\documentclass{situationist-preprint}

\usepackage{longtable}
\usepackage{array}

\title{On Personal Identity}
\msid{PID-01-a}
\msfamily{PID-01}
\msdiscipline{Personal Identity and the Self}
\mstype{Principal manuscript}
\mslean{lean/Philosophy/PID/PID\_01.lean}
\mskeywords{personal identity; psychological continuity; branching; fission; teletransportation; tie-breaker; phase sortal; animalism; exhaustive enumeration}

% Generated by generate.py; do not edit.
\newcommand{\stcases}{11{,}009}
\newcommand{\ststages}{43{,}137}
\newcommand{\sthone}{26}
\newcommand{\sthtwo}{5{,}437}
\newcommand{\sththree}{5{,}546}
\newcommand{\stclsN}{4{,}445}
\newcommand{\stclsB}{2{,}388}
\newcommand{\stclsE}{2{,}076}
\newcommand{\stclsF}{624}
\newcommand{\stclsBE}{228}
\newcommand{\stclsBF}{1{,}248}
\newcommand{\stclsEF}{0}
\newcommand{\stclsBEF}{0}
\newcommand{\stmachine}{6{,}162}
\newcommand{\stfusionstruct}{7{,}092}
\newcommand{\stvI}{14{,}664}
\newcommand{\stvL}{4{,}152}
\newcommand{\stvE}{5{,}880}
\newcommand{\stvD}{3{,}448}
\newcommand{\stvU}{7{,}776}
\newcommand{\stvN}{7{,}217}
\newcommand{\stordstages}{2{,}634}
\newcommand{\stdisC}{4{,}338}
\newcommand{\stdisP}{4{,}064}
\newcommand{\stdisW}{7{,}792}
\newcommand{\stwidenonlocal}{2{,}430}
\newcommand{\stbiodis}{6{,}411}
\newcommand{\stbiodisN}{3{,}951}
\newcommand{\stnonbranching}{4{,}445}
\newcommand{\stpbdefect}{4{,}064}
\newcommand{\stcontestB}{4{,}152}
\newcommand{\stcontestE}{2{,}304}
\newcommand{\stcontestF}{1{,}872}
\newcommand{\caseOrdinary}{A0005}
\newcommand{\caseDisrupted}{A0001}
\newcommand{\caseTeleport}{A0004}
\newcommand{\caseBranchLine}{A0016}
\newcommand{\caseBranchLineLate}{B0150}
\newcommand{\caseFission}{A0015}
\newcommand{\caseFissionUneven}{A0012}
\newcommand{\caseBranchLineWeak}{A0014}
\newcommand{\caseUploadD}{A0007}
\newcommand{\caseUploadN}{A0026}
\newcommand{\caseUploadDownload}{B0039}
\newcommand{\caseTeleBranch}{B0140}
\newcommand{\caseFusionLife}{B0292}
\newcommand{\caseFusionNone}{B0291}
\newcommand{\caseTransplant}{A0009}
\newcommand{\caseBranchLineTele}{B0921}
\newcommand{\caseFissionLater}{B0892}
\newcommand{\caseBranchLineThenTele}{B0925}
\newcommand{\caseSleepLong}{C0113}


\newenvironment{casetable}{\par\scriptsize\setlength{\tabcolsep}{3.5pt}\setlength{\LTleft}{0pt}\setlength{\LTright}{\fill}%
  \begin{longtable}{@{}>{\ttfamily}l>{\ttfamily}l>{\ttfamily}l>{\ttfamily}l>{\ttfamily}l>{\ttfamily}l@{}}
  \textbf{ID} & \textbf{structure} & \textbf{PID-01} & \textbf{bio} & \textbf{cl} & \textbf{rivals}\\ \midrule\endhead}%
  {\end{longtable}}

\begin{document}
\maketitle

\begin{abstract}
What makes a person at one time the same person as a person at another? We argue that the answer is psychological continuity, with bodily continuity as a tie-breaker in cases of branching. \enquote{Person} is a phase sortal of the organism. A later person-stage belongs to the same person as an earlier one if it is psychologically continuous with it, by a non-branching process, through a cause of the right kind, with a realizer of the original's kind. Where continuity branches, the tie is broken by the continuation of the life of the original organism, not by the quantity of matter retained: a branch is the original only if it alone continues that life, and otherwise neither is. We prove that the criterion respects the logic of identity, never making two simultaneous stages stages of the original; that in ordinary cases it agrees with the biological and bodily criteria; that its verdict on a stage never depends on what happens later; and that it yields the verdicts on teletransportation, fission and uploading. An appendix computes the verdicts for all \stcases{} case structures within a stated bound and checks each result. The uniqueness theorem is formalized in Lean.
\end{abstract}

\section{The problem}\label{s:problem}

A woman walks into a room at noon and a woman walks out at one. What makes it the case, when it is the case, that the woman who walks out is the woman who walked in? The question is not how we know it. We know it, in the cases that occur, by seeing that the same body went in and came out, by recognizing a face and a voice, and occasionally by asking what she remembers. The question is what the fact consists in: what relation between the person at noon and the person at one is necessary and sufficient for their being one and the same person.

There are three families of answer. The \emph{psychological} answer says that the fact consists in relations of memory, intention, belief, desire, value and character between the earlier and the later person, or in the continuity that overlapping chains of such relations make up. The \emph{physical} answer says that it consists in the continuity of something material: the body, the brain, or the life of the organism. The \emph{simple} answer says that it consists in nothing further: personal identity is a primitive fact, not made up of any other relations, and the relations of the first two answers are at best evidence of it. Each answer has an ancestry going back at least to the seventeenth century, and each has been refined in response to a few imagined cases that the others handle better.

We defend a version of the psychological answer. Its central claim is that a person persists by psychological continuity. It has three further features, each of which responds to a difficulty that the history of the debate has made vivid. First, it says what persons are, so that the criterion is not left hanging: \enquote{person} is a phase sortal of the organism, true of an organism during the phase in which it has the capacities for rational and self-conscious thought, as \enquote{child} is true of it during an earlier phase. Secondly, it requires that the continuity be non-branching and that it proceed through a cause of the right kind, where a cause is of the right kind only if the later stages are realized in a thing of the same kind as the earlier ones. Thirdly, it says what happens when continuity branches. The tie is then broken by bodily continuity, read as the continuation of the life of the original organism: a branch is the original if it alone continues that life, and otherwise no branch is.

The answer has to meet four tests. It must deliver the verdicts on the imagined cases that have shaped the debate: the transplant of a cerebrum, teletransportation with and without destruction of the original, fission, fusion and uploading. It must respect the logic of identity, which is one-one and transitive while psychological continuity can hold one-many. It must agree with the biological and bodily criteria in the ordinary cases in which almost every judgment of identity is made, so that the evidence on which we rely is evidence of what identity consists in. And it must not make the identity of a person at a time depend on what happens later. We prove that the criterion meets the last three tests in general, and that it meets the first for every case in a large finite class. The class is that of all case structures, as Definition~\ref{d:case} will make precise, with at most four person-stages after the original, at most two at each of at most three later times, at most two immediate predecessors for each stage and at most one stage realized in a thing of another kind. There are \stcases{} such structures up to isomorphism. Appendix~\ref{s:appendix} lists every one, with the verdict of the criterion on every stage and the verdicts of four rival criteria, and the program that produced it checks each of the numbered results of Sections~\ref{s:criterion}--\ref{s:ordinary} on each case.

The plan is this. Section~\ref{s:history} reviews the debate from Locke to the present and draws from it the constraints that a criterion must meet. Section~\ref{s:framework} sets out the framework: what a person is, what a case structure is, and the three principles of the criterion. Section~\ref{s:criterion} states the criterion and proves that it respects the logic of identity. Section~\ref{s:branching} treats branching and the tie-breaker, and derives the verdicts of PID-02, PID-03 and PID-14. Section~\ref{s:ordinary} treats ordinary cases and temporal locality. Section~\ref{s:cases} discusses particular entries of the enumeration. Section~\ref{s:objections} answers objections, Section~\ref{s:limits} records what the result does not establish, and Section~\ref{s:formal} describes the formalization.

\section{The debate}\label{s:history}

\subsection{Locke}

The modern debate begins with the chapter on identity and diversity that Locke added to the second edition of the \emph{Essay} \citep[II.xxvii]{Locke1694}. Locke's first move was to insist that identity is relative to a sort: whether something at one time is the same as something at another depends on what kind of thing is in question. The same mass of matter is the same particles; the same oak or the same horse is the same continued life, communicated to successive particles organized in the right way; the same man is the same living organized body. The same person, Locke held, is none of these. A person is a thinking intelligent being that has reason and reflection and can consider itself as itself, the same thinking thing, in different times and places (II.xxvii.9), and the identity of a person reaches as far back as its consciousness can be extended to past actions and thoughts.

Two of Locke's examples set the agenda. If the soul of a prince, carrying the consciousness of the prince's past life, entered the body of a cobbler, the cobbler's body would afterwards house the prince: everyone would say the same man, but the person would be the prince (II.xxvii.15). And if the same man had distinct and incommunicable consciousnesses by day and by night, the day man and the night man would be distinct persons (II.xxvii.23). The first example is the ancestor of every transplant case; the second, of every case of divided consciousness. Locke also gave a reason for caring about the distinction: \enquote{person} is a forensic term, appropriating actions and their merit, and so belongs to intelligent agents capable of law, happiness and misery (II.xxvii.26). Whom we may reward and punish for an act should depend on whether the agent can own the act, which on Locke's view is a matter of consciousness.

Locke's view contains, in compressed form, both of the ideas that the present manuscript combines. It contains the psychological criterion: personal identity follows consciousness. And it contains the doctrine that identity is sortal-relative, with the life of the organism as the principle of identity for the man. Locke treated the man and the person as two things that may coincide, but he did not ask what happens when consciousness is extended to two later beings, nor how the person is related to the man when they coincide.

\subsection{Hume}

Hume gave the problem a skeptical turn. Looking within, he found no impression of a self that persists unchanged through a life, only a succession of perceptions, related by resemblance and causation, which the imagination, passing easily along them, takes for one continued thing \citep[I.iv.6]{Hume1739}. The identity we ascribe to the mind is therefore, on his view, a fiction of the same kind as the identity we ascribe to a plant or an animal whose parts are replaced, and disputes about it are at bottom verbal, except so far as the relation of parts gives rise to some fiction or imaginary principle of union. Two features of Hume's discussion survive in the modern debate. His relations of resemblance and causation are the ancestors of psychological connectedness, and his comparison with plants and animals anticipates the role that the continuation of a life, a process that replaces its materials while keeping its organization, plays in the tie-breaker of Section~\ref{s:branching}. His conclusion that the identity is fictitious we do not accept: PID-04 holds that the self is a construction, and that constructions are real, and the criterion below specifies the relations on which the construction rests.

\subsection{Butler and Reid}

Two objections to Locke became canonical. \citet{Butler1736} argued that consciousness of personal identity presupposes personal identity and so cannot constitute it: I can remember only my own experiences, so that a genuine memory of an experience already implies that the rememberer is the one who had it. The objection does not refute the claim that identity consists in memory; it shows that the memory relation cannot be specified, in the criterion, as one that already involves identity. \citet[III.6]{Reid1785} gave the second objection. A brave officer remembers being flogged as a boy for robbing an orchard; as a general he remembers taking a standard as an officer, but has entirely lost the memory of the flogging. By Locke's criterion the general is the officer and the officer is the boy, but the general is not the boy, which violates the transitivity of identity.

Reid himself held that personal identity is simple and indivisible, that it admits of no degrees and is not made up of anything else; memory gives us certain knowledge of our own past, and the testimony of others and the evidence of the body give us knowledge of the identity of others, but none of these constitutes identity \citep[III.4]{Reid1785}. Butler held a similar view. The simple view has had defenders since, most prominently \citet{Swinburne1984}, who argued that the possibility of a person's continuing in either of two fission products, with no fact about the products that settles which, shows that personal identity is a further fact over and above physical and psychological continuity.

The two objections shaped the later psychological tradition. The transitivity objection was answered by replacing direct memory connections with continuity: overlapping chains of direct connections, which is transitive by construction. The circularity objection was answered, as we shall see, by Shoemaker's notion of quasi-memory. Neither answer touches the third difficulty the tradition later discovered, which is that continuity, unlike identity, can hold one-many.

\subsection{Williams}

Williams opened the twentieth-century debate with a dilemma for the psychological criterion \citep{Williams1957}. Suppose a man, Charles, wakes one morning claiming to remember the life of Guy Fawkes, and that his claims are in every checkable detail correct. Should we say that he is Guy Fawkes? If it is logically possible that Charles undergoes this change, it is logically possible that his brother Robert undergoes it at the same time. Both cannot be Guy Fawkes, since they are not one another; it would be arbitrary to say that one of them is; so neither is. But if neither is Guy Fawkes when both have the memories, then Charles is not Guy Fawkes when he alone has them, since whether Charles is Guy Fawkes cannot depend on what happens to Robert. This is the \emph{reduplication argument}. It concludes that psychological similarity, however complete, cannot constitute personal identity, and that the identity of persons requires bodily continuity.

The argument rests on a premise that the debate came to call the only-$x$-and-$y$ principle: whether $x$ is $y$ depends only on facts about $x$, $y$ and the relations between them \citep{Wiggins1980,Noonan1985}. Every non-branching criterion denies this premise, and so the reduplication argument was the first form of what we shall treat as the problem of branching. Williams later gave a second argument \citep{Williams1970}. Two people, A and B, are told that their bodies will undergo an exchange of psychological states, after which one of the resulting persons will be tortured and the other rewarded. Told the story in these terms, it is natural for A to want the B-body-person rewarded, since that person will have A's memories, character and concerns. Told instead that he will be tortured after his memories have been removed and replaced by others, A will naturally fear the torture. Since the two stories describe the same events, Williams concluded that our attitudes do not reliably follow psychological continuity, and that the case against bodily criteria is weaker than it looks. We shall return to the lesson in Section~\ref{s:objections}: it concerns what matters in survival as much as identity.

\subsection{Shoemaker}

\citet{Shoemaker1963} described the case that made the psychological criterion the default view for a generation. Brown's brain is removed and placed in the head of Robinson, whose own brain has been removed; the resulting person, Brownson, has Robinson's body and Brown's memories and character. Most who consider the case judge that Brownson is Brown. The case is Locke's prince and cobbler with the soul replaced by a brain, and it shows that the bodily criterion conflicts with the judgment that a person goes where her psychology goes, even when the psychology is carried by a physical organ.

\citet{Shoemaker1970} then answered Butler. A \emph{quasi-memory} is a state that is like a memory except that it does not imply that the subject had the remembered experience: it requires only that someone had it and that the state is caused by the experience in the right way. Memory is quasi-memory of one's own experiences. The criterion can be stated in terms of quasi-memory, and of the analogous quasi-intentions and the persistence of beliefs and traits, without presupposing identity. Shoemaker also gave the psychological criterion its causal form: the later states must be caused by the earlier ones, and the causal chain must be of an appropriate kind. In the form defended by \citet{Shoemaker1984}, the theory is functionalist: mental states are individuated by their causal roles, and personal identity consists in the continuity of a psychology, so understood, by a non-branching process. Shoemaker accepted that the relevant causation need not run through a single brain; he considered a device that transfers the states of one brain to another, and argued that a person could survive by such a transfer. The device is instructive for the kind condition defended below. It moves a psychology from one human brain to another, so that the later stages are realized in a thing of the original's kind and evolve under the same dynamics; on our criterion the transfer preserves identity when it is uncontested, as teletransportation does. Shoemaker also answered the reduplication argument in the way that became standard: the criterion is non-branching, so that when Charles alone has Guy Fawkes's psychology, by a cause of the right kind, he is Guy Fawkes, and when Charles and Robert both have it, neither is. The answer accepts that whether Charles is Guy Fawkes depends on what happens to Robert, which is the premise of Williams's argument; it rejects the inference from that dependence to the conclusion that psychology cannot constitute identity. We return to this point in reply to Objection~\ref{o:onlyxy}.

\subsection{Parfit}

\citet{Parfit1971,Parfit1984} gave the psychological criterion its classical statement and drew from it the most radical consequences. His Relation R is psychological connectedness and continuity, with any cause; personal identity consists in R holding uniquely, so that the criterion is non-branching by stipulation. Parfit distinguished three versions according to the cause required: on the narrow version the cause must be the normal one, the continued existence of the brain; on the wide version any reliable cause will do; on the widest, any cause \citep[ch.~10]{Parfit1984}. He described the cases that the present manuscript calls by his names. In \emph{Simple Teletransportation} a scanner records the state of every cell of a person's body, destroying it, and a replicator on Mars builds an exact copy from the record. In the \emph{Branch-Line} case a new scanner records the information without destroying the original, and the original, on Earth, is told that the scan has damaged her heart and that she will die within days. In \emph{My Division} the two halves of a person's brain, each of which carries her whole psychology, are transplanted into two bodies.

Parfit's conclusion was that identity is not what matters in survival. In My Division, each successor stands to the original in Relation R, and the original cannot be both; to say that she is one would be arbitrary, and to say that she is neither would be to say that a relation that would have been survival had one transplant failed becomes death because both succeed. Parfit concluded that the question whether she survives is empty, in the sense that once we know the facts about R we know everything there is to know, and that what matters is R, not identity. He supported the conclusion with the spectrum cases of \citep[ch.~11]{Parfit1984}, in which a person's cells or psychological states are replaced in degrees, and in which there is no sharp line between survival and replacement. In the \emph{combined spectrum}, the cases range from one in which nothing is changed to one in which every cell of a person's brain and body is replaced and her psychology is replaced by that of someone else; in between, the replacement is partial in both respects. Parfit argued that there is no plausible place at which to draw a line between the cases in which the original survives and those in which she is replaced, that a sharp line at an arbitrary point would make a trivial difference into the difference between life and death, and that the question of identity in the middle cases is therefore empty.

Parfit's view is a form of \emph{reductionism}: the fact of personal identity consists in nothing more than the holding of physical and psychological continuities, which can be described without presupposing that the subjects whose states they relate are persons. Reductionism, so stated, is common ground between Parfit and the present manuscript. What we do not accept is his further claim that, because the fact reduces, the question of identity can be empty in cases where all the continuities are known. Our criterion gives a determinate answer in every case of the enumeration, including the Branch-Line case and fission, and it locates whatever indeterminacy remains in the vagueness of the predicates the criterion uses, which by MET-18 is not indeterminacy in the facts. \citet{Johnston1997} argued, against Parfit, that even if reductionism is true it does not support the revision of our concerns that Parfit proposed, since the ordinary concern for one's own future does not presuppose a further fact. PID-06 accepts a moderate form of Parfit's revision: what matters is carried mostly by connectedness, with identity retaining a small residual weight.

Two alternatives to Parfit's non-branching clause were proposed at about the same time. \citet{Nozick1981} proposed the closest continuer theory: a later being is the original if it is the closest of her continuers and close enough, where closeness weighs psychological and physical continuity together; if two continuers are equally close, neither is the original. \citet{Lewis1976} proposed that in fission there were two persons all along, sharing their stages before the fission, so that identity is not violated and nothing that matters is lost; Lewis's account requires the perdurantist view that persons are sums of temporal stages, which PID-11 defends. \citet{Perry1972} considered fission at length and proposed that the identity of a person may be relative to the time at which she is counted. We return to all three proposals in Section~\ref{s:branching}.

\subsection{Olson, animalism and the human being}

The last major development is animalism. \citet{Snowdon1990}, \citet{Olson1997,Olson2007} and, in a different metaphysical setting, \citet{vanInwagen1990} argued that each of us is a human organism, an animal of the species \emph{Homo sapiens}, and that an organism persists through the continuity of its biological life, not of its psychology. Olson's argument is the thinking animal argument: there is a human animal where you are; it thinks your thoughts; and if you were not it, there would be two thinkers of your thoughts, one of which would be mistaken about which thing it is. Olson argued that the psychological criterion gives the wrong persistence conditions for the thing that thinks: it implies that you were never a fetus, since a fetus has no psychology, and that in a cerebrum transplant you would leave behind a living, empty-headed animal that had until then been thinking your thoughts. PID-05 defends animalism in this form. Two further problems sharpen the challenge. The \emph{fetus problem} is that each of us was once a fetus, which had no psychology, so that if we persist by psychological continuity we were never fetuses, and some other thing, the organism, was the fetus and later came to think our thoughts. The \emph{too many thinkers} problem is that, if we are not organisms, then both we and the organisms that coincide with us think our thoughts, and we can have no reason to believe that we are the persons rather than the organisms \citep{Olson2007}. Both problems are answered by the distinction between the person, unified by psychological continuity, and the organism, unified by its life, if the person is a phase of the organism. The organism was a fetus, and is what each of us is; the person is the organism in its personal phase; there is one thinker at each time, counted by its present stage. \citet{Johnston2007} pressed against animalism the cases in which an organism is reduced to a living brain or head: the animalist must say either that the organism survives as the reduced remnant, against the biological criterion, or that a new thinking thing comes into existence when the reduction is completed, which is hard to believe. PID-05-a answers the objection for the brain in the vat; our criterion adds that, whatever the organism does, the person continues in the remnant, by uncontested continuity of the right kind.

\citet{Johnston1987} took a related but distinct position. He argued that we are human beings, and that the method of deciding the question by consulting our responses to imagined cases is less probative than its users assumed, since our responses may reflect an implicit theory that does not settle the matter. Johnston allowed that a human being could survive as a functioning brain, so that the transplant case does not by itself favor the psychological criterion over a criterion based on the human being. In later work \citep{Johnston2010}, he argued that our persistence conditions are not fixed independently of our dispositions of self-identification, a view that makes the self, in his term, protean, and he drew from it consequences for the significance of death. The position challenges every criterion that, like ours, claims to state the persistence conditions of persons independently of how persons regard themselves.

The debate since 2000 has developed along these lines. \citet{Baker2000} argued that persons are constituted by organisms without being identical with them. \citet{McMahan2002} argued that we are embodied minds, essentially the parts of brains that sustain consciousness, and \citet{Parfit2012} came to a similar view, arguing that we are the thinking parts of animals. \citet{Unger1990} defended a physically based psychological criterion, on which what persists must be the physical realizer of a core psychology. \citet{Schechtman1996} argued that the question of reidentification has been confused with the question of what makes a life one's own, which a narrative account answers, and \citet{Schechtman2014} developed a \enquote{person life} view on which the unit of persistence is a characteristically human life, integrating biological, psychological and social functions. \citet{Snowdon2014} restated animalism. \citet{Chalmers2010} raised the question of uploading, arguing that gradual uploading offers the best prospect of survival; PID-14 answers it.

\subsection{Recent views and the present criterion}

The present criterion shares something with each of the recent views and differs from each at an identifiable point. The embodied mind view of \citet{McMahan2002} and the view of \citet{Parfit2012} that we are the thinking parts of animals agree with us on the transplant: the person goes with the cerebrum. They disagree on teletransportation. If a person is the part of a brain that thinks, she does not survive the destruction of that part, and the replica is someone else; our criterion says that the person, in the sense of the psychological criterion, survives, and PID-05 adds that the organism does not. The disagreement reflects a choice about what the thing whose persistence is in question is: these views take it to be a physical part, we take it to be the person as a phase of the organism, unified by a relation that can be carried by a reconstruction of a realizer of the right kind.

The physically based criterion of \citet{Unger1990} requires that the physical realizer of a core psychology continue. It agrees with our kind condition in excluding uploads and in counting the gradual replacement of a brain's matter as survival so long as the realizer's capacity is continuously preserved. It differs from ours on teletransportation, which on Unger's view breaks the continuity of the realizer, and so it gives the verdict that the narrow version of the psychological criterion gives. Our kind condition requires sameness of the kind of realizer and of its dynamics, not of the realizer itself.

The constitution view of \citet{Baker2000} holds that a person is constituted by an organism, as a statue is constituted by a lump of marble, without being identical with it, and that what makes a thing a person is a capacity for a first-person perspective. It agrees with us that the person and the organism can come apart. It differs in making the person a second thing that coincides with the organism in ordinary cases; on our view the person in ordinary cases is the organism in a phase, and the second sum that the psychological criterion picks out diverges from the organism only when the psychology leaves it. The person life view of \citet{Schechtman2014} comes closest to the phase sortal: a person's persistence is that of a characteristically human life in which biological, psychological and social functions are integrated. Our tie-breaker can be read as a precise statement of one element of that view, that when psychology branches the life in question is the biological life of which the psychology was a phase.

Animalism \citep{Olson1997,Snowdon2014} is not a rival of our criterion but a thesis about a different question, as PID-05-a argues. The protean view of \citet{Johnston2010} is a rival, since it denies that persistence conditions are fixed independently of identification, and we answer it in Section~\ref{s:objections}.

\subsection{What a criterion must do}

The history leaves five constraints.

\emph{The transplant.} A criterion must say that a person goes where her psychology goes when it is moved, intact and uniquely, by a process that keeps it realized in a thing of the right kind: Locke's cobbler, Shoemaker's Brownson. A criterion that says otherwise of persons must explain the near-unanimity of the judgment; animalism does so by distinguishing the person from the organism, and PID-05 adopts that distinction.

\emph{The Branch-Line case.} A criterion must not say that a person who has been scanned, without harm, ceases to exist because a replica has been made elsewhere. A pure non-branching criterion says exactly this.

\emph{The logic of identity.} Identity is one-one, transitive and absolute (LOG-12). A criterion must never make one earlier person identical with two later persons who exist at the same time, and must not make identity depend on a relation that is not transitive.

\emph{Ordinary cases.} In the cases in which psychology continues in a living human organism without interruption and without branching, a criterion must agree with the biological and bodily criteria. Otherwise ordinary bodily evidence would not be evidence for the relation that constitutes identity, and almost all our judgments of identity would be unwarranted.

\emph{Temporal locality.} Whether a later stage is a stage of the original should be settled by what has happened up to that stage's time, not by what happens afterwards. Williams's reduplication argument turns on the intuition that whether Charles is Guy Fawkes cannot depend on what happens to Robert; we accept that identity can depend on what happens elsewhere at the same time (PID-02-a, Corollary~3.4), but a criterion that made a person's identity depend on what will happen later would make it unknowable in principle until the end of time.

The criterion of this manuscript meets all five. The first two are verdicts on cases, which Proposition~\ref{p:verdicts} derives; the third is Theorem~\ref{t:unique}; the fourth is Theorem~\ref{t:ordinary}; the fifth is Theorem~\ref{t:local}.

\section{Framework}\label{s:framework}

\subsection{Persons as a phase of the organism}

\begin{definition}[Person]\label{d:person}
A thing is a \emph{person} at a time if and only if it has at that time the capacities for rational thought and for thinking of itself as itself, the same thinking thing in different times and places. \enquote{Person} is a \emph{phase sortal} of the organism: it is true of a human organism during the phase of its life in which it has these capacities, as \enquote{child} and \enquote{adult} are true of it during other phases \citep{Wiggins1980}. A \emph{person-stage} is a temporal stage of a thing at a time at which the thing is a person.
\end{definition}

Definition~\ref{d:person} takes Locke's definition of a person and Wiggins's distinction between substance sortals, which give the persistence conditions of the things that fall under them, and phase sortals, which apply to such things for part of their existence. It agrees with PID-05-a, Definition~2.2, and with the reconciliation in PID-05-a, Proposition~4.1. The organism is what each of us fundamentally is. The person is the organism in its personal phase, in every case in which the organism's psychology continues in the organism. The criterion of this manuscript says when person-stages are stages of one person; since \enquote{person} is a phase sortal, it does not say what the thinker of a person's thoughts fundamentally is. In the cases in which a person's psychology leaves the organism, the person, in the sense of the criterion, and the organism come apart, and PID-05-a explains why \enquote{you} then refers to the organism. Following PID-11, we read persons as perdurers: the person is a sum of person-stages unified by the relation the criterion specifies, and the organism a sum of organism-stages unified by the continuation of a life.

\subsection{Case structures}

\begin{definition}[Case structure]\label{d:case}
A \emph{case structure} consists of finitely many person-stages, each at one of the times $0, 1, \dots, h$, with a single \emph{original} stage $o$ at time $0$, and a set of \emph{edges}, each from a stage at a time $t$ to a stage at time $t+1$, such that every stage other than $o$ has at least one incoming edge. Each stage is realized in a thing of the original's kind or in a thing of another kind. Each edge carries a degree of psychological connectedness (Definition~\ref{d:conn}) and records whether the later stage continues the life of the organism that realizes the earlier (Definition~\ref{d:life}). An edge from $x$ to $y$ means that $y$'s psychological states, or $y$'s organism, or both, are produced by those of $x$; we call $x$ a \emph{parent} of $y$.
\end{definition}

A case structure is the skeleton of a case described in the literature, stripped of everything except the relations on which the rival criteria disagree. The times are not clock times but the successive phases of a process: the scan, the transmission, the activation of a replica, the death of an original. A stage with no outgoing edge is a stage after which that line of the process ends, by death or destruction. A stage with two outgoing edges is one from which two later stages descend, as in Branch-Line and fission. A stage with two incoming edges is the product of a fusion, or of a transplant of a psychology into an existing organism. Every link spans one phase; Section~\ref{s:limits} records what this excludes.

\begin{definition}[Connectedness]\label{d:conn}
A \emph{direct psychological connection} between a stage $x$ and a later stage $y$ is the holding at $y$ of a quasi-memory of an experience at $x$, the carrying out at $y$ of an intention formed at $x$, or the persistence from $x$ to $y$ of a belief, desire, value or trait of character, or of another psychological capacity. Connections are for the most part \emph{dispositional}: they consist in the persistence of states and capacities, not in their exercise. The \emph{degree} of connectedness along an edge is $2$ if the connections hold to the extent that they hold over a day of an ordinary life, $1$ if they hold strongly but in part, and $0$ if they do not hold strongly.
\end{definition}

The scale is coarse by design. Its three values are what the rival criteria need: the psychological criterion needs the distinction between strong connection and its absence, and the closest continuer theory needs, in addition, the distinction between more and less connection. The threshold between degree $0$ and degree $1$ is vague, as are the predicates of everyday thought. By MET-18 vagueness lies in language and knowledge, so that there is a sharp threshold whose location we cannot know; the criterion inherits whatever vagueness the notion of strong connection has, and the enumeration treats the degree as given.

Definition~\ref{d:conn} counts connections of character and value, and of psychological capacities generally, not only of memory. A person who loses her autobiographical memory but keeps her beliefs about the world, her values, her skills and her character remains strongly connected to her earlier self, as PID-08 argues. The definition also makes connection mainly dispositional, which settles the old worry about interruption: in dreamless sleep, under general anesthesia or in a coma from which a person recovers, there may be no occurrent mental states, but the beliefs, intentions, values and traits persist in the brain that realizes them, and memories formed before are retrievable after. An edge spanning such an interval has degree $2$.

\begin{definition}[Continuity]\label{d:cont}
A later stage $y$ is \emph{psychologically continuous} with an earlier stage $x$ if there is a path of edges from $x$ to $y$ each of degree at least $1$. It is psychologically continuous with $x$ \emph{through a cause of the right kind} if, in addition, every stage on some such path is realized in a thing of $x$'s kind.
\end{definition}

Continuity is the ancestral of strong connection, and so it is transitive, which answers Reid. It requires a path through actual intermediate stages, each strongly connected with the next, and so a general who remembers nothing of his boyhood is nonetheless continuous with the boy, through the officer.

\begin{definition}[Life]\label{d:life}
A later stage $y$ \emph{continues the life} of an earlier stage $x$ if $y$ is realized in an organism whose life, the self-maintaining metabolic process by which an organism takes in matter, regulates its own conditions and restores its organization after perturbation (PID-05-a, Definition~2.1), is the continuation of the life of the organism that realizes $x$. An edge along which this holds is a \emph{life edge}.
\end{definition}

The notion of a life is Locke's principle of the identity of an oak or a horse, and the principle of persistence for organisms that the biological criterion states (PID-05-a, Assumption~2.7). A life continues through the replacement of the matter it organizes, through the loss of limbs, through the transplant of organs into and out of the body and through dreamless sleep and anesthesia. It does not continue into a replica built from a record, which begins a new life, nor into a brain removed from the body and sustained by a machine, which is an organ kept alive and not a living organism (PID-05-a, Assumption~2.8), nor into a computer. A life edge joins two stages realized in organisms; a stage realized in a thing of another kind neither continues nor is continued by a life.

\begin{assumption}[Lives do not divide or merge]\label{a:life}
The life of a human organism at a time is continued by at most one organism at a later time, and the life of a human organism at a time is the continuation of the life of at most one organism at an earlier time.
\end{assumption}

Assumption~\ref{a:life} is a biological fact about human beings, not a stipulation. Some organisms divide: an amoeba's life ends, or continues twice, when it splits. A human organism after its early embryonic stages does not, and the cases in which two human organisms come to share a body, such as conjoined twinning, are cases in which two lives are integrated to some degree, not one in which a life has become two. Nothing in what follows depends on how such cases are classified, since none of them involves a person-stage that continues two lives.

\subsection{The beginning and end of the personal phase}

Since \enquote{person} is a phase sortal, the person in the sense of Definition~\ref{d:person} begins when the organism acquires the capacities for rational and self-conscious thought, and may end before the organism does. The criterion concerns stages within the personal phase, and three consequences follow.

First, the person was never a fetus, though the organism was, and so, by PID-05, were you. The criterion is silent on the identity of the organism before its personal phase, which is a question for the biological criterion; it says only that the first person-stages of an organism are unified with its later ones by psychological continuity. The acquisition of the capacities is gradual, and the beginning of the personal phase is vague in the way that the threshold of strong connection is vague.

Secondly, if an organism permanently loses the capacities for rational and self-conscious thought, as in a persistent vegetative state, the person has no further stages, though the organism lives on. The biological criterion and the criterion then come apart, since the organism continues and the person does not; but the later stages of the organism are not person-stages, and so are not stages of any person. In the encoding of Appendix~\ref{s:appendix} such stages do not appear, since every stage of a case structure is a person-stage. What does appear is the case in which the organism goes on as a person but with none of the original's psychology, as when a disease or an injury destroys the capacities on which strong connection depends and a new psychology later develops in the same organism; the encoding represents it by a life edge of degree $0$.

Thirdly, a progressive dementia that erodes memory but leaves other capacities intact for a long time does not end the person at the first loss. By Definition~\ref{d:conn}, connections of belief, value, character and skill continue to hold, and strong connection, which is the condition of each step, is a matter of the overall extent of connection, not of any one kind. The person ends, if she does, when the overall extent of connection over a short interval falls below the threshold, and not before.

\subsection{The three principles}

\begin{principle}[Psychological criterion]\label{p:psych}
A later person-stage is a stage of the same person as an earlier one if it is psychologically continuous with it through a cause of the right kind, by a process that does not branch; and it is a stage of the same person only if it is psychologically continuous with it through a cause of the right kind.
\end{principle}

\begin{principle}[Kind condition]\label{p:kind}
Psychological continuity is through a cause of the right kind only if the later stages are realized in a thing of the same kind as the earlier ones. For a human person this kind is living human neural tissue whose activity evolves under the dynamics of a human brain: in every ordinary case and in teletransportation a living human organism, and in the case of a brain kept alive in a vat the living brain of such an organism.
\end{principle}

Principle~\ref{p:kind} is the kind condition of PID-14-a, Principle~P2, which states it for the cases in which the realizer is a whole organism; we state it so that it applies also to the brain in the vat, which PID-05-a, Proposition~4.1, treats as a case in which the person, in the sense of the psychological criterion, continues. Its support is the composition principle of MET-06-a, applied to persons by PID-11: only stages that are linked by the relevant dynamics, in that later stages are what earlier ones evolve into under the laws that govern them, compose a persisting thing. A human person's psychology is realized in, and evolves under, the dynamics of a living human brain. A teletransporter's replica, a transplanted cerebrum and a brain in a vat continue those dynamics; a computer that runs an emulation of the brain does not, since its states evolve by the execution of a program on hardware, whatever their functional organization. The condition is narrower than Parfit's wide criterion and wider than his narrow one: it does not require the normal cause, but it does require a cause that reconstructs or preserves a realizer of the original's kind.

\begin{principle}[Tie-breaker]\label{p:tie}
Where the process branches, a branch is the original if and only if it, and no other branch, continues the life of the original's organism. If no branch alone continues that life, neither branch is the original.
\end{principle}

Principle~\ref{p:psych} is the non-branching psychological criterion of \citet{Shoemaker1984} and \citet{Parfit1984}, narrowed by Principle~\ref{p:kind}; Principle~\ref{p:tie} is the bodily tie-breaker, read as PID-01-b reads it, as the continuation of a life rather than the quantity of matter retained. The two principles are stated loosely here, with \enquote{the process branches} and \enquote{the original's organism} left at the level of the literature. Section~\ref{s:criterion} makes them precise, and the precision matters: the obvious reading of \enquote{the process branches}, on which a process branches if at any time two later stages are continuous with the original, makes the criterion violate temporal locality (Proposition~\ref{p:global}), and the obvious reading of \enquote{the original's organism}, as the organism that realized the original at the beginning of the case, gives the wrong verdict whenever a person who has survived teletransportation later undergoes a Branch-Line scan.

\section{The criterion}\label{s:criterion}

The two refinements are that branching is assessed at each step, relative to the latest stage of the original, and that the life whose continuation breaks a tie is the life of the organism that realizes that latest stage. Both follow PID-01-b, Definition~2.1, which defines the continuation of a person's life relative to an earlier person $x$, as the continuation of the life of the organism in which $x$ was realized.

\begin{definition}[Candidate]\label{d:cand}
A stage $y$ is a \emph{candidate} of a stage $x$ if there is an edge from $x$ to $y$ of degree at least $1$ and both $x$ and $y$ are realized in a thing of the original's kind. A candidate is a one-step continuation through a cause of the right kind.
\end{definition}

\begin{definition}[Contested step]\label{d:contested}
The step from $x$ to a candidate $y$ of $x$ is \emph{contested} if $x$ has another candidate (\emph{fission}), or if $y$ is a candidate of another stage $z \neq x$ (\emph{fusion}). Otherwise it is \emph{uncontested}.
\end{definition}

\begin{definition}[The criterion]\label{d:crit}
The stages of the original person are given inductively. The original stage $o$ is a stage of the original. If $x$ is a stage of the original and $y$ a candidate of $x$, then $y$ is a stage of the original if and only if the step from $x$ to $y$ is uncontested, or it is contested and the edge from $x$ to $y$ is a life edge. Nothing else is a stage of the original. A later stage that is a stage of the original is \emph{identical with the original}, in the sense that it is a stage of the person whose stage $o$ is.
\end{definition}

Definition~\ref{d:crit} implements Principles~\ref{p:psych}--\ref{p:tie}. A step that preserves strong connection through a realizer of the right kind, and that is uncontested, is non-branching continuity, and by Principle~\ref{p:psych} it preserves identity. A contested step is a branching, and by Principle~\ref{p:tie} it preserves identity just when it continues the life of the organism that realizes the latest stage of the original. The clause \enquote{nothing else} is the necessity half of Principle~\ref{p:psych}, with Principle~\ref{p:kind}: a stage is a stage of the original only if it is reached from $o$ by right-kind continuity.

In the enumeration of Appendix~\ref{s:appendix}, each later stage receives one of six verdicts. The verdict \texttt{I} says that the stage is identical with the original. The verdicts \texttt{L} and \texttt{E} are given to candidates of a stage of the original that the criterion excludes in a contested step: \texttt{L} when another candidate of that stage continues the life and is the original, \texttt{E} when no candidate of that stage is the original, the \enquote{neither} verdict of Principle~\ref{p:tie}. The verdict \texttt{D} is given to a stage continuous with the original through a cause of the right kind that descends from a stage that is not the original, such as the later stages of a replica in Branch-Line. The verdict \texttt{U} is given to a stage that is psychologically continuous with the original only through a realizer of another kind, such as an upload, and the verdict \texttt{N} to a stage that is not psychologically continuous with the original at all. Only \texttt{I} is identity; \texttt{L}, \texttt{E}, \texttt{D} and \texttt{U} mark different kinds of successor, which is the term PID-14-a, Definition~2.2, uses for a later person psychologically continuous with the original who is not she.

\begin{proposition}[Path characterization]\label{p:path}
A stage $y$ is identical with the original if and only if there is a path from $o$ to $y$ every step of which is from a stage to a candidate of it, and is either uncontested or along a life edge.
\end{proposition}

\begin{proof}
If $y$ is a stage of the original, the inductive clauses of Definition~\ref{d:crit} that make it so trace such a path back to $o$. Conversely, given such a path, induction along it shows that each stage on it is a stage of the original, since each step meets the condition of the inductive clause.
\end{proof}

Proposition~\ref{p:path} gives a second way of computing the criterion, by searching for paths, and the program of Appendix~\ref{s:appendix} computes it in both ways and checks that they agree on every stage of every case.

\begin{lemma}\label{l:parent}
Every stage of the original other than $o$ is a candidate of a stage of the original one time earlier.
\end{lemma}

\begin{proof}
By the clause \enquote{nothing else} of Definition~\ref{d:crit}, a stage of the original other than $o$ is made so by the inductive clause, which makes it a candidate of a stage of the original; and edges join stages at consecutive times (Definition~\ref{d:case}).
\end{proof}

\begin{theorem}[Uniqueness]\label{t:unique}
At no time are two distinct stages identical with the original. The criterion therefore never makes one person identical with two simultaneous persons, and respects the logic of identity.
\end{theorem}

\begin{proof}
By induction on time. At time $0$ only $o$ is a stage of the original, since by Lemma~\ref{l:parent} every other stage of the original is at a time after some other stage. Suppose that at time $t$ at most one stage is a stage of the original, and that $y$ and $y'$ are stages of the original at time $t+1$. By Lemma~\ref{l:parent}, $y$ is a candidate of a stage $x$ of the original at $t$ and $y'$ of a stage $x'$ of the original at $t$, and by the hypothesis $x = x'$. If $y \neq y'$, then $x$ has two candidates, so the steps from $x$ to $y$ and from $x$ to $y'$ are both contested (Definition~\ref{d:contested}). By Definition~\ref{d:crit}, since both $y$ and $y'$ are stages of the original, both edges are life edges. Then the life of the organism that realizes $x$ is continued by two organisms at $t+1$, contrary to Assumption~\ref{a:life}. So $y = y'$.
\end{proof}

The proof shows where the logic of identity is secured. It is not secured by the non-branching clause alone, since the tie-breaker can make a branch the original; it is secured by the conjunction of the tie-breaker with the biological fact that lives do not divide. A tie-breaker based on any relation that can itself hold one-many, such as degree of material continuity, would need a further clause to break its own ties. The theorem is the main result of the formalization (Section~\ref{s:formal}), where the dependence on Assumption~\ref{a:life} appears as the dependence on the axiom \texttt{life\_unique}.

\begin{corollary}[One chain]\label{c:chain}
The stages of the original form a single unbroken chain $o = x_0, x_1, \dots, x_k$, with $x_i$ at time $i$ and each $x_{i+1}$ a candidate of $x_i$. The relation of being stages of the original person is reflexive, symmetric and transitive on the stages it relates, and the original person, once she has ceased to have a stage, has no later stage.
\end{corollary}

\begin{proof}
By Theorem~\ref{t:unique} there is at most one stage of the original at each time, and by Lemma~\ref{l:parent} each such stage after $o$ is a candidate of the one at the preceding time; so if there is a stage of the original at time $t$ there is one at every earlier time. The relation of being stages of the original is that of belonging to one chain, which is an equivalence relation.
\end{proof}

Corollary~\ref{c:chain} has a consequence for the debate between Lewis and his critics. On the perdurantist reading of PID-11, a person is a sum of stages unified by a persistence relation, and Lewis held that in fission there are two such sums sharing the pre-fission stages \citep{Lewis1976}. Our criterion picks out, for a given original stage, a single sum: the chain of Corollary~\ref{c:chain}. When the chain ends at a fission, the successors on each side begin chains of their own. Whether one also counts, as Lewis did, the sum of the original's chain with a successor's chain as a person is a question about which sums of stages the persistence relation unifies, and PID-03-a, Remark~3.5, shows that nothing about what matters depends on the answer. What the criterion settles is that the original, counted by her stage before the fission, is not identical with either successor.

\begin{remark}\label{r:reid}
Corollary~\ref{c:chain} answers Reid. The general is identical with the boy, because there is a chain of uncontested steps of strong connection from the boy to the officer and from the officer to the general, though the general has no memory of the boy's flogging. The criterion requires continuity, not connectedness between the end points. PID-07 draws the consequence that a person is identical with her childhood self even where the connections between them are weak.
\end{remark}

\section{Branching and the tie-breaker}\label{s:branching}

\begin{definition}[Branching]\label{d:branch}
A case \emph{branches at} a stage $x$ of the original if some step from $x$ to a candidate of $x$ is contested. The candidates of $x$, and in a fusion the stage into which $x$'s psychology is merged, are the \emph{branches} at $x$.
\end{definition}

Definition~\ref{d:branch} makes precise the definition of PID-01-b, Definition~1.1, and of PID-02-a, Definition~2.2, on which a process branches if two or more later persons are each psychologically continuous with an earlier person. The difference is that branching is assessed at the latest stage of the original, not over the whole history of the case. The two readings agree on every case discussed in PID-02-a and PID-03-a, and they disagree only where a person who has already survived a branching, or a teletransportation, branches again, or where branching occurs among the successors of a person who is no longer the original. We show in Section~\ref{s:ordinary} that the global reading violates temporal locality and that ours does not.

\begin{lemma}\label{l:onelife}
At any stage of the original, at most one candidate continues the life.
\end{lemma}

\begin{proof}
By Assumption~\ref{a:life}, the life of the organism that realizes a stage is continued by at most one later stage.
\end{proof}

\begin{theorem}[Tie-breaker]\label{t:tie}
Let the case branch at a stage $x$ of the original. (i) If a candidate of $x$ continues the life of $x$'s organism, that candidate is the original, and every other candidate of $x$ is not. (ii) If no candidate of $x$ continues that life and $x$ has two or more candidates, none of them is the original. (iii) If $x$ has a single candidate $y$ and the step is contested by fusion, $y$ is the original if and only if the edge from $x$ to $y$ is a life edge.
\end{theorem}

\begin{proof}
In (i), the life-continuing candidate is the original by Definition~\ref{d:crit}, since its step is contested and along a life edge; any other candidate's step is contested, since $x$ has two candidates, and is not along a life edge, by Lemma~\ref{l:onelife}, so that candidate is not the original. In (ii), each candidate's step is contested and none is along a life edge. Clause (iii) is Definition~\ref{d:crit} for a step contested only by fusion.
\end{proof}

\begin{corollary}[Neither]\label{c:neither}
In fission, in which each of two candidates of a stage of the original is realized in an organism that does not continue the original's life, the original is identical with neither.
\end{corollary}

\begin{proposition}[Verdicts]\label{p:verdicts}
The criterion yields the following verdicts.
\begin{enumerate}[label=(\roman*)]
\item In Simple Teletransportation, and in the single transplant of PID-03-a, Assumption~2.5, the person survives as the replica or the recipient.
\item In Branch-Line, the original is the person on Earth and not the replica.
\item If the original is destroyed after the replica is activated, she is the person on Earth until her destruction and the replica is a distinct person; if she is destroyed before it is activated, she survives as the replica.
\item In fission she is neither successor.
\item In destructive and nondestructive uploading, the upload is not the original; in nondestructive uploading she is the person who remains in her body.
\item In the brain-in-a-vat case, the person in the vat is the original person, though not the original organism.
\end{enumerate}
\end{proposition}

\begin{proof}
(i) The replica is a candidate of the original, since it is strongly connected with her and realized in a living human organism, and the step is uncontested, since the original has no other candidate. So it is the original by Definition~\ref{d:crit}. The same holds of the recipient of a single transplanted hemisphere. (ii) The original has two candidates, the person on Earth and the replica; the case branches at the original; the person on Earth continues the life of her organism and the replica does not, so by Theorem~\ref{t:tie}(i) the person on Earth is the original and the replica is not. (iii) If the original exists when the replica is activated, the case branches at the original as in (ii), and by Corollary~\ref{c:chain} the replica, having once failed to be the original, never becomes her; if the original is destroyed before activation, the original stage has a single candidate, and (i) applies. (iv) is Corollary~\ref{c:neither}: in the fission of PID-03-a, Definition~2.1, each hemisphere is housed in a body of its own and the original organism does not survive. (v) An upload is realized in a thing of another kind, so it is not a candidate (Definition~\ref{d:cand}); it is not a stage of the original by the clause \enquote{nothing else}. In nondestructive uploading the person on Earth is the only candidate of the original, so her step is uncontested and she is the original. (vi) The vat brain is realized in a thing of the original's kind (Principle~\ref{p:kind}), and its step is uncontested; it is therefore the original in the sense of Definition~\ref{d:crit}. By PID-05-a, Theorem~3.2, it is not the original organism.
\end{proof}

Clause (iii) is PID-02-a, Theorem~3.3, including the temporal condition of its revision, and clauses (i) and (ii) are its Propositions~3.1 and~3.2. Clause (iv) is PID-03-a, Lemma~3.1, clause (v) is PID-14-a, Theorem~3.1 and Proposition~3.4, and clause (vi) is the reading of the vat case in PID-05-a, Proposition~4.1. The proof of clause (v) shows something PID-14-a states separately: nondestructive uploading does not even make the case branch at the original, since the upload is not a candidate. Its failure to be the original is overdetermined: it fails Principle~\ref{p:kind}, and it would fail Principle~\ref{p:tie} if Principle~\ref{p:kind} were dropped (PID-14-a, Proposition~3.4).

\begin{proposition}[Kind]\label{p:kindprop}
Every stage of the original is realized in a thing of the original's kind; in particular no stage that is psychologically continuous with the original only through a realizer of another kind is the original, even if it is itself realized in a thing of the original's kind.
\end{proposition}

\begin{proof}
By Lemma~\ref{l:parent} and Definition~\ref{d:cand}, every stage of the original after $o$ is realized in a thing of the original's kind, and so is every stage on the chain of Corollary~\ref{c:chain}. A stage that is continuous with the original only through a realizer of another kind is reached by no chain of candidates, and so is not a stage of the original.
\end{proof}

The second clause of Proposition~\ref{p:kindprop} concerns a case that the literature has not much discussed: a person is uploaded, and the emulation is later downloaded into a newly grown human body. The final stage is realized in a thing of the original's kind, but its continuity with the original passes through the computer, under a different dynamics. On Principle~\ref{p:kind} the chain of right-kind continuity is broken at the upload and is not restored by the download. The downloaded person is a successor of the original, as the upload was, and her relation to the original has, by PID-06, much of what matters.

\subsection{Against the alternatives}

There are four other ways of handling branching. The enumeration computes the first three on every case, and Appendix~\ref{s:appendix} reports the cases on which they differ from ours.

\paragraph{No tie-breaker.} The pure non-branching criterion says that in a contested step neither branch is the original. In Branch-Line it says that the woman on Earth, who has been read by a scanner and has otherwise undergone nothing, ceased to exist when the replica was activated. This is not a verdict that the theory of personal identity can accept. The person on Earth has the same body, the same uninterrupted stream of consciousness and the same organism as before; nothing has happened to her except that a record has been made and used elsewhere. On the criterion without a tie-breaker, any person can be killed by being copied without her knowledge. The defect appears in every case in which a life-continuing branch exists; the enumeration finds it in \stdisP{} cases.

\paragraph{Degree of psychological continuity.} The closest continuer theory of \citet{Nozick1981}, in the form that weighs psychological continuity only, says that in a contested step the candidate most closely connected with the original is the original, and that if two are equally close neither is. In Branch-Line the replica is exactly similar, so the degrees are equal and the theory gives the same intolerable verdict as the pure non-branching criterion. Where the degrees differ, it makes identity turn on the difference between a memory kept and a memory lost: if the scan slightly disturbs the original's memory, the theory makes the replica the original and the woman on Earth, with her body, life and unbroken consciousness, someone else. The enumeration exhibits the case as \caseBranchLineWeak. A difference in the degree of connection is not the sort of fact that could make the difference between a person's survival and her death, as PID-01-b argues. The theory differs from ours in \stdisC{} cases.

\paragraph{Degree of bodily continuity.} A tie-breaker that weighs the quantity of body or brain retained, so that the successor with more of the original matter is the original, has the same defect in material form. When a brain's two hemispheres are divided between two bodies, nothing about the person, her thought or her life distinguishes the successor with fifty-one per cent of the neurons from the one with forty-nine in a way that could ground identity with one and not the other. Our case structures do not represent quantities of matter, and so the enumeration does not compute this rival; but its defect does not need computing. A life, by contrast, is all or nothing in the cases at issue: an organism whose metabolism continues in a body is alive, and a hemisphere transplanted into a new body continues no life but that body's. This is why Principle~\ref{p:tie} looks to the continuation of a life.

\paragraph{The global reading.} The criterion that \citet{Parfit1984} states says that personal identity consists in Relation R holding uniquely, and on a natural reading uniqueness is assessed over the whole history of the case: if at any time two later stages are each R-related to the original, identity holds with neither. This is the \emph{wide global} criterion of the enumeration, wide because it allows any reliable cause, global because it assesses branching over the whole case. It agrees with ours on many cases, but differs in two ways. Without Principle~\ref{p:kind} it counts an upload as a branch, and so says that in nondestructive uploading the person on Earth, who has merely been scanned, ceased to be when the upload was started. And reading branching globally it violates temporal locality, as Proposition~\ref{p:global} will show. It differs from ours in \stdisW{} cases.

\paragraph{Convention.} One might hold that when psychology branches there is no fact about which branch is the original until we decide, and that a convention may settle it. A convention that settled the matter by stipulation would not explain why the Branch-Line verdict is not open to decision. No community could reasonably decide that the woman on Earth had ceased to exist; and if the convention that it would adopt is fixed by the facts, then the facts already settle the matter, and the convention adds nothing. Perry's proposal that identity be relative to the time of counting \citep{Perry1972} is a sophisticated form of this view, and it gives up the absoluteness of identity, which LOG-12 defends.

\subsection{Why the tie-breaker is not ad hoc}

A clause is ad hoc if its only support is that it saves a theory from counterexamples. Principle~\ref{p:tie} has independent support.

The first source is Definition~\ref{d:person}. A person, in the ordinary case, is an organism during the phase in which it has the capacities for rational and self-conscious thought, and the person's stages are unified by psychological continuity because those capacities are exercised by one organism (PID-05-a, Proposition~4.1). When psychology branches, the organism is not divided: it remains where its life continues. The person phase of that organism continues with it, so long as the organism keeps its psychology. Principle~\ref{p:tie} says no more than this. The psychological criterion departs from this default only when the organism ceases, or loses the psychology, and the psychology continues uniquely elsewhere, as in teletransportation and the transplant.

The second source is Principle~\ref{p:kind}. The kind condition already gives the realizer a role in identity: continuity sustains the persistence of a person only when the later stages are realized in a thing of the same kind as the earlier ones. In a branching each branch may meet this condition, and the question is which of two realizers of the right kind is the one the person had. The continuation of one life is the natural answer, because a life is what makes successive stages of an organism stages of one realizer.

The third source is the structure of the logic of identity. Theorem~\ref{t:unique} requires a tie-breaker based on a relation that cannot itself branch. Lives are the only relation in the vicinity that satisfies this condition as a matter of biological fact, rather than by a stipulated metric; quantities of matter and degrees of connection can tie, and lives cannot fork.

\subsection{Fusion}

In fusion, two persons' psychologies are merged into one later person who is strongly connected with both. The case is the mirror of fission, and Definition~\ref{d:contested} treats it as a contest. Clause (iii) of Theorem~\ref{t:tie} says that the fusion product is the original if and only if it continues the life of the original's organism. If the merged psychology is realized in the original organism, the original survives, enriched by the psychology of the other person, who has ceased to be. If it is realized in the other person's organism, that person survives and the original has ceased to be. If it is realized in a new organism, neither survives. The verdicts follow the same rule as fission, and they are required by the logic of identity for the same reason: the fusion product cannot be identical with two distinct earlier persons. The enumeration contains \stfusionstruct{} structures with a stage that has two parents, and the program checks clause (iii) on every one in which a step from a stage of the original is contested by fusion.

\section{Ordinary cases and temporal locality}\label{s:ordinary}

\begin{definition}[Ordinary stage]\label{d:ordinary}
A stage $y$ is \emph{ordinary} if it is reached from $o$ by a path of uncontested steps, each from a stage to a candidate of it along a life edge.
\end{definition}

Definition~\ref{d:ordinary} restates, in the terms of the framework, the ordinary case of PID-01-c, Definition~1.1: both stages belong to one living human organism, the organism has the capacities for rational and self-conscious thought throughout, and no other person is psychologically continuous with the earlier stage. The requirement that each step be of degree at least $1$ is what PID-01-c, Lemma~2.1, derives from the normal operation of one brain in one organism; in the framework it is part of what makes the case ordinary.

\begin{theorem}[Agreement in ordinary cases]\label{t:ordinary}
At every ordinary stage the criterion, the biological criterion and the criterion of bodily continuity all give identity with the original.
\end{theorem}

\begin{proof}
By induction along the path of Definition~\ref{d:ordinary}. Each step is from a stage of the original to a candidate of it, and is uncontested, so by Definition~\ref{d:crit} the later stage is a stage of the original. Each step is along a life edge, so the later stage's organism continues the life of the original organism, which is the biological criterion of PID-05-a, Assumption~2.7, and the body that has that life is the original body continued.
\end{proof}

In the encoding of case structures, bodily continuity and the continuation of life coincide, since a body that ceases to be alive is not the realizer of a person-stage; in the column \texttt{bio} of the tables, \texttt{+} marks stages that both criteria count as the original, reached from $o$ by life edges.

\begin{corollary}[Determinacy]\label{c:determinate}
In ordinary cases personal identity is determinate: every criterion in contention gives the same verdict, so that no indeterminacy in which criterion is correct can make the verdict indeterminate.
\end{corollary}

Corollary~\ref{c:determinate} is the conclusion of PID-01-c in its strongest form. Ordinary cases differ from the puzzle cases in kind, not degree: in every puzzle case some relation that normally accompanies psychological continuity is absent, whether the organism, the body, uniqueness or the same kind of realizer. Whatever indeterminacy the puzzles reveal arises from the separation of relations that ordinary cases keep together, and so cannot spread to them. It follows, as PID-01-c, Corollary~2.3, argues, that the bodily and behavioral evidence on which we rely in reidentifying persons is evidence of the relation that constitutes their identity. Williams's complaint against psychological criteria, that we do not reidentify persons by their memories and could not \citep{Williams1957}, mistakes a claim about what identity consists in for a claim about how it is known.

\begin{proposition}[Where the criteria diverge]\label{p:diverge}
In a case that does not branch at any stage of the original, the criterion and the biological criterion disagree about a stage only if the stage is not ordinary: either the psychology has at some step left the organism, by a step that is not along a life edge, or the organism has at some step lost the psychology, by a life edge of degree $0$.
\end{proposition}

\begin{proof}
By Theorem~\ref{t:ordinary}, the criteria agree at every ordinary stage. A stage that is not ordinary is reached from $o$ by no path of uncontested life edges of degree at least $1$; in a case that does not branch at any stage of the original, no step from a stage of the original is contested, so on any path from $o$ to the stage some step either is not a life edge or has degree $0$.
\end{proof}

The disagreements Proposition~\ref{p:diverge} allows are the ones PID-05-a predicts. When the psychology leaves the organism, as in teletransportation, the transplant and the vat, the criterion follows the psychology and the biological criterion follows the organism: the person continues, and the organism, which is what you are, does not. When the organism loses the psychology, through injury or disease that destroys the capacities on which strong connection depends, the biological criterion says that the organism continues and the criterion says that the person whose stages were unified by psychological continuity has ceased to have stages. The two sums, in the language of PID-11, share stages in ordinary cases and diverge in these. The enumeration contains \stbiodisN{} cases that do not branch at any stage of the original and in which the criteria disagree about some stage; the program checks that every disagreement is at a stage that is not ordinary.

\begin{proposition}[Dispositional continuity]\label{p:sleep}
Dreamless sleep, general anesthesia and comparable interruptions of consciousness do not interrupt personal identity: a stage after such an interruption is ordinary if the stage before it is.
\end{proposition}

\begin{proof}
By Definition~\ref{d:conn}, the degree of connectedness across the interruption is fixed by the persistence of dispositional states, which persist in the brain that realizes them, and so is $2$. The organism's life continues, so the edge is a life edge, and no other stage is a candidate of the stage before, so the step is uncontested. By Definition~\ref{d:ordinary} the later stage is ordinary if the earlier is.
\end{proof}

\begin{theorem}[Temporal locality]\label{t:local}
Whether a stage at time $t$ is identical with the original depends only on the stages and edges at times up to $t$.
\end{theorem}

\begin{proof}
By Definition~\ref{d:crit}, whether a stage $y$ at $t$ is a stage of the original depends on whether its parent $x$ at $t-1$ is a stage of the original, on whether $y$ is a candidate of $x$, on whether the step is contested and on whether it is along a life edge. Whether $y$ is a candidate of $x$ and whether the edge is a life edge are features of the edge between them. Whether the step is contested depends on the other candidates of $x$, which are at $t$, and on the other parents of $y$, which are at $t-1$. By induction on $t$, whether $x$ is a stage of the original depends only on the case up to $t-1$.
\end{proof}

Theorem~\ref{t:local} is what PID-02-a needs in its reply to the closest continuer theory: our verdict does not make the identity of a person at a time depend on what happens later. Identity on our view is extrinsic in the sense of PID-02-a, Corollary~3.4: whether the person on Earth in Branch-Line is the original depends on whether a replica was activated, and whether a replica is the original depends on whether the original survived its activation. But all the facts on which identity depends are facts about the present and the past of the stage in question. The global reading of branching does not have this property.

\begin{proposition}[The global reading is not local]\label{p:global}
On the wide global criterion, whether a stage at time $t$ is the original can depend on events after $t$.
\end{proposition}

\begin{proof}
In case \caseTeleBranch{} of the enumeration, the original is teletransported, and the person who results is later scanned without harm, a replica of her being made. Truncated after the teletransportation, the case does not branch, and the wide global criterion says that the person who results is the original. In the full case, two later stages are each continuous with the original at the second time, so on the global reading the case branches and the wide global criterion says that no later stage, including the person who resulted from the teletransportation, is the original.
\end{proof}

The enumeration finds \stwidenonlocal{} cases in which the wide global criterion's verdict on some stage changes when the case is extended, and none in which ours does. The difference matters. On the global reading, a person who has survived teletransportation, and who has lived for years as the original, would turn out never to have been the original if a copy of her is later made from a scan. Whether she is the person who signed a contract, or who is owed a debt, would be unsettled while she lives. This is the defect Williams's reduplication argument was designed to expose, and it is a defect of the global reading of the non-branching clause, not of the psychological criterion.

\section{Further cases}\label{s:further}

The criterion was shaped by the cases of Section~\ref{s:history}. Its value is tested by cases it was not built for. We take nine, and state the verdict of the criterion on each, with the principles from which it follows.

\paragraph{Hemispherectomy.} A person's left hemisphere is removed to stop intractable seizures, and she continues with the right hemisphere alone, in her own body. Her psychology is diminished in some respects, but her beliefs, values, character and much of her memory persist, and her organism lives on. The step is along a life edge, of degree at least $1$, and uncontested, since the removed hemisphere is destroyed. The stage after the operation is ordinary (Definition~\ref{d:ordinary}), and by Theorem~\ref{t:ordinary} every criterion agrees that she survives. The case shows that the person does not consist of the matter that realizes her, nor of half of it: what persists is the organism's psychology, realized in what remains.

\paragraph{The asymmetric division.} Suppose instead that the removed hemisphere, which is as capable as the remaining one of sustaining the whole psychology, is transplanted into a new body. The original organism continues with one hemisphere and a second organism houses the other. The case has the structure of Branch-Line: two candidates, one continuing the life. By Theorem~\ref{t:tie}(i) the person who stays in the original body is the original, and the recipient is a successor. Those who find this verdict arbitrary, since the two hemispheres are equal, should observe that it is the verdict the transplant intuition itself delivers when a single hemisphere is removed and destroyed: the person who remains in the body is the original. The second hemisphere's being kept alive elsewhere does not change what happened to her.

\paragraph{The body exchange.} In Williams's case, the psychologies of A and B are exchanged, so that the person in B's body has A's psychology and conversely \citep{Williams1970}. For each original, the case has one candidate, the person in the other body, and the step is uncontested; so A survives in B's body and B in A's, as the first presentation of the story suggests. The second presentation, in which A is told that he will lose his memories and be given others before being tortured, describes the same events as happening to the organism A, and the fear it elicits is fear on behalf of the organism, which PID-05 says is what A is. The two presentations track two sums of stages that share their earlier stages and diverge: the person and the organism. Williams's case therefore does not show that our attitudes are unreliable; it shows that they have two objects.

\paragraph{The combined spectrum.} In the middle of Parfit's spectrum, part of a person's brain is replaced by cells that carry part of someone else's psychology. Our criterion asks of each step whether the later stage is strongly connected with the earlier, realized in a thing of the same kind and, if the step is contested, whether it continues the life. In the spectrum the organism's life continues throughout, so there is no contest, and identity holds exactly when the connection with the original is strong. The threshold is vague, and by MET-18 it is sharp but unknowable. The question is therefore not empty, but its answer cannot be known in the middle of the spectrum. What matters, by PID-06, varies smoothly with connectedness and does not depend on where the threshold falls.

\paragraph{Neural prosthesis.} Neurons are replaced one by one with devices that are supplied by the brain's circulation and integrated into its physiology. By the reasoning of PID-14-a, Section~4, and MET-06-a, the devices become parts of the living organism, the psychology remains realized in a living human organism, and each step is ordinary. The person survives. If instead the functions are transferred step by step to an external computer and the brain is dismantled, the realizer is at the end of another kind, and by Proposition~\ref{p:kindprop} the end point is not the original.

\paragraph{Total amnesia.} A person loses all episodic and autobiographical memory, retrograde and anterograde, while retaining her beliefs about the world, her skills, her values and her character. By Definition~\ref{d:conn}, connections of these kinds are direct psychological connections, and their persistence across the onset of the amnesia keeps the step strong. She survives, as PID-08 argues. The criterion does not need memory, only some sufficient extent of connection of some kinds.

\paragraph{Split brains.} After commissurotomy, the two hemispheres of a patient's brain sometimes support experiences and actions that are not integrated with one another, as in the experiments in which each hand acts on information presented to one visual field. The case is not branching in the sense of Definition~\ref{d:branch}: there is one organism, one life and, by MND-15, a consciousness that is partially unified rather than divided into two streams; PID-15 reaches the same conclusion for dissociative conditions. Every step is along a life edge and uncontested, so the person persists. Branching requires two later persons, each continuous with the original, and partial disunity within one person does not supply them.

\paragraph{Conjoined twins.} In dicephalic twinning, two heads, each with its own brain and its own psychology, share one body below the neck. If the twins form one organism, two persons are phases of one organism, each unified with her own earlier stages by psychological continuity. The criterion does not require that each person be the only person realized in her organism; it requires that each be psychologically continuous with her own earlier stages, and the tie-breaker is never called on, since no psychology branches. Whether the twins are one organism or two is a question for biology that the criterion does not need to settle.

\paragraph{Teletransportation with delay.} The original is destroyed, and the record of her state is stored for a year before a replica is activated. The step from the original's last stage to the replica spans the interval, and it is uncontested, so she survives, scattered in time, as PID-11 allows and PID-02-a's temporal condition requires. If a second replica is activated from the same record after the first, the second replica is continuous with the original through the record but is not a candidate of the latest stage of the original, which by then is the first replica; the second is a successor. If the two are activated at the same moment, the case is fission, and neither is the original.

\section{Consequences}\label{s:consequences}

Locke called \enquote{person} a forensic term, and a criterion of personal identity is tested by what it implies for the practices in which identity is used. We draw out four consequences.

\paragraph{Responsibility.} A person is responsible for an act only if she is the person who did it, and on our criterion she is so exactly when a chain of uncontested steps of strong connection, or of life-continuing steps where the chain branches, runs from the agent to her. Identity is not, however, the whole of the matter. PID-07 argues that the obligations a person incurred as a child are weak because her connectedness with her childhood self is weak, though identity holds; the same reasoning applies to the distant past of an adult, and to a person whose character has been transformed by injury. The criterion fixes who is answerable; connectedness, by PID-06, bears on how much. In Branch-Line, the person on Earth is answerable for what the original did, and the replica, who is not the original, is not, though she remembers doing it and is as connected with the original as the person on Earth. This is the verdict that the forensic use of the concept requires: two people cannot both be answerable as the agent of one act, though both may have reason to feel as if they were.

\paragraph{Prudence.} The original has reason to care about her future in the way she cares about her own life, and the criterion says whose future is hers. But the extent of that reason is fixed mostly by connectedness, not identity (PID-06). In fission, the original has nearly as much reason to care about each successor as about her own future, though she is neither (PID-03); in Branch-Line, she has reason to care about the replica, who will remember her life and carry on her projects, though the replica is not she. A person deciding whether to step into a nondestructive scanner should therefore not reason that, since she will not be the replica, the replica's future is no concern of hers; and a person deciding whether to step into a destructive teletransporter should not reason that, since she will survive, nothing of value is at stake in the timing of the destruction. The first is a mistake about what matters, and the second a mistake about identity, which by PID-02-a depends on that timing.

\paragraph{Dementia and advance decisions.} A person in the late stages of a dementia is, on our criterion, the person who earlier made decisions about her care, so long as the chain of strong connection is unbroken, and in most cases it is unbroken, since the decline is gradual and each step preserves much. Her earlier decisions are therefore decisions about her own future, not about a successor's. But her connectedness with her earlier self may by then be weak, and the weight of the earlier decisions, which derives in part from the fact that they express the values of the person whose life is in question, weakens with it. The criterion settles the question of identity on which the authority of advance decisions is often taken to depend, and leaves to the theory of what matters the question how much that authority is worth when connectedness has faded.

\paragraph{Records and copies.} Theorem~\ref{t:local} has a practical consequence for any future in which scanning and replication are possible. Whether a person who has survived a teletransportation is the original is settled when she is activated, and does not depend on whether her record is later used again. Institutions that must decide who holds a person's property, or who is party to her contracts, can therefore rely on the order of events: the first activation of a record after the original's destruction yields the original, and every later use of the record, while she exists, yields a successor. A global reading of the non-branching clause would leave every such question open for as long as a record survives.

\section{What the enumeration shows}\label{s:cases}

Appendix~\ref{s:appendix} lists \stcases{} case structures, all those within the bound stated at the end of Section~\ref{s:problem}, with \ststages{} later stages among them. The program that produced it computes the criterion of Definition~\ref{d:crit} on every stage, and checks Theorem~\ref{t:unique}, Corollary~\ref{c:chain}, Proposition~\ref{p:path}, Theorem~\ref{t:tie}, Proposition~\ref{p:kindprop}, Theorem~\ref{t:ordinary}, Proposition~\ref{p:diverge} and Theorem~\ref{t:local} on every case. Every check succeeds. The checks are not independent evidence that the criterion is true; they are evidence that the criterion as stated has the properties claimed for it, in every case of the class, and that the general proofs of Sections~\ref{s:criterion}--\ref{s:ordinary} contain no error that a case in the class would expose. The appendix's preface describes the encoding; here we discuss particular entries.

\paragraph{Ordinary survival and sleep.} Case \caseOrdinary{} is the simplest case: one stage after the original, strongly connected to degree $2$ and continuing the life. All criteria give identity. Case \caseSleepLong{} is a person who passes, over three phases, through two intervals in which connectedness is partial, as across a long illness or in the slow change of character that PID-07 discusses. Every stage is ordinary, every criterion gives identity, and Proposition~\ref{p:sleep} accounts for interruptions of consciousness in the same way.

\paragraph{Teletransportation and the vat.} Case \caseTeleport{} is a single step of degree $2$ to a stage of the original's kind that does not continue the life. It represents Simple Teletransportation, the transplant of a cerebrum into a body whose own brain has been removed, the single transplant of PID-03-a and the brain in a vat: the encoding does not distinguish them, because they do not differ in any relation on which the criteria disagree. The criterion gives identity and the biological criterion denies it, which is the divergence of Proposition~\ref{p:diverge}.

\paragraph{Branch-Line.} Case \caseBranchLine{} is Branch-Line: two candidates of the original, one continuing the life. The row reads \texttt{LI}: the replica, \texttt{1a}, loses the tie, and the person on Earth, \texttt{1b}, is the original. The rival column shows that the closest continuer theory, the criterion without a tie-breaker and the wide global criterion all deny that the person on Earth is the original; this is the intolerable verdict described above, and it is the reason a tie-breaker is needed. Case \caseBranchLineLate{} adds a further phase in which the replica lives on and the original has died: the replica's later stage, \texttt{2a}, receives \texttt{D}, a successor descending from a stage that was not the original. This is the case added in the revision of PID-02-a: the original was destroyed after the replica was activated, and she does not survive as the replica. Case \caseBranchLineThenTele{} continues Branch-Line differently: the person on Earth is later teletransported while the replica lives on. The criterion says that she survives the teletransportation, since the step from her stage on Earth is uncontested, and the wide global criterion says that she does not, because at the later time two stages are continuous with the original.

\paragraph{Branch-Line with a disturbed original.} Case \caseBranchLineWeak{} is Branch-Line in which the scan has partly disrupted the original's psychology, so that the person on Earth is connected to degree $1$ and the replica to degree $2$. The criterion gives identity to the person on Earth, who continues the life. The closest continuer theory gives it to the replica. It seems to us plain that a scan that disturbs some of a woman's memories does not thereby transfer her existence to Mars.

\paragraph{Fission.} Case \caseFission{} is symmetric fission: two candidates of degree $2$, neither continuing the life. The row reads \texttt{EE}, and every rival agrees. Case \caseFissionUneven{} is uneven fission, with one successor more closely connected than the other; the criterion still gives \texttt{EE}, since neither continues the life, while the closest continuer theory makes the better-connected successor the original. Case \caseFissionLater{} lets both successors live on; their later stages receive \texttt{D}. Each successor is a person, as PID-03-a, Theorem~3.2, holds, and each has much of what matters (PID-03-a, Theorem~3.3; PID-06).

\paragraph{Uploading.} Case \caseUploadD{} is destructive uploading; the upload receives \texttt{U}, and only the wide global criterion, which drops Principle~\ref{p:kind}, makes it the original. Case \caseUploadN{} is nondestructive uploading: the row reads \texttt{IU}, the person on Earth being the original and the upload a successor of another kind; the wide global criterion, counting the upload as a branch, says that the person on Earth ceased to be. Case \caseUploadDownload{} is an upload later downloaded into a new human body; both later stages receive \texttt{U}, as Proposition~\ref{p:kindprop} requires.

\paragraph{Fusion.} Cases \caseFusionLife{} and \caseFusionNone{} begin with Branch-Line and then merge the person on Earth and the replica into one person. In \caseFusionLife{} the merged person continues the life of the original's organism and is the original; in \caseFusionNone{} the merged person is realized in a new organism, and by Theorem~\ref{t:tie}(iii) she is neither the original nor the replica, receiving \texttt{E}.

\paragraph{Teletransportation followed by branching.} Case \caseTeleBranch{} is the case of Proposition~\ref{p:global}. The original is teletransported (\texttt{1a}, \texttt{I}), and the resulting person, whose new organism has a life of its own, is later scanned without harm (\texttt{2b} continues that life, \texttt{2a} is the new replica). The criterion gives \texttt{I/LI}: the teletransported person is the original, and remains so on Earth after the second scan, because the tie at \texttt{1a} is broken by the life of \texttt{1a}'s organism, which is the replica's organism of the first teletransportation, not the organism that realized $o$. A tie-breaker that looked to the life of the organism that realized $o$ would give \texttt{I/EE}, killing a person who has been scanned without harm, and this is why Definition~\ref{d:crit} relativizes the tie-breaker to the latest stage of the original.

\paragraph{Totals.} Of the \stcases{} cases, \stclsN{} do not branch at any stage of the original. In the others, the program finds \stcontestB{} contested steps resolved by a life-continuing candidate, \stcontestE{} contested steps in fission with no life-continuing candidate, resolved as \enquote{neither}, and \stcontestF{} steps contested by fusion. Among the \ststages{} later stages, \stvI{} are identical with the original; \stvL{} lose a tie to a life-continuing branch; \stvE{} are branches of which neither is the original; \stvD{} descend from a stage that is not the original; \stvU{} are continuous with the original only through a realizer of another kind; and \stvN{} are not psychologically continuous with the original. There are \stordstages{} ordinary stages, and at every one all criteria agree.

\section{Objections and replies}\label{s:objections}

\begin{objection}[The only-$x$-and-$y$ principle]\label{o:onlyxy}
Whether the person on Earth in Branch-Line is the original depends, on this criterion, on whether a replica was made on Mars; whether a replica is the original depends on whether the original survived. Identity has been made to depend on facts about third things. But whether $x$ is $y$ should depend only on $x$, $y$ and the relations between them \citep{Wiggins1980,Noonan1985}. The criterion therefore fails, and the reduplication argument of \citet{Williams1957} shows why.
\end{objection}

\begin{reply}
The extrinsicness is real, and PID-02-a accepts it, in its Corollary~3.4. It is the price of any account on which identity is one-one and the relation that grounds it can hold one-many; and the relation that grounds the persistence of persons can hold one-many, as Branch-Line shows. The only-$x$-and-$y$ principle holds for kinds whose persistence relation cannot branch, such as artefacts, whose persistence is grounded in a maintenance process aimed at one artefact (MET-07), and fails for kinds whose persistence relation can, as PID-02-a argues. Two features of our criterion limit the damage. First, by Theorem~\ref{t:local} the extrinsic facts on which identity depends are never later than the stage in question, so the reduplication argument's sharpest form, on which Charles's identity with Guy Fawkes depends on what will later happen to Robert, does not apply. Secondly, the extrinsic dependence is confined to branching cases: in every case that does not branch at a stage of the original, the verdict on a stage depends only on the chain of steps leading to it. And in branching cases the tie-breaker restores an intrinsic relation as the decisive one: whether a branch continues the life of the original's organism is a relation between the branch and the original. What is extrinsic is only whether the tie-breaker is called on.
\end{reply}

\begin{objection}[Collapse]\label{o:collapse}
If bodily continuity can decide identity where psychology branches, it should decide it everywhere. A criterion that lets the body decide in Branch-Line, but not in teletransportation or the transplant, has two criteria and no principle for choosing between them. Either the psychological criterion is the criterion, and the tie-breaker is an arbitrary addition, or the bodily criterion is, and the psychological criterion is a mistake.
\end{objection}

\begin{reply}
The tie-breaker operates only among branches, each of which is psychologically continuous with the original through a cause of the right kind. It never makes a person identical with a stage that has lost her psychology: a life edge of degree $0$ gives a stage that is not a candidate and so cannot be the original, as case \caseDisrupted{} shows. Nor does it deny survival where psychology continues uniquely without the original organism, as in teletransportation. A criterion on which bodily continuity were necessary would give the opposite verdict there, and the enumeration contains \stbiodisN{} non-branching cases in which our criterion and the biological criterion disagree. The principle that decides between the two relations is the one stated in Section~\ref{s:branching}: the person is the organism in its personal phase, the organism's psychology unifies the phase, and when the psychology divides the organism does not. The psychological criterion says what unifies a person's stages; the tie-breaker says which of several equally good unifiers belongs to the person the stages were, and it answers by appeal to the organism the person was a phase of. That is one principle, not two.
\end{reply}

\begin{objection}[The thinking animal]\label{o:animal}
If \enquote{person} is a phase sortal of the organism, then a person is an organism in a certain phase, and must persist as organisms do, by the continuation of a life. The criterion says that the person survives teletransportation, though the organism does not. Either the person is not a phase of the organism, or the criterion is false. And if the person in the vat is a person distinct from the organism, there are two thinkers of the vat person's thoughts before the operation, the person and the organism \citep{Olson1997}.
\end{objection}

\begin{reply}
The objection runs together two questions, as PID-05-a, Proposition~4.1, shows: what unifies person-stages into one person, and what the thinker of a person's thoughts fundamentally is. On the perdurantist account of PID-11 there are two sums, the organism, unified by the continuation of a life, and the person, unified by the relation of Definition~\ref{d:crit}. In every ordinary case they share all their stages during the organism's personal phase; that is the sense in which \enquote{person} is a phase sortal of the organism, and it is what Theorem~\ref{t:ordinary} states. In teletransportation and in the vat they share their earlier stages and diverge. There is, at any time, one thinker of the thoughts then occurring, counted by its present stage; PID-05-a, Section~4, reads the thinking animal argument so, and shows that the first-person pronoun refers to the organism, the continuant whose persistence the substance sortal fixes. Our criterion does not say that you survive teletransportation in the sense in which PID-05 denies that you survive in a vat; it says that the person continues, and PID-05-a agrees.
\end{reply}

\begin{objection}[Double success]\label{o:double}
In fission each successor has everything that would have made her the original had the other transplant failed. On the criterion, the original then ceases to exist. But a double success cannot be a failure \citep[ch.~12]{Parfit1984}: if one successor would have been survival, two cannot be death. The \enquote{neither} verdict is therefore absurd, and so is any criterion that yields it.
\end{objection}

\begin{reply}
The verdict is not that fission is as bad as death, but that it is not identity. PID-03-a argues that fission produces two persons, each of whom has what matters in survival, and that identity is therefore not what matters; PID-06-a argues that connectedness, not identity, carries most of what matters, so that the original has nearly as much reason to care about each successor as about her own future. The objection assumes that identity and what matters coincide. The \enquote{neither} verdict is forced by the logic of identity, by Theorem~\ref{t:unique}, once a successor is not singled out by the continuation of the original's life, and the alternative verdicts, that she is one successor or both, are worse: the first is arbitrary, and the second gives up the absoluteness of identity or posits, with \citet{Lewis1976}, a second person who was present all along. We accept that the double success is not survival in the strict sense, and that it is nearly as good.
\end{reply}

\begin{objection}[Arbitrary kinds]\label{o:kind}
The kind condition draws a sharp line through a continuum. A brain can be replaced neuron by neuron with functionally equivalent devices; at what point does the realizer cease to be of the human kind? Principle~\ref{p:kind} has no answer, and so the criterion is either indeterminate or arbitrary in the cases that matter most to the future of persons \citep{Chalmers2010}.
\end{objection}

\begin{reply}
PID-14-a distinguishes two kinds of gradual replacement. If the devices are integrated into a living brain in a living organism, supplied by its circulation and regulated by its physiology, they become parts of the organism, as pacemakers and prosthetic joints do; the psychology remains realized in a living human organism, every step is along a life edge, and the person survives, by Proposition~\ref{p:sleep}'s reasoning applied to each step. If the functions are transferred step by step to an external system and the brain is dismantled, the end point is an upload, and by Proposition~\ref{p:kindprop} it is not the original. That there is no sharp step at which the original ceases shows that the transition is vague, not that she survives. By MET-18 the vagueness lies in the predicate \enquote{realized in a thing of the original's kind}, which has a sharp but unknowable boundary; it is not a vagueness in the world, and it does not make the criterion arbitrary. Every criterion that says anything about the persistence of persons faces the same structure: a vague relation, a sharp but unknowable threshold, and a clear verdict on either side.
\end{reply}

\begin{objection}[Life as an accident]\label{o:accident}
Whether a branch continues the life of the original's organism can depend on surgical accidents of no psychological significance: whether the surgeon leaves the right hemisphere or the left in the skull, or whether the replica is built in a new body or in the original's revived one. A criterion of personal identity should not make a person's survival turn on such facts.
\end{objection}

\begin{reply}
The facts in question are not accidents from the point of view of the thing whose survival is at issue. They decide which of two equally good psychological continuers is realized in the organism the person was a phase of, and that organism is, by PID-05, what the person fundamentally is. That it is a matter of surgical choice which hemisphere stays is no more troubling than that it is a matter of surgical choice which hemisphere is destroyed in a single transplant; in both cases the choice decides what happens to a thing, and the criterion reports it. What would be troubling is a criterion that made survival turn on facts that do not concern the person or her organism at all, such as the degree of a memory's preservation in a successor; the closest continuer theory does that, as case \caseBranchLineWeak{} shows, and Principle~\ref{p:tie} does not. Where no branch continues the life, the criterion gives the \enquote{neither} verdict rather than inventing a significance for the accident.
\end{reply}

\begin{objection}[The method of cases]\label{o:method}
The criterion is built on responses to imagined cases that will never occur, and \citet{Johnston1987} has argued that such responses are not reliable guides to the persistence conditions of persons. Worse, if our persistence conditions are partly fixed by our dispositions of identification \citep{Johnston2010}, there is no single answer to the question the criterion purports to answer, since different persons, or different cultures, could identify differently.
\end{objection}

\begin{reply}
The criterion does not rest on puzzle-case responses alone. Its central clauses are supported by the ordinary cases, where by Theorem~\ref{t:ordinary} it agrees with every rival, and by the logic of identity, which by Theorem~\ref{t:unique} it respects; the puzzle cases are used to choose among criteria that agree on ordinary cases, and the responses to them on which we rely, that the transplant recipient and the person on Earth in Branch-Line are the original, are among the most stable in the literature. Against the protean view, we grant that a person's dispositions of identification shape what she has reason to care about, which PID-06 traces to connectedness, and how she understands her life, which PID-09 calls narrative identity, co-authored by others and distinct from numerical identity. They do not fix numerical identity. If they did, a person could survive a fission as the successor she chose to identify with, and identity would follow choices that could themselves branch, as two successors can identify differently; the logic of identity would then require a tie-breaker of its own. That the continuation of a life is not subject to choice is part of what fits it for the role.
\end{reply}

\section{What the result does not establish}\label{s:limits}

The criterion says what makes person-stages stages of one person. It does not say what each of us fundamentally is, which PID-05 answers; it does not say what matters in survival, which PID-03 and PID-06 answer; and it does not say whether disembodied persons are possible, which PID-10 holds they are in principle, since a phase sortal can apply to things of more than one kind. If there were disembodied persons, Principle~\ref{p:kind} would require a separate statement of their kind, and the tie-breaker would need a counterpart of a life; we take no position on either.

The enumeration is exhaustive only up to its bound. It does not represent records that are stored and used more than once, since every link spans one phase. In such a case, a second replica activated later from the same record is continuous with the original through the record, but is not a candidate of the latest stage of the original, which by then is the first replica; on the natural extension of Definition~\ref{d:crit} it is a successor, not the original, which agrees with PID-02-a, Theorem~3.3, and preserves Theorem~\ref{t:local}. It does not represent quantities of matter, and so does not compute the material tie-breaker discussed in Section~\ref{s:branching}. It does not represent cases with more than two stages at a time, more than four stages in all or more than one stage realized in another kind. The general proofs of Theorems~\ref{t:unique}, \ref{t:tie}, \ref{t:ordinary} and~\ref{t:local} are not so limited; the enumeration checks them, and the particular verdicts, on the class it covers.

The criterion does not settle where the threshold of strong connection lies, nor where the boundary of the human kind lies, nor whether a given organism is alive at a given time. Each is a question for which the criterion takes an answer as input. It does not settle questions of responsibility, compensation or inheritance, which depend on what matters as well as on identity, and in which, as PID-07 observes of the childhood self, weak connectedness can weaken obligations that identity would otherwise sustain.

\section{Formalization}\label{s:formal}

The module \nolinkurl{Philosophy/PID/PID_01.lean} formalizes the logical skeleton of Sections~\ref{s:framework}--\ref{s:ordinary}. Stages, times, the original stage, edges, strong connection, sameness of kind and life continuation are primitives, introduced as axioms with docstrings citing Definitions~\ref{d:case}--\ref{d:life} and Principle~\ref{p:kind}. Assumption~\ref{a:life} is the axiom \texttt{life\_unique}. Candidates and contested steps are defined (Definitions~\ref{d:cand} and~\ref{d:contested}). Being a stage of the original is a primitive predicate \texttt{Same}, and the criterion of Definition~\ref{d:crit} is stated by four axioms: \texttt{same\_origin}, \texttt{principle\_P1} for the sufficiency half of Principle~\ref{p:psych}, \texttt{principle\_P3} for Principle~\ref{p:tie}, and \texttt{necessity} for the clause \enquote{nothing else}, which states the necessity half of Principle~\ref{p:psych} together with Principle~\ref{p:kind}. Stating the criterion by axioms rather than by an inductive definition keeps its philosophical content visible in the axiom list of each theorem. The module proves Lemma~\ref{l:parent}, Theorem~\ref{t:unique}, Corollary~\ref{c:chain}, Theorem~\ref{t:tie}, clauses (i) and (ii), Proposition~\ref{p:kindprop} and Theorem~\ref{t:ordinary}. The main theorem is:

\begin{lstlisting}
theorem uniqueness : ∀ n (y y' : Stage), Same y → Same y' →
    time y = n → time y' = n → y = y'
\end{lstlisting}

The checker reports (namespace prefix \texttt{Philosophy.PID\_01} omitted):

\begin{lstlisting}
'uniqueness' depends on axioms: [propext, Classical.choice,
 Quot.sound, Edge, Life, Same, SameKind, Stage, Strong,
 edge_time, life_unique, necessity, origin, origin_time,
 principle_P3, time]
\end{lstlisting}

\noindent and for Theorem~\ref{t:ordinary}:

\begin{lstlisting}
'theorem_6_2' depends on axioms: [Edge, Life, Same, SameKind,
 Stage, Strong, origin, principle_P1, same_origin]
\end{lstlisting}

The axiom lists show the structure of the argument. The uniqueness theorem does not depend on the sufficiency half of Principle~\ref{p:psych} (\texttt{principle\_P1}): it would hold for any criterion that makes a stage of the original a candidate of an earlier one and breaks ties by life. It depends essentially on \texttt{life\_unique}, which is Assumption~\ref{a:life}, and on \texttt{principle\_P3}, the tie-breaker: without them the theorem fails, since nothing would prevent both branches of a fission from being the original. The appeal to classical logic is the case distinction on whether two stages are identical. Theorem~\ref{t:ordinary} depends on the sufficiency half and not on the tie-breaker, as it should, since in ordinary cases the tie-breaker is idle; it is constructive. The particular verdicts of Proposition~\ref{p:verdicts} and of Appendix~\ref{s:appendix} are not formalized in Lean; they are computed by the program described in the appendix, which checks the theorems on every case of the enumerated class.

\section{Conclusion}\label{s:conclusion}

A person persists by psychological continuity through a cause of the right kind, by a non-branching process; where the process branches at a stage of the original, the branch that alone continues the life of that stage's organism is the original, and if none does, neither is. \enquote{Person} is a phase sortal of the organism, and the tie-breaker follows from that status. The criterion respects the logic of identity, agrees with the biological and bodily criteria in ordinary cases, never makes identity depend on what happens later, and yields the verdicts of PID-02, PID-03 and PID-14. An exhaustive enumeration of the \stcases{} case structures within the stated bound confirms each of these results case by case.

\begin{thebibliography}{99}\small
\bibitem[Baker(2000)]{Baker2000} Baker, L.~R. (2000). \emph{Persons and Bodies: A Constitution View}. Cambridge University Press.
\bibitem[Butler(1736)]{Butler1736} Butler, J. (1736). Of personal identity. Dissertation I appended to \emph{The Analogy of Religion, Natural and Revealed}. London.
\bibitem[Chalmers(2010)]{Chalmers2010} Chalmers, D.~J. (2010). The singularity: A philosophical analysis. \emph{Journal of Consciousness Studies}, 17(9--10), 7--65.
\bibitem[Johnston(1987)]{Johnston1987} Johnston, M. (1987). Human beings. \emph{The Journal of Philosophy}, 84(2), 59--83.
\bibitem[Hume(1739)]{Hume1739} Hume, D. (1739). \emph{A Treatise of Human Nature}. London. Cited by book, part and section.
\bibitem[Johnston(1997)]{Johnston1997} Johnston, M. (1997). Human concerns without superlative selves. In J. Dancy (ed.), \emph{Reading Parfit}. Blackwell.
\bibitem[Johnston(2007)]{Johnston2007} Johnston, M. (2007). \enquote{Human beings} revisited: My body is not an animal. In D. Zimmerman (ed.), \emph{Oxford Studies in Metaphysics}, Vol.~3. Oxford University Press.
\bibitem[Johnston(2010)]{Johnston2010} Johnston, M. (2010). \emph{Surviving Death}. Princeton University Press.
\bibitem[Lewis(1976)]{Lewis1976} Lewis, D. (1976). Survival and identity. In A.~O. Rorty (ed.), \emph{The Identities of Persons} (pp.~17--40). University of California Press.
\bibitem[Locke(1694)]{Locke1694} Locke, J. (1694). \emph{An Essay Concerning Human Understanding}, 2nd edn. London.
\bibitem[McMahan(2002)]{McMahan2002} McMahan, J. (2002). \emph{The Ethics of Killing: Problems at the Margins of Life}. Oxford University Press.
\bibitem[Noonan(1985)]{Noonan1985} Noonan, H.~W. (1985). The only $x$ and $y$ principle. \emph{Analysis}, 45(2), 79--83.
\bibitem[Nozick(1981)]{Nozick1981} Nozick, R. (1981). \emph{Philosophical Explanations}. Harvard University Press.
\bibitem[Olson(1997)]{Olson1997} Olson, E.~T. (1997). \emph{The Human Animal: Personal Identity without Psychology}. Oxford University Press.
\bibitem[Olson(2007)]{Olson2007} Olson, E.~T. (2007). \emph{What Are We? A Study in Personal Ontology}. Oxford University Press.
\bibitem[Parfit(1971)]{Parfit1971} Parfit, D. (1971). Personal identity. \emph{The Philosophical Review}, 80(1), 3--27.
\bibitem[Parfit(1984)]{Parfit1984} Parfit, D. (1984). \emph{Reasons and Persons}. Oxford University Press.
\bibitem[Parfit(2012)]{Parfit2012} Parfit, D. (2012). We are not human beings. \emph{Philosophy}, 87(1), 5--28.
\bibitem[Perry(1972)]{Perry1972} Perry, J. (1972). Can the self divide? \emph{The Journal of Philosophy}, 69(16), 463--488.
\bibitem[Reid(1785)]{Reid1785} Reid, T. (1785). \emph{Essays on the Intellectual Powers of Man}. Edinburgh.
\bibitem[Schechtman(1996)]{Schechtman1996} Schechtman, M. (1996). \emph{The Constitution of Selves}. Cornell University Press.
\bibitem[Schechtman(2014)]{Schechtman2014} Schechtman, M. (2014). \emph{Staying Alive: Personal Identity, Practical Concerns, and the Unity of a Life}. Oxford University Press.
\bibitem[Shoemaker(1963)]{Shoemaker1963} Shoemaker, S. (1963). \emph{Self-Knowledge and Self-Identity}. Cornell University Press.
\bibitem[Shoemaker(1970)]{Shoemaker1970} Shoemaker, S. (1970). Persons and their pasts. \emph{American Philosophical Quarterly}, 7(4), 269--285.
\bibitem[Shoemaker(1984)]{Shoemaker1984} Shoemaker, S. (1984). Personal identity: A materialist's account. In S. Shoemaker and R. Swinburne, \emph{Personal Identity} (pp.~67--132). Blackwell.
\bibitem[Snowdon(1990)]{Snowdon1990} Snowdon, P.~F. (1990). Persons, animals, and ourselves. In C. Gill (ed.), \emph{The Person and the Human Mind} (pp.~83--107). Clarendon Press.
\bibitem[Snowdon(2014)]{Snowdon2014} Snowdon, P.~F. (2014). \emph{Persons, Animals, Ourselves}. Oxford University Press.
\bibitem[Swinburne(1984)]{Swinburne1984} Swinburne, R. (1984). Personal identity: The dualist theory. In S. Shoemaker and R. Swinburne, \emph{Personal Identity}. Blackwell.
\bibitem[Unger(1990)]{Unger1990} Unger, P. (1990). \emph{Identity, Consciousness and Value}. Oxford University Press.
\bibitem[van Inwagen(1990)]{vanInwagen1990} van Inwagen, P. (1990). \emph{Material Beings}. Cornell University Press.
\bibitem[Wiggins(1980)]{Wiggins1980} Wiggins, D. (1980). \emph{Sameness and Substance}. Blackwell.
\bibitem[Williams(1957)]{Williams1957} Williams, B. (1957). Personal identity and individuation. \emph{Proceedings of the Aristotelian Society}, 57, 229--252.
\bibitem[Williams(1970)]{Williams1970} Williams, B. (1970). The self and the future. \emph{The Philosophical Review}, 79(2), 161--180.
\end{thebibliography}

\clearpage
\appendix
\section{The enumeration of case structures}\label{s:appendix}

\subsection*{What is enumerated}

The tables list every case structure in the sense of Definition~\ref{d:case} that satisfies the following bound: at most three times after the original; at most two stages at each time; at most four stages after the original in all; at most two parents for each stage; and at most one stage realized in a thing of another kind than the original's. Each stage is realized in a thing of the original's kind, written as a plain stage name, or of another kind, marked by an asterisk. Each edge carries a degree of connectedness $0$, $1$ or $2$ (Definition~\ref{d:conn}) and may be a life edge, written \texttt{L}. Three constraints apply. A life edge joins two stages of the original's kind, since a thing of another kind has no life (Definition~\ref{d:life}). An edge of degree $0$ is a life edge, since an edge records that the later stage's psychology or organism, or both, are produced by the earlier stage's, and an edge that carried neither would record nothing. And, by Assumption~\ref{a:life}, a stage has at most one incoming and at most one outgoing life edge. Every stage at a time after the original has at least one parent at the preceding time, so that every case issues from the one original stage.

\subsection*{Encoding and identifiers}

Stages at time $t$ are named \texttt{$t$a} and \texttt{$t$b}; the original is \texttt{o}. A stage is written with its parents and edge labels after the sign \texttt{<}: thus \texttt{1b<o2L} is a stage at time $1$ whose parent is the original, connected to degree $2$ along a life edge, and \texttt{2a<1a2,1b2L} is a stage at time $2$ with two parents, a fusion. Times are separated by a slash. Two structures are isomorphic if a relabelling of stages that preserves times carries the one onto the other, with its kinds and edge labels. Each isomorphism class is represented by its canonical form, the representative whose level-by-level encoding, with the stages at each time in sorted order, is lexicographically least; the program computes it by a search over the orderings of each level that is exact, since it keeps every ordering that attains the least encoding at each level. Cases are listed by height (the last time at which there is a stage), then by the number of stages, then by the number of stages of another kind and the number of edges, then by canonical form. The identifier is a letter for the height, \texttt{A}, \texttt{B} or \texttt{C}, followed by a serial number within the height.

\subsection*{What is computed}

The column \texttt{PID-01} gives the verdict of Definition~\ref{d:crit} on each later stage, in the order of the structure, with times separated by a slash: \texttt{I} identical with the original; \texttt{L} a candidate of a stage of the original that loses a contested step to a life-continuing candidate; \texttt{E} a candidate of a stage of the original in a contested step that no candidate wins (\enquote{neither}); \texttt{D} continuous with the original through a cause of the right kind, but descended from a stage that is not the original; \texttt{U} continuous with the original only through a realizer of another kind; \texttt{N} not psychologically continuous with the original. The column \texttt{bio} gives the verdict of the biological criterion, which in this encoding is also that of bodily continuity: \texttt{+} if the stage is reached from the original by life edges, \texttt{-} otherwise. The column \texttt{cl} classifies the case: \texttt{N} if it does not branch at any stage of the original; otherwise the letters \texttt{B} (a contested fission step resolved by a life-continuing candidate), \texttt{E} (a contested fission step with no life-continuing candidate) and \texttt{F} (a step contested by fusion) for the kinds of contest that occur. The column \texttt{rivals} lists, for each of three rival criteria, the stages on which its verdict about identity with the original differs from ours: \texttt{C} the closest continuer theory by degree of connectedness (in a contested step, the candidate strictly more closely connected than every rival wins, and otherwise none); \texttt{P} the pure non-branching criterion with the kind condition and without a tie-breaker; \texttt{W} the wide global criterion (continuity by any cause, with branching assessed over the whole case). Rivals that differ on the same stages are grouped, so that \texttt{CPW:1b} says that all three differ from ours on stage \texttt{1b} and only there. An empty entry means that all three agree with ours on every stage.

\subsection*{The program and the checks}

The program \texttt{gen/generate.py}, in the manuscript's directory, uses only the Python standard library and runs in a few seconds; its output is deterministic. It generates case structures level by level: from each canonical case of height $h$ it forms every admissible next level, with every choice of kinds, parents and edge labels, and reduces the result to canonical form. Every case of height $h+1$ within the bound arises in this way from its truncation at $h$, which is itself a case within the bound, and an isomorphism of the larger case restricts to an isomorphism of the truncations; so extending one representative of each class at height $h$ reaches every class at height $h+1$, and the enumeration is exhaustive up to the bound.

For every case the program computes the criterion of Definition~\ref{d:crit} forward in time, and asserts: (1) that at most one stage at each time is the original (Theorem~\ref{t:unique}); (2) that every stage of the original after $o$ has a parent that is a stage of the original, and that the times at which the original has a stage form an initial segment (Corollary~\ref{c:chain}); (3) that the forward computation agrees with a separate search for admissible paths (Proposition~\ref{p:path}); (4) that in every contested step at a stage of the original, a life-continuing candidate is the original and every other candidate is not, that if there is none no contested candidate is the original, and that a step contested only by fusion preserves identity exactly when it is a life edge (Theorem~\ref{t:tie}); (5) that every stage of the original is of the original's kind and continuous with it through a cause of the right kind (Proposition~\ref{p:kindprop}); (6) that at every ordinary stage the criterion and the biological criterion both give identity (Theorem~\ref{t:ordinary}); (7) that in a case that does not branch at any stage of the original, every disagreement between the criterion and the biological criterion is at a stage that is not ordinary (Proposition~\ref{p:diverge}); and (8) that truncating the case at any earlier time leaves the verdicts on all earlier stages unchanged (Theorem~\ref{t:local}). All eight assertions hold in every case.

\subsection*{Summary}

The enumeration contains \stcases{} cases: \sthone{} of height $1$, \sthtwo{} of height $2$ and \sththree{} of height $3$, with \ststages{} later stages in all. \stmachine{} cases contain a stage of another kind and \stfusionstruct{} contain a stage with two parents. By class: \stclsN{} do not branch at any stage of the original (\texttt{N}); \stclsB{} contain only contests resolved by a life-continuing candidate (\texttt{B}); \stclsE{} only fission contests resolved as \enquote{neither} (\texttt{E}); \stclsF{} only contests by fusion (\texttt{F}); \stclsBE{} both of the first two kinds (\texttt{BE}); and \stclsBF{} contests of the first and third kinds (\texttt{BF}). Verdicts on stages: \stvI{} \texttt{I}, \stvL{} \texttt{L}, \stvE{} \texttt{E}, \stvD{} \texttt{D}, \stvU{} \texttt{U} and \stvN{} \texttt{N}; \stordstages{} stages are ordinary. The closest continuer theory differs from the criterion in \stdisC{} cases, the criterion without a tie-breaker in \stdisP{}, and the wide global criterion in \stdisW{}; the wide global criterion violates temporal locality in \stwidenonlocal{} cases. The biological criterion differs from the criterion on some stage in \stbiodis{} cases, of which \stbiodisN{} do not branch at any stage of the original; in all of them the stages of disagreement are not ordinary. The cases discussed in Section~\ref{s:cases} are \caseOrdinary, \caseSleepLong, \caseTeleport, \caseBranchLine, \caseBranchLineLate, \caseBranchLineThenTele, \caseBranchLineWeak, \caseFission, \caseFissionUneven, \caseFissionLater, \caseUploadD, \caseUploadN, \caseUploadDownload, \caseFusionLife, \caseFusionNone{} and \caseTeleBranch; case \caseDisrupted{} is discussed in the reply to Objection~\ref{o:collapse}.

% Generated by generate.py; do not edit.
\subsection*{Height 1, 1 stage after the original (7 cases)}
\begin{casetable}
A0001 & 1a<o0L & N & + & N &  \\
A0002 & 1a<o1 & I & {-} & N &  \\
A0003 & 1a<o1L & I & + & N &  \\
A0004 & 1a<o2 & I & {-} & N &  \\
A0005 & 1a<o2L & I & + & N &  \\
A0006 & 1a*<o1 & U & {-} & N & W:1a \\
A0007 & 1a*<o2 & U & {-} & N & W:1a \\
\end{casetable}
\subsection*{Height 1, 2 stages after the original (19 cases)}
\begin{casetable}
A0008 & 1a<o0L 1b<o1 & NI & +{-} & N &  \\
A0009 & 1a<o0L 1b<o2 & NI & +{-} & N &  \\
A0010 & 1a<o1 1b<o1 & EE & {-}{-} & E &  \\
A0011 & 1a<o1 1b<o1L & LI & {-}+ & B & CPW:1b \\
A0012 & 1a<o1 1b<o2 & EE & {-}{-} & E & C:1b \\
A0013 & 1a<o1 1b<o2L & LI & {-}+ & B & PW:1b \\
A0014 & 1a<o1L 1b<o2 & IL & +{-} & B & C:1a,1b PW:1a \\
A0015 & 1a<o2 1b<o2 & EE & {-}{-} & E &  \\
A0016 & 1a<o2 1b<o2L & LI & {-}+ & B & CPW:1b \\
A0017 & 1a<o0L 1b*<o1 & NU & +{-} & N & W:1b \\
A0018 & 1a<o0L 1b*<o2 & NU & +{-} & N & W:1b \\
A0019 & 1a<o1 1b*<o1 & IU & {-}{-} & N & W:1a \\
A0020 & 1a<o1 1b*<o2 & IU & {-}{-} & N & W:1a \\
A0021 & 1a<o1L 1b*<o1 & IU & +{-} & N & W:1a \\
A0022 & 1a<o1L 1b*<o2 & IU & +{-} & N & W:1a \\
A0023 & 1a<o2 1b*<o1 & IU & {-}{-} & N & W:1a \\
A0024 & 1a<o2 1b*<o2 & IU & {-}{-} & N & W:1a \\
A0025 & 1a<o2L 1b*<o1 & IU & +{-} & N & W:1a \\
A0026 & 1a<o2L 1b*<o2 & IU & +{-} & N & W:1a \\
\end{casetable}
\subsection*{Height 2, 2 stages after the original (39 cases)}
\begin{casetable}
B0001 & 1a<o0L / 2a<1a0L & N/N & +/+ & N &  \\
B0002 & 1a<o0L / 2a<1a1 & N/N & +/{-} & N &  \\
B0003 & 1a<o0L / 2a<1a1L & N/N & +/+ & N &  \\
B0004 & 1a<o0L / 2a<1a2 & N/N & +/{-} & N &  \\
B0005 & 1a<o0L / 2a<1a2L & N/N & +/+ & N &  \\
B0006 & 1a<o1 / 2a<1a0L & I/N & {-}/{-} & N &  \\
B0007 & 1a<o1 / 2a<1a1 & I/I & {-}/{-} & N &  \\
B0008 & 1a<o1 / 2a<1a1L & I/I & {-}/{-} & N &  \\
B0009 & 1a<o1 / 2a<1a2 & I/I & {-}/{-} & N &  \\
B0010 & 1a<o1 / 2a<1a2L & I/I & {-}/{-} & N &  \\
B0011 & 1a<o1L / 2a<1a0L & I/N & +/+ & N &  \\
B0012 & 1a<o1L / 2a<1a1 & I/I & +/{-} & N &  \\
B0013 & 1a<o1L / 2a<1a1L & I/I & +/+ & N &  \\
B0014 & 1a<o1L / 2a<1a2 & I/I & +/{-} & N &  \\
B0015 & 1a<o1L / 2a<1a2L & I/I & +/+ & N &  \\
B0016 & 1a<o2 / 2a<1a0L & I/N & {-}/{-} & N &  \\
B0017 & 1a<o2 / 2a<1a1 & I/I & {-}/{-} & N &  \\
B0018 & 1a<o2 / 2a<1a1L & I/I & {-}/{-} & N &  \\
B0019 & 1a<o2 / 2a<1a2 & I/I & {-}/{-} & N &  \\
B0020 & 1a<o2 / 2a<1a2L & I/I & {-}/{-} & N &  \\
B0021 & 1a<o2L / 2a<1a0L & I/N & +/+ & N &  \\
B0022 & 1a<o2L / 2a<1a1 & I/I & +/{-} & N &  \\
B0023 & 1a<o2L / 2a<1a1L & I/I & +/+ & N &  \\
B0024 & 1a<o2L / 2a<1a2 & I/I & +/{-} & N &  \\
B0025 & 1a<o2L / 2a<1a2L & I/I & +/+ & N &  \\
B0026 & 1a<o0L / 2a*<1a1 & N/N & +/{-} & N &  \\
B0027 & 1a<o0L / 2a*<1a2 & N/N & +/{-} & N &  \\
B0028 & 1a<o1 / 2a*<1a1 & I/U & {-}/{-} & N & W:2a \\
B0029 & 1a<o1 / 2a*<1a2 & I/U & {-}/{-} & N & W:2a \\
B0030 & 1a<o1L / 2a*<1a1 & I/U & +/{-} & N & W:2a \\
B0031 & 1a<o1L / 2a*<1a2 & I/U & +/{-} & N & W:2a \\
B0032 & 1a<o2 / 2a*<1a1 & I/U & {-}/{-} & N & W:2a \\
B0033 & 1a<o2 / 2a*<1a2 & I/U & {-}/{-} & N & W:2a \\
B0034 & 1a<o2L / 2a*<1a1 & I/U & +/{-} & N & W:2a \\
B0035 & 1a<o2L / 2a*<1a2 & I/U & +/{-} & N & W:2a \\
B0036 & 1a*<o1 / 2a<1a1 & U/U & {-}/{-} & N & W:1a,2a \\
B0037 & 1a*<o1 / 2a<1a2 & U/U & {-}/{-} & N & W:1a,2a \\
B0038 & 1a*<o2 / 2a<1a1 & U/U & {-}/{-} & N & W:1a,2a \\
B0039 & 1a*<o2 / 2a<1a2 & U/U & {-}/{-} & N & W:1a,2a \\
\end{casetable}
\subsection*{Height 2, 3 stages after the original (547 cases)}
\begin{casetable}
B0040 & 1a<o0L / 2a<1a0L 2b<1a1 & N/NN & +/+{-} & N &  \\
B0041 & 1a<o0L / 2a<1a0L 2b<1a2 & N/NN & +/+{-} & N &  \\
B0042 & 1a<o0L / 2a<1a1 2b<1a1 & N/NN & +/{-}{-} & N &  \\
B0043 & 1a<o0L / 2a<1a1 2b<1a1L & N/NN & +/{-}+ & N &  \\
B0044 & 1a<o0L / 2a<1a1 2b<1a2 & N/NN & +/{-}{-} & N &  \\
B0045 & 1a<o0L / 2a<1a1 2b<1a2L & N/NN & +/{-}+ & N &  \\
B0046 & 1a<o0L / 2a<1a1L 2b<1a2 & N/NN & +/+{-} & N &  \\
B0047 & 1a<o0L / 2a<1a2 2b<1a2 & N/NN & +/{-}{-} & N &  \\
B0048 & 1a<o0L / 2a<1a2 2b<1a2L & N/NN & +/{-}+ & N &  \\
B0049 & 1a<o0L 1b<o1 / 2a<1a0L & NI/N & +{-}/+ & N &  \\
B0050 & 1a<o0L 1b<o1 / 2a<1a1 & NI/N & +{-}/{-} & N &  \\
B0051 & 1a<o0L 1b<o1 / 2a<1a1L & NI/N & +{-}/+ & N &  \\
B0052 & 1a<o0L 1b<o1 / 2a<1a2 & NI/N & +{-}/{-} & N &  \\
B0053 & 1a<o0L 1b<o1 / 2a<1a2L & NI/N & +{-}/+ & N &  \\
B0054 & 1a<o0L 1b<o1 / 2a<1b0L & NI/N & +{-}/{-} & N &  \\
B0055 & 1a<o0L 1b<o1 / 2a<1b1 & NI/I & +{-}/{-} & N &  \\
B0056 & 1a<o0L 1b<o1 / 2a<1b1L & NI/I & +{-}/{-} & N &  \\
B0057 & 1a<o0L 1b<o1 / 2a<1b2 & NI/I & +{-}/{-} & N &  \\
B0058 & 1a<o0L 1b<o1 / 2a<1b2L & NI/I & +{-}/{-} & N &  \\
B0059 & 1a<o0L 1b<o2 / 2a<1a0L & NI/N & +{-}/+ & N &  \\
B0060 & 1a<o0L 1b<o2 / 2a<1a1 & NI/N & +{-}/{-} & N &  \\
B0061 & 1a<o0L 1b<o2 / 2a<1a1L & NI/N & +{-}/+ & N &  \\
B0062 & 1a<o0L 1b<o2 / 2a<1a2 & NI/N & +{-}/{-} & N &  \\
B0063 & 1a<o0L 1b<o2 / 2a<1a2L & NI/N & +{-}/+ & N &  \\
B0064 & 1a<o0L 1b<o2 / 2a<1b0L & NI/N & +{-}/{-} & N &  \\
B0065 & 1a<o0L 1b<o2 / 2a<1b1 & NI/I & +{-}/{-} & N &  \\
B0066 & 1a<o0L 1b<o2 / 2a<1b1L & NI/I & +{-}/{-} & N &  \\
B0067 & 1a<o0L 1b<o2 / 2a<1b2 & NI/I & +{-}/{-} & N &  \\
B0068 & 1a<o0L 1b<o2 / 2a<1b2L & NI/I & +{-}/{-} & N &  \\
B0069 & 1a<o1 / 2a<1a0L 2b<1a1 & I/NI & {-}/{-}{-} & N &  \\
B0070 & 1a<o1 / 2a<1a0L 2b<1a2 & I/NI & {-}/{-}{-} & N &  \\
B0071 & 1a<o1 / 2a<1a1 2b<1a1 & I/EE & {-}/{-}{-} & E & W:1a \\
B0072 & 1a<o1 / 2a<1a1 2b<1a1L & I/LI & {-}/{-}{-} & B & CP:2b W:1a,2b \\
B0073 & 1a<o1 / 2a<1a1 2b<1a2 & I/EE & {-}/{-}{-} & E & C:2b W:1a \\
B0074 & 1a<o1 / 2a<1a1 2b<1a2L & I/LI & {-}/{-}{-} & B & P:2b W:1a,2b \\
B0075 & 1a<o1 / 2a<1a1L 2b<1a2 & I/IL & {-}/{-}{-} & B & C:2a,2b P:2a W:1a,2a \\
B0076 & 1a<o1 / 2a<1a2 2b<1a2 & I/EE & {-}/{-}{-} & E & W:1a \\
B0077 & 1a<o1 / 2a<1a2 2b<1a2L & I/LI & {-}/{-}{-} & B & CP:2b W:1a,2b \\
B0078 & 1a<o1 1b<o1 / 2a<1a0L & EE/N & {-}{-}/{-} & E &  \\
B0079 & 1a<o1 1b<o1 / 2a<1a1 & EE/D & {-}{-}/{-} & E &  \\
B0080 & 1a<o1 1b<o1 / 2a<1a1L & EE/D & {-}{-}/{-} & E &  \\
B0081 & 1a<o1 1b<o1 / 2a<1a2 & EE/D & {-}{-}/{-} & E &  \\
B0082 & 1a<o1 1b<o1 / 2a<1a2L & EE/D & {-}{-}/{-} & E &  \\
B0083 & 1a<o1 1b<o1L / 2a<1a0L & LI/N & {-}+/{-} & B & CPW:1b \\
B0084 & 1a<o1 1b<o1L / 2a<1a1 & LI/D & {-}+/{-} & B & CPW:1b \\
B0085 & 1a<o1 1b<o1L / 2a<1a1L & LI/D & {-}+/{-} & B & CPW:1b \\
B0086 & 1a<o1 1b<o1L / 2a<1a2 & LI/D & {-}+/{-} & B & CPW:1b \\
B0087 & 1a<o1 1b<o1L / 2a<1a2L & LI/D & {-}+/{-} & B & CPW:1b \\
B0088 & 1a<o1 1b<o1L / 2a<1b0L & LI/N & {-}+/+ & B & CPW:1b \\
B0089 & 1a<o1 1b<o1L / 2a<1b1 & LI/I & {-}+/{-} & B & CPW:1b,2a \\
B0090 & 1a<o1 1b<o1L / 2a<1b1L & LI/I & {-}+/+ & B & CPW:1b,2a \\
B0091 & 1a<o1 1b<o1L / 2a<1b2 & LI/I & {-}+/{-} & B & CPW:1b,2a \\
B0092 & 1a<o1 1b<o1L / 2a<1b2L & LI/I & {-}+/+ & B & CPW:1b,2a \\
B0093 & 1a<o1 1b<o2 / 2a<1a0L & EE/N & {-}{-}/{-} & E & C:1b \\
B0094 & 1a<o1 1b<o2 / 2a<1a1 & EE/D & {-}{-}/{-} & E & C:1b \\
B0095 & 1a<o1 1b<o2 / 2a<1a1L & EE/D & {-}{-}/{-} & E & C:1b \\
B0096 & 1a<o1 1b<o2 / 2a<1a2 & EE/D & {-}{-}/{-} & E & C:1b \\
B0097 & 1a<o1 1b<o2 / 2a<1a2L & EE/D & {-}{-}/{-} & E & C:1b \\
B0098 & 1a<o1 1b<o2 / 2a<1b0L & EE/N & {-}{-}/{-} & E & C:1b \\
B0099 & 1a<o1 1b<o2 / 2a<1b1 & EE/D & {-}{-}/{-} & E & C:1b,2a \\
B0100 & 1a<o1 1b<o2 / 2a<1b1L & EE/D & {-}{-}/{-} & E & C:1b,2a \\
B0101 & 1a<o1 1b<o2 / 2a<1b2 & EE/D & {-}{-}/{-} & E & C:1b,2a \\
B0102 & 1a<o1 1b<o2 / 2a<1b2L & EE/D & {-}{-}/{-} & E & C:1b,2a \\
B0103 & 1a<o1 1b<o2L / 2a<1a0L & LI/N & {-}+/{-} & B & PW:1b \\
B0104 & 1a<o1 1b<o2L / 2a<1a1 & LI/D & {-}+/{-} & B & PW:1b \\
B0105 & 1a<o1 1b<o2L / 2a<1a1L & LI/D & {-}+/{-} & B & PW:1b \\
B0106 & 1a<o1 1b<o2L / 2a<1a2 & LI/D & {-}+/{-} & B & PW:1b \\
B0107 & 1a<o1 1b<o2L / 2a<1a2L & LI/D & {-}+/{-} & B & PW:1b \\
B0108 & 1a<o1 1b<o2L / 2a<1b0L & LI/N & {-}+/+ & B & PW:1b \\
B0109 & 1a<o1 1b<o2L / 2a<1b1 & LI/I & {-}+/{-} & B & PW:1b,2a \\
B0110 & 1a<o1 1b<o2L / 2a<1b1L & LI/I & {-}+/+ & B & PW:1b,2a \\
B0111 & 1a<o1 1b<o2L / 2a<1b2 & LI/I & {-}+/{-} & B & PW:1b,2a \\
B0112 & 1a<o1 1b<o2L / 2a<1b2L & LI/I & {-}+/+ & B & PW:1b,2a \\
B0113 & 1a<o1L / 2a<1a0L 2b<1a1 & I/NI & +/+{-} & N &  \\
B0114 & 1a<o1L / 2a<1a0L 2b<1a2 & I/NI & +/+{-} & N &  \\
B0115 & 1a<o1L / 2a<1a1 2b<1a1 & I/EE & +/{-}{-} & E & W:1a \\
B0116 & 1a<o1L / 2a<1a1 2b<1a1L & I/LI & +/{-}+ & B & CP:2b W:1a,2b \\
B0117 & 1a<o1L / 2a<1a1 2b<1a2 & I/EE & +/{-}{-} & E & C:2b W:1a \\
B0118 & 1a<o1L / 2a<1a1 2b<1a2L & I/LI & +/{-}+ & B & P:2b W:1a,2b \\
B0119 & 1a<o1L / 2a<1a1L 2b<1a2 & I/IL & +/+{-} & B & C:2a,2b P:2a W:1a,2a \\
B0120 & 1a<o1L / 2a<1a2 2b<1a2 & I/EE & +/{-}{-} & E & W:1a \\
B0121 & 1a<o1L / 2a<1a2 2b<1a2L & I/LI & +/{-}+ & B & CP:2b W:1a,2b \\
B0122 & 1a<o1L 1b<o2 / 2a<1a0L & IL/N & +{-}/+ & B & C:1a,1b PW:1a \\
B0123 & 1a<o1L 1b<o2 / 2a<1a1 & IL/I & +{-}/{-} & B & C:1a,1b,2a PW:1a,2a \\
B0124 & 1a<o1L 1b<o2 / 2a<1a1L & IL/I & +{-}/+ & B & C:1a,1b,2a PW:1a,2a \\
B0125 & 1a<o1L 1b<o2 / 2a<1a2 & IL/I & +{-}/{-} & B & C:1a,1b,2a PW:1a,2a \\
B0126 & 1a<o1L 1b<o2 / 2a<1a2L & IL/I & +{-}/+ & B & C:1a,1b,2a PW:1a,2a \\
B0127 & 1a<o1L 1b<o2 / 2a<1b0L & IL/N & +{-}/{-} & B & C:1a,1b PW:1a \\
B0128 & 1a<o1L 1b<o2 / 2a<1b1 & IL/D & +{-}/{-} & B & C:1a,1b,2a PW:1a \\
B0129 & 1a<o1L 1b<o2 / 2a<1b1L & IL/D & +{-}/{-} & B & C:1a,1b,2a PW:1a \\
B0130 & 1a<o1L 1b<o2 / 2a<1b2 & IL/D & +{-}/{-} & B & C:1a,1b,2a PW:1a \\
B0131 & 1a<o1L 1b<o2 / 2a<1b2L & IL/D & +{-}/{-} & B & C:1a,1b,2a PW:1a \\
B0132 & 1a<o2 / 2a<1a0L 2b<1a1 & I/NI & {-}/{-}{-} & N &  \\
B0133 & 1a<o2 / 2a<1a0L 2b<1a2 & I/NI & {-}/{-}{-} & N &  \\
B0134 & 1a<o2 / 2a<1a1 2b<1a1 & I/EE & {-}/{-}{-} & E & W:1a \\
B0135 & 1a<o2 / 2a<1a1 2b<1a1L & I/LI & {-}/{-}{-} & B & CP:2b W:1a,2b \\
B0136 & 1a<o2 / 2a<1a1 2b<1a2 & I/EE & {-}/{-}{-} & E & C:2b W:1a \\
B0137 & 1a<o2 / 2a<1a1 2b<1a2L & I/LI & {-}/{-}{-} & B & P:2b W:1a,2b \\
B0138 & 1a<o2 / 2a<1a1L 2b<1a2 & I/IL & {-}/{-}{-} & B & C:2a,2b P:2a W:1a,2a \\
B0139 & 1a<o2 / 2a<1a2 2b<1a2 & I/EE & {-}/{-}{-} & E & W:1a \\
B0140 & 1a<o2 / 2a<1a2 2b<1a2L & I/LI & {-}/{-}{-} & B & CP:2b W:1a,2b \\
B0141 & 1a<o2 1b<o2 / 2a<1a0L & EE/N & {-}{-}/{-} & E &  \\
B0142 & 1a<o2 1b<o2 / 2a<1a1 & EE/D & {-}{-}/{-} & E &  \\
B0143 & 1a<o2 1b<o2 / 2a<1a1L & EE/D & {-}{-}/{-} & E &  \\
B0144 & 1a<o2 1b<o2 / 2a<1a2 & EE/D & {-}{-}/{-} & E &  \\
B0145 & 1a<o2 1b<o2 / 2a<1a2L & EE/D & {-}{-}/{-} & E &  \\
B0146 & 1a<o2 1b<o2L / 2a<1a0L & LI/N & {-}+/{-} & B & CPW:1b \\
B0147 & 1a<o2 1b<o2L / 2a<1a1 & LI/D & {-}+/{-} & B & CPW:1b \\
B0148 & 1a<o2 1b<o2L / 2a<1a1L & LI/D & {-}+/{-} & B & CPW:1b \\
B0149 & 1a<o2 1b<o2L / 2a<1a2 & LI/D & {-}+/{-} & B & CPW:1b \\
B0150 & 1a<o2 1b<o2L / 2a<1a2L & LI/D & {-}+/{-} & B & CPW:1b \\
B0151 & 1a<o2 1b<o2L / 2a<1b0L & LI/N & {-}+/+ & B & CPW:1b \\
B0152 & 1a<o2 1b<o2L / 2a<1b1 & LI/I & {-}+/{-} & B & CPW:1b,2a \\
B0153 & 1a<o2 1b<o2L / 2a<1b1L & LI/I & {-}+/+ & B & CPW:1b,2a \\
B0154 & 1a<o2 1b<o2L / 2a<1b2 & LI/I & {-}+/{-} & B & CPW:1b,2a \\
B0155 & 1a<o2 1b<o2L / 2a<1b2L & LI/I & {-}+/+ & B & CPW:1b,2a \\
B0156 & 1a<o2L / 2a<1a0L 2b<1a1 & I/NI & +/+{-} & N &  \\
B0157 & 1a<o2L / 2a<1a0L 2b<1a2 & I/NI & +/+{-} & N &  \\
B0158 & 1a<o2L / 2a<1a1 2b<1a1 & I/EE & +/{-}{-} & E & W:1a \\
B0159 & 1a<o2L / 2a<1a1 2b<1a1L & I/LI & +/{-}+ & B & CP:2b W:1a,2b \\
B0160 & 1a<o2L / 2a<1a1 2b<1a2 & I/EE & +/{-}{-} & E & C:2b W:1a \\
B0161 & 1a<o2L / 2a<1a1 2b<1a2L & I/LI & +/{-}+ & B & P:2b W:1a,2b \\
B0162 & 1a<o2L / 2a<1a1L 2b<1a2 & I/IL & +/+{-} & B & C:2a,2b P:2a W:1a,2a \\
B0163 & 1a<o2L / 2a<1a2 2b<1a2 & I/EE & +/{-}{-} & E & W:1a \\
B0164 & 1a<o2L / 2a<1a2 2b<1a2L & I/LI & +/{-}+ & B & CP:2b W:1a,2b \\
B0165 & 1a<o0L 1b<o1 / 2a<1a0L,1b1 & NI/I & +{-}/+ & N &  \\
B0166 & 1a<o0L 1b<o1 / 2a<1a0L,1b2 & NI/I & +{-}/+ & N &  \\
B0167 & 1a<o0L 1b<o1 / 2a<1a1,1b0L & NI/N & +{-}/{-} & N &  \\
B0168 & 1a<o0L 1b<o1 / 2a<1a1,1b1 & NI/E & +{-}/{-} & F & W:2a \\
B0169 & 1a<o0L 1b<o1 / 2a<1a1,1b1L & NI/I & +{-}/{-} & F & CP:2a \\
B0170 & 1a<o0L 1b<o1 / 2a<1a1,1b2 & NI/E & +{-}/{-} & F & CW:2a \\
B0171 & 1a<o0L 1b<o1 / 2a<1a1,1b2L & NI/I & +{-}/{-} & F & P:2a \\
B0172 & 1a<o0L 1b<o1 / 2a<1a1L,1b1 & NI/E & +{-}/+ & F & W:2a \\
B0173 & 1a<o0L 1b<o1 / 2a<1a1L,1b2 & NI/E & +{-}/+ & F & CW:2a \\
B0174 & 1a<o0L 1b<o1 / 2a<1a2,1b0L & NI/N & +{-}/{-} & N &  \\
B0175 & 1a<o0L 1b<o1 / 2a<1a2,1b1 & NI/E & +{-}/{-} & F & W:2a \\
B0176 & 1a<o0L 1b<o1 / 2a<1a2,1b1L & NI/I & +{-}/{-} & F & CP:2a \\
B0177 & 1a<o0L 1b<o1 / 2a<1a2,1b2 & NI/E & +{-}/{-} & F & W:2a \\
B0178 & 1a<o0L 1b<o1 / 2a<1a2,1b2L & NI/I & +{-}/{-} & F & CP:2a \\
B0179 & 1a<o0L 1b<o1 / 2a<1a2L,1b1 & NI/E & +{-}/+ & F & W:2a \\
B0180 & 1a<o0L 1b<o1 / 2a<1a2L,1b2 & NI/E & +{-}/+ & F & W:2a \\
B0181 & 1a<o0L 1b<o2 / 2a<1a0L,1b1 & NI/I & +{-}/+ & N &  \\
B0182 & 1a<o0L 1b<o2 / 2a<1a0L,1b2 & NI/I & +{-}/+ & N &  \\
B0183 & 1a<o0L 1b<o2 / 2a<1a1,1b0L & NI/N & +{-}/{-} & N &  \\
B0184 & 1a<o0L 1b<o2 / 2a<1a1,1b1 & NI/E & +{-}/{-} & F & W:2a \\
B0185 & 1a<o0L 1b<o2 / 2a<1a1,1b1L & NI/I & +{-}/{-} & F & CP:2a \\
B0186 & 1a<o0L 1b<o2 / 2a<1a1,1b2 & NI/E & +{-}/{-} & F & CW:2a \\
B0187 & 1a<o0L 1b<o2 / 2a<1a1,1b2L & NI/I & +{-}/{-} & F & P:2a \\
B0188 & 1a<o0L 1b<o2 / 2a<1a1L,1b1 & NI/E & +{-}/+ & F & W:2a \\
B0189 & 1a<o0L 1b<o2 / 2a<1a1L,1b2 & NI/E & +{-}/+ & F & CW:2a \\
B0190 & 1a<o0L 1b<o2 / 2a<1a2,1b0L & NI/N & +{-}/{-} & N &  \\
B0191 & 1a<o0L 1b<o2 / 2a<1a2,1b1 & NI/E & +{-}/{-} & F & W:2a \\
B0192 & 1a<o0L 1b<o2 / 2a<1a2,1b1L & NI/I & +{-}/{-} & F & CP:2a \\
B0193 & 1a<o0L 1b<o2 / 2a<1a2,1b2 & NI/E & +{-}/{-} & F & W:2a \\
B0194 & 1a<o0L 1b<o2 / 2a<1a2,1b2L & NI/I & +{-}/{-} & F & CP:2a \\
B0195 & 1a<o0L 1b<o2 / 2a<1a2L,1b1 & NI/E & +{-}/+ & F & W:2a \\
B0196 & 1a<o0L 1b<o2 / 2a<1a2L,1b2 & NI/E & +{-}/+ & F & W:2a \\
B0197 & 1a<o1 1b<o1 / 2a<1a0L,1b1 & EE/D & {-}{-}/{-} & E &  \\
B0198 & 1a<o1 1b<o1 / 2a<1a0L,1b2 & EE/D & {-}{-}/{-} & E &  \\
B0199 & 1a<o1 1b<o1 / 2a<1a1,1b1 & EE/D & {-}{-}/{-} & E &  \\
B0200 & 1a<o1 1b<o1 / 2a<1a1,1b1L & EE/D & {-}{-}/{-} & E &  \\
B0201 & 1a<o1 1b<o1 / 2a<1a1,1b2 & EE/D & {-}{-}/{-} & E &  \\
B0202 & 1a<o1 1b<o1 / 2a<1a1,1b2L & EE/D & {-}{-}/{-} & E &  \\
B0203 & 1a<o1 1b<o1 / 2a<1a1L,1b2 & EE/D & {-}{-}/{-} & E &  \\
B0204 & 1a<o1 1b<o1 / 2a<1a2,1b2 & EE/D & {-}{-}/{-} & E &  \\
B0205 & 1a<o1 1b<o1 / 2a<1a2,1b2L & EE/D & {-}{-}/{-} & E &  \\
B0206 & 1a<o1 1b<o1L / 2a<1a0L,1b1 & LI/I & {-}+/{-} & B & CPW:1b,2a \\
B0207 & 1a<o1 1b<o1L / 2a<1a0L,1b2 & LI/I & {-}+/{-} & B & CPW:1b,2a \\
B0208 & 1a<o1 1b<o1L / 2a<1a1,1b0L & LI/D & {-}+/+ & B & CPW:1b \\
B0209 & 1a<o1 1b<o1L / 2a<1a1,1b1 & LI/E & {-}+/{-} & BF & CPW:1b \\
B0210 & 1a<o1 1b<o1L / 2a<1a1,1b1L & LI/I & {-}+/+ & BF & CPW:1b,2a \\
B0211 & 1a<o1 1b<o1L / 2a<1a1,1b2 & LI/E & {-}+/{-} & BF & CPW:1b \\
B0212 & 1a<o1 1b<o1L / 2a<1a1,1b2L & LI/I & {-}+/+ & BF & CPW:1b,2a \\
B0213 & 1a<o1 1b<o1L / 2a<1a1L,1b1 & LI/E & {-}+/{-} & BF & CPW:1b \\
B0214 & 1a<o1 1b<o1L / 2a<1a1L,1b2 & LI/E & {-}+/{-} & BF & CPW:1b \\
B0215 & 1a<o1 1b<o1L / 2a<1a2,1b0L & LI/D & {-}+/+ & B & CPW:1b \\
B0216 & 1a<o1 1b<o1L / 2a<1a2,1b1 & LI/E & {-}+/{-} & BF & CPW:1b \\
B0217 & 1a<o1 1b<o1L / 2a<1a2,1b1L & LI/I & {-}+/+ & BF & CPW:1b,2a \\
B0218 & 1a<o1 1b<o1L / 2a<1a2,1b2 & LI/E & {-}+/{-} & BF & CPW:1b \\
B0219 & 1a<o1 1b<o1L / 2a<1a2,1b2L & LI/I & {-}+/+ & BF & CPW:1b,2a \\
B0220 & 1a<o1 1b<o1L / 2a<1a2L,1b1 & LI/E & {-}+/{-} & BF & CPW:1b \\
B0221 & 1a<o1 1b<o1L / 2a<1a2L,1b2 & LI/E & {-}+/{-} & BF & CPW:1b \\
B0222 & 1a<o1 1b<o2 / 2a<1a0L,1b1 & EE/D & {-}{-}/{-} & E & C:1b,2a \\
B0223 & 1a<o1 1b<o2 / 2a<1a0L,1b2 & EE/D & {-}{-}/{-} & E & C:1b,2a \\
B0224 & 1a<o1 1b<o2 / 2a<1a1,1b0L & EE/D & {-}{-}/{-} & E & C:1b \\
B0225 & 1a<o1 1b<o2 / 2a<1a1,1b1 & EE/D & {-}{-}/{-} & E & C:1b \\
B0226 & 1a<o1 1b<o2 / 2a<1a1,1b1L & EE/D & {-}{-}/{-} & E & C:1b \\
B0227 & 1a<o1 1b<o2 / 2a<1a1,1b2 & EE/D & {-}{-}/{-} & E & C:1b,2a \\
B0228 & 1a<o1 1b<o2 / 2a<1a1,1b2L & EE/D & {-}{-}/{-} & E & C:1b,2a \\
B0229 & 1a<o1 1b<o2 / 2a<1a1L,1b1 & EE/D & {-}{-}/{-} & E & C:1b \\
B0230 & 1a<o1 1b<o2 / 2a<1a1L,1b2 & EE/D & {-}{-}/{-} & E & C:1b,2a \\
B0231 & 1a<o1 1b<o2 / 2a<1a2,1b0L & EE/D & {-}{-}/{-} & E & C:1b \\
B0232 & 1a<o1 1b<o2 / 2a<1a2,1b1 & EE/D & {-}{-}/{-} & E & C:1b \\
B0233 & 1a<o1 1b<o2 / 2a<1a2,1b1L & EE/D & {-}{-}/{-} & E & C:1b \\
B0234 & 1a<o1 1b<o2 / 2a<1a2,1b2 & EE/D & {-}{-}/{-} & E & C:1b \\
B0235 & 1a<o1 1b<o2 / 2a<1a2,1b2L & EE/D & {-}{-}/{-} & E & C:1b \\
B0236 & 1a<o1 1b<o2 / 2a<1a2L,1b1 & EE/D & {-}{-}/{-} & E & C:1b \\
B0237 & 1a<o1 1b<o2 / 2a<1a2L,1b2 & EE/D & {-}{-}/{-} & E & C:1b \\
B0238 & 1a<o1 1b<o2L / 2a<1a0L,1b1 & LI/I & {-}+/{-} & B & PW:1b,2a \\
B0239 & 1a<o1 1b<o2L / 2a<1a0L,1b2 & LI/I & {-}+/{-} & B & PW:1b,2a \\
B0240 & 1a<o1 1b<o2L / 2a<1a1,1b0L & LI/D & {-}+/+ & B & PW:1b \\
B0241 & 1a<o1 1b<o2L / 2a<1a1,1b1 & LI/E & {-}+/{-} & BF & PW:1b \\
B0242 & 1a<o1 1b<o2L / 2a<1a1,1b1L & LI/I & {-}+/+ & BF & C:2a PW:1b,2a \\
B0243 & 1a<o1 1b<o2L / 2a<1a1,1b2 & LI/E & {-}+/{-} & BF & C:2a PW:1b \\
B0244 & 1a<o1 1b<o2L / 2a<1a1,1b2L & LI/I & {-}+/+ & BF & PW:1b,2a \\
B0245 & 1a<o1 1b<o2L / 2a<1a1L,1b1 & LI/E & {-}+/{-} & BF & PW:1b \\
B0246 & 1a<o1 1b<o2L / 2a<1a1L,1b2 & LI/E & {-}+/{-} & BF & C:2a PW:1b \\
B0247 & 1a<o1 1b<o2L / 2a<1a2,1b0L & LI/D & {-}+/+ & B & PW:1b \\
B0248 & 1a<o1 1b<o2L / 2a<1a2,1b1 & LI/E & {-}+/{-} & BF & PW:1b \\
B0249 & 1a<o1 1b<o2L / 2a<1a2,1b1L & LI/I & {-}+/+ & BF & C:2a PW:1b,2a \\
B0250 & 1a<o1 1b<o2L / 2a<1a2,1b2 & LI/E & {-}+/{-} & BF & PW:1b \\
B0251 & 1a<o1 1b<o2L / 2a<1a2,1b2L & LI/I & {-}+/+ & BF & C:2a PW:1b,2a \\
B0252 & 1a<o1 1b<o2L / 2a<1a2L,1b1 & LI/E & {-}+/{-} & BF & PW:1b \\
B0253 & 1a<o1 1b<o2L / 2a<1a2L,1b2 & LI/E & {-}+/{-} & BF & PW:1b \\
B0254 & 1a<o1L 1b<o2 / 2a<1a0L,1b1 & IL/D & +{-}/+ & B & C:1a,1b,2a PW:1a \\
B0255 & 1a<o1L 1b<o2 / 2a<1a0L,1b2 & IL/D & +{-}/+ & B & C:1a,1b,2a PW:1a \\
B0256 & 1a<o1L 1b<o2 / 2a<1a1,1b0L & IL/I & +{-}/{-} & B & C:1a,1b,2a PW:1a,2a \\
B0257 & 1a<o1L 1b<o2 / 2a<1a1,1b1 & IL/E & +{-}/{-} & BF & C:1a,1b PW:1a \\
B0258 & 1a<o1L 1b<o2 / 2a<1a1,1b1L & IL/E & +{-}/{-} & BF & C:1a,1b PW:1a \\
B0259 & 1a<o1L 1b<o2 / 2a<1a1,1b2 & IL/E & +{-}/{-} & BF & C:1a,1b,2a PW:1a \\
B0260 & 1a<o1L 1b<o2 / 2a<1a1,1b2L & IL/E & +{-}/{-} & BF & C:1a,1b,2a PW:1a \\
B0261 & 1a<o1L 1b<o2 / 2a<1a1L,1b1 & IL/I & +{-}/+ & BF & C:1a,1b,2a PW:1a,2a \\
B0262 & 1a<o1L 1b<o2 / 2a<1a1L,1b2 & IL/I & +{-}/+ & BF & C:1a,1b PW:1a,2a \\
B0263 & 1a<o1L 1b<o2 / 2a<1a2,1b0L & IL/I & +{-}/{-} & B & C:1a,1b,2a PW:1a,2a \\
B0264 & 1a<o1L 1b<o2 / 2a<1a2,1b1 & IL/E & +{-}/{-} & BF & C:1a,1b PW:1a \\
B0265 & 1a<o1L 1b<o2 / 2a<1a2,1b1L & IL/E & +{-}/{-} & BF & C:1a,1b PW:1a \\
B0266 & 1a<o1L 1b<o2 / 2a<1a2,1b2 & IL/E & +{-}/{-} & BF & C:1a,1b PW:1a \\
B0267 & 1a<o1L 1b<o2 / 2a<1a2,1b2L & IL/E & +{-}/{-} & BF & C:1a,1b PW:1a \\
B0268 & 1a<o1L 1b<o2 / 2a<1a2L,1b1 & IL/I & +{-}/+ & BF & C:1a,1b,2a PW:1a,2a \\
B0269 & 1a<o1L 1b<o2 / 2a<1a2L,1b2 & IL/I & +{-}/+ & BF & C:1a,1b,2a PW:1a,2a \\
B0270 & 1a<o2 1b<o2 / 2a<1a0L,1b1 & EE/D & {-}{-}/{-} & E &  \\
B0271 & 1a<o2 1b<o2 / 2a<1a0L,1b2 & EE/D & {-}{-}/{-} & E &  \\
B0272 & 1a<o2 1b<o2 / 2a<1a1,1b1 & EE/D & {-}{-}/{-} & E &  \\
B0273 & 1a<o2 1b<o2 / 2a<1a1,1b1L & EE/D & {-}{-}/{-} & E &  \\
B0274 & 1a<o2 1b<o2 / 2a<1a1,1b2 & EE/D & {-}{-}/{-} & E &  \\
B0275 & 1a<o2 1b<o2 / 2a<1a1,1b2L & EE/D & {-}{-}/{-} & E &  \\
B0276 & 1a<o2 1b<o2 / 2a<1a1L,1b2 & EE/D & {-}{-}/{-} & E &  \\
B0277 & 1a<o2 1b<o2 / 2a<1a2,1b2 & EE/D & {-}{-}/{-} & E &  \\
B0278 & 1a<o2 1b<o2 / 2a<1a2,1b2L & EE/D & {-}{-}/{-} & E &  \\
B0279 & 1a<o2 1b<o2L / 2a<1a0L,1b1 & LI/I & {-}+/{-} & B & CPW:1b,2a \\
B0280 & 1a<o2 1b<o2L / 2a<1a0L,1b2 & LI/I & {-}+/{-} & B & CPW:1b,2a \\
B0281 & 1a<o2 1b<o2L / 2a<1a1,1b0L & LI/D & {-}+/+ & B & CPW:1b \\
B0282 & 1a<o2 1b<o2L / 2a<1a1,1b1 & LI/E & {-}+/{-} & BF & CPW:1b \\
B0283 & 1a<o2 1b<o2L / 2a<1a1,1b1L & LI/I & {-}+/+ & BF & CPW:1b,2a \\
B0284 & 1a<o2 1b<o2L / 2a<1a1,1b2 & LI/E & {-}+/{-} & BF & CPW:1b \\
B0285 & 1a<o2 1b<o2L / 2a<1a1,1b2L & LI/I & {-}+/+ & BF & CPW:1b,2a \\
B0286 & 1a<o2 1b<o2L / 2a<1a1L,1b1 & LI/E & {-}+/{-} & BF & CPW:1b \\
B0287 & 1a<o2 1b<o2L / 2a<1a1L,1b2 & LI/E & {-}+/{-} & BF & CPW:1b \\
B0288 & 1a<o2 1b<o2L / 2a<1a2,1b0L & LI/D & {-}+/+ & B & CPW:1b \\
B0289 & 1a<o2 1b<o2L / 2a<1a2,1b1 & LI/E & {-}+/{-} & BF & CPW:1b \\
B0290 & 1a<o2 1b<o2L / 2a<1a2,1b1L & LI/I & {-}+/+ & BF & CPW:1b,2a \\
B0291 & 1a<o2 1b<o2L / 2a<1a2,1b2 & LI/E & {-}+/{-} & BF & CPW:1b \\
B0292 & 1a<o2 1b<o2L / 2a<1a2,1b2L & LI/I & {-}+/+ & BF & CPW:1b,2a \\
B0293 & 1a<o2 1b<o2L / 2a<1a2L,1b1 & LI/E & {-}+/{-} & BF & CPW:1b \\
B0294 & 1a<o2 1b<o2L / 2a<1a2L,1b2 & LI/E & {-}+/{-} & BF & CPW:1b \\
B0295 & 1a<o0L / 2a<1a0L 2b*<1a1 & N/NN & +/+{-} & N &  \\
B0296 & 1a<o0L / 2a<1a0L 2b*<1a2 & N/NN & +/+{-} & N &  \\
B0297 & 1a<o0L / 2a<1a1 2b*<1a1 & N/NN & +/{-}{-} & N &  \\
B0298 & 1a<o0L / 2a<1a1 2b*<1a2 & N/NN & +/{-}{-} & N &  \\
B0299 & 1a<o0L / 2a<1a1L 2b*<1a1 & N/NN & +/+{-} & N &  \\
B0300 & 1a<o0L / 2a<1a1L 2b*<1a2 & N/NN & +/+{-} & N &  \\
B0301 & 1a<o0L / 2a<1a2 2b*<1a1 & N/NN & +/{-}{-} & N &  \\
B0302 & 1a<o0L / 2a<1a2 2b*<1a2 & N/NN & +/{-}{-} & N &  \\
B0303 & 1a<o0L / 2a<1a2L 2b*<1a1 & N/NN & +/+{-} & N &  \\
B0304 & 1a<o0L / 2a<1a2L 2b*<1a2 & N/NN & +/+{-} & N &  \\
B0305 & 1a<o0L 1b<o1 / 2a*<1a1 & NI/N & +{-}/{-} & N &  \\
B0306 & 1a<o0L 1b<o1 / 2a*<1a2 & NI/N & +{-}/{-} & N &  \\
B0307 & 1a<o0L 1b<o1 / 2a*<1b1 & NI/U & +{-}/{-} & N & W:2a \\
B0308 & 1a<o0L 1b<o1 / 2a*<1b2 & NI/U & +{-}/{-} & N & W:2a \\
B0309 & 1a<o0L 1b<o2 / 2a*<1a1 & NI/N & +{-}/{-} & N &  \\
B0310 & 1a<o0L 1b<o2 / 2a*<1a2 & NI/N & +{-}/{-} & N &  \\
B0311 & 1a<o0L 1b<o2 / 2a*<1b1 & NI/U & +{-}/{-} & N & W:2a \\
B0312 & 1a<o0L 1b<o2 / 2a*<1b2 & NI/U & +{-}/{-} & N & W:2a \\
B0313 & 1a<o0L 1b*<o1 / 2a<1a0L & NU/N & +{-}/+ & N & W:1b \\
B0314 & 1a<o0L 1b*<o1 / 2a<1a1 & NU/N & +{-}/{-} & N & W:1b \\
B0315 & 1a<o0L 1b*<o1 / 2a<1a1L & NU/N & +{-}/+ & N & W:1b \\
B0316 & 1a<o0L 1b*<o1 / 2a<1a2 & NU/N & +{-}/{-} & N & W:1b \\
B0317 & 1a<o0L 1b*<o1 / 2a<1a2L & NU/N & +{-}/+ & N & W:1b \\
B0318 & 1a<o0L 1b*<o1 / 2a<1b1 & NU/U & +{-}/{-} & N & W:1b,2a \\
B0319 & 1a<o0L 1b*<o1 / 2a<1b2 & NU/U & +{-}/{-} & N & W:1b,2a \\
B0320 & 1a<o0L 1b*<o2 / 2a<1a0L & NU/N & +{-}/+ & N & W:1b \\
B0321 & 1a<o0L 1b*<o2 / 2a<1a1 & NU/N & +{-}/{-} & N & W:1b \\
B0322 & 1a<o0L 1b*<o2 / 2a<1a1L & NU/N & +{-}/+ & N & W:1b \\
B0323 & 1a<o0L 1b*<o2 / 2a<1a2 & NU/N & +{-}/{-} & N & W:1b \\
B0324 & 1a<o0L 1b*<o2 / 2a<1a2L & NU/N & +{-}/+ & N & W:1b \\
B0325 & 1a<o0L 1b*<o2 / 2a<1b1 & NU/U & +{-}/{-} & N & W:1b,2a \\
B0326 & 1a<o0L 1b*<o2 / 2a<1b2 & NU/U & +{-}/{-} & N & W:1b,2a \\
B0327 & 1a<o1 / 2a<1a0L 2b*<1a1 & I/NU & {-}/{-}{-} & N & W:2b \\
B0328 & 1a<o1 / 2a<1a0L 2b*<1a2 & I/NU & {-}/{-}{-} & N & W:2b \\
B0329 & 1a<o1 / 2a<1a1 2b*<1a1 & I/IU & {-}/{-}{-} & N & W:1a,2a \\
B0330 & 1a<o1 / 2a<1a1 2b*<1a2 & I/IU & {-}/{-}{-} & N & W:1a,2a \\
B0331 & 1a<o1 / 2a<1a1L 2b*<1a1 & I/IU & {-}/{-}{-} & N & W:1a,2a \\
B0332 & 1a<o1 / 2a<1a1L 2b*<1a2 & I/IU & {-}/{-}{-} & N & W:1a,2a \\
B0333 & 1a<o1 / 2a<1a2 2b*<1a1 & I/IU & {-}/{-}{-} & N & W:1a,2a \\
B0334 & 1a<o1 / 2a<1a2 2b*<1a2 & I/IU & {-}/{-}{-} & N & W:1a,2a \\
B0335 & 1a<o1 / 2a<1a2L 2b*<1a1 & I/IU & {-}/{-}{-} & N & W:1a,2a \\
B0336 & 1a<o1 / 2a<1a2L 2b*<1a2 & I/IU & {-}/{-}{-} & N & W:1a,2a \\
B0337 & 1a<o1 1b<o1 / 2a*<1a1 & EE/U & {-}{-}/{-} & E &  \\
B0338 & 1a<o1 1b<o1 / 2a*<1a2 & EE/U & {-}{-}/{-} & E &  \\
B0339 & 1a<o1 1b<o1L / 2a*<1a1 & LI/U & {-}+/{-} & B & CPW:1b \\
B0340 & 1a<o1 1b<o1L / 2a*<1a2 & LI/U & {-}+/{-} & B & CPW:1b \\
B0341 & 1a<o1 1b<o1L / 2a*<1b1 & LI/U & {-}+/{-} & B & CPW:1b \\
B0342 & 1a<o1 1b<o1L / 2a*<1b2 & LI/U & {-}+/{-} & B & CPW:1b \\
B0343 & 1a<o1 1b<o2 / 2a*<1a1 & EE/U & {-}{-}/{-} & E & C:1b \\
B0344 & 1a<o1 1b<o2 / 2a*<1a2 & EE/U & {-}{-}/{-} & E & C:1b \\
B0345 & 1a<o1 1b<o2 / 2a*<1b1 & EE/U & {-}{-}/{-} & E & C:1b \\
B0346 & 1a<o1 1b<o2 / 2a*<1b2 & EE/U & {-}{-}/{-} & E & C:1b \\
B0347 & 1a<o1 1b<o2L / 2a*<1a1 & LI/U & {-}+/{-} & B & PW:1b \\
B0348 & 1a<o1 1b<o2L / 2a*<1a2 & LI/U & {-}+/{-} & B & PW:1b \\
B0349 & 1a<o1 1b<o2L / 2a*<1b1 & LI/U & {-}+/{-} & B & PW:1b \\
B0350 & 1a<o1 1b<o2L / 2a*<1b2 & LI/U & {-}+/{-} & B & PW:1b \\
B0351 & 1a<o1 1b*<o1 / 2a<1a0L & IU/N & {-}{-}/{-} & N & W:1a \\
B0352 & 1a<o1 1b*<o1 / 2a<1a1 & IU/I & {-}{-}/{-} & N & W:1a,2a \\
B0353 & 1a<o1 1b*<o1 / 2a<1a1L & IU/I & {-}{-}/{-} & N & W:1a,2a \\
B0354 & 1a<o1 1b*<o1 / 2a<1a2 & IU/I & {-}{-}/{-} & N & W:1a,2a \\
B0355 & 1a<o1 1b*<o1 / 2a<1a2L & IU/I & {-}{-}/{-} & N & W:1a,2a \\
B0356 & 1a<o1 1b*<o1 / 2a<1b1 & IU/U & {-}{-}/{-} & N & W:1a \\
B0357 & 1a<o1 1b*<o1 / 2a<1b2 & IU/U & {-}{-}/{-} & N & W:1a \\
B0358 & 1a<o1 1b*<o2 / 2a<1a0L & IU/N & {-}{-}/{-} & N & W:1a \\
B0359 & 1a<o1 1b*<o2 / 2a<1a1 & IU/I & {-}{-}/{-} & N & W:1a,2a \\
B0360 & 1a<o1 1b*<o2 / 2a<1a1L & IU/I & {-}{-}/{-} & N & W:1a,2a \\
B0361 & 1a<o1 1b*<o2 / 2a<1a2 & IU/I & {-}{-}/{-} & N & W:1a,2a \\
B0362 & 1a<o1 1b*<o2 / 2a<1a2L & IU/I & {-}{-}/{-} & N & W:1a,2a \\
B0363 & 1a<o1 1b*<o2 / 2a<1b1 & IU/U & {-}{-}/{-} & N & W:1a \\
B0364 & 1a<o1 1b*<o2 / 2a<1b2 & IU/U & {-}{-}/{-} & N & W:1a \\
B0365 & 1a<o1L / 2a<1a0L 2b*<1a1 & I/NU & +/+{-} & N & W:2b \\
B0366 & 1a<o1L / 2a<1a0L 2b*<1a2 & I/NU & +/+{-} & N & W:2b \\
B0367 & 1a<o1L / 2a<1a1 2b*<1a1 & I/IU & +/{-}{-} & N & W:1a,2a \\
B0368 & 1a<o1L / 2a<1a1 2b*<1a2 & I/IU & +/{-}{-} & N & W:1a,2a \\
B0369 & 1a<o1L / 2a<1a1L 2b*<1a1 & I/IU & +/+{-} & N & W:1a,2a \\
B0370 & 1a<o1L / 2a<1a1L 2b*<1a2 & I/IU & +/+{-} & N & W:1a,2a \\
B0371 & 1a<o1L / 2a<1a2 2b*<1a1 & I/IU & +/{-}{-} & N & W:1a,2a \\
B0372 & 1a<o1L / 2a<1a2 2b*<1a2 & I/IU & +/{-}{-} & N & W:1a,2a \\
B0373 & 1a<o1L / 2a<1a2L 2b*<1a1 & I/IU & +/+{-} & N & W:1a,2a \\
B0374 & 1a<o1L / 2a<1a2L 2b*<1a2 & I/IU & +/+{-} & N & W:1a,2a \\
B0375 & 1a<o1L 1b<o2 / 2a*<1a1 & IL/U & +{-}/{-} & B & C:1a,1b PW:1a \\
B0376 & 1a<o1L 1b<o2 / 2a*<1a2 & IL/U & +{-}/{-} & B & C:1a,1b PW:1a \\
B0377 & 1a<o1L 1b<o2 / 2a*<1b1 & IL/U & +{-}/{-} & B & C:1a,1b PW:1a \\
B0378 & 1a<o1L 1b<o2 / 2a*<1b2 & IL/U & +{-}/{-} & B & C:1a,1b PW:1a \\
B0379 & 1a<o1L 1b*<o1 / 2a<1a0L & IU/N & +{-}/+ & N & W:1a \\
B0380 & 1a<o1L 1b*<o1 / 2a<1a1 & IU/I & +{-}/{-} & N & W:1a,2a \\
B0381 & 1a<o1L 1b*<o1 / 2a<1a1L & IU/I & +{-}/+ & N & W:1a,2a \\
B0382 & 1a<o1L 1b*<o1 / 2a<1a2 & IU/I & +{-}/{-} & N & W:1a,2a \\
B0383 & 1a<o1L 1b*<o1 / 2a<1a2L & IU/I & +{-}/+ & N & W:1a,2a \\
B0384 & 1a<o1L 1b*<o1 / 2a<1b1 & IU/U & +{-}/{-} & N & W:1a \\
B0385 & 1a<o1L 1b*<o1 / 2a<1b2 & IU/U & +{-}/{-} & N & W:1a \\
B0386 & 1a<o1L 1b*<o2 / 2a<1a0L & IU/N & +{-}/+ & N & W:1a \\
B0387 & 1a<o1L 1b*<o2 / 2a<1a1 & IU/I & +{-}/{-} & N & W:1a,2a \\
B0388 & 1a<o1L 1b*<o2 / 2a<1a1L & IU/I & +{-}/+ & N & W:1a,2a \\
B0389 & 1a<o1L 1b*<o2 / 2a<1a2 & IU/I & +{-}/{-} & N & W:1a,2a \\
B0390 & 1a<o1L 1b*<o2 / 2a<1a2L & IU/I & +{-}/+ & N & W:1a,2a \\
B0391 & 1a<o1L 1b*<o2 / 2a<1b1 & IU/U & +{-}/{-} & N & W:1a \\
B0392 & 1a<o1L 1b*<o2 / 2a<1b2 & IU/U & +{-}/{-} & N & W:1a \\
B0393 & 1a<o2 / 2a<1a0L 2b*<1a1 & I/NU & {-}/{-}{-} & N & W:2b \\
B0394 & 1a<o2 / 2a<1a0L 2b*<1a2 & I/NU & {-}/{-}{-} & N & W:2b \\
B0395 & 1a<o2 / 2a<1a1 2b*<1a1 & I/IU & {-}/{-}{-} & N & W:1a,2a \\
B0396 & 1a<o2 / 2a<1a1 2b*<1a2 & I/IU & {-}/{-}{-} & N & W:1a,2a \\
B0397 & 1a<o2 / 2a<1a1L 2b*<1a1 & I/IU & {-}/{-}{-} & N & W:1a,2a \\
B0398 & 1a<o2 / 2a<1a1L 2b*<1a2 & I/IU & {-}/{-}{-} & N & W:1a,2a \\
B0399 & 1a<o2 / 2a<1a2 2b*<1a1 & I/IU & {-}/{-}{-} & N & W:1a,2a \\
B0400 & 1a<o2 / 2a<1a2 2b*<1a2 & I/IU & {-}/{-}{-} & N & W:1a,2a \\
B0401 & 1a<o2 / 2a<1a2L 2b*<1a1 & I/IU & {-}/{-}{-} & N & W:1a,2a \\
B0402 & 1a<o2 / 2a<1a2L 2b*<1a2 & I/IU & {-}/{-}{-} & N & W:1a,2a \\
B0403 & 1a<o2 1b<o2 / 2a*<1a1 & EE/U & {-}{-}/{-} & E &  \\
B0404 & 1a<o2 1b<o2 / 2a*<1a2 & EE/U & {-}{-}/{-} & E &  \\
B0405 & 1a<o2 1b<o2L / 2a*<1a1 & LI/U & {-}+/{-} & B & CPW:1b \\
B0406 & 1a<o2 1b<o2L / 2a*<1a2 & LI/U & {-}+/{-} & B & CPW:1b \\
B0407 & 1a<o2 1b<o2L / 2a*<1b1 & LI/U & {-}+/{-} & B & CPW:1b \\
B0408 & 1a<o2 1b<o2L / 2a*<1b2 & LI/U & {-}+/{-} & B & CPW:1b \\
B0409 & 1a<o2 1b*<o1 / 2a<1a0L & IU/N & {-}{-}/{-} & N & W:1a \\
B0410 & 1a<o2 1b*<o1 / 2a<1a1 & IU/I & {-}{-}/{-} & N & W:1a,2a \\
B0411 & 1a<o2 1b*<o1 / 2a<1a1L & IU/I & {-}{-}/{-} & N & W:1a,2a \\
B0412 & 1a<o2 1b*<o1 / 2a<1a2 & IU/I & {-}{-}/{-} & N & W:1a,2a \\
B0413 & 1a<o2 1b*<o1 / 2a<1a2L & IU/I & {-}{-}/{-} & N & W:1a,2a \\
B0414 & 1a<o2 1b*<o1 / 2a<1b1 & IU/U & {-}{-}/{-} & N & W:1a \\
B0415 & 1a<o2 1b*<o1 / 2a<1b2 & IU/U & {-}{-}/{-} & N & W:1a \\
B0416 & 1a<o2 1b*<o2 / 2a<1a0L & IU/N & {-}{-}/{-} & N & W:1a \\
B0417 & 1a<o2 1b*<o2 / 2a<1a1 & IU/I & {-}{-}/{-} & N & W:1a,2a \\
B0418 & 1a<o2 1b*<o2 / 2a<1a1L & IU/I & {-}{-}/{-} & N & W:1a,2a \\
B0419 & 1a<o2 1b*<o2 / 2a<1a2 & IU/I & {-}{-}/{-} & N & W:1a,2a \\
B0420 & 1a<o2 1b*<o2 / 2a<1a2L & IU/I & {-}{-}/{-} & N & W:1a,2a \\
B0421 & 1a<o2 1b*<o2 / 2a<1b1 & IU/U & {-}{-}/{-} & N & W:1a \\
B0422 & 1a<o2 1b*<o2 / 2a<1b2 & IU/U & {-}{-}/{-} & N & W:1a \\
B0423 & 1a<o2L / 2a<1a0L 2b*<1a1 & I/NU & +/+{-} & N & W:2b \\
B0424 & 1a<o2L / 2a<1a0L 2b*<1a2 & I/NU & +/+{-} & N & W:2b \\
B0425 & 1a<o2L / 2a<1a1 2b*<1a1 & I/IU & +/{-}{-} & N & W:1a,2a \\
B0426 & 1a<o2L / 2a<1a1 2b*<1a2 & I/IU & +/{-}{-} & N & W:1a,2a \\
B0427 & 1a<o2L / 2a<1a1L 2b*<1a1 & I/IU & +/+{-} & N & W:1a,2a \\
B0428 & 1a<o2L / 2a<1a1L 2b*<1a2 & I/IU & +/+{-} & N & W:1a,2a \\
B0429 & 1a<o2L / 2a<1a2 2b*<1a1 & I/IU & +/{-}{-} & N & W:1a,2a \\
B0430 & 1a<o2L / 2a<1a2 2b*<1a2 & I/IU & +/{-}{-} & N & W:1a,2a \\
B0431 & 1a<o2L / 2a<1a2L 2b*<1a1 & I/IU & +/+{-} & N & W:1a,2a \\
B0432 & 1a<o2L / 2a<1a2L 2b*<1a2 & I/IU & +/+{-} & N & W:1a,2a \\
B0433 & 1a<o2L 1b*<o1 / 2a<1a0L & IU/N & +{-}/+ & N & W:1a \\
B0434 & 1a<o2L 1b*<o1 / 2a<1a1 & IU/I & +{-}/{-} & N & W:1a,2a \\
B0435 & 1a<o2L 1b*<o1 / 2a<1a1L & IU/I & +{-}/+ & N & W:1a,2a \\
B0436 & 1a<o2L 1b*<o1 / 2a<1a2 & IU/I & +{-}/{-} & N & W:1a,2a \\
B0437 & 1a<o2L 1b*<o1 / 2a<1a2L & IU/I & +{-}/+ & N & W:1a,2a \\
B0438 & 1a<o2L 1b*<o1 / 2a<1b1 & IU/U & +{-}/{-} & N & W:1a \\
B0439 & 1a<o2L 1b*<o1 / 2a<1b2 & IU/U & +{-}/{-} & N & W:1a \\
B0440 & 1a<o2L 1b*<o2 / 2a<1a0L & IU/N & +{-}/+ & N & W:1a \\
B0441 & 1a<o2L 1b*<o2 / 2a<1a1 & IU/I & +{-}/{-} & N & W:1a,2a \\
B0442 & 1a<o2L 1b*<o2 / 2a<1a1L & IU/I & +{-}/+ & N & W:1a,2a \\
B0443 & 1a<o2L 1b*<o2 / 2a<1a2 & IU/I & +{-}/{-} & N & W:1a,2a \\
B0444 & 1a<o2L 1b*<o2 / 2a<1a2L & IU/I & +{-}/+ & N & W:1a,2a \\
B0445 & 1a<o2L 1b*<o2 / 2a<1b1 & IU/U & +{-}/{-} & N & W:1a \\
B0446 & 1a<o2L 1b*<o2 / 2a<1b2 & IU/U & +{-}/{-} & N & W:1a \\
B0447 & 1a*<o1 / 2a<1a1 2b<1a1 & U/UU & {-}/{-}{-} & N &  \\
B0448 & 1a*<o1 / 2a<1a1 2b<1a2 & U/UU & {-}/{-}{-} & N &  \\
B0449 & 1a*<o1 / 2a<1a2 2b<1a2 & U/UU & {-}/{-}{-} & N &  \\
B0450 & 1a*<o2 / 2a<1a1 2b<1a1 & U/UU & {-}/{-}{-} & N &  \\
B0451 & 1a*<o2 / 2a<1a1 2b<1a2 & U/UU & {-}/{-}{-} & N &  \\
B0452 & 1a*<o2 / 2a<1a2 2b<1a2 & U/UU & {-}/{-}{-} & N &  \\
B0453 & 1a<o0L 1b<o1 / 2a*<1a1,1b1 & NI/U & +{-}/{-} & N & W:2a \\
B0454 & 1a<o0L 1b<o1 / 2a*<1a1,1b2 & NI/U & +{-}/{-} & N & W:2a \\
B0455 & 1a<o0L 1b<o1 / 2a*<1a2,1b1 & NI/U & +{-}/{-} & N & W:2a \\
B0456 & 1a<o0L 1b<o1 / 2a*<1a2,1b2 & NI/U & +{-}/{-} & N & W:2a \\
B0457 & 1a<o0L 1b<o2 / 2a*<1a1,1b1 & NI/U & +{-}/{-} & N & W:2a \\
B0458 & 1a<o0L 1b<o2 / 2a*<1a1,1b2 & NI/U & +{-}/{-} & N & W:2a \\
B0459 & 1a<o0L 1b<o2 / 2a*<1a2,1b1 & NI/U & +{-}/{-} & N & W:2a \\
B0460 & 1a<o0L 1b<o2 / 2a*<1a2,1b2 & NI/U & +{-}/{-} & N & W:2a \\
B0461 & 1a<o0L 1b*<o1 / 2a<1a0L,1b1 & NU/U & +{-}/+ & N & W:1b,2a \\
B0462 & 1a<o0L 1b*<o1 / 2a<1a0L,1b2 & NU/U & +{-}/+ & N & W:1b,2a \\
B0463 & 1a<o0L 1b*<o1 / 2a<1a1,1b1 & NU/U & +{-}/{-} & N & W:1b,2a \\
B0464 & 1a<o0L 1b*<o1 / 2a<1a1,1b2 & NU/U & +{-}/{-} & N & W:1b,2a \\
B0465 & 1a<o0L 1b*<o1 / 2a<1a1L,1b1 & NU/U & +{-}/+ & N & W:1b,2a \\
B0466 & 1a<o0L 1b*<o1 / 2a<1a1L,1b2 & NU/U & +{-}/+ & N & W:1b,2a \\
B0467 & 1a<o0L 1b*<o1 / 2a<1a2,1b1 & NU/U & +{-}/{-} & N & W:1b,2a \\
B0468 & 1a<o0L 1b*<o1 / 2a<1a2,1b2 & NU/U & +{-}/{-} & N & W:1b,2a \\
B0469 & 1a<o0L 1b*<o1 / 2a<1a2L,1b1 & NU/U & +{-}/+ & N & W:1b,2a \\
B0470 & 1a<o0L 1b*<o1 / 2a<1a2L,1b2 & NU/U & +{-}/+ & N & W:1b,2a \\
B0471 & 1a<o0L 1b*<o2 / 2a<1a0L,1b1 & NU/U & +{-}/+ & N & W:1b,2a \\
B0472 & 1a<o0L 1b*<o2 / 2a<1a0L,1b2 & NU/U & +{-}/+ & N & W:1b,2a \\
B0473 & 1a<o0L 1b*<o2 / 2a<1a1,1b1 & NU/U & +{-}/{-} & N & W:1b,2a \\
B0474 & 1a<o0L 1b*<o2 / 2a<1a1,1b2 & NU/U & +{-}/{-} & N & W:1b,2a \\
B0475 & 1a<o0L 1b*<o2 / 2a<1a1L,1b1 & NU/U & +{-}/+ & N & W:1b,2a \\
B0476 & 1a<o0L 1b*<o2 / 2a<1a1L,1b2 & NU/U & +{-}/+ & N & W:1b,2a \\
B0477 & 1a<o0L 1b*<o2 / 2a<1a2,1b1 & NU/U & +{-}/{-} & N & W:1b,2a \\
B0478 & 1a<o0L 1b*<o2 / 2a<1a2,1b2 & NU/U & +{-}/{-} & N & W:1b,2a \\
B0479 & 1a<o0L 1b*<o2 / 2a<1a2L,1b1 & NU/U & +{-}/+ & N & W:1b,2a \\
B0480 & 1a<o0L 1b*<o2 / 2a<1a2L,1b2 & NU/U & +{-}/+ & N & W:1b,2a \\
B0481 & 1a<o1 1b<o1 / 2a*<1a1,1b1 & EE/U & {-}{-}/{-} & E &  \\
B0482 & 1a<o1 1b<o1 / 2a*<1a1,1b2 & EE/U & {-}{-}/{-} & E &  \\
B0483 & 1a<o1 1b<o1 / 2a*<1a2,1b2 & EE/U & {-}{-}/{-} & E &  \\
B0484 & 1a<o1 1b<o1L / 2a*<1a1,1b1 & LI/U & {-}+/{-} & B & CPW:1b \\
B0485 & 1a<o1 1b<o1L / 2a*<1a1,1b2 & LI/U & {-}+/{-} & B & CPW:1b \\
B0486 & 1a<o1 1b<o1L / 2a*<1a2,1b1 & LI/U & {-}+/{-} & B & CPW:1b \\
B0487 & 1a<o1 1b<o1L / 2a*<1a2,1b2 & LI/U & {-}+/{-} & B & CPW:1b \\
B0488 & 1a<o1 1b<o2 / 2a*<1a1,1b1 & EE/U & {-}{-}/{-} & E & C:1b \\
B0489 & 1a<o1 1b<o2 / 2a*<1a1,1b2 & EE/U & {-}{-}/{-} & E & C:1b \\
B0490 & 1a<o1 1b<o2 / 2a*<1a2,1b1 & EE/U & {-}{-}/{-} & E & C:1b \\
B0491 & 1a<o1 1b<o2 / 2a*<1a2,1b2 & EE/U & {-}{-}/{-} & E & C:1b \\
B0492 & 1a<o1 1b<o2L / 2a*<1a1,1b1 & LI/U & {-}+/{-} & B & PW:1b \\
B0493 & 1a<o1 1b<o2L / 2a*<1a1,1b2 & LI/U & {-}+/{-} & B & PW:1b \\
B0494 & 1a<o1 1b<o2L / 2a*<1a2,1b1 & LI/U & {-}+/{-} & B & PW:1b \\
B0495 & 1a<o1 1b<o2L / 2a*<1a2,1b2 & LI/U & {-}+/{-} & B & PW:1b \\
B0496 & 1a<o1 1b*<o1 / 2a<1a0L,1b1 & IU/U & {-}{-}/{-} & N & W:1a \\
B0497 & 1a<o1 1b*<o1 / 2a<1a0L,1b2 & IU/U & {-}{-}/{-} & N & W:1a \\
B0498 & 1a<o1 1b*<o1 / 2a<1a1,1b1 & IU/I & {-}{-}/{-} & N & W:1a,2a \\
B0499 & 1a<o1 1b*<o1 / 2a<1a1,1b2 & IU/I & {-}{-}/{-} & N & W:1a,2a \\
B0500 & 1a<o1 1b*<o1 / 2a<1a1L,1b1 & IU/I & {-}{-}/{-} & N & W:1a,2a \\
B0501 & 1a<o1 1b*<o1 / 2a<1a1L,1b2 & IU/I & {-}{-}/{-} & N & W:1a,2a \\
B0502 & 1a<o1 1b*<o1 / 2a<1a2,1b1 & IU/I & {-}{-}/{-} & N & W:1a,2a \\
B0503 & 1a<o1 1b*<o1 / 2a<1a2,1b2 & IU/I & {-}{-}/{-} & N & W:1a,2a \\
B0504 & 1a<o1 1b*<o1 / 2a<1a2L,1b1 & IU/I & {-}{-}/{-} & N & W:1a,2a \\
B0505 & 1a<o1 1b*<o1 / 2a<1a2L,1b2 & IU/I & {-}{-}/{-} & N & W:1a,2a \\
B0506 & 1a<o1 1b*<o2 / 2a<1a0L,1b1 & IU/U & {-}{-}/{-} & N & W:1a \\
B0507 & 1a<o1 1b*<o2 / 2a<1a0L,1b2 & IU/U & {-}{-}/{-} & N & W:1a \\
B0508 & 1a<o1 1b*<o2 / 2a<1a1,1b1 & IU/I & {-}{-}/{-} & N & W:1a,2a \\
B0509 & 1a<o1 1b*<o2 / 2a<1a1,1b2 & IU/I & {-}{-}/{-} & N & W:1a,2a \\
B0510 & 1a<o1 1b*<o2 / 2a<1a1L,1b1 & IU/I & {-}{-}/{-} & N & W:1a,2a \\
B0511 & 1a<o1 1b*<o2 / 2a<1a1L,1b2 & IU/I & {-}{-}/{-} & N & W:1a,2a \\
B0512 & 1a<o1 1b*<o2 / 2a<1a2,1b1 & IU/I & {-}{-}/{-} & N & W:1a,2a \\
B0513 & 1a<o1 1b*<o2 / 2a<1a2,1b2 & IU/I & {-}{-}/{-} & N & W:1a,2a \\
B0514 & 1a<o1 1b*<o2 / 2a<1a2L,1b1 & IU/I & {-}{-}/{-} & N & W:1a,2a \\
B0515 & 1a<o1 1b*<o2 / 2a<1a2L,1b2 & IU/I & {-}{-}/{-} & N & W:1a,2a \\
B0516 & 1a<o1L 1b<o2 / 2a*<1a1,1b1 & IL/U & +{-}/{-} & B & C:1a,1b PW:1a \\
B0517 & 1a<o1L 1b<o2 / 2a*<1a1,1b2 & IL/U & +{-}/{-} & B & C:1a,1b PW:1a \\
B0518 & 1a<o1L 1b<o2 / 2a*<1a2,1b1 & IL/U & +{-}/{-} & B & C:1a,1b PW:1a \\
B0519 & 1a<o1L 1b<o2 / 2a*<1a2,1b2 & IL/U & +{-}/{-} & B & C:1a,1b PW:1a \\
B0520 & 1a<o1L 1b*<o1 / 2a<1a0L,1b1 & IU/U & +{-}/+ & N & W:1a \\
B0521 & 1a<o1L 1b*<o1 / 2a<1a0L,1b2 & IU/U & +{-}/+ & N & W:1a \\
B0522 & 1a<o1L 1b*<o1 / 2a<1a1,1b1 & IU/I & +{-}/{-} & N & W:1a,2a \\
B0523 & 1a<o1L 1b*<o1 / 2a<1a1,1b2 & IU/I & +{-}/{-} & N & W:1a,2a \\
B0524 & 1a<o1L 1b*<o1 / 2a<1a1L,1b1 & IU/I & +{-}/+ & N & W:1a,2a \\
B0525 & 1a<o1L 1b*<o1 / 2a<1a1L,1b2 & IU/I & +{-}/+ & N & W:1a,2a \\
B0526 & 1a<o1L 1b*<o1 / 2a<1a2,1b1 & IU/I & +{-}/{-} & N & W:1a,2a \\
B0527 & 1a<o1L 1b*<o1 / 2a<1a2,1b2 & IU/I & +{-}/{-} & N & W:1a,2a \\
B0528 & 1a<o1L 1b*<o1 / 2a<1a2L,1b1 & IU/I & +{-}/+ & N & W:1a,2a \\
B0529 & 1a<o1L 1b*<o1 / 2a<1a2L,1b2 & IU/I & +{-}/+ & N & W:1a,2a \\
B0530 & 1a<o1L 1b*<o2 / 2a<1a0L,1b1 & IU/U & +{-}/+ & N & W:1a \\
B0531 & 1a<o1L 1b*<o2 / 2a<1a0L,1b2 & IU/U & +{-}/+ & N & W:1a \\
B0532 & 1a<o1L 1b*<o2 / 2a<1a1,1b1 & IU/I & +{-}/{-} & N & W:1a,2a \\
B0533 & 1a<o1L 1b*<o2 / 2a<1a1,1b2 & IU/I & +{-}/{-} & N & W:1a,2a \\
B0534 & 1a<o1L 1b*<o2 / 2a<1a1L,1b1 & IU/I & +{-}/+ & N & W:1a,2a \\
B0535 & 1a<o1L 1b*<o2 / 2a<1a1L,1b2 & IU/I & +{-}/+ & N & W:1a,2a \\
B0536 & 1a<o1L 1b*<o2 / 2a<1a2,1b1 & IU/I & +{-}/{-} & N & W:1a,2a \\
B0537 & 1a<o1L 1b*<o2 / 2a<1a2,1b2 & IU/I & +{-}/{-} & N & W:1a,2a \\
B0538 & 1a<o1L 1b*<o2 / 2a<1a2L,1b1 & IU/I & +{-}/+ & N & W:1a,2a \\
B0539 & 1a<o1L 1b*<o2 / 2a<1a2L,1b2 & IU/I & +{-}/+ & N & W:1a,2a \\
B0540 & 1a<o2 1b<o2 / 2a*<1a1,1b1 & EE/U & {-}{-}/{-} & E &  \\
B0541 & 1a<o2 1b<o2 / 2a*<1a1,1b2 & EE/U & {-}{-}/{-} & E &  \\
B0542 & 1a<o2 1b<o2 / 2a*<1a2,1b2 & EE/U & {-}{-}/{-} & E &  \\
B0543 & 1a<o2 1b<o2L / 2a*<1a1,1b1 & LI/U & {-}+/{-} & B & CPW:1b \\
B0544 & 1a<o2 1b<o2L / 2a*<1a1,1b2 & LI/U & {-}+/{-} & B & CPW:1b \\
B0545 & 1a<o2 1b<o2L / 2a*<1a2,1b1 & LI/U & {-}+/{-} & B & CPW:1b \\
B0546 & 1a<o2 1b<o2L / 2a*<1a2,1b2 & LI/U & {-}+/{-} & B & CPW:1b \\
B0547 & 1a<o2 1b*<o1 / 2a<1a0L,1b1 & IU/U & {-}{-}/{-} & N & W:1a \\
B0548 & 1a<o2 1b*<o1 / 2a<1a0L,1b2 & IU/U & {-}{-}/{-} & N & W:1a \\
B0549 & 1a<o2 1b*<o1 / 2a<1a1,1b1 & IU/I & {-}{-}/{-} & N & W:1a,2a \\
B0550 & 1a<o2 1b*<o1 / 2a<1a1,1b2 & IU/I & {-}{-}/{-} & N & W:1a,2a \\
B0551 & 1a<o2 1b*<o1 / 2a<1a1L,1b1 & IU/I & {-}{-}/{-} & N & W:1a,2a \\
B0552 & 1a<o2 1b*<o1 / 2a<1a1L,1b2 & IU/I & {-}{-}/{-} & N & W:1a,2a \\
B0553 & 1a<o2 1b*<o1 / 2a<1a2,1b1 & IU/I & {-}{-}/{-} & N & W:1a,2a \\
B0554 & 1a<o2 1b*<o1 / 2a<1a2,1b2 & IU/I & {-}{-}/{-} & N & W:1a,2a \\
B0555 & 1a<o2 1b*<o1 / 2a<1a2L,1b1 & IU/I & {-}{-}/{-} & N & W:1a,2a \\
B0556 & 1a<o2 1b*<o1 / 2a<1a2L,1b2 & IU/I & {-}{-}/{-} & N & W:1a,2a \\
B0557 & 1a<o2 1b*<o2 / 2a<1a0L,1b1 & IU/U & {-}{-}/{-} & N & W:1a \\
B0558 & 1a<o2 1b*<o2 / 2a<1a0L,1b2 & IU/U & {-}{-}/{-} & N & W:1a \\
B0559 & 1a<o2 1b*<o2 / 2a<1a1,1b1 & IU/I & {-}{-}/{-} & N & W:1a,2a \\
B0560 & 1a<o2 1b*<o2 / 2a<1a1,1b2 & IU/I & {-}{-}/{-} & N & W:1a,2a \\
B0561 & 1a<o2 1b*<o2 / 2a<1a1L,1b1 & IU/I & {-}{-}/{-} & N & W:1a,2a \\
B0562 & 1a<o2 1b*<o2 / 2a<1a1L,1b2 & IU/I & {-}{-}/{-} & N & W:1a,2a \\
B0563 & 1a<o2 1b*<o2 / 2a<1a2,1b1 & IU/I & {-}{-}/{-} & N & W:1a,2a \\
B0564 & 1a<o2 1b*<o2 / 2a<1a2,1b2 & IU/I & {-}{-}/{-} & N & W:1a,2a \\
B0565 & 1a<o2 1b*<o2 / 2a<1a2L,1b1 & IU/I & {-}{-}/{-} & N & W:1a,2a \\
B0566 & 1a<o2 1b*<o2 / 2a<1a2L,1b2 & IU/I & {-}{-}/{-} & N & W:1a,2a \\
B0567 & 1a<o2L 1b*<o1 / 2a<1a0L,1b1 & IU/U & +{-}/+ & N & W:1a \\
B0568 & 1a<o2L 1b*<o1 / 2a<1a0L,1b2 & IU/U & +{-}/+ & N & W:1a \\
B0569 & 1a<o2L 1b*<o1 / 2a<1a1,1b1 & IU/I & +{-}/{-} & N & W:1a,2a \\
B0570 & 1a<o2L 1b*<o1 / 2a<1a1,1b2 & IU/I & +{-}/{-} & N & W:1a,2a \\
B0571 & 1a<o2L 1b*<o1 / 2a<1a1L,1b1 & IU/I & +{-}/+ & N & W:1a,2a \\
B0572 & 1a<o2L 1b*<o1 / 2a<1a1L,1b2 & IU/I & +{-}/+ & N & W:1a,2a \\
B0573 & 1a<o2L 1b*<o1 / 2a<1a2,1b1 & IU/I & +{-}/{-} & N & W:1a,2a \\
B0574 & 1a<o2L 1b*<o1 / 2a<1a2,1b2 & IU/I & +{-}/{-} & N & W:1a,2a \\
B0575 & 1a<o2L 1b*<o1 / 2a<1a2L,1b1 & IU/I & +{-}/+ & N & W:1a,2a \\
B0576 & 1a<o2L 1b*<o1 / 2a<1a2L,1b2 & IU/I & +{-}/+ & N & W:1a,2a \\
B0577 & 1a<o2L 1b*<o2 / 2a<1a0L,1b1 & IU/U & +{-}/+ & N & W:1a \\
B0578 & 1a<o2L 1b*<o2 / 2a<1a0L,1b2 & IU/U & +{-}/+ & N & W:1a \\
B0579 & 1a<o2L 1b*<o2 / 2a<1a1,1b1 & IU/I & +{-}/{-} & N & W:1a,2a \\
B0580 & 1a<o2L 1b*<o2 / 2a<1a1,1b2 & IU/I & +{-}/{-} & N & W:1a,2a \\
B0581 & 1a<o2L 1b*<o2 / 2a<1a1L,1b1 & IU/I & +{-}/+ & N & W:1a,2a \\
B0582 & 1a<o2L 1b*<o2 / 2a<1a1L,1b2 & IU/I & +{-}/+ & N & W:1a,2a \\
B0583 & 1a<o2L 1b*<o2 / 2a<1a2,1b1 & IU/I & +{-}/{-} & N & W:1a,2a \\
B0584 & 1a<o2L 1b*<o2 / 2a<1a2,1b2 & IU/I & +{-}/{-} & N & W:1a,2a \\
B0585 & 1a<o2L 1b*<o2 / 2a<1a2L,1b1 & IU/I & +{-}/+ & N & W:1a,2a \\
B0586 & 1a<o2L 1b*<o2 / 2a<1a2L,1b2 & IU/I & +{-}/+ & N & W:1a,2a \\
\end{casetable}
\subsection*{Height 2, 4 stages after the original (4851 cases)}
\begin{casetable}
B0587 & 1a<o0L 1b<o1 / 2a<1a0L 2b<1a1 & NI/NN & +{-}/+{-} & N &  \\
B0588 & 1a<o0L 1b<o1 / 2a<1a0L 2b<1a2 & NI/NN & +{-}/+{-} & N &  \\
B0589 & 1a<o0L 1b<o1 / 2a<1a0L 2b<1b0L & NI/NN & +{-}/+{-} & N &  \\
B0590 & 1a<o0L 1b<o1 / 2a<1a0L 2b<1b1 & NI/NI & +{-}/+{-} & N &  \\
B0591 & 1a<o0L 1b<o1 / 2a<1a0L 2b<1b1L & NI/NI & +{-}/+{-} & N &  \\
B0592 & 1a<o0L 1b<o1 / 2a<1a0L 2b<1b2 & NI/NI & +{-}/+{-} & N &  \\
B0593 & 1a<o0L 1b<o1 / 2a<1a0L 2b<1b2L & NI/NI & +{-}/+{-} & N &  \\
B0594 & 1a<o0L 1b<o1 / 2a<1a1 2b<1a1 & NI/NN & +{-}/{-}{-} & N &  \\
B0595 & 1a<o0L 1b<o1 / 2a<1a1 2b<1a1L & NI/NN & +{-}/{-}+ & N &  \\
B0596 & 1a<o0L 1b<o1 / 2a<1a1 2b<1a2 & NI/NN & +{-}/{-}{-} & N &  \\
B0597 & 1a<o0L 1b<o1 / 2a<1a1 2b<1a2L & NI/NN & +{-}/{-}+ & N &  \\
B0598 & 1a<o0L 1b<o1 / 2a<1a1 2b<1b0L & NI/NN & +{-}/{-}{-} & N &  \\
B0599 & 1a<o0L 1b<o1 / 2a<1a1 2b<1b1 & NI/NI & +{-}/{-}{-} & N &  \\
B0600 & 1a<o0L 1b<o1 / 2a<1a1 2b<1b1L & NI/NI & +{-}/{-}{-} & N &  \\
B0601 & 1a<o0L 1b<o1 / 2a<1a1 2b<1b2 & NI/NI & +{-}/{-}{-} & N &  \\
B0602 & 1a<o0L 1b<o1 / 2a<1a1 2b<1b2L & NI/NI & +{-}/{-}{-} & N &  \\
B0603 & 1a<o0L 1b<o1 / 2a<1a1L 2b<1a2 & NI/NN & +{-}/+{-} & N &  \\
B0604 & 1a<o0L 1b<o1 / 2a<1a1L 2b<1b0L & NI/NN & +{-}/+{-} & N &  \\
B0605 & 1a<o0L 1b<o1 / 2a<1a1L 2b<1b1 & NI/NI & +{-}/+{-} & N &  \\
B0606 & 1a<o0L 1b<o1 / 2a<1a1L 2b<1b1L & NI/NI & +{-}/+{-} & N &  \\
B0607 & 1a<o0L 1b<o1 / 2a<1a1L 2b<1b2 & NI/NI & +{-}/+{-} & N &  \\
B0608 & 1a<o0L 1b<o1 / 2a<1a1L 2b<1b2L & NI/NI & +{-}/+{-} & N &  \\
B0609 & 1a<o0L 1b<o1 / 2a<1a2 2b<1a2 & NI/NN & +{-}/{-}{-} & N &  \\
B0610 & 1a<o0L 1b<o1 / 2a<1a2 2b<1a2L & NI/NN & +{-}/{-}+ & N &  \\
B0611 & 1a<o0L 1b<o1 / 2a<1a2 2b<1b0L & NI/NN & +{-}/{-}{-} & N &  \\
B0612 & 1a<o0L 1b<o1 / 2a<1a2 2b<1b1 & NI/NI & +{-}/{-}{-} & N &  \\
B0613 & 1a<o0L 1b<o1 / 2a<1a2 2b<1b1L & NI/NI & +{-}/{-}{-} & N &  \\
B0614 & 1a<o0L 1b<o1 / 2a<1a2 2b<1b2 & NI/NI & +{-}/{-}{-} & N &  \\
B0615 & 1a<o0L 1b<o1 / 2a<1a2 2b<1b2L & NI/NI & +{-}/{-}{-} & N &  \\
B0616 & 1a<o0L 1b<o1 / 2a<1a2L 2b<1b0L & NI/NN & +{-}/+{-} & N &  \\
B0617 & 1a<o0L 1b<o1 / 2a<1a2L 2b<1b1 & NI/NI & +{-}/+{-} & N &  \\
B0618 & 1a<o0L 1b<o1 / 2a<1a2L 2b<1b1L & NI/NI & +{-}/+{-} & N &  \\
B0619 & 1a<o0L 1b<o1 / 2a<1a2L 2b<1b2 & NI/NI & +{-}/+{-} & N &  \\
B0620 & 1a<o0L 1b<o1 / 2a<1a2L 2b<1b2L & NI/NI & +{-}/+{-} & N &  \\
B0621 & 1a<o0L 1b<o1 / 2a<1b0L 2b<1b1 & NI/NI & +{-}/{-}{-} & N &  \\
B0622 & 1a<o0L 1b<o1 / 2a<1b0L 2b<1b2 & NI/NI & +{-}/{-}{-} & N &  \\
B0623 & 1a<o0L 1b<o1 / 2a<1b1 2b<1b1 & NI/EE & +{-}/{-}{-} & E & W:1b \\
B0624 & 1a<o0L 1b<o1 / 2a<1b1 2b<1b1L & NI/LI & +{-}/{-}{-} & B & CP:2b W:1b,2b \\
B0625 & 1a<o0L 1b<o1 / 2a<1b1 2b<1b2 & NI/EE & +{-}/{-}{-} & E & C:2b W:1b \\
B0626 & 1a<o0L 1b<o1 / 2a<1b1 2b<1b2L & NI/LI & +{-}/{-}{-} & B & P:2b W:1b,2b \\
B0627 & 1a<o0L 1b<o1 / 2a<1b1L 2b<1b2 & NI/IL & +{-}/{-}{-} & B & C:2a,2b P:2a W:1b,2a \\
B0628 & 1a<o0L 1b<o1 / 2a<1b2 2b<1b2 & NI/EE & +{-}/{-}{-} & E & W:1b \\
B0629 & 1a<o0L 1b<o1 / 2a<1b2 2b<1b2L & NI/LI & +{-}/{-}{-} & B & CP:2b W:1b,2b \\
B0630 & 1a<o0L 1b<o2 / 2a<1a0L 2b<1a1 & NI/NN & +{-}/+{-} & N &  \\
B0631 & 1a<o0L 1b<o2 / 2a<1a0L 2b<1a2 & NI/NN & +{-}/+{-} & N &  \\
B0632 & 1a<o0L 1b<o2 / 2a<1a0L 2b<1b0L & NI/NN & +{-}/+{-} & N &  \\
B0633 & 1a<o0L 1b<o2 / 2a<1a0L 2b<1b1 & NI/NI & +{-}/+{-} & N &  \\
B0634 & 1a<o0L 1b<o2 / 2a<1a0L 2b<1b1L & NI/NI & +{-}/+{-} & N &  \\
B0635 & 1a<o0L 1b<o2 / 2a<1a0L 2b<1b2 & NI/NI & +{-}/+{-} & N &  \\
B0636 & 1a<o0L 1b<o2 / 2a<1a0L 2b<1b2L & NI/NI & +{-}/+{-} & N &  \\
B0637 & 1a<o0L 1b<o2 / 2a<1a1 2b<1a1 & NI/NN & +{-}/{-}{-} & N &  \\
B0638 & 1a<o0L 1b<o2 / 2a<1a1 2b<1a1L & NI/NN & +{-}/{-}+ & N &  \\
B0639 & 1a<o0L 1b<o2 / 2a<1a1 2b<1a2 & NI/NN & +{-}/{-}{-} & N &  \\
B0640 & 1a<o0L 1b<o2 / 2a<1a1 2b<1a2L & NI/NN & +{-}/{-}+ & N &  \\
B0641 & 1a<o0L 1b<o2 / 2a<1a1 2b<1b0L & NI/NN & +{-}/{-}{-} & N &  \\
B0642 & 1a<o0L 1b<o2 / 2a<1a1 2b<1b1 & NI/NI & +{-}/{-}{-} & N &  \\
B0643 & 1a<o0L 1b<o2 / 2a<1a1 2b<1b1L & NI/NI & +{-}/{-}{-} & N &  \\
B0644 & 1a<o0L 1b<o2 / 2a<1a1 2b<1b2 & NI/NI & +{-}/{-}{-} & N &  \\
B0645 & 1a<o0L 1b<o2 / 2a<1a1 2b<1b2L & NI/NI & +{-}/{-}{-} & N &  \\
B0646 & 1a<o0L 1b<o2 / 2a<1a1L 2b<1a2 & NI/NN & +{-}/+{-} & N &  \\
B0647 & 1a<o0L 1b<o2 / 2a<1a1L 2b<1b0L & NI/NN & +{-}/+{-} & N &  \\
B0648 & 1a<o0L 1b<o2 / 2a<1a1L 2b<1b1 & NI/NI & +{-}/+{-} & N &  \\
B0649 & 1a<o0L 1b<o2 / 2a<1a1L 2b<1b1L & NI/NI & +{-}/+{-} & N &  \\
B0650 & 1a<o0L 1b<o2 / 2a<1a1L 2b<1b2 & NI/NI & +{-}/+{-} & N &  \\
B0651 & 1a<o0L 1b<o2 / 2a<1a1L 2b<1b2L & NI/NI & +{-}/+{-} & N &  \\
B0652 & 1a<o0L 1b<o2 / 2a<1a2 2b<1a2 & NI/NN & +{-}/{-}{-} & N &  \\
B0653 & 1a<o0L 1b<o2 / 2a<1a2 2b<1a2L & NI/NN & +{-}/{-}+ & N &  \\
B0654 & 1a<o0L 1b<o2 / 2a<1a2 2b<1b0L & NI/NN & +{-}/{-}{-} & N &  \\
B0655 & 1a<o0L 1b<o2 / 2a<1a2 2b<1b1 & NI/NI & +{-}/{-}{-} & N &  \\
B0656 & 1a<o0L 1b<o2 / 2a<1a2 2b<1b1L & NI/NI & +{-}/{-}{-} & N &  \\
B0657 & 1a<o0L 1b<o2 / 2a<1a2 2b<1b2 & NI/NI & +{-}/{-}{-} & N &  \\
B0658 & 1a<o0L 1b<o2 / 2a<1a2 2b<1b2L & NI/NI & +{-}/{-}{-} & N &  \\
B0659 & 1a<o0L 1b<o2 / 2a<1a2L 2b<1b0L & NI/NN & +{-}/+{-} & N &  \\
B0660 & 1a<o0L 1b<o2 / 2a<1a2L 2b<1b1 & NI/NI & +{-}/+{-} & N &  \\
B0661 & 1a<o0L 1b<o2 / 2a<1a2L 2b<1b1L & NI/NI & +{-}/+{-} & N &  \\
B0662 & 1a<o0L 1b<o2 / 2a<1a2L 2b<1b2 & NI/NI & +{-}/+{-} & N &  \\
B0663 & 1a<o0L 1b<o2 / 2a<1a2L 2b<1b2L & NI/NI & +{-}/+{-} & N &  \\
B0664 & 1a<o0L 1b<o2 / 2a<1b0L 2b<1b1 & NI/NI & +{-}/{-}{-} & N &  \\
B0665 & 1a<o0L 1b<o2 / 2a<1b0L 2b<1b2 & NI/NI & +{-}/{-}{-} & N &  \\
B0666 & 1a<o0L 1b<o2 / 2a<1b1 2b<1b1 & NI/EE & +{-}/{-}{-} & E & W:1b \\
B0667 & 1a<o0L 1b<o2 / 2a<1b1 2b<1b1L & NI/LI & +{-}/{-}{-} & B & CP:2b W:1b,2b \\
B0668 & 1a<o0L 1b<o2 / 2a<1b1 2b<1b2 & NI/EE & +{-}/{-}{-} & E & C:2b W:1b \\
B0669 & 1a<o0L 1b<o2 / 2a<1b1 2b<1b2L & NI/LI & +{-}/{-}{-} & B & P:2b W:1b,2b \\
B0670 & 1a<o0L 1b<o2 / 2a<1b1L 2b<1b2 & NI/IL & +{-}/{-}{-} & B & C:2a,2b P:2a W:1b,2a \\
B0671 & 1a<o0L 1b<o2 / 2a<1b2 2b<1b2 & NI/EE & +{-}/{-}{-} & E & W:1b \\
B0672 & 1a<o0L 1b<o2 / 2a<1b2 2b<1b2L & NI/LI & +{-}/{-}{-} & B & CP:2b W:1b,2b \\
B0673 & 1a<o1 1b<o1 / 2a<1a0L 2b<1a1 & EE/ND & {-}{-}/{-}{-} & E &  \\
B0674 & 1a<o1 1b<o1 / 2a<1a0L 2b<1a2 & EE/ND & {-}{-}/{-}{-} & E &  \\
B0675 & 1a<o1 1b<o1 / 2a<1a0L 2b<1b0L & EE/NN & {-}{-}/{-}{-} & E &  \\
B0676 & 1a<o1 1b<o1 / 2a<1a0L 2b<1b1 & EE/ND & {-}{-}/{-}{-} & E &  \\
B0677 & 1a<o1 1b<o1 / 2a<1a0L 2b<1b1L & EE/ND & {-}{-}/{-}{-} & E &  \\
B0678 & 1a<o1 1b<o1 / 2a<1a0L 2b<1b2 & EE/ND & {-}{-}/{-}{-} & E &  \\
B0679 & 1a<o1 1b<o1 / 2a<1a0L 2b<1b2L & EE/ND & {-}{-}/{-}{-} & E &  \\
B0680 & 1a<o1 1b<o1 / 2a<1a1 2b<1a1 & EE/DD & {-}{-}/{-}{-} & E &  \\
B0681 & 1a<o1 1b<o1 / 2a<1a1 2b<1a1L & EE/DD & {-}{-}/{-}{-} & E &  \\
B0682 & 1a<o1 1b<o1 / 2a<1a1 2b<1a2 & EE/DD & {-}{-}/{-}{-} & E &  \\
B0683 & 1a<o1 1b<o1 / 2a<1a1 2b<1a2L & EE/DD & {-}{-}/{-}{-} & E &  \\
B0684 & 1a<o1 1b<o1 / 2a<1a1 2b<1b1 & EE/DD & {-}{-}/{-}{-} & E &  \\
B0685 & 1a<o1 1b<o1 / 2a<1a1 2b<1b1L & EE/DD & {-}{-}/{-}{-} & E &  \\
B0686 & 1a<o1 1b<o1 / 2a<1a1 2b<1b2 & EE/DD & {-}{-}/{-}{-} & E &  \\
B0687 & 1a<o1 1b<o1 / 2a<1a1 2b<1b2L & EE/DD & {-}{-}/{-}{-} & E &  \\
B0688 & 1a<o1 1b<o1 / 2a<1a1L 2b<1a2 & EE/DD & {-}{-}/{-}{-} & E &  \\
B0689 & 1a<o1 1b<o1 / 2a<1a1L 2b<1b1L & EE/DD & {-}{-}/{-}{-} & E &  \\
B0690 & 1a<o1 1b<o1 / 2a<1a1L 2b<1b2 & EE/DD & {-}{-}/{-}{-} & E &  \\
B0691 & 1a<o1 1b<o1 / 2a<1a1L 2b<1b2L & EE/DD & {-}{-}/{-}{-} & E &  \\
B0692 & 1a<o1 1b<o1 / 2a<1a2 2b<1a2 & EE/DD & {-}{-}/{-}{-} & E &  \\
B0693 & 1a<o1 1b<o1 / 2a<1a2 2b<1a2L & EE/DD & {-}{-}/{-}{-} & E &  \\
B0694 & 1a<o1 1b<o1 / 2a<1a2 2b<1b2 & EE/DD & {-}{-}/{-}{-} & E &  \\
B0695 & 1a<o1 1b<o1 / 2a<1a2 2b<1b2L & EE/DD & {-}{-}/{-}{-} & E &  \\
B0696 & 1a<o1 1b<o1 / 2a<1a2L 2b<1b2L & EE/DD & {-}{-}/{-}{-} & E &  \\
B0697 & 1a<o1 1b<o1L / 2a<1a0L 2b<1a1 & LI/ND & {-}+/{-}{-} & B & CPW:1b \\
B0698 & 1a<o1 1b<o1L / 2a<1a0L 2b<1a2 & LI/ND & {-}+/{-}{-} & B & CPW:1b \\
B0699 & 1a<o1 1b<o1L / 2a<1a0L 2b<1b0L & LI/NN & {-}+/{-}+ & B & CPW:1b \\
B0700 & 1a<o1 1b<o1L / 2a<1a0L 2b<1b1 & LI/NI & {-}+/{-}{-} & B & CPW:1b,2b \\
B0701 & 1a<o1 1b<o1L / 2a<1a0L 2b<1b1L & LI/NI & {-}+/{-}+ & B & CPW:1b,2b \\
B0702 & 1a<o1 1b<o1L / 2a<1a0L 2b<1b2 & LI/NI & {-}+/{-}{-} & B & CPW:1b,2b \\
B0703 & 1a<o1 1b<o1L / 2a<1a0L 2b<1b2L & LI/NI & {-}+/{-}+ & B & CPW:1b,2b \\
B0704 & 1a<o1 1b<o1L / 2a<1a1 2b<1a1 & LI/DD & {-}+/{-}{-} & B & CPW:1b \\
B0705 & 1a<o1 1b<o1L / 2a<1a1 2b<1a1L & LI/DD & {-}+/{-}{-} & B & CPW:1b \\
B0706 & 1a<o1 1b<o1L / 2a<1a1 2b<1a2 & LI/DD & {-}+/{-}{-} & B & CPW:1b \\
B0707 & 1a<o1 1b<o1L / 2a<1a1 2b<1a2L & LI/DD & {-}+/{-}{-} & B & CPW:1b \\
B0708 & 1a<o1 1b<o1L / 2a<1a1 2b<1b0L & LI/DN & {-}+/{-}+ & B & CPW:1b \\
B0709 & 1a<o1 1b<o1L / 2a<1a1 2b<1b1 & LI/DI & {-}+/{-}{-} & B & CPW:1b,2b \\
B0710 & 1a<o1 1b<o1L / 2a<1a1 2b<1b1L & LI/DI & {-}+/{-}+ & B & CPW:1b,2b \\
B0711 & 1a<o1 1b<o1L / 2a<1a1 2b<1b2 & LI/DI & {-}+/{-}{-} & B & CPW:1b,2b \\
B0712 & 1a<o1 1b<o1L / 2a<1a1 2b<1b2L & LI/DI & {-}+/{-}+ & B & CPW:1b,2b \\
B0713 & 1a<o1 1b<o1L / 2a<1a1L 2b<1a2 & LI/DD & {-}+/{-}{-} & B & CPW:1b \\
B0714 & 1a<o1 1b<o1L / 2a<1a1L 2b<1b0L & LI/DN & {-}+/{-}+ & B & CPW:1b \\
B0715 & 1a<o1 1b<o1L / 2a<1a1L 2b<1b1 & LI/DI & {-}+/{-}{-} & B & CPW:1b,2b \\
B0716 & 1a<o1 1b<o1L / 2a<1a1L 2b<1b1L & LI/DI & {-}+/{-}+ & B & CPW:1b,2b \\
B0717 & 1a<o1 1b<o1L / 2a<1a1L 2b<1b2 & LI/DI & {-}+/{-}{-} & B & CPW:1b,2b \\
B0718 & 1a<o1 1b<o1L / 2a<1a1L 2b<1b2L & LI/DI & {-}+/{-}+ & B & CPW:1b,2b \\
B0719 & 1a<o1 1b<o1L / 2a<1a2 2b<1a2 & LI/DD & {-}+/{-}{-} & B & CPW:1b \\
B0720 & 1a<o1 1b<o1L / 2a<1a2 2b<1a2L & LI/DD & {-}+/{-}{-} & B & CPW:1b \\
B0721 & 1a<o1 1b<o1L / 2a<1a2 2b<1b0L & LI/DN & {-}+/{-}+ & B & CPW:1b \\
B0722 & 1a<o1 1b<o1L / 2a<1a2 2b<1b1 & LI/DI & {-}+/{-}{-} & B & CPW:1b,2b \\
B0723 & 1a<o1 1b<o1L / 2a<1a2 2b<1b1L & LI/DI & {-}+/{-}+ & B & CPW:1b,2b \\
B0724 & 1a<o1 1b<o1L / 2a<1a2 2b<1b2 & LI/DI & {-}+/{-}{-} & B & CPW:1b,2b \\
B0725 & 1a<o1 1b<o1L / 2a<1a2 2b<1b2L & LI/DI & {-}+/{-}+ & B & CPW:1b,2b \\
B0726 & 1a<o1 1b<o1L / 2a<1a2L 2b<1b0L & LI/DN & {-}+/{-}+ & B & CPW:1b \\
B0727 & 1a<o1 1b<o1L / 2a<1a2L 2b<1b1 & LI/DI & {-}+/{-}{-} & B & CPW:1b,2b \\
B0728 & 1a<o1 1b<o1L / 2a<1a2L 2b<1b1L & LI/DI & {-}+/{-}+ & B & CPW:1b,2b \\
B0729 & 1a<o1 1b<o1L / 2a<1a2L 2b<1b2 & LI/DI & {-}+/{-}{-} & B & CPW:1b,2b \\
B0730 & 1a<o1 1b<o1L / 2a<1a2L 2b<1b2L & LI/DI & {-}+/{-}+ & B & CPW:1b,2b \\
B0731 & 1a<o1 1b<o1L / 2a<1b0L 2b<1b1 & LI/NI & {-}+/+{-} & B & CPW:1b,2b \\
B0732 & 1a<o1 1b<o1L / 2a<1b0L 2b<1b2 & LI/NI & {-}+/+{-} & B & CPW:1b,2b \\
B0733 & 1a<o1 1b<o1L / 2a<1b1 2b<1b1 & LI/EE & {-}+/{-}{-} & BE & CPW:1b \\
B0734 & 1a<o1 1b<o1L / 2a<1b1 2b<1b1L & LI/LI & {-}+/{-}+ & B & CPW:1b,2b \\
B0735 & 1a<o1 1b<o1L / 2a<1b1 2b<1b2 & LI/EE & {-}+/{-}{-} & BE & CPW:1b \\
B0736 & 1a<o1 1b<o1L / 2a<1b1 2b<1b2L & LI/LI & {-}+/{-}+ & B & CPW:1b,2b \\
B0737 & 1a<o1 1b<o1L / 2a<1b1L 2b<1b2 & LI/IL & {-}+/+{-} & B & CPW:1b,2a \\
B0738 & 1a<o1 1b<o1L / 2a<1b2 2b<1b2 & LI/EE & {-}+/{-}{-} & BE & CPW:1b \\
B0739 & 1a<o1 1b<o1L / 2a<1b2 2b<1b2L & LI/LI & {-}+/{-}+ & B & CPW:1b,2b \\
B0740 & 1a<o1 1b<o2 / 2a<1a0L 2b<1a1 & EE/ND & {-}{-}/{-}{-} & E & C:1b \\
B0741 & 1a<o1 1b<o2 / 2a<1a0L 2b<1a2 & EE/ND & {-}{-}/{-}{-} & E & C:1b \\
B0742 & 1a<o1 1b<o2 / 2a<1a0L 2b<1b0L & EE/NN & {-}{-}/{-}{-} & E & C:1b \\
B0743 & 1a<o1 1b<o2 / 2a<1a0L 2b<1b1 & EE/ND & {-}{-}/{-}{-} & E & C:1b,2b \\
B0744 & 1a<o1 1b<o2 / 2a<1a0L 2b<1b1L & EE/ND & {-}{-}/{-}{-} & E & C:1b,2b \\
B0745 & 1a<o1 1b<o2 / 2a<1a0L 2b<1b2 & EE/ND & {-}{-}/{-}{-} & E & C:1b,2b \\
B0746 & 1a<o1 1b<o2 / 2a<1a0L 2b<1b2L & EE/ND & {-}{-}/{-}{-} & E & C:1b,2b \\
B0747 & 1a<o1 1b<o2 / 2a<1a1 2b<1a1 & EE/DD & {-}{-}/{-}{-} & E & C:1b \\
B0748 & 1a<o1 1b<o2 / 2a<1a1 2b<1a1L & EE/DD & {-}{-}/{-}{-} & E & C:1b \\
B0749 & 1a<o1 1b<o2 / 2a<1a1 2b<1a2 & EE/DD & {-}{-}/{-}{-} & E & C:1b \\
B0750 & 1a<o1 1b<o2 / 2a<1a1 2b<1a2L & EE/DD & {-}{-}/{-}{-} & E & C:1b \\
B0751 & 1a<o1 1b<o2 / 2a<1a1 2b<1b0L & EE/DN & {-}{-}/{-}{-} & E & C:1b \\
B0752 & 1a<o1 1b<o2 / 2a<1a1 2b<1b1 & EE/DD & {-}{-}/{-}{-} & E & C:1b,2b \\
B0753 & 1a<o1 1b<o2 / 2a<1a1 2b<1b1L & EE/DD & {-}{-}/{-}{-} & E & C:1b,2b \\
B0754 & 1a<o1 1b<o2 / 2a<1a1 2b<1b2 & EE/DD & {-}{-}/{-}{-} & E & C:1b,2b \\
B0755 & 1a<o1 1b<o2 / 2a<1a1 2b<1b2L & EE/DD & {-}{-}/{-}{-} & E & C:1b,2b \\
B0756 & 1a<o1 1b<o2 / 2a<1a1L 2b<1a2 & EE/DD & {-}{-}/{-}{-} & E & C:1b \\
B0757 & 1a<o1 1b<o2 / 2a<1a1L 2b<1b0L & EE/DN & {-}{-}/{-}{-} & E & C:1b \\
B0758 & 1a<o1 1b<o2 / 2a<1a1L 2b<1b1 & EE/DD & {-}{-}/{-}{-} & E & C:1b,2b \\
B0759 & 1a<o1 1b<o2 / 2a<1a1L 2b<1b1L & EE/DD & {-}{-}/{-}{-} & E & C:1b,2b \\
B0760 & 1a<o1 1b<o2 / 2a<1a1L 2b<1b2 & EE/DD & {-}{-}/{-}{-} & E & C:1b,2b \\
B0761 & 1a<o1 1b<o2 / 2a<1a1L 2b<1b2L & EE/DD & {-}{-}/{-}{-} & E & C:1b,2b \\
B0762 & 1a<o1 1b<o2 / 2a<1a2 2b<1a2 & EE/DD & {-}{-}/{-}{-} & E & C:1b \\
B0763 & 1a<o1 1b<o2 / 2a<1a2 2b<1a2L & EE/DD & {-}{-}/{-}{-} & E & C:1b \\
B0764 & 1a<o1 1b<o2 / 2a<1a2 2b<1b0L & EE/DN & {-}{-}/{-}{-} & E & C:1b \\
B0765 & 1a<o1 1b<o2 / 2a<1a2 2b<1b1 & EE/DD & {-}{-}/{-}{-} & E & C:1b,2b \\
B0766 & 1a<o1 1b<o2 / 2a<1a2 2b<1b1L & EE/DD & {-}{-}/{-}{-} & E & C:1b,2b \\
B0767 & 1a<o1 1b<o2 / 2a<1a2 2b<1b2 & EE/DD & {-}{-}/{-}{-} & E & C:1b,2b \\
B0768 & 1a<o1 1b<o2 / 2a<1a2 2b<1b2L & EE/DD & {-}{-}/{-}{-} & E & C:1b,2b \\
B0769 & 1a<o1 1b<o2 / 2a<1a2L 2b<1b0L & EE/DN & {-}{-}/{-}{-} & E & C:1b \\
B0770 & 1a<o1 1b<o2 / 2a<1a2L 2b<1b1 & EE/DD & {-}{-}/{-}{-} & E & C:1b,2b \\
B0771 & 1a<o1 1b<o2 / 2a<1a2L 2b<1b1L & EE/DD & {-}{-}/{-}{-} & E & C:1b,2b \\
B0772 & 1a<o1 1b<o2 / 2a<1a2L 2b<1b2 & EE/DD & {-}{-}/{-}{-} & E & C:1b,2b \\
B0773 & 1a<o1 1b<o2 / 2a<1a2L 2b<1b2L & EE/DD & {-}{-}/{-}{-} & E & C:1b,2b \\
B0774 & 1a<o1 1b<o2 / 2a<1b0L 2b<1b1 & EE/ND & {-}{-}/{-}{-} & E & C:1b,2b \\
B0775 & 1a<o1 1b<o2 / 2a<1b0L 2b<1b2 & EE/ND & {-}{-}/{-}{-} & E & C:1b,2b \\
B0776 & 1a<o1 1b<o2 / 2a<1b1 2b<1b1 & EE/DD & {-}{-}/{-}{-} & E & C:1b \\
B0777 & 1a<o1 1b<o2 / 2a<1b1 2b<1b1L & EE/DD & {-}{-}/{-}{-} & E & C:1b \\
B0778 & 1a<o1 1b<o2 / 2a<1b1 2b<1b2 & EE/DD & {-}{-}/{-}{-} & E & C:1b,2b \\
B0779 & 1a<o1 1b<o2 / 2a<1b1 2b<1b2L & EE/DD & {-}{-}/{-}{-} & E & C:1b,2b \\
B0780 & 1a<o1 1b<o2 / 2a<1b1L 2b<1b2 & EE/DD & {-}{-}/{-}{-} & E & C:1b,2b \\
B0781 & 1a<o1 1b<o2 / 2a<1b2 2b<1b2 & EE/DD & {-}{-}/{-}{-} & E & C:1b \\
B0782 & 1a<o1 1b<o2 / 2a<1b2 2b<1b2L & EE/DD & {-}{-}/{-}{-} & E & C:1b \\
B0783 & 1a<o1 1b<o2L / 2a<1a0L 2b<1a1 & LI/ND & {-}+/{-}{-} & B & PW:1b \\
B0784 & 1a<o1 1b<o2L / 2a<1a0L 2b<1a2 & LI/ND & {-}+/{-}{-} & B & PW:1b \\
B0785 & 1a<o1 1b<o2L / 2a<1a0L 2b<1b0L & LI/NN & {-}+/{-}+ & B & PW:1b \\
B0786 & 1a<o1 1b<o2L / 2a<1a0L 2b<1b1 & LI/NI & {-}+/{-}{-} & B & PW:1b,2b \\
B0787 & 1a<o1 1b<o2L / 2a<1a0L 2b<1b1L & LI/NI & {-}+/{-}+ & B & PW:1b,2b \\
B0788 & 1a<o1 1b<o2L / 2a<1a0L 2b<1b2 & LI/NI & {-}+/{-}{-} & B & PW:1b,2b \\
B0789 & 1a<o1 1b<o2L / 2a<1a0L 2b<1b2L & LI/NI & {-}+/{-}+ & B & PW:1b,2b \\
B0790 & 1a<o1 1b<o2L / 2a<1a1 2b<1a1 & LI/DD & {-}+/{-}{-} & B & PW:1b \\
B0791 & 1a<o1 1b<o2L / 2a<1a1 2b<1a1L & LI/DD & {-}+/{-}{-} & B & PW:1b \\
B0792 & 1a<o1 1b<o2L / 2a<1a1 2b<1a2 & LI/DD & {-}+/{-}{-} & B & PW:1b \\
B0793 & 1a<o1 1b<o2L / 2a<1a1 2b<1a2L & LI/DD & {-}+/{-}{-} & B & PW:1b \\
B0794 & 1a<o1 1b<o2L / 2a<1a1 2b<1b0L & LI/DN & {-}+/{-}+ & B & PW:1b \\
B0795 & 1a<o1 1b<o2L / 2a<1a1 2b<1b1 & LI/DI & {-}+/{-}{-} & B & PW:1b,2b \\
B0796 & 1a<o1 1b<o2L / 2a<1a1 2b<1b1L & LI/DI & {-}+/{-}+ & B & PW:1b,2b \\
B0797 & 1a<o1 1b<o2L / 2a<1a1 2b<1b2 & LI/DI & {-}+/{-}{-} & B & PW:1b,2b \\
B0798 & 1a<o1 1b<o2L / 2a<1a1 2b<1b2L & LI/DI & {-}+/{-}+ & B & PW:1b,2b \\
B0799 & 1a<o1 1b<o2L / 2a<1a1L 2b<1a2 & LI/DD & {-}+/{-}{-} & B & PW:1b \\
B0800 & 1a<o1 1b<o2L / 2a<1a1L 2b<1b0L & LI/DN & {-}+/{-}+ & B & PW:1b \\
B0801 & 1a<o1 1b<o2L / 2a<1a1L 2b<1b1 & LI/DI & {-}+/{-}{-} & B & PW:1b,2b \\
B0802 & 1a<o1 1b<o2L / 2a<1a1L 2b<1b1L & LI/DI & {-}+/{-}+ & B & PW:1b,2b \\
B0803 & 1a<o1 1b<o2L / 2a<1a1L 2b<1b2 & LI/DI & {-}+/{-}{-} & B & PW:1b,2b \\
B0804 & 1a<o1 1b<o2L / 2a<1a1L 2b<1b2L & LI/DI & {-}+/{-}+ & B & PW:1b,2b \\
B0805 & 1a<o1 1b<o2L / 2a<1a2 2b<1a2 & LI/DD & {-}+/{-}{-} & B & PW:1b \\
B0806 & 1a<o1 1b<o2L / 2a<1a2 2b<1a2L & LI/DD & {-}+/{-}{-} & B & PW:1b \\
B0807 & 1a<o1 1b<o2L / 2a<1a2 2b<1b0L & LI/DN & {-}+/{-}+ & B & PW:1b \\
B0808 & 1a<o1 1b<o2L / 2a<1a2 2b<1b1 & LI/DI & {-}+/{-}{-} & B & PW:1b,2b \\
B0809 & 1a<o1 1b<o2L / 2a<1a2 2b<1b1L & LI/DI & {-}+/{-}+ & B & PW:1b,2b \\
B0810 & 1a<o1 1b<o2L / 2a<1a2 2b<1b2 & LI/DI & {-}+/{-}{-} & B & PW:1b,2b \\
B0811 & 1a<o1 1b<o2L / 2a<1a2 2b<1b2L & LI/DI & {-}+/{-}+ & B & PW:1b,2b \\
B0812 & 1a<o1 1b<o2L / 2a<1a2L 2b<1b0L & LI/DN & {-}+/{-}+ & B & PW:1b \\
B0813 & 1a<o1 1b<o2L / 2a<1a2L 2b<1b1 & LI/DI & {-}+/{-}{-} & B & PW:1b,2b \\
B0814 & 1a<o1 1b<o2L / 2a<1a2L 2b<1b1L & LI/DI & {-}+/{-}+ & B & PW:1b,2b \\
B0815 & 1a<o1 1b<o2L / 2a<1a2L 2b<1b2 & LI/DI & {-}+/{-}{-} & B & PW:1b,2b \\
B0816 & 1a<o1 1b<o2L / 2a<1a2L 2b<1b2L & LI/DI & {-}+/{-}+ & B & PW:1b,2b \\
B0817 & 1a<o1 1b<o2L / 2a<1b0L 2b<1b1 & LI/NI & {-}+/+{-} & B & PW:1b,2b \\
B0818 & 1a<o1 1b<o2L / 2a<1b0L 2b<1b2 & LI/NI & {-}+/+{-} & B & PW:1b,2b \\
B0819 & 1a<o1 1b<o2L / 2a<1b1 2b<1b1 & LI/EE & {-}+/{-}{-} & BE & PW:1b \\
B0820 & 1a<o1 1b<o2L / 2a<1b1 2b<1b1L & LI/LI & {-}+/{-}+ & B & C:2b PW:1b,2b \\
B0821 & 1a<o1 1b<o2L / 2a<1b1 2b<1b2 & LI/EE & {-}+/{-}{-} & BE & C:2b PW:1b \\
B0822 & 1a<o1 1b<o2L / 2a<1b1 2b<1b2L & LI/LI & {-}+/{-}+ & B & PW:1b,2b \\
B0823 & 1a<o1 1b<o2L / 2a<1b1L 2b<1b2 & LI/IL & {-}+/+{-} & B & C:2a,2b PW:1b,2a \\
B0824 & 1a<o1 1b<o2L / 2a<1b2 2b<1b2 & LI/EE & {-}+/{-}{-} & BE & PW:1b \\
B0825 & 1a<o1 1b<o2L / 2a<1b2 2b<1b2L & LI/LI & {-}+/{-}+ & B & C:2b PW:1b,2b \\
B0826 & 1a<o1L 1b<o2 / 2a<1a0L 2b<1a1 & IL/NI & +{-}/+{-} & B & C:1a,1b,2b PW:1a,2b \\
B0827 & 1a<o1L 1b<o2 / 2a<1a0L 2b<1a2 & IL/NI & +{-}/+{-} & B & C:1a,1b,2b PW:1a,2b \\
B0828 & 1a<o1L 1b<o2 / 2a<1a0L 2b<1b0L & IL/NN & +{-}/+{-} & B & C:1a,1b PW:1a \\
B0829 & 1a<o1L 1b<o2 / 2a<1a0L 2b<1b1 & IL/ND & +{-}/+{-} & B & C:1a,1b,2b PW:1a \\
B0830 & 1a<o1L 1b<o2 / 2a<1a0L 2b<1b1L & IL/ND & +{-}/+{-} & B & C:1a,1b,2b PW:1a \\
B0831 & 1a<o1L 1b<o2 / 2a<1a0L 2b<1b2 & IL/ND & +{-}/+{-} & B & C:1a,1b,2b PW:1a \\
B0832 & 1a<o1L 1b<o2 / 2a<1a0L 2b<1b2L & IL/ND & +{-}/+{-} & B & C:1a,1b,2b PW:1a \\
B0833 & 1a<o1L 1b<o2 / 2a<1a1 2b<1a1 & IL/EE & +{-}/{-}{-} & BE & C:1a,1b PW:1a \\
B0834 & 1a<o1L 1b<o2 / 2a<1a1 2b<1a1L & IL/LI & +{-}/{-}+ & B & C:1a,1b,2b PW:1a,2b \\
B0835 & 1a<o1L 1b<o2 / 2a<1a1 2b<1a2 & IL/EE & +{-}/{-}{-} & BE & C:1a,1b PW:1a \\
B0836 & 1a<o1L 1b<o2 / 2a<1a1 2b<1a2L & IL/LI & +{-}/{-}+ & B & C:1a,1b,2b PW:1a,2b \\
B0837 & 1a<o1L 1b<o2 / 2a<1a1 2b<1b0L & IL/IN & +{-}/{-}{-} & B & C:1a,1b,2a PW:1a,2a \\
B0838 & 1a<o1L 1b<o2 / 2a<1a1 2b<1b1 & IL/ID & +{-}/{-}{-} & B & C:1a,1b,2a,2b PW:1a,2a \\
B0839 & 1a<o1L 1b<o2 / 2a<1a1 2b<1b1L & IL/ID & +{-}/{-}{-} & B & C:1a,1b,2a,2b PW:1a,2a \\
B0840 & 1a<o1L 1b<o2 / 2a<1a1 2b<1b2 & IL/ID & +{-}/{-}{-} & B & C:1a,1b,2a,2b PW:1a,2a \\
B0841 & 1a<o1L 1b<o2 / 2a<1a1 2b<1b2L & IL/ID & +{-}/{-}{-} & B & C:1a,1b,2a,2b PW:1a,2a \\
B0842 & 1a<o1L 1b<o2 / 2a<1a1L 2b<1a2 & IL/IL & +{-}/+{-} & B & C:1a,1b,2a PW:1a,2a \\
B0843 & 1a<o1L 1b<o2 / 2a<1a1L 2b<1b0L & IL/IN & +{-}/+{-} & B & C:1a,1b,2a PW:1a,2a \\
B0844 & 1a<o1L 1b<o2 / 2a<1a1L 2b<1b1 & IL/ID & +{-}/+{-} & B & C:1a,1b,2a,2b PW:1a,2a \\
B0845 & 1a<o1L 1b<o2 / 2a<1a1L 2b<1b1L & IL/ID & +{-}/+{-} & B & C:1a,1b,2a,2b PW:1a,2a \\
B0846 & 1a<o1L 1b<o2 / 2a<1a1L 2b<1b2 & IL/ID & +{-}/+{-} & B & C:1a,1b,2a,2b PW:1a,2a \\
B0847 & 1a<o1L 1b<o2 / 2a<1a1L 2b<1b2L & IL/ID & +{-}/+{-} & B & C:1a,1b,2a,2b PW:1a,2a \\
B0848 & 1a<o1L 1b<o2 / 2a<1a2 2b<1a2 & IL/EE & +{-}/{-}{-} & BE & C:1a,1b PW:1a \\
B0849 & 1a<o1L 1b<o2 / 2a<1a2 2b<1a2L & IL/LI & +{-}/{-}+ & B & C:1a,1b,2b PW:1a,2b \\
B0850 & 1a<o1L 1b<o2 / 2a<1a2 2b<1b0L & IL/IN & +{-}/{-}{-} & B & C:1a,1b,2a PW:1a,2a \\
B0851 & 1a<o1L 1b<o2 / 2a<1a2 2b<1b1 & IL/ID & +{-}/{-}{-} & B & C:1a,1b,2a,2b PW:1a,2a \\
B0852 & 1a<o1L 1b<o2 / 2a<1a2 2b<1b1L & IL/ID & +{-}/{-}{-} & B & C:1a,1b,2a,2b PW:1a,2a \\
B0853 & 1a<o1L 1b<o2 / 2a<1a2 2b<1b2 & IL/ID & +{-}/{-}{-} & B & C:1a,1b,2a,2b PW:1a,2a \\
B0854 & 1a<o1L 1b<o2 / 2a<1a2 2b<1b2L & IL/ID & +{-}/{-}{-} & B & C:1a,1b,2a,2b PW:1a,2a \\
B0855 & 1a<o1L 1b<o2 / 2a<1a2L 2b<1b0L & IL/IN & +{-}/+{-} & B & C:1a,1b,2a PW:1a,2a \\
B0856 & 1a<o1L 1b<o2 / 2a<1a2L 2b<1b1 & IL/ID & +{-}/+{-} & B & C:1a,1b,2a,2b PW:1a,2a \\
B0857 & 1a<o1L 1b<o2 / 2a<1a2L 2b<1b1L & IL/ID & +{-}/+{-} & B & C:1a,1b,2a,2b PW:1a,2a \\
B0858 & 1a<o1L 1b<o2 / 2a<1a2L 2b<1b2 & IL/ID & +{-}/+{-} & B & C:1a,1b,2a,2b PW:1a,2a \\
B0859 & 1a<o1L 1b<o2 / 2a<1a2L 2b<1b2L & IL/ID & +{-}/+{-} & B & C:1a,1b,2a,2b PW:1a,2a \\
B0860 & 1a<o1L 1b<o2 / 2a<1b0L 2b<1b1 & IL/ND & +{-}/{-}{-} & B & C:1a,1b,2b PW:1a \\
B0861 & 1a<o1L 1b<o2 / 2a<1b0L 2b<1b2 & IL/ND & +{-}/{-}{-} & B & C:1a,1b,2b PW:1a \\
B0862 & 1a<o1L 1b<o2 / 2a<1b1 2b<1b1 & IL/DD & +{-}/{-}{-} & B & C:1a,1b PW:1a \\
B0863 & 1a<o1L 1b<o2 / 2a<1b1 2b<1b1L & IL/DD & +{-}/{-}{-} & B & C:1a,1b PW:1a \\
B0864 & 1a<o1L 1b<o2 / 2a<1b1 2b<1b2 & IL/DD & +{-}/{-}{-} & B & C:1a,1b,2b PW:1a \\
B0865 & 1a<o1L 1b<o2 / 2a<1b1 2b<1b2L & IL/DD & +{-}/{-}{-} & B & C:1a,1b,2b PW:1a \\
B0866 & 1a<o1L 1b<o2 / 2a<1b1L 2b<1b2 & IL/DD & +{-}/{-}{-} & B & C:1a,1b,2b PW:1a \\
B0867 & 1a<o1L 1b<o2 / 2a<1b2 2b<1b2 & IL/DD & +{-}/{-}{-} & B & C:1a,1b PW:1a \\
B0868 & 1a<o1L 1b<o2 / 2a<1b2 2b<1b2L & IL/DD & +{-}/{-}{-} & B & C:1a,1b PW:1a \\
B0869 & 1a<o2 1b<o2 / 2a<1a0L 2b<1a1 & EE/ND & {-}{-}/{-}{-} & E &  \\
B0870 & 1a<o2 1b<o2 / 2a<1a0L 2b<1a2 & EE/ND & {-}{-}/{-}{-} & E &  \\
B0871 & 1a<o2 1b<o2 / 2a<1a0L 2b<1b0L & EE/NN & {-}{-}/{-}{-} & E &  \\
B0872 & 1a<o2 1b<o2 / 2a<1a0L 2b<1b1 & EE/ND & {-}{-}/{-}{-} & E &  \\
B0873 & 1a<o2 1b<o2 / 2a<1a0L 2b<1b1L & EE/ND & {-}{-}/{-}{-} & E &  \\
B0874 & 1a<o2 1b<o2 / 2a<1a0L 2b<1b2 & EE/ND & {-}{-}/{-}{-} & E &  \\
B0875 & 1a<o2 1b<o2 / 2a<1a0L 2b<1b2L & EE/ND & {-}{-}/{-}{-} & E &  \\
B0876 & 1a<o2 1b<o2 / 2a<1a1 2b<1a1 & EE/DD & {-}{-}/{-}{-} & E &  \\
B0877 & 1a<o2 1b<o2 / 2a<1a1 2b<1a1L & EE/DD & {-}{-}/{-}{-} & E &  \\
B0878 & 1a<o2 1b<o2 / 2a<1a1 2b<1a2 & EE/DD & {-}{-}/{-}{-} & E &  \\
B0879 & 1a<o2 1b<o2 / 2a<1a1 2b<1a2L & EE/DD & {-}{-}/{-}{-} & E &  \\
B0880 & 1a<o2 1b<o2 / 2a<1a1 2b<1b1 & EE/DD & {-}{-}/{-}{-} & E &  \\
B0881 & 1a<o2 1b<o2 / 2a<1a1 2b<1b1L & EE/DD & {-}{-}/{-}{-} & E &  \\
B0882 & 1a<o2 1b<o2 / 2a<1a1 2b<1b2 & EE/DD & {-}{-}/{-}{-} & E &  \\
B0883 & 1a<o2 1b<o2 / 2a<1a1 2b<1b2L & EE/DD & {-}{-}/{-}{-} & E &  \\
B0884 & 1a<o2 1b<o2 / 2a<1a1L 2b<1a2 & EE/DD & {-}{-}/{-}{-} & E &  \\
B0885 & 1a<o2 1b<o2 / 2a<1a1L 2b<1b1L & EE/DD & {-}{-}/{-}{-} & E &  \\
B0886 & 1a<o2 1b<o2 / 2a<1a1L 2b<1b2 & EE/DD & {-}{-}/{-}{-} & E &  \\
B0887 & 1a<o2 1b<o2 / 2a<1a1L 2b<1b2L & EE/DD & {-}{-}/{-}{-} & E &  \\
B0888 & 1a<o2 1b<o2 / 2a<1a2 2b<1a2 & EE/DD & {-}{-}/{-}{-} & E &  \\
B0889 & 1a<o2 1b<o2 / 2a<1a2 2b<1a2L & EE/DD & {-}{-}/{-}{-} & E &  \\
B0890 & 1a<o2 1b<o2 / 2a<1a2 2b<1b2 & EE/DD & {-}{-}/{-}{-} & E &  \\
B0891 & 1a<o2 1b<o2 / 2a<1a2 2b<1b2L & EE/DD & {-}{-}/{-}{-} & E &  \\
B0892 & 1a<o2 1b<o2 / 2a<1a2L 2b<1b2L & EE/DD & {-}{-}/{-}{-} & E &  \\
B0893 & 1a<o2 1b<o2L / 2a<1a0L 2b<1a1 & LI/ND & {-}+/{-}{-} & B & CPW:1b \\
B0894 & 1a<o2 1b<o2L / 2a<1a0L 2b<1a2 & LI/ND & {-}+/{-}{-} & B & CPW:1b \\
B0895 & 1a<o2 1b<o2L / 2a<1a0L 2b<1b0L & LI/NN & {-}+/{-}+ & B & CPW:1b \\
B0896 & 1a<o2 1b<o2L / 2a<1a0L 2b<1b1 & LI/NI & {-}+/{-}{-} & B & CPW:1b,2b \\
B0897 & 1a<o2 1b<o2L / 2a<1a0L 2b<1b1L & LI/NI & {-}+/{-}+ & B & CPW:1b,2b \\
B0898 & 1a<o2 1b<o2L / 2a<1a0L 2b<1b2 & LI/NI & {-}+/{-}{-} & B & CPW:1b,2b \\
B0899 & 1a<o2 1b<o2L / 2a<1a0L 2b<1b2L & LI/NI & {-}+/{-}+ & B & CPW:1b,2b \\
B0900 & 1a<o2 1b<o2L / 2a<1a1 2b<1a1 & LI/DD & {-}+/{-}{-} & B & CPW:1b \\
B0901 & 1a<o2 1b<o2L / 2a<1a1 2b<1a1L & LI/DD & {-}+/{-}{-} & B & CPW:1b \\
B0902 & 1a<o2 1b<o2L / 2a<1a1 2b<1a2 & LI/DD & {-}+/{-}{-} & B & CPW:1b \\
B0903 & 1a<o2 1b<o2L / 2a<1a1 2b<1a2L & LI/DD & {-}+/{-}{-} & B & CPW:1b \\
B0904 & 1a<o2 1b<o2L / 2a<1a1 2b<1b0L & LI/DN & {-}+/{-}+ & B & CPW:1b \\
B0905 & 1a<o2 1b<o2L / 2a<1a1 2b<1b1 & LI/DI & {-}+/{-}{-} & B & CPW:1b,2b \\
B0906 & 1a<o2 1b<o2L / 2a<1a1 2b<1b1L & LI/DI & {-}+/{-}+ & B & CPW:1b,2b \\
B0907 & 1a<o2 1b<o2L / 2a<1a1 2b<1b2 & LI/DI & {-}+/{-}{-} & B & CPW:1b,2b \\
B0908 & 1a<o2 1b<o2L / 2a<1a1 2b<1b2L & LI/DI & {-}+/{-}+ & B & CPW:1b,2b \\
B0909 & 1a<o2 1b<o2L / 2a<1a1L 2b<1a2 & LI/DD & {-}+/{-}{-} & B & CPW:1b \\
B0910 & 1a<o2 1b<o2L / 2a<1a1L 2b<1b0L & LI/DN & {-}+/{-}+ & B & CPW:1b \\
B0911 & 1a<o2 1b<o2L / 2a<1a1L 2b<1b1 & LI/DI & {-}+/{-}{-} & B & CPW:1b,2b \\
B0912 & 1a<o2 1b<o2L / 2a<1a1L 2b<1b1L & LI/DI & {-}+/{-}+ & B & CPW:1b,2b \\
B0913 & 1a<o2 1b<o2L / 2a<1a1L 2b<1b2 & LI/DI & {-}+/{-}{-} & B & CPW:1b,2b \\
B0914 & 1a<o2 1b<o2L / 2a<1a1L 2b<1b2L & LI/DI & {-}+/{-}+ & B & CPW:1b,2b \\
B0915 & 1a<o2 1b<o2L / 2a<1a2 2b<1a2 & LI/DD & {-}+/{-}{-} & B & CPW:1b \\
B0916 & 1a<o2 1b<o2L / 2a<1a2 2b<1a2L & LI/DD & {-}+/{-}{-} & B & CPW:1b \\
B0917 & 1a<o2 1b<o2L / 2a<1a2 2b<1b0L & LI/DN & {-}+/{-}+ & B & CPW:1b \\
B0918 & 1a<o2 1b<o2L / 2a<1a2 2b<1b1 & LI/DI & {-}+/{-}{-} & B & CPW:1b,2b \\
B0919 & 1a<o2 1b<o2L / 2a<1a2 2b<1b1L & LI/DI & {-}+/{-}+ & B & CPW:1b,2b \\
B0920 & 1a<o2 1b<o2L / 2a<1a2 2b<1b2 & LI/DI & {-}+/{-}{-} & B & CPW:1b,2b \\
B0921 & 1a<o2 1b<o2L / 2a<1a2 2b<1b2L & LI/DI & {-}+/{-}+ & B & CPW:1b,2b \\
B0922 & 1a<o2 1b<o2L / 2a<1a2L 2b<1b0L & LI/DN & {-}+/{-}+ & B & CPW:1b \\
B0923 & 1a<o2 1b<o2L / 2a<1a2L 2b<1b1 & LI/DI & {-}+/{-}{-} & B & CPW:1b,2b \\
B0924 & 1a<o2 1b<o2L / 2a<1a2L 2b<1b1L & LI/DI & {-}+/{-}+ & B & CPW:1b,2b \\
B0925 & 1a<o2 1b<o2L / 2a<1a2L 2b<1b2 & LI/DI & {-}+/{-}{-} & B & CPW:1b,2b \\
B0926 & 1a<o2 1b<o2L / 2a<1a2L 2b<1b2L & LI/DI & {-}+/{-}+ & B & CPW:1b,2b \\
B0927 & 1a<o2 1b<o2L / 2a<1b0L 2b<1b1 & LI/NI & {-}+/+{-} & B & CPW:1b,2b \\
B0928 & 1a<o2 1b<o2L / 2a<1b0L 2b<1b2 & LI/NI & {-}+/+{-} & B & CPW:1b,2b \\
B0929 & 1a<o2 1b<o2L / 2a<1b1 2b<1b1 & LI/EE & {-}+/{-}{-} & BE & CPW:1b \\
B0930 & 1a<o2 1b<o2L / 2a<1b1 2b<1b1L & LI/LI & {-}+/{-}+ & B & CPW:1b,2b \\
B0931 & 1a<o2 1b<o2L / 2a<1b1 2b<1b2 & LI/EE & {-}+/{-}{-} & BE & CPW:1b \\
B0932 & 1a<o2 1b<o2L / 2a<1b1 2b<1b2L & LI/LI & {-}+/{-}+ & B & CPW:1b,2b \\
B0933 & 1a<o2 1b<o2L / 2a<1b1L 2b<1b2 & LI/IL & {-}+/+{-} & B & CPW:1b,2a \\
B0934 & 1a<o2 1b<o2L / 2a<1b2 2b<1b2 & LI/EE & {-}+/{-}{-} & BE & CPW:1b \\
B0935 & 1a<o2 1b<o2L / 2a<1b2 2b<1b2L & LI/LI & {-}+/{-}+ & B & CPW:1b,2b \\
B0936 & 1a<o0L 1b<o1 / 2a<1a0L 2b<1a1,1b0L & NI/NN & +{-}/+{-} & N &  \\
B0937 & 1a<o0L 1b<o1 / 2a<1a0L 2b<1a1,1b1 & NI/NE & +{-}/+{-} & F & W:2b \\
B0938 & 1a<o0L 1b<o1 / 2a<1a0L 2b<1a1,1b1L & NI/NI & +{-}/+{-} & F & CP:2b \\
B0939 & 1a<o0L 1b<o1 / 2a<1a0L 2b<1a1,1b2 & NI/NE & +{-}/+{-} & F & CW:2b \\
B0940 & 1a<o0L 1b<o1 / 2a<1a0L 2b<1a1,1b2L & NI/NI & +{-}/+{-} & F & P:2b \\
B0941 & 1a<o0L 1b<o1 / 2a<1a0L 2b<1a2,1b0L & NI/NN & +{-}/+{-} & N &  \\
B0942 & 1a<o0L 1b<o1 / 2a<1a0L 2b<1a2,1b1 & NI/NE & +{-}/+{-} & F & W:2b \\
B0943 & 1a<o0L 1b<o1 / 2a<1a0L 2b<1a2,1b1L & NI/NI & +{-}/+{-} & F & CP:2b \\
B0944 & 1a<o0L 1b<o1 / 2a<1a0L 2b<1a2,1b2 & NI/NE & +{-}/+{-} & F & W:2b \\
B0945 & 1a<o0L 1b<o1 / 2a<1a0L 2b<1a2,1b2L & NI/NI & +{-}/+{-} & F & CP:2b \\
B0946 & 1a<o0L 1b<o1 / 2a<1a0L,1b1 2b<1a1 & NI/IN & +{-}/+{-} & N &  \\
B0947 & 1a<o0L 1b<o1 / 2a<1a0L,1b1 2b<1a2 & NI/IN & +{-}/+{-} & N &  \\
B0948 & 1a<o0L 1b<o1 / 2a<1a0L,1b1 2b<1b0L & NI/IN & +{-}/+{-} & N &  \\
B0949 & 1a<o0L 1b<o1 / 2a<1a0L,1b1 2b<1b1 & NI/EE & +{-}/+{-} & E & W:1b \\
B0950 & 1a<o0L 1b<o1 / 2a<1a0L,1b1 2b<1b1L & NI/LI & +{-}/+{-} & B & CP:2b W:1b,2b \\
B0951 & 1a<o0L 1b<o1 / 2a<1a0L,1b1 2b<1b2 & NI/EE & +{-}/+{-} & E & C:2b W:1b \\
B0952 & 1a<o0L 1b<o1 / 2a<1a0L,1b1 2b<1b2L & NI/LI & +{-}/+{-} & B & P:2b W:1b,2b \\
B0953 & 1a<o0L 1b<o1 / 2a<1a0L,1b2 2b<1a1 & NI/IN & +{-}/+{-} & N &  \\
B0954 & 1a<o0L 1b<o1 / 2a<1a0L,1b2 2b<1a2 & NI/IN & +{-}/+{-} & N &  \\
B0955 & 1a<o0L 1b<o1 / 2a<1a0L,1b2 2b<1b0L & NI/IN & +{-}/+{-} & N &  \\
B0956 & 1a<o0L 1b<o1 / 2a<1a0L,1b2 2b<1b1 & NI/EE & +{-}/+{-} & E & C:2a W:1b \\
B0957 & 1a<o0L 1b<o1 / 2a<1a0L,1b2 2b<1b1L & NI/LI & +{-}/+{-} & B & C:2a,2b P:2b W:1b,2b \\
B0958 & 1a<o0L 1b<o1 / 2a<1a0L,1b2 2b<1b2 & NI/EE & +{-}/+{-} & E & W:1b \\
B0959 & 1a<o0L 1b<o1 / 2a<1a0L,1b2 2b<1b2L & NI/LI & +{-}/+{-} & B & CP:2b W:1b,2b \\
B0960 & 1a<o0L 1b<o1 / 2a<1a1 2b<1a1,1b0L & NI/NN & +{-}/{-}{-} & N &  \\
B0961 & 1a<o0L 1b<o1 / 2a<1a1 2b<1a1,1b1 & NI/NE & +{-}/{-}{-} & F & W:2b \\
B0962 & 1a<o0L 1b<o1 / 2a<1a1 2b<1a1,1b1L & NI/NI & +{-}/{-}{-} & F & CP:2b \\
B0963 & 1a<o0L 1b<o1 / 2a<1a1 2b<1a1,1b2 & NI/NE & +{-}/{-}{-} & F & CW:2b \\
B0964 & 1a<o0L 1b<o1 / 2a<1a1 2b<1a1,1b2L & NI/NI & +{-}/{-}{-} & F & P:2b \\
B0965 & 1a<o0L 1b<o1 / 2a<1a1 2b<1a1L,1b1 & NI/NE & +{-}/{-}+ & F & W:2b \\
B0966 & 1a<o0L 1b<o1 / 2a<1a1 2b<1a1L,1b2 & NI/NE & +{-}/{-}+ & F & CW:2b \\
B0967 & 1a<o0L 1b<o1 / 2a<1a1 2b<1a2,1b0L & NI/NN & +{-}/{-}{-} & N &  \\
B0968 & 1a<o0L 1b<o1 / 2a<1a1 2b<1a2,1b1 & NI/NE & +{-}/{-}{-} & F & W:2b \\
B0969 & 1a<o0L 1b<o1 / 2a<1a1 2b<1a2,1b1L & NI/NI & +{-}/{-}{-} & F & CP:2b \\
B0970 & 1a<o0L 1b<o1 / 2a<1a1 2b<1a2,1b2 & NI/NE & +{-}/{-}{-} & F & W:2b \\
B0971 & 1a<o0L 1b<o1 / 2a<1a1 2b<1a2,1b2L & NI/NI & +{-}/{-}{-} & F & CP:2b \\
B0972 & 1a<o0L 1b<o1 / 2a<1a1 2b<1a2L,1b1 & NI/NE & +{-}/{-}+ & F & W:2b \\
B0973 & 1a<o0L 1b<o1 / 2a<1a1 2b<1a2L,1b2 & NI/NE & +{-}/{-}+ & F & W:2b \\
B0974 & 1a<o0L 1b<o1 / 2a<1a1,1b0L 2b<1a1L & NI/NN & +{-}/{-}+ & N &  \\
B0975 & 1a<o0L 1b<o1 / 2a<1a1,1b0L 2b<1a2 & NI/NN & +{-}/{-}{-} & N &  \\
B0976 & 1a<o0L 1b<o1 / 2a<1a1,1b0L 2b<1a2L & NI/NN & +{-}/{-}+ & N &  \\
B0977 & 1a<o0L 1b<o1 / 2a<1a1,1b0L 2b<1b1 & NI/NI & +{-}/{-}{-} & N &  \\
B0978 & 1a<o0L 1b<o1 / 2a<1a1,1b0L 2b<1b2 & NI/NI & +{-}/{-}{-} & N &  \\
B0979 & 1a<o0L 1b<o1 / 2a<1a1,1b1 2b<1a1L & NI/EN & +{-}/{-}+ & F & W:2a \\
B0980 & 1a<o0L 1b<o1 / 2a<1a1,1b1 2b<1a2 & NI/EN & +{-}/{-}{-} & F & W:2a \\
B0981 & 1a<o0L 1b<o1 / 2a<1a1,1b1 2b<1a2L & NI/EN & +{-}/{-}+ & F & W:2a \\
B0982 & 1a<o0L 1b<o1 / 2a<1a1,1b1 2b<1b0L & NI/EN & +{-}/{-}{-} & F & W:2a \\
B0983 & 1a<o0L 1b<o1 / 2a<1a1,1b1 2b<1b1 & NI/EE & +{-}/{-}{-} & E & W:1b \\
B0984 & 1a<o0L 1b<o1 / 2a<1a1,1b1 2b<1b1L & NI/LI & +{-}/{-}{-} & B & CP:2b W:1b,2b \\
B0985 & 1a<o0L 1b<o1 / 2a<1a1,1b1 2b<1b2 & NI/EE & +{-}/{-}{-} & E & C:2b W:1b \\
B0986 & 1a<o0L 1b<o1 / 2a<1a1,1b1 2b<1b2L & NI/LI & +{-}/{-}{-} & B & P:2b W:1b,2b \\
B0987 & 1a<o0L 1b<o1 / 2a<1a1,1b1L 2b<1a1L & NI/IN & +{-}/{-}+ & F & CP:2a \\
B0988 & 1a<o0L 1b<o1 / 2a<1a1,1b1L 2b<1a2 & NI/IN & +{-}/{-}{-} & F & CP:2a \\
B0989 & 1a<o0L 1b<o1 / 2a<1a1,1b1L 2b<1a2L & NI/IN & +{-}/{-}+ & F & CP:2a \\
B0990 & 1a<o0L 1b<o1 / 2a<1a1,1b1L 2b<1b1 & NI/IL & +{-}/{-}{-} & B & CP:2a W:1b,2a \\
B0991 & 1a<o0L 1b<o1 / 2a<1a1,1b1L 2b<1b2 & NI/IL & +{-}/{-}{-} & B & C:2a,2b P:2a W:1b,2a \\
B0992 & 1a<o0L 1b<o1 / 2a<1a1,1b2 2b<1a1L & NI/EN & +{-}/{-}+ & F & CW:2a \\
B0993 & 1a<o0L 1b<o1 / 2a<1a1,1b2 2b<1a2 & NI/EN & +{-}/{-}{-} & F & CW:2a \\
B0994 & 1a<o0L 1b<o1 / 2a<1a1,1b2 2b<1a2L & NI/EN & +{-}/{-}+ & F & CW:2a \\
B0995 & 1a<o0L 1b<o1 / 2a<1a1,1b2 2b<1b0L & NI/EN & +{-}/{-}{-} & F & CW:2a \\
B0996 & 1a<o0L 1b<o1 / 2a<1a1,1b2 2b<1b1 & NI/EE & +{-}/{-}{-} & E & C:2a W:1b \\
B0997 & 1a<o0L 1b<o1 / 2a<1a1,1b2 2b<1b1L & NI/LI & +{-}/{-}{-} & B & C:2a,2b P:2b W:1b,2b \\
B0998 & 1a<o0L 1b<o1 / 2a<1a1,1b2 2b<1b2 & NI/EE & +{-}/{-}{-} & E & W:1b \\
B0999 & 1a<o0L 1b<o1 / 2a<1a1,1b2 2b<1b2L & NI/LI & +{-}/{-}{-} & B & CP:2b W:1b,2b \\
B1000 & 1a<o0L 1b<o1 / 2a<1a1,1b2L 2b<1a1L & NI/IN & +{-}/{-}+ & F & P:2a \\
B1001 & 1a<o0L 1b<o1 / 2a<1a1,1b2L 2b<1a2 & NI/IN & +{-}/{-}{-} & F & P:2a \\
B1002 & 1a<o0L 1b<o1 / 2a<1a1,1b2L 2b<1a2L & NI/IN & +{-}/{-}+ & F & P:2a \\
B1003 & 1a<o0L 1b<o1 / 2a<1a1,1b2L 2b<1b1 & NI/IL & +{-}/{-}{-} & B & P:2a W:1b,2a \\
B1004 & 1a<o0L 1b<o1 / 2a<1a1,1b2L 2b<1b2 & NI/IL & +{-}/{-}{-} & B & CP:2a W:1b,2a \\
B1005 & 1a<o0L 1b<o1 / 2a<1a1L 2b<1a2,1b0L & NI/NN & +{-}/+{-} & N &  \\
B1006 & 1a<o0L 1b<o1 / 2a<1a1L 2b<1a2,1b1 & NI/NE & +{-}/+{-} & F & W:2b \\
B1007 & 1a<o0L 1b<o1 / 2a<1a1L 2b<1a2,1b1L & NI/NI & +{-}/+{-} & F & CP:2b \\
B1008 & 1a<o0L 1b<o1 / 2a<1a1L 2b<1a2,1b2 & NI/NE & +{-}/+{-} & F & W:2b \\
B1009 & 1a<o0L 1b<o1 / 2a<1a1L 2b<1a2,1b2L & NI/NI & +{-}/+{-} & F & CP:2b \\
B1010 & 1a<o0L 1b<o1 / 2a<1a1L,1b1 2b<1a2 & NI/EN & +{-}/+{-} & F & W:2a \\
B1011 & 1a<o0L 1b<o1 / 2a<1a1L,1b1 2b<1b0L & NI/EN & +{-}/+{-} & F & W:2a \\
B1012 & 1a<o0L 1b<o1 / 2a<1a1L,1b1 2b<1b1 & NI/EE & +{-}/+{-} & E & W:1b \\
B1013 & 1a<o0L 1b<o1 / 2a<1a1L,1b1 2b<1b1L & NI/LI & +{-}/+{-} & B & CP:2b W:1b,2b \\
B1014 & 1a<o0L 1b<o1 / 2a<1a1L,1b1 2b<1b2 & NI/EE & +{-}/+{-} & E & C:2b W:1b \\
B1015 & 1a<o0L 1b<o1 / 2a<1a1L,1b1 2b<1b2L & NI/LI & +{-}/+{-} & B & P:2b W:1b,2b \\
B1016 & 1a<o0L 1b<o1 / 2a<1a1L,1b2 2b<1a2 & NI/EN & +{-}/+{-} & F & CW:2a \\
B1017 & 1a<o0L 1b<o1 / 2a<1a1L,1b2 2b<1b0L & NI/EN & +{-}/+{-} & F & CW:2a \\
B1018 & 1a<o0L 1b<o1 / 2a<1a1L,1b2 2b<1b1 & NI/EE & +{-}/+{-} & E & C:2a W:1b \\
B1019 & 1a<o0L 1b<o1 / 2a<1a1L,1b2 2b<1b1L & NI/LI & +{-}/+{-} & B & C:2a,2b P:2b W:1b,2b \\
B1020 & 1a<o0L 1b<o1 / 2a<1a1L,1b2 2b<1b2 & NI/EE & +{-}/+{-} & E & W:1b \\
B1021 & 1a<o0L 1b<o1 / 2a<1a1L,1b2 2b<1b2L & NI/LI & +{-}/+{-} & B & CP:2b W:1b,2b \\
B1022 & 1a<o0L 1b<o1 / 2a<1a2 2b<1a2,1b0L & NI/NN & +{-}/{-}{-} & N &  \\
B1023 & 1a<o0L 1b<o1 / 2a<1a2 2b<1a2,1b1 & NI/NE & +{-}/{-}{-} & F & W:2b \\
B1024 & 1a<o0L 1b<o1 / 2a<1a2 2b<1a2,1b1L & NI/NI & +{-}/{-}{-} & F & CP:2b \\
B1025 & 1a<o0L 1b<o1 / 2a<1a2 2b<1a2,1b2 & NI/NE & +{-}/{-}{-} & F & W:2b \\
B1026 & 1a<o0L 1b<o1 / 2a<1a2 2b<1a2,1b2L & NI/NI & +{-}/{-}{-} & F & CP:2b \\
B1027 & 1a<o0L 1b<o1 / 2a<1a2 2b<1a2L,1b1 & NI/NE & +{-}/{-}+ & F & W:2b \\
B1028 & 1a<o0L 1b<o1 / 2a<1a2 2b<1a2L,1b2 & NI/NE & +{-}/{-}+ & F & W:2b \\
B1029 & 1a<o0L 1b<o1 / 2a<1a2,1b0L 2b<1a2L & NI/NN & +{-}/{-}+ & N &  \\
B1030 & 1a<o0L 1b<o1 / 2a<1a2,1b0L 2b<1b1 & NI/NI & +{-}/{-}{-} & N &  \\
B1031 & 1a<o0L 1b<o1 / 2a<1a2,1b0L 2b<1b2 & NI/NI & +{-}/{-}{-} & N &  \\
B1032 & 1a<o0L 1b<o1 / 2a<1a2,1b1 2b<1a2L & NI/EN & +{-}/{-}+ & F & W:2a \\
B1033 & 1a<o0L 1b<o1 / 2a<1a2,1b1 2b<1b0L & NI/EN & +{-}/{-}{-} & F & W:2a \\
B1034 & 1a<o0L 1b<o1 / 2a<1a2,1b1 2b<1b1 & NI/EE & +{-}/{-}{-} & E & W:1b \\
B1035 & 1a<o0L 1b<o1 / 2a<1a2,1b1 2b<1b1L & NI/LI & +{-}/{-}{-} & B & CP:2b W:1b,2b \\
B1036 & 1a<o0L 1b<o1 / 2a<1a2,1b1 2b<1b2 & NI/EE & +{-}/{-}{-} & E & C:2b W:1b \\
B1037 & 1a<o0L 1b<o1 / 2a<1a2,1b1 2b<1b2L & NI/LI & +{-}/{-}{-} & B & P:2b W:1b,2b \\
B1038 & 1a<o0L 1b<o1 / 2a<1a2,1b1L 2b<1a2L & NI/IN & +{-}/{-}+ & F & CP:2a \\
B1039 & 1a<o0L 1b<o1 / 2a<1a2,1b1L 2b<1b1 & NI/IL & +{-}/{-}{-} & B & CP:2a W:1b,2a \\
B1040 & 1a<o0L 1b<o1 / 2a<1a2,1b1L 2b<1b2 & NI/IL & +{-}/{-}{-} & B & C:2a,2b P:2a W:1b,2a \\
B1041 & 1a<o0L 1b<o1 / 2a<1a2,1b2 2b<1a2L & NI/EN & +{-}/{-}+ & F & W:2a \\
B1042 & 1a<o0L 1b<o1 / 2a<1a2,1b2 2b<1b0L & NI/EN & +{-}/{-}{-} & F & W:2a \\
B1043 & 1a<o0L 1b<o1 / 2a<1a2,1b2 2b<1b1 & NI/EE & +{-}/{-}{-} & E & W:1b \\
B1044 & 1a<o0L 1b<o1 / 2a<1a2,1b2 2b<1b1L & NI/LI & +{-}/{-}{-} & B & CP:2b W:1b,2b \\
B1045 & 1a<o0L 1b<o1 / 2a<1a2,1b2 2b<1b2 & NI/EE & +{-}/{-}{-} & E & W:1b \\
B1046 & 1a<o0L 1b<o1 / 2a<1a2,1b2 2b<1b2L & NI/LI & +{-}/{-}{-} & B & CP:2b W:1b,2b \\
B1047 & 1a<o0L 1b<o1 / 2a<1a2,1b2L 2b<1a2L & NI/IN & +{-}/{-}+ & F & CP:2a \\
B1048 & 1a<o0L 1b<o1 / 2a<1a2,1b2L 2b<1b1 & NI/IL & +{-}/{-}{-} & B & CP:2a W:1b,2a \\
B1049 & 1a<o0L 1b<o1 / 2a<1a2,1b2L 2b<1b2 & NI/IL & +{-}/{-}{-} & B & CP:2a W:1b,2a \\
B1050 & 1a<o0L 1b<o1 / 2a<1a2L,1b1 2b<1b0L & NI/EN & +{-}/+{-} & F & W:2a \\
B1051 & 1a<o0L 1b<o1 / 2a<1a2L,1b1 2b<1b1 & NI/EE & +{-}/+{-} & E & W:1b \\
B1052 & 1a<o0L 1b<o1 / 2a<1a2L,1b1 2b<1b1L & NI/LI & +{-}/+{-} & B & CP:2b W:1b,2b \\
B1053 & 1a<o0L 1b<o1 / 2a<1a2L,1b1 2b<1b2 & NI/EE & +{-}/+{-} & E & C:2b W:1b \\
B1054 & 1a<o0L 1b<o1 / 2a<1a2L,1b1 2b<1b2L & NI/LI & +{-}/+{-} & B & P:2b W:1b,2b \\
B1055 & 1a<o0L 1b<o1 / 2a<1a2L,1b2 2b<1b0L & NI/EN & +{-}/+{-} & F & W:2a \\
B1056 & 1a<o0L 1b<o1 / 2a<1a2L,1b2 2b<1b1 & NI/EE & +{-}/+{-} & E & W:1b \\
B1057 & 1a<o0L 1b<o1 / 2a<1a2L,1b2 2b<1b1L & NI/LI & +{-}/+{-} & B & CP:2b W:1b,2b \\
B1058 & 1a<o0L 1b<o1 / 2a<1a2L,1b2 2b<1b2 & NI/EE & +{-}/+{-} & E & W:1b \\
B1059 & 1a<o0L 1b<o1 / 2a<1a2L,1b2 2b<1b2L & NI/LI & +{-}/+{-} & B & CP:2b W:1b,2b \\
B1060 & 1a<o0L 1b<o2 / 2a<1a0L 2b<1a1,1b0L & NI/NN & +{-}/+{-} & N &  \\
B1061 & 1a<o0L 1b<o2 / 2a<1a0L 2b<1a1,1b1 & NI/NE & +{-}/+{-} & F & W:2b \\
B1062 & 1a<o0L 1b<o2 / 2a<1a0L 2b<1a1,1b1L & NI/NI & +{-}/+{-} & F & CP:2b \\
B1063 & 1a<o0L 1b<o2 / 2a<1a0L 2b<1a1,1b2 & NI/NE & +{-}/+{-} & F & CW:2b \\
B1064 & 1a<o0L 1b<o2 / 2a<1a0L 2b<1a1,1b2L & NI/NI & +{-}/+{-} & F & P:2b \\
B1065 & 1a<o0L 1b<o2 / 2a<1a0L 2b<1a2,1b0L & NI/NN & +{-}/+{-} & N &  \\
B1066 & 1a<o0L 1b<o2 / 2a<1a0L 2b<1a2,1b1 & NI/NE & +{-}/+{-} & F & W:2b \\
B1067 & 1a<o0L 1b<o2 / 2a<1a0L 2b<1a2,1b1L & NI/NI & +{-}/+{-} & F & CP:2b \\
B1068 & 1a<o0L 1b<o2 / 2a<1a0L 2b<1a2,1b2 & NI/NE & +{-}/+{-} & F & W:2b \\
B1069 & 1a<o0L 1b<o2 / 2a<1a0L 2b<1a2,1b2L & NI/NI & +{-}/+{-} & F & CP:2b \\
B1070 & 1a<o0L 1b<o2 / 2a<1a0L,1b1 2b<1a1 & NI/IN & +{-}/+{-} & N &  \\
B1071 & 1a<o0L 1b<o2 / 2a<1a0L,1b1 2b<1a2 & NI/IN & +{-}/+{-} & N &  \\
B1072 & 1a<o0L 1b<o2 / 2a<1a0L,1b1 2b<1b0L & NI/IN & +{-}/+{-} & N &  \\
B1073 & 1a<o0L 1b<o2 / 2a<1a0L,1b1 2b<1b1 & NI/EE & +{-}/+{-} & E & W:1b \\
B1074 & 1a<o0L 1b<o2 / 2a<1a0L,1b1 2b<1b1L & NI/LI & +{-}/+{-} & B & CP:2b W:1b,2b \\
B1075 & 1a<o0L 1b<o2 / 2a<1a0L,1b1 2b<1b2 & NI/EE & +{-}/+{-} & E & C:2b W:1b \\
B1076 & 1a<o0L 1b<o2 / 2a<1a0L,1b1 2b<1b2L & NI/LI & +{-}/+{-} & B & P:2b W:1b,2b \\
B1077 & 1a<o0L 1b<o2 / 2a<1a0L,1b2 2b<1a1 & NI/IN & +{-}/+{-} & N &  \\
B1078 & 1a<o0L 1b<o2 / 2a<1a0L,1b2 2b<1a2 & NI/IN & +{-}/+{-} & N &  \\
B1079 & 1a<o0L 1b<o2 / 2a<1a0L,1b2 2b<1b0L & NI/IN & +{-}/+{-} & N &  \\
B1080 & 1a<o0L 1b<o2 / 2a<1a0L,1b2 2b<1b1 & NI/EE & +{-}/+{-} & E & C:2a W:1b \\
B1081 & 1a<o0L 1b<o2 / 2a<1a0L,1b2 2b<1b1L & NI/LI & +{-}/+{-} & B & C:2a,2b P:2b W:1b,2b \\
B1082 & 1a<o0L 1b<o2 / 2a<1a0L,1b2 2b<1b2 & NI/EE & +{-}/+{-} & E & W:1b \\
B1083 & 1a<o0L 1b<o2 / 2a<1a0L,1b2 2b<1b2L & NI/LI & +{-}/+{-} & B & CP:2b W:1b,2b \\
B1084 & 1a<o0L 1b<o2 / 2a<1a1 2b<1a1,1b0L & NI/NN & +{-}/{-}{-} & N &  \\
B1085 & 1a<o0L 1b<o2 / 2a<1a1 2b<1a1,1b1 & NI/NE & +{-}/{-}{-} & F & W:2b \\
B1086 & 1a<o0L 1b<o2 / 2a<1a1 2b<1a1,1b1L & NI/NI & +{-}/{-}{-} & F & CP:2b \\
B1087 & 1a<o0L 1b<o2 / 2a<1a1 2b<1a1,1b2 & NI/NE & +{-}/{-}{-} & F & CW:2b \\
B1088 & 1a<o0L 1b<o2 / 2a<1a1 2b<1a1,1b2L & NI/NI & +{-}/{-}{-} & F & P:2b \\
B1089 & 1a<o0L 1b<o2 / 2a<1a1 2b<1a1L,1b1 & NI/NE & +{-}/{-}+ & F & W:2b \\
B1090 & 1a<o0L 1b<o2 / 2a<1a1 2b<1a1L,1b2 & NI/NE & +{-}/{-}+ & F & CW:2b \\
B1091 & 1a<o0L 1b<o2 / 2a<1a1 2b<1a2,1b0L & NI/NN & +{-}/{-}{-} & N &  \\
B1092 & 1a<o0L 1b<o2 / 2a<1a1 2b<1a2,1b1 & NI/NE & +{-}/{-}{-} & F & W:2b \\
B1093 & 1a<o0L 1b<o2 / 2a<1a1 2b<1a2,1b1L & NI/NI & +{-}/{-}{-} & F & CP:2b \\
B1094 & 1a<o0L 1b<o2 / 2a<1a1 2b<1a2,1b2 & NI/NE & +{-}/{-}{-} & F & W:2b \\
B1095 & 1a<o0L 1b<o2 / 2a<1a1 2b<1a2,1b2L & NI/NI & +{-}/{-}{-} & F & CP:2b \\
B1096 & 1a<o0L 1b<o2 / 2a<1a1 2b<1a2L,1b1 & NI/NE & +{-}/{-}+ & F & W:2b \\
B1097 & 1a<o0L 1b<o2 / 2a<1a1 2b<1a2L,1b2 & NI/NE & +{-}/{-}+ & F & W:2b \\
B1098 & 1a<o0L 1b<o2 / 2a<1a1,1b0L 2b<1a1L & NI/NN & +{-}/{-}+ & N &  \\
B1099 & 1a<o0L 1b<o2 / 2a<1a1,1b0L 2b<1a2 & NI/NN & +{-}/{-}{-} & N &  \\
B1100 & 1a<o0L 1b<o2 / 2a<1a1,1b0L 2b<1a2L & NI/NN & +{-}/{-}+ & N &  \\
B1101 & 1a<o0L 1b<o2 / 2a<1a1,1b0L 2b<1b1 & NI/NI & +{-}/{-}{-} & N &  \\
B1102 & 1a<o0L 1b<o2 / 2a<1a1,1b0L 2b<1b2 & NI/NI & +{-}/{-}{-} & N &  \\
B1103 & 1a<o0L 1b<o2 / 2a<1a1,1b1 2b<1a1L & NI/EN & +{-}/{-}+ & F & W:2a \\
B1104 & 1a<o0L 1b<o2 / 2a<1a1,1b1 2b<1a2 & NI/EN & +{-}/{-}{-} & F & W:2a \\
B1105 & 1a<o0L 1b<o2 / 2a<1a1,1b1 2b<1a2L & NI/EN & +{-}/{-}+ & F & W:2a \\
B1106 & 1a<o0L 1b<o2 / 2a<1a1,1b1 2b<1b0L & NI/EN & +{-}/{-}{-} & F & W:2a \\
B1107 & 1a<o0L 1b<o2 / 2a<1a1,1b1 2b<1b1 & NI/EE & +{-}/{-}{-} & E & W:1b \\
B1108 & 1a<o0L 1b<o2 / 2a<1a1,1b1 2b<1b1L & NI/LI & +{-}/{-}{-} & B & CP:2b W:1b,2b \\
B1109 & 1a<o0L 1b<o2 / 2a<1a1,1b1 2b<1b2 & NI/EE & +{-}/{-}{-} & E & C:2b W:1b \\
B1110 & 1a<o0L 1b<o2 / 2a<1a1,1b1 2b<1b2L & NI/LI & +{-}/{-}{-} & B & P:2b W:1b,2b \\
B1111 & 1a<o0L 1b<o2 / 2a<1a1,1b1L 2b<1a1L & NI/IN & +{-}/{-}+ & F & CP:2a \\
B1112 & 1a<o0L 1b<o2 / 2a<1a1,1b1L 2b<1a2 & NI/IN & +{-}/{-}{-} & F & CP:2a \\
B1113 & 1a<o0L 1b<o2 / 2a<1a1,1b1L 2b<1a2L & NI/IN & +{-}/{-}+ & F & CP:2a \\
B1114 & 1a<o0L 1b<o2 / 2a<1a1,1b1L 2b<1b1 & NI/IL & +{-}/{-}{-} & B & CP:2a W:1b,2a \\
B1115 & 1a<o0L 1b<o2 / 2a<1a1,1b1L 2b<1b2 & NI/IL & +{-}/{-}{-} & B & C:2a,2b P:2a W:1b,2a \\
B1116 & 1a<o0L 1b<o2 / 2a<1a1,1b2 2b<1a1L & NI/EN & +{-}/{-}+ & F & CW:2a \\
B1117 & 1a<o0L 1b<o2 / 2a<1a1,1b2 2b<1a2 & NI/EN & +{-}/{-}{-} & F & CW:2a \\
B1118 & 1a<o0L 1b<o2 / 2a<1a1,1b2 2b<1a2L & NI/EN & +{-}/{-}+ & F & CW:2a \\
B1119 & 1a<o0L 1b<o2 / 2a<1a1,1b2 2b<1b0L & NI/EN & +{-}/{-}{-} & F & CW:2a \\
B1120 & 1a<o0L 1b<o2 / 2a<1a1,1b2 2b<1b1 & NI/EE & +{-}/{-}{-} & E & C:2a W:1b \\
B1121 & 1a<o0L 1b<o2 / 2a<1a1,1b2 2b<1b1L & NI/LI & +{-}/{-}{-} & B & C:2a,2b P:2b W:1b,2b \\
B1122 & 1a<o0L 1b<o2 / 2a<1a1,1b2 2b<1b2 & NI/EE & +{-}/{-}{-} & E & W:1b \\
B1123 & 1a<o0L 1b<o2 / 2a<1a1,1b2 2b<1b2L & NI/LI & +{-}/{-}{-} & B & CP:2b W:1b,2b \\
B1124 & 1a<o0L 1b<o2 / 2a<1a1,1b2L 2b<1a1L & NI/IN & +{-}/{-}+ & F & P:2a \\
B1125 & 1a<o0L 1b<o2 / 2a<1a1,1b2L 2b<1a2 & NI/IN & +{-}/{-}{-} & F & P:2a \\
B1126 & 1a<o0L 1b<o2 / 2a<1a1,1b2L 2b<1a2L & NI/IN & +{-}/{-}+ & F & P:2a \\
B1127 & 1a<o0L 1b<o2 / 2a<1a1,1b2L 2b<1b1 & NI/IL & +{-}/{-}{-} & B & P:2a W:1b,2a \\
B1128 & 1a<o0L 1b<o2 / 2a<1a1,1b2L 2b<1b2 & NI/IL & +{-}/{-}{-} & B & CP:2a W:1b,2a \\
B1129 & 1a<o0L 1b<o2 / 2a<1a1L 2b<1a2,1b0L & NI/NN & +{-}/+{-} & N &  \\
B1130 & 1a<o0L 1b<o2 / 2a<1a1L 2b<1a2,1b1 & NI/NE & +{-}/+{-} & F & W:2b \\
B1131 & 1a<o0L 1b<o2 / 2a<1a1L 2b<1a2,1b1L & NI/NI & +{-}/+{-} & F & CP:2b \\
B1132 & 1a<o0L 1b<o2 / 2a<1a1L 2b<1a2,1b2 & NI/NE & +{-}/+{-} & F & W:2b \\
B1133 & 1a<o0L 1b<o2 / 2a<1a1L 2b<1a2,1b2L & NI/NI & +{-}/+{-} & F & CP:2b \\
B1134 & 1a<o0L 1b<o2 / 2a<1a1L,1b1 2b<1a2 & NI/EN & +{-}/+{-} & F & W:2a \\
B1135 & 1a<o0L 1b<o2 / 2a<1a1L,1b1 2b<1b0L & NI/EN & +{-}/+{-} & F & W:2a \\
B1136 & 1a<o0L 1b<o2 / 2a<1a1L,1b1 2b<1b1 & NI/EE & +{-}/+{-} & E & W:1b \\
B1137 & 1a<o0L 1b<o2 / 2a<1a1L,1b1 2b<1b1L & NI/LI & +{-}/+{-} & B & CP:2b W:1b,2b \\
B1138 & 1a<o0L 1b<o2 / 2a<1a1L,1b1 2b<1b2 & NI/EE & +{-}/+{-} & E & C:2b W:1b \\
B1139 & 1a<o0L 1b<o2 / 2a<1a1L,1b1 2b<1b2L & NI/LI & +{-}/+{-} & B & P:2b W:1b,2b \\
B1140 & 1a<o0L 1b<o2 / 2a<1a1L,1b2 2b<1a2 & NI/EN & +{-}/+{-} & F & CW:2a \\
B1141 & 1a<o0L 1b<o2 / 2a<1a1L,1b2 2b<1b0L & NI/EN & +{-}/+{-} & F & CW:2a \\
B1142 & 1a<o0L 1b<o2 / 2a<1a1L,1b2 2b<1b1 & NI/EE & +{-}/+{-} & E & C:2a W:1b \\
B1143 & 1a<o0L 1b<o2 / 2a<1a1L,1b2 2b<1b1L & NI/LI & +{-}/+{-} & B & C:2a,2b P:2b W:1b,2b \\
B1144 & 1a<o0L 1b<o2 / 2a<1a1L,1b2 2b<1b2 & NI/EE & +{-}/+{-} & E & W:1b \\
B1145 & 1a<o0L 1b<o2 / 2a<1a1L,1b2 2b<1b2L & NI/LI & +{-}/+{-} & B & CP:2b W:1b,2b \\
B1146 & 1a<o0L 1b<o2 / 2a<1a2 2b<1a2,1b0L & NI/NN & +{-}/{-}{-} & N &  \\
B1147 & 1a<o0L 1b<o2 / 2a<1a2 2b<1a2,1b1 & NI/NE & +{-}/{-}{-} & F & W:2b \\
B1148 & 1a<o0L 1b<o2 / 2a<1a2 2b<1a2,1b1L & NI/NI & +{-}/{-}{-} & F & CP:2b \\
B1149 & 1a<o0L 1b<o2 / 2a<1a2 2b<1a2,1b2 & NI/NE & +{-}/{-}{-} & F & W:2b \\
B1150 & 1a<o0L 1b<o2 / 2a<1a2 2b<1a2,1b2L & NI/NI & +{-}/{-}{-} & F & CP:2b \\
B1151 & 1a<o0L 1b<o2 / 2a<1a2 2b<1a2L,1b1 & NI/NE & +{-}/{-}+ & F & W:2b \\
B1152 & 1a<o0L 1b<o2 / 2a<1a2 2b<1a2L,1b2 & NI/NE & +{-}/{-}+ & F & W:2b \\
B1153 & 1a<o0L 1b<o2 / 2a<1a2,1b0L 2b<1a2L & NI/NN & +{-}/{-}+ & N &  \\
B1154 & 1a<o0L 1b<o2 / 2a<1a2,1b0L 2b<1b1 & NI/NI & +{-}/{-}{-} & N &  \\
B1155 & 1a<o0L 1b<o2 / 2a<1a2,1b0L 2b<1b2 & NI/NI & +{-}/{-}{-} & N &  \\
B1156 & 1a<o0L 1b<o2 / 2a<1a2,1b1 2b<1a2L & NI/EN & +{-}/{-}+ & F & W:2a \\
B1157 & 1a<o0L 1b<o2 / 2a<1a2,1b1 2b<1b0L & NI/EN & +{-}/{-}{-} & F & W:2a \\
B1158 & 1a<o0L 1b<o2 / 2a<1a2,1b1 2b<1b1 & NI/EE & +{-}/{-}{-} & E & W:1b \\
B1159 & 1a<o0L 1b<o2 / 2a<1a2,1b1 2b<1b1L & NI/LI & +{-}/{-}{-} & B & CP:2b W:1b,2b \\
B1160 & 1a<o0L 1b<o2 / 2a<1a2,1b1 2b<1b2 & NI/EE & +{-}/{-}{-} & E & C:2b W:1b \\
B1161 & 1a<o0L 1b<o2 / 2a<1a2,1b1 2b<1b2L & NI/LI & +{-}/{-}{-} & B & P:2b W:1b,2b \\
B1162 & 1a<o0L 1b<o2 / 2a<1a2,1b1L 2b<1a2L & NI/IN & +{-}/{-}+ & F & CP:2a \\
B1163 & 1a<o0L 1b<o2 / 2a<1a2,1b1L 2b<1b1 & NI/IL & +{-}/{-}{-} & B & CP:2a W:1b,2a \\
B1164 & 1a<o0L 1b<o2 / 2a<1a2,1b1L 2b<1b2 & NI/IL & +{-}/{-}{-} & B & C:2a,2b P:2a W:1b,2a \\
B1165 & 1a<o0L 1b<o2 / 2a<1a2,1b2 2b<1a2L & NI/EN & +{-}/{-}+ & F & W:2a \\
B1166 & 1a<o0L 1b<o2 / 2a<1a2,1b2 2b<1b0L & NI/EN & +{-}/{-}{-} & F & W:2a \\
B1167 & 1a<o0L 1b<o2 / 2a<1a2,1b2 2b<1b1 & NI/EE & +{-}/{-}{-} & E & W:1b \\
B1168 & 1a<o0L 1b<o2 / 2a<1a2,1b2 2b<1b1L & NI/LI & +{-}/{-}{-} & B & CP:2b W:1b,2b \\
B1169 & 1a<o0L 1b<o2 / 2a<1a2,1b2 2b<1b2 & NI/EE & +{-}/{-}{-} & E & W:1b \\
B1170 & 1a<o0L 1b<o2 / 2a<1a2,1b2 2b<1b2L & NI/LI & +{-}/{-}{-} & B & CP:2b W:1b,2b \\
B1171 & 1a<o0L 1b<o2 / 2a<1a2,1b2L 2b<1a2L & NI/IN & +{-}/{-}+ & F & CP:2a \\
B1172 & 1a<o0L 1b<o2 / 2a<1a2,1b2L 2b<1b1 & NI/IL & +{-}/{-}{-} & B & CP:2a W:1b,2a \\
B1173 & 1a<o0L 1b<o2 / 2a<1a2,1b2L 2b<1b2 & NI/IL & +{-}/{-}{-} & B & CP:2a W:1b,2a \\
B1174 & 1a<o0L 1b<o2 / 2a<1a2L,1b1 2b<1b0L & NI/EN & +{-}/+{-} & F & W:2a \\
B1175 & 1a<o0L 1b<o2 / 2a<1a2L,1b1 2b<1b1 & NI/EE & +{-}/+{-} & E & W:1b \\
B1176 & 1a<o0L 1b<o2 / 2a<1a2L,1b1 2b<1b1L & NI/LI & +{-}/+{-} & B & CP:2b W:1b,2b \\
B1177 & 1a<o0L 1b<o2 / 2a<1a2L,1b1 2b<1b2 & NI/EE & +{-}/+{-} & E & C:2b W:1b \\
B1178 & 1a<o0L 1b<o2 / 2a<1a2L,1b1 2b<1b2L & NI/LI & +{-}/+{-} & B & P:2b W:1b,2b \\
B1179 & 1a<o0L 1b<o2 / 2a<1a2L,1b2 2b<1b0L & NI/EN & +{-}/+{-} & F & W:2a \\
B1180 & 1a<o0L 1b<o2 / 2a<1a2L,1b2 2b<1b1 & NI/EE & +{-}/+{-} & E & W:1b \\
B1181 & 1a<o0L 1b<o2 / 2a<1a2L,1b2 2b<1b1L & NI/LI & +{-}/+{-} & B & CP:2b W:1b,2b \\
B1182 & 1a<o0L 1b<o2 / 2a<1a2L,1b2 2b<1b2 & NI/EE & +{-}/+{-} & E & W:1b \\
B1183 & 1a<o0L 1b<o2 / 2a<1a2L,1b2 2b<1b2L & NI/LI & +{-}/+{-} & B & CP:2b W:1b,2b \\
B1184 & 1a<o1 1b<o1 / 2a<1a0L 2b<1a1,1b0L & EE/ND & {-}{-}/{-}{-} & E &  \\
B1185 & 1a<o1 1b<o1 / 2a<1a0L 2b<1a1,1b1 & EE/ND & {-}{-}/{-}{-} & E &  \\
B1186 & 1a<o1 1b<o1 / 2a<1a0L 2b<1a1,1b1L & EE/ND & {-}{-}/{-}{-} & E &  \\
B1187 & 1a<o1 1b<o1 / 2a<1a0L 2b<1a1,1b2 & EE/ND & {-}{-}/{-}{-} & E &  \\
B1188 & 1a<o1 1b<o1 / 2a<1a0L 2b<1a1,1b2L & EE/ND & {-}{-}/{-}{-} & E &  \\
B1189 & 1a<o1 1b<o1 / 2a<1a0L 2b<1a2,1b0L & EE/ND & {-}{-}/{-}{-} & E &  \\
B1190 & 1a<o1 1b<o1 / 2a<1a0L 2b<1a2,1b1 & EE/ND & {-}{-}/{-}{-} & E &  \\
B1191 & 1a<o1 1b<o1 / 2a<1a0L 2b<1a2,1b1L & EE/ND & {-}{-}/{-}{-} & E &  \\
B1192 & 1a<o1 1b<o1 / 2a<1a0L 2b<1a2,1b2 & EE/ND & {-}{-}/{-}{-} & E &  \\
B1193 & 1a<o1 1b<o1 / 2a<1a0L 2b<1a2,1b2L & EE/ND & {-}{-}/{-}{-} & E &  \\
B1194 & 1a<o1 1b<o1 / 2a<1a0L,1b1 2b<1a1 & EE/DD & {-}{-}/{-}{-} & E &  \\
B1195 & 1a<o1 1b<o1 / 2a<1a0L,1b1 2b<1a2 & EE/DD & {-}{-}/{-}{-} & E &  \\
B1196 & 1a<o1 1b<o1 / 2a<1a0L,1b1 2b<1b1 & EE/DD & {-}{-}/{-}{-} & E &  \\
B1197 & 1a<o1 1b<o1 / 2a<1a0L,1b1 2b<1b1L & EE/DD & {-}{-}/{-}{-} & E &  \\
B1198 & 1a<o1 1b<o1 / 2a<1a0L,1b1 2b<1b2 & EE/DD & {-}{-}/{-}{-} & E &  \\
B1199 & 1a<o1 1b<o1 / 2a<1a0L,1b1 2b<1b2L & EE/DD & {-}{-}/{-}{-} & E &  \\
B1200 & 1a<o1 1b<o1 / 2a<1a0L,1b2 2b<1a1 & EE/DD & {-}{-}/{-}{-} & E &  \\
B1201 & 1a<o1 1b<o1 / 2a<1a0L,1b2 2b<1a2 & EE/DD & {-}{-}/{-}{-} & E &  \\
B1202 & 1a<o1 1b<o1 / 2a<1a0L,1b2 2b<1b1 & EE/DD & {-}{-}/{-}{-} & E &  \\
B1203 & 1a<o1 1b<o1 / 2a<1a0L,1b2 2b<1b1L & EE/DD & {-}{-}/{-}{-} & E &  \\
B1204 & 1a<o1 1b<o1 / 2a<1a0L,1b2 2b<1b2 & EE/DD & {-}{-}/{-}{-} & E &  \\
B1205 & 1a<o1 1b<o1 / 2a<1a0L,1b2 2b<1b2L & EE/DD & {-}{-}/{-}{-} & E &  \\
B1206 & 1a<o1 1b<o1 / 2a<1a1 2b<1a1,1b1 & EE/DD & {-}{-}/{-}{-} & E &  \\
B1207 & 1a<o1 1b<o1 / 2a<1a1 2b<1a1,1b1L & EE/DD & {-}{-}/{-}{-} & E &  \\
B1208 & 1a<o1 1b<o1 / 2a<1a1 2b<1a1,1b2 & EE/DD & {-}{-}/{-}{-} & E &  \\
B1209 & 1a<o1 1b<o1 / 2a<1a1 2b<1a1,1b2L & EE/DD & {-}{-}/{-}{-} & E &  \\
B1210 & 1a<o1 1b<o1 / 2a<1a1 2b<1a1L,1b1 & EE/DD & {-}{-}/{-}{-} & E &  \\
B1211 & 1a<o1 1b<o1 / 2a<1a1 2b<1a1L,1b2 & EE/DD & {-}{-}/{-}{-} & E &  \\
B1212 & 1a<o1 1b<o1 / 2a<1a1 2b<1a2,1b1 & EE/DD & {-}{-}/{-}{-} & E &  \\
B1213 & 1a<o1 1b<o1 / 2a<1a1 2b<1a2,1b1L & EE/DD & {-}{-}/{-}{-} & E &  \\
B1214 & 1a<o1 1b<o1 / 2a<1a1 2b<1a2,1b2 & EE/DD & {-}{-}/{-}{-} & E &  \\
B1215 & 1a<o1 1b<o1 / 2a<1a1 2b<1a2,1b2L & EE/DD & {-}{-}/{-}{-} & E &  \\
B1216 & 1a<o1 1b<o1 / 2a<1a1 2b<1a2L,1b1 & EE/DD & {-}{-}/{-}{-} & E &  \\
B1217 & 1a<o1 1b<o1 / 2a<1a1 2b<1a2L,1b2 & EE/DD & {-}{-}/{-}{-} & E &  \\
B1218 & 1a<o1 1b<o1 / 2a<1a1,1b1 2b<1a1L & EE/DD & {-}{-}/{-}{-} & E &  \\
B1219 & 1a<o1 1b<o1 / 2a<1a1,1b1 2b<1a2 & EE/DD & {-}{-}/{-}{-} & E &  \\
B1220 & 1a<o1 1b<o1 / 2a<1a1,1b1 2b<1a2L & EE/DD & {-}{-}/{-}{-} & E &  \\
B1221 & 1a<o1 1b<o1 / 2a<1a1,1b1L 2b<1a1L & EE/DD & {-}{-}/{-}{-} & E &  \\
B1222 & 1a<o1 1b<o1 / 2a<1a1,1b1L 2b<1a2 & EE/DD & {-}{-}/{-}{-} & E &  \\
B1223 & 1a<o1 1b<o1 / 2a<1a1,1b1L 2b<1a2L & EE/DD & {-}{-}/{-}{-} & E &  \\
B1224 & 1a<o1 1b<o1 / 2a<1a1,1b1L 2b<1b2 & EE/DD & {-}{-}/{-}{-} & E &  \\
B1225 & 1a<o1 1b<o1 / 2a<1a1,1b2 2b<1a1L & EE/DD & {-}{-}/{-}{-} & E &  \\
B1226 & 1a<o1 1b<o1 / 2a<1a1,1b2 2b<1a2 & EE/DD & {-}{-}/{-}{-} & E &  \\
B1227 & 1a<o1 1b<o1 / 2a<1a1,1b2 2b<1a2L & EE/DD & {-}{-}/{-}{-} & E &  \\
B1228 & 1a<o1 1b<o1 / 2a<1a1,1b2 2b<1b1L & EE/DD & {-}{-}/{-}{-} & E &  \\
B1229 & 1a<o1 1b<o1 / 2a<1a1,1b2 2b<1b2 & EE/DD & {-}{-}/{-}{-} & E &  \\
B1230 & 1a<o1 1b<o1 / 2a<1a1,1b2 2b<1b2L & EE/DD & {-}{-}/{-}{-} & E &  \\
B1231 & 1a<o1 1b<o1 / 2a<1a1,1b2L 2b<1a1L & EE/DD & {-}{-}/{-}{-} & E &  \\
B1232 & 1a<o1 1b<o1 / 2a<1a1,1b2L 2b<1a2 & EE/DD & {-}{-}/{-}{-} & E &  \\
B1233 & 1a<o1 1b<o1 / 2a<1a1,1b2L 2b<1a2L & EE/DD & {-}{-}/{-}{-} & E &  \\
B1234 & 1a<o1 1b<o1 / 2a<1a1,1b2L 2b<1b2 & EE/DD & {-}{-}/{-}{-} & E &  \\
B1235 & 1a<o1 1b<o1 / 2a<1a1L 2b<1a2,1b1L & EE/DD & {-}{-}/{-}{-} & E &  \\
B1236 & 1a<o1 1b<o1 / 2a<1a1L 2b<1a2,1b2 & EE/DD & {-}{-}/{-}{-} & E &  \\
B1237 & 1a<o1 1b<o1 / 2a<1a1L 2b<1a2,1b2L & EE/DD & {-}{-}/{-}{-} & E &  \\
B1238 & 1a<o1 1b<o1 / 2a<1a1L,1b2 2b<1a2 & EE/DD & {-}{-}/{-}{-} & E &  \\
B1239 & 1a<o1 1b<o1 / 2a<1a1L,1b2 2b<1b2 & EE/DD & {-}{-}/{-}{-} & E &  \\
B1240 & 1a<o1 1b<o1 / 2a<1a1L,1b2 2b<1b2L & EE/DD & {-}{-}/{-}{-} & E &  \\
B1241 & 1a<o1 1b<o1 / 2a<1a2 2b<1a2,1b2 & EE/DD & {-}{-}/{-}{-} & E &  \\
B1242 & 1a<o1 1b<o1 / 2a<1a2 2b<1a2,1b2L & EE/DD & {-}{-}/{-}{-} & E &  \\
B1243 & 1a<o1 1b<o1 / 2a<1a2 2b<1a2L,1b2 & EE/DD & {-}{-}/{-}{-} & E &  \\
B1244 & 1a<o1 1b<o1 / 2a<1a2,1b2 2b<1a2L & EE/DD & {-}{-}/{-}{-} & E &  \\
B1245 & 1a<o1 1b<o1 / 2a<1a2,1b2L 2b<1a2L & EE/DD & {-}{-}/{-}{-} & E &  \\
B1246 & 1a<o1 1b<o1L / 2a<1a0L 2b<1a1,1b0L & LI/ND & {-}+/{-}+ & B & CPW:1b \\
B1247 & 1a<o1 1b<o1L / 2a<1a0L 2b<1a1,1b1 & LI/NE & {-}+/{-}{-} & BF & CPW:1b \\
B1248 & 1a<o1 1b<o1L / 2a<1a0L 2b<1a1,1b1L & LI/NI & {-}+/{-}+ & BF & CPW:1b,2b \\
B1249 & 1a<o1 1b<o1L / 2a<1a0L 2b<1a1,1b2 & LI/NE & {-}+/{-}{-} & BF & CPW:1b \\
B1250 & 1a<o1 1b<o1L / 2a<1a0L 2b<1a1,1b2L & LI/NI & {-}+/{-}+ & BF & CPW:1b,2b \\
B1251 & 1a<o1 1b<o1L / 2a<1a0L 2b<1a2,1b0L & LI/ND & {-}+/{-}+ & B & CPW:1b \\
B1252 & 1a<o1 1b<o1L / 2a<1a0L 2b<1a2,1b1 & LI/NE & {-}+/{-}{-} & BF & CPW:1b \\
B1253 & 1a<o1 1b<o1L / 2a<1a0L 2b<1a2,1b1L & LI/NI & {-}+/{-}+ & BF & CPW:1b,2b \\
B1254 & 1a<o1 1b<o1L / 2a<1a0L 2b<1a2,1b2 & LI/NE & {-}+/{-}{-} & BF & CPW:1b \\
B1255 & 1a<o1 1b<o1L / 2a<1a0L 2b<1a2,1b2L & LI/NI & {-}+/{-}+ & BF & CPW:1b,2b \\
B1256 & 1a<o1 1b<o1L / 2a<1a0L,1b1 2b<1a1 & LI/ID & {-}+/{-}{-} & B & CPW:1b,2a \\
B1257 & 1a<o1 1b<o1L / 2a<1a0L,1b1 2b<1a2 & LI/ID & {-}+/{-}{-} & B & CPW:1b,2a \\
B1258 & 1a<o1 1b<o1L / 2a<1a0L,1b1 2b<1b0L & LI/IN & {-}+/{-}+ & B & CPW:1b,2a \\
B1259 & 1a<o1 1b<o1L / 2a<1a0L,1b1 2b<1b1 & LI/EE & {-}+/{-}{-} & BE & CPW:1b \\
B1260 & 1a<o1 1b<o1L / 2a<1a0L,1b1 2b<1b1L & LI/LI & {-}+/{-}+ & B & CPW:1b,2b \\
B1261 & 1a<o1 1b<o1L / 2a<1a0L,1b1 2b<1b2 & LI/EE & {-}+/{-}{-} & BE & CPW:1b \\
B1262 & 1a<o1 1b<o1L / 2a<1a0L,1b1 2b<1b2L & LI/LI & {-}+/{-}+ & B & CPW:1b,2b \\
B1263 & 1a<o1 1b<o1L / 2a<1a0L,1b2 2b<1a1 & LI/ID & {-}+/{-}{-} & B & CPW:1b,2a \\
B1264 & 1a<o1 1b<o1L / 2a<1a0L,1b2 2b<1a2 & LI/ID & {-}+/{-}{-} & B & CPW:1b,2a \\
B1265 & 1a<o1 1b<o1L / 2a<1a0L,1b2 2b<1b0L & LI/IN & {-}+/{-}+ & B & CPW:1b,2a \\
B1266 & 1a<o1 1b<o1L / 2a<1a0L,1b2 2b<1b1 & LI/EE & {-}+/{-}{-} & BE & CPW:1b \\
B1267 & 1a<o1 1b<o1L / 2a<1a0L,1b2 2b<1b1L & LI/LI & {-}+/{-}+ & B & CPW:1b,2b \\
B1268 & 1a<o1 1b<o1L / 2a<1a0L,1b2 2b<1b2 & LI/EE & {-}+/{-}{-} & BE & CPW:1b \\
B1269 & 1a<o1 1b<o1L / 2a<1a0L,1b2 2b<1b2L & LI/LI & {-}+/{-}+ & B & CPW:1b,2b \\
B1270 & 1a<o1 1b<o1L / 2a<1a1 2b<1a1,1b0L & LI/DD & {-}+/{-}+ & B & CPW:1b \\
B1271 & 1a<o1 1b<o1L / 2a<1a1 2b<1a1,1b1 & LI/DE & {-}+/{-}{-} & BF & CPW:1b \\
B1272 & 1a<o1 1b<o1L / 2a<1a1 2b<1a1,1b1L & LI/DI & {-}+/{-}+ & BF & CPW:1b,2b \\
B1273 & 1a<o1 1b<o1L / 2a<1a1 2b<1a1,1b2 & LI/DE & {-}+/{-}{-} & BF & CPW:1b \\
B1274 & 1a<o1 1b<o1L / 2a<1a1 2b<1a1,1b2L & LI/DI & {-}+/{-}+ & BF & CPW:1b,2b \\
B1275 & 1a<o1 1b<o1L / 2a<1a1 2b<1a1L,1b1 & LI/DE & {-}+/{-}{-} & BF & CPW:1b \\
B1276 & 1a<o1 1b<o1L / 2a<1a1 2b<1a1L,1b2 & LI/DE & {-}+/{-}{-} & BF & CPW:1b \\
B1277 & 1a<o1 1b<o1L / 2a<1a1 2b<1a2,1b0L & LI/DD & {-}+/{-}+ & B & CPW:1b \\
B1278 & 1a<o1 1b<o1L / 2a<1a1 2b<1a2,1b1 & LI/DE & {-}+/{-}{-} & BF & CPW:1b \\
B1279 & 1a<o1 1b<o1L / 2a<1a1 2b<1a2,1b1L & LI/DI & {-}+/{-}+ & BF & CPW:1b,2b \\
B1280 & 1a<o1 1b<o1L / 2a<1a1 2b<1a2,1b2 & LI/DE & {-}+/{-}{-} & BF & CPW:1b \\
B1281 & 1a<o1 1b<o1L / 2a<1a1 2b<1a2,1b2L & LI/DI & {-}+/{-}+ & BF & CPW:1b,2b \\
B1282 & 1a<o1 1b<o1L / 2a<1a1 2b<1a2L,1b1 & LI/DE & {-}+/{-}{-} & BF & CPW:1b \\
B1283 & 1a<o1 1b<o1L / 2a<1a1 2b<1a2L,1b2 & LI/DE & {-}+/{-}{-} & BF & CPW:1b \\
B1284 & 1a<o1 1b<o1L / 2a<1a1,1b0L 2b<1a1L & LI/DD & {-}+/+{-} & B & CPW:1b \\
B1285 & 1a<o1 1b<o1L / 2a<1a1,1b0L 2b<1a2 & LI/DD & {-}+/+{-} & B & CPW:1b \\
B1286 & 1a<o1 1b<o1L / 2a<1a1,1b0L 2b<1a2L & LI/DD & {-}+/+{-} & B & CPW:1b \\
B1287 & 1a<o1 1b<o1L / 2a<1a1,1b0L 2b<1b1 & LI/DI & {-}+/+{-} & B & CPW:1b,2b \\
B1288 & 1a<o1 1b<o1L / 2a<1a1,1b0L 2b<1b2 & LI/DI & {-}+/+{-} & B & CPW:1b,2b \\
B1289 & 1a<o1 1b<o1L / 2a<1a1,1b1 2b<1a1L & LI/ED & {-}+/{-}{-} & BF & CPW:1b \\
B1290 & 1a<o1 1b<o1L / 2a<1a1,1b1 2b<1a2 & LI/ED & {-}+/{-}{-} & BF & CPW:1b \\
B1291 & 1a<o1 1b<o1L / 2a<1a1,1b1 2b<1a2L & LI/ED & {-}+/{-}{-} & BF & CPW:1b \\
B1292 & 1a<o1 1b<o1L / 2a<1a1,1b1 2b<1b0L & LI/EN & {-}+/{-}+ & BF & CPW:1b \\
B1293 & 1a<o1 1b<o1L / 2a<1a1,1b1 2b<1b1 & LI/EE & {-}+/{-}{-} & BE & CPW:1b \\
B1294 & 1a<o1 1b<o1L / 2a<1a1,1b1 2b<1b1L & LI/LI & {-}+/{-}+ & B & CPW:1b,2b \\
B1295 & 1a<o1 1b<o1L / 2a<1a1,1b1 2b<1b2 & LI/EE & {-}+/{-}{-} & BE & CPW:1b \\
B1296 & 1a<o1 1b<o1L / 2a<1a1,1b1 2b<1b2L & LI/LI & {-}+/{-}+ & B & CPW:1b,2b \\
B1297 & 1a<o1 1b<o1L / 2a<1a1,1b1L 2b<1a1L & LI/ID & {-}+/+{-} & BF & CPW:1b,2a \\
B1298 & 1a<o1 1b<o1L / 2a<1a1,1b1L 2b<1a2 & LI/ID & {-}+/+{-} & BF & CPW:1b,2a \\
B1299 & 1a<o1 1b<o1L / 2a<1a1,1b1L 2b<1a2L & LI/ID & {-}+/+{-} & BF & CPW:1b,2a \\
B1300 & 1a<o1 1b<o1L / 2a<1a1,1b1L 2b<1b1 & LI/IL & {-}+/+{-} & B & CPW:1b,2a \\
B1301 & 1a<o1 1b<o1L / 2a<1a1,1b1L 2b<1b2 & LI/IL & {-}+/+{-} & B & CPW:1b,2a \\
B1302 & 1a<o1 1b<o1L / 2a<1a1,1b2 2b<1a1L & LI/ED & {-}+/{-}{-} & BF & CPW:1b \\
B1303 & 1a<o1 1b<o1L / 2a<1a1,1b2 2b<1a2 & LI/ED & {-}+/{-}{-} & BF & CPW:1b \\
B1304 & 1a<o1 1b<o1L / 2a<1a1,1b2 2b<1a2L & LI/ED & {-}+/{-}{-} & BF & CPW:1b \\
B1305 & 1a<o1 1b<o1L / 2a<1a1,1b2 2b<1b0L & LI/EN & {-}+/{-}+ & BF & CPW:1b \\
B1306 & 1a<o1 1b<o1L / 2a<1a1,1b2 2b<1b1 & LI/EE & {-}+/{-}{-} & BE & CPW:1b \\
B1307 & 1a<o1 1b<o1L / 2a<1a1,1b2 2b<1b1L & LI/LI & {-}+/{-}+ & B & CPW:1b,2b \\
B1308 & 1a<o1 1b<o1L / 2a<1a1,1b2 2b<1b2 & LI/EE & {-}+/{-}{-} & BE & CPW:1b \\
B1309 & 1a<o1 1b<o1L / 2a<1a1,1b2 2b<1b2L & LI/LI & {-}+/{-}+ & B & CPW:1b,2b \\
B1310 & 1a<o1 1b<o1L / 2a<1a1,1b2L 2b<1a1L & LI/ID & {-}+/+{-} & BF & CPW:1b,2a \\
B1311 & 1a<o1 1b<o1L / 2a<1a1,1b2L 2b<1a2 & LI/ID & {-}+/+{-} & BF & CPW:1b,2a \\
B1312 & 1a<o1 1b<o1L / 2a<1a1,1b2L 2b<1a2L & LI/ID & {-}+/+{-} & BF & CPW:1b,2a \\
B1313 & 1a<o1 1b<o1L / 2a<1a1,1b2L 2b<1b1 & LI/IL & {-}+/+{-} & B & CPW:1b,2a \\
B1314 & 1a<o1 1b<o1L / 2a<1a1,1b2L 2b<1b2 & LI/IL & {-}+/+{-} & B & CPW:1b,2a \\
B1315 & 1a<o1 1b<o1L / 2a<1a1L 2b<1a2,1b0L & LI/DD & {-}+/{-}+ & B & CPW:1b \\
B1316 & 1a<o1 1b<o1L / 2a<1a1L 2b<1a2,1b1 & LI/DE & {-}+/{-}{-} & BF & CPW:1b \\
B1317 & 1a<o1 1b<o1L / 2a<1a1L 2b<1a2,1b1L & LI/DI & {-}+/{-}+ & BF & CPW:1b,2b \\
B1318 & 1a<o1 1b<o1L / 2a<1a1L 2b<1a2,1b2 & LI/DE & {-}+/{-}{-} & BF & CPW:1b \\
B1319 & 1a<o1 1b<o1L / 2a<1a1L 2b<1a2,1b2L & LI/DI & {-}+/{-}+ & BF & CPW:1b,2b \\
B1320 & 1a<o1 1b<o1L / 2a<1a1L,1b1 2b<1a2 & LI/ED & {-}+/{-}{-} & BF & CPW:1b \\
B1321 & 1a<o1 1b<o1L / 2a<1a1L,1b1 2b<1b0L & LI/EN & {-}+/{-}+ & BF & CPW:1b \\
B1322 & 1a<o1 1b<o1L / 2a<1a1L,1b1 2b<1b1 & LI/EE & {-}+/{-}{-} & BE & CPW:1b \\
B1323 & 1a<o1 1b<o1L / 2a<1a1L,1b1 2b<1b1L & LI/LI & {-}+/{-}+ & B & CPW:1b,2b \\
B1324 & 1a<o1 1b<o1L / 2a<1a1L,1b1 2b<1b2 & LI/EE & {-}+/{-}{-} & BE & CPW:1b \\
B1325 & 1a<o1 1b<o1L / 2a<1a1L,1b1 2b<1b2L & LI/LI & {-}+/{-}+ & B & CPW:1b,2b \\
B1326 & 1a<o1 1b<o1L / 2a<1a1L,1b2 2b<1a2 & LI/ED & {-}+/{-}{-} & BF & CPW:1b \\
B1327 & 1a<o1 1b<o1L / 2a<1a1L,1b2 2b<1b0L & LI/EN & {-}+/{-}+ & BF & CPW:1b \\
B1328 & 1a<o1 1b<o1L / 2a<1a1L,1b2 2b<1b1 & LI/EE & {-}+/{-}{-} & BE & CPW:1b \\
B1329 & 1a<o1 1b<o1L / 2a<1a1L,1b2 2b<1b1L & LI/LI & {-}+/{-}+ & B & CPW:1b,2b \\
B1330 & 1a<o1 1b<o1L / 2a<1a1L,1b2 2b<1b2 & LI/EE & {-}+/{-}{-} & BE & CPW:1b \\
B1331 & 1a<o1 1b<o1L / 2a<1a1L,1b2 2b<1b2L & LI/LI & {-}+/{-}+ & B & CPW:1b,2b \\
B1332 & 1a<o1 1b<o1L / 2a<1a2 2b<1a2,1b0L & LI/DD & {-}+/{-}+ & B & CPW:1b \\
B1333 & 1a<o1 1b<o1L / 2a<1a2 2b<1a2,1b1 & LI/DE & {-}+/{-}{-} & BF & CPW:1b \\
B1334 & 1a<o1 1b<o1L / 2a<1a2 2b<1a2,1b1L & LI/DI & {-}+/{-}+ & BF & CPW:1b,2b \\
B1335 & 1a<o1 1b<o1L / 2a<1a2 2b<1a2,1b2 & LI/DE & {-}+/{-}{-} & BF & CPW:1b \\
B1336 & 1a<o1 1b<o1L / 2a<1a2 2b<1a2,1b2L & LI/DI & {-}+/{-}+ & BF & CPW:1b,2b \\
B1337 & 1a<o1 1b<o1L / 2a<1a2 2b<1a2L,1b1 & LI/DE & {-}+/{-}{-} & BF & CPW:1b \\
B1338 & 1a<o1 1b<o1L / 2a<1a2 2b<1a2L,1b2 & LI/DE & {-}+/{-}{-} & BF & CPW:1b \\
B1339 & 1a<o1 1b<o1L / 2a<1a2,1b0L 2b<1a2L & LI/DD & {-}+/+{-} & B & CPW:1b \\
B1340 & 1a<o1 1b<o1L / 2a<1a2,1b0L 2b<1b1 & LI/DI & {-}+/+{-} & B & CPW:1b,2b \\
B1341 & 1a<o1 1b<o1L / 2a<1a2,1b0L 2b<1b2 & LI/DI & {-}+/+{-} & B & CPW:1b,2b \\
B1342 & 1a<o1 1b<o1L / 2a<1a2,1b1 2b<1a2L & LI/ED & {-}+/{-}{-} & BF & CPW:1b \\
B1343 & 1a<o1 1b<o1L / 2a<1a2,1b1 2b<1b0L & LI/EN & {-}+/{-}+ & BF & CPW:1b \\
B1344 & 1a<o1 1b<o1L / 2a<1a2,1b1 2b<1b1 & LI/EE & {-}+/{-}{-} & BE & CPW:1b \\
B1345 & 1a<o1 1b<o1L / 2a<1a2,1b1 2b<1b1L & LI/LI & {-}+/{-}+ & B & CPW:1b,2b \\
B1346 & 1a<o1 1b<o1L / 2a<1a2,1b1 2b<1b2 & LI/EE & {-}+/{-}{-} & BE & CPW:1b \\
B1347 & 1a<o1 1b<o1L / 2a<1a2,1b1 2b<1b2L & LI/LI & {-}+/{-}+ & B & CPW:1b,2b \\
B1348 & 1a<o1 1b<o1L / 2a<1a2,1b1L 2b<1a2L & LI/ID & {-}+/+{-} & BF & CPW:1b,2a \\
B1349 & 1a<o1 1b<o1L / 2a<1a2,1b1L 2b<1b1 & LI/IL & {-}+/+{-} & B & CPW:1b,2a \\
B1350 & 1a<o1 1b<o1L / 2a<1a2,1b1L 2b<1b2 & LI/IL & {-}+/+{-} & B & CPW:1b,2a \\
B1351 & 1a<o1 1b<o1L / 2a<1a2,1b2 2b<1a2L & LI/ED & {-}+/{-}{-} & BF & CPW:1b \\
B1352 & 1a<o1 1b<o1L / 2a<1a2,1b2 2b<1b0L & LI/EN & {-}+/{-}+ & BF & CPW:1b \\
B1353 & 1a<o1 1b<o1L / 2a<1a2,1b2 2b<1b1 & LI/EE & {-}+/{-}{-} & BE & CPW:1b \\
B1354 & 1a<o1 1b<o1L / 2a<1a2,1b2 2b<1b1L & LI/LI & {-}+/{-}+ & B & CPW:1b,2b \\
B1355 & 1a<o1 1b<o1L / 2a<1a2,1b2 2b<1b2 & LI/EE & {-}+/{-}{-} & BE & CPW:1b \\
B1356 & 1a<o1 1b<o1L / 2a<1a2,1b2 2b<1b2L & LI/LI & {-}+/{-}+ & B & CPW:1b,2b \\
B1357 & 1a<o1 1b<o1L / 2a<1a2,1b2L 2b<1a2L & LI/ID & {-}+/+{-} & BF & CPW:1b,2a \\
B1358 & 1a<o1 1b<o1L / 2a<1a2,1b2L 2b<1b1 & LI/IL & {-}+/+{-} & B & CPW:1b,2a \\
B1359 & 1a<o1 1b<o1L / 2a<1a2,1b2L 2b<1b2 & LI/IL & {-}+/+{-} & B & CPW:1b,2a \\
B1360 & 1a<o1 1b<o1L / 2a<1a2L,1b1 2b<1b0L & LI/EN & {-}+/{-}+ & BF & CPW:1b \\
B1361 & 1a<o1 1b<o1L / 2a<1a2L,1b1 2b<1b1 & LI/EE & {-}+/{-}{-} & BE & CPW:1b \\
B1362 & 1a<o1 1b<o1L / 2a<1a2L,1b1 2b<1b1L & LI/LI & {-}+/{-}+ & B & CPW:1b,2b \\
B1363 & 1a<o1 1b<o1L / 2a<1a2L,1b1 2b<1b2 & LI/EE & {-}+/{-}{-} & BE & CPW:1b \\
B1364 & 1a<o1 1b<o1L / 2a<1a2L,1b1 2b<1b2L & LI/LI & {-}+/{-}+ & B & CPW:1b,2b \\
B1365 & 1a<o1 1b<o1L / 2a<1a2L,1b2 2b<1b0L & LI/EN & {-}+/{-}+ & BF & CPW:1b \\
B1366 & 1a<o1 1b<o1L / 2a<1a2L,1b2 2b<1b1 & LI/EE & {-}+/{-}{-} & BE & CPW:1b \\
B1367 & 1a<o1 1b<o1L / 2a<1a2L,1b2 2b<1b1L & LI/LI & {-}+/{-}+ & B & CPW:1b,2b \\
B1368 & 1a<o1 1b<o1L / 2a<1a2L,1b2 2b<1b2 & LI/EE & {-}+/{-}{-} & BE & CPW:1b \\
B1369 & 1a<o1 1b<o1L / 2a<1a2L,1b2 2b<1b2L & LI/LI & {-}+/{-}+ & B & CPW:1b,2b \\
B1370 & 1a<o1 1b<o2 / 2a<1a0L 2b<1a1,1b0L & EE/ND & {-}{-}/{-}{-} & E & C:1b \\
B1371 & 1a<o1 1b<o2 / 2a<1a0L 2b<1a1,1b1 & EE/ND & {-}{-}/{-}{-} & E & C:1b \\
B1372 & 1a<o1 1b<o2 / 2a<1a0L 2b<1a1,1b1L & EE/ND & {-}{-}/{-}{-} & E & C:1b \\
B1373 & 1a<o1 1b<o2 / 2a<1a0L 2b<1a1,1b2 & EE/ND & {-}{-}/{-}{-} & E & C:1b,2b \\
B1374 & 1a<o1 1b<o2 / 2a<1a0L 2b<1a1,1b2L & EE/ND & {-}{-}/{-}{-} & E & C:1b,2b \\
B1375 & 1a<o1 1b<o2 / 2a<1a0L 2b<1a2,1b0L & EE/ND & {-}{-}/{-}{-} & E & C:1b \\
B1376 & 1a<o1 1b<o2 / 2a<1a0L 2b<1a2,1b1 & EE/ND & {-}{-}/{-}{-} & E & C:1b \\
B1377 & 1a<o1 1b<o2 / 2a<1a0L 2b<1a2,1b1L & EE/ND & {-}{-}/{-}{-} & E & C:1b \\
B1378 & 1a<o1 1b<o2 / 2a<1a0L 2b<1a2,1b2 & EE/ND & {-}{-}/{-}{-} & E & C:1b \\
B1379 & 1a<o1 1b<o2 / 2a<1a0L 2b<1a2,1b2L & EE/ND & {-}{-}/{-}{-} & E & C:1b \\
B1380 & 1a<o1 1b<o2 / 2a<1a0L,1b1 2b<1a1 & EE/DD & {-}{-}/{-}{-} & E & C:1b,2a \\
B1381 & 1a<o1 1b<o2 / 2a<1a0L,1b1 2b<1a2 & EE/DD & {-}{-}/{-}{-} & E & C:1b,2a \\
B1382 & 1a<o1 1b<o2 / 2a<1a0L,1b1 2b<1b0L & EE/DN & {-}{-}/{-}{-} & E & C:1b,2a \\
B1383 & 1a<o1 1b<o2 / 2a<1a0L,1b1 2b<1b1 & EE/DD & {-}{-}/{-}{-} & E & C:1b \\
B1384 & 1a<o1 1b<o2 / 2a<1a0L,1b1 2b<1b1L & EE/DD & {-}{-}/{-}{-} & E & C:1b \\
B1385 & 1a<o1 1b<o2 / 2a<1a0L,1b1 2b<1b2 & EE/DD & {-}{-}/{-}{-} & E & C:1b,2b \\
B1386 & 1a<o1 1b<o2 / 2a<1a0L,1b1 2b<1b2L & EE/DD & {-}{-}/{-}{-} & E & C:1b,2b \\
B1387 & 1a<o1 1b<o2 / 2a<1a0L,1b2 2b<1a1 & EE/DD & {-}{-}/{-}{-} & E & C:1b,2a \\
B1388 & 1a<o1 1b<o2 / 2a<1a0L,1b2 2b<1a2 & EE/DD & {-}{-}/{-}{-} & E & C:1b,2a \\
B1389 & 1a<o1 1b<o2 / 2a<1a0L,1b2 2b<1b0L & EE/DN & {-}{-}/{-}{-} & E & C:1b,2a \\
B1390 & 1a<o1 1b<o2 / 2a<1a0L,1b2 2b<1b1 & EE/DD & {-}{-}/{-}{-} & E & C:1b,2a \\
B1391 & 1a<o1 1b<o2 / 2a<1a0L,1b2 2b<1b1L & EE/DD & {-}{-}/{-}{-} & E & C:1b,2a \\
B1392 & 1a<o1 1b<o2 / 2a<1a0L,1b2 2b<1b2 & EE/DD & {-}{-}/{-}{-} & E & C:1b \\
B1393 & 1a<o1 1b<o2 / 2a<1a0L,1b2 2b<1b2L & EE/DD & {-}{-}/{-}{-} & E & C:1b \\
B1394 & 1a<o1 1b<o2 / 2a<1a1 2b<1a1,1b0L & EE/DD & {-}{-}/{-}{-} & E & C:1b \\
B1395 & 1a<o1 1b<o2 / 2a<1a1 2b<1a1,1b1 & EE/DD & {-}{-}/{-}{-} & E & C:1b \\
B1396 & 1a<o1 1b<o2 / 2a<1a1 2b<1a1,1b1L & EE/DD & {-}{-}/{-}{-} & E & C:1b \\
B1397 & 1a<o1 1b<o2 / 2a<1a1 2b<1a1,1b2 & EE/DD & {-}{-}/{-}{-} & E & C:1b,2b \\
B1398 & 1a<o1 1b<o2 / 2a<1a1 2b<1a1,1b2L & EE/DD & {-}{-}/{-}{-} & E & C:1b,2b \\
B1399 & 1a<o1 1b<o2 / 2a<1a1 2b<1a1L,1b1 & EE/DD & {-}{-}/{-}{-} & E & C:1b \\
B1400 & 1a<o1 1b<o2 / 2a<1a1 2b<1a1L,1b2 & EE/DD & {-}{-}/{-}{-} & E & C:1b,2b \\
B1401 & 1a<o1 1b<o2 / 2a<1a1 2b<1a2,1b0L & EE/DD & {-}{-}/{-}{-} & E & C:1b \\
B1402 & 1a<o1 1b<o2 / 2a<1a1 2b<1a2,1b1 & EE/DD & {-}{-}/{-}{-} & E & C:1b \\
B1403 & 1a<o1 1b<o2 / 2a<1a1 2b<1a2,1b1L & EE/DD & {-}{-}/{-}{-} & E & C:1b \\
B1404 & 1a<o1 1b<o2 / 2a<1a1 2b<1a2,1b2 & EE/DD & {-}{-}/{-}{-} & E & C:1b \\
B1405 & 1a<o1 1b<o2 / 2a<1a1 2b<1a2,1b2L & EE/DD & {-}{-}/{-}{-} & E & C:1b \\
B1406 & 1a<o1 1b<o2 / 2a<1a1 2b<1a2L,1b1 & EE/DD & {-}{-}/{-}{-} & E & C:1b \\
B1407 & 1a<o1 1b<o2 / 2a<1a1 2b<1a2L,1b2 & EE/DD & {-}{-}/{-}{-} & E & C:1b \\
B1408 & 1a<o1 1b<o2 / 2a<1a1,1b0L 2b<1a1L & EE/DD & {-}{-}/{-}{-} & E & C:1b \\
B1409 & 1a<o1 1b<o2 / 2a<1a1,1b0L 2b<1a2 & EE/DD & {-}{-}/{-}{-} & E & C:1b \\
B1410 & 1a<o1 1b<o2 / 2a<1a1,1b0L 2b<1a2L & EE/DD & {-}{-}/{-}{-} & E & C:1b \\
B1411 & 1a<o1 1b<o2 / 2a<1a1,1b0L 2b<1b1 & EE/DD & {-}{-}/{-}{-} & E & C:1b,2b \\
B1412 & 1a<o1 1b<o2 / 2a<1a1,1b0L 2b<1b2 & EE/DD & {-}{-}/{-}{-} & E & C:1b,2b \\
B1413 & 1a<o1 1b<o2 / 2a<1a1,1b1 2b<1a1L & EE/DD & {-}{-}/{-}{-} & E & C:1b \\
B1414 & 1a<o1 1b<o2 / 2a<1a1,1b1 2b<1a2 & EE/DD & {-}{-}/{-}{-} & E & C:1b \\
B1415 & 1a<o1 1b<o2 / 2a<1a1,1b1 2b<1a2L & EE/DD & {-}{-}/{-}{-} & E & C:1b \\
B1416 & 1a<o1 1b<o2 / 2a<1a1,1b1 2b<1b0L & EE/DN & {-}{-}/{-}{-} & E & C:1b \\
B1417 & 1a<o1 1b<o2 / 2a<1a1,1b1 2b<1b1 & EE/DD & {-}{-}/{-}{-} & E & C:1b \\
B1418 & 1a<o1 1b<o2 / 2a<1a1,1b1 2b<1b1L & EE/DD & {-}{-}/{-}{-} & E & C:1b \\
B1419 & 1a<o1 1b<o2 / 2a<1a1,1b1 2b<1b2 & EE/DD & {-}{-}/{-}{-} & E & C:1b,2b \\
B1420 & 1a<o1 1b<o2 / 2a<1a1,1b1 2b<1b2L & EE/DD & {-}{-}/{-}{-} & E & C:1b,2b \\
B1421 & 1a<o1 1b<o2 / 2a<1a1,1b1L 2b<1a1L & EE/DD & {-}{-}/{-}{-} & E & C:1b \\
B1422 & 1a<o1 1b<o2 / 2a<1a1,1b1L 2b<1a2 & EE/DD & {-}{-}/{-}{-} & E & C:1b \\
B1423 & 1a<o1 1b<o2 / 2a<1a1,1b1L 2b<1a2L & EE/DD & {-}{-}/{-}{-} & E & C:1b \\
B1424 & 1a<o1 1b<o2 / 2a<1a1,1b1L 2b<1b1 & EE/DD & {-}{-}/{-}{-} & E & C:1b \\
B1425 & 1a<o1 1b<o2 / 2a<1a1,1b1L 2b<1b2 & EE/DD & {-}{-}/{-}{-} & E & C:1b,2b \\
B1426 & 1a<o1 1b<o2 / 2a<1a1,1b2 2b<1a1L & EE/DD & {-}{-}/{-}{-} & E & C:1b,2a \\
B1427 & 1a<o1 1b<o2 / 2a<1a1,1b2 2b<1a2 & EE/DD & {-}{-}/{-}{-} & E & C:1b,2a \\
B1428 & 1a<o1 1b<o2 / 2a<1a1,1b2 2b<1a2L & EE/DD & {-}{-}/{-}{-} & E & C:1b,2a \\
B1429 & 1a<o1 1b<o2 / 2a<1a1,1b2 2b<1b0L & EE/DN & {-}{-}/{-}{-} & E & C:1b,2a \\
B1430 & 1a<o1 1b<o2 / 2a<1a1,1b2 2b<1b1 & EE/DD & {-}{-}/{-}{-} & E & C:1b,2a \\
B1431 & 1a<o1 1b<o2 / 2a<1a1,1b2 2b<1b1L & EE/DD & {-}{-}/{-}{-} & E & C:1b,2a \\
B1432 & 1a<o1 1b<o2 / 2a<1a1,1b2 2b<1b2 & EE/DD & {-}{-}/{-}{-} & E & C:1b \\
B1433 & 1a<o1 1b<o2 / 2a<1a1,1b2 2b<1b2L & EE/DD & {-}{-}/{-}{-} & E & C:1b \\
B1434 & 1a<o1 1b<o2 / 2a<1a1,1b2L 2b<1a1L & EE/DD & {-}{-}/{-}{-} & E & C:1b,2a \\
B1435 & 1a<o1 1b<o2 / 2a<1a1,1b2L 2b<1a2 & EE/DD & {-}{-}/{-}{-} & E & C:1b,2a \\
B1436 & 1a<o1 1b<o2 / 2a<1a1,1b2L 2b<1a2L & EE/DD & {-}{-}/{-}{-} & E & C:1b,2a \\
B1437 & 1a<o1 1b<o2 / 2a<1a1,1b2L 2b<1b1 & EE/DD & {-}{-}/{-}{-} & E & C:1b,2a \\
B1438 & 1a<o1 1b<o2 / 2a<1a1,1b2L 2b<1b2 & EE/DD & {-}{-}/{-}{-} & E & C:1b \\
B1439 & 1a<o1 1b<o2 / 2a<1a1L 2b<1a2,1b0L & EE/DD & {-}{-}/{-}{-} & E & C:1b \\
B1440 & 1a<o1 1b<o2 / 2a<1a1L 2b<1a2,1b1 & EE/DD & {-}{-}/{-}{-} & E & C:1b \\
B1441 & 1a<o1 1b<o2 / 2a<1a1L 2b<1a2,1b1L & EE/DD & {-}{-}/{-}{-} & E & C:1b \\
B1442 & 1a<o1 1b<o2 / 2a<1a1L 2b<1a2,1b2 & EE/DD & {-}{-}/{-}{-} & E & C:1b \\
B1443 & 1a<o1 1b<o2 / 2a<1a1L 2b<1a2,1b2L & EE/DD & {-}{-}/{-}{-} & E & C:1b \\
B1444 & 1a<o1 1b<o2 / 2a<1a1L,1b1 2b<1a2 & EE/DD & {-}{-}/{-}{-} & E & C:1b \\
B1445 & 1a<o1 1b<o2 / 2a<1a1L,1b1 2b<1b0L & EE/DN & {-}{-}/{-}{-} & E & C:1b \\
B1446 & 1a<o1 1b<o2 / 2a<1a1L,1b1 2b<1b1 & EE/DD & {-}{-}/{-}{-} & E & C:1b \\
B1447 & 1a<o1 1b<o2 / 2a<1a1L,1b1 2b<1b1L & EE/DD & {-}{-}/{-}{-} & E & C:1b \\
B1448 & 1a<o1 1b<o2 / 2a<1a1L,1b1 2b<1b2 & EE/DD & {-}{-}/{-}{-} & E & C:1b,2b \\
B1449 & 1a<o1 1b<o2 / 2a<1a1L,1b1 2b<1b2L & EE/DD & {-}{-}/{-}{-} & E & C:1b,2b \\
B1450 & 1a<o1 1b<o2 / 2a<1a1L,1b2 2b<1a2 & EE/DD & {-}{-}/{-}{-} & E & C:1b,2a \\
B1451 & 1a<o1 1b<o2 / 2a<1a1L,1b2 2b<1b0L & EE/DN & {-}{-}/{-}{-} & E & C:1b,2a \\
B1452 & 1a<o1 1b<o2 / 2a<1a1L,1b2 2b<1b1 & EE/DD & {-}{-}/{-}{-} & E & C:1b,2a \\
B1453 & 1a<o1 1b<o2 / 2a<1a1L,1b2 2b<1b1L & EE/DD & {-}{-}/{-}{-} & E & C:1b,2a \\
B1454 & 1a<o1 1b<o2 / 2a<1a1L,1b2 2b<1b2 & EE/DD & {-}{-}/{-}{-} & E & C:1b \\
B1455 & 1a<o1 1b<o2 / 2a<1a1L,1b2 2b<1b2L & EE/DD & {-}{-}/{-}{-} & E & C:1b \\
B1456 & 1a<o1 1b<o2 / 2a<1a2 2b<1a2,1b0L & EE/DD & {-}{-}/{-}{-} & E & C:1b \\
B1457 & 1a<o1 1b<o2 / 2a<1a2 2b<1a2,1b1 & EE/DD & {-}{-}/{-}{-} & E & C:1b \\
B1458 & 1a<o1 1b<o2 / 2a<1a2 2b<1a2,1b1L & EE/DD & {-}{-}/{-}{-} & E & C:1b \\
B1459 & 1a<o1 1b<o2 / 2a<1a2 2b<1a2,1b2 & EE/DD & {-}{-}/{-}{-} & E & C:1b \\
B1460 & 1a<o1 1b<o2 / 2a<1a2 2b<1a2,1b2L & EE/DD & {-}{-}/{-}{-} & E & C:1b \\
B1461 & 1a<o1 1b<o2 / 2a<1a2 2b<1a2L,1b1 & EE/DD & {-}{-}/{-}{-} & E & C:1b \\
B1462 & 1a<o1 1b<o2 / 2a<1a2 2b<1a2L,1b2 & EE/DD & {-}{-}/{-}{-} & E & C:1b \\
B1463 & 1a<o1 1b<o2 / 2a<1a2,1b0L 2b<1a2L & EE/DD & {-}{-}/{-}{-} & E & C:1b \\
B1464 & 1a<o1 1b<o2 / 2a<1a2,1b0L 2b<1b1 & EE/DD & {-}{-}/{-}{-} & E & C:1b,2b \\
B1465 & 1a<o1 1b<o2 / 2a<1a2,1b0L 2b<1b2 & EE/DD & {-}{-}/{-}{-} & E & C:1b,2b \\
B1466 & 1a<o1 1b<o2 / 2a<1a2,1b1 2b<1a2L & EE/DD & {-}{-}/{-}{-} & E & C:1b \\
B1467 & 1a<o1 1b<o2 / 2a<1a2,1b1 2b<1b0L & EE/DN & {-}{-}/{-}{-} & E & C:1b \\
B1468 & 1a<o1 1b<o2 / 2a<1a2,1b1 2b<1b1 & EE/DD & {-}{-}/{-}{-} & E & C:1b \\
B1469 & 1a<o1 1b<o2 / 2a<1a2,1b1 2b<1b1L & EE/DD & {-}{-}/{-}{-} & E & C:1b \\
B1470 & 1a<o1 1b<o2 / 2a<1a2,1b1 2b<1b2 & EE/DD & {-}{-}/{-}{-} & E & C:1b,2b \\
B1471 & 1a<o1 1b<o2 / 2a<1a2,1b1 2b<1b2L & EE/DD & {-}{-}/{-}{-} & E & C:1b,2b \\
B1472 & 1a<o1 1b<o2 / 2a<1a2,1b1L 2b<1a2L & EE/DD & {-}{-}/{-}{-} & E & C:1b \\
B1473 & 1a<o1 1b<o2 / 2a<1a2,1b1L 2b<1b1 & EE/DD & {-}{-}/{-}{-} & E & C:1b \\
B1474 & 1a<o1 1b<o2 / 2a<1a2,1b1L 2b<1b2 & EE/DD & {-}{-}/{-}{-} & E & C:1b,2b \\
B1475 & 1a<o1 1b<o2 / 2a<1a2,1b2 2b<1a2L & EE/DD & {-}{-}/{-}{-} & E & C:1b \\
B1476 & 1a<o1 1b<o2 / 2a<1a2,1b2 2b<1b0L & EE/DN & {-}{-}/{-}{-} & E & C:1b \\
B1477 & 1a<o1 1b<o2 / 2a<1a2,1b2 2b<1b1 & EE/DD & {-}{-}/{-}{-} & E & C:1b \\
B1478 & 1a<o1 1b<o2 / 2a<1a2,1b2 2b<1b1L & EE/DD & {-}{-}/{-}{-} & E & C:1b \\
B1479 & 1a<o1 1b<o2 / 2a<1a2,1b2 2b<1b2 & EE/DD & {-}{-}/{-}{-} & E & C:1b \\
B1480 & 1a<o1 1b<o2 / 2a<1a2,1b2 2b<1b2L & EE/DD & {-}{-}/{-}{-} & E & C:1b \\
B1481 & 1a<o1 1b<o2 / 2a<1a2,1b2L 2b<1a2L & EE/DD & {-}{-}/{-}{-} & E & C:1b \\
B1482 & 1a<o1 1b<o2 / 2a<1a2,1b2L 2b<1b1 & EE/DD & {-}{-}/{-}{-} & E & C:1b \\
B1483 & 1a<o1 1b<o2 / 2a<1a2,1b2L 2b<1b2 & EE/DD & {-}{-}/{-}{-} & E & C:1b \\
B1484 & 1a<o1 1b<o2 / 2a<1a2L,1b1 2b<1b0L & EE/DN & {-}{-}/{-}{-} & E & C:1b \\
B1485 & 1a<o1 1b<o2 / 2a<1a2L,1b1 2b<1b1 & EE/DD & {-}{-}/{-}{-} & E & C:1b \\
B1486 & 1a<o1 1b<o2 / 2a<1a2L,1b1 2b<1b1L & EE/DD & {-}{-}/{-}{-} & E & C:1b \\
B1487 & 1a<o1 1b<o2 / 2a<1a2L,1b1 2b<1b2 & EE/DD & {-}{-}/{-}{-} & E & C:1b,2b \\
B1488 & 1a<o1 1b<o2 / 2a<1a2L,1b1 2b<1b2L & EE/DD & {-}{-}/{-}{-} & E & C:1b,2b \\
B1489 & 1a<o1 1b<o2 / 2a<1a2L,1b2 2b<1b0L & EE/DN & {-}{-}/{-}{-} & E & C:1b \\
B1490 & 1a<o1 1b<o2 / 2a<1a2L,1b2 2b<1b1 & EE/DD & {-}{-}/{-}{-} & E & C:1b \\
B1491 & 1a<o1 1b<o2 / 2a<1a2L,1b2 2b<1b1L & EE/DD & {-}{-}/{-}{-} & E & C:1b \\
B1492 & 1a<o1 1b<o2 / 2a<1a2L,1b2 2b<1b2 & EE/DD & {-}{-}/{-}{-} & E & C:1b \\
B1493 & 1a<o1 1b<o2 / 2a<1a2L,1b2 2b<1b2L & EE/DD & {-}{-}/{-}{-} & E & C:1b \\
B1494 & 1a<o1 1b<o2L / 2a<1a0L 2b<1a1,1b0L & LI/ND & {-}+/{-}+ & B & PW:1b \\
B1495 & 1a<o1 1b<o2L / 2a<1a0L 2b<1a1,1b1 & LI/NE & {-}+/{-}{-} & BF & PW:1b \\
B1496 & 1a<o1 1b<o2L / 2a<1a0L 2b<1a1,1b1L & LI/NI & {-}+/{-}+ & BF & C:2b PW:1b,2b \\
B1497 & 1a<o1 1b<o2L / 2a<1a0L 2b<1a1,1b2 & LI/NE & {-}+/{-}{-} & BF & C:2b PW:1b \\
B1498 & 1a<o1 1b<o2L / 2a<1a0L 2b<1a1,1b2L & LI/NI & {-}+/{-}+ & BF & PW:1b,2b \\
B1499 & 1a<o1 1b<o2L / 2a<1a0L 2b<1a2,1b0L & LI/ND & {-}+/{-}+ & B & PW:1b \\
B1500 & 1a<o1 1b<o2L / 2a<1a0L 2b<1a2,1b1 & LI/NE & {-}+/{-}{-} & BF & PW:1b \\
B1501 & 1a<o1 1b<o2L / 2a<1a0L 2b<1a2,1b1L & LI/NI & {-}+/{-}+ & BF & C:2b PW:1b,2b \\
B1502 & 1a<o1 1b<o2L / 2a<1a0L 2b<1a2,1b2 & LI/NE & {-}+/{-}{-} & BF & PW:1b \\
B1503 & 1a<o1 1b<o2L / 2a<1a0L 2b<1a2,1b2L & LI/NI & {-}+/{-}+ & BF & C:2b PW:1b,2b \\
B1504 & 1a<o1 1b<o2L / 2a<1a0L,1b1 2b<1a1 & LI/ID & {-}+/{-}{-} & B & PW:1b,2a \\
B1505 & 1a<o1 1b<o2L / 2a<1a0L,1b1 2b<1a2 & LI/ID & {-}+/{-}{-} & B & PW:1b,2a \\
B1506 & 1a<o1 1b<o2L / 2a<1a0L,1b1 2b<1b0L & LI/IN & {-}+/{-}+ & B & PW:1b,2a \\
B1507 & 1a<o1 1b<o2L / 2a<1a0L,1b1 2b<1b1 & LI/EE & {-}+/{-}{-} & BE & PW:1b \\
B1508 & 1a<o1 1b<o2L / 2a<1a0L,1b1 2b<1b1L & LI/LI & {-}+/{-}+ & B & C:2b PW:1b,2b \\
B1509 & 1a<o1 1b<o2L / 2a<1a0L,1b1 2b<1b2 & LI/EE & {-}+/{-}{-} & BE & C:2b PW:1b \\
B1510 & 1a<o1 1b<o2L / 2a<1a0L,1b1 2b<1b2L & LI/LI & {-}+/{-}+ & B & PW:1b,2b \\
B1511 & 1a<o1 1b<o2L / 2a<1a0L,1b2 2b<1a1 & LI/ID & {-}+/{-}{-} & B & PW:1b,2a \\
B1512 & 1a<o1 1b<o2L / 2a<1a0L,1b2 2b<1a2 & LI/ID & {-}+/{-}{-} & B & PW:1b,2a \\
B1513 & 1a<o1 1b<o2L / 2a<1a0L,1b2 2b<1b0L & LI/IN & {-}+/{-}+ & B & PW:1b,2a \\
B1514 & 1a<o1 1b<o2L / 2a<1a0L,1b2 2b<1b1 & LI/EE & {-}+/{-}{-} & BE & C:2a PW:1b \\
B1515 & 1a<o1 1b<o2L / 2a<1a0L,1b2 2b<1b1L & LI/LI & {-}+/{-}+ & B & C:2a,2b PW:1b,2b \\
B1516 & 1a<o1 1b<o2L / 2a<1a0L,1b2 2b<1b2 & LI/EE & {-}+/{-}{-} & BE & PW:1b \\
B1517 & 1a<o1 1b<o2L / 2a<1a0L,1b2 2b<1b2L & LI/LI & {-}+/{-}+ & B & C:2b PW:1b,2b \\
B1518 & 1a<o1 1b<o2L / 2a<1a1 2b<1a1,1b0L & LI/DD & {-}+/{-}+ & B & PW:1b \\
B1519 & 1a<o1 1b<o2L / 2a<1a1 2b<1a1,1b1 & LI/DE & {-}+/{-}{-} & BF & PW:1b \\
B1520 & 1a<o1 1b<o2L / 2a<1a1 2b<1a1,1b1L & LI/DI & {-}+/{-}+ & BF & C:2b PW:1b,2b \\
B1521 & 1a<o1 1b<o2L / 2a<1a1 2b<1a1,1b2 & LI/DE & {-}+/{-}{-} & BF & C:2b PW:1b \\
B1522 & 1a<o1 1b<o2L / 2a<1a1 2b<1a1,1b2L & LI/DI & {-}+/{-}+ & BF & PW:1b,2b \\
B1523 & 1a<o1 1b<o2L / 2a<1a1 2b<1a1L,1b1 & LI/DE & {-}+/{-}{-} & BF & PW:1b \\
B1524 & 1a<o1 1b<o2L / 2a<1a1 2b<1a1L,1b2 & LI/DE & {-}+/{-}{-} & BF & C:2b PW:1b \\
B1525 & 1a<o1 1b<o2L / 2a<1a1 2b<1a2,1b0L & LI/DD & {-}+/{-}+ & B & PW:1b \\
B1526 & 1a<o1 1b<o2L / 2a<1a1 2b<1a2,1b1 & LI/DE & {-}+/{-}{-} & BF & PW:1b \\
B1527 & 1a<o1 1b<o2L / 2a<1a1 2b<1a2,1b1L & LI/DI & {-}+/{-}+ & BF & C:2b PW:1b,2b \\
B1528 & 1a<o1 1b<o2L / 2a<1a1 2b<1a2,1b2 & LI/DE & {-}+/{-}{-} & BF & PW:1b \\
B1529 & 1a<o1 1b<o2L / 2a<1a1 2b<1a2,1b2L & LI/DI & {-}+/{-}+ & BF & C:2b PW:1b,2b \\
B1530 & 1a<o1 1b<o2L / 2a<1a1 2b<1a2L,1b1 & LI/DE & {-}+/{-}{-} & BF & PW:1b \\
B1531 & 1a<o1 1b<o2L / 2a<1a1 2b<1a2L,1b2 & LI/DE & {-}+/{-}{-} & BF & PW:1b \\
B1532 & 1a<o1 1b<o2L / 2a<1a1,1b0L 2b<1a1L & LI/DD & {-}+/+{-} & B & PW:1b \\
B1533 & 1a<o1 1b<o2L / 2a<1a1,1b0L 2b<1a2 & LI/DD & {-}+/+{-} & B & PW:1b \\
B1534 & 1a<o1 1b<o2L / 2a<1a1,1b0L 2b<1a2L & LI/DD & {-}+/+{-} & B & PW:1b \\
B1535 & 1a<o1 1b<o2L / 2a<1a1,1b0L 2b<1b1 & LI/DI & {-}+/+{-} & B & PW:1b,2b \\
B1536 & 1a<o1 1b<o2L / 2a<1a1,1b0L 2b<1b2 & LI/DI & {-}+/+{-} & B & PW:1b,2b \\
B1537 & 1a<o1 1b<o2L / 2a<1a1,1b1 2b<1a1L & LI/ED & {-}+/{-}{-} & BF & PW:1b \\
B1538 & 1a<o1 1b<o2L / 2a<1a1,1b1 2b<1a2 & LI/ED & {-}+/{-}{-} & BF & PW:1b \\
B1539 & 1a<o1 1b<o2L / 2a<1a1,1b1 2b<1a2L & LI/ED & {-}+/{-}{-} & BF & PW:1b \\
B1540 & 1a<o1 1b<o2L / 2a<1a1,1b1 2b<1b0L & LI/EN & {-}+/{-}+ & BF & PW:1b \\
B1541 & 1a<o1 1b<o2L / 2a<1a1,1b1 2b<1b1 & LI/EE & {-}+/{-}{-} & BE & PW:1b \\
B1542 & 1a<o1 1b<o2L / 2a<1a1,1b1 2b<1b1L & LI/LI & {-}+/{-}+ & B & C:2b PW:1b,2b \\
B1543 & 1a<o1 1b<o2L / 2a<1a1,1b1 2b<1b2 & LI/EE & {-}+/{-}{-} & BE & C:2b PW:1b \\
B1544 & 1a<o1 1b<o2L / 2a<1a1,1b1 2b<1b2L & LI/LI & {-}+/{-}+ & B & PW:1b,2b \\
B1545 & 1a<o1 1b<o2L / 2a<1a1,1b1L 2b<1a1L & LI/ID & {-}+/+{-} & BF & C:2a PW:1b,2a \\
B1546 & 1a<o1 1b<o2L / 2a<1a1,1b1L 2b<1a2 & LI/ID & {-}+/+{-} & BF & C:2a PW:1b,2a \\
B1547 & 1a<o1 1b<o2L / 2a<1a1,1b1L 2b<1a2L & LI/ID & {-}+/+{-} & BF & C:2a PW:1b,2a \\
B1548 & 1a<o1 1b<o2L / 2a<1a1,1b1L 2b<1b1 & LI/IL & {-}+/+{-} & B & C:2a PW:1b,2a \\
B1549 & 1a<o1 1b<o2L / 2a<1a1,1b1L 2b<1b2 & LI/IL & {-}+/+{-} & B & C:2a,2b PW:1b,2a \\
B1550 & 1a<o1 1b<o2L / 2a<1a1,1b2 2b<1a1L & LI/ED & {-}+/{-}{-} & BF & C:2a PW:1b \\
B1551 & 1a<o1 1b<o2L / 2a<1a1,1b2 2b<1a2 & LI/ED & {-}+/{-}{-} & BF & C:2a PW:1b \\
B1552 & 1a<o1 1b<o2L / 2a<1a1,1b2 2b<1a2L & LI/ED & {-}+/{-}{-} & BF & C:2a PW:1b \\
B1553 & 1a<o1 1b<o2L / 2a<1a1,1b2 2b<1b0L & LI/EN & {-}+/{-}+ & BF & C:2a PW:1b \\
B1554 & 1a<o1 1b<o2L / 2a<1a1,1b2 2b<1b1 & LI/EE & {-}+/{-}{-} & BE & C:2a PW:1b \\
B1555 & 1a<o1 1b<o2L / 2a<1a1,1b2 2b<1b1L & LI/LI & {-}+/{-}+ & B & C:2a,2b PW:1b,2b \\
B1556 & 1a<o1 1b<o2L / 2a<1a1,1b2 2b<1b2 & LI/EE & {-}+/{-}{-} & BE & PW:1b \\
B1557 & 1a<o1 1b<o2L / 2a<1a1,1b2 2b<1b2L & LI/LI & {-}+/{-}+ & B & C:2b PW:1b,2b \\
B1558 & 1a<o1 1b<o2L / 2a<1a1,1b2L 2b<1a1L & LI/ID & {-}+/+{-} & BF & PW:1b,2a \\
B1559 & 1a<o1 1b<o2L / 2a<1a1,1b2L 2b<1a2 & LI/ID & {-}+/+{-} & BF & PW:1b,2a \\
B1560 & 1a<o1 1b<o2L / 2a<1a1,1b2L 2b<1a2L & LI/ID & {-}+/+{-} & BF & PW:1b,2a \\
B1561 & 1a<o1 1b<o2L / 2a<1a1,1b2L 2b<1b1 & LI/IL & {-}+/+{-} & B & PW:1b,2a \\
B1562 & 1a<o1 1b<o2L / 2a<1a1,1b2L 2b<1b2 & LI/IL & {-}+/+{-} & B & C:2a PW:1b,2a \\
B1563 & 1a<o1 1b<o2L / 2a<1a1L 2b<1a2,1b0L & LI/DD & {-}+/{-}+ & B & PW:1b \\
B1564 & 1a<o1 1b<o2L / 2a<1a1L 2b<1a2,1b1 & LI/DE & {-}+/{-}{-} & BF & PW:1b \\
B1565 & 1a<o1 1b<o2L / 2a<1a1L 2b<1a2,1b1L & LI/DI & {-}+/{-}+ & BF & C:2b PW:1b,2b \\
B1566 & 1a<o1 1b<o2L / 2a<1a1L 2b<1a2,1b2 & LI/DE & {-}+/{-}{-} & BF & PW:1b \\
B1567 & 1a<o1 1b<o2L / 2a<1a1L 2b<1a2,1b2L & LI/DI & {-}+/{-}+ & BF & C:2b PW:1b,2b \\
B1568 & 1a<o1 1b<o2L / 2a<1a1L,1b1 2b<1a2 & LI/ED & {-}+/{-}{-} & BF & PW:1b \\
B1569 & 1a<o1 1b<o2L / 2a<1a1L,1b1 2b<1b0L & LI/EN & {-}+/{-}+ & BF & PW:1b \\
B1570 & 1a<o1 1b<o2L / 2a<1a1L,1b1 2b<1b1 & LI/EE & {-}+/{-}{-} & BE & PW:1b \\
B1571 & 1a<o1 1b<o2L / 2a<1a1L,1b1 2b<1b1L & LI/LI & {-}+/{-}+ & B & C:2b PW:1b,2b \\
B1572 & 1a<o1 1b<o2L / 2a<1a1L,1b1 2b<1b2 & LI/EE & {-}+/{-}{-} & BE & C:2b PW:1b \\
B1573 & 1a<o1 1b<o2L / 2a<1a1L,1b1 2b<1b2L & LI/LI & {-}+/{-}+ & B & PW:1b,2b \\
B1574 & 1a<o1 1b<o2L / 2a<1a1L,1b2 2b<1a2 & LI/ED & {-}+/{-}{-} & BF & C:2a PW:1b \\
B1575 & 1a<o1 1b<o2L / 2a<1a1L,1b2 2b<1b0L & LI/EN & {-}+/{-}+ & BF & C:2a PW:1b \\
B1576 & 1a<o1 1b<o2L / 2a<1a1L,1b2 2b<1b1 & LI/EE & {-}+/{-}{-} & BE & C:2a PW:1b \\
B1577 & 1a<o1 1b<o2L / 2a<1a1L,1b2 2b<1b1L & LI/LI & {-}+/{-}+ & B & C:2a,2b PW:1b,2b \\
B1578 & 1a<o1 1b<o2L / 2a<1a1L,1b2 2b<1b2 & LI/EE & {-}+/{-}{-} & BE & PW:1b \\
B1579 & 1a<o1 1b<o2L / 2a<1a1L,1b2 2b<1b2L & LI/LI & {-}+/{-}+ & B & C:2b PW:1b,2b \\
B1580 & 1a<o1 1b<o2L / 2a<1a2 2b<1a2,1b0L & LI/DD & {-}+/{-}+ & B & PW:1b \\
B1581 & 1a<o1 1b<o2L / 2a<1a2 2b<1a2,1b1 & LI/DE & {-}+/{-}{-} & BF & PW:1b \\
B1582 & 1a<o1 1b<o2L / 2a<1a2 2b<1a2,1b1L & LI/DI & {-}+/{-}+ & BF & C:2b PW:1b,2b \\
B1583 & 1a<o1 1b<o2L / 2a<1a2 2b<1a2,1b2 & LI/DE & {-}+/{-}{-} & BF & PW:1b \\
B1584 & 1a<o1 1b<o2L / 2a<1a2 2b<1a2,1b2L & LI/DI & {-}+/{-}+ & BF & C:2b PW:1b,2b \\
B1585 & 1a<o1 1b<o2L / 2a<1a2 2b<1a2L,1b1 & LI/DE & {-}+/{-}{-} & BF & PW:1b \\
B1586 & 1a<o1 1b<o2L / 2a<1a2 2b<1a2L,1b2 & LI/DE & {-}+/{-}{-} & BF & PW:1b \\
B1587 & 1a<o1 1b<o2L / 2a<1a2,1b0L 2b<1a2L & LI/DD & {-}+/+{-} & B & PW:1b \\
B1588 & 1a<o1 1b<o2L / 2a<1a2,1b0L 2b<1b1 & LI/DI & {-}+/+{-} & B & PW:1b,2b \\
B1589 & 1a<o1 1b<o2L / 2a<1a2,1b0L 2b<1b2 & LI/DI & {-}+/+{-} & B & PW:1b,2b \\
B1590 & 1a<o1 1b<o2L / 2a<1a2,1b1 2b<1a2L & LI/ED & {-}+/{-}{-} & BF & PW:1b \\
B1591 & 1a<o1 1b<o2L / 2a<1a2,1b1 2b<1b0L & LI/EN & {-}+/{-}+ & BF & PW:1b \\
B1592 & 1a<o1 1b<o2L / 2a<1a2,1b1 2b<1b1 & LI/EE & {-}+/{-}{-} & BE & PW:1b \\
B1593 & 1a<o1 1b<o2L / 2a<1a2,1b1 2b<1b1L & LI/LI & {-}+/{-}+ & B & C:2b PW:1b,2b \\
B1594 & 1a<o1 1b<o2L / 2a<1a2,1b1 2b<1b2 & LI/EE & {-}+/{-}{-} & BE & C:2b PW:1b \\
B1595 & 1a<o1 1b<o2L / 2a<1a2,1b1 2b<1b2L & LI/LI & {-}+/{-}+ & B & PW:1b,2b \\
B1596 & 1a<o1 1b<o2L / 2a<1a2,1b1L 2b<1a2L & LI/ID & {-}+/+{-} & BF & C:2a PW:1b,2a \\
B1597 & 1a<o1 1b<o2L / 2a<1a2,1b1L 2b<1b1 & LI/IL & {-}+/+{-} & B & C:2a PW:1b,2a \\
B1598 & 1a<o1 1b<o2L / 2a<1a2,1b1L 2b<1b2 & LI/IL & {-}+/+{-} & B & C:2a,2b PW:1b,2a \\
B1599 & 1a<o1 1b<o2L / 2a<1a2,1b2 2b<1a2L & LI/ED & {-}+/{-}{-} & BF & PW:1b \\
B1600 & 1a<o1 1b<o2L / 2a<1a2,1b2 2b<1b0L & LI/EN & {-}+/{-}+ & BF & PW:1b \\
B1601 & 1a<o1 1b<o2L / 2a<1a2,1b2 2b<1b1 & LI/EE & {-}+/{-}{-} & BE & PW:1b \\
B1602 & 1a<o1 1b<o2L / 2a<1a2,1b2 2b<1b1L & LI/LI & {-}+/{-}+ & B & C:2b PW:1b,2b \\
B1603 & 1a<o1 1b<o2L / 2a<1a2,1b2 2b<1b2 & LI/EE & {-}+/{-}{-} & BE & PW:1b \\
B1604 & 1a<o1 1b<o2L / 2a<1a2,1b2 2b<1b2L & LI/LI & {-}+/{-}+ & B & C:2b PW:1b,2b \\
B1605 & 1a<o1 1b<o2L / 2a<1a2,1b2L 2b<1a2L & LI/ID & {-}+/+{-} & BF & C:2a PW:1b,2a \\
B1606 & 1a<o1 1b<o2L / 2a<1a2,1b2L 2b<1b1 & LI/IL & {-}+/+{-} & B & C:2a PW:1b,2a \\
B1607 & 1a<o1 1b<o2L / 2a<1a2,1b2L 2b<1b2 & LI/IL & {-}+/+{-} & B & C:2a PW:1b,2a \\
B1608 & 1a<o1 1b<o2L / 2a<1a2L,1b1 2b<1b0L & LI/EN & {-}+/{-}+ & BF & PW:1b \\
B1609 & 1a<o1 1b<o2L / 2a<1a2L,1b1 2b<1b1 & LI/EE & {-}+/{-}{-} & BE & PW:1b \\
B1610 & 1a<o1 1b<o2L / 2a<1a2L,1b1 2b<1b1L & LI/LI & {-}+/{-}+ & B & C:2b PW:1b,2b \\
B1611 & 1a<o1 1b<o2L / 2a<1a2L,1b1 2b<1b2 & LI/EE & {-}+/{-}{-} & BE & C:2b PW:1b \\
B1612 & 1a<o1 1b<o2L / 2a<1a2L,1b1 2b<1b2L & LI/LI & {-}+/{-}+ & B & PW:1b,2b \\
B1613 & 1a<o1 1b<o2L / 2a<1a2L,1b2 2b<1b0L & LI/EN & {-}+/{-}+ & BF & PW:1b \\
B1614 & 1a<o1 1b<o2L / 2a<1a2L,1b2 2b<1b1 & LI/EE & {-}+/{-}{-} & BE & PW:1b \\
B1615 & 1a<o1 1b<o2L / 2a<1a2L,1b2 2b<1b1L & LI/LI & {-}+/{-}+ & B & C:2b PW:1b,2b \\
B1616 & 1a<o1 1b<o2L / 2a<1a2L,1b2 2b<1b2 & LI/EE & {-}+/{-}{-} & BE & PW:1b \\
B1617 & 1a<o1 1b<o2L / 2a<1a2L,1b2 2b<1b2L & LI/LI & {-}+/{-}+ & B & C:2b PW:1b,2b \\
B1618 & 1a<o1L 1b<o2 / 2a<1a0L 2b<1a1,1b0L & IL/NI & +{-}/+{-} & B & C:1a,1b,2b PW:1a,2b \\
B1619 & 1a<o1L 1b<o2 / 2a<1a0L 2b<1a1,1b1 & IL/NE & +{-}/+{-} & BF & C:1a,1b PW:1a \\
B1620 & 1a<o1L 1b<o2 / 2a<1a0L 2b<1a1,1b1L & IL/NE & +{-}/+{-} & BF & C:1a,1b PW:1a \\
B1621 & 1a<o1L 1b<o2 / 2a<1a0L 2b<1a1,1b2 & IL/NE & +{-}/+{-} & BF & C:1a,1b,2b PW:1a \\
B1622 & 1a<o1L 1b<o2 / 2a<1a0L 2b<1a1,1b2L & IL/NE & +{-}/+{-} & BF & C:1a,1b,2b PW:1a \\
B1623 & 1a<o1L 1b<o2 / 2a<1a0L 2b<1a2,1b0L & IL/NI & +{-}/+{-} & B & C:1a,1b,2b PW:1a,2b \\
B1624 & 1a<o1L 1b<o2 / 2a<1a0L 2b<1a2,1b1 & IL/NE & +{-}/+{-} & BF & C:1a,1b PW:1a \\
B1625 & 1a<o1L 1b<o2 / 2a<1a0L 2b<1a2,1b1L & IL/NE & +{-}/+{-} & BF & C:1a,1b PW:1a \\
B1626 & 1a<o1L 1b<o2 / 2a<1a0L 2b<1a2,1b2 & IL/NE & +{-}/+{-} & BF & C:1a,1b PW:1a \\
B1627 & 1a<o1L 1b<o2 / 2a<1a0L 2b<1a2,1b2L & IL/NE & +{-}/+{-} & BF & C:1a,1b PW:1a \\
B1628 & 1a<o1L 1b<o2 / 2a<1a0L,1b1 2b<1a1 & IL/DI & +{-}/+{-} & B & C:1a,1b,2a,2b PW:1a,2b \\
B1629 & 1a<o1L 1b<o2 / 2a<1a0L,1b1 2b<1a2 & IL/DI & +{-}/+{-} & B & C:1a,1b,2a,2b PW:1a,2b \\
B1630 & 1a<o1L 1b<o2 / 2a<1a0L,1b1 2b<1b0L & IL/DN & +{-}/+{-} & B & C:1a,1b,2a PW:1a \\
B1631 & 1a<o1L 1b<o2 / 2a<1a0L,1b1 2b<1b1 & IL/DD & +{-}/+{-} & B & C:1a,1b PW:1a \\
B1632 & 1a<o1L 1b<o2 / 2a<1a0L,1b1 2b<1b1L & IL/DD & +{-}/+{-} & B & C:1a,1b PW:1a \\
B1633 & 1a<o1L 1b<o2 / 2a<1a0L,1b1 2b<1b2 & IL/DD & +{-}/+{-} & B & C:1a,1b,2b PW:1a \\
B1634 & 1a<o1L 1b<o2 / 2a<1a0L,1b1 2b<1b2L & IL/DD & +{-}/+{-} & B & C:1a,1b,2b PW:1a \\
B1635 & 1a<o1L 1b<o2 / 2a<1a0L,1b2 2b<1a1 & IL/DI & +{-}/+{-} & B & C:1a,1b,2a,2b PW:1a,2b \\
B1636 & 1a<o1L 1b<o2 / 2a<1a0L,1b2 2b<1a2 & IL/DI & +{-}/+{-} & B & C:1a,1b,2a,2b PW:1a,2b \\
B1637 & 1a<o1L 1b<o2 / 2a<1a0L,1b2 2b<1b0L & IL/DN & +{-}/+{-} & B & C:1a,1b,2a PW:1a \\
B1638 & 1a<o1L 1b<o2 / 2a<1a0L,1b2 2b<1b1 & IL/DD & +{-}/+{-} & B & C:1a,1b,2a PW:1a \\
B1639 & 1a<o1L 1b<o2 / 2a<1a0L,1b2 2b<1b1L & IL/DD & +{-}/+{-} & B & C:1a,1b,2a PW:1a \\
B1640 & 1a<o1L 1b<o2 / 2a<1a0L,1b2 2b<1b2 & IL/DD & +{-}/+{-} & B & C:1a,1b PW:1a \\
B1641 & 1a<o1L 1b<o2 / 2a<1a0L,1b2 2b<1b2L & IL/DD & +{-}/+{-} & B & C:1a,1b PW:1a \\
B1642 & 1a<o1L 1b<o2 / 2a<1a1 2b<1a1,1b0L & IL/EE & +{-}/{-}{-} & BE & C:1a,1b PW:1a \\
B1643 & 1a<o1L 1b<o2 / 2a<1a1 2b<1a1,1b1 & IL/EE & +{-}/{-}{-} & BE & C:1a,1b PW:1a \\
B1644 & 1a<o1L 1b<o2 / 2a<1a1 2b<1a1,1b1L & IL/EE & +{-}/{-}{-} & BE & C:1a,1b PW:1a \\
B1645 & 1a<o1L 1b<o2 / 2a<1a1 2b<1a1,1b2 & IL/EE & +{-}/{-}{-} & BE & C:1a,1b,2b PW:1a \\
B1646 & 1a<o1L 1b<o2 / 2a<1a1 2b<1a1,1b2L & IL/EE & +{-}/{-}{-} & BE & C:1a,1b,2b PW:1a \\
B1647 & 1a<o1L 1b<o2 / 2a<1a1 2b<1a1L,1b1 & IL/LI & +{-}/{-}+ & B & C:1a,1b,2b PW:1a,2b \\
B1648 & 1a<o1L 1b<o2 / 2a<1a1 2b<1a1L,1b2 & IL/LI & +{-}/{-}+ & B & C:1a,1b PW:1a,2b \\
B1649 & 1a<o1L 1b<o2 / 2a<1a1 2b<1a2,1b0L & IL/EE & +{-}/{-}{-} & BE & C:1a,1b PW:1a \\
B1650 & 1a<o1L 1b<o2 / 2a<1a1 2b<1a2,1b1 & IL/EE & +{-}/{-}{-} & BE & C:1a,1b PW:1a \\
B1651 & 1a<o1L 1b<o2 / 2a<1a1 2b<1a2,1b1L & IL/EE & +{-}/{-}{-} & BE & C:1a,1b PW:1a \\
B1652 & 1a<o1L 1b<o2 / 2a<1a1 2b<1a2,1b2 & IL/EE & +{-}/{-}{-} & BE & C:1a,1b PW:1a \\
B1653 & 1a<o1L 1b<o2 / 2a<1a1 2b<1a2,1b2L & IL/EE & +{-}/{-}{-} & BE & C:1a,1b PW:1a \\
B1654 & 1a<o1L 1b<o2 / 2a<1a1 2b<1a2L,1b1 & IL/LI & +{-}/{-}+ & B & C:1a,1b,2b PW:1a,2b \\
B1655 & 1a<o1L 1b<o2 / 2a<1a1 2b<1a2L,1b2 & IL/LI & +{-}/{-}+ & B & C:1a,1b,2b PW:1a,2b \\
B1656 & 1a<o1L 1b<o2 / 2a<1a1,1b0L 2b<1a1L & IL/LI & +{-}/{-}+ & B & C:1a,1b,2b PW:1a,2b \\
B1657 & 1a<o1L 1b<o2 / 2a<1a1,1b0L 2b<1a2 & IL/EE & +{-}/{-}{-} & BE & C:1a,1b PW:1a \\
B1658 & 1a<o1L 1b<o2 / 2a<1a1,1b0L 2b<1a2L & IL/LI & +{-}/{-}+ & B & C:1a,1b,2b PW:1a,2b \\
B1659 & 1a<o1L 1b<o2 / 2a<1a1,1b0L 2b<1b1 & IL/ID & +{-}/{-}{-} & B & C:1a,1b,2a,2b PW:1a,2a \\
B1660 & 1a<o1L 1b<o2 / 2a<1a1,1b0L 2b<1b2 & IL/ID & +{-}/{-}{-} & B & C:1a,1b,2a,2b PW:1a,2a \\
B1661 & 1a<o1L 1b<o2 / 2a<1a1,1b1 2b<1a1L & IL/LI & +{-}/{-}+ & B & C:1a,1b,2b PW:1a,2b \\
B1662 & 1a<o1L 1b<o2 / 2a<1a1,1b1 2b<1a2 & IL/EE & +{-}/{-}{-} & BE & C:1a,1b PW:1a \\
B1663 & 1a<o1L 1b<o2 / 2a<1a1,1b1 2b<1a2L & IL/LI & +{-}/{-}+ & B & C:1a,1b,2b PW:1a,2b \\
B1664 & 1a<o1L 1b<o2 / 2a<1a1,1b1 2b<1b0L & IL/EN & +{-}/{-}{-} & BF & C:1a,1b PW:1a \\
B1665 & 1a<o1L 1b<o2 / 2a<1a1,1b1 2b<1b1 & IL/ED & +{-}/{-}{-} & BF & C:1a,1b PW:1a \\
B1666 & 1a<o1L 1b<o2 / 2a<1a1,1b1 2b<1b1L & IL/ED & +{-}/{-}{-} & BF & C:1a,1b PW:1a \\
B1667 & 1a<o1L 1b<o2 / 2a<1a1,1b1 2b<1b2 & IL/ED & +{-}/{-}{-} & BF & C:1a,1b,2b PW:1a \\
B1668 & 1a<o1L 1b<o2 / 2a<1a1,1b1 2b<1b2L & IL/ED & +{-}/{-}{-} & BF & C:1a,1b,2b PW:1a \\
B1669 & 1a<o1L 1b<o2 / 2a<1a1,1b1L 2b<1a1L & IL/LI & +{-}/{-}+ & B & C:1a,1b,2b PW:1a,2b \\
B1670 & 1a<o1L 1b<o2 / 2a<1a1,1b1L 2b<1a2 & IL/EE & +{-}/{-}{-} & BE & C:1a,1b PW:1a \\
B1671 & 1a<o1L 1b<o2 / 2a<1a1,1b1L 2b<1a2L & IL/LI & +{-}/{-}+ & B & C:1a,1b,2b PW:1a,2b \\
B1672 & 1a<o1L 1b<o2 / 2a<1a1,1b1L 2b<1b1 & IL/ED & +{-}/{-}{-} & BF & C:1a,1b PW:1a \\
B1673 & 1a<o1L 1b<o2 / 2a<1a1,1b1L 2b<1b2 & IL/ED & +{-}/{-}{-} & BF & C:1a,1b,2b PW:1a \\
B1674 & 1a<o1L 1b<o2 / 2a<1a1,1b2 2b<1a1L & IL/LI & +{-}/{-}+ & B & C:1a,1b,2a,2b PW:1a,2b \\
B1675 & 1a<o1L 1b<o2 / 2a<1a1,1b2 2b<1a2 & IL/EE & +{-}/{-}{-} & BE & C:1a,1b,2a PW:1a \\
B1676 & 1a<o1L 1b<o2 / 2a<1a1,1b2 2b<1a2L & IL/LI & +{-}/{-}+ & B & C:1a,1b,2a,2b PW:1a,2b \\
B1677 & 1a<o1L 1b<o2 / 2a<1a1,1b2 2b<1b0L & IL/EN & +{-}/{-}{-} & BF & C:1a,1b,2a PW:1a \\
B1678 & 1a<o1L 1b<o2 / 2a<1a1,1b2 2b<1b1 & IL/ED & +{-}/{-}{-} & BF & C:1a,1b,2a PW:1a \\
B1679 & 1a<o1L 1b<o2 / 2a<1a1,1b2 2b<1b1L & IL/ED & +{-}/{-}{-} & BF & C:1a,1b,2a PW:1a \\
B1680 & 1a<o1L 1b<o2 / 2a<1a1,1b2 2b<1b2 & IL/ED & +{-}/{-}{-} & BF & C:1a,1b PW:1a \\
B1681 & 1a<o1L 1b<o2 / 2a<1a1,1b2 2b<1b2L & IL/ED & +{-}/{-}{-} & BF & C:1a,1b PW:1a \\
B1682 & 1a<o1L 1b<o2 / 2a<1a1,1b2L 2b<1a1L & IL/LI & +{-}/{-}+ & B & C:1a,1b,2a,2b PW:1a,2b \\
B1683 & 1a<o1L 1b<o2 / 2a<1a1,1b2L 2b<1a2 & IL/EE & +{-}/{-}{-} & BE & C:1a,1b,2a PW:1a \\
B1684 & 1a<o1L 1b<o2 / 2a<1a1,1b2L 2b<1a2L & IL/LI & +{-}/{-}+ & B & C:1a,1b,2a,2b PW:1a,2b \\
B1685 & 1a<o1L 1b<o2 / 2a<1a1,1b2L 2b<1b1 & IL/ED & +{-}/{-}{-} & BF & C:1a,1b,2a PW:1a \\
B1686 & 1a<o1L 1b<o2 / 2a<1a1,1b2L 2b<1b2 & IL/ED & +{-}/{-}{-} & BF & C:1a,1b PW:1a \\
B1687 & 1a<o1L 1b<o2 / 2a<1a1L 2b<1a2,1b0L & IL/IL & +{-}/+{-} & B & C:1a,1b,2a PW:1a,2a \\
B1688 & 1a<o1L 1b<o2 / 2a<1a1L 2b<1a2,1b1 & IL/IL & +{-}/+{-} & B & C:1a,1b,2a PW:1a,2a \\
B1689 & 1a<o1L 1b<o2 / 2a<1a1L 2b<1a2,1b1L & IL/IL & +{-}/+{-} & B & C:1a,1b,2a PW:1a,2a \\
B1690 & 1a<o1L 1b<o2 / 2a<1a1L 2b<1a2,1b2 & IL/IL & +{-}/+{-} & B & C:1a,1b,2a PW:1a,2a \\
B1691 & 1a<o1L 1b<o2 / 2a<1a1L 2b<1a2,1b2L & IL/IL & +{-}/+{-} & B & C:1a,1b,2a PW:1a,2a \\
B1692 & 1a<o1L 1b<o2 / 2a<1a1L,1b1 2b<1a2 & IL/IL & +{-}/+{-} & B & C:1a,1b,2a PW:1a,2a \\
B1693 & 1a<o1L 1b<o2 / 2a<1a1L,1b1 2b<1b0L & IL/IN & +{-}/+{-} & BF & C:1a,1b,2a PW:1a,2a \\
B1694 & 1a<o1L 1b<o2 / 2a<1a1L,1b1 2b<1b1 & IL/ID & +{-}/+{-} & BF & C:1a,1b,2a PW:1a,2a \\
B1695 & 1a<o1L 1b<o2 / 2a<1a1L,1b1 2b<1b1L & IL/ID & +{-}/+{-} & BF & C:1a,1b,2a PW:1a,2a \\
B1696 & 1a<o1L 1b<o2 / 2a<1a1L,1b1 2b<1b2 & IL/ID & +{-}/+{-} & BF & C:1a,1b,2a,2b PW:1a,2a \\
B1697 & 1a<o1L 1b<o2 / 2a<1a1L,1b1 2b<1b2L & IL/ID & +{-}/+{-} & BF & C:1a,1b,2a,2b PW:1a,2a \\
B1698 & 1a<o1L 1b<o2 / 2a<1a1L,1b2 2b<1a2 & IL/IL & +{-}/+{-} & B & C:1a,1b PW:1a,2a \\
B1699 & 1a<o1L 1b<o2 / 2a<1a1L,1b2 2b<1b0L & IL/IN & +{-}/+{-} & BF & C:1a,1b PW:1a,2a \\
B1700 & 1a<o1L 1b<o2 / 2a<1a1L,1b2 2b<1b1 & IL/ID & +{-}/+{-} & BF & C:1a,1b PW:1a,2a \\
B1701 & 1a<o1L 1b<o2 / 2a<1a1L,1b2 2b<1b1L & IL/ID & +{-}/+{-} & BF & C:1a,1b PW:1a,2a \\
B1702 & 1a<o1L 1b<o2 / 2a<1a1L,1b2 2b<1b2 & IL/ID & +{-}/+{-} & BF & C:1a,1b,2a PW:1a,2a \\
B1703 & 1a<o1L 1b<o2 / 2a<1a1L,1b2 2b<1b2L & IL/ID & +{-}/+{-} & BF & C:1a,1b,2a PW:1a,2a \\
B1704 & 1a<o1L 1b<o2 / 2a<1a2 2b<1a2,1b0L & IL/EE & +{-}/{-}{-} & BE & C:1a,1b PW:1a \\
B1705 & 1a<o1L 1b<o2 / 2a<1a2 2b<1a2,1b1 & IL/EE & +{-}/{-}{-} & BE & C:1a,1b PW:1a \\
B1706 & 1a<o1L 1b<o2 / 2a<1a2 2b<1a2,1b1L & IL/EE & +{-}/{-}{-} & BE & C:1a,1b PW:1a \\
B1707 & 1a<o1L 1b<o2 / 2a<1a2 2b<1a2,1b2 & IL/EE & +{-}/{-}{-} & BE & C:1a,1b PW:1a \\
B1708 & 1a<o1L 1b<o2 / 2a<1a2 2b<1a2,1b2L & IL/EE & +{-}/{-}{-} & BE & C:1a,1b PW:1a \\
B1709 & 1a<o1L 1b<o2 / 2a<1a2 2b<1a2L,1b1 & IL/LI & +{-}/{-}+ & B & C:1a,1b,2b PW:1a,2b \\
B1710 & 1a<o1L 1b<o2 / 2a<1a2 2b<1a2L,1b2 & IL/LI & +{-}/{-}+ & B & C:1a,1b,2b PW:1a,2b \\
B1711 & 1a<o1L 1b<o2 / 2a<1a2,1b0L 2b<1a2L & IL/LI & +{-}/{-}+ & B & C:1a,1b,2b PW:1a,2b \\
B1712 & 1a<o1L 1b<o2 / 2a<1a2,1b0L 2b<1b1 & IL/ID & +{-}/{-}{-} & B & C:1a,1b,2a,2b PW:1a,2a \\
B1713 & 1a<o1L 1b<o2 / 2a<1a2,1b0L 2b<1b2 & IL/ID & +{-}/{-}{-} & B & C:1a,1b,2a,2b PW:1a,2a \\
B1714 & 1a<o1L 1b<o2 / 2a<1a2,1b1 2b<1a2L & IL/LI & +{-}/{-}+ & B & C:1a,1b,2b PW:1a,2b \\
B1715 & 1a<o1L 1b<o2 / 2a<1a2,1b1 2b<1b0L & IL/EN & +{-}/{-}{-} & BF & C:1a,1b PW:1a \\
B1716 & 1a<o1L 1b<o2 / 2a<1a2,1b1 2b<1b1 & IL/ED & +{-}/{-}{-} & BF & C:1a,1b PW:1a \\
B1717 & 1a<o1L 1b<o2 / 2a<1a2,1b1 2b<1b1L & IL/ED & +{-}/{-}{-} & BF & C:1a,1b PW:1a \\
B1718 & 1a<o1L 1b<o2 / 2a<1a2,1b1 2b<1b2 & IL/ED & +{-}/{-}{-} & BF & C:1a,1b,2b PW:1a \\
B1719 & 1a<o1L 1b<o2 / 2a<1a2,1b1 2b<1b2L & IL/ED & +{-}/{-}{-} & BF & C:1a,1b,2b PW:1a \\
B1720 & 1a<o1L 1b<o2 / 2a<1a2,1b1L 2b<1a2L & IL/LI & +{-}/{-}+ & B & C:1a,1b,2b PW:1a,2b \\
B1721 & 1a<o1L 1b<o2 / 2a<1a2,1b1L 2b<1b1 & IL/ED & +{-}/{-}{-} & BF & C:1a,1b PW:1a \\
B1722 & 1a<o1L 1b<o2 / 2a<1a2,1b1L 2b<1b2 & IL/ED & +{-}/{-}{-} & BF & C:1a,1b,2b PW:1a \\
B1723 & 1a<o1L 1b<o2 / 2a<1a2,1b2 2b<1a2L & IL/LI & +{-}/{-}+ & B & C:1a,1b,2b PW:1a,2b \\
B1724 & 1a<o1L 1b<o2 / 2a<1a2,1b2 2b<1b0L & IL/EN & +{-}/{-}{-} & BF & C:1a,1b PW:1a \\
B1725 & 1a<o1L 1b<o2 / 2a<1a2,1b2 2b<1b1 & IL/ED & +{-}/{-}{-} & BF & C:1a,1b PW:1a \\
B1726 & 1a<o1L 1b<o2 / 2a<1a2,1b2 2b<1b1L & IL/ED & +{-}/{-}{-} & BF & C:1a,1b PW:1a \\
B1727 & 1a<o1L 1b<o2 / 2a<1a2,1b2 2b<1b2 & IL/ED & +{-}/{-}{-} & BF & C:1a,1b PW:1a \\
B1728 & 1a<o1L 1b<o2 / 2a<1a2,1b2 2b<1b2L & IL/ED & +{-}/{-}{-} & BF & C:1a,1b PW:1a \\
B1729 & 1a<o1L 1b<o2 / 2a<1a2,1b2L 2b<1a2L & IL/LI & +{-}/{-}+ & B & C:1a,1b,2b PW:1a,2b \\
B1730 & 1a<o1L 1b<o2 / 2a<1a2,1b2L 2b<1b1 & IL/ED & +{-}/{-}{-} & BF & C:1a,1b PW:1a \\
B1731 & 1a<o1L 1b<o2 / 2a<1a2,1b2L 2b<1b2 & IL/ED & +{-}/{-}{-} & BF & C:1a,1b PW:1a \\
B1732 & 1a<o1L 1b<o2 / 2a<1a2L,1b1 2b<1b0L & IL/IN & +{-}/+{-} & BF & C:1a,1b,2a PW:1a,2a \\
B1733 & 1a<o1L 1b<o2 / 2a<1a2L,1b1 2b<1b1 & IL/ID & +{-}/+{-} & BF & C:1a,1b,2a PW:1a,2a \\
B1734 & 1a<o1L 1b<o2 / 2a<1a2L,1b1 2b<1b1L & IL/ID & +{-}/+{-} & BF & C:1a,1b,2a PW:1a,2a \\
B1735 & 1a<o1L 1b<o2 / 2a<1a2L,1b1 2b<1b2 & IL/ID & +{-}/+{-} & BF & C:1a,1b,2a,2b PW:1a,2a \\
B1736 & 1a<o1L 1b<o2 / 2a<1a2L,1b1 2b<1b2L & IL/ID & +{-}/+{-} & BF & C:1a,1b,2a,2b PW:1a,2a \\
B1737 & 1a<o1L 1b<o2 / 2a<1a2L,1b2 2b<1b0L & IL/IN & +{-}/+{-} & BF & C:1a,1b,2a PW:1a,2a \\
B1738 & 1a<o1L 1b<o2 / 2a<1a2L,1b2 2b<1b1 & IL/ID & +{-}/+{-} & BF & C:1a,1b,2a PW:1a,2a \\
B1739 & 1a<o1L 1b<o2 / 2a<1a2L,1b2 2b<1b1L & IL/ID & +{-}/+{-} & BF & C:1a,1b,2a PW:1a,2a \\
B1740 & 1a<o1L 1b<o2 / 2a<1a2L,1b2 2b<1b2 & IL/ID & +{-}/+{-} & BF & C:1a,1b,2a PW:1a,2a \\
B1741 & 1a<o1L 1b<o2 / 2a<1a2L,1b2 2b<1b2L & IL/ID & +{-}/+{-} & BF & C:1a,1b,2a PW:1a,2a \\
B1742 & 1a<o2 1b<o2 / 2a<1a0L 2b<1a1,1b0L & EE/ND & {-}{-}/{-}{-} & E &  \\
B1743 & 1a<o2 1b<o2 / 2a<1a0L 2b<1a1,1b1 & EE/ND & {-}{-}/{-}{-} & E &  \\
B1744 & 1a<o2 1b<o2 / 2a<1a0L 2b<1a1,1b1L & EE/ND & {-}{-}/{-}{-} & E &  \\
B1745 & 1a<o2 1b<o2 / 2a<1a0L 2b<1a1,1b2 & EE/ND & {-}{-}/{-}{-} & E &  \\
B1746 & 1a<o2 1b<o2 / 2a<1a0L 2b<1a1,1b2L & EE/ND & {-}{-}/{-}{-} & E &  \\
B1747 & 1a<o2 1b<o2 / 2a<1a0L 2b<1a2,1b0L & EE/ND & {-}{-}/{-}{-} & E &  \\
B1748 & 1a<o2 1b<o2 / 2a<1a0L 2b<1a2,1b1 & EE/ND & {-}{-}/{-}{-} & E &  \\
B1749 & 1a<o2 1b<o2 / 2a<1a0L 2b<1a2,1b1L & EE/ND & {-}{-}/{-}{-} & E &  \\
B1750 & 1a<o2 1b<o2 / 2a<1a0L 2b<1a2,1b2 & EE/ND & {-}{-}/{-}{-} & E &  \\
B1751 & 1a<o2 1b<o2 / 2a<1a0L 2b<1a2,1b2L & EE/ND & {-}{-}/{-}{-} & E &  \\
B1752 & 1a<o2 1b<o2 / 2a<1a0L,1b1 2b<1a1 & EE/DD & {-}{-}/{-}{-} & E &  \\
B1753 & 1a<o2 1b<o2 / 2a<1a0L,1b1 2b<1a2 & EE/DD & {-}{-}/{-}{-} & E &  \\
B1754 & 1a<o2 1b<o2 / 2a<1a0L,1b1 2b<1b1 & EE/DD & {-}{-}/{-}{-} & E &  \\
B1755 & 1a<o2 1b<o2 / 2a<1a0L,1b1 2b<1b1L & EE/DD & {-}{-}/{-}{-} & E &  \\
B1756 & 1a<o2 1b<o2 / 2a<1a0L,1b1 2b<1b2 & EE/DD & {-}{-}/{-}{-} & E &  \\
B1757 & 1a<o2 1b<o2 / 2a<1a0L,1b1 2b<1b2L & EE/DD & {-}{-}/{-}{-} & E &  \\
B1758 & 1a<o2 1b<o2 / 2a<1a0L,1b2 2b<1a1 & EE/DD & {-}{-}/{-}{-} & E &  \\
B1759 & 1a<o2 1b<o2 / 2a<1a0L,1b2 2b<1a2 & EE/DD & {-}{-}/{-}{-} & E &  \\
B1760 & 1a<o2 1b<o2 / 2a<1a0L,1b2 2b<1b1 & EE/DD & {-}{-}/{-}{-} & E &  \\
B1761 & 1a<o2 1b<o2 / 2a<1a0L,1b2 2b<1b1L & EE/DD & {-}{-}/{-}{-} & E &  \\
B1762 & 1a<o2 1b<o2 / 2a<1a0L,1b2 2b<1b2 & EE/DD & {-}{-}/{-}{-} & E &  \\
B1763 & 1a<o2 1b<o2 / 2a<1a0L,1b2 2b<1b2L & EE/DD & {-}{-}/{-}{-} & E &  \\
B1764 & 1a<o2 1b<o2 / 2a<1a1 2b<1a1,1b1 & EE/DD & {-}{-}/{-}{-} & E &  \\
B1765 & 1a<o2 1b<o2 / 2a<1a1 2b<1a1,1b1L & EE/DD & {-}{-}/{-}{-} & E &  \\
B1766 & 1a<o2 1b<o2 / 2a<1a1 2b<1a1,1b2 & EE/DD & {-}{-}/{-}{-} & E &  \\
B1767 & 1a<o2 1b<o2 / 2a<1a1 2b<1a1,1b2L & EE/DD & {-}{-}/{-}{-} & E &  \\
B1768 & 1a<o2 1b<o2 / 2a<1a1 2b<1a1L,1b1 & EE/DD & {-}{-}/{-}{-} & E &  \\
B1769 & 1a<o2 1b<o2 / 2a<1a1 2b<1a1L,1b2 & EE/DD & {-}{-}/{-}{-} & E &  \\
B1770 & 1a<o2 1b<o2 / 2a<1a1 2b<1a2,1b1 & EE/DD & {-}{-}/{-}{-} & E &  \\
B1771 & 1a<o2 1b<o2 / 2a<1a1 2b<1a2,1b1L & EE/DD & {-}{-}/{-}{-} & E &  \\
B1772 & 1a<o2 1b<o2 / 2a<1a1 2b<1a2,1b2 & EE/DD & {-}{-}/{-}{-} & E &  \\
B1773 & 1a<o2 1b<o2 / 2a<1a1 2b<1a2,1b2L & EE/DD & {-}{-}/{-}{-} & E &  \\
B1774 & 1a<o2 1b<o2 / 2a<1a1 2b<1a2L,1b1 & EE/DD & {-}{-}/{-}{-} & E &  \\
B1775 & 1a<o2 1b<o2 / 2a<1a1 2b<1a2L,1b2 & EE/DD & {-}{-}/{-}{-} & E &  \\
B1776 & 1a<o2 1b<o2 / 2a<1a1,1b1 2b<1a1L & EE/DD & {-}{-}/{-}{-} & E &  \\
B1777 & 1a<o2 1b<o2 / 2a<1a1,1b1 2b<1a2 & EE/DD & {-}{-}/{-}{-} & E &  \\
B1778 & 1a<o2 1b<o2 / 2a<1a1,1b1 2b<1a2L & EE/DD & {-}{-}/{-}{-} & E &  \\
B1779 & 1a<o2 1b<o2 / 2a<1a1,1b1L 2b<1a1L & EE/DD & {-}{-}/{-}{-} & E &  \\
B1780 & 1a<o2 1b<o2 / 2a<1a1,1b1L 2b<1a2 & EE/DD & {-}{-}/{-}{-} & E &  \\
B1781 & 1a<o2 1b<o2 / 2a<1a1,1b1L 2b<1a2L & EE/DD & {-}{-}/{-}{-} & E &  \\
B1782 & 1a<o2 1b<o2 / 2a<1a1,1b1L 2b<1b2 & EE/DD & {-}{-}/{-}{-} & E &  \\
B1783 & 1a<o2 1b<o2 / 2a<1a1,1b2 2b<1a1L & EE/DD & {-}{-}/{-}{-} & E &  \\
B1784 & 1a<o2 1b<o2 / 2a<1a1,1b2 2b<1a2 & EE/DD & {-}{-}/{-}{-} & E &  \\
B1785 & 1a<o2 1b<o2 / 2a<1a1,1b2 2b<1a2L & EE/DD & {-}{-}/{-}{-} & E &  \\
B1786 & 1a<o2 1b<o2 / 2a<1a1,1b2 2b<1b1L & EE/DD & {-}{-}/{-}{-} & E &  \\
B1787 & 1a<o2 1b<o2 / 2a<1a1,1b2 2b<1b2 & EE/DD & {-}{-}/{-}{-} & E &  \\
B1788 & 1a<o2 1b<o2 / 2a<1a1,1b2 2b<1b2L & EE/DD & {-}{-}/{-}{-} & E &  \\
B1789 & 1a<o2 1b<o2 / 2a<1a1,1b2L 2b<1a1L & EE/DD & {-}{-}/{-}{-} & E &  \\
B1790 & 1a<o2 1b<o2 / 2a<1a1,1b2L 2b<1a2 & EE/DD & {-}{-}/{-}{-} & E &  \\
B1791 & 1a<o2 1b<o2 / 2a<1a1,1b2L 2b<1a2L & EE/DD & {-}{-}/{-}{-} & E &  \\
B1792 & 1a<o2 1b<o2 / 2a<1a1,1b2L 2b<1b2 & EE/DD & {-}{-}/{-}{-} & E &  \\
B1793 & 1a<o2 1b<o2 / 2a<1a1L 2b<1a2,1b1L & EE/DD & {-}{-}/{-}{-} & E &  \\
B1794 & 1a<o2 1b<o2 / 2a<1a1L 2b<1a2,1b2 & EE/DD & {-}{-}/{-}{-} & E &  \\
B1795 & 1a<o2 1b<o2 / 2a<1a1L 2b<1a2,1b2L & EE/DD & {-}{-}/{-}{-} & E &  \\
B1796 & 1a<o2 1b<o2 / 2a<1a1L,1b2 2b<1a2 & EE/DD & {-}{-}/{-}{-} & E &  \\
B1797 & 1a<o2 1b<o2 / 2a<1a1L,1b2 2b<1b2 & EE/DD & {-}{-}/{-}{-} & E &  \\
B1798 & 1a<o2 1b<o2 / 2a<1a1L,1b2 2b<1b2L & EE/DD & {-}{-}/{-}{-} & E &  \\
B1799 & 1a<o2 1b<o2 / 2a<1a2 2b<1a2,1b2 & EE/DD & {-}{-}/{-}{-} & E &  \\
B1800 & 1a<o2 1b<o2 / 2a<1a2 2b<1a2,1b2L & EE/DD & {-}{-}/{-}{-} & E &  \\
B1801 & 1a<o2 1b<o2 / 2a<1a2 2b<1a2L,1b2 & EE/DD & {-}{-}/{-}{-} & E &  \\
B1802 & 1a<o2 1b<o2 / 2a<1a2,1b2 2b<1a2L & EE/DD & {-}{-}/{-}{-} & E &  \\
B1803 & 1a<o2 1b<o2 / 2a<1a2,1b2L 2b<1a2L & EE/DD & {-}{-}/{-}{-} & E &  \\
B1804 & 1a<o2 1b<o2L / 2a<1a0L 2b<1a1,1b0L & LI/ND & {-}+/{-}+ & B & CPW:1b \\
B1805 & 1a<o2 1b<o2L / 2a<1a0L 2b<1a1,1b1 & LI/NE & {-}+/{-}{-} & BF & CPW:1b \\
B1806 & 1a<o2 1b<o2L / 2a<1a0L 2b<1a1,1b1L & LI/NI & {-}+/{-}+ & BF & CPW:1b,2b \\
B1807 & 1a<o2 1b<o2L / 2a<1a0L 2b<1a1,1b2 & LI/NE & {-}+/{-}{-} & BF & CPW:1b \\
B1808 & 1a<o2 1b<o2L / 2a<1a0L 2b<1a1,1b2L & LI/NI & {-}+/{-}+ & BF & CPW:1b,2b \\
B1809 & 1a<o2 1b<o2L / 2a<1a0L 2b<1a2,1b0L & LI/ND & {-}+/{-}+ & B & CPW:1b \\
B1810 & 1a<o2 1b<o2L / 2a<1a0L 2b<1a2,1b1 & LI/NE & {-}+/{-}{-} & BF & CPW:1b \\
B1811 & 1a<o2 1b<o2L / 2a<1a0L 2b<1a2,1b1L & LI/NI & {-}+/{-}+ & BF & CPW:1b,2b \\
B1812 & 1a<o2 1b<o2L / 2a<1a0L 2b<1a2,1b2 & LI/NE & {-}+/{-}{-} & BF & CPW:1b \\
B1813 & 1a<o2 1b<o2L / 2a<1a0L 2b<1a2,1b2L & LI/NI & {-}+/{-}+ & BF & CPW:1b,2b \\
B1814 & 1a<o2 1b<o2L / 2a<1a0L,1b1 2b<1a1 & LI/ID & {-}+/{-}{-} & B & CPW:1b,2a \\
B1815 & 1a<o2 1b<o2L / 2a<1a0L,1b1 2b<1a2 & LI/ID & {-}+/{-}{-} & B & CPW:1b,2a \\
B1816 & 1a<o2 1b<o2L / 2a<1a0L,1b1 2b<1b0L & LI/IN & {-}+/{-}+ & B & CPW:1b,2a \\
B1817 & 1a<o2 1b<o2L / 2a<1a0L,1b1 2b<1b1 & LI/EE & {-}+/{-}{-} & BE & CPW:1b \\
B1818 & 1a<o2 1b<o2L / 2a<1a0L,1b1 2b<1b1L & LI/LI & {-}+/{-}+ & B & CPW:1b,2b \\
B1819 & 1a<o2 1b<o2L / 2a<1a0L,1b1 2b<1b2 & LI/EE & {-}+/{-}{-} & BE & CPW:1b \\
B1820 & 1a<o2 1b<o2L / 2a<1a0L,1b1 2b<1b2L & LI/LI & {-}+/{-}+ & B & CPW:1b,2b \\
B1821 & 1a<o2 1b<o2L / 2a<1a0L,1b2 2b<1a1 & LI/ID & {-}+/{-}{-} & B & CPW:1b,2a \\
B1822 & 1a<o2 1b<o2L / 2a<1a0L,1b2 2b<1a2 & LI/ID & {-}+/{-}{-} & B & CPW:1b,2a \\
B1823 & 1a<o2 1b<o2L / 2a<1a0L,1b2 2b<1b0L & LI/IN & {-}+/{-}+ & B & CPW:1b,2a \\
B1824 & 1a<o2 1b<o2L / 2a<1a0L,1b2 2b<1b1 & LI/EE & {-}+/{-}{-} & BE & CPW:1b \\
B1825 & 1a<o2 1b<o2L / 2a<1a0L,1b2 2b<1b1L & LI/LI & {-}+/{-}+ & B & CPW:1b,2b \\
B1826 & 1a<o2 1b<o2L / 2a<1a0L,1b2 2b<1b2 & LI/EE & {-}+/{-}{-} & BE & CPW:1b \\
B1827 & 1a<o2 1b<o2L / 2a<1a0L,1b2 2b<1b2L & LI/LI & {-}+/{-}+ & B & CPW:1b,2b \\
B1828 & 1a<o2 1b<o2L / 2a<1a1 2b<1a1,1b0L & LI/DD & {-}+/{-}+ & B & CPW:1b \\
B1829 & 1a<o2 1b<o2L / 2a<1a1 2b<1a1,1b1 & LI/DE & {-}+/{-}{-} & BF & CPW:1b \\
B1830 & 1a<o2 1b<o2L / 2a<1a1 2b<1a1,1b1L & LI/DI & {-}+/{-}+ & BF & CPW:1b,2b \\
B1831 & 1a<o2 1b<o2L / 2a<1a1 2b<1a1,1b2 & LI/DE & {-}+/{-}{-} & BF & CPW:1b \\
B1832 & 1a<o2 1b<o2L / 2a<1a1 2b<1a1,1b2L & LI/DI & {-}+/{-}+ & BF & CPW:1b,2b \\
B1833 & 1a<o2 1b<o2L / 2a<1a1 2b<1a1L,1b1 & LI/DE & {-}+/{-}{-} & BF & CPW:1b \\
B1834 & 1a<o2 1b<o2L / 2a<1a1 2b<1a1L,1b2 & LI/DE & {-}+/{-}{-} & BF & CPW:1b \\
B1835 & 1a<o2 1b<o2L / 2a<1a1 2b<1a2,1b0L & LI/DD & {-}+/{-}+ & B & CPW:1b \\
B1836 & 1a<o2 1b<o2L / 2a<1a1 2b<1a2,1b1 & LI/DE & {-}+/{-}{-} & BF & CPW:1b \\
B1837 & 1a<o2 1b<o2L / 2a<1a1 2b<1a2,1b1L & LI/DI & {-}+/{-}+ & BF & CPW:1b,2b \\
B1838 & 1a<o2 1b<o2L / 2a<1a1 2b<1a2,1b2 & LI/DE & {-}+/{-}{-} & BF & CPW:1b \\
B1839 & 1a<o2 1b<o2L / 2a<1a1 2b<1a2,1b2L & LI/DI & {-}+/{-}+ & BF & CPW:1b,2b \\
B1840 & 1a<o2 1b<o2L / 2a<1a1 2b<1a2L,1b1 & LI/DE & {-}+/{-}{-} & BF & CPW:1b \\
B1841 & 1a<o2 1b<o2L / 2a<1a1 2b<1a2L,1b2 & LI/DE & {-}+/{-}{-} & BF & CPW:1b \\
B1842 & 1a<o2 1b<o2L / 2a<1a1,1b0L 2b<1a1L & LI/DD & {-}+/+{-} & B & CPW:1b \\
B1843 & 1a<o2 1b<o2L / 2a<1a1,1b0L 2b<1a2 & LI/DD & {-}+/+{-} & B & CPW:1b \\
B1844 & 1a<o2 1b<o2L / 2a<1a1,1b0L 2b<1a2L & LI/DD & {-}+/+{-} & B & CPW:1b \\
B1845 & 1a<o2 1b<o2L / 2a<1a1,1b0L 2b<1b1 & LI/DI & {-}+/+{-} & B & CPW:1b,2b \\
B1846 & 1a<o2 1b<o2L / 2a<1a1,1b0L 2b<1b2 & LI/DI & {-}+/+{-} & B & CPW:1b,2b \\
B1847 & 1a<o2 1b<o2L / 2a<1a1,1b1 2b<1a1L & LI/ED & {-}+/{-}{-} & BF & CPW:1b \\
B1848 & 1a<o2 1b<o2L / 2a<1a1,1b1 2b<1a2 & LI/ED & {-}+/{-}{-} & BF & CPW:1b \\
B1849 & 1a<o2 1b<o2L / 2a<1a1,1b1 2b<1a2L & LI/ED & {-}+/{-}{-} & BF & CPW:1b \\
B1850 & 1a<o2 1b<o2L / 2a<1a1,1b1 2b<1b0L & LI/EN & {-}+/{-}+ & BF & CPW:1b \\
B1851 & 1a<o2 1b<o2L / 2a<1a1,1b1 2b<1b1 & LI/EE & {-}+/{-}{-} & BE & CPW:1b \\
B1852 & 1a<o2 1b<o2L / 2a<1a1,1b1 2b<1b1L & LI/LI & {-}+/{-}+ & B & CPW:1b,2b \\
B1853 & 1a<o2 1b<o2L / 2a<1a1,1b1 2b<1b2 & LI/EE & {-}+/{-}{-} & BE & CPW:1b \\
B1854 & 1a<o2 1b<o2L / 2a<1a1,1b1 2b<1b2L & LI/LI & {-}+/{-}+ & B & CPW:1b,2b \\
B1855 & 1a<o2 1b<o2L / 2a<1a1,1b1L 2b<1a1L & LI/ID & {-}+/+{-} & BF & CPW:1b,2a \\
B1856 & 1a<o2 1b<o2L / 2a<1a1,1b1L 2b<1a2 & LI/ID & {-}+/+{-} & BF & CPW:1b,2a \\
B1857 & 1a<o2 1b<o2L / 2a<1a1,1b1L 2b<1a2L & LI/ID & {-}+/+{-} & BF & CPW:1b,2a \\
B1858 & 1a<o2 1b<o2L / 2a<1a1,1b1L 2b<1b1 & LI/IL & {-}+/+{-} & B & CPW:1b,2a \\
B1859 & 1a<o2 1b<o2L / 2a<1a1,1b1L 2b<1b2 & LI/IL & {-}+/+{-} & B & CPW:1b,2a \\
B1860 & 1a<o2 1b<o2L / 2a<1a1,1b2 2b<1a1L & LI/ED & {-}+/{-}{-} & BF & CPW:1b \\
B1861 & 1a<o2 1b<o2L / 2a<1a1,1b2 2b<1a2 & LI/ED & {-}+/{-}{-} & BF & CPW:1b \\
B1862 & 1a<o2 1b<o2L / 2a<1a1,1b2 2b<1a2L & LI/ED & {-}+/{-}{-} & BF & CPW:1b \\
B1863 & 1a<o2 1b<o2L / 2a<1a1,1b2 2b<1b0L & LI/EN & {-}+/{-}+ & BF & CPW:1b \\
B1864 & 1a<o2 1b<o2L / 2a<1a1,1b2 2b<1b1 & LI/EE & {-}+/{-}{-} & BE & CPW:1b \\
B1865 & 1a<o2 1b<o2L / 2a<1a1,1b2 2b<1b1L & LI/LI & {-}+/{-}+ & B & CPW:1b,2b \\
B1866 & 1a<o2 1b<o2L / 2a<1a1,1b2 2b<1b2 & LI/EE & {-}+/{-}{-} & BE & CPW:1b \\
B1867 & 1a<o2 1b<o2L / 2a<1a1,1b2 2b<1b2L & LI/LI & {-}+/{-}+ & B & CPW:1b,2b \\
B1868 & 1a<o2 1b<o2L / 2a<1a1,1b2L 2b<1a1L & LI/ID & {-}+/+{-} & BF & CPW:1b,2a \\
B1869 & 1a<o2 1b<o2L / 2a<1a1,1b2L 2b<1a2 & LI/ID & {-}+/+{-} & BF & CPW:1b,2a \\
B1870 & 1a<o2 1b<o2L / 2a<1a1,1b2L 2b<1a2L & LI/ID & {-}+/+{-} & BF & CPW:1b,2a \\
B1871 & 1a<o2 1b<o2L / 2a<1a1,1b2L 2b<1b1 & LI/IL & {-}+/+{-} & B & CPW:1b,2a \\
B1872 & 1a<o2 1b<o2L / 2a<1a1,1b2L 2b<1b2 & LI/IL & {-}+/+{-} & B & CPW:1b,2a \\
B1873 & 1a<o2 1b<o2L / 2a<1a1L 2b<1a2,1b0L & LI/DD & {-}+/{-}+ & B & CPW:1b \\
B1874 & 1a<o2 1b<o2L / 2a<1a1L 2b<1a2,1b1 & LI/DE & {-}+/{-}{-} & BF & CPW:1b \\
B1875 & 1a<o2 1b<o2L / 2a<1a1L 2b<1a2,1b1L & LI/DI & {-}+/{-}+ & BF & CPW:1b,2b \\
B1876 & 1a<o2 1b<o2L / 2a<1a1L 2b<1a2,1b2 & LI/DE & {-}+/{-}{-} & BF & CPW:1b \\
B1877 & 1a<o2 1b<o2L / 2a<1a1L 2b<1a2,1b2L & LI/DI & {-}+/{-}+ & BF & CPW:1b,2b \\
B1878 & 1a<o2 1b<o2L / 2a<1a1L,1b1 2b<1a2 & LI/ED & {-}+/{-}{-} & BF & CPW:1b \\
B1879 & 1a<o2 1b<o2L / 2a<1a1L,1b1 2b<1b0L & LI/EN & {-}+/{-}+ & BF & CPW:1b \\
B1880 & 1a<o2 1b<o2L / 2a<1a1L,1b1 2b<1b1 & LI/EE & {-}+/{-}{-} & BE & CPW:1b \\
B1881 & 1a<o2 1b<o2L / 2a<1a1L,1b1 2b<1b1L & LI/LI & {-}+/{-}+ & B & CPW:1b,2b \\
B1882 & 1a<o2 1b<o2L / 2a<1a1L,1b1 2b<1b2 & LI/EE & {-}+/{-}{-} & BE & CPW:1b \\
B1883 & 1a<o2 1b<o2L / 2a<1a1L,1b1 2b<1b2L & LI/LI & {-}+/{-}+ & B & CPW:1b,2b \\
B1884 & 1a<o2 1b<o2L / 2a<1a1L,1b2 2b<1a2 & LI/ED & {-}+/{-}{-} & BF & CPW:1b \\
B1885 & 1a<o2 1b<o2L / 2a<1a1L,1b2 2b<1b0L & LI/EN & {-}+/{-}+ & BF & CPW:1b \\
B1886 & 1a<o2 1b<o2L / 2a<1a1L,1b2 2b<1b1 & LI/EE & {-}+/{-}{-} & BE & CPW:1b \\
B1887 & 1a<o2 1b<o2L / 2a<1a1L,1b2 2b<1b1L & LI/LI & {-}+/{-}+ & B & CPW:1b,2b \\
B1888 & 1a<o2 1b<o2L / 2a<1a1L,1b2 2b<1b2 & LI/EE & {-}+/{-}{-} & BE & CPW:1b \\
B1889 & 1a<o2 1b<o2L / 2a<1a1L,1b2 2b<1b2L & LI/LI & {-}+/{-}+ & B & CPW:1b,2b \\
B1890 & 1a<o2 1b<o2L / 2a<1a2 2b<1a2,1b0L & LI/DD & {-}+/{-}+ & B & CPW:1b \\
B1891 & 1a<o2 1b<o2L / 2a<1a2 2b<1a2,1b1 & LI/DE & {-}+/{-}{-} & BF & CPW:1b \\
B1892 & 1a<o2 1b<o2L / 2a<1a2 2b<1a2,1b1L & LI/DI & {-}+/{-}+ & BF & CPW:1b,2b \\
B1893 & 1a<o2 1b<o2L / 2a<1a2 2b<1a2,1b2 & LI/DE & {-}+/{-}{-} & BF & CPW:1b \\
B1894 & 1a<o2 1b<o2L / 2a<1a2 2b<1a2,1b2L & LI/DI & {-}+/{-}+ & BF & CPW:1b,2b \\
B1895 & 1a<o2 1b<o2L / 2a<1a2 2b<1a2L,1b1 & LI/DE & {-}+/{-}{-} & BF & CPW:1b \\
B1896 & 1a<o2 1b<o2L / 2a<1a2 2b<1a2L,1b2 & LI/DE & {-}+/{-}{-} & BF & CPW:1b \\
B1897 & 1a<o2 1b<o2L / 2a<1a2,1b0L 2b<1a2L & LI/DD & {-}+/+{-} & B & CPW:1b \\
B1898 & 1a<o2 1b<o2L / 2a<1a2,1b0L 2b<1b1 & LI/DI & {-}+/+{-} & B & CPW:1b,2b \\
B1899 & 1a<o2 1b<o2L / 2a<1a2,1b0L 2b<1b2 & LI/DI & {-}+/+{-} & B & CPW:1b,2b \\
B1900 & 1a<o2 1b<o2L / 2a<1a2,1b1 2b<1a2L & LI/ED & {-}+/{-}{-} & BF & CPW:1b \\
B1901 & 1a<o2 1b<o2L / 2a<1a2,1b1 2b<1b0L & LI/EN & {-}+/{-}+ & BF & CPW:1b \\
B1902 & 1a<o2 1b<o2L / 2a<1a2,1b1 2b<1b1 & LI/EE & {-}+/{-}{-} & BE & CPW:1b \\
B1903 & 1a<o2 1b<o2L / 2a<1a2,1b1 2b<1b1L & LI/LI & {-}+/{-}+ & B & CPW:1b,2b \\
B1904 & 1a<o2 1b<o2L / 2a<1a2,1b1 2b<1b2 & LI/EE & {-}+/{-}{-} & BE & CPW:1b \\
B1905 & 1a<o2 1b<o2L / 2a<1a2,1b1 2b<1b2L & LI/LI & {-}+/{-}+ & B & CPW:1b,2b \\
B1906 & 1a<o2 1b<o2L / 2a<1a2,1b1L 2b<1a2L & LI/ID & {-}+/+{-} & BF & CPW:1b,2a \\
B1907 & 1a<o2 1b<o2L / 2a<1a2,1b1L 2b<1b1 & LI/IL & {-}+/+{-} & B & CPW:1b,2a \\
B1908 & 1a<o2 1b<o2L / 2a<1a2,1b1L 2b<1b2 & LI/IL & {-}+/+{-} & B & CPW:1b,2a \\
B1909 & 1a<o2 1b<o2L / 2a<1a2,1b2 2b<1a2L & LI/ED & {-}+/{-}{-} & BF & CPW:1b \\
B1910 & 1a<o2 1b<o2L / 2a<1a2,1b2 2b<1b0L & LI/EN & {-}+/{-}+ & BF & CPW:1b \\
B1911 & 1a<o2 1b<o2L / 2a<1a2,1b2 2b<1b1 & LI/EE & {-}+/{-}{-} & BE & CPW:1b \\
B1912 & 1a<o2 1b<o2L / 2a<1a2,1b2 2b<1b1L & LI/LI & {-}+/{-}+ & B & CPW:1b,2b \\
B1913 & 1a<o2 1b<o2L / 2a<1a2,1b2 2b<1b2 & LI/EE & {-}+/{-}{-} & BE & CPW:1b \\
B1914 & 1a<o2 1b<o2L / 2a<1a2,1b2 2b<1b2L & LI/LI & {-}+/{-}+ & B & CPW:1b,2b \\
B1915 & 1a<o2 1b<o2L / 2a<1a2,1b2L 2b<1a2L & LI/ID & {-}+/+{-} & BF & CPW:1b,2a \\
B1916 & 1a<o2 1b<o2L / 2a<1a2,1b2L 2b<1b1 & LI/IL & {-}+/+{-} & B & CPW:1b,2a \\
B1917 & 1a<o2 1b<o2L / 2a<1a2,1b2L 2b<1b2 & LI/IL & {-}+/+{-} & B & CPW:1b,2a \\
B1918 & 1a<o2 1b<o2L / 2a<1a2L,1b1 2b<1b0L & LI/EN & {-}+/{-}+ & BF & CPW:1b \\
B1919 & 1a<o2 1b<o2L / 2a<1a2L,1b1 2b<1b1 & LI/EE & {-}+/{-}{-} & BE & CPW:1b \\
B1920 & 1a<o2 1b<o2L / 2a<1a2L,1b1 2b<1b1L & LI/LI & {-}+/{-}+ & B & CPW:1b,2b \\
B1921 & 1a<o2 1b<o2L / 2a<1a2L,1b1 2b<1b2 & LI/EE & {-}+/{-}{-} & BE & CPW:1b \\
B1922 & 1a<o2 1b<o2L / 2a<1a2L,1b1 2b<1b2L & LI/LI & {-}+/{-}+ & B & CPW:1b,2b \\
B1923 & 1a<o2 1b<o2L / 2a<1a2L,1b2 2b<1b0L & LI/EN & {-}+/{-}+ & BF & CPW:1b \\
B1924 & 1a<o2 1b<o2L / 2a<1a2L,1b2 2b<1b1 & LI/EE & {-}+/{-}{-} & BE & CPW:1b \\
B1925 & 1a<o2 1b<o2L / 2a<1a2L,1b2 2b<1b1L & LI/LI & {-}+/{-}+ & B & CPW:1b,2b \\
B1926 & 1a<o2 1b<o2L / 2a<1a2L,1b2 2b<1b2 & LI/EE & {-}+/{-}{-} & BE & CPW:1b \\
B1927 & 1a<o2 1b<o2L / 2a<1a2L,1b2 2b<1b2L & LI/LI & {-}+/{-}+ & B & CPW:1b,2b \\
B1928 & 1a<o0L 1b<o1 / 2a<1a0L,1b1 2b<1a1,1b0L & NI/IN & +{-}/+{-} & N &  \\
B1929 & 1a<o0L 1b<o1 / 2a<1a0L,1b1 2b<1a1,1b1 & NI/EE & +{-}/+{-} & E & W:1b \\
B1930 & 1a<o0L 1b<o1 / 2a<1a0L,1b1 2b<1a1,1b1L & NI/LI & +{-}/+{-} & B & CP:2b W:1b,2b \\
B1931 & 1a<o0L 1b<o1 / 2a<1a0L,1b1 2b<1a1,1b2 & NI/EE & +{-}/+{-} & E & C:2b W:1b \\
B1932 & 1a<o0L 1b<o1 / 2a<1a0L,1b1 2b<1a1,1b2L & NI/LI & +{-}/+{-} & B & P:2b W:1b,2b \\
B1933 & 1a<o0L 1b<o1 / 2a<1a0L,1b1 2b<1a2,1b0L & NI/IN & +{-}/+{-} & N &  \\
B1934 & 1a<o0L 1b<o1 / 2a<1a0L,1b1 2b<1a2,1b1 & NI/EE & +{-}/+{-} & E & W:1b \\
B1935 & 1a<o0L 1b<o1 / 2a<1a0L,1b1 2b<1a2,1b1L & NI/LI & +{-}/+{-} & B & CP:2b W:1b,2b \\
B1936 & 1a<o0L 1b<o1 / 2a<1a0L,1b1 2b<1a2,1b2 & NI/EE & +{-}/+{-} & E & W:1b \\
B1937 & 1a<o0L 1b<o1 / 2a<1a0L,1b1 2b<1a2,1b2L & NI/LI & +{-}/+{-} & B & CP:2b W:1b,2b \\
B1938 & 1a<o0L 1b<o1 / 2a<1a0L,1b2 2b<1a1,1b0L & NI/IN & +{-}/+{-} & N &  \\
B1939 & 1a<o0L 1b<o1 / 2a<1a0L,1b2 2b<1a1,1b1 & NI/EE & +{-}/+{-} & E & C:2a W:1b \\
B1940 & 1a<o0L 1b<o1 / 2a<1a0L,1b2 2b<1a1,1b1L & NI/LI & +{-}/+{-} & B & C:2a,2b P:2b W:1b,2b \\
B1941 & 1a<o0L 1b<o1 / 2a<1a0L,1b2 2b<1a1,1b2 & NI/EE & +{-}/+{-} & E & W:1b \\
B1942 & 1a<o0L 1b<o1 / 2a<1a0L,1b2 2b<1a1,1b2L & NI/LI & +{-}/+{-} & B & CP:2b W:1b,2b \\
B1943 & 1a<o0L 1b<o1 / 2a<1a0L,1b2 2b<1a2,1b0L & NI/IN & +{-}/+{-} & N &  \\
B1944 & 1a<o0L 1b<o1 / 2a<1a0L,1b2 2b<1a2,1b1 & NI/EE & +{-}/+{-} & E & C:2a W:1b \\
B1945 & 1a<o0L 1b<o1 / 2a<1a0L,1b2 2b<1a2,1b1L & NI/LI & +{-}/+{-} & B & C:2a,2b P:2b W:1b,2b \\
B1946 & 1a<o0L 1b<o1 / 2a<1a0L,1b2 2b<1a2,1b2 & NI/EE & +{-}/+{-} & E & W:1b \\
B1947 & 1a<o0L 1b<o1 / 2a<1a0L,1b2 2b<1a2,1b2L & NI/LI & +{-}/+{-} & B & CP:2b W:1b,2b \\
B1948 & 1a<o0L 1b<o1 / 2a<1a1,1b0L 2b<1a1,1b1 & NI/NE & +{-}/{-}{-} & F & W:2b \\
B1949 & 1a<o0L 1b<o1 / 2a<1a1,1b0L 2b<1a1,1b2 & NI/NE & +{-}/{-}{-} & F & CW:2b \\
B1950 & 1a<o0L 1b<o1 / 2a<1a1,1b0L 2b<1a1L,1b1 & NI/NE & +{-}/{-}+ & F & W:2b \\
B1951 & 1a<o0L 1b<o1 / 2a<1a1,1b0L 2b<1a1L,1b2 & NI/NE & +{-}/{-}+ & F & CW:2b \\
B1952 & 1a<o0L 1b<o1 / 2a<1a1,1b0L 2b<1a2,1b1 & NI/NE & +{-}/{-}{-} & F & W:2b \\
B1953 & 1a<o0L 1b<o1 / 2a<1a1,1b0L 2b<1a2,1b2 & NI/NE & +{-}/{-}{-} & F & W:2b \\
B1954 & 1a<o0L 1b<o1 / 2a<1a1,1b0L 2b<1a2L,1b1 & NI/NE & +{-}/{-}+ & F & W:2b \\
B1955 & 1a<o0L 1b<o1 / 2a<1a1,1b0L 2b<1a2L,1b2 & NI/NE & +{-}/{-}+ & F & W:2b \\
B1956 & 1a<o0L 1b<o1 / 2a<1a1,1b1 2b<1a1,1b1 & NI/EE & +{-}/{-}{-} & E & W:1b \\
B1957 & 1a<o0L 1b<o1 / 2a<1a1,1b1 2b<1a1,1b1L & NI/LI & +{-}/{-}{-} & B & CP:2b W:1b,2b \\
B1958 & 1a<o0L 1b<o1 / 2a<1a1,1b1 2b<1a1,1b2 & NI/EE & +{-}/{-}{-} & E & C:2b W:1b \\
B1959 & 1a<o0L 1b<o1 / 2a<1a1,1b1 2b<1a1,1b2L & NI/LI & +{-}/{-}{-} & B & P:2b W:1b,2b \\
B1960 & 1a<o0L 1b<o1 / 2a<1a1,1b1 2b<1a1L,1b1 & NI/EE & +{-}/{-}+ & E & W:1b \\
B1961 & 1a<o0L 1b<o1 / 2a<1a1,1b1 2b<1a1L,1b2 & NI/EE & +{-}/{-}+ & E & C:2b W:1b \\
B1962 & 1a<o0L 1b<o1 / 2a<1a1,1b1 2b<1a2,1b0L & NI/EN & +{-}/{-}{-} & F & W:2a \\
B1963 & 1a<o0L 1b<o1 / 2a<1a1,1b1 2b<1a2,1b1 & NI/EE & +{-}/{-}{-} & E & W:1b \\
B1964 & 1a<o0L 1b<o1 / 2a<1a1,1b1 2b<1a2,1b1L & NI/LI & +{-}/{-}{-} & B & CP:2b W:1b,2b \\
B1965 & 1a<o0L 1b<o1 / 2a<1a1,1b1 2b<1a2,1b2 & NI/EE & +{-}/{-}{-} & E & W:1b \\
B1966 & 1a<o0L 1b<o1 / 2a<1a1,1b1 2b<1a2,1b2L & NI/LI & +{-}/{-}{-} & B & CP:2b W:1b,2b \\
B1967 & 1a<o0L 1b<o1 / 2a<1a1,1b1 2b<1a2L,1b1 & NI/EE & +{-}/{-}+ & E & W:1b \\
B1968 & 1a<o0L 1b<o1 / 2a<1a1,1b1 2b<1a2L,1b2 & NI/EE & +{-}/{-}+ & E & W:1b \\
B1969 & 1a<o0L 1b<o1 / 2a<1a1,1b1L 2b<1a1,1b2 & NI/IL & +{-}/{-}{-} & B & C:2a,2b P:2a W:1b,2a \\
B1970 & 1a<o0L 1b<o1 / 2a<1a1,1b1L 2b<1a1L,1b1 & NI/IL & +{-}/{-}+ & B & CP:2a W:1b,2a \\
B1971 & 1a<o0L 1b<o1 / 2a<1a1,1b1L 2b<1a1L,1b2 & NI/IL & +{-}/{-}+ & B & C:2a,2b P:2a W:1b,2a \\
B1972 & 1a<o0L 1b<o1 / 2a<1a1,1b1L 2b<1a2,1b1 & NI/IL & +{-}/{-}{-} & B & CP:2a W:1b,2a \\
B1973 & 1a<o0L 1b<o1 / 2a<1a1,1b1L 2b<1a2,1b2 & NI/IL & +{-}/{-}{-} & B & CP:2a W:1b,2a \\
B1974 & 1a<o0L 1b<o1 / 2a<1a1,1b1L 2b<1a2L,1b1 & NI/IL & +{-}/{-}+ & B & CP:2a W:1b,2a \\
B1975 & 1a<o0L 1b<o1 / 2a<1a1,1b1L 2b<1a2L,1b2 & NI/IL & +{-}/{-}+ & B & CP:2a W:1b,2a \\
B1976 & 1a<o0L 1b<o1 / 2a<1a1,1b2 2b<1a1,1b2 & NI/EE & +{-}/{-}{-} & E & W:1b \\
B1977 & 1a<o0L 1b<o1 / 2a<1a1,1b2 2b<1a1,1b2L & NI/LI & +{-}/{-}{-} & B & CP:2b W:1b,2b \\
B1978 & 1a<o0L 1b<o1 / 2a<1a1,1b2 2b<1a1L,1b1 & NI/EE & +{-}/{-}+ & E & C:2a W:1b \\
B1979 & 1a<o0L 1b<o1 / 2a<1a1,1b2 2b<1a1L,1b2 & NI/EE & +{-}/{-}+ & E & W:1b \\
B1980 & 1a<o0L 1b<o1 / 2a<1a1,1b2 2b<1a2,1b0L & NI/EN & +{-}/{-}{-} & F & CW:2a \\
B1981 & 1a<o0L 1b<o1 / 2a<1a1,1b2 2b<1a2,1b1 & NI/EE & +{-}/{-}{-} & E & C:2a W:1b \\
B1982 & 1a<o0L 1b<o1 / 2a<1a1,1b2 2b<1a2,1b1L & NI/LI & +{-}/{-}{-} & B & C:2a,2b P:2b W:1b,2b \\
B1983 & 1a<o0L 1b<o1 / 2a<1a1,1b2 2b<1a2,1b2 & NI/EE & +{-}/{-}{-} & E & W:1b \\
B1984 & 1a<o0L 1b<o1 / 2a<1a1,1b2 2b<1a2,1b2L & NI/LI & +{-}/{-}{-} & B & CP:2b W:1b,2b \\
B1985 & 1a<o0L 1b<o1 / 2a<1a1,1b2 2b<1a2L,1b1 & NI/EE & +{-}/{-}+ & E & C:2a W:1b \\
B1986 & 1a<o0L 1b<o1 / 2a<1a1,1b2 2b<1a2L,1b2 & NI/EE & +{-}/{-}+ & E & W:1b \\
B1987 & 1a<o0L 1b<o1 / 2a<1a1,1b2L 2b<1a1L,1b1 & NI/IL & +{-}/{-}+ & B & P:2a W:1b,2a \\
B1988 & 1a<o0L 1b<o1 / 2a<1a1,1b2L 2b<1a1L,1b2 & NI/IL & +{-}/{-}+ & B & CP:2a W:1b,2a \\
B1989 & 1a<o0L 1b<o1 / 2a<1a1,1b2L 2b<1a2,1b1 & NI/IL & +{-}/{-}{-} & B & P:2a W:1b,2a \\
B1990 & 1a<o0L 1b<o1 / 2a<1a1,1b2L 2b<1a2,1b2 & NI/IL & +{-}/{-}{-} & B & CP:2a W:1b,2a \\
B1991 & 1a<o0L 1b<o1 / 2a<1a1,1b2L 2b<1a2L,1b1 & NI/IL & +{-}/{-}+ & B & P:2a W:1b,2a \\
B1992 & 1a<o0L 1b<o1 / 2a<1a1,1b2L 2b<1a2L,1b2 & NI/IL & +{-}/{-}+ & B & CP:2a W:1b,2a \\
B1993 & 1a<o0L 1b<o1 / 2a<1a1L,1b1 2b<1a2,1b0L & NI/EN & +{-}/+{-} & F & W:2a \\
B1994 & 1a<o0L 1b<o1 / 2a<1a1L,1b1 2b<1a2,1b1 & NI/EE & +{-}/+{-} & E & W:1b \\
B1995 & 1a<o0L 1b<o1 / 2a<1a1L,1b1 2b<1a2,1b1L & NI/LI & +{-}/+{-} & B & CP:2b W:1b,2b \\
B1996 & 1a<o0L 1b<o1 / 2a<1a1L,1b1 2b<1a2,1b2 & NI/EE & +{-}/+{-} & E & W:1b \\
B1997 & 1a<o0L 1b<o1 / 2a<1a1L,1b1 2b<1a2,1b2L & NI/LI & +{-}/+{-} & B & CP:2b W:1b,2b \\
B1998 & 1a<o0L 1b<o1 / 2a<1a1L,1b2 2b<1a2,1b0L & NI/EN & +{-}/+{-} & F & CW:2a \\
B1999 & 1a<o0L 1b<o1 / 2a<1a1L,1b2 2b<1a2,1b1 & NI/EE & +{-}/+{-} & E & C:2a W:1b \\
B2000 & 1a<o0L 1b<o1 / 2a<1a1L,1b2 2b<1a2,1b1L & NI/LI & +{-}/+{-} & B & C:2a,2b P:2b W:1b,2b \\
B2001 & 1a<o0L 1b<o1 / 2a<1a1L,1b2 2b<1a2,1b2 & NI/EE & +{-}/+{-} & E & W:1b \\
B2002 & 1a<o0L 1b<o1 / 2a<1a1L,1b2 2b<1a2,1b2L & NI/LI & +{-}/+{-} & B & CP:2b W:1b,2b \\
B2003 & 1a<o0L 1b<o1 / 2a<1a2,1b0L 2b<1a2,1b1 & NI/NE & +{-}/{-}{-} & F & W:2b \\
B2004 & 1a<o0L 1b<o1 / 2a<1a2,1b0L 2b<1a2,1b2 & NI/NE & +{-}/{-}{-} & F & W:2b \\
B2005 & 1a<o0L 1b<o1 / 2a<1a2,1b0L 2b<1a2L,1b1 & NI/NE & +{-}/{-}+ & F & W:2b \\
B2006 & 1a<o0L 1b<o1 / 2a<1a2,1b0L 2b<1a2L,1b2 & NI/NE & +{-}/{-}+ & F & W:2b \\
B2007 & 1a<o0L 1b<o1 / 2a<1a2,1b1 2b<1a2,1b1 & NI/EE & +{-}/{-}{-} & E & W:1b \\
B2008 & 1a<o0L 1b<o1 / 2a<1a2,1b1 2b<1a2,1b1L & NI/LI & +{-}/{-}{-} & B & CP:2b W:1b,2b \\
B2009 & 1a<o0L 1b<o1 / 2a<1a2,1b1 2b<1a2,1b2 & NI/EE & +{-}/{-}{-} & E & W:1b \\
B2010 & 1a<o0L 1b<o1 / 2a<1a2,1b1 2b<1a2,1b2L & NI/LI & +{-}/{-}{-} & B & CP:2b W:1b,2b \\
B2011 & 1a<o0L 1b<o1 / 2a<1a2,1b1 2b<1a2L,1b1 & NI/EE & +{-}/{-}+ & E & W:1b \\
B2012 & 1a<o0L 1b<o1 / 2a<1a2,1b1 2b<1a2L,1b2 & NI/EE & +{-}/{-}+ & E & W:1b \\
B2013 & 1a<o0L 1b<o1 / 2a<1a2,1b1L 2b<1a2,1b2 & NI/IL & +{-}/{-}{-} & B & CP:2a W:1b,2a \\
B2014 & 1a<o0L 1b<o1 / 2a<1a2,1b1L 2b<1a2L,1b1 & NI/IL & +{-}/{-}+ & B & CP:2a W:1b,2a \\
B2015 & 1a<o0L 1b<o1 / 2a<1a2,1b1L 2b<1a2L,1b2 & NI/IL & +{-}/{-}+ & B & CP:2a W:1b,2a \\
B2016 & 1a<o0L 1b<o1 / 2a<1a2,1b2 2b<1a2,1b2 & NI/EE & +{-}/{-}{-} & E & W:1b \\
B2017 & 1a<o0L 1b<o1 / 2a<1a2,1b2 2b<1a2,1b2L & NI/LI & +{-}/{-}{-} & B & CP:2b W:1b,2b \\
B2018 & 1a<o0L 1b<o1 / 2a<1a2,1b2 2b<1a2L,1b1 & NI/EE & +{-}/{-}+ & E & W:1b \\
B2019 & 1a<o0L 1b<o1 / 2a<1a2,1b2 2b<1a2L,1b2 & NI/EE & +{-}/{-}+ & E & W:1b \\
B2020 & 1a<o0L 1b<o1 / 2a<1a2,1b2L 2b<1a2L,1b1 & NI/IL & +{-}/{-}+ & B & CP:2a W:1b,2a \\
B2021 & 1a<o0L 1b<o1 / 2a<1a2,1b2L 2b<1a2L,1b2 & NI/IL & +{-}/{-}+ & B & CP:2a W:1b,2a \\
B2022 & 1a<o0L 1b<o2 / 2a<1a0L,1b1 2b<1a1,1b0L & NI/IN & +{-}/+{-} & N &  \\
B2023 & 1a<o0L 1b<o2 / 2a<1a0L,1b1 2b<1a1,1b1 & NI/EE & +{-}/+{-} & E & W:1b \\
B2024 & 1a<o0L 1b<o2 / 2a<1a0L,1b1 2b<1a1,1b1L & NI/LI & +{-}/+{-} & B & CP:2b W:1b,2b \\
B2025 & 1a<o0L 1b<o2 / 2a<1a0L,1b1 2b<1a1,1b2 & NI/EE & +{-}/+{-} & E & C:2b W:1b \\
B2026 & 1a<o0L 1b<o2 / 2a<1a0L,1b1 2b<1a1,1b2L & NI/LI & +{-}/+{-} & B & P:2b W:1b,2b \\
B2027 & 1a<o0L 1b<o2 / 2a<1a0L,1b1 2b<1a2,1b0L & NI/IN & +{-}/+{-} & N &  \\
B2028 & 1a<o0L 1b<o2 / 2a<1a0L,1b1 2b<1a2,1b1 & NI/EE & +{-}/+{-} & E & W:1b \\
B2029 & 1a<o0L 1b<o2 / 2a<1a0L,1b1 2b<1a2,1b1L & NI/LI & +{-}/+{-} & B & CP:2b W:1b,2b \\
B2030 & 1a<o0L 1b<o2 / 2a<1a0L,1b1 2b<1a2,1b2 & NI/EE & +{-}/+{-} & E & W:1b \\
B2031 & 1a<o0L 1b<o2 / 2a<1a0L,1b1 2b<1a2,1b2L & NI/LI & +{-}/+{-} & B & CP:2b W:1b,2b \\
B2032 & 1a<o0L 1b<o2 / 2a<1a0L,1b2 2b<1a1,1b0L & NI/IN & +{-}/+{-} & N &  \\
B2033 & 1a<o0L 1b<o2 / 2a<1a0L,1b2 2b<1a1,1b1 & NI/EE & +{-}/+{-} & E & C:2a W:1b \\
B2034 & 1a<o0L 1b<o2 / 2a<1a0L,1b2 2b<1a1,1b1L & NI/LI & +{-}/+{-} & B & C:2a,2b P:2b W:1b,2b \\
B2035 & 1a<o0L 1b<o2 / 2a<1a0L,1b2 2b<1a1,1b2 & NI/EE & +{-}/+{-} & E & W:1b \\
B2036 & 1a<o0L 1b<o2 / 2a<1a0L,1b2 2b<1a1,1b2L & NI/LI & +{-}/+{-} & B & CP:2b W:1b,2b \\
B2037 & 1a<o0L 1b<o2 / 2a<1a0L,1b2 2b<1a2,1b0L & NI/IN & +{-}/+{-} & N &  \\
B2038 & 1a<o0L 1b<o2 / 2a<1a0L,1b2 2b<1a2,1b1 & NI/EE & +{-}/+{-} & E & C:2a W:1b \\
B2039 & 1a<o0L 1b<o2 / 2a<1a0L,1b2 2b<1a2,1b1L & NI/LI & +{-}/+{-} & B & C:2a,2b P:2b W:1b,2b \\
B2040 & 1a<o0L 1b<o2 / 2a<1a0L,1b2 2b<1a2,1b2 & NI/EE & +{-}/+{-} & E & W:1b \\
B2041 & 1a<o0L 1b<o2 / 2a<1a0L,1b2 2b<1a2,1b2L & NI/LI & +{-}/+{-} & B & CP:2b W:1b,2b \\
B2042 & 1a<o0L 1b<o2 / 2a<1a1,1b0L 2b<1a1,1b1 & NI/NE & +{-}/{-}{-} & F & W:2b \\
B2043 & 1a<o0L 1b<o2 / 2a<1a1,1b0L 2b<1a1,1b2 & NI/NE & +{-}/{-}{-} & F & CW:2b \\
B2044 & 1a<o0L 1b<o2 / 2a<1a1,1b0L 2b<1a1L,1b1 & NI/NE & +{-}/{-}+ & F & W:2b \\
B2045 & 1a<o0L 1b<o2 / 2a<1a1,1b0L 2b<1a1L,1b2 & NI/NE & +{-}/{-}+ & F & CW:2b \\
B2046 & 1a<o0L 1b<o2 / 2a<1a1,1b0L 2b<1a2,1b1 & NI/NE & +{-}/{-}{-} & F & W:2b \\
B2047 & 1a<o0L 1b<o2 / 2a<1a1,1b0L 2b<1a2,1b2 & NI/NE & +{-}/{-}{-} & F & W:2b \\
B2048 & 1a<o0L 1b<o2 / 2a<1a1,1b0L 2b<1a2L,1b1 & NI/NE & +{-}/{-}+ & F & W:2b \\
B2049 & 1a<o0L 1b<o2 / 2a<1a1,1b0L 2b<1a2L,1b2 & NI/NE & +{-}/{-}+ & F & W:2b \\
B2050 & 1a<o0L 1b<o2 / 2a<1a1,1b1 2b<1a1,1b1 & NI/EE & +{-}/{-}{-} & E & W:1b \\
B2051 & 1a<o0L 1b<o2 / 2a<1a1,1b1 2b<1a1,1b1L & NI/LI & +{-}/{-}{-} & B & CP:2b W:1b,2b \\
B2052 & 1a<o0L 1b<o2 / 2a<1a1,1b1 2b<1a1,1b2 & NI/EE & +{-}/{-}{-} & E & C:2b W:1b \\
B2053 & 1a<o0L 1b<o2 / 2a<1a1,1b1 2b<1a1,1b2L & NI/LI & +{-}/{-}{-} & B & P:2b W:1b,2b \\
B2054 & 1a<o0L 1b<o2 / 2a<1a1,1b1 2b<1a1L,1b1 & NI/EE & +{-}/{-}+ & E & W:1b \\
B2055 & 1a<o0L 1b<o2 / 2a<1a1,1b1 2b<1a1L,1b2 & NI/EE & +{-}/{-}+ & E & C:2b W:1b \\
B2056 & 1a<o0L 1b<o2 / 2a<1a1,1b1 2b<1a2,1b0L & NI/EN & +{-}/{-}{-} & F & W:2a \\
B2057 & 1a<o0L 1b<o2 / 2a<1a1,1b1 2b<1a2,1b1 & NI/EE & +{-}/{-}{-} & E & W:1b \\
B2058 & 1a<o0L 1b<o2 / 2a<1a1,1b1 2b<1a2,1b1L & NI/LI & +{-}/{-}{-} & B & CP:2b W:1b,2b \\
B2059 & 1a<o0L 1b<o2 / 2a<1a1,1b1 2b<1a2,1b2 & NI/EE & +{-}/{-}{-} & E & W:1b \\
B2060 & 1a<o0L 1b<o2 / 2a<1a1,1b1 2b<1a2,1b2L & NI/LI & +{-}/{-}{-} & B & CP:2b W:1b,2b \\
B2061 & 1a<o0L 1b<o2 / 2a<1a1,1b1 2b<1a2L,1b1 & NI/EE & +{-}/{-}+ & E & W:1b \\
B2062 & 1a<o0L 1b<o2 / 2a<1a1,1b1 2b<1a2L,1b2 & NI/EE & +{-}/{-}+ & E & W:1b \\
B2063 & 1a<o0L 1b<o2 / 2a<1a1,1b1L 2b<1a1,1b2 & NI/IL & +{-}/{-}{-} & B & C:2a,2b P:2a W:1b,2a \\
B2064 & 1a<o0L 1b<o2 / 2a<1a1,1b1L 2b<1a1L,1b1 & NI/IL & +{-}/{-}+ & B & CP:2a W:1b,2a \\
B2065 & 1a<o0L 1b<o2 / 2a<1a1,1b1L 2b<1a1L,1b2 & NI/IL & +{-}/{-}+ & B & C:2a,2b P:2a W:1b,2a \\
B2066 & 1a<o0L 1b<o2 / 2a<1a1,1b1L 2b<1a2,1b1 & NI/IL & +{-}/{-}{-} & B & CP:2a W:1b,2a \\
B2067 & 1a<o0L 1b<o2 / 2a<1a1,1b1L 2b<1a2,1b2 & NI/IL & +{-}/{-}{-} & B & CP:2a W:1b,2a \\
B2068 & 1a<o0L 1b<o2 / 2a<1a1,1b1L 2b<1a2L,1b1 & NI/IL & +{-}/{-}+ & B & CP:2a W:1b,2a \\
B2069 & 1a<o0L 1b<o2 / 2a<1a1,1b1L 2b<1a2L,1b2 & NI/IL & +{-}/{-}+ & B & CP:2a W:1b,2a \\
B2070 & 1a<o0L 1b<o2 / 2a<1a1,1b2 2b<1a1,1b2 & NI/EE & +{-}/{-}{-} & E & W:1b \\
B2071 & 1a<o0L 1b<o2 / 2a<1a1,1b2 2b<1a1,1b2L & NI/LI & +{-}/{-}{-} & B & CP:2b W:1b,2b \\
B2072 & 1a<o0L 1b<o2 / 2a<1a1,1b2 2b<1a1L,1b1 & NI/EE & +{-}/{-}+ & E & C:2a W:1b \\
B2073 & 1a<o0L 1b<o2 / 2a<1a1,1b2 2b<1a1L,1b2 & NI/EE & +{-}/{-}+ & E & W:1b \\
B2074 & 1a<o0L 1b<o2 / 2a<1a1,1b2 2b<1a2,1b0L & NI/EN & +{-}/{-}{-} & F & CW:2a \\
B2075 & 1a<o0L 1b<o2 / 2a<1a1,1b2 2b<1a2,1b1 & NI/EE & +{-}/{-}{-} & E & C:2a W:1b \\
B2076 & 1a<o0L 1b<o2 / 2a<1a1,1b2 2b<1a2,1b1L & NI/LI & +{-}/{-}{-} & B & C:2a,2b P:2b W:1b,2b \\
B2077 & 1a<o0L 1b<o2 / 2a<1a1,1b2 2b<1a2,1b2 & NI/EE & +{-}/{-}{-} & E & W:1b \\
B2078 & 1a<o0L 1b<o2 / 2a<1a1,1b2 2b<1a2,1b2L & NI/LI & +{-}/{-}{-} & B & CP:2b W:1b,2b \\
B2079 & 1a<o0L 1b<o2 / 2a<1a1,1b2 2b<1a2L,1b1 & NI/EE & +{-}/{-}+ & E & C:2a W:1b \\
B2080 & 1a<o0L 1b<o2 / 2a<1a1,1b2 2b<1a2L,1b2 & NI/EE & +{-}/{-}+ & E & W:1b \\
B2081 & 1a<o0L 1b<o2 / 2a<1a1,1b2L 2b<1a1L,1b1 & NI/IL & +{-}/{-}+ & B & P:2a W:1b,2a \\
B2082 & 1a<o0L 1b<o2 / 2a<1a1,1b2L 2b<1a1L,1b2 & NI/IL & +{-}/{-}+ & B & CP:2a W:1b,2a \\
B2083 & 1a<o0L 1b<o2 / 2a<1a1,1b2L 2b<1a2,1b1 & NI/IL & +{-}/{-}{-} & B & P:2a W:1b,2a \\
B2084 & 1a<o0L 1b<o2 / 2a<1a1,1b2L 2b<1a2,1b2 & NI/IL & +{-}/{-}{-} & B & CP:2a W:1b,2a \\
B2085 & 1a<o0L 1b<o2 / 2a<1a1,1b2L 2b<1a2L,1b1 & NI/IL & +{-}/{-}+ & B & P:2a W:1b,2a \\
B2086 & 1a<o0L 1b<o2 / 2a<1a1,1b2L 2b<1a2L,1b2 & NI/IL & +{-}/{-}+ & B & CP:2a W:1b,2a \\
B2087 & 1a<o0L 1b<o2 / 2a<1a1L,1b1 2b<1a2,1b0L & NI/EN & +{-}/+{-} & F & W:2a \\
B2088 & 1a<o0L 1b<o2 / 2a<1a1L,1b1 2b<1a2,1b1 & NI/EE & +{-}/+{-} & E & W:1b \\
B2089 & 1a<o0L 1b<o2 / 2a<1a1L,1b1 2b<1a2,1b1L & NI/LI & +{-}/+{-} & B & CP:2b W:1b,2b \\
B2090 & 1a<o0L 1b<o2 / 2a<1a1L,1b1 2b<1a2,1b2 & NI/EE & +{-}/+{-} & E & W:1b \\
B2091 & 1a<o0L 1b<o2 / 2a<1a1L,1b1 2b<1a2,1b2L & NI/LI & +{-}/+{-} & B & CP:2b W:1b,2b \\
B2092 & 1a<o0L 1b<o2 / 2a<1a1L,1b2 2b<1a2,1b0L & NI/EN & +{-}/+{-} & F & CW:2a \\
B2093 & 1a<o0L 1b<o2 / 2a<1a1L,1b2 2b<1a2,1b1 & NI/EE & +{-}/+{-} & E & C:2a W:1b \\
B2094 & 1a<o0L 1b<o2 / 2a<1a1L,1b2 2b<1a2,1b1L & NI/LI & +{-}/+{-} & B & C:2a,2b P:2b W:1b,2b \\
B2095 & 1a<o0L 1b<o2 / 2a<1a1L,1b2 2b<1a2,1b2 & NI/EE & +{-}/+{-} & E & W:1b \\
B2096 & 1a<o0L 1b<o2 / 2a<1a1L,1b2 2b<1a2,1b2L & NI/LI & +{-}/+{-} & B & CP:2b W:1b,2b \\
B2097 & 1a<o0L 1b<o2 / 2a<1a2,1b0L 2b<1a2,1b1 & NI/NE & +{-}/{-}{-} & F & W:2b \\
B2098 & 1a<o0L 1b<o2 / 2a<1a2,1b0L 2b<1a2,1b2 & NI/NE & +{-}/{-}{-} & F & W:2b \\
B2099 & 1a<o0L 1b<o2 / 2a<1a2,1b0L 2b<1a2L,1b1 & NI/NE & +{-}/{-}+ & F & W:2b \\
B2100 & 1a<o0L 1b<o2 / 2a<1a2,1b0L 2b<1a2L,1b2 & NI/NE & +{-}/{-}+ & F & W:2b \\
B2101 & 1a<o0L 1b<o2 / 2a<1a2,1b1 2b<1a2,1b1 & NI/EE & +{-}/{-}{-} & E & W:1b \\
B2102 & 1a<o0L 1b<o2 / 2a<1a2,1b1 2b<1a2,1b1L & NI/LI & +{-}/{-}{-} & B & CP:2b W:1b,2b \\
B2103 & 1a<o0L 1b<o2 / 2a<1a2,1b1 2b<1a2,1b2 & NI/EE & +{-}/{-}{-} & E & W:1b \\
B2104 & 1a<o0L 1b<o2 / 2a<1a2,1b1 2b<1a2,1b2L & NI/LI & +{-}/{-}{-} & B & CP:2b W:1b,2b \\
B2105 & 1a<o0L 1b<o2 / 2a<1a2,1b1 2b<1a2L,1b1 & NI/EE & +{-}/{-}+ & E & W:1b \\
B2106 & 1a<o0L 1b<o2 / 2a<1a2,1b1 2b<1a2L,1b2 & NI/EE & +{-}/{-}+ & E & W:1b \\
B2107 & 1a<o0L 1b<o2 / 2a<1a2,1b1L 2b<1a2,1b2 & NI/IL & +{-}/{-}{-} & B & CP:2a W:1b,2a \\
B2108 & 1a<o0L 1b<o2 / 2a<1a2,1b1L 2b<1a2L,1b1 & NI/IL & +{-}/{-}+ & B & CP:2a W:1b,2a \\
B2109 & 1a<o0L 1b<o2 / 2a<1a2,1b1L 2b<1a2L,1b2 & NI/IL & +{-}/{-}+ & B & CP:2a W:1b,2a \\
B2110 & 1a<o0L 1b<o2 / 2a<1a2,1b2 2b<1a2,1b2 & NI/EE & +{-}/{-}{-} & E & W:1b \\
B2111 & 1a<o0L 1b<o2 / 2a<1a2,1b2 2b<1a2,1b2L & NI/LI & +{-}/{-}{-} & B & CP:2b W:1b,2b \\
B2112 & 1a<o0L 1b<o2 / 2a<1a2,1b2 2b<1a2L,1b1 & NI/EE & +{-}/{-}+ & E & W:1b \\
B2113 & 1a<o0L 1b<o2 / 2a<1a2,1b2 2b<1a2L,1b2 & NI/EE & +{-}/{-}+ & E & W:1b \\
B2114 & 1a<o0L 1b<o2 / 2a<1a2,1b2L 2b<1a2L,1b1 & NI/IL & +{-}/{-}+ & B & CP:2a W:1b,2a \\
B2115 & 1a<o0L 1b<o2 / 2a<1a2,1b2L 2b<1a2L,1b2 & NI/IL & +{-}/{-}+ & B & CP:2a W:1b,2a \\
B2116 & 1a<o1 1b<o1 / 2a<1a0L,1b1 2b<1a1,1b0L & EE/DD & {-}{-}/{-}{-} & E &  \\
B2117 & 1a<o1 1b<o1 / 2a<1a0L,1b1 2b<1a1,1b1 & EE/DD & {-}{-}/{-}{-} & E &  \\
B2118 & 1a<o1 1b<o1 / 2a<1a0L,1b1 2b<1a1,1b1L & EE/DD & {-}{-}/{-}{-} & E &  \\
B2119 & 1a<o1 1b<o1 / 2a<1a0L,1b1 2b<1a1,1b2 & EE/DD & {-}{-}/{-}{-} & E &  \\
B2120 & 1a<o1 1b<o1 / 2a<1a0L,1b1 2b<1a1,1b2L & EE/DD & {-}{-}/{-}{-} & E &  \\
B2121 & 1a<o1 1b<o1 / 2a<1a0L,1b1 2b<1a2,1b0L & EE/DD & {-}{-}/{-}{-} & E &  \\
B2122 & 1a<o1 1b<o1 / 2a<1a0L,1b1 2b<1a2,1b1 & EE/DD & {-}{-}/{-}{-} & E &  \\
B2123 & 1a<o1 1b<o1 / 2a<1a0L,1b1 2b<1a2,1b1L & EE/DD & {-}{-}/{-}{-} & E &  \\
B2124 & 1a<o1 1b<o1 / 2a<1a0L,1b1 2b<1a2,1b2 & EE/DD & {-}{-}/{-}{-} & E &  \\
B2125 & 1a<o1 1b<o1 / 2a<1a0L,1b1 2b<1a2,1b2L & EE/DD & {-}{-}/{-}{-} & E &  \\
B2126 & 1a<o1 1b<o1 / 2a<1a0L,1b2 2b<1a1,1b1 & EE/DD & {-}{-}/{-}{-} & E &  \\
B2127 & 1a<o1 1b<o1 / 2a<1a0L,1b2 2b<1a1,1b1L & EE/DD & {-}{-}/{-}{-} & E &  \\
B2128 & 1a<o1 1b<o1 / 2a<1a0L,1b2 2b<1a1,1b2 & EE/DD & {-}{-}/{-}{-} & E &  \\
B2129 & 1a<o1 1b<o1 / 2a<1a0L,1b2 2b<1a1,1b2L & EE/DD & {-}{-}/{-}{-} & E &  \\
B2130 & 1a<o1 1b<o1 / 2a<1a0L,1b2 2b<1a2,1b0L & EE/DD & {-}{-}/{-}{-} & E &  \\
B2131 & 1a<o1 1b<o1 / 2a<1a0L,1b2 2b<1a2,1b1 & EE/DD & {-}{-}/{-}{-} & E &  \\
B2132 & 1a<o1 1b<o1 / 2a<1a0L,1b2 2b<1a2,1b1L & EE/DD & {-}{-}/{-}{-} & E &  \\
B2133 & 1a<o1 1b<o1 / 2a<1a0L,1b2 2b<1a2,1b2 & EE/DD & {-}{-}/{-}{-} & E &  \\
B2134 & 1a<o1 1b<o1 / 2a<1a0L,1b2 2b<1a2,1b2L & EE/DD & {-}{-}/{-}{-} & E &  \\
B2135 & 1a<o1 1b<o1 / 2a<1a1,1b1 2b<1a1,1b1 & EE/DD & {-}{-}/{-}{-} & E &  \\
B2136 & 1a<o1 1b<o1 / 2a<1a1,1b1 2b<1a1,1b1L & EE/DD & {-}{-}/{-}{-} & E &  \\
B2137 & 1a<o1 1b<o1 / 2a<1a1,1b1 2b<1a1,1b2 & EE/DD & {-}{-}/{-}{-} & E &  \\
B2138 & 1a<o1 1b<o1 / 2a<1a1,1b1 2b<1a1,1b2L & EE/DD & {-}{-}/{-}{-} & E &  \\
B2139 & 1a<o1 1b<o1 / 2a<1a1,1b1 2b<1a1L,1b2 & EE/DD & {-}{-}/{-}{-} & E &  \\
B2140 & 1a<o1 1b<o1 / 2a<1a1,1b1 2b<1a2,1b2 & EE/DD & {-}{-}/{-}{-} & E &  \\
B2141 & 1a<o1 1b<o1 / 2a<1a1,1b1 2b<1a2,1b2L & EE/DD & {-}{-}/{-}{-} & E &  \\
B2142 & 1a<o1 1b<o1 / 2a<1a1,1b1L 2b<1a1,1b2 & EE/DD & {-}{-}/{-}{-} & E &  \\
B2143 & 1a<o1 1b<o1 / 2a<1a1,1b1L 2b<1a1L,1b1 & EE/DD & {-}{-}/{-}{-} & E &  \\
B2144 & 1a<o1 1b<o1 / 2a<1a1,1b1L 2b<1a1L,1b2 & EE/DD & {-}{-}/{-}{-} & E &  \\
B2145 & 1a<o1 1b<o1 / 2a<1a1,1b1L 2b<1a2,1b1 & EE/DD & {-}{-}/{-}{-} & E &  \\
B2146 & 1a<o1 1b<o1 / 2a<1a1,1b1L 2b<1a2,1b2 & EE/DD & {-}{-}/{-}{-} & E &  \\
B2147 & 1a<o1 1b<o1 / 2a<1a1,1b1L 2b<1a2L,1b1 & EE/DD & {-}{-}/{-}{-} & E &  \\
B2148 & 1a<o1 1b<o1 / 2a<1a1,1b1L 2b<1a2L,1b2 & EE/DD & {-}{-}/{-}{-} & E &  \\
B2149 & 1a<o1 1b<o1 / 2a<1a1,1b2 2b<1a1,1b2 & EE/DD & {-}{-}/{-}{-} & E &  \\
B2150 & 1a<o1 1b<o1 / 2a<1a1,1b2 2b<1a1,1b2L & EE/DD & {-}{-}/{-}{-} & E &  \\
B2151 & 1a<o1 1b<o1 / 2a<1a1,1b2 2b<1a1L,1b2 & EE/DD & {-}{-}/{-}{-} & E &  \\
B2152 & 1a<o1 1b<o1 / 2a<1a1,1b2 2b<1a2,1b1 & EE/DD & {-}{-}/{-}{-} & E &  \\
B2153 & 1a<o1 1b<o1 / 2a<1a1,1b2 2b<1a2,1b1L & EE/DD & {-}{-}/{-}{-} & E &  \\
B2154 & 1a<o1 1b<o1 / 2a<1a1,1b2 2b<1a2,1b2 & EE/DD & {-}{-}/{-}{-} & E &  \\
B2155 & 1a<o1 1b<o1 / 2a<1a1,1b2 2b<1a2,1b2L & EE/DD & {-}{-}/{-}{-} & E &  \\
B2156 & 1a<o1 1b<o1 / 2a<1a1,1b2 2b<1a2L,1b1 & EE/DD & {-}{-}/{-}{-} & E &  \\
B2157 & 1a<o1 1b<o1 / 2a<1a1,1b2 2b<1a2L,1b2 & EE/DD & {-}{-}/{-}{-} & E &  \\
B2158 & 1a<o1 1b<o1 / 2a<1a1,1b2L 2b<1a1L,1b2 & EE/DD & {-}{-}/{-}{-} & E &  \\
B2159 & 1a<o1 1b<o1 / 2a<1a1,1b2L 2b<1a2,1b2 & EE/DD & {-}{-}/{-}{-} & E &  \\
B2160 & 1a<o1 1b<o1 / 2a<1a1,1b2L 2b<1a2L,1b1 & EE/DD & {-}{-}/{-}{-} & E &  \\
B2161 & 1a<o1 1b<o1 / 2a<1a1,1b2L 2b<1a2L,1b2 & EE/DD & {-}{-}/{-}{-} & E &  \\
B2162 & 1a<o1 1b<o1 / 2a<1a1L,1b2 2b<1a2,1b1L & EE/DD & {-}{-}/{-}{-} & E &  \\
B2163 & 1a<o1 1b<o1 / 2a<1a1L,1b2 2b<1a2,1b2 & EE/DD & {-}{-}/{-}{-} & E &  \\
B2164 & 1a<o1 1b<o1 / 2a<1a1L,1b2 2b<1a2,1b2L & EE/DD & {-}{-}/{-}{-} & E &  \\
B2165 & 1a<o1 1b<o1 / 2a<1a2,1b2 2b<1a2,1b2 & EE/DD & {-}{-}/{-}{-} & E &  \\
B2166 & 1a<o1 1b<o1 / 2a<1a2,1b2 2b<1a2,1b2L & EE/DD & {-}{-}/{-}{-} & E &  \\
B2167 & 1a<o1 1b<o1 / 2a<1a2,1b2L 2b<1a2L,1b2 & EE/DD & {-}{-}/{-}{-} & E &  \\
B2168 & 1a<o1 1b<o1L / 2a<1a0L,1b1 2b<1a1,1b0L & LI/ID & {-}+/{-}+ & B & CPW:1b,2a \\
B2169 & 1a<o1 1b<o1L / 2a<1a0L,1b1 2b<1a1,1b1 & LI/EE & {-}+/{-}{-} & BE & CPW:1b \\
B2170 & 1a<o1 1b<o1L / 2a<1a0L,1b1 2b<1a1,1b1L & LI/LI & {-}+/{-}+ & B & CPW:1b,2b \\
B2171 & 1a<o1 1b<o1L / 2a<1a0L,1b1 2b<1a1,1b2 & LI/EE & {-}+/{-}{-} & BE & CPW:1b \\
B2172 & 1a<o1 1b<o1L / 2a<1a0L,1b1 2b<1a1,1b2L & LI/LI & {-}+/{-}+ & B & CPW:1b,2b \\
B2173 & 1a<o1 1b<o1L / 2a<1a0L,1b1 2b<1a2,1b0L & LI/ID & {-}+/{-}+ & B & CPW:1b,2a \\
B2174 & 1a<o1 1b<o1L / 2a<1a0L,1b1 2b<1a2,1b1 & LI/EE & {-}+/{-}{-} & BE & CPW:1b \\
B2175 & 1a<o1 1b<o1L / 2a<1a0L,1b1 2b<1a2,1b1L & LI/LI & {-}+/{-}+ & B & CPW:1b,2b \\
B2176 & 1a<o1 1b<o1L / 2a<1a0L,1b1 2b<1a2,1b2 & LI/EE & {-}+/{-}{-} & BE & CPW:1b \\
B2177 & 1a<o1 1b<o1L / 2a<1a0L,1b1 2b<1a2,1b2L & LI/LI & {-}+/{-}+ & B & CPW:1b,2b \\
B2178 & 1a<o1 1b<o1L / 2a<1a0L,1b2 2b<1a1,1b0L & LI/ID & {-}+/{-}+ & B & CPW:1b,2a \\
B2179 & 1a<o1 1b<o1L / 2a<1a0L,1b2 2b<1a1,1b1 & LI/EE & {-}+/{-}{-} & BE & CPW:1b \\
B2180 & 1a<o1 1b<o1L / 2a<1a0L,1b2 2b<1a1,1b1L & LI/LI & {-}+/{-}+ & B & CPW:1b,2b \\
B2181 & 1a<o1 1b<o1L / 2a<1a0L,1b2 2b<1a1,1b2 & LI/EE & {-}+/{-}{-} & BE & CPW:1b \\
B2182 & 1a<o1 1b<o1L / 2a<1a0L,1b2 2b<1a1,1b2L & LI/LI & {-}+/{-}+ & B & CPW:1b,2b \\
B2183 & 1a<o1 1b<o1L / 2a<1a0L,1b2 2b<1a2,1b0L & LI/ID & {-}+/{-}+ & B & CPW:1b,2a \\
B2184 & 1a<o1 1b<o1L / 2a<1a0L,1b2 2b<1a2,1b1 & LI/EE & {-}+/{-}{-} & BE & CPW:1b \\
B2185 & 1a<o1 1b<o1L / 2a<1a0L,1b2 2b<1a2,1b1L & LI/LI & {-}+/{-}+ & B & CPW:1b,2b \\
B2186 & 1a<o1 1b<o1L / 2a<1a0L,1b2 2b<1a2,1b2 & LI/EE & {-}+/{-}{-} & BE & CPW:1b \\
B2187 & 1a<o1 1b<o1L / 2a<1a0L,1b2 2b<1a2,1b2L & LI/LI & {-}+/{-}+ & B & CPW:1b,2b \\
B2188 & 1a<o1 1b<o1L / 2a<1a1,1b0L 2b<1a1,1b1 & LI/DE & {-}+/+{-} & BF & CPW:1b \\
B2189 & 1a<o1 1b<o1L / 2a<1a1,1b0L 2b<1a1,1b2 & LI/DE & {-}+/+{-} & BF & CPW:1b \\
B2190 & 1a<o1 1b<o1L / 2a<1a1,1b0L 2b<1a1L,1b1 & LI/DE & {-}+/+{-} & BF & CPW:1b \\
B2191 & 1a<o1 1b<o1L / 2a<1a1,1b0L 2b<1a1L,1b2 & LI/DE & {-}+/+{-} & BF & CPW:1b \\
B2192 & 1a<o1 1b<o1L / 2a<1a1,1b0L 2b<1a2,1b1 & LI/DE & {-}+/+{-} & BF & CPW:1b \\
B2193 & 1a<o1 1b<o1L / 2a<1a1,1b0L 2b<1a2,1b2 & LI/DE & {-}+/+{-} & BF & CPW:1b \\
B2194 & 1a<o1 1b<o1L / 2a<1a1,1b0L 2b<1a2L,1b1 & LI/DE & {-}+/+{-} & BF & CPW:1b \\
B2195 & 1a<o1 1b<o1L / 2a<1a1,1b0L 2b<1a2L,1b2 & LI/DE & {-}+/+{-} & BF & CPW:1b \\
B2196 & 1a<o1 1b<o1L / 2a<1a1,1b1 2b<1a1,1b1 & LI/EE & {-}+/{-}{-} & BE & CPW:1b \\
B2197 & 1a<o1 1b<o1L / 2a<1a1,1b1 2b<1a1,1b1L & LI/LI & {-}+/{-}+ & B & CPW:1b,2b \\
B2198 & 1a<o1 1b<o1L / 2a<1a1,1b1 2b<1a1,1b2 & LI/EE & {-}+/{-}{-} & BE & CPW:1b \\
B2199 & 1a<o1 1b<o1L / 2a<1a1,1b1 2b<1a1,1b2L & LI/LI & {-}+/{-}+ & B & CPW:1b,2b \\
B2200 & 1a<o1 1b<o1L / 2a<1a1,1b1 2b<1a1L,1b1 & LI/EE & {-}+/{-}{-} & BE & CPW:1b \\
B2201 & 1a<o1 1b<o1L / 2a<1a1,1b1 2b<1a1L,1b2 & LI/EE & {-}+/{-}{-} & BE & CPW:1b \\
B2202 & 1a<o1 1b<o1L / 2a<1a1,1b1 2b<1a2,1b0L & LI/ED & {-}+/{-}+ & BF & CPW:1b \\
B2203 & 1a<o1 1b<o1L / 2a<1a1,1b1 2b<1a2,1b1 & LI/EE & {-}+/{-}{-} & BE & CPW:1b \\
B2204 & 1a<o1 1b<o1L / 2a<1a1,1b1 2b<1a2,1b1L & LI/LI & {-}+/{-}+ & B & CPW:1b,2b \\
B2205 & 1a<o1 1b<o1L / 2a<1a1,1b1 2b<1a2,1b2 & LI/EE & {-}+/{-}{-} & BE & CPW:1b \\
B2206 & 1a<o1 1b<o1L / 2a<1a1,1b1 2b<1a2,1b2L & LI/LI & {-}+/{-}+ & B & CPW:1b,2b \\
B2207 & 1a<o1 1b<o1L / 2a<1a1,1b1 2b<1a2L,1b1 & LI/EE & {-}+/{-}{-} & BE & CPW:1b \\
B2208 & 1a<o1 1b<o1L / 2a<1a1,1b1 2b<1a2L,1b2 & LI/EE & {-}+/{-}{-} & BE & CPW:1b \\
B2209 & 1a<o1 1b<o1L / 2a<1a1,1b1L 2b<1a1,1b2 & LI/IL & {-}+/+{-} & B & CPW:1b,2a \\
B2210 & 1a<o1 1b<o1L / 2a<1a1,1b1L 2b<1a1L,1b1 & LI/IL & {-}+/+{-} & B & CPW:1b,2a \\
B2211 & 1a<o1 1b<o1L / 2a<1a1,1b1L 2b<1a1L,1b2 & LI/IL & {-}+/+{-} & B & CPW:1b,2a \\
B2212 & 1a<o1 1b<o1L / 2a<1a1,1b1L 2b<1a2,1b1 & LI/IL & {-}+/+{-} & B & CPW:1b,2a \\
B2213 & 1a<o1 1b<o1L / 2a<1a1,1b1L 2b<1a2,1b2 & LI/IL & {-}+/+{-} & B & CPW:1b,2a \\
B2214 & 1a<o1 1b<o1L / 2a<1a1,1b1L 2b<1a2L,1b1 & LI/IL & {-}+/+{-} & B & CPW:1b,2a \\
B2215 & 1a<o1 1b<o1L / 2a<1a1,1b1L 2b<1a2L,1b2 & LI/IL & {-}+/+{-} & B & CPW:1b,2a \\
B2216 & 1a<o1 1b<o1L / 2a<1a1,1b2 2b<1a1,1b2 & LI/EE & {-}+/{-}{-} & BE & CPW:1b \\
B2217 & 1a<o1 1b<o1L / 2a<1a1,1b2 2b<1a1,1b2L & LI/LI & {-}+/{-}+ & B & CPW:1b,2b \\
B2218 & 1a<o1 1b<o1L / 2a<1a1,1b2 2b<1a1L,1b1 & LI/EE & {-}+/{-}{-} & BE & CPW:1b \\
B2219 & 1a<o1 1b<o1L / 2a<1a1,1b2 2b<1a1L,1b2 & LI/EE & {-}+/{-}{-} & BE & CPW:1b \\
B2220 & 1a<o1 1b<o1L / 2a<1a1,1b2 2b<1a2,1b0L & LI/ED & {-}+/{-}+ & BF & CPW:1b \\
B2221 & 1a<o1 1b<o1L / 2a<1a1,1b2 2b<1a2,1b1 & LI/EE & {-}+/{-}{-} & BE & CPW:1b \\
B2222 & 1a<o1 1b<o1L / 2a<1a1,1b2 2b<1a2,1b1L & LI/LI & {-}+/{-}+ & B & CPW:1b,2b \\
B2223 & 1a<o1 1b<o1L / 2a<1a1,1b2 2b<1a2,1b2 & LI/EE & {-}+/{-}{-} & BE & CPW:1b \\
B2224 & 1a<o1 1b<o1L / 2a<1a1,1b2 2b<1a2,1b2L & LI/LI & {-}+/{-}+ & B & CPW:1b,2b \\
B2225 & 1a<o1 1b<o1L / 2a<1a1,1b2 2b<1a2L,1b1 & LI/EE & {-}+/{-}{-} & BE & CPW:1b \\
B2226 & 1a<o1 1b<o1L / 2a<1a1,1b2 2b<1a2L,1b2 & LI/EE & {-}+/{-}{-} & BE & CPW:1b \\
B2227 & 1a<o1 1b<o1L / 2a<1a1,1b2L 2b<1a1L,1b1 & LI/IL & {-}+/+{-} & B & CPW:1b,2a \\
B2228 & 1a<o1 1b<o1L / 2a<1a1,1b2L 2b<1a1L,1b2 & LI/IL & {-}+/+{-} & B & CPW:1b,2a \\
B2229 & 1a<o1 1b<o1L / 2a<1a1,1b2L 2b<1a2,1b1 & LI/IL & {-}+/+{-} & B & CPW:1b,2a \\
B2230 & 1a<o1 1b<o1L / 2a<1a1,1b2L 2b<1a2,1b2 & LI/IL & {-}+/+{-} & B & CPW:1b,2a \\
B2231 & 1a<o1 1b<o1L / 2a<1a1,1b2L 2b<1a2L,1b1 & LI/IL & {-}+/+{-} & B & CPW:1b,2a \\
B2232 & 1a<o1 1b<o1L / 2a<1a1,1b2L 2b<1a2L,1b2 & LI/IL & {-}+/+{-} & B & CPW:1b,2a \\
B2233 & 1a<o1 1b<o1L / 2a<1a1L,1b1 2b<1a2,1b0L & LI/ED & {-}+/{-}+ & BF & CPW:1b \\
B2234 & 1a<o1 1b<o1L / 2a<1a1L,1b1 2b<1a2,1b1 & LI/EE & {-}+/{-}{-} & BE & CPW:1b \\
B2235 & 1a<o1 1b<o1L / 2a<1a1L,1b1 2b<1a2,1b1L & LI/LI & {-}+/{-}+ & B & CPW:1b,2b \\
B2236 & 1a<o1 1b<o1L / 2a<1a1L,1b1 2b<1a2,1b2 & LI/EE & {-}+/{-}{-} & BE & CPW:1b \\
B2237 & 1a<o1 1b<o1L / 2a<1a1L,1b1 2b<1a2,1b2L & LI/LI & {-}+/{-}+ & B & CPW:1b,2b \\
B2238 & 1a<o1 1b<o1L / 2a<1a1L,1b2 2b<1a2,1b0L & LI/ED & {-}+/{-}+ & BF & CPW:1b \\
B2239 & 1a<o1 1b<o1L / 2a<1a1L,1b2 2b<1a2,1b1 & LI/EE & {-}+/{-}{-} & BE & CPW:1b \\
B2240 & 1a<o1 1b<o1L / 2a<1a1L,1b2 2b<1a2,1b1L & LI/LI & {-}+/{-}+ & B & CPW:1b,2b \\
B2241 & 1a<o1 1b<o1L / 2a<1a1L,1b2 2b<1a2,1b2 & LI/EE & {-}+/{-}{-} & BE & CPW:1b \\
B2242 & 1a<o1 1b<o1L / 2a<1a1L,1b2 2b<1a2,1b2L & LI/LI & {-}+/{-}+ & B & CPW:1b,2b \\
B2243 & 1a<o1 1b<o1L / 2a<1a2,1b0L 2b<1a2,1b1 & LI/DE & {-}+/+{-} & BF & CPW:1b \\
B2244 & 1a<o1 1b<o1L / 2a<1a2,1b0L 2b<1a2,1b2 & LI/DE & {-}+/+{-} & BF & CPW:1b \\
B2245 & 1a<o1 1b<o1L / 2a<1a2,1b0L 2b<1a2L,1b1 & LI/DE & {-}+/+{-} & BF & CPW:1b \\
B2246 & 1a<o1 1b<o1L / 2a<1a2,1b0L 2b<1a2L,1b2 & LI/DE & {-}+/+{-} & BF & CPW:1b \\
B2247 & 1a<o1 1b<o1L / 2a<1a2,1b1 2b<1a2,1b1 & LI/EE & {-}+/{-}{-} & BE & CPW:1b \\
B2248 & 1a<o1 1b<o1L / 2a<1a2,1b1 2b<1a2,1b1L & LI/LI & {-}+/{-}+ & B & CPW:1b,2b \\
B2249 & 1a<o1 1b<o1L / 2a<1a2,1b1 2b<1a2,1b2 & LI/EE & {-}+/{-}{-} & BE & CPW:1b \\
B2250 & 1a<o1 1b<o1L / 2a<1a2,1b1 2b<1a2,1b2L & LI/LI & {-}+/{-}+ & B & CPW:1b,2b \\
B2251 & 1a<o1 1b<o1L / 2a<1a2,1b1 2b<1a2L,1b1 & LI/EE & {-}+/{-}{-} & BE & CPW:1b \\
B2252 & 1a<o1 1b<o1L / 2a<1a2,1b1 2b<1a2L,1b2 & LI/EE & {-}+/{-}{-} & BE & CPW:1b \\
B2253 & 1a<o1 1b<o1L / 2a<1a2,1b1L 2b<1a2,1b2 & LI/IL & {-}+/+{-} & B & CPW:1b,2a \\
B2254 & 1a<o1 1b<o1L / 2a<1a2,1b1L 2b<1a2L,1b1 & LI/IL & {-}+/+{-} & B & CPW:1b,2a \\
B2255 & 1a<o1 1b<o1L / 2a<1a2,1b1L 2b<1a2L,1b2 & LI/IL & {-}+/+{-} & B & CPW:1b,2a \\
B2256 & 1a<o1 1b<o1L / 2a<1a2,1b2 2b<1a2,1b2 & LI/EE & {-}+/{-}{-} & BE & CPW:1b \\
B2257 & 1a<o1 1b<o1L / 2a<1a2,1b2 2b<1a2,1b2L & LI/LI & {-}+/{-}+ & B & CPW:1b,2b \\
B2258 & 1a<o1 1b<o1L / 2a<1a2,1b2 2b<1a2L,1b1 & LI/EE & {-}+/{-}{-} & BE & CPW:1b \\
B2259 & 1a<o1 1b<o1L / 2a<1a2,1b2 2b<1a2L,1b2 & LI/EE & {-}+/{-}{-} & BE & CPW:1b \\
B2260 & 1a<o1 1b<o1L / 2a<1a2,1b2L 2b<1a2L,1b1 & LI/IL & {-}+/+{-} & B & CPW:1b,2a \\
B2261 & 1a<o1 1b<o1L / 2a<1a2,1b2L 2b<1a2L,1b2 & LI/IL & {-}+/+{-} & B & CPW:1b,2a \\
B2262 & 1a<o1 1b<o2 / 2a<1a0L,1b1 2b<1a1,1b0L & EE/DD & {-}{-}/{-}{-} & E & C:1b,2a \\
B2263 & 1a<o1 1b<o2 / 2a<1a0L,1b1 2b<1a1,1b1 & EE/DD & {-}{-}/{-}{-} & E & C:1b \\
B2264 & 1a<o1 1b<o2 / 2a<1a0L,1b1 2b<1a1,1b1L & EE/DD & {-}{-}/{-}{-} & E & C:1b \\
B2265 & 1a<o1 1b<o2 / 2a<1a0L,1b1 2b<1a1,1b2 & EE/DD & {-}{-}/{-}{-} & E & C:1b,2b \\
B2266 & 1a<o1 1b<o2 / 2a<1a0L,1b1 2b<1a1,1b2L & EE/DD & {-}{-}/{-}{-} & E & C:1b,2b \\
B2267 & 1a<o1 1b<o2 / 2a<1a0L,1b1 2b<1a2,1b0L & EE/DD & {-}{-}/{-}{-} & E & C:1b,2a \\
B2268 & 1a<o1 1b<o2 / 2a<1a0L,1b1 2b<1a2,1b1 & EE/DD & {-}{-}/{-}{-} & E & C:1b \\
B2269 & 1a<o1 1b<o2 / 2a<1a0L,1b1 2b<1a2,1b1L & EE/DD & {-}{-}/{-}{-} & E & C:1b \\
B2270 & 1a<o1 1b<o2 / 2a<1a0L,1b1 2b<1a2,1b2 & EE/DD & {-}{-}/{-}{-} & E & C:1b \\
B2271 & 1a<o1 1b<o2 / 2a<1a0L,1b1 2b<1a2,1b2L & EE/DD & {-}{-}/{-}{-} & E & C:1b \\
B2272 & 1a<o1 1b<o2 / 2a<1a0L,1b2 2b<1a1,1b0L & EE/DD & {-}{-}/{-}{-} & E & C:1b,2a \\
B2273 & 1a<o1 1b<o2 / 2a<1a0L,1b2 2b<1a1,1b1 & EE/DD & {-}{-}/{-}{-} & E & C:1b,2a \\
B2274 & 1a<o1 1b<o2 / 2a<1a0L,1b2 2b<1a1,1b1L & EE/DD & {-}{-}/{-}{-} & E & C:1b,2a \\
B2275 & 1a<o1 1b<o2 / 2a<1a0L,1b2 2b<1a1,1b2 & EE/DD & {-}{-}/{-}{-} & E & C:1b \\
B2276 & 1a<o1 1b<o2 / 2a<1a0L,1b2 2b<1a1,1b2L & EE/DD & {-}{-}/{-}{-} & E & C:1b \\
B2277 & 1a<o1 1b<o2 / 2a<1a0L,1b2 2b<1a2,1b0L & EE/DD & {-}{-}/{-}{-} & E & C:1b,2a \\
B2278 & 1a<o1 1b<o2 / 2a<1a0L,1b2 2b<1a2,1b1 & EE/DD & {-}{-}/{-}{-} & E & C:1b,2a \\
B2279 & 1a<o1 1b<o2 / 2a<1a0L,1b2 2b<1a2,1b1L & EE/DD & {-}{-}/{-}{-} & E & C:1b,2a \\
B2280 & 1a<o1 1b<o2 / 2a<1a0L,1b2 2b<1a2,1b2 & EE/DD & {-}{-}/{-}{-} & E & C:1b \\
B2281 & 1a<o1 1b<o2 / 2a<1a0L,1b2 2b<1a2,1b2L & EE/DD & {-}{-}/{-}{-} & E & C:1b \\
B2282 & 1a<o1 1b<o2 / 2a<1a1,1b0L 2b<1a1,1b1 & EE/DD & {-}{-}/{-}{-} & E & C:1b \\
B2283 & 1a<o1 1b<o2 / 2a<1a1,1b0L 2b<1a1,1b2 & EE/DD & {-}{-}/{-}{-} & E & C:1b,2b \\
B2284 & 1a<o1 1b<o2 / 2a<1a1,1b0L 2b<1a1L,1b1 & EE/DD & {-}{-}/{-}{-} & E & C:1b \\
B2285 & 1a<o1 1b<o2 / 2a<1a1,1b0L 2b<1a1L,1b2 & EE/DD & {-}{-}/{-}{-} & E & C:1b,2b \\
B2286 & 1a<o1 1b<o2 / 2a<1a1,1b0L 2b<1a2,1b1 & EE/DD & {-}{-}/{-}{-} & E & C:1b \\
B2287 & 1a<o1 1b<o2 / 2a<1a1,1b0L 2b<1a2,1b2 & EE/DD & {-}{-}/{-}{-} & E & C:1b \\
B2288 & 1a<o1 1b<o2 / 2a<1a1,1b0L 2b<1a2L,1b1 & EE/DD & {-}{-}/{-}{-} & E & C:1b \\
B2289 & 1a<o1 1b<o2 / 2a<1a1,1b0L 2b<1a2L,1b2 & EE/DD & {-}{-}/{-}{-} & E & C:1b \\
B2290 & 1a<o1 1b<o2 / 2a<1a1,1b1 2b<1a1,1b1 & EE/DD & {-}{-}/{-}{-} & E & C:1b \\
B2291 & 1a<o1 1b<o2 / 2a<1a1,1b1 2b<1a1,1b1L & EE/DD & {-}{-}/{-}{-} & E & C:1b \\
B2292 & 1a<o1 1b<o2 / 2a<1a1,1b1 2b<1a1,1b2 & EE/DD & {-}{-}/{-}{-} & E & C:1b,2b \\
B2293 & 1a<o1 1b<o2 / 2a<1a1,1b1 2b<1a1,1b2L & EE/DD & {-}{-}/{-}{-} & E & C:1b,2b \\
B2294 & 1a<o1 1b<o2 / 2a<1a1,1b1 2b<1a1L,1b1 & EE/DD & {-}{-}/{-}{-} & E & C:1b \\
B2295 & 1a<o1 1b<o2 / 2a<1a1,1b1 2b<1a1L,1b2 & EE/DD & {-}{-}/{-}{-} & E & C:1b,2b \\
B2296 & 1a<o1 1b<o2 / 2a<1a1,1b1 2b<1a2,1b0L & EE/DD & {-}{-}/{-}{-} & E & C:1b \\
B2297 & 1a<o1 1b<o2 / 2a<1a1,1b1 2b<1a2,1b1 & EE/DD & {-}{-}/{-}{-} & E & C:1b \\
B2298 & 1a<o1 1b<o2 / 2a<1a1,1b1 2b<1a2,1b1L & EE/DD & {-}{-}/{-}{-} & E & C:1b \\
B2299 & 1a<o1 1b<o2 / 2a<1a1,1b1 2b<1a2,1b2 & EE/DD & {-}{-}/{-}{-} & E & C:1b \\
B2300 & 1a<o1 1b<o2 / 2a<1a1,1b1 2b<1a2,1b2L & EE/DD & {-}{-}/{-}{-} & E & C:1b \\
B2301 & 1a<o1 1b<o2 / 2a<1a1,1b1 2b<1a2L,1b1 & EE/DD & {-}{-}/{-}{-} & E & C:1b \\
B2302 & 1a<o1 1b<o2 / 2a<1a1,1b1 2b<1a2L,1b2 & EE/DD & {-}{-}/{-}{-} & E & C:1b \\
B2303 & 1a<o1 1b<o2 / 2a<1a1,1b1L 2b<1a1,1b2 & EE/DD & {-}{-}/{-}{-} & E & C:1b,2b \\
B2304 & 1a<o1 1b<o2 / 2a<1a1,1b1L 2b<1a1L,1b1 & EE/DD & {-}{-}/{-}{-} & E & C:1b \\
B2305 & 1a<o1 1b<o2 / 2a<1a1,1b1L 2b<1a1L,1b2 & EE/DD & {-}{-}/{-}{-} & E & C:1b,2b \\
B2306 & 1a<o1 1b<o2 / 2a<1a1,1b1L 2b<1a2,1b1 & EE/DD & {-}{-}/{-}{-} & E & C:1b \\
B2307 & 1a<o1 1b<o2 / 2a<1a1,1b1L 2b<1a2,1b2 & EE/DD & {-}{-}/{-}{-} & E & C:1b \\
B2308 & 1a<o1 1b<o2 / 2a<1a1,1b1L 2b<1a2L,1b1 & EE/DD & {-}{-}/{-}{-} & E & C:1b \\
B2309 & 1a<o1 1b<o2 / 2a<1a1,1b1L 2b<1a2L,1b2 & EE/DD & {-}{-}/{-}{-} & E & C:1b \\
B2310 & 1a<o1 1b<o2 / 2a<1a1,1b2 2b<1a1,1b2 & EE/DD & {-}{-}/{-}{-} & E & C:1b \\
B2311 & 1a<o1 1b<o2 / 2a<1a1,1b2 2b<1a1,1b2L & EE/DD & {-}{-}/{-}{-} & E & C:1b \\
B2312 & 1a<o1 1b<o2 / 2a<1a1,1b2 2b<1a1L,1b1 & EE/DD & {-}{-}/{-}{-} & E & C:1b,2a \\
B2313 & 1a<o1 1b<o2 / 2a<1a1,1b2 2b<1a1L,1b2 & EE/DD & {-}{-}/{-}{-} & E & C:1b \\
B2314 & 1a<o1 1b<o2 / 2a<1a1,1b2 2b<1a2,1b0L & EE/DD & {-}{-}/{-}{-} & E & C:1b,2a \\
B2315 & 1a<o1 1b<o2 / 2a<1a1,1b2 2b<1a2,1b1 & EE/DD & {-}{-}/{-}{-} & E & C:1b,2a \\
B2316 & 1a<o1 1b<o2 / 2a<1a1,1b2 2b<1a2,1b1L & EE/DD & {-}{-}/{-}{-} & E & C:1b,2a \\
B2317 & 1a<o1 1b<o2 / 2a<1a1,1b2 2b<1a2,1b2 & EE/DD & {-}{-}/{-}{-} & E & C:1b \\
B2318 & 1a<o1 1b<o2 / 2a<1a1,1b2 2b<1a2,1b2L & EE/DD & {-}{-}/{-}{-} & E & C:1b \\
B2319 & 1a<o1 1b<o2 / 2a<1a1,1b2 2b<1a2L,1b1 & EE/DD & {-}{-}/{-}{-} & E & C:1b,2a \\
B2320 & 1a<o1 1b<o2 / 2a<1a1,1b2 2b<1a2L,1b2 & EE/DD & {-}{-}/{-}{-} & E & C:1b \\
B2321 & 1a<o1 1b<o2 / 2a<1a1,1b2L 2b<1a1L,1b1 & EE/DD & {-}{-}/{-}{-} & E & C:1b,2a \\
B2322 & 1a<o1 1b<o2 / 2a<1a1,1b2L 2b<1a1L,1b2 & EE/DD & {-}{-}/{-}{-} & E & C:1b \\
B2323 & 1a<o1 1b<o2 / 2a<1a1,1b2L 2b<1a2,1b1 & EE/DD & {-}{-}/{-}{-} & E & C:1b,2a \\
B2324 & 1a<o1 1b<o2 / 2a<1a1,1b2L 2b<1a2,1b2 & EE/DD & {-}{-}/{-}{-} & E & C:1b \\
B2325 & 1a<o1 1b<o2 / 2a<1a1,1b2L 2b<1a2L,1b1 & EE/DD & {-}{-}/{-}{-} & E & C:1b,2a \\
B2326 & 1a<o1 1b<o2 / 2a<1a1,1b2L 2b<1a2L,1b2 & EE/DD & {-}{-}/{-}{-} & E & C:1b \\
B2327 & 1a<o1 1b<o2 / 2a<1a1L,1b1 2b<1a2,1b0L & EE/DD & {-}{-}/{-}{-} & E & C:1b \\
B2328 & 1a<o1 1b<o2 / 2a<1a1L,1b1 2b<1a2,1b1 & EE/DD & {-}{-}/{-}{-} & E & C:1b \\
B2329 & 1a<o1 1b<o2 / 2a<1a1L,1b1 2b<1a2,1b1L & EE/DD & {-}{-}/{-}{-} & E & C:1b \\
B2330 & 1a<o1 1b<o2 / 2a<1a1L,1b1 2b<1a2,1b2 & EE/DD & {-}{-}/{-}{-} & E & C:1b \\
B2331 & 1a<o1 1b<o2 / 2a<1a1L,1b1 2b<1a2,1b2L & EE/DD & {-}{-}/{-}{-} & E & C:1b \\
B2332 & 1a<o1 1b<o2 / 2a<1a1L,1b2 2b<1a2,1b0L & EE/DD & {-}{-}/{-}{-} & E & C:1b,2a \\
B2333 & 1a<o1 1b<o2 / 2a<1a1L,1b2 2b<1a2,1b1 & EE/DD & {-}{-}/{-}{-} & E & C:1b,2a \\
B2334 & 1a<o1 1b<o2 / 2a<1a1L,1b2 2b<1a2,1b1L & EE/DD & {-}{-}/{-}{-} & E & C:1b,2a \\
B2335 & 1a<o1 1b<o2 / 2a<1a1L,1b2 2b<1a2,1b2 & EE/DD & {-}{-}/{-}{-} & E & C:1b \\
B2336 & 1a<o1 1b<o2 / 2a<1a1L,1b2 2b<1a2,1b2L & EE/DD & {-}{-}/{-}{-} & E & C:1b \\
B2337 & 1a<o1 1b<o2 / 2a<1a2,1b0L 2b<1a2,1b1 & EE/DD & {-}{-}/{-}{-} & E & C:1b \\
B2338 & 1a<o1 1b<o2 / 2a<1a2,1b0L 2b<1a2,1b2 & EE/DD & {-}{-}/{-}{-} & E & C:1b \\
B2339 & 1a<o1 1b<o2 / 2a<1a2,1b0L 2b<1a2L,1b1 & EE/DD & {-}{-}/{-}{-} & E & C:1b \\
B2340 & 1a<o1 1b<o2 / 2a<1a2,1b0L 2b<1a2L,1b2 & EE/DD & {-}{-}/{-}{-} & E & C:1b \\
B2341 & 1a<o1 1b<o2 / 2a<1a2,1b1 2b<1a2,1b1 & EE/DD & {-}{-}/{-}{-} & E & C:1b \\
B2342 & 1a<o1 1b<o2 / 2a<1a2,1b1 2b<1a2,1b1L & EE/DD & {-}{-}/{-}{-} & E & C:1b \\
B2343 & 1a<o1 1b<o2 / 2a<1a2,1b1 2b<1a2,1b2 & EE/DD & {-}{-}/{-}{-} & E & C:1b \\
B2344 & 1a<o1 1b<o2 / 2a<1a2,1b1 2b<1a2,1b2L & EE/DD & {-}{-}/{-}{-} & E & C:1b \\
B2345 & 1a<o1 1b<o2 / 2a<1a2,1b1 2b<1a2L,1b1 & EE/DD & {-}{-}/{-}{-} & E & C:1b \\
B2346 & 1a<o1 1b<o2 / 2a<1a2,1b1 2b<1a2L,1b2 & EE/DD & {-}{-}/{-}{-} & E & C:1b \\
B2347 & 1a<o1 1b<o2 / 2a<1a2,1b1L 2b<1a2,1b2 & EE/DD & {-}{-}/{-}{-} & E & C:1b \\
B2348 & 1a<o1 1b<o2 / 2a<1a2,1b1L 2b<1a2L,1b1 & EE/DD & {-}{-}/{-}{-} & E & C:1b \\
B2349 & 1a<o1 1b<o2 / 2a<1a2,1b1L 2b<1a2L,1b2 & EE/DD & {-}{-}/{-}{-} & E & C:1b \\
B2350 & 1a<o1 1b<o2 / 2a<1a2,1b2 2b<1a2,1b2 & EE/DD & {-}{-}/{-}{-} & E & C:1b \\
B2351 & 1a<o1 1b<o2 / 2a<1a2,1b2 2b<1a2,1b2L & EE/DD & {-}{-}/{-}{-} & E & C:1b \\
B2352 & 1a<o1 1b<o2 / 2a<1a2,1b2 2b<1a2L,1b1 & EE/DD & {-}{-}/{-}{-} & E & C:1b \\
B2353 & 1a<o1 1b<o2 / 2a<1a2,1b2 2b<1a2L,1b2 & EE/DD & {-}{-}/{-}{-} & E & C:1b \\
B2354 & 1a<o1 1b<o2 / 2a<1a2,1b2L 2b<1a2L,1b1 & EE/DD & {-}{-}/{-}{-} & E & C:1b \\
B2355 & 1a<o1 1b<o2 / 2a<1a2,1b2L 2b<1a2L,1b2 & EE/DD & {-}{-}/{-}{-} & E & C:1b \\
B2356 & 1a<o1 1b<o2L / 2a<1a0L,1b1 2b<1a1,1b0L & LI/ID & {-}+/{-}+ & B & PW:1b,2a \\
B2357 & 1a<o1 1b<o2L / 2a<1a0L,1b1 2b<1a1,1b1 & LI/EE & {-}+/{-}{-} & BE & PW:1b \\
B2358 & 1a<o1 1b<o2L / 2a<1a0L,1b1 2b<1a1,1b1L & LI/LI & {-}+/{-}+ & B & C:2b PW:1b,2b \\
B2359 & 1a<o1 1b<o2L / 2a<1a0L,1b1 2b<1a1,1b2 & LI/EE & {-}+/{-}{-} & BE & C:2b PW:1b \\
B2360 & 1a<o1 1b<o2L / 2a<1a0L,1b1 2b<1a1,1b2L & LI/LI & {-}+/{-}+ & B & PW:1b,2b \\
B2361 & 1a<o1 1b<o2L / 2a<1a0L,1b1 2b<1a2,1b0L & LI/ID & {-}+/{-}+ & B & PW:1b,2a \\
B2362 & 1a<o1 1b<o2L / 2a<1a0L,1b1 2b<1a2,1b1 & LI/EE & {-}+/{-}{-} & BE & PW:1b \\
B2363 & 1a<o1 1b<o2L / 2a<1a0L,1b1 2b<1a2,1b1L & LI/LI & {-}+/{-}+ & B & C:2b PW:1b,2b \\
B2364 & 1a<o1 1b<o2L / 2a<1a0L,1b1 2b<1a2,1b2 & LI/EE & {-}+/{-}{-} & BE & PW:1b \\
B2365 & 1a<o1 1b<o2L / 2a<1a0L,1b1 2b<1a2,1b2L & LI/LI & {-}+/{-}+ & B & C:2b PW:1b,2b \\
B2366 & 1a<o1 1b<o2L / 2a<1a0L,1b2 2b<1a1,1b0L & LI/ID & {-}+/{-}+ & B & PW:1b,2a \\
B2367 & 1a<o1 1b<o2L / 2a<1a0L,1b2 2b<1a1,1b1 & LI/EE & {-}+/{-}{-} & BE & C:2a PW:1b \\
B2368 & 1a<o1 1b<o2L / 2a<1a0L,1b2 2b<1a1,1b1L & LI/LI & {-}+/{-}+ & B & C:2a,2b PW:1b,2b \\
B2369 & 1a<o1 1b<o2L / 2a<1a0L,1b2 2b<1a1,1b2 & LI/EE & {-}+/{-}{-} & BE & PW:1b \\
B2370 & 1a<o1 1b<o2L / 2a<1a0L,1b2 2b<1a1,1b2L & LI/LI & {-}+/{-}+ & B & C:2b PW:1b,2b \\
B2371 & 1a<o1 1b<o2L / 2a<1a0L,1b2 2b<1a2,1b0L & LI/ID & {-}+/{-}+ & B & PW:1b,2a \\
B2372 & 1a<o1 1b<o2L / 2a<1a0L,1b2 2b<1a2,1b1 & LI/EE & {-}+/{-}{-} & BE & C:2a PW:1b \\
B2373 & 1a<o1 1b<o2L / 2a<1a0L,1b2 2b<1a2,1b1L & LI/LI & {-}+/{-}+ & B & C:2a,2b PW:1b,2b \\
B2374 & 1a<o1 1b<o2L / 2a<1a0L,1b2 2b<1a2,1b2 & LI/EE & {-}+/{-}{-} & BE & PW:1b \\
B2375 & 1a<o1 1b<o2L / 2a<1a0L,1b2 2b<1a2,1b2L & LI/LI & {-}+/{-}+ & B & C:2b PW:1b,2b \\
B2376 & 1a<o1 1b<o2L / 2a<1a1,1b0L 2b<1a1,1b1 & LI/DE & {-}+/+{-} & BF & PW:1b \\
B2377 & 1a<o1 1b<o2L / 2a<1a1,1b0L 2b<1a1,1b2 & LI/DE & {-}+/+{-} & BF & C:2b PW:1b \\
B2378 & 1a<o1 1b<o2L / 2a<1a1,1b0L 2b<1a1L,1b1 & LI/DE & {-}+/+{-} & BF & PW:1b \\
B2379 & 1a<o1 1b<o2L / 2a<1a1,1b0L 2b<1a1L,1b2 & LI/DE & {-}+/+{-} & BF & C:2b PW:1b \\
B2380 & 1a<o1 1b<o2L / 2a<1a1,1b0L 2b<1a2,1b1 & LI/DE & {-}+/+{-} & BF & PW:1b \\
B2381 & 1a<o1 1b<o2L / 2a<1a1,1b0L 2b<1a2,1b2 & LI/DE & {-}+/+{-} & BF & PW:1b \\
B2382 & 1a<o1 1b<o2L / 2a<1a1,1b0L 2b<1a2L,1b1 & LI/DE & {-}+/+{-} & BF & PW:1b \\
B2383 & 1a<o1 1b<o2L / 2a<1a1,1b0L 2b<1a2L,1b2 & LI/DE & {-}+/+{-} & BF & PW:1b \\
B2384 & 1a<o1 1b<o2L / 2a<1a1,1b1 2b<1a1,1b1 & LI/EE & {-}+/{-}{-} & BE & PW:1b \\
B2385 & 1a<o1 1b<o2L / 2a<1a1,1b1 2b<1a1,1b1L & LI/LI & {-}+/{-}+ & B & C:2b PW:1b,2b \\
B2386 & 1a<o1 1b<o2L / 2a<1a1,1b1 2b<1a1,1b2 & LI/EE & {-}+/{-}{-} & BE & C:2b PW:1b \\
B2387 & 1a<o1 1b<o2L / 2a<1a1,1b1 2b<1a1,1b2L & LI/LI & {-}+/{-}+ & B & PW:1b,2b \\
B2388 & 1a<o1 1b<o2L / 2a<1a1,1b1 2b<1a1L,1b1 & LI/EE & {-}+/{-}{-} & BE & PW:1b \\
B2389 & 1a<o1 1b<o2L / 2a<1a1,1b1 2b<1a1L,1b2 & LI/EE & {-}+/{-}{-} & BE & C:2b PW:1b \\
B2390 & 1a<o1 1b<o2L / 2a<1a1,1b1 2b<1a2,1b0L & LI/ED & {-}+/{-}+ & BF & PW:1b \\
B2391 & 1a<o1 1b<o2L / 2a<1a1,1b1 2b<1a2,1b1 & LI/EE & {-}+/{-}{-} & BE & PW:1b \\
B2392 & 1a<o1 1b<o2L / 2a<1a1,1b1 2b<1a2,1b1L & LI/LI & {-}+/{-}+ & B & C:2b PW:1b,2b \\
B2393 & 1a<o1 1b<o2L / 2a<1a1,1b1 2b<1a2,1b2 & LI/EE & {-}+/{-}{-} & BE & PW:1b \\
B2394 & 1a<o1 1b<o2L / 2a<1a1,1b1 2b<1a2,1b2L & LI/LI & {-}+/{-}+ & B & C:2b PW:1b,2b \\
B2395 & 1a<o1 1b<o2L / 2a<1a1,1b1 2b<1a2L,1b1 & LI/EE & {-}+/{-}{-} & BE & PW:1b \\
B2396 & 1a<o1 1b<o2L / 2a<1a1,1b1 2b<1a2L,1b2 & LI/EE & {-}+/{-}{-} & BE & PW:1b \\
B2397 & 1a<o1 1b<o2L / 2a<1a1,1b1L 2b<1a1,1b2 & LI/IL & {-}+/+{-} & B & C:2a,2b PW:1b,2a \\
B2398 & 1a<o1 1b<o2L / 2a<1a1,1b1L 2b<1a1L,1b1 & LI/IL & {-}+/+{-} & B & C:2a PW:1b,2a \\
B2399 & 1a<o1 1b<o2L / 2a<1a1,1b1L 2b<1a1L,1b2 & LI/IL & {-}+/+{-} & B & C:2a,2b PW:1b,2a \\
B2400 & 1a<o1 1b<o2L / 2a<1a1,1b1L 2b<1a2,1b1 & LI/IL & {-}+/+{-} & B & C:2a PW:1b,2a \\
B2401 & 1a<o1 1b<o2L / 2a<1a1,1b1L 2b<1a2,1b2 & LI/IL & {-}+/+{-} & B & C:2a PW:1b,2a \\
B2402 & 1a<o1 1b<o2L / 2a<1a1,1b1L 2b<1a2L,1b1 & LI/IL & {-}+/+{-} & B & C:2a PW:1b,2a \\
B2403 & 1a<o1 1b<o2L / 2a<1a1,1b1L 2b<1a2L,1b2 & LI/IL & {-}+/+{-} & B & C:2a PW:1b,2a \\
B2404 & 1a<o1 1b<o2L / 2a<1a1,1b2 2b<1a1,1b2 & LI/EE & {-}+/{-}{-} & BE & PW:1b \\
B2405 & 1a<o1 1b<o2L / 2a<1a1,1b2 2b<1a1,1b2L & LI/LI & {-}+/{-}+ & B & C:2b PW:1b,2b \\
B2406 & 1a<o1 1b<o2L / 2a<1a1,1b2 2b<1a1L,1b1 & LI/EE & {-}+/{-}{-} & BE & C:2a PW:1b \\
B2407 & 1a<o1 1b<o2L / 2a<1a1,1b2 2b<1a1L,1b2 & LI/EE & {-}+/{-}{-} & BE & PW:1b \\
B2408 & 1a<o1 1b<o2L / 2a<1a1,1b2 2b<1a2,1b0L & LI/ED & {-}+/{-}+ & BF & C:2a PW:1b \\
B2409 & 1a<o1 1b<o2L / 2a<1a1,1b2 2b<1a2,1b1 & LI/EE & {-}+/{-}{-} & BE & C:2a PW:1b \\
B2410 & 1a<o1 1b<o2L / 2a<1a1,1b2 2b<1a2,1b1L & LI/LI & {-}+/{-}+ & B & C:2a,2b PW:1b,2b \\
B2411 & 1a<o1 1b<o2L / 2a<1a1,1b2 2b<1a2,1b2 & LI/EE & {-}+/{-}{-} & BE & PW:1b \\
B2412 & 1a<o1 1b<o2L / 2a<1a1,1b2 2b<1a2,1b2L & LI/LI & {-}+/{-}+ & B & C:2b PW:1b,2b \\
B2413 & 1a<o1 1b<o2L / 2a<1a1,1b2 2b<1a2L,1b1 & LI/EE & {-}+/{-}{-} & BE & C:2a PW:1b \\
B2414 & 1a<o1 1b<o2L / 2a<1a1,1b2 2b<1a2L,1b2 & LI/EE & {-}+/{-}{-} & BE & PW:1b \\
B2415 & 1a<o1 1b<o2L / 2a<1a1,1b2L 2b<1a1L,1b1 & LI/IL & {-}+/+{-} & B & PW:1b,2a \\
B2416 & 1a<o1 1b<o2L / 2a<1a1,1b2L 2b<1a1L,1b2 & LI/IL & {-}+/+{-} & B & C:2a PW:1b,2a \\
B2417 & 1a<o1 1b<o2L / 2a<1a1,1b2L 2b<1a2,1b1 & LI/IL & {-}+/+{-} & B & PW:1b,2a \\
B2418 & 1a<o1 1b<o2L / 2a<1a1,1b2L 2b<1a2,1b2 & LI/IL & {-}+/+{-} & B & C:2a PW:1b,2a \\
B2419 & 1a<o1 1b<o2L / 2a<1a1,1b2L 2b<1a2L,1b1 & LI/IL & {-}+/+{-} & B & PW:1b,2a \\
B2420 & 1a<o1 1b<o2L / 2a<1a1,1b2L 2b<1a2L,1b2 & LI/IL & {-}+/+{-} & B & C:2a PW:1b,2a \\
B2421 & 1a<o1 1b<o2L / 2a<1a1L,1b1 2b<1a2,1b0L & LI/ED & {-}+/{-}+ & BF & PW:1b \\
B2422 & 1a<o1 1b<o2L / 2a<1a1L,1b1 2b<1a2,1b1 & LI/EE & {-}+/{-}{-} & BE & PW:1b \\
B2423 & 1a<o1 1b<o2L / 2a<1a1L,1b1 2b<1a2,1b1L & LI/LI & {-}+/{-}+ & B & C:2b PW:1b,2b \\
B2424 & 1a<o1 1b<o2L / 2a<1a1L,1b1 2b<1a2,1b2 & LI/EE & {-}+/{-}{-} & BE & PW:1b \\
B2425 & 1a<o1 1b<o2L / 2a<1a1L,1b1 2b<1a2,1b2L & LI/LI & {-}+/{-}+ & B & C:2b PW:1b,2b \\
B2426 & 1a<o1 1b<o2L / 2a<1a1L,1b2 2b<1a2,1b0L & LI/ED & {-}+/{-}+ & BF & C:2a PW:1b \\
B2427 & 1a<o1 1b<o2L / 2a<1a1L,1b2 2b<1a2,1b1 & LI/EE & {-}+/{-}{-} & BE & C:2a PW:1b \\
B2428 & 1a<o1 1b<o2L / 2a<1a1L,1b2 2b<1a2,1b1L & LI/LI & {-}+/{-}+ & B & C:2a,2b PW:1b,2b \\
B2429 & 1a<o1 1b<o2L / 2a<1a1L,1b2 2b<1a2,1b2 & LI/EE & {-}+/{-}{-} & BE & PW:1b \\
B2430 & 1a<o1 1b<o2L / 2a<1a1L,1b2 2b<1a2,1b2L & LI/LI & {-}+/{-}+ & B & C:2b PW:1b,2b \\
B2431 & 1a<o1 1b<o2L / 2a<1a2,1b0L 2b<1a2,1b1 & LI/DE & {-}+/+{-} & BF & PW:1b \\
B2432 & 1a<o1 1b<o2L / 2a<1a2,1b0L 2b<1a2,1b2 & LI/DE & {-}+/+{-} & BF & PW:1b \\
B2433 & 1a<o1 1b<o2L / 2a<1a2,1b0L 2b<1a2L,1b1 & LI/DE & {-}+/+{-} & BF & PW:1b \\
B2434 & 1a<o1 1b<o2L / 2a<1a2,1b0L 2b<1a2L,1b2 & LI/DE & {-}+/+{-} & BF & PW:1b \\
B2435 & 1a<o1 1b<o2L / 2a<1a2,1b1 2b<1a2,1b1 & LI/EE & {-}+/{-}{-} & BE & PW:1b \\
B2436 & 1a<o1 1b<o2L / 2a<1a2,1b1 2b<1a2,1b1L & LI/LI & {-}+/{-}+ & B & C:2b PW:1b,2b \\
B2437 & 1a<o1 1b<o2L / 2a<1a2,1b1 2b<1a2,1b2 & LI/EE & {-}+/{-}{-} & BE & PW:1b \\
B2438 & 1a<o1 1b<o2L / 2a<1a2,1b1 2b<1a2,1b2L & LI/LI & {-}+/{-}+ & B & C:2b PW:1b,2b \\
B2439 & 1a<o1 1b<o2L / 2a<1a2,1b1 2b<1a2L,1b1 & LI/EE & {-}+/{-}{-} & BE & PW:1b \\
B2440 & 1a<o1 1b<o2L / 2a<1a2,1b1 2b<1a2L,1b2 & LI/EE & {-}+/{-}{-} & BE & PW:1b \\
B2441 & 1a<o1 1b<o2L / 2a<1a2,1b1L 2b<1a2,1b2 & LI/IL & {-}+/+{-} & B & C:2a PW:1b,2a \\
B2442 & 1a<o1 1b<o2L / 2a<1a2,1b1L 2b<1a2L,1b1 & LI/IL & {-}+/+{-} & B & C:2a PW:1b,2a \\
B2443 & 1a<o1 1b<o2L / 2a<1a2,1b1L 2b<1a2L,1b2 & LI/IL & {-}+/+{-} & B & C:2a PW:1b,2a \\
B2444 & 1a<o1 1b<o2L / 2a<1a2,1b2 2b<1a2,1b2 & LI/EE & {-}+/{-}{-} & BE & PW:1b \\
B2445 & 1a<o1 1b<o2L / 2a<1a2,1b2 2b<1a2,1b2L & LI/LI & {-}+/{-}+ & B & C:2b PW:1b,2b \\
B2446 & 1a<o1 1b<o2L / 2a<1a2,1b2 2b<1a2L,1b1 & LI/EE & {-}+/{-}{-} & BE & PW:1b \\
B2447 & 1a<o1 1b<o2L / 2a<1a2,1b2 2b<1a2L,1b2 & LI/EE & {-}+/{-}{-} & BE & PW:1b \\
B2448 & 1a<o1 1b<o2L / 2a<1a2,1b2L 2b<1a2L,1b1 & LI/IL & {-}+/+{-} & B & C:2a PW:1b,2a \\
B2449 & 1a<o1 1b<o2L / 2a<1a2,1b2L 2b<1a2L,1b2 & LI/IL & {-}+/+{-} & B & C:2a PW:1b,2a \\
B2450 & 1a<o1L 1b<o2 / 2a<1a0L,1b1 2b<1a1,1b0L & IL/DI & +{-}/+{-} & B & C:1a,1b,2a,2b PW:1a,2b \\
B2451 & 1a<o1L 1b<o2 / 2a<1a0L,1b1 2b<1a1,1b1 & IL/DE & +{-}/+{-} & BF & C:1a,1b PW:1a \\
B2452 & 1a<o1L 1b<o2 / 2a<1a0L,1b1 2b<1a1,1b1L & IL/DE & +{-}/+{-} & BF & C:1a,1b PW:1a \\
B2453 & 1a<o1L 1b<o2 / 2a<1a0L,1b1 2b<1a1,1b2 & IL/DE & +{-}/+{-} & BF & C:1a,1b,2b PW:1a \\
B2454 & 1a<o1L 1b<o2 / 2a<1a0L,1b1 2b<1a1,1b2L & IL/DE & +{-}/+{-} & BF & C:1a,1b,2b PW:1a \\
B2455 & 1a<o1L 1b<o2 / 2a<1a0L,1b1 2b<1a2,1b0L & IL/DI & +{-}/+{-} & B & C:1a,1b,2a,2b PW:1a,2b \\
B2456 & 1a<o1L 1b<o2 / 2a<1a0L,1b1 2b<1a2,1b1 & IL/DE & +{-}/+{-} & BF & C:1a,1b PW:1a \\
B2457 & 1a<o1L 1b<o2 / 2a<1a0L,1b1 2b<1a2,1b1L & IL/DE & +{-}/+{-} & BF & C:1a,1b PW:1a \\
B2458 & 1a<o1L 1b<o2 / 2a<1a0L,1b1 2b<1a2,1b2 & IL/DE & +{-}/+{-} & BF & C:1a,1b PW:1a \\
B2459 & 1a<o1L 1b<o2 / 2a<1a0L,1b1 2b<1a2,1b2L & IL/DE & +{-}/+{-} & BF & C:1a,1b PW:1a \\
B2460 & 1a<o1L 1b<o2 / 2a<1a0L,1b2 2b<1a1,1b0L & IL/DI & +{-}/+{-} & B & C:1a,1b,2a,2b PW:1a,2b \\
B2461 & 1a<o1L 1b<o2 / 2a<1a0L,1b2 2b<1a1,1b1 & IL/DE & +{-}/+{-} & BF & C:1a,1b,2a PW:1a \\
B2462 & 1a<o1L 1b<o2 / 2a<1a0L,1b2 2b<1a1,1b1L & IL/DE & +{-}/+{-} & BF & C:1a,1b,2a PW:1a \\
B2463 & 1a<o1L 1b<o2 / 2a<1a0L,1b2 2b<1a1,1b2 & IL/DE & +{-}/+{-} & BF & C:1a,1b PW:1a \\
B2464 & 1a<o1L 1b<o2 / 2a<1a0L,1b2 2b<1a1,1b2L & IL/DE & +{-}/+{-} & BF & C:1a,1b PW:1a \\
B2465 & 1a<o1L 1b<o2 / 2a<1a0L,1b2 2b<1a2,1b0L & IL/DI & +{-}/+{-} & B & C:1a,1b,2a,2b PW:1a,2b \\
B2466 & 1a<o1L 1b<o2 / 2a<1a0L,1b2 2b<1a2,1b1 & IL/DE & +{-}/+{-} & BF & C:1a,1b,2a PW:1a \\
B2467 & 1a<o1L 1b<o2 / 2a<1a0L,1b2 2b<1a2,1b1L & IL/DE & +{-}/+{-} & BF & C:1a,1b,2a PW:1a \\
B2468 & 1a<o1L 1b<o2 / 2a<1a0L,1b2 2b<1a2,1b2 & IL/DE & +{-}/+{-} & BF & C:1a,1b PW:1a \\
B2469 & 1a<o1L 1b<o2 / 2a<1a0L,1b2 2b<1a2,1b2L & IL/DE & +{-}/+{-} & BF & C:1a,1b PW:1a \\
B2470 & 1a<o1L 1b<o2 / 2a<1a1,1b0L 2b<1a1,1b1 & IL/EE & +{-}/{-}{-} & BE & C:1a,1b PW:1a \\
B2471 & 1a<o1L 1b<o2 / 2a<1a1,1b0L 2b<1a1,1b2 & IL/EE & +{-}/{-}{-} & BE & C:1a,1b,2b PW:1a \\
B2472 & 1a<o1L 1b<o2 / 2a<1a1,1b0L 2b<1a1L,1b1 & IL/LI & +{-}/{-}+ & B & C:1a,1b,2b PW:1a,2b \\
B2473 & 1a<o1L 1b<o2 / 2a<1a1,1b0L 2b<1a1L,1b2 & IL/LI & +{-}/{-}+ & B & C:1a,1b PW:1a,2b \\
B2474 & 1a<o1L 1b<o2 / 2a<1a1,1b0L 2b<1a2,1b1 & IL/EE & +{-}/{-}{-} & BE & C:1a,1b PW:1a \\
B2475 & 1a<o1L 1b<o2 / 2a<1a1,1b0L 2b<1a2,1b2 & IL/EE & +{-}/{-}{-} & BE & C:1a,1b PW:1a \\
B2476 & 1a<o1L 1b<o2 / 2a<1a1,1b0L 2b<1a2L,1b1 & IL/LI & +{-}/{-}+ & B & C:1a,1b,2b PW:1a,2b \\
B2477 & 1a<o1L 1b<o2 / 2a<1a1,1b0L 2b<1a2L,1b2 & IL/LI & +{-}/{-}+ & B & C:1a,1b,2b PW:1a,2b \\
B2478 & 1a<o1L 1b<o2 / 2a<1a1,1b1 2b<1a1,1b1 & IL/EE & +{-}/{-}{-} & BE & C:1a,1b PW:1a \\
B2479 & 1a<o1L 1b<o2 / 2a<1a1,1b1 2b<1a1,1b1L & IL/EE & +{-}/{-}{-} & BE & C:1a,1b PW:1a \\
B2480 & 1a<o1L 1b<o2 / 2a<1a1,1b1 2b<1a1,1b2 & IL/EE & +{-}/{-}{-} & BE & C:1a,1b,2b PW:1a \\
B2481 & 1a<o1L 1b<o2 / 2a<1a1,1b1 2b<1a1,1b2L & IL/EE & +{-}/{-}{-} & BE & C:1a,1b,2b PW:1a \\
B2482 & 1a<o1L 1b<o2 / 2a<1a1,1b1 2b<1a1L,1b1 & IL/LI & +{-}/{-}+ & B & C:1a,1b,2b PW:1a,2b \\
B2483 & 1a<o1L 1b<o2 / 2a<1a1,1b1 2b<1a1L,1b2 & IL/LI & +{-}/{-}+ & B & C:1a,1b PW:1a,2b \\
B2484 & 1a<o1L 1b<o2 / 2a<1a1,1b1 2b<1a2,1b0L & IL/EE & +{-}/{-}{-} & BE & C:1a,1b PW:1a \\
B2485 & 1a<o1L 1b<o2 / 2a<1a1,1b1 2b<1a2,1b1 & IL/EE & +{-}/{-}{-} & BE & C:1a,1b PW:1a \\
B2486 & 1a<o1L 1b<o2 / 2a<1a1,1b1 2b<1a2,1b1L & IL/EE & +{-}/{-}{-} & BE & C:1a,1b PW:1a \\
B2487 & 1a<o1L 1b<o2 / 2a<1a1,1b1 2b<1a2,1b2 & IL/EE & +{-}/{-}{-} & BE & C:1a,1b PW:1a \\
B2488 & 1a<o1L 1b<o2 / 2a<1a1,1b1 2b<1a2,1b2L & IL/EE & +{-}/{-}{-} & BE & C:1a,1b PW:1a \\
B2489 & 1a<o1L 1b<o2 / 2a<1a1,1b1 2b<1a2L,1b1 & IL/LI & +{-}/{-}+ & B & C:1a,1b,2b PW:1a,2b \\
B2490 & 1a<o1L 1b<o2 / 2a<1a1,1b1 2b<1a2L,1b2 & IL/LI & +{-}/{-}+ & B & C:1a,1b,2b PW:1a,2b \\
B2491 & 1a<o1L 1b<o2 / 2a<1a1,1b1L 2b<1a1,1b2 & IL/EE & +{-}/{-}{-} & BE & C:1a,1b,2b PW:1a \\
B2492 & 1a<o1L 1b<o2 / 2a<1a1,1b1L 2b<1a1L,1b1 & IL/LI & +{-}/{-}+ & B & C:1a,1b,2b PW:1a,2b \\
B2493 & 1a<o1L 1b<o2 / 2a<1a1,1b1L 2b<1a1L,1b2 & IL/LI & +{-}/{-}+ & B & C:1a,1b PW:1a,2b \\
B2494 & 1a<o1L 1b<o2 / 2a<1a1,1b1L 2b<1a2,1b1 & IL/EE & +{-}/{-}{-} & BE & C:1a,1b PW:1a \\
B2495 & 1a<o1L 1b<o2 / 2a<1a1,1b1L 2b<1a2,1b2 & IL/EE & +{-}/{-}{-} & BE & C:1a,1b PW:1a \\
B2496 & 1a<o1L 1b<o2 / 2a<1a1,1b1L 2b<1a2L,1b1 & IL/LI & +{-}/{-}+ & B & C:1a,1b,2b PW:1a,2b \\
B2497 & 1a<o1L 1b<o2 / 2a<1a1,1b1L 2b<1a2L,1b2 & IL/LI & +{-}/{-}+ & B & C:1a,1b,2b PW:1a,2b \\
B2498 & 1a<o1L 1b<o2 / 2a<1a1,1b2 2b<1a1,1b2 & IL/EE & +{-}/{-}{-} & BE & C:1a,1b PW:1a \\
B2499 & 1a<o1L 1b<o2 / 2a<1a1,1b2 2b<1a1,1b2L & IL/EE & +{-}/{-}{-} & BE & C:1a,1b PW:1a \\
B2500 & 1a<o1L 1b<o2 / 2a<1a1,1b2 2b<1a1L,1b1 & IL/LI & +{-}/{-}+ & B & C:1a,1b,2a,2b PW:1a,2b \\
B2501 & 1a<o1L 1b<o2 / 2a<1a1,1b2 2b<1a1L,1b2 & IL/LI & +{-}/{-}+ & B & C:1a,1b,2b PW:1a,2b \\
B2502 & 1a<o1L 1b<o2 / 2a<1a1,1b2 2b<1a2,1b0L & IL/EE & +{-}/{-}{-} & BE & C:1a,1b,2a PW:1a \\
B2503 & 1a<o1L 1b<o2 / 2a<1a1,1b2 2b<1a2,1b1 & IL/EE & +{-}/{-}{-} & BE & C:1a,1b,2a PW:1a \\
B2504 & 1a<o1L 1b<o2 / 2a<1a1,1b2 2b<1a2,1b1L & IL/EE & +{-}/{-}{-} & BE & C:1a,1b,2a PW:1a \\
B2505 & 1a<o1L 1b<o2 / 2a<1a1,1b2 2b<1a2,1b2 & IL/EE & +{-}/{-}{-} & BE & C:1a,1b PW:1a \\
B2506 & 1a<o1L 1b<o2 / 2a<1a1,1b2 2b<1a2,1b2L & IL/EE & +{-}/{-}{-} & BE & C:1a,1b PW:1a \\
B2507 & 1a<o1L 1b<o2 / 2a<1a1,1b2 2b<1a2L,1b1 & IL/LI & +{-}/{-}+ & B & C:1a,1b,2a,2b PW:1a,2b \\
B2508 & 1a<o1L 1b<o2 / 2a<1a1,1b2 2b<1a2L,1b2 & IL/LI & +{-}/{-}+ & B & C:1a,1b,2b PW:1a,2b \\
B2509 & 1a<o1L 1b<o2 / 2a<1a1,1b2L 2b<1a1L,1b1 & IL/LI & +{-}/{-}+ & B & C:1a,1b,2a,2b PW:1a,2b \\
B2510 & 1a<o1L 1b<o2 / 2a<1a1,1b2L 2b<1a1L,1b2 & IL/LI & +{-}/{-}+ & B & C:1a,1b,2b PW:1a,2b \\
B2511 & 1a<o1L 1b<o2 / 2a<1a1,1b2L 2b<1a2,1b1 & IL/EE & +{-}/{-}{-} & BE & C:1a,1b,2a PW:1a \\
B2512 & 1a<o1L 1b<o2 / 2a<1a1,1b2L 2b<1a2,1b2 & IL/EE & +{-}/{-}{-} & BE & C:1a,1b PW:1a \\
B2513 & 1a<o1L 1b<o2 / 2a<1a1,1b2L 2b<1a2L,1b1 & IL/LI & +{-}/{-}+ & B & C:1a,1b,2a,2b PW:1a,2b \\
B2514 & 1a<o1L 1b<o2 / 2a<1a1,1b2L 2b<1a2L,1b2 & IL/LI & +{-}/{-}+ & B & C:1a,1b,2b PW:1a,2b \\
B2515 & 1a<o1L 1b<o2 / 2a<1a1L,1b1 2b<1a2,1b0L & IL/IL & +{-}/+{-} & B & C:1a,1b,2a PW:1a,2a \\
B2516 & 1a<o1L 1b<o2 / 2a<1a1L,1b1 2b<1a2,1b1 & IL/IL & +{-}/+{-} & B & C:1a,1b,2a PW:1a,2a \\
B2517 & 1a<o1L 1b<o2 / 2a<1a1L,1b1 2b<1a2,1b1L & IL/IL & +{-}/+{-} & B & C:1a,1b,2a PW:1a,2a \\
B2518 & 1a<o1L 1b<o2 / 2a<1a1L,1b1 2b<1a2,1b2 & IL/IL & +{-}/+{-} & B & C:1a,1b,2a PW:1a,2a \\
B2519 & 1a<o1L 1b<o2 / 2a<1a1L,1b1 2b<1a2,1b2L & IL/IL & +{-}/+{-} & B & C:1a,1b,2a PW:1a,2a \\
B2520 & 1a<o1L 1b<o2 / 2a<1a1L,1b2 2b<1a2,1b0L & IL/IL & +{-}/+{-} & B & C:1a,1b PW:1a,2a \\
B2521 & 1a<o1L 1b<o2 / 2a<1a1L,1b2 2b<1a2,1b1 & IL/IL & +{-}/+{-} & B & C:1a,1b PW:1a,2a \\
B2522 & 1a<o1L 1b<o2 / 2a<1a1L,1b2 2b<1a2,1b1L & IL/IL & +{-}/+{-} & B & C:1a,1b PW:1a,2a \\
B2523 & 1a<o1L 1b<o2 / 2a<1a1L,1b2 2b<1a2,1b2 & IL/IL & +{-}/+{-} & B & C:1a,1b,2a PW:1a,2a \\
B2524 & 1a<o1L 1b<o2 / 2a<1a1L,1b2 2b<1a2,1b2L & IL/IL & +{-}/+{-} & B & C:1a,1b,2a PW:1a,2a \\
B2525 & 1a<o1L 1b<o2 / 2a<1a2,1b0L 2b<1a2,1b1 & IL/EE & +{-}/{-}{-} & BE & C:1a,1b PW:1a \\
B2526 & 1a<o1L 1b<o2 / 2a<1a2,1b0L 2b<1a2,1b2 & IL/EE & +{-}/{-}{-} & BE & C:1a,1b PW:1a \\
B2527 & 1a<o1L 1b<o2 / 2a<1a2,1b0L 2b<1a2L,1b1 & IL/LI & +{-}/{-}+ & B & C:1a,1b,2b PW:1a,2b \\
B2528 & 1a<o1L 1b<o2 / 2a<1a2,1b0L 2b<1a2L,1b2 & IL/LI & +{-}/{-}+ & B & C:1a,1b,2b PW:1a,2b \\
B2529 & 1a<o1L 1b<o2 / 2a<1a2,1b1 2b<1a2,1b1 & IL/EE & +{-}/{-}{-} & BE & C:1a,1b PW:1a \\
B2530 & 1a<o1L 1b<o2 / 2a<1a2,1b1 2b<1a2,1b1L & IL/EE & +{-}/{-}{-} & BE & C:1a,1b PW:1a \\
B2531 & 1a<o1L 1b<o2 / 2a<1a2,1b1 2b<1a2,1b2 & IL/EE & +{-}/{-}{-} & BE & C:1a,1b PW:1a \\
B2532 & 1a<o1L 1b<o2 / 2a<1a2,1b1 2b<1a2,1b2L & IL/EE & +{-}/{-}{-} & BE & C:1a,1b PW:1a \\
B2533 & 1a<o1L 1b<o2 / 2a<1a2,1b1 2b<1a2L,1b1 & IL/LI & +{-}/{-}+ & B & C:1a,1b,2b PW:1a,2b \\
B2534 & 1a<o1L 1b<o2 / 2a<1a2,1b1 2b<1a2L,1b2 & IL/LI & +{-}/{-}+ & B & C:1a,1b,2b PW:1a,2b \\
B2535 & 1a<o1L 1b<o2 / 2a<1a2,1b1L 2b<1a2,1b2 & IL/EE & +{-}/{-}{-} & BE & C:1a,1b PW:1a \\
B2536 & 1a<o1L 1b<o2 / 2a<1a2,1b1L 2b<1a2L,1b1 & IL/LI & +{-}/{-}+ & B & C:1a,1b,2b PW:1a,2b \\
B2537 & 1a<o1L 1b<o2 / 2a<1a2,1b1L 2b<1a2L,1b2 & IL/LI & +{-}/{-}+ & B & C:1a,1b,2b PW:1a,2b \\
B2538 & 1a<o1L 1b<o2 / 2a<1a2,1b2 2b<1a2,1b2 & IL/EE & +{-}/{-}{-} & BE & C:1a,1b PW:1a \\
B2539 & 1a<o1L 1b<o2 / 2a<1a2,1b2 2b<1a2,1b2L & IL/EE & +{-}/{-}{-} & BE & C:1a,1b PW:1a \\
B2540 & 1a<o1L 1b<o2 / 2a<1a2,1b2 2b<1a2L,1b1 & IL/LI & +{-}/{-}+ & B & C:1a,1b,2b PW:1a,2b \\
B2541 & 1a<o1L 1b<o2 / 2a<1a2,1b2 2b<1a2L,1b2 & IL/LI & +{-}/{-}+ & B & C:1a,1b,2b PW:1a,2b \\
B2542 & 1a<o1L 1b<o2 / 2a<1a2,1b2L 2b<1a2L,1b1 & IL/LI & +{-}/{-}+ & B & C:1a,1b,2b PW:1a,2b \\
B2543 & 1a<o1L 1b<o2 / 2a<1a2,1b2L 2b<1a2L,1b2 & IL/LI & +{-}/{-}+ & B & C:1a,1b,2b PW:1a,2b \\
B2544 & 1a<o2 1b<o2 / 2a<1a0L,1b1 2b<1a1,1b0L & EE/DD & {-}{-}/{-}{-} & E &  \\
B2545 & 1a<o2 1b<o2 / 2a<1a0L,1b1 2b<1a1,1b1 & EE/DD & {-}{-}/{-}{-} & E &  \\
B2546 & 1a<o2 1b<o2 / 2a<1a0L,1b1 2b<1a1,1b1L & EE/DD & {-}{-}/{-}{-} & E &  \\
B2547 & 1a<o2 1b<o2 / 2a<1a0L,1b1 2b<1a1,1b2 & EE/DD & {-}{-}/{-}{-} & E &  \\
B2548 & 1a<o2 1b<o2 / 2a<1a0L,1b1 2b<1a1,1b2L & EE/DD & {-}{-}/{-}{-} & E &  \\
B2549 & 1a<o2 1b<o2 / 2a<1a0L,1b1 2b<1a2,1b0L & EE/DD & {-}{-}/{-}{-} & E &  \\
B2550 & 1a<o2 1b<o2 / 2a<1a0L,1b1 2b<1a2,1b1 & EE/DD & {-}{-}/{-}{-} & E &  \\
B2551 & 1a<o2 1b<o2 / 2a<1a0L,1b1 2b<1a2,1b1L & EE/DD & {-}{-}/{-}{-} & E &  \\
B2552 & 1a<o2 1b<o2 / 2a<1a0L,1b1 2b<1a2,1b2 & EE/DD & {-}{-}/{-}{-} & E &  \\
B2553 & 1a<o2 1b<o2 / 2a<1a0L,1b1 2b<1a2,1b2L & EE/DD & {-}{-}/{-}{-} & E &  \\
B2554 & 1a<o2 1b<o2 / 2a<1a0L,1b2 2b<1a1,1b1 & EE/DD & {-}{-}/{-}{-} & E &  \\
B2555 & 1a<o2 1b<o2 / 2a<1a0L,1b2 2b<1a1,1b1L & EE/DD & {-}{-}/{-}{-} & E &  \\
B2556 & 1a<o2 1b<o2 / 2a<1a0L,1b2 2b<1a1,1b2 & EE/DD & {-}{-}/{-}{-} & E &  \\
B2557 & 1a<o2 1b<o2 / 2a<1a0L,1b2 2b<1a1,1b2L & EE/DD & {-}{-}/{-}{-} & E &  \\
B2558 & 1a<o2 1b<o2 / 2a<1a0L,1b2 2b<1a2,1b0L & EE/DD & {-}{-}/{-}{-} & E &  \\
B2559 & 1a<o2 1b<o2 / 2a<1a0L,1b2 2b<1a2,1b1 & EE/DD & {-}{-}/{-}{-} & E &  \\
B2560 & 1a<o2 1b<o2 / 2a<1a0L,1b2 2b<1a2,1b1L & EE/DD & {-}{-}/{-}{-} & E &  \\
B2561 & 1a<o2 1b<o2 / 2a<1a0L,1b2 2b<1a2,1b2 & EE/DD & {-}{-}/{-}{-} & E &  \\
B2562 & 1a<o2 1b<o2 / 2a<1a0L,1b2 2b<1a2,1b2L & EE/DD & {-}{-}/{-}{-} & E &  \\
B2563 & 1a<o2 1b<o2 / 2a<1a1,1b1 2b<1a1,1b1 & EE/DD & {-}{-}/{-}{-} & E &  \\
B2564 & 1a<o2 1b<o2 / 2a<1a1,1b1 2b<1a1,1b1L & EE/DD & {-}{-}/{-}{-} & E &  \\
B2565 & 1a<o2 1b<o2 / 2a<1a1,1b1 2b<1a1,1b2 & EE/DD & {-}{-}/{-}{-} & E &  \\
B2566 & 1a<o2 1b<o2 / 2a<1a1,1b1 2b<1a1,1b2L & EE/DD & {-}{-}/{-}{-} & E &  \\
B2567 & 1a<o2 1b<o2 / 2a<1a1,1b1 2b<1a1L,1b2 & EE/DD & {-}{-}/{-}{-} & E &  \\
B2568 & 1a<o2 1b<o2 / 2a<1a1,1b1 2b<1a2,1b2 & EE/DD & {-}{-}/{-}{-} & E &  \\
B2569 & 1a<o2 1b<o2 / 2a<1a1,1b1 2b<1a2,1b2L & EE/DD & {-}{-}/{-}{-} & E &  \\
B2570 & 1a<o2 1b<o2 / 2a<1a1,1b1L 2b<1a1,1b2 & EE/DD & {-}{-}/{-}{-} & E &  \\
B2571 & 1a<o2 1b<o2 / 2a<1a1,1b1L 2b<1a1L,1b1 & EE/DD & {-}{-}/{-}{-} & E &  \\
B2572 & 1a<o2 1b<o2 / 2a<1a1,1b1L 2b<1a1L,1b2 & EE/DD & {-}{-}/{-}{-} & E &  \\
B2573 & 1a<o2 1b<o2 / 2a<1a1,1b1L 2b<1a2,1b1 & EE/DD & {-}{-}/{-}{-} & E &  \\
B2574 & 1a<o2 1b<o2 / 2a<1a1,1b1L 2b<1a2,1b2 & EE/DD & {-}{-}/{-}{-} & E &  \\
B2575 & 1a<o2 1b<o2 / 2a<1a1,1b1L 2b<1a2L,1b1 & EE/DD & {-}{-}/{-}{-} & E &  \\
B2576 & 1a<o2 1b<o2 / 2a<1a1,1b1L 2b<1a2L,1b2 & EE/DD & {-}{-}/{-}{-} & E &  \\
B2577 & 1a<o2 1b<o2 / 2a<1a1,1b2 2b<1a1,1b2 & EE/DD & {-}{-}/{-}{-} & E &  \\
B2578 & 1a<o2 1b<o2 / 2a<1a1,1b2 2b<1a1,1b2L & EE/DD & {-}{-}/{-}{-} & E &  \\
B2579 & 1a<o2 1b<o2 / 2a<1a1,1b2 2b<1a1L,1b2 & EE/DD & {-}{-}/{-}{-} & E &  \\
B2580 & 1a<o2 1b<o2 / 2a<1a1,1b2 2b<1a2,1b1 & EE/DD & {-}{-}/{-}{-} & E &  \\
B2581 & 1a<o2 1b<o2 / 2a<1a1,1b2 2b<1a2,1b1L & EE/DD & {-}{-}/{-}{-} & E &  \\
B2582 & 1a<o2 1b<o2 / 2a<1a1,1b2 2b<1a2,1b2 & EE/DD & {-}{-}/{-}{-} & E &  \\
B2583 & 1a<o2 1b<o2 / 2a<1a1,1b2 2b<1a2,1b2L & EE/DD & {-}{-}/{-}{-} & E &  \\
B2584 & 1a<o2 1b<o2 / 2a<1a1,1b2 2b<1a2L,1b1 & EE/DD & {-}{-}/{-}{-} & E &  \\
B2585 & 1a<o2 1b<o2 / 2a<1a1,1b2 2b<1a2L,1b2 & EE/DD & {-}{-}/{-}{-} & E &  \\
B2586 & 1a<o2 1b<o2 / 2a<1a1,1b2L 2b<1a1L,1b2 & EE/DD & {-}{-}/{-}{-} & E &  \\
B2587 & 1a<o2 1b<o2 / 2a<1a1,1b2L 2b<1a2,1b2 & EE/DD & {-}{-}/{-}{-} & E &  \\
B2588 & 1a<o2 1b<o2 / 2a<1a1,1b2L 2b<1a2L,1b1 & EE/DD & {-}{-}/{-}{-} & E &  \\
B2589 & 1a<o2 1b<o2 / 2a<1a1,1b2L 2b<1a2L,1b2 & EE/DD & {-}{-}/{-}{-} & E &  \\
B2590 & 1a<o2 1b<o2 / 2a<1a1L,1b2 2b<1a2,1b1L & EE/DD & {-}{-}/{-}{-} & E &  \\
B2591 & 1a<o2 1b<o2 / 2a<1a1L,1b2 2b<1a2,1b2 & EE/DD & {-}{-}/{-}{-} & E &  \\
B2592 & 1a<o2 1b<o2 / 2a<1a1L,1b2 2b<1a2,1b2L & EE/DD & {-}{-}/{-}{-} & E &  \\
B2593 & 1a<o2 1b<o2 / 2a<1a2,1b2 2b<1a2,1b2 & EE/DD & {-}{-}/{-}{-} & E &  \\
B2594 & 1a<o2 1b<o2 / 2a<1a2,1b2 2b<1a2,1b2L & EE/DD & {-}{-}/{-}{-} & E &  \\
B2595 & 1a<o2 1b<o2 / 2a<1a2,1b2L 2b<1a2L,1b2 & EE/DD & {-}{-}/{-}{-} & E &  \\
B2596 & 1a<o2 1b<o2L / 2a<1a0L,1b1 2b<1a1,1b0L & LI/ID & {-}+/{-}+ & B & CPW:1b,2a \\
B2597 & 1a<o2 1b<o2L / 2a<1a0L,1b1 2b<1a1,1b1 & LI/EE & {-}+/{-}{-} & BE & CPW:1b \\
B2598 & 1a<o2 1b<o2L / 2a<1a0L,1b1 2b<1a1,1b1L & LI/LI & {-}+/{-}+ & B & CPW:1b,2b \\
B2599 & 1a<o2 1b<o2L / 2a<1a0L,1b1 2b<1a1,1b2 & LI/EE & {-}+/{-}{-} & BE & CPW:1b \\
B2600 & 1a<o2 1b<o2L / 2a<1a0L,1b1 2b<1a1,1b2L & LI/LI & {-}+/{-}+ & B & CPW:1b,2b \\
B2601 & 1a<o2 1b<o2L / 2a<1a0L,1b1 2b<1a2,1b0L & LI/ID & {-}+/{-}+ & B & CPW:1b,2a \\
B2602 & 1a<o2 1b<o2L / 2a<1a0L,1b1 2b<1a2,1b1 & LI/EE & {-}+/{-}{-} & BE & CPW:1b \\
B2603 & 1a<o2 1b<o2L / 2a<1a0L,1b1 2b<1a2,1b1L & LI/LI & {-}+/{-}+ & B & CPW:1b,2b \\
B2604 & 1a<o2 1b<o2L / 2a<1a0L,1b1 2b<1a2,1b2 & LI/EE & {-}+/{-}{-} & BE & CPW:1b \\
B2605 & 1a<o2 1b<o2L / 2a<1a0L,1b1 2b<1a2,1b2L & LI/LI & {-}+/{-}+ & B & CPW:1b,2b \\
B2606 & 1a<o2 1b<o2L / 2a<1a0L,1b2 2b<1a1,1b0L & LI/ID & {-}+/{-}+ & B & CPW:1b,2a \\
B2607 & 1a<o2 1b<o2L / 2a<1a0L,1b2 2b<1a1,1b1 & LI/EE & {-}+/{-}{-} & BE & CPW:1b \\
B2608 & 1a<o2 1b<o2L / 2a<1a0L,1b2 2b<1a1,1b1L & LI/LI & {-}+/{-}+ & B & CPW:1b,2b \\
B2609 & 1a<o2 1b<o2L / 2a<1a0L,1b2 2b<1a1,1b2 & LI/EE & {-}+/{-}{-} & BE & CPW:1b \\
B2610 & 1a<o2 1b<o2L / 2a<1a0L,1b2 2b<1a1,1b2L & LI/LI & {-}+/{-}+ & B & CPW:1b,2b \\
B2611 & 1a<o2 1b<o2L / 2a<1a0L,1b2 2b<1a2,1b0L & LI/ID & {-}+/{-}+ & B & CPW:1b,2a \\
B2612 & 1a<o2 1b<o2L / 2a<1a0L,1b2 2b<1a2,1b1 & LI/EE & {-}+/{-}{-} & BE & CPW:1b \\
B2613 & 1a<o2 1b<o2L / 2a<1a0L,1b2 2b<1a2,1b1L & LI/LI & {-}+/{-}+ & B & CPW:1b,2b \\
B2614 & 1a<o2 1b<o2L / 2a<1a0L,1b2 2b<1a2,1b2 & LI/EE & {-}+/{-}{-} & BE & CPW:1b \\
B2615 & 1a<o2 1b<o2L / 2a<1a0L,1b2 2b<1a2,1b2L & LI/LI & {-}+/{-}+ & B & CPW:1b,2b \\
B2616 & 1a<o2 1b<o2L / 2a<1a1,1b0L 2b<1a1,1b1 & LI/DE & {-}+/+{-} & BF & CPW:1b \\
B2617 & 1a<o2 1b<o2L / 2a<1a1,1b0L 2b<1a1,1b2 & LI/DE & {-}+/+{-} & BF & CPW:1b \\
B2618 & 1a<o2 1b<o2L / 2a<1a1,1b0L 2b<1a1L,1b1 & LI/DE & {-}+/+{-} & BF & CPW:1b \\
B2619 & 1a<o2 1b<o2L / 2a<1a1,1b0L 2b<1a1L,1b2 & LI/DE & {-}+/+{-} & BF & CPW:1b \\
B2620 & 1a<o2 1b<o2L / 2a<1a1,1b0L 2b<1a2,1b1 & LI/DE & {-}+/+{-} & BF & CPW:1b \\
B2621 & 1a<o2 1b<o2L / 2a<1a1,1b0L 2b<1a2,1b2 & LI/DE & {-}+/+{-} & BF & CPW:1b \\
B2622 & 1a<o2 1b<o2L / 2a<1a1,1b0L 2b<1a2L,1b1 & LI/DE & {-}+/+{-} & BF & CPW:1b \\
B2623 & 1a<o2 1b<o2L / 2a<1a1,1b0L 2b<1a2L,1b2 & LI/DE & {-}+/+{-} & BF & CPW:1b \\
B2624 & 1a<o2 1b<o2L / 2a<1a1,1b1 2b<1a1,1b1 & LI/EE & {-}+/{-}{-} & BE & CPW:1b \\
B2625 & 1a<o2 1b<o2L / 2a<1a1,1b1 2b<1a1,1b1L & LI/LI & {-}+/{-}+ & B & CPW:1b,2b \\
B2626 & 1a<o2 1b<o2L / 2a<1a1,1b1 2b<1a1,1b2 & LI/EE & {-}+/{-}{-} & BE & CPW:1b \\
B2627 & 1a<o2 1b<o2L / 2a<1a1,1b1 2b<1a1,1b2L & LI/LI & {-}+/{-}+ & B & CPW:1b,2b \\
B2628 & 1a<o2 1b<o2L / 2a<1a1,1b1 2b<1a1L,1b1 & LI/EE & {-}+/{-}{-} & BE & CPW:1b \\
B2629 & 1a<o2 1b<o2L / 2a<1a1,1b1 2b<1a1L,1b2 & LI/EE & {-}+/{-}{-} & BE & CPW:1b \\
B2630 & 1a<o2 1b<o2L / 2a<1a1,1b1 2b<1a2,1b0L & LI/ED & {-}+/{-}+ & BF & CPW:1b \\
B2631 & 1a<o2 1b<o2L / 2a<1a1,1b1 2b<1a2,1b1 & LI/EE & {-}+/{-}{-} & BE & CPW:1b \\
B2632 & 1a<o2 1b<o2L / 2a<1a1,1b1 2b<1a2,1b1L & LI/LI & {-}+/{-}+ & B & CPW:1b,2b \\
B2633 & 1a<o2 1b<o2L / 2a<1a1,1b1 2b<1a2,1b2 & LI/EE & {-}+/{-}{-} & BE & CPW:1b \\
B2634 & 1a<o2 1b<o2L / 2a<1a1,1b1 2b<1a2,1b2L & LI/LI & {-}+/{-}+ & B & CPW:1b,2b \\
B2635 & 1a<o2 1b<o2L / 2a<1a1,1b1 2b<1a2L,1b1 & LI/EE & {-}+/{-}{-} & BE & CPW:1b \\
B2636 & 1a<o2 1b<o2L / 2a<1a1,1b1 2b<1a2L,1b2 & LI/EE & {-}+/{-}{-} & BE & CPW:1b \\
B2637 & 1a<o2 1b<o2L / 2a<1a1,1b1L 2b<1a1,1b2 & LI/IL & {-}+/+{-} & B & CPW:1b,2a \\
B2638 & 1a<o2 1b<o2L / 2a<1a1,1b1L 2b<1a1L,1b1 & LI/IL & {-}+/+{-} & B & CPW:1b,2a \\
B2639 & 1a<o2 1b<o2L / 2a<1a1,1b1L 2b<1a1L,1b2 & LI/IL & {-}+/+{-} & B & CPW:1b,2a \\
B2640 & 1a<o2 1b<o2L / 2a<1a1,1b1L 2b<1a2,1b1 & LI/IL & {-}+/+{-} & B & CPW:1b,2a \\
B2641 & 1a<o2 1b<o2L / 2a<1a1,1b1L 2b<1a2,1b2 & LI/IL & {-}+/+{-} & B & CPW:1b,2a \\
B2642 & 1a<o2 1b<o2L / 2a<1a1,1b1L 2b<1a2L,1b1 & LI/IL & {-}+/+{-} & B & CPW:1b,2a \\
B2643 & 1a<o2 1b<o2L / 2a<1a1,1b1L 2b<1a2L,1b2 & LI/IL & {-}+/+{-} & B & CPW:1b,2a \\
B2644 & 1a<o2 1b<o2L / 2a<1a1,1b2 2b<1a1,1b2 & LI/EE & {-}+/{-}{-} & BE & CPW:1b \\
B2645 & 1a<o2 1b<o2L / 2a<1a1,1b2 2b<1a1,1b2L & LI/LI & {-}+/{-}+ & B & CPW:1b,2b \\
B2646 & 1a<o2 1b<o2L / 2a<1a1,1b2 2b<1a1L,1b1 & LI/EE & {-}+/{-}{-} & BE & CPW:1b \\
B2647 & 1a<o2 1b<o2L / 2a<1a1,1b2 2b<1a1L,1b2 & LI/EE & {-}+/{-}{-} & BE & CPW:1b \\
B2648 & 1a<o2 1b<o2L / 2a<1a1,1b2 2b<1a2,1b0L & LI/ED & {-}+/{-}+ & BF & CPW:1b \\
B2649 & 1a<o2 1b<o2L / 2a<1a1,1b2 2b<1a2,1b1 & LI/EE & {-}+/{-}{-} & BE & CPW:1b \\
B2650 & 1a<o2 1b<o2L / 2a<1a1,1b2 2b<1a2,1b1L & LI/LI & {-}+/{-}+ & B & CPW:1b,2b \\
B2651 & 1a<o2 1b<o2L / 2a<1a1,1b2 2b<1a2,1b2 & LI/EE & {-}+/{-}{-} & BE & CPW:1b \\
B2652 & 1a<o2 1b<o2L / 2a<1a1,1b2 2b<1a2,1b2L & LI/LI & {-}+/{-}+ & B & CPW:1b,2b \\
B2653 & 1a<o2 1b<o2L / 2a<1a1,1b2 2b<1a2L,1b1 & LI/EE & {-}+/{-}{-} & BE & CPW:1b \\
B2654 & 1a<o2 1b<o2L / 2a<1a1,1b2 2b<1a2L,1b2 & LI/EE & {-}+/{-}{-} & BE & CPW:1b \\
B2655 & 1a<o2 1b<o2L / 2a<1a1,1b2L 2b<1a1L,1b1 & LI/IL & {-}+/+{-} & B & CPW:1b,2a \\
B2656 & 1a<o2 1b<o2L / 2a<1a1,1b2L 2b<1a1L,1b2 & LI/IL & {-}+/+{-} & B & CPW:1b,2a \\
B2657 & 1a<o2 1b<o2L / 2a<1a1,1b2L 2b<1a2,1b1 & LI/IL & {-}+/+{-} & B & CPW:1b,2a \\
B2658 & 1a<o2 1b<o2L / 2a<1a1,1b2L 2b<1a2,1b2 & LI/IL & {-}+/+{-} & B & CPW:1b,2a \\
B2659 & 1a<o2 1b<o2L / 2a<1a1,1b2L 2b<1a2L,1b1 & LI/IL & {-}+/+{-} & B & CPW:1b,2a \\
B2660 & 1a<o2 1b<o2L / 2a<1a1,1b2L 2b<1a2L,1b2 & LI/IL & {-}+/+{-} & B & CPW:1b,2a \\
B2661 & 1a<o2 1b<o2L / 2a<1a1L,1b1 2b<1a2,1b0L & LI/ED & {-}+/{-}+ & BF & CPW:1b \\
B2662 & 1a<o2 1b<o2L / 2a<1a1L,1b1 2b<1a2,1b1 & LI/EE & {-}+/{-}{-} & BE & CPW:1b \\
B2663 & 1a<o2 1b<o2L / 2a<1a1L,1b1 2b<1a2,1b1L & LI/LI & {-}+/{-}+ & B & CPW:1b,2b \\
B2664 & 1a<o2 1b<o2L / 2a<1a1L,1b1 2b<1a2,1b2 & LI/EE & {-}+/{-}{-} & BE & CPW:1b \\
B2665 & 1a<o2 1b<o2L / 2a<1a1L,1b1 2b<1a2,1b2L & LI/LI & {-}+/{-}+ & B & CPW:1b,2b \\
B2666 & 1a<o2 1b<o2L / 2a<1a1L,1b2 2b<1a2,1b0L & LI/ED & {-}+/{-}+ & BF & CPW:1b \\
B2667 & 1a<o2 1b<o2L / 2a<1a1L,1b2 2b<1a2,1b1 & LI/EE & {-}+/{-}{-} & BE & CPW:1b \\
B2668 & 1a<o2 1b<o2L / 2a<1a1L,1b2 2b<1a2,1b1L & LI/LI & {-}+/{-}+ & B & CPW:1b,2b \\
B2669 & 1a<o2 1b<o2L / 2a<1a1L,1b2 2b<1a2,1b2 & LI/EE & {-}+/{-}{-} & BE & CPW:1b \\
B2670 & 1a<o2 1b<o2L / 2a<1a1L,1b2 2b<1a2,1b2L & LI/LI & {-}+/{-}+ & B & CPW:1b,2b \\
B2671 & 1a<o2 1b<o2L / 2a<1a2,1b0L 2b<1a2,1b1 & LI/DE & {-}+/+{-} & BF & CPW:1b \\
B2672 & 1a<o2 1b<o2L / 2a<1a2,1b0L 2b<1a2,1b2 & LI/DE & {-}+/+{-} & BF & CPW:1b \\
B2673 & 1a<o2 1b<o2L / 2a<1a2,1b0L 2b<1a2L,1b1 & LI/DE & {-}+/+{-} & BF & CPW:1b \\
B2674 & 1a<o2 1b<o2L / 2a<1a2,1b0L 2b<1a2L,1b2 & LI/DE & {-}+/+{-} & BF & CPW:1b \\
B2675 & 1a<o2 1b<o2L / 2a<1a2,1b1 2b<1a2,1b1 & LI/EE & {-}+/{-}{-} & BE & CPW:1b \\
B2676 & 1a<o2 1b<o2L / 2a<1a2,1b1 2b<1a2,1b1L & LI/LI & {-}+/{-}+ & B & CPW:1b,2b \\
B2677 & 1a<o2 1b<o2L / 2a<1a2,1b1 2b<1a2,1b2 & LI/EE & {-}+/{-}{-} & BE & CPW:1b \\
B2678 & 1a<o2 1b<o2L / 2a<1a2,1b1 2b<1a2,1b2L & LI/LI & {-}+/{-}+ & B & CPW:1b,2b \\
B2679 & 1a<o2 1b<o2L / 2a<1a2,1b1 2b<1a2L,1b1 & LI/EE & {-}+/{-}{-} & BE & CPW:1b \\
B2680 & 1a<o2 1b<o2L / 2a<1a2,1b1 2b<1a2L,1b2 & LI/EE & {-}+/{-}{-} & BE & CPW:1b \\
B2681 & 1a<o2 1b<o2L / 2a<1a2,1b1L 2b<1a2,1b2 & LI/IL & {-}+/+{-} & B & CPW:1b,2a \\
B2682 & 1a<o2 1b<o2L / 2a<1a2,1b1L 2b<1a2L,1b1 & LI/IL & {-}+/+{-} & B & CPW:1b,2a \\
B2683 & 1a<o2 1b<o2L / 2a<1a2,1b1L 2b<1a2L,1b2 & LI/IL & {-}+/+{-} & B & CPW:1b,2a \\
B2684 & 1a<o2 1b<o2L / 2a<1a2,1b2 2b<1a2,1b2 & LI/EE & {-}+/{-}{-} & BE & CPW:1b \\
B2685 & 1a<o2 1b<o2L / 2a<1a2,1b2 2b<1a2,1b2L & LI/LI & {-}+/{-}+ & B & CPW:1b,2b \\
B2686 & 1a<o2 1b<o2L / 2a<1a2,1b2 2b<1a2L,1b1 & LI/EE & {-}+/{-}{-} & BE & CPW:1b \\
B2687 & 1a<o2 1b<o2L / 2a<1a2,1b2 2b<1a2L,1b2 & LI/EE & {-}+/{-}{-} & BE & CPW:1b \\
B2688 & 1a<o2 1b<o2L / 2a<1a2,1b2L 2b<1a2L,1b1 & LI/IL & {-}+/+{-} & B & CPW:1b,2a \\
B2689 & 1a<o2 1b<o2L / 2a<1a2,1b2L 2b<1a2L,1b2 & LI/IL & {-}+/+{-} & B & CPW:1b,2a \\
B2690 & 1a<o0L 1b<o1 / 2a<1a0L 2b*<1a1 & NI/NN & +{-}/+{-} & N &  \\
B2691 & 1a<o0L 1b<o1 / 2a<1a0L 2b*<1a2 & NI/NN & +{-}/+{-} & N &  \\
B2692 & 1a<o0L 1b<o1 / 2a<1a0L 2b*<1b1 & NI/NU & +{-}/+{-} & N & W:2b \\
B2693 & 1a<o0L 1b<o1 / 2a<1a0L 2b*<1b2 & NI/NU & +{-}/+{-} & N & W:2b \\
B2694 & 1a<o0L 1b<o1 / 2a<1a1 2b*<1a1 & NI/NN & +{-}/{-}{-} & N &  \\
B2695 & 1a<o0L 1b<o1 / 2a<1a1 2b*<1a2 & NI/NN & +{-}/{-}{-} & N &  \\
B2696 & 1a<o0L 1b<o1 / 2a<1a1 2b*<1b1 & NI/NU & +{-}/{-}{-} & N & W:2b \\
B2697 & 1a<o0L 1b<o1 / 2a<1a1 2b*<1b2 & NI/NU & +{-}/{-}{-} & N & W:2b \\
B2698 & 1a<o0L 1b<o1 / 2a<1a1L 2b*<1a1 & NI/NN & +{-}/+{-} & N &  \\
B2699 & 1a<o0L 1b<o1 / 2a<1a1L 2b*<1a2 & NI/NN & +{-}/+{-} & N &  \\
B2700 & 1a<o0L 1b<o1 / 2a<1a1L 2b*<1b1 & NI/NU & +{-}/+{-} & N & W:2b \\
B2701 & 1a<o0L 1b<o1 / 2a<1a1L 2b*<1b2 & NI/NU & +{-}/+{-} & N & W:2b \\
B2702 & 1a<o0L 1b<o1 / 2a<1a2 2b*<1a1 & NI/NN & +{-}/{-}{-} & N &  \\
B2703 & 1a<o0L 1b<o1 / 2a<1a2 2b*<1a2 & NI/NN & +{-}/{-}{-} & N &  \\
B2704 & 1a<o0L 1b<o1 / 2a<1a2 2b*<1b1 & NI/NU & +{-}/{-}{-} & N & W:2b \\
B2705 & 1a<o0L 1b<o1 / 2a<1a2 2b*<1b2 & NI/NU & +{-}/{-}{-} & N & W:2b \\
B2706 & 1a<o0L 1b<o1 / 2a<1a2L 2b*<1a1 & NI/NN & +{-}/+{-} & N &  \\
B2707 & 1a<o0L 1b<o1 / 2a<1a2L 2b*<1a2 & NI/NN & +{-}/+{-} & N &  \\
B2708 & 1a<o0L 1b<o1 / 2a<1a2L 2b*<1b1 & NI/NU & +{-}/+{-} & N & W:2b \\
B2709 & 1a<o0L 1b<o1 / 2a<1a2L 2b*<1b2 & NI/NU & +{-}/+{-} & N & W:2b \\
B2710 & 1a<o0L 1b<o1 / 2a<1b0L 2b*<1a1 & NI/NN & +{-}/{-}{-} & N &  \\
B2711 & 1a<o0L 1b<o1 / 2a<1b0L 2b*<1a2 & NI/NN & +{-}/{-}{-} & N &  \\
B2712 & 1a<o0L 1b<o1 / 2a<1b0L 2b*<1b1 & NI/NU & +{-}/{-}{-} & N & W:2b \\
B2713 & 1a<o0L 1b<o1 / 2a<1b0L 2b*<1b2 & NI/NU & +{-}/{-}{-} & N & W:2b \\
B2714 & 1a<o0L 1b<o1 / 2a<1b1 2b*<1a1 & NI/IN & +{-}/{-}{-} & N &  \\
B2715 & 1a<o0L 1b<o1 / 2a<1b1 2b*<1a2 & NI/IN & +{-}/{-}{-} & N &  \\
B2716 & 1a<o0L 1b<o1 / 2a<1b1 2b*<1b1 & NI/IU & +{-}/{-}{-} & N & W:1b,2a \\
B2717 & 1a<o0L 1b<o1 / 2a<1b1 2b*<1b2 & NI/IU & +{-}/{-}{-} & N & W:1b,2a \\
B2718 & 1a<o0L 1b<o1 / 2a<1b1L 2b*<1a1 & NI/IN & +{-}/{-}{-} & N &  \\
B2719 & 1a<o0L 1b<o1 / 2a<1b1L 2b*<1a2 & NI/IN & +{-}/{-}{-} & N &  \\
B2720 & 1a<o0L 1b<o1 / 2a<1b1L 2b*<1b1 & NI/IU & +{-}/{-}{-} & N & W:1b,2a \\
B2721 & 1a<o0L 1b<o1 / 2a<1b1L 2b*<1b2 & NI/IU & +{-}/{-}{-} & N & W:1b,2a \\
B2722 & 1a<o0L 1b<o1 / 2a<1b2 2b*<1a1 & NI/IN & +{-}/{-}{-} & N &  \\
B2723 & 1a<o0L 1b<o1 / 2a<1b2 2b*<1a2 & NI/IN & +{-}/{-}{-} & N &  \\
B2724 & 1a<o0L 1b<o1 / 2a<1b2 2b*<1b1 & NI/IU & +{-}/{-}{-} & N & W:1b,2a \\
B2725 & 1a<o0L 1b<o1 / 2a<1b2 2b*<1b2 & NI/IU & +{-}/{-}{-} & N & W:1b,2a \\
B2726 & 1a<o0L 1b<o1 / 2a<1b2L 2b*<1a1 & NI/IN & +{-}/{-}{-} & N &  \\
B2727 & 1a<o0L 1b<o1 / 2a<1b2L 2b*<1a2 & NI/IN & +{-}/{-}{-} & N &  \\
B2728 & 1a<o0L 1b<o1 / 2a<1b2L 2b*<1b1 & NI/IU & +{-}/{-}{-} & N & W:1b,2a \\
B2729 & 1a<o0L 1b<o1 / 2a<1b2L 2b*<1b2 & NI/IU & +{-}/{-}{-} & N & W:1b,2a \\
B2730 & 1a<o0L 1b<o2 / 2a<1a0L 2b*<1a1 & NI/NN & +{-}/+{-} & N &  \\
B2731 & 1a<o0L 1b<o2 / 2a<1a0L 2b*<1a2 & NI/NN & +{-}/+{-} & N &  \\
B2732 & 1a<o0L 1b<o2 / 2a<1a0L 2b*<1b1 & NI/NU & +{-}/+{-} & N & W:2b \\
B2733 & 1a<o0L 1b<o2 / 2a<1a0L 2b*<1b2 & NI/NU & +{-}/+{-} & N & W:2b \\
B2734 & 1a<o0L 1b<o2 / 2a<1a1 2b*<1a1 & NI/NN & +{-}/{-}{-} & N &  \\
B2735 & 1a<o0L 1b<o2 / 2a<1a1 2b*<1a2 & NI/NN & +{-}/{-}{-} & N &  \\
B2736 & 1a<o0L 1b<o2 / 2a<1a1 2b*<1b1 & NI/NU & +{-}/{-}{-} & N & W:2b \\
B2737 & 1a<o0L 1b<o2 / 2a<1a1 2b*<1b2 & NI/NU & +{-}/{-}{-} & N & W:2b \\
B2738 & 1a<o0L 1b<o2 / 2a<1a1L 2b*<1a1 & NI/NN & +{-}/+{-} & N &  \\
B2739 & 1a<o0L 1b<o2 / 2a<1a1L 2b*<1a2 & NI/NN & +{-}/+{-} & N &  \\
B2740 & 1a<o0L 1b<o2 / 2a<1a1L 2b*<1b1 & NI/NU & +{-}/+{-} & N & W:2b \\
B2741 & 1a<o0L 1b<o2 / 2a<1a1L 2b*<1b2 & NI/NU & +{-}/+{-} & N & W:2b \\
B2742 & 1a<o0L 1b<o2 / 2a<1a2 2b*<1a1 & NI/NN & +{-}/{-}{-} & N &  \\
B2743 & 1a<o0L 1b<o2 / 2a<1a2 2b*<1a2 & NI/NN & +{-}/{-}{-} & N &  \\
B2744 & 1a<o0L 1b<o2 / 2a<1a2 2b*<1b1 & NI/NU & +{-}/{-}{-} & N & W:2b \\
B2745 & 1a<o0L 1b<o2 / 2a<1a2 2b*<1b2 & NI/NU & +{-}/{-}{-} & N & W:2b \\
B2746 & 1a<o0L 1b<o2 / 2a<1a2L 2b*<1a1 & NI/NN & +{-}/+{-} & N &  \\
B2747 & 1a<o0L 1b<o2 / 2a<1a2L 2b*<1a2 & NI/NN & +{-}/+{-} & N &  \\
B2748 & 1a<o0L 1b<o2 / 2a<1a2L 2b*<1b1 & NI/NU & +{-}/+{-} & N & W:2b \\
B2749 & 1a<o0L 1b<o2 / 2a<1a2L 2b*<1b2 & NI/NU & +{-}/+{-} & N & W:2b \\
B2750 & 1a<o0L 1b<o2 / 2a<1b0L 2b*<1a1 & NI/NN & +{-}/{-}{-} & N &  \\
B2751 & 1a<o0L 1b<o2 / 2a<1b0L 2b*<1a2 & NI/NN & +{-}/{-}{-} & N &  \\
B2752 & 1a<o0L 1b<o2 / 2a<1b0L 2b*<1b1 & NI/NU & +{-}/{-}{-} & N & W:2b \\
B2753 & 1a<o0L 1b<o2 / 2a<1b0L 2b*<1b2 & NI/NU & +{-}/{-}{-} & N & W:2b \\
B2754 & 1a<o0L 1b<o2 / 2a<1b1 2b*<1a1 & NI/IN & +{-}/{-}{-} & N &  \\
B2755 & 1a<o0L 1b<o2 / 2a<1b1 2b*<1a2 & NI/IN & +{-}/{-}{-} & N &  \\
B2756 & 1a<o0L 1b<o2 / 2a<1b1 2b*<1b1 & NI/IU & +{-}/{-}{-} & N & W:1b,2a \\
B2757 & 1a<o0L 1b<o2 / 2a<1b1 2b*<1b2 & NI/IU & +{-}/{-}{-} & N & W:1b,2a \\
B2758 & 1a<o0L 1b<o2 / 2a<1b1L 2b*<1a1 & NI/IN & +{-}/{-}{-} & N &  \\
B2759 & 1a<o0L 1b<o2 / 2a<1b1L 2b*<1a2 & NI/IN & +{-}/{-}{-} & N &  \\
B2760 & 1a<o0L 1b<o2 / 2a<1b1L 2b*<1b1 & NI/IU & +{-}/{-}{-} & N & W:1b,2a \\
B2761 & 1a<o0L 1b<o2 / 2a<1b1L 2b*<1b2 & NI/IU & +{-}/{-}{-} & N & W:1b,2a \\
B2762 & 1a<o0L 1b<o2 / 2a<1b2 2b*<1a1 & NI/IN & +{-}/{-}{-} & N &  \\
B2763 & 1a<o0L 1b<o2 / 2a<1b2 2b*<1a2 & NI/IN & +{-}/{-}{-} & N &  \\
B2764 & 1a<o0L 1b<o2 / 2a<1b2 2b*<1b1 & NI/IU & +{-}/{-}{-} & N & W:1b,2a \\
B2765 & 1a<o0L 1b<o2 / 2a<1b2 2b*<1b2 & NI/IU & +{-}/{-}{-} & N & W:1b,2a \\
B2766 & 1a<o0L 1b<o2 / 2a<1b2L 2b*<1a1 & NI/IN & +{-}/{-}{-} & N &  \\
B2767 & 1a<o0L 1b<o2 / 2a<1b2L 2b*<1a2 & NI/IN & +{-}/{-}{-} & N &  \\
B2768 & 1a<o0L 1b<o2 / 2a<1b2L 2b*<1b1 & NI/IU & +{-}/{-}{-} & N & W:1b,2a \\
B2769 & 1a<o0L 1b<o2 / 2a<1b2L 2b*<1b2 & NI/IU & +{-}/{-}{-} & N & W:1b,2a \\
B2770 & 1a<o0L 1b*<o1 / 2a<1a0L 2b<1a1 & NU/NN & +{-}/+{-} & N & W:1b \\
B2771 & 1a<o0L 1b*<o1 / 2a<1a0L 2b<1a2 & NU/NN & +{-}/+{-} & N & W:1b \\
B2772 & 1a<o0L 1b*<o1 / 2a<1a0L 2b<1b1 & NU/NU & +{-}/+{-} & N & W:1b,2b \\
B2773 & 1a<o0L 1b*<o1 / 2a<1a0L 2b<1b2 & NU/NU & +{-}/+{-} & N & W:1b,2b \\
B2774 & 1a<o0L 1b*<o1 / 2a<1a1 2b<1a1 & NU/NN & +{-}/{-}{-} & N & W:1b \\
B2775 & 1a<o0L 1b*<o1 / 2a<1a1 2b<1a1L & NU/NN & +{-}/{-}+ & N & W:1b \\
B2776 & 1a<o0L 1b*<o1 / 2a<1a1 2b<1a2 & NU/NN & +{-}/{-}{-} & N & W:1b \\
B2777 & 1a<o0L 1b*<o1 / 2a<1a1 2b<1a2L & NU/NN & +{-}/{-}+ & N & W:1b \\
B2778 & 1a<o0L 1b*<o1 / 2a<1a1 2b<1b1 & NU/NU & +{-}/{-}{-} & N & W:1b,2b \\
B2779 & 1a<o0L 1b*<o1 / 2a<1a1 2b<1b2 & NU/NU & +{-}/{-}{-} & N & W:1b,2b \\
B2780 & 1a<o0L 1b*<o1 / 2a<1a1L 2b<1a2 & NU/NN & +{-}/+{-} & N & W:1b \\
B2781 & 1a<o0L 1b*<o1 / 2a<1a1L 2b<1b1 & NU/NU & +{-}/+{-} & N & W:1b,2b \\
B2782 & 1a<o0L 1b*<o1 / 2a<1a1L 2b<1b2 & NU/NU & +{-}/+{-} & N & W:1b,2b \\
B2783 & 1a<o0L 1b*<o1 / 2a<1a2 2b<1a2 & NU/NN & +{-}/{-}{-} & N & W:1b \\
B2784 & 1a<o0L 1b*<o1 / 2a<1a2 2b<1a2L & NU/NN & +{-}/{-}+ & N & W:1b \\
B2785 & 1a<o0L 1b*<o1 / 2a<1a2 2b<1b1 & NU/NU & +{-}/{-}{-} & N & W:1b,2b \\
B2786 & 1a<o0L 1b*<o1 / 2a<1a2 2b<1b2 & NU/NU & +{-}/{-}{-} & N & W:1b,2b \\
B2787 & 1a<o0L 1b*<o1 / 2a<1a2L 2b<1b1 & NU/NU & +{-}/+{-} & N & W:1b,2b \\
B2788 & 1a<o0L 1b*<o1 / 2a<1a2L 2b<1b2 & NU/NU & +{-}/+{-} & N & W:1b,2b \\
B2789 & 1a<o0L 1b*<o1 / 2a<1b1 2b<1b1 & NU/UU & +{-}/{-}{-} & N &  \\
B2790 & 1a<o0L 1b*<o1 / 2a<1b1 2b<1b2 & NU/UU & +{-}/{-}{-} & N &  \\
B2791 & 1a<o0L 1b*<o1 / 2a<1b2 2b<1b2 & NU/UU & +{-}/{-}{-} & N &  \\
B2792 & 1a<o0L 1b*<o2 / 2a<1a0L 2b<1a1 & NU/NN & +{-}/+{-} & N & W:1b \\
B2793 & 1a<o0L 1b*<o2 / 2a<1a0L 2b<1a2 & NU/NN & +{-}/+{-} & N & W:1b \\
B2794 & 1a<o0L 1b*<o2 / 2a<1a0L 2b<1b1 & NU/NU & +{-}/+{-} & N & W:1b,2b \\
B2795 & 1a<o0L 1b*<o2 / 2a<1a0L 2b<1b2 & NU/NU & +{-}/+{-} & N & W:1b,2b \\
B2796 & 1a<o0L 1b*<o2 / 2a<1a1 2b<1a1 & NU/NN & +{-}/{-}{-} & N & W:1b \\
B2797 & 1a<o0L 1b*<o2 / 2a<1a1 2b<1a1L & NU/NN & +{-}/{-}+ & N & W:1b \\
B2798 & 1a<o0L 1b*<o2 / 2a<1a1 2b<1a2 & NU/NN & +{-}/{-}{-} & N & W:1b \\
B2799 & 1a<o0L 1b*<o2 / 2a<1a1 2b<1a2L & NU/NN & +{-}/{-}+ & N & W:1b \\
B2800 & 1a<o0L 1b*<o2 / 2a<1a1 2b<1b1 & NU/NU & +{-}/{-}{-} & N & W:1b,2b \\
B2801 & 1a<o0L 1b*<o2 / 2a<1a1 2b<1b2 & NU/NU & +{-}/{-}{-} & N & W:1b,2b \\
B2802 & 1a<o0L 1b*<o2 / 2a<1a1L 2b<1a2 & NU/NN & +{-}/+{-} & N & W:1b \\
B2803 & 1a<o0L 1b*<o2 / 2a<1a1L 2b<1b1 & NU/NU & +{-}/+{-} & N & W:1b,2b \\
B2804 & 1a<o0L 1b*<o2 / 2a<1a1L 2b<1b2 & NU/NU & +{-}/+{-} & N & W:1b,2b \\
B2805 & 1a<o0L 1b*<o2 / 2a<1a2 2b<1a2 & NU/NN & +{-}/{-}{-} & N & W:1b \\
B2806 & 1a<o0L 1b*<o2 / 2a<1a2 2b<1a2L & NU/NN & +{-}/{-}+ & N & W:1b \\
B2807 & 1a<o0L 1b*<o2 / 2a<1a2 2b<1b1 & NU/NU & +{-}/{-}{-} & N & W:1b,2b \\
B2808 & 1a<o0L 1b*<o2 / 2a<1a2 2b<1b2 & NU/NU & +{-}/{-}{-} & N & W:1b,2b \\
B2809 & 1a<o0L 1b*<o2 / 2a<1a2L 2b<1b1 & NU/NU & +{-}/+{-} & N & W:1b,2b \\
B2810 & 1a<o0L 1b*<o2 / 2a<1a2L 2b<1b2 & NU/NU & +{-}/+{-} & N & W:1b,2b \\
B2811 & 1a<o0L 1b*<o2 / 2a<1b1 2b<1b1 & NU/UU & +{-}/{-}{-} & N &  \\
B2812 & 1a<o0L 1b*<o2 / 2a<1b1 2b<1b2 & NU/UU & +{-}/{-}{-} & N &  \\
B2813 & 1a<o0L 1b*<o2 / 2a<1b2 2b<1b2 & NU/UU & +{-}/{-}{-} & N &  \\
B2814 & 1a<o1 1b<o1 / 2a<1a0L 2b*<1a1 & EE/NU & {-}{-}/{-}{-} & E &  \\
B2815 & 1a<o1 1b<o1 / 2a<1a0L 2b*<1a2 & EE/NU & {-}{-}/{-}{-} & E &  \\
B2816 & 1a<o1 1b<o1 / 2a<1a0L 2b*<1b1 & EE/NU & {-}{-}/{-}{-} & E &  \\
B2817 & 1a<o1 1b<o1 / 2a<1a0L 2b*<1b2 & EE/NU & {-}{-}/{-}{-} & E &  \\
B2818 & 1a<o1 1b<o1 / 2a<1a1 2b*<1a1 & EE/DU & {-}{-}/{-}{-} & E &  \\
B2819 & 1a<o1 1b<o1 / 2a<1a1 2b*<1a2 & EE/DU & {-}{-}/{-}{-} & E &  \\
B2820 & 1a<o1 1b<o1 / 2a<1a1 2b*<1b1 & EE/DU & {-}{-}/{-}{-} & E &  \\
B2821 & 1a<o1 1b<o1 / 2a<1a1 2b*<1b2 & EE/DU & {-}{-}/{-}{-} & E &  \\
B2822 & 1a<o1 1b<o1 / 2a<1a1L 2b*<1a1 & EE/DU & {-}{-}/{-}{-} & E &  \\
B2823 & 1a<o1 1b<o1 / 2a<1a1L 2b*<1a2 & EE/DU & {-}{-}/{-}{-} & E &  \\
B2824 & 1a<o1 1b<o1 / 2a<1a1L 2b*<1b1 & EE/DU & {-}{-}/{-}{-} & E &  \\
B2825 & 1a<o1 1b<o1 / 2a<1a1L 2b*<1b2 & EE/DU & {-}{-}/{-}{-} & E &  \\
B2826 & 1a<o1 1b<o1 / 2a<1a2 2b*<1a1 & EE/DU & {-}{-}/{-}{-} & E &  \\
B2827 & 1a<o1 1b<o1 / 2a<1a2 2b*<1a2 & EE/DU & {-}{-}/{-}{-} & E &  \\
B2828 & 1a<o1 1b<o1 / 2a<1a2 2b*<1b1 & EE/DU & {-}{-}/{-}{-} & E &  \\
B2829 & 1a<o1 1b<o1 / 2a<1a2 2b*<1b2 & EE/DU & {-}{-}/{-}{-} & E &  \\
B2830 & 1a<o1 1b<o1 / 2a<1a2L 2b*<1a1 & EE/DU & {-}{-}/{-}{-} & E &  \\
B2831 & 1a<o1 1b<o1 / 2a<1a2L 2b*<1a2 & EE/DU & {-}{-}/{-}{-} & E &  \\
B2832 & 1a<o1 1b<o1 / 2a<1a2L 2b*<1b1 & EE/DU & {-}{-}/{-}{-} & E &  \\
B2833 & 1a<o1 1b<o1 / 2a<1a2L 2b*<1b2 & EE/DU & {-}{-}/{-}{-} & E &  \\
B2834 & 1a<o1 1b<o1L / 2a<1a0L 2b*<1a1 & LI/NU & {-}+/{-}{-} & B & CPW:1b \\
B2835 & 1a<o1 1b<o1L / 2a<1a0L 2b*<1a2 & LI/NU & {-}+/{-}{-} & B & CPW:1b \\
B2836 & 1a<o1 1b<o1L / 2a<1a0L 2b*<1b1 & LI/NU & {-}+/{-}{-} & B & CPW:1b \\
B2837 & 1a<o1 1b<o1L / 2a<1a0L 2b*<1b2 & LI/NU & {-}+/{-}{-} & B & CPW:1b \\
B2838 & 1a<o1 1b<o1L / 2a<1a1 2b*<1a1 & LI/DU & {-}+/{-}{-} & B & CPW:1b \\
B2839 & 1a<o1 1b<o1L / 2a<1a1 2b*<1a2 & LI/DU & {-}+/{-}{-} & B & CPW:1b \\
B2840 & 1a<o1 1b<o1L / 2a<1a1 2b*<1b1 & LI/DU & {-}+/{-}{-} & B & CPW:1b \\
B2841 & 1a<o1 1b<o1L / 2a<1a1 2b*<1b2 & LI/DU & {-}+/{-}{-} & B & CPW:1b \\
B2842 & 1a<o1 1b<o1L / 2a<1a1L 2b*<1a1 & LI/DU & {-}+/{-}{-} & B & CPW:1b \\
B2843 & 1a<o1 1b<o1L / 2a<1a1L 2b*<1a2 & LI/DU & {-}+/{-}{-} & B & CPW:1b \\
B2844 & 1a<o1 1b<o1L / 2a<1a1L 2b*<1b1 & LI/DU & {-}+/{-}{-} & B & CPW:1b \\
B2845 & 1a<o1 1b<o1L / 2a<1a1L 2b*<1b2 & LI/DU & {-}+/{-}{-} & B & CPW:1b \\
B2846 & 1a<o1 1b<o1L / 2a<1a2 2b*<1a1 & LI/DU & {-}+/{-}{-} & B & CPW:1b \\
B2847 & 1a<o1 1b<o1L / 2a<1a2 2b*<1a2 & LI/DU & {-}+/{-}{-} & B & CPW:1b \\
B2848 & 1a<o1 1b<o1L / 2a<1a2 2b*<1b1 & LI/DU & {-}+/{-}{-} & B & CPW:1b \\
B2849 & 1a<o1 1b<o1L / 2a<1a2 2b*<1b2 & LI/DU & {-}+/{-}{-} & B & CPW:1b \\
B2850 & 1a<o1 1b<o1L / 2a<1a2L 2b*<1a1 & LI/DU & {-}+/{-}{-} & B & CPW:1b \\
B2851 & 1a<o1 1b<o1L / 2a<1a2L 2b*<1a2 & LI/DU & {-}+/{-}{-} & B & CPW:1b \\
B2852 & 1a<o1 1b<o1L / 2a<1a2L 2b*<1b1 & LI/DU & {-}+/{-}{-} & B & CPW:1b \\
B2853 & 1a<o1 1b<o1L / 2a<1a2L 2b*<1b2 & LI/DU & {-}+/{-}{-} & B & CPW:1b \\
B2854 & 1a<o1 1b<o1L / 2a<1b0L 2b*<1a1 & LI/NU & {-}+/+{-} & B & CPW:1b \\
B2855 & 1a<o1 1b<o1L / 2a<1b0L 2b*<1a2 & LI/NU & {-}+/+{-} & B & CPW:1b \\
B2856 & 1a<o1 1b<o1L / 2a<1b0L 2b*<1b1 & LI/NU & {-}+/+{-} & B & CPW:1b \\
B2857 & 1a<o1 1b<o1L / 2a<1b0L 2b*<1b2 & LI/NU & {-}+/+{-} & B & CPW:1b \\
B2858 & 1a<o1 1b<o1L / 2a<1b1 2b*<1a1 & LI/IU & {-}+/{-}{-} & B & CPW:1b,2a \\
B2859 & 1a<o1 1b<o1L / 2a<1b1 2b*<1a2 & LI/IU & {-}+/{-}{-} & B & CPW:1b,2a \\
B2860 & 1a<o1 1b<o1L / 2a<1b1 2b*<1b1 & LI/IU & {-}+/{-}{-} & B & CPW:1b,2a \\
B2861 & 1a<o1 1b<o1L / 2a<1b1 2b*<1b2 & LI/IU & {-}+/{-}{-} & B & CPW:1b,2a \\
B2862 & 1a<o1 1b<o1L / 2a<1b1L 2b*<1a1 & LI/IU & {-}+/+{-} & B & CPW:1b,2a \\
B2863 & 1a<o1 1b<o1L / 2a<1b1L 2b*<1a2 & LI/IU & {-}+/+{-} & B & CPW:1b,2a \\
B2864 & 1a<o1 1b<o1L / 2a<1b1L 2b*<1b1 & LI/IU & {-}+/+{-} & B & CPW:1b,2a \\
B2865 & 1a<o1 1b<o1L / 2a<1b1L 2b*<1b2 & LI/IU & {-}+/+{-} & B & CPW:1b,2a \\
B2866 & 1a<o1 1b<o1L / 2a<1b2 2b*<1a1 & LI/IU & {-}+/{-}{-} & B & CPW:1b,2a \\
B2867 & 1a<o1 1b<o1L / 2a<1b2 2b*<1a2 & LI/IU & {-}+/{-}{-} & B & CPW:1b,2a \\
B2868 & 1a<o1 1b<o1L / 2a<1b2 2b*<1b1 & LI/IU & {-}+/{-}{-} & B & CPW:1b,2a \\
B2869 & 1a<o1 1b<o1L / 2a<1b2 2b*<1b2 & LI/IU & {-}+/{-}{-} & B & CPW:1b,2a \\
B2870 & 1a<o1 1b<o1L / 2a<1b2L 2b*<1a1 & LI/IU & {-}+/+{-} & B & CPW:1b,2a \\
B2871 & 1a<o1 1b<o1L / 2a<1b2L 2b*<1a2 & LI/IU & {-}+/+{-} & B & CPW:1b,2a \\
B2872 & 1a<o1 1b<o1L / 2a<1b2L 2b*<1b1 & LI/IU & {-}+/+{-} & B & CPW:1b,2a \\
B2873 & 1a<o1 1b<o1L / 2a<1b2L 2b*<1b2 & LI/IU & {-}+/+{-} & B & CPW:1b,2a \\
B2874 & 1a<o1 1b<o2 / 2a<1a0L 2b*<1a1 & EE/NU & {-}{-}/{-}{-} & E & C:1b \\
B2875 & 1a<o1 1b<o2 / 2a<1a0L 2b*<1a2 & EE/NU & {-}{-}/{-}{-} & E & C:1b \\
B2876 & 1a<o1 1b<o2 / 2a<1a0L 2b*<1b1 & EE/NU & {-}{-}/{-}{-} & E & C:1b \\
B2877 & 1a<o1 1b<o2 / 2a<1a0L 2b*<1b2 & EE/NU & {-}{-}/{-}{-} & E & C:1b \\
B2878 & 1a<o1 1b<o2 / 2a<1a1 2b*<1a1 & EE/DU & {-}{-}/{-}{-} & E & C:1b \\
B2879 & 1a<o1 1b<o2 / 2a<1a1 2b*<1a2 & EE/DU & {-}{-}/{-}{-} & E & C:1b \\
B2880 & 1a<o1 1b<o2 / 2a<1a1 2b*<1b1 & EE/DU & {-}{-}/{-}{-} & E & C:1b \\
B2881 & 1a<o1 1b<o2 / 2a<1a1 2b*<1b2 & EE/DU & {-}{-}/{-}{-} & E & C:1b \\
B2882 & 1a<o1 1b<o2 / 2a<1a1L 2b*<1a1 & EE/DU & {-}{-}/{-}{-} & E & C:1b \\
B2883 & 1a<o1 1b<o2 / 2a<1a1L 2b*<1a2 & EE/DU & {-}{-}/{-}{-} & E & C:1b \\
B2884 & 1a<o1 1b<o2 / 2a<1a1L 2b*<1b1 & EE/DU & {-}{-}/{-}{-} & E & C:1b \\
B2885 & 1a<o1 1b<o2 / 2a<1a1L 2b*<1b2 & EE/DU & {-}{-}/{-}{-} & E & C:1b \\
B2886 & 1a<o1 1b<o2 / 2a<1a2 2b*<1a1 & EE/DU & {-}{-}/{-}{-} & E & C:1b \\
B2887 & 1a<o1 1b<o2 / 2a<1a2 2b*<1a2 & EE/DU & {-}{-}/{-}{-} & E & C:1b \\
B2888 & 1a<o1 1b<o2 / 2a<1a2 2b*<1b1 & EE/DU & {-}{-}/{-}{-} & E & C:1b \\
B2889 & 1a<o1 1b<o2 / 2a<1a2 2b*<1b2 & EE/DU & {-}{-}/{-}{-} & E & C:1b \\
B2890 & 1a<o1 1b<o2 / 2a<1a2L 2b*<1a1 & EE/DU & {-}{-}/{-}{-} & E & C:1b \\
B2891 & 1a<o1 1b<o2 / 2a<1a2L 2b*<1a2 & EE/DU & {-}{-}/{-}{-} & E & C:1b \\
B2892 & 1a<o1 1b<o2 / 2a<1a2L 2b*<1b1 & EE/DU & {-}{-}/{-}{-} & E & C:1b \\
B2893 & 1a<o1 1b<o2 / 2a<1a2L 2b*<1b2 & EE/DU & {-}{-}/{-}{-} & E & C:1b \\
B2894 & 1a<o1 1b<o2 / 2a<1b0L 2b*<1a1 & EE/NU & {-}{-}/{-}{-} & E & C:1b \\
B2895 & 1a<o1 1b<o2 / 2a<1b0L 2b*<1a2 & EE/NU & {-}{-}/{-}{-} & E & C:1b \\
B2896 & 1a<o1 1b<o2 / 2a<1b0L 2b*<1b1 & EE/NU & {-}{-}/{-}{-} & E & C:1b \\
B2897 & 1a<o1 1b<o2 / 2a<1b0L 2b*<1b2 & EE/NU & {-}{-}/{-}{-} & E & C:1b \\
B2898 & 1a<o1 1b<o2 / 2a<1b1 2b*<1a1 & EE/DU & {-}{-}/{-}{-} & E & C:1b,2a \\
B2899 & 1a<o1 1b<o2 / 2a<1b1 2b*<1a2 & EE/DU & {-}{-}/{-}{-} & E & C:1b,2a \\
B2900 & 1a<o1 1b<o2 / 2a<1b1 2b*<1b1 & EE/DU & {-}{-}/{-}{-} & E & C:1b,2a \\
B2901 & 1a<o1 1b<o2 / 2a<1b1 2b*<1b2 & EE/DU & {-}{-}/{-}{-} & E & C:1b,2a \\
B2902 & 1a<o1 1b<o2 / 2a<1b1L 2b*<1a1 & EE/DU & {-}{-}/{-}{-} & E & C:1b,2a \\
B2903 & 1a<o1 1b<o2 / 2a<1b1L 2b*<1a2 & EE/DU & {-}{-}/{-}{-} & E & C:1b,2a \\
B2904 & 1a<o1 1b<o2 / 2a<1b1L 2b*<1b1 & EE/DU & {-}{-}/{-}{-} & E & C:1b,2a \\
B2905 & 1a<o1 1b<o2 / 2a<1b1L 2b*<1b2 & EE/DU & {-}{-}/{-}{-} & E & C:1b,2a \\
B2906 & 1a<o1 1b<o2 / 2a<1b2 2b*<1a1 & EE/DU & {-}{-}/{-}{-} & E & C:1b,2a \\
B2907 & 1a<o1 1b<o2 / 2a<1b2 2b*<1a2 & EE/DU & {-}{-}/{-}{-} & E & C:1b,2a \\
B2908 & 1a<o1 1b<o2 / 2a<1b2 2b*<1b1 & EE/DU & {-}{-}/{-}{-} & E & C:1b,2a \\
B2909 & 1a<o1 1b<o2 / 2a<1b2 2b*<1b2 & EE/DU & {-}{-}/{-}{-} & E & C:1b,2a \\
B2910 & 1a<o1 1b<o2 / 2a<1b2L 2b*<1a1 & EE/DU & {-}{-}/{-}{-} & E & C:1b,2a \\
B2911 & 1a<o1 1b<o2 / 2a<1b2L 2b*<1a2 & EE/DU & {-}{-}/{-}{-} & E & C:1b,2a \\
B2912 & 1a<o1 1b<o2 / 2a<1b2L 2b*<1b1 & EE/DU & {-}{-}/{-}{-} & E & C:1b,2a \\
B2913 & 1a<o1 1b<o2 / 2a<1b2L 2b*<1b2 & EE/DU & {-}{-}/{-}{-} & E & C:1b,2a \\
B2914 & 1a<o1 1b<o2L / 2a<1a0L 2b*<1a1 & LI/NU & {-}+/{-}{-} & B & PW:1b \\
B2915 & 1a<o1 1b<o2L / 2a<1a0L 2b*<1a2 & LI/NU & {-}+/{-}{-} & B & PW:1b \\
B2916 & 1a<o1 1b<o2L / 2a<1a0L 2b*<1b1 & LI/NU & {-}+/{-}{-} & B & PW:1b \\
B2917 & 1a<o1 1b<o2L / 2a<1a0L 2b*<1b2 & LI/NU & {-}+/{-}{-} & B & PW:1b \\
B2918 & 1a<o1 1b<o2L / 2a<1a1 2b*<1a1 & LI/DU & {-}+/{-}{-} & B & PW:1b \\
B2919 & 1a<o1 1b<o2L / 2a<1a1 2b*<1a2 & LI/DU & {-}+/{-}{-} & B & PW:1b \\
B2920 & 1a<o1 1b<o2L / 2a<1a1 2b*<1b1 & LI/DU & {-}+/{-}{-} & B & PW:1b \\
B2921 & 1a<o1 1b<o2L / 2a<1a1 2b*<1b2 & LI/DU & {-}+/{-}{-} & B & PW:1b \\
B2922 & 1a<o1 1b<o2L / 2a<1a1L 2b*<1a1 & LI/DU & {-}+/{-}{-} & B & PW:1b \\
B2923 & 1a<o1 1b<o2L / 2a<1a1L 2b*<1a2 & LI/DU & {-}+/{-}{-} & B & PW:1b \\
B2924 & 1a<o1 1b<o2L / 2a<1a1L 2b*<1b1 & LI/DU & {-}+/{-}{-} & B & PW:1b \\
B2925 & 1a<o1 1b<o2L / 2a<1a1L 2b*<1b2 & LI/DU & {-}+/{-}{-} & B & PW:1b \\
B2926 & 1a<o1 1b<o2L / 2a<1a2 2b*<1a1 & LI/DU & {-}+/{-}{-} & B & PW:1b \\
B2927 & 1a<o1 1b<o2L / 2a<1a2 2b*<1a2 & LI/DU & {-}+/{-}{-} & B & PW:1b \\
B2928 & 1a<o1 1b<o2L / 2a<1a2 2b*<1b1 & LI/DU & {-}+/{-}{-} & B & PW:1b \\
B2929 & 1a<o1 1b<o2L / 2a<1a2 2b*<1b2 & LI/DU & {-}+/{-}{-} & B & PW:1b \\
B2930 & 1a<o1 1b<o2L / 2a<1a2L 2b*<1a1 & LI/DU & {-}+/{-}{-} & B & PW:1b \\
B2931 & 1a<o1 1b<o2L / 2a<1a2L 2b*<1a2 & LI/DU & {-}+/{-}{-} & B & PW:1b \\
B2932 & 1a<o1 1b<o2L / 2a<1a2L 2b*<1b1 & LI/DU & {-}+/{-}{-} & B & PW:1b \\
B2933 & 1a<o1 1b<o2L / 2a<1a2L 2b*<1b2 & LI/DU & {-}+/{-}{-} & B & PW:1b \\
B2934 & 1a<o1 1b<o2L / 2a<1b0L 2b*<1a1 & LI/NU & {-}+/+{-} & B & PW:1b \\
B2935 & 1a<o1 1b<o2L / 2a<1b0L 2b*<1a2 & LI/NU & {-}+/+{-} & B & PW:1b \\
B2936 & 1a<o1 1b<o2L / 2a<1b0L 2b*<1b1 & LI/NU & {-}+/+{-} & B & PW:1b \\
B2937 & 1a<o1 1b<o2L / 2a<1b0L 2b*<1b2 & LI/NU & {-}+/+{-} & B & PW:1b \\
B2938 & 1a<o1 1b<o2L / 2a<1b1 2b*<1a1 & LI/IU & {-}+/{-}{-} & B & PW:1b,2a \\
B2939 & 1a<o1 1b<o2L / 2a<1b1 2b*<1a2 & LI/IU & {-}+/{-}{-} & B & PW:1b,2a \\
B2940 & 1a<o1 1b<o2L / 2a<1b1 2b*<1b1 & LI/IU & {-}+/{-}{-} & B & PW:1b,2a \\
B2941 & 1a<o1 1b<o2L / 2a<1b1 2b*<1b2 & LI/IU & {-}+/{-}{-} & B & PW:1b,2a \\
B2942 & 1a<o1 1b<o2L / 2a<1b1L 2b*<1a1 & LI/IU & {-}+/+{-} & B & PW:1b,2a \\
B2943 & 1a<o1 1b<o2L / 2a<1b1L 2b*<1a2 & LI/IU & {-}+/+{-} & B & PW:1b,2a \\
B2944 & 1a<o1 1b<o2L / 2a<1b1L 2b*<1b1 & LI/IU & {-}+/+{-} & B & PW:1b,2a \\
B2945 & 1a<o1 1b<o2L / 2a<1b1L 2b*<1b2 & LI/IU & {-}+/+{-} & B & PW:1b,2a \\
B2946 & 1a<o1 1b<o2L / 2a<1b2 2b*<1a1 & LI/IU & {-}+/{-}{-} & B & PW:1b,2a \\
B2947 & 1a<o1 1b<o2L / 2a<1b2 2b*<1a2 & LI/IU & {-}+/{-}{-} & B & PW:1b,2a \\
B2948 & 1a<o1 1b<o2L / 2a<1b2 2b*<1b1 & LI/IU & {-}+/{-}{-} & B & PW:1b,2a \\
B2949 & 1a<o1 1b<o2L / 2a<1b2 2b*<1b2 & LI/IU & {-}+/{-}{-} & B & PW:1b,2a \\
B2950 & 1a<o1 1b<o2L / 2a<1b2L 2b*<1a1 & LI/IU & {-}+/+{-} & B & PW:1b,2a \\
B2951 & 1a<o1 1b<o2L / 2a<1b2L 2b*<1a2 & LI/IU & {-}+/+{-} & B & PW:1b,2a \\
B2952 & 1a<o1 1b<o2L / 2a<1b2L 2b*<1b1 & LI/IU & {-}+/+{-} & B & PW:1b,2a \\
B2953 & 1a<o1 1b<o2L / 2a<1b2L 2b*<1b2 & LI/IU & {-}+/+{-} & B & PW:1b,2a \\
B2954 & 1a<o1 1b*<o1 / 2a<1a0L 2b<1a1 & IU/NI & {-}{-}/{-}{-} & N & W:1a,2b \\
B2955 & 1a<o1 1b*<o1 / 2a<1a0L 2b<1a2 & IU/NI & {-}{-}/{-}{-} & N & W:1a,2b \\
B2956 & 1a<o1 1b*<o1 / 2a<1a0L 2b<1b1 & IU/NU & {-}{-}/{-}{-} & N & W:1a \\
B2957 & 1a<o1 1b*<o1 / 2a<1a0L 2b<1b2 & IU/NU & {-}{-}/{-}{-} & N & W:1a \\
B2958 & 1a<o1 1b*<o1 / 2a<1a1 2b<1a1 & IU/EE & {-}{-}/{-}{-} & E & W:1a \\
B2959 & 1a<o1 1b*<o1 / 2a<1a1 2b<1a1L & IU/LI & {-}{-}/{-}{-} & B & CP:2b W:1a,2b \\
B2960 & 1a<o1 1b*<o1 / 2a<1a1 2b<1a2 & IU/EE & {-}{-}/{-}{-} & E & C:2b W:1a \\
B2961 & 1a<o1 1b*<o1 / 2a<1a1 2b<1a2L & IU/LI & {-}{-}/{-}{-} & B & P:2b W:1a,2b \\
B2962 & 1a<o1 1b*<o1 / 2a<1a1 2b<1b1 & IU/IU & {-}{-}/{-}{-} & N & W:1a,2a \\
B2963 & 1a<o1 1b*<o1 / 2a<1a1 2b<1b2 & IU/IU & {-}{-}/{-}{-} & N & W:1a,2a \\
B2964 & 1a<o1 1b*<o1 / 2a<1a1L 2b<1a2 & IU/IL & {-}{-}/{-}{-} & B & C:2a,2b P:2a W:1a,2a \\
B2965 & 1a<o1 1b*<o1 / 2a<1a1L 2b<1b1 & IU/IU & {-}{-}/{-}{-} & N & W:1a,2a \\
B2966 & 1a<o1 1b*<o1 / 2a<1a1L 2b<1b2 & IU/IU & {-}{-}/{-}{-} & N & W:1a,2a \\
B2967 & 1a<o1 1b*<o1 / 2a<1a2 2b<1a2 & IU/EE & {-}{-}/{-}{-} & E & W:1a \\
B2968 & 1a<o1 1b*<o1 / 2a<1a2 2b<1a2L & IU/LI & {-}{-}/{-}{-} & B & CP:2b W:1a,2b \\
B2969 & 1a<o1 1b*<o1 / 2a<1a2 2b<1b1 & IU/IU & {-}{-}/{-}{-} & N & W:1a,2a \\
B2970 & 1a<o1 1b*<o1 / 2a<1a2 2b<1b2 & IU/IU & {-}{-}/{-}{-} & N & W:1a,2a \\
B2971 & 1a<o1 1b*<o1 / 2a<1a2L 2b<1b1 & IU/IU & {-}{-}/{-}{-} & N & W:1a,2a \\
B2972 & 1a<o1 1b*<o1 / 2a<1a2L 2b<1b2 & IU/IU & {-}{-}/{-}{-} & N & W:1a,2a \\
B2973 & 1a<o1 1b*<o1 / 2a<1b1 2b<1b1 & IU/UU & {-}{-}/{-}{-} & N & W:1a \\
B2974 & 1a<o1 1b*<o1 / 2a<1b1 2b<1b2 & IU/UU & {-}{-}/{-}{-} & N & W:1a \\
B2975 & 1a<o1 1b*<o1 / 2a<1b2 2b<1b2 & IU/UU & {-}{-}/{-}{-} & N & W:1a \\
B2976 & 1a<o1 1b*<o2 / 2a<1a0L 2b<1a1 & IU/NI & {-}{-}/{-}{-} & N & W:1a,2b \\
B2977 & 1a<o1 1b*<o2 / 2a<1a0L 2b<1a2 & IU/NI & {-}{-}/{-}{-} & N & W:1a,2b \\
B2978 & 1a<o1 1b*<o2 / 2a<1a0L 2b<1b1 & IU/NU & {-}{-}/{-}{-} & N & W:1a \\
B2979 & 1a<o1 1b*<o2 / 2a<1a0L 2b<1b2 & IU/NU & {-}{-}/{-}{-} & N & W:1a \\
B2980 & 1a<o1 1b*<o2 / 2a<1a1 2b<1a1 & IU/EE & {-}{-}/{-}{-} & E & W:1a \\
B2981 & 1a<o1 1b*<o2 / 2a<1a1 2b<1a1L & IU/LI & {-}{-}/{-}{-} & B & CP:2b W:1a,2b \\
B2982 & 1a<o1 1b*<o2 / 2a<1a1 2b<1a2 & IU/EE & {-}{-}/{-}{-} & E & C:2b W:1a \\
B2983 & 1a<o1 1b*<o2 / 2a<1a1 2b<1a2L & IU/LI & {-}{-}/{-}{-} & B & P:2b W:1a,2b \\
B2984 & 1a<o1 1b*<o2 / 2a<1a1 2b<1b1 & IU/IU & {-}{-}/{-}{-} & N & W:1a,2a \\
B2985 & 1a<o1 1b*<o2 / 2a<1a1 2b<1b2 & IU/IU & {-}{-}/{-}{-} & N & W:1a,2a \\
B2986 & 1a<o1 1b*<o2 / 2a<1a1L 2b<1a2 & IU/IL & {-}{-}/{-}{-} & B & C:2a,2b P:2a W:1a,2a \\
B2987 & 1a<o1 1b*<o2 / 2a<1a1L 2b<1b1 & IU/IU & {-}{-}/{-}{-} & N & W:1a,2a \\
B2988 & 1a<o1 1b*<o2 / 2a<1a1L 2b<1b2 & IU/IU & {-}{-}/{-}{-} & N & W:1a,2a \\
B2989 & 1a<o1 1b*<o2 / 2a<1a2 2b<1a2 & IU/EE & {-}{-}/{-}{-} & E & W:1a \\
B2990 & 1a<o1 1b*<o2 / 2a<1a2 2b<1a2L & IU/LI & {-}{-}/{-}{-} & B & CP:2b W:1a,2b \\
B2991 & 1a<o1 1b*<o2 / 2a<1a2 2b<1b1 & IU/IU & {-}{-}/{-}{-} & N & W:1a,2a \\
B2992 & 1a<o1 1b*<o2 / 2a<1a2 2b<1b2 & IU/IU & {-}{-}/{-}{-} & N & W:1a,2a \\
B2993 & 1a<o1 1b*<o2 / 2a<1a2L 2b<1b1 & IU/IU & {-}{-}/{-}{-} & N & W:1a,2a \\
B2994 & 1a<o1 1b*<o2 / 2a<1a2L 2b<1b2 & IU/IU & {-}{-}/{-}{-} & N & W:1a,2a \\
B2995 & 1a<o1 1b*<o2 / 2a<1b1 2b<1b1 & IU/UU & {-}{-}/{-}{-} & N & W:1a \\
B2996 & 1a<o1 1b*<o2 / 2a<1b1 2b<1b2 & IU/UU & {-}{-}/{-}{-} & N & W:1a \\
B2997 & 1a<o1 1b*<o2 / 2a<1b2 2b<1b2 & IU/UU & {-}{-}/{-}{-} & N & W:1a \\
B2998 & 1a<o1L 1b<o2 / 2a<1a0L 2b*<1a1 & IL/NU & +{-}/+{-} & B & C:1a,1b PW:1a \\
B2999 & 1a<o1L 1b<o2 / 2a<1a0L 2b*<1a2 & IL/NU & +{-}/+{-} & B & C:1a,1b PW:1a \\
B3000 & 1a<o1L 1b<o2 / 2a<1a0L 2b*<1b1 & IL/NU & +{-}/+{-} & B & C:1a,1b PW:1a \\
B3001 & 1a<o1L 1b<o2 / 2a<1a0L 2b*<1b2 & IL/NU & +{-}/+{-} & B & C:1a,1b PW:1a \\
B3002 & 1a<o1L 1b<o2 / 2a<1a1 2b*<1a1 & IL/IU & +{-}/{-}{-} & B & C:1a,1b,2a PW:1a,2a \\
B3003 & 1a<o1L 1b<o2 / 2a<1a1 2b*<1a2 & IL/IU & +{-}/{-}{-} & B & C:1a,1b,2a PW:1a,2a \\
B3004 & 1a<o1L 1b<o2 / 2a<1a1 2b*<1b1 & IL/IU & +{-}/{-}{-} & B & C:1a,1b,2a PW:1a,2a \\
B3005 & 1a<o1L 1b<o2 / 2a<1a1 2b*<1b2 & IL/IU & +{-}/{-}{-} & B & C:1a,1b,2a PW:1a,2a \\
B3006 & 1a<o1L 1b<o2 / 2a<1a1L 2b*<1a1 & IL/IU & +{-}/+{-} & B & C:1a,1b,2a PW:1a,2a \\
B3007 & 1a<o1L 1b<o2 / 2a<1a1L 2b*<1a2 & IL/IU & +{-}/+{-} & B & C:1a,1b,2a PW:1a,2a \\
B3008 & 1a<o1L 1b<o2 / 2a<1a1L 2b*<1b1 & IL/IU & +{-}/+{-} & B & C:1a,1b,2a PW:1a,2a \\
B3009 & 1a<o1L 1b<o2 / 2a<1a1L 2b*<1b2 & IL/IU & +{-}/+{-} & B & C:1a,1b,2a PW:1a,2a \\
B3010 & 1a<o1L 1b<o2 / 2a<1a2 2b*<1a1 & IL/IU & +{-}/{-}{-} & B & C:1a,1b,2a PW:1a,2a \\
B3011 & 1a<o1L 1b<o2 / 2a<1a2 2b*<1a2 & IL/IU & +{-}/{-}{-} & B & C:1a,1b,2a PW:1a,2a \\
B3012 & 1a<o1L 1b<o2 / 2a<1a2 2b*<1b1 & IL/IU & +{-}/{-}{-} & B & C:1a,1b,2a PW:1a,2a \\
B3013 & 1a<o1L 1b<o2 / 2a<1a2 2b*<1b2 & IL/IU & +{-}/{-}{-} & B & C:1a,1b,2a PW:1a,2a \\
B3014 & 1a<o1L 1b<o2 / 2a<1a2L 2b*<1a1 & IL/IU & +{-}/+{-} & B & C:1a,1b,2a PW:1a,2a \\
B3015 & 1a<o1L 1b<o2 / 2a<1a2L 2b*<1a2 & IL/IU & +{-}/+{-} & B & C:1a,1b,2a PW:1a,2a \\
B3016 & 1a<o1L 1b<o2 / 2a<1a2L 2b*<1b1 & IL/IU & +{-}/+{-} & B & C:1a,1b,2a PW:1a,2a \\
B3017 & 1a<o1L 1b<o2 / 2a<1a2L 2b*<1b2 & IL/IU & +{-}/+{-} & B & C:1a,1b,2a PW:1a,2a \\
B3018 & 1a<o1L 1b<o2 / 2a<1b0L 2b*<1a1 & IL/NU & +{-}/{-}{-} & B & C:1a,1b PW:1a \\
B3019 & 1a<o1L 1b<o2 / 2a<1b0L 2b*<1a2 & IL/NU & +{-}/{-}{-} & B & C:1a,1b PW:1a \\
B3020 & 1a<o1L 1b<o2 / 2a<1b0L 2b*<1b1 & IL/NU & +{-}/{-}{-} & B & C:1a,1b PW:1a \\
B3021 & 1a<o1L 1b<o2 / 2a<1b0L 2b*<1b2 & IL/NU & +{-}/{-}{-} & B & C:1a,1b PW:1a \\
B3022 & 1a<o1L 1b<o2 / 2a<1b1 2b*<1a1 & IL/DU & +{-}/{-}{-} & B & C:1a,1b,2a PW:1a \\
B3023 & 1a<o1L 1b<o2 / 2a<1b1 2b*<1a2 & IL/DU & +{-}/{-}{-} & B & C:1a,1b,2a PW:1a \\
B3024 & 1a<o1L 1b<o2 / 2a<1b1 2b*<1b1 & IL/DU & +{-}/{-}{-} & B & C:1a,1b,2a PW:1a \\
B3025 & 1a<o1L 1b<o2 / 2a<1b1 2b*<1b2 & IL/DU & +{-}/{-}{-} & B & C:1a,1b,2a PW:1a \\
B3026 & 1a<o1L 1b<o2 / 2a<1b1L 2b*<1a1 & IL/DU & +{-}/{-}{-} & B & C:1a,1b,2a PW:1a \\
B3027 & 1a<o1L 1b<o2 / 2a<1b1L 2b*<1a2 & IL/DU & +{-}/{-}{-} & B & C:1a,1b,2a PW:1a \\
B3028 & 1a<o1L 1b<o2 / 2a<1b1L 2b*<1b1 & IL/DU & +{-}/{-}{-} & B & C:1a,1b,2a PW:1a \\
B3029 & 1a<o1L 1b<o2 / 2a<1b1L 2b*<1b2 & IL/DU & +{-}/{-}{-} & B & C:1a,1b,2a PW:1a \\
B3030 & 1a<o1L 1b<o2 / 2a<1b2 2b*<1a1 & IL/DU & +{-}/{-}{-} & B & C:1a,1b,2a PW:1a \\
B3031 & 1a<o1L 1b<o2 / 2a<1b2 2b*<1a2 & IL/DU & +{-}/{-}{-} & B & C:1a,1b,2a PW:1a \\
B3032 & 1a<o1L 1b<o2 / 2a<1b2 2b*<1b1 & IL/DU & +{-}/{-}{-} & B & C:1a,1b,2a PW:1a \\
B3033 & 1a<o1L 1b<o2 / 2a<1b2 2b*<1b2 & IL/DU & +{-}/{-}{-} & B & C:1a,1b,2a PW:1a \\
B3034 & 1a<o1L 1b<o2 / 2a<1b2L 2b*<1a1 & IL/DU & +{-}/{-}{-} & B & C:1a,1b,2a PW:1a \\
B3035 & 1a<o1L 1b<o2 / 2a<1b2L 2b*<1a2 & IL/DU & +{-}/{-}{-} & B & C:1a,1b,2a PW:1a \\
B3036 & 1a<o1L 1b<o2 / 2a<1b2L 2b*<1b1 & IL/DU & +{-}/{-}{-} & B & C:1a,1b,2a PW:1a \\
B3037 & 1a<o1L 1b<o2 / 2a<1b2L 2b*<1b2 & IL/DU & +{-}/{-}{-} & B & C:1a,1b,2a PW:1a \\
B3038 & 1a<o1L 1b*<o1 / 2a<1a0L 2b<1a1 & IU/NI & +{-}/+{-} & N & W:1a,2b \\
B3039 & 1a<o1L 1b*<o1 / 2a<1a0L 2b<1a2 & IU/NI & +{-}/+{-} & N & W:1a,2b \\
B3040 & 1a<o1L 1b*<o1 / 2a<1a0L 2b<1b1 & IU/NU & +{-}/+{-} & N & W:1a \\
B3041 & 1a<o1L 1b*<o1 / 2a<1a0L 2b<1b2 & IU/NU & +{-}/+{-} & N & W:1a \\
B3042 & 1a<o1L 1b*<o1 / 2a<1a1 2b<1a1 & IU/EE & +{-}/{-}{-} & E & W:1a \\
B3043 & 1a<o1L 1b*<o1 / 2a<1a1 2b<1a1L & IU/LI & +{-}/{-}+ & B & CP:2b W:1a,2b \\
B3044 & 1a<o1L 1b*<o1 / 2a<1a1 2b<1a2 & IU/EE & +{-}/{-}{-} & E & C:2b W:1a \\
B3045 & 1a<o1L 1b*<o1 / 2a<1a1 2b<1a2L & IU/LI & +{-}/{-}+ & B & P:2b W:1a,2b \\
B3046 & 1a<o1L 1b*<o1 / 2a<1a1 2b<1b1 & IU/IU & +{-}/{-}{-} & N & W:1a,2a \\
B3047 & 1a<o1L 1b*<o1 / 2a<1a1 2b<1b2 & IU/IU & +{-}/{-}{-} & N & W:1a,2a \\
B3048 & 1a<o1L 1b*<o1 / 2a<1a1L 2b<1a2 & IU/IL & +{-}/+{-} & B & C:2a,2b P:2a W:1a,2a \\
B3049 & 1a<o1L 1b*<o1 / 2a<1a1L 2b<1b1 & IU/IU & +{-}/+{-} & N & W:1a,2a \\
B3050 & 1a<o1L 1b*<o1 / 2a<1a1L 2b<1b2 & IU/IU & +{-}/+{-} & N & W:1a,2a \\
B3051 & 1a<o1L 1b*<o1 / 2a<1a2 2b<1a2 & IU/EE & +{-}/{-}{-} & E & W:1a \\
B3052 & 1a<o1L 1b*<o1 / 2a<1a2 2b<1a2L & IU/LI & +{-}/{-}+ & B & CP:2b W:1a,2b \\
B3053 & 1a<o1L 1b*<o1 / 2a<1a2 2b<1b1 & IU/IU & +{-}/{-}{-} & N & W:1a,2a \\
B3054 & 1a<o1L 1b*<o1 / 2a<1a2 2b<1b2 & IU/IU & +{-}/{-}{-} & N & W:1a,2a \\
B3055 & 1a<o1L 1b*<o1 / 2a<1a2L 2b<1b1 & IU/IU & +{-}/+{-} & N & W:1a,2a \\
B3056 & 1a<o1L 1b*<o1 / 2a<1a2L 2b<1b2 & IU/IU & +{-}/+{-} & N & W:1a,2a \\
B3057 & 1a<o1L 1b*<o1 / 2a<1b1 2b<1b1 & IU/UU & +{-}/{-}{-} & N & W:1a \\
B3058 & 1a<o1L 1b*<o1 / 2a<1b1 2b<1b2 & IU/UU & +{-}/{-}{-} & N & W:1a \\
B3059 & 1a<o1L 1b*<o1 / 2a<1b2 2b<1b2 & IU/UU & +{-}/{-}{-} & N & W:1a \\
B3060 & 1a<o1L 1b*<o2 / 2a<1a0L 2b<1a1 & IU/NI & +{-}/+{-} & N & W:1a,2b \\
B3061 & 1a<o1L 1b*<o2 / 2a<1a0L 2b<1a2 & IU/NI & +{-}/+{-} & N & W:1a,2b \\
B3062 & 1a<o1L 1b*<o2 / 2a<1a0L 2b<1b1 & IU/NU & +{-}/+{-} & N & W:1a \\
B3063 & 1a<o1L 1b*<o2 / 2a<1a0L 2b<1b2 & IU/NU & +{-}/+{-} & N & W:1a \\
B3064 & 1a<o1L 1b*<o2 / 2a<1a1 2b<1a1 & IU/EE & +{-}/{-}{-} & E & W:1a \\
B3065 & 1a<o1L 1b*<o2 / 2a<1a1 2b<1a1L & IU/LI & +{-}/{-}+ & B & CP:2b W:1a,2b \\
B3066 & 1a<o1L 1b*<o2 / 2a<1a1 2b<1a2 & IU/EE & +{-}/{-}{-} & E & C:2b W:1a \\
B3067 & 1a<o1L 1b*<o2 / 2a<1a1 2b<1a2L & IU/LI & +{-}/{-}+ & B & P:2b W:1a,2b \\
B3068 & 1a<o1L 1b*<o2 / 2a<1a1 2b<1b1 & IU/IU & +{-}/{-}{-} & N & W:1a,2a \\
B3069 & 1a<o1L 1b*<o2 / 2a<1a1 2b<1b2 & IU/IU & +{-}/{-}{-} & N & W:1a,2a \\
B3070 & 1a<o1L 1b*<o2 / 2a<1a1L 2b<1a2 & IU/IL & +{-}/+{-} & B & C:2a,2b P:2a W:1a,2a \\
B3071 & 1a<o1L 1b*<o2 / 2a<1a1L 2b<1b1 & IU/IU & +{-}/+{-} & N & W:1a,2a \\
B3072 & 1a<o1L 1b*<o2 / 2a<1a1L 2b<1b2 & IU/IU & +{-}/+{-} & N & W:1a,2a \\
B3073 & 1a<o1L 1b*<o2 / 2a<1a2 2b<1a2 & IU/EE & +{-}/{-}{-} & E & W:1a \\
B3074 & 1a<o1L 1b*<o2 / 2a<1a2 2b<1a2L & IU/LI & +{-}/{-}+ & B & CP:2b W:1a,2b \\
B3075 & 1a<o1L 1b*<o2 / 2a<1a2 2b<1b1 & IU/IU & +{-}/{-}{-} & N & W:1a,2a \\
B3076 & 1a<o1L 1b*<o2 / 2a<1a2 2b<1b2 & IU/IU & +{-}/{-}{-} & N & W:1a,2a \\
B3077 & 1a<o1L 1b*<o2 / 2a<1a2L 2b<1b1 & IU/IU & +{-}/+{-} & N & W:1a,2a \\
B3078 & 1a<o1L 1b*<o2 / 2a<1a2L 2b<1b2 & IU/IU & +{-}/+{-} & N & W:1a,2a \\
B3079 & 1a<o1L 1b*<o2 / 2a<1b1 2b<1b1 & IU/UU & +{-}/{-}{-} & N & W:1a \\
B3080 & 1a<o1L 1b*<o2 / 2a<1b1 2b<1b2 & IU/UU & +{-}/{-}{-} & N & W:1a \\
B3081 & 1a<o1L 1b*<o2 / 2a<1b2 2b<1b2 & IU/UU & +{-}/{-}{-} & N & W:1a \\
B3082 & 1a<o2 1b<o2 / 2a<1a0L 2b*<1a1 & EE/NU & {-}{-}/{-}{-} & E &  \\
B3083 & 1a<o2 1b<o2 / 2a<1a0L 2b*<1a2 & EE/NU & {-}{-}/{-}{-} & E &  \\
B3084 & 1a<o2 1b<o2 / 2a<1a0L 2b*<1b1 & EE/NU & {-}{-}/{-}{-} & E &  \\
B3085 & 1a<o2 1b<o2 / 2a<1a0L 2b*<1b2 & EE/NU & {-}{-}/{-}{-} & E &  \\
B3086 & 1a<o2 1b<o2 / 2a<1a1 2b*<1a1 & EE/DU & {-}{-}/{-}{-} & E &  \\
B3087 & 1a<o2 1b<o2 / 2a<1a1 2b*<1a2 & EE/DU & {-}{-}/{-}{-} & E &  \\
B3088 & 1a<o2 1b<o2 / 2a<1a1 2b*<1b1 & EE/DU & {-}{-}/{-}{-} & E &  \\
B3089 & 1a<o2 1b<o2 / 2a<1a1 2b*<1b2 & EE/DU & {-}{-}/{-}{-} & E &  \\
B3090 & 1a<o2 1b<o2 / 2a<1a1L 2b*<1a1 & EE/DU & {-}{-}/{-}{-} & E &  \\
B3091 & 1a<o2 1b<o2 / 2a<1a1L 2b*<1a2 & EE/DU & {-}{-}/{-}{-} & E &  \\
B3092 & 1a<o2 1b<o2 / 2a<1a1L 2b*<1b1 & EE/DU & {-}{-}/{-}{-} & E &  \\
B3093 & 1a<o2 1b<o2 / 2a<1a1L 2b*<1b2 & EE/DU & {-}{-}/{-}{-} & E &  \\
B3094 & 1a<o2 1b<o2 / 2a<1a2 2b*<1a1 & EE/DU & {-}{-}/{-}{-} & E &  \\
B3095 & 1a<o2 1b<o2 / 2a<1a2 2b*<1a2 & EE/DU & {-}{-}/{-}{-} & E &  \\
B3096 & 1a<o2 1b<o2 / 2a<1a2 2b*<1b1 & EE/DU & {-}{-}/{-}{-} & E &  \\
B3097 & 1a<o2 1b<o2 / 2a<1a2 2b*<1b2 & EE/DU & {-}{-}/{-}{-} & E &  \\
B3098 & 1a<o2 1b<o2 / 2a<1a2L 2b*<1a1 & EE/DU & {-}{-}/{-}{-} & E &  \\
B3099 & 1a<o2 1b<o2 / 2a<1a2L 2b*<1a2 & EE/DU & {-}{-}/{-}{-} & E &  \\
B3100 & 1a<o2 1b<o2 / 2a<1a2L 2b*<1b1 & EE/DU & {-}{-}/{-}{-} & E &  \\
B3101 & 1a<o2 1b<o2 / 2a<1a2L 2b*<1b2 & EE/DU & {-}{-}/{-}{-} & E &  \\
B3102 & 1a<o2 1b<o2L / 2a<1a0L 2b*<1a1 & LI/NU & {-}+/{-}{-} & B & CPW:1b \\
B3103 & 1a<o2 1b<o2L / 2a<1a0L 2b*<1a2 & LI/NU & {-}+/{-}{-} & B & CPW:1b \\
B3104 & 1a<o2 1b<o2L / 2a<1a0L 2b*<1b1 & LI/NU & {-}+/{-}{-} & B & CPW:1b \\
B3105 & 1a<o2 1b<o2L / 2a<1a0L 2b*<1b2 & LI/NU & {-}+/{-}{-} & B & CPW:1b \\
B3106 & 1a<o2 1b<o2L / 2a<1a1 2b*<1a1 & LI/DU & {-}+/{-}{-} & B & CPW:1b \\
B3107 & 1a<o2 1b<o2L / 2a<1a1 2b*<1a2 & LI/DU & {-}+/{-}{-} & B & CPW:1b \\
B3108 & 1a<o2 1b<o2L / 2a<1a1 2b*<1b1 & LI/DU & {-}+/{-}{-} & B & CPW:1b \\
B3109 & 1a<o2 1b<o2L / 2a<1a1 2b*<1b2 & LI/DU & {-}+/{-}{-} & B & CPW:1b \\
B3110 & 1a<o2 1b<o2L / 2a<1a1L 2b*<1a1 & LI/DU & {-}+/{-}{-} & B & CPW:1b \\
B3111 & 1a<o2 1b<o2L / 2a<1a1L 2b*<1a2 & LI/DU & {-}+/{-}{-} & B & CPW:1b \\
B3112 & 1a<o2 1b<o2L / 2a<1a1L 2b*<1b1 & LI/DU & {-}+/{-}{-} & B & CPW:1b \\
B3113 & 1a<o2 1b<o2L / 2a<1a1L 2b*<1b2 & LI/DU & {-}+/{-}{-} & B & CPW:1b \\
B3114 & 1a<o2 1b<o2L / 2a<1a2 2b*<1a1 & LI/DU & {-}+/{-}{-} & B & CPW:1b \\
B3115 & 1a<o2 1b<o2L / 2a<1a2 2b*<1a2 & LI/DU & {-}+/{-}{-} & B & CPW:1b \\
B3116 & 1a<o2 1b<o2L / 2a<1a2 2b*<1b1 & LI/DU & {-}+/{-}{-} & B & CPW:1b \\
B3117 & 1a<o2 1b<o2L / 2a<1a2 2b*<1b2 & LI/DU & {-}+/{-}{-} & B & CPW:1b \\
B3118 & 1a<o2 1b<o2L / 2a<1a2L 2b*<1a1 & LI/DU & {-}+/{-}{-} & B & CPW:1b \\
B3119 & 1a<o2 1b<o2L / 2a<1a2L 2b*<1a2 & LI/DU & {-}+/{-}{-} & B & CPW:1b \\
B3120 & 1a<o2 1b<o2L / 2a<1a2L 2b*<1b1 & LI/DU & {-}+/{-}{-} & B & CPW:1b \\
B3121 & 1a<o2 1b<o2L / 2a<1a2L 2b*<1b2 & LI/DU & {-}+/{-}{-} & B & CPW:1b \\
B3122 & 1a<o2 1b<o2L / 2a<1b0L 2b*<1a1 & LI/NU & {-}+/+{-} & B & CPW:1b \\
B3123 & 1a<o2 1b<o2L / 2a<1b0L 2b*<1a2 & LI/NU & {-}+/+{-} & B & CPW:1b \\
B3124 & 1a<o2 1b<o2L / 2a<1b0L 2b*<1b1 & LI/NU & {-}+/+{-} & B & CPW:1b \\
B3125 & 1a<o2 1b<o2L / 2a<1b0L 2b*<1b2 & LI/NU & {-}+/+{-} & B & CPW:1b \\
B3126 & 1a<o2 1b<o2L / 2a<1b1 2b*<1a1 & LI/IU & {-}+/{-}{-} & B & CPW:1b,2a \\
B3127 & 1a<o2 1b<o2L / 2a<1b1 2b*<1a2 & LI/IU & {-}+/{-}{-} & B & CPW:1b,2a \\
B3128 & 1a<o2 1b<o2L / 2a<1b1 2b*<1b1 & LI/IU & {-}+/{-}{-} & B & CPW:1b,2a \\
B3129 & 1a<o2 1b<o2L / 2a<1b1 2b*<1b2 & LI/IU & {-}+/{-}{-} & B & CPW:1b,2a \\
B3130 & 1a<o2 1b<o2L / 2a<1b1L 2b*<1a1 & LI/IU & {-}+/+{-} & B & CPW:1b,2a \\
B3131 & 1a<o2 1b<o2L / 2a<1b1L 2b*<1a2 & LI/IU & {-}+/+{-} & B & CPW:1b,2a \\
B3132 & 1a<o2 1b<o2L / 2a<1b1L 2b*<1b1 & LI/IU & {-}+/+{-} & B & CPW:1b,2a \\
B3133 & 1a<o2 1b<o2L / 2a<1b1L 2b*<1b2 & LI/IU & {-}+/+{-} & B & CPW:1b,2a \\
B3134 & 1a<o2 1b<o2L / 2a<1b2 2b*<1a1 & LI/IU & {-}+/{-}{-} & B & CPW:1b,2a \\
B3135 & 1a<o2 1b<o2L / 2a<1b2 2b*<1a2 & LI/IU & {-}+/{-}{-} & B & CPW:1b,2a \\
B3136 & 1a<o2 1b<o2L / 2a<1b2 2b*<1b1 & LI/IU & {-}+/{-}{-} & B & CPW:1b,2a \\
B3137 & 1a<o2 1b<o2L / 2a<1b2 2b*<1b2 & LI/IU & {-}+/{-}{-} & B & CPW:1b,2a \\
B3138 & 1a<o2 1b<o2L / 2a<1b2L 2b*<1a1 & LI/IU & {-}+/+{-} & B & CPW:1b,2a \\
B3139 & 1a<o2 1b<o2L / 2a<1b2L 2b*<1a2 & LI/IU & {-}+/+{-} & B & CPW:1b,2a \\
B3140 & 1a<o2 1b<o2L / 2a<1b2L 2b*<1b1 & LI/IU & {-}+/+{-} & B & CPW:1b,2a \\
B3141 & 1a<o2 1b<o2L / 2a<1b2L 2b*<1b2 & LI/IU & {-}+/+{-} & B & CPW:1b,2a \\
B3142 & 1a<o2 1b*<o1 / 2a<1a0L 2b<1a1 & IU/NI & {-}{-}/{-}{-} & N & W:1a,2b \\
B3143 & 1a<o2 1b*<o1 / 2a<1a0L 2b<1a2 & IU/NI & {-}{-}/{-}{-} & N & W:1a,2b \\
B3144 & 1a<o2 1b*<o1 / 2a<1a0L 2b<1b1 & IU/NU & {-}{-}/{-}{-} & N & W:1a \\
B3145 & 1a<o2 1b*<o1 / 2a<1a0L 2b<1b2 & IU/NU & {-}{-}/{-}{-} & N & W:1a \\
B3146 & 1a<o2 1b*<o1 / 2a<1a1 2b<1a1 & IU/EE & {-}{-}/{-}{-} & E & W:1a \\
B3147 & 1a<o2 1b*<o1 / 2a<1a1 2b<1a1L & IU/LI & {-}{-}/{-}{-} & B & CP:2b W:1a,2b \\
B3148 & 1a<o2 1b*<o1 / 2a<1a1 2b<1a2 & IU/EE & {-}{-}/{-}{-} & E & C:2b W:1a \\
B3149 & 1a<o2 1b*<o1 / 2a<1a1 2b<1a2L & IU/LI & {-}{-}/{-}{-} & B & P:2b W:1a,2b \\
B3150 & 1a<o2 1b*<o1 / 2a<1a1 2b<1b1 & IU/IU & {-}{-}/{-}{-} & N & W:1a,2a \\
B3151 & 1a<o2 1b*<o1 / 2a<1a1 2b<1b2 & IU/IU & {-}{-}/{-}{-} & N & W:1a,2a \\
B3152 & 1a<o2 1b*<o1 / 2a<1a1L 2b<1a2 & IU/IL & {-}{-}/{-}{-} & B & C:2a,2b P:2a W:1a,2a \\
B3153 & 1a<o2 1b*<o1 / 2a<1a1L 2b<1b1 & IU/IU & {-}{-}/{-}{-} & N & W:1a,2a \\
B3154 & 1a<o2 1b*<o1 / 2a<1a1L 2b<1b2 & IU/IU & {-}{-}/{-}{-} & N & W:1a,2a \\
B3155 & 1a<o2 1b*<o1 / 2a<1a2 2b<1a2 & IU/EE & {-}{-}/{-}{-} & E & W:1a \\
B3156 & 1a<o2 1b*<o1 / 2a<1a2 2b<1a2L & IU/LI & {-}{-}/{-}{-} & B & CP:2b W:1a,2b \\
B3157 & 1a<o2 1b*<o1 / 2a<1a2 2b<1b1 & IU/IU & {-}{-}/{-}{-} & N & W:1a,2a \\
B3158 & 1a<o2 1b*<o1 / 2a<1a2 2b<1b2 & IU/IU & {-}{-}/{-}{-} & N & W:1a,2a \\
B3159 & 1a<o2 1b*<o1 / 2a<1a2L 2b<1b1 & IU/IU & {-}{-}/{-}{-} & N & W:1a,2a \\
B3160 & 1a<o2 1b*<o1 / 2a<1a2L 2b<1b2 & IU/IU & {-}{-}/{-}{-} & N & W:1a,2a \\
B3161 & 1a<o2 1b*<o1 / 2a<1b1 2b<1b1 & IU/UU & {-}{-}/{-}{-} & N & W:1a \\
B3162 & 1a<o2 1b*<o1 / 2a<1b1 2b<1b2 & IU/UU & {-}{-}/{-}{-} & N & W:1a \\
B3163 & 1a<o2 1b*<o1 / 2a<1b2 2b<1b2 & IU/UU & {-}{-}/{-}{-} & N & W:1a \\
B3164 & 1a<o2 1b*<o2 / 2a<1a0L 2b<1a1 & IU/NI & {-}{-}/{-}{-} & N & W:1a,2b \\
B3165 & 1a<o2 1b*<o2 / 2a<1a0L 2b<1a2 & IU/NI & {-}{-}/{-}{-} & N & W:1a,2b \\
B3166 & 1a<o2 1b*<o2 / 2a<1a0L 2b<1b1 & IU/NU & {-}{-}/{-}{-} & N & W:1a \\
B3167 & 1a<o2 1b*<o2 / 2a<1a0L 2b<1b2 & IU/NU & {-}{-}/{-}{-} & N & W:1a \\
B3168 & 1a<o2 1b*<o2 / 2a<1a1 2b<1a1 & IU/EE & {-}{-}/{-}{-} & E & W:1a \\
B3169 & 1a<o2 1b*<o2 / 2a<1a1 2b<1a1L & IU/LI & {-}{-}/{-}{-} & B & CP:2b W:1a,2b \\
B3170 & 1a<o2 1b*<o2 / 2a<1a1 2b<1a2 & IU/EE & {-}{-}/{-}{-} & E & C:2b W:1a \\
B3171 & 1a<o2 1b*<o2 / 2a<1a1 2b<1a2L & IU/LI & {-}{-}/{-}{-} & B & P:2b W:1a,2b \\
B3172 & 1a<o2 1b*<o2 / 2a<1a1 2b<1b1 & IU/IU & {-}{-}/{-}{-} & N & W:1a,2a \\
B3173 & 1a<o2 1b*<o2 / 2a<1a1 2b<1b2 & IU/IU & {-}{-}/{-}{-} & N & W:1a,2a \\
B3174 & 1a<o2 1b*<o2 / 2a<1a1L 2b<1a2 & IU/IL & {-}{-}/{-}{-} & B & C:2a,2b P:2a W:1a,2a \\
B3175 & 1a<o2 1b*<o2 / 2a<1a1L 2b<1b1 & IU/IU & {-}{-}/{-}{-} & N & W:1a,2a \\
B3176 & 1a<o2 1b*<o2 / 2a<1a1L 2b<1b2 & IU/IU & {-}{-}/{-}{-} & N & W:1a,2a \\
B3177 & 1a<o2 1b*<o2 / 2a<1a2 2b<1a2 & IU/EE & {-}{-}/{-}{-} & E & W:1a \\
B3178 & 1a<o2 1b*<o2 / 2a<1a2 2b<1a2L & IU/LI & {-}{-}/{-}{-} & B & CP:2b W:1a,2b \\
B3179 & 1a<o2 1b*<o2 / 2a<1a2 2b<1b1 & IU/IU & {-}{-}/{-}{-} & N & W:1a,2a \\
B3180 & 1a<o2 1b*<o2 / 2a<1a2 2b<1b2 & IU/IU & {-}{-}/{-}{-} & N & W:1a,2a \\
B3181 & 1a<o2 1b*<o2 / 2a<1a2L 2b<1b1 & IU/IU & {-}{-}/{-}{-} & N & W:1a,2a \\
B3182 & 1a<o2 1b*<o2 / 2a<1a2L 2b<1b2 & IU/IU & {-}{-}/{-}{-} & N & W:1a,2a \\
B3183 & 1a<o2 1b*<o2 / 2a<1b1 2b<1b1 & IU/UU & {-}{-}/{-}{-} & N & W:1a \\
B3184 & 1a<o2 1b*<o2 / 2a<1b1 2b<1b2 & IU/UU & {-}{-}/{-}{-} & N & W:1a \\
B3185 & 1a<o2 1b*<o2 / 2a<1b2 2b<1b2 & IU/UU & {-}{-}/{-}{-} & N & W:1a \\
B3186 & 1a<o2L 1b*<o1 / 2a<1a0L 2b<1a1 & IU/NI & +{-}/+{-} & N & W:1a,2b \\
B3187 & 1a<o2L 1b*<o1 / 2a<1a0L 2b<1a2 & IU/NI & +{-}/+{-} & N & W:1a,2b \\
B3188 & 1a<o2L 1b*<o1 / 2a<1a0L 2b<1b1 & IU/NU & +{-}/+{-} & N & W:1a \\
B3189 & 1a<o2L 1b*<o1 / 2a<1a0L 2b<1b2 & IU/NU & +{-}/+{-} & N & W:1a \\
B3190 & 1a<o2L 1b*<o1 / 2a<1a1 2b<1a1 & IU/EE & +{-}/{-}{-} & E & W:1a \\
B3191 & 1a<o2L 1b*<o1 / 2a<1a1 2b<1a1L & IU/LI & +{-}/{-}+ & B & CP:2b W:1a,2b \\
B3192 & 1a<o2L 1b*<o1 / 2a<1a1 2b<1a2 & IU/EE & +{-}/{-}{-} & E & C:2b W:1a \\
B3193 & 1a<o2L 1b*<o1 / 2a<1a1 2b<1a2L & IU/LI & +{-}/{-}+ & B & P:2b W:1a,2b \\
B3194 & 1a<o2L 1b*<o1 / 2a<1a1 2b<1b1 & IU/IU & +{-}/{-}{-} & N & W:1a,2a \\
B3195 & 1a<o2L 1b*<o1 / 2a<1a1 2b<1b2 & IU/IU & +{-}/{-}{-} & N & W:1a,2a \\
B3196 & 1a<o2L 1b*<o1 / 2a<1a1L 2b<1a2 & IU/IL & +{-}/+{-} & B & C:2a,2b P:2a W:1a,2a \\
B3197 & 1a<o2L 1b*<o1 / 2a<1a1L 2b<1b1 & IU/IU & +{-}/+{-} & N & W:1a,2a \\
B3198 & 1a<o2L 1b*<o1 / 2a<1a1L 2b<1b2 & IU/IU & +{-}/+{-} & N & W:1a,2a \\
B3199 & 1a<o2L 1b*<o1 / 2a<1a2 2b<1a2 & IU/EE & +{-}/{-}{-} & E & W:1a \\
B3200 & 1a<o2L 1b*<o1 / 2a<1a2 2b<1a2L & IU/LI & +{-}/{-}+ & B & CP:2b W:1a,2b \\
B3201 & 1a<o2L 1b*<o1 / 2a<1a2 2b<1b1 & IU/IU & +{-}/{-}{-} & N & W:1a,2a \\
B3202 & 1a<o2L 1b*<o1 / 2a<1a2 2b<1b2 & IU/IU & +{-}/{-}{-} & N & W:1a,2a \\
B3203 & 1a<o2L 1b*<o1 / 2a<1a2L 2b<1b1 & IU/IU & +{-}/+{-} & N & W:1a,2a \\
B3204 & 1a<o2L 1b*<o1 / 2a<1a2L 2b<1b2 & IU/IU & +{-}/+{-} & N & W:1a,2a \\
B3205 & 1a<o2L 1b*<o1 / 2a<1b1 2b<1b1 & IU/UU & +{-}/{-}{-} & N & W:1a \\
B3206 & 1a<o2L 1b*<o1 / 2a<1b1 2b<1b2 & IU/UU & +{-}/{-}{-} & N & W:1a \\
B3207 & 1a<o2L 1b*<o1 / 2a<1b2 2b<1b2 & IU/UU & +{-}/{-}{-} & N & W:1a \\
B3208 & 1a<o2L 1b*<o2 / 2a<1a0L 2b<1a1 & IU/NI & +{-}/+{-} & N & W:1a,2b \\
B3209 & 1a<o2L 1b*<o2 / 2a<1a0L 2b<1a2 & IU/NI & +{-}/+{-} & N & W:1a,2b \\
B3210 & 1a<o2L 1b*<o2 / 2a<1a0L 2b<1b1 & IU/NU & +{-}/+{-} & N & W:1a \\
B3211 & 1a<o2L 1b*<o2 / 2a<1a0L 2b<1b2 & IU/NU & +{-}/+{-} & N & W:1a \\
B3212 & 1a<o2L 1b*<o2 / 2a<1a1 2b<1a1 & IU/EE & +{-}/{-}{-} & E & W:1a \\
B3213 & 1a<o2L 1b*<o2 / 2a<1a1 2b<1a1L & IU/LI & +{-}/{-}+ & B & CP:2b W:1a,2b \\
B3214 & 1a<o2L 1b*<o2 / 2a<1a1 2b<1a2 & IU/EE & +{-}/{-}{-} & E & C:2b W:1a \\
B3215 & 1a<o2L 1b*<o2 / 2a<1a1 2b<1a2L & IU/LI & +{-}/{-}+ & B & P:2b W:1a,2b \\
B3216 & 1a<o2L 1b*<o2 / 2a<1a1 2b<1b1 & IU/IU & +{-}/{-}{-} & N & W:1a,2a \\
B3217 & 1a<o2L 1b*<o2 / 2a<1a1 2b<1b2 & IU/IU & +{-}/{-}{-} & N & W:1a,2a \\
B3218 & 1a<o2L 1b*<o2 / 2a<1a1L 2b<1a2 & IU/IL & +{-}/+{-} & B & C:2a,2b P:2a W:1a,2a \\
B3219 & 1a<o2L 1b*<o2 / 2a<1a1L 2b<1b1 & IU/IU & +{-}/+{-} & N & W:1a,2a \\
B3220 & 1a<o2L 1b*<o2 / 2a<1a1L 2b<1b2 & IU/IU & +{-}/+{-} & N & W:1a,2a \\
B3221 & 1a<o2L 1b*<o2 / 2a<1a2 2b<1a2 & IU/EE & +{-}/{-}{-} & E & W:1a \\
B3222 & 1a<o2L 1b*<o2 / 2a<1a2 2b<1a2L & IU/LI & +{-}/{-}+ & B & CP:2b W:1a,2b \\
B3223 & 1a<o2L 1b*<o2 / 2a<1a2 2b<1b1 & IU/IU & +{-}/{-}{-} & N & W:1a,2a \\
B3224 & 1a<o2L 1b*<o2 / 2a<1a2 2b<1b2 & IU/IU & +{-}/{-}{-} & N & W:1a,2a \\
B3225 & 1a<o2L 1b*<o2 / 2a<1a2L 2b<1b1 & IU/IU & +{-}/+{-} & N & W:1a,2a \\
B3226 & 1a<o2L 1b*<o2 / 2a<1a2L 2b<1b2 & IU/IU & +{-}/+{-} & N & W:1a,2a \\
B3227 & 1a<o2L 1b*<o2 / 2a<1b1 2b<1b1 & IU/UU & +{-}/{-}{-} & N & W:1a \\
B3228 & 1a<o2L 1b*<o2 / 2a<1b1 2b<1b2 & IU/UU & +{-}/{-}{-} & N & W:1a \\
B3229 & 1a<o2L 1b*<o2 / 2a<1b2 2b<1b2 & IU/UU & +{-}/{-}{-} & N & W:1a \\
B3230 & 1a<o0L 1b<o1 / 2a<1a0L 2b*<1a1,1b1 & NI/NU & +{-}/+{-} & N & W:2b \\
B3231 & 1a<o0L 1b<o1 / 2a<1a0L 2b*<1a1,1b2 & NI/NU & +{-}/+{-} & N & W:2b \\
B3232 & 1a<o0L 1b<o1 / 2a<1a0L 2b*<1a2,1b1 & NI/NU & +{-}/+{-} & N & W:2b \\
B3233 & 1a<o0L 1b<o1 / 2a<1a0L 2b*<1a2,1b2 & NI/NU & +{-}/+{-} & N & W:2b \\
B3234 & 1a<o0L 1b<o1 / 2a<1a0L,1b1 2b*<1a1 & NI/IN & +{-}/+{-} & N &  \\
B3235 & 1a<o0L 1b<o1 / 2a<1a0L,1b1 2b*<1a2 & NI/IN & +{-}/+{-} & N &  \\
B3236 & 1a<o0L 1b<o1 / 2a<1a0L,1b1 2b*<1b1 & NI/IU & +{-}/+{-} & N & W:1b,2a \\
B3237 & 1a<o0L 1b<o1 / 2a<1a0L,1b1 2b*<1b2 & NI/IU & +{-}/+{-} & N & W:1b,2a \\
B3238 & 1a<o0L 1b<o1 / 2a<1a0L,1b2 2b*<1a1 & NI/IN & +{-}/+{-} & N &  \\
B3239 & 1a<o0L 1b<o1 / 2a<1a0L,1b2 2b*<1a2 & NI/IN & +{-}/+{-} & N &  \\
B3240 & 1a<o0L 1b<o1 / 2a<1a0L,1b2 2b*<1b1 & NI/IU & +{-}/+{-} & N & W:1b,2a \\
B3241 & 1a<o0L 1b<o1 / 2a<1a0L,1b2 2b*<1b2 & NI/IU & +{-}/+{-} & N & W:1b,2a \\
B3242 & 1a<o0L 1b<o1 / 2a<1a1 2b*<1a1,1b1 & NI/NU & +{-}/{-}{-} & N & W:2b \\
B3243 & 1a<o0L 1b<o1 / 2a<1a1 2b*<1a1,1b2 & NI/NU & +{-}/{-}{-} & N & W:2b \\
B3244 & 1a<o0L 1b<o1 / 2a<1a1 2b*<1a2,1b1 & NI/NU & +{-}/{-}{-} & N & W:2b \\
B3245 & 1a<o0L 1b<o1 / 2a<1a1 2b*<1a2,1b2 & NI/NU & +{-}/{-}{-} & N & W:2b \\
B3246 & 1a<o0L 1b<o1 / 2a<1a1,1b0L 2b*<1a1 & NI/NN & +{-}/{-}{-} & N &  \\
B3247 & 1a<o0L 1b<o1 / 2a<1a1,1b0L 2b*<1a2 & NI/NN & +{-}/{-}{-} & N &  \\
B3248 & 1a<o0L 1b<o1 / 2a<1a1,1b0L 2b*<1b1 & NI/NU & +{-}/{-}{-} & N & W:2b \\
B3249 & 1a<o0L 1b<o1 / 2a<1a1,1b0L 2b*<1b2 & NI/NU & +{-}/{-}{-} & N & W:2b \\
B3250 & 1a<o0L 1b<o1 / 2a<1a1,1b1 2b*<1a1 & NI/EN & +{-}/{-}{-} & F & W:2a \\
B3251 & 1a<o0L 1b<o1 / 2a<1a1,1b1 2b*<1a2 & NI/EN & +{-}/{-}{-} & F & W:2a \\
B3252 & 1a<o0L 1b<o1 / 2a<1a1,1b1 2b*<1b1 & NI/EU & +{-}/{-}{-} & F & W:1b \\
B3253 & 1a<o0L 1b<o1 / 2a<1a1,1b1 2b*<1b2 & NI/EU & +{-}/{-}{-} & F & W:1b \\
B3254 & 1a<o0L 1b<o1 / 2a<1a1,1b1L 2b*<1a1 & NI/IN & +{-}/{-}{-} & F & CP:2a \\
B3255 & 1a<o0L 1b<o1 / 2a<1a1,1b1L 2b*<1a2 & NI/IN & +{-}/{-}{-} & F & CP:2a \\
B3256 & 1a<o0L 1b<o1 / 2a<1a1,1b1L 2b*<1b1 & NI/IU & +{-}/{-}{-} & F & CP:2a W:1b,2a \\
B3257 & 1a<o0L 1b<o1 / 2a<1a1,1b1L 2b*<1b2 & NI/IU & +{-}/{-}{-} & F & CP:2a W:1b,2a \\
B3258 & 1a<o0L 1b<o1 / 2a<1a1,1b2 2b*<1a1 & NI/EN & +{-}/{-}{-} & F & CW:2a \\
B3259 & 1a<o0L 1b<o1 / 2a<1a1,1b2 2b*<1a2 & NI/EN & +{-}/{-}{-} & F & CW:2a \\
B3260 & 1a<o0L 1b<o1 / 2a<1a1,1b2 2b*<1b1 & NI/EU & +{-}/{-}{-} & F & C:2a W:1b \\
B3261 & 1a<o0L 1b<o1 / 2a<1a1,1b2 2b*<1b2 & NI/EU & +{-}/{-}{-} & F & C:2a W:1b \\
B3262 & 1a<o0L 1b<o1 / 2a<1a1,1b2L 2b*<1a1 & NI/IN & +{-}/{-}{-} & F & P:2a \\
B3263 & 1a<o0L 1b<o1 / 2a<1a1,1b2L 2b*<1a2 & NI/IN & +{-}/{-}{-} & F & P:2a \\
B3264 & 1a<o0L 1b<o1 / 2a<1a1,1b2L 2b*<1b1 & NI/IU & +{-}/{-}{-} & F & P:2a W:1b,2a \\
B3265 & 1a<o0L 1b<o1 / 2a<1a1,1b2L 2b*<1b2 & NI/IU & +{-}/{-}{-} & F & P:2a W:1b,2a \\
B3266 & 1a<o0L 1b<o1 / 2a<1a1L 2b*<1a1,1b1 & NI/NU & +{-}/+{-} & N & W:2b \\
B3267 & 1a<o0L 1b<o1 / 2a<1a1L 2b*<1a1,1b2 & NI/NU & +{-}/+{-} & N & W:2b \\
B3268 & 1a<o0L 1b<o1 / 2a<1a1L 2b*<1a2,1b1 & NI/NU & +{-}/+{-} & N & W:2b \\
B3269 & 1a<o0L 1b<o1 / 2a<1a1L 2b*<1a2,1b2 & NI/NU & +{-}/+{-} & N & W:2b \\
B3270 & 1a<o0L 1b<o1 / 2a<1a1L,1b1 2b*<1a1 & NI/EN & +{-}/+{-} & F & W:2a \\
B3271 & 1a<o0L 1b<o1 / 2a<1a1L,1b1 2b*<1a2 & NI/EN & +{-}/+{-} & F & W:2a \\
B3272 & 1a<o0L 1b<o1 / 2a<1a1L,1b1 2b*<1b1 & NI/EU & +{-}/+{-} & F & W:1b \\
B3273 & 1a<o0L 1b<o1 / 2a<1a1L,1b1 2b*<1b2 & NI/EU & +{-}/+{-} & F & W:1b \\
B3274 & 1a<o0L 1b<o1 / 2a<1a1L,1b2 2b*<1a1 & NI/EN & +{-}/+{-} & F & CW:2a \\
B3275 & 1a<o0L 1b<o1 / 2a<1a1L,1b2 2b*<1a2 & NI/EN & +{-}/+{-} & F & CW:2a \\
B3276 & 1a<o0L 1b<o1 / 2a<1a1L,1b2 2b*<1b1 & NI/EU & +{-}/+{-} & F & C:2a W:1b \\
B3277 & 1a<o0L 1b<o1 / 2a<1a1L,1b2 2b*<1b2 & NI/EU & +{-}/+{-} & F & C:2a W:1b \\
B3278 & 1a<o0L 1b<o1 / 2a<1a2 2b*<1a1,1b1 & NI/NU & +{-}/{-}{-} & N & W:2b \\
B3279 & 1a<o0L 1b<o1 / 2a<1a2 2b*<1a1,1b2 & NI/NU & +{-}/{-}{-} & N & W:2b \\
B3280 & 1a<o0L 1b<o1 / 2a<1a2 2b*<1a2,1b1 & NI/NU & +{-}/{-}{-} & N & W:2b \\
B3281 & 1a<o0L 1b<o1 / 2a<1a2 2b*<1a2,1b2 & NI/NU & +{-}/{-}{-} & N & W:2b \\
B3282 & 1a<o0L 1b<o1 / 2a<1a2,1b0L 2b*<1a1 & NI/NN & +{-}/{-}{-} & N &  \\
B3283 & 1a<o0L 1b<o1 / 2a<1a2,1b0L 2b*<1a2 & NI/NN & +{-}/{-}{-} & N &  \\
B3284 & 1a<o0L 1b<o1 / 2a<1a2,1b0L 2b*<1b1 & NI/NU & +{-}/{-}{-} & N & W:2b \\
B3285 & 1a<o0L 1b<o1 / 2a<1a2,1b0L 2b*<1b2 & NI/NU & +{-}/{-}{-} & N & W:2b \\
B3286 & 1a<o0L 1b<o1 / 2a<1a2,1b1 2b*<1a1 & NI/EN & +{-}/{-}{-} & F & W:2a \\
B3287 & 1a<o0L 1b<o1 / 2a<1a2,1b1 2b*<1a2 & NI/EN & +{-}/{-}{-} & F & W:2a \\
B3288 & 1a<o0L 1b<o1 / 2a<1a2,1b1 2b*<1b1 & NI/EU & +{-}/{-}{-} & F & W:1b \\
B3289 & 1a<o0L 1b<o1 / 2a<1a2,1b1 2b*<1b2 & NI/EU & +{-}/{-}{-} & F & W:1b \\
B3290 & 1a<o0L 1b<o1 / 2a<1a2,1b1L 2b*<1a1 & NI/IN & +{-}/{-}{-} & F & CP:2a \\
B3291 & 1a<o0L 1b<o1 / 2a<1a2,1b1L 2b*<1a2 & NI/IN & +{-}/{-}{-} & F & CP:2a \\
B3292 & 1a<o0L 1b<o1 / 2a<1a2,1b1L 2b*<1b1 & NI/IU & +{-}/{-}{-} & F & CP:2a W:1b,2a \\
B3293 & 1a<o0L 1b<o1 / 2a<1a2,1b1L 2b*<1b2 & NI/IU & +{-}/{-}{-} & F & CP:2a W:1b,2a \\
B3294 & 1a<o0L 1b<o1 / 2a<1a2,1b2 2b*<1a1 & NI/EN & +{-}/{-}{-} & F & W:2a \\
B3295 & 1a<o0L 1b<o1 / 2a<1a2,1b2 2b*<1a2 & NI/EN & +{-}/{-}{-} & F & W:2a \\
B3296 & 1a<o0L 1b<o1 / 2a<1a2,1b2 2b*<1b1 & NI/EU & +{-}/{-}{-} & F & W:1b \\
B3297 & 1a<o0L 1b<o1 / 2a<1a2,1b2 2b*<1b2 & NI/EU & +{-}/{-}{-} & F & W:1b \\
B3298 & 1a<o0L 1b<o1 / 2a<1a2,1b2L 2b*<1a1 & NI/IN & +{-}/{-}{-} & F & CP:2a \\
B3299 & 1a<o0L 1b<o1 / 2a<1a2,1b2L 2b*<1a2 & NI/IN & +{-}/{-}{-} & F & CP:2a \\
B3300 & 1a<o0L 1b<o1 / 2a<1a2,1b2L 2b*<1b1 & NI/IU & +{-}/{-}{-} & F & CP:2a W:1b,2a \\
B3301 & 1a<o0L 1b<o1 / 2a<1a2,1b2L 2b*<1b2 & NI/IU & +{-}/{-}{-} & F & CP:2a W:1b,2a \\
B3302 & 1a<o0L 1b<o1 / 2a<1a2L 2b*<1a1,1b1 & NI/NU & +{-}/+{-} & N & W:2b \\
B3303 & 1a<o0L 1b<o1 / 2a<1a2L 2b*<1a1,1b2 & NI/NU & +{-}/+{-} & N & W:2b \\
B3304 & 1a<o0L 1b<o1 / 2a<1a2L 2b*<1a2,1b1 & NI/NU & +{-}/+{-} & N & W:2b \\
B3305 & 1a<o0L 1b<o1 / 2a<1a2L 2b*<1a2,1b2 & NI/NU & +{-}/+{-} & N & W:2b \\
B3306 & 1a<o0L 1b<o1 / 2a<1a2L,1b1 2b*<1a1 & NI/EN & +{-}/+{-} & F & W:2a \\
B3307 & 1a<o0L 1b<o1 / 2a<1a2L,1b1 2b*<1a2 & NI/EN & +{-}/+{-} & F & W:2a \\
B3308 & 1a<o0L 1b<o1 / 2a<1a2L,1b1 2b*<1b1 & NI/EU & +{-}/+{-} & F & W:1b \\
B3309 & 1a<o0L 1b<o1 / 2a<1a2L,1b1 2b*<1b2 & NI/EU & +{-}/+{-} & F & W:1b \\
B3310 & 1a<o0L 1b<o1 / 2a<1a2L,1b2 2b*<1a1 & NI/EN & +{-}/+{-} & F & W:2a \\
B3311 & 1a<o0L 1b<o1 / 2a<1a2L,1b2 2b*<1a2 & NI/EN & +{-}/+{-} & F & W:2a \\
B3312 & 1a<o0L 1b<o1 / 2a<1a2L,1b2 2b*<1b1 & NI/EU & +{-}/+{-} & F & W:1b \\
B3313 & 1a<o0L 1b<o1 / 2a<1a2L,1b2 2b*<1b2 & NI/EU & +{-}/+{-} & F & W:1b \\
B3314 & 1a<o0L 1b<o1 / 2a<1b0L 2b*<1a1,1b1 & NI/NU & +{-}/{-}{-} & N & W:2b \\
B3315 & 1a<o0L 1b<o1 / 2a<1b0L 2b*<1a1,1b2 & NI/NU & +{-}/{-}{-} & N & W:2b \\
B3316 & 1a<o0L 1b<o1 / 2a<1b0L 2b*<1a2,1b1 & NI/NU & +{-}/{-}{-} & N & W:2b \\
B3317 & 1a<o0L 1b<o1 / 2a<1b0L 2b*<1a2,1b2 & NI/NU & +{-}/{-}{-} & N & W:2b \\
B3318 & 1a<o0L 1b<o1 / 2a<1b1 2b*<1a1,1b1 & NI/IU & +{-}/{-}{-} & N & W:1b,2a \\
B3319 & 1a<o0L 1b<o1 / 2a<1b1 2b*<1a1,1b2 & NI/IU & +{-}/{-}{-} & N & W:1b,2a \\
B3320 & 1a<o0L 1b<o1 / 2a<1b1 2b*<1a2,1b1 & NI/IU & +{-}/{-}{-} & N & W:1b,2a \\
B3321 & 1a<o0L 1b<o1 / 2a<1b1 2b*<1a2,1b2 & NI/IU & +{-}/{-}{-} & N & W:1b,2a \\
B3322 & 1a<o0L 1b<o1 / 2a<1b1L 2b*<1a1,1b1 & NI/IU & +{-}/{-}{-} & N & W:1b,2a \\
B3323 & 1a<o0L 1b<o1 / 2a<1b1L 2b*<1a1,1b2 & NI/IU & +{-}/{-}{-} & N & W:1b,2a \\
B3324 & 1a<o0L 1b<o1 / 2a<1b1L 2b*<1a2,1b1 & NI/IU & +{-}/{-}{-} & N & W:1b,2a \\
B3325 & 1a<o0L 1b<o1 / 2a<1b1L 2b*<1a2,1b2 & NI/IU & +{-}/{-}{-} & N & W:1b,2a \\
B3326 & 1a<o0L 1b<o1 / 2a<1b2 2b*<1a1,1b1 & NI/IU & +{-}/{-}{-} & N & W:1b,2a \\
B3327 & 1a<o0L 1b<o1 / 2a<1b2 2b*<1a1,1b2 & NI/IU & +{-}/{-}{-} & N & W:1b,2a \\
B3328 & 1a<o0L 1b<o1 / 2a<1b2 2b*<1a2,1b1 & NI/IU & +{-}/{-}{-} & N & W:1b,2a \\
B3329 & 1a<o0L 1b<o1 / 2a<1b2 2b*<1a2,1b2 & NI/IU & +{-}/{-}{-} & N & W:1b,2a \\
B3330 & 1a<o0L 1b<o1 / 2a<1b2L 2b*<1a1,1b1 & NI/IU & +{-}/{-}{-} & N & W:1b,2a \\
B3331 & 1a<o0L 1b<o1 / 2a<1b2L 2b*<1a1,1b2 & NI/IU & +{-}/{-}{-} & N & W:1b,2a \\
B3332 & 1a<o0L 1b<o1 / 2a<1b2L 2b*<1a2,1b1 & NI/IU & +{-}/{-}{-} & N & W:1b,2a \\
B3333 & 1a<o0L 1b<o1 / 2a<1b2L 2b*<1a2,1b2 & NI/IU & +{-}/{-}{-} & N & W:1b,2a \\
B3334 & 1a<o0L 1b<o2 / 2a<1a0L 2b*<1a1,1b1 & NI/NU & +{-}/+{-} & N & W:2b \\
B3335 & 1a<o0L 1b<o2 / 2a<1a0L 2b*<1a1,1b2 & NI/NU & +{-}/+{-} & N & W:2b \\
B3336 & 1a<o0L 1b<o2 / 2a<1a0L 2b*<1a2,1b1 & NI/NU & +{-}/+{-} & N & W:2b \\
B3337 & 1a<o0L 1b<o2 / 2a<1a0L 2b*<1a2,1b2 & NI/NU & +{-}/+{-} & N & W:2b \\
B3338 & 1a<o0L 1b<o2 / 2a<1a0L,1b1 2b*<1a1 & NI/IN & +{-}/+{-} & N &  \\
B3339 & 1a<o0L 1b<o2 / 2a<1a0L,1b1 2b*<1a2 & NI/IN & +{-}/+{-} & N &  \\
B3340 & 1a<o0L 1b<o2 / 2a<1a0L,1b1 2b*<1b1 & NI/IU & +{-}/+{-} & N & W:1b,2a \\
B3341 & 1a<o0L 1b<o2 / 2a<1a0L,1b1 2b*<1b2 & NI/IU & +{-}/+{-} & N & W:1b,2a \\
B3342 & 1a<o0L 1b<o2 / 2a<1a0L,1b2 2b*<1a1 & NI/IN & +{-}/+{-} & N &  \\
B3343 & 1a<o0L 1b<o2 / 2a<1a0L,1b2 2b*<1a2 & NI/IN & +{-}/+{-} & N &  \\
B3344 & 1a<o0L 1b<o2 / 2a<1a0L,1b2 2b*<1b1 & NI/IU & +{-}/+{-} & N & W:1b,2a \\
B3345 & 1a<o0L 1b<o2 / 2a<1a0L,1b2 2b*<1b2 & NI/IU & +{-}/+{-} & N & W:1b,2a \\
B3346 & 1a<o0L 1b<o2 / 2a<1a1 2b*<1a1,1b1 & NI/NU & +{-}/{-}{-} & N & W:2b \\
B3347 & 1a<o0L 1b<o2 / 2a<1a1 2b*<1a1,1b2 & NI/NU & +{-}/{-}{-} & N & W:2b \\
B3348 & 1a<o0L 1b<o2 / 2a<1a1 2b*<1a2,1b1 & NI/NU & +{-}/{-}{-} & N & W:2b \\
B3349 & 1a<o0L 1b<o2 / 2a<1a1 2b*<1a2,1b2 & NI/NU & +{-}/{-}{-} & N & W:2b \\
B3350 & 1a<o0L 1b<o2 / 2a<1a1,1b0L 2b*<1a1 & NI/NN & +{-}/{-}{-} & N &  \\
B3351 & 1a<o0L 1b<o2 / 2a<1a1,1b0L 2b*<1a2 & NI/NN & +{-}/{-}{-} & N &  \\
B3352 & 1a<o0L 1b<o2 / 2a<1a1,1b0L 2b*<1b1 & NI/NU & +{-}/{-}{-} & N & W:2b \\
B3353 & 1a<o0L 1b<o2 / 2a<1a1,1b0L 2b*<1b2 & NI/NU & +{-}/{-}{-} & N & W:2b \\
B3354 & 1a<o0L 1b<o2 / 2a<1a1,1b1 2b*<1a1 & NI/EN & +{-}/{-}{-} & F & W:2a \\
B3355 & 1a<o0L 1b<o2 / 2a<1a1,1b1 2b*<1a2 & NI/EN & +{-}/{-}{-} & F & W:2a \\
B3356 & 1a<o0L 1b<o2 / 2a<1a1,1b1 2b*<1b1 & NI/EU & +{-}/{-}{-} & F & W:1b \\
B3357 & 1a<o0L 1b<o2 / 2a<1a1,1b1 2b*<1b2 & NI/EU & +{-}/{-}{-} & F & W:1b \\
B3358 & 1a<o0L 1b<o2 / 2a<1a1,1b1L 2b*<1a1 & NI/IN & +{-}/{-}{-} & F & CP:2a \\
B3359 & 1a<o0L 1b<o2 / 2a<1a1,1b1L 2b*<1a2 & NI/IN & +{-}/{-}{-} & F & CP:2a \\
B3360 & 1a<o0L 1b<o2 / 2a<1a1,1b1L 2b*<1b1 & NI/IU & +{-}/{-}{-} & F & CP:2a W:1b,2a \\
B3361 & 1a<o0L 1b<o2 / 2a<1a1,1b1L 2b*<1b2 & NI/IU & +{-}/{-}{-} & F & CP:2a W:1b,2a \\
B3362 & 1a<o0L 1b<o2 / 2a<1a1,1b2 2b*<1a1 & NI/EN & +{-}/{-}{-} & F & CW:2a \\
B3363 & 1a<o0L 1b<o2 / 2a<1a1,1b2 2b*<1a2 & NI/EN & +{-}/{-}{-} & F & CW:2a \\
B3364 & 1a<o0L 1b<o2 / 2a<1a1,1b2 2b*<1b1 & NI/EU & +{-}/{-}{-} & F & C:2a W:1b \\
B3365 & 1a<o0L 1b<o2 / 2a<1a1,1b2 2b*<1b2 & NI/EU & +{-}/{-}{-} & F & C:2a W:1b \\
B3366 & 1a<o0L 1b<o2 / 2a<1a1,1b2L 2b*<1a1 & NI/IN & +{-}/{-}{-} & F & P:2a \\
B3367 & 1a<o0L 1b<o2 / 2a<1a1,1b2L 2b*<1a2 & NI/IN & +{-}/{-}{-} & F & P:2a \\
B3368 & 1a<o0L 1b<o2 / 2a<1a1,1b2L 2b*<1b1 & NI/IU & +{-}/{-}{-} & F & P:2a W:1b,2a \\
B3369 & 1a<o0L 1b<o2 / 2a<1a1,1b2L 2b*<1b2 & NI/IU & +{-}/{-}{-} & F & P:2a W:1b,2a \\
B3370 & 1a<o0L 1b<o2 / 2a<1a1L 2b*<1a1,1b1 & NI/NU & +{-}/+{-} & N & W:2b \\
B3371 & 1a<o0L 1b<o2 / 2a<1a1L 2b*<1a1,1b2 & NI/NU & +{-}/+{-} & N & W:2b \\
B3372 & 1a<o0L 1b<o2 / 2a<1a1L 2b*<1a2,1b1 & NI/NU & +{-}/+{-} & N & W:2b \\
B3373 & 1a<o0L 1b<o2 / 2a<1a1L 2b*<1a2,1b2 & NI/NU & +{-}/+{-} & N & W:2b \\
B3374 & 1a<o0L 1b<o2 / 2a<1a1L,1b1 2b*<1a1 & NI/EN & +{-}/+{-} & F & W:2a \\
B3375 & 1a<o0L 1b<o2 / 2a<1a1L,1b1 2b*<1a2 & NI/EN & +{-}/+{-} & F & W:2a \\
B3376 & 1a<o0L 1b<o2 / 2a<1a1L,1b1 2b*<1b1 & NI/EU & +{-}/+{-} & F & W:1b \\
B3377 & 1a<o0L 1b<o2 / 2a<1a1L,1b1 2b*<1b2 & NI/EU & +{-}/+{-} & F & W:1b \\
B3378 & 1a<o0L 1b<o2 / 2a<1a1L,1b2 2b*<1a1 & NI/EN & +{-}/+{-} & F & CW:2a \\
B3379 & 1a<o0L 1b<o2 / 2a<1a1L,1b2 2b*<1a2 & NI/EN & +{-}/+{-} & F & CW:2a \\
B3380 & 1a<o0L 1b<o2 / 2a<1a1L,1b2 2b*<1b1 & NI/EU & +{-}/+{-} & F & C:2a W:1b \\
B3381 & 1a<o0L 1b<o2 / 2a<1a1L,1b2 2b*<1b2 & NI/EU & +{-}/+{-} & F & C:2a W:1b \\
B3382 & 1a<o0L 1b<o2 / 2a<1a2 2b*<1a1,1b1 & NI/NU & +{-}/{-}{-} & N & W:2b \\
B3383 & 1a<o0L 1b<o2 / 2a<1a2 2b*<1a1,1b2 & NI/NU & +{-}/{-}{-} & N & W:2b \\
B3384 & 1a<o0L 1b<o2 / 2a<1a2 2b*<1a2,1b1 & NI/NU & +{-}/{-}{-} & N & W:2b \\
B3385 & 1a<o0L 1b<o2 / 2a<1a2 2b*<1a2,1b2 & NI/NU & +{-}/{-}{-} & N & W:2b \\
B3386 & 1a<o0L 1b<o2 / 2a<1a2,1b0L 2b*<1a1 & NI/NN & +{-}/{-}{-} & N &  \\
B3387 & 1a<o0L 1b<o2 / 2a<1a2,1b0L 2b*<1a2 & NI/NN & +{-}/{-}{-} & N &  \\
B3388 & 1a<o0L 1b<o2 / 2a<1a2,1b0L 2b*<1b1 & NI/NU & +{-}/{-}{-} & N & W:2b \\
B3389 & 1a<o0L 1b<o2 / 2a<1a2,1b0L 2b*<1b2 & NI/NU & +{-}/{-}{-} & N & W:2b \\
B3390 & 1a<o0L 1b<o2 / 2a<1a2,1b1 2b*<1a1 & NI/EN & +{-}/{-}{-} & F & W:2a \\
B3391 & 1a<o0L 1b<o2 / 2a<1a2,1b1 2b*<1a2 & NI/EN & +{-}/{-}{-} & F & W:2a \\
B3392 & 1a<o0L 1b<o2 / 2a<1a2,1b1 2b*<1b1 & NI/EU & +{-}/{-}{-} & F & W:1b \\
B3393 & 1a<o0L 1b<o2 / 2a<1a2,1b1 2b*<1b2 & NI/EU & +{-}/{-}{-} & F & W:1b \\
B3394 & 1a<o0L 1b<o2 / 2a<1a2,1b1L 2b*<1a1 & NI/IN & +{-}/{-}{-} & F & CP:2a \\
B3395 & 1a<o0L 1b<o2 / 2a<1a2,1b1L 2b*<1a2 & NI/IN & +{-}/{-}{-} & F & CP:2a \\
B3396 & 1a<o0L 1b<o2 / 2a<1a2,1b1L 2b*<1b1 & NI/IU & +{-}/{-}{-} & F & CP:2a W:1b,2a \\
B3397 & 1a<o0L 1b<o2 / 2a<1a2,1b1L 2b*<1b2 & NI/IU & +{-}/{-}{-} & F & CP:2a W:1b,2a \\
B3398 & 1a<o0L 1b<o2 / 2a<1a2,1b2 2b*<1a1 & NI/EN & +{-}/{-}{-} & F & W:2a \\
B3399 & 1a<o0L 1b<o2 / 2a<1a2,1b2 2b*<1a2 & NI/EN & +{-}/{-}{-} & F & W:2a \\
B3400 & 1a<o0L 1b<o2 / 2a<1a2,1b2 2b*<1b1 & NI/EU & +{-}/{-}{-} & F & W:1b \\
B3401 & 1a<o0L 1b<o2 / 2a<1a2,1b2 2b*<1b2 & NI/EU & +{-}/{-}{-} & F & W:1b \\
B3402 & 1a<o0L 1b<o2 / 2a<1a2,1b2L 2b*<1a1 & NI/IN & +{-}/{-}{-} & F & CP:2a \\
B3403 & 1a<o0L 1b<o2 / 2a<1a2,1b2L 2b*<1a2 & NI/IN & +{-}/{-}{-} & F & CP:2a \\
B3404 & 1a<o0L 1b<o2 / 2a<1a2,1b2L 2b*<1b1 & NI/IU & +{-}/{-}{-} & F & CP:2a W:1b,2a \\
B3405 & 1a<o0L 1b<o2 / 2a<1a2,1b2L 2b*<1b2 & NI/IU & +{-}/{-}{-} & F & CP:2a W:1b,2a \\
B3406 & 1a<o0L 1b<o2 / 2a<1a2L 2b*<1a1,1b1 & NI/NU & +{-}/+{-} & N & W:2b \\
B3407 & 1a<o0L 1b<o2 / 2a<1a2L 2b*<1a1,1b2 & NI/NU & +{-}/+{-} & N & W:2b \\
B3408 & 1a<o0L 1b<o2 / 2a<1a2L 2b*<1a2,1b1 & NI/NU & +{-}/+{-} & N & W:2b \\
B3409 & 1a<o0L 1b<o2 / 2a<1a2L 2b*<1a2,1b2 & NI/NU & +{-}/+{-} & N & W:2b \\
B3410 & 1a<o0L 1b<o2 / 2a<1a2L,1b1 2b*<1a1 & NI/EN & +{-}/+{-} & F & W:2a \\
B3411 & 1a<o0L 1b<o2 / 2a<1a2L,1b1 2b*<1a2 & NI/EN & +{-}/+{-} & F & W:2a \\
B3412 & 1a<o0L 1b<o2 / 2a<1a2L,1b1 2b*<1b1 & NI/EU & +{-}/+{-} & F & W:1b \\
B3413 & 1a<o0L 1b<o2 / 2a<1a2L,1b1 2b*<1b2 & NI/EU & +{-}/+{-} & F & W:1b \\
B3414 & 1a<o0L 1b<o2 / 2a<1a2L,1b2 2b*<1a1 & NI/EN & +{-}/+{-} & F & W:2a \\
B3415 & 1a<o0L 1b<o2 / 2a<1a2L,1b2 2b*<1a2 & NI/EN & +{-}/+{-} & F & W:2a \\
B3416 & 1a<o0L 1b<o2 / 2a<1a2L,1b2 2b*<1b1 & NI/EU & +{-}/+{-} & F & W:1b \\
B3417 & 1a<o0L 1b<o2 / 2a<1a2L,1b2 2b*<1b2 & NI/EU & +{-}/+{-} & F & W:1b \\
B3418 & 1a<o0L 1b<o2 / 2a<1b0L 2b*<1a1,1b1 & NI/NU & +{-}/{-}{-} & N & W:2b \\
B3419 & 1a<o0L 1b<o2 / 2a<1b0L 2b*<1a1,1b2 & NI/NU & +{-}/{-}{-} & N & W:2b \\
B3420 & 1a<o0L 1b<o2 / 2a<1b0L 2b*<1a2,1b1 & NI/NU & +{-}/{-}{-} & N & W:2b \\
B3421 & 1a<o0L 1b<o2 / 2a<1b0L 2b*<1a2,1b2 & NI/NU & +{-}/{-}{-} & N & W:2b \\
B3422 & 1a<o0L 1b<o2 / 2a<1b1 2b*<1a1,1b1 & NI/IU & +{-}/{-}{-} & N & W:1b,2a \\
B3423 & 1a<o0L 1b<o2 / 2a<1b1 2b*<1a1,1b2 & NI/IU & +{-}/{-}{-} & N & W:1b,2a \\
B3424 & 1a<o0L 1b<o2 / 2a<1b1 2b*<1a2,1b1 & NI/IU & +{-}/{-}{-} & N & W:1b,2a \\
B3425 & 1a<o0L 1b<o2 / 2a<1b1 2b*<1a2,1b2 & NI/IU & +{-}/{-}{-} & N & W:1b,2a \\
B3426 & 1a<o0L 1b<o2 / 2a<1b1L 2b*<1a1,1b1 & NI/IU & +{-}/{-}{-} & N & W:1b,2a \\
B3427 & 1a<o0L 1b<o2 / 2a<1b1L 2b*<1a1,1b2 & NI/IU & +{-}/{-}{-} & N & W:1b,2a \\
B3428 & 1a<o0L 1b<o2 / 2a<1b1L 2b*<1a2,1b1 & NI/IU & +{-}/{-}{-} & N & W:1b,2a \\
B3429 & 1a<o0L 1b<o2 / 2a<1b1L 2b*<1a2,1b2 & NI/IU & +{-}/{-}{-} & N & W:1b,2a \\
B3430 & 1a<o0L 1b<o2 / 2a<1b2 2b*<1a1,1b1 & NI/IU & +{-}/{-}{-} & N & W:1b,2a \\
B3431 & 1a<o0L 1b<o2 / 2a<1b2 2b*<1a1,1b2 & NI/IU & +{-}/{-}{-} & N & W:1b,2a \\
B3432 & 1a<o0L 1b<o2 / 2a<1b2 2b*<1a2,1b1 & NI/IU & +{-}/{-}{-} & N & W:1b,2a \\
B3433 & 1a<o0L 1b<o2 / 2a<1b2 2b*<1a2,1b2 & NI/IU & +{-}/{-}{-} & N & W:1b,2a \\
B3434 & 1a<o0L 1b<o2 / 2a<1b2L 2b*<1a1,1b1 & NI/IU & +{-}/{-}{-} & N & W:1b,2a \\
B3435 & 1a<o0L 1b<o2 / 2a<1b2L 2b*<1a1,1b2 & NI/IU & +{-}/{-}{-} & N & W:1b,2a \\
B3436 & 1a<o0L 1b<o2 / 2a<1b2L 2b*<1a2,1b1 & NI/IU & +{-}/{-}{-} & N & W:1b,2a \\
B3437 & 1a<o0L 1b<o2 / 2a<1b2L 2b*<1a2,1b2 & NI/IU & +{-}/{-}{-} & N & W:1b,2a \\
B3438 & 1a<o0L 1b*<o1 / 2a<1a0L 2b<1a1,1b1 & NU/NU & +{-}/+{-} & N & W:1b,2b \\
B3439 & 1a<o0L 1b*<o1 / 2a<1a0L 2b<1a1,1b2 & NU/NU & +{-}/+{-} & N & W:1b,2b \\
B3440 & 1a<o0L 1b*<o1 / 2a<1a0L 2b<1a2,1b1 & NU/NU & +{-}/+{-} & N & W:1b,2b \\
B3441 & 1a<o0L 1b*<o1 / 2a<1a0L 2b<1a2,1b2 & NU/NU & +{-}/+{-} & N & W:1b,2b \\
B3442 & 1a<o0L 1b*<o1 / 2a<1a0L,1b1 2b<1a1 & NU/UN & +{-}/+{-} & N & W:1b,2a \\
B3443 & 1a<o0L 1b*<o1 / 2a<1a0L,1b1 2b<1a2 & NU/UN & +{-}/+{-} & N & W:1b,2a \\
B3444 & 1a<o0L 1b*<o1 / 2a<1a0L,1b1 2b<1b1 & NU/UU & +{-}/+{-} & N &  \\
B3445 & 1a<o0L 1b*<o1 / 2a<1a0L,1b1 2b<1b2 & NU/UU & +{-}/+{-} & N &  \\
B3446 & 1a<o0L 1b*<o1 / 2a<1a0L,1b2 2b<1a1 & NU/UN & +{-}/+{-} & N & W:1b,2a \\
B3447 & 1a<o0L 1b*<o1 / 2a<1a0L,1b2 2b<1a2 & NU/UN & +{-}/+{-} & N & W:1b,2a \\
B3448 & 1a<o0L 1b*<o1 / 2a<1a0L,1b2 2b<1b1 & NU/UU & +{-}/+{-} & N &  \\
B3449 & 1a<o0L 1b*<o1 / 2a<1a0L,1b2 2b<1b2 & NU/UU & +{-}/+{-} & N &  \\
B3450 & 1a<o0L 1b*<o1 / 2a<1a1 2b<1a1,1b1 & NU/NU & +{-}/{-}{-} & N & W:1b,2b \\
B3451 & 1a<o0L 1b*<o1 / 2a<1a1 2b<1a1,1b2 & NU/NU & +{-}/{-}{-} & N & W:1b,2b \\
B3452 & 1a<o0L 1b*<o1 / 2a<1a1 2b<1a1L,1b1 & NU/NU & +{-}/{-}+ & N & W:1b,2b \\
B3453 & 1a<o0L 1b*<o1 / 2a<1a1 2b<1a1L,1b2 & NU/NU & +{-}/{-}+ & N & W:1b,2b \\
B3454 & 1a<o0L 1b*<o1 / 2a<1a1 2b<1a2,1b1 & NU/NU & +{-}/{-}{-} & N & W:1b,2b \\
B3455 & 1a<o0L 1b*<o1 / 2a<1a1 2b<1a2,1b2 & NU/NU & +{-}/{-}{-} & N & W:1b,2b \\
B3456 & 1a<o0L 1b*<o1 / 2a<1a1 2b<1a2L,1b1 & NU/NU & +{-}/{-}+ & N & W:1b,2b \\
B3457 & 1a<o0L 1b*<o1 / 2a<1a1 2b<1a2L,1b2 & NU/NU & +{-}/{-}+ & N & W:1b,2b \\
B3458 & 1a<o0L 1b*<o1 / 2a<1a1,1b1 2b<1a1L & NU/UN & +{-}/{-}+ & N & W:1b,2a \\
B3459 & 1a<o0L 1b*<o1 / 2a<1a1,1b1 2b<1a2 & NU/UN & +{-}/{-}{-} & N & W:1b,2a \\
B3460 & 1a<o0L 1b*<o1 / 2a<1a1,1b1 2b<1a2L & NU/UN & +{-}/{-}+ & N & W:1b,2a \\
B3461 & 1a<o0L 1b*<o1 / 2a<1a1,1b1 2b<1b1 & NU/UU & +{-}/{-}{-} & N &  \\
B3462 & 1a<o0L 1b*<o1 / 2a<1a1,1b1 2b<1b2 & NU/UU & +{-}/{-}{-} & N &  \\
B3463 & 1a<o0L 1b*<o1 / 2a<1a1,1b2 2b<1a1L & NU/UN & +{-}/{-}+ & N & W:1b,2a \\
B3464 & 1a<o0L 1b*<o1 / 2a<1a1,1b2 2b<1a2 & NU/UN & +{-}/{-}{-} & N & W:1b,2a \\
B3465 & 1a<o0L 1b*<o1 / 2a<1a1,1b2 2b<1a2L & NU/UN & +{-}/{-}+ & N & W:1b,2a \\
B3466 & 1a<o0L 1b*<o1 / 2a<1a1,1b2 2b<1b1 & NU/UU & +{-}/{-}{-} & N &  \\
B3467 & 1a<o0L 1b*<o1 / 2a<1a1,1b2 2b<1b2 & NU/UU & +{-}/{-}{-} & N &  \\
B3468 & 1a<o0L 1b*<o1 / 2a<1a1L 2b<1a2,1b1 & NU/NU & +{-}/+{-} & N & W:1b,2b \\
B3469 & 1a<o0L 1b*<o1 / 2a<1a1L 2b<1a2,1b2 & NU/NU & +{-}/+{-} & N & W:1b,2b \\
B3470 & 1a<o0L 1b*<o1 / 2a<1a1L,1b1 2b<1a2 & NU/UN & +{-}/+{-} & N & W:1b,2a \\
B3471 & 1a<o0L 1b*<o1 / 2a<1a1L,1b1 2b<1b1 & NU/UU & +{-}/+{-} & N &  \\
B3472 & 1a<o0L 1b*<o1 / 2a<1a1L,1b1 2b<1b2 & NU/UU & +{-}/+{-} & N &  \\
B3473 & 1a<o0L 1b*<o1 / 2a<1a1L,1b2 2b<1a2 & NU/UN & +{-}/+{-} & N & W:1b,2a \\
B3474 & 1a<o0L 1b*<o1 / 2a<1a1L,1b2 2b<1b1 & NU/UU & +{-}/+{-} & N &  \\
B3475 & 1a<o0L 1b*<o1 / 2a<1a1L,1b2 2b<1b2 & NU/UU & +{-}/+{-} & N &  \\
B3476 & 1a<o0L 1b*<o1 / 2a<1a2 2b<1a2,1b1 & NU/NU & +{-}/{-}{-} & N & W:1b,2b \\
B3477 & 1a<o0L 1b*<o1 / 2a<1a2 2b<1a2,1b2 & NU/NU & +{-}/{-}{-} & N & W:1b,2b \\
B3478 & 1a<o0L 1b*<o1 / 2a<1a2 2b<1a2L,1b1 & NU/NU & +{-}/{-}+ & N & W:1b,2b \\
B3479 & 1a<o0L 1b*<o1 / 2a<1a2 2b<1a2L,1b2 & NU/NU & +{-}/{-}+ & N & W:1b,2b \\
B3480 & 1a<o0L 1b*<o1 / 2a<1a2,1b1 2b<1a2L & NU/UN & +{-}/{-}+ & N & W:1b,2a \\
B3481 & 1a<o0L 1b*<o1 / 2a<1a2,1b1 2b<1b1 & NU/UU & +{-}/{-}{-} & N &  \\
B3482 & 1a<o0L 1b*<o1 / 2a<1a2,1b1 2b<1b2 & NU/UU & +{-}/{-}{-} & N &  \\
B3483 & 1a<o0L 1b*<o1 / 2a<1a2,1b2 2b<1a2L & NU/UN & +{-}/{-}+ & N & W:1b,2a \\
B3484 & 1a<o0L 1b*<o1 / 2a<1a2,1b2 2b<1b1 & NU/UU & +{-}/{-}{-} & N &  \\
B3485 & 1a<o0L 1b*<o1 / 2a<1a2,1b2 2b<1b2 & NU/UU & +{-}/{-}{-} & N &  \\
B3486 & 1a<o0L 1b*<o1 / 2a<1a2L,1b1 2b<1b1 & NU/UU & +{-}/+{-} & N &  \\
B3487 & 1a<o0L 1b*<o1 / 2a<1a2L,1b1 2b<1b2 & NU/UU & +{-}/+{-} & N &  \\
B3488 & 1a<o0L 1b*<o1 / 2a<1a2L,1b2 2b<1b1 & NU/UU & +{-}/+{-} & N &  \\
B3489 & 1a<o0L 1b*<o1 / 2a<1a2L,1b2 2b<1b2 & NU/UU & +{-}/+{-} & N &  \\
B3490 & 1a<o0L 1b*<o2 / 2a<1a0L 2b<1a1,1b1 & NU/NU & +{-}/+{-} & N & W:1b,2b \\
B3491 & 1a<o0L 1b*<o2 / 2a<1a0L 2b<1a1,1b2 & NU/NU & +{-}/+{-} & N & W:1b,2b \\
B3492 & 1a<o0L 1b*<o2 / 2a<1a0L 2b<1a2,1b1 & NU/NU & +{-}/+{-} & N & W:1b,2b \\
B3493 & 1a<o0L 1b*<o2 / 2a<1a0L 2b<1a2,1b2 & NU/NU & +{-}/+{-} & N & W:1b,2b \\
B3494 & 1a<o0L 1b*<o2 / 2a<1a0L,1b1 2b<1a1 & NU/UN & +{-}/+{-} & N & W:1b,2a \\
B3495 & 1a<o0L 1b*<o2 / 2a<1a0L,1b1 2b<1a2 & NU/UN & +{-}/+{-} & N & W:1b,2a \\
B3496 & 1a<o0L 1b*<o2 / 2a<1a0L,1b1 2b<1b1 & NU/UU & +{-}/+{-} & N &  \\
B3497 & 1a<o0L 1b*<o2 / 2a<1a0L,1b1 2b<1b2 & NU/UU & +{-}/+{-} & N &  \\
B3498 & 1a<o0L 1b*<o2 / 2a<1a0L,1b2 2b<1a1 & NU/UN & +{-}/+{-} & N & W:1b,2a \\
B3499 & 1a<o0L 1b*<o2 / 2a<1a0L,1b2 2b<1a2 & NU/UN & +{-}/+{-} & N & W:1b,2a \\
B3500 & 1a<o0L 1b*<o2 / 2a<1a0L,1b2 2b<1b1 & NU/UU & +{-}/+{-} & N &  \\
B3501 & 1a<o0L 1b*<o2 / 2a<1a0L,1b2 2b<1b2 & NU/UU & +{-}/+{-} & N &  \\
B3502 & 1a<o0L 1b*<o2 / 2a<1a1 2b<1a1,1b1 & NU/NU & +{-}/{-}{-} & N & W:1b,2b \\
B3503 & 1a<o0L 1b*<o2 / 2a<1a1 2b<1a1,1b2 & NU/NU & +{-}/{-}{-} & N & W:1b,2b \\
B3504 & 1a<o0L 1b*<o2 / 2a<1a1 2b<1a1L,1b1 & NU/NU & +{-}/{-}+ & N & W:1b,2b \\
B3505 & 1a<o0L 1b*<o2 / 2a<1a1 2b<1a1L,1b2 & NU/NU & +{-}/{-}+ & N & W:1b,2b \\
B3506 & 1a<o0L 1b*<o2 / 2a<1a1 2b<1a2,1b1 & NU/NU & +{-}/{-}{-} & N & W:1b,2b \\
B3507 & 1a<o0L 1b*<o2 / 2a<1a1 2b<1a2,1b2 & NU/NU & +{-}/{-}{-} & N & W:1b,2b \\
B3508 & 1a<o0L 1b*<o2 / 2a<1a1 2b<1a2L,1b1 & NU/NU & +{-}/{-}+ & N & W:1b,2b \\
B3509 & 1a<o0L 1b*<o2 / 2a<1a1 2b<1a2L,1b2 & NU/NU & +{-}/{-}+ & N & W:1b,2b \\
B3510 & 1a<o0L 1b*<o2 / 2a<1a1,1b1 2b<1a1L & NU/UN & +{-}/{-}+ & N & W:1b,2a \\
B3511 & 1a<o0L 1b*<o2 / 2a<1a1,1b1 2b<1a2 & NU/UN & +{-}/{-}{-} & N & W:1b,2a \\
B3512 & 1a<o0L 1b*<o2 / 2a<1a1,1b1 2b<1a2L & NU/UN & +{-}/{-}+ & N & W:1b,2a \\
B3513 & 1a<o0L 1b*<o2 / 2a<1a1,1b1 2b<1b1 & NU/UU & +{-}/{-}{-} & N &  \\
B3514 & 1a<o0L 1b*<o2 / 2a<1a1,1b1 2b<1b2 & NU/UU & +{-}/{-}{-} & N &  \\
B3515 & 1a<o0L 1b*<o2 / 2a<1a1,1b2 2b<1a1L & NU/UN & +{-}/{-}+ & N & W:1b,2a \\
B3516 & 1a<o0L 1b*<o2 / 2a<1a1,1b2 2b<1a2 & NU/UN & +{-}/{-}{-} & N & W:1b,2a \\
B3517 & 1a<o0L 1b*<o2 / 2a<1a1,1b2 2b<1a2L & NU/UN & +{-}/{-}+ & N & W:1b,2a \\
B3518 & 1a<o0L 1b*<o2 / 2a<1a1,1b2 2b<1b1 & NU/UU & +{-}/{-}{-} & N &  \\
B3519 & 1a<o0L 1b*<o2 / 2a<1a1,1b2 2b<1b2 & NU/UU & +{-}/{-}{-} & N &  \\
B3520 & 1a<o0L 1b*<o2 / 2a<1a1L 2b<1a2,1b1 & NU/NU & +{-}/+{-} & N & W:1b,2b \\
B3521 & 1a<o0L 1b*<o2 / 2a<1a1L 2b<1a2,1b2 & NU/NU & +{-}/+{-} & N & W:1b,2b \\
B3522 & 1a<o0L 1b*<o2 / 2a<1a1L,1b1 2b<1a2 & NU/UN & +{-}/+{-} & N & W:1b,2a \\
B3523 & 1a<o0L 1b*<o2 / 2a<1a1L,1b1 2b<1b1 & NU/UU & +{-}/+{-} & N &  \\
B3524 & 1a<o0L 1b*<o2 / 2a<1a1L,1b1 2b<1b2 & NU/UU & +{-}/+{-} & N &  \\
B3525 & 1a<o0L 1b*<o2 / 2a<1a1L,1b2 2b<1a2 & NU/UN & +{-}/+{-} & N & W:1b,2a \\
B3526 & 1a<o0L 1b*<o2 / 2a<1a1L,1b2 2b<1b1 & NU/UU & +{-}/+{-} & N &  \\
B3527 & 1a<o0L 1b*<o2 / 2a<1a1L,1b2 2b<1b2 & NU/UU & +{-}/+{-} & N &  \\
B3528 & 1a<o0L 1b*<o2 / 2a<1a2 2b<1a2,1b1 & NU/NU & +{-}/{-}{-} & N & W:1b,2b \\
B3529 & 1a<o0L 1b*<o2 / 2a<1a2 2b<1a2,1b2 & NU/NU & +{-}/{-}{-} & N & W:1b,2b \\
B3530 & 1a<o0L 1b*<o2 / 2a<1a2 2b<1a2L,1b1 & NU/NU & +{-}/{-}+ & N & W:1b,2b \\
B3531 & 1a<o0L 1b*<o2 / 2a<1a2 2b<1a2L,1b2 & NU/NU & +{-}/{-}+ & N & W:1b,2b \\
B3532 & 1a<o0L 1b*<o2 / 2a<1a2,1b1 2b<1a2L & NU/UN & +{-}/{-}+ & N & W:1b,2a \\
B3533 & 1a<o0L 1b*<o2 / 2a<1a2,1b1 2b<1b1 & NU/UU & +{-}/{-}{-} & N &  \\
B3534 & 1a<o0L 1b*<o2 / 2a<1a2,1b1 2b<1b2 & NU/UU & +{-}/{-}{-} & N &  \\
B3535 & 1a<o0L 1b*<o2 / 2a<1a2,1b2 2b<1a2L & NU/UN & +{-}/{-}+ & N & W:1b,2a \\
B3536 & 1a<o0L 1b*<o2 / 2a<1a2,1b2 2b<1b1 & NU/UU & +{-}/{-}{-} & N &  \\
B3537 & 1a<o0L 1b*<o2 / 2a<1a2,1b2 2b<1b2 & NU/UU & +{-}/{-}{-} & N &  \\
B3538 & 1a<o0L 1b*<o2 / 2a<1a2L,1b1 2b<1b1 & NU/UU & +{-}/+{-} & N &  \\
B3539 & 1a<o0L 1b*<o2 / 2a<1a2L,1b1 2b<1b2 & NU/UU & +{-}/+{-} & N &  \\
B3540 & 1a<o0L 1b*<o2 / 2a<1a2L,1b2 2b<1b1 & NU/UU & +{-}/+{-} & N &  \\
B3541 & 1a<o0L 1b*<o2 / 2a<1a2L,1b2 2b<1b2 & NU/UU & +{-}/+{-} & N &  \\
B3542 & 1a<o1 1b<o1 / 2a<1a0L 2b*<1a1,1b1 & EE/NU & {-}{-}/{-}{-} & E &  \\
B3543 & 1a<o1 1b<o1 / 2a<1a0L 2b*<1a1,1b2 & EE/NU & {-}{-}/{-}{-} & E &  \\
B3544 & 1a<o1 1b<o1 / 2a<1a0L 2b*<1a2,1b1 & EE/NU & {-}{-}/{-}{-} & E &  \\
B3545 & 1a<o1 1b<o1 / 2a<1a0L 2b*<1a2,1b2 & EE/NU & {-}{-}/{-}{-} & E &  \\
B3546 & 1a<o1 1b<o1 / 2a<1a0L,1b1 2b*<1a1 & EE/DU & {-}{-}/{-}{-} & E &  \\
B3547 & 1a<o1 1b<o1 / 2a<1a0L,1b1 2b*<1a2 & EE/DU & {-}{-}/{-}{-} & E &  \\
B3548 & 1a<o1 1b<o1 / 2a<1a0L,1b1 2b*<1b1 & EE/DU & {-}{-}/{-}{-} & E &  \\
B3549 & 1a<o1 1b<o1 / 2a<1a0L,1b1 2b*<1b2 & EE/DU & {-}{-}/{-}{-} & E &  \\
B3550 & 1a<o1 1b<o1 / 2a<1a0L,1b2 2b*<1a1 & EE/DU & {-}{-}/{-}{-} & E &  \\
B3551 & 1a<o1 1b<o1 / 2a<1a0L,1b2 2b*<1a2 & EE/DU & {-}{-}/{-}{-} & E &  \\
B3552 & 1a<o1 1b<o1 / 2a<1a0L,1b2 2b*<1b1 & EE/DU & {-}{-}/{-}{-} & E &  \\
B3553 & 1a<o1 1b<o1 / 2a<1a0L,1b2 2b*<1b2 & EE/DU & {-}{-}/{-}{-} & E &  \\
B3554 & 1a<o1 1b<o1 / 2a<1a1 2b*<1a1,1b1 & EE/DU & {-}{-}/{-}{-} & E &  \\
B3555 & 1a<o1 1b<o1 / 2a<1a1 2b*<1a1,1b2 & EE/DU & {-}{-}/{-}{-} & E &  \\
B3556 & 1a<o1 1b<o1 / 2a<1a1 2b*<1a2,1b1 & EE/DU & {-}{-}/{-}{-} & E &  \\
B3557 & 1a<o1 1b<o1 / 2a<1a1 2b*<1a2,1b2 & EE/DU & {-}{-}/{-}{-} & E &  \\
B3558 & 1a<o1 1b<o1 / 2a<1a1,1b1 2b*<1a1 & EE/DU & {-}{-}/{-}{-} & E &  \\
B3559 & 1a<o1 1b<o1 / 2a<1a1,1b1 2b*<1a2 & EE/DU & {-}{-}/{-}{-} & E &  \\
B3560 & 1a<o1 1b<o1 / 2a<1a1,1b1L 2b*<1a1 & EE/DU & {-}{-}/{-}{-} & E &  \\
B3561 & 1a<o1 1b<o1 / 2a<1a1,1b1L 2b*<1a2 & EE/DU & {-}{-}/{-}{-} & E &  \\
B3562 & 1a<o1 1b<o1 / 2a<1a1,1b1L 2b*<1b1 & EE/DU & {-}{-}/{-}{-} & E &  \\
B3563 & 1a<o1 1b<o1 / 2a<1a1,1b1L 2b*<1b2 & EE/DU & {-}{-}/{-}{-} & E &  \\
B3564 & 1a<o1 1b<o1 / 2a<1a1,1b2 2b*<1a1 & EE/DU & {-}{-}/{-}{-} & E &  \\
B3565 & 1a<o1 1b<o1 / 2a<1a1,1b2 2b*<1a2 & EE/DU & {-}{-}/{-}{-} & E &  \\
B3566 & 1a<o1 1b<o1 / 2a<1a1,1b2 2b*<1b1 & EE/DU & {-}{-}/{-}{-} & E &  \\
B3567 & 1a<o1 1b<o1 / 2a<1a1,1b2 2b*<1b2 & EE/DU & {-}{-}/{-}{-} & E &  \\
B3568 & 1a<o1 1b<o1 / 2a<1a1,1b2L 2b*<1a1 & EE/DU & {-}{-}/{-}{-} & E &  \\
B3569 & 1a<o1 1b<o1 / 2a<1a1,1b2L 2b*<1a2 & EE/DU & {-}{-}/{-}{-} & E &  \\
B3570 & 1a<o1 1b<o1 / 2a<1a1,1b2L 2b*<1b1 & EE/DU & {-}{-}/{-}{-} & E &  \\
B3571 & 1a<o1 1b<o1 / 2a<1a1,1b2L 2b*<1b2 & EE/DU & {-}{-}/{-}{-} & E &  \\
B3572 & 1a<o1 1b<o1 / 2a<1a1L 2b*<1a1,1b1 & EE/DU & {-}{-}/{-}{-} & E &  \\
B3573 & 1a<o1 1b<o1 / 2a<1a1L 2b*<1a1,1b2 & EE/DU & {-}{-}/{-}{-} & E &  \\
B3574 & 1a<o1 1b<o1 / 2a<1a1L 2b*<1a2,1b1 & EE/DU & {-}{-}/{-}{-} & E &  \\
B3575 & 1a<o1 1b<o1 / 2a<1a1L 2b*<1a2,1b2 & EE/DU & {-}{-}/{-}{-} & E &  \\
B3576 & 1a<o1 1b<o1 / 2a<1a1L,1b2 2b*<1a1 & EE/DU & {-}{-}/{-}{-} & E &  \\
B3577 & 1a<o1 1b<o1 / 2a<1a1L,1b2 2b*<1a2 & EE/DU & {-}{-}/{-}{-} & E &  \\
B3578 & 1a<o1 1b<o1 / 2a<1a1L,1b2 2b*<1b1 & EE/DU & {-}{-}/{-}{-} & E &  \\
B3579 & 1a<o1 1b<o1 / 2a<1a1L,1b2 2b*<1b2 & EE/DU & {-}{-}/{-}{-} & E &  \\
B3580 & 1a<o1 1b<o1 / 2a<1a2 2b*<1a1,1b1 & EE/DU & {-}{-}/{-}{-} & E &  \\
B3581 & 1a<o1 1b<o1 / 2a<1a2 2b*<1a1,1b2 & EE/DU & {-}{-}/{-}{-} & E &  \\
B3582 & 1a<o1 1b<o1 / 2a<1a2 2b*<1a2,1b1 & EE/DU & {-}{-}/{-}{-} & E &  \\
B3583 & 1a<o1 1b<o1 / 2a<1a2 2b*<1a2,1b2 & EE/DU & {-}{-}/{-}{-} & E &  \\
B3584 & 1a<o1 1b<o1 / 2a<1a2,1b2 2b*<1a1 & EE/DU & {-}{-}/{-}{-} & E &  \\
B3585 & 1a<o1 1b<o1 / 2a<1a2,1b2 2b*<1a2 & EE/DU & {-}{-}/{-}{-} & E &  \\
B3586 & 1a<o1 1b<o1 / 2a<1a2,1b2L 2b*<1a1 & EE/DU & {-}{-}/{-}{-} & E &  \\
B3587 & 1a<o1 1b<o1 / 2a<1a2,1b2L 2b*<1a2 & EE/DU & {-}{-}/{-}{-} & E &  \\
B3588 & 1a<o1 1b<o1 / 2a<1a2,1b2L 2b*<1b1 & EE/DU & {-}{-}/{-}{-} & E &  \\
B3589 & 1a<o1 1b<o1 / 2a<1a2,1b2L 2b*<1b2 & EE/DU & {-}{-}/{-}{-} & E &  \\
B3590 & 1a<o1 1b<o1 / 2a<1a2L 2b*<1a1,1b1 & EE/DU & {-}{-}/{-}{-} & E &  \\
B3591 & 1a<o1 1b<o1 / 2a<1a2L 2b*<1a1,1b2 & EE/DU & {-}{-}/{-}{-} & E &  \\
B3592 & 1a<o1 1b<o1 / 2a<1a2L 2b*<1a2,1b1 & EE/DU & {-}{-}/{-}{-} & E &  \\
B3593 & 1a<o1 1b<o1 / 2a<1a2L 2b*<1a2,1b2 & EE/DU & {-}{-}/{-}{-} & E &  \\
B3594 & 1a<o1 1b<o1L / 2a<1a0L 2b*<1a1,1b1 & LI/NU & {-}+/{-}{-} & B & CPW:1b \\
B3595 & 1a<o1 1b<o1L / 2a<1a0L 2b*<1a1,1b2 & LI/NU & {-}+/{-}{-} & B & CPW:1b \\
B3596 & 1a<o1 1b<o1L / 2a<1a0L 2b*<1a2,1b1 & LI/NU & {-}+/{-}{-} & B & CPW:1b \\
B3597 & 1a<o1 1b<o1L / 2a<1a0L 2b*<1a2,1b2 & LI/NU & {-}+/{-}{-} & B & CPW:1b \\
B3598 & 1a<o1 1b<o1L / 2a<1a0L,1b1 2b*<1a1 & LI/IU & {-}+/{-}{-} & B & CPW:1b,2a \\
B3599 & 1a<o1 1b<o1L / 2a<1a0L,1b1 2b*<1a2 & LI/IU & {-}+/{-}{-} & B & CPW:1b,2a \\
B3600 & 1a<o1 1b<o1L / 2a<1a0L,1b1 2b*<1b1 & LI/IU & {-}+/{-}{-} & B & CPW:1b,2a \\
B3601 & 1a<o1 1b<o1L / 2a<1a0L,1b1 2b*<1b2 & LI/IU & {-}+/{-}{-} & B & CPW:1b,2a \\
B3602 & 1a<o1 1b<o1L / 2a<1a0L,1b2 2b*<1a1 & LI/IU & {-}+/{-}{-} & B & CPW:1b,2a \\
B3603 & 1a<o1 1b<o1L / 2a<1a0L,1b2 2b*<1a2 & LI/IU & {-}+/{-}{-} & B & CPW:1b,2a \\
B3604 & 1a<o1 1b<o1L / 2a<1a0L,1b2 2b*<1b1 & LI/IU & {-}+/{-}{-} & B & CPW:1b,2a \\
B3605 & 1a<o1 1b<o1L / 2a<1a0L,1b2 2b*<1b2 & LI/IU & {-}+/{-}{-} & B & CPW:1b,2a \\
B3606 & 1a<o1 1b<o1L / 2a<1a1 2b*<1a1,1b1 & LI/DU & {-}+/{-}{-} & B & CPW:1b \\
B3607 & 1a<o1 1b<o1L / 2a<1a1 2b*<1a1,1b2 & LI/DU & {-}+/{-}{-} & B & CPW:1b \\
B3608 & 1a<o1 1b<o1L / 2a<1a1 2b*<1a2,1b1 & LI/DU & {-}+/{-}{-} & B & CPW:1b \\
B3609 & 1a<o1 1b<o1L / 2a<1a1 2b*<1a2,1b2 & LI/DU & {-}+/{-}{-} & B & CPW:1b \\
B3610 & 1a<o1 1b<o1L / 2a<1a1,1b0L 2b*<1a1 & LI/DU & {-}+/+{-} & B & CPW:1b \\
B3611 & 1a<o1 1b<o1L / 2a<1a1,1b0L 2b*<1a2 & LI/DU & {-}+/+{-} & B & CPW:1b \\
B3612 & 1a<o1 1b<o1L / 2a<1a1,1b0L 2b*<1b1 & LI/DU & {-}+/+{-} & B & CPW:1b \\
B3613 & 1a<o1 1b<o1L / 2a<1a1,1b0L 2b*<1b2 & LI/DU & {-}+/+{-} & B & CPW:1b \\
B3614 & 1a<o1 1b<o1L / 2a<1a1,1b1 2b*<1a1 & LI/EU & {-}+/{-}{-} & BF & CPW:1b \\
B3615 & 1a<o1 1b<o1L / 2a<1a1,1b1 2b*<1a2 & LI/EU & {-}+/{-}{-} & BF & CPW:1b \\
B3616 & 1a<o1 1b<o1L / 2a<1a1,1b1 2b*<1b1 & LI/EU & {-}+/{-}{-} & BF & CPW:1b \\
B3617 & 1a<o1 1b<o1L / 2a<1a1,1b1 2b*<1b2 & LI/EU & {-}+/{-}{-} & BF & CPW:1b \\
B3618 & 1a<o1 1b<o1L / 2a<1a1,1b1L 2b*<1a1 & LI/IU & {-}+/+{-} & BF & CPW:1b,2a \\
B3619 & 1a<o1 1b<o1L / 2a<1a1,1b1L 2b*<1a2 & LI/IU & {-}+/+{-} & BF & CPW:1b,2a \\
B3620 & 1a<o1 1b<o1L / 2a<1a1,1b1L 2b*<1b1 & LI/IU & {-}+/+{-} & BF & CPW:1b,2a \\
B3621 & 1a<o1 1b<o1L / 2a<1a1,1b1L 2b*<1b2 & LI/IU & {-}+/+{-} & BF & CPW:1b,2a \\
B3622 & 1a<o1 1b<o1L / 2a<1a1,1b2 2b*<1a1 & LI/EU & {-}+/{-}{-} & BF & CPW:1b \\
B3623 & 1a<o1 1b<o1L / 2a<1a1,1b2 2b*<1a2 & LI/EU & {-}+/{-}{-} & BF & CPW:1b \\
B3624 & 1a<o1 1b<o1L / 2a<1a1,1b2 2b*<1b1 & LI/EU & {-}+/{-}{-} & BF & CPW:1b \\
B3625 & 1a<o1 1b<o1L / 2a<1a1,1b2 2b*<1b2 & LI/EU & {-}+/{-}{-} & BF & CPW:1b \\
B3626 & 1a<o1 1b<o1L / 2a<1a1,1b2L 2b*<1a1 & LI/IU & {-}+/+{-} & BF & CPW:1b,2a \\
B3627 & 1a<o1 1b<o1L / 2a<1a1,1b2L 2b*<1a2 & LI/IU & {-}+/+{-} & BF & CPW:1b,2a \\
B3628 & 1a<o1 1b<o1L / 2a<1a1,1b2L 2b*<1b1 & LI/IU & {-}+/+{-} & BF & CPW:1b,2a \\
B3629 & 1a<o1 1b<o1L / 2a<1a1,1b2L 2b*<1b2 & LI/IU & {-}+/+{-} & BF & CPW:1b,2a \\
B3630 & 1a<o1 1b<o1L / 2a<1a1L 2b*<1a1,1b1 & LI/DU & {-}+/{-}{-} & B & CPW:1b \\
B3631 & 1a<o1 1b<o1L / 2a<1a1L 2b*<1a1,1b2 & LI/DU & {-}+/{-}{-} & B & CPW:1b \\
B3632 & 1a<o1 1b<o1L / 2a<1a1L 2b*<1a2,1b1 & LI/DU & {-}+/{-}{-} & B & CPW:1b \\
B3633 & 1a<o1 1b<o1L / 2a<1a1L 2b*<1a2,1b2 & LI/DU & {-}+/{-}{-} & B & CPW:1b \\
B3634 & 1a<o1 1b<o1L / 2a<1a1L,1b1 2b*<1a1 & LI/EU & {-}+/{-}{-} & BF & CPW:1b \\
B3635 & 1a<o1 1b<o1L / 2a<1a1L,1b1 2b*<1a2 & LI/EU & {-}+/{-}{-} & BF & CPW:1b \\
B3636 & 1a<o1 1b<o1L / 2a<1a1L,1b1 2b*<1b1 & LI/EU & {-}+/{-}{-} & BF & CPW:1b \\
B3637 & 1a<o1 1b<o1L / 2a<1a1L,1b1 2b*<1b2 & LI/EU & {-}+/{-}{-} & BF & CPW:1b \\
B3638 & 1a<o1 1b<o1L / 2a<1a1L,1b2 2b*<1a1 & LI/EU & {-}+/{-}{-} & BF & CPW:1b \\
B3639 & 1a<o1 1b<o1L / 2a<1a1L,1b2 2b*<1a2 & LI/EU & {-}+/{-}{-} & BF & CPW:1b \\
B3640 & 1a<o1 1b<o1L / 2a<1a1L,1b2 2b*<1b1 & LI/EU & {-}+/{-}{-} & BF & CPW:1b \\
B3641 & 1a<o1 1b<o1L / 2a<1a1L,1b2 2b*<1b2 & LI/EU & {-}+/{-}{-} & BF & CPW:1b \\
B3642 & 1a<o1 1b<o1L / 2a<1a2 2b*<1a1,1b1 & LI/DU & {-}+/{-}{-} & B & CPW:1b \\
B3643 & 1a<o1 1b<o1L / 2a<1a2 2b*<1a1,1b2 & LI/DU & {-}+/{-}{-} & B & CPW:1b \\
B3644 & 1a<o1 1b<o1L / 2a<1a2 2b*<1a2,1b1 & LI/DU & {-}+/{-}{-} & B & CPW:1b \\
B3645 & 1a<o1 1b<o1L / 2a<1a2 2b*<1a2,1b2 & LI/DU & {-}+/{-}{-} & B & CPW:1b \\
B3646 & 1a<o1 1b<o1L / 2a<1a2,1b0L 2b*<1a1 & LI/DU & {-}+/+{-} & B & CPW:1b \\
B3647 & 1a<o1 1b<o1L / 2a<1a2,1b0L 2b*<1a2 & LI/DU & {-}+/+{-} & B & CPW:1b \\
B3648 & 1a<o1 1b<o1L / 2a<1a2,1b0L 2b*<1b1 & LI/DU & {-}+/+{-} & B & CPW:1b \\
B3649 & 1a<o1 1b<o1L / 2a<1a2,1b0L 2b*<1b2 & LI/DU & {-}+/+{-} & B & CPW:1b \\
B3650 & 1a<o1 1b<o1L / 2a<1a2,1b1 2b*<1a1 & LI/EU & {-}+/{-}{-} & BF & CPW:1b \\
B3651 & 1a<o1 1b<o1L / 2a<1a2,1b1 2b*<1a2 & LI/EU & {-}+/{-}{-} & BF & CPW:1b \\
B3652 & 1a<o1 1b<o1L / 2a<1a2,1b1 2b*<1b1 & LI/EU & {-}+/{-}{-} & BF & CPW:1b \\
B3653 & 1a<o1 1b<o1L / 2a<1a2,1b1 2b*<1b2 & LI/EU & {-}+/{-}{-} & BF & CPW:1b \\
B3654 & 1a<o1 1b<o1L / 2a<1a2,1b1L 2b*<1a1 & LI/IU & {-}+/+{-} & BF & CPW:1b,2a \\
B3655 & 1a<o1 1b<o1L / 2a<1a2,1b1L 2b*<1a2 & LI/IU & {-}+/+{-} & BF & CPW:1b,2a \\
B3656 & 1a<o1 1b<o1L / 2a<1a2,1b1L 2b*<1b1 & LI/IU & {-}+/+{-} & BF & CPW:1b,2a \\
B3657 & 1a<o1 1b<o1L / 2a<1a2,1b1L 2b*<1b2 & LI/IU & {-}+/+{-} & BF & CPW:1b,2a \\
B3658 & 1a<o1 1b<o1L / 2a<1a2,1b2 2b*<1a1 & LI/EU & {-}+/{-}{-} & BF & CPW:1b \\
B3659 & 1a<o1 1b<o1L / 2a<1a2,1b2 2b*<1a2 & LI/EU & {-}+/{-}{-} & BF & CPW:1b \\
B3660 & 1a<o1 1b<o1L / 2a<1a2,1b2 2b*<1b1 & LI/EU & {-}+/{-}{-} & BF & CPW:1b \\
B3661 & 1a<o1 1b<o1L / 2a<1a2,1b2 2b*<1b2 & LI/EU & {-}+/{-}{-} & BF & CPW:1b \\
B3662 & 1a<o1 1b<o1L / 2a<1a2,1b2L 2b*<1a1 & LI/IU & {-}+/+{-} & BF & CPW:1b,2a \\
B3663 & 1a<o1 1b<o1L / 2a<1a2,1b2L 2b*<1a2 & LI/IU & {-}+/+{-} & BF & CPW:1b,2a \\
B3664 & 1a<o1 1b<o1L / 2a<1a2,1b2L 2b*<1b1 & LI/IU & {-}+/+{-} & BF & CPW:1b,2a \\
B3665 & 1a<o1 1b<o1L / 2a<1a2,1b2L 2b*<1b2 & LI/IU & {-}+/+{-} & BF & CPW:1b,2a \\
B3666 & 1a<o1 1b<o1L / 2a<1a2L 2b*<1a1,1b1 & LI/DU & {-}+/{-}{-} & B & CPW:1b \\
B3667 & 1a<o1 1b<o1L / 2a<1a2L 2b*<1a1,1b2 & LI/DU & {-}+/{-}{-} & B & CPW:1b \\
B3668 & 1a<o1 1b<o1L / 2a<1a2L 2b*<1a2,1b1 & LI/DU & {-}+/{-}{-} & B & CPW:1b \\
B3669 & 1a<o1 1b<o1L / 2a<1a2L 2b*<1a2,1b2 & LI/DU & {-}+/{-}{-} & B & CPW:1b \\
B3670 & 1a<o1 1b<o1L / 2a<1a2L,1b1 2b*<1a1 & LI/EU & {-}+/{-}{-} & BF & CPW:1b \\
B3671 & 1a<o1 1b<o1L / 2a<1a2L,1b1 2b*<1a2 & LI/EU & {-}+/{-}{-} & BF & CPW:1b \\
B3672 & 1a<o1 1b<o1L / 2a<1a2L,1b1 2b*<1b1 & LI/EU & {-}+/{-}{-} & BF & CPW:1b \\
B3673 & 1a<o1 1b<o1L / 2a<1a2L,1b1 2b*<1b2 & LI/EU & {-}+/{-}{-} & BF & CPW:1b \\
B3674 & 1a<o1 1b<o1L / 2a<1a2L,1b2 2b*<1a1 & LI/EU & {-}+/{-}{-} & BF & CPW:1b \\
B3675 & 1a<o1 1b<o1L / 2a<1a2L,1b2 2b*<1a2 & LI/EU & {-}+/{-}{-} & BF & CPW:1b \\
B3676 & 1a<o1 1b<o1L / 2a<1a2L,1b2 2b*<1b1 & LI/EU & {-}+/{-}{-} & BF & CPW:1b \\
B3677 & 1a<o1 1b<o1L / 2a<1a2L,1b2 2b*<1b2 & LI/EU & {-}+/{-}{-} & BF & CPW:1b \\
B3678 & 1a<o1 1b<o1L / 2a<1b0L 2b*<1a1,1b1 & LI/NU & {-}+/+{-} & B & CPW:1b \\
B3679 & 1a<o1 1b<o1L / 2a<1b0L 2b*<1a1,1b2 & LI/NU & {-}+/+{-} & B & CPW:1b \\
B3680 & 1a<o1 1b<o1L / 2a<1b0L 2b*<1a2,1b1 & LI/NU & {-}+/+{-} & B & CPW:1b \\
B3681 & 1a<o1 1b<o1L / 2a<1b0L 2b*<1a2,1b2 & LI/NU & {-}+/+{-} & B & CPW:1b \\
B3682 & 1a<o1 1b<o1L / 2a<1b1 2b*<1a1,1b1 & LI/IU & {-}+/{-}{-} & B & CPW:1b,2a \\
B3683 & 1a<o1 1b<o1L / 2a<1b1 2b*<1a1,1b2 & LI/IU & {-}+/{-}{-} & B & CPW:1b,2a \\
B3684 & 1a<o1 1b<o1L / 2a<1b1 2b*<1a2,1b1 & LI/IU & {-}+/{-}{-} & B & CPW:1b,2a \\
B3685 & 1a<o1 1b<o1L / 2a<1b1 2b*<1a2,1b2 & LI/IU & {-}+/{-}{-} & B & CPW:1b,2a \\
B3686 & 1a<o1 1b<o1L / 2a<1b1L 2b*<1a1,1b1 & LI/IU & {-}+/+{-} & B & CPW:1b,2a \\
B3687 & 1a<o1 1b<o1L / 2a<1b1L 2b*<1a1,1b2 & LI/IU & {-}+/+{-} & B & CPW:1b,2a \\
B3688 & 1a<o1 1b<o1L / 2a<1b1L 2b*<1a2,1b1 & LI/IU & {-}+/+{-} & B & CPW:1b,2a \\
B3689 & 1a<o1 1b<o1L / 2a<1b1L 2b*<1a2,1b2 & LI/IU & {-}+/+{-} & B & CPW:1b,2a \\
B3690 & 1a<o1 1b<o1L / 2a<1b2 2b*<1a1,1b1 & LI/IU & {-}+/{-}{-} & B & CPW:1b,2a \\
B3691 & 1a<o1 1b<o1L / 2a<1b2 2b*<1a1,1b2 & LI/IU & {-}+/{-}{-} & B & CPW:1b,2a \\
B3692 & 1a<o1 1b<o1L / 2a<1b2 2b*<1a2,1b1 & LI/IU & {-}+/{-}{-} & B & CPW:1b,2a \\
B3693 & 1a<o1 1b<o1L / 2a<1b2 2b*<1a2,1b2 & LI/IU & {-}+/{-}{-} & B & CPW:1b,2a \\
B3694 & 1a<o1 1b<o1L / 2a<1b2L 2b*<1a1,1b1 & LI/IU & {-}+/+{-} & B & CPW:1b,2a \\
B3695 & 1a<o1 1b<o1L / 2a<1b2L 2b*<1a1,1b2 & LI/IU & {-}+/+{-} & B & CPW:1b,2a \\
B3696 & 1a<o1 1b<o1L / 2a<1b2L 2b*<1a2,1b1 & LI/IU & {-}+/+{-} & B & CPW:1b,2a \\
B3697 & 1a<o1 1b<o1L / 2a<1b2L 2b*<1a2,1b2 & LI/IU & {-}+/+{-} & B & CPW:1b,2a \\
B3698 & 1a<o1 1b<o2 / 2a<1a0L 2b*<1a1,1b1 & EE/NU & {-}{-}/{-}{-} & E & C:1b \\
B3699 & 1a<o1 1b<o2 / 2a<1a0L 2b*<1a1,1b2 & EE/NU & {-}{-}/{-}{-} & E & C:1b \\
B3700 & 1a<o1 1b<o2 / 2a<1a0L 2b*<1a2,1b1 & EE/NU & {-}{-}/{-}{-} & E & C:1b \\
B3701 & 1a<o1 1b<o2 / 2a<1a0L 2b*<1a2,1b2 & EE/NU & {-}{-}/{-}{-} & E & C:1b \\
B3702 & 1a<o1 1b<o2 / 2a<1a0L,1b1 2b*<1a1 & EE/DU & {-}{-}/{-}{-} & E & C:1b,2a \\
B3703 & 1a<o1 1b<o2 / 2a<1a0L,1b1 2b*<1a2 & EE/DU & {-}{-}/{-}{-} & E & C:1b,2a \\
B3704 & 1a<o1 1b<o2 / 2a<1a0L,1b1 2b*<1b1 & EE/DU & {-}{-}/{-}{-} & E & C:1b,2a \\
B3705 & 1a<o1 1b<o2 / 2a<1a0L,1b1 2b*<1b2 & EE/DU & {-}{-}/{-}{-} & E & C:1b,2a \\
B3706 & 1a<o1 1b<o2 / 2a<1a0L,1b2 2b*<1a1 & EE/DU & {-}{-}/{-}{-} & E & C:1b,2a \\
B3707 & 1a<o1 1b<o2 / 2a<1a0L,1b2 2b*<1a2 & EE/DU & {-}{-}/{-}{-} & E & C:1b,2a \\
B3708 & 1a<o1 1b<o2 / 2a<1a0L,1b2 2b*<1b1 & EE/DU & {-}{-}/{-}{-} & E & C:1b,2a \\
B3709 & 1a<o1 1b<o2 / 2a<1a0L,1b2 2b*<1b2 & EE/DU & {-}{-}/{-}{-} & E & C:1b,2a \\
B3710 & 1a<o1 1b<o2 / 2a<1a1 2b*<1a1,1b1 & EE/DU & {-}{-}/{-}{-} & E & C:1b \\
B3711 & 1a<o1 1b<o2 / 2a<1a1 2b*<1a1,1b2 & EE/DU & {-}{-}/{-}{-} & E & C:1b \\
B3712 & 1a<o1 1b<o2 / 2a<1a1 2b*<1a2,1b1 & EE/DU & {-}{-}/{-}{-} & E & C:1b \\
B3713 & 1a<o1 1b<o2 / 2a<1a1 2b*<1a2,1b2 & EE/DU & {-}{-}/{-}{-} & E & C:1b \\
B3714 & 1a<o1 1b<o2 / 2a<1a1,1b0L 2b*<1a1 & EE/DU & {-}{-}/{-}{-} & E & C:1b \\
B3715 & 1a<o1 1b<o2 / 2a<1a1,1b0L 2b*<1a2 & EE/DU & {-}{-}/{-}{-} & E & C:1b \\
B3716 & 1a<o1 1b<o2 / 2a<1a1,1b0L 2b*<1b1 & EE/DU & {-}{-}/{-}{-} & E & C:1b \\
B3717 & 1a<o1 1b<o2 / 2a<1a1,1b0L 2b*<1b2 & EE/DU & {-}{-}/{-}{-} & E & C:1b \\
B3718 & 1a<o1 1b<o2 / 2a<1a1,1b1 2b*<1a1 & EE/DU & {-}{-}/{-}{-} & E & C:1b \\
B3719 & 1a<o1 1b<o2 / 2a<1a1,1b1 2b*<1a2 & EE/DU & {-}{-}/{-}{-} & E & C:1b \\
B3720 & 1a<o1 1b<o2 / 2a<1a1,1b1 2b*<1b1 & EE/DU & {-}{-}/{-}{-} & E & C:1b \\
B3721 & 1a<o1 1b<o2 / 2a<1a1,1b1 2b*<1b2 & EE/DU & {-}{-}/{-}{-} & E & C:1b \\
B3722 & 1a<o1 1b<o2 / 2a<1a1,1b1L 2b*<1a1 & EE/DU & {-}{-}/{-}{-} & E & C:1b \\
B3723 & 1a<o1 1b<o2 / 2a<1a1,1b1L 2b*<1a2 & EE/DU & {-}{-}/{-}{-} & E & C:1b \\
B3724 & 1a<o1 1b<o2 / 2a<1a1,1b1L 2b*<1b1 & EE/DU & {-}{-}/{-}{-} & E & C:1b \\
B3725 & 1a<o1 1b<o2 / 2a<1a1,1b1L 2b*<1b2 & EE/DU & {-}{-}/{-}{-} & E & C:1b \\
B3726 & 1a<o1 1b<o2 / 2a<1a1,1b2 2b*<1a1 & EE/DU & {-}{-}/{-}{-} & E & C:1b,2a \\
B3727 & 1a<o1 1b<o2 / 2a<1a1,1b2 2b*<1a2 & EE/DU & {-}{-}/{-}{-} & E & C:1b,2a \\
B3728 & 1a<o1 1b<o2 / 2a<1a1,1b2 2b*<1b1 & EE/DU & {-}{-}/{-}{-} & E & C:1b,2a \\
B3729 & 1a<o1 1b<o2 / 2a<1a1,1b2 2b*<1b2 & EE/DU & {-}{-}/{-}{-} & E & C:1b,2a \\
B3730 & 1a<o1 1b<o2 / 2a<1a1,1b2L 2b*<1a1 & EE/DU & {-}{-}/{-}{-} & E & C:1b,2a \\
B3731 & 1a<o1 1b<o2 / 2a<1a1,1b2L 2b*<1a2 & EE/DU & {-}{-}/{-}{-} & E & C:1b,2a \\
B3732 & 1a<o1 1b<o2 / 2a<1a1,1b2L 2b*<1b1 & EE/DU & {-}{-}/{-}{-} & E & C:1b,2a \\
B3733 & 1a<o1 1b<o2 / 2a<1a1,1b2L 2b*<1b2 & EE/DU & {-}{-}/{-}{-} & E & C:1b,2a \\
B3734 & 1a<o1 1b<o2 / 2a<1a1L 2b*<1a1,1b1 & EE/DU & {-}{-}/{-}{-} & E & C:1b \\
B3735 & 1a<o1 1b<o2 / 2a<1a1L 2b*<1a1,1b2 & EE/DU & {-}{-}/{-}{-} & E & C:1b \\
B3736 & 1a<o1 1b<o2 / 2a<1a1L 2b*<1a2,1b1 & EE/DU & {-}{-}/{-}{-} & E & C:1b \\
B3737 & 1a<o1 1b<o2 / 2a<1a1L 2b*<1a2,1b2 & EE/DU & {-}{-}/{-}{-} & E & C:1b \\
B3738 & 1a<o1 1b<o2 / 2a<1a1L,1b1 2b*<1a1 & EE/DU & {-}{-}/{-}{-} & E & C:1b \\
B3739 & 1a<o1 1b<o2 / 2a<1a1L,1b1 2b*<1a2 & EE/DU & {-}{-}/{-}{-} & E & C:1b \\
B3740 & 1a<o1 1b<o2 / 2a<1a1L,1b1 2b*<1b1 & EE/DU & {-}{-}/{-}{-} & E & C:1b \\
B3741 & 1a<o1 1b<o2 / 2a<1a1L,1b1 2b*<1b2 & EE/DU & {-}{-}/{-}{-} & E & C:1b \\
B3742 & 1a<o1 1b<o2 / 2a<1a1L,1b2 2b*<1a1 & EE/DU & {-}{-}/{-}{-} & E & C:1b,2a \\
B3743 & 1a<o1 1b<o2 / 2a<1a1L,1b2 2b*<1a2 & EE/DU & {-}{-}/{-}{-} & E & C:1b,2a \\
B3744 & 1a<o1 1b<o2 / 2a<1a1L,1b2 2b*<1b1 & EE/DU & {-}{-}/{-}{-} & E & C:1b,2a \\
B3745 & 1a<o1 1b<o2 / 2a<1a1L,1b2 2b*<1b2 & EE/DU & {-}{-}/{-}{-} & E & C:1b,2a \\
B3746 & 1a<o1 1b<o2 / 2a<1a2 2b*<1a1,1b1 & EE/DU & {-}{-}/{-}{-} & E & C:1b \\
B3747 & 1a<o1 1b<o2 / 2a<1a2 2b*<1a1,1b2 & EE/DU & {-}{-}/{-}{-} & E & C:1b \\
B3748 & 1a<o1 1b<o2 / 2a<1a2 2b*<1a2,1b1 & EE/DU & {-}{-}/{-}{-} & E & C:1b \\
B3749 & 1a<o1 1b<o2 / 2a<1a2 2b*<1a2,1b2 & EE/DU & {-}{-}/{-}{-} & E & C:1b \\
B3750 & 1a<o1 1b<o2 / 2a<1a2,1b0L 2b*<1a1 & EE/DU & {-}{-}/{-}{-} & E & C:1b \\
B3751 & 1a<o1 1b<o2 / 2a<1a2,1b0L 2b*<1a2 & EE/DU & {-}{-}/{-}{-} & E & C:1b \\
B3752 & 1a<o1 1b<o2 / 2a<1a2,1b0L 2b*<1b1 & EE/DU & {-}{-}/{-}{-} & E & C:1b \\
B3753 & 1a<o1 1b<o2 / 2a<1a2,1b0L 2b*<1b2 & EE/DU & {-}{-}/{-}{-} & E & C:1b \\
B3754 & 1a<o1 1b<o2 / 2a<1a2,1b1 2b*<1a1 & EE/DU & {-}{-}/{-}{-} & E & C:1b \\
B3755 & 1a<o1 1b<o2 / 2a<1a2,1b1 2b*<1a2 & EE/DU & {-}{-}/{-}{-} & E & C:1b \\
B3756 & 1a<o1 1b<o2 / 2a<1a2,1b1 2b*<1b1 & EE/DU & {-}{-}/{-}{-} & E & C:1b \\
B3757 & 1a<o1 1b<o2 / 2a<1a2,1b1 2b*<1b2 & EE/DU & {-}{-}/{-}{-} & E & C:1b \\
B3758 & 1a<o1 1b<o2 / 2a<1a2,1b1L 2b*<1a1 & EE/DU & {-}{-}/{-}{-} & E & C:1b \\
B3759 & 1a<o1 1b<o2 / 2a<1a2,1b1L 2b*<1a2 & EE/DU & {-}{-}/{-}{-} & E & C:1b \\
B3760 & 1a<o1 1b<o2 / 2a<1a2,1b1L 2b*<1b1 & EE/DU & {-}{-}/{-}{-} & E & C:1b \\
B3761 & 1a<o1 1b<o2 / 2a<1a2,1b1L 2b*<1b2 & EE/DU & {-}{-}/{-}{-} & E & C:1b \\
B3762 & 1a<o1 1b<o2 / 2a<1a2,1b2 2b*<1a1 & EE/DU & {-}{-}/{-}{-} & E & C:1b \\
B3763 & 1a<o1 1b<o2 / 2a<1a2,1b2 2b*<1a2 & EE/DU & {-}{-}/{-}{-} & E & C:1b \\
B3764 & 1a<o1 1b<o2 / 2a<1a2,1b2 2b*<1b1 & EE/DU & {-}{-}/{-}{-} & E & C:1b \\
B3765 & 1a<o1 1b<o2 / 2a<1a2,1b2 2b*<1b2 & EE/DU & {-}{-}/{-}{-} & E & C:1b \\
B3766 & 1a<o1 1b<o2 / 2a<1a2,1b2L 2b*<1a1 & EE/DU & {-}{-}/{-}{-} & E & C:1b \\
B3767 & 1a<o1 1b<o2 / 2a<1a2,1b2L 2b*<1a2 & EE/DU & {-}{-}/{-}{-} & E & C:1b \\
B3768 & 1a<o1 1b<o2 / 2a<1a2,1b2L 2b*<1b1 & EE/DU & {-}{-}/{-}{-} & E & C:1b \\
B3769 & 1a<o1 1b<o2 / 2a<1a2,1b2L 2b*<1b2 & EE/DU & {-}{-}/{-}{-} & E & C:1b \\
B3770 & 1a<o1 1b<o2 / 2a<1a2L 2b*<1a1,1b1 & EE/DU & {-}{-}/{-}{-} & E & C:1b \\
B3771 & 1a<o1 1b<o2 / 2a<1a2L 2b*<1a1,1b2 & EE/DU & {-}{-}/{-}{-} & E & C:1b \\
B3772 & 1a<o1 1b<o2 / 2a<1a2L 2b*<1a2,1b1 & EE/DU & {-}{-}/{-}{-} & E & C:1b \\
B3773 & 1a<o1 1b<o2 / 2a<1a2L 2b*<1a2,1b2 & EE/DU & {-}{-}/{-}{-} & E & C:1b \\
B3774 & 1a<o1 1b<o2 / 2a<1a2L,1b1 2b*<1a1 & EE/DU & {-}{-}/{-}{-} & E & C:1b \\
B3775 & 1a<o1 1b<o2 / 2a<1a2L,1b1 2b*<1a2 & EE/DU & {-}{-}/{-}{-} & E & C:1b \\
B3776 & 1a<o1 1b<o2 / 2a<1a2L,1b1 2b*<1b1 & EE/DU & {-}{-}/{-}{-} & E & C:1b \\
B3777 & 1a<o1 1b<o2 / 2a<1a2L,1b1 2b*<1b2 & EE/DU & {-}{-}/{-}{-} & E & C:1b \\
B3778 & 1a<o1 1b<o2 / 2a<1a2L,1b2 2b*<1a1 & EE/DU & {-}{-}/{-}{-} & E & C:1b \\
B3779 & 1a<o1 1b<o2 / 2a<1a2L,1b2 2b*<1a2 & EE/DU & {-}{-}/{-}{-} & E & C:1b \\
B3780 & 1a<o1 1b<o2 / 2a<1a2L,1b2 2b*<1b1 & EE/DU & {-}{-}/{-}{-} & E & C:1b \\
B3781 & 1a<o1 1b<o2 / 2a<1a2L,1b2 2b*<1b2 & EE/DU & {-}{-}/{-}{-} & E & C:1b \\
B3782 & 1a<o1 1b<o2 / 2a<1b0L 2b*<1a1,1b1 & EE/NU & {-}{-}/{-}{-} & E & C:1b \\
B3783 & 1a<o1 1b<o2 / 2a<1b0L 2b*<1a1,1b2 & EE/NU & {-}{-}/{-}{-} & E & C:1b \\
B3784 & 1a<o1 1b<o2 / 2a<1b0L 2b*<1a2,1b1 & EE/NU & {-}{-}/{-}{-} & E & C:1b \\
B3785 & 1a<o1 1b<o2 / 2a<1b0L 2b*<1a2,1b2 & EE/NU & {-}{-}/{-}{-} & E & C:1b \\
B3786 & 1a<o1 1b<o2 / 2a<1b1 2b*<1a1,1b1 & EE/DU & {-}{-}/{-}{-} & E & C:1b,2a \\
B3787 & 1a<o1 1b<o2 / 2a<1b1 2b*<1a1,1b2 & EE/DU & {-}{-}/{-}{-} & E & C:1b,2a \\
B3788 & 1a<o1 1b<o2 / 2a<1b1 2b*<1a2,1b1 & EE/DU & {-}{-}/{-}{-} & E & C:1b,2a \\
B3789 & 1a<o1 1b<o2 / 2a<1b1 2b*<1a2,1b2 & EE/DU & {-}{-}/{-}{-} & E & C:1b,2a \\
B3790 & 1a<o1 1b<o2 / 2a<1b1L 2b*<1a1,1b1 & EE/DU & {-}{-}/{-}{-} & E & C:1b,2a \\
B3791 & 1a<o1 1b<o2 / 2a<1b1L 2b*<1a1,1b2 & EE/DU & {-}{-}/{-}{-} & E & C:1b,2a \\
B3792 & 1a<o1 1b<o2 / 2a<1b1L 2b*<1a2,1b1 & EE/DU & {-}{-}/{-}{-} & E & C:1b,2a \\
B3793 & 1a<o1 1b<o2 / 2a<1b1L 2b*<1a2,1b2 & EE/DU & {-}{-}/{-}{-} & E & C:1b,2a \\
B3794 & 1a<o1 1b<o2 / 2a<1b2 2b*<1a1,1b1 & EE/DU & {-}{-}/{-}{-} & E & C:1b,2a \\
B3795 & 1a<o1 1b<o2 / 2a<1b2 2b*<1a1,1b2 & EE/DU & {-}{-}/{-}{-} & E & C:1b,2a \\
B3796 & 1a<o1 1b<o2 / 2a<1b2 2b*<1a2,1b1 & EE/DU & {-}{-}/{-}{-} & E & C:1b,2a \\
B3797 & 1a<o1 1b<o2 / 2a<1b2 2b*<1a2,1b2 & EE/DU & {-}{-}/{-}{-} & E & C:1b,2a \\
B3798 & 1a<o1 1b<o2 / 2a<1b2L 2b*<1a1,1b1 & EE/DU & {-}{-}/{-}{-} & E & C:1b,2a \\
B3799 & 1a<o1 1b<o2 / 2a<1b2L 2b*<1a1,1b2 & EE/DU & {-}{-}/{-}{-} & E & C:1b,2a \\
B3800 & 1a<o1 1b<o2 / 2a<1b2L 2b*<1a2,1b1 & EE/DU & {-}{-}/{-}{-} & E & C:1b,2a \\
B3801 & 1a<o1 1b<o2 / 2a<1b2L 2b*<1a2,1b2 & EE/DU & {-}{-}/{-}{-} & E & C:1b,2a \\
B3802 & 1a<o1 1b<o2L / 2a<1a0L 2b*<1a1,1b1 & LI/NU & {-}+/{-}{-} & B & PW:1b \\
B3803 & 1a<o1 1b<o2L / 2a<1a0L 2b*<1a1,1b2 & LI/NU & {-}+/{-}{-} & B & PW:1b \\
B3804 & 1a<o1 1b<o2L / 2a<1a0L 2b*<1a2,1b1 & LI/NU & {-}+/{-}{-} & B & PW:1b \\
B3805 & 1a<o1 1b<o2L / 2a<1a0L 2b*<1a2,1b2 & LI/NU & {-}+/{-}{-} & B & PW:1b \\
B3806 & 1a<o1 1b<o2L / 2a<1a0L,1b1 2b*<1a1 & LI/IU & {-}+/{-}{-} & B & PW:1b,2a \\
B3807 & 1a<o1 1b<o2L / 2a<1a0L,1b1 2b*<1a2 & LI/IU & {-}+/{-}{-} & B & PW:1b,2a \\
B3808 & 1a<o1 1b<o2L / 2a<1a0L,1b1 2b*<1b1 & LI/IU & {-}+/{-}{-} & B & PW:1b,2a \\
B3809 & 1a<o1 1b<o2L / 2a<1a0L,1b1 2b*<1b2 & LI/IU & {-}+/{-}{-} & B & PW:1b,2a \\
B3810 & 1a<o1 1b<o2L / 2a<1a0L,1b2 2b*<1a1 & LI/IU & {-}+/{-}{-} & B & PW:1b,2a \\
B3811 & 1a<o1 1b<o2L / 2a<1a0L,1b2 2b*<1a2 & LI/IU & {-}+/{-}{-} & B & PW:1b,2a \\
B3812 & 1a<o1 1b<o2L / 2a<1a0L,1b2 2b*<1b1 & LI/IU & {-}+/{-}{-} & B & PW:1b,2a \\
B3813 & 1a<o1 1b<o2L / 2a<1a0L,1b2 2b*<1b2 & LI/IU & {-}+/{-}{-} & B & PW:1b,2a \\
B3814 & 1a<o1 1b<o2L / 2a<1a1 2b*<1a1,1b1 & LI/DU & {-}+/{-}{-} & B & PW:1b \\
B3815 & 1a<o1 1b<o2L / 2a<1a1 2b*<1a1,1b2 & LI/DU & {-}+/{-}{-} & B & PW:1b \\
B3816 & 1a<o1 1b<o2L / 2a<1a1 2b*<1a2,1b1 & LI/DU & {-}+/{-}{-} & B & PW:1b \\
B3817 & 1a<o1 1b<o2L / 2a<1a1 2b*<1a2,1b2 & LI/DU & {-}+/{-}{-} & B & PW:1b \\
B3818 & 1a<o1 1b<o2L / 2a<1a1,1b0L 2b*<1a1 & LI/DU & {-}+/+{-} & B & PW:1b \\
B3819 & 1a<o1 1b<o2L / 2a<1a1,1b0L 2b*<1a2 & LI/DU & {-}+/+{-} & B & PW:1b \\
B3820 & 1a<o1 1b<o2L / 2a<1a1,1b0L 2b*<1b1 & LI/DU & {-}+/+{-} & B & PW:1b \\
B3821 & 1a<o1 1b<o2L / 2a<1a1,1b0L 2b*<1b2 & LI/DU & {-}+/+{-} & B & PW:1b \\
B3822 & 1a<o1 1b<o2L / 2a<1a1,1b1 2b*<1a1 & LI/EU & {-}+/{-}{-} & BF & PW:1b \\
B3823 & 1a<o1 1b<o2L / 2a<1a1,1b1 2b*<1a2 & LI/EU & {-}+/{-}{-} & BF & PW:1b \\
B3824 & 1a<o1 1b<o2L / 2a<1a1,1b1 2b*<1b1 & LI/EU & {-}+/{-}{-} & BF & PW:1b \\
B3825 & 1a<o1 1b<o2L / 2a<1a1,1b1 2b*<1b2 & LI/EU & {-}+/{-}{-} & BF & PW:1b \\
B3826 & 1a<o1 1b<o2L / 2a<1a1,1b1L 2b*<1a1 & LI/IU & {-}+/+{-} & BF & C:2a PW:1b,2a \\
B3827 & 1a<o1 1b<o2L / 2a<1a1,1b1L 2b*<1a2 & LI/IU & {-}+/+{-} & BF & C:2a PW:1b,2a \\
B3828 & 1a<o1 1b<o2L / 2a<1a1,1b1L 2b*<1b1 & LI/IU & {-}+/+{-} & BF & C:2a PW:1b,2a \\
B3829 & 1a<o1 1b<o2L / 2a<1a1,1b1L 2b*<1b2 & LI/IU & {-}+/+{-} & BF & C:2a PW:1b,2a \\
B3830 & 1a<o1 1b<o2L / 2a<1a1,1b2 2b*<1a1 & LI/EU & {-}+/{-}{-} & BF & C:2a PW:1b \\
B3831 & 1a<o1 1b<o2L / 2a<1a1,1b2 2b*<1a2 & LI/EU & {-}+/{-}{-} & BF & C:2a PW:1b \\
B3832 & 1a<o1 1b<o2L / 2a<1a1,1b2 2b*<1b1 & LI/EU & {-}+/{-}{-} & BF & C:2a PW:1b \\
B3833 & 1a<o1 1b<o2L / 2a<1a1,1b2 2b*<1b2 & LI/EU & {-}+/{-}{-} & BF & C:2a PW:1b \\
B3834 & 1a<o1 1b<o2L / 2a<1a1,1b2L 2b*<1a1 & LI/IU & {-}+/+{-} & BF & PW:1b,2a \\
B3835 & 1a<o1 1b<o2L / 2a<1a1,1b2L 2b*<1a2 & LI/IU & {-}+/+{-} & BF & PW:1b,2a \\
B3836 & 1a<o1 1b<o2L / 2a<1a1,1b2L 2b*<1b1 & LI/IU & {-}+/+{-} & BF & PW:1b,2a \\
B3837 & 1a<o1 1b<o2L / 2a<1a1,1b2L 2b*<1b2 & LI/IU & {-}+/+{-} & BF & PW:1b,2a \\
B3838 & 1a<o1 1b<o2L / 2a<1a1L 2b*<1a1,1b1 & LI/DU & {-}+/{-}{-} & B & PW:1b \\
B3839 & 1a<o1 1b<o2L / 2a<1a1L 2b*<1a1,1b2 & LI/DU & {-}+/{-}{-} & B & PW:1b \\
B3840 & 1a<o1 1b<o2L / 2a<1a1L 2b*<1a2,1b1 & LI/DU & {-}+/{-}{-} & B & PW:1b \\
B3841 & 1a<o1 1b<o2L / 2a<1a1L 2b*<1a2,1b2 & LI/DU & {-}+/{-}{-} & B & PW:1b \\
B3842 & 1a<o1 1b<o2L / 2a<1a1L,1b1 2b*<1a1 & LI/EU & {-}+/{-}{-} & BF & PW:1b \\
B3843 & 1a<o1 1b<o2L / 2a<1a1L,1b1 2b*<1a2 & LI/EU & {-}+/{-}{-} & BF & PW:1b \\
B3844 & 1a<o1 1b<o2L / 2a<1a1L,1b1 2b*<1b1 & LI/EU & {-}+/{-}{-} & BF & PW:1b \\
B3845 & 1a<o1 1b<o2L / 2a<1a1L,1b1 2b*<1b2 & LI/EU & {-}+/{-}{-} & BF & PW:1b \\
B3846 & 1a<o1 1b<o2L / 2a<1a1L,1b2 2b*<1a1 & LI/EU & {-}+/{-}{-} & BF & C:2a PW:1b \\
B3847 & 1a<o1 1b<o2L / 2a<1a1L,1b2 2b*<1a2 & LI/EU & {-}+/{-}{-} & BF & C:2a PW:1b \\
B3848 & 1a<o1 1b<o2L / 2a<1a1L,1b2 2b*<1b1 & LI/EU & {-}+/{-}{-} & BF & C:2a PW:1b \\
B3849 & 1a<o1 1b<o2L / 2a<1a1L,1b2 2b*<1b2 & LI/EU & {-}+/{-}{-} & BF & C:2a PW:1b \\
B3850 & 1a<o1 1b<o2L / 2a<1a2 2b*<1a1,1b1 & LI/DU & {-}+/{-}{-} & B & PW:1b \\
B3851 & 1a<o1 1b<o2L / 2a<1a2 2b*<1a1,1b2 & LI/DU & {-}+/{-}{-} & B & PW:1b \\
B3852 & 1a<o1 1b<o2L / 2a<1a2 2b*<1a2,1b1 & LI/DU & {-}+/{-}{-} & B & PW:1b \\
B3853 & 1a<o1 1b<o2L / 2a<1a2 2b*<1a2,1b2 & LI/DU & {-}+/{-}{-} & B & PW:1b \\
B3854 & 1a<o1 1b<o2L / 2a<1a2,1b0L 2b*<1a1 & LI/DU & {-}+/+{-} & B & PW:1b \\
B3855 & 1a<o1 1b<o2L / 2a<1a2,1b0L 2b*<1a2 & LI/DU & {-}+/+{-} & B & PW:1b \\
B3856 & 1a<o1 1b<o2L / 2a<1a2,1b0L 2b*<1b1 & LI/DU & {-}+/+{-} & B & PW:1b \\
B3857 & 1a<o1 1b<o2L / 2a<1a2,1b0L 2b*<1b2 & LI/DU & {-}+/+{-} & B & PW:1b \\
B3858 & 1a<o1 1b<o2L / 2a<1a2,1b1 2b*<1a1 & LI/EU & {-}+/{-}{-} & BF & PW:1b \\
B3859 & 1a<o1 1b<o2L / 2a<1a2,1b1 2b*<1a2 & LI/EU & {-}+/{-}{-} & BF & PW:1b \\
B3860 & 1a<o1 1b<o2L / 2a<1a2,1b1 2b*<1b1 & LI/EU & {-}+/{-}{-} & BF & PW:1b \\
B3861 & 1a<o1 1b<o2L / 2a<1a2,1b1 2b*<1b2 & LI/EU & {-}+/{-}{-} & BF & PW:1b \\
B3862 & 1a<o1 1b<o2L / 2a<1a2,1b1L 2b*<1a1 & LI/IU & {-}+/+{-} & BF & C:2a PW:1b,2a \\
B3863 & 1a<o1 1b<o2L / 2a<1a2,1b1L 2b*<1a2 & LI/IU & {-}+/+{-} & BF & C:2a PW:1b,2a \\
B3864 & 1a<o1 1b<o2L / 2a<1a2,1b1L 2b*<1b1 & LI/IU & {-}+/+{-} & BF & C:2a PW:1b,2a \\
B3865 & 1a<o1 1b<o2L / 2a<1a2,1b1L 2b*<1b2 & LI/IU & {-}+/+{-} & BF & C:2a PW:1b,2a \\
B3866 & 1a<o1 1b<o2L / 2a<1a2,1b2 2b*<1a1 & LI/EU & {-}+/{-}{-} & BF & PW:1b \\
B3867 & 1a<o1 1b<o2L / 2a<1a2,1b2 2b*<1a2 & LI/EU & {-}+/{-}{-} & BF & PW:1b \\
B3868 & 1a<o1 1b<o2L / 2a<1a2,1b2 2b*<1b1 & LI/EU & {-}+/{-}{-} & BF & PW:1b \\
B3869 & 1a<o1 1b<o2L / 2a<1a2,1b2 2b*<1b2 & LI/EU & {-}+/{-}{-} & BF & PW:1b \\
B3870 & 1a<o1 1b<o2L / 2a<1a2,1b2L 2b*<1a1 & LI/IU & {-}+/+{-} & BF & C:2a PW:1b,2a \\
B3871 & 1a<o1 1b<o2L / 2a<1a2,1b2L 2b*<1a2 & LI/IU & {-}+/+{-} & BF & C:2a PW:1b,2a \\
B3872 & 1a<o1 1b<o2L / 2a<1a2,1b2L 2b*<1b1 & LI/IU & {-}+/+{-} & BF & C:2a PW:1b,2a \\
B3873 & 1a<o1 1b<o2L / 2a<1a2,1b2L 2b*<1b2 & LI/IU & {-}+/+{-} & BF & C:2a PW:1b,2a \\
B3874 & 1a<o1 1b<o2L / 2a<1a2L 2b*<1a1,1b1 & LI/DU & {-}+/{-}{-} & B & PW:1b \\
B3875 & 1a<o1 1b<o2L / 2a<1a2L 2b*<1a1,1b2 & LI/DU & {-}+/{-}{-} & B & PW:1b \\
B3876 & 1a<o1 1b<o2L / 2a<1a2L 2b*<1a2,1b1 & LI/DU & {-}+/{-}{-} & B & PW:1b \\
B3877 & 1a<o1 1b<o2L / 2a<1a2L 2b*<1a2,1b2 & LI/DU & {-}+/{-}{-} & B & PW:1b \\
B3878 & 1a<o1 1b<o2L / 2a<1a2L,1b1 2b*<1a1 & LI/EU & {-}+/{-}{-} & BF & PW:1b \\
B3879 & 1a<o1 1b<o2L / 2a<1a2L,1b1 2b*<1a2 & LI/EU & {-}+/{-}{-} & BF & PW:1b \\
B3880 & 1a<o1 1b<o2L / 2a<1a2L,1b1 2b*<1b1 & LI/EU & {-}+/{-}{-} & BF & PW:1b \\
B3881 & 1a<o1 1b<o2L / 2a<1a2L,1b1 2b*<1b2 & LI/EU & {-}+/{-}{-} & BF & PW:1b \\
B3882 & 1a<o1 1b<o2L / 2a<1a2L,1b2 2b*<1a1 & LI/EU & {-}+/{-}{-} & BF & PW:1b \\
B3883 & 1a<o1 1b<o2L / 2a<1a2L,1b2 2b*<1a2 & LI/EU & {-}+/{-}{-} & BF & PW:1b \\
B3884 & 1a<o1 1b<o2L / 2a<1a2L,1b2 2b*<1b1 & LI/EU & {-}+/{-}{-} & BF & PW:1b \\
B3885 & 1a<o1 1b<o2L / 2a<1a2L,1b2 2b*<1b2 & LI/EU & {-}+/{-}{-} & BF & PW:1b \\
B3886 & 1a<o1 1b<o2L / 2a<1b0L 2b*<1a1,1b1 & LI/NU & {-}+/+{-} & B & PW:1b \\
B3887 & 1a<o1 1b<o2L / 2a<1b0L 2b*<1a1,1b2 & LI/NU & {-}+/+{-} & B & PW:1b \\
B3888 & 1a<o1 1b<o2L / 2a<1b0L 2b*<1a2,1b1 & LI/NU & {-}+/+{-} & B & PW:1b \\
B3889 & 1a<o1 1b<o2L / 2a<1b0L 2b*<1a2,1b2 & LI/NU & {-}+/+{-} & B & PW:1b \\
B3890 & 1a<o1 1b<o2L / 2a<1b1 2b*<1a1,1b1 & LI/IU & {-}+/{-}{-} & B & PW:1b,2a \\
B3891 & 1a<o1 1b<o2L / 2a<1b1 2b*<1a1,1b2 & LI/IU & {-}+/{-}{-} & B & PW:1b,2a \\
B3892 & 1a<o1 1b<o2L / 2a<1b1 2b*<1a2,1b1 & LI/IU & {-}+/{-}{-} & B & PW:1b,2a \\
B3893 & 1a<o1 1b<o2L / 2a<1b1 2b*<1a2,1b2 & LI/IU & {-}+/{-}{-} & B & PW:1b,2a \\
B3894 & 1a<o1 1b<o2L / 2a<1b1L 2b*<1a1,1b1 & LI/IU & {-}+/+{-} & B & PW:1b,2a \\
B3895 & 1a<o1 1b<o2L / 2a<1b1L 2b*<1a1,1b2 & LI/IU & {-}+/+{-} & B & PW:1b,2a \\
B3896 & 1a<o1 1b<o2L / 2a<1b1L 2b*<1a2,1b1 & LI/IU & {-}+/+{-} & B & PW:1b,2a \\
B3897 & 1a<o1 1b<o2L / 2a<1b1L 2b*<1a2,1b2 & LI/IU & {-}+/+{-} & B & PW:1b,2a \\
B3898 & 1a<o1 1b<o2L / 2a<1b2 2b*<1a1,1b1 & LI/IU & {-}+/{-}{-} & B & PW:1b,2a \\
B3899 & 1a<o1 1b<o2L / 2a<1b2 2b*<1a1,1b2 & LI/IU & {-}+/{-}{-} & B & PW:1b,2a \\
B3900 & 1a<o1 1b<o2L / 2a<1b2 2b*<1a2,1b1 & LI/IU & {-}+/{-}{-} & B & PW:1b,2a \\
B3901 & 1a<o1 1b<o2L / 2a<1b2 2b*<1a2,1b2 & LI/IU & {-}+/{-}{-} & B & PW:1b,2a \\
B3902 & 1a<o1 1b<o2L / 2a<1b2L 2b*<1a1,1b1 & LI/IU & {-}+/+{-} & B & PW:1b,2a \\
B3903 & 1a<o1 1b<o2L / 2a<1b2L 2b*<1a1,1b2 & LI/IU & {-}+/+{-} & B & PW:1b,2a \\
B3904 & 1a<o1 1b<o2L / 2a<1b2L 2b*<1a2,1b1 & LI/IU & {-}+/+{-} & B & PW:1b,2a \\
B3905 & 1a<o1 1b<o2L / 2a<1b2L 2b*<1a2,1b2 & LI/IU & {-}+/+{-} & B & PW:1b,2a \\
B3906 & 1a<o1 1b*<o1 / 2a<1a0L 2b<1a1,1b1 & IU/NI & {-}{-}/{-}{-} & N & W:1a,2b \\
B3907 & 1a<o1 1b*<o1 / 2a<1a0L 2b<1a1,1b2 & IU/NI & {-}{-}/{-}{-} & N & W:1a,2b \\
B3908 & 1a<o1 1b*<o1 / 2a<1a0L 2b<1a2,1b1 & IU/NI & {-}{-}/{-}{-} & N & W:1a,2b \\
B3909 & 1a<o1 1b*<o1 / 2a<1a0L 2b<1a2,1b2 & IU/NI & {-}{-}/{-}{-} & N & W:1a,2b \\
B3910 & 1a<o1 1b*<o1 / 2a<1a0L,1b1 2b<1a1 & IU/UI & {-}{-}/{-}{-} & N & W:1a,2b \\
B3911 & 1a<o1 1b*<o1 / 2a<1a0L,1b1 2b<1a2 & IU/UI & {-}{-}/{-}{-} & N & W:1a,2b \\
B3912 & 1a<o1 1b*<o1 / 2a<1a0L,1b1 2b<1b1 & IU/UU & {-}{-}/{-}{-} & N & W:1a \\
B3913 & 1a<o1 1b*<o1 / 2a<1a0L,1b1 2b<1b2 & IU/UU & {-}{-}/{-}{-} & N & W:1a \\
B3914 & 1a<o1 1b*<o1 / 2a<1a0L,1b2 2b<1a1 & IU/UI & {-}{-}/{-}{-} & N & W:1a,2b \\
B3915 & 1a<o1 1b*<o1 / 2a<1a0L,1b2 2b<1a2 & IU/UI & {-}{-}/{-}{-} & N & W:1a,2b \\
B3916 & 1a<o1 1b*<o1 / 2a<1a0L,1b2 2b<1b1 & IU/UU & {-}{-}/{-}{-} & N & W:1a \\
B3917 & 1a<o1 1b*<o1 / 2a<1a0L,1b2 2b<1b2 & IU/UU & {-}{-}/{-}{-} & N & W:1a \\
B3918 & 1a<o1 1b*<o1 / 2a<1a1 2b<1a1,1b1 & IU/EE & {-}{-}/{-}{-} & E & W:1a \\
B3919 & 1a<o1 1b*<o1 / 2a<1a1 2b<1a1,1b2 & IU/EE & {-}{-}/{-}{-} & E & W:1a \\
B3920 & 1a<o1 1b*<o1 / 2a<1a1 2b<1a1L,1b1 & IU/LI & {-}{-}/{-}{-} & B & CP:2b W:1a,2b \\
B3921 & 1a<o1 1b*<o1 / 2a<1a1 2b<1a1L,1b2 & IU/LI & {-}{-}/{-}{-} & B & CP:2b W:1a,2b \\
B3922 & 1a<o1 1b*<o1 / 2a<1a1 2b<1a2,1b1 & IU/EE & {-}{-}/{-}{-} & E & C:2b W:1a \\
B3923 & 1a<o1 1b*<o1 / 2a<1a1 2b<1a2,1b2 & IU/EE & {-}{-}/{-}{-} & E & C:2b W:1a \\
B3924 & 1a<o1 1b*<o1 / 2a<1a1 2b<1a2L,1b1 & IU/LI & {-}{-}/{-}{-} & B & P:2b W:1a,2b \\
B3925 & 1a<o1 1b*<o1 / 2a<1a1 2b<1a2L,1b2 & IU/LI & {-}{-}/{-}{-} & B & P:2b W:1a,2b \\
B3926 & 1a<o1 1b*<o1 / 2a<1a1,1b1 2b<1a1L & IU/LI & {-}{-}/{-}{-} & B & CP:2b W:1a,2b \\
B3927 & 1a<o1 1b*<o1 / 2a<1a1,1b1 2b<1a2 & IU/EE & {-}{-}/{-}{-} & E & C:2b W:1a \\
B3928 & 1a<o1 1b*<o1 / 2a<1a1,1b1 2b<1a2L & IU/LI & {-}{-}/{-}{-} & B & P:2b W:1a,2b \\
B3929 & 1a<o1 1b*<o1 / 2a<1a1,1b1 2b<1b1 & IU/IU & {-}{-}/{-}{-} & N & W:1a,2a \\
B3930 & 1a<o1 1b*<o1 / 2a<1a1,1b1 2b<1b2 & IU/IU & {-}{-}/{-}{-} & N & W:1a,2a \\
B3931 & 1a<o1 1b*<o1 / 2a<1a1,1b2 2b<1a1L & IU/LI & {-}{-}/{-}{-} & B & CP:2b W:1a,2b \\
B3932 & 1a<o1 1b*<o1 / 2a<1a1,1b2 2b<1a2 & IU/EE & {-}{-}/{-}{-} & E & C:2b W:1a \\
B3933 & 1a<o1 1b*<o1 / 2a<1a1,1b2 2b<1a2L & IU/LI & {-}{-}/{-}{-} & B & P:2b W:1a,2b \\
B3934 & 1a<o1 1b*<o1 / 2a<1a1,1b2 2b<1b1 & IU/IU & {-}{-}/{-}{-} & N & W:1a,2a \\
B3935 & 1a<o1 1b*<o1 / 2a<1a1,1b2 2b<1b2 & IU/IU & {-}{-}/{-}{-} & N & W:1a,2a \\
B3936 & 1a<o1 1b*<o1 / 2a<1a1L 2b<1a2,1b1 & IU/IL & {-}{-}/{-}{-} & B & C:2a,2b P:2a W:1a,2a \\
B3937 & 1a<o1 1b*<o1 / 2a<1a1L 2b<1a2,1b2 & IU/IL & {-}{-}/{-}{-} & B & C:2a,2b P:2a W:1a,2a \\
B3938 & 1a<o1 1b*<o1 / 2a<1a1L,1b1 2b<1a2 & IU/IL & {-}{-}/{-}{-} & B & C:2a,2b P:2a W:1a,2a \\
B3939 & 1a<o1 1b*<o1 / 2a<1a1L,1b1 2b<1b1 & IU/IU & {-}{-}/{-}{-} & N & W:1a,2a \\
B3940 & 1a<o1 1b*<o1 / 2a<1a1L,1b1 2b<1b2 & IU/IU & {-}{-}/{-}{-} & N & W:1a,2a \\
B3941 & 1a<o1 1b*<o1 / 2a<1a1L,1b2 2b<1a2 & IU/IL & {-}{-}/{-}{-} & B & C:2a,2b P:2a W:1a,2a \\
B3942 & 1a<o1 1b*<o1 / 2a<1a1L,1b2 2b<1b1 & IU/IU & {-}{-}/{-}{-} & N & W:1a,2a \\
B3943 & 1a<o1 1b*<o1 / 2a<1a1L,1b2 2b<1b2 & IU/IU & {-}{-}/{-}{-} & N & W:1a,2a \\
B3944 & 1a<o1 1b*<o1 / 2a<1a2 2b<1a2,1b1 & IU/EE & {-}{-}/{-}{-} & E & W:1a \\
B3945 & 1a<o1 1b*<o1 / 2a<1a2 2b<1a2,1b2 & IU/EE & {-}{-}/{-}{-} & E & W:1a \\
B3946 & 1a<o1 1b*<o1 / 2a<1a2 2b<1a2L,1b1 & IU/LI & {-}{-}/{-}{-} & B & CP:2b W:1a,2b \\
B3947 & 1a<o1 1b*<o1 / 2a<1a2 2b<1a2L,1b2 & IU/LI & {-}{-}/{-}{-} & B & CP:2b W:1a,2b \\
B3948 & 1a<o1 1b*<o1 / 2a<1a2,1b1 2b<1a2L & IU/LI & {-}{-}/{-}{-} & B & CP:2b W:1a,2b \\
B3949 & 1a<o1 1b*<o1 / 2a<1a2,1b1 2b<1b1 & IU/IU & {-}{-}/{-}{-} & N & W:1a,2a \\
B3950 & 1a<o1 1b*<o1 / 2a<1a2,1b1 2b<1b2 & IU/IU & {-}{-}/{-}{-} & N & W:1a,2a \\
B3951 & 1a<o1 1b*<o1 / 2a<1a2,1b2 2b<1a2L & IU/LI & {-}{-}/{-}{-} & B & CP:2b W:1a,2b \\
B3952 & 1a<o1 1b*<o1 / 2a<1a2,1b2 2b<1b1 & IU/IU & {-}{-}/{-}{-} & N & W:1a,2a \\
B3953 & 1a<o1 1b*<o1 / 2a<1a2,1b2 2b<1b2 & IU/IU & {-}{-}/{-}{-} & N & W:1a,2a \\
B3954 & 1a<o1 1b*<o1 / 2a<1a2L,1b1 2b<1b1 & IU/IU & {-}{-}/{-}{-} & N & W:1a,2a \\
B3955 & 1a<o1 1b*<o1 / 2a<1a2L,1b1 2b<1b2 & IU/IU & {-}{-}/{-}{-} & N & W:1a,2a \\
B3956 & 1a<o1 1b*<o1 / 2a<1a2L,1b2 2b<1b1 & IU/IU & {-}{-}/{-}{-} & N & W:1a,2a \\
B3957 & 1a<o1 1b*<o1 / 2a<1a2L,1b2 2b<1b2 & IU/IU & {-}{-}/{-}{-} & N & W:1a,2a \\
B3958 & 1a<o1 1b*<o2 / 2a<1a0L 2b<1a1,1b1 & IU/NI & {-}{-}/{-}{-} & N & W:1a,2b \\
B3959 & 1a<o1 1b*<o2 / 2a<1a0L 2b<1a1,1b2 & IU/NI & {-}{-}/{-}{-} & N & W:1a,2b \\
B3960 & 1a<o1 1b*<o2 / 2a<1a0L 2b<1a2,1b1 & IU/NI & {-}{-}/{-}{-} & N & W:1a,2b \\
B3961 & 1a<o1 1b*<o2 / 2a<1a0L 2b<1a2,1b2 & IU/NI & {-}{-}/{-}{-} & N & W:1a,2b \\
B3962 & 1a<o1 1b*<o2 / 2a<1a0L,1b1 2b<1a1 & IU/UI & {-}{-}/{-}{-} & N & W:1a,2b \\
B3963 & 1a<o1 1b*<o2 / 2a<1a0L,1b1 2b<1a2 & IU/UI & {-}{-}/{-}{-} & N & W:1a,2b \\
B3964 & 1a<o1 1b*<o2 / 2a<1a0L,1b1 2b<1b1 & IU/UU & {-}{-}/{-}{-} & N & W:1a \\
B3965 & 1a<o1 1b*<o2 / 2a<1a0L,1b1 2b<1b2 & IU/UU & {-}{-}/{-}{-} & N & W:1a \\
B3966 & 1a<o1 1b*<o2 / 2a<1a0L,1b2 2b<1a1 & IU/UI & {-}{-}/{-}{-} & N & W:1a,2b \\
B3967 & 1a<o1 1b*<o2 / 2a<1a0L,1b2 2b<1a2 & IU/UI & {-}{-}/{-}{-} & N & W:1a,2b \\
B3968 & 1a<o1 1b*<o2 / 2a<1a0L,1b2 2b<1b1 & IU/UU & {-}{-}/{-}{-} & N & W:1a \\
B3969 & 1a<o1 1b*<o2 / 2a<1a0L,1b2 2b<1b2 & IU/UU & {-}{-}/{-}{-} & N & W:1a \\
B3970 & 1a<o1 1b*<o2 / 2a<1a1 2b<1a1,1b1 & IU/EE & {-}{-}/{-}{-} & E & W:1a \\
B3971 & 1a<o1 1b*<o2 / 2a<1a1 2b<1a1,1b2 & IU/EE & {-}{-}/{-}{-} & E & W:1a \\
B3972 & 1a<o1 1b*<o2 / 2a<1a1 2b<1a1L,1b1 & IU/LI & {-}{-}/{-}{-} & B & CP:2b W:1a,2b \\
B3973 & 1a<o1 1b*<o2 / 2a<1a1 2b<1a1L,1b2 & IU/LI & {-}{-}/{-}{-} & B & CP:2b W:1a,2b \\
B3974 & 1a<o1 1b*<o2 / 2a<1a1 2b<1a2,1b1 & IU/EE & {-}{-}/{-}{-} & E & C:2b W:1a \\
B3975 & 1a<o1 1b*<o2 / 2a<1a1 2b<1a2,1b2 & IU/EE & {-}{-}/{-}{-} & E & C:2b W:1a \\
B3976 & 1a<o1 1b*<o2 / 2a<1a1 2b<1a2L,1b1 & IU/LI & {-}{-}/{-}{-} & B & P:2b W:1a,2b \\
B3977 & 1a<o1 1b*<o2 / 2a<1a1 2b<1a2L,1b2 & IU/LI & {-}{-}/{-}{-} & B & P:2b W:1a,2b \\
B3978 & 1a<o1 1b*<o2 / 2a<1a1,1b1 2b<1a1L & IU/LI & {-}{-}/{-}{-} & B & CP:2b W:1a,2b \\
B3979 & 1a<o1 1b*<o2 / 2a<1a1,1b1 2b<1a2 & IU/EE & {-}{-}/{-}{-} & E & C:2b W:1a \\
B3980 & 1a<o1 1b*<o2 / 2a<1a1,1b1 2b<1a2L & IU/LI & {-}{-}/{-}{-} & B & P:2b W:1a,2b \\
B3981 & 1a<o1 1b*<o2 / 2a<1a1,1b1 2b<1b1 & IU/IU & {-}{-}/{-}{-} & N & W:1a,2a \\
B3982 & 1a<o1 1b*<o2 / 2a<1a1,1b1 2b<1b2 & IU/IU & {-}{-}/{-}{-} & N & W:1a,2a \\
B3983 & 1a<o1 1b*<o2 / 2a<1a1,1b2 2b<1a1L & IU/LI & {-}{-}/{-}{-} & B & CP:2b W:1a,2b \\
B3984 & 1a<o1 1b*<o2 / 2a<1a1,1b2 2b<1a2 & IU/EE & {-}{-}/{-}{-} & E & C:2b W:1a \\
B3985 & 1a<o1 1b*<o2 / 2a<1a1,1b2 2b<1a2L & IU/LI & {-}{-}/{-}{-} & B & P:2b W:1a,2b \\
B3986 & 1a<o1 1b*<o2 / 2a<1a1,1b2 2b<1b1 & IU/IU & {-}{-}/{-}{-} & N & W:1a,2a \\
B3987 & 1a<o1 1b*<o2 / 2a<1a1,1b2 2b<1b2 & IU/IU & {-}{-}/{-}{-} & N & W:1a,2a \\
B3988 & 1a<o1 1b*<o2 / 2a<1a1L 2b<1a2,1b1 & IU/IL & {-}{-}/{-}{-} & B & C:2a,2b P:2a W:1a,2a \\
B3989 & 1a<o1 1b*<o2 / 2a<1a1L 2b<1a2,1b2 & IU/IL & {-}{-}/{-}{-} & B & C:2a,2b P:2a W:1a,2a \\
B3990 & 1a<o1 1b*<o2 / 2a<1a1L,1b1 2b<1a2 & IU/IL & {-}{-}/{-}{-} & B & C:2a,2b P:2a W:1a,2a \\
B3991 & 1a<o1 1b*<o2 / 2a<1a1L,1b1 2b<1b1 & IU/IU & {-}{-}/{-}{-} & N & W:1a,2a \\
B3992 & 1a<o1 1b*<o2 / 2a<1a1L,1b1 2b<1b2 & IU/IU & {-}{-}/{-}{-} & N & W:1a,2a \\
B3993 & 1a<o1 1b*<o2 / 2a<1a1L,1b2 2b<1a2 & IU/IL & {-}{-}/{-}{-} & B & C:2a,2b P:2a W:1a,2a \\
B3994 & 1a<o1 1b*<o2 / 2a<1a1L,1b2 2b<1b1 & IU/IU & {-}{-}/{-}{-} & N & W:1a,2a \\
B3995 & 1a<o1 1b*<o2 / 2a<1a1L,1b2 2b<1b2 & IU/IU & {-}{-}/{-}{-} & N & W:1a,2a \\
B3996 & 1a<o1 1b*<o2 / 2a<1a2 2b<1a2,1b1 & IU/EE & {-}{-}/{-}{-} & E & W:1a \\
B3997 & 1a<o1 1b*<o2 / 2a<1a2 2b<1a2,1b2 & IU/EE & {-}{-}/{-}{-} & E & W:1a \\
B3998 & 1a<o1 1b*<o2 / 2a<1a2 2b<1a2L,1b1 & IU/LI & {-}{-}/{-}{-} & B & CP:2b W:1a,2b \\
B3999 & 1a<o1 1b*<o2 / 2a<1a2 2b<1a2L,1b2 & IU/LI & {-}{-}/{-}{-} & B & CP:2b W:1a,2b \\
B4000 & 1a<o1 1b*<o2 / 2a<1a2,1b1 2b<1a2L & IU/LI & {-}{-}/{-}{-} & B & CP:2b W:1a,2b \\
B4001 & 1a<o1 1b*<o2 / 2a<1a2,1b1 2b<1b1 & IU/IU & {-}{-}/{-}{-} & N & W:1a,2a \\
B4002 & 1a<o1 1b*<o2 / 2a<1a2,1b1 2b<1b2 & IU/IU & {-}{-}/{-}{-} & N & W:1a,2a \\
B4003 & 1a<o1 1b*<o2 / 2a<1a2,1b2 2b<1a2L & IU/LI & {-}{-}/{-}{-} & B & CP:2b W:1a,2b \\
B4004 & 1a<o1 1b*<o2 / 2a<1a2,1b2 2b<1b1 & IU/IU & {-}{-}/{-}{-} & N & W:1a,2a \\
B4005 & 1a<o1 1b*<o2 / 2a<1a2,1b2 2b<1b2 & IU/IU & {-}{-}/{-}{-} & N & W:1a,2a \\
B4006 & 1a<o1 1b*<o2 / 2a<1a2L,1b1 2b<1b1 & IU/IU & {-}{-}/{-}{-} & N & W:1a,2a \\
B4007 & 1a<o1 1b*<o2 / 2a<1a2L,1b1 2b<1b2 & IU/IU & {-}{-}/{-}{-} & N & W:1a,2a \\
B4008 & 1a<o1 1b*<o2 / 2a<1a2L,1b2 2b<1b1 & IU/IU & {-}{-}/{-}{-} & N & W:1a,2a \\
B4009 & 1a<o1 1b*<o2 / 2a<1a2L,1b2 2b<1b2 & IU/IU & {-}{-}/{-}{-} & N & W:1a,2a \\
B4010 & 1a<o1L 1b<o2 / 2a<1a0L 2b*<1a1,1b1 & IL/NU & +{-}/+{-} & B & C:1a,1b PW:1a \\
B4011 & 1a<o1L 1b<o2 / 2a<1a0L 2b*<1a1,1b2 & IL/NU & +{-}/+{-} & B & C:1a,1b PW:1a \\
B4012 & 1a<o1L 1b<o2 / 2a<1a0L 2b*<1a2,1b1 & IL/NU & +{-}/+{-} & B & C:1a,1b PW:1a \\
B4013 & 1a<o1L 1b<o2 / 2a<1a0L 2b*<1a2,1b2 & IL/NU & +{-}/+{-} & B & C:1a,1b PW:1a \\
B4014 & 1a<o1L 1b<o2 / 2a<1a0L,1b1 2b*<1a1 & IL/DU & +{-}/+{-} & B & C:1a,1b,2a PW:1a \\
B4015 & 1a<o1L 1b<o2 / 2a<1a0L,1b1 2b*<1a2 & IL/DU & +{-}/+{-} & B & C:1a,1b,2a PW:1a \\
B4016 & 1a<o1L 1b<o2 / 2a<1a0L,1b1 2b*<1b1 & IL/DU & +{-}/+{-} & B & C:1a,1b,2a PW:1a \\
B4017 & 1a<o1L 1b<o2 / 2a<1a0L,1b1 2b*<1b2 & IL/DU & +{-}/+{-} & B & C:1a,1b,2a PW:1a \\
B4018 & 1a<o1L 1b<o2 / 2a<1a0L,1b2 2b*<1a1 & IL/DU & +{-}/+{-} & B & C:1a,1b,2a PW:1a \\
B4019 & 1a<o1L 1b<o2 / 2a<1a0L,1b2 2b*<1a2 & IL/DU & +{-}/+{-} & B & C:1a,1b,2a PW:1a \\
B4020 & 1a<o1L 1b<o2 / 2a<1a0L,1b2 2b*<1b1 & IL/DU & +{-}/+{-} & B & C:1a,1b,2a PW:1a \\
B4021 & 1a<o1L 1b<o2 / 2a<1a0L,1b2 2b*<1b2 & IL/DU & +{-}/+{-} & B & C:1a,1b,2a PW:1a \\
B4022 & 1a<o1L 1b<o2 / 2a<1a1 2b*<1a1,1b1 & IL/IU & +{-}/{-}{-} & B & C:1a,1b,2a PW:1a,2a \\
B4023 & 1a<o1L 1b<o2 / 2a<1a1 2b*<1a1,1b2 & IL/IU & +{-}/{-}{-} & B & C:1a,1b,2a PW:1a,2a \\
B4024 & 1a<o1L 1b<o2 / 2a<1a1 2b*<1a2,1b1 & IL/IU & +{-}/{-}{-} & B & C:1a,1b,2a PW:1a,2a \\
B4025 & 1a<o1L 1b<o2 / 2a<1a1 2b*<1a2,1b2 & IL/IU & +{-}/{-}{-} & B & C:1a,1b,2a PW:1a,2a \\
B4026 & 1a<o1L 1b<o2 / 2a<1a1,1b0L 2b*<1a1 & IL/IU & +{-}/{-}{-} & B & C:1a,1b,2a PW:1a,2a \\
B4027 & 1a<o1L 1b<o2 / 2a<1a1,1b0L 2b*<1a2 & IL/IU & +{-}/{-}{-} & B & C:1a,1b,2a PW:1a,2a \\
B4028 & 1a<o1L 1b<o2 / 2a<1a1,1b0L 2b*<1b1 & IL/IU & +{-}/{-}{-} & B & C:1a,1b,2a PW:1a,2a \\
B4029 & 1a<o1L 1b<o2 / 2a<1a1,1b0L 2b*<1b2 & IL/IU & +{-}/{-}{-} & B & C:1a,1b,2a PW:1a,2a \\
B4030 & 1a<o1L 1b<o2 / 2a<1a1,1b1 2b*<1a1 & IL/EU & +{-}/{-}{-} & BF & C:1a,1b PW:1a \\
B4031 & 1a<o1L 1b<o2 / 2a<1a1,1b1 2b*<1a2 & IL/EU & +{-}/{-}{-} & BF & C:1a,1b PW:1a \\
B4032 & 1a<o1L 1b<o2 / 2a<1a1,1b1 2b*<1b1 & IL/EU & +{-}/{-}{-} & BF & C:1a,1b PW:1a \\
B4033 & 1a<o1L 1b<o2 / 2a<1a1,1b1 2b*<1b2 & IL/EU & +{-}/{-}{-} & BF & C:1a,1b PW:1a \\
B4034 & 1a<o1L 1b<o2 / 2a<1a1,1b1L 2b*<1a1 & IL/EU & +{-}/{-}{-} & BF & C:1a,1b PW:1a \\
B4035 & 1a<o1L 1b<o2 / 2a<1a1,1b1L 2b*<1a2 & IL/EU & +{-}/{-}{-} & BF & C:1a,1b PW:1a \\
B4036 & 1a<o1L 1b<o2 / 2a<1a1,1b1L 2b*<1b1 & IL/EU & +{-}/{-}{-} & BF & C:1a,1b PW:1a \\
B4037 & 1a<o1L 1b<o2 / 2a<1a1,1b1L 2b*<1b2 & IL/EU & +{-}/{-}{-} & BF & C:1a,1b PW:1a \\
B4038 & 1a<o1L 1b<o2 / 2a<1a1,1b2 2b*<1a1 & IL/EU & +{-}/{-}{-} & BF & C:1a,1b,2a PW:1a \\
B4039 & 1a<o1L 1b<o2 / 2a<1a1,1b2 2b*<1a2 & IL/EU & +{-}/{-}{-} & BF & C:1a,1b,2a PW:1a \\
B4040 & 1a<o1L 1b<o2 / 2a<1a1,1b2 2b*<1b1 & IL/EU & +{-}/{-}{-} & BF & C:1a,1b,2a PW:1a \\
B4041 & 1a<o1L 1b<o2 / 2a<1a1,1b2 2b*<1b2 & IL/EU & +{-}/{-}{-} & BF & C:1a,1b,2a PW:1a \\
B4042 & 1a<o1L 1b<o2 / 2a<1a1,1b2L 2b*<1a1 & IL/EU & +{-}/{-}{-} & BF & C:1a,1b,2a PW:1a \\
B4043 & 1a<o1L 1b<o2 / 2a<1a1,1b2L 2b*<1a2 & IL/EU & +{-}/{-}{-} & BF & C:1a,1b,2a PW:1a \\
B4044 & 1a<o1L 1b<o2 / 2a<1a1,1b2L 2b*<1b1 & IL/EU & +{-}/{-}{-} & BF & C:1a,1b,2a PW:1a \\
B4045 & 1a<o1L 1b<o2 / 2a<1a1,1b2L 2b*<1b2 & IL/EU & +{-}/{-}{-} & BF & C:1a,1b,2a PW:1a \\
B4046 & 1a<o1L 1b<o2 / 2a<1a1L 2b*<1a1,1b1 & IL/IU & +{-}/+{-} & B & C:1a,1b,2a PW:1a,2a \\
B4047 & 1a<o1L 1b<o2 / 2a<1a1L 2b*<1a1,1b2 & IL/IU & +{-}/+{-} & B & C:1a,1b,2a PW:1a,2a \\
B4048 & 1a<o1L 1b<o2 / 2a<1a1L 2b*<1a2,1b1 & IL/IU & +{-}/+{-} & B & C:1a,1b,2a PW:1a,2a \\
B4049 & 1a<o1L 1b<o2 / 2a<1a1L 2b*<1a2,1b2 & IL/IU & +{-}/+{-} & B & C:1a,1b,2a PW:1a,2a \\
B4050 & 1a<o1L 1b<o2 / 2a<1a1L,1b1 2b*<1a1 & IL/IU & +{-}/+{-} & BF & C:1a,1b,2a PW:1a,2a \\
B4051 & 1a<o1L 1b<o2 / 2a<1a1L,1b1 2b*<1a2 & IL/IU & +{-}/+{-} & BF & C:1a,1b,2a PW:1a,2a \\
B4052 & 1a<o1L 1b<o2 / 2a<1a1L,1b1 2b*<1b1 & IL/IU & +{-}/+{-} & BF & C:1a,1b,2a PW:1a,2a \\
B4053 & 1a<o1L 1b<o2 / 2a<1a1L,1b1 2b*<1b2 & IL/IU & +{-}/+{-} & BF & C:1a,1b,2a PW:1a,2a \\
B4054 & 1a<o1L 1b<o2 / 2a<1a1L,1b2 2b*<1a1 & IL/IU & +{-}/+{-} & BF & C:1a,1b PW:1a,2a \\
B4055 & 1a<o1L 1b<o2 / 2a<1a1L,1b2 2b*<1a2 & IL/IU & +{-}/+{-} & BF & C:1a,1b PW:1a,2a \\
B4056 & 1a<o1L 1b<o2 / 2a<1a1L,1b2 2b*<1b1 & IL/IU & +{-}/+{-} & BF & C:1a,1b PW:1a,2a \\
B4057 & 1a<o1L 1b<o2 / 2a<1a1L,1b2 2b*<1b2 & IL/IU & +{-}/+{-} & BF & C:1a,1b PW:1a,2a \\
B4058 & 1a<o1L 1b<o2 / 2a<1a2 2b*<1a1,1b1 & IL/IU & +{-}/{-}{-} & B & C:1a,1b,2a PW:1a,2a \\
B4059 & 1a<o1L 1b<o2 / 2a<1a2 2b*<1a1,1b2 & IL/IU & +{-}/{-}{-} & B & C:1a,1b,2a PW:1a,2a \\
B4060 & 1a<o1L 1b<o2 / 2a<1a2 2b*<1a2,1b1 & IL/IU & +{-}/{-}{-} & B & C:1a,1b,2a PW:1a,2a \\
B4061 & 1a<o1L 1b<o2 / 2a<1a2 2b*<1a2,1b2 & IL/IU & +{-}/{-}{-} & B & C:1a,1b,2a PW:1a,2a \\
B4062 & 1a<o1L 1b<o2 / 2a<1a2,1b0L 2b*<1a1 & IL/IU & +{-}/{-}{-} & B & C:1a,1b,2a PW:1a,2a \\
B4063 & 1a<o1L 1b<o2 / 2a<1a2,1b0L 2b*<1a2 & IL/IU & +{-}/{-}{-} & B & C:1a,1b,2a PW:1a,2a \\
B4064 & 1a<o1L 1b<o2 / 2a<1a2,1b0L 2b*<1b1 & IL/IU & +{-}/{-}{-} & B & C:1a,1b,2a PW:1a,2a \\
B4065 & 1a<o1L 1b<o2 / 2a<1a2,1b0L 2b*<1b2 & IL/IU & +{-}/{-}{-} & B & C:1a,1b,2a PW:1a,2a \\
B4066 & 1a<o1L 1b<o2 / 2a<1a2,1b1 2b*<1a1 & IL/EU & +{-}/{-}{-} & BF & C:1a,1b PW:1a \\
B4067 & 1a<o1L 1b<o2 / 2a<1a2,1b1 2b*<1a2 & IL/EU & +{-}/{-}{-} & BF & C:1a,1b PW:1a \\
B4068 & 1a<o1L 1b<o2 / 2a<1a2,1b1 2b*<1b1 & IL/EU & +{-}/{-}{-} & BF & C:1a,1b PW:1a \\
B4069 & 1a<o1L 1b<o2 / 2a<1a2,1b1 2b*<1b2 & IL/EU & +{-}/{-}{-} & BF & C:1a,1b PW:1a \\
B4070 & 1a<o1L 1b<o2 / 2a<1a2,1b1L 2b*<1a1 & IL/EU & +{-}/{-}{-} & BF & C:1a,1b PW:1a \\
B4071 & 1a<o1L 1b<o2 / 2a<1a2,1b1L 2b*<1a2 & IL/EU & +{-}/{-}{-} & BF & C:1a,1b PW:1a \\
B4072 & 1a<o1L 1b<o2 / 2a<1a2,1b1L 2b*<1b1 & IL/EU & +{-}/{-}{-} & BF & C:1a,1b PW:1a \\
B4073 & 1a<o1L 1b<o2 / 2a<1a2,1b1L 2b*<1b2 & IL/EU & +{-}/{-}{-} & BF & C:1a,1b PW:1a \\
B4074 & 1a<o1L 1b<o2 / 2a<1a2,1b2 2b*<1a1 & IL/EU & +{-}/{-}{-} & BF & C:1a,1b PW:1a \\
B4075 & 1a<o1L 1b<o2 / 2a<1a2,1b2 2b*<1a2 & IL/EU & +{-}/{-}{-} & BF & C:1a,1b PW:1a \\
B4076 & 1a<o1L 1b<o2 / 2a<1a2,1b2 2b*<1b1 & IL/EU & +{-}/{-}{-} & BF & C:1a,1b PW:1a \\
B4077 & 1a<o1L 1b<o2 / 2a<1a2,1b2 2b*<1b2 & IL/EU & +{-}/{-}{-} & BF & C:1a,1b PW:1a \\
B4078 & 1a<o1L 1b<o2 / 2a<1a2,1b2L 2b*<1a1 & IL/EU & +{-}/{-}{-} & BF & C:1a,1b PW:1a \\
B4079 & 1a<o1L 1b<o2 / 2a<1a2,1b2L 2b*<1a2 & IL/EU & +{-}/{-}{-} & BF & C:1a,1b PW:1a \\
B4080 & 1a<o1L 1b<o2 / 2a<1a2,1b2L 2b*<1b1 & IL/EU & +{-}/{-}{-} & BF & C:1a,1b PW:1a \\
B4081 & 1a<o1L 1b<o2 / 2a<1a2,1b2L 2b*<1b2 & IL/EU & +{-}/{-}{-} & BF & C:1a,1b PW:1a \\
B4082 & 1a<o1L 1b<o2 / 2a<1a2L 2b*<1a1,1b1 & IL/IU & +{-}/+{-} & B & C:1a,1b,2a PW:1a,2a \\
B4083 & 1a<o1L 1b<o2 / 2a<1a2L 2b*<1a1,1b2 & IL/IU & +{-}/+{-} & B & C:1a,1b,2a PW:1a,2a \\
B4084 & 1a<o1L 1b<o2 / 2a<1a2L 2b*<1a2,1b1 & IL/IU & +{-}/+{-} & B & C:1a,1b,2a PW:1a,2a \\
B4085 & 1a<o1L 1b<o2 / 2a<1a2L 2b*<1a2,1b2 & IL/IU & +{-}/+{-} & B & C:1a,1b,2a PW:1a,2a \\
B4086 & 1a<o1L 1b<o2 / 2a<1a2L,1b1 2b*<1a1 & IL/IU & +{-}/+{-} & BF & C:1a,1b,2a PW:1a,2a \\
B4087 & 1a<o1L 1b<o2 / 2a<1a2L,1b1 2b*<1a2 & IL/IU & +{-}/+{-} & BF & C:1a,1b,2a PW:1a,2a \\
B4088 & 1a<o1L 1b<o2 / 2a<1a2L,1b1 2b*<1b1 & IL/IU & +{-}/+{-} & BF & C:1a,1b,2a PW:1a,2a \\
B4089 & 1a<o1L 1b<o2 / 2a<1a2L,1b1 2b*<1b2 & IL/IU & +{-}/+{-} & BF & C:1a,1b,2a PW:1a,2a \\
B4090 & 1a<o1L 1b<o2 / 2a<1a2L,1b2 2b*<1a1 & IL/IU & +{-}/+{-} & BF & C:1a,1b,2a PW:1a,2a \\
B4091 & 1a<o1L 1b<o2 / 2a<1a2L,1b2 2b*<1a2 & IL/IU & +{-}/+{-} & BF & C:1a,1b,2a PW:1a,2a \\
B4092 & 1a<o1L 1b<o2 / 2a<1a2L,1b2 2b*<1b1 & IL/IU & +{-}/+{-} & BF & C:1a,1b,2a PW:1a,2a \\
B4093 & 1a<o1L 1b<o2 / 2a<1a2L,1b2 2b*<1b2 & IL/IU & +{-}/+{-} & BF & C:1a,1b,2a PW:1a,2a \\
B4094 & 1a<o1L 1b<o2 / 2a<1b0L 2b*<1a1,1b1 & IL/NU & +{-}/{-}{-} & B & C:1a,1b PW:1a \\
B4095 & 1a<o1L 1b<o2 / 2a<1b0L 2b*<1a1,1b2 & IL/NU & +{-}/{-}{-} & B & C:1a,1b PW:1a \\
B4096 & 1a<o1L 1b<o2 / 2a<1b0L 2b*<1a2,1b1 & IL/NU & +{-}/{-}{-} & B & C:1a,1b PW:1a \\
B4097 & 1a<o1L 1b<o2 / 2a<1b0L 2b*<1a2,1b2 & IL/NU & +{-}/{-}{-} & B & C:1a,1b PW:1a \\
B4098 & 1a<o1L 1b<o2 / 2a<1b1 2b*<1a1,1b1 & IL/DU & +{-}/{-}{-} & B & C:1a,1b,2a PW:1a \\
B4099 & 1a<o1L 1b<o2 / 2a<1b1 2b*<1a1,1b2 & IL/DU & +{-}/{-}{-} & B & C:1a,1b,2a PW:1a \\
B4100 & 1a<o1L 1b<o2 / 2a<1b1 2b*<1a2,1b1 & IL/DU & +{-}/{-}{-} & B & C:1a,1b,2a PW:1a \\
B4101 & 1a<o1L 1b<o2 / 2a<1b1 2b*<1a2,1b2 & IL/DU & +{-}/{-}{-} & B & C:1a,1b,2a PW:1a \\
B4102 & 1a<o1L 1b<o2 / 2a<1b1L 2b*<1a1,1b1 & IL/DU & +{-}/{-}{-} & B & C:1a,1b,2a PW:1a \\
B4103 & 1a<o1L 1b<o2 / 2a<1b1L 2b*<1a1,1b2 & IL/DU & +{-}/{-}{-} & B & C:1a,1b,2a PW:1a \\
B4104 & 1a<o1L 1b<o2 / 2a<1b1L 2b*<1a2,1b1 & IL/DU & +{-}/{-}{-} & B & C:1a,1b,2a PW:1a \\
B4105 & 1a<o1L 1b<o2 / 2a<1b1L 2b*<1a2,1b2 & IL/DU & +{-}/{-}{-} & B & C:1a,1b,2a PW:1a \\
B4106 & 1a<o1L 1b<o2 / 2a<1b2 2b*<1a1,1b1 & IL/DU & +{-}/{-}{-} & B & C:1a,1b,2a PW:1a \\
B4107 & 1a<o1L 1b<o2 / 2a<1b2 2b*<1a1,1b2 & IL/DU & +{-}/{-}{-} & B & C:1a,1b,2a PW:1a \\
B4108 & 1a<o1L 1b<o2 / 2a<1b2 2b*<1a2,1b1 & IL/DU & +{-}/{-}{-} & B & C:1a,1b,2a PW:1a \\
B4109 & 1a<o1L 1b<o2 / 2a<1b2 2b*<1a2,1b2 & IL/DU & +{-}/{-}{-} & B & C:1a,1b,2a PW:1a \\
B4110 & 1a<o1L 1b<o2 / 2a<1b2L 2b*<1a1,1b1 & IL/DU & +{-}/{-}{-} & B & C:1a,1b,2a PW:1a \\
B4111 & 1a<o1L 1b<o2 / 2a<1b2L 2b*<1a1,1b2 & IL/DU & +{-}/{-}{-} & B & C:1a,1b,2a PW:1a \\
B4112 & 1a<o1L 1b<o2 / 2a<1b2L 2b*<1a2,1b1 & IL/DU & +{-}/{-}{-} & B & C:1a,1b,2a PW:1a \\
B4113 & 1a<o1L 1b<o2 / 2a<1b2L 2b*<1a2,1b2 & IL/DU & +{-}/{-}{-} & B & C:1a,1b,2a PW:1a \\
B4114 & 1a<o1L 1b*<o1 / 2a<1a0L 2b<1a1,1b1 & IU/NI & +{-}/+{-} & N & W:1a,2b \\
B4115 & 1a<o1L 1b*<o1 / 2a<1a0L 2b<1a1,1b2 & IU/NI & +{-}/+{-} & N & W:1a,2b \\
B4116 & 1a<o1L 1b*<o1 / 2a<1a0L 2b<1a2,1b1 & IU/NI & +{-}/+{-} & N & W:1a,2b \\
B4117 & 1a<o1L 1b*<o1 / 2a<1a0L 2b<1a2,1b2 & IU/NI & +{-}/+{-} & N & W:1a,2b \\
B4118 & 1a<o1L 1b*<o1 / 2a<1a0L,1b1 2b<1a1 & IU/UI & +{-}/+{-} & N & W:1a,2b \\
B4119 & 1a<o1L 1b*<o1 / 2a<1a0L,1b1 2b<1a2 & IU/UI & +{-}/+{-} & N & W:1a,2b \\
B4120 & 1a<o1L 1b*<o1 / 2a<1a0L,1b1 2b<1b1 & IU/UU & +{-}/+{-} & N & W:1a \\
B4121 & 1a<o1L 1b*<o1 / 2a<1a0L,1b1 2b<1b2 & IU/UU & +{-}/+{-} & N & W:1a \\
B4122 & 1a<o1L 1b*<o1 / 2a<1a0L,1b2 2b<1a1 & IU/UI & +{-}/+{-} & N & W:1a,2b \\
B4123 & 1a<o1L 1b*<o1 / 2a<1a0L,1b2 2b<1a2 & IU/UI & +{-}/+{-} & N & W:1a,2b \\
B4124 & 1a<o1L 1b*<o1 / 2a<1a0L,1b2 2b<1b1 & IU/UU & +{-}/+{-} & N & W:1a \\
B4125 & 1a<o1L 1b*<o1 / 2a<1a0L,1b2 2b<1b2 & IU/UU & +{-}/+{-} & N & W:1a \\
B4126 & 1a<o1L 1b*<o1 / 2a<1a1 2b<1a1,1b1 & IU/EE & +{-}/{-}{-} & E & W:1a \\
B4127 & 1a<o1L 1b*<o1 / 2a<1a1 2b<1a1,1b2 & IU/EE & +{-}/{-}{-} & E & W:1a \\
B4128 & 1a<o1L 1b*<o1 / 2a<1a1 2b<1a1L,1b1 & IU/LI & +{-}/{-}+ & B & CP:2b W:1a,2b \\
B4129 & 1a<o1L 1b*<o1 / 2a<1a1 2b<1a1L,1b2 & IU/LI & +{-}/{-}+ & B & CP:2b W:1a,2b \\
B4130 & 1a<o1L 1b*<o1 / 2a<1a1 2b<1a2,1b1 & IU/EE & +{-}/{-}{-} & E & C:2b W:1a \\
B4131 & 1a<o1L 1b*<o1 / 2a<1a1 2b<1a2,1b2 & IU/EE & +{-}/{-}{-} & E & C:2b W:1a \\
B4132 & 1a<o1L 1b*<o1 / 2a<1a1 2b<1a2L,1b1 & IU/LI & +{-}/{-}+ & B & P:2b W:1a,2b \\
B4133 & 1a<o1L 1b*<o1 / 2a<1a1 2b<1a2L,1b2 & IU/LI & +{-}/{-}+ & B & P:2b W:1a,2b \\
B4134 & 1a<o1L 1b*<o1 / 2a<1a1,1b1 2b<1a1L & IU/LI & +{-}/{-}+ & B & CP:2b W:1a,2b \\
B4135 & 1a<o1L 1b*<o1 / 2a<1a1,1b1 2b<1a2 & IU/EE & +{-}/{-}{-} & E & C:2b W:1a \\
B4136 & 1a<o1L 1b*<o1 / 2a<1a1,1b1 2b<1a2L & IU/LI & +{-}/{-}+ & B & P:2b W:1a,2b \\
B4137 & 1a<o1L 1b*<o1 / 2a<1a1,1b1 2b<1b1 & IU/IU & +{-}/{-}{-} & N & W:1a,2a \\
B4138 & 1a<o1L 1b*<o1 / 2a<1a1,1b1 2b<1b2 & IU/IU & +{-}/{-}{-} & N & W:1a,2a \\
B4139 & 1a<o1L 1b*<o1 / 2a<1a1,1b2 2b<1a1L & IU/LI & +{-}/{-}+ & B & CP:2b W:1a,2b \\
B4140 & 1a<o1L 1b*<o1 / 2a<1a1,1b2 2b<1a2 & IU/EE & +{-}/{-}{-} & E & C:2b W:1a \\
B4141 & 1a<o1L 1b*<o1 / 2a<1a1,1b2 2b<1a2L & IU/LI & +{-}/{-}+ & B & P:2b W:1a,2b \\
B4142 & 1a<o1L 1b*<o1 / 2a<1a1,1b2 2b<1b1 & IU/IU & +{-}/{-}{-} & N & W:1a,2a \\
B4143 & 1a<o1L 1b*<o1 / 2a<1a1,1b2 2b<1b2 & IU/IU & +{-}/{-}{-} & N & W:1a,2a \\
B4144 & 1a<o1L 1b*<o1 / 2a<1a1L 2b<1a2,1b1 & IU/IL & +{-}/+{-} & B & C:2a,2b P:2a W:1a,2a \\
B4145 & 1a<o1L 1b*<o1 / 2a<1a1L 2b<1a2,1b2 & IU/IL & +{-}/+{-} & B & C:2a,2b P:2a W:1a,2a \\
B4146 & 1a<o1L 1b*<o1 / 2a<1a1L,1b1 2b<1a2 & IU/IL & +{-}/+{-} & B & C:2a,2b P:2a W:1a,2a \\
B4147 & 1a<o1L 1b*<o1 / 2a<1a1L,1b1 2b<1b1 & IU/IU & +{-}/+{-} & N & W:1a,2a \\
B4148 & 1a<o1L 1b*<o1 / 2a<1a1L,1b1 2b<1b2 & IU/IU & +{-}/+{-} & N & W:1a,2a \\
B4149 & 1a<o1L 1b*<o1 / 2a<1a1L,1b2 2b<1a2 & IU/IL & +{-}/+{-} & B & C:2a,2b P:2a W:1a,2a \\
B4150 & 1a<o1L 1b*<o1 / 2a<1a1L,1b2 2b<1b1 & IU/IU & +{-}/+{-} & N & W:1a,2a \\
B4151 & 1a<o1L 1b*<o1 / 2a<1a1L,1b2 2b<1b2 & IU/IU & +{-}/+{-} & N & W:1a,2a \\
B4152 & 1a<o1L 1b*<o1 / 2a<1a2 2b<1a2,1b1 & IU/EE & +{-}/{-}{-} & E & W:1a \\
B4153 & 1a<o1L 1b*<o1 / 2a<1a2 2b<1a2,1b2 & IU/EE & +{-}/{-}{-} & E & W:1a \\
B4154 & 1a<o1L 1b*<o1 / 2a<1a2 2b<1a2L,1b1 & IU/LI & +{-}/{-}+ & B & CP:2b W:1a,2b \\
B4155 & 1a<o1L 1b*<o1 / 2a<1a2 2b<1a2L,1b2 & IU/LI & +{-}/{-}+ & B & CP:2b W:1a,2b \\
B4156 & 1a<o1L 1b*<o1 / 2a<1a2,1b1 2b<1a2L & IU/LI & +{-}/{-}+ & B & CP:2b W:1a,2b \\
B4157 & 1a<o1L 1b*<o1 / 2a<1a2,1b1 2b<1b1 & IU/IU & +{-}/{-}{-} & N & W:1a,2a \\
B4158 & 1a<o1L 1b*<o1 / 2a<1a2,1b1 2b<1b2 & IU/IU & +{-}/{-}{-} & N & W:1a,2a \\
B4159 & 1a<o1L 1b*<o1 / 2a<1a2,1b2 2b<1a2L & IU/LI & +{-}/{-}+ & B & CP:2b W:1a,2b \\
B4160 & 1a<o1L 1b*<o1 / 2a<1a2,1b2 2b<1b1 & IU/IU & +{-}/{-}{-} & N & W:1a,2a \\
B4161 & 1a<o1L 1b*<o1 / 2a<1a2,1b2 2b<1b2 & IU/IU & +{-}/{-}{-} & N & W:1a,2a \\
B4162 & 1a<o1L 1b*<o1 / 2a<1a2L,1b1 2b<1b1 & IU/IU & +{-}/+{-} & N & W:1a,2a \\
B4163 & 1a<o1L 1b*<o1 / 2a<1a2L,1b1 2b<1b2 & IU/IU & +{-}/+{-} & N & W:1a,2a \\
B4164 & 1a<o1L 1b*<o1 / 2a<1a2L,1b2 2b<1b1 & IU/IU & +{-}/+{-} & N & W:1a,2a \\
B4165 & 1a<o1L 1b*<o1 / 2a<1a2L,1b2 2b<1b2 & IU/IU & +{-}/+{-} & N & W:1a,2a \\
B4166 & 1a<o1L 1b*<o2 / 2a<1a0L 2b<1a1,1b1 & IU/NI & +{-}/+{-} & N & W:1a,2b \\
B4167 & 1a<o1L 1b*<o2 / 2a<1a0L 2b<1a1,1b2 & IU/NI & +{-}/+{-} & N & W:1a,2b \\
B4168 & 1a<o1L 1b*<o2 / 2a<1a0L 2b<1a2,1b1 & IU/NI & +{-}/+{-} & N & W:1a,2b \\
B4169 & 1a<o1L 1b*<o2 / 2a<1a0L 2b<1a2,1b2 & IU/NI & +{-}/+{-} & N & W:1a,2b \\
B4170 & 1a<o1L 1b*<o2 / 2a<1a0L,1b1 2b<1a1 & IU/UI & +{-}/+{-} & N & W:1a,2b \\
B4171 & 1a<o1L 1b*<o2 / 2a<1a0L,1b1 2b<1a2 & IU/UI & +{-}/+{-} & N & W:1a,2b \\
B4172 & 1a<o1L 1b*<o2 / 2a<1a0L,1b1 2b<1b1 & IU/UU & +{-}/+{-} & N & W:1a \\
B4173 & 1a<o1L 1b*<o2 / 2a<1a0L,1b1 2b<1b2 & IU/UU & +{-}/+{-} & N & W:1a \\
B4174 & 1a<o1L 1b*<o2 / 2a<1a0L,1b2 2b<1a1 & IU/UI & +{-}/+{-} & N & W:1a,2b \\
B4175 & 1a<o1L 1b*<o2 / 2a<1a0L,1b2 2b<1a2 & IU/UI & +{-}/+{-} & N & W:1a,2b \\
B4176 & 1a<o1L 1b*<o2 / 2a<1a0L,1b2 2b<1b1 & IU/UU & +{-}/+{-} & N & W:1a \\
B4177 & 1a<o1L 1b*<o2 / 2a<1a0L,1b2 2b<1b2 & IU/UU & +{-}/+{-} & N & W:1a \\
B4178 & 1a<o1L 1b*<o2 / 2a<1a1 2b<1a1,1b1 & IU/EE & +{-}/{-}{-} & E & W:1a \\
B4179 & 1a<o1L 1b*<o2 / 2a<1a1 2b<1a1,1b2 & IU/EE & +{-}/{-}{-} & E & W:1a \\
B4180 & 1a<o1L 1b*<o2 / 2a<1a1 2b<1a1L,1b1 & IU/LI & +{-}/{-}+ & B & CP:2b W:1a,2b \\
B4181 & 1a<o1L 1b*<o2 / 2a<1a1 2b<1a1L,1b2 & IU/LI & +{-}/{-}+ & B & CP:2b W:1a,2b \\
B4182 & 1a<o1L 1b*<o2 / 2a<1a1 2b<1a2,1b1 & IU/EE & +{-}/{-}{-} & E & C:2b W:1a \\
B4183 & 1a<o1L 1b*<o2 / 2a<1a1 2b<1a2,1b2 & IU/EE & +{-}/{-}{-} & E & C:2b W:1a \\
B4184 & 1a<o1L 1b*<o2 / 2a<1a1 2b<1a2L,1b1 & IU/LI & +{-}/{-}+ & B & P:2b W:1a,2b \\
B4185 & 1a<o1L 1b*<o2 / 2a<1a1 2b<1a2L,1b2 & IU/LI & +{-}/{-}+ & B & P:2b W:1a,2b \\
B4186 & 1a<o1L 1b*<o2 / 2a<1a1,1b1 2b<1a1L & IU/LI & +{-}/{-}+ & B & CP:2b W:1a,2b \\
B4187 & 1a<o1L 1b*<o2 / 2a<1a1,1b1 2b<1a2 & IU/EE & +{-}/{-}{-} & E & C:2b W:1a \\
B4188 & 1a<o1L 1b*<o2 / 2a<1a1,1b1 2b<1a2L & IU/LI & +{-}/{-}+ & B & P:2b W:1a,2b \\
B4189 & 1a<o1L 1b*<o2 / 2a<1a1,1b1 2b<1b1 & IU/IU & +{-}/{-}{-} & N & W:1a,2a \\
B4190 & 1a<o1L 1b*<o2 / 2a<1a1,1b1 2b<1b2 & IU/IU & +{-}/{-}{-} & N & W:1a,2a \\
B4191 & 1a<o1L 1b*<o2 / 2a<1a1,1b2 2b<1a1L & IU/LI & +{-}/{-}+ & B & CP:2b W:1a,2b \\
B4192 & 1a<o1L 1b*<o2 / 2a<1a1,1b2 2b<1a2 & IU/EE & +{-}/{-}{-} & E & C:2b W:1a \\
B4193 & 1a<o1L 1b*<o2 / 2a<1a1,1b2 2b<1a2L & IU/LI & +{-}/{-}+ & B & P:2b W:1a,2b \\
B4194 & 1a<o1L 1b*<o2 / 2a<1a1,1b2 2b<1b1 & IU/IU & +{-}/{-}{-} & N & W:1a,2a \\
B4195 & 1a<o1L 1b*<o2 / 2a<1a1,1b2 2b<1b2 & IU/IU & +{-}/{-}{-} & N & W:1a,2a \\
B4196 & 1a<o1L 1b*<o2 / 2a<1a1L 2b<1a2,1b1 & IU/IL & +{-}/+{-} & B & C:2a,2b P:2a W:1a,2a \\
B4197 & 1a<o1L 1b*<o2 / 2a<1a1L 2b<1a2,1b2 & IU/IL & +{-}/+{-} & B & C:2a,2b P:2a W:1a,2a \\
B4198 & 1a<o1L 1b*<o2 / 2a<1a1L,1b1 2b<1a2 & IU/IL & +{-}/+{-} & B & C:2a,2b P:2a W:1a,2a \\
B4199 & 1a<o1L 1b*<o2 / 2a<1a1L,1b1 2b<1b1 & IU/IU & +{-}/+{-} & N & W:1a,2a \\
B4200 & 1a<o1L 1b*<o2 / 2a<1a1L,1b1 2b<1b2 & IU/IU & +{-}/+{-} & N & W:1a,2a \\
B4201 & 1a<o1L 1b*<o2 / 2a<1a1L,1b2 2b<1a2 & IU/IL & +{-}/+{-} & B & C:2a,2b P:2a W:1a,2a \\
B4202 & 1a<o1L 1b*<o2 / 2a<1a1L,1b2 2b<1b1 & IU/IU & +{-}/+{-} & N & W:1a,2a \\
B4203 & 1a<o1L 1b*<o2 / 2a<1a1L,1b2 2b<1b2 & IU/IU & +{-}/+{-} & N & W:1a,2a \\
B4204 & 1a<o1L 1b*<o2 / 2a<1a2 2b<1a2,1b1 & IU/EE & +{-}/{-}{-} & E & W:1a \\
B4205 & 1a<o1L 1b*<o2 / 2a<1a2 2b<1a2,1b2 & IU/EE & +{-}/{-}{-} & E & W:1a \\
B4206 & 1a<o1L 1b*<o2 / 2a<1a2 2b<1a2L,1b1 & IU/LI & +{-}/{-}+ & B & CP:2b W:1a,2b \\
B4207 & 1a<o1L 1b*<o2 / 2a<1a2 2b<1a2L,1b2 & IU/LI & +{-}/{-}+ & B & CP:2b W:1a,2b \\
B4208 & 1a<o1L 1b*<o2 / 2a<1a2,1b1 2b<1a2L & IU/LI & +{-}/{-}+ & B & CP:2b W:1a,2b \\
B4209 & 1a<o1L 1b*<o2 / 2a<1a2,1b1 2b<1b1 & IU/IU & +{-}/{-}{-} & N & W:1a,2a \\
B4210 & 1a<o1L 1b*<o2 / 2a<1a2,1b1 2b<1b2 & IU/IU & +{-}/{-}{-} & N & W:1a,2a \\
B4211 & 1a<o1L 1b*<o2 / 2a<1a2,1b2 2b<1a2L & IU/LI & +{-}/{-}+ & B & CP:2b W:1a,2b \\
B4212 & 1a<o1L 1b*<o2 / 2a<1a2,1b2 2b<1b1 & IU/IU & +{-}/{-}{-} & N & W:1a,2a \\
B4213 & 1a<o1L 1b*<o2 / 2a<1a2,1b2 2b<1b2 & IU/IU & +{-}/{-}{-} & N & W:1a,2a \\
B4214 & 1a<o1L 1b*<o2 / 2a<1a2L,1b1 2b<1b1 & IU/IU & +{-}/+{-} & N & W:1a,2a \\
B4215 & 1a<o1L 1b*<o2 / 2a<1a2L,1b1 2b<1b2 & IU/IU & +{-}/+{-} & N & W:1a,2a \\
B4216 & 1a<o1L 1b*<o2 / 2a<1a2L,1b2 2b<1b1 & IU/IU & +{-}/+{-} & N & W:1a,2a \\
B4217 & 1a<o1L 1b*<o2 / 2a<1a2L,1b2 2b<1b2 & IU/IU & +{-}/+{-} & N & W:1a,2a \\
B4218 & 1a<o2 1b<o2 / 2a<1a0L 2b*<1a1,1b1 & EE/NU & {-}{-}/{-}{-} & E &  \\
B4219 & 1a<o2 1b<o2 / 2a<1a0L 2b*<1a1,1b2 & EE/NU & {-}{-}/{-}{-} & E &  \\
B4220 & 1a<o2 1b<o2 / 2a<1a0L 2b*<1a2,1b1 & EE/NU & {-}{-}/{-}{-} & E &  \\
B4221 & 1a<o2 1b<o2 / 2a<1a0L 2b*<1a2,1b2 & EE/NU & {-}{-}/{-}{-} & E &  \\
B4222 & 1a<o2 1b<o2 / 2a<1a0L,1b1 2b*<1a1 & EE/DU & {-}{-}/{-}{-} & E &  \\
B4223 & 1a<o2 1b<o2 / 2a<1a0L,1b1 2b*<1a2 & EE/DU & {-}{-}/{-}{-} & E &  \\
B4224 & 1a<o2 1b<o2 / 2a<1a0L,1b1 2b*<1b1 & EE/DU & {-}{-}/{-}{-} & E &  \\
B4225 & 1a<o2 1b<o2 / 2a<1a0L,1b1 2b*<1b2 & EE/DU & {-}{-}/{-}{-} & E &  \\
B4226 & 1a<o2 1b<o2 / 2a<1a0L,1b2 2b*<1a1 & EE/DU & {-}{-}/{-}{-} & E &  \\
B4227 & 1a<o2 1b<o2 / 2a<1a0L,1b2 2b*<1a2 & EE/DU & {-}{-}/{-}{-} & E &  \\
B4228 & 1a<o2 1b<o2 / 2a<1a0L,1b2 2b*<1b1 & EE/DU & {-}{-}/{-}{-} & E &  \\
B4229 & 1a<o2 1b<o2 / 2a<1a0L,1b2 2b*<1b2 & EE/DU & {-}{-}/{-}{-} & E &  \\
B4230 & 1a<o2 1b<o2 / 2a<1a1 2b*<1a1,1b1 & EE/DU & {-}{-}/{-}{-} & E &  \\
B4231 & 1a<o2 1b<o2 / 2a<1a1 2b*<1a1,1b2 & EE/DU & {-}{-}/{-}{-} & E &  \\
B4232 & 1a<o2 1b<o2 / 2a<1a1 2b*<1a2,1b1 & EE/DU & {-}{-}/{-}{-} & E &  \\
B4233 & 1a<o2 1b<o2 / 2a<1a1 2b*<1a2,1b2 & EE/DU & {-}{-}/{-}{-} & E &  \\
B4234 & 1a<o2 1b<o2 / 2a<1a1,1b1 2b*<1a1 & EE/DU & {-}{-}/{-}{-} & E &  \\
B4235 & 1a<o2 1b<o2 / 2a<1a1,1b1 2b*<1a2 & EE/DU & {-}{-}/{-}{-} & E &  \\
B4236 & 1a<o2 1b<o2 / 2a<1a1,1b1L 2b*<1a1 & EE/DU & {-}{-}/{-}{-} & E &  \\
B4237 & 1a<o2 1b<o2 / 2a<1a1,1b1L 2b*<1a2 & EE/DU & {-}{-}/{-}{-} & E &  \\
B4238 & 1a<o2 1b<o2 / 2a<1a1,1b1L 2b*<1b1 & EE/DU & {-}{-}/{-}{-} & E &  \\
B4239 & 1a<o2 1b<o2 / 2a<1a1,1b1L 2b*<1b2 & EE/DU & {-}{-}/{-}{-} & E &  \\
B4240 & 1a<o2 1b<o2 / 2a<1a1,1b2 2b*<1a1 & EE/DU & {-}{-}/{-}{-} & E &  \\
B4241 & 1a<o2 1b<o2 / 2a<1a1,1b2 2b*<1a2 & EE/DU & {-}{-}/{-}{-} & E &  \\
B4242 & 1a<o2 1b<o2 / 2a<1a1,1b2 2b*<1b1 & EE/DU & {-}{-}/{-}{-} & E &  \\
B4243 & 1a<o2 1b<o2 / 2a<1a1,1b2 2b*<1b2 & EE/DU & {-}{-}/{-}{-} & E &  \\
B4244 & 1a<o2 1b<o2 / 2a<1a1,1b2L 2b*<1a1 & EE/DU & {-}{-}/{-}{-} & E &  \\
B4245 & 1a<o2 1b<o2 / 2a<1a1,1b2L 2b*<1a2 & EE/DU & {-}{-}/{-}{-} & E &  \\
B4246 & 1a<o2 1b<o2 / 2a<1a1,1b2L 2b*<1b1 & EE/DU & {-}{-}/{-}{-} & E &  \\
B4247 & 1a<o2 1b<o2 / 2a<1a1,1b2L 2b*<1b2 & EE/DU & {-}{-}/{-}{-} & E &  \\
B4248 & 1a<o2 1b<o2 / 2a<1a1L 2b*<1a1,1b1 & EE/DU & {-}{-}/{-}{-} & E &  \\
B4249 & 1a<o2 1b<o2 / 2a<1a1L 2b*<1a1,1b2 & EE/DU & {-}{-}/{-}{-} & E &  \\
B4250 & 1a<o2 1b<o2 / 2a<1a1L 2b*<1a2,1b1 & EE/DU & {-}{-}/{-}{-} & E &  \\
B4251 & 1a<o2 1b<o2 / 2a<1a1L 2b*<1a2,1b2 & EE/DU & {-}{-}/{-}{-} & E &  \\
B4252 & 1a<o2 1b<o2 / 2a<1a1L,1b2 2b*<1a1 & EE/DU & {-}{-}/{-}{-} & E &  \\
B4253 & 1a<o2 1b<o2 / 2a<1a1L,1b2 2b*<1a2 & EE/DU & {-}{-}/{-}{-} & E &  \\
B4254 & 1a<o2 1b<o2 / 2a<1a1L,1b2 2b*<1b1 & EE/DU & {-}{-}/{-}{-} & E &  \\
B4255 & 1a<o2 1b<o2 / 2a<1a1L,1b2 2b*<1b2 & EE/DU & {-}{-}/{-}{-} & E &  \\
B4256 & 1a<o2 1b<o2 / 2a<1a2 2b*<1a1,1b1 & EE/DU & {-}{-}/{-}{-} & E &  \\
B4257 & 1a<o2 1b<o2 / 2a<1a2 2b*<1a1,1b2 & EE/DU & {-}{-}/{-}{-} & E &  \\
B4258 & 1a<o2 1b<o2 / 2a<1a2 2b*<1a2,1b1 & EE/DU & {-}{-}/{-}{-} & E &  \\
B4259 & 1a<o2 1b<o2 / 2a<1a2 2b*<1a2,1b2 & EE/DU & {-}{-}/{-}{-} & E &  \\
B4260 & 1a<o2 1b<o2 / 2a<1a2,1b2 2b*<1a1 & EE/DU & {-}{-}/{-}{-} & E &  \\
B4261 & 1a<o2 1b<o2 / 2a<1a2,1b2 2b*<1a2 & EE/DU & {-}{-}/{-}{-} & E &  \\
B4262 & 1a<o2 1b<o2 / 2a<1a2,1b2L 2b*<1a1 & EE/DU & {-}{-}/{-}{-} & E &  \\
B4263 & 1a<o2 1b<o2 / 2a<1a2,1b2L 2b*<1a2 & EE/DU & {-}{-}/{-}{-} & E &  \\
B4264 & 1a<o2 1b<o2 / 2a<1a2,1b2L 2b*<1b1 & EE/DU & {-}{-}/{-}{-} & E &  \\
B4265 & 1a<o2 1b<o2 / 2a<1a2,1b2L 2b*<1b2 & EE/DU & {-}{-}/{-}{-} & E &  \\
B4266 & 1a<o2 1b<o2 / 2a<1a2L 2b*<1a1,1b1 & EE/DU & {-}{-}/{-}{-} & E &  \\
B4267 & 1a<o2 1b<o2 / 2a<1a2L 2b*<1a1,1b2 & EE/DU & {-}{-}/{-}{-} & E &  \\
B4268 & 1a<o2 1b<o2 / 2a<1a2L 2b*<1a2,1b1 & EE/DU & {-}{-}/{-}{-} & E &  \\
B4269 & 1a<o2 1b<o2 / 2a<1a2L 2b*<1a2,1b2 & EE/DU & {-}{-}/{-}{-} & E &  \\
B4270 & 1a<o2 1b<o2L / 2a<1a0L 2b*<1a1,1b1 & LI/NU & {-}+/{-}{-} & B & CPW:1b \\
B4271 & 1a<o2 1b<o2L / 2a<1a0L 2b*<1a1,1b2 & LI/NU & {-}+/{-}{-} & B & CPW:1b \\
B4272 & 1a<o2 1b<o2L / 2a<1a0L 2b*<1a2,1b1 & LI/NU & {-}+/{-}{-} & B & CPW:1b \\
B4273 & 1a<o2 1b<o2L / 2a<1a0L 2b*<1a2,1b2 & LI/NU & {-}+/{-}{-} & B & CPW:1b \\
B4274 & 1a<o2 1b<o2L / 2a<1a0L,1b1 2b*<1a1 & LI/IU & {-}+/{-}{-} & B & CPW:1b,2a \\
B4275 & 1a<o2 1b<o2L / 2a<1a0L,1b1 2b*<1a2 & LI/IU & {-}+/{-}{-} & B & CPW:1b,2a \\
B4276 & 1a<o2 1b<o2L / 2a<1a0L,1b1 2b*<1b1 & LI/IU & {-}+/{-}{-} & B & CPW:1b,2a \\
B4277 & 1a<o2 1b<o2L / 2a<1a0L,1b1 2b*<1b2 & LI/IU & {-}+/{-}{-} & B & CPW:1b,2a \\
B4278 & 1a<o2 1b<o2L / 2a<1a0L,1b2 2b*<1a1 & LI/IU & {-}+/{-}{-} & B & CPW:1b,2a \\
B4279 & 1a<o2 1b<o2L / 2a<1a0L,1b2 2b*<1a2 & LI/IU & {-}+/{-}{-} & B & CPW:1b,2a \\
B4280 & 1a<o2 1b<o2L / 2a<1a0L,1b2 2b*<1b1 & LI/IU & {-}+/{-}{-} & B & CPW:1b,2a \\
B4281 & 1a<o2 1b<o2L / 2a<1a0L,1b2 2b*<1b2 & LI/IU & {-}+/{-}{-} & B & CPW:1b,2a \\
B4282 & 1a<o2 1b<o2L / 2a<1a1 2b*<1a1,1b1 & LI/DU & {-}+/{-}{-} & B & CPW:1b \\
B4283 & 1a<o2 1b<o2L / 2a<1a1 2b*<1a1,1b2 & LI/DU & {-}+/{-}{-} & B & CPW:1b \\
B4284 & 1a<o2 1b<o2L / 2a<1a1 2b*<1a2,1b1 & LI/DU & {-}+/{-}{-} & B & CPW:1b \\
B4285 & 1a<o2 1b<o2L / 2a<1a1 2b*<1a2,1b2 & LI/DU & {-}+/{-}{-} & B & CPW:1b \\
B4286 & 1a<o2 1b<o2L / 2a<1a1,1b0L 2b*<1a1 & LI/DU & {-}+/+{-} & B & CPW:1b \\
B4287 & 1a<o2 1b<o2L / 2a<1a1,1b0L 2b*<1a2 & LI/DU & {-}+/+{-} & B & CPW:1b \\
B4288 & 1a<o2 1b<o2L / 2a<1a1,1b0L 2b*<1b1 & LI/DU & {-}+/+{-} & B & CPW:1b \\
B4289 & 1a<o2 1b<o2L / 2a<1a1,1b0L 2b*<1b2 & LI/DU & {-}+/+{-} & B & CPW:1b \\
B4290 & 1a<o2 1b<o2L / 2a<1a1,1b1 2b*<1a1 & LI/EU & {-}+/{-}{-} & BF & CPW:1b \\
B4291 & 1a<o2 1b<o2L / 2a<1a1,1b1 2b*<1a2 & LI/EU & {-}+/{-}{-} & BF & CPW:1b \\
B4292 & 1a<o2 1b<o2L / 2a<1a1,1b1 2b*<1b1 & LI/EU & {-}+/{-}{-} & BF & CPW:1b \\
B4293 & 1a<o2 1b<o2L / 2a<1a1,1b1 2b*<1b2 & LI/EU & {-}+/{-}{-} & BF & CPW:1b \\
B4294 & 1a<o2 1b<o2L / 2a<1a1,1b1L 2b*<1a1 & LI/IU & {-}+/+{-} & BF & CPW:1b,2a \\
B4295 & 1a<o2 1b<o2L / 2a<1a1,1b1L 2b*<1a2 & LI/IU & {-}+/+{-} & BF & CPW:1b,2a \\
B4296 & 1a<o2 1b<o2L / 2a<1a1,1b1L 2b*<1b1 & LI/IU & {-}+/+{-} & BF & CPW:1b,2a \\
B4297 & 1a<o2 1b<o2L / 2a<1a1,1b1L 2b*<1b2 & LI/IU & {-}+/+{-} & BF & CPW:1b,2a \\
B4298 & 1a<o2 1b<o2L / 2a<1a1,1b2 2b*<1a1 & LI/EU & {-}+/{-}{-} & BF & CPW:1b \\
B4299 & 1a<o2 1b<o2L / 2a<1a1,1b2 2b*<1a2 & LI/EU & {-}+/{-}{-} & BF & CPW:1b \\
B4300 & 1a<o2 1b<o2L / 2a<1a1,1b2 2b*<1b1 & LI/EU & {-}+/{-}{-} & BF & CPW:1b \\
B4301 & 1a<o2 1b<o2L / 2a<1a1,1b2 2b*<1b2 & LI/EU & {-}+/{-}{-} & BF & CPW:1b \\
B4302 & 1a<o2 1b<o2L / 2a<1a1,1b2L 2b*<1a1 & LI/IU & {-}+/+{-} & BF & CPW:1b,2a \\
B4303 & 1a<o2 1b<o2L / 2a<1a1,1b2L 2b*<1a2 & LI/IU & {-}+/+{-} & BF & CPW:1b,2a \\
B4304 & 1a<o2 1b<o2L / 2a<1a1,1b2L 2b*<1b1 & LI/IU & {-}+/+{-} & BF & CPW:1b,2a \\
B4305 & 1a<o2 1b<o2L / 2a<1a1,1b2L 2b*<1b2 & LI/IU & {-}+/+{-} & BF & CPW:1b,2a \\
B4306 & 1a<o2 1b<o2L / 2a<1a1L 2b*<1a1,1b1 & LI/DU & {-}+/{-}{-} & B & CPW:1b \\
B4307 & 1a<o2 1b<o2L / 2a<1a1L 2b*<1a1,1b2 & LI/DU & {-}+/{-}{-} & B & CPW:1b \\
B4308 & 1a<o2 1b<o2L / 2a<1a1L 2b*<1a2,1b1 & LI/DU & {-}+/{-}{-} & B & CPW:1b \\
B4309 & 1a<o2 1b<o2L / 2a<1a1L 2b*<1a2,1b2 & LI/DU & {-}+/{-}{-} & B & CPW:1b \\
B4310 & 1a<o2 1b<o2L / 2a<1a1L,1b1 2b*<1a1 & LI/EU & {-}+/{-}{-} & BF & CPW:1b \\
B4311 & 1a<o2 1b<o2L / 2a<1a1L,1b1 2b*<1a2 & LI/EU & {-}+/{-}{-} & BF & CPW:1b \\
B4312 & 1a<o2 1b<o2L / 2a<1a1L,1b1 2b*<1b1 & LI/EU & {-}+/{-}{-} & BF & CPW:1b \\
B4313 & 1a<o2 1b<o2L / 2a<1a1L,1b1 2b*<1b2 & LI/EU & {-}+/{-}{-} & BF & CPW:1b \\
B4314 & 1a<o2 1b<o2L / 2a<1a1L,1b2 2b*<1a1 & LI/EU & {-}+/{-}{-} & BF & CPW:1b \\
B4315 & 1a<o2 1b<o2L / 2a<1a1L,1b2 2b*<1a2 & LI/EU & {-}+/{-}{-} & BF & CPW:1b \\
B4316 & 1a<o2 1b<o2L / 2a<1a1L,1b2 2b*<1b1 & LI/EU & {-}+/{-}{-} & BF & CPW:1b \\
B4317 & 1a<o2 1b<o2L / 2a<1a1L,1b2 2b*<1b2 & LI/EU & {-}+/{-}{-} & BF & CPW:1b \\
B4318 & 1a<o2 1b<o2L / 2a<1a2 2b*<1a1,1b1 & LI/DU & {-}+/{-}{-} & B & CPW:1b \\
B4319 & 1a<o2 1b<o2L / 2a<1a2 2b*<1a1,1b2 & LI/DU & {-}+/{-}{-} & B & CPW:1b \\
B4320 & 1a<o2 1b<o2L / 2a<1a2 2b*<1a2,1b1 & LI/DU & {-}+/{-}{-} & B & CPW:1b \\
B4321 & 1a<o2 1b<o2L / 2a<1a2 2b*<1a2,1b2 & LI/DU & {-}+/{-}{-} & B & CPW:1b \\
B4322 & 1a<o2 1b<o2L / 2a<1a2,1b0L 2b*<1a1 & LI/DU & {-}+/+{-} & B & CPW:1b \\
B4323 & 1a<o2 1b<o2L / 2a<1a2,1b0L 2b*<1a2 & LI/DU & {-}+/+{-} & B & CPW:1b \\
B4324 & 1a<o2 1b<o2L / 2a<1a2,1b0L 2b*<1b1 & LI/DU & {-}+/+{-} & B & CPW:1b \\
B4325 & 1a<o2 1b<o2L / 2a<1a2,1b0L 2b*<1b2 & LI/DU & {-}+/+{-} & B & CPW:1b \\
B4326 & 1a<o2 1b<o2L / 2a<1a2,1b1 2b*<1a1 & LI/EU & {-}+/{-}{-} & BF & CPW:1b \\
B4327 & 1a<o2 1b<o2L / 2a<1a2,1b1 2b*<1a2 & LI/EU & {-}+/{-}{-} & BF & CPW:1b \\
B4328 & 1a<o2 1b<o2L / 2a<1a2,1b1 2b*<1b1 & LI/EU & {-}+/{-}{-} & BF & CPW:1b \\
B4329 & 1a<o2 1b<o2L / 2a<1a2,1b1 2b*<1b2 & LI/EU & {-}+/{-}{-} & BF & CPW:1b \\
B4330 & 1a<o2 1b<o2L / 2a<1a2,1b1L 2b*<1a1 & LI/IU & {-}+/+{-} & BF & CPW:1b,2a \\
B4331 & 1a<o2 1b<o2L / 2a<1a2,1b1L 2b*<1a2 & LI/IU & {-}+/+{-} & BF & CPW:1b,2a \\
B4332 & 1a<o2 1b<o2L / 2a<1a2,1b1L 2b*<1b1 & LI/IU & {-}+/+{-} & BF & CPW:1b,2a \\
B4333 & 1a<o2 1b<o2L / 2a<1a2,1b1L 2b*<1b2 & LI/IU & {-}+/+{-} & BF & CPW:1b,2a \\
B4334 & 1a<o2 1b<o2L / 2a<1a2,1b2 2b*<1a1 & LI/EU & {-}+/{-}{-} & BF & CPW:1b \\
B4335 & 1a<o2 1b<o2L / 2a<1a2,1b2 2b*<1a2 & LI/EU & {-}+/{-}{-} & BF & CPW:1b \\
B4336 & 1a<o2 1b<o2L / 2a<1a2,1b2 2b*<1b1 & LI/EU & {-}+/{-}{-} & BF & CPW:1b \\
B4337 & 1a<o2 1b<o2L / 2a<1a2,1b2 2b*<1b2 & LI/EU & {-}+/{-}{-} & BF & CPW:1b \\
B4338 & 1a<o2 1b<o2L / 2a<1a2,1b2L 2b*<1a1 & LI/IU & {-}+/+{-} & BF & CPW:1b,2a \\
B4339 & 1a<o2 1b<o2L / 2a<1a2,1b2L 2b*<1a2 & LI/IU & {-}+/+{-} & BF & CPW:1b,2a \\
B4340 & 1a<o2 1b<o2L / 2a<1a2,1b2L 2b*<1b1 & LI/IU & {-}+/+{-} & BF & CPW:1b,2a \\
B4341 & 1a<o2 1b<o2L / 2a<1a2,1b2L 2b*<1b2 & LI/IU & {-}+/+{-} & BF & CPW:1b,2a \\
B4342 & 1a<o2 1b<o2L / 2a<1a2L 2b*<1a1,1b1 & LI/DU & {-}+/{-}{-} & B & CPW:1b \\
B4343 & 1a<o2 1b<o2L / 2a<1a2L 2b*<1a1,1b2 & LI/DU & {-}+/{-}{-} & B & CPW:1b \\
B4344 & 1a<o2 1b<o2L / 2a<1a2L 2b*<1a2,1b1 & LI/DU & {-}+/{-}{-} & B & CPW:1b \\
B4345 & 1a<o2 1b<o2L / 2a<1a2L 2b*<1a2,1b2 & LI/DU & {-}+/{-}{-} & B & CPW:1b \\
B4346 & 1a<o2 1b<o2L / 2a<1a2L,1b1 2b*<1a1 & LI/EU & {-}+/{-}{-} & BF & CPW:1b \\
B4347 & 1a<o2 1b<o2L / 2a<1a2L,1b1 2b*<1a2 & LI/EU & {-}+/{-}{-} & BF & CPW:1b \\
B4348 & 1a<o2 1b<o2L / 2a<1a2L,1b1 2b*<1b1 & LI/EU & {-}+/{-}{-} & BF & CPW:1b \\
B4349 & 1a<o2 1b<o2L / 2a<1a2L,1b1 2b*<1b2 & LI/EU & {-}+/{-}{-} & BF & CPW:1b \\
B4350 & 1a<o2 1b<o2L / 2a<1a2L,1b2 2b*<1a1 & LI/EU & {-}+/{-}{-} & BF & CPW:1b \\
B4351 & 1a<o2 1b<o2L / 2a<1a2L,1b2 2b*<1a2 & LI/EU & {-}+/{-}{-} & BF & CPW:1b \\
B4352 & 1a<o2 1b<o2L / 2a<1a2L,1b2 2b*<1b1 & LI/EU & {-}+/{-}{-} & BF & CPW:1b \\
B4353 & 1a<o2 1b<o2L / 2a<1a2L,1b2 2b*<1b2 & LI/EU & {-}+/{-}{-} & BF & CPW:1b \\
B4354 & 1a<o2 1b<o2L / 2a<1b0L 2b*<1a1,1b1 & LI/NU & {-}+/+{-} & B & CPW:1b \\
B4355 & 1a<o2 1b<o2L / 2a<1b0L 2b*<1a1,1b2 & LI/NU & {-}+/+{-} & B & CPW:1b \\
B4356 & 1a<o2 1b<o2L / 2a<1b0L 2b*<1a2,1b1 & LI/NU & {-}+/+{-} & B & CPW:1b \\
B4357 & 1a<o2 1b<o2L / 2a<1b0L 2b*<1a2,1b2 & LI/NU & {-}+/+{-} & B & CPW:1b \\
B4358 & 1a<o2 1b<o2L / 2a<1b1 2b*<1a1,1b1 & LI/IU & {-}+/{-}{-} & B & CPW:1b,2a \\
B4359 & 1a<o2 1b<o2L / 2a<1b1 2b*<1a1,1b2 & LI/IU & {-}+/{-}{-} & B & CPW:1b,2a \\
B4360 & 1a<o2 1b<o2L / 2a<1b1 2b*<1a2,1b1 & LI/IU & {-}+/{-}{-} & B & CPW:1b,2a \\
B4361 & 1a<o2 1b<o2L / 2a<1b1 2b*<1a2,1b2 & LI/IU & {-}+/{-}{-} & B & CPW:1b,2a \\
B4362 & 1a<o2 1b<o2L / 2a<1b1L 2b*<1a1,1b1 & LI/IU & {-}+/+{-} & B & CPW:1b,2a \\
B4363 & 1a<o2 1b<o2L / 2a<1b1L 2b*<1a1,1b2 & LI/IU & {-}+/+{-} & B & CPW:1b,2a \\
B4364 & 1a<o2 1b<o2L / 2a<1b1L 2b*<1a2,1b1 & LI/IU & {-}+/+{-} & B & CPW:1b,2a \\
B4365 & 1a<o2 1b<o2L / 2a<1b1L 2b*<1a2,1b2 & LI/IU & {-}+/+{-} & B & CPW:1b,2a \\
B4366 & 1a<o2 1b<o2L / 2a<1b2 2b*<1a1,1b1 & LI/IU & {-}+/{-}{-} & B & CPW:1b,2a \\
B4367 & 1a<o2 1b<o2L / 2a<1b2 2b*<1a1,1b2 & LI/IU & {-}+/{-}{-} & B & CPW:1b,2a \\
B4368 & 1a<o2 1b<o2L / 2a<1b2 2b*<1a2,1b1 & LI/IU & {-}+/{-}{-} & B & CPW:1b,2a \\
B4369 & 1a<o2 1b<o2L / 2a<1b2 2b*<1a2,1b2 & LI/IU & {-}+/{-}{-} & B & CPW:1b,2a \\
B4370 & 1a<o2 1b<o2L / 2a<1b2L 2b*<1a1,1b1 & LI/IU & {-}+/+{-} & B & CPW:1b,2a \\
B4371 & 1a<o2 1b<o2L / 2a<1b2L 2b*<1a1,1b2 & LI/IU & {-}+/+{-} & B & CPW:1b,2a \\
B4372 & 1a<o2 1b<o2L / 2a<1b2L 2b*<1a2,1b1 & LI/IU & {-}+/+{-} & B & CPW:1b,2a \\
B4373 & 1a<o2 1b<o2L / 2a<1b2L 2b*<1a2,1b2 & LI/IU & {-}+/+{-} & B & CPW:1b,2a \\
B4374 & 1a<o2 1b*<o1 / 2a<1a0L 2b<1a1,1b1 & IU/NI & {-}{-}/{-}{-} & N & W:1a,2b \\
B4375 & 1a<o2 1b*<o1 / 2a<1a0L 2b<1a1,1b2 & IU/NI & {-}{-}/{-}{-} & N & W:1a,2b \\
B4376 & 1a<o2 1b*<o1 / 2a<1a0L 2b<1a2,1b1 & IU/NI & {-}{-}/{-}{-} & N & W:1a,2b \\
B4377 & 1a<o2 1b*<o1 / 2a<1a0L 2b<1a2,1b2 & IU/NI & {-}{-}/{-}{-} & N & W:1a,2b \\
B4378 & 1a<o2 1b*<o1 / 2a<1a0L,1b1 2b<1a1 & IU/UI & {-}{-}/{-}{-} & N & W:1a,2b \\
B4379 & 1a<o2 1b*<o1 / 2a<1a0L,1b1 2b<1a2 & IU/UI & {-}{-}/{-}{-} & N & W:1a,2b \\
B4380 & 1a<o2 1b*<o1 / 2a<1a0L,1b1 2b<1b1 & IU/UU & {-}{-}/{-}{-} & N & W:1a \\
B4381 & 1a<o2 1b*<o1 / 2a<1a0L,1b1 2b<1b2 & IU/UU & {-}{-}/{-}{-} & N & W:1a \\
B4382 & 1a<o2 1b*<o1 / 2a<1a0L,1b2 2b<1a1 & IU/UI & {-}{-}/{-}{-} & N & W:1a,2b \\
B4383 & 1a<o2 1b*<o1 / 2a<1a0L,1b2 2b<1a2 & IU/UI & {-}{-}/{-}{-} & N & W:1a,2b \\
B4384 & 1a<o2 1b*<o1 / 2a<1a0L,1b2 2b<1b1 & IU/UU & {-}{-}/{-}{-} & N & W:1a \\
B4385 & 1a<o2 1b*<o1 / 2a<1a0L,1b2 2b<1b2 & IU/UU & {-}{-}/{-}{-} & N & W:1a \\
B4386 & 1a<o2 1b*<o1 / 2a<1a1 2b<1a1,1b1 & IU/EE & {-}{-}/{-}{-} & E & W:1a \\
B4387 & 1a<o2 1b*<o1 / 2a<1a1 2b<1a1,1b2 & IU/EE & {-}{-}/{-}{-} & E & W:1a \\
B4388 & 1a<o2 1b*<o1 / 2a<1a1 2b<1a1L,1b1 & IU/LI & {-}{-}/{-}{-} & B & CP:2b W:1a,2b \\
B4389 & 1a<o2 1b*<o1 / 2a<1a1 2b<1a1L,1b2 & IU/LI & {-}{-}/{-}{-} & B & CP:2b W:1a,2b \\
B4390 & 1a<o2 1b*<o1 / 2a<1a1 2b<1a2,1b1 & IU/EE & {-}{-}/{-}{-} & E & C:2b W:1a \\
B4391 & 1a<o2 1b*<o1 / 2a<1a1 2b<1a2,1b2 & IU/EE & {-}{-}/{-}{-} & E & C:2b W:1a \\
B4392 & 1a<o2 1b*<o1 / 2a<1a1 2b<1a2L,1b1 & IU/LI & {-}{-}/{-}{-} & B & P:2b W:1a,2b \\
B4393 & 1a<o2 1b*<o1 / 2a<1a1 2b<1a2L,1b2 & IU/LI & {-}{-}/{-}{-} & B & P:2b W:1a,2b \\
B4394 & 1a<o2 1b*<o1 / 2a<1a1,1b1 2b<1a1L & IU/LI & {-}{-}/{-}{-} & B & CP:2b W:1a,2b \\
B4395 & 1a<o2 1b*<o1 / 2a<1a1,1b1 2b<1a2 & IU/EE & {-}{-}/{-}{-} & E & C:2b W:1a \\
B4396 & 1a<o2 1b*<o1 / 2a<1a1,1b1 2b<1a2L & IU/LI & {-}{-}/{-}{-} & B & P:2b W:1a,2b \\
B4397 & 1a<o2 1b*<o1 / 2a<1a1,1b1 2b<1b1 & IU/IU & {-}{-}/{-}{-} & N & W:1a,2a \\
B4398 & 1a<o2 1b*<o1 / 2a<1a1,1b1 2b<1b2 & IU/IU & {-}{-}/{-}{-} & N & W:1a,2a \\
B4399 & 1a<o2 1b*<o1 / 2a<1a1,1b2 2b<1a1L & IU/LI & {-}{-}/{-}{-} & B & CP:2b W:1a,2b \\
B4400 & 1a<o2 1b*<o1 / 2a<1a1,1b2 2b<1a2 & IU/EE & {-}{-}/{-}{-} & E & C:2b W:1a \\
B4401 & 1a<o2 1b*<o1 / 2a<1a1,1b2 2b<1a2L & IU/LI & {-}{-}/{-}{-} & B & P:2b W:1a,2b \\
B4402 & 1a<o2 1b*<o1 / 2a<1a1,1b2 2b<1b1 & IU/IU & {-}{-}/{-}{-} & N & W:1a,2a \\
B4403 & 1a<o2 1b*<o1 / 2a<1a1,1b2 2b<1b2 & IU/IU & {-}{-}/{-}{-} & N & W:1a,2a \\
B4404 & 1a<o2 1b*<o1 / 2a<1a1L 2b<1a2,1b1 & IU/IL & {-}{-}/{-}{-} & B & C:2a,2b P:2a W:1a,2a \\
B4405 & 1a<o2 1b*<o1 / 2a<1a1L 2b<1a2,1b2 & IU/IL & {-}{-}/{-}{-} & B & C:2a,2b P:2a W:1a,2a \\
B4406 & 1a<o2 1b*<o1 / 2a<1a1L,1b1 2b<1a2 & IU/IL & {-}{-}/{-}{-} & B & C:2a,2b P:2a W:1a,2a \\
B4407 & 1a<o2 1b*<o1 / 2a<1a1L,1b1 2b<1b1 & IU/IU & {-}{-}/{-}{-} & N & W:1a,2a \\
B4408 & 1a<o2 1b*<o1 / 2a<1a1L,1b1 2b<1b2 & IU/IU & {-}{-}/{-}{-} & N & W:1a,2a \\
B4409 & 1a<o2 1b*<o1 / 2a<1a1L,1b2 2b<1a2 & IU/IL & {-}{-}/{-}{-} & B & C:2a,2b P:2a W:1a,2a \\
B4410 & 1a<o2 1b*<o1 / 2a<1a1L,1b2 2b<1b1 & IU/IU & {-}{-}/{-}{-} & N & W:1a,2a \\
B4411 & 1a<o2 1b*<o1 / 2a<1a1L,1b2 2b<1b2 & IU/IU & {-}{-}/{-}{-} & N & W:1a,2a \\
B4412 & 1a<o2 1b*<o1 / 2a<1a2 2b<1a2,1b1 & IU/EE & {-}{-}/{-}{-} & E & W:1a \\
B4413 & 1a<o2 1b*<o1 / 2a<1a2 2b<1a2,1b2 & IU/EE & {-}{-}/{-}{-} & E & W:1a \\
B4414 & 1a<o2 1b*<o1 / 2a<1a2 2b<1a2L,1b1 & IU/LI & {-}{-}/{-}{-} & B & CP:2b W:1a,2b \\
B4415 & 1a<o2 1b*<o1 / 2a<1a2 2b<1a2L,1b2 & IU/LI & {-}{-}/{-}{-} & B & CP:2b W:1a,2b \\
B4416 & 1a<o2 1b*<o1 / 2a<1a2,1b1 2b<1a2L & IU/LI & {-}{-}/{-}{-} & B & CP:2b W:1a,2b \\
B4417 & 1a<o2 1b*<o1 / 2a<1a2,1b1 2b<1b1 & IU/IU & {-}{-}/{-}{-} & N & W:1a,2a \\
B4418 & 1a<o2 1b*<o1 / 2a<1a2,1b1 2b<1b2 & IU/IU & {-}{-}/{-}{-} & N & W:1a,2a \\
B4419 & 1a<o2 1b*<o1 / 2a<1a2,1b2 2b<1a2L & IU/LI & {-}{-}/{-}{-} & B & CP:2b W:1a,2b \\
B4420 & 1a<o2 1b*<o1 / 2a<1a2,1b2 2b<1b1 & IU/IU & {-}{-}/{-}{-} & N & W:1a,2a \\
B4421 & 1a<o2 1b*<o1 / 2a<1a2,1b2 2b<1b2 & IU/IU & {-}{-}/{-}{-} & N & W:1a,2a \\
B4422 & 1a<o2 1b*<o1 / 2a<1a2L,1b1 2b<1b1 & IU/IU & {-}{-}/{-}{-} & N & W:1a,2a \\
B4423 & 1a<o2 1b*<o1 / 2a<1a2L,1b1 2b<1b2 & IU/IU & {-}{-}/{-}{-} & N & W:1a,2a \\
B4424 & 1a<o2 1b*<o1 / 2a<1a2L,1b2 2b<1b1 & IU/IU & {-}{-}/{-}{-} & N & W:1a,2a \\
B4425 & 1a<o2 1b*<o1 / 2a<1a2L,1b2 2b<1b2 & IU/IU & {-}{-}/{-}{-} & N & W:1a,2a \\
B4426 & 1a<o2 1b*<o2 / 2a<1a0L 2b<1a1,1b1 & IU/NI & {-}{-}/{-}{-} & N & W:1a,2b \\
B4427 & 1a<o2 1b*<o2 / 2a<1a0L 2b<1a1,1b2 & IU/NI & {-}{-}/{-}{-} & N & W:1a,2b \\
B4428 & 1a<o2 1b*<o2 / 2a<1a0L 2b<1a2,1b1 & IU/NI & {-}{-}/{-}{-} & N & W:1a,2b \\
B4429 & 1a<o2 1b*<o2 / 2a<1a0L 2b<1a2,1b2 & IU/NI & {-}{-}/{-}{-} & N & W:1a,2b \\
B4430 & 1a<o2 1b*<o2 / 2a<1a0L,1b1 2b<1a1 & IU/UI & {-}{-}/{-}{-} & N & W:1a,2b \\
B4431 & 1a<o2 1b*<o2 / 2a<1a0L,1b1 2b<1a2 & IU/UI & {-}{-}/{-}{-} & N & W:1a,2b \\
B4432 & 1a<o2 1b*<o2 / 2a<1a0L,1b1 2b<1b1 & IU/UU & {-}{-}/{-}{-} & N & W:1a \\
B4433 & 1a<o2 1b*<o2 / 2a<1a0L,1b1 2b<1b2 & IU/UU & {-}{-}/{-}{-} & N & W:1a \\
B4434 & 1a<o2 1b*<o2 / 2a<1a0L,1b2 2b<1a1 & IU/UI & {-}{-}/{-}{-} & N & W:1a,2b \\
B4435 & 1a<o2 1b*<o2 / 2a<1a0L,1b2 2b<1a2 & IU/UI & {-}{-}/{-}{-} & N & W:1a,2b \\
B4436 & 1a<o2 1b*<o2 / 2a<1a0L,1b2 2b<1b1 & IU/UU & {-}{-}/{-}{-} & N & W:1a \\
B4437 & 1a<o2 1b*<o2 / 2a<1a0L,1b2 2b<1b2 & IU/UU & {-}{-}/{-}{-} & N & W:1a \\
B4438 & 1a<o2 1b*<o2 / 2a<1a1 2b<1a1,1b1 & IU/EE & {-}{-}/{-}{-} & E & W:1a \\
B4439 & 1a<o2 1b*<o2 / 2a<1a1 2b<1a1,1b2 & IU/EE & {-}{-}/{-}{-} & E & W:1a \\
B4440 & 1a<o2 1b*<o2 / 2a<1a1 2b<1a1L,1b1 & IU/LI & {-}{-}/{-}{-} & B & CP:2b W:1a,2b \\
B4441 & 1a<o2 1b*<o2 / 2a<1a1 2b<1a1L,1b2 & IU/LI & {-}{-}/{-}{-} & B & CP:2b W:1a,2b \\
B4442 & 1a<o2 1b*<o2 / 2a<1a1 2b<1a2,1b1 & IU/EE & {-}{-}/{-}{-} & E & C:2b W:1a \\
B4443 & 1a<o2 1b*<o2 / 2a<1a1 2b<1a2,1b2 & IU/EE & {-}{-}/{-}{-} & E & C:2b W:1a \\
B4444 & 1a<o2 1b*<o2 / 2a<1a1 2b<1a2L,1b1 & IU/LI & {-}{-}/{-}{-} & B & P:2b W:1a,2b \\
B4445 & 1a<o2 1b*<o2 / 2a<1a1 2b<1a2L,1b2 & IU/LI & {-}{-}/{-}{-} & B & P:2b W:1a,2b \\
B4446 & 1a<o2 1b*<o2 / 2a<1a1,1b1 2b<1a1L & IU/LI & {-}{-}/{-}{-} & B & CP:2b W:1a,2b \\
B4447 & 1a<o2 1b*<o2 / 2a<1a1,1b1 2b<1a2 & IU/EE & {-}{-}/{-}{-} & E & C:2b W:1a \\
B4448 & 1a<o2 1b*<o2 / 2a<1a1,1b1 2b<1a2L & IU/LI & {-}{-}/{-}{-} & B & P:2b W:1a,2b \\
B4449 & 1a<o2 1b*<o2 / 2a<1a1,1b1 2b<1b1 & IU/IU & {-}{-}/{-}{-} & N & W:1a,2a \\
B4450 & 1a<o2 1b*<o2 / 2a<1a1,1b1 2b<1b2 & IU/IU & {-}{-}/{-}{-} & N & W:1a,2a \\
B4451 & 1a<o2 1b*<o2 / 2a<1a1,1b2 2b<1a1L & IU/LI & {-}{-}/{-}{-} & B & CP:2b W:1a,2b \\
B4452 & 1a<o2 1b*<o2 / 2a<1a1,1b2 2b<1a2 & IU/EE & {-}{-}/{-}{-} & E & C:2b W:1a \\
B4453 & 1a<o2 1b*<o2 / 2a<1a1,1b2 2b<1a2L & IU/LI & {-}{-}/{-}{-} & B & P:2b W:1a,2b \\
B4454 & 1a<o2 1b*<o2 / 2a<1a1,1b2 2b<1b1 & IU/IU & {-}{-}/{-}{-} & N & W:1a,2a \\
B4455 & 1a<o2 1b*<o2 / 2a<1a1,1b2 2b<1b2 & IU/IU & {-}{-}/{-}{-} & N & W:1a,2a \\
B4456 & 1a<o2 1b*<o2 / 2a<1a1L 2b<1a2,1b1 & IU/IL & {-}{-}/{-}{-} & B & C:2a,2b P:2a W:1a,2a \\
B4457 & 1a<o2 1b*<o2 / 2a<1a1L 2b<1a2,1b2 & IU/IL & {-}{-}/{-}{-} & B & C:2a,2b P:2a W:1a,2a \\
B4458 & 1a<o2 1b*<o2 / 2a<1a1L,1b1 2b<1a2 & IU/IL & {-}{-}/{-}{-} & B & C:2a,2b P:2a W:1a,2a \\
B4459 & 1a<o2 1b*<o2 / 2a<1a1L,1b1 2b<1b1 & IU/IU & {-}{-}/{-}{-} & N & W:1a,2a \\
B4460 & 1a<o2 1b*<o2 / 2a<1a1L,1b1 2b<1b2 & IU/IU & {-}{-}/{-}{-} & N & W:1a,2a \\
B4461 & 1a<o2 1b*<o2 / 2a<1a1L,1b2 2b<1a2 & IU/IL & {-}{-}/{-}{-} & B & C:2a,2b P:2a W:1a,2a \\
B4462 & 1a<o2 1b*<o2 / 2a<1a1L,1b2 2b<1b1 & IU/IU & {-}{-}/{-}{-} & N & W:1a,2a \\
B4463 & 1a<o2 1b*<o2 / 2a<1a1L,1b2 2b<1b2 & IU/IU & {-}{-}/{-}{-} & N & W:1a,2a \\
B4464 & 1a<o2 1b*<o2 / 2a<1a2 2b<1a2,1b1 & IU/EE & {-}{-}/{-}{-} & E & W:1a \\
B4465 & 1a<o2 1b*<o2 / 2a<1a2 2b<1a2,1b2 & IU/EE & {-}{-}/{-}{-} & E & W:1a \\
B4466 & 1a<o2 1b*<o2 / 2a<1a2 2b<1a2L,1b1 & IU/LI & {-}{-}/{-}{-} & B & CP:2b W:1a,2b \\
B4467 & 1a<o2 1b*<o2 / 2a<1a2 2b<1a2L,1b2 & IU/LI & {-}{-}/{-}{-} & B & CP:2b W:1a,2b \\
B4468 & 1a<o2 1b*<o2 / 2a<1a2,1b1 2b<1a2L & IU/LI & {-}{-}/{-}{-} & B & CP:2b W:1a,2b \\
B4469 & 1a<o2 1b*<o2 / 2a<1a2,1b1 2b<1b1 & IU/IU & {-}{-}/{-}{-} & N & W:1a,2a \\
B4470 & 1a<o2 1b*<o2 / 2a<1a2,1b1 2b<1b2 & IU/IU & {-}{-}/{-}{-} & N & W:1a,2a \\
B4471 & 1a<o2 1b*<o2 / 2a<1a2,1b2 2b<1a2L & IU/LI & {-}{-}/{-}{-} & B & CP:2b W:1a,2b \\
B4472 & 1a<o2 1b*<o2 / 2a<1a2,1b2 2b<1b1 & IU/IU & {-}{-}/{-}{-} & N & W:1a,2a \\
B4473 & 1a<o2 1b*<o2 / 2a<1a2,1b2 2b<1b2 & IU/IU & {-}{-}/{-}{-} & N & W:1a,2a \\
B4474 & 1a<o2 1b*<o2 / 2a<1a2L,1b1 2b<1b1 & IU/IU & {-}{-}/{-}{-} & N & W:1a,2a \\
B4475 & 1a<o2 1b*<o2 / 2a<1a2L,1b1 2b<1b2 & IU/IU & {-}{-}/{-}{-} & N & W:1a,2a \\
B4476 & 1a<o2 1b*<o2 / 2a<1a2L,1b2 2b<1b1 & IU/IU & {-}{-}/{-}{-} & N & W:1a,2a \\
B4477 & 1a<o2 1b*<o2 / 2a<1a2L,1b2 2b<1b2 & IU/IU & {-}{-}/{-}{-} & N & W:1a,2a \\
B4478 & 1a<o2L 1b*<o1 / 2a<1a0L 2b<1a1,1b1 & IU/NI & +{-}/+{-} & N & W:1a,2b \\
B4479 & 1a<o2L 1b*<o1 / 2a<1a0L 2b<1a1,1b2 & IU/NI & +{-}/+{-} & N & W:1a,2b \\
B4480 & 1a<o2L 1b*<o1 / 2a<1a0L 2b<1a2,1b1 & IU/NI & +{-}/+{-} & N & W:1a,2b \\
B4481 & 1a<o2L 1b*<o1 / 2a<1a0L 2b<1a2,1b2 & IU/NI & +{-}/+{-} & N & W:1a,2b \\
B4482 & 1a<o2L 1b*<o1 / 2a<1a0L,1b1 2b<1a1 & IU/UI & +{-}/+{-} & N & W:1a,2b \\
B4483 & 1a<o2L 1b*<o1 / 2a<1a0L,1b1 2b<1a2 & IU/UI & +{-}/+{-} & N & W:1a,2b \\
B4484 & 1a<o2L 1b*<o1 / 2a<1a0L,1b1 2b<1b1 & IU/UU & +{-}/+{-} & N & W:1a \\
B4485 & 1a<o2L 1b*<o1 / 2a<1a0L,1b1 2b<1b2 & IU/UU & +{-}/+{-} & N & W:1a \\
B4486 & 1a<o2L 1b*<o1 / 2a<1a0L,1b2 2b<1a1 & IU/UI & +{-}/+{-} & N & W:1a,2b \\
B4487 & 1a<o2L 1b*<o1 / 2a<1a0L,1b2 2b<1a2 & IU/UI & +{-}/+{-} & N & W:1a,2b \\
B4488 & 1a<o2L 1b*<o1 / 2a<1a0L,1b2 2b<1b1 & IU/UU & +{-}/+{-} & N & W:1a \\
B4489 & 1a<o2L 1b*<o1 / 2a<1a0L,1b2 2b<1b2 & IU/UU & +{-}/+{-} & N & W:1a \\
B4490 & 1a<o2L 1b*<o1 / 2a<1a1 2b<1a1,1b1 & IU/EE & +{-}/{-}{-} & E & W:1a \\
B4491 & 1a<o2L 1b*<o1 / 2a<1a1 2b<1a1,1b2 & IU/EE & +{-}/{-}{-} & E & W:1a \\
B4492 & 1a<o2L 1b*<o1 / 2a<1a1 2b<1a1L,1b1 & IU/LI & +{-}/{-}+ & B & CP:2b W:1a,2b \\
B4493 & 1a<o2L 1b*<o1 / 2a<1a1 2b<1a1L,1b2 & IU/LI & +{-}/{-}+ & B & CP:2b W:1a,2b \\
B4494 & 1a<o2L 1b*<o1 / 2a<1a1 2b<1a2,1b1 & IU/EE & +{-}/{-}{-} & E & C:2b W:1a \\
B4495 & 1a<o2L 1b*<o1 / 2a<1a1 2b<1a2,1b2 & IU/EE & +{-}/{-}{-} & E & C:2b W:1a \\
B4496 & 1a<o2L 1b*<o1 / 2a<1a1 2b<1a2L,1b1 & IU/LI & +{-}/{-}+ & B & P:2b W:1a,2b \\
B4497 & 1a<o2L 1b*<o1 / 2a<1a1 2b<1a2L,1b2 & IU/LI & +{-}/{-}+ & B & P:2b W:1a,2b \\
B4498 & 1a<o2L 1b*<o1 / 2a<1a1,1b1 2b<1a1L & IU/LI & +{-}/{-}+ & B & CP:2b W:1a,2b \\
B4499 & 1a<o2L 1b*<o1 / 2a<1a1,1b1 2b<1a2 & IU/EE & +{-}/{-}{-} & E & C:2b W:1a \\
B4500 & 1a<o2L 1b*<o1 / 2a<1a1,1b1 2b<1a2L & IU/LI & +{-}/{-}+ & B & P:2b W:1a,2b \\
B4501 & 1a<o2L 1b*<o1 / 2a<1a1,1b1 2b<1b1 & IU/IU & +{-}/{-}{-} & N & W:1a,2a \\
B4502 & 1a<o2L 1b*<o1 / 2a<1a1,1b1 2b<1b2 & IU/IU & +{-}/{-}{-} & N & W:1a,2a \\
B4503 & 1a<o2L 1b*<o1 / 2a<1a1,1b2 2b<1a1L & IU/LI & +{-}/{-}+ & B & CP:2b W:1a,2b \\
B4504 & 1a<o2L 1b*<o1 / 2a<1a1,1b2 2b<1a2 & IU/EE & +{-}/{-}{-} & E & C:2b W:1a \\
B4505 & 1a<o2L 1b*<o1 / 2a<1a1,1b2 2b<1a2L & IU/LI & +{-}/{-}+ & B & P:2b W:1a,2b \\
B4506 & 1a<o2L 1b*<o1 / 2a<1a1,1b2 2b<1b1 & IU/IU & +{-}/{-}{-} & N & W:1a,2a \\
B4507 & 1a<o2L 1b*<o1 / 2a<1a1,1b2 2b<1b2 & IU/IU & +{-}/{-}{-} & N & W:1a,2a \\
B4508 & 1a<o2L 1b*<o1 / 2a<1a1L 2b<1a2,1b1 & IU/IL & +{-}/+{-} & B & C:2a,2b P:2a W:1a,2a \\
B4509 & 1a<o2L 1b*<o1 / 2a<1a1L 2b<1a2,1b2 & IU/IL & +{-}/+{-} & B & C:2a,2b P:2a W:1a,2a \\
B4510 & 1a<o2L 1b*<o1 / 2a<1a1L,1b1 2b<1a2 & IU/IL & +{-}/+{-} & B & C:2a,2b P:2a W:1a,2a \\
B4511 & 1a<o2L 1b*<o1 / 2a<1a1L,1b1 2b<1b1 & IU/IU & +{-}/+{-} & N & W:1a,2a \\
B4512 & 1a<o2L 1b*<o1 / 2a<1a1L,1b1 2b<1b2 & IU/IU & +{-}/+{-} & N & W:1a,2a \\
B4513 & 1a<o2L 1b*<o1 / 2a<1a1L,1b2 2b<1a2 & IU/IL & +{-}/+{-} & B & C:2a,2b P:2a W:1a,2a \\
B4514 & 1a<o2L 1b*<o1 / 2a<1a1L,1b2 2b<1b1 & IU/IU & +{-}/+{-} & N & W:1a,2a \\
B4515 & 1a<o2L 1b*<o1 / 2a<1a1L,1b2 2b<1b2 & IU/IU & +{-}/+{-} & N & W:1a,2a \\
B4516 & 1a<o2L 1b*<o1 / 2a<1a2 2b<1a2,1b1 & IU/EE & +{-}/{-}{-} & E & W:1a \\
B4517 & 1a<o2L 1b*<o1 / 2a<1a2 2b<1a2,1b2 & IU/EE & +{-}/{-}{-} & E & W:1a \\
B4518 & 1a<o2L 1b*<o1 / 2a<1a2 2b<1a2L,1b1 & IU/LI & +{-}/{-}+ & B & CP:2b W:1a,2b \\
B4519 & 1a<o2L 1b*<o1 / 2a<1a2 2b<1a2L,1b2 & IU/LI & +{-}/{-}+ & B & CP:2b W:1a,2b \\
B4520 & 1a<o2L 1b*<o1 / 2a<1a2,1b1 2b<1a2L & IU/LI & +{-}/{-}+ & B & CP:2b W:1a,2b \\
B4521 & 1a<o2L 1b*<o1 / 2a<1a2,1b1 2b<1b1 & IU/IU & +{-}/{-}{-} & N & W:1a,2a \\
B4522 & 1a<o2L 1b*<o1 / 2a<1a2,1b1 2b<1b2 & IU/IU & +{-}/{-}{-} & N & W:1a,2a \\
B4523 & 1a<o2L 1b*<o1 / 2a<1a2,1b2 2b<1a2L & IU/LI & +{-}/{-}+ & B & CP:2b W:1a,2b \\
B4524 & 1a<o2L 1b*<o1 / 2a<1a2,1b2 2b<1b1 & IU/IU & +{-}/{-}{-} & N & W:1a,2a \\
B4525 & 1a<o2L 1b*<o1 / 2a<1a2,1b2 2b<1b2 & IU/IU & +{-}/{-}{-} & N & W:1a,2a \\
B4526 & 1a<o2L 1b*<o1 / 2a<1a2L,1b1 2b<1b1 & IU/IU & +{-}/+{-} & N & W:1a,2a \\
B4527 & 1a<o2L 1b*<o1 / 2a<1a2L,1b1 2b<1b2 & IU/IU & +{-}/+{-} & N & W:1a,2a \\
B4528 & 1a<o2L 1b*<o1 / 2a<1a2L,1b2 2b<1b1 & IU/IU & +{-}/+{-} & N & W:1a,2a \\
B4529 & 1a<o2L 1b*<o1 / 2a<1a2L,1b2 2b<1b2 & IU/IU & +{-}/+{-} & N & W:1a,2a \\
B4530 & 1a<o2L 1b*<o2 / 2a<1a0L 2b<1a1,1b1 & IU/NI & +{-}/+{-} & N & W:1a,2b \\
B4531 & 1a<o2L 1b*<o2 / 2a<1a0L 2b<1a1,1b2 & IU/NI & +{-}/+{-} & N & W:1a,2b \\
B4532 & 1a<o2L 1b*<o2 / 2a<1a0L 2b<1a2,1b1 & IU/NI & +{-}/+{-} & N & W:1a,2b \\
B4533 & 1a<o2L 1b*<o2 / 2a<1a0L 2b<1a2,1b2 & IU/NI & +{-}/+{-} & N & W:1a,2b \\
B4534 & 1a<o2L 1b*<o2 / 2a<1a0L,1b1 2b<1a1 & IU/UI & +{-}/+{-} & N & W:1a,2b \\
B4535 & 1a<o2L 1b*<o2 / 2a<1a0L,1b1 2b<1a2 & IU/UI & +{-}/+{-} & N & W:1a,2b \\
B4536 & 1a<o2L 1b*<o2 / 2a<1a0L,1b1 2b<1b1 & IU/UU & +{-}/+{-} & N & W:1a \\
B4537 & 1a<o2L 1b*<o2 / 2a<1a0L,1b1 2b<1b2 & IU/UU & +{-}/+{-} & N & W:1a \\
B4538 & 1a<o2L 1b*<o2 / 2a<1a0L,1b2 2b<1a1 & IU/UI & +{-}/+{-} & N & W:1a,2b \\
B4539 & 1a<o2L 1b*<o2 / 2a<1a0L,1b2 2b<1a2 & IU/UI & +{-}/+{-} & N & W:1a,2b \\
B4540 & 1a<o2L 1b*<o2 / 2a<1a0L,1b2 2b<1b1 & IU/UU & +{-}/+{-} & N & W:1a \\
B4541 & 1a<o2L 1b*<o2 / 2a<1a0L,1b2 2b<1b2 & IU/UU & +{-}/+{-} & N & W:1a \\
B4542 & 1a<o2L 1b*<o2 / 2a<1a1 2b<1a1,1b1 & IU/EE & +{-}/{-}{-} & E & W:1a \\
B4543 & 1a<o2L 1b*<o2 / 2a<1a1 2b<1a1,1b2 & IU/EE & +{-}/{-}{-} & E & W:1a \\
B4544 & 1a<o2L 1b*<o2 / 2a<1a1 2b<1a1L,1b1 & IU/LI & +{-}/{-}+ & B & CP:2b W:1a,2b \\
B4545 & 1a<o2L 1b*<o2 / 2a<1a1 2b<1a1L,1b2 & IU/LI & +{-}/{-}+ & B & CP:2b W:1a,2b \\
B4546 & 1a<o2L 1b*<o2 / 2a<1a1 2b<1a2,1b1 & IU/EE & +{-}/{-}{-} & E & C:2b W:1a \\
B4547 & 1a<o2L 1b*<o2 / 2a<1a1 2b<1a2,1b2 & IU/EE & +{-}/{-}{-} & E & C:2b W:1a \\
B4548 & 1a<o2L 1b*<o2 / 2a<1a1 2b<1a2L,1b1 & IU/LI & +{-}/{-}+ & B & P:2b W:1a,2b \\
B4549 & 1a<o2L 1b*<o2 / 2a<1a1 2b<1a2L,1b2 & IU/LI & +{-}/{-}+ & B & P:2b W:1a,2b \\
B4550 & 1a<o2L 1b*<o2 / 2a<1a1,1b1 2b<1a1L & IU/LI & +{-}/{-}+ & B & CP:2b W:1a,2b \\
B4551 & 1a<o2L 1b*<o2 / 2a<1a1,1b1 2b<1a2 & IU/EE & +{-}/{-}{-} & E & C:2b W:1a \\
B4552 & 1a<o2L 1b*<o2 / 2a<1a1,1b1 2b<1a2L & IU/LI & +{-}/{-}+ & B & P:2b W:1a,2b \\
B4553 & 1a<o2L 1b*<o2 / 2a<1a1,1b1 2b<1b1 & IU/IU & +{-}/{-}{-} & N & W:1a,2a \\
B4554 & 1a<o2L 1b*<o2 / 2a<1a1,1b1 2b<1b2 & IU/IU & +{-}/{-}{-} & N & W:1a,2a \\
B4555 & 1a<o2L 1b*<o2 / 2a<1a1,1b2 2b<1a1L & IU/LI & +{-}/{-}+ & B & CP:2b W:1a,2b \\
B4556 & 1a<o2L 1b*<o2 / 2a<1a1,1b2 2b<1a2 & IU/EE & +{-}/{-}{-} & E & C:2b W:1a \\
B4557 & 1a<o2L 1b*<o2 / 2a<1a1,1b2 2b<1a2L & IU/LI & +{-}/{-}+ & B & P:2b W:1a,2b \\
B4558 & 1a<o2L 1b*<o2 / 2a<1a1,1b2 2b<1b1 & IU/IU & +{-}/{-}{-} & N & W:1a,2a \\
B4559 & 1a<o2L 1b*<o2 / 2a<1a1,1b2 2b<1b2 & IU/IU & +{-}/{-}{-} & N & W:1a,2a \\
B4560 & 1a<o2L 1b*<o2 / 2a<1a1L 2b<1a2,1b1 & IU/IL & +{-}/+{-} & B & C:2a,2b P:2a W:1a,2a \\
B4561 & 1a<o2L 1b*<o2 / 2a<1a1L 2b<1a2,1b2 & IU/IL & +{-}/+{-} & B & C:2a,2b P:2a W:1a,2a \\
B4562 & 1a<o2L 1b*<o2 / 2a<1a1L,1b1 2b<1a2 & IU/IL & +{-}/+{-} & B & C:2a,2b P:2a W:1a,2a \\
B4563 & 1a<o2L 1b*<o2 / 2a<1a1L,1b1 2b<1b1 & IU/IU & +{-}/+{-} & N & W:1a,2a \\
B4564 & 1a<o2L 1b*<o2 / 2a<1a1L,1b1 2b<1b2 & IU/IU & +{-}/+{-} & N & W:1a,2a \\
B4565 & 1a<o2L 1b*<o2 / 2a<1a1L,1b2 2b<1a2 & IU/IL & +{-}/+{-} & B & C:2a,2b P:2a W:1a,2a \\
B4566 & 1a<o2L 1b*<o2 / 2a<1a1L,1b2 2b<1b1 & IU/IU & +{-}/+{-} & N & W:1a,2a \\
B4567 & 1a<o2L 1b*<o2 / 2a<1a1L,1b2 2b<1b2 & IU/IU & +{-}/+{-} & N & W:1a,2a \\
B4568 & 1a<o2L 1b*<o2 / 2a<1a2 2b<1a2,1b1 & IU/EE & +{-}/{-}{-} & E & W:1a \\
B4569 & 1a<o2L 1b*<o2 / 2a<1a2 2b<1a2,1b2 & IU/EE & +{-}/{-}{-} & E & W:1a \\
B4570 & 1a<o2L 1b*<o2 / 2a<1a2 2b<1a2L,1b1 & IU/LI & +{-}/{-}+ & B & CP:2b W:1a,2b \\
B4571 & 1a<o2L 1b*<o2 / 2a<1a2 2b<1a2L,1b2 & IU/LI & +{-}/{-}+ & B & CP:2b W:1a,2b \\
B4572 & 1a<o2L 1b*<o2 / 2a<1a2,1b1 2b<1a2L & IU/LI & +{-}/{-}+ & B & CP:2b W:1a,2b \\
B4573 & 1a<o2L 1b*<o2 / 2a<1a2,1b1 2b<1b1 & IU/IU & +{-}/{-}{-} & N & W:1a,2a \\
B4574 & 1a<o2L 1b*<o2 / 2a<1a2,1b1 2b<1b2 & IU/IU & +{-}/{-}{-} & N & W:1a,2a \\
B4575 & 1a<o2L 1b*<o2 / 2a<1a2,1b2 2b<1a2L & IU/LI & +{-}/{-}+ & B & CP:2b W:1a,2b \\
B4576 & 1a<o2L 1b*<o2 / 2a<1a2,1b2 2b<1b1 & IU/IU & +{-}/{-}{-} & N & W:1a,2a \\
B4577 & 1a<o2L 1b*<o2 / 2a<1a2,1b2 2b<1b2 & IU/IU & +{-}/{-}{-} & N & W:1a,2a \\
B4578 & 1a<o2L 1b*<o2 / 2a<1a2L,1b1 2b<1b1 & IU/IU & +{-}/+{-} & N & W:1a,2a \\
B4579 & 1a<o2L 1b*<o2 / 2a<1a2L,1b1 2b<1b2 & IU/IU & +{-}/+{-} & N & W:1a,2a \\
B4580 & 1a<o2L 1b*<o2 / 2a<1a2L,1b2 2b<1b1 & IU/IU & +{-}/+{-} & N & W:1a,2a \\
B4581 & 1a<o2L 1b*<o2 / 2a<1a2L,1b2 2b<1b2 & IU/IU & +{-}/+{-} & N & W:1a,2a \\
B4582 & 1a<o0L 1b<o1 / 2a<1a0L,1b1 2b*<1a1,1b1 & NI/IU & +{-}/+{-} & N & W:1b,2a \\
B4583 & 1a<o0L 1b<o1 / 2a<1a0L,1b1 2b*<1a1,1b2 & NI/IU & +{-}/+{-} & N & W:1b,2a \\
B4584 & 1a<o0L 1b<o1 / 2a<1a0L,1b1 2b*<1a2,1b1 & NI/IU & +{-}/+{-} & N & W:1b,2a \\
B4585 & 1a<o0L 1b<o1 / 2a<1a0L,1b1 2b*<1a2,1b2 & NI/IU & +{-}/+{-} & N & W:1b,2a \\
B4586 & 1a<o0L 1b<o1 / 2a<1a0L,1b2 2b*<1a1,1b1 & NI/IU & +{-}/+{-} & N & W:1b,2a \\
B4587 & 1a<o0L 1b<o1 / 2a<1a0L,1b2 2b*<1a1,1b2 & NI/IU & +{-}/+{-} & N & W:1b,2a \\
B4588 & 1a<o0L 1b<o1 / 2a<1a0L,1b2 2b*<1a2,1b1 & NI/IU & +{-}/+{-} & N & W:1b,2a \\
B4589 & 1a<o0L 1b<o1 / 2a<1a0L,1b2 2b*<1a2,1b2 & NI/IU & +{-}/+{-} & N & W:1b,2a \\
B4590 & 1a<o0L 1b<o1 / 2a<1a1,1b0L 2b*<1a1,1b1 & NI/NU & +{-}/{-}{-} & N & W:2b \\
B4591 & 1a<o0L 1b<o1 / 2a<1a1,1b0L 2b*<1a1,1b2 & NI/NU & +{-}/{-}{-} & N & W:2b \\
B4592 & 1a<o0L 1b<o1 / 2a<1a1,1b0L 2b*<1a2,1b1 & NI/NU & +{-}/{-}{-} & N & W:2b \\
B4593 & 1a<o0L 1b<o1 / 2a<1a1,1b0L 2b*<1a2,1b2 & NI/NU & +{-}/{-}{-} & N & W:2b \\
B4594 & 1a<o0L 1b<o1 / 2a<1a1,1b1 2b*<1a1,1b1 & NI/EU & +{-}/{-}{-} & F & W:1b \\
B4595 & 1a<o0L 1b<o1 / 2a<1a1,1b1 2b*<1a1,1b2 & NI/EU & +{-}/{-}{-} & F & W:1b \\
B4596 & 1a<o0L 1b<o1 / 2a<1a1,1b1 2b*<1a2,1b1 & NI/EU & +{-}/{-}{-} & F & W:1b \\
B4597 & 1a<o0L 1b<o1 / 2a<1a1,1b1 2b*<1a2,1b2 & NI/EU & +{-}/{-}{-} & F & W:1b \\
B4598 & 1a<o0L 1b<o1 / 2a<1a1,1b1L 2b*<1a1,1b1 & NI/IU & +{-}/{-}{-} & F & CP:2a W:1b,2a \\
B4599 & 1a<o0L 1b<o1 / 2a<1a1,1b1L 2b*<1a1,1b2 & NI/IU & +{-}/{-}{-} & F & CP:2a W:1b,2a \\
B4600 & 1a<o0L 1b<o1 / 2a<1a1,1b1L 2b*<1a2,1b1 & NI/IU & +{-}/{-}{-} & F & CP:2a W:1b,2a \\
B4601 & 1a<o0L 1b<o1 / 2a<1a1,1b1L 2b*<1a2,1b2 & NI/IU & +{-}/{-}{-} & F & CP:2a W:1b,2a \\
B4602 & 1a<o0L 1b<o1 / 2a<1a1,1b2 2b*<1a1,1b1 & NI/EU & +{-}/{-}{-} & F & C:2a W:1b \\
B4603 & 1a<o0L 1b<o1 / 2a<1a1,1b2 2b*<1a1,1b2 & NI/EU & +{-}/{-}{-} & F & C:2a W:1b \\
B4604 & 1a<o0L 1b<o1 / 2a<1a1,1b2 2b*<1a2,1b1 & NI/EU & +{-}/{-}{-} & F & C:2a W:1b \\
B4605 & 1a<o0L 1b<o1 / 2a<1a1,1b2 2b*<1a2,1b2 & NI/EU & +{-}/{-}{-} & F & C:2a W:1b \\
B4606 & 1a<o0L 1b<o1 / 2a<1a1,1b2L 2b*<1a1,1b1 & NI/IU & +{-}/{-}{-} & F & P:2a W:1b,2a \\
B4607 & 1a<o0L 1b<o1 / 2a<1a1,1b2L 2b*<1a1,1b2 & NI/IU & +{-}/{-}{-} & F & P:2a W:1b,2a \\
B4608 & 1a<o0L 1b<o1 / 2a<1a1,1b2L 2b*<1a2,1b1 & NI/IU & +{-}/{-}{-} & F & P:2a W:1b,2a \\
B4609 & 1a<o0L 1b<o1 / 2a<1a1,1b2L 2b*<1a2,1b2 & NI/IU & +{-}/{-}{-} & F & P:2a W:1b,2a \\
B4610 & 1a<o0L 1b<o1 / 2a<1a1L,1b1 2b*<1a1,1b1 & NI/EU & +{-}/+{-} & F & W:1b \\
B4611 & 1a<o0L 1b<o1 / 2a<1a1L,1b1 2b*<1a1,1b2 & NI/EU & +{-}/+{-} & F & W:1b \\
B4612 & 1a<o0L 1b<o1 / 2a<1a1L,1b1 2b*<1a2,1b1 & NI/EU & +{-}/+{-} & F & W:1b \\
B4613 & 1a<o0L 1b<o1 / 2a<1a1L,1b1 2b*<1a2,1b2 & NI/EU & +{-}/+{-} & F & W:1b \\
B4614 & 1a<o0L 1b<o1 / 2a<1a1L,1b2 2b*<1a1,1b1 & NI/EU & +{-}/+{-} & F & C:2a W:1b \\
B4615 & 1a<o0L 1b<o1 / 2a<1a1L,1b2 2b*<1a1,1b2 & NI/EU & +{-}/+{-} & F & C:2a W:1b \\
B4616 & 1a<o0L 1b<o1 / 2a<1a1L,1b2 2b*<1a2,1b1 & NI/EU & +{-}/+{-} & F & C:2a W:1b \\
B4617 & 1a<o0L 1b<o1 / 2a<1a1L,1b2 2b*<1a2,1b2 & NI/EU & +{-}/+{-} & F & C:2a W:1b \\
B4618 & 1a<o0L 1b<o1 / 2a<1a2,1b0L 2b*<1a1,1b1 & NI/NU & +{-}/{-}{-} & N & W:2b \\
B4619 & 1a<o0L 1b<o1 / 2a<1a2,1b0L 2b*<1a1,1b2 & NI/NU & +{-}/{-}{-} & N & W:2b \\
B4620 & 1a<o0L 1b<o1 / 2a<1a2,1b0L 2b*<1a2,1b1 & NI/NU & +{-}/{-}{-} & N & W:2b \\
B4621 & 1a<o0L 1b<o1 / 2a<1a2,1b0L 2b*<1a2,1b2 & NI/NU & +{-}/{-}{-} & N & W:2b \\
B4622 & 1a<o0L 1b<o1 / 2a<1a2,1b1 2b*<1a1,1b1 & NI/EU & +{-}/{-}{-} & F & W:1b \\
B4623 & 1a<o0L 1b<o1 / 2a<1a2,1b1 2b*<1a1,1b2 & NI/EU & +{-}/{-}{-} & F & W:1b \\
B4624 & 1a<o0L 1b<o1 / 2a<1a2,1b1 2b*<1a2,1b1 & NI/EU & +{-}/{-}{-} & F & W:1b \\
B4625 & 1a<o0L 1b<o1 / 2a<1a2,1b1 2b*<1a2,1b2 & NI/EU & +{-}/{-}{-} & F & W:1b \\
B4626 & 1a<o0L 1b<o1 / 2a<1a2,1b1L 2b*<1a1,1b1 & NI/IU & +{-}/{-}{-} & F & CP:2a W:1b,2a \\
B4627 & 1a<o0L 1b<o1 / 2a<1a2,1b1L 2b*<1a1,1b2 & NI/IU & +{-}/{-}{-} & F & CP:2a W:1b,2a \\
B4628 & 1a<o0L 1b<o1 / 2a<1a2,1b1L 2b*<1a2,1b1 & NI/IU & +{-}/{-}{-} & F & CP:2a W:1b,2a \\
B4629 & 1a<o0L 1b<o1 / 2a<1a2,1b1L 2b*<1a2,1b2 & NI/IU & +{-}/{-}{-} & F & CP:2a W:1b,2a \\
B4630 & 1a<o0L 1b<o1 / 2a<1a2,1b2 2b*<1a1,1b1 & NI/EU & +{-}/{-}{-} & F & W:1b \\
B4631 & 1a<o0L 1b<o1 / 2a<1a2,1b2 2b*<1a1,1b2 & NI/EU & +{-}/{-}{-} & F & W:1b \\
B4632 & 1a<o0L 1b<o1 / 2a<1a2,1b2 2b*<1a2,1b1 & NI/EU & +{-}/{-}{-} & F & W:1b \\
B4633 & 1a<o0L 1b<o1 / 2a<1a2,1b2 2b*<1a2,1b2 & NI/EU & +{-}/{-}{-} & F & W:1b \\
B4634 & 1a<o0L 1b<o1 / 2a<1a2,1b2L 2b*<1a1,1b1 & NI/IU & +{-}/{-}{-} & F & CP:2a W:1b,2a \\
B4635 & 1a<o0L 1b<o1 / 2a<1a2,1b2L 2b*<1a1,1b2 & NI/IU & +{-}/{-}{-} & F & CP:2a W:1b,2a \\
B4636 & 1a<o0L 1b<o1 / 2a<1a2,1b2L 2b*<1a2,1b1 & NI/IU & +{-}/{-}{-} & F & CP:2a W:1b,2a \\
B4637 & 1a<o0L 1b<o1 / 2a<1a2,1b2L 2b*<1a2,1b2 & NI/IU & +{-}/{-}{-} & F & CP:2a W:1b,2a \\
B4638 & 1a<o0L 1b<o1 / 2a<1a2L,1b1 2b*<1a1,1b1 & NI/EU & +{-}/+{-} & F & W:1b \\
B4639 & 1a<o0L 1b<o1 / 2a<1a2L,1b1 2b*<1a1,1b2 & NI/EU & +{-}/+{-} & F & W:1b \\
B4640 & 1a<o0L 1b<o1 / 2a<1a2L,1b1 2b*<1a2,1b1 & NI/EU & +{-}/+{-} & F & W:1b \\
B4641 & 1a<o0L 1b<o1 / 2a<1a2L,1b1 2b*<1a2,1b2 & NI/EU & +{-}/+{-} & F & W:1b \\
B4642 & 1a<o0L 1b<o1 / 2a<1a2L,1b2 2b*<1a1,1b1 & NI/EU & +{-}/+{-} & F & W:1b \\
B4643 & 1a<o0L 1b<o1 / 2a<1a2L,1b2 2b*<1a1,1b2 & NI/EU & +{-}/+{-} & F & W:1b \\
B4644 & 1a<o0L 1b<o1 / 2a<1a2L,1b2 2b*<1a2,1b1 & NI/EU & +{-}/+{-} & F & W:1b \\
B4645 & 1a<o0L 1b<o1 / 2a<1a2L,1b2 2b*<1a2,1b2 & NI/EU & +{-}/+{-} & F & W:1b \\
B4646 & 1a<o0L 1b<o2 / 2a<1a0L,1b1 2b*<1a1,1b1 & NI/IU & +{-}/+{-} & N & W:1b,2a \\
B4647 & 1a<o0L 1b<o2 / 2a<1a0L,1b1 2b*<1a1,1b2 & NI/IU & +{-}/+{-} & N & W:1b,2a \\
B4648 & 1a<o0L 1b<o2 / 2a<1a0L,1b1 2b*<1a2,1b1 & NI/IU & +{-}/+{-} & N & W:1b,2a \\
B4649 & 1a<o0L 1b<o2 / 2a<1a0L,1b1 2b*<1a2,1b2 & NI/IU & +{-}/+{-} & N & W:1b,2a \\
B4650 & 1a<o0L 1b<o2 / 2a<1a0L,1b2 2b*<1a1,1b1 & NI/IU & +{-}/+{-} & N & W:1b,2a \\
B4651 & 1a<o0L 1b<o2 / 2a<1a0L,1b2 2b*<1a1,1b2 & NI/IU & +{-}/+{-} & N & W:1b,2a \\
B4652 & 1a<o0L 1b<o2 / 2a<1a0L,1b2 2b*<1a2,1b1 & NI/IU & +{-}/+{-} & N & W:1b,2a \\
B4653 & 1a<o0L 1b<o2 / 2a<1a0L,1b2 2b*<1a2,1b2 & NI/IU & +{-}/+{-} & N & W:1b,2a \\
B4654 & 1a<o0L 1b<o2 / 2a<1a1,1b0L 2b*<1a1,1b1 & NI/NU & +{-}/{-}{-} & N & W:2b \\
B4655 & 1a<o0L 1b<o2 / 2a<1a1,1b0L 2b*<1a1,1b2 & NI/NU & +{-}/{-}{-} & N & W:2b \\
B4656 & 1a<o0L 1b<o2 / 2a<1a1,1b0L 2b*<1a2,1b1 & NI/NU & +{-}/{-}{-} & N & W:2b \\
B4657 & 1a<o0L 1b<o2 / 2a<1a1,1b0L 2b*<1a2,1b2 & NI/NU & +{-}/{-}{-} & N & W:2b \\
B4658 & 1a<o0L 1b<o2 / 2a<1a1,1b1 2b*<1a1,1b1 & NI/EU & +{-}/{-}{-} & F & W:1b \\
B4659 & 1a<o0L 1b<o2 / 2a<1a1,1b1 2b*<1a1,1b2 & NI/EU & +{-}/{-}{-} & F & W:1b \\
B4660 & 1a<o0L 1b<o2 / 2a<1a1,1b1 2b*<1a2,1b1 & NI/EU & +{-}/{-}{-} & F & W:1b \\
B4661 & 1a<o0L 1b<o2 / 2a<1a1,1b1 2b*<1a2,1b2 & NI/EU & +{-}/{-}{-} & F & W:1b \\
B4662 & 1a<o0L 1b<o2 / 2a<1a1,1b1L 2b*<1a1,1b1 & NI/IU & +{-}/{-}{-} & F & CP:2a W:1b,2a \\
B4663 & 1a<o0L 1b<o2 / 2a<1a1,1b1L 2b*<1a1,1b2 & NI/IU & +{-}/{-}{-} & F & CP:2a W:1b,2a \\
B4664 & 1a<o0L 1b<o2 / 2a<1a1,1b1L 2b*<1a2,1b1 & NI/IU & +{-}/{-}{-} & F & CP:2a W:1b,2a \\
B4665 & 1a<o0L 1b<o2 / 2a<1a1,1b1L 2b*<1a2,1b2 & NI/IU & +{-}/{-}{-} & F & CP:2a W:1b,2a \\
B4666 & 1a<o0L 1b<o2 / 2a<1a1,1b2 2b*<1a1,1b1 & NI/EU & +{-}/{-}{-} & F & C:2a W:1b \\
B4667 & 1a<o0L 1b<o2 / 2a<1a1,1b2 2b*<1a1,1b2 & NI/EU & +{-}/{-}{-} & F & C:2a W:1b \\
B4668 & 1a<o0L 1b<o2 / 2a<1a1,1b2 2b*<1a2,1b1 & NI/EU & +{-}/{-}{-} & F & C:2a W:1b \\
B4669 & 1a<o0L 1b<o2 / 2a<1a1,1b2 2b*<1a2,1b2 & NI/EU & +{-}/{-}{-} & F & C:2a W:1b \\
B4670 & 1a<o0L 1b<o2 / 2a<1a1,1b2L 2b*<1a1,1b1 & NI/IU & +{-}/{-}{-} & F & P:2a W:1b,2a \\
B4671 & 1a<o0L 1b<o2 / 2a<1a1,1b2L 2b*<1a1,1b2 & NI/IU & +{-}/{-}{-} & F & P:2a W:1b,2a \\
B4672 & 1a<o0L 1b<o2 / 2a<1a1,1b2L 2b*<1a2,1b1 & NI/IU & +{-}/{-}{-} & F & P:2a W:1b,2a \\
B4673 & 1a<o0L 1b<o2 / 2a<1a1,1b2L 2b*<1a2,1b2 & NI/IU & +{-}/{-}{-} & F & P:2a W:1b,2a \\
B4674 & 1a<o0L 1b<o2 / 2a<1a1L,1b1 2b*<1a1,1b1 & NI/EU & +{-}/+{-} & F & W:1b \\
B4675 & 1a<o0L 1b<o2 / 2a<1a1L,1b1 2b*<1a1,1b2 & NI/EU & +{-}/+{-} & F & W:1b \\
B4676 & 1a<o0L 1b<o2 / 2a<1a1L,1b1 2b*<1a2,1b1 & NI/EU & +{-}/+{-} & F & W:1b \\
B4677 & 1a<o0L 1b<o2 / 2a<1a1L,1b1 2b*<1a2,1b2 & NI/EU & +{-}/+{-} & F & W:1b \\
B4678 & 1a<o0L 1b<o2 / 2a<1a1L,1b2 2b*<1a1,1b1 & NI/EU & +{-}/+{-} & F & C:2a W:1b \\
B4679 & 1a<o0L 1b<o2 / 2a<1a1L,1b2 2b*<1a1,1b2 & NI/EU & +{-}/+{-} & F & C:2a W:1b \\
B4680 & 1a<o0L 1b<o2 / 2a<1a1L,1b2 2b*<1a2,1b1 & NI/EU & +{-}/+{-} & F & C:2a W:1b \\
B4681 & 1a<o0L 1b<o2 / 2a<1a1L,1b2 2b*<1a2,1b2 & NI/EU & +{-}/+{-} & F & C:2a W:1b \\
B4682 & 1a<o0L 1b<o2 / 2a<1a2,1b0L 2b*<1a1,1b1 & NI/NU & +{-}/{-}{-} & N & W:2b \\
B4683 & 1a<o0L 1b<o2 / 2a<1a2,1b0L 2b*<1a1,1b2 & NI/NU & +{-}/{-}{-} & N & W:2b \\
B4684 & 1a<o0L 1b<o2 / 2a<1a2,1b0L 2b*<1a2,1b1 & NI/NU & +{-}/{-}{-} & N & W:2b \\
B4685 & 1a<o0L 1b<o2 / 2a<1a2,1b0L 2b*<1a2,1b2 & NI/NU & +{-}/{-}{-} & N & W:2b \\
B4686 & 1a<o0L 1b<o2 / 2a<1a2,1b1 2b*<1a1,1b1 & NI/EU & +{-}/{-}{-} & F & W:1b \\
B4687 & 1a<o0L 1b<o2 / 2a<1a2,1b1 2b*<1a1,1b2 & NI/EU & +{-}/{-}{-} & F & W:1b \\
B4688 & 1a<o0L 1b<o2 / 2a<1a2,1b1 2b*<1a2,1b1 & NI/EU & +{-}/{-}{-} & F & W:1b \\
B4689 & 1a<o0L 1b<o2 / 2a<1a2,1b1 2b*<1a2,1b2 & NI/EU & +{-}/{-}{-} & F & W:1b \\
B4690 & 1a<o0L 1b<o2 / 2a<1a2,1b1L 2b*<1a1,1b1 & NI/IU & +{-}/{-}{-} & F & CP:2a W:1b,2a \\
B4691 & 1a<o0L 1b<o2 / 2a<1a2,1b1L 2b*<1a1,1b2 & NI/IU & +{-}/{-}{-} & F & CP:2a W:1b,2a \\
B4692 & 1a<o0L 1b<o2 / 2a<1a2,1b1L 2b*<1a2,1b1 & NI/IU & +{-}/{-}{-} & F & CP:2a W:1b,2a \\
B4693 & 1a<o0L 1b<o2 / 2a<1a2,1b1L 2b*<1a2,1b2 & NI/IU & +{-}/{-}{-} & F & CP:2a W:1b,2a \\
B4694 & 1a<o0L 1b<o2 / 2a<1a2,1b2 2b*<1a1,1b1 & NI/EU & +{-}/{-}{-} & F & W:1b \\
B4695 & 1a<o0L 1b<o2 / 2a<1a2,1b2 2b*<1a1,1b2 & NI/EU & +{-}/{-}{-} & F & W:1b \\
B4696 & 1a<o0L 1b<o2 / 2a<1a2,1b2 2b*<1a2,1b1 & NI/EU & +{-}/{-}{-} & F & W:1b \\
B4697 & 1a<o0L 1b<o2 / 2a<1a2,1b2 2b*<1a2,1b2 & NI/EU & +{-}/{-}{-} & F & W:1b \\
B4698 & 1a<o0L 1b<o2 / 2a<1a2,1b2L 2b*<1a1,1b1 & NI/IU & +{-}/{-}{-} & F & CP:2a W:1b,2a \\
B4699 & 1a<o0L 1b<o2 / 2a<1a2,1b2L 2b*<1a1,1b2 & NI/IU & +{-}/{-}{-} & F & CP:2a W:1b,2a \\
B4700 & 1a<o0L 1b<o2 / 2a<1a2,1b2L 2b*<1a2,1b1 & NI/IU & +{-}/{-}{-} & F & CP:2a W:1b,2a \\
B4701 & 1a<o0L 1b<o2 / 2a<1a2,1b2L 2b*<1a2,1b2 & NI/IU & +{-}/{-}{-} & F & CP:2a W:1b,2a \\
B4702 & 1a<o0L 1b<o2 / 2a<1a2L,1b1 2b*<1a1,1b1 & NI/EU & +{-}/+{-} & F & W:1b \\
B4703 & 1a<o0L 1b<o2 / 2a<1a2L,1b1 2b*<1a1,1b2 & NI/EU & +{-}/+{-} & F & W:1b \\
B4704 & 1a<o0L 1b<o2 / 2a<1a2L,1b1 2b*<1a2,1b1 & NI/EU & +{-}/+{-} & F & W:1b \\
B4705 & 1a<o0L 1b<o2 / 2a<1a2L,1b1 2b*<1a2,1b2 & NI/EU & +{-}/+{-} & F & W:1b \\
B4706 & 1a<o0L 1b<o2 / 2a<1a2L,1b2 2b*<1a1,1b1 & NI/EU & +{-}/+{-} & F & W:1b \\
B4707 & 1a<o0L 1b<o2 / 2a<1a2L,1b2 2b*<1a1,1b2 & NI/EU & +{-}/+{-} & F & W:1b \\
B4708 & 1a<o0L 1b<o2 / 2a<1a2L,1b2 2b*<1a2,1b1 & NI/EU & +{-}/+{-} & F & W:1b \\
B4709 & 1a<o0L 1b<o2 / 2a<1a2L,1b2 2b*<1a2,1b2 & NI/EU & +{-}/+{-} & F & W:1b \\
B4710 & 1a<o0L 1b*<o1 / 2a<1a0L,1b1 2b<1a1,1b1 & NU/UU & +{-}/+{-} & N &  \\
B4711 & 1a<o0L 1b*<o1 / 2a<1a0L,1b1 2b<1a1,1b2 & NU/UU & +{-}/+{-} & N &  \\
B4712 & 1a<o0L 1b*<o1 / 2a<1a0L,1b1 2b<1a2,1b1 & NU/UU & +{-}/+{-} & N &  \\
B4713 & 1a<o0L 1b*<o1 / 2a<1a0L,1b1 2b<1a2,1b2 & NU/UU & +{-}/+{-} & N &  \\
B4714 & 1a<o0L 1b*<o1 / 2a<1a0L,1b2 2b<1a1,1b1 & NU/UU & +{-}/+{-} & N &  \\
B4715 & 1a<o0L 1b*<o1 / 2a<1a0L,1b2 2b<1a1,1b2 & NU/UU & +{-}/+{-} & N &  \\
B4716 & 1a<o0L 1b*<o1 / 2a<1a0L,1b2 2b<1a2,1b1 & NU/UU & +{-}/+{-} & N &  \\
B4717 & 1a<o0L 1b*<o1 / 2a<1a0L,1b2 2b<1a2,1b2 & NU/UU & +{-}/+{-} & N &  \\
B4718 & 1a<o0L 1b*<o1 / 2a<1a1,1b1 2b<1a1,1b1 & NU/UU & +{-}/{-}{-} & N &  \\
B4719 & 1a<o0L 1b*<o1 / 2a<1a1,1b1 2b<1a1,1b2 & NU/UU & +{-}/{-}{-} & N &  \\
B4720 & 1a<o0L 1b*<o1 / 2a<1a1,1b1 2b<1a1L,1b1 & NU/UU & +{-}/{-}+ & N &  \\
B4721 & 1a<o0L 1b*<o1 / 2a<1a1,1b1 2b<1a1L,1b2 & NU/UU & +{-}/{-}+ & N &  \\
B4722 & 1a<o0L 1b*<o1 / 2a<1a1,1b1 2b<1a2,1b1 & NU/UU & +{-}/{-}{-} & N &  \\
B4723 & 1a<o0L 1b*<o1 / 2a<1a1,1b1 2b<1a2,1b2 & NU/UU & +{-}/{-}{-} & N &  \\
B4724 & 1a<o0L 1b*<o1 / 2a<1a1,1b1 2b<1a2L,1b1 & NU/UU & +{-}/{-}+ & N &  \\
B4725 & 1a<o0L 1b*<o1 / 2a<1a1,1b1 2b<1a2L,1b2 & NU/UU & +{-}/{-}+ & N &  \\
B4726 & 1a<o0L 1b*<o1 / 2a<1a1,1b2 2b<1a1,1b2 & NU/UU & +{-}/{-}{-} & N &  \\
B4727 & 1a<o0L 1b*<o1 / 2a<1a1,1b2 2b<1a1L,1b1 & NU/UU & +{-}/{-}+ & N &  \\
B4728 & 1a<o0L 1b*<o1 / 2a<1a1,1b2 2b<1a1L,1b2 & NU/UU & +{-}/{-}+ & N &  \\
B4729 & 1a<o0L 1b*<o1 / 2a<1a1,1b2 2b<1a2,1b1 & NU/UU & +{-}/{-}{-} & N &  \\
B4730 & 1a<o0L 1b*<o1 / 2a<1a1,1b2 2b<1a2,1b2 & NU/UU & +{-}/{-}{-} & N &  \\
B4731 & 1a<o0L 1b*<o1 / 2a<1a1,1b2 2b<1a2L,1b1 & NU/UU & +{-}/{-}+ & N &  \\
B4732 & 1a<o0L 1b*<o1 / 2a<1a1,1b2 2b<1a2L,1b2 & NU/UU & +{-}/{-}+ & N &  \\
B4733 & 1a<o0L 1b*<o1 / 2a<1a1L,1b1 2b<1a2,1b1 & NU/UU & +{-}/+{-} & N &  \\
B4734 & 1a<o0L 1b*<o1 / 2a<1a1L,1b1 2b<1a2,1b2 & NU/UU & +{-}/+{-} & N &  \\
B4735 & 1a<o0L 1b*<o1 / 2a<1a1L,1b2 2b<1a2,1b1 & NU/UU & +{-}/+{-} & N &  \\
B4736 & 1a<o0L 1b*<o1 / 2a<1a1L,1b2 2b<1a2,1b2 & NU/UU & +{-}/+{-} & N &  \\
B4737 & 1a<o0L 1b*<o1 / 2a<1a2,1b1 2b<1a2,1b1 & NU/UU & +{-}/{-}{-} & N &  \\
B4738 & 1a<o0L 1b*<o1 / 2a<1a2,1b1 2b<1a2,1b2 & NU/UU & +{-}/{-}{-} & N &  \\
B4739 & 1a<o0L 1b*<o1 / 2a<1a2,1b1 2b<1a2L,1b1 & NU/UU & +{-}/{-}+ & N &  \\
B4740 & 1a<o0L 1b*<o1 / 2a<1a2,1b1 2b<1a2L,1b2 & NU/UU & +{-}/{-}+ & N &  \\
B4741 & 1a<o0L 1b*<o1 / 2a<1a2,1b2 2b<1a2,1b2 & NU/UU & +{-}/{-}{-} & N &  \\
B4742 & 1a<o0L 1b*<o1 / 2a<1a2,1b2 2b<1a2L,1b1 & NU/UU & +{-}/{-}+ & N &  \\
B4743 & 1a<o0L 1b*<o1 / 2a<1a2,1b2 2b<1a2L,1b2 & NU/UU & +{-}/{-}+ & N &  \\
B4744 & 1a<o0L 1b*<o2 / 2a<1a0L,1b1 2b<1a1,1b1 & NU/UU & +{-}/+{-} & N &  \\
B4745 & 1a<o0L 1b*<o2 / 2a<1a0L,1b1 2b<1a1,1b2 & NU/UU & +{-}/+{-} & N &  \\
B4746 & 1a<o0L 1b*<o2 / 2a<1a0L,1b1 2b<1a2,1b1 & NU/UU & +{-}/+{-} & N &  \\
B4747 & 1a<o0L 1b*<o2 / 2a<1a0L,1b1 2b<1a2,1b2 & NU/UU & +{-}/+{-} & N &  \\
B4748 & 1a<o0L 1b*<o2 / 2a<1a0L,1b2 2b<1a1,1b1 & NU/UU & +{-}/+{-} & N &  \\
B4749 & 1a<o0L 1b*<o2 / 2a<1a0L,1b2 2b<1a1,1b2 & NU/UU & +{-}/+{-} & N &  \\
B4750 & 1a<o0L 1b*<o2 / 2a<1a0L,1b2 2b<1a2,1b1 & NU/UU & +{-}/+{-} & N &  \\
B4751 & 1a<o0L 1b*<o2 / 2a<1a0L,1b2 2b<1a2,1b2 & NU/UU & +{-}/+{-} & N &  \\
B4752 & 1a<o0L 1b*<o2 / 2a<1a1,1b1 2b<1a1,1b1 & NU/UU & +{-}/{-}{-} & N &  \\
B4753 & 1a<o0L 1b*<o2 / 2a<1a1,1b1 2b<1a1,1b2 & NU/UU & +{-}/{-}{-} & N &  \\
B4754 & 1a<o0L 1b*<o2 / 2a<1a1,1b1 2b<1a1L,1b1 & NU/UU & +{-}/{-}+ & N &  \\
B4755 & 1a<o0L 1b*<o2 / 2a<1a1,1b1 2b<1a1L,1b2 & NU/UU & +{-}/{-}+ & N &  \\
B4756 & 1a<o0L 1b*<o2 / 2a<1a1,1b1 2b<1a2,1b1 & NU/UU & +{-}/{-}{-} & N &  \\
B4757 & 1a<o0L 1b*<o2 / 2a<1a1,1b1 2b<1a2,1b2 & NU/UU & +{-}/{-}{-} & N &  \\
B4758 & 1a<o0L 1b*<o2 / 2a<1a1,1b1 2b<1a2L,1b1 & NU/UU & +{-}/{-}+ & N &  \\
B4759 & 1a<o0L 1b*<o2 / 2a<1a1,1b1 2b<1a2L,1b2 & NU/UU & +{-}/{-}+ & N &  \\
B4760 & 1a<o0L 1b*<o2 / 2a<1a1,1b2 2b<1a1,1b2 & NU/UU & +{-}/{-}{-} & N &  \\
B4761 & 1a<o0L 1b*<o2 / 2a<1a1,1b2 2b<1a1L,1b1 & NU/UU & +{-}/{-}+ & N &  \\
B4762 & 1a<o0L 1b*<o2 / 2a<1a1,1b2 2b<1a1L,1b2 & NU/UU & +{-}/{-}+ & N &  \\
B4763 & 1a<o0L 1b*<o2 / 2a<1a1,1b2 2b<1a2,1b1 & NU/UU & +{-}/{-}{-} & N &  \\
B4764 & 1a<o0L 1b*<o2 / 2a<1a1,1b2 2b<1a2,1b2 & NU/UU & +{-}/{-}{-} & N &  \\
B4765 & 1a<o0L 1b*<o2 / 2a<1a1,1b2 2b<1a2L,1b1 & NU/UU & +{-}/{-}+ & N &  \\
B4766 & 1a<o0L 1b*<o2 / 2a<1a1,1b2 2b<1a2L,1b2 & NU/UU & +{-}/{-}+ & N &  \\
B4767 & 1a<o0L 1b*<o2 / 2a<1a1L,1b1 2b<1a2,1b1 & NU/UU & +{-}/+{-} & N &  \\
B4768 & 1a<o0L 1b*<o2 / 2a<1a1L,1b1 2b<1a2,1b2 & NU/UU & +{-}/+{-} & N &  \\
B4769 & 1a<o0L 1b*<o2 / 2a<1a1L,1b2 2b<1a2,1b1 & NU/UU & +{-}/+{-} & N &  \\
B4770 & 1a<o0L 1b*<o2 / 2a<1a1L,1b2 2b<1a2,1b2 & NU/UU & +{-}/+{-} & N &  \\
B4771 & 1a<o0L 1b*<o2 / 2a<1a2,1b1 2b<1a2,1b1 & NU/UU & +{-}/{-}{-} & N &  \\
B4772 & 1a<o0L 1b*<o2 / 2a<1a2,1b1 2b<1a2,1b2 & NU/UU & +{-}/{-}{-} & N &  \\
B4773 & 1a<o0L 1b*<o2 / 2a<1a2,1b1 2b<1a2L,1b1 & NU/UU & +{-}/{-}+ & N &  \\
B4774 & 1a<o0L 1b*<o2 / 2a<1a2,1b1 2b<1a2L,1b2 & NU/UU & +{-}/{-}+ & N &  \\
B4775 & 1a<o0L 1b*<o2 / 2a<1a2,1b2 2b<1a2,1b2 & NU/UU & +{-}/{-}{-} & N &  \\
B4776 & 1a<o0L 1b*<o2 / 2a<1a2,1b2 2b<1a2L,1b1 & NU/UU & +{-}/{-}+ & N &  \\
B4777 & 1a<o0L 1b*<o2 / 2a<1a2,1b2 2b<1a2L,1b2 & NU/UU & +{-}/{-}+ & N &  \\
B4778 & 1a<o1 1b<o1 / 2a<1a0L,1b1 2b*<1a1,1b1 & EE/DU & {-}{-}/{-}{-} & E &  \\
B4779 & 1a<o1 1b<o1 / 2a<1a0L,1b1 2b*<1a1,1b2 & EE/DU & {-}{-}/{-}{-} & E &  \\
B4780 & 1a<o1 1b<o1 / 2a<1a0L,1b1 2b*<1a2,1b1 & EE/DU & {-}{-}/{-}{-} & E &  \\
B4781 & 1a<o1 1b<o1 / 2a<1a0L,1b1 2b*<1a2,1b2 & EE/DU & {-}{-}/{-}{-} & E &  \\
B4782 & 1a<o1 1b<o1 / 2a<1a0L,1b2 2b*<1a1,1b1 & EE/DU & {-}{-}/{-}{-} & E &  \\
B4783 & 1a<o1 1b<o1 / 2a<1a0L,1b2 2b*<1a1,1b2 & EE/DU & {-}{-}/{-}{-} & E &  \\
B4784 & 1a<o1 1b<o1 / 2a<1a0L,1b2 2b*<1a2,1b1 & EE/DU & {-}{-}/{-}{-} & E &  \\
B4785 & 1a<o1 1b<o1 / 2a<1a0L,1b2 2b*<1a2,1b2 & EE/DU & {-}{-}/{-}{-} & E &  \\
B4786 & 1a<o1 1b<o1 / 2a<1a1,1b1 2b*<1a1,1b1 & EE/DU & {-}{-}/{-}{-} & E &  \\
B4787 & 1a<o1 1b<o1 / 2a<1a1,1b1 2b*<1a1,1b2 & EE/DU & {-}{-}/{-}{-} & E &  \\
B4788 & 1a<o1 1b<o1 / 2a<1a1,1b1 2b*<1a2,1b2 & EE/DU & {-}{-}/{-}{-} & E &  \\
B4789 & 1a<o1 1b<o1 / 2a<1a1,1b1L 2b*<1a1,1b1 & EE/DU & {-}{-}/{-}{-} & E &  \\
B4790 & 1a<o1 1b<o1 / 2a<1a1,1b1L 2b*<1a1,1b2 & EE/DU & {-}{-}/{-}{-} & E &  \\
B4791 & 1a<o1 1b<o1 / 2a<1a1,1b1L 2b*<1a2,1b1 & EE/DU & {-}{-}/{-}{-} & E &  \\
B4792 & 1a<o1 1b<o1 / 2a<1a1,1b1L 2b*<1a2,1b2 & EE/DU & {-}{-}/{-}{-} & E &  \\
B4793 & 1a<o1 1b<o1 / 2a<1a1,1b2 2b*<1a1,1b1 & EE/DU & {-}{-}/{-}{-} & E &  \\
B4794 & 1a<o1 1b<o1 / 2a<1a1,1b2 2b*<1a1,1b2 & EE/DU & {-}{-}/{-}{-} & E &  \\
B4795 & 1a<o1 1b<o1 / 2a<1a1,1b2 2b*<1a2,1b1 & EE/DU & {-}{-}/{-}{-} & E &  \\
B4796 & 1a<o1 1b<o1 / 2a<1a1,1b2 2b*<1a2,1b2 & EE/DU & {-}{-}/{-}{-} & E &  \\
B4797 & 1a<o1 1b<o1 / 2a<1a1,1b2L 2b*<1a1,1b1 & EE/DU & {-}{-}/{-}{-} & E &  \\
B4798 & 1a<o1 1b<o1 / 2a<1a1,1b2L 2b*<1a1,1b2 & EE/DU & {-}{-}/{-}{-} & E &  \\
B4799 & 1a<o1 1b<o1 / 2a<1a1,1b2L 2b*<1a2,1b1 & EE/DU & {-}{-}/{-}{-} & E &  \\
B4800 & 1a<o1 1b<o1 / 2a<1a1,1b2L 2b*<1a2,1b2 & EE/DU & {-}{-}/{-}{-} & E &  \\
B4801 & 1a<o1 1b<o1 / 2a<1a1L,1b2 2b*<1a1,1b1 & EE/DU & {-}{-}/{-}{-} & E &  \\
B4802 & 1a<o1 1b<o1 / 2a<1a1L,1b2 2b*<1a1,1b2 & EE/DU & {-}{-}/{-}{-} & E &  \\
B4803 & 1a<o1 1b<o1 / 2a<1a1L,1b2 2b*<1a2,1b1 & EE/DU & {-}{-}/{-}{-} & E &  \\
B4804 & 1a<o1 1b<o1 / 2a<1a1L,1b2 2b*<1a2,1b2 & EE/DU & {-}{-}/{-}{-} & E &  \\
B4805 & 1a<o1 1b<o1 / 2a<1a2,1b2 2b*<1a1,1b1 & EE/DU & {-}{-}/{-}{-} & E &  \\
B4806 & 1a<o1 1b<o1 / 2a<1a2,1b2 2b*<1a1,1b2 & EE/DU & {-}{-}/{-}{-} & E &  \\
B4807 & 1a<o1 1b<o1 / 2a<1a2,1b2 2b*<1a2,1b2 & EE/DU & {-}{-}/{-}{-} & E &  \\
B4808 & 1a<o1 1b<o1 / 2a<1a2,1b2L 2b*<1a1,1b1 & EE/DU & {-}{-}/{-}{-} & E &  \\
B4809 & 1a<o1 1b<o1 / 2a<1a2,1b2L 2b*<1a1,1b2 & EE/DU & {-}{-}/{-}{-} & E &  \\
B4810 & 1a<o1 1b<o1 / 2a<1a2,1b2L 2b*<1a2,1b1 & EE/DU & {-}{-}/{-}{-} & E &  \\
B4811 & 1a<o1 1b<o1 / 2a<1a2,1b2L 2b*<1a2,1b2 & EE/DU & {-}{-}/{-}{-} & E &  \\
B4812 & 1a<o1 1b<o1L / 2a<1a0L,1b1 2b*<1a1,1b1 & LI/IU & {-}+/{-}{-} & B & CPW:1b,2a \\
B4813 & 1a<o1 1b<o1L / 2a<1a0L,1b1 2b*<1a1,1b2 & LI/IU & {-}+/{-}{-} & B & CPW:1b,2a \\
B4814 & 1a<o1 1b<o1L / 2a<1a0L,1b1 2b*<1a2,1b1 & LI/IU & {-}+/{-}{-} & B & CPW:1b,2a \\
B4815 & 1a<o1 1b<o1L / 2a<1a0L,1b1 2b*<1a2,1b2 & LI/IU & {-}+/{-}{-} & B & CPW:1b,2a \\
B4816 & 1a<o1 1b<o1L / 2a<1a0L,1b2 2b*<1a1,1b1 & LI/IU & {-}+/{-}{-} & B & CPW:1b,2a \\
B4817 & 1a<o1 1b<o1L / 2a<1a0L,1b2 2b*<1a1,1b2 & LI/IU & {-}+/{-}{-} & B & CPW:1b,2a \\
B4818 & 1a<o1 1b<o1L / 2a<1a0L,1b2 2b*<1a2,1b1 & LI/IU & {-}+/{-}{-} & B & CPW:1b,2a \\
B4819 & 1a<o1 1b<o1L / 2a<1a0L,1b2 2b*<1a2,1b2 & LI/IU & {-}+/{-}{-} & B & CPW:1b,2a \\
B4820 & 1a<o1 1b<o1L / 2a<1a1,1b0L 2b*<1a1,1b1 & LI/DU & {-}+/+{-} & B & CPW:1b \\
B4821 & 1a<o1 1b<o1L / 2a<1a1,1b0L 2b*<1a1,1b2 & LI/DU & {-}+/+{-} & B & CPW:1b \\
B4822 & 1a<o1 1b<o1L / 2a<1a1,1b0L 2b*<1a2,1b1 & LI/DU & {-}+/+{-} & B & CPW:1b \\
B4823 & 1a<o1 1b<o1L / 2a<1a1,1b0L 2b*<1a2,1b2 & LI/DU & {-}+/+{-} & B & CPW:1b \\
B4824 & 1a<o1 1b<o1L / 2a<1a1,1b1 2b*<1a1,1b1 & LI/EU & {-}+/{-}{-} & BF & CPW:1b \\
B4825 & 1a<o1 1b<o1L / 2a<1a1,1b1 2b*<1a1,1b2 & LI/EU & {-}+/{-}{-} & BF & CPW:1b \\
B4826 & 1a<o1 1b<o1L / 2a<1a1,1b1 2b*<1a2,1b1 & LI/EU & {-}+/{-}{-} & BF & CPW:1b \\
B4827 & 1a<o1 1b<o1L / 2a<1a1,1b1 2b*<1a2,1b2 & LI/EU & {-}+/{-}{-} & BF & CPW:1b \\
B4828 & 1a<o1 1b<o1L / 2a<1a1,1b1L 2b*<1a1,1b1 & LI/IU & {-}+/+{-} & BF & CPW:1b,2a \\
B4829 & 1a<o1 1b<o1L / 2a<1a1,1b1L 2b*<1a1,1b2 & LI/IU & {-}+/+{-} & BF & CPW:1b,2a \\
B4830 & 1a<o1 1b<o1L / 2a<1a1,1b1L 2b*<1a2,1b1 & LI/IU & {-}+/+{-} & BF & CPW:1b,2a \\
B4831 & 1a<o1 1b<o1L / 2a<1a1,1b1L 2b*<1a2,1b2 & LI/IU & {-}+/+{-} & BF & CPW:1b,2a \\
B4832 & 1a<o1 1b<o1L / 2a<1a1,1b2 2b*<1a1,1b1 & LI/EU & {-}+/{-}{-} & BF & CPW:1b \\
B4833 & 1a<o1 1b<o1L / 2a<1a1,1b2 2b*<1a1,1b2 & LI/EU & {-}+/{-}{-} & BF & CPW:1b \\
B4834 & 1a<o1 1b<o1L / 2a<1a1,1b2 2b*<1a2,1b1 & LI/EU & {-}+/{-}{-} & BF & CPW:1b \\
B4835 & 1a<o1 1b<o1L / 2a<1a1,1b2 2b*<1a2,1b2 & LI/EU & {-}+/{-}{-} & BF & CPW:1b \\
B4836 & 1a<o1 1b<o1L / 2a<1a1,1b2L 2b*<1a1,1b1 & LI/IU & {-}+/+{-} & BF & CPW:1b,2a \\
B4837 & 1a<o1 1b<o1L / 2a<1a1,1b2L 2b*<1a1,1b2 & LI/IU & {-}+/+{-} & BF & CPW:1b,2a \\
B4838 & 1a<o1 1b<o1L / 2a<1a1,1b2L 2b*<1a2,1b1 & LI/IU & {-}+/+{-} & BF & CPW:1b,2a \\
B4839 & 1a<o1 1b<o1L / 2a<1a1,1b2L 2b*<1a2,1b2 & LI/IU & {-}+/+{-} & BF & CPW:1b,2a \\
B4840 & 1a<o1 1b<o1L / 2a<1a1L,1b1 2b*<1a1,1b1 & LI/EU & {-}+/{-}{-} & BF & CPW:1b \\
B4841 & 1a<o1 1b<o1L / 2a<1a1L,1b1 2b*<1a1,1b2 & LI/EU & {-}+/{-}{-} & BF & CPW:1b \\
B4842 & 1a<o1 1b<o1L / 2a<1a1L,1b1 2b*<1a2,1b1 & LI/EU & {-}+/{-}{-} & BF & CPW:1b \\
B4843 & 1a<o1 1b<o1L / 2a<1a1L,1b1 2b*<1a2,1b2 & LI/EU & {-}+/{-}{-} & BF & CPW:1b \\
B4844 & 1a<o1 1b<o1L / 2a<1a1L,1b2 2b*<1a1,1b1 & LI/EU & {-}+/{-}{-} & BF & CPW:1b \\
B4845 & 1a<o1 1b<o1L / 2a<1a1L,1b2 2b*<1a1,1b2 & LI/EU & {-}+/{-}{-} & BF & CPW:1b \\
B4846 & 1a<o1 1b<o1L / 2a<1a1L,1b2 2b*<1a2,1b1 & LI/EU & {-}+/{-}{-} & BF & CPW:1b \\
B4847 & 1a<o1 1b<o1L / 2a<1a1L,1b2 2b*<1a2,1b2 & LI/EU & {-}+/{-}{-} & BF & CPW:1b \\
B4848 & 1a<o1 1b<o1L / 2a<1a2,1b0L 2b*<1a1,1b1 & LI/DU & {-}+/+{-} & B & CPW:1b \\
B4849 & 1a<o1 1b<o1L / 2a<1a2,1b0L 2b*<1a1,1b2 & LI/DU & {-}+/+{-} & B & CPW:1b \\
B4850 & 1a<o1 1b<o1L / 2a<1a2,1b0L 2b*<1a2,1b1 & LI/DU & {-}+/+{-} & B & CPW:1b \\
B4851 & 1a<o1 1b<o1L / 2a<1a2,1b0L 2b*<1a2,1b2 & LI/DU & {-}+/+{-} & B & CPW:1b \\
B4852 & 1a<o1 1b<o1L / 2a<1a2,1b1 2b*<1a1,1b1 & LI/EU & {-}+/{-}{-} & BF & CPW:1b \\
B4853 & 1a<o1 1b<o1L / 2a<1a2,1b1 2b*<1a1,1b2 & LI/EU & {-}+/{-}{-} & BF & CPW:1b \\
B4854 & 1a<o1 1b<o1L / 2a<1a2,1b1 2b*<1a2,1b1 & LI/EU & {-}+/{-}{-} & BF & CPW:1b \\
B4855 & 1a<o1 1b<o1L / 2a<1a2,1b1 2b*<1a2,1b2 & LI/EU & {-}+/{-}{-} & BF & CPW:1b \\
B4856 & 1a<o1 1b<o1L / 2a<1a2,1b1L 2b*<1a1,1b1 & LI/IU & {-}+/+{-} & BF & CPW:1b,2a \\
B4857 & 1a<o1 1b<o1L / 2a<1a2,1b1L 2b*<1a1,1b2 & LI/IU & {-}+/+{-} & BF & CPW:1b,2a \\
B4858 & 1a<o1 1b<o1L / 2a<1a2,1b1L 2b*<1a2,1b1 & LI/IU & {-}+/+{-} & BF & CPW:1b,2a \\
B4859 & 1a<o1 1b<o1L / 2a<1a2,1b1L 2b*<1a2,1b2 & LI/IU & {-}+/+{-} & BF & CPW:1b,2a \\
B4860 & 1a<o1 1b<o1L / 2a<1a2,1b2 2b*<1a1,1b1 & LI/EU & {-}+/{-}{-} & BF & CPW:1b \\
B4861 & 1a<o1 1b<o1L / 2a<1a2,1b2 2b*<1a1,1b2 & LI/EU & {-}+/{-}{-} & BF & CPW:1b \\
B4862 & 1a<o1 1b<o1L / 2a<1a2,1b2 2b*<1a2,1b1 & LI/EU & {-}+/{-}{-} & BF & CPW:1b \\
B4863 & 1a<o1 1b<o1L / 2a<1a2,1b2 2b*<1a2,1b2 & LI/EU & {-}+/{-}{-} & BF & CPW:1b \\
B4864 & 1a<o1 1b<o1L / 2a<1a2,1b2L 2b*<1a1,1b1 & LI/IU & {-}+/+{-} & BF & CPW:1b,2a \\
B4865 & 1a<o1 1b<o1L / 2a<1a2,1b2L 2b*<1a1,1b2 & LI/IU & {-}+/+{-} & BF & CPW:1b,2a \\
B4866 & 1a<o1 1b<o1L / 2a<1a2,1b2L 2b*<1a2,1b1 & LI/IU & {-}+/+{-} & BF & CPW:1b,2a \\
B4867 & 1a<o1 1b<o1L / 2a<1a2,1b2L 2b*<1a2,1b2 & LI/IU & {-}+/+{-} & BF & CPW:1b,2a \\
B4868 & 1a<o1 1b<o1L / 2a<1a2L,1b1 2b*<1a1,1b1 & LI/EU & {-}+/{-}{-} & BF & CPW:1b \\
B4869 & 1a<o1 1b<o1L / 2a<1a2L,1b1 2b*<1a1,1b2 & LI/EU & {-}+/{-}{-} & BF & CPW:1b \\
B4870 & 1a<o1 1b<o1L / 2a<1a2L,1b1 2b*<1a2,1b1 & LI/EU & {-}+/{-}{-} & BF & CPW:1b \\
B4871 & 1a<o1 1b<o1L / 2a<1a2L,1b1 2b*<1a2,1b2 & LI/EU & {-}+/{-}{-} & BF & CPW:1b \\
B4872 & 1a<o1 1b<o1L / 2a<1a2L,1b2 2b*<1a1,1b1 & LI/EU & {-}+/{-}{-} & BF & CPW:1b \\
B4873 & 1a<o1 1b<o1L / 2a<1a2L,1b2 2b*<1a1,1b2 & LI/EU & {-}+/{-}{-} & BF & CPW:1b \\
B4874 & 1a<o1 1b<o1L / 2a<1a2L,1b2 2b*<1a2,1b1 & LI/EU & {-}+/{-}{-} & BF & CPW:1b \\
B4875 & 1a<o1 1b<o1L / 2a<1a2L,1b2 2b*<1a2,1b2 & LI/EU & {-}+/{-}{-} & BF & CPW:1b \\
B4876 & 1a<o1 1b<o2 / 2a<1a0L,1b1 2b*<1a1,1b1 & EE/DU & {-}{-}/{-}{-} & E & C:1b,2a \\
B4877 & 1a<o1 1b<o2 / 2a<1a0L,1b1 2b*<1a1,1b2 & EE/DU & {-}{-}/{-}{-} & E & C:1b,2a \\
B4878 & 1a<o1 1b<o2 / 2a<1a0L,1b1 2b*<1a2,1b1 & EE/DU & {-}{-}/{-}{-} & E & C:1b,2a \\
B4879 & 1a<o1 1b<o2 / 2a<1a0L,1b1 2b*<1a2,1b2 & EE/DU & {-}{-}/{-}{-} & E & C:1b,2a \\
B4880 & 1a<o1 1b<o2 / 2a<1a0L,1b2 2b*<1a1,1b1 & EE/DU & {-}{-}/{-}{-} & E & C:1b,2a \\
B4881 & 1a<o1 1b<o2 / 2a<1a0L,1b2 2b*<1a1,1b2 & EE/DU & {-}{-}/{-}{-} & E & C:1b,2a \\
B4882 & 1a<o1 1b<o2 / 2a<1a0L,1b2 2b*<1a2,1b1 & EE/DU & {-}{-}/{-}{-} & E & C:1b,2a \\
B4883 & 1a<o1 1b<o2 / 2a<1a0L,1b2 2b*<1a2,1b2 & EE/DU & {-}{-}/{-}{-} & E & C:1b,2a \\
B4884 & 1a<o1 1b<o2 / 2a<1a1,1b0L 2b*<1a1,1b1 & EE/DU & {-}{-}/{-}{-} & E & C:1b \\
B4885 & 1a<o1 1b<o2 / 2a<1a1,1b0L 2b*<1a1,1b2 & EE/DU & {-}{-}/{-}{-} & E & C:1b \\
B4886 & 1a<o1 1b<o2 / 2a<1a1,1b0L 2b*<1a2,1b1 & EE/DU & {-}{-}/{-}{-} & E & C:1b \\
B4887 & 1a<o1 1b<o2 / 2a<1a1,1b0L 2b*<1a2,1b2 & EE/DU & {-}{-}/{-}{-} & E & C:1b \\
B4888 & 1a<o1 1b<o2 / 2a<1a1,1b1 2b*<1a1,1b1 & EE/DU & {-}{-}/{-}{-} & E & C:1b \\
B4889 & 1a<o1 1b<o2 / 2a<1a1,1b1 2b*<1a1,1b2 & EE/DU & {-}{-}/{-}{-} & E & C:1b \\
B4890 & 1a<o1 1b<o2 / 2a<1a1,1b1 2b*<1a2,1b1 & EE/DU & {-}{-}/{-}{-} & E & C:1b \\
B4891 & 1a<o1 1b<o2 / 2a<1a1,1b1 2b*<1a2,1b2 & EE/DU & {-}{-}/{-}{-} & E & C:1b \\
B4892 & 1a<o1 1b<o2 / 2a<1a1,1b1L 2b*<1a1,1b1 & EE/DU & {-}{-}/{-}{-} & E & C:1b \\
B4893 & 1a<o1 1b<o2 / 2a<1a1,1b1L 2b*<1a1,1b2 & EE/DU & {-}{-}/{-}{-} & E & C:1b \\
B4894 & 1a<o1 1b<o2 / 2a<1a1,1b1L 2b*<1a2,1b1 & EE/DU & {-}{-}/{-}{-} & E & C:1b \\
B4895 & 1a<o1 1b<o2 / 2a<1a1,1b1L 2b*<1a2,1b2 & EE/DU & {-}{-}/{-}{-} & E & C:1b \\
B4896 & 1a<o1 1b<o2 / 2a<1a1,1b2 2b*<1a1,1b1 & EE/DU & {-}{-}/{-}{-} & E & C:1b,2a \\
B4897 & 1a<o1 1b<o2 / 2a<1a1,1b2 2b*<1a1,1b2 & EE/DU & {-}{-}/{-}{-} & E & C:1b,2a \\
B4898 & 1a<o1 1b<o2 / 2a<1a1,1b2 2b*<1a2,1b1 & EE/DU & {-}{-}/{-}{-} & E & C:1b,2a \\
B4899 & 1a<o1 1b<o2 / 2a<1a1,1b2 2b*<1a2,1b2 & EE/DU & {-}{-}/{-}{-} & E & C:1b,2a \\
B4900 & 1a<o1 1b<o2 / 2a<1a1,1b2L 2b*<1a1,1b1 & EE/DU & {-}{-}/{-}{-} & E & C:1b,2a \\
B4901 & 1a<o1 1b<o2 / 2a<1a1,1b2L 2b*<1a1,1b2 & EE/DU & {-}{-}/{-}{-} & E & C:1b,2a \\
B4902 & 1a<o1 1b<o2 / 2a<1a1,1b2L 2b*<1a2,1b1 & EE/DU & {-}{-}/{-}{-} & E & C:1b,2a \\
B4903 & 1a<o1 1b<o2 / 2a<1a1,1b2L 2b*<1a2,1b2 & EE/DU & {-}{-}/{-}{-} & E & C:1b,2a \\
B4904 & 1a<o1 1b<o2 / 2a<1a1L,1b1 2b*<1a1,1b1 & EE/DU & {-}{-}/{-}{-} & E & C:1b \\
B4905 & 1a<o1 1b<o2 / 2a<1a1L,1b1 2b*<1a1,1b2 & EE/DU & {-}{-}/{-}{-} & E & C:1b \\
B4906 & 1a<o1 1b<o2 / 2a<1a1L,1b1 2b*<1a2,1b1 & EE/DU & {-}{-}/{-}{-} & E & C:1b \\
B4907 & 1a<o1 1b<o2 / 2a<1a1L,1b1 2b*<1a2,1b2 & EE/DU & {-}{-}/{-}{-} & E & C:1b \\
B4908 & 1a<o1 1b<o2 / 2a<1a1L,1b2 2b*<1a1,1b1 & EE/DU & {-}{-}/{-}{-} & E & C:1b,2a \\
B4909 & 1a<o1 1b<o2 / 2a<1a1L,1b2 2b*<1a1,1b2 & EE/DU & {-}{-}/{-}{-} & E & C:1b,2a \\
B4910 & 1a<o1 1b<o2 / 2a<1a1L,1b2 2b*<1a2,1b1 & EE/DU & {-}{-}/{-}{-} & E & C:1b,2a \\
B4911 & 1a<o1 1b<o2 / 2a<1a1L,1b2 2b*<1a2,1b2 & EE/DU & {-}{-}/{-}{-} & E & C:1b,2a \\
B4912 & 1a<o1 1b<o2 / 2a<1a2,1b0L 2b*<1a1,1b1 & EE/DU & {-}{-}/{-}{-} & E & C:1b \\
B4913 & 1a<o1 1b<o2 / 2a<1a2,1b0L 2b*<1a1,1b2 & EE/DU & {-}{-}/{-}{-} & E & C:1b \\
B4914 & 1a<o1 1b<o2 / 2a<1a2,1b0L 2b*<1a2,1b1 & EE/DU & {-}{-}/{-}{-} & E & C:1b \\
B4915 & 1a<o1 1b<o2 / 2a<1a2,1b0L 2b*<1a2,1b2 & EE/DU & {-}{-}/{-}{-} & E & C:1b \\
B4916 & 1a<o1 1b<o2 / 2a<1a2,1b1 2b*<1a1,1b1 & EE/DU & {-}{-}/{-}{-} & E & C:1b \\
B4917 & 1a<o1 1b<o2 / 2a<1a2,1b1 2b*<1a1,1b2 & EE/DU & {-}{-}/{-}{-} & E & C:1b \\
B4918 & 1a<o1 1b<o2 / 2a<1a2,1b1 2b*<1a2,1b1 & EE/DU & {-}{-}/{-}{-} & E & C:1b \\
B4919 & 1a<o1 1b<o2 / 2a<1a2,1b1 2b*<1a2,1b2 & EE/DU & {-}{-}/{-}{-} & E & C:1b \\
B4920 & 1a<o1 1b<o2 / 2a<1a2,1b1L 2b*<1a1,1b1 & EE/DU & {-}{-}/{-}{-} & E & C:1b \\
B4921 & 1a<o1 1b<o2 / 2a<1a2,1b1L 2b*<1a1,1b2 & EE/DU & {-}{-}/{-}{-} & E & C:1b \\
B4922 & 1a<o1 1b<o2 / 2a<1a2,1b1L 2b*<1a2,1b1 & EE/DU & {-}{-}/{-}{-} & E & C:1b \\
B4923 & 1a<o1 1b<o2 / 2a<1a2,1b1L 2b*<1a2,1b2 & EE/DU & {-}{-}/{-}{-} & E & C:1b \\
B4924 & 1a<o1 1b<o2 / 2a<1a2,1b2 2b*<1a1,1b1 & EE/DU & {-}{-}/{-}{-} & E & C:1b \\
B4925 & 1a<o1 1b<o2 / 2a<1a2,1b2 2b*<1a1,1b2 & EE/DU & {-}{-}/{-}{-} & E & C:1b \\
B4926 & 1a<o1 1b<o2 / 2a<1a2,1b2 2b*<1a2,1b1 & EE/DU & {-}{-}/{-}{-} & E & C:1b \\
B4927 & 1a<o1 1b<o2 / 2a<1a2,1b2 2b*<1a2,1b2 & EE/DU & {-}{-}/{-}{-} & E & C:1b \\
B4928 & 1a<o1 1b<o2 / 2a<1a2,1b2L 2b*<1a1,1b1 & EE/DU & {-}{-}/{-}{-} & E & C:1b \\
B4929 & 1a<o1 1b<o2 / 2a<1a2,1b2L 2b*<1a1,1b2 & EE/DU & {-}{-}/{-}{-} & E & C:1b \\
B4930 & 1a<o1 1b<o2 / 2a<1a2,1b2L 2b*<1a2,1b1 & EE/DU & {-}{-}/{-}{-} & E & C:1b \\
B4931 & 1a<o1 1b<o2 / 2a<1a2,1b2L 2b*<1a2,1b2 & EE/DU & {-}{-}/{-}{-} & E & C:1b \\
B4932 & 1a<o1 1b<o2 / 2a<1a2L,1b1 2b*<1a1,1b1 & EE/DU & {-}{-}/{-}{-} & E & C:1b \\
B4933 & 1a<o1 1b<o2 / 2a<1a2L,1b1 2b*<1a1,1b2 & EE/DU & {-}{-}/{-}{-} & E & C:1b \\
B4934 & 1a<o1 1b<o2 / 2a<1a2L,1b1 2b*<1a2,1b1 & EE/DU & {-}{-}/{-}{-} & E & C:1b \\
B4935 & 1a<o1 1b<o2 / 2a<1a2L,1b1 2b*<1a2,1b2 & EE/DU & {-}{-}/{-}{-} & E & C:1b \\
B4936 & 1a<o1 1b<o2 / 2a<1a2L,1b2 2b*<1a1,1b1 & EE/DU & {-}{-}/{-}{-} & E & C:1b \\
B4937 & 1a<o1 1b<o2 / 2a<1a2L,1b2 2b*<1a1,1b2 & EE/DU & {-}{-}/{-}{-} & E & C:1b \\
B4938 & 1a<o1 1b<o2 / 2a<1a2L,1b2 2b*<1a2,1b1 & EE/DU & {-}{-}/{-}{-} & E & C:1b \\
B4939 & 1a<o1 1b<o2 / 2a<1a2L,1b2 2b*<1a2,1b2 & EE/DU & {-}{-}/{-}{-} & E & C:1b \\
B4940 & 1a<o1 1b<o2L / 2a<1a0L,1b1 2b*<1a1,1b1 & LI/IU & {-}+/{-}{-} & B & PW:1b,2a \\
B4941 & 1a<o1 1b<o2L / 2a<1a0L,1b1 2b*<1a1,1b2 & LI/IU & {-}+/{-}{-} & B & PW:1b,2a \\
B4942 & 1a<o1 1b<o2L / 2a<1a0L,1b1 2b*<1a2,1b1 & LI/IU & {-}+/{-}{-} & B & PW:1b,2a \\
B4943 & 1a<o1 1b<o2L / 2a<1a0L,1b1 2b*<1a2,1b2 & LI/IU & {-}+/{-}{-} & B & PW:1b,2a \\
B4944 & 1a<o1 1b<o2L / 2a<1a0L,1b2 2b*<1a1,1b1 & LI/IU & {-}+/{-}{-} & B & PW:1b,2a \\
B4945 & 1a<o1 1b<o2L / 2a<1a0L,1b2 2b*<1a1,1b2 & LI/IU & {-}+/{-}{-} & B & PW:1b,2a \\
B4946 & 1a<o1 1b<o2L / 2a<1a0L,1b2 2b*<1a2,1b1 & LI/IU & {-}+/{-}{-} & B & PW:1b,2a \\
B4947 & 1a<o1 1b<o2L / 2a<1a0L,1b2 2b*<1a2,1b2 & LI/IU & {-}+/{-}{-} & B & PW:1b,2a \\
B4948 & 1a<o1 1b<o2L / 2a<1a1,1b0L 2b*<1a1,1b1 & LI/DU & {-}+/+{-} & B & PW:1b \\
B4949 & 1a<o1 1b<o2L / 2a<1a1,1b0L 2b*<1a1,1b2 & LI/DU & {-}+/+{-} & B & PW:1b \\
B4950 & 1a<o1 1b<o2L / 2a<1a1,1b0L 2b*<1a2,1b1 & LI/DU & {-}+/+{-} & B & PW:1b \\
B4951 & 1a<o1 1b<o2L / 2a<1a1,1b0L 2b*<1a2,1b2 & LI/DU & {-}+/+{-} & B & PW:1b \\
B4952 & 1a<o1 1b<o2L / 2a<1a1,1b1 2b*<1a1,1b1 & LI/EU & {-}+/{-}{-} & BF & PW:1b \\
B4953 & 1a<o1 1b<o2L / 2a<1a1,1b1 2b*<1a1,1b2 & LI/EU & {-}+/{-}{-} & BF & PW:1b \\
B4954 & 1a<o1 1b<o2L / 2a<1a1,1b1 2b*<1a2,1b1 & LI/EU & {-}+/{-}{-} & BF & PW:1b \\
B4955 & 1a<o1 1b<o2L / 2a<1a1,1b1 2b*<1a2,1b2 & LI/EU & {-}+/{-}{-} & BF & PW:1b \\
B4956 & 1a<o1 1b<o2L / 2a<1a1,1b1L 2b*<1a1,1b1 & LI/IU & {-}+/+{-} & BF & C:2a PW:1b,2a \\
B4957 & 1a<o1 1b<o2L / 2a<1a1,1b1L 2b*<1a1,1b2 & LI/IU & {-}+/+{-} & BF & C:2a PW:1b,2a \\
B4958 & 1a<o1 1b<o2L / 2a<1a1,1b1L 2b*<1a2,1b1 & LI/IU & {-}+/+{-} & BF & C:2a PW:1b,2a \\
B4959 & 1a<o1 1b<o2L / 2a<1a1,1b1L 2b*<1a2,1b2 & LI/IU & {-}+/+{-} & BF & C:2a PW:1b,2a \\
B4960 & 1a<o1 1b<o2L / 2a<1a1,1b2 2b*<1a1,1b1 & LI/EU & {-}+/{-}{-} & BF & C:2a PW:1b \\
B4961 & 1a<o1 1b<o2L / 2a<1a1,1b2 2b*<1a1,1b2 & LI/EU & {-}+/{-}{-} & BF & C:2a PW:1b \\
B4962 & 1a<o1 1b<o2L / 2a<1a1,1b2 2b*<1a2,1b1 & LI/EU & {-}+/{-}{-} & BF & C:2a PW:1b \\
B4963 & 1a<o1 1b<o2L / 2a<1a1,1b2 2b*<1a2,1b2 & LI/EU & {-}+/{-}{-} & BF & C:2a PW:1b \\
B4964 & 1a<o1 1b<o2L / 2a<1a1,1b2L 2b*<1a1,1b1 & LI/IU & {-}+/+{-} & BF & PW:1b,2a \\
B4965 & 1a<o1 1b<o2L / 2a<1a1,1b2L 2b*<1a1,1b2 & LI/IU & {-}+/+{-} & BF & PW:1b,2a \\
B4966 & 1a<o1 1b<o2L / 2a<1a1,1b2L 2b*<1a2,1b1 & LI/IU & {-}+/+{-} & BF & PW:1b,2a \\
B4967 & 1a<o1 1b<o2L / 2a<1a1,1b2L 2b*<1a2,1b2 & LI/IU & {-}+/+{-} & BF & PW:1b,2a \\
B4968 & 1a<o1 1b<o2L / 2a<1a1L,1b1 2b*<1a1,1b1 & LI/EU & {-}+/{-}{-} & BF & PW:1b \\
B4969 & 1a<o1 1b<o2L / 2a<1a1L,1b1 2b*<1a1,1b2 & LI/EU & {-}+/{-}{-} & BF & PW:1b \\
B4970 & 1a<o1 1b<o2L / 2a<1a1L,1b1 2b*<1a2,1b1 & LI/EU & {-}+/{-}{-} & BF & PW:1b \\
B4971 & 1a<o1 1b<o2L / 2a<1a1L,1b1 2b*<1a2,1b2 & LI/EU & {-}+/{-}{-} & BF & PW:1b \\
B4972 & 1a<o1 1b<o2L / 2a<1a1L,1b2 2b*<1a1,1b1 & LI/EU & {-}+/{-}{-} & BF & C:2a PW:1b \\
B4973 & 1a<o1 1b<o2L / 2a<1a1L,1b2 2b*<1a1,1b2 & LI/EU & {-}+/{-}{-} & BF & C:2a PW:1b \\
B4974 & 1a<o1 1b<o2L / 2a<1a1L,1b2 2b*<1a2,1b1 & LI/EU & {-}+/{-}{-} & BF & C:2a PW:1b \\
B4975 & 1a<o1 1b<o2L / 2a<1a1L,1b2 2b*<1a2,1b2 & LI/EU & {-}+/{-}{-} & BF & C:2a PW:1b \\
B4976 & 1a<o1 1b<o2L / 2a<1a2,1b0L 2b*<1a1,1b1 & LI/DU & {-}+/+{-} & B & PW:1b \\
B4977 & 1a<o1 1b<o2L / 2a<1a2,1b0L 2b*<1a1,1b2 & LI/DU & {-}+/+{-} & B & PW:1b \\
B4978 & 1a<o1 1b<o2L / 2a<1a2,1b0L 2b*<1a2,1b1 & LI/DU & {-}+/+{-} & B & PW:1b \\
B4979 & 1a<o1 1b<o2L / 2a<1a2,1b0L 2b*<1a2,1b2 & LI/DU & {-}+/+{-} & B & PW:1b \\
B4980 & 1a<o1 1b<o2L / 2a<1a2,1b1 2b*<1a1,1b1 & LI/EU & {-}+/{-}{-} & BF & PW:1b \\
B4981 & 1a<o1 1b<o2L / 2a<1a2,1b1 2b*<1a1,1b2 & LI/EU & {-}+/{-}{-} & BF & PW:1b \\
B4982 & 1a<o1 1b<o2L / 2a<1a2,1b1 2b*<1a2,1b1 & LI/EU & {-}+/{-}{-} & BF & PW:1b \\
B4983 & 1a<o1 1b<o2L / 2a<1a2,1b1 2b*<1a2,1b2 & LI/EU & {-}+/{-}{-} & BF & PW:1b \\
B4984 & 1a<o1 1b<o2L / 2a<1a2,1b1L 2b*<1a1,1b1 & LI/IU & {-}+/+{-} & BF & C:2a PW:1b,2a \\
B4985 & 1a<o1 1b<o2L / 2a<1a2,1b1L 2b*<1a1,1b2 & LI/IU & {-}+/+{-} & BF & C:2a PW:1b,2a \\
B4986 & 1a<o1 1b<o2L / 2a<1a2,1b1L 2b*<1a2,1b1 & LI/IU & {-}+/+{-} & BF & C:2a PW:1b,2a \\
B4987 & 1a<o1 1b<o2L / 2a<1a2,1b1L 2b*<1a2,1b2 & LI/IU & {-}+/+{-} & BF & C:2a PW:1b,2a \\
B4988 & 1a<o1 1b<o2L / 2a<1a2,1b2 2b*<1a1,1b1 & LI/EU & {-}+/{-}{-} & BF & PW:1b \\
B4989 & 1a<o1 1b<o2L / 2a<1a2,1b2 2b*<1a1,1b2 & LI/EU & {-}+/{-}{-} & BF & PW:1b \\
B4990 & 1a<o1 1b<o2L / 2a<1a2,1b2 2b*<1a2,1b1 & LI/EU & {-}+/{-}{-} & BF & PW:1b \\
B4991 & 1a<o1 1b<o2L / 2a<1a2,1b2 2b*<1a2,1b2 & LI/EU & {-}+/{-}{-} & BF & PW:1b \\
B4992 & 1a<o1 1b<o2L / 2a<1a2,1b2L 2b*<1a1,1b1 & LI/IU & {-}+/+{-} & BF & C:2a PW:1b,2a \\
B4993 & 1a<o1 1b<o2L / 2a<1a2,1b2L 2b*<1a1,1b2 & LI/IU & {-}+/+{-} & BF & C:2a PW:1b,2a \\
B4994 & 1a<o1 1b<o2L / 2a<1a2,1b2L 2b*<1a2,1b1 & LI/IU & {-}+/+{-} & BF & C:2a PW:1b,2a \\
B4995 & 1a<o1 1b<o2L / 2a<1a2,1b2L 2b*<1a2,1b2 & LI/IU & {-}+/+{-} & BF & C:2a PW:1b,2a \\
B4996 & 1a<o1 1b<o2L / 2a<1a2L,1b1 2b*<1a1,1b1 & LI/EU & {-}+/{-}{-} & BF & PW:1b \\
B4997 & 1a<o1 1b<o2L / 2a<1a2L,1b1 2b*<1a1,1b2 & LI/EU & {-}+/{-}{-} & BF & PW:1b \\
B4998 & 1a<o1 1b<o2L / 2a<1a2L,1b1 2b*<1a2,1b1 & LI/EU & {-}+/{-}{-} & BF & PW:1b \\
B4999 & 1a<o1 1b<o2L / 2a<1a2L,1b1 2b*<1a2,1b2 & LI/EU & {-}+/{-}{-} & BF & PW:1b \\
B5000 & 1a<o1 1b<o2L / 2a<1a2L,1b2 2b*<1a1,1b1 & LI/EU & {-}+/{-}{-} & BF & PW:1b \\
B5001 & 1a<o1 1b<o2L / 2a<1a2L,1b2 2b*<1a1,1b2 & LI/EU & {-}+/{-}{-} & BF & PW:1b \\
B5002 & 1a<o1 1b<o2L / 2a<1a2L,1b2 2b*<1a2,1b1 & LI/EU & {-}+/{-}{-} & BF & PW:1b \\
B5003 & 1a<o1 1b<o2L / 2a<1a2L,1b2 2b*<1a2,1b2 & LI/EU & {-}+/{-}{-} & BF & PW:1b \\
B5004 & 1a<o1 1b*<o1 / 2a<1a0L,1b1 2b<1a1,1b1 & IU/UI & {-}{-}/{-}{-} & N & W:1a,2b \\
B5005 & 1a<o1 1b*<o1 / 2a<1a0L,1b1 2b<1a1,1b2 & IU/UI & {-}{-}/{-}{-} & N & W:1a,2b \\
B5006 & 1a<o1 1b*<o1 / 2a<1a0L,1b1 2b<1a2,1b1 & IU/UI & {-}{-}/{-}{-} & N & W:1a,2b \\
B5007 & 1a<o1 1b*<o1 / 2a<1a0L,1b1 2b<1a2,1b2 & IU/UI & {-}{-}/{-}{-} & N & W:1a,2b \\
B5008 & 1a<o1 1b*<o1 / 2a<1a0L,1b2 2b<1a1,1b1 & IU/UI & {-}{-}/{-}{-} & N & W:1a,2b \\
B5009 & 1a<o1 1b*<o1 / 2a<1a0L,1b2 2b<1a1,1b2 & IU/UI & {-}{-}/{-}{-} & N & W:1a,2b \\
B5010 & 1a<o1 1b*<o1 / 2a<1a0L,1b2 2b<1a2,1b1 & IU/UI & {-}{-}/{-}{-} & N & W:1a,2b \\
B5011 & 1a<o1 1b*<o1 / 2a<1a0L,1b2 2b<1a2,1b2 & IU/UI & {-}{-}/{-}{-} & N & W:1a,2b \\
B5012 & 1a<o1 1b*<o1 / 2a<1a1,1b1 2b<1a1,1b1 & IU/EE & {-}{-}/{-}{-} & E & W:1a \\
B5013 & 1a<o1 1b*<o1 / 2a<1a1,1b1 2b<1a1,1b2 & IU/EE & {-}{-}/{-}{-} & E & W:1a \\
B5014 & 1a<o1 1b*<o1 / 2a<1a1,1b1 2b<1a1L,1b1 & IU/LI & {-}{-}/{-}{-} & B & CP:2b W:1a,2b \\
B5015 & 1a<o1 1b*<o1 / 2a<1a1,1b1 2b<1a1L,1b2 & IU/LI & {-}{-}/{-}{-} & B & CP:2b W:1a,2b \\
B5016 & 1a<o1 1b*<o1 / 2a<1a1,1b1 2b<1a2,1b1 & IU/EE & {-}{-}/{-}{-} & E & C:2b W:1a \\
B5017 & 1a<o1 1b*<o1 / 2a<1a1,1b1 2b<1a2,1b2 & IU/EE & {-}{-}/{-}{-} & E & C:2b W:1a \\
B5018 & 1a<o1 1b*<o1 / 2a<1a1,1b1 2b<1a2L,1b1 & IU/LI & {-}{-}/{-}{-} & B & P:2b W:1a,2b \\
B5019 & 1a<o1 1b*<o1 / 2a<1a1,1b1 2b<1a2L,1b2 & IU/LI & {-}{-}/{-}{-} & B & P:2b W:1a,2b \\
B5020 & 1a<o1 1b*<o1 / 2a<1a1,1b2 2b<1a1,1b2 & IU/EE & {-}{-}/{-}{-} & E & W:1a \\
B5021 & 1a<o1 1b*<o1 / 2a<1a1,1b2 2b<1a1L,1b1 & IU/LI & {-}{-}/{-}{-} & B & CP:2b W:1a,2b \\
B5022 & 1a<o1 1b*<o1 / 2a<1a1,1b2 2b<1a1L,1b2 & IU/LI & {-}{-}/{-}{-} & B & CP:2b W:1a,2b \\
B5023 & 1a<o1 1b*<o1 / 2a<1a1,1b2 2b<1a2,1b1 & IU/EE & {-}{-}/{-}{-} & E & C:2b W:1a \\
B5024 & 1a<o1 1b*<o1 / 2a<1a1,1b2 2b<1a2,1b2 & IU/EE & {-}{-}/{-}{-} & E & C:2b W:1a \\
B5025 & 1a<o1 1b*<o1 / 2a<1a1,1b2 2b<1a2L,1b1 & IU/LI & {-}{-}/{-}{-} & B & P:2b W:1a,2b \\
B5026 & 1a<o1 1b*<o1 / 2a<1a1,1b2 2b<1a2L,1b2 & IU/LI & {-}{-}/{-}{-} & B & P:2b W:1a,2b \\
B5027 & 1a<o1 1b*<o1 / 2a<1a1L,1b1 2b<1a2,1b1 & IU/IL & {-}{-}/{-}{-} & B & C:2a,2b P:2a W:1a,2a \\
B5028 & 1a<o1 1b*<o1 / 2a<1a1L,1b1 2b<1a2,1b2 & IU/IL & {-}{-}/{-}{-} & B & C:2a,2b P:2a W:1a,2a \\
B5029 & 1a<o1 1b*<o1 / 2a<1a1L,1b2 2b<1a2,1b1 & IU/IL & {-}{-}/{-}{-} & B & C:2a,2b P:2a W:1a,2a \\
B5030 & 1a<o1 1b*<o1 / 2a<1a1L,1b2 2b<1a2,1b2 & IU/IL & {-}{-}/{-}{-} & B & C:2a,2b P:2a W:1a,2a \\
B5031 & 1a<o1 1b*<o1 / 2a<1a2,1b1 2b<1a2,1b1 & IU/EE & {-}{-}/{-}{-} & E & W:1a \\
B5032 & 1a<o1 1b*<o1 / 2a<1a2,1b1 2b<1a2,1b2 & IU/EE & {-}{-}/{-}{-} & E & W:1a \\
B5033 & 1a<o1 1b*<o1 / 2a<1a2,1b1 2b<1a2L,1b1 & IU/LI & {-}{-}/{-}{-} & B & CP:2b W:1a,2b \\
B5034 & 1a<o1 1b*<o1 / 2a<1a2,1b1 2b<1a2L,1b2 & IU/LI & {-}{-}/{-}{-} & B & CP:2b W:1a,2b \\
B5035 & 1a<o1 1b*<o1 / 2a<1a2,1b2 2b<1a2,1b2 & IU/EE & {-}{-}/{-}{-} & E & W:1a \\
B5036 & 1a<o1 1b*<o1 / 2a<1a2,1b2 2b<1a2L,1b1 & IU/LI & {-}{-}/{-}{-} & B & CP:2b W:1a,2b \\
B5037 & 1a<o1 1b*<o1 / 2a<1a2,1b2 2b<1a2L,1b2 & IU/LI & {-}{-}/{-}{-} & B & CP:2b W:1a,2b \\
B5038 & 1a<o1 1b*<o2 / 2a<1a0L,1b1 2b<1a1,1b1 & IU/UI & {-}{-}/{-}{-} & N & W:1a,2b \\
B5039 & 1a<o1 1b*<o2 / 2a<1a0L,1b1 2b<1a1,1b2 & IU/UI & {-}{-}/{-}{-} & N & W:1a,2b \\
B5040 & 1a<o1 1b*<o2 / 2a<1a0L,1b1 2b<1a2,1b1 & IU/UI & {-}{-}/{-}{-} & N & W:1a,2b \\
B5041 & 1a<o1 1b*<o2 / 2a<1a0L,1b1 2b<1a2,1b2 & IU/UI & {-}{-}/{-}{-} & N & W:1a,2b \\
B5042 & 1a<o1 1b*<o2 / 2a<1a0L,1b2 2b<1a1,1b1 & IU/UI & {-}{-}/{-}{-} & N & W:1a,2b \\
B5043 & 1a<o1 1b*<o2 / 2a<1a0L,1b2 2b<1a1,1b2 & IU/UI & {-}{-}/{-}{-} & N & W:1a,2b \\
B5044 & 1a<o1 1b*<o2 / 2a<1a0L,1b2 2b<1a2,1b1 & IU/UI & {-}{-}/{-}{-} & N & W:1a,2b \\
B5045 & 1a<o1 1b*<o2 / 2a<1a0L,1b2 2b<1a2,1b2 & IU/UI & {-}{-}/{-}{-} & N & W:1a,2b \\
B5046 & 1a<o1 1b*<o2 / 2a<1a1,1b1 2b<1a1,1b1 & IU/EE & {-}{-}/{-}{-} & E & W:1a \\
B5047 & 1a<o1 1b*<o2 / 2a<1a1,1b1 2b<1a1,1b2 & IU/EE & {-}{-}/{-}{-} & E & W:1a \\
B5048 & 1a<o1 1b*<o2 / 2a<1a1,1b1 2b<1a1L,1b1 & IU/LI & {-}{-}/{-}{-} & B & CP:2b W:1a,2b \\
B5049 & 1a<o1 1b*<o2 / 2a<1a1,1b1 2b<1a1L,1b2 & IU/LI & {-}{-}/{-}{-} & B & CP:2b W:1a,2b \\
B5050 & 1a<o1 1b*<o2 / 2a<1a1,1b1 2b<1a2,1b1 & IU/EE & {-}{-}/{-}{-} & E & C:2b W:1a \\
B5051 & 1a<o1 1b*<o2 / 2a<1a1,1b1 2b<1a2,1b2 & IU/EE & {-}{-}/{-}{-} & E & C:2b W:1a \\
B5052 & 1a<o1 1b*<o2 / 2a<1a1,1b1 2b<1a2L,1b1 & IU/LI & {-}{-}/{-}{-} & B & P:2b W:1a,2b \\
B5053 & 1a<o1 1b*<o2 / 2a<1a1,1b1 2b<1a2L,1b2 & IU/LI & {-}{-}/{-}{-} & B & P:2b W:1a,2b \\
B5054 & 1a<o1 1b*<o2 / 2a<1a1,1b2 2b<1a1,1b2 & IU/EE & {-}{-}/{-}{-} & E & W:1a \\
B5055 & 1a<o1 1b*<o2 / 2a<1a1,1b2 2b<1a1L,1b1 & IU/LI & {-}{-}/{-}{-} & B & CP:2b W:1a,2b \\
B5056 & 1a<o1 1b*<o2 / 2a<1a1,1b2 2b<1a1L,1b2 & IU/LI & {-}{-}/{-}{-} & B & CP:2b W:1a,2b \\
B5057 & 1a<o1 1b*<o2 / 2a<1a1,1b2 2b<1a2,1b1 & IU/EE & {-}{-}/{-}{-} & E & C:2b W:1a \\
B5058 & 1a<o1 1b*<o2 / 2a<1a1,1b2 2b<1a2,1b2 & IU/EE & {-}{-}/{-}{-} & E & C:2b W:1a \\
B5059 & 1a<o1 1b*<o2 / 2a<1a1,1b2 2b<1a2L,1b1 & IU/LI & {-}{-}/{-}{-} & B & P:2b W:1a,2b \\
B5060 & 1a<o1 1b*<o2 / 2a<1a1,1b2 2b<1a2L,1b2 & IU/LI & {-}{-}/{-}{-} & B & P:2b W:1a,2b \\
B5061 & 1a<o1 1b*<o2 / 2a<1a1L,1b1 2b<1a2,1b1 & IU/IL & {-}{-}/{-}{-} & B & C:2a,2b P:2a W:1a,2a \\
B5062 & 1a<o1 1b*<o2 / 2a<1a1L,1b1 2b<1a2,1b2 & IU/IL & {-}{-}/{-}{-} & B & C:2a,2b P:2a W:1a,2a \\
B5063 & 1a<o1 1b*<o2 / 2a<1a1L,1b2 2b<1a2,1b1 & IU/IL & {-}{-}/{-}{-} & B & C:2a,2b P:2a W:1a,2a \\
B5064 & 1a<o1 1b*<o2 / 2a<1a1L,1b2 2b<1a2,1b2 & IU/IL & {-}{-}/{-}{-} & B & C:2a,2b P:2a W:1a,2a \\
B5065 & 1a<o1 1b*<o2 / 2a<1a2,1b1 2b<1a2,1b1 & IU/EE & {-}{-}/{-}{-} & E & W:1a \\
B5066 & 1a<o1 1b*<o2 / 2a<1a2,1b1 2b<1a2,1b2 & IU/EE & {-}{-}/{-}{-} & E & W:1a \\
B5067 & 1a<o1 1b*<o2 / 2a<1a2,1b1 2b<1a2L,1b1 & IU/LI & {-}{-}/{-}{-} & B & CP:2b W:1a,2b \\
B5068 & 1a<o1 1b*<o2 / 2a<1a2,1b1 2b<1a2L,1b2 & IU/LI & {-}{-}/{-}{-} & B & CP:2b W:1a,2b \\
B5069 & 1a<o1 1b*<o2 / 2a<1a2,1b2 2b<1a2,1b2 & IU/EE & {-}{-}/{-}{-} & E & W:1a \\
B5070 & 1a<o1 1b*<o2 / 2a<1a2,1b2 2b<1a2L,1b1 & IU/LI & {-}{-}/{-}{-} & B & CP:2b W:1a,2b \\
B5071 & 1a<o1 1b*<o2 / 2a<1a2,1b2 2b<1a2L,1b2 & IU/LI & {-}{-}/{-}{-} & B & CP:2b W:1a,2b \\
B5072 & 1a<o1L 1b<o2 / 2a<1a0L,1b1 2b*<1a1,1b1 & IL/DU & +{-}/+{-} & B & C:1a,1b,2a PW:1a \\
B5073 & 1a<o1L 1b<o2 / 2a<1a0L,1b1 2b*<1a1,1b2 & IL/DU & +{-}/+{-} & B & C:1a,1b,2a PW:1a \\
B5074 & 1a<o1L 1b<o2 / 2a<1a0L,1b1 2b*<1a2,1b1 & IL/DU & +{-}/+{-} & B & C:1a,1b,2a PW:1a \\
B5075 & 1a<o1L 1b<o2 / 2a<1a0L,1b1 2b*<1a2,1b2 & IL/DU & +{-}/+{-} & B & C:1a,1b,2a PW:1a \\
B5076 & 1a<o1L 1b<o2 / 2a<1a0L,1b2 2b*<1a1,1b1 & IL/DU & +{-}/+{-} & B & C:1a,1b,2a PW:1a \\
B5077 & 1a<o1L 1b<o2 / 2a<1a0L,1b2 2b*<1a1,1b2 & IL/DU & +{-}/+{-} & B & C:1a,1b,2a PW:1a \\
B5078 & 1a<o1L 1b<o2 / 2a<1a0L,1b2 2b*<1a2,1b1 & IL/DU & +{-}/+{-} & B & C:1a,1b,2a PW:1a \\
B5079 & 1a<o1L 1b<o2 / 2a<1a0L,1b2 2b*<1a2,1b2 & IL/DU & +{-}/+{-} & B & C:1a,1b,2a PW:1a \\
B5080 & 1a<o1L 1b<o2 / 2a<1a1,1b0L 2b*<1a1,1b1 & IL/IU & +{-}/{-}{-} & B & C:1a,1b,2a PW:1a,2a \\
B5081 & 1a<o1L 1b<o2 / 2a<1a1,1b0L 2b*<1a1,1b2 & IL/IU & +{-}/{-}{-} & B & C:1a,1b,2a PW:1a,2a \\
B5082 & 1a<o1L 1b<o2 / 2a<1a1,1b0L 2b*<1a2,1b1 & IL/IU & +{-}/{-}{-} & B & C:1a,1b,2a PW:1a,2a \\
B5083 & 1a<o1L 1b<o2 / 2a<1a1,1b0L 2b*<1a2,1b2 & IL/IU & +{-}/{-}{-} & B & C:1a,1b,2a PW:1a,2a \\
B5084 & 1a<o1L 1b<o2 / 2a<1a1,1b1 2b*<1a1,1b1 & IL/EU & +{-}/{-}{-} & BF & C:1a,1b PW:1a \\
B5085 & 1a<o1L 1b<o2 / 2a<1a1,1b1 2b*<1a1,1b2 & IL/EU & +{-}/{-}{-} & BF & C:1a,1b PW:1a \\
B5086 & 1a<o1L 1b<o2 / 2a<1a1,1b1 2b*<1a2,1b1 & IL/EU & +{-}/{-}{-} & BF & C:1a,1b PW:1a \\
B5087 & 1a<o1L 1b<o2 / 2a<1a1,1b1 2b*<1a2,1b2 & IL/EU & +{-}/{-}{-} & BF & C:1a,1b PW:1a \\
B5088 & 1a<o1L 1b<o2 / 2a<1a1,1b1L 2b*<1a1,1b1 & IL/EU & +{-}/{-}{-} & BF & C:1a,1b PW:1a \\
B5089 & 1a<o1L 1b<o2 / 2a<1a1,1b1L 2b*<1a1,1b2 & IL/EU & +{-}/{-}{-} & BF & C:1a,1b PW:1a \\
B5090 & 1a<o1L 1b<o2 / 2a<1a1,1b1L 2b*<1a2,1b1 & IL/EU & +{-}/{-}{-} & BF & C:1a,1b PW:1a \\
B5091 & 1a<o1L 1b<o2 / 2a<1a1,1b1L 2b*<1a2,1b2 & IL/EU & +{-}/{-}{-} & BF & C:1a,1b PW:1a \\
B5092 & 1a<o1L 1b<o2 / 2a<1a1,1b2 2b*<1a1,1b1 & IL/EU & +{-}/{-}{-} & BF & C:1a,1b,2a PW:1a \\
B5093 & 1a<o1L 1b<o2 / 2a<1a1,1b2 2b*<1a1,1b2 & IL/EU & +{-}/{-}{-} & BF & C:1a,1b,2a PW:1a \\
B5094 & 1a<o1L 1b<o2 / 2a<1a1,1b2 2b*<1a2,1b1 & IL/EU & +{-}/{-}{-} & BF & C:1a,1b,2a PW:1a \\
B5095 & 1a<o1L 1b<o2 / 2a<1a1,1b2 2b*<1a2,1b2 & IL/EU & +{-}/{-}{-} & BF & C:1a,1b,2a PW:1a \\
B5096 & 1a<o1L 1b<o2 / 2a<1a1,1b2L 2b*<1a1,1b1 & IL/EU & +{-}/{-}{-} & BF & C:1a,1b,2a PW:1a \\
B5097 & 1a<o1L 1b<o2 / 2a<1a1,1b2L 2b*<1a1,1b2 & IL/EU & +{-}/{-}{-} & BF & C:1a,1b,2a PW:1a \\
B5098 & 1a<o1L 1b<o2 / 2a<1a1,1b2L 2b*<1a2,1b1 & IL/EU & +{-}/{-}{-} & BF & C:1a,1b,2a PW:1a \\
B5099 & 1a<o1L 1b<o2 / 2a<1a1,1b2L 2b*<1a2,1b2 & IL/EU & +{-}/{-}{-} & BF & C:1a,1b,2a PW:1a \\
B5100 & 1a<o1L 1b<o2 / 2a<1a1L,1b1 2b*<1a1,1b1 & IL/IU & +{-}/+{-} & BF & C:1a,1b,2a PW:1a,2a \\
B5101 & 1a<o1L 1b<o2 / 2a<1a1L,1b1 2b*<1a1,1b2 & IL/IU & +{-}/+{-} & BF & C:1a,1b,2a PW:1a,2a \\
B5102 & 1a<o1L 1b<o2 / 2a<1a1L,1b1 2b*<1a2,1b1 & IL/IU & +{-}/+{-} & BF & C:1a,1b,2a PW:1a,2a \\
B5103 & 1a<o1L 1b<o2 / 2a<1a1L,1b1 2b*<1a2,1b2 & IL/IU & +{-}/+{-} & BF & C:1a,1b,2a PW:1a,2a \\
B5104 & 1a<o1L 1b<o2 / 2a<1a1L,1b2 2b*<1a1,1b1 & IL/IU & +{-}/+{-} & BF & C:1a,1b PW:1a,2a \\
B5105 & 1a<o1L 1b<o2 / 2a<1a1L,1b2 2b*<1a1,1b2 & IL/IU & +{-}/+{-} & BF & C:1a,1b PW:1a,2a \\
B5106 & 1a<o1L 1b<o2 / 2a<1a1L,1b2 2b*<1a2,1b1 & IL/IU & +{-}/+{-} & BF & C:1a,1b PW:1a,2a \\
B5107 & 1a<o1L 1b<o2 / 2a<1a1L,1b2 2b*<1a2,1b2 & IL/IU & +{-}/+{-} & BF & C:1a,1b PW:1a,2a \\
B5108 & 1a<o1L 1b<o2 / 2a<1a2,1b0L 2b*<1a1,1b1 & IL/IU & +{-}/{-}{-} & B & C:1a,1b,2a PW:1a,2a \\
B5109 & 1a<o1L 1b<o2 / 2a<1a2,1b0L 2b*<1a1,1b2 & IL/IU & +{-}/{-}{-} & B & C:1a,1b,2a PW:1a,2a \\
B5110 & 1a<o1L 1b<o2 / 2a<1a2,1b0L 2b*<1a2,1b1 & IL/IU & +{-}/{-}{-} & B & C:1a,1b,2a PW:1a,2a \\
B5111 & 1a<o1L 1b<o2 / 2a<1a2,1b0L 2b*<1a2,1b2 & IL/IU & +{-}/{-}{-} & B & C:1a,1b,2a PW:1a,2a \\
B5112 & 1a<o1L 1b<o2 / 2a<1a2,1b1 2b*<1a1,1b1 & IL/EU & +{-}/{-}{-} & BF & C:1a,1b PW:1a \\
B5113 & 1a<o1L 1b<o2 / 2a<1a2,1b1 2b*<1a1,1b2 & IL/EU & +{-}/{-}{-} & BF & C:1a,1b PW:1a \\
B5114 & 1a<o1L 1b<o2 / 2a<1a2,1b1 2b*<1a2,1b1 & IL/EU & +{-}/{-}{-} & BF & C:1a,1b PW:1a \\
B5115 & 1a<o1L 1b<o2 / 2a<1a2,1b1 2b*<1a2,1b2 & IL/EU & +{-}/{-}{-} & BF & C:1a,1b PW:1a \\
B5116 & 1a<o1L 1b<o2 / 2a<1a2,1b1L 2b*<1a1,1b1 & IL/EU & +{-}/{-}{-} & BF & C:1a,1b PW:1a \\
B5117 & 1a<o1L 1b<o2 / 2a<1a2,1b1L 2b*<1a1,1b2 & IL/EU & +{-}/{-}{-} & BF & C:1a,1b PW:1a \\
B5118 & 1a<o1L 1b<o2 / 2a<1a2,1b1L 2b*<1a2,1b1 & IL/EU & +{-}/{-}{-} & BF & C:1a,1b PW:1a \\
B5119 & 1a<o1L 1b<o2 / 2a<1a2,1b1L 2b*<1a2,1b2 & IL/EU & +{-}/{-}{-} & BF & C:1a,1b PW:1a \\
B5120 & 1a<o1L 1b<o2 / 2a<1a2,1b2 2b*<1a1,1b1 & IL/EU & +{-}/{-}{-} & BF & C:1a,1b PW:1a \\
B5121 & 1a<o1L 1b<o2 / 2a<1a2,1b2 2b*<1a1,1b2 & IL/EU & +{-}/{-}{-} & BF & C:1a,1b PW:1a \\
B5122 & 1a<o1L 1b<o2 / 2a<1a2,1b2 2b*<1a2,1b1 & IL/EU & +{-}/{-}{-} & BF & C:1a,1b PW:1a \\
B5123 & 1a<o1L 1b<o2 / 2a<1a2,1b2 2b*<1a2,1b2 & IL/EU & +{-}/{-}{-} & BF & C:1a,1b PW:1a \\
B5124 & 1a<o1L 1b<o2 / 2a<1a2,1b2L 2b*<1a1,1b1 & IL/EU & +{-}/{-}{-} & BF & C:1a,1b PW:1a \\
B5125 & 1a<o1L 1b<o2 / 2a<1a2,1b2L 2b*<1a1,1b2 & IL/EU & +{-}/{-}{-} & BF & C:1a,1b PW:1a \\
B5126 & 1a<o1L 1b<o2 / 2a<1a2,1b2L 2b*<1a2,1b1 & IL/EU & +{-}/{-}{-} & BF & C:1a,1b PW:1a \\
B5127 & 1a<o1L 1b<o2 / 2a<1a2,1b2L 2b*<1a2,1b2 & IL/EU & +{-}/{-}{-} & BF & C:1a,1b PW:1a \\
B5128 & 1a<o1L 1b<o2 / 2a<1a2L,1b1 2b*<1a1,1b1 & IL/IU & +{-}/+{-} & BF & C:1a,1b,2a PW:1a,2a \\
B5129 & 1a<o1L 1b<o2 / 2a<1a2L,1b1 2b*<1a1,1b2 & IL/IU & +{-}/+{-} & BF & C:1a,1b,2a PW:1a,2a \\
B5130 & 1a<o1L 1b<o2 / 2a<1a2L,1b1 2b*<1a2,1b1 & IL/IU & +{-}/+{-} & BF & C:1a,1b,2a PW:1a,2a \\
B5131 & 1a<o1L 1b<o2 / 2a<1a2L,1b1 2b*<1a2,1b2 & IL/IU & +{-}/+{-} & BF & C:1a,1b,2a PW:1a,2a \\
B5132 & 1a<o1L 1b<o2 / 2a<1a2L,1b2 2b*<1a1,1b1 & IL/IU & +{-}/+{-} & BF & C:1a,1b,2a PW:1a,2a \\
B5133 & 1a<o1L 1b<o2 / 2a<1a2L,1b2 2b*<1a1,1b2 & IL/IU & +{-}/+{-} & BF & C:1a,1b,2a PW:1a,2a \\
B5134 & 1a<o1L 1b<o2 / 2a<1a2L,1b2 2b*<1a2,1b1 & IL/IU & +{-}/+{-} & BF & C:1a,1b,2a PW:1a,2a \\
B5135 & 1a<o1L 1b<o2 / 2a<1a2L,1b2 2b*<1a2,1b2 & IL/IU & +{-}/+{-} & BF & C:1a,1b,2a PW:1a,2a \\
B5136 & 1a<o1L 1b*<o1 / 2a<1a0L,1b1 2b<1a1,1b1 & IU/UI & +{-}/+{-} & N & W:1a,2b \\
B5137 & 1a<o1L 1b*<o1 / 2a<1a0L,1b1 2b<1a1,1b2 & IU/UI & +{-}/+{-} & N & W:1a,2b \\
B5138 & 1a<o1L 1b*<o1 / 2a<1a0L,1b1 2b<1a2,1b1 & IU/UI & +{-}/+{-} & N & W:1a,2b \\
B5139 & 1a<o1L 1b*<o1 / 2a<1a0L,1b1 2b<1a2,1b2 & IU/UI & +{-}/+{-} & N & W:1a,2b \\
B5140 & 1a<o1L 1b*<o1 / 2a<1a0L,1b2 2b<1a1,1b1 & IU/UI & +{-}/+{-} & N & W:1a,2b \\
B5141 & 1a<o1L 1b*<o1 / 2a<1a0L,1b2 2b<1a1,1b2 & IU/UI & +{-}/+{-} & N & W:1a,2b \\
B5142 & 1a<o1L 1b*<o1 / 2a<1a0L,1b2 2b<1a2,1b1 & IU/UI & +{-}/+{-} & N & W:1a,2b \\
B5143 & 1a<o1L 1b*<o1 / 2a<1a0L,1b2 2b<1a2,1b2 & IU/UI & +{-}/+{-} & N & W:1a,2b \\
B5144 & 1a<o1L 1b*<o1 / 2a<1a1,1b1 2b<1a1,1b1 & IU/EE & +{-}/{-}{-} & E & W:1a \\
B5145 & 1a<o1L 1b*<o1 / 2a<1a1,1b1 2b<1a1,1b2 & IU/EE & +{-}/{-}{-} & E & W:1a \\
B5146 & 1a<o1L 1b*<o1 / 2a<1a1,1b1 2b<1a1L,1b1 & IU/LI & +{-}/{-}+ & B & CP:2b W:1a,2b \\
B5147 & 1a<o1L 1b*<o1 / 2a<1a1,1b1 2b<1a1L,1b2 & IU/LI & +{-}/{-}+ & B & CP:2b W:1a,2b \\
B5148 & 1a<o1L 1b*<o1 / 2a<1a1,1b1 2b<1a2,1b1 & IU/EE & +{-}/{-}{-} & E & C:2b W:1a \\
B5149 & 1a<o1L 1b*<o1 / 2a<1a1,1b1 2b<1a2,1b2 & IU/EE & +{-}/{-}{-} & E & C:2b W:1a \\
B5150 & 1a<o1L 1b*<o1 / 2a<1a1,1b1 2b<1a2L,1b1 & IU/LI & +{-}/{-}+ & B & P:2b W:1a,2b \\
B5151 & 1a<o1L 1b*<o1 / 2a<1a1,1b1 2b<1a2L,1b2 & IU/LI & +{-}/{-}+ & B & P:2b W:1a,2b \\
B5152 & 1a<o1L 1b*<o1 / 2a<1a1,1b2 2b<1a1,1b2 & IU/EE & +{-}/{-}{-} & E & W:1a \\
B5153 & 1a<o1L 1b*<o1 / 2a<1a1,1b2 2b<1a1L,1b1 & IU/LI & +{-}/{-}+ & B & CP:2b W:1a,2b \\
B5154 & 1a<o1L 1b*<o1 / 2a<1a1,1b2 2b<1a1L,1b2 & IU/LI & +{-}/{-}+ & B & CP:2b W:1a,2b \\
B5155 & 1a<o1L 1b*<o1 / 2a<1a1,1b2 2b<1a2,1b1 & IU/EE & +{-}/{-}{-} & E & C:2b W:1a \\
B5156 & 1a<o1L 1b*<o1 / 2a<1a1,1b2 2b<1a2,1b2 & IU/EE & +{-}/{-}{-} & E & C:2b W:1a \\
B5157 & 1a<o1L 1b*<o1 / 2a<1a1,1b2 2b<1a2L,1b1 & IU/LI & +{-}/{-}+ & B & P:2b W:1a,2b \\
B5158 & 1a<o1L 1b*<o1 / 2a<1a1,1b2 2b<1a2L,1b2 & IU/LI & +{-}/{-}+ & B & P:2b W:1a,2b \\
B5159 & 1a<o1L 1b*<o1 / 2a<1a1L,1b1 2b<1a2,1b1 & IU/IL & +{-}/+{-} & B & C:2a,2b P:2a W:1a,2a \\
B5160 & 1a<o1L 1b*<o1 / 2a<1a1L,1b1 2b<1a2,1b2 & IU/IL & +{-}/+{-} & B & C:2a,2b P:2a W:1a,2a \\
B5161 & 1a<o1L 1b*<o1 / 2a<1a1L,1b2 2b<1a2,1b1 & IU/IL & +{-}/+{-} & B & C:2a,2b P:2a W:1a,2a \\
B5162 & 1a<o1L 1b*<o1 / 2a<1a1L,1b2 2b<1a2,1b2 & IU/IL & +{-}/+{-} & B & C:2a,2b P:2a W:1a,2a \\
B5163 & 1a<o1L 1b*<o1 / 2a<1a2,1b1 2b<1a2,1b1 & IU/EE & +{-}/{-}{-} & E & W:1a \\
B5164 & 1a<o1L 1b*<o1 / 2a<1a2,1b1 2b<1a2,1b2 & IU/EE & +{-}/{-}{-} & E & W:1a \\
B5165 & 1a<o1L 1b*<o1 / 2a<1a2,1b1 2b<1a2L,1b1 & IU/LI & +{-}/{-}+ & B & CP:2b W:1a,2b \\
B5166 & 1a<o1L 1b*<o1 / 2a<1a2,1b1 2b<1a2L,1b2 & IU/LI & +{-}/{-}+ & B & CP:2b W:1a,2b \\
B5167 & 1a<o1L 1b*<o1 / 2a<1a2,1b2 2b<1a2,1b2 & IU/EE & +{-}/{-}{-} & E & W:1a \\
B5168 & 1a<o1L 1b*<o1 / 2a<1a2,1b2 2b<1a2L,1b1 & IU/LI & +{-}/{-}+ & B & CP:2b W:1a,2b \\
B5169 & 1a<o1L 1b*<o1 / 2a<1a2,1b2 2b<1a2L,1b2 & IU/LI & +{-}/{-}+ & B & CP:2b W:1a,2b \\
B5170 & 1a<o1L 1b*<o2 / 2a<1a0L,1b1 2b<1a1,1b1 & IU/UI & +{-}/+{-} & N & W:1a,2b \\
B5171 & 1a<o1L 1b*<o2 / 2a<1a0L,1b1 2b<1a1,1b2 & IU/UI & +{-}/+{-} & N & W:1a,2b \\
B5172 & 1a<o1L 1b*<o2 / 2a<1a0L,1b1 2b<1a2,1b1 & IU/UI & +{-}/+{-} & N & W:1a,2b \\
B5173 & 1a<o1L 1b*<o2 / 2a<1a0L,1b1 2b<1a2,1b2 & IU/UI & +{-}/+{-} & N & W:1a,2b \\
B5174 & 1a<o1L 1b*<o2 / 2a<1a0L,1b2 2b<1a1,1b1 & IU/UI & +{-}/+{-} & N & W:1a,2b \\
B5175 & 1a<o1L 1b*<o2 / 2a<1a0L,1b2 2b<1a1,1b2 & IU/UI & +{-}/+{-} & N & W:1a,2b \\
B5176 & 1a<o1L 1b*<o2 / 2a<1a0L,1b2 2b<1a2,1b1 & IU/UI & +{-}/+{-} & N & W:1a,2b \\
B5177 & 1a<o1L 1b*<o2 / 2a<1a0L,1b2 2b<1a2,1b2 & IU/UI & +{-}/+{-} & N & W:1a,2b \\
B5178 & 1a<o1L 1b*<o2 / 2a<1a1,1b1 2b<1a1,1b1 & IU/EE & +{-}/{-}{-} & E & W:1a \\
B5179 & 1a<o1L 1b*<o2 / 2a<1a1,1b1 2b<1a1,1b2 & IU/EE & +{-}/{-}{-} & E & W:1a \\
B5180 & 1a<o1L 1b*<o2 / 2a<1a1,1b1 2b<1a1L,1b1 & IU/LI & +{-}/{-}+ & B & CP:2b W:1a,2b \\
B5181 & 1a<o1L 1b*<o2 / 2a<1a1,1b1 2b<1a1L,1b2 & IU/LI & +{-}/{-}+ & B & CP:2b W:1a,2b \\
B5182 & 1a<o1L 1b*<o2 / 2a<1a1,1b1 2b<1a2,1b1 & IU/EE & +{-}/{-}{-} & E & C:2b W:1a \\
B5183 & 1a<o1L 1b*<o2 / 2a<1a1,1b1 2b<1a2,1b2 & IU/EE & +{-}/{-}{-} & E & C:2b W:1a \\
B5184 & 1a<o1L 1b*<o2 / 2a<1a1,1b1 2b<1a2L,1b1 & IU/LI & +{-}/{-}+ & B & P:2b W:1a,2b \\
B5185 & 1a<o1L 1b*<o2 / 2a<1a1,1b1 2b<1a2L,1b2 & IU/LI & +{-}/{-}+ & B & P:2b W:1a,2b \\
B5186 & 1a<o1L 1b*<o2 / 2a<1a1,1b2 2b<1a1,1b2 & IU/EE & +{-}/{-}{-} & E & W:1a \\
B5187 & 1a<o1L 1b*<o2 / 2a<1a1,1b2 2b<1a1L,1b1 & IU/LI & +{-}/{-}+ & B & CP:2b W:1a,2b \\
B5188 & 1a<o1L 1b*<o2 / 2a<1a1,1b2 2b<1a1L,1b2 & IU/LI & +{-}/{-}+ & B & CP:2b W:1a,2b \\
B5189 & 1a<o1L 1b*<o2 / 2a<1a1,1b2 2b<1a2,1b1 & IU/EE & +{-}/{-}{-} & E & C:2b W:1a \\
B5190 & 1a<o1L 1b*<o2 / 2a<1a1,1b2 2b<1a2,1b2 & IU/EE & +{-}/{-}{-} & E & C:2b W:1a \\
B5191 & 1a<o1L 1b*<o2 / 2a<1a1,1b2 2b<1a2L,1b1 & IU/LI & +{-}/{-}+ & B & P:2b W:1a,2b \\
B5192 & 1a<o1L 1b*<o2 / 2a<1a1,1b2 2b<1a2L,1b2 & IU/LI & +{-}/{-}+ & B & P:2b W:1a,2b \\
B5193 & 1a<o1L 1b*<o2 / 2a<1a1L,1b1 2b<1a2,1b1 & IU/IL & +{-}/+{-} & B & C:2a,2b P:2a W:1a,2a \\
B5194 & 1a<o1L 1b*<o2 / 2a<1a1L,1b1 2b<1a2,1b2 & IU/IL & +{-}/+{-} & B & C:2a,2b P:2a W:1a,2a \\
B5195 & 1a<o1L 1b*<o2 / 2a<1a1L,1b2 2b<1a2,1b1 & IU/IL & +{-}/+{-} & B & C:2a,2b P:2a W:1a,2a \\
B5196 & 1a<o1L 1b*<o2 / 2a<1a1L,1b2 2b<1a2,1b2 & IU/IL & +{-}/+{-} & B & C:2a,2b P:2a W:1a,2a \\
B5197 & 1a<o1L 1b*<o2 / 2a<1a2,1b1 2b<1a2,1b1 & IU/EE & +{-}/{-}{-} & E & W:1a \\
B5198 & 1a<o1L 1b*<o2 / 2a<1a2,1b1 2b<1a2,1b2 & IU/EE & +{-}/{-}{-} & E & W:1a \\
B5199 & 1a<o1L 1b*<o2 / 2a<1a2,1b1 2b<1a2L,1b1 & IU/LI & +{-}/{-}+ & B & CP:2b W:1a,2b \\
B5200 & 1a<o1L 1b*<o2 / 2a<1a2,1b1 2b<1a2L,1b2 & IU/LI & +{-}/{-}+ & B & CP:2b W:1a,2b \\
B5201 & 1a<o1L 1b*<o2 / 2a<1a2,1b2 2b<1a2,1b2 & IU/EE & +{-}/{-}{-} & E & W:1a \\
B5202 & 1a<o1L 1b*<o2 / 2a<1a2,1b2 2b<1a2L,1b1 & IU/LI & +{-}/{-}+ & B & CP:2b W:1a,2b \\
B5203 & 1a<o1L 1b*<o2 / 2a<1a2,1b2 2b<1a2L,1b2 & IU/LI & +{-}/{-}+ & B & CP:2b W:1a,2b \\
B5204 & 1a<o2 1b<o2 / 2a<1a0L,1b1 2b*<1a1,1b1 & EE/DU & {-}{-}/{-}{-} & E &  \\
B5205 & 1a<o2 1b<o2 / 2a<1a0L,1b1 2b*<1a1,1b2 & EE/DU & {-}{-}/{-}{-} & E &  \\
B5206 & 1a<o2 1b<o2 / 2a<1a0L,1b1 2b*<1a2,1b1 & EE/DU & {-}{-}/{-}{-} & E &  \\
B5207 & 1a<o2 1b<o2 / 2a<1a0L,1b1 2b*<1a2,1b2 & EE/DU & {-}{-}/{-}{-} & E &  \\
B5208 & 1a<o2 1b<o2 / 2a<1a0L,1b2 2b*<1a1,1b1 & EE/DU & {-}{-}/{-}{-} & E &  \\
B5209 & 1a<o2 1b<o2 / 2a<1a0L,1b2 2b*<1a1,1b2 & EE/DU & {-}{-}/{-}{-} & E &  \\
B5210 & 1a<o2 1b<o2 / 2a<1a0L,1b2 2b*<1a2,1b1 & EE/DU & {-}{-}/{-}{-} & E &  \\
B5211 & 1a<o2 1b<o2 / 2a<1a0L,1b2 2b*<1a2,1b2 & EE/DU & {-}{-}/{-}{-} & E &  \\
B5212 & 1a<o2 1b<o2 / 2a<1a1,1b1 2b*<1a1,1b1 & EE/DU & {-}{-}/{-}{-} & E &  \\
B5213 & 1a<o2 1b<o2 / 2a<1a1,1b1 2b*<1a1,1b2 & EE/DU & {-}{-}/{-}{-} & E &  \\
B5214 & 1a<o2 1b<o2 / 2a<1a1,1b1 2b*<1a2,1b2 & EE/DU & {-}{-}/{-}{-} & E &  \\
B5215 & 1a<o2 1b<o2 / 2a<1a1,1b1L 2b*<1a1,1b1 & EE/DU & {-}{-}/{-}{-} & E &  \\
B5216 & 1a<o2 1b<o2 / 2a<1a1,1b1L 2b*<1a1,1b2 & EE/DU & {-}{-}/{-}{-} & E &  \\
B5217 & 1a<o2 1b<o2 / 2a<1a1,1b1L 2b*<1a2,1b1 & EE/DU & {-}{-}/{-}{-} & E &  \\
B5218 & 1a<o2 1b<o2 / 2a<1a1,1b1L 2b*<1a2,1b2 & EE/DU & {-}{-}/{-}{-} & E &  \\
B5219 & 1a<o2 1b<o2 / 2a<1a1,1b2 2b*<1a1,1b1 & EE/DU & {-}{-}/{-}{-} & E &  \\
B5220 & 1a<o2 1b<o2 / 2a<1a1,1b2 2b*<1a1,1b2 & EE/DU & {-}{-}/{-}{-} & E &  \\
B5221 & 1a<o2 1b<o2 / 2a<1a1,1b2 2b*<1a2,1b1 & EE/DU & {-}{-}/{-}{-} & E &  \\
B5222 & 1a<o2 1b<o2 / 2a<1a1,1b2 2b*<1a2,1b2 & EE/DU & {-}{-}/{-}{-} & E &  \\
B5223 & 1a<o2 1b<o2 / 2a<1a1,1b2L 2b*<1a1,1b1 & EE/DU & {-}{-}/{-}{-} & E &  \\
B5224 & 1a<o2 1b<o2 / 2a<1a1,1b2L 2b*<1a1,1b2 & EE/DU & {-}{-}/{-}{-} & E &  \\
B5225 & 1a<o2 1b<o2 / 2a<1a1,1b2L 2b*<1a2,1b1 & EE/DU & {-}{-}/{-}{-} & E &  \\
B5226 & 1a<o2 1b<o2 / 2a<1a1,1b2L 2b*<1a2,1b2 & EE/DU & {-}{-}/{-}{-} & E &  \\
B5227 & 1a<o2 1b<o2 / 2a<1a1L,1b2 2b*<1a1,1b1 & EE/DU & {-}{-}/{-}{-} & E &  \\
B5228 & 1a<o2 1b<o2 / 2a<1a1L,1b2 2b*<1a1,1b2 & EE/DU & {-}{-}/{-}{-} & E &  \\
B5229 & 1a<o2 1b<o2 / 2a<1a1L,1b2 2b*<1a2,1b1 & EE/DU & {-}{-}/{-}{-} & E &  \\
B5230 & 1a<o2 1b<o2 / 2a<1a1L,1b2 2b*<1a2,1b2 & EE/DU & {-}{-}/{-}{-} & E &  \\
B5231 & 1a<o2 1b<o2 / 2a<1a2,1b2 2b*<1a1,1b1 & EE/DU & {-}{-}/{-}{-} & E &  \\
B5232 & 1a<o2 1b<o2 / 2a<1a2,1b2 2b*<1a1,1b2 & EE/DU & {-}{-}/{-}{-} & E &  \\
B5233 & 1a<o2 1b<o2 / 2a<1a2,1b2 2b*<1a2,1b2 & EE/DU & {-}{-}/{-}{-} & E &  \\
B5234 & 1a<o2 1b<o2 / 2a<1a2,1b2L 2b*<1a1,1b1 & EE/DU & {-}{-}/{-}{-} & E &  \\
B5235 & 1a<o2 1b<o2 / 2a<1a2,1b2L 2b*<1a1,1b2 & EE/DU & {-}{-}/{-}{-} & E &  \\
B5236 & 1a<o2 1b<o2 / 2a<1a2,1b2L 2b*<1a2,1b1 & EE/DU & {-}{-}/{-}{-} & E &  \\
B5237 & 1a<o2 1b<o2 / 2a<1a2,1b2L 2b*<1a2,1b2 & EE/DU & {-}{-}/{-}{-} & E &  \\
B5238 & 1a<o2 1b<o2L / 2a<1a0L,1b1 2b*<1a1,1b1 & LI/IU & {-}+/{-}{-} & B & CPW:1b,2a \\
B5239 & 1a<o2 1b<o2L / 2a<1a0L,1b1 2b*<1a1,1b2 & LI/IU & {-}+/{-}{-} & B & CPW:1b,2a \\
B5240 & 1a<o2 1b<o2L / 2a<1a0L,1b1 2b*<1a2,1b1 & LI/IU & {-}+/{-}{-} & B & CPW:1b,2a \\
B5241 & 1a<o2 1b<o2L / 2a<1a0L,1b1 2b*<1a2,1b2 & LI/IU & {-}+/{-}{-} & B & CPW:1b,2a \\
B5242 & 1a<o2 1b<o2L / 2a<1a0L,1b2 2b*<1a1,1b1 & LI/IU & {-}+/{-}{-} & B & CPW:1b,2a \\
B5243 & 1a<o2 1b<o2L / 2a<1a0L,1b2 2b*<1a1,1b2 & LI/IU & {-}+/{-}{-} & B & CPW:1b,2a \\
B5244 & 1a<o2 1b<o2L / 2a<1a0L,1b2 2b*<1a2,1b1 & LI/IU & {-}+/{-}{-} & B & CPW:1b,2a \\
B5245 & 1a<o2 1b<o2L / 2a<1a0L,1b2 2b*<1a2,1b2 & LI/IU & {-}+/{-}{-} & B & CPW:1b,2a \\
B5246 & 1a<o2 1b<o2L / 2a<1a1,1b0L 2b*<1a1,1b1 & LI/DU & {-}+/+{-} & B & CPW:1b \\
B5247 & 1a<o2 1b<o2L / 2a<1a1,1b0L 2b*<1a1,1b2 & LI/DU & {-}+/+{-} & B & CPW:1b \\
B5248 & 1a<o2 1b<o2L / 2a<1a1,1b0L 2b*<1a2,1b1 & LI/DU & {-}+/+{-} & B & CPW:1b \\
B5249 & 1a<o2 1b<o2L / 2a<1a1,1b0L 2b*<1a2,1b2 & LI/DU & {-}+/+{-} & B & CPW:1b \\
B5250 & 1a<o2 1b<o2L / 2a<1a1,1b1 2b*<1a1,1b1 & LI/EU & {-}+/{-}{-} & BF & CPW:1b \\
B5251 & 1a<o2 1b<o2L / 2a<1a1,1b1 2b*<1a1,1b2 & LI/EU & {-}+/{-}{-} & BF & CPW:1b \\
B5252 & 1a<o2 1b<o2L / 2a<1a1,1b1 2b*<1a2,1b1 & LI/EU & {-}+/{-}{-} & BF & CPW:1b \\
B5253 & 1a<o2 1b<o2L / 2a<1a1,1b1 2b*<1a2,1b2 & LI/EU & {-}+/{-}{-} & BF & CPW:1b \\
B5254 & 1a<o2 1b<o2L / 2a<1a1,1b1L 2b*<1a1,1b1 & LI/IU & {-}+/+{-} & BF & CPW:1b,2a \\
B5255 & 1a<o2 1b<o2L / 2a<1a1,1b1L 2b*<1a1,1b2 & LI/IU & {-}+/+{-} & BF & CPW:1b,2a \\
B5256 & 1a<o2 1b<o2L / 2a<1a1,1b1L 2b*<1a2,1b1 & LI/IU & {-}+/+{-} & BF & CPW:1b,2a \\
B5257 & 1a<o2 1b<o2L / 2a<1a1,1b1L 2b*<1a2,1b2 & LI/IU & {-}+/+{-} & BF & CPW:1b,2a \\
B5258 & 1a<o2 1b<o2L / 2a<1a1,1b2 2b*<1a1,1b1 & LI/EU & {-}+/{-}{-} & BF & CPW:1b \\
B5259 & 1a<o2 1b<o2L / 2a<1a1,1b2 2b*<1a1,1b2 & LI/EU & {-}+/{-}{-} & BF & CPW:1b \\
B5260 & 1a<o2 1b<o2L / 2a<1a1,1b2 2b*<1a2,1b1 & LI/EU & {-}+/{-}{-} & BF & CPW:1b \\
B5261 & 1a<o2 1b<o2L / 2a<1a1,1b2 2b*<1a2,1b2 & LI/EU & {-}+/{-}{-} & BF & CPW:1b \\
B5262 & 1a<o2 1b<o2L / 2a<1a1,1b2L 2b*<1a1,1b1 & LI/IU & {-}+/+{-} & BF & CPW:1b,2a \\
B5263 & 1a<o2 1b<o2L / 2a<1a1,1b2L 2b*<1a1,1b2 & LI/IU & {-}+/+{-} & BF & CPW:1b,2a \\
B5264 & 1a<o2 1b<o2L / 2a<1a1,1b2L 2b*<1a2,1b1 & LI/IU & {-}+/+{-} & BF & CPW:1b,2a \\
B5265 & 1a<o2 1b<o2L / 2a<1a1,1b2L 2b*<1a2,1b2 & LI/IU & {-}+/+{-} & BF & CPW:1b,2a \\
B5266 & 1a<o2 1b<o2L / 2a<1a1L,1b1 2b*<1a1,1b1 & LI/EU & {-}+/{-}{-} & BF & CPW:1b \\
B5267 & 1a<o2 1b<o2L / 2a<1a1L,1b1 2b*<1a1,1b2 & LI/EU & {-}+/{-}{-} & BF & CPW:1b \\
B5268 & 1a<o2 1b<o2L / 2a<1a1L,1b1 2b*<1a2,1b1 & LI/EU & {-}+/{-}{-} & BF & CPW:1b \\
B5269 & 1a<o2 1b<o2L / 2a<1a1L,1b1 2b*<1a2,1b2 & LI/EU & {-}+/{-}{-} & BF & CPW:1b \\
B5270 & 1a<o2 1b<o2L / 2a<1a1L,1b2 2b*<1a1,1b1 & LI/EU & {-}+/{-}{-} & BF & CPW:1b \\
B5271 & 1a<o2 1b<o2L / 2a<1a1L,1b2 2b*<1a1,1b2 & LI/EU & {-}+/{-}{-} & BF & CPW:1b \\
B5272 & 1a<o2 1b<o2L / 2a<1a1L,1b2 2b*<1a2,1b1 & LI/EU & {-}+/{-}{-} & BF & CPW:1b \\
B5273 & 1a<o2 1b<o2L / 2a<1a1L,1b2 2b*<1a2,1b2 & LI/EU & {-}+/{-}{-} & BF & CPW:1b \\
B5274 & 1a<o2 1b<o2L / 2a<1a2,1b0L 2b*<1a1,1b1 & LI/DU & {-}+/+{-} & B & CPW:1b \\
B5275 & 1a<o2 1b<o2L / 2a<1a2,1b0L 2b*<1a1,1b2 & LI/DU & {-}+/+{-} & B & CPW:1b \\
B5276 & 1a<o2 1b<o2L / 2a<1a2,1b0L 2b*<1a2,1b1 & LI/DU & {-}+/+{-} & B & CPW:1b \\
B5277 & 1a<o2 1b<o2L / 2a<1a2,1b0L 2b*<1a2,1b2 & LI/DU & {-}+/+{-} & B & CPW:1b \\
B5278 & 1a<o2 1b<o2L / 2a<1a2,1b1 2b*<1a1,1b1 & LI/EU & {-}+/{-}{-} & BF & CPW:1b \\
B5279 & 1a<o2 1b<o2L / 2a<1a2,1b1 2b*<1a1,1b2 & LI/EU & {-}+/{-}{-} & BF & CPW:1b \\
B5280 & 1a<o2 1b<o2L / 2a<1a2,1b1 2b*<1a2,1b1 & LI/EU & {-}+/{-}{-} & BF & CPW:1b \\
B5281 & 1a<o2 1b<o2L / 2a<1a2,1b1 2b*<1a2,1b2 & LI/EU & {-}+/{-}{-} & BF & CPW:1b \\
B5282 & 1a<o2 1b<o2L / 2a<1a2,1b1L 2b*<1a1,1b1 & LI/IU & {-}+/+{-} & BF & CPW:1b,2a \\
B5283 & 1a<o2 1b<o2L / 2a<1a2,1b1L 2b*<1a1,1b2 & LI/IU & {-}+/+{-} & BF & CPW:1b,2a \\
B5284 & 1a<o2 1b<o2L / 2a<1a2,1b1L 2b*<1a2,1b1 & LI/IU & {-}+/+{-} & BF & CPW:1b,2a \\
B5285 & 1a<o2 1b<o2L / 2a<1a2,1b1L 2b*<1a2,1b2 & LI/IU & {-}+/+{-} & BF & CPW:1b,2a \\
B5286 & 1a<o2 1b<o2L / 2a<1a2,1b2 2b*<1a1,1b1 & LI/EU & {-}+/{-}{-} & BF & CPW:1b \\
B5287 & 1a<o2 1b<o2L / 2a<1a2,1b2 2b*<1a1,1b2 & LI/EU & {-}+/{-}{-} & BF & CPW:1b \\
B5288 & 1a<o2 1b<o2L / 2a<1a2,1b2 2b*<1a2,1b1 & LI/EU & {-}+/{-}{-} & BF & CPW:1b \\
B5289 & 1a<o2 1b<o2L / 2a<1a2,1b2 2b*<1a2,1b2 & LI/EU & {-}+/{-}{-} & BF & CPW:1b \\
B5290 & 1a<o2 1b<o2L / 2a<1a2,1b2L 2b*<1a1,1b1 & LI/IU & {-}+/+{-} & BF & CPW:1b,2a \\
B5291 & 1a<o2 1b<o2L / 2a<1a2,1b2L 2b*<1a1,1b2 & LI/IU & {-}+/+{-} & BF & CPW:1b,2a \\
B5292 & 1a<o2 1b<o2L / 2a<1a2,1b2L 2b*<1a2,1b1 & LI/IU & {-}+/+{-} & BF & CPW:1b,2a \\
B5293 & 1a<o2 1b<o2L / 2a<1a2,1b2L 2b*<1a2,1b2 & LI/IU & {-}+/+{-} & BF & CPW:1b,2a \\
B5294 & 1a<o2 1b<o2L / 2a<1a2L,1b1 2b*<1a1,1b1 & LI/EU & {-}+/{-}{-} & BF & CPW:1b \\
B5295 & 1a<o2 1b<o2L / 2a<1a2L,1b1 2b*<1a1,1b2 & LI/EU & {-}+/{-}{-} & BF & CPW:1b \\
B5296 & 1a<o2 1b<o2L / 2a<1a2L,1b1 2b*<1a2,1b1 & LI/EU & {-}+/{-}{-} & BF & CPW:1b \\
B5297 & 1a<o2 1b<o2L / 2a<1a2L,1b1 2b*<1a2,1b2 & LI/EU & {-}+/{-}{-} & BF & CPW:1b \\
B5298 & 1a<o2 1b<o2L / 2a<1a2L,1b2 2b*<1a1,1b1 & LI/EU & {-}+/{-}{-} & BF & CPW:1b \\
B5299 & 1a<o2 1b<o2L / 2a<1a2L,1b2 2b*<1a1,1b2 & LI/EU & {-}+/{-}{-} & BF & CPW:1b \\
B5300 & 1a<o2 1b<o2L / 2a<1a2L,1b2 2b*<1a2,1b1 & LI/EU & {-}+/{-}{-} & BF & CPW:1b \\
B5301 & 1a<o2 1b<o2L / 2a<1a2L,1b2 2b*<1a2,1b2 & LI/EU & {-}+/{-}{-} & BF & CPW:1b \\
B5302 & 1a<o2 1b*<o1 / 2a<1a0L,1b1 2b<1a1,1b1 & IU/UI & {-}{-}/{-}{-} & N & W:1a,2b \\
B5303 & 1a<o2 1b*<o1 / 2a<1a0L,1b1 2b<1a1,1b2 & IU/UI & {-}{-}/{-}{-} & N & W:1a,2b \\
B5304 & 1a<o2 1b*<o1 / 2a<1a0L,1b1 2b<1a2,1b1 & IU/UI & {-}{-}/{-}{-} & N & W:1a,2b \\
B5305 & 1a<o2 1b*<o1 / 2a<1a0L,1b1 2b<1a2,1b2 & IU/UI & {-}{-}/{-}{-} & N & W:1a,2b \\
B5306 & 1a<o2 1b*<o1 / 2a<1a0L,1b2 2b<1a1,1b1 & IU/UI & {-}{-}/{-}{-} & N & W:1a,2b \\
B5307 & 1a<o2 1b*<o1 / 2a<1a0L,1b2 2b<1a1,1b2 & IU/UI & {-}{-}/{-}{-} & N & W:1a,2b \\
B5308 & 1a<o2 1b*<o1 / 2a<1a0L,1b2 2b<1a2,1b1 & IU/UI & {-}{-}/{-}{-} & N & W:1a,2b \\
B5309 & 1a<o2 1b*<o1 / 2a<1a0L,1b2 2b<1a2,1b2 & IU/UI & {-}{-}/{-}{-} & N & W:1a,2b \\
B5310 & 1a<o2 1b*<o1 / 2a<1a1,1b1 2b<1a1,1b1 & IU/EE & {-}{-}/{-}{-} & E & W:1a \\
B5311 & 1a<o2 1b*<o1 / 2a<1a1,1b1 2b<1a1,1b2 & IU/EE & {-}{-}/{-}{-} & E & W:1a \\
B5312 & 1a<o2 1b*<o1 / 2a<1a1,1b1 2b<1a1L,1b1 & IU/LI & {-}{-}/{-}{-} & B & CP:2b W:1a,2b \\
B5313 & 1a<o2 1b*<o1 / 2a<1a1,1b1 2b<1a1L,1b2 & IU/LI & {-}{-}/{-}{-} & B & CP:2b W:1a,2b \\
B5314 & 1a<o2 1b*<o1 / 2a<1a1,1b1 2b<1a2,1b1 & IU/EE & {-}{-}/{-}{-} & E & C:2b W:1a \\
B5315 & 1a<o2 1b*<o1 / 2a<1a1,1b1 2b<1a2,1b2 & IU/EE & {-}{-}/{-}{-} & E & C:2b W:1a \\
B5316 & 1a<o2 1b*<o1 / 2a<1a1,1b1 2b<1a2L,1b1 & IU/LI & {-}{-}/{-}{-} & B & P:2b W:1a,2b \\
B5317 & 1a<o2 1b*<o1 / 2a<1a1,1b1 2b<1a2L,1b2 & IU/LI & {-}{-}/{-}{-} & B & P:2b W:1a,2b \\
B5318 & 1a<o2 1b*<o1 / 2a<1a1,1b2 2b<1a1,1b2 & IU/EE & {-}{-}/{-}{-} & E & W:1a \\
B5319 & 1a<o2 1b*<o1 / 2a<1a1,1b2 2b<1a1L,1b1 & IU/LI & {-}{-}/{-}{-} & B & CP:2b W:1a,2b \\
B5320 & 1a<o2 1b*<o1 / 2a<1a1,1b2 2b<1a1L,1b2 & IU/LI & {-}{-}/{-}{-} & B & CP:2b W:1a,2b \\
B5321 & 1a<o2 1b*<o1 / 2a<1a1,1b2 2b<1a2,1b1 & IU/EE & {-}{-}/{-}{-} & E & C:2b W:1a \\
B5322 & 1a<o2 1b*<o1 / 2a<1a1,1b2 2b<1a2,1b2 & IU/EE & {-}{-}/{-}{-} & E & C:2b W:1a \\
B5323 & 1a<o2 1b*<o1 / 2a<1a1,1b2 2b<1a2L,1b1 & IU/LI & {-}{-}/{-}{-} & B & P:2b W:1a,2b \\
B5324 & 1a<o2 1b*<o1 / 2a<1a1,1b2 2b<1a2L,1b2 & IU/LI & {-}{-}/{-}{-} & B & P:2b W:1a,2b \\
B5325 & 1a<o2 1b*<o1 / 2a<1a1L,1b1 2b<1a2,1b1 & IU/IL & {-}{-}/{-}{-} & B & C:2a,2b P:2a W:1a,2a \\
B5326 & 1a<o2 1b*<o1 / 2a<1a1L,1b1 2b<1a2,1b2 & IU/IL & {-}{-}/{-}{-} & B & C:2a,2b P:2a W:1a,2a \\
B5327 & 1a<o2 1b*<o1 / 2a<1a1L,1b2 2b<1a2,1b1 & IU/IL & {-}{-}/{-}{-} & B & C:2a,2b P:2a W:1a,2a \\
B5328 & 1a<o2 1b*<o1 / 2a<1a1L,1b2 2b<1a2,1b2 & IU/IL & {-}{-}/{-}{-} & B & C:2a,2b P:2a W:1a,2a \\
B5329 & 1a<o2 1b*<o1 / 2a<1a2,1b1 2b<1a2,1b1 & IU/EE & {-}{-}/{-}{-} & E & W:1a \\
B5330 & 1a<o2 1b*<o1 / 2a<1a2,1b1 2b<1a2,1b2 & IU/EE & {-}{-}/{-}{-} & E & W:1a \\
B5331 & 1a<o2 1b*<o1 / 2a<1a2,1b1 2b<1a2L,1b1 & IU/LI & {-}{-}/{-}{-} & B & CP:2b W:1a,2b \\
B5332 & 1a<o2 1b*<o1 / 2a<1a2,1b1 2b<1a2L,1b2 & IU/LI & {-}{-}/{-}{-} & B & CP:2b W:1a,2b \\
B5333 & 1a<o2 1b*<o1 / 2a<1a2,1b2 2b<1a2,1b2 & IU/EE & {-}{-}/{-}{-} & E & W:1a \\
B5334 & 1a<o2 1b*<o1 / 2a<1a2,1b2 2b<1a2L,1b1 & IU/LI & {-}{-}/{-}{-} & B & CP:2b W:1a,2b \\
B5335 & 1a<o2 1b*<o1 / 2a<1a2,1b2 2b<1a2L,1b2 & IU/LI & {-}{-}/{-}{-} & B & CP:2b W:1a,2b \\
B5336 & 1a<o2 1b*<o2 / 2a<1a0L,1b1 2b<1a1,1b1 & IU/UI & {-}{-}/{-}{-} & N & W:1a,2b \\
B5337 & 1a<o2 1b*<o2 / 2a<1a0L,1b1 2b<1a1,1b2 & IU/UI & {-}{-}/{-}{-} & N & W:1a,2b \\
B5338 & 1a<o2 1b*<o2 / 2a<1a0L,1b1 2b<1a2,1b1 & IU/UI & {-}{-}/{-}{-} & N & W:1a,2b \\
B5339 & 1a<o2 1b*<o2 / 2a<1a0L,1b1 2b<1a2,1b2 & IU/UI & {-}{-}/{-}{-} & N & W:1a,2b \\
B5340 & 1a<o2 1b*<o2 / 2a<1a0L,1b2 2b<1a1,1b1 & IU/UI & {-}{-}/{-}{-} & N & W:1a,2b \\
B5341 & 1a<o2 1b*<o2 / 2a<1a0L,1b2 2b<1a1,1b2 & IU/UI & {-}{-}/{-}{-} & N & W:1a,2b \\
B5342 & 1a<o2 1b*<o2 / 2a<1a0L,1b2 2b<1a2,1b1 & IU/UI & {-}{-}/{-}{-} & N & W:1a,2b \\
B5343 & 1a<o2 1b*<o2 / 2a<1a0L,1b2 2b<1a2,1b2 & IU/UI & {-}{-}/{-}{-} & N & W:1a,2b \\
B5344 & 1a<o2 1b*<o2 / 2a<1a1,1b1 2b<1a1,1b1 & IU/EE & {-}{-}/{-}{-} & E & W:1a \\
B5345 & 1a<o2 1b*<o2 / 2a<1a1,1b1 2b<1a1,1b2 & IU/EE & {-}{-}/{-}{-} & E & W:1a \\
B5346 & 1a<o2 1b*<o2 / 2a<1a1,1b1 2b<1a1L,1b1 & IU/LI & {-}{-}/{-}{-} & B & CP:2b W:1a,2b \\
B5347 & 1a<o2 1b*<o2 / 2a<1a1,1b1 2b<1a1L,1b2 & IU/LI & {-}{-}/{-}{-} & B & CP:2b W:1a,2b \\
B5348 & 1a<o2 1b*<o2 / 2a<1a1,1b1 2b<1a2,1b1 & IU/EE & {-}{-}/{-}{-} & E & C:2b W:1a \\
B5349 & 1a<o2 1b*<o2 / 2a<1a1,1b1 2b<1a2,1b2 & IU/EE & {-}{-}/{-}{-} & E & C:2b W:1a \\
B5350 & 1a<o2 1b*<o2 / 2a<1a1,1b1 2b<1a2L,1b1 & IU/LI & {-}{-}/{-}{-} & B & P:2b W:1a,2b \\
B5351 & 1a<o2 1b*<o2 / 2a<1a1,1b1 2b<1a2L,1b2 & IU/LI & {-}{-}/{-}{-} & B & P:2b W:1a,2b \\
B5352 & 1a<o2 1b*<o2 / 2a<1a1,1b2 2b<1a1,1b2 & IU/EE & {-}{-}/{-}{-} & E & W:1a \\
B5353 & 1a<o2 1b*<o2 / 2a<1a1,1b2 2b<1a1L,1b1 & IU/LI & {-}{-}/{-}{-} & B & CP:2b W:1a,2b \\
B5354 & 1a<o2 1b*<o2 / 2a<1a1,1b2 2b<1a1L,1b2 & IU/LI & {-}{-}/{-}{-} & B & CP:2b W:1a,2b \\
B5355 & 1a<o2 1b*<o2 / 2a<1a1,1b2 2b<1a2,1b1 & IU/EE & {-}{-}/{-}{-} & E & C:2b W:1a \\
B5356 & 1a<o2 1b*<o2 / 2a<1a1,1b2 2b<1a2,1b2 & IU/EE & {-}{-}/{-}{-} & E & C:2b W:1a \\
B5357 & 1a<o2 1b*<o2 / 2a<1a1,1b2 2b<1a2L,1b1 & IU/LI & {-}{-}/{-}{-} & B & P:2b W:1a,2b \\
B5358 & 1a<o2 1b*<o2 / 2a<1a1,1b2 2b<1a2L,1b2 & IU/LI & {-}{-}/{-}{-} & B & P:2b W:1a,2b \\
B5359 & 1a<o2 1b*<o2 / 2a<1a1L,1b1 2b<1a2,1b1 & IU/IL & {-}{-}/{-}{-} & B & C:2a,2b P:2a W:1a,2a \\
B5360 & 1a<o2 1b*<o2 / 2a<1a1L,1b1 2b<1a2,1b2 & IU/IL & {-}{-}/{-}{-} & B & C:2a,2b P:2a W:1a,2a \\
B5361 & 1a<o2 1b*<o2 / 2a<1a1L,1b2 2b<1a2,1b1 & IU/IL & {-}{-}/{-}{-} & B & C:2a,2b P:2a W:1a,2a \\
B5362 & 1a<o2 1b*<o2 / 2a<1a1L,1b2 2b<1a2,1b2 & IU/IL & {-}{-}/{-}{-} & B & C:2a,2b P:2a W:1a,2a \\
B5363 & 1a<o2 1b*<o2 / 2a<1a2,1b1 2b<1a2,1b1 & IU/EE & {-}{-}/{-}{-} & E & W:1a \\
B5364 & 1a<o2 1b*<o2 / 2a<1a2,1b1 2b<1a2,1b2 & IU/EE & {-}{-}/{-}{-} & E & W:1a \\
B5365 & 1a<o2 1b*<o2 / 2a<1a2,1b1 2b<1a2L,1b1 & IU/LI & {-}{-}/{-}{-} & B & CP:2b W:1a,2b \\
B5366 & 1a<o2 1b*<o2 / 2a<1a2,1b1 2b<1a2L,1b2 & IU/LI & {-}{-}/{-}{-} & B & CP:2b W:1a,2b \\
B5367 & 1a<o2 1b*<o2 / 2a<1a2,1b2 2b<1a2,1b2 & IU/EE & {-}{-}/{-}{-} & E & W:1a \\
B5368 & 1a<o2 1b*<o2 / 2a<1a2,1b2 2b<1a2L,1b1 & IU/LI & {-}{-}/{-}{-} & B & CP:2b W:1a,2b \\
B5369 & 1a<o2 1b*<o2 / 2a<1a2,1b2 2b<1a2L,1b2 & IU/LI & {-}{-}/{-}{-} & B & CP:2b W:1a,2b \\
B5370 & 1a<o2L 1b*<o1 / 2a<1a0L,1b1 2b<1a1,1b1 & IU/UI & +{-}/+{-} & N & W:1a,2b \\
B5371 & 1a<o2L 1b*<o1 / 2a<1a0L,1b1 2b<1a1,1b2 & IU/UI & +{-}/+{-} & N & W:1a,2b \\
B5372 & 1a<o2L 1b*<o1 / 2a<1a0L,1b1 2b<1a2,1b1 & IU/UI & +{-}/+{-} & N & W:1a,2b \\
B5373 & 1a<o2L 1b*<o1 / 2a<1a0L,1b1 2b<1a2,1b2 & IU/UI & +{-}/+{-} & N & W:1a,2b \\
B5374 & 1a<o2L 1b*<o1 / 2a<1a0L,1b2 2b<1a1,1b1 & IU/UI & +{-}/+{-} & N & W:1a,2b \\
B5375 & 1a<o2L 1b*<o1 / 2a<1a0L,1b2 2b<1a1,1b2 & IU/UI & +{-}/+{-} & N & W:1a,2b \\
B5376 & 1a<o2L 1b*<o1 / 2a<1a0L,1b2 2b<1a2,1b1 & IU/UI & +{-}/+{-} & N & W:1a,2b \\
B5377 & 1a<o2L 1b*<o1 / 2a<1a0L,1b2 2b<1a2,1b2 & IU/UI & +{-}/+{-} & N & W:1a,2b \\
B5378 & 1a<o2L 1b*<o1 / 2a<1a1,1b1 2b<1a1,1b1 & IU/EE & +{-}/{-}{-} & E & W:1a \\
B5379 & 1a<o2L 1b*<o1 / 2a<1a1,1b1 2b<1a1,1b2 & IU/EE & +{-}/{-}{-} & E & W:1a \\
B5380 & 1a<o2L 1b*<o1 / 2a<1a1,1b1 2b<1a1L,1b1 & IU/LI & +{-}/{-}+ & B & CP:2b W:1a,2b \\
B5381 & 1a<o2L 1b*<o1 / 2a<1a1,1b1 2b<1a1L,1b2 & IU/LI & +{-}/{-}+ & B & CP:2b W:1a,2b \\
B5382 & 1a<o2L 1b*<o1 / 2a<1a1,1b1 2b<1a2,1b1 & IU/EE & +{-}/{-}{-} & E & C:2b W:1a \\
B5383 & 1a<o2L 1b*<o1 / 2a<1a1,1b1 2b<1a2,1b2 & IU/EE & +{-}/{-}{-} & E & C:2b W:1a \\
B5384 & 1a<o2L 1b*<o1 / 2a<1a1,1b1 2b<1a2L,1b1 & IU/LI & +{-}/{-}+ & B & P:2b W:1a,2b \\
B5385 & 1a<o2L 1b*<o1 / 2a<1a1,1b1 2b<1a2L,1b2 & IU/LI & +{-}/{-}+ & B & P:2b W:1a,2b \\
B5386 & 1a<o2L 1b*<o1 / 2a<1a1,1b2 2b<1a1,1b2 & IU/EE & +{-}/{-}{-} & E & W:1a \\
B5387 & 1a<o2L 1b*<o1 / 2a<1a1,1b2 2b<1a1L,1b1 & IU/LI & +{-}/{-}+ & B & CP:2b W:1a,2b \\
B5388 & 1a<o2L 1b*<o1 / 2a<1a1,1b2 2b<1a1L,1b2 & IU/LI & +{-}/{-}+ & B & CP:2b W:1a,2b \\
B5389 & 1a<o2L 1b*<o1 / 2a<1a1,1b2 2b<1a2,1b1 & IU/EE & +{-}/{-}{-} & E & C:2b W:1a \\
B5390 & 1a<o2L 1b*<o1 / 2a<1a1,1b2 2b<1a2,1b2 & IU/EE & +{-}/{-}{-} & E & C:2b W:1a \\
B5391 & 1a<o2L 1b*<o1 / 2a<1a1,1b2 2b<1a2L,1b1 & IU/LI & +{-}/{-}+ & B & P:2b W:1a,2b \\
B5392 & 1a<o2L 1b*<o1 / 2a<1a1,1b2 2b<1a2L,1b2 & IU/LI & +{-}/{-}+ & B & P:2b W:1a,2b \\
B5393 & 1a<o2L 1b*<o1 / 2a<1a1L,1b1 2b<1a2,1b1 & IU/IL & +{-}/+{-} & B & C:2a,2b P:2a W:1a,2a \\
B5394 & 1a<o2L 1b*<o1 / 2a<1a1L,1b1 2b<1a2,1b2 & IU/IL & +{-}/+{-} & B & C:2a,2b P:2a W:1a,2a \\
B5395 & 1a<o2L 1b*<o1 / 2a<1a1L,1b2 2b<1a2,1b1 & IU/IL & +{-}/+{-} & B & C:2a,2b P:2a W:1a,2a \\
B5396 & 1a<o2L 1b*<o1 / 2a<1a1L,1b2 2b<1a2,1b2 & IU/IL & +{-}/+{-} & B & C:2a,2b P:2a W:1a,2a \\
B5397 & 1a<o2L 1b*<o1 / 2a<1a2,1b1 2b<1a2,1b1 & IU/EE & +{-}/{-}{-} & E & W:1a \\
B5398 & 1a<o2L 1b*<o1 / 2a<1a2,1b1 2b<1a2,1b2 & IU/EE & +{-}/{-}{-} & E & W:1a \\
B5399 & 1a<o2L 1b*<o1 / 2a<1a2,1b1 2b<1a2L,1b1 & IU/LI & +{-}/{-}+ & B & CP:2b W:1a,2b \\
B5400 & 1a<o2L 1b*<o1 / 2a<1a2,1b1 2b<1a2L,1b2 & IU/LI & +{-}/{-}+ & B & CP:2b W:1a,2b \\
B5401 & 1a<o2L 1b*<o1 / 2a<1a2,1b2 2b<1a2,1b2 & IU/EE & +{-}/{-}{-} & E & W:1a \\
B5402 & 1a<o2L 1b*<o1 / 2a<1a2,1b2 2b<1a2L,1b1 & IU/LI & +{-}/{-}+ & B & CP:2b W:1a,2b \\
B5403 & 1a<o2L 1b*<o1 / 2a<1a2,1b2 2b<1a2L,1b2 & IU/LI & +{-}/{-}+ & B & CP:2b W:1a,2b \\
B5404 & 1a<o2L 1b*<o2 / 2a<1a0L,1b1 2b<1a1,1b1 & IU/UI & +{-}/+{-} & N & W:1a,2b \\
B5405 & 1a<o2L 1b*<o2 / 2a<1a0L,1b1 2b<1a1,1b2 & IU/UI & +{-}/+{-} & N & W:1a,2b \\
B5406 & 1a<o2L 1b*<o2 / 2a<1a0L,1b1 2b<1a2,1b1 & IU/UI & +{-}/+{-} & N & W:1a,2b \\
B5407 & 1a<o2L 1b*<o2 / 2a<1a0L,1b1 2b<1a2,1b2 & IU/UI & +{-}/+{-} & N & W:1a,2b \\
B5408 & 1a<o2L 1b*<o2 / 2a<1a0L,1b2 2b<1a1,1b1 & IU/UI & +{-}/+{-} & N & W:1a,2b \\
B5409 & 1a<o2L 1b*<o2 / 2a<1a0L,1b2 2b<1a1,1b2 & IU/UI & +{-}/+{-} & N & W:1a,2b \\
B5410 & 1a<o2L 1b*<o2 / 2a<1a0L,1b2 2b<1a2,1b1 & IU/UI & +{-}/+{-} & N & W:1a,2b \\
B5411 & 1a<o2L 1b*<o2 / 2a<1a0L,1b2 2b<1a2,1b2 & IU/UI & +{-}/+{-} & N & W:1a,2b \\
B5412 & 1a<o2L 1b*<o2 / 2a<1a1,1b1 2b<1a1,1b1 & IU/EE & +{-}/{-}{-} & E & W:1a \\
B5413 & 1a<o2L 1b*<o2 / 2a<1a1,1b1 2b<1a1,1b2 & IU/EE & +{-}/{-}{-} & E & W:1a \\
B5414 & 1a<o2L 1b*<o2 / 2a<1a1,1b1 2b<1a1L,1b1 & IU/LI & +{-}/{-}+ & B & CP:2b W:1a,2b \\
B5415 & 1a<o2L 1b*<o2 / 2a<1a1,1b1 2b<1a1L,1b2 & IU/LI & +{-}/{-}+ & B & CP:2b W:1a,2b \\
B5416 & 1a<o2L 1b*<o2 / 2a<1a1,1b1 2b<1a2,1b1 & IU/EE & +{-}/{-}{-} & E & C:2b W:1a \\
B5417 & 1a<o2L 1b*<o2 / 2a<1a1,1b1 2b<1a2,1b2 & IU/EE & +{-}/{-}{-} & E & C:2b W:1a \\
B5418 & 1a<o2L 1b*<o2 / 2a<1a1,1b1 2b<1a2L,1b1 & IU/LI & +{-}/{-}+ & B & P:2b W:1a,2b \\
B5419 & 1a<o2L 1b*<o2 / 2a<1a1,1b1 2b<1a2L,1b2 & IU/LI & +{-}/{-}+ & B & P:2b W:1a,2b \\
B5420 & 1a<o2L 1b*<o2 / 2a<1a1,1b2 2b<1a1,1b2 & IU/EE & +{-}/{-}{-} & E & W:1a \\
B5421 & 1a<o2L 1b*<o2 / 2a<1a1,1b2 2b<1a1L,1b1 & IU/LI & +{-}/{-}+ & B & CP:2b W:1a,2b \\
B5422 & 1a<o2L 1b*<o2 / 2a<1a1,1b2 2b<1a1L,1b2 & IU/LI & +{-}/{-}+ & B & CP:2b W:1a,2b \\
B5423 & 1a<o2L 1b*<o2 / 2a<1a1,1b2 2b<1a2,1b1 & IU/EE & +{-}/{-}{-} & E & C:2b W:1a \\
B5424 & 1a<o2L 1b*<o2 / 2a<1a1,1b2 2b<1a2,1b2 & IU/EE & +{-}/{-}{-} & E & C:2b W:1a \\
B5425 & 1a<o2L 1b*<o2 / 2a<1a1,1b2 2b<1a2L,1b1 & IU/LI & +{-}/{-}+ & B & P:2b W:1a,2b \\
B5426 & 1a<o2L 1b*<o2 / 2a<1a1,1b2 2b<1a2L,1b2 & IU/LI & +{-}/{-}+ & B & P:2b W:1a,2b \\
B5427 & 1a<o2L 1b*<o2 / 2a<1a1L,1b1 2b<1a2,1b1 & IU/IL & +{-}/+{-} & B & C:2a,2b P:2a W:1a,2a \\
B5428 & 1a<o2L 1b*<o2 / 2a<1a1L,1b1 2b<1a2,1b2 & IU/IL & +{-}/+{-} & B & C:2a,2b P:2a W:1a,2a \\
B5429 & 1a<o2L 1b*<o2 / 2a<1a1L,1b2 2b<1a2,1b1 & IU/IL & +{-}/+{-} & B & C:2a,2b P:2a W:1a,2a \\
B5430 & 1a<o2L 1b*<o2 / 2a<1a1L,1b2 2b<1a2,1b2 & IU/IL & +{-}/+{-} & B & C:2a,2b P:2a W:1a,2a \\
B5431 & 1a<o2L 1b*<o2 / 2a<1a2,1b1 2b<1a2,1b1 & IU/EE & +{-}/{-}{-} & E & W:1a \\
B5432 & 1a<o2L 1b*<o2 / 2a<1a2,1b1 2b<1a2,1b2 & IU/EE & +{-}/{-}{-} & E & W:1a \\
B5433 & 1a<o2L 1b*<o2 / 2a<1a2,1b1 2b<1a2L,1b1 & IU/LI & +{-}/{-}+ & B & CP:2b W:1a,2b \\
B5434 & 1a<o2L 1b*<o2 / 2a<1a2,1b1 2b<1a2L,1b2 & IU/LI & +{-}/{-}+ & B & CP:2b W:1a,2b \\
B5435 & 1a<o2L 1b*<o2 / 2a<1a2,1b2 2b<1a2,1b2 & IU/EE & +{-}/{-}{-} & E & W:1a \\
B5436 & 1a<o2L 1b*<o2 / 2a<1a2,1b2 2b<1a2L,1b1 & IU/LI & +{-}/{-}+ & B & CP:2b W:1a,2b \\
B5437 & 1a<o2L 1b*<o2 / 2a<1a2,1b2 2b<1a2L,1b2 & IU/LI & +{-}/{-}+ & B & CP:2b W:1a,2b \\
\end{casetable}
\subsection*{Height 3, 3 stages after the original (215 cases)}
\begin{casetable}
C0001 & 1a<o0L / 2a<1a0L / 3a<2a0L & N/N/N & +/+/+ & N &  \\
C0002 & 1a<o0L / 2a<1a0L / 3a<2a1 & N/N/N & +/+/{-} & N &  \\
C0003 & 1a<o0L / 2a<1a0L / 3a<2a1L & N/N/N & +/+/+ & N &  \\
C0004 & 1a<o0L / 2a<1a0L / 3a<2a2 & N/N/N & +/+/{-} & N &  \\
C0005 & 1a<o0L / 2a<1a0L / 3a<2a2L & N/N/N & +/+/+ & N &  \\
C0006 & 1a<o0L / 2a<1a1 / 3a<2a0L & N/N/N & +/{-}/{-} & N &  \\
C0007 & 1a<o0L / 2a<1a1 / 3a<2a1 & N/N/N & +/{-}/{-} & N &  \\
C0008 & 1a<o0L / 2a<1a1 / 3a<2a1L & N/N/N & +/{-}/{-} & N &  \\
C0009 & 1a<o0L / 2a<1a1 / 3a<2a2 & N/N/N & +/{-}/{-} & N &  \\
C0010 & 1a<o0L / 2a<1a1 / 3a<2a2L & N/N/N & +/{-}/{-} & N &  \\
C0011 & 1a<o0L / 2a<1a1L / 3a<2a0L & N/N/N & +/+/+ & N &  \\
C0012 & 1a<o0L / 2a<1a1L / 3a<2a1 & N/N/N & +/+/{-} & N &  \\
C0013 & 1a<o0L / 2a<1a1L / 3a<2a1L & N/N/N & +/+/+ & N &  \\
C0014 & 1a<o0L / 2a<1a1L / 3a<2a2 & N/N/N & +/+/{-} & N &  \\
C0015 & 1a<o0L / 2a<1a1L / 3a<2a2L & N/N/N & +/+/+ & N &  \\
C0016 & 1a<o0L / 2a<1a2 / 3a<2a0L & N/N/N & +/{-}/{-} & N &  \\
C0017 & 1a<o0L / 2a<1a2 / 3a<2a1 & N/N/N & +/{-}/{-} & N &  \\
C0018 & 1a<o0L / 2a<1a2 / 3a<2a1L & N/N/N & +/{-}/{-} & N &  \\
C0019 & 1a<o0L / 2a<1a2 / 3a<2a2 & N/N/N & +/{-}/{-} & N &  \\
C0020 & 1a<o0L / 2a<1a2 / 3a<2a2L & N/N/N & +/{-}/{-} & N &  \\
C0021 & 1a<o0L / 2a<1a2L / 3a<2a0L & N/N/N & +/+/+ & N &  \\
C0022 & 1a<o0L / 2a<1a2L / 3a<2a1 & N/N/N & +/+/{-} & N &  \\
C0023 & 1a<o0L / 2a<1a2L / 3a<2a1L & N/N/N & +/+/+ & N &  \\
C0024 & 1a<o0L / 2a<1a2L / 3a<2a2 & N/N/N & +/+/{-} & N &  \\
C0025 & 1a<o0L / 2a<1a2L / 3a<2a2L & N/N/N & +/+/+ & N &  \\
C0026 & 1a<o1 / 2a<1a0L / 3a<2a0L & I/N/N & {-}/{-}/{-} & N &  \\
C0027 & 1a<o1 / 2a<1a0L / 3a<2a1 & I/N/N & {-}/{-}/{-} & N &  \\
C0028 & 1a<o1 / 2a<1a0L / 3a<2a1L & I/N/N & {-}/{-}/{-} & N &  \\
C0029 & 1a<o1 / 2a<1a0L / 3a<2a2 & I/N/N & {-}/{-}/{-} & N &  \\
C0030 & 1a<o1 / 2a<1a0L / 3a<2a2L & I/N/N & {-}/{-}/{-} & N &  \\
C0031 & 1a<o1 / 2a<1a1 / 3a<2a0L & I/I/N & {-}/{-}/{-} & N &  \\
C0032 & 1a<o1 / 2a<1a1 / 3a<2a1 & I/I/I & {-}/{-}/{-} & N &  \\
C0033 & 1a<o1 / 2a<1a1 / 3a<2a1L & I/I/I & {-}/{-}/{-} & N &  \\
C0034 & 1a<o1 / 2a<1a1 / 3a<2a2 & I/I/I & {-}/{-}/{-} & N &  \\
C0035 & 1a<o1 / 2a<1a1 / 3a<2a2L & I/I/I & {-}/{-}/{-} & N &  \\
C0036 & 1a<o1 / 2a<1a1L / 3a<2a0L & I/I/N & {-}/{-}/{-} & N &  \\
C0037 & 1a<o1 / 2a<1a1L / 3a<2a1 & I/I/I & {-}/{-}/{-} & N &  \\
C0038 & 1a<o1 / 2a<1a1L / 3a<2a1L & I/I/I & {-}/{-}/{-} & N &  \\
C0039 & 1a<o1 / 2a<1a1L / 3a<2a2 & I/I/I & {-}/{-}/{-} & N &  \\
C0040 & 1a<o1 / 2a<1a1L / 3a<2a2L & I/I/I & {-}/{-}/{-} & N &  \\
C0041 & 1a<o1 / 2a<1a2 / 3a<2a0L & I/I/N & {-}/{-}/{-} & N &  \\
C0042 & 1a<o1 / 2a<1a2 / 3a<2a1 & I/I/I & {-}/{-}/{-} & N &  \\
C0043 & 1a<o1 / 2a<1a2 / 3a<2a1L & I/I/I & {-}/{-}/{-} & N &  \\
C0044 & 1a<o1 / 2a<1a2 / 3a<2a2 & I/I/I & {-}/{-}/{-} & N &  \\
C0045 & 1a<o1 / 2a<1a2 / 3a<2a2L & I/I/I & {-}/{-}/{-} & N &  \\
C0046 & 1a<o1 / 2a<1a2L / 3a<2a0L & I/I/N & {-}/{-}/{-} & N &  \\
C0047 & 1a<o1 / 2a<1a2L / 3a<2a1 & I/I/I & {-}/{-}/{-} & N &  \\
C0048 & 1a<o1 / 2a<1a2L / 3a<2a1L & I/I/I & {-}/{-}/{-} & N &  \\
C0049 & 1a<o1 / 2a<1a2L / 3a<2a2 & I/I/I & {-}/{-}/{-} & N &  \\
C0050 & 1a<o1 / 2a<1a2L / 3a<2a2L & I/I/I & {-}/{-}/{-} & N &  \\
C0051 & 1a<o1L / 2a<1a0L / 3a<2a0L & I/N/N & +/+/+ & N &  \\
C0052 & 1a<o1L / 2a<1a0L / 3a<2a1 & I/N/N & +/+/{-} & N &  \\
C0053 & 1a<o1L / 2a<1a0L / 3a<2a1L & I/N/N & +/+/+ & N &  \\
C0054 & 1a<o1L / 2a<1a0L / 3a<2a2 & I/N/N & +/+/{-} & N &  \\
C0055 & 1a<o1L / 2a<1a0L / 3a<2a2L & I/N/N & +/+/+ & N &  \\
C0056 & 1a<o1L / 2a<1a1 / 3a<2a0L & I/I/N & +/{-}/{-} & N &  \\
C0057 & 1a<o1L / 2a<1a1 / 3a<2a1 & I/I/I & +/{-}/{-} & N &  \\
C0058 & 1a<o1L / 2a<1a1 / 3a<2a1L & I/I/I & +/{-}/{-} & N &  \\
C0059 & 1a<o1L / 2a<1a1 / 3a<2a2 & I/I/I & +/{-}/{-} & N &  \\
C0060 & 1a<o1L / 2a<1a1 / 3a<2a2L & I/I/I & +/{-}/{-} & N &  \\
C0061 & 1a<o1L / 2a<1a1L / 3a<2a0L & I/I/N & +/+/+ & N &  \\
C0062 & 1a<o1L / 2a<1a1L / 3a<2a1 & I/I/I & +/+/{-} & N &  \\
C0063 & 1a<o1L / 2a<1a1L / 3a<2a1L & I/I/I & +/+/+ & N &  \\
C0064 & 1a<o1L / 2a<1a1L / 3a<2a2 & I/I/I & +/+/{-} & N &  \\
C0065 & 1a<o1L / 2a<1a1L / 3a<2a2L & I/I/I & +/+/+ & N &  \\
C0066 & 1a<o1L / 2a<1a2 / 3a<2a0L & I/I/N & +/{-}/{-} & N &  \\
C0067 & 1a<o1L / 2a<1a2 / 3a<2a1 & I/I/I & +/{-}/{-} & N &  \\
C0068 & 1a<o1L / 2a<1a2 / 3a<2a1L & I/I/I & +/{-}/{-} & N &  \\
C0069 & 1a<o1L / 2a<1a2 / 3a<2a2 & I/I/I & +/{-}/{-} & N &  \\
C0070 & 1a<o1L / 2a<1a2 / 3a<2a2L & I/I/I & +/{-}/{-} & N &  \\
C0071 & 1a<o1L / 2a<1a2L / 3a<2a0L & I/I/N & +/+/+ & N &  \\
C0072 & 1a<o1L / 2a<1a2L / 3a<2a1 & I/I/I & +/+/{-} & N &  \\
C0073 & 1a<o1L / 2a<1a2L / 3a<2a1L & I/I/I & +/+/+ & N &  \\
C0074 & 1a<o1L / 2a<1a2L / 3a<2a2 & I/I/I & +/+/{-} & N &  \\
C0075 & 1a<o1L / 2a<1a2L / 3a<2a2L & I/I/I & +/+/+ & N &  \\
C0076 & 1a<o2 / 2a<1a0L / 3a<2a0L & I/N/N & {-}/{-}/{-} & N &  \\
C0077 & 1a<o2 / 2a<1a0L / 3a<2a1 & I/N/N & {-}/{-}/{-} & N &  \\
C0078 & 1a<o2 / 2a<1a0L / 3a<2a1L & I/N/N & {-}/{-}/{-} & N &  \\
C0079 & 1a<o2 / 2a<1a0L / 3a<2a2 & I/N/N & {-}/{-}/{-} & N &  \\
C0080 & 1a<o2 / 2a<1a0L / 3a<2a2L & I/N/N & {-}/{-}/{-} & N &  \\
C0081 & 1a<o2 / 2a<1a1 / 3a<2a0L & I/I/N & {-}/{-}/{-} & N &  \\
C0082 & 1a<o2 / 2a<1a1 / 3a<2a1 & I/I/I & {-}/{-}/{-} & N &  \\
C0083 & 1a<o2 / 2a<1a1 / 3a<2a1L & I/I/I & {-}/{-}/{-} & N &  \\
C0084 & 1a<o2 / 2a<1a1 / 3a<2a2 & I/I/I & {-}/{-}/{-} & N &  \\
C0085 & 1a<o2 / 2a<1a1 / 3a<2a2L & I/I/I & {-}/{-}/{-} & N &  \\
C0086 & 1a<o2 / 2a<1a1L / 3a<2a0L & I/I/N & {-}/{-}/{-} & N &  \\
C0087 & 1a<o2 / 2a<1a1L / 3a<2a1 & I/I/I & {-}/{-}/{-} & N &  \\
C0088 & 1a<o2 / 2a<1a1L / 3a<2a1L & I/I/I & {-}/{-}/{-} & N &  \\
C0089 & 1a<o2 / 2a<1a1L / 3a<2a2 & I/I/I & {-}/{-}/{-} & N &  \\
C0090 & 1a<o2 / 2a<1a1L / 3a<2a2L & I/I/I & {-}/{-}/{-} & N &  \\
C0091 & 1a<o2 / 2a<1a2 / 3a<2a0L & I/I/N & {-}/{-}/{-} & N &  \\
C0092 & 1a<o2 / 2a<1a2 / 3a<2a1 & I/I/I & {-}/{-}/{-} & N &  \\
C0093 & 1a<o2 / 2a<1a2 / 3a<2a1L & I/I/I & {-}/{-}/{-} & N &  \\
C0094 & 1a<o2 / 2a<1a2 / 3a<2a2 & I/I/I & {-}/{-}/{-} & N &  \\
C0095 & 1a<o2 / 2a<1a2 / 3a<2a2L & I/I/I & {-}/{-}/{-} & N &  \\
C0096 & 1a<o2 / 2a<1a2L / 3a<2a0L & I/I/N & {-}/{-}/{-} & N &  \\
C0097 & 1a<o2 / 2a<1a2L / 3a<2a1 & I/I/I & {-}/{-}/{-} & N &  \\
C0098 & 1a<o2 / 2a<1a2L / 3a<2a1L & I/I/I & {-}/{-}/{-} & N &  \\
C0099 & 1a<o2 / 2a<1a2L / 3a<2a2 & I/I/I & {-}/{-}/{-} & N &  \\
C0100 & 1a<o2 / 2a<1a2L / 3a<2a2L & I/I/I & {-}/{-}/{-} & N &  \\
C0101 & 1a<o2L / 2a<1a0L / 3a<2a0L & I/N/N & +/+/+ & N &  \\
C0102 & 1a<o2L / 2a<1a0L / 3a<2a1 & I/N/N & +/+/{-} & N &  \\
C0103 & 1a<o2L / 2a<1a0L / 3a<2a1L & I/N/N & +/+/+ & N &  \\
C0104 & 1a<o2L / 2a<1a0L / 3a<2a2 & I/N/N & +/+/{-} & N &  \\
C0105 & 1a<o2L / 2a<1a0L / 3a<2a2L & I/N/N & +/+/+ & N &  \\
C0106 & 1a<o2L / 2a<1a1 / 3a<2a0L & I/I/N & +/{-}/{-} & N &  \\
C0107 & 1a<o2L / 2a<1a1 / 3a<2a1 & I/I/I & +/{-}/{-} & N &  \\
C0108 & 1a<o2L / 2a<1a1 / 3a<2a1L & I/I/I & +/{-}/{-} & N &  \\
C0109 & 1a<o2L / 2a<1a1 / 3a<2a2 & I/I/I & +/{-}/{-} & N &  \\
C0110 & 1a<o2L / 2a<1a1 / 3a<2a2L & I/I/I & +/{-}/{-} & N &  \\
C0111 & 1a<o2L / 2a<1a1L / 3a<2a0L & I/I/N & +/+/+ & N &  \\
C0112 & 1a<o2L / 2a<1a1L / 3a<2a1 & I/I/I & +/+/{-} & N &  \\
C0113 & 1a<o2L / 2a<1a1L / 3a<2a1L & I/I/I & +/+/+ & N &  \\
C0114 & 1a<o2L / 2a<1a1L / 3a<2a2 & I/I/I & +/+/{-} & N &  \\
C0115 & 1a<o2L / 2a<1a1L / 3a<2a2L & I/I/I & +/+/+ & N &  \\
C0116 & 1a<o2L / 2a<1a2 / 3a<2a0L & I/I/N & +/{-}/{-} & N &  \\
C0117 & 1a<o2L / 2a<1a2 / 3a<2a1 & I/I/I & +/{-}/{-} & N &  \\
C0118 & 1a<o2L / 2a<1a2 / 3a<2a1L & I/I/I & +/{-}/{-} & N &  \\
C0119 & 1a<o2L / 2a<1a2 / 3a<2a2 & I/I/I & +/{-}/{-} & N &  \\
C0120 & 1a<o2L / 2a<1a2 / 3a<2a2L & I/I/I & +/{-}/{-} & N &  \\
C0121 & 1a<o2L / 2a<1a2L / 3a<2a0L & I/I/N & +/+/+ & N &  \\
C0122 & 1a<o2L / 2a<1a2L / 3a<2a1 & I/I/I & +/+/{-} & N &  \\
C0123 & 1a<o2L / 2a<1a2L / 3a<2a1L & I/I/I & +/+/+ & N &  \\
C0124 & 1a<o2L / 2a<1a2L / 3a<2a2 & I/I/I & +/+/{-} & N &  \\
C0125 & 1a<o2L / 2a<1a2L / 3a<2a2L & I/I/I & +/+/+ & N &  \\
C0126 & 1a<o0L / 2a<1a0L / 3a*<2a1 & N/N/N & +/+/{-} & N &  \\
C0127 & 1a<o0L / 2a<1a0L / 3a*<2a2 & N/N/N & +/+/{-} & N &  \\
C0128 & 1a<o0L / 2a<1a1 / 3a*<2a1 & N/N/N & +/{-}/{-} & N &  \\
C0129 & 1a<o0L / 2a<1a1 / 3a*<2a2 & N/N/N & +/{-}/{-} & N &  \\
C0130 & 1a<o0L / 2a<1a1L / 3a*<2a1 & N/N/N & +/+/{-} & N &  \\
C0131 & 1a<o0L / 2a<1a1L / 3a*<2a2 & N/N/N & +/+/{-} & N &  \\
C0132 & 1a<o0L / 2a<1a2 / 3a*<2a1 & N/N/N & +/{-}/{-} & N &  \\
C0133 & 1a<o0L / 2a<1a2 / 3a*<2a2 & N/N/N & +/{-}/{-} & N &  \\
C0134 & 1a<o0L / 2a<1a2L / 3a*<2a1 & N/N/N & +/+/{-} & N &  \\
C0135 & 1a<o0L / 2a<1a2L / 3a*<2a2 & N/N/N & +/+/{-} & N &  \\
C0136 & 1a<o0L / 2a*<1a1 / 3a<2a1 & N/N/N & +/{-}/{-} & N &  \\
C0137 & 1a<o0L / 2a*<1a1 / 3a<2a2 & N/N/N & +/{-}/{-} & N &  \\
C0138 & 1a<o0L / 2a*<1a2 / 3a<2a1 & N/N/N & +/{-}/{-} & N &  \\
C0139 & 1a<o0L / 2a*<1a2 / 3a<2a2 & N/N/N & +/{-}/{-} & N &  \\
C0140 & 1a<o1 / 2a<1a0L / 3a*<2a1 & I/N/N & {-}/{-}/{-} & N &  \\
C0141 & 1a<o1 / 2a<1a0L / 3a*<2a2 & I/N/N & {-}/{-}/{-} & N &  \\
C0142 & 1a<o1 / 2a<1a1 / 3a*<2a1 & I/I/U & {-}/{-}/{-} & N & W:3a \\
C0143 & 1a<o1 / 2a<1a1 / 3a*<2a2 & I/I/U & {-}/{-}/{-} & N & W:3a \\
C0144 & 1a<o1 / 2a<1a1L / 3a*<2a1 & I/I/U & {-}/{-}/{-} & N & W:3a \\
C0145 & 1a<o1 / 2a<1a1L / 3a*<2a2 & I/I/U & {-}/{-}/{-} & N & W:3a \\
C0146 & 1a<o1 / 2a<1a2 / 3a*<2a1 & I/I/U & {-}/{-}/{-} & N & W:3a \\
C0147 & 1a<o1 / 2a<1a2 / 3a*<2a2 & I/I/U & {-}/{-}/{-} & N & W:3a \\
C0148 & 1a<o1 / 2a<1a2L / 3a*<2a1 & I/I/U & {-}/{-}/{-} & N & W:3a \\
C0149 & 1a<o1 / 2a<1a2L / 3a*<2a2 & I/I/U & {-}/{-}/{-} & N & W:3a \\
C0150 & 1a<o1 / 2a*<1a1 / 3a<2a1 & I/U/U & {-}/{-}/{-} & N & W:2a,3a \\
C0151 & 1a<o1 / 2a*<1a1 / 3a<2a2 & I/U/U & {-}/{-}/{-} & N & W:2a,3a \\
C0152 & 1a<o1 / 2a*<1a2 / 3a<2a1 & I/U/U & {-}/{-}/{-} & N & W:2a,3a \\
C0153 & 1a<o1 / 2a*<1a2 / 3a<2a2 & I/U/U & {-}/{-}/{-} & N & W:2a,3a \\
C0154 & 1a<o1L / 2a<1a0L / 3a*<2a1 & I/N/N & +/+/{-} & N &  \\
C0155 & 1a<o1L / 2a<1a0L / 3a*<2a2 & I/N/N & +/+/{-} & N &  \\
C0156 & 1a<o1L / 2a<1a1 / 3a*<2a1 & I/I/U & +/{-}/{-} & N & W:3a \\
C0157 & 1a<o1L / 2a<1a1 / 3a*<2a2 & I/I/U & +/{-}/{-} & N & W:3a \\
C0158 & 1a<o1L / 2a<1a1L / 3a*<2a1 & I/I/U & +/+/{-} & N & W:3a \\
C0159 & 1a<o1L / 2a<1a1L / 3a*<2a2 & I/I/U & +/+/{-} & N & W:3a \\
C0160 & 1a<o1L / 2a<1a2 / 3a*<2a1 & I/I/U & +/{-}/{-} & N & W:3a \\
C0161 & 1a<o1L / 2a<1a2 / 3a*<2a2 & I/I/U & +/{-}/{-} & N & W:3a \\
C0162 & 1a<o1L / 2a<1a2L / 3a*<2a1 & I/I/U & +/+/{-} & N & W:3a \\
C0163 & 1a<o1L / 2a<1a2L / 3a*<2a2 & I/I/U & +/+/{-} & N & W:3a \\
C0164 & 1a<o1L / 2a*<1a1 / 3a<2a1 & I/U/U & +/{-}/{-} & N & W:2a,3a \\
C0165 & 1a<o1L / 2a*<1a1 / 3a<2a2 & I/U/U & +/{-}/{-} & N & W:2a,3a \\
C0166 & 1a<o1L / 2a*<1a2 / 3a<2a1 & I/U/U & +/{-}/{-} & N & W:2a,3a \\
C0167 & 1a<o1L / 2a*<1a2 / 3a<2a2 & I/U/U & +/{-}/{-} & N & W:2a,3a \\
C0168 & 1a<o2 / 2a<1a0L / 3a*<2a1 & I/N/N & {-}/{-}/{-} & N &  \\
C0169 & 1a<o2 / 2a<1a0L / 3a*<2a2 & I/N/N & {-}/{-}/{-} & N &  \\
C0170 & 1a<o2 / 2a<1a1 / 3a*<2a1 & I/I/U & {-}/{-}/{-} & N & W:3a \\
C0171 & 1a<o2 / 2a<1a1 / 3a*<2a2 & I/I/U & {-}/{-}/{-} & N & W:3a \\
C0172 & 1a<o2 / 2a<1a1L / 3a*<2a1 & I/I/U & {-}/{-}/{-} & N & W:3a \\
C0173 & 1a<o2 / 2a<1a1L / 3a*<2a2 & I/I/U & {-}/{-}/{-} & N & W:3a \\
C0174 & 1a<o2 / 2a<1a2 / 3a*<2a1 & I/I/U & {-}/{-}/{-} & N & W:3a \\
C0175 & 1a<o2 / 2a<1a2 / 3a*<2a2 & I/I/U & {-}/{-}/{-} & N & W:3a \\
C0176 & 1a<o2 / 2a<1a2L / 3a*<2a1 & I/I/U & {-}/{-}/{-} & N & W:3a \\
C0177 & 1a<o2 / 2a<1a2L / 3a*<2a2 & I/I/U & {-}/{-}/{-} & N & W:3a \\
C0178 & 1a<o2 / 2a*<1a1 / 3a<2a1 & I/U/U & {-}/{-}/{-} & N & W:2a,3a \\
C0179 & 1a<o2 / 2a*<1a1 / 3a<2a2 & I/U/U & {-}/{-}/{-} & N & W:2a,3a \\
C0180 & 1a<o2 / 2a*<1a2 / 3a<2a1 & I/U/U & {-}/{-}/{-} & N & W:2a,3a \\
C0181 & 1a<o2 / 2a*<1a2 / 3a<2a2 & I/U/U & {-}/{-}/{-} & N & W:2a,3a \\
C0182 & 1a<o2L / 2a<1a0L / 3a*<2a1 & I/N/N & +/+/{-} & N &  \\
C0183 & 1a<o2L / 2a<1a0L / 3a*<2a2 & I/N/N & +/+/{-} & N &  \\
C0184 & 1a<o2L / 2a<1a1 / 3a*<2a1 & I/I/U & +/{-}/{-} & N & W:3a \\
C0185 & 1a<o2L / 2a<1a1 / 3a*<2a2 & I/I/U & +/{-}/{-} & N & W:3a \\
C0186 & 1a<o2L / 2a<1a1L / 3a*<2a1 & I/I/U & +/+/{-} & N & W:3a \\
C0187 & 1a<o2L / 2a<1a1L / 3a*<2a2 & I/I/U & +/+/{-} & N & W:3a \\
C0188 & 1a<o2L / 2a<1a2 / 3a*<2a1 & I/I/U & +/{-}/{-} & N & W:3a \\
C0189 & 1a<o2L / 2a<1a2 / 3a*<2a2 & I/I/U & +/{-}/{-} & N & W:3a \\
C0190 & 1a<o2L / 2a<1a2L / 3a*<2a1 & I/I/U & +/+/{-} & N & W:3a \\
C0191 & 1a<o2L / 2a<1a2L / 3a*<2a2 & I/I/U & +/+/{-} & N & W:3a \\
C0192 & 1a<o2L / 2a*<1a1 / 3a<2a1 & I/U/U & +/{-}/{-} & N & W:2a,3a \\
C0193 & 1a<o2L / 2a*<1a1 / 3a<2a2 & I/U/U & +/{-}/{-} & N & W:2a,3a \\
C0194 & 1a<o2L / 2a*<1a2 / 3a<2a1 & I/U/U & +/{-}/{-} & N & W:2a,3a \\
C0195 & 1a<o2L / 2a*<1a2 / 3a<2a2 & I/U/U & +/{-}/{-} & N & W:2a,3a \\
C0196 & 1a*<o1 / 2a<1a1 / 3a<2a0L & U/U/N & {-}/{-}/{-} & N & W:1a,2a \\
C0197 & 1a*<o1 / 2a<1a1 / 3a<2a1 & U/U/U & {-}/{-}/{-} & N & W:1a,2a,3a \\
C0198 & 1a*<o1 / 2a<1a1 / 3a<2a1L & U/U/U & {-}/{-}/{-} & N & W:1a,2a,3a \\
C0199 & 1a*<o1 / 2a<1a1 / 3a<2a2 & U/U/U & {-}/{-}/{-} & N & W:1a,2a,3a \\
C0200 & 1a*<o1 / 2a<1a1 / 3a<2a2L & U/U/U & {-}/{-}/{-} & N & W:1a,2a,3a \\
C0201 & 1a*<o1 / 2a<1a2 / 3a<2a0L & U/U/N & {-}/{-}/{-} & N & W:1a,2a \\
C0202 & 1a*<o1 / 2a<1a2 / 3a<2a1 & U/U/U & {-}/{-}/{-} & N & W:1a,2a,3a \\
C0203 & 1a*<o1 / 2a<1a2 / 3a<2a1L & U/U/U & {-}/{-}/{-} & N & W:1a,2a,3a \\
C0204 & 1a*<o1 / 2a<1a2 / 3a<2a2 & U/U/U & {-}/{-}/{-} & N & W:1a,2a,3a \\
C0205 & 1a*<o1 / 2a<1a2 / 3a<2a2L & U/U/U & {-}/{-}/{-} & N & W:1a,2a,3a \\
C0206 & 1a*<o2 / 2a<1a1 / 3a<2a0L & U/U/N & {-}/{-}/{-} & N & W:1a,2a \\
C0207 & 1a*<o2 / 2a<1a1 / 3a<2a1 & U/U/U & {-}/{-}/{-} & N & W:1a,2a,3a \\
C0208 & 1a*<o2 / 2a<1a1 / 3a<2a1L & U/U/U & {-}/{-}/{-} & N & W:1a,2a,3a \\
C0209 & 1a*<o2 / 2a<1a1 / 3a<2a2 & U/U/U & {-}/{-}/{-} & N & W:1a,2a,3a \\
C0210 & 1a*<o2 / 2a<1a1 / 3a<2a2L & U/U/U & {-}/{-}/{-} & N & W:1a,2a,3a \\
C0211 & 1a*<o2 / 2a<1a2 / 3a<2a0L & U/U/N & {-}/{-}/{-} & N & W:1a,2a \\
C0212 & 1a*<o2 / 2a<1a2 / 3a<2a1 & U/U/U & {-}/{-}/{-} & N & W:1a,2a,3a \\
C0213 & 1a*<o2 / 2a<1a2 / 3a<2a1L & U/U/U & {-}/{-}/{-} & N & W:1a,2a,3a \\
C0214 & 1a*<o2 / 2a<1a2 / 3a<2a2 & U/U/U & {-}/{-}/{-} & N & W:1a,2a,3a \\
C0215 & 1a*<o2 / 2a<1a2 / 3a<2a2L & U/U/U & {-}/{-}/{-} & N & W:1a,2a,3a \\
\end{casetable}
\subsection*{Height 3, 4 stages after the original (5331 cases)}
\begin{casetable}
C0216 & 1a<o0L / 2a<1a0L / 3a<2a0L 3b<2a1 & N/N/NN & +/+/+{-} & N &  \\
C0217 & 1a<o0L / 2a<1a0L / 3a<2a0L 3b<2a2 & N/N/NN & +/+/+{-} & N &  \\
C0218 & 1a<o0L / 2a<1a0L / 3a<2a1 3b<2a1 & N/N/NN & +/+/{-}{-} & N &  \\
C0219 & 1a<o0L / 2a<1a0L / 3a<2a1 3b<2a1L & N/N/NN & +/+/{-}+ & N &  \\
C0220 & 1a<o0L / 2a<1a0L / 3a<2a1 3b<2a2 & N/N/NN & +/+/{-}{-} & N &  \\
C0221 & 1a<o0L / 2a<1a0L / 3a<2a1 3b<2a2L & N/N/NN & +/+/{-}+ & N &  \\
C0222 & 1a<o0L / 2a<1a0L / 3a<2a1L 3b<2a2 & N/N/NN & +/+/+{-} & N &  \\
C0223 & 1a<o0L / 2a<1a0L / 3a<2a2 3b<2a2 & N/N/NN & +/+/{-}{-} & N &  \\
C0224 & 1a<o0L / 2a<1a0L / 3a<2a2 3b<2a2L & N/N/NN & +/+/{-}+ & N &  \\
C0225 & 1a<o0L / 2a<1a0L 2b<1a1 / 3a<2a0L & N/NN/N & +/+{-}/+ & N &  \\
C0226 & 1a<o0L / 2a<1a0L 2b<1a1 / 3a<2a1 & N/NN/N & +/+{-}/{-} & N &  \\
C0227 & 1a<o0L / 2a<1a0L 2b<1a1 / 3a<2a1L & N/NN/N & +/+{-}/+ & N &  \\
C0228 & 1a<o0L / 2a<1a0L 2b<1a1 / 3a<2a2 & N/NN/N & +/+{-}/{-} & N &  \\
C0229 & 1a<o0L / 2a<1a0L 2b<1a1 / 3a<2a2L & N/NN/N & +/+{-}/+ & N &  \\
C0230 & 1a<o0L / 2a<1a0L 2b<1a1 / 3a<2b0L & N/NN/N & +/+{-}/{-} & N &  \\
C0231 & 1a<o0L / 2a<1a0L 2b<1a1 / 3a<2b1 & N/NN/N & +/+{-}/{-} & N &  \\
C0232 & 1a<o0L / 2a<1a0L 2b<1a1 / 3a<2b1L & N/NN/N & +/+{-}/{-} & N &  \\
C0233 & 1a<o0L / 2a<1a0L 2b<1a1 / 3a<2b2 & N/NN/N & +/+{-}/{-} & N &  \\
C0234 & 1a<o0L / 2a<1a0L 2b<1a1 / 3a<2b2L & N/NN/N & +/+{-}/{-} & N &  \\
C0235 & 1a<o0L / 2a<1a0L 2b<1a2 / 3a<2a0L & N/NN/N & +/+{-}/+ & N &  \\
C0236 & 1a<o0L / 2a<1a0L 2b<1a2 / 3a<2a1 & N/NN/N & +/+{-}/{-} & N &  \\
C0237 & 1a<o0L / 2a<1a0L 2b<1a2 / 3a<2a1L & N/NN/N & +/+{-}/+ & N &  \\
C0238 & 1a<o0L / 2a<1a0L 2b<1a2 / 3a<2a2 & N/NN/N & +/+{-}/{-} & N &  \\
C0239 & 1a<o0L / 2a<1a0L 2b<1a2 / 3a<2a2L & N/NN/N & +/+{-}/+ & N &  \\
C0240 & 1a<o0L / 2a<1a0L 2b<1a2 / 3a<2b0L & N/NN/N & +/+{-}/{-} & N &  \\
C0241 & 1a<o0L / 2a<1a0L 2b<1a2 / 3a<2b1 & N/NN/N & +/+{-}/{-} & N &  \\
C0242 & 1a<o0L / 2a<1a0L 2b<1a2 / 3a<2b1L & N/NN/N & +/+{-}/{-} & N &  \\
C0243 & 1a<o0L / 2a<1a0L 2b<1a2 / 3a<2b2 & N/NN/N & +/+{-}/{-} & N &  \\
C0244 & 1a<o0L / 2a<1a0L 2b<1a2 / 3a<2b2L & N/NN/N & +/+{-}/{-} & N &  \\
C0245 & 1a<o0L / 2a<1a1 / 3a<2a0L 3b<2a1 & N/N/NN & +/{-}/{-}{-} & N &  \\
C0246 & 1a<o0L / 2a<1a1 / 3a<2a0L 3b<2a2 & N/N/NN & +/{-}/{-}{-} & N &  \\
C0247 & 1a<o0L / 2a<1a1 / 3a<2a1 3b<2a1 & N/N/NN & +/{-}/{-}{-} & N &  \\
C0248 & 1a<o0L / 2a<1a1 / 3a<2a1 3b<2a1L & N/N/NN & +/{-}/{-}{-} & N &  \\
C0249 & 1a<o0L / 2a<1a1 / 3a<2a1 3b<2a2 & N/N/NN & +/{-}/{-}{-} & N &  \\
C0250 & 1a<o0L / 2a<1a1 / 3a<2a1 3b<2a2L & N/N/NN & +/{-}/{-}{-} & N &  \\
C0251 & 1a<o0L / 2a<1a1 / 3a<2a1L 3b<2a2 & N/N/NN & +/{-}/{-}{-} & N &  \\
C0252 & 1a<o0L / 2a<1a1 / 3a<2a2 3b<2a2 & N/N/NN & +/{-}/{-}{-} & N &  \\
C0253 & 1a<o0L / 2a<1a1 / 3a<2a2 3b<2a2L & N/N/NN & +/{-}/{-}{-} & N &  \\
C0254 & 1a<o0L / 2a<1a1 2b<1a1 / 3a<2a0L & N/NN/N & +/{-}{-}/{-} & N &  \\
C0255 & 1a<o0L / 2a<1a1 2b<1a1 / 3a<2a1 & N/NN/N & +/{-}{-}/{-} & N &  \\
C0256 & 1a<o0L / 2a<1a1 2b<1a1 / 3a<2a1L & N/NN/N & +/{-}{-}/{-} & N &  \\
C0257 & 1a<o0L / 2a<1a1 2b<1a1 / 3a<2a2 & N/NN/N & +/{-}{-}/{-} & N &  \\
C0258 & 1a<o0L / 2a<1a1 2b<1a1 / 3a<2a2L & N/NN/N & +/{-}{-}/{-} & N &  \\
C0259 & 1a<o0L / 2a<1a1 2b<1a1L / 3a<2a0L & N/NN/N & +/{-}+/{-} & N &  \\
C0260 & 1a<o0L / 2a<1a1 2b<1a1L / 3a<2a1 & N/NN/N & +/{-}+/{-} & N &  \\
C0261 & 1a<o0L / 2a<1a1 2b<1a1L / 3a<2a1L & N/NN/N & +/{-}+/{-} & N &  \\
C0262 & 1a<o0L / 2a<1a1 2b<1a1L / 3a<2a2 & N/NN/N & +/{-}+/{-} & N &  \\
C0263 & 1a<o0L / 2a<1a1 2b<1a1L / 3a<2a2L & N/NN/N & +/{-}+/{-} & N &  \\
C0264 & 1a<o0L / 2a<1a1 2b<1a1L / 3a<2b0L & N/NN/N & +/{-}+/+ & N &  \\
C0265 & 1a<o0L / 2a<1a1 2b<1a1L / 3a<2b1 & N/NN/N & +/{-}+/{-} & N &  \\
C0266 & 1a<o0L / 2a<1a1 2b<1a1L / 3a<2b1L & N/NN/N & +/{-}+/+ & N &  \\
C0267 & 1a<o0L / 2a<1a1 2b<1a1L / 3a<2b2 & N/NN/N & +/{-}+/{-} & N &  \\
C0268 & 1a<o0L / 2a<1a1 2b<1a1L / 3a<2b2L & N/NN/N & +/{-}+/+ & N &  \\
C0269 & 1a<o0L / 2a<1a1 2b<1a2 / 3a<2a0L & N/NN/N & +/{-}{-}/{-} & N &  \\
C0270 & 1a<o0L / 2a<1a1 2b<1a2 / 3a<2a1 & N/NN/N & +/{-}{-}/{-} & N &  \\
C0271 & 1a<o0L / 2a<1a1 2b<1a2 / 3a<2a1L & N/NN/N & +/{-}{-}/{-} & N &  \\
C0272 & 1a<o0L / 2a<1a1 2b<1a2 / 3a<2a2 & N/NN/N & +/{-}{-}/{-} & N &  \\
C0273 & 1a<o0L / 2a<1a1 2b<1a2 / 3a<2a2L & N/NN/N & +/{-}{-}/{-} & N &  \\
C0274 & 1a<o0L / 2a<1a1 2b<1a2 / 3a<2b0L & N/NN/N & +/{-}{-}/{-} & N &  \\
C0275 & 1a<o0L / 2a<1a1 2b<1a2 / 3a<2b1 & N/NN/N & +/{-}{-}/{-} & N &  \\
C0276 & 1a<o0L / 2a<1a1 2b<1a2 / 3a<2b1L & N/NN/N & +/{-}{-}/{-} & N &  \\
C0277 & 1a<o0L / 2a<1a1 2b<1a2 / 3a<2b2 & N/NN/N & +/{-}{-}/{-} & N &  \\
C0278 & 1a<o0L / 2a<1a1 2b<1a2 / 3a<2b2L & N/NN/N & +/{-}{-}/{-} & N &  \\
C0279 & 1a<o0L / 2a<1a1 2b<1a2L / 3a<2a0L & N/NN/N & +/{-}+/{-} & N &  \\
C0280 & 1a<o0L / 2a<1a1 2b<1a2L / 3a<2a1 & N/NN/N & +/{-}+/{-} & N &  \\
C0281 & 1a<o0L / 2a<1a1 2b<1a2L / 3a<2a1L & N/NN/N & +/{-}+/{-} & N &  \\
C0282 & 1a<o0L / 2a<1a1 2b<1a2L / 3a<2a2 & N/NN/N & +/{-}+/{-} & N &  \\
C0283 & 1a<o0L / 2a<1a1 2b<1a2L / 3a<2a2L & N/NN/N & +/{-}+/{-} & N &  \\
C0284 & 1a<o0L / 2a<1a1 2b<1a2L / 3a<2b0L & N/NN/N & +/{-}+/+ & N &  \\
C0285 & 1a<o0L / 2a<1a1 2b<1a2L / 3a<2b1 & N/NN/N & +/{-}+/{-} & N &  \\
C0286 & 1a<o0L / 2a<1a1 2b<1a2L / 3a<2b1L & N/NN/N & +/{-}+/+ & N &  \\
C0287 & 1a<o0L / 2a<1a1 2b<1a2L / 3a<2b2 & N/NN/N & +/{-}+/{-} & N &  \\
C0288 & 1a<o0L / 2a<1a1 2b<1a2L / 3a<2b2L & N/NN/N & +/{-}+/+ & N &  \\
C0289 & 1a<o0L / 2a<1a1L / 3a<2a0L 3b<2a1 & N/N/NN & +/+/+{-} & N &  \\
C0290 & 1a<o0L / 2a<1a1L / 3a<2a0L 3b<2a2 & N/N/NN & +/+/+{-} & N &  \\
C0291 & 1a<o0L / 2a<1a1L / 3a<2a1 3b<2a1 & N/N/NN & +/+/{-}{-} & N &  \\
C0292 & 1a<o0L / 2a<1a1L / 3a<2a1 3b<2a1L & N/N/NN & +/+/{-}+ & N &  \\
C0293 & 1a<o0L / 2a<1a1L / 3a<2a1 3b<2a2 & N/N/NN & +/+/{-}{-} & N &  \\
C0294 & 1a<o0L / 2a<1a1L / 3a<2a1 3b<2a2L & N/N/NN & +/+/{-}+ & N &  \\
C0295 & 1a<o0L / 2a<1a1L / 3a<2a1L 3b<2a2 & N/N/NN & +/+/+{-} & N &  \\
C0296 & 1a<o0L / 2a<1a1L / 3a<2a2 3b<2a2 & N/N/NN & +/+/{-}{-} & N &  \\
C0297 & 1a<o0L / 2a<1a1L / 3a<2a2 3b<2a2L & N/N/NN & +/+/{-}+ & N &  \\
C0298 & 1a<o0L / 2a<1a1L 2b<1a2 / 3a<2a0L & N/NN/N & +/+{-}/+ & N &  \\
C0299 & 1a<o0L / 2a<1a1L 2b<1a2 / 3a<2a1 & N/NN/N & +/+{-}/{-} & N &  \\
C0300 & 1a<o0L / 2a<1a1L 2b<1a2 / 3a<2a1L & N/NN/N & +/+{-}/+ & N &  \\
C0301 & 1a<o0L / 2a<1a1L 2b<1a2 / 3a<2a2 & N/NN/N & +/+{-}/{-} & N &  \\
C0302 & 1a<o0L / 2a<1a1L 2b<1a2 / 3a<2a2L & N/NN/N & +/+{-}/+ & N &  \\
C0303 & 1a<o0L / 2a<1a1L 2b<1a2 / 3a<2b0L & N/NN/N & +/+{-}/{-} & N &  \\
C0304 & 1a<o0L / 2a<1a1L 2b<1a2 / 3a<2b1 & N/NN/N & +/+{-}/{-} & N &  \\
C0305 & 1a<o0L / 2a<1a1L 2b<1a2 / 3a<2b1L & N/NN/N & +/+{-}/{-} & N &  \\
C0306 & 1a<o0L / 2a<1a1L 2b<1a2 / 3a<2b2 & N/NN/N & +/+{-}/{-} & N &  \\
C0307 & 1a<o0L / 2a<1a1L 2b<1a2 / 3a<2b2L & N/NN/N & +/+{-}/{-} & N &  \\
C0308 & 1a<o0L / 2a<1a2 / 3a<2a0L 3b<2a1 & N/N/NN & +/{-}/{-}{-} & N &  \\
C0309 & 1a<o0L / 2a<1a2 / 3a<2a0L 3b<2a2 & N/N/NN & +/{-}/{-}{-} & N &  \\
C0310 & 1a<o0L / 2a<1a2 / 3a<2a1 3b<2a1 & N/N/NN & +/{-}/{-}{-} & N &  \\
C0311 & 1a<o0L / 2a<1a2 / 3a<2a1 3b<2a1L & N/N/NN & +/{-}/{-}{-} & N &  \\
C0312 & 1a<o0L / 2a<1a2 / 3a<2a1 3b<2a2 & N/N/NN & +/{-}/{-}{-} & N &  \\
C0313 & 1a<o0L / 2a<1a2 / 3a<2a1 3b<2a2L & N/N/NN & +/{-}/{-}{-} & N &  \\
C0314 & 1a<o0L / 2a<1a2 / 3a<2a1L 3b<2a2 & N/N/NN & +/{-}/{-}{-} & N &  \\
C0315 & 1a<o0L / 2a<1a2 / 3a<2a2 3b<2a2 & N/N/NN & +/{-}/{-}{-} & N &  \\
C0316 & 1a<o0L / 2a<1a2 / 3a<2a2 3b<2a2L & N/N/NN & +/{-}/{-}{-} & N &  \\
C0317 & 1a<o0L / 2a<1a2 2b<1a2 / 3a<2a0L & N/NN/N & +/{-}{-}/{-} & N &  \\
C0318 & 1a<o0L / 2a<1a2 2b<1a2 / 3a<2a1 & N/NN/N & +/{-}{-}/{-} & N &  \\
C0319 & 1a<o0L / 2a<1a2 2b<1a2 / 3a<2a1L & N/NN/N & +/{-}{-}/{-} & N &  \\
C0320 & 1a<o0L / 2a<1a2 2b<1a2 / 3a<2a2 & N/NN/N & +/{-}{-}/{-} & N &  \\
C0321 & 1a<o0L / 2a<1a2 2b<1a2 / 3a<2a2L & N/NN/N & +/{-}{-}/{-} & N &  \\
C0322 & 1a<o0L / 2a<1a2 2b<1a2L / 3a<2a0L & N/NN/N & +/{-}+/{-} & N &  \\
C0323 & 1a<o0L / 2a<1a2 2b<1a2L / 3a<2a1 & N/NN/N & +/{-}+/{-} & N &  \\
C0324 & 1a<o0L / 2a<1a2 2b<1a2L / 3a<2a1L & N/NN/N & +/{-}+/{-} & N &  \\
C0325 & 1a<o0L / 2a<1a2 2b<1a2L / 3a<2a2 & N/NN/N & +/{-}+/{-} & N &  \\
C0326 & 1a<o0L / 2a<1a2 2b<1a2L / 3a<2a2L & N/NN/N & +/{-}+/{-} & N &  \\
C0327 & 1a<o0L / 2a<1a2 2b<1a2L / 3a<2b0L & N/NN/N & +/{-}+/+ & N &  \\
C0328 & 1a<o0L / 2a<1a2 2b<1a2L / 3a<2b1 & N/NN/N & +/{-}+/{-} & N &  \\
C0329 & 1a<o0L / 2a<1a2 2b<1a2L / 3a<2b1L & N/NN/N & +/{-}+/+ & N &  \\
C0330 & 1a<o0L / 2a<1a2 2b<1a2L / 3a<2b2 & N/NN/N & +/{-}+/{-} & N &  \\
C0331 & 1a<o0L / 2a<1a2 2b<1a2L / 3a<2b2L & N/NN/N & +/{-}+/+ & N &  \\
C0332 & 1a<o0L / 2a<1a2L / 3a<2a0L 3b<2a1 & N/N/NN & +/+/+{-} & N &  \\
C0333 & 1a<o0L / 2a<1a2L / 3a<2a0L 3b<2a2 & N/N/NN & +/+/+{-} & N &  \\
C0334 & 1a<o0L / 2a<1a2L / 3a<2a1 3b<2a1 & N/N/NN & +/+/{-}{-} & N &  \\
C0335 & 1a<o0L / 2a<1a2L / 3a<2a1 3b<2a1L & N/N/NN & +/+/{-}+ & N &  \\
C0336 & 1a<o0L / 2a<1a2L / 3a<2a1 3b<2a2 & N/N/NN & +/+/{-}{-} & N &  \\
C0337 & 1a<o0L / 2a<1a2L / 3a<2a1 3b<2a2L & N/N/NN & +/+/{-}+ & N &  \\
C0338 & 1a<o0L / 2a<1a2L / 3a<2a1L 3b<2a2 & N/N/NN & +/+/+{-} & N &  \\
C0339 & 1a<o0L / 2a<1a2L / 3a<2a2 3b<2a2 & N/N/NN & +/+/{-}{-} & N &  \\
C0340 & 1a<o0L / 2a<1a2L / 3a<2a2 3b<2a2L & N/N/NN & +/+/{-}+ & N &  \\
C0341 & 1a<o0L 1b<o1 / 2a<1a0L / 3a<2a0L & NI/N/N & +{-}/+/+ & N &  \\
C0342 & 1a<o0L 1b<o1 / 2a<1a0L / 3a<2a1 & NI/N/N & +{-}/+/{-} & N &  \\
C0343 & 1a<o0L 1b<o1 / 2a<1a0L / 3a<2a1L & NI/N/N & +{-}/+/+ & N &  \\
C0344 & 1a<o0L 1b<o1 / 2a<1a0L / 3a<2a2 & NI/N/N & +{-}/+/{-} & N &  \\
C0345 & 1a<o0L 1b<o1 / 2a<1a0L / 3a<2a2L & NI/N/N & +{-}/+/+ & N &  \\
C0346 & 1a<o0L 1b<o1 / 2a<1a1 / 3a<2a0L & NI/N/N & +{-}/{-}/{-} & N &  \\
C0347 & 1a<o0L 1b<o1 / 2a<1a1 / 3a<2a1 & NI/N/N & +{-}/{-}/{-} & N &  \\
C0348 & 1a<o0L 1b<o1 / 2a<1a1 / 3a<2a1L & NI/N/N & +{-}/{-}/{-} & N &  \\
C0349 & 1a<o0L 1b<o1 / 2a<1a1 / 3a<2a2 & NI/N/N & +{-}/{-}/{-} & N &  \\
C0350 & 1a<o0L 1b<o1 / 2a<1a1 / 3a<2a2L & NI/N/N & +{-}/{-}/{-} & N &  \\
C0351 & 1a<o0L 1b<o1 / 2a<1a1L / 3a<2a0L & NI/N/N & +{-}/+/+ & N &  \\
C0352 & 1a<o0L 1b<o1 / 2a<1a1L / 3a<2a1 & NI/N/N & +{-}/+/{-} & N &  \\
C0353 & 1a<o0L 1b<o1 / 2a<1a1L / 3a<2a1L & NI/N/N & +{-}/+/+ & N &  \\
C0354 & 1a<o0L 1b<o1 / 2a<1a1L / 3a<2a2 & NI/N/N & +{-}/+/{-} & N &  \\
C0355 & 1a<o0L 1b<o1 / 2a<1a1L / 3a<2a2L & NI/N/N & +{-}/+/+ & N &  \\
C0356 & 1a<o0L 1b<o1 / 2a<1a2 / 3a<2a0L & NI/N/N & +{-}/{-}/{-} & N &  \\
C0357 & 1a<o0L 1b<o1 / 2a<1a2 / 3a<2a1 & NI/N/N & +{-}/{-}/{-} & N &  \\
C0358 & 1a<o0L 1b<o1 / 2a<1a2 / 3a<2a1L & NI/N/N & +{-}/{-}/{-} & N &  \\
C0359 & 1a<o0L 1b<o1 / 2a<1a2 / 3a<2a2 & NI/N/N & +{-}/{-}/{-} & N &  \\
C0360 & 1a<o0L 1b<o1 / 2a<1a2 / 3a<2a2L & NI/N/N & +{-}/{-}/{-} & N &  \\
C0361 & 1a<o0L 1b<o1 / 2a<1a2L / 3a<2a0L & NI/N/N & +{-}/+/+ & N &  \\
C0362 & 1a<o0L 1b<o1 / 2a<1a2L / 3a<2a1 & NI/N/N & +{-}/+/{-} & N &  \\
C0363 & 1a<o0L 1b<o1 / 2a<1a2L / 3a<2a1L & NI/N/N & +{-}/+/+ & N &  \\
C0364 & 1a<o0L 1b<o1 / 2a<1a2L / 3a<2a2 & NI/N/N & +{-}/+/{-} & N &  \\
C0365 & 1a<o0L 1b<o1 / 2a<1a2L / 3a<2a2L & NI/N/N & +{-}/+/+ & N &  \\
C0366 & 1a<o0L 1b<o1 / 2a<1b0L / 3a<2a0L & NI/N/N & +{-}/{-}/{-} & N &  \\
C0367 & 1a<o0L 1b<o1 / 2a<1b0L / 3a<2a1 & NI/N/N & +{-}/{-}/{-} & N &  \\
C0368 & 1a<o0L 1b<o1 / 2a<1b0L / 3a<2a1L & NI/N/N & +{-}/{-}/{-} & N &  \\
C0369 & 1a<o0L 1b<o1 / 2a<1b0L / 3a<2a2 & NI/N/N & +{-}/{-}/{-} & N &  \\
C0370 & 1a<o0L 1b<o1 / 2a<1b0L / 3a<2a2L & NI/N/N & +{-}/{-}/{-} & N &  \\
C0371 & 1a<o0L 1b<o1 / 2a<1b1 / 3a<2a0L & NI/I/N & +{-}/{-}/{-} & N &  \\
C0372 & 1a<o0L 1b<o1 / 2a<1b1 / 3a<2a1 & NI/I/I & +{-}/{-}/{-} & N &  \\
C0373 & 1a<o0L 1b<o1 / 2a<1b1 / 3a<2a1L & NI/I/I & +{-}/{-}/{-} & N &  \\
C0374 & 1a<o0L 1b<o1 / 2a<1b1 / 3a<2a2 & NI/I/I & +{-}/{-}/{-} & N &  \\
C0375 & 1a<o0L 1b<o1 / 2a<1b1 / 3a<2a2L & NI/I/I & +{-}/{-}/{-} & N &  \\
C0376 & 1a<o0L 1b<o1 / 2a<1b1L / 3a<2a0L & NI/I/N & +{-}/{-}/{-} & N &  \\
C0377 & 1a<o0L 1b<o1 / 2a<1b1L / 3a<2a1 & NI/I/I & +{-}/{-}/{-} & N &  \\
C0378 & 1a<o0L 1b<o1 / 2a<1b1L / 3a<2a1L & NI/I/I & +{-}/{-}/{-} & N &  \\
C0379 & 1a<o0L 1b<o1 / 2a<1b1L / 3a<2a2 & NI/I/I & +{-}/{-}/{-} & N &  \\
C0380 & 1a<o0L 1b<o1 / 2a<1b1L / 3a<2a2L & NI/I/I & +{-}/{-}/{-} & N &  \\
C0381 & 1a<o0L 1b<o1 / 2a<1b2 / 3a<2a0L & NI/I/N & +{-}/{-}/{-} & N &  \\
C0382 & 1a<o0L 1b<o1 / 2a<1b2 / 3a<2a1 & NI/I/I & +{-}/{-}/{-} & N &  \\
C0383 & 1a<o0L 1b<o1 / 2a<1b2 / 3a<2a1L & NI/I/I & +{-}/{-}/{-} & N &  \\
C0384 & 1a<o0L 1b<o1 / 2a<1b2 / 3a<2a2 & NI/I/I & +{-}/{-}/{-} & N &  \\
C0385 & 1a<o0L 1b<o1 / 2a<1b2 / 3a<2a2L & NI/I/I & +{-}/{-}/{-} & N &  \\
C0386 & 1a<o0L 1b<o1 / 2a<1b2L / 3a<2a0L & NI/I/N & +{-}/{-}/{-} & N &  \\
C0387 & 1a<o0L 1b<o1 / 2a<1b2L / 3a<2a1 & NI/I/I & +{-}/{-}/{-} & N &  \\
C0388 & 1a<o0L 1b<o1 / 2a<1b2L / 3a<2a1L & NI/I/I & +{-}/{-}/{-} & N &  \\
C0389 & 1a<o0L 1b<o1 / 2a<1b2L / 3a<2a2 & NI/I/I & +{-}/{-}/{-} & N &  \\
C0390 & 1a<o0L 1b<o1 / 2a<1b2L / 3a<2a2L & NI/I/I & +{-}/{-}/{-} & N &  \\
C0391 & 1a<o0L 1b<o2 / 2a<1a0L / 3a<2a0L & NI/N/N & +{-}/+/+ & N &  \\
C0392 & 1a<o0L 1b<o2 / 2a<1a0L / 3a<2a1 & NI/N/N & +{-}/+/{-} & N &  \\
C0393 & 1a<o0L 1b<o2 / 2a<1a0L / 3a<2a1L & NI/N/N & +{-}/+/+ & N &  \\
C0394 & 1a<o0L 1b<o2 / 2a<1a0L / 3a<2a2 & NI/N/N & +{-}/+/{-} & N &  \\
C0395 & 1a<o0L 1b<o2 / 2a<1a0L / 3a<2a2L & NI/N/N & +{-}/+/+ & N &  \\
C0396 & 1a<o0L 1b<o2 / 2a<1a1 / 3a<2a0L & NI/N/N & +{-}/{-}/{-} & N &  \\
C0397 & 1a<o0L 1b<o2 / 2a<1a1 / 3a<2a1 & NI/N/N & +{-}/{-}/{-} & N &  \\
C0398 & 1a<o0L 1b<o2 / 2a<1a1 / 3a<2a1L & NI/N/N & +{-}/{-}/{-} & N &  \\
C0399 & 1a<o0L 1b<o2 / 2a<1a1 / 3a<2a2 & NI/N/N & +{-}/{-}/{-} & N &  \\
C0400 & 1a<o0L 1b<o2 / 2a<1a1 / 3a<2a2L & NI/N/N & +{-}/{-}/{-} & N &  \\
C0401 & 1a<o0L 1b<o2 / 2a<1a1L / 3a<2a0L & NI/N/N & +{-}/+/+ & N &  \\
C0402 & 1a<o0L 1b<o2 / 2a<1a1L / 3a<2a1 & NI/N/N & +{-}/+/{-} & N &  \\
C0403 & 1a<o0L 1b<o2 / 2a<1a1L / 3a<2a1L & NI/N/N & +{-}/+/+ & N &  \\
C0404 & 1a<o0L 1b<o2 / 2a<1a1L / 3a<2a2 & NI/N/N & +{-}/+/{-} & N &  \\
C0405 & 1a<o0L 1b<o2 / 2a<1a1L / 3a<2a2L & NI/N/N & +{-}/+/+ & N &  \\
C0406 & 1a<o0L 1b<o2 / 2a<1a2 / 3a<2a0L & NI/N/N & +{-}/{-}/{-} & N &  \\
C0407 & 1a<o0L 1b<o2 / 2a<1a2 / 3a<2a1 & NI/N/N & +{-}/{-}/{-} & N &  \\
C0408 & 1a<o0L 1b<o2 / 2a<1a2 / 3a<2a1L & NI/N/N & +{-}/{-}/{-} & N &  \\
C0409 & 1a<o0L 1b<o2 / 2a<1a2 / 3a<2a2 & NI/N/N & +{-}/{-}/{-} & N &  \\
C0410 & 1a<o0L 1b<o2 / 2a<1a2 / 3a<2a2L & NI/N/N & +{-}/{-}/{-} & N &  \\
C0411 & 1a<o0L 1b<o2 / 2a<1a2L / 3a<2a0L & NI/N/N & +{-}/+/+ & N &  \\
C0412 & 1a<o0L 1b<o2 / 2a<1a2L / 3a<2a1 & NI/N/N & +{-}/+/{-} & N &  \\
C0413 & 1a<o0L 1b<o2 / 2a<1a2L / 3a<2a1L & NI/N/N & +{-}/+/+ & N &  \\
C0414 & 1a<o0L 1b<o2 / 2a<1a2L / 3a<2a2 & NI/N/N & +{-}/+/{-} & N &  \\
C0415 & 1a<o0L 1b<o2 / 2a<1a2L / 3a<2a2L & NI/N/N & +{-}/+/+ & N &  \\
C0416 & 1a<o0L 1b<o2 / 2a<1b0L / 3a<2a0L & NI/N/N & +{-}/{-}/{-} & N &  \\
C0417 & 1a<o0L 1b<o2 / 2a<1b0L / 3a<2a1 & NI/N/N & +{-}/{-}/{-} & N &  \\
C0418 & 1a<o0L 1b<o2 / 2a<1b0L / 3a<2a1L & NI/N/N & +{-}/{-}/{-} & N &  \\
C0419 & 1a<o0L 1b<o2 / 2a<1b0L / 3a<2a2 & NI/N/N & +{-}/{-}/{-} & N &  \\
C0420 & 1a<o0L 1b<o2 / 2a<1b0L / 3a<2a2L & NI/N/N & +{-}/{-}/{-} & N &  \\
C0421 & 1a<o0L 1b<o2 / 2a<1b1 / 3a<2a0L & NI/I/N & +{-}/{-}/{-} & N &  \\
C0422 & 1a<o0L 1b<o2 / 2a<1b1 / 3a<2a1 & NI/I/I & +{-}/{-}/{-} & N &  \\
C0423 & 1a<o0L 1b<o2 / 2a<1b1 / 3a<2a1L & NI/I/I & +{-}/{-}/{-} & N &  \\
C0424 & 1a<o0L 1b<o2 / 2a<1b1 / 3a<2a2 & NI/I/I & +{-}/{-}/{-} & N &  \\
C0425 & 1a<o0L 1b<o2 / 2a<1b1 / 3a<2a2L & NI/I/I & +{-}/{-}/{-} & N &  \\
C0426 & 1a<o0L 1b<o2 / 2a<1b1L / 3a<2a0L & NI/I/N & +{-}/{-}/{-} & N &  \\
C0427 & 1a<o0L 1b<o2 / 2a<1b1L / 3a<2a1 & NI/I/I & +{-}/{-}/{-} & N &  \\
C0428 & 1a<o0L 1b<o2 / 2a<1b1L / 3a<2a1L & NI/I/I & +{-}/{-}/{-} & N &  \\
C0429 & 1a<o0L 1b<o2 / 2a<1b1L / 3a<2a2 & NI/I/I & +{-}/{-}/{-} & N &  \\
C0430 & 1a<o0L 1b<o2 / 2a<1b1L / 3a<2a2L & NI/I/I & +{-}/{-}/{-} & N &  \\
C0431 & 1a<o0L 1b<o2 / 2a<1b2 / 3a<2a0L & NI/I/N & +{-}/{-}/{-} & N &  \\
C0432 & 1a<o0L 1b<o2 / 2a<1b2 / 3a<2a1 & NI/I/I & +{-}/{-}/{-} & N &  \\
C0433 & 1a<o0L 1b<o2 / 2a<1b2 / 3a<2a1L & NI/I/I & +{-}/{-}/{-} & N &  \\
C0434 & 1a<o0L 1b<o2 / 2a<1b2 / 3a<2a2 & NI/I/I & +{-}/{-}/{-} & N &  \\
C0435 & 1a<o0L 1b<o2 / 2a<1b2 / 3a<2a2L & NI/I/I & +{-}/{-}/{-} & N &  \\
C0436 & 1a<o0L 1b<o2 / 2a<1b2L / 3a<2a0L & NI/I/N & +{-}/{-}/{-} & N &  \\
C0437 & 1a<o0L 1b<o2 / 2a<1b2L / 3a<2a1 & NI/I/I & +{-}/{-}/{-} & N &  \\
C0438 & 1a<o0L 1b<o2 / 2a<1b2L / 3a<2a1L & NI/I/I & +{-}/{-}/{-} & N &  \\
C0439 & 1a<o0L 1b<o2 / 2a<1b2L / 3a<2a2 & NI/I/I & +{-}/{-}/{-} & N &  \\
C0440 & 1a<o0L 1b<o2 / 2a<1b2L / 3a<2a2L & NI/I/I & +{-}/{-}/{-} & N &  \\
C0441 & 1a<o1 / 2a<1a0L / 3a<2a0L 3b<2a1 & I/N/NN & {-}/{-}/{-}{-} & N &  \\
C0442 & 1a<o1 / 2a<1a0L / 3a<2a0L 3b<2a2 & I/N/NN & {-}/{-}/{-}{-} & N &  \\
C0443 & 1a<o1 / 2a<1a0L / 3a<2a1 3b<2a1 & I/N/NN & {-}/{-}/{-}{-} & N &  \\
C0444 & 1a<o1 / 2a<1a0L / 3a<2a1 3b<2a1L & I/N/NN & {-}/{-}/{-}{-} & N &  \\
C0445 & 1a<o1 / 2a<1a0L / 3a<2a1 3b<2a2 & I/N/NN & {-}/{-}/{-}{-} & N &  \\
C0446 & 1a<o1 / 2a<1a0L / 3a<2a1 3b<2a2L & I/N/NN & {-}/{-}/{-}{-} & N &  \\
C0447 & 1a<o1 / 2a<1a0L / 3a<2a1L 3b<2a2 & I/N/NN & {-}/{-}/{-}{-} & N &  \\
C0448 & 1a<o1 / 2a<1a0L / 3a<2a2 3b<2a2 & I/N/NN & {-}/{-}/{-}{-} & N &  \\
C0449 & 1a<o1 / 2a<1a0L / 3a<2a2 3b<2a2L & I/N/NN & {-}/{-}/{-}{-} & N &  \\
C0450 & 1a<o1 / 2a<1a0L 2b<1a1 / 3a<2a0L & I/NI/N & {-}/{-}{-}/{-} & N &  \\
C0451 & 1a<o1 / 2a<1a0L 2b<1a1 / 3a<2a1 & I/NI/N & {-}/{-}{-}/{-} & N &  \\
C0452 & 1a<o1 / 2a<1a0L 2b<1a1 / 3a<2a1L & I/NI/N & {-}/{-}{-}/{-} & N &  \\
C0453 & 1a<o1 / 2a<1a0L 2b<1a1 / 3a<2a2 & I/NI/N & {-}/{-}{-}/{-} & N &  \\
C0454 & 1a<o1 / 2a<1a0L 2b<1a1 / 3a<2a2L & I/NI/N & {-}/{-}{-}/{-} & N &  \\
C0455 & 1a<o1 / 2a<1a0L 2b<1a1 / 3a<2b0L & I/NI/N & {-}/{-}{-}/{-} & N &  \\
C0456 & 1a<o1 / 2a<1a0L 2b<1a1 / 3a<2b1 & I/NI/I & {-}/{-}{-}/{-} & N &  \\
C0457 & 1a<o1 / 2a<1a0L 2b<1a1 / 3a<2b1L & I/NI/I & {-}/{-}{-}/{-} & N &  \\
C0458 & 1a<o1 / 2a<1a0L 2b<1a1 / 3a<2b2 & I/NI/I & {-}/{-}{-}/{-} & N &  \\
C0459 & 1a<o1 / 2a<1a0L 2b<1a1 / 3a<2b2L & I/NI/I & {-}/{-}{-}/{-} & N &  \\
C0460 & 1a<o1 / 2a<1a0L 2b<1a2 / 3a<2a0L & I/NI/N & {-}/{-}{-}/{-} & N &  \\
C0461 & 1a<o1 / 2a<1a0L 2b<1a2 / 3a<2a1 & I/NI/N & {-}/{-}{-}/{-} & N &  \\
C0462 & 1a<o1 / 2a<1a0L 2b<1a2 / 3a<2a1L & I/NI/N & {-}/{-}{-}/{-} & N &  \\
C0463 & 1a<o1 / 2a<1a0L 2b<1a2 / 3a<2a2 & I/NI/N & {-}/{-}{-}/{-} & N &  \\
C0464 & 1a<o1 / 2a<1a0L 2b<1a2 / 3a<2a2L & I/NI/N & {-}/{-}{-}/{-} & N &  \\
C0465 & 1a<o1 / 2a<1a0L 2b<1a2 / 3a<2b0L & I/NI/N & {-}/{-}{-}/{-} & N &  \\
C0466 & 1a<o1 / 2a<1a0L 2b<1a2 / 3a<2b1 & I/NI/I & {-}/{-}{-}/{-} & N &  \\
C0467 & 1a<o1 / 2a<1a0L 2b<1a2 / 3a<2b1L & I/NI/I & {-}/{-}{-}/{-} & N &  \\
C0468 & 1a<o1 / 2a<1a0L 2b<1a2 / 3a<2b2 & I/NI/I & {-}/{-}{-}/{-} & N &  \\
C0469 & 1a<o1 / 2a<1a0L 2b<1a2 / 3a<2b2L & I/NI/I & {-}/{-}{-}/{-} & N &  \\
C0470 & 1a<o1 / 2a<1a1 / 3a<2a0L 3b<2a1 & I/I/NI & {-}/{-}/{-}{-} & N &  \\
C0471 & 1a<o1 / 2a<1a1 / 3a<2a0L 3b<2a2 & I/I/NI & {-}/{-}/{-}{-} & N &  \\
C0472 & 1a<o1 / 2a<1a1 / 3a<2a1 3b<2a1 & I/I/EE & {-}/{-}/{-}{-} & E & W:1a,2a \\
C0473 & 1a<o1 / 2a<1a1 / 3a<2a1 3b<2a1L & I/I/LI & {-}/{-}/{-}{-} & B & CP:3b W:1a,2a,3b \\
C0474 & 1a<o1 / 2a<1a1 / 3a<2a1 3b<2a2 & I/I/EE & {-}/{-}/{-}{-} & E & C:3b W:1a,2a \\
C0475 & 1a<o1 / 2a<1a1 / 3a<2a1 3b<2a2L & I/I/LI & {-}/{-}/{-}{-} & B & P:3b W:1a,2a,3b \\
C0476 & 1a<o1 / 2a<1a1 / 3a<2a1L 3b<2a2 & I/I/IL & {-}/{-}/{-}{-} & B & C:3a,3b P:3a W:1a,2a,3a \\
C0477 & 1a<o1 / 2a<1a1 / 3a<2a2 3b<2a2 & I/I/EE & {-}/{-}/{-}{-} & E & W:1a,2a \\
C0478 & 1a<o1 / 2a<1a1 / 3a<2a2 3b<2a2L & I/I/LI & {-}/{-}/{-}{-} & B & CP:3b W:1a,2a,3b \\
C0479 & 1a<o1 / 2a<1a1 2b<1a1 / 3a<2a0L & I/EE/N & {-}/{-}{-}/{-} & E & W:1a \\
C0480 & 1a<o1 / 2a<1a1 2b<1a1 / 3a<2a1 & I/EE/D & {-}/{-}{-}/{-} & E & W:1a \\
C0481 & 1a<o1 / 2a<1a1 2b<1a1 / 3a<2a1L & I/EE/D & {-}/{-}{-}/{-} & E & W:1a \\
C0482 & 1a<o1 / 2a<1a1 2b<1a1 / 3a<2a2 & I/EE/D & {-}/{-}{-}/{-} & E & W:1a \\
C0483 & 1a<o1 / 2a<1a1 2b<1a1 / 3a<2a2L & I/EE/D & {-}/{-}{-}/{-} & E & W:1a \\
C0484 & 1a<o1 / 2a<1a1 2b<1a1L / 3a<2a0L & I/LI/N & {-}/{-}{-}/{-} & B & CP:2b W:1a,2b \\
C0485 & 1a<o1 / 2a<1a1 2b<1a1L / 3a<2a1 & I/LI/D & {-}/{-}{-}/{-} & B & CP:2b W:1a,2b \\
C0486 & 1a<o1 / 2a<1a1 2b<1a1L / 3a<2a1L & I/LI/D & {-}/{-}{-}/{-} & B & CP:2b W:1a,2b \\
C0487 & 1a<o1 / 2a<1a1 2b<1a1L / 3a<2a2 & I/LI/D & {-}/{-}{-}/{-} & B & CP:2b W:1a,2b \\
C0488 & 1a<o1 / 2a<1a1 2b<1a1L / 3a<2a2L & I/LI/D & {-}/{-}{-}/{-} & B & CP:2b W:1a,2b \\
C0489 & 1a<o1 / 2a<1a1 2b<1a1L / 3a<2b0L & I/LI/N & {-}/{-}{-}/{-} & B & CP:2b W:1a,2b \\
C0490 & 1a<o1 / 2a<1a1 2b<1a1L / 3a<2b1 & I/LI/I & {-}/{-}{-}/{-} & B & CP:2b,3a W:1a,2b,3a \\
C0491 & 1a<o1 / 2a<1a1 2b<1a1L / 3a<2b1L & I/LI/I & {-}/{-}{-}/{-} & B & CP:2b,3a W:1a,2b,3a \\
C0492 & 1a<o1 / 2a<1a1 2b<1a1L / 3a<2b2 & I/LI/I & {-}/{-}{-}/{-} & B & CP:2b,3a W:1a,2b,3a \\
C0493 & 1a<o1 / 2a<1a1 2b<1a1L / 3a<2b2L & I/LI/I & {-}/{-}{-}/{-} & B & CP:2b,3a W:1a,2b,3a \\
C0494 & 1a<o1 / 2a<1a1 2b<1a2 / 3a<2a0L & I/EE/N & {-}/{-}{-}/{-} & E & C:2b W:1a \\
C0495 & 1a<o1 / 2a<1a1 2b<1a2 / 3a<2a1 & I/EE/D & {-}/{-}{-}/{-} & E & C:2b W:1a \\
C0496 & 1a<o1 / 2a<1a1 2b<1a2 / 3a<2a1L & I/EE/D & {-}/{-}{-}/{-} & E & C:2b W:1a \\
C0497 & 1a<o1 / 2a<1a1 2b<1a2 / 3a<2a2 & I/EE/D & {-}/{-}{-}/{-} & E & C:2b W:1a \\
C0498 & 1a<o1 / 2a<1a1 2b<1a2 / 3a<2a2L & I/EE/D & {-}/{-}{-}/{-} & E & C:2b W:1a \\
C0499 & 1a<o1 / 2a<1a1 2b<1a2 / 3a<2b0L & I/EE/N & {-}/{-}{-}/{-} & E & C:2b W:1a \\
C0500 & 1a<o1 / 2a<1a1 2b<1a2 / 3a<2b1 & I/EE/D & {-}/{-}{-}/{-} & E & C:2b,3a W:1a \\
C0501 & 1a<o1 / 2a<1a1 2b<1a2 / 3a<2b1L & I/EE/D & {-}/{-}{-}/{-} & E & C:2b,3a W:1a \\
C0502 & 1a<o1 / 2a<1a1 2b<1a2 / 3a<2b2 & I/EE/D & {-}/{-}{-}/{-} & E & C:2b,3a W:1a \\
C0503 & 1a<o1 / 2a<1a1 2b<1a2 / 3a<2b2L & I/EE/D & {-}/{-}{-}/{-} & E & C:2b,3a W:1a \\
C0504 & 1a<o1 / 2a<1a1 2b<1a2L / 3a<2a0L & I/LI/N & {-}/{-}{-}/{-} & B & P:2b W:1a,2b \\
C0505 & 1a<o1 / 2a<1a1 2b<1a2L / 3a<2a1 & I/LI/D & {-}/{-}{-}/{-} & B & P:2b W:1a,2b \\
C0506 & 1a<o1 / 2a<1a1 2b<1a2L / 3a<2a1L & I/LI/D & {-}/{-}{-}/{-} & B & P:2b W:1a,2b \\
C0507 & 1a<o1 / 2a<1a1 2b<1a2L / 3a<2a2 & I/LI/D & {-}/{-}{-}/{-} & B & P:2b W:1a,2b \\
C0508 & 1a<o1 / 2a<1a1 2b<1a2L / 3a<2a2L & I/LI/D & {-}/{-}{-}/{-} & B & P:2b W:1a,2b \\
C0509 & 1a<o1 / 2a<1a1 2b<1a2L / 3a<2b0L & I/LI/N & {-}/{-}{-}/{-} & B & P:2b W:1a,2b \\
C0510 & 1a<o1 / 2a<1a1 2b<1a2L / 3a<2b1 & I/LI/I & {-}/{-}{-}/{-} & B & P:2b,3a W:1a,2b,3a \\
C0511 & 1a<o1 / 2a<1a1 2b<1a2L / 3a<2b1L & I/LI/I & {-}/{-}{-}/{-} & B & P:2b,3a W:1a,2b,3a \\
C0512 & 1a<o1 / 2a<1a1 2b<1a2L / 3a<2b2 & I/LI/I & {-}/{-}{-}/{-} & B & P:2b,3a W:1a,2b,3a \\
C0513 & 1a<o1 / 2a<1a1 2b<1a2L / 3a<2b2L & I/LI/I & {-}/{-}{-}/{-} & B & P:2b,3a W:1a,2b,3a \\
C0514 & 1a<o1 / 2a<1a1L / 3a<2a0L 3b<2a1 & I/I/NI & {-}/{-}/{-}{-} & N &  \\
C0515 & 1a<o1 / 2a<1a1L / 3a<2a0L 3b<2a2 & I/I/NI & {-}/{-}/{-}{-} & N &  \\
C0516 & 1a<o1 / 2a<1a1L / 3a<2a1 3b<2a1 & I/I/EE & {-}/{-}/{-}{-} & E & W:1a,2a \\
C0517 & 1a<o1 / 2a<1a1L / 3a<2a1 3b<2a1L & I/I/LI & {-}/{-}/{-}{-} & B & CP:3b W:1a,2a,3b \\
C0518 & 1a<o1 / 2a<1a1L / 3a<2a1 3b<2a2 & I/I/EE & {-}/{-}/{-}{-} & E & C:3b W:1a,2a \\
C0519 & 1a<o1 / 2a<1a1L / 3a<2a1 3b<2a2L & I/I/LI & {-}/{-}/{-}{-} & B & P:3b W:1a,2a,3b \\
C0520 & 1a<o1 / 2a<1a1L / 3a<2a1L 3b<2a2 & I/I/IL & {-}/{-}/{-}{-} & B & C:3a,3b P:3a W:1a,2a,3a \\
C0521 & 1a<o1 / 2a<1a1L / 3a<2a2 3b<2a2 & I/I/EE & {-}/{-}/{-}{-} & E & W:1a,2a \\
C0522 & 1a<o1 / 2a<1a1L / 3a<2a2 3b<2a2L & I/I/LI & {-}/{-}/{-}{-} & B & CP:3b W:1a,2a,3b \\
C0523 & 1a<o1 / 2a<1a1L 2b<1a2 / 3a<2a0L & I/IL/N & {-}/{-}{-}/{-} & B & C:2a,2b P:2a W:1a,2a \\
C0524 & 1a<o1 / 2a<1a1L 2b<1a2 / 3a<2a1 & I/IL/I & {-}/{-}{-}/{-} & B & C:2a,2b,3a P:2a,3a W:1a,2a,3a \\
C0525 & 1a<o1 / 2a<1a1L 2b<1a2 / 3a<2a1L & I/IL/I & {-}/{-}{-}/{-} & B & C:2a,2b,3a P:2a,3a W:1a,2a,3a \\
C0526 & 1a<o1 / 2a<1a1L 2b<1a2 / 3a<2a2 & I/IL/I & {-}/{-}{-}/{-} & B & C:2a,2b,3a P:2a,3a W:1a,2a,3a \\
C0527 & 1a<o1 / 2a<1a1L 2b<1a2 / 3a<2a2L & I/IL/I & {-}/{-}{-}/{-} & B & C:2a,2b,3a P:2a,3a W:1a,2a,3a \\
C0528 & 1a<o1 / 2a<1a1L 2b<1a2 / 3a<2b0L & I/IL/N & {-}/{-}{-}/{-} & B & C:2a,2b P:2a W:1a,2a \\
C0529 & 1a<o1 / 2a<1a1L 2b<1a2 / 3a<2b1 & I/IL/D & {-}/{-}{-}/{-} & B & C:2a,2b,3a P:2a W:1a,2a \\
C0530 & 1a<o1 / 2a<1a1L 2b<1a2 / 3a<2b1L & I/IL/D & {-}/{-}{-}/{-} & B & C:2a,2b,3a P:2a W:1a,2a \\
C0531 & 1a<o1 / 2a<1a1L 2b<1a2 / 3a<2b2 & I/IL/D & {-}/{-}{-}/{-} & B & C:2a,2b,3a P:2a W:1a,2a \\
C0532 & 1a<o1 / 2a<1a1L 2b<1a2 / 3a<2b2L & I/IL/D & {-}/{-}{-}/{-} & B & C:2a,2b,3a P:2a W:1a,2a \\
C0533 & 1a<o1 / 2a<1a2 / 3a<2a0L 3b<2a1 & I/I/NI & {-}/{-}/{-}{-} & N &  \\
C0534 & 1a<o1 / 2a<1a2 / 3a<2a0L 3b<2a2 & I/I/NI & {-}/{-}/{-}{-} & N &  \\
C0535 & 1a<o1 / 2a<1a2 / 3a<2a1 3b<2a1 & I/I/EE & {-}/{-}/{-}{-} & E & W:1a,2a \\
C0536 & 1a<o1 / 2a<1a2 / 3a<2a1 3b<2a1L & I/I/LI & {-}/{-}/{-}{-} & B & CP:3b W:1a,2a,3b \\
C0537 & 1a<o1 / 2a<1a2 / 3a<2a1 3b<2a2 & I/I/EE & {-}/{-}/{-}{-} & E & C:3b W:1a,2a \\
C0538 & 1a<o1 / 2a<1a2 / 3a<2a1 3b<2a2L & I/I/LI & {-}/{-}/{-}{-} & B & P:3b W:1a,2a,3b \\
C0539 & 1a<o1 / 2a<1a2 / 3a<2a1L 3b<2a2 & I/I/IL & {-}/{-}/{-}{-} & B & C:3a,3b P:3a W:1a,2a,3a \\
C0540 & 1a<o1 / 2a<1a2 / 3a<2a2 3b<2a2 & I/I/EE & {-}/{-}/{-}{-} & E & W:1a,2a \\
C0541 & 1a<o1 / 2a<1a2 / 3a<2a2 3b<2a2L & I/I/LI & {-}/{-}/{-}{-} & B & CP:3b W:1a,2a,3b \\
C0542 & 1a<o1 / 2a<1a2 2b<1a2 / 3a<2a0L & I/EE/N & {-}/{-}{-}/{-} & E & W:1a \\
C0543 & 1a<o1 / 2a<1a2 2b<1a2 / 3a<2a1 & I/EE/D & {-}/{-}{-}/{-} & E & W:1a \\
C0544 & 1a<o1 / 2a<1a2 2b<1a2 / 3a<2a1L & I/EE/D & {-}/{-}{-}/{-} & E & W:1a \\
C0545 & 1a<o1 / 2a<1a2 2b<1a2 / 3a<2a2 & I/EE/D & {-}/{-}{-}/{-} & E & W:1a \\
C0546 & 1a<o1 / 2a<1a2 2b<1a2 / 3a<2a2L & I/EE/D & {-}/{-}{-}/{-} & E & W:1a \\
C0547 & 1a<o1 / 2a<1a2 2b<1a2L / 3a<2a0L & I/LI/N & {-}/{-}{-}/{-} & B & CP:2b W:1a,2b \\
C0548 & 1a<o1 / 2a<1a2 2b<1a2L / 3a<2a1 & I/LI/D & {-}/{-}{-}/{-} & B & CP:2b W:1a,2b \\
C0549 & 1a<o1 / 2a<1a2 2b<1a2L / 3a<2a1L & I/LI/D & {-}/{-}{-}/{-} & B & CP:2b W:1a,2b \\
C0550 & 1a<o1 / 2a<1a2 2b<1a2L / 3a<2a2 & I/LI/D & {-}/{-}{-}/{-} & B & CP:2b W:1a,2b \\
C0551 & 1a<o1 / 2a<1a2 2b<1a2L / 3a<2a2L & I/LI/D & {-}/{-}{-}/{-} & B & CP:2b W:1a,2b \\
C0552 & 1a<o1 / 2a<1a2 2b<1a2L / 3a<2b0L & I/LI/N & {-}/{-}{-}/{-} & B & CP:2b W:1a,2b \\
C0553 & 1a<o1 / 2a<1a2 2b<1a2L / 3a<2b1 & I/LI/I & {-}/{-}{-}/{-} & B & CP:2b,3a W:1a,2b,3a \\
C0554 & 1a<o1 / 2a<1a2 2b<1a2L / 3a<2b1L & I/LI/I & {-}/{-}{-}/{-} & B & CP:2b,3a W:1a,2b,3a \\
C0555 & 1a<o1 / 2a<1a2 2b<1a2L / 3a<2b2 & I/LI/I & {-}/{-}{-}/{-} & B & CP:2b,3a W:1a,2b,3a \\
C0556 & 1a<o1 / 2a<1a2 2b<1a2L / 3a<2b2L & I/LI/I & {-}/{-}{-}/{-} & B & CP:2b,3a W:1a,2b,3a \\
C0557 & 1a<o1 / 2a<1a2L / 3a<2a0L 3b<2a1 & I/I/NI & {-}/{-}/{-}{-} & N &  \\
C0558 & 1a<o1 / 2a<1a2L / 3a<2a0L 3b<2a2 & I/I/NI & {-}/{-}/{-}{-} & N &  \\
C0559 & 1a<o1 / 2a<1a2L / 3a<2a1 3b<2a1 & I/I/EE & {-}/{-}/{-}{-} & E & W:1a,2a \\
C0560 & 1a<o1 / 2a<1a2L / 3a<2a1 3b<2a1L & I/I/LI & {-}/{-}/{-}{-} & B & CP:3b W:1a,2a,3b \\
C0561 & 1a<o1 / 2a<1a2L / 3a<2a1 3b<2a2 & I/I/EE & {-}/{-}/{-}{-} & E & C:3b W:1a,2a \\
C0562 & 1a<o1 / 2a<1a2L / 3a<2a1 3b<2a2L & I/I/LI & {-}/{-}/{-}{-} & B & P:3b W:1a,2a,3b \\
C0563 & 1a<o1 / 2a<1a2L / 3a<2a1L 3b<2a2 & I/I/IL & {-}/{-}/{-}{-} & B & C:3a,3b P:3a W:1a,2a,3a \\
C0564 & 1a<o1 / 2a<1a2L / 3a<2a2 3b<2a2 & I/I/EE & {-}/{-}/{-}{-} & E & W:1a,2a \\
C0565 & 1a<o1 / 2a<1a2L / 3a<2a2 3b<2a2L & I/I/LI & {-}/{-}/{-}{-} & B & CP:3b W:1a,2a,3b \\
C0566 & 1a<o1 1b<o1 / 2a<1a0L / 3a<2a0L & EE/N/N & {-}{-}/{-}/{-} & E &  \\
C0567 & 1a<o1 1b<o1 / 2a<1a0L / 3a<2a1 & EE/N/N & {-}{-}/{-}/{-} & E &  \\
C0568 & 1a<o1 1b<o1 / 2a<1a0L / 3a<2a1L & EE/N/N & {-}{-}/{-}/{-} & E &  \\
C0569 & 1a<o1 1b<o1 / 2a<1a0L / 3a<2a2 & EE/N/N & {-}{-}/{-}/{-} & E &  \\
C0570 & 1a<o1 1b<o1 / 2a<1a0L / 3a<2a2L & EE/N/N & {-}{-}/{-}/{-} & E &  \\
C0571 & 1a<o1 1b<o1 / 2a<1a1 / 3a<2a0L & EE/D/N & {-}{-}/{-}/{-} & E &  \\
C0572 & 1a<o1 1b<o1 / 2a<1a1 / 3a<2a1 & EE/D/D & {-}{-}/{-}/{-} & E &  \\
C0573 & 1a<o1 1b<o1 / 2a<1a1 / 3a<2a1L & EE/D/D & {-}{-}/{-}/{-} & E &  \\
C0574 & 1a<o1 1b<o1 / 2a<1a1 / 3a<2a2 & EE/D/D & {-}{-}/{-}/{-} & E &  \\
C0575 & 1a<o1 1b<o1 / 2a<1a1 / 3a<2a2L & EE/D/D & {-}{-}/{-}/{-} & E &  \\
C0576 & 1a<o1 1b<o1 / 2a<1a1L / 3a<2a0L & EE/D/N & {-}{-}/{-}/{-} & E &  \\
C0577 & 1a<o1 1b<o1 / 2a<1a1L / 3a<2a1 & EE/D/D & {-}{-}/{-}/{-} & E &  \\
C0578 & 1a<o1 1b<o1 / 2a<1a1L / 3a<2a1L & EE/D/D & {-}{-}/{-}/{-} & E &  \\
C0579 & 1a<o1 1b<o1 / 2a<1a1L / 3a<2a2 & EE/D/D & {-}{-}/{-}/{-} & E &  \\
C0580 & 1a<o1 1b<o1 / 2a<1a1L / 3a<2a2L & EE/D/D & {-}{-}/{-}/{-} & E &  \\
C0581 & 1a<o1 1b<o1 / 2a<1a2 / 3a<2a0L & EE/D/N & {-}{-}/{-}/{-} & E &  \\
C0582 & 1a<o1 1b<o1 / 2a<1a2 / 3a<2a1 & EE/D/D & {-}{-}/{-}/{-} & E &  \\
C0583 & 1a<o1 1b<o1 / 2a<1a2 / 3a<2a1L & EE/D/D & {-}{-}/{-}/{-} & E &  \\
C0584 & 1a<o1 1b<o1 / 2a<1a2 / 3a<2a2 & EE/D/D & {-}{-}/{-}/{-} & E &  \\
C0585 & 1a<o1 1b<o1 / 2a<1a2 / 3a<2a2L & EE/D/D & {-}{-}/{-}/{-} & E &  \\
C0586 & 1a<o1 1b<o1 / 2a<1a2L / 3a<2a0L & EE/D/N & {-}{-}/{-}/{-} & E &  \\
C0587 & 1a<o1 1b<o1 / 2a<1a2L / 3a<2a1 & EE/D/D & {-}{-}/{-}/{-} & E &  \\
C0588 & 1a<o1 1b<o1 / 2a<1a2L / 3a<2a1L & EE/D/D & {-}{-}/{-}/{-} & E &  \\
C0589 & 1a<o1 1b<o1 / 2a<1a2L / 3a<2a2 & EE/D/D & {-}{-}/{-}/{-} & E &  \\
C0590 & 1a<o1 1b<o1 / 2a<1a2L / 3a<2a2L & EE/D/D & {-}{-}/{-}/{-} & E &  \\
C0591 & 1a<o1 1b<o1L / 2a<1a0L / 3a<2a0L & LI/N/N & {-}+/{-}/{-} & B & CPW:1b \\
C0592 & 1a<o1 1b<o1L / 2a<1a0L / 3a<2a1 & LI/N/N & {-}+/{-}/{-} & B & CPW:1b \\
C0593 & 1a<o1 1b<o1L / 2a<1a0L / 3a<2a1L & LI/N/N & {-}+/{-}/{-} & B & CPW:1b \\
C0594 & 1a<o1 1b<o1L / 2a<1a0L / 3a<2a2 & LI/N/N & {-}+/{-}/{-} & B & CPW:1b \\
C0595 & 1a<o1 1b<o1L / 2a<1a0L / 3a<2a2L & LI/N/N & {-}+/{-}/{-} & B & CPW:1b \\
C0596 & 1a<o1 1b<o1L / 2a<1a1 / 3a<2a0L & LI/D/N & {-}+/{-}/{-} & B & CPW:1b \\
C0597 & 1a<o1 1b<o1L / 2a<1a1 / 3a<2a1 & LI/D/D & {-}+/{-}/{-} & B & CPW:1b \\
C0598 & 1a<o1 1b<o1L / 2a<1a1 / 3a<2a1L & LI/D/D & {-}+/{-}/{-} & B & CPW:1b \\
C0599 & 1a<o1 1b<o1L / 2a<1a1 / 3a<2a2 & LI/D/D & {-}+/{-}/{-} & B & CPW:1b \\
C0600 & 1a<o1 1b<o1L / 2a<1a1 / 3a<2a2L & LI/D/D & {-}+/{-}/{-} & B & CPW:1b \\
C0601 & 1a<o1 1b<o1L / 2a<1a1L / 3a<2a0L & LI/D/N & {-}+/{-}/{-} & B & CPW:1b \\
C0602 & 1a<o1 1b<o1L / 2a<1a1L / 3a<2a1 & LI/D/D & {-}+/{-}/{-} & B & CPW:1b \\
C0603 & 1a<o1 1b<o1L / 2a<1a1L / 3a<2a1L & LI/D/D & {-}+/{-}/{-} & B & CPW:1b \\
C0604 & 1a<o1 1b<o1L / 2a<1a1L / 3a<2a2 & LI/D/D & {-}+/{-}/{-} & B & CPW:1b \\
C0605 & 1a<o1 1b<o1L / 2a<1a1L / 3a<2a2L & LI/D/D & {-}+/{-}/{-} & B & CPW:1b \\
C0606 & 1a<o1 1b<o1L / 2a<1a2 / 3a<2a0L & LI/D/N & {-}+/{-}/{-} & B & CPW:1b \\
C0607 & 1a<o1 1b<o1L / 2a<1a2 / 3a<2a1 & LI/D/D & {-}+/{-}/{-} & B & CPW:1b \\
C0608 & 1a<o1 1b<o1L / 2a<1a2 / 3a<2a1L & LI/D/D & {-}+/{-}/{-} & B & CPW:1b \\
C0609 & 1a<o1 1b<o1L / 2a<1a2 / 3a<2a2 & LI/D/D & {-}+/{-}/{-} & B & CPW:1b \\
C0610 & 1a<o1 1b<o1L / 2a<1a2 / 3a<2a2L & LI/D/D & {-}+/{-}/{-} & B & CPW:1b \\
C0611 & 1a<o1 1b<o1L / 2a<1a2L / 3a<2a0L & LI/D/N & {-}+/{-}/{-} & B & CPW:1b \\
C0612 & 1a<o1 1b<o1L / 2a<1a2L / 3a<2a1 & LI/D/D & {-}+/{-}/{-} & B & CPW:1b \\
C0613 & 1a<o1 1b<o1L / 2a<1a2L / 3a<2a1L & LI/D/D & {-}+/{-}/{-} & B & CPW:1b \\
C0614 & 1a<o1 1b<o1L / 2a<1a2L / 3a<2a2 & LI/D/D & {-}+/{-}/{-} & B & CPW:1b \\
C0615 & 1a<o1 1b<o1L / 2a<1a2L / 3a<2a2L & LI/D/D & {-}+/{-}/{-} & B & CPW:1b \\
C0616 & 1a<o1 1b<o1L / 2a<1b0L / 3a<2a0L & LI/N/N & {-}+/+/+ & B & CPW:1b \\
C0617 & 1a<o1 1b<o1L / 2a<1b0L / 3a<2a1 & LI/N/N & {-}+/+/{-} & B & CPW:1b \\
C0618 & 1a<o1 1b<o1L / 2a<1b0L / 3a<2a1L & LI/N/N & {-}+/+/+ & B & CPW:1b \\
C0619 & 1a<o1 1b<o1L / 2a<1b0L / 3a<2a2 & LI/N/N & {-}+/+/{-} & B & CPW:1b \\
C0620 & 1a<o1 1b<o1L / 2a<1b0L / 3a<2a2L & LI/N/N & {-}+/+/+ & B & CPW:1b \\
C0621 & 1a<o1 1b<o1L / 2a<1b1 / 3a<2a0L & LI/I/N & {-}+/{-}/{-} & B & CPW:1b,2a \\
C0622 & 1a<o1 1b<o1L / 2a<1b1 / 3a<2a1 & LI/I/I & {-}+/{-}/{-} & B & CPW:1b,2a,3a \\
C0623 & 1a<o1 1b<o1L / 2a<1b1 / 3a<2a1L & LI/I/I & {-}+/{-}/{-} & B & CPW:1b,2a,3a \\
C0624 & 1a<o1 1b<o1L / 2a<1b1 / 3a<2a2 & LI/I/I & {-}+/{-}/{-} & B & CPW:1b,2a,3a \\
C0625 & 1a<o1 1b<o1L / 2a<1b1 / 3a<2a2L & LI/I/I & {-}+/{-}/{-} & B & CPW:1b,2a,3a \\
C0626 & 1a<o1 1b<o1L / 2a<1b1L / 3a<2a0L & LI/I/N & {-}+/+/+ & B & CPW:1b,2a \\
C0627 & 1a<o1 1b<o1L / 2a<1b1L / 3a<2a1 & LI/I/I & {-}+/+/{-} & B & CPW:1b,2a,3a \\
C0628 & 1a<o1 1b<o1L / 2a<1b1L / 3a<2a1L & LI/I/I & {-}+/+/+ & B & CPW:1b,2a,3a \\
C0629 & 1a<o1 1b<o1L / 2a<1b1L / 3a<2a2 & LI/I/I & {-}+/+/{-} & B & CPW:1b,2a,3a \\
C0630 & 1a<o1 1b<o1L / 2a<1b1L / 3a<2a2L & LI/I/I & {-}+/+/+ & B & CPW:1b,2a,3a \\
C0631 & 1a<o1 1b<o1L / 2a<1b2 / 3a<2a0L & LI/I/N & {-}+/{-}/{-} & B & CPW:1b,2a \\
C0632 & 1a<o1 1b<o1L / 2a<1b2 / 3a<2a1 & LI/I/I & {-}+/{-}/{-} & B & CPW:1b,2a,3a \\
C0633 & 1a<o1 1b<o1L / 2a<1b2 / 3a<2a1L & LI/I/I & {-}+/{-}/{-} & B & CPW:1b,2a,3a \\
C0634 & 1a<o1 1b<o1L / 2a<1b2 / 3a<2a2 & LI/I/I & {-}+/{-}/{-} & B & CPW:1b,2a,3a \\
C0635 & 1a<o1 1b<o1L / 2a<1b2 / 3a<2a2L & LI/I/I & {-}+/{-}/{-} & B & CPW:1b,2a,3a \\
C0636 & 1a<o1 1b<o1L / 2a<1b2L / 3a<2a0L & LI/I/N & {-}+/+/+ & B & CPW:1b,2a \\
C0637 & 1a<o1 1b<o1L / 2a<1b2L / 3a<2a1 & LI/I/I & {-}+/+/{-} & B & CPW:1b,2a,3a \\
C0638 & 1a<o1 1b<o1L / 2a<1b2L / 3a<2a1L & LI/I/I & {-}+/+/+ & B & CPW:1b,2a,3a \\
C0639 & 1a<o1 1b<o1L / 2a<1b2L / 3a<2a2 & LI/I/I & {-}+/+/{-} & B & CPW:1b,2a,3a \\
C0640 & 1a<o1 1b<o1L / 2a<1b2L / 3a<2a2L & LI/I/I & {-}+/+/+ & B & CPW:1b,2a,3a \\
C0641 & 1a<o1 1b<o2 / 2a<1a0L / 3a<2a0L & EE/N/N & {-}{-}/{-}/{-} & E & C:1b \\
C0642 & 1a<o1 1b<o2 / 2a<1a0L / 3a<2a1 & EE/N/N & {-}{-}/{-}/{-} & E & C:1b \\
C0643 & 1a<o1 1b<o2 / 2a<1a0L / 3a<2a1L & EE/N/N & {-}{-}/{-}/{-} & E & C:1b \\
C0644 & 1a<o1 1b<o2 / 2a<1a0L / 3a<2a2 & EE/N/N & {-}{-}/{-}/{-} & E & C:1b \\
C0645 & 1a<o1 1b<o2 / 2a<1a0L / 3a<2a2L & EE/N/N & {-}{-}/{-}/{-} & E & C:1b \\
C0646 & 1a<o1 1b<o2 / 2a<1a1 / 3a<2a0L & EE/D/N & {-}{-}/{-}/{-} & E & C:1b \\
C0647 & 1a<o1 1b<o2 / 2a<1a1 / 3a<2a1 & EE/D/D & {-}{-}/{-}/{-} & E & C:1b \\
C0648 & 1a<o1 1b<o2 / 2a<1a1 / 3a<2a1L & EE/D/D & {-}{-}/{-}/{-} & E & C:1b \\
C0649 & 1a<o1 1b<o2 / 2a<1a1 / 3a<2a2 & EE/D/D & {-}{-}/{-}/{-} & E & C:1b \\
C0650 & 1a<o1 1b<o2 / 2a<1a1 / 3a<2a2L & EE/D/D & {-}{-}/{-}/{-} & E & C:1b \\
C0651 & 1a<o1 1b<o2 / 2a<1a1L / 3a<2a0L & EE/D/N & {-}{-}/{-}/{-} & E & C:1b \\
C0652 & 1a<o1 1b<o2 / 2a<1a1L / 3a<2a1 & EE/D/D & {-}{-}/{-}/{-} & E & C:1b \\
C0653 & 1a<o1 1b<o2 / 2a<1a1L / 3a<2a1L & EE/D/D & {-}{-}/{-}/{-} & E & C:1b \\
C0654 & 1a<o1 1b<o2 / 2a<1a1L / 3a<2a2 & EE/D/D & {-}{-}/{-}/{-} & E & C:1b \\
C0655 & 1a<o1 1b<o2 / 2a<1a1L / 3a<2a2L & EE/D/D & {-}{-}/{-}/{-} & E & C:1b \\
C0656 & 1a<o1 1b<o2 / 2a<1a2 / 3a<2a0L & EE/D/N & {-}{-}/{-}/{-} & E & C:1b \\
C0657 & 1a<o1 1b<o2 / 2a<1a2 / 3a<2a1 & EE/D/D & {-}{-}/{-}/{-} & E & C:1b \\
C0658 & 1a<o1 1b<o2 / 2a<1a2 / 3a<2a1L & EE/D/D & {-}{-}/{-}/{-} & E & C:1b \\
C0659 & 1a<o1 1b<o2 / 2a<1a2 / 3a<2a2 & EE/D/D & {-}{-}/{-}/{-} & E & C:1b \\
C0660 & 1a<o1 1b<o2 / 2a<1a2 / 3a<2a2L & EE/D/D & {-}{-}/{-}/{-} & E & C:1b \\
C0661 & 1a<o1 1b<o2 / 2a<1a2L / 3a<2a0L & EE/D/N & {-}{-}/{-}/{-} & E & C:1b \\
C0662 & 1a<o1 1b<o2 / 2a<1a2L / 3a<2a1 & EE/D/D & {-}{-}/{-}/{-} & E & C:1b \\
C0663 & 1a<o1 1b<o2 / 2a<1a2L / 3a<2a1L & EE/D/D & {-}{-}/{-}/{-} & E & C:1b \\
C0664 & 1a<o1 1b<o2 / 2a<1a2L / 3a<2a2 & EE/D/D & {-}{-}/{-}/{-} & E & C:1b \\
C0665 & 1a<o1 1b<o2 / 2a<1a2L / 3a<2a2L & EE/D/D & {-}{-}/{-}/{-} & E & C:1b \\
C0666 & 1a<o1 1b<o2 / 2a<1b0L / 3a<2a0L & EE/N/N & {-}{-}/{-}/{-} & E & C:1b \\
C0667 & 1a<o1 1b<o2 / 2a<1b0L / 3a<2a1 & EE/N/N & {-}{-}/{-}/{-} & E & C:1b \\
C0668 & 1a<o1 1b<o2 / 2a<1b0L / 3a<2a1L & EE/N/N & {-}{-}/{-}/{-} & E & C:1b \\
C0669 & 1a<o1 1b<o2 / 2a<1b0L / 3a<2a2 & EE/N/N & {-}{-}/{-}/{-} & E & C:1b \\
C0670 & 1a<o1 1b<o2 / 2a<1b0L / 3a<2a2L & EE/N/N & {-}{-}/{-}/{-} & E & C:1b \\
C0671 & 1a<o1 1b<o2 / 2a<1b1 / 3a<2a0L & EE/D/N & {-}{-}/{-}/{-} & E & C:1b,2a \\
C0672 & 1a<o1 1b<o2 / 2a<1b1 / 3a<2a1 & EE/D/D & {-}{-}/{-}/{-} & E & C:1b,2a,3a \\
C0673 & 1a<o1 1b<o2 / 2a<1b1 / 3a<2a1L & EE/D/D & {-}{-}/{-}/{-} & E & C:1b,2a,3a \\
C0674 & 1a<o1 1b<o2 / 2a<1b1 / 3a<2a2 & EE/D/D & {-}{-}/{-}/{-} & E & C:1b,2a,3a \\
C0675 & 1a<o1 1b<o2 / 2a<1b1 / 3a<2a2L & EE/D/D & {-}{-}/{-}/{-} & E & C:1b,2a,3a \\
C0676 & 1a<o1 1b<o2 / 2a<1b1L / 3a<2a0L & EE/D/N & {-}{-}/{-}/{-} & E & C:1b,2a \\
C0677 & 1a<o1 1b<o2 / 2a<1b1L / 3a<2a1 & EE/D/D & {-}{-}/{-}/{-} & E & C:1b,2a,3a \\
C0678 & 1a<o1 1b<o2 / 2a<1b1L / 3a<2a1L & EE/D/D & {-}{-}/{-}/{-} & E & C:1b,2a,3a \\
C0679 & 1a<o1 1b<o2 / 2a<1b1L / 3a<2a2 & EE/D/D & {-}{-}/{-}/{-} & E & C:1b,2a,3a \\
C0680 & 1a<o1 1b<o2 / 2a<1b1L / 3a<2a2L & EE/D/D & {-}{-}/{-}/{-} & E & C:1b,2a,3a \\
C0681 & 1a<o1 1b<o2 / 2a<1b2 / 3a<2a0L & EE/D/N & {-}{-}/{-}/{-} & E & C:1b,2a \\
C0682 & 1a<o1 1b<o2 / 2a<1b2 / 3a<2a1 & EE/D/D & {-}{-}/{-}/{-} & E & C:1b,2a,3a \\
C0683 & 1a<o1 1b<o2 / 2a<1b2 / 3a<2a1L & EE/D/D & {-}{-}/{-}/{-} & E & C:1b,2a,3a \\
C0684 & 1a<o1 1b<o2 / 2a<1b2 / 3a<2a2 & EE/D/D & {-}{-}/{-}/{-} & E & C:1b,2a,3a \\
C0685 & 1a<o1 1b<o2 / 2a<1b2 / 3a<2a2L & EE/D/D & {-}{-}/{-}/{-} & E & C:1b,2a,3a \\
C0686 & 1a<o1 1b<o2 / 2a<1b2L / 3a<2a0L & EE/D/N & {-}{-}/{-}/{-} & E & C:1b,2a \\
C0687 & 1a<o1 1b<o2 / 2a<1b2L / 3a<2a1 & EE/D/D & {-}{-}/{-}/{-} & E & C:1b,2a,3a \\
C0688 & 1a<o1 1b<o2 / 2a<1b2L / 3a<2a1L & EE/D/D & {-}{-}/{-}/{-} & E & C:1b,2a,3a \\
C0689 & 1a<o1 1b<o2 / 2a<1b2L / 3a<2a2 & EE/D/D & {-}{-}/{-}/{-} & E & C:1b,2a,3a \\
C0690 & 1a<o1 1b<o2 / 2a<1b2L / 3a<2a2L & EE/D/D & {-}{-}/{-}/{-} & E & C:1b,2a,3a \\
C0691 & 1a<o1 1b<o2L / 2a<1a0L / 3a<2a0L & LI/N/N & {-}+/{-}/{-} & B & PW:1b \\
C0692 & 1a<o1 1b<o2L / 2a<1a0L / 3a<2a1 & LI/N/N & {-}+/{-}/{-} & B & PW:1b \\
C0693 & 1a<o1 1b<o2L / 2a<1a0L / 3a<2a1L & LI/N/N & {-}+/{-}/{-} & B & PW:1b \\
C0694 & 1a<o1 1b<o2L / 2a<1a0L / 3a<2a2 & LI/N/N & {-}+/{-}/{-} & B & PW:1b \\
C0695 & 1a<o1 1b<o2L / 2a<1a0L / 3a<2a2L & LI/N/N & {-}+/{-}/{-} & B & PW:1b \\
C0696 & 1a<o1 1b<o2L / 2a<1a1 / 3a<2a0L & LI/D/N & {-}+/{-}/{-} & B & PW:1b \\
C0697 & 1a<o1 1b<o2L / 2a<1a1 / 3a<2a1 & LI/D/D & {-}+/{-}/{-} & B & PW:1b \\
C0698 & 1a<o1 1b<o2L / 2a<1a1 / 3a<2a1L & LI/D/D & {-}+/{-}/{-} & B & PW:1b \\
C0699 & 1a<o1 1b<o2L / 2a<1a1 / 3a<2a2 & LI/D/D & {-}+/{-}/{-} & B & PW:1b \\
C0700 & 1a<o1 1b<o2L / 2a<1a1 / 3a<2a2L & LI/D/D & {-}+/{-}/{-} & B & PW:1b \\
C0701 & 1a<o1 1b<o2L / 2a<1a1L / 3a<2a0L & LI/D/N & {-}+/{-}/{-} & B & PW:1b \\
C0702 & 1a<o1 1b<o2L / 2a<1a1L / 3a<2a1 & LI/D/D & {-}+/{-}/{-} & B & PW:1b \\
C0703 & 1a<o1 1b<o2L / 2a<1a1L / 3a<2a1L & LI/D/D & {-}+/{-}/{-} & B & PW:1b \\
C0704 & 1a<o1 1b<o2L / 2a<1a1L / 3a<2a2 & LI/D/D & {-}+/{-}/{-} & B & PW:1b \\
C0705 & 1a<o1 1b<o2L / 2a<1a1L / 3a<2a2L & LI/D/D & {-}+/{-}/{-} & B & PW:1b \\
C0706 & 1a<o1 1b<o2L / 2a<1a2 / 3a<2a0L & LI/D/N & {-}+/{-}/{-} & B & PW:1b \\
C0707 & 1a<o1 1b<o2L / 2a<1a2 / 3a<2a1 & LI/D/D & {-}+/{-}/{-} & B & PW:1b \\
C0708 & 1a<o1 1b<o2L / 2a<1a2 / 3a<2a1L & LI/D/D & {-}+/{-}/{-} & B & PW:1b \\
C0709 & 1a<o1 1b<o2L / 2a<1a2 / 3a<2a2 & LI/D/D & {-}+/{-}/{-} & B & PW:1b \\
C0710 & 1a<o1 1b<o2L / 2a<1a2 / 3a<2a2L & LI/D/D & {-}+/{-}/{-} & B & PW:1b \\
C0711 & 1a<o1 1b<o2L / 2a<1a2L / 3a<2a0L & LI/D/N & {-}+/{-}/{-} & B & PW:1b \\
C0712 & 1a<o1 1b<o2L / 2a<1a2L / 3a<2a1 & LI/D/D & {-}+/{-}/{-} & B & PW:1b \\
C0713 & 1a<o1 1b<o2L / 2a<1a2L / 3a<2a1L & LI/D/D & {-}+/{-}/{-} & B & PW:1b \\
C0714 & 1a<o1 1b<o2L / 2a<1a2L / 3a<2a2 & LI/D/D & {-}+/{-}/{-} & B & PW:1b \\
C0715 & 1a<o1 1b<o2L / 2a<1a2L / 3a<2a2L & LI/D/D & {-}+/{-}/{-} & B & PW:1b \\
C0716 & 1a<o1 1b<o2L / 2a<1b0L / 3a<2a0L & LI/N/N & {-}+/+/+ & B & PW:1b \\
C0717 & 1a<o1 1b<o2L / 2a<1b0L / 3a<2a1 & LI/N/N & {-}+/+/{-} & B & PW:1b \\
C0718 & 1a<o1 1b<o2L / 2a<1b0L / 3a<2a1L & LI/N/N & {-}+/+/+ & B & PW:1b \\
C0719 & 1a<o1 1b<o2L / 2a<1b0L / 3a<2a2 & LI/N/N & {-}+/+/{-} & B & PW:1b \\
C0720 & 1a<o1 1b<o2L / 2a<1b0L / 3a<2a2L & LI/N/N & {-}+/+/+ & B & PW:1b \\
C0721 & 1a<o1 1b<o2L / 2a<1b1 / 3a<2a0L & LI/I/N & {-}+/{-}/{-} & B & PW:1b,2a \\
C0722 & 1a<o1 1b<o2L / 2a<1b1 / 3a<2a1 & LI/I/I & {-}+/{-}/{-} & B & PW:1b,2a,3a \\
C0723 & 1a<o1 1b<o2L / 2a<1b1 / 3a<2a1L & LI/I/I & {-}+/{-}/{-} & B & PW:1b,2a,3a \\
C0724 & 1a<o1 1b<o2L / 2a<1b1 / 3a<2a2 & LI/I/I & {-}+/{-}/{-} & B & PW:1b,2a,3a \\
C0725 & 1a<o1 1b<o2L / 2a<1b1 / 3a<2a2L & LI/I/I & {-}+/{-}/{-} & B & PW:1b,2a,3a \\
C0726 & 1a<o1 1b<o2L / 2a<1b1L / 3a<2a0L & LI/I/N & {-}+/+/+ & B & PW:1b,2a \\
C0727 & 1a<o1 1b<o2L / 2a<1b1L / 3a<2a1 & LI/I/I & {-}+/+/{-} & B & PW:1b,2a,3a \\
C0728 & 1a<o1 1b<o2L / 2a<1b1L / 3a<2a1L & LI/I/I & {-}+/+/+ & B & PW:1b,2a,3a \\
C0729 & 1a<o1 1b<o2L / 2a<1b1L / 3a<2a2 & LI/I/I & {-}+/+/{-} & B & PW:1b,2a,3a \\
C0730 & 1a<o1 1b<o2L / 2a<1b1L / 3a<2a2L & LI/I/I & {-}+/+/+ & B & PW:1b,2a,3a \\
C0731 & 1a<o1 1b<o2L / 2a<1b2 / 3a<2a0L & LI/I/N & {-}+/{-}/{-} & B & PW:1b,2a \\
C0732 & 1a<o1 1b<o2L / 2a<1b2 / 3a<2a1 & LI/I/I & {-}+/{-}/{-} & B & PW:1b,2a,3a \\
C0733 & 1a<o1 1b<o2L / 2a<1b2 / 3a<2a1L & LI/I/I & {-}+/{-}/{-} & B & PW:1b,2a,3a \\
C0734 & 1a<o1 1b<o2L / 2a<1b2 / 3a<2a2 & LI/I/I & {-}+/{-}/{-} & B & PW:1b,2a,3a \\
C0735 & 1a<o1 1b<o2L / 2a<1b2 / 3a<2a2L & LI/I/I & {-}+/{-}/{-} & B & PW:1b,2a,3a \\
C0736 & 1a<o1 1b<o2L / 2a<1b2L / 3a<2a0L & LI/I/N & {-}+/+/+ & B & PW:1b,2a \\
C0737 & 1a<o1 1b<o2L / 2a<1b2L / 3a<2a1 & LI/I/I & {-}+/+/{-} & B & PW:1b,2a,3a \\
C0738 & 1a<o1 1b<o2L / 2a<1b2L / 3a<2a1L & LI/I/I & {-}+/+/+ & B & PW:1b,2a,3a \\
C0739 & 1a<o1 1b<o2L / 2a<1b2L / 3a<2a2 & LI/I/I & {-}+/+/{-} & B & PW:1b,2a,3a \\
C0740 & 1a<o1 1b<o2L / 2a<1b2L / 3a<2a2L & LI/I/I & {-}+/+/+ & B & PW:1b,2a,3a \\
C0741 & 1a<o1L / 2a<1a0L / 3a<2a0L 3b<2a1 & I/N/NN & +/+/+{-} & N &  \\
C0742 & 1a<o1L / 2a<1a0L / 3a<2a0L 3b<2a2 & I/N/NN & +/+/+{-} & N &  \\
C0743 & 1a<o1L / 2a<1a0L / 3a<2a1 3b<2a1 & I/N/NN & +/+/{-}{-} & N &  \\
C0744 & 1a<o1L / 2a<1a0L / 3a<2a1 3b<2a1L & I/N/NN & +/+/{-}+ & N &  \\
C0745 & 1a<o1L / 2a<1a0L / 3a<2a1 3b<2a2 & I/N/NN & +/+/{-}{-} & N &  \\
C0746 & 1a<o1L / 2a<1a0L / 3a<2a1 3b<2a2L & I/N/NN & +/+/{-}+ & N &  \\
C0747 & 1a<o1L / 2a<1a0L / 3a<2a1L 3b<2a2 & I/N/NN & +/+/+{-} & N &  \\
C0748 & 1a<o1L / 2a<1a0L / 3a<2a2 3b<2a2 & I/N/NN & +/+/{-}{-} & N &  \\
C0749 & 1a<o1L / 2a<1a0L / 3a<2a2 3b<2a2L & I/N/NN & +/+/{-}+ & N &  \\
C0750 & 1a<o1L / 2a<1a0L 2b<1a1 / 3a<2a0L & I/NI/N & +/+{-}/+ & N &  \\
C0751 & 1a<o1L / 2a<1a0L 2b<1a1 / 3a<2a1 & I/NI/N & +/+{-}/{-} & N &  \\
C0752 & 1a<o1L / 2a<1a0L 2b<1a1 / 3a<2a1L & I/NI/N & +/+{-}/+ & N &  \\
C0753 & 1a<o1L / 2a<1a0L 2b<1a1 / 3a<2a2 & I/NI/N & +/+{-}/{-} & N &  \\
C0754 & 1a<o1L / 2a<1a0L 2b<1a1 / 3a<2a2L & I/NI/N & +/+{-}/+ & N &  \\
C0755 & 1a<o1L / 2a<1a0L 2b<1a1 / 3a<2b0L & I/NI/N & +/+{-}/{-} & N &  \\
C0756 & 1a<o1L / 2a<1a0L 2b<1a1 / 3a<2b1 & I/NI/I & +/+{-}/{-} & N &  \\
C0757 & 1a<o1L / 2a<1a0L 2b<1a1 / 3a<2b1L & I/NI/I & +/+{-}/{-} & N &  \\
C0758 & 1a<o1L / 2a<1a0L 2b<1a1 / 3a<2b2 & I/NI/I & +/+{-}/{-} & N &  \\
C0759 & 1a<o1L / 2a<1a0L 2b<1a1 / 3a<2b2L & I/NI/I & +/+{-}/{-} & N &  \\
C0760 & 1a<o1L / 2a<1a0L 2b<1a2 / 3a<2a0L & I/NI/N & +/+{-}/+ & N &  \\
C0761 & 1a<o1L / 2a<1a0L 2b<1a2 / 3a<2a1 & I/NI/N & +/+{-}/{-} & N &  \\
C0762 & 1a<o1L / 2a<1a0L 2b<1a2 / 3a<2a1L & I/NI/N & +/+{-}/+ & N &  \\
C0763 & 1a<o1L / 2a<1a0L 2b<1a2 / 3a<2a2 & I/NI/N & +/+{-}/{-} & N &  \\
C0764 & 1a<o1L / 2a<1a0L 2b<1a2 / 3a<2a2L & I/NI/N & +/+{-}/+ & N &  \\
C0765 & 1a<o1L / 2a<1a0L 2b<1a2 / 3a<2b0L & I/NI/N & +/+{-}/{-} & N &  \\
C0766 & 1a<o1L / 2a<1a0L 2b<1a2 / 3a<2b1 & I/NI/I & +/+{-}/{-} & N &  \\
C0767 & 1a<o1L / 2a<1a0L 2b<1a2 / 3a<2b1L & I/NI/I & +/+{-}/{-} & N &  \\
C0768 & 1a<o1L / 2a<1a0L 2b<1a2 / 3a<2b2 & I/NI/I & +/+{-}/{-} & N &  \\
C0769 & 1a<o1L / 2a<1a0L 2b<1a2 / 3a<2b2L & I/NI/I & +/+{-}/{-} & N &  \\
C0770 & 1a<o1L / 2a<1a1 / 3a<2a0L 3b<2a1 & I/I/NI & +/{-}/{-}{-} & N &  \\
C0771 & 1a<o1L / 2a<1a1 / 3a<2a0L 3b<2a2 & I/I/NI & +/{-}/{-}{-} & N &  \\
C0772 & 1a<o1L / 2a<1a1 / 3a<2a1 3b<2a1 & I/I/EE & +/{-}/{-}{-} & E & W:1a,2a \\
C0773 & 1a<o1L / 2a<1a1 / 3a<2a1 3b<2a1L & I/I/LI & +/{-}/{-}{-} & B & CP:3b W:1a,2a,3b \\
C0774 & 1a<o1L / 2a<1a1 / 3a<2a1 3b<2a2 & I/I/EE & +/{-}/{-}{-} & E & C:3b W:1a,2a \\
C0775 & 1a<o1L / 2a<1a1 / 3a<2a1 3b<2a2L & I/I/LI & +/{-}/{-}{-} & B & P:3b W:1a,2a,3b \\
C0776 & 1a<o1L / 2a<1a1 / 3a<2a1L 3b<2a2 & I/I/IL & +/{-}/{-}{-} & B & C:3a,3b P:3a W:1a,2a,3a \\
C0777 & 1a<o1L / 2a<1a1 / 3a<2a2 3b<2a2 & I/I/EE & +/{-}/{-}{-} & E & W:1a,2a \\
C0778 & 1a<o1L / 2a<1a1 / 3a<2a2 3b<2a2L & I/I/LI & +/{-}/{-}{-} & B & CP:3b W:1a,2a,3b \\
C0779 & 1a<o1L / 2a<1a1 2b<1a1 / 3a<2a0L & I/EE/N & +/{-}{-}/{-} & E & W:1a \\
C0780 & 1a<o1L / 2a<1a1 2b<1a1 / 3a<2a1 & I/EE/D & +/{-}{-}/{-} & E & W:1a \\
C0781 & 1a<o1L / 2a<1a1 2b<1a1 / 3a<2a1L & I/EE/D & +/{-}{-}/{-} & E & W:1a \\
C0782 & 1a<o1L / 2a<1a1 2b<1a1 / 3a<2a2 & I/EE/D & +/{-}{-}/{-} & E & W:1a \\
C0783 & 1a<o1L / 2a<1a1 2b<1a1 / 3a<2a2L & I/EE/D & +/{-}{-}/{-} & E & W:1a \\
C0784 & 1a<o1L / 2a<1a1 2b<1a1L / 3a<2a0L & I/LI/N & +/{-}+/{-} & B & CP:2b W:1a,2b \\
C0785 & 1a<o1L / 2a<1a1 2b<1a1L / 3a<2a1 & I/LI/D & +/{-}+/{-} & B & CP:2b W:1a,2b \\
C0786 & 1a<o1L / 2a<1a1 2b<1a1L / 3a<2a1L & I/LI/D & +/{-}+/{-} & B & CP:2b W:1a,2b \\
C0787 & 1a<o1L / 2a<1a1 2b<1a1L / 3a<2a2 & I/LI/D & +/{-}+/{-} & B & CP:2b W:1a,2b \\
C0788 & 1a<o1L / 2a<1a1 2b<1a1L / 3a<2a2L & I/LI/D & +/{-}+/{-} & B & CP:2b W:1a,2b \\
C0789 & 1a<o1L / 2a<1a1 2b<1a1L / 3a<2b0L & I/LI/N & +/{-}+/+ & B & CP:2b W:1a,2b \\
C0790 & 1a<o1L / 2a<1a1 2b<1a1L / 3a<2b1 & I/LI/I & +/{-}+/{-} & B & CP:2b,3a W:1a,2b,3a \\
C0791 & 1a<o1L / 2a<1a1 2b<1a1L / 3a<2b1L & I/LI/I & +/{-}+/+ & B & CP:2b,3a W:1a,2b,3a \\
C0792 & 1a<o1L / 2a<1a1 2b<1a1L / 3a<2b2 & I/LI/I & +/{-}+/{-} & B & CP:2b,3a W:1a,2b,3a \\
C0793 & 1a<o1L / 2a<1a1 2b<1a1L / 3a<2b2L & I/LI/I & +/{-}+/+ & B & CP:2b,3a W:1a,2b,3a \\
C0794 & 1a<o1L / 2a<1a1 2b<1a2 / 3a<2a0L & I/EE/N & +/{-}{-}/{-} & E & C:2b W:1a \\
C0795 & 1a<o1L / 2a<1a1 2b<1a2 / 3a<2a1 & I/EE/D & +/{-}{-}/{-} & E & C:2b W:1a \\
C0796 & 1a<o1L / 2a<1a1 2b<1a2 / 3a<2a1L & I/EE/D & +/{-}{-}/{-} & E & C:2b W:1a \\
C0797 & 1a<o1L / 2a<1a1 2b<1a2 / 3a<2a2 & I/EE/D & +/{-}{-}/{-} & E & C:2b W:1a \\
C0798 & 1a<o1L / 2a<1a1 2b<1a2 / 3a<2a2L & I/EE/D & +/{-}{-}/{-} & E & C:2b W:1a \\
C0799 & 1a<o1L / 2a<1a1 2b<1a2 / 3a<2b0L & I/EE/N & +/{-}{-}/{-} & E & C:2b W:1a \\
C0800 & 1a<o1L / 2a<1a1 2b<1a2 / 3a<2b1 & I/EE/D & +/{-}{-}/{-} & E & C:2b,3a W:1a \\
C0801 & 1a<o1L / 2a<1a1 2b<1a2 / 3a<2b1L & I/EE/D & +/{-}{-}/{-} & E & C:2b,3a W:1a \\
C0802 & 1a<o1L / 2a<1a1 2b<1a2 / 3a<2b2 & I/EE/D & +/{-}{-}/{-} & E & C:2b,3a W:1a \\
C0803 & 1a<o1L / 2a<1a1 2b<1a2 / 3a<2b2L & I/EE/D & +/{-}{-}/{-} & E & C:2b,3a W:1a \\
C0804 & 1a<o1L / 2a<1a1 2b<1a2L / 3a<2a0L & I/LI/N & +/{-}+/{-} & B & P:2b W:1a,2b \\
C0805 & 1a<o1L / 2a<1a1 2b<1a2L / 3a<2a1 & I/LI/D & +/{-}+/{-} & B & P:2b W:1a,2b \\
C0806 & 1a<o1L / 2a<1a1 2b<1a2L / 3a<2a1L & I/LI/D & +/{-}+/{-} & B & P:2b W:1a,2b \\
C0807 & 1a<o1L / 2a<1a1 2b<1a2L / 3a<2a2 & I/LI/D & +/{-}+/{-} & B & P:2b W:1a,2b \\
C0808 & 1a<o1L / 2a<1a1 2b<1a2L / 3a<2a2L & I/LI/D & +/{-}+/{-} & B & P:2b W:1a,2b \\
C0809 & 1a<o1L / 2a<1a1 2b<1a2L / 3a<2b0L & I/LI/N & +/{-}+/+ & B & P:2b W:1a,2b \\
C0810 & 1a<o1L / 2a<1a1 2b<1a2L / 3a<2b1 & I/LI/I & +/{-}+/{-} & B & P:2b,3a W:1a,2b,3a \\
C0811 & 1a<o1L / 2a<1a1 2b<1a2L / 3a<2b1L & I/LI/I & +/{-}+/+ & B & P:2b,3a W:1a,2b,3a \\
C0812 & 1a<o1L / 2a<1a1 2b<1a2L / 3a<2b2 & I/LI/I & +/{-}+/{-} & B & P:2b,3a W:1a,2b,3a \\
C0813 & 1a<o1L / 2a<1a1 2b<1a2L / 3a<2b2L & I/LI/I & +/{-}+/+ & B & P:2b,3a W:1a,2b,3a \\
C0814 & 1a<o1L / 2a<1a1L / 3a<2a0L 3b<2a1 & I/I/NI & +/+/+{-} & N &  \\
C0815 & 1a<o1L / 2a<1a1L / 3a<2a0L 3b<2a2 & I/I/NI & +/+/+{-} & N &  \\
C0816 & 1a<o1L / 2a<1a1L / 3a<2a1 3b<2a1 & I/I/EE & +/+/{-}{-} & E & W:1a,2a \\
C0817 & 1a<o1L / 2a<1a1L / 3a<2a1 3b<2a1L & I/I/LI & +/+/{-}+ & B & CP:3b W:1a,2a,3b \\
C0818 & 1a<o1L / 2a<1a1L / 3a<2a1 3b<2a2 & I/I/EE & +/+/{-}{-} & E & C:3b W:1a,2a \\
C0819 & 1a<o1L / 2a<1a1L / 3a<2a1 3b<2a2L & I/I/LI & +/+/{-}+ & B & P:3b W:1a,2a,3b \\
C0820 & 1a<o1L / 2a<1a1L / 3a<2a1L 3b<2a2 & I/I/IL & +/+/+{-} & B & C:3a,3b P:3a W:1a,2a,3a \\
C0821 & 1a<o1L / 2a<1a1L / 3a<2a2 3b<2a2 & I/I/EE & +/+/{-}{-} & E & W:1a,2a \\
C0822 & 1a<o1L / 2a<1a1L / 3a<2a2 3b<2a2L & I/I/LI & +/+/{-}+ & B & CP:3b W:1a,2a,3b \\
C0823 & 1a<o1L / 2a<1a1L 2b<1a2 / 3a<2a0L & I/IL/N & +/+{-}/+ & B & C:2a,2b P:2a W:1a,2a \\
C0824 & 1a<o1L / 2a<1a1L 2b<1a2 / 3a<2a1 & I/IL/I & +/+{-}/{-} & B & C:2a,2b,3a P:2a,3a W:1a,2a,3a \\
C0825 & 1a<o1L / 2a<1a1L 2b<1a2 / 3a<2a1L & I/IL/I & +/+{-}/+ & B & C:2a,2b,3a P:2a,3a W:1a,2a,3a \\
C0826 & 1a<o1L / 2a<1a1L 2b<1a2 / 3a<2a2 & I/IL/I & +/+{-}/{-} & B & C:2a,2b,3a P:2a,3a W:1a,2a,3a \\
C0827 & 1a<o1L / 2a<1a1L 2b<1a2 / 3a<2a2L & I/IL/I & +/+{-}/+ & B & C:2a,2b,3a P:2a,3a W:1a,2a,3a \\
C0828 & 1a<o1L / 2a<1a1L 2b<1a2 / 3a<2b0L & I/IL/N & +/+{-}/{-} & B & C:2a,2b P:2a W:1a,2a \\
C0829 & 1a<o1L / 2a<1a1L 2b<1a2 / 3a<2b1 & I/IL/D & +/+{-}/{-} & B & C:2a,2b,3a P:2a W:1a,2a \\
C0830 & 1a<o1L / 2a<1a1L 2b<1a2 / 3a<2b1L & I/IL/D & +/+{-}/{-} & B & C:2a,2b,3a P:2a W:1a,2a \\
C0831 & 1a<o1L / 2a<1a1L 2b<1a2 / 3a<2b2 & I/IL/D & +/+{-}/{-} & B & C:2a,2b,3a P:2a W:1a,2a \\
C0832 & 1a<o1L / 2a<1a1L 2b<1a2 / 3a<2b2L & I/IL/D & +/+{-}/{-} & B & C:2a,2b,3a P:2a W:1a,2a \\
C0833 & 1a<o1L / 2a<1a2 / 3a<2a0L 3b<2a1 & I/I/NI & +/{-}/{-}{-} & N &  \\
C0834 & 1a<o1L / 2a<1a2 / 3a<2a0L 3b<2a2 & I/I/NI & +/{-}/{-}{-} & N &  \\
C0835 & 1a<o1L / 2a<1a2 / 3a<2a1 3b<2a1 & I/I/EE & +/{-}/{-}{-} & E & W:1a,2a \\
C0836 & 1a<o1L / 2a<1a2 / 3a<2a1 3b<2a1L & I/I/LI & +/{-}/{-}{-} & B & CP:3b W:1a,2a,3b \\
C0837 & 1a<o1L / 2a<1a2 / 3a<2a1 3b<2a2 & I/I/EE & +/{-}/{-}{-} & E & C:3b W:1a,2a \\
C0838 & 1a<o1L / 2a<1a2 / 3a<2a1 3b<2a2L & I/I/LI & +/{-}/{-}{-} & B & P:3b W:1a,2a,3b \\
C0839 & 1a<o1L / 2a<1a2 / 3a<2a1L 3b<2a2 & I/I/IL & +/{-}/{-}{-} & B & C:3a,3b P:3a W:1a,2a,3a \\
C0840 & 1a<o1L / 2a<1a2 / 3a<2a2 3b<2a2 & I/I/EE & +/{-}/{-}{-} & E & W:1a,2a \\
C0841 & 1a<o1L / 2a<1a2 / 3a<2a2 3b<2a2L & I/I/LI & +/{-}/{-}{-} & B & CP:3b W:1a,2a,3b \\
C0842 & 1a<o1L / 2a<1a2 2b<1a2 / 3a<2a0L & I/EE/N & +/{-}{-}/{-} & E & W:1a \\
C0843 & 1a<o1L / 2a<1a2 2b<1a2 / 3a<2a1 & I/EE/D & +/{-}{-}/{-} & E & W:1a \\
C0844 & 1a<o1L / 2a<1a2 2b<1a2 / 3a<2a1L & I/EE/D & +/{-}{-}/{-} & E & W:1a \\
C0845 & 1a<o1L / 2a<1a2 2b<1a2 / 3a<2a2 & I/EE/D & +/{-}{-}/{-} & E & W:1a \\
C0846 & 1a<o1L / 2a<1a2 2b<1a2 / 3a<2a2L & I/EE/D & +/{-}{-}/{-} & E & W:1a \\
C0847 & 1a<o1L / 2a<1a2 2b<1a2L / 3a<2a0L & I/LI/N & +/{-}+/{-} & B & CP:2b W:1a,2b \\
C0848 & 1a<o1L / 2a<1a2 2b<1a2L / 3a<2a1 & I/LI/D & +/{-}+/{-} & B & CP:2b W:1a,2b \\
C0849 & 1a<o1L / 2a<1a2 2b<1a2L / 3a<2a1L & I/LI/D & +/{-}+/{-} & B & CP:2b W:1a,2b \\
C0850 & 1a<o1L / 2a<1a2 2b<1a2L / 3a<2a2 & I/LI/D & +/{-}+/{-} & B & CP:2b W:1a,2b \\
C0851 & 1a<o1L / 2a<1a2 2b<1a2L / 3a<2a2L & I/LI/D & +/{-}+/{-} & B & CP:2b W:1a,2b \\
C0852 & 1a<o1L / 2a<1a2 2b<1a2L / 3a<2b0L & I/LI/N & +/{-}+/+ & B & CP:2b W:1a,2b \\
C0853 & 1a<o1L / 2a<1a2 2b<1a2L / 3a<2b1 & I/LI/I & +/{-}+/{-} & B & CP:2b,3a W:1a,2b,3a \\
C0854 & 1a<o1L / 2a<1a2 2b<1a2L / 3a<2b1L & I/LI/I & +/{-}+/+ & B & CP:2b,3a W:1a,2b,3a \\
C0855 & 1a<o1L / 2a<1a2 2b<1a2L / 3a<2b2 & I/LI/I & +/{-}+/{-} & B & CP:2b,3a W:1a,2b,3a \\
C0856 & 1a<o1L / 2a<1a2 2b<1a2L / 3a<2b2L & I/LI/I & +/{-}+/+ & B & CP:2b,3a W:1a,2b,3a \\
C0857 & 1a<o1L / 2a<1a2L / 3a<2a0L 3b<2a1 & I/I/NI & +/+/+{-} & N &  \\
C0858 & 1a<o1L / 2a<1a2L / 3a<2a0L 3b<2a2 & I/I/NI & +/+/+{-} & N &  \\
C0859 & 1a<o1L / 2a<1a2L / 3a<2a1 3b<2a1 & I/I/EE & +/+/{-}{-} & E & W:1a,2a \\
C0860 & 1a<o1L / 2a<1a2L / 3a<2a1 3b<2a1L & I/I/LI & +/+/{-}+ & B & CP:3b W:1a,2a,3b \\
C0861 & 1a<o1L / 2a<1a2L / 3a<2a1 3b<2a2 & I/I/EE & +/+/{-}{-} & E & C:3b W:1a,2a \\
C0862 & 1a<o1L / 2a<1a2L / 3a<2a1 3b<2a2L & I/I/LI & +/+/{-}+ & B & P:3b W:1a,2a,3b \\
C0863 & 1a<o1L / 2a<1a2L / 3a<2a1L 3b<2a2 & I/I/IL & +/+/+{-} & B & C:3a,3b P:3a W:1a,2a,3a \\
C0864 & 1a<o1L / 2a<1a2L / 3a<2a2 3b<2a2 & I/I/EE & +/+/{-}{-} & E & W:1a,2a \\
C0865 & 1a<o1L / 2a<1a2L / 3a<2a2 3b<2a2L & I/I/LI & +/+/{-}+ & B & CP:3b W:1a,2a,3b \\
C0866 & 1a<o1L 1b<o2 / 2a<1a0L / 3a<2a0L & IL/N/N & +{-}/+/+ & B & C:1a,1b PW:1a \\
C0867 & 1a<o1L 1b<o2 / 2a<1a0L / 3a<2a1 & IL/N/N & +{-}/+/{-} & B & C:1a,1b PW:1a \\
C0868 & 1a<o1L 1b<o2 / 2a<1a0L / 3a<2a1L & IL/N/N & +{-}/+/+ & B & C:1a,1b PW:1a \\
C0869 & 1a<o1L 1b<o2 / 2a<1a0L / 3a<2a2 & IL/N/N & +{-}/+/{-} & B & C:1a,1b PW:1a \\
C0870 & 1a<o1L 1b<o2 / 2a<1a0L / 3a<2a2L & IL/N/N & +{-}/+/+ & B & C:1a,1b PW:1a \\
C0871 & 1a<o1L 1b<o2 / 2a<1a1 / 3a<2a0L & IL/I/N & +{-}/{-}/{-} & B & C:1a,1b,2a PW:1a,2a \\
C0872 & 1a<o1L 1b<o2 / 2a<1a1 / 3a<2a1 & IL/I/I & +{-}/{-}/{-} & B & C:1a,1b,2a,3a PW:1a,2a,3a \\
C0873 & 1a<o1L 1b<o2 / 2a<1a1 / 3a<2a1L & IL/I/I & +{-}/{-}/{-} & B & C:1a,1b,2a,3a PW:1a,2a,3a \\
C0874 & 1a<o1L 1b<o2 / 2a<1a1 / 3a<2a2 & IL/I/I & +{-}/{-}/{-} & B & C:1a,1b,2a,3a PW:1a,2a,3a \\
C0875 & 1a<o1L 1b<o2 / 2a<1a1 / 3a<2a2L & IL/I/I & +{-}/{-}/{-} & B & C:1a,1b,2a,3a PW:1a,2a,3a \\
C0876 & 1a<o1L 1b<o2 / 2a<1a1L / 3a<2a0L & IL/I/N & +{-}/+/+ & B & C:1a,1b,2a PW:1a,2a \\
C0877 & 1a<o1L 1b<o2 / 2a<1a1L / 3a<2a1 & IL/I/I & +{-}/+/{-} & B & C:1a,1b,2a,3a PW:1a,2a,3a \\
C0878 & 1a<o1L 1b<o2 / 2a<1a1L / 3a<2a1L & IL/I/I & +{-}/+/+ & B & C:1a,1b,2a,3a PW:1a,2a,3a \\
C0879 & 1a<o1L 1b<o2 / 2a<1a1L / 3a<2a2 & IL/I/I & +{-}/+/{-} & B & C:1a,1b,2a,3a PW:1a,2a,3a \\
C0880 & 1a<o1L 1b<o2 / 2a<1a1L / 3a<2a2L & IL/I/I & +{-}/+/+ & B & C:1a,1b,2a,3a PW:1a,2a,3a \\
C0881 & 1a<o1L 1b<o2 / 2a<1a2 / 3a<2a0L & IL/I/N & +{-}/{-}/{-} & B & C:1a,1b,2a PW:1a,2a \\
C0882 & 1a<o1L 1b<o2 / 2a<1a2 / 3a<2a1 & IL/I/I & +{-}/{-}/{-} & B & C:1a,1b,2a,3a PW:1a,2a,3a \\
C0883 & 1a<o1L 1b<o2 / 2a<1a2 / 3a<2a1L & IL/I/I & +{-}/{-}/{-} & B & C:1a,1b,2a,3a PW:1a,2a,3a \\
C0884 & 1a<o1L 1b<o2 / 2a<1a2 / 3a<2a2 & IL/I/I & +{-}/{-}/{-} & B & C:1a,1b,2a,3a PW:1a,2a,3a \\
C0885 & 1a<o1L 1b<o2 / 2a<1a2 / 3a<2a2L & IL/I/I & +{-}/{-}/{-} & B & C:1a,1b,2a,3a PW:1a,2a,3a \\
C0886 & 1a<o1L 1b<o2 / 2a<1a2L / 3a<2a0L & IL/I/N & +{-}/+/+ & B & C:1a,1b,2a PW:1a,2a \\
C0887 & 1a<o1L 1b<o2 / 2a<1a2L / 3a<2a1 & IL/I/I & +{-}/+/{-} & B & C:1a,1b,2a,3a PW:1a,2a,3a \\
C0888 & 1a<o1L 1b<o2 / 2a<1a2L / 3a<2a1L & IL/I/I & +{-}/+/+ & B & C:1a,1b,2a,3a PW:1a,2a,3a \\
C0889 & 1a<o1L 1b<o2 / 2a<1a2L / 3a<2a2 & IL/I/I & +{-}/+/{-} & B & C:1a,1b,2a,3a PW:1a,2a,3a \\
C0890 & 1a<o1L 1b<o2 / 2a<1a2L / 3a<2a2L & IL/I/I & +{-}/+/+ & B & C:1a,1b,2a,3a PW:1a,2a,3a \\
C0891 & 1a<o1L 1b<o2 / 2a<1b0L / 3a<2a0L & IL/N/N & +{-}/{-}/{-} & B & C:1a,1b PW:1a \\
C0892 & 1a<o1L 1b<o2 / 2a<1b0L / 3a<2a1 & IL/N/N & +{-}/{-}/{-} & B & C:1a,1b PW:1a \\
C0893 & 1a<o1L 1b<o2 / 2a<1b0L / 3a<2a1L & IL/N/N & +{-}/{-}/{-} & B & C:1a,1b PW:1a \\
C0894 & 1a<o1L 1b<o2 / 2a<1b0L / 3a<2a2 & IL/N/N & +{-}/{-}/{-} & B & C:1a,1b PW:1a \\
C0895 & 1a<o1L 1b<o2 / 2a<1b0L / 3a<2a2L & IL/N/N & +{-}/{-}/{-} & B & C:1a,1b PW:1a \\
C0896 & 1a<o1L 1b<o2 / 2a<1b1 / 3a<2a0L & IL/D/N & +{-}/{-}/{-} & B & C:1a,1b,2a PW:1a \\
C0897 & 1a<o1L 1b<o2 / 2a<1b1 / 3a<2a1 & IL/D/D & +{-}/{-}/{-} & B & C:1a,1b,2a,3a PW:1a \\
C0898 & 1a<o1L 1b<o2 / 2a<1b1 / 3a<2a1L & IL/D/D & +{-}/{-}/{-} & B & C:1a,1b,2a,3a PW:1a \\
C0899 & 1a<o1L 1b<o2 / 2a<1b1 / 3a<2a2 & IL/D/D & +{-}/{-}/{-} & B & C:1a,1b,2a,3a PW:1a \\
C0900 & 1a<o1L 1b<o2 / 2a<1b1 / 3a<2a2L & IL/D/D & +{-}/{-}/{-} & B & C:1a,1b,2a,3a PW:1a \\
C0901 & 1a<o1L 1b<o2 / 2a<1b1L / 3a<2a0L & IL/D/N & +{-}/{-}/{-} & B & C:1a,1b,2a PW:1a \\
C0902 & 1a<o1L 1b<o2 / 2a<1b1L / 3a<2a1 & IL/D/D & +{-}/{-}/{-} & B & C:1a,1b,2a,3a PW:1a \\
C0903 & 1a<o1L 1b<o2 / 2a<1b1L / 3a<2a1L & IL/D/D & +{-}/{-}/{-} & B & C:1a,1b,2a,3a PW:1a \\
C0904 & 1a<o1L 1b<o2 / 2a<1b1L / 3a<2a2 & IL/D/D & +{-}/{-}/{-} & B & C:1a,1b,2a,3a PW:1a \\
C0905 & 1a<o1L 1b<o2 / 2a<1b1L / 3a<2a2L & IL/D/D & +{-}/{-}/{-} & B & C:1a,1b,2a,3a PW:1a \\
C0906 & 1a<o1L 1b<o2 / 2a<1b2 / 3a<2a0L & IL/D/N & +{-}/{-}/{-} & B & C:1a,1b,2a PW:1a \\
C0907 & 1a<o1L 1b<o2 / 2a<1b2 / 3a<2a1 & IL/D/D & +{-}/{-}/{-} & B & C:1a,1b,2a,3a PW:1a \\
C0908 & 1a<o1L 1b<o2 / 2a<1b2 / 3a<2a1L & IL/D/D & +{-}/{-}/{-} & B & C:1a,1b,2a,3a PW:1a \\
C0909 & 1a<o1L 1b<o2 / 2a<1b2 / 3a<2a2 & IL/D/D & +{-}/{-}/{-} & B & C:1a,1b,2a,3a PW:1a \\
C0910 & 1a<o1L 1b<o2 / 2a<1b2 / 3a<2a2L & IL/D/D & +{-}/{-}/{-} & B & C:1a,1b,2a,3a PW:1a \\
C0911 & 1a<o1L 1b<o2 / 2a<1b2L / 3a<2a0L & IL/D/N & +{-}/{-}/{-} & B & C:1a,1b,2a PW:1a \\
C0912 & 1a<o1L 1b<o2 / 2a<1b2L / 3a<2a1 & IL/D/D & +{-}/{-}/{-} & B & C:1a,1b,2a,3a PW:1a \\
C0913 & 1a<o1L 1b<o2 / 2a<1b2L / 3a<2a1L & IL/D/D & +{-}/{-}/{-} & B & C:1a,1b,2a,3a PW:1a \\
C0914 & 1a<o1L 1b<o2 / 2a<1b2L / 3a<2a2 & IL/D/D & +{-}/{-}/{-} & B & C:1a,1b,2a,3a PW:1a \\
C0915 & 1a<o1L 1b<o2 / 2a<1b2L / 3a<2a2L & IL/D/D & +{-}/{-}/{-} & B & C:1a,1b,2a,3a PW:1a \\
C0916 & 1a<o2 / 2a<1a0L / 3a<2a0L 3b<2a1 & I/N/NN & {-}/{-}/{-}{-} & N &  \\
C0917 & 1a<o2 / 2a<1a0L / 3a<2a0L 3b<2a2 & I/N/NN & {-}/{-}/{-}{-} & N &  \\
C0918 & 1a<o2 / 2a<1a0L / 3a<2a1 3b<2a1 & I/N/NN & {-}/{-}/{-}{-} & N &  \\
C0919 & 1a<o2 / 2a<1a0L / 3a<2a1 3b<2a1L & I/N/NN & {-}/{-}/{-}{-} & N &  \\
C0920 & 1a<o2 / 2a<1a0L / 3a<2a1 3b<2a2 & I/N/NN & {-}/{-}/{-}{-} & N &  \\
C0921 & 1a<o2 / 2a<1a0L / 3a<2a1 3b<2a2L & I/N/NN & {-}/{-}/{-}{-} & N &  \\
C0922 & 1a<o2 / 2a<1a0L / 3a<2a1L 3b<2a2 & I/N/NN & {-}/{-}/{-}{-} & N &  \\
C0923 & 1a<o2 / 2a<1a0L / 3a<2a2 3b<2a2 & I/N/NN & {-}/{-}/{-}{-} & N &  \\
C0924 & 1a<o2 / 2a<1a0L / 3a<2a2 3b<2a2L & I/N/NN & {-}/{-}/{-}{-} & N &  \\
C0925 & 1a<o2 / 2a<1a0L 2b<1a1 / 3a<2a0L & I/NI/N & {-}/{-}{-}/{-} & N &  \\
C0926 & 1a<o2 / 2a<1a0L 2b<1a1 / 3a<2a1 & I/NI/N & {-}/{-}{-}/{-} & N &  \\
C0927 & 1a<o2 / 2a<1a0L 2b<1a1 / 3a<2a1L & I/NI/N & {-}/{-}{-}/{-} & N &  \\
C0928 & 1a<o2 / 2a<1a0L 2b<1a1 / 3a<2a2 & I/NI/N & {-}/{-}{-}/{-} & N &  \\
C0929 & 1a<o2 / 2a<1a0L 2b<1a1 / 3a<2a2L & I/NI/N & {-}/{-}{-}/{-} & N &  \\
C0930 & 1a<o2 / 2a<1a0L 2b<1a1 / 3a<2b0L & I/NI/N & {-}/{-}{-}/{-} & N &  \\
C0931 & 1a<o2 / 2a<1a0L 2b<1a1 / 3a<2b1 & I/NI/I & {-}/{-}{-}/{-} & N &  \\
C0932 & 1a<o2 / 2a<1a0L 2b<1a1 / 3a<2b1L & I/NI/I & {-}/{-}{-}/{-} & N &  \\
C0933 & 1a<o2 / 2a<1a0L 2b<1a1 / 3a<2b2 & I/NI/I & {-}/{-}{-}/{-} & N &  \\
C0934 & 1a<o2 / 2a<1a0L 2b<1a1 / 3a<2b2L & I/NI/I & {-}/{-}{-}/{-} & N &  \\
C0935 & 1a<o2 / 2a<1a0L 2b<1a2 / 3a<2a0L & I/NI/N & {-}/{-}{-}/{-} & N &  \\
C0936 & 1a<o2 / 2a<1a0L 2b<1a2 / 3a<2a1 & I/NI/N & {-}/{-}{-}/{-} & N &  \\
C0937 & 1a<o2 / 2a<1a0L 2b<1a2 / 3a<2a1L & I/NI/N & {-}/{-}{-}/{-} & N &  \\
C0938 & 1a<o2 / 2a<1a0L 2b<1a2 / 3a<2a2 & I/NI/N & {-}/{-}{-}/{-} & N &  \\
C0939 & 1a<o2 / 2a<1a0L 2b<1a2 / 3a<2a2L & I/NI/N & {-}/{-}{-}/{-} & N &  \\
C0940 & 1a<o2 / 2a<1a0L 2b<1a2 / 3a<2b0L & I/NI/N & {-}/{-}{-}/{-} & N &  \\
C0941 & 1a<o2 / 2a<1a0L 2b<1a2 / 3a<2b1 & I/NI/I & {-}/{-}{-}/{-} & N &  \\
C0942 & 1a<o2 / 2a<1a0L 2b<1a2 / 3a<2b1L & I/NI/I & {-}/{-}{-}/{-} & N &  \\
C0943 & 1a<o2 / 2a<1a0L 2b<1a2 / 3a<2b2 & I/NI/I & {-}/{-}{-}/{-} & N &  \\
C0944 & 1a<o2 / 2a<1a0L 2b<1a2 / 3a<2b2L & I/NI/I & {-}/{-}{-}/{-} & N &  \\
C0945 & 1a<o2 / 2a<1a1 / 3a<2a0L 3b<2a1 & I/I/NI & {-}/{-}/{-}{-} & N &  \\
C0946 & 1a<o2 / 2a<1a1 / 3a<2a0L 3b<2a2 & I/I/NI & {-}/{-}/{-}{-} & N &  \\
C0947 & 1a<o2 / 2a<1a1 / 3a<2a1 3b<2a1 & I/I/EE & {-}/{-}/{-}{-} & E & W:1a,2a \\
C0948 & 1a<o2 / 2a<1a1 / 3a<2a1 3b<2a1L & I/I/LI & {-}/{-}/{-}{-} & B & CP:3b W:1a,2a,3b \\
C0949 & 1a<o2 / 2a<1a1 / 3a<2a1 3b<2a2 & I/I/EE & {-}/{-}/{-}{-} & E & C:3b W:1a,2a \\
C0950 & 1a<o2 / 2a<1a1 / 3a<2a1 3b<2a2L & I/I/LI & {-}/{-}/{-}{-} & B & P:3b W:1a,2a,3b \\
C0951 & 1a<o2 / 2a<1a1 / 3a<2a1L 3b<2a2 & I/I/IL & {-}/{-}/{-}{-} & B & C:3a,3b P:3a W:1a,2a,3a \\
C0952 & 1a<o2 / 2a<1a1 / 3a<2a2 3b<2a2 & I/I/EE & {-}/{-}/{-}{-} & E & W:1a,2a \\
C0953 & 1a<o2 / 2a<1a1 / 3a<2a2 3b<2a2L & I/I/LI & {-}/{-}/{-}{-} & B & CP:3b W:1a,2a,3b \\
C0954 & 1a<o2 / 2a<1a1 2b<1a1 / 3a<2a0L & I/EE/N & {-}/{-}{-}/{-} & E & W:1a \\
C0955 & 1a<o2 / 2a<1a1 2b<1a1 / 3a<2a1 & I/EE/D & {-}/{-}{-}/{-} & E & W:1a \\
C0956 & 1a<o2 / 2a<1a1 2b<1a1 / 3a<2a1L & I/EE/D & {-}/{-}{-}/{-} & E & W:1a \\
C0957 & 1a<o2 / 2a<1a1 2b<1a1 / 3a<2a2 & I/EE/D & {-}/{-}{-}/{-} & E & W:1a \\
C0958 & 1a<o2 / 2a<1a1 2b<1a1 / 3a<2a2L & I/EE/D & {-}/{-}{-}/{-} & E & W:1a \\
C0959 & 1a<o2 / 2a<1a1 2b<1a1L / 3a<2a0L & I/LI/N & {-}/{-}{-}/{-} & B & CP:2b W:1a,2b \\
C0960 & 1a<o2 / 2a<1a1 2b<1a1L / 3a<2a1 & I/LI/D & {-}/{-}{-}/{-} & B & CP:2b W:1a,2b \\
C0961 & 1a<o2 / 2a<1a1 2b<1a1L / 3a<2a1L & I/LI/D & {-}/{-}{-}/{-} & B & CP:2b W:1a,2b \\
C0962 & 1a<o2 / 2a<1a1 2b<1a1L / 3a<2a2 & I/LI/D & {-}/{-}{-}/{-} & B & CP:2b W:1a,2b \\
C0963 & 1a<o2 / 2a<1a1 2b<1a1L / 3a<2a2L & I/LI/D & {-}/{-}{-}/{-} & B & CP:2b W:1a,2b \\
C0964 & 1a<o2 / 2a<1a1 2b<1a1L / 3a<2b0L & I/LI/N & {-}/{-}{-}/{-} & B & CP:2b W:1a,2b \\
C0965 & 1a<o2 / 2a<1a1 2b<1a1L / 3a<2b1 & I/LI/I & {-}/{-}{-}/{-} & B & CP:2b,3a W:1a,2b,3a \\
C0966 & 1a<o2 / 2a<1a1 2b<1a1L / 3a<2b1L & I/LI/I & {-}/{-}{-}/{-} & B & CP:2b,3a W:1a,2b,3a \\
C0967 & 1a<o2 / 2a<1a1 2b<1a1L / 3a<2b2 & I/LI/I & {-}/{-}{-}/{-} & B & CP:2b,3a W:1a,2b,3a \\
C0968 & 1a<o2 / 2a<1a1 2b<1a1L / 3a<2b2L & I/LI/I & {-}/{-}{-}/{-} & B & CP:2b,3a W:1a,2b,3a \\
C0969 & 1a<o2 / 2a<1a1 2b<1a2 / 3a<2a0L & I/EE/N & {-}/{-}{-}/{-} & E & C:2b W:1a \\
C0970 & 1a<o2 / 2a<1a1 2b<1a2 / 3a<2a1 & I/EE/D & {-}/{-}{-}/{-} & E & C:2b W:1a \\
C0971 & 1a<o2 / 2a<1a1 2b<1a2 / 3a<2a1L & I/EE/D & {-}/{-}{-}/{-} & E & C:2b W:1a \\
C0972 & 1a<o2 / 2a<1a1 2b<1a2 / 3a<2a2 & I/EE/D & {-}/{-}{-}/{-} & E & C:2b W:1a \\
C0973 & 1a<o2 / 2a<1a1 2b<1a2 / 3a<2a2L & I/EE/D & {-}/{-}{-}/{-} & E & C:2b W:1a \\
C0974 & 1a<o2 / 2a<1a1 2b<1a2 / 3a<2b0L & I/EE/N & {-}/{-}{-}/{-} & E & C:2b W:1a \\
C0975 & 1a<o2 / 2a<1a1 2b<1a2 / 3a<2b1 & I/EE/D & {-}/{-}{-}/{-} & E & C:2b,3a W:1a \\
C0976 & 1a<o2 / 2a<1a1 2b<1a2 / 3a<2b1L & I/EE/D & {-}/{-}{-}/{-} & E & C:2b,3a W:1a \\
C0977 & 1a<o2 / 2a<1a1 2b<1a2 / 3a<2b2 & I/EE/D & {-}/{-}{-}/{-} & E & C:2b,3a W:1a \\
C0978 & 1a<o2 / 2a<1a1 2b<1a2 / 3a<2b2L & I/EE/D & {-}/{-}{-}/{-} & E & C:2b,3a W:1a \\
C0979 & 1a<o2 / 2a<1a1 2b<1a2L / 3a<2a0L & I/LI/N & {-}/{-}{-}/{-} & B & P:2b W:1a,2b \\
C0980 & 1a<o2 / 2a<1a1 2b<1a2L / 3a<2a1 & I/LI/D & {-}/{-}{-}/{-} & B & P:2b W:1a,2b \\
C0981 & 1a<o2 / 2a<1a1 2b<1a2L / 3a<2a1L & I/LI/D & {-}/{-}{-}/{-} & B & P:2b W:1a,2b \\
C0982 & 1a<o2 / 2a<1a1 2b<1a2L / 3a<2a2 & I/LI/D & {-}/{-}{-}/{-} & B & P:2b W:1a,2b \\
C0983 & 1a<o2 / 2a<1a1 2b<1a2L / 3a<2a2L & I/LI/D & {-}/{-}{-}/{-} & B & P:2b W:1a,2b \\
C0984 & 1a<o2 / 2a<1a1 2b<1a2L / 3a<2b0L & I/LI/N & {-}/{-}{-}/{-} & B & P:2b W:1a,2b \\
C0985 & 1a<o2 / 2a<1a1 2b<1a2L / 3a<2b1 & I/LI/I & {-}/{-}{-}/{-} & B & P:2b,3a W:1a,2b,3a \\
C0986 & 1a<o2 / 2a<1a1 2b<1a2L / 3a<2b1L & I/LI/I & {-}/{-}{-}/{-} & B & P:2b,3a W:1a,2b,3a \\
C0987 & 1a<o2 / 2a<1a1 2b<1a2L / 3a<2b2 & I/LI/I & {-}/{-}{-}/{-} & B & P:2b,3a W:1a,2b,3a \\
C0988 & 1a<o2 / 2a<1a1 2b<1a2L / 3a<2b2L & I/LI/I & {-}/{-}{-}/{-} & B & P:2b,3a W:1a,2b,3a \\
C0989 & 1a<o2 / 2a<1a1L / 3a<2a0L 3b<2a1 & I/I/NI & {-}/{-}/{-}{-} & N &  \\
C0990 & 1a<o2 / 2a<1a1L / 3a<2a0L 3b<2a2 & I/I/NI & {-}/{-}/{-}{-} & N &  \\
C0991 & 1a<o2 / 2a<1a1L / 3a<2a1 3b<2a1 & I/I/EE & {-}/{-}/{-}{-} & E & W:1a,2a \\
C0992 & 1a<o2 / 2a<1a1L / 3a<2a1 3b<2a1L & I/I/LI & {-}/{-}/{-}{-} & B & CP:3b W:1a,2a,3b \\
C0993 & 1a<o2 / 2a<1a1L / 3a<2a1 3b<2a2 & I/I/EE & {-}/{-}/{-}{-} & E & C:3b W:1a,2a \\
C0994 & 1a<o2 / 2a<1a1L / 3a<2a1 3b<2a2L & I/I/LI & {-}/{-}/{-}{-} & B & P:3b W:1a,2a,3b \\
C0995 & 1a<o2 / 2a<1a1L / 3a<2a1L 3b<2a2 & I/I/IL & {-}/{-}/{-}{-} & B & C:3a,3b P:3a W:1a,2a,3a \\
C0996 & 1a<o2 / 2a<1a1L / 3a<2a2 3b<2a2 & I/I/EE & {-}/{-}/{-}{-} & E & W:1a,2a \\
C0997 & 1a<o2 / 2a<1a1L / 3a<2a2 3b<2a2L & I/I/LI & {-}/{-}/{-}{-} & B & CP:3b W:1a,2a,3b \\
C0998 & 1a<o2 / 2a<1a1L 2b<1a2 / 3a<2a0L & I/IL/N & {-}/{-}{-}/{-} & B & C:2a,2b P:2a W:1a,2a \\
C0999 & 1a<o2 / 2a<1a1L 2b<1a2 / 3a<2a1 & I/IL/I & {-}/{-}{-}/{-} & B & C:2a,2b,3a P:2a,3a W:1a,2a,3a \\
C1000 & 1a<o2 / 2a<1a1L 2b<1a2 / 3a<2a1L & I/IL/I & {-}/{-}{-}/{-} & B & C:2a,2b,3a P:2a,3a W:1a,2a,3a \\
C1001 & 1a<o2 / 2a<1a1L 2b<1a2 / 3a<2a2 & I/IL/I & {-}/{-}{-}/{-} & B & C:2a,2b,3a P:2a,3a W:1a,2a,3a \\
C1002 & 1a<o2 / 2a<1a1L 2b<1a2 / 3a<2a2L & I/IL/I & {-}/{-}{-}/{-} & B & C:2a,2b,3a P:2a,3a W:1a,2a,3a \\
C1003 & 1a<o2 / 2a<1a1L 2b<1a2 / 3a<2b0L & I/IL/N & {-}/{-}{-}/{-} & B & C:2a,2b P:2a W:1a,2a \\
C1004 & 1a<o2 / 2a<1a1L 2b<1a2 / 3a<2b1 & I/IL/D & {-}/{-}{-}/{-} & B & C:2a,2b,3a P:2a W:1a,2a \\
C1005 & 1a<o2 / 2a<1a1L 2b<1a2 / 3a<2b1L & I/IL/D & {-}/{-}{-}/{-} & B & C:2a,2b,3a P:2a W:1a,2a \\
C1006 & 1a<o2 / 2a<1a1L 2b<1a2 / 3a<2b2 & I/IL/D & {-}/{-}{-}/{-} & B & C:2a,2b,3a P:2a W:1a,2a \\
C1007 & 1a<o2 / 2a<1a1L 2b<1a2 / 3a<2b2L & I/IL/D & {-}/{-}{-}/{-} & B & C:2a,2b,3a P:2a W:1a,2a \\
C1008 & 1a<o2 / 2a<1a2 / 3a<2a0L 3b<2a1 & I/I/NI & {-}/{-}/{-}{-} & N &  \\
C1009 & 1a<o2 / 2a<1a2 / 3a<2a0L 3b<2a2 & I/I/NI & {-}/{-}/{-}{-} & N &  \\
C1010 & 1a<o2 / 2a<1a2 / 3a<2a1 3b<2a1 & I/I/EE & {-}/{-}/{-}{-} & E & W:1a,2a \\
C1011 & 1a<o2 / 2a<1a2 / 3a<2a1 3b<2a1L & I/I/LI & {-}/{-}/{-}{-} & B & CP:3b W:1a,2a,3b \\
C1012 & 1a<o2 / 2a<1a2 / 3a<2a1 3b<2a2 & I/I/EE & {-}/{-}/{-}{-} & E & C:3b W:1a,2a \\
C1013 & 1a<o2 / 2a<1a2 / 3a<2a1 3b<2a2L & I/I/LI & {-}/{-}/{-}{-} & B & P:3b W:1a,2a,3b \\
C1014 & 1a<o2 / 2a<1a2 / 3a<2a1L 3b<2a2 & I/I/IL & {-}/{-}/{-}{-} & B & C:3a,3b P:3a W:1a,2a,3a \\
C1015 & 1a<o2 / 2a<1a2 / 3a<2a2 3b<2a2 & I/I/EE & {-}/{-}/{-}{-} & E & W:1a,2a \\
C1016 & 1a<o2 / 2a<1a2 / 3a<2a2 3b<2a2L & I/I/LI & {-}/{-}/{-}{-} & B & CP:3b W:1a,2a,3b \\
C1017 & 1a<o2 / 2a<1a2 2b<1a2 / 3a<2a0L & I/EE/N & {-}/{-}{-}/{-} & E & W:1a \\
C1018 & 1a<o2 / 2a<1a2 2b<1a2 / 3a<2a1 & I/EE/D & {-}/{-}{-}/{-} & E & W:1a \\
C1019 & 1a<o2 / 2a<1a2 2b<1a2 / 3a<2a1L & I/EE/D & {-}/{-}{-}/{-} & E & W:1a \\
C1020 & 1a<o2 / 2a<1a2 2b<1a2 / 3a<2a2 & I/EE/D & {-}/{-}{-}/{-} & E & W:1a \\
C1021 & 1a<o2 / 2a<1a2 2b<1a2 / 3a<2a2L & I/EE/D & {-}/{-}{-}/{-} & E & W:1a \\
C1022 & 1a<o2 / 2a<1a2 2b<1a2L / 3a<2a0L & I/LI/N & {-}/{-}{-}/{-} & B & CP:2b W:1a,2b \\
C1023 & 1a<o2 / 2a<1a2 2b<1a2L / 3a<2a1 & I/LI/D & {-}/{-}{-}/{-} & B & CP:2b W:1a,2b \\
C1024 & 1a<o2 / 2a<1a2 2b<1a2L / 3a<2a1L & I/LI/D & {-}/{-}{-}/{-} & B & CP:2b W:1a,2b \\
C1025 & 1a<o2 / 2a<1a2 2b<1a2L / 3a<2a2 & I/LI/D & {-}/{-}{-}/{-} & B & CP:2b W:1a,2b \\
C1026 & 1a<o2 / 2a<1a2 2b<1a2L / 3a<2a2L & I/LI/D & {-}/{-}{-}/{-} & B & CP:2b W:1a,2b \\
C1027 & 1a<o2 / 2a<1a2 2b<1a2L / 3a<2b0L & I/LI/N & {-}/{-}{-}/{-} & B & CP:2b W:1a,2b \\
C1028 & 1a<o2 / 2a<1a2 2b<1a2L / 3a<2b1 & I/LI/I & {-}/{-}{-}/{-} & B & CP:2b,3a W:1a,2b,3a \\
C1029 & 1a<o2 / 2a<1a2 2b<1a2L / 3a<2b1L & I/LI/I & {-}/{-}{-}/{-} & B & CP:2b,3a W:1a,2b,3a \\
C1030 & 1a<o2 / 2a<1a2 2b<1a2L / 3a<2b2 & I/LI/I & {-}/{-}{-}/{-} & B & CP:2b,3a W:1a,2b,3a \\
C1031 & 1a<o2 / 2a<1a2 2b<1a2L / 3a<2b2L & I/LI/I & {-}/{-}{-}/{-} & B & CP:2b,3a W:1a,2b,3a \\
C1032 & 1a<o2 / 2a<1a2L / 3a<2a0L 3b<2a1 & I/I/NI & {-}/{-}/{-}{-} & N &  \\
C1033 & 1a<o2 / 2a<1a2L / 3a<2a0L 3b<2a2 & I/I/NI & {-}/{-}/{-}{-} & N &  \\
C1034 & 1a<o2 / 2a<1a2L / 3a<2a1 3b<2a1 & I/I/EE & {-}/{-}/{-}{-} & E & W:1a,2a \\
C1035 & 1a<o2 / 2a<1a2L / 3a<2a1 3b<2a1L & I/I/LI & {-}/{-}/{-}{-} & B & CP:3b W:1a,2a,3b \\
C1036 & 1a<o2 / 2a<1a2L / 3a<2a1 3b<2a2 & I/I/EE & {-}/{-}/{-}{-} & E & C:3b W:1a,2a \\
C1037 & 1a<o2 / 2a<1a2L / 3a<2a1 3b<2a2L & I/I/LI & {-}/{-}/{-}{-} & B & P:3b W:1a,2a,3b \\
C1038 & 1a<o2 / 2a<1a2L / 3a<2a1L 3b<2a2 & I/I/IL & {-}/{-}/{-}{-} & B & C:3a,3b P:3a W:1a,2a,3a \\
C1039 & 1a<o2 / 2a<1a2L / 3a<2a2 3b<2a2 & I/I/EE & {-}/{-}/{-}{-} & E & W:1a,2a \\
C1040 & 1a<o2 / 2a<1a2L / 3a<2a2 3b<2a2L & I/I/LI & {-}/{-}/{-}{-} & B & CP:3b W:1a,2a,3b \\
C1041 & 1a<o2 1b<o2 / 2a<1a0L / 3a<2a0L & EE/N/N & {-}{-}/{-}/{-} & E &  \\
C1042 & 1a<o2 1b<o2 / 2a<1a0L / 3a<2a1 & EE/N/N & {-}{-}/{-}/{-} & E &  \\
C1043 & 1a<o2 1b<o2 / 2a<1a0L / 3a<2a1L & EE/N/N & {-}{-}/{-}/{-} & E &  \\
C1044 & 1a<o2 1b<o2 / 2a<1a0L / 3a<2a2 & EE/N/N & {-}{-}/{-}/{-} & E &  \\
C1045 & 1a<o2 1b<o2 / 2a<1a0L / 3a<2a2L & EE/N/N & {-}{-}/{-}/{-} & E &  \\
C1046 & 1a<o2 1b<o2 / 2a<1a1 / 3a<2a0L & EE/D/N & {-}{-}/{-}/{-} & E &  \\
C1047 & 1a<o2 1b<o2 / 2a<1a1 / 3a<2a1 & EE/D/D & {-}{-}/{-}/{-} & E &  \\
C1048 & 1a<o2 1b<o2 / 2a<1a1 / 3a<2a1L & EE/D/D & {-}{-}/{-}/{-} & E &  \\
C1049 & 1a<o2 1b<o2 / 2a<1a1 / 3a<2a2 & EE/D/D & {-}{-}/{-}/{-} & E &  \\
C1050 & 1a<o2 1b<o2 / 2a<1a1 / 3a<2a2L & EE/D/D & {-}{-}/{-}/{-} & E &  \\
C1051 & 1a<o2 1b<o2 / 2a<1a1L / 3a<2a0L & EE/D/N & {-}{-}/{-}/{-} & E &  \\
C1052 & 1a<o2 1b<o2 / 2a<1a1L / 3a<2a1 & EE/D/D & {-}{-}/{-}/{-} & E &  \\
C1053 & 1a<o2 1b<o2 / 2a<1a1L / 3a<2a1L & EE/D/D & {-}{-}/{-}/{-} & E &  \\
C1054 & 1a<o2 1b<o2 / 2a<1a1L / 3a<2a2 & EE/D/D & {-}{-}/{-}/{-} & E &  \\
C1055 & 1a<o2 1b<o2 / 2a<1a1L / 3a<2a2L & EE/D/D & {-}{-}/{-}/{-} & E &  \\
C1056 & 1a<o2 1b<o2 / 2a<1a2 / 3a<2a0L & EE/D/N & {-}{-}/{-}/{-} & E &  \\
C1057 & 1a<o2 1b<o2 / 2a<1a2 / 3a<2a1 & EE/D/D & {-}{-}/{-}/{-} & E &  \\
C1058 & 1a<o2 1b<o2 / 2a<1a2 / 3a<2a1L & EE/D/D & {-}{-}/{-}/{-} & E &  \\
C1059 & 1a<o2 1b<o2 / 2a<1a2 / 3a<2a2 & EE/D/D & {-}{-}/{-}/{-} & E &  \\
C1060 & 1a<o2 1b<o2 / 2a<1a2 / 3a<2a2L & EE/D/D & {-}{-}/{-}/{-} & E &  \\
C1061 & 1a<o2 1b<o2 / 2a<1a2L / 3a<2a0L & EE/D/N & {-}{-}/{-}/{-} & E &  \\
C1062 & 1a<o2 1b<o2 / 2a<1a2L / 3a<2a1 & EE/D/D & {-}{-}/{-}/{-} & E &  \\
C1063 & 1a<o2 1b<o2 / 2a<1a2L / 3a<2a1L & EE/D/D & {-}{-}/{-}/{-} & E &  \\
C1064 & 1a<o2 1b<o2 / 2a<1a2L / 3a<2a2 & EE/D/D & {-}{-}/{-}/{-} & E &  \\
C1065 & 1a<o2 1b<o2 / 2a<1a2L / 3a<2a2L & EE/D/D & {-}{-}/{-}/{-} & E &  \\
C1066 & 1a<o2 1b<o2L / 2a<1a0L / 3a<2a0L & LI/N/N & {-}+/{-}/{-} & B & CPW:1b \\
C1067 & 1a<o2 1b<o2L / 2a<1a0L / 3a<2a1 & LI/N/N & {-}+/{-}/{-} & B & CPW:1b \\
C1068 & 1a<o2 1b<o2L / 2a<1a0L / 3a<2a1L & LI/N/N & {-}+/{-}/{-} & B & CPW:1b \\
C1069 & 1a<o2 1b<o2L / 2a<1a0L / 3a<2a2 & LI/N/N & {-}+/{-}/{-} & B & CPW:1b \\
C1070 & 1a<o2 1b<o2L / 2a<1a0L / 3a<2a2L & LI/N/N & {-}+/{-}/{-} & B & CPW:1b \\
C1071 & 1a<o2 1b<o2L / 2a<1a1 / 3a<2a0L & LI/D/N & {-}+/{-}/{-} & B & CPW:1b \\
C1072 & 1a<o2 1b<o2L / 2a<1a1 / 3a<2a1 & LI/D/D & {-}+/{-}/{-} & B & CPW:1b \\
C1073 & 1a<o2 1b<o2L / 2a<1a1 / 3a<2a1L & LI/D/D & {-}+/{-}/{-} & B & CPW:1b \\
C1074 & 1a<o2 1b<o2L / 2a<1a1 / 3a<2a2 & LI/D/D & {-}+/{-}/{-} & B & CPW:1b \\
C1075 & 1a<o2 1b<o2L / 2a<1a1 / 3a<2a2L & LI/D/D & {-}+/{-}/{-} & B & CPW:1b \\
C1076 & 1a<o2 1b<o2L / 2a<1a1L / 3a<2a0L & LI/D/N & {-}+/{-}/{-} & B & CPW:1b \\
C1077 & 1a<o2 1b<o2L / 2a<1a1L / 3a<2a1 & LI/D/D & {-}+/{-}/{-} & B & CPW:1b \\
C1078 & 1a<o2 1b<o2L / 2a<1a1L / 3a<2a1L & LI/D/D & {-}+/{-}/{-} & B & CPW:1b \\
C1079 & 1a<o2 1b<o2L / 2a<1a1L / 3a<2a2 & LI/D/D & {-}+/{-}/{-} & B & CPW:1b \\
C1080 & 1a<o2 1b<o2L / 2a<1a1L / 3a<2a2L & LI/D/D & {-}+/{-}/{-} & B & CPW:1b \\
C1081 & 1a<o2 1b<o2L / 2a<1a2 / 3a<2a0L & LI/D/N & {-}+/{-}/{-} & B & CPW:1b \\
C1082 & 1a<o2 1b<o2L / 2a<1a2 / 3a<2a1 & LI/D/D & {-}+/{-}/{-} & B & CPW:1b \\
C1083 & 1a<o2 1b<o2L / 2a<1a2 / 3a<2a1L & LI/D/D & {-}+/{-}/{-} & B & CPW:1b \\
C1084 & 1a<o2 1b<o2L / 2a<1a2 / 3a<2a2 & LI/D/D & {-}+/{-}/{-} & B & CPW:1b \\
C1085 & 1a<o2 1b<o2L / 2a<1a2 / 3a<2a2L & LI/D/D & {-}+/{-}/{-} & B & CPW:1b \\
C1086 & 1a<o2 1b<o2L / 2a<1a2L / 3a<2a0L & LI/D/N & {-}+/{-}/{-} & B & CPW:1b \\
C1087 & 1a<o2 1b<o2L / 2a<1a2L / 3a<2a1 & LI/D/D & {-}+/{-}/{-} & B & CPW:1b \\
C1088 & 1a<o2 1b<o2L / 2a<1a2L / 3a<2a1L & LI/D/D & {-}+/{-}/{-} & B & CPW:1b \\
C1089 & 1a<o2 1b<o2L / 2a<1a2L / 3a<2a2 & LI/D/D & {-}+/{-}/{-} & B & CPW:1b \\
C1090 & 1a<o2 1b<o2L / 2a<1a2L / 3a<2a2L & LI/D/D & {-}+/{-}/{-} & B & CPW:1b \\
C1091 & 1a<o2 1b<o2L / 2a<1b0L / 3a<2a0L & LI/N/N & {-}+/+/+ & B & CPW:1b \\
C1092 & 1a<o2 1b<o2L / 2a<1b0L / 3a<2a1 & LI/N/N & {-}+/+/{-} & B & CPW:1b \\
C1093 & 1a<o2 1b<o2L / 2a<1b0L / 3a<2a1L & LI/N/N & {-}+/+/+ & B & CPW:1b \\
C1094 & 1a<o2 1b<o2L / 2a<1b0L / 3a<2a2 & LI/N/N & {-}+/+/{-} & B & CPW:1b \\
C1095 & 1a<o2 1b<o2L / 2a<1b0L / 3a<2a2L & LI/N/N & {-}+/+/+ & B & CPW:1b \\
C1096 & 1a<o2 1b<o2L / 2a<1b1 / 3a<2a0L & LI/I/N & {-}+/{-}/{-} & B & CPW:1b,2a \\
C1097 & 1a<o2 1b<o2L / 2a<1b1 / 3a<2a1 & LI/I/I & {-}+/{-}/{-} & B & CPW:1b,2a,3a \\
C1098 & 1a<o2 1b<o2L / 2a<1b1 / 3a<2a1L & LI/I/I & {-}+/{-}/{-} & B & CPW:1b,2a,3a \\
C1099 & 1a<o2 1b<o2L / 2a<1b1 / 3a<2a2 & LI/I/I & {-}+/{-}/{-} & B & CPW:1b,2a,3a \\
C1100 & 1a<o2 1b<o2L / 2a<1b1 / 3a<2a2L & LI/I/I & {-}+/{-}/{-} & B & CPW:1b,2a,3a \\
C1101 & 1a<o2 1b<o2L / 2a<1b1L / 3a<2a0L & LI/I/N & {-}+/+/+ & B & CPW:1b,2a \\
C1102 & 1a<o2 1b<o2L / 2a<1b1L / 3a<2a1 & LI/I/I & {-}+/+/{-} & B & CPW:1b,2a,3a \\
C1103 & 1a<o2 1b<o2L / 2a<1b1L / 3a<2a1L & LI/I/I & {-}+/+/+ & B & CPW:1b,2a,3a \\
C1104 & 1a<o2 1b<o2L / 2a<1b1L / 3a<2a2 & LI/I/I & {-}+/+/{-} & B & CPW:1b,2a,3a \\
C1105 & 1a<o2 1b<o2L / 2a<1b1L / 3a<2a2L & LI/I/I & {-}+/+/+ & B & CPW:1b,2a,3a \\
C1106 & 1a<o2 1b<o2L / 2a<1b2 / 3a<2a0L & LI/I/N & {-}+/{-}/{-} & B & CPW:1b,2a \\
C1107 & 1a<o2 1b<o2L / 2a<1b2 / 3a<2a1 & LI/I/I & {-}+/{-}/{-} & B & CPW:1b,2a,3a \\
C1108 & 1a<o2 1b<o2L / 2a<1b2 / 3a<2a1L & LI/I/I & {-}+/{-}/{-} & B & CPW:1b,2a,3a \\
C1109 & 1a<o2 1b<o2L / 2a<1b2 / 3a<2a2 & LI/I/I & {-}+/{-}/{-} & B & CPW:1b,2a,3a \\
C1110 & 1a<o2 1b<o2L / 2a<1b2 / 3a<2a2L & LI/I/I & {-}+/{-}/{-} & B & CPW:1b,2a,3a \\
C1111 & 1a<o2 1b<o2L / 2a<1b2L / 3a<2a0L & LI/I/N & {-}+/+/+ & B & CPW:1b,2a \\
C1112 & 1a<o2 1b<o2L / 2a<1b2L / 3a<2a1 & LI/I/I & {-}+/+/{-} & B & CPW:1b,2a,3a \\
C1113 & 1a<o2 1b<o2L / 2a<1b2L / 3a<2a1L & LI/I/I & {-}+/+/+ & B & CPW:1b,2a,3a \\
C1114 & 1a<o2 1b<o2L / 2a<1b2L / 3a<2a2 & LI/I/I & {-}+/+/{-} & B & CPW:1b,2a,3a \\
C1115 & 1a<o2 1b<o2L / 2a<1b2L / 3a<2a2L & LI/I/I & {-}+/+/+ & B & CPW:1b,2a,3a \\
C1116 & 1a<o2L / 2a<1a0L / 3a<2a0L 3b<2a1 & I/N/NN & +/+/+{-} & N &  \\
C1117 & 1a<o2L / 2a<1a0L / 3a<2a0L 3b<2a2 & I/N/NN & +/+/+{-} & N &  \\
C1118 & 1a<o2L / 2a<1a0L / 3a<2a1 3b<2a1 & I/N/NN & +/+/{-}{-} & N &  \\
C1119 & 1a<o2L / 2a<1a0L / 3a<2a1 3b<2a1L & I/N/NN & +/+/{-}+ & N &  \\
C1120 & 1a<o2L / 2a<1a0L / 3a<2a1 3b<2a2 & I/N/NN & +/+/{-}{-} & N &  \\
C1121 & 1a<o2L / 2a<1a0L / 3a<2a1 3b<2a2L & I/N/NN & +/+/{-}+ & N &  \\
C1122 & 1a<o2L / 2a<1a0L / 3a<2a1L 3b<2a2 & I/N/NN & +/+/+{-} & N &  \\
C1123 & 1a<o2L / 2a<1a0L / 3a<2a2 3b<2a2 & I/N/NN & +/+/{-}{-} & N &  \\
C1124 & 1a<o2L / 2a<1a0L / 3a<2a2 3b<2a2L & I/N/NN & +/+/{-}+ & N &  \\
C1125 & 1a<o2L / 2a<1a0L 2b<1a1 / 3a<2a0L & I/NI/N & +/+{-}/+ & N &  \\
C1126 & 1a<o2L / 2a<1a0L 2b<1a1 / 3a<2a1 & I/NI/N & +/+{-}/{-} & N &  \\
C1127 & 1a<o2L / 2a<1a0L 2b<1a1 / 3a<2a1L & I/NI/N & +/+{-}/+ & N &  \\
C1128 & 1a<o2L / 2a<1a0L 2b<1a1 / 3a<2a2 & I/NI/N & +/+{-}/{-} & N &  \\
C1129 & 1a<o2L / 2a<1a0L 2b<1a1 / 3a<2a2L & I/NI/N & +/+{-}/+ & N &  \\
C1130 & 1a<o2L / 2a<1a0L 2b<1a1 / 3a<2b0L & I/NI/N & +/+{-}/{-} & N &  \\
C1131 & 1a<o2L / 2a<1a0L 2b<1a1 / 3a<2b1 & I/NI/I & +/+{-}/{-} & N &  \\
C1132 & 1a<o2L / 2a<1a0L 2b<1a1 / 3a<2b1L & I/NI/I & +/+{-}/{-} & N &  \\
C1133 & 1a<o2L / 2a<1a0L 2b<1a1 / 3a<2b2 & I/NI/I & +/+{-}/{-} & N &  \\
C1134 & 1a<o2L / 2a<1a0L 2b<1a1 / 3a<2b2L & I/NI/I & +/+{-}/{-} & N &  \\
C1135 & 1a<o2L / 2a<1a0L 2b<1a2 / 3a<2a0L & I/NI/N & +/+{-}/+ & N &  \\
C1136 & 1a<o2L / 2a<1a0L 2b<1a2 / 3a<2a1 & I/NI/N & +/+{-}/{-} & N &  \\
C1137 & 1a<o2L / 2a<1a0L 2b<1a2 / 3a<2a1L & I/NI/N & +/+{-}/+ & N &  \\
C1138 & 1a<o2L / 2a<1a0L 2b<1a2 / 3a<2a2 & I/NI/N & +/+{-}/{-} & N &  \\
C1139 & 1a<o2L / 2a<1a0L 2b<1a2 / 3a<2a2L & I/NI/N & +/+{-}/+ & N &  \\
C1140 & 1a<o2L / 2a<1a0L 2b<1a2 / 3a<2b0L & I/NI/N & +/+{-}/{-} & N &  \\
C1141 & 1a<o2L / 2a<1a0L 2b<1a2 / 3a<2b1 & I/NI/I & +/+{-}/{-} & N &  \\
C1142 & 1a<o2L / 2a<1a0L 2b<1a2 / 3a<2b1L & I/NI/I & +/+{-}/{-} & N &  \\
C1143 & 1a<o2L / 2a<1a0L 2b<1a2 / 3a<2b2 & I/NI/I & +/+{-}/{-} & N &  \\
C1144 & 1a<o2L / 2a<1a0L 2b<1a2 / 3a<2b2L & I/NI/I & +/+{-}/{-} & N &  \\
C1145 & 1a<o2L / 2a<1a1 / 3a<2a0L 3b<2a1 & I/I/NI & +/{-}/{-}{-} & N &  \\
C1146 & 1a<o2L / 2a<1a1 / 3a<2a0L 3b<2a2 & I/I/NI & +/{-}/{-}{-} & N &  \\
C1147 & 1a<o2L / 2a<1a1 / 3a<2a1 3b<2a1 & I/I/EE & +/{-}/{-}{-} & E & W:1a,2a \\
C1148 & 1a<o2L / 2a<1a1 / 3a<2a1 3b<2a1L & I/I/LI & +/{-}/{-}{-} & B & CP:3b W:1a,2a,3b \\
C1149 & 1a<o2L / 2a<1a1 / 3a<2a1 3b<2a2 & I/I/EE & +/{-}/{-}{-} & E & C:3b W:1a,2a \\
C1150 & 1a<o2L / 2a<1a1 / 3a<2a1 3b<2a2L & I/I/LI & +/{-}/{-}{-} & B & P:3b W:1a,2a,3b \\
C1151 & 1a<o2L / 2a<1a1 / 3a<2a1L 3b<2a2 & I/I/IL & +/{-}/{-}{-} & B & C:3a,3b P:3a W:1a,2a,3a \\
C1152 & 1a<o2L / 2a<1a1 / 3a<2a2 3b<2a2 & I/I/EE & +/{-}/{-}{-} & E & W:1a,2a \\
C1153 & 1a<o2L / 2a<1a1 / 3a<2a2 3b<2a2L & I/I/LI & +/{-}/{-}{-} & B & CP:3b W:1a,2a,3b \\
C1154 & 1a<o2L / 2a<1a1 2b<1a1 / 3a<2a0L & I/EE/N & +/{-}{-}/{-} & E & W:1a \\
C1155 & 1a<o2L / 2a<1a1 2b<1a1 / 3a<2a1 & I/EE/D & +/{-}{-}/{-} & E & W:1a \\
C1156 & 1a<o2L / 2a<1a1 2b<1a1 / 3a<2a1L & I/EE/D & +/{-}{-}/{-} & E & W:1a \\
C1157 & 1a<o2L / 2a<1a1 2b<1a1 / 3a<2a2 & I/EE/D & +/{-}{-}/{-} & E & W:1a \\
C1158 & 1a<o2L / 2a<1a1 2b<1a1 / 3a<2a2L & I/EE/D & +/{-}{-}/{-} & E & W:1a \\
C1159 & 1a<o2L / 2a<1a1 2b<1a1L / 3a<2a0L & I/LI/N & +/{-}+/{-} & B & CP:2b W:1a,2b \\
C1160 & 1a<o2L / 2a<1a1 2b<1a1L / 3a<2a1 & I/LI/D & +/{-}+/{-} & B & CP:2b W:1a,2b \\
C1161 & 1a<o2L / 2a<1a1 2b<1a1L / 3a<2a1L & I/LI/D & +/{-}+/{-} & B & CP:2b W:1a,2b \\
C1162 & 1a<o2L / 2a<1a1 2b<1a1L / 3a<2a2 & I/LI/D & +/{-}+/{-} & B & CP:2b W:1a,2b \\
C1163 & 1a<o2L / 2a<1a1 2b<1a1L / 3a<2a2L & I/LI/D & +/{-}+/{-} & B & CP:2b W:1a,2b \\
C1164 & 1a<o2L / 2a<1a1 2b<1a1L / 3a<2b0L & I/LI/N & +/{-}+/+ & B & CP:2b W:1a,2b \\
C1165 & 1a<o2L / 2a<1a1 2b<1a1L / 3a<2b1 & I/LI/I & +/{-}+/{-} & B & CP:2b,3a W:1a,2b,3a \\
C1166 & 1a<o2L / 2a<1a1 2b<1a1L / 3a<2b1L & I/LI/I & +/{-}+/+ & B & CP:2b,3a W:1a,2b,3a \\
C1167 & 1a<o2L / 2a<1a1 2b<1a1L / 3a<2b2 & I/LI/I & +/{-}+/{-} & B & CP:2b,3a W:1a,2b,3a \\
C1168 & 1a<o2L / 2a<1a1 2b<1a1L / 3a<2b2L & I/LI/I & +/{-}+/+ & B & CP:2b,3a W:1a,2b,3a \\
C1169 & 1a<o2L / 2a<1a1 2b<1a2 / 3a<2a0L & I/EE/N & +/{-}{-}/{-} & E & C:2b W:1a \\
C1170 & 1a<o2L / 2a<1a1 2b<1a2 / 3a<2a1 & I/EE/D & +/{-}{-}/{-} & E & C:2b W:1a \\
C1171 & 1a<o2L / 2a<1a1 2b<1a2 / 3a<2a1L & I/EE/D & +/{-}{-}/{-} & E & C:2b W:1a \\
C1172 & 1a<o2L / 2a<1a1 2b<1a2 / 3a<2a2 & I/EE/D & +/{-}{-}/{-} & E & C:2b W:1a \\
C1173 & 1a<o2L / 2a<1a1 2b<1a2 / 3a<2a2L & I/EE/D & +/{-}{-}/{-} & E & C:2b W:1a \\
C1174 & 1a<o2L / 2a<1a1 2b<1a2 / 3a<2b0L & I/EE/N & +/{-}{-}/{-} & E & C:2b W:1a \\
C1175 & 1a<o2L / 2a<1a1 2b<1a2 / 3a<2b1 & I/EE/D & +/{-}{-}/{-} & E & C:2b,3a W:1a \\
C1176 & 1a<o2L / 2a<1a1 2b<1a2 / 3a<2b1L & I/EE/D & +/{-}{-}/{-} & E & C:2b,3a W:1a \\
C1177 & 1a<o2L / 2a<1a1 2b<1a2 / 3a<2b2 & I/EE/D & +/{-}{-}/{-} & E & C:2b,3a W:1a \\
C1178 & 1a<o2L / 2a<1a1 2b<1a2 / 3a<2b2L & I/EE/D & +/{-}{-}/{-} & E & C:2b,3a W:1a \\
C1179 & 1a<o2L / 2a<1a1 2b<1a2L / 3a<2a0L & I/LI/N & +/{-}+/{-} & B & P:2b W:1a,2b \\
C1180 & 1a<o2L / 2a<1a1 2b<1a2L / 3a<2a1 & I/LI/D & +/{-}+/{-} & B & P:2b W:1a,2b \\
C1181 & 1a<o2L / 2a<1a1 2b<1a2L / 3a<2a1L & I/LI/D & +/{-}+/{-} & B & P:2b W:1a,2b \\
C1182 & 1a<o2L / 2a<1a1 2b<1a2L / 3a<2a2 & I/LI/D & +/{-}+/{-} & B & P:2b W:1a,2b \\
C1183 & 1a<o2L / 2a<1a1 2b<1a2L / 3a<2a2L & I/LI/D & +/{-}+/{-} & B & P:2b W:1a,2b \\
C1184 & 1a<o2L / 2a<1a1 2b<1a2L / 3a<2b0L & I/LI/N & +/{-}+/+ & B & P:2b W:1a,2b \\
C1185 & 1a<o2L / 2a<1a1 2b<1a2L / 3a<2b1 & I/LI/I & +/{-}+/{-} & B & P:2b,3a W:1a,2b,3a \\
C1186 & 1a<o2L / 2a<1a1 2b<1a2L / 3a<2b1L & I/LI/I & +/{-}+/+ & B & P:2b,3a W:1a,2b,3a \\
C1187 & 1a<o2L / 2a<1a1 2b<1a2L / 3a<2b2 & I/LI/I & +/{-}+/{-} & B & P:2b,3a W:1a,2b,3a \\
C1188 & 1a<o2L / 2a<1a1 2b<1a2L / 3a<2b2L & I/LI/I & +/{-}+/+ & B & P:2b,3a W:1a,2b,3a \\
C1189 & 1a<o2L / 2a<1a1L / 3a<2a0L 3b<2a1 & I/I/NI & +/+/+{-} & N &  \\
C1190 & 1a<o2L / 2a<1a1L / 3a<2a0L 3b<2a2 & I/I/NI & +/+/+{-} & N &  \\
C1191 & 1a<o2L / 2a<1a1L / 3a<2a1 3b<2a1 & I/I/EE & +/+/{-}{-} & E & W:1a,2a \\
C1192 & 1a<o2L / 2a<1a1L / 3a<2a1 3b<2a1L & I/I/LI & +/+/{-}+ & B & CP:3b W:1a,2a,3b \\
C1193 & 1a<o2L / 2a<1a1L / 3a<2a1 3b<2a2 & I/I/EE & +/+/{-}{-} & E & C:3b W:1a,2a \\
C1194 & 1a<o2L / 2a<1a1L / 3a<2a1 3b<2a2L & I/I/LI & +/+/{-}+ & B & P:3b W:1a,2a,3b \\
C1195 & 1a<o2L / 2a<1a1L / 3a<2a1L 3b<2a2 & I/I/IL & +/+/+{-} & B & C:3a,3b P:3a W:1a,2a,3a \\
C1196 & 1a<o2L / 2a<1a1L / 3a<2a2 3b<2a2 & I/I/EE & +/+/{-}{-} & E & W:1a,2a \\
C1197 & 1a<o2L / 2a<1a1L / 3a<2a2 3b<2a2L & I/I/LI & +/+/{-}+ & B & CP:3b W:1a,2a,3b \\
C1198 & 1a<o2L / 2a<1a1L 2b<1a2 / 3a<2a0L & I/IL/N & +/+{-}/+ & B & C:2a,2b P:2a W:1a,2a \\
C1199 & 1a<o2L / 2a<1a1L 2b<1a2 / 3a<2a1 & I/IL/I & +/+{-}/{-} & B & C:2a,2b,3a P:2a,3a W:1a,2a,3a \\
C1200 & 1a<o2L / 2a<1a1L 2b<1a2 / 3a<2a1L & I/IL/I & +/+{-}/+ & B & C:2a,2b,3a P:2a,3a W:1a,2a,3a \\
C1201 & 1a<o2L / 2a<1a1L 2b<1a2 / 3a<2a2 & I/IL/I & +/+{-}/{-} & B & C:2a,2b,3a P:2a,3a W:1a,2a,3a \\
C1202 & 1a<o2L / 2a<1a1L 2b<1a2 / 3a<2a2L & I/IL/I & +/+{-}/+ & B & C:2a,2b,3a P:2a,3a W:1a,2a,3a \\
C1203 & 1a<o2L / 2a<1a1L 2b<1a2 / 3a<2b0L & I/IL/N & +/+{-}/{-} & B & C:2a,2b P:2a W:1a,2a \\
C1204 & 1a<o2L / 2a<1a1L 2b<1a2 / 3a<2b1 & I/IL/D & +/+{-}/{-} & B & C:2a,2b,3a P:2a W:1a,2a \\
C1205 & 1a<o2L / 2a<1a1L 2b<1a2 / 3a<2b1L & I/IL/D & +/+{-}/{-} & B & C:2a,2b,3a P:2a W:1a,2a \\
C1206 & 1a<o2L / 2a<1a1L 2b<1a2 / 3a<2b2 & I/IL/D & +/+{-}/{-} & B & C:2a,2b,3a P:2a W:1a,2a \\
C1207 & 1a<o2L / 2a<1a1L 2b<1a2 / 3a<2b2L & I/IL/D & +/+{-}/{-} & B & C:2a,2b,3a P:2a W:1a,2a \\
C1208 & 1a<o2L / 2a<1a2 / 3a<2a0L 3b<2a1 & I/I/NI & +/{-}/{-}{-} & N &  \\
C1209 & 1a<o2L / 2a<1a2 / 3a<2a0L 3b<2a2 & I/I/NI & +/{-}/{-}{-} & N &  \\
C1210 & 1a<o2L / 2a<1a2 / 3a<2a1 3b<2a1 & I/I/EE & +/{-}/{-}{-} & E & W:1a,2a \\
C1211 & 1a<o2L / 2a<1a2 / 3a<2a1 3b<2a1L & I/I/LI & +/{-}/{-}{-} & B & CP:3b W:1a,2a,3b \\
C1212 & 1a<o2L / 2a<1a2 / 3a<2a1 3b<2a2 & I/I/EE & +/{-}/{-}{-} & E & C:3b W:1a,2a \\
C1213 & 1a<o2L / 2a<1a2 / 3a<2a1 3b<2a2L & I/I/LI & +/{-}/{-}{-} & B & P:3b W:1a,2a,3b \\
C1214 & 1a<o2L / 2a<1a2 / 3a<2a1L 3b<2a2 & I/I/IL & +/{-}/{-}{-} & B & C:3a,3b P:3a W:1a,2a,3a \\
C1215 & 1a<o2L / 2a<1a2 / 3a<2a2 3b<2a2 & I/I/EE & +/{-}/{-}{-} & E & W:1a,2a \\
C1216 & 1a<o2L / 2a<1a2 / 3a<2a2 3b<2a2L & I/I/LI & +/{-}/{-}{-} & B & CP:3b W:1a,2a,3b \\
C1217 & 1a<o2L / 2a<1a2 2b<1a2 / 3a<2a0L & I/EE/N & +/{-}{-}/{-} & E & W:1a \\
C1218 & 1a<o2L / 2a<1a2 2b<1a2 / 3a<2a1 & I/EE/D & +/{-}{-}/{-} & E & W:1a \\
C1219 & 1a<o2L / 2a<1a2 2b<1a2 / 3a<2a1L & I/EE/D & +/{-}{-}/{-} & E & W:1a \\
C1220 & 1a<o2L / 2a<1a2 2b<1a2 / 3a<2a2 & I/EE/D & +/{-}{-}/{-} & E & W:1a \\
C1221 & 1a<o2L / 2a<1a2 2b<1a2 / 3a<2a2L & I/EE/D & +/{-}{-}/{-} & E & W:1a \\
C1222 & 1a<o2L / 2a<1a2 2b<1a2L / 3a<2a0L & I/LI/N & +/{-}+/{-} & B & CP:2b W:1a,2b \\
C1223 & 1a<o2L / 2a<1a2 2b<1a2L / 3a<2a1 & I/LI/D & +/{-}+/{-} & B & CP:2b W:1a,2b \\
C1224 & 1a<o2L / 2a<1a2 2b<1a2L / 3a<2a1L & I/LI/D & +/{-}+/{-} & B & CP:2b W:1a,2b \\
C1225 & 1a<o2L / 2a<1a2 2b<1a2L / 3a<2a2 & I/LI/D & +/{-}+/{-} & B & CP:2b W:1a,2b \\
C1226 & 1a<o2L / 2a<1a2 2b<1a2L / 3a<2a2L & I/LI/D & +/{-}+/{-} & B & CP:2b W:1a,2b \\
C1227 & 1a<o2L / 2a<1a2 2b<1a2L / 3a<2b0L & I/LI/N & +/{-}+/+ & B & CP:2b W:1a,2b \\
C1228 & 1a<o2L / 2a<1a2 2b<1a2L / 3a<2b1 & I/LI/I & +/{-}+/{-} & B & CP:2b,3a W:1a,2b,3a \\
C1229 & 1a<o2L / 2a<1a2 2b<1a2L / 3a<2b1L & I/LI/I & +/{-}+/+ & B & CP:2b,3a W:1a,2b,3a \\
C1230 & 1a<o2L / 2a<1a2 2b<1a2L / 3a<2b2 & I/LI/I & +/{-}+/{-} & B & CP:2b,3a W:1a,2b,3a \\
C1231 & 1a<o2L / 2a<1a2 2b<1a2L / 3a<2b2L & I/LI/I & +/{-}+/+ & B & CP:2b,3a W:1a,2b,3a \\
C1232 & 1a<o2L / 2a<1a2L / 3a<2a0L 3b<2a1 & I/I/NI & +/+/+{-} & N &  \\
C1233 & 1a<o2L / 2a<1a2L / 3a<2a0L 3b<2a2 & I/I/NI & +/+/+{-} & N &  \\
C1234 & 1a<o2L / 2a<1a2L / 3a<2a1 3b<2a1 & I/I/EE & +/+/{-}{-} & E & W:1a,2a \\
C1235 & 1a<o2L / 2a<1a2L / 3a<2a1 3b<2a1L & I/I/LI & +/+/{-}+ & B & CP:3b W:1a,2a,3b \\
C1236 & 1a<o2L / 2a<1a2L / 3a<2a1 3b<2a2 & I/I/EE & +/+/{-}{-} & E & C:3b W:1a,2a \\
C1237 & 1a<o2L / 2a<1a2L / 3a<2a1 3b<2a2L & I/I/LI & +/+/{-}+ & B & P:3b W:1a,2a,3b \\
C1238 & 1a<o2L / 2a<1a2L / 3a<2a1L 3b<2a2 & I/I/IL & +/+/+{-} & B & C:3a,3b P:3a W:1a,2a,3a \\
C1239 & 1a<o2L / 2a<1a2L / 3a<2a2 3b<2a2 & I/I/EE & +/+/{-}{-} & E & W:1a,2a \\
C1240 & 1a<o2L / 2a<1a2L / 3a<2a2 3b<2a2L & I/I/LI & +/+/{-}+ & B & CP:3b W:1a,2a,3b \\
C1241 & 1a<o0L / 2a<1a0L 2b<1a1 / 3a<2a0L,2b1 & N/NN/N & +/+{-}/+ & N &  \\
C1242 & 1a<o0L / 2a<1a0L 2b<1a1 / 3a<2a0L,2b2 & N/NN/N & +/+{-}/+ & N &  \\
C1243 & 1a<o0L / 2a<1a0L 2b<1a1 / 3a<2a1,2b0L & N/NN/N & +/+{-}/{-} & N &  \\
C1244 & 1a<o0L / 2a<1a0L 2b<1a1 / 3a<2a1,2b1 & N/NN/N & +/+{-}/{-} & N &  \\
C1245 & 1a<o0L / 2a<1a0L 2b<1a1 / 3a<2a1,2b1L & N/NN/N & +/+{-}/{-} & N &  \\
C1246 & 1a<o0L / 2a<1a0L 2b<1a1 / 3a<2a1,2b2 & N/NN/N & +/+{-}/{-} & N &  \\
C1247 & 1a<o0L / 2a<1a0L 2b<1a1 / 3a<2a1,2b2L & N/NN/N & +/+{-}/{-} & N &  \\
C1248 & 1a<o0L / 2a<1a0L 2b<1a1 / 3a<2a1L,2b1 & N/NN/N & +/+{-}/+ & N &  \\
C1249 & 1a<o0L / 2a<1a0L 2b<1a1 / 3a<2a1L,2b2 & N/NN/N & +/+{-}/+ & N &  \\
C1250 & 1a<o0L / 2a<1a0L 2b<1a1 / 3a<2a2,2b0L & N/NN/N & +/+{-}/{-} & N &  \\
C1251 & 1a<o0L / 2a<1a0L 2b<1a1 / 3a<2a2,2b1 & N/NN/N & +/+{-}/{-} & N &  \\
C1252 & 1a<o0L / 2a<1a0L 2b<1a1 / 3a<2a2,2b1L & N/NN/N & +/+{-}/{-} & N &  \\
C1253 & 1a<o0L / 2a<1a0L 2b<1a1 / 3a<2a2,2b2 & N/NN/N & +/+{-}/{-} & N &  \\
C1254 & 1a<o0L / 2a<1a0L 2b<1a1 / 3a<2a2,2b2L & N/NN/N & +/+{-}/{-} & N &  \\
C1255 & 1a<o0L / 2a<1a0L 2b<1a1 / 3a<2a2L,2b1 & N/NN/N & +/+{-}/+ & N &  \\
C1256 & 1a<o0L / 2a<1a0L 2b<1a1 / 3a<2a2L,2b2 & N/NN/N & +/+{-}/+ & N &  \\
C1257 & 1a<o0L / 2a<1a0L 2b<1a2 / 3a<2a0L,2b1 & N/NN/N & +/+{-}/+ & N &  \\
C1258 & 1a<o0L / 2a<1a0L 2b<1a2 / 3a<2a0L,2b2 & N/NN/N & +/+{-}/+ & N &  \\
C1259 & 1a<o0L / 2a<1a0L 2b<1a2 / 3a<2a1,2b0L & N/NN/N & +/+{-}/{-} & N &  \\
C1260 & 1a<o0L / 2a<1a0L 2b<1a2 / 3a<2a1,2b1 & N/NN/N & +/+{-}/{-} & N &  \\
C1261 & 1a<o0L / 2a<1a0L 2b<1a2 / 3a<2a1,2b1L & N/NN/N & +/+{-}/{-} & N &  \\
C1262 & 1a<o0L / 2a<1a0L 2b<1a2 / 3a<2a1,2b2 & N/NN/N & +/+{-}/{-} & N &  \\
C1263 & 1a<o0L / 2a<1a0L 2b<1a2 / 3a<2a1,2b2L & N/NN/N & +/+{-}/{-} & N &  \\
C1264 & 1a<o0L / 2a<1a0L 2b<1a2 / 3a<2a1L,2b1 & N/NN/N & +/+{-}/+ & N &  \\
C1265 & 1a<o0L / 2a<1a0L 2b<1a2 / 3a<2a1L,2b2 & N/NN/N & +/+{-}/+ & N &  \\
C1266 & 1a<o0L / 2a<1a0L 2b<1a2 / 3a<2a2,2b0L & N/NN/N & +/+{-}/{-} & N &  \\
C1267 & 1a<o0L / 2a<1a0L 2b<1a2 / 3a<2a2,2b1 & N/NN/N & +/+{-}/{-} & N &  \\
C1268 & 1a<o0L / 2a<1a0L 2b<1a2 / 3a<2a2,2b1L & N/NN/N & +/+{-}/{-} & N &  \\
C1269 & 1a<o0L / 2a<1a0L 2b<1a2 / 3a<2a2,2b2 & N/NN/N & +/+{-}/{-} & N &  \\
C1270 & 1a<o0L / 2a<1a0L 2b<1a2 / 3a<2a2,2b2L & N/NN/N & +/+{-}/{-} & N &  \\
C1271 & 1a<o0L / 2a<1a0L 2b<1a2 / 3a<2a2L,2b1 & N/NN/N & +/+{-}/+ & N &  \\
C1272 & 1a<o0L / 2a<1a0L 2b<1a2 / 3a<2a2L,2b2 & N/NN/N & +/+{-}/+ & N &  \\
C1273 & 1a<o0L / 2a<1a1 2b<1a1 / 3a<2a0L,2b1 & N/NN/N & +/{-}{-}/{-} & N &  \\
C1274 & 1a<o0L / 2a<1a1 2b<1a1 / 3a<2a0L,2b2 & N/NN/N & +/{-}{-}/{-} & N &  \\
C1275 & 1a<o0L / 2a<1a1 2b<1a1 / 3a<2a1,2b1 & N/NN/N & +/{-}{-}/{-} & N &  \\
C1276 & 1a<o0L / 2a<1a1 2b<1a1 / 3a<2a1,2b1L & N/NN/N & +/{-}{-}/{-} & N &  \\
C1277 & 1a<o0L / 2a<1a1 2b<1a1 / 3a<2a1,2b2 & N/NN/N & +/{-}{-}/{-} & N &  \\
C1278 & 1a<o0L / 2a<1a1 2b<1a1 / 3a<2a1,2b2L & N/NN/N & +/{-}{-}/{-} & N &  \\
C1279 & 1a<o0L / 2a<1a1 2b<1a1 / 3a<2a1L,2b2 & N/NN/N & +/{-}{-}/{-} & N &  \\
C1280 & 1a<o0L / 2a<1a1 2b<1a1 / 3a<2a2,2b2 & N/NN/N & +/{-}{-}/{-} & N &  \\
C1281 & 1a<o0L / 2a<1a1 2b<1a1 / 3a<2a2,2b2L & N/NN/N & +/{-}{-}/{-} & N &  \\
C1282 & 1a<o0L / 2a<1a1 2b<1a1L / 3a<2a0L,2b1 & N/NN/N & +/{-}+/{-} & N &  \\
C1283 & 1a<o0L / 2a<1a1 2b<1a1L / 3a<2a0L,2b2 & N/NN/N & +/{-}+/{-} & N &  \\
C1284 & 1a<o0L / 2a<1a1 2b<1a1L / 3a<2a1,2b0L & N/NN/N & +/{-}+/+ & N &  \\
C1285 & 1a<o0L / 2a<1a1 2b<1a1L / 3a<2a1,2b1 & N/NN/N & +/{-}+/{-} & N &  \\
C1286 & 1a<o0L / 2a<1a1 2b<1a1L / 3a<2a1,2b1L & N/NN/N & +/{-}+/+ & N &  \\
C1287 & 1a<o0L / 2a<1a1 2b<1a1L / 3a<2a1,2b2 & N/NN/N & +/{-}+/{-} & N &  \\
C1288 & 1a<o0L / 2a<1a1 2b<1a1L / 3a<2a1,2b2L & N/NN/N & +/{-}+/+ & N &  \\
C1289 & 1a<o0L / 2a<1a1 2b<1a1L / 3a<2a1L,2b1 & N/NN/N & +/{-}+/{-} & N &  \\
C1290 & 1a<o0L / 2a<1a1 2b<1a1L / 3a<2a1L,2b2 & N/NN/N & +/{-}+/{-} & N &  \\
C1291 & 1a<o0L / 2a<1a1 2b<1a1L / 3a<2a2,2b0L & N/NN/N & +/{-}+/+ & N &  \\
C1292 & 1a<o0L / 2a<1a1 2b<1a1L / 3a<2a2,2b1 & N/NN/N & +/{-}+/{-} & N &  \\
C1293 & 1a<o0L / 2a<1a1 2b<1a1L / 3a<2a2,2b1L & N/NN/N & +/{-}+/+ & N &  \\
C1294 & 1a<o0L / 2a<1a1 2b<1a1L / 3a<2a2,2b2 & N/NN/N & +/{-}+/{-} & N &  \\
C1295 & 1a<o0L / 2a<1a1 2b<1a1L / 3a<2a2,2b2L & N/NN/N & +/{-}+/+ & N &  \\
C1296 & 1a<o0L / 2a<1a1 2b<1a1L / 3a<2a2L,2b1 & N/NN/N & +/{-}+/{-} & N &  \\
C1297 & 1a<o0L / 2a<1a1 2b<1a1L / 3a<2a2L,2b2 & N/NN/N & +/{-}+/{-} & N &  \\
C1298 & 1a<o0L / 2a<1a1 2b<1a2 / 3a<2a0L,2b1 & N/NN/N & +/{-}{-}/{-} & N &  \\
C1299 & 1a<o0L / 2a<1a1 2b<1a2 / 3a<2a0L,2b2 & N/NN/N & +/{-}{-}/{-} & N &  \\
C1300 & 1a<o0L / 2a<1a1 2b<1a2 / 3a<2a1,2b0L & N/NN/N & +/{-}{-}/{-} & N &  \\
C1301 & 1a<o0L / 2a<1a1 2b<1a2 / 3a<2a1,2b1 & N/NN/N & +/{-}{-}/{-} & N &  \\
C1302 & 1a<o0L / 2a<1a1 2b<1a2 / 3a<2a1,2b1L & N/NN/N & +/{-}{-}/{-} & N &  \\
C1303 & 1a<o0L / 2a<1a1 2b<1a2 / 3a<2a1,2b2 & N/NN/N & +/{-}{-}/{-} & N &  \\
C1304 & 1a<o0L / 2a<1a1 2b<1a2 / 3a<2a1,2b2L & N/NN/N & +/{-}{-}/{-} & N &  \\
C1305 & 1a<o0L / 2a<1a1 2b<1a2 / 3a<2a1L,2b1 & N/NN/N & +/{-}{-}/{-} & N &  \\
C1306 & 1a<o0L / 2a<1a1 2b<1a2 / 3a<2a1L,2b2 & N/NN/N & +/{-}{-}/{-} & N &  \\
C1307 & 1a<o0L / 2a<1a1 2b<1a2 / 3a<2a2,2b0L & N/NN/N & +/{-}{-}/{-} & N &  \\
C1308 & 1a<o0L / 2a<1a1 2b<1a2 / 3a<2a2,2b1 & N/NN/N & +/{-}{-}/{-} & N &  \\
C1309 & 1a<o0L / 2a<1a1 2b<1a2 / 3a<2a2,2b1L & N/NN/N & +/{-}{-}/{-} & N &  \\
C1310 & 1a<o0L / 2a<1a1 2b<1a2 / 3a<2a2,2b2 & N/NN/N & +/{-}{-}/{-} & N &  \\
C1311 & 1a<o0L / 2a<1a1 2b<1a2 / 3a<2a2,2b2L & N/NN/N & +/{-}{-}/{-} & N &  \\
C1312 & 1a<o0L / 2a<1a1 2b<1a2 / 3a<2a2L,2b1 & N/NN/N & +/{-}{-}/{-} & N &  \\
C1313 & 1a<o0L / 2a<1a1 2b<1a2 / 3a<2a2L,2b2 & N/NN/N & +/{-}{-}/{-} & N &  \\
C1314 & 1a<o0L / 2a<1a1 2b<1a2L / 3a<2a0L,2b1 & N/NN/N & +/{-}+/{-} & N &  \\
C1315 & 1a<o0L / 2a<1a1 2b<1a2L / 3a<2a0L,2b2 & N/NN/N & +/{-}+/{-} & N &  \\
C1316 & 1a<o0L / 2a<1a1 2b<1a2L / 3a<2a1,2b0L & N/NN/N & +/{-}+/+ & N &  \\
C1317 & 1a<o0L / 2a<1a1 2b<1a2L / 3a<2a1,2b1 & N/NN/N & +/{-}+/{-} & N &  \\
C1318 & 1a<o0L / 2a<1a1 2b<1a2L / 3a<2a1,2b1L & N/NN/N & +/{-}+/+ & N &  \\
C1319 & 1a<o0L / 2a<1a1 2b<1a2L / 3a<2a1,2b2 & N/NN/N & +/{-}+/{-} & N &  \\
C1320 & 1a<o0L / 2a<1a1 2b<1a2L / 3a<2a1,2b2L & N/NN/N & +/{-}+/+ & N &  \\
C1321 & 1a<o0L / 2a<1a1 2b<1a2L / 3a<2a1L,2b1 & N/NN/N & +/{-}+/{-} & N &  \\
C1322 & 1a<o0L / 2a<1a1 2b<1a2L / 3a<2a1L,2b2 & N/NN/N & +/{-}+/{-} & N &  \\
C1323 & 1a<o0L / 2a<1a1 2b<1a2L / 3a<2a2,2b0L & N/NN/N & +/{-}+/+ & N &  \\
C1324 & 1a<o0L / 2a<1a1 2b<1a2L / 3a<2a2,2b1 & N/NN/N & +/{-}+/{-} & N &  \\
C1325 & 1a<o0L / 2a<1a1 2b<1a2L / 3a<2a2,2b1L & N/NN/N & +/{-}+/+ & N &  \\
C1326 & 1a<o0L / 2a<1a1 2b<1a2L / 3a<2a2,2b2 & N/NN/N & +/{-}+/{-} & N &  \\
C1327 & 1a<o0L / 2a<1a1 2b<1a2L / 3a<2a2,2b2L & N/NN/N & +/{-}+/+ & N &  \\
C1328 & 1a<o0L / 2a<1a1 2b<1a2L / 3a<2a2L,2b1 & N/NN/N & +/{-}+/{-} & N &  \\
C1329 & 1a<o0L / 2a<1a1 2b<1a2L / 3a<2a2L,2b2 & N/NN/N & +/{-}+/{-} & N &  \\
C1330 & 1a<o0L / 2a<1a1L 2b<1a2 / 3a<2a0L,2b1 & N/NN/N & +/+{-}/+ & N &  \\
C1331 & 1a<o0L / 2a<1a1L 2b<1a2 / 3a<2a0L,2b2 & N/NN/N & +/+{-}/+ & N &  \\
C1332 & 1a<o0L / 2a<1a1L 2b<1a2 / 3a<2a1,2b0L & N/NN/N & +/+{-}/{-} & N &  \\
C1333 & 1a<o0L / 2a<1a1L 2b<1a2 / 3a<2a1,2b1 & N/NN/N & +/+{-}/{-} & N &  \\
C1334 & 1a<o0L / 2a<1a1L 2b<1a2 / 3a<2a1,2b1L & N/NN/N & +/+{-}/{-} & N &  \\
C1335 & 1a<o0L / 2a<1a1L 2b<1a2 / 3a<2a1,2b2 & N/NN/N & +/+{-}/{-} & N &  \\
C1336 & 1a<o0L / 2a<1a1L 2b<1a2 / 3a<2a1,2b2L & N/NN/N & +/+{-}/{-} & N &  \\
C1337 & 1a<o0L / 2a<1a1L 2b<1a2 / 3a<2a1L,2b1 & N/NN/N & +/+{-}/+ & N &  \\
C1338 & 1a<o0L / 2a<1a1L 2b<1a2 / 3a<2a1L,2b2 & N/NN/N & +/+{-}/+ & N &  \\
C1339 & 1a<o0L / 2a<1a1L 2b<1a2 / 3a<2a2,2b0L & N/NN/N & +/+{-}/{-} & N &  \\
C1340 & 1a<o0L / 2a<1a1L 2b<1a2 / 3a<2a2,2b1 & N/NN/N & +/+{-}/{-} & N &  \\
C1341 & 1a<o0L / 2a<1a1L 2b<1a2 / 3a<2a2,2b1L & N/NN/N & +/+{-}/{-} & N &  \\
C1342 & 1a<o0L / 2a<1a1L 2b<1a2 / 3a<2a2,2b2 & N/NN/N & +/+{-}/{-} & N &  \\
C1343 & 1a<o0L / 2a<1a1L 2b<1a2 / 3a<2a2,2b2L & N/NN/N & +/+{-}/{-} & N &  \\
C1344 & 1a<o0L / 2a<1a1L 2b<1a2 / 3a<2a2L,2b1 & N/NN/N & +/+{-}/+ & N &  \\
C1345 & 1a<o0L / 2a<1a1L 2b<1a2 / 3a<2a2L,2b2 & N/NN/N & +/+{-}/+ & N &  \\
C1346 & 1a<o0L / 2a<1a2 2b<1a2 / 3a<2a0L,2b1 & N/NN/N & +/{-}{-}/{-} & N &  \\
C1347 & 1a<o0L / 2a<1a2 2b<1a2 / 3a<2a0L,2b2 & N/NN/N & +/{-}{-}/{-} & N &  \\
C1348 & 1a<o0L / 2a<1a2 2b<1a2 / 3a<2a1,2b1 & N/NN/N & +/{-}{-}/{-} & N &  \\
C1349 & 1a<o0L / 2a<1a2 2b<1a2 / 3a<2a1,2b1L & N/NN/N & +/{-}{-}/{-} & N &  \\
C1350 & 1a<o0L / 2a<1a2 2b<1a2 / 3a<2a1,2b2 & N/NN/N & +/{-}{-}/{-} & N &  \\
C1351 & 1a<o0L / 2a<1a2 2b<1a2 / 3a<2a1,2b2L & N/NN/N & +/{-}{-}/{-} & N &  \\
C1352 & 1a<o0L / 2a<1a2 2b<1a2 / 3a<2a1L,2b2 & N/NN/N & +/{-}{-}/{-} & N &  \\
C1353 & 1a<o0L / 2a<1a2 2b<1a2 / 3a<2a2,2b2 & N/NN/N & +/{-}{-}/{-} & N &  \\
C1354 & 1a<o0L / 2a<1a2 2b<1a2 / 3a<2a2,2b2L & N/NN/N & +/{-}{-}/{-} & N &  \\
C1355 & 1a<o0L / 2a<1a2 2b<1a2L / 3a<2a0L,2b1 & N/NN/N & +/{-}+/{-} & N &  \\
C1356 & 1a<o0L / 2a<1a2 2b<1a2L / 3a<2a0L,2b2 & N/NN/N & +/{-}+/{-} & N &  \\
C1357 & 1a<o0L / 2a<1a2 2b<1a2L / 3a<2a1,2b0L & N/NN/N & +/{-}+/+ & N &  \\
C1358 & 1a<o0L / 2a<1a2 2b<1a2L / 3a<2a1,2b1 & N/NN/N & +/{-}+/{-} & N &  \\
C1359 & 1a<o0L / 2a<1a2 2b<1a2L / 3a<2a1,2b1L & N/NN/N & +/{-}+/+ & N &  \\
C1360 & 1a<o0L / 2a<1a2 2b<1a2L / 3a<2a1,2b2 & N/NN/N & +/{-}+/{-} & N &  \\
C1361 & 1a<o0L / 2a<1a2 2b<1a2L / 3a<2a1,2b2L & N/NN/N & +/{-}+/+ & N &  \\
C1362 & 1a<o0L / 2a<1a2 2b<1a2L / 3a<2a1L,2b1 & N/NN/N & +/{-}+/{-} & N &  \\
C1363 & 1a<o0L / 2a<1a2 2b<1a2L / 3a<2a1L,2b2 & N/NN/N & +/{-}+/{-} & N &  \\
C1364 & 1a<o0L / 2a<1a2 2b<1a2L / 3a<2a2,2b0L & N/NN/N & +/{-}+/+ & N &  \\
C1365 & 1a<o0L / 2a<1a2 2b<1a2L / 3a<2a2,2b1 & N/NN/N & +/{-}+/{-} & N &  \\
C1366 & 1a<o0L / 2a<1a2 2b<1a2L / 3a<2a2,2b1L & N/NN/N & +/{-}+/+ & N &  \\
C1367 & 1a<o0L / 2a<1a2 2b<1a2L / 3a<2a2,2b2 & N/NN/N & +/{-}+/{-} & N &  \\
C1368 & 1a<o0L / 2a<1a2 2b<1a2L / 3a<2a2,2b2L & N/NN/N & +/{-}+/+ & N &  \\
C1369 & 1a<o0L / 2a<1a2 2b<1a2L / 3a<2a2L,2b1 & N/NN/N & +/{-}+/{-} & N &  \\
C1370 & 1a<o0L / 2a<1a2 2b<1a2L / 3a<2a2L,2b2 & N/NN/N & +/{-}+/{-} & N &  \\
C1371 & 1a<o0L 1b<o1 / 2a<1a0L,1b1 / 3a<2a0L & NI/I/N & +{-}/+/+ & N &  \\
C1372 & 1a<o0L 1b<o1 / 2a<1a0L,1b1 / 3a<2a1 & NI/I/I & +{-}/+/{-} & N &  \\
C1373 & 1a<o0L 1b<o1 / 2a<1a0L,1b1 / 3a<2a1L & NI/I/I & +{-}/+/+ & N &  \\
C1374 & 1a<o0L 1b<o1 / 2a<1a0L,1b1 / 3a<2a2 & NI/I/I & +{-}/+/{-} & N &  \\
C1375 & 1a<o0L 1b<o1 / 2a<1a0L,1b1 / 3a<2a2L & NI/I/I & +{-}/+/+ & N &  \\
C1376 & 1a<o0L 1b<o1 / 2a<1a0L,1b2 / 3a<2a0L & NI/I/N & +{-}/+/+ & N &  \\
C1377 & 1a<o0L 1b<o1 / 2a<1a0L,1b2 / 3a<2a1 & NI/I/I & +{-}/+/{-} & N &  \\
C1378 & 1a<o0L 1b<o1 / 2a<1a0L,1b2 / 3a<2a1L & NI/I/I & +{-}/+/+ & N &  \\
C1379 & 1a<o0L 1b<o1 / 2a<1a0L,1b2 / 3a<2a2 & NI/I/I & +{-}/+/{-} & N &  \\
C1380 & 1a<o0L 1b<o1 / 2a<1a0L,1b2 / 3a<2a2L & NI/I/I & +{-}/+/+ & N &  \\
C1381 & 1a<o0L 1b<o1 / 2a<1a1,1b0L / 3a<2a0L & NI/N/N & +{-}/{-}/{-} & N &  \\
C1382 & 1a<o0L 1b<o1 / 2a<1a1,1b0L / 3a<2a1 & NI/N/N & +{-}/{-}/{-} & N &  \\
C1383 & 1a<o0L 1b<o1 / 2a<1a1,1b0L / 3a<2a1L & NI/N/N & +{-}/{-}/{-} & N &  \\
C1384 & 1a<o0L 1b<o1 / 2a<1a1,1b0L / 3a<2a2 & NI/N/N & +{-}/{-}/{-} & N &  \\
C1385 & 1a<o0L 1b<o1 / 2a<1a1,1b0L / 3a<2a2L & NI/N/N & +{-}/{-}/{-} & N &  \\
C1386 & 1a<o0L 1b<o1 / 2a<1a1,1b1 / 3a<2a0L & NI/E/N & +{-}/{-}/{-} & F & W:2a \\
C1387 & 1a<o0L 1b<o1 / 2a<1a1,1b1 / 3a<2a1 & NI/E/D & +{-}/{-}/{-} & F & W:2a,3a \\
C1388 & 1a<o0L 1b<o1 / 2a<1a1,1b1 / 3a<2a1L & NI/E/D & +{-}/{-}/{-} & F & W:2a,3a \\
C1389 & 1a<o0L 1b<o1 / 2a<1a1,1b1 / 3a<2a2 & NI/E/D & +{-}/{-}/{-} & F & W:2a,3a \\
C1390 & 1a<o0L 1b<o1 / 2a<1a1,1b1 / 3a<2a2L & NI/E/D & +{-}/{-}/{-} & F & W:2a,3a \\
C1391 & 1a<o0L 1b<o1 / 2a<1a1,1b1L / 3a<2a0L & NI/I/N & +{-}/{-}/{-} & F & CP:2a \\
C1392 & 1a<o0L 1b<o1 / 2a<1a1,1b1L / 3a<2a1 & NI/I/I & +{-}/{-}/{-} & F & CP:2a,3a \\
C1393 & 1a<o0L 1b<o1 / 2a<1a1,1b1L / 3a<2a1L & NI/I/I & +{-}/{-}/{-} & F & CP:2a,3a \\
C1394 & 1a<o0L 1b<o1 / 2a<1a1,1b1L / 3a<2a2 & NI/I/I & +{-}/{-}/{-} & F & CP:2a,3a \\
C1395 & 1a<o0L 1b<o1 / 2a<1a1,1b1L / 3a<2a2L & NI/I/I & +{-}/{-}/{-} & F & CP:2a,3a \\
C1396 & 1a<o0L 1b<o1 / 2a<1a1,1b2 / 3a<2a0L & NI/E/N & +{-}/{-}/{-} & F & CW:2a \\
C1397 & 1a<o0L 1b<o1 / 2a<1a1,1b2 / 3a<2a1 & NI/E/D & +{-}/{-}/{-} & F & CW:2a,3a \\
C1398 & 1a<o0L 1b<o1 / 2a<1a1,1b2 / 3a<2a1L & NI/E/D & +{-}/{-}/{-} & F & CW:2a,3a \\
C1399 & 1a<o0L 1b<o1 / 2a<1a1,1b2 / 3a<2a2 & NI/E/D & +{-}/{-}/{-} & F & CW:2a,3a \\
C1400 & 1a<o0L 1b<o1 / 2a<1a1,1b2 / 3a<2a2L & NI/E/D & +{-}/{-}/{-} & F & CW:2a,3a \\
C1401 & 1a<o0L 1b<o1 / 2a<1a1,1b2L / 3a<2a0L & NI/I/N & +{-}/{-}/{-} & F & P:2a \\
C1402 & 1a<o0L 1b<o1 / 2a<1a1,1b2L / 3a<2a1 & NI/I/I & +{-}/{-}/{-} & F & P:2a,3a \\
C1403 & 1a<o0L 1b<o1 / 2a<1a1,1b2L / 3a<2a1L & NI/I/I & +{-}/{-}/{-} & F & P:2a,3a \\
C1404 & 1a<o0L 1b<o1 / 2a<1a1,1b2L / 3a<2a2 & NI/I/I & +{-}/{-}/{-} & F & P:2a,3a \\
C1405 & 1a<o0L 1b<o1 / 2a<1a1,1b2L / 3a<2a2L & NI/I/I & +{-}/{-}/{-} & F & P:2a,3a \\
C1406 & 1a<o0L 1b<o1 / 2a<1a1L,1b1 / 3a<2a0L & NI/E/N & +{-}/+/+ & F & W:2a \\
C1407 & 1a<o0L 1b<o1 / 2a<1a1L,1b1 / 3a<2a1 & NI/E/D & +{-}/+/{-} & F & W:2a,3a \\
C1408 & 1a<o0L 1b<o1 / 2a<1a1L,1b1 / 3a<2a1L & NI/E/D & +{-}/+/+ & F & W:2a,3a \\
C1409 & 1a<o0L 1b<o1 / 2a<1a1L,1b1 / 3a<2a2 & NI/E/D & +{-}/+/{-} & F & W:2a,3a \\
C1410 & 1a<o0L 1b<o1 / 2a<1a1L,1b1 / 3a<2a2L & NI/E/D & +{-}/+/+ & F & W:2a,3a \\
C1411 & 1a<o0L 1b<o1 / 2a<1a1L,1b2 / 3a<2a0L & NI/E/N & +{-}/+/+ & F & CW:2a \\
C1412 & 1a<o0L 1b<o1 / 2a<1a1L,1b2 / 3a<2a1 & NI/E/D & +{-}/+/{-} & F & CW:2a,3a \\
C1413 & 1a<o0L 1b<o1 / 2a<1a1L,1b2 / 3a<2a1L & NI/E/D & +{-}/+/+ & F & CW:2a,3a \\
C1414 & 1a<o0L 1b<o1 / 2a<1a1L,1b2 / 3a<2a2 & NI/E/D & +{-}/+/{-} & F & CW:2a,3a \\
C1415 & 1a<o0L 1b<o1 / 2a<1a1L,1b2 / 3a<2a2L & NI/E/D & +{-}/+/+ & F & CW:2a,3a \\
C1416 & 1a<o0L 1b<o1 / 2a<1a2,1b0L / 3a<2a0L & NI/N/N & +{-}/{-}/{-} & N &  \\
C1417 & 1a<o0L 1b<o1 / 2a<1a2,1b0L / 3a<2a1 & NI/N/N & +{-}/{-}/{-} & N &  \\
C1418 & 1a<o0L 1b<o1 / 2a<1a2,1b0L / 3a<2a1L & NI/N/N & +{-}/{-}/{-} & N &  \\
C1419 & 1a<o0L 1b<o1 / 2a<1a2,1b0L / 3a<2a2 & NI/N/N & +{-}/{-}/{-} & N &  \\
C1420 & 1a<o0L 1b<o1 / 2a<1a2,1b0L / 3a<2a2L & NI/N/N & +{-}/{-}/{-} & N &  \\
C1421 & 1a<o0L 1b<o1 / 2a<1a2,1b1 / 3a<2a0L & NI/E/N & +{-}/{-}/{-} & F & W:2a \\
C1422 & 1a<o0L 1b<o1 / 2a<1a2,1b1 / 3a<2a1 & NI/E/D & +{-}/{-}/{-} & F & W:2a,3a \\
C1423 & 1a<o0L 1b<o1 / 2a<1a2,1b1 / 3a<2a1L & NI/E/D & +{-}/{-}/{-} & F & W:2a,3a \\
C1424 & 1a<o0L 1b<o1 / 2a<1a2,1b1 / 3a<2a2 & NI/E/D & +{-}/{-}/{-} & F & W:2a,3a \\
C1425 & 1a<o0L 1b<o1 / 2a<1a2,1b1 / 3a<2a2L & NI/E/D & +{-}/{-}/{-} & F & W:2a,3a \\
C1426 & 1a<o0L 1b<o1 / 2a<1a2,1b1L / 3a<2a0L & NI/I/N & +{-}/{-}/{-} & F & CP:2a \\
C1427 & 1a<o0L 1b<o1 / 2a<1a2,1b1L / 3a<2a1 & NI/I/I & +{-}/{-}/{-} & F & CP:2a,3a \\
C1428 & 1a<o0L 1b<o1 / 2a<1a2,1b1L / 3a<2a1L & NI/I/I & +{-}/{-}/{-} & F & CP:2a,3a \\
C1429 & 1a<o0L 1b<o1 / 2a<1a2,1b1L / 3a<2a2 & NI/I/I & +{-}/{-}/{-} & F & CP:2a,3a \\
C1430 & 1a<o0L 1b<o1 / 2a<1a2,1b1L / 3a<2a2L & NI/I/I & +{-}/{-}/{-} & F & CP:2a,3a \\
C1431 & 1a<o0L 1b<o1 / 2a<1a2,1b2 / 3a<2a0L & NI/E/N & +{-}/{-}/{-} & F & W:2a \\
C1432 & 1a<o0L 1b<o1 / 2a<1a2,1b2 / 3a<2a1 & NI/E/D & +{-}/{-}/{-} & F & W:2a,3a \\
C1433 & 1a<o0L 1b<o1 / 2a<1a2,1b2 / 3a<2a1L & NI/E/D & +{-}/{-}/{-} & F & W:2a,3a \\
C1434 & 1a<o0L 1b<o1 / 2a<1a2,1b2 / 3a<2a2 & NI/E/D & +{-}/{-}/{-} & F & W:2a,3a \\
C1435 & 1a<o0L 1b<o1 / 2a<1a2,1b2 / 3a<2a2L & NI/E/D & +{-}/{-}/{-} & F & W:2a,3a \\
C1436 & 1a<o0L 1b<o1 / 2a<1a2,1b2L / 3a<2a0L & NI/I/N & +{-}/{-}/{-} & F & CP:2a \\
C1437 & 1a<o0L 1b<o1 / 2a<1a2,1b2L / 3a<2a1 & NI/I/I & +{-}/{-}/{-} & F & CP:2a,3a \\
C1438 & 1a<o0L 1b<o1 / 2a<1a2,1b2L / 3a<2a1L & NI/I/I & +{-}/{-}/{-} & F & CP:2a,3a \\
C1439 & 1a<o0L 1b<o1 / 2a<1a2,1b2L / 3a<2a2 & NI/I/I & +{-}/{-}/{-} & F & CP:2a,3a \\
C1440 & 1a<o0L 1b<o1 / 2a<1a2,1b2L / 3a<2a2L & NI/I/I & +{-}/{-}/{-} & F & CP:2a,3a \\
C1441 & 1a<o0L 1b<o1 / 2a<1a2L,1b1 / 3a<2a0L & NI/E/N & +{-}/+/+ & F & W:2a \\
C1442 & 1a<o0L 1b<o1 / 2a<1a2L,1b1 / 3a<2a1 & NI/E/D & +{-}/+/{-} & F & W:2a,3a \\
C1443 & 1a<o0L 1b<o1 / 2a<1a2L,1b1 / 3a<2a1L & NI/E/D & +{-}/+/+ & F & W:2a,3a \\
C1444 & 1a<o0L 1b<o1 / 2a<1a2L,1b1 / 3a<2a2 & NI/E/D & +{-}/+/{-} & F & W:2a,3a \\
C1445 & 1a<o0L 1b<o1 / 2a<1a2L,1b1 / 3a<2a2L & NI/E/D & +{-}/+/+ & F & W:2a,3a \\
C1446 & 1a<o0L 1b<o1 / 2a<1a2L,1b2 / 3a<2a0L & NI/E/N & +{-}/+/+ & F & W:2a \\
C1447 & 1a<o0L 1b<o1 / 2a<1a2L,1b2 / 3a<2a1 & NI/E/D & +{-}/+/{-} & F & W:2a,3a \\
C1448 & 1a<o0L 1b<o1 / 2a<1a2L,1b2 / 3a<2a1L & NI/E/D & +{-}/+/+ & F & W:2a,3a \\
C1449 & 1a<o0L 1b<o1 / 2a<1a2L,1b2 / 3a<2a2 & NI/E/D & +{-}/+/{-} & F & W:2a,3a \\
C1450 & 1a<o0L 1b<o1 / 2a<1a2L,1b2 / 3a<2a2L & NI/E/D & +{-}/+/+ & F & W:2a,3a \\
C1451 & 1a<o0L 1b<o2 / 2a<1a0L,1b1 / 3a<2a0L & NI/I/N & +{-}/+/+ & N &  \\
C1452 & 1a<o0L 1b<o2 / 2a<1a0L,1b1 / 3a<2a1 & NI/I/I & +{-}/+/{-} & N &  \\
C1453 & 1a<o0L 1b<o2 / 2a<1a0L,1b1 / 3a<2a1L & NI/I/I & +{-}/+/+ & N &  \\
C1454 & 1a<o0L 1b<o2 / 2a<1a0L,1b1 / 3a<2a2 & NI/I/I & +{-}/+/{-} & N &  \\
C1455 & 1a<o0L 1b<o2 / 2a<1a0L,1b1 / 3a<2a2L & NI/I/I & +{-}/+/+ & N &  \\
C1456 & 1a<o0L 1b<o2 / 2a<1a0L,1b2 / 3a<2a0L & NI/I/N & +{-}/+/+ & N &  \\
C1457 & 1a<o0L 1b<o2 / 2a<1a0L,1b2 / 3a<2a1 & NI/I/I & +{-}/+/{-} & N &  \\
C1458 & 1a<o0L 1b<o2 / 2a<1a0L,1b2 / 3a<2a1L & NI/I/I & +{-}/+/+ & N &  \\
C1459 & 1a<o0L 1b<o2 / 2a<1a0L,1b2 / 3a<2a2 & NI/I/I & +{-}/+/{-} & N &  \\
C1460 & 1a<o0L 1b<o2 / 2a<1a0L,1b2 / 3a<2a2L & NI/I/I & +{-}/+/+ & N &  \\
C1461 & 1a<o0L 1b<o2 / 2a<1a1,1b0L / 3a<2a0L & NI/N/N & +{-}/{-}/{-} & N &  \\
C1462 & 1a<o0L 1b<o2 / 2a<1a1,1b0L / 3a<2a1 & NI/N/N & +{-}/{-}/{-} & N &  \\
C1463 & 1a<o0L 1b<o2 / 2a<1a1,1b0L / 3a<2a1L & NI/N/N & +{-}/{-}/{-} & N &  \\
C1464 & 1a<o0L 1b<o2 / 2a<1a1,1b0L / 3a<2a2 & NI/N/N & +{-}/{-}/{-} & N &  \\
C1465 & 1a<o0L 1b<o2 / 2a<1a1,1b0L / 3a<2a2L & NI/N/N & +{-}/{-}/{-} & N &  \\
C1466 & 1a<o0L 1b<o2 / 2a<1a1,1b1 / 3a<2a0L & NI/E/N & +{-}/{-}/{-} & F & W:2a \\
C1467 & 1a<o0L 1b<o2 / 2a<1a1,1b1 / 3a<2a1 & NI/E/D & +{-}/{-}/{-} & F & W:2a,3a \\
C1468 & 1a<o0L 1b<o2 / 2a<1a1,1b1 / 3a<2a1L & NI/E/D & +{-}/{-}/{-} & F & W:2a,3a \\
C1469 & 1a<o0L 1b<o2 / 2a<1a1,1b1 / 3a<2a2 & NI/E/D & +{-}/{-}/{-} & F & W:2a,3a \\
C1470 & 1a<o0L 1b<o2 / 2a<1a1,1b1 / 3a<2a2L & NI/E/D & +{-}/{-}/{-} & F & W:2a,3a \\
C1471 & 1a<o0L 1b<o2 / 2a<1a1,1b1L / 3a<2a0L & NI/I/N & +{-}/{-}/{-} & F & CP:2a \\
C1472 & 1a<o0L 1b<o2 / 2a<1a1,1b1L / 3a<2a1 & NI/I/I & +{-}/{-}/{-} & F & CP:2a,3a \\
C1473 & 1a<o0L 1b<o2 / 2a<1a1,1b1L / 3a<2a1L & NI/I/I & +{-}/{-}/{-} & F & CP:2a,3a \\
C1474 & 1a<o0L 1b<o2 / 2a<1a1,1b1L / 3a<2a2 & NI/I/I & +{-}/{-}/{-} & F & CP:2a,3a \\
C1475 & 1a<o0L 1b<o2 / 2a<1a1,1b1L / 3a<2a2L & NI/I/I & +{-}/{-}/{-} & F & CP:2a,3a \\
C1476 & 1a<o0L 1b<o2 / 2a<1a1,1b2 / 3a<2a0L & NI/E/N & +{-}/{-}/{-} & F & CW:2a \\
C1477 & 1a<o0L 1b<o2 / 2a<1a1,1b2 / 3a<2a1 & NI/E/D & +{-}/{-}/{-} & F & CW:2a,3a \\
C1478 & 1a<o0L 1b<o2 / 2a<1a1,1b2 / 3a<2a1L & NI/E/D & +{-}/{-}/{-} & F & CW:2a,3a \\
C1479 & 1a<o0L 1b<o2 / 2a<1a1,1b2 / 3a<2a2 & NI/E/D & +{-}/{-}/{-} & F & CW:2a,3a \\
C1480 & 1a<o0L 1b<o2 / 2a<1a1,1b2 / 3a<2a2L & NI/E/D & +{-}/{-}/{-} & F & CW:2a,3a \\
C1481 & 1a<o0L 1b<o2 / 2a<1a1,1b2L / 3a<2a0L & NI/I/N & +{-}/{-}/{-} & F & P:2a \\
C1482 & 1a<o0L 1b<o2 / 2a<1a1,1b2L / 3a<2a1 & NI/I/I & +{-}/{-}/{-} & F & P:2a,3a \\
C1483 & 1a<o0L 1b<o2 / 2a<1a1,1b2L / 3a<2a1L & NI/I/I & +{-}/{-}/{-} & F & P:2a,3a \\
C1484 & 1a<o0L 1b<o2 / 2a<1a1,1b2L / 3a<2a2 & NI/I/I & +{-}/{-}/{-} & F & P:2a,3a \\
C1485 & 1a<o0L 1b<o2 / 2a<1a1,1b2L / 3a<2a2L & NI/I/I & +{-}/{-}/{-} & F & P:2a,3a \\
C1486 & 1a<o0L 1b<o2 / 2a<1a1L,1b1 / 3a<2a0L & NI/E/N & +{-}/+/+ & F & W:2a \\
C1487 & 1a<o0L 1b<o2 / 2a<1a1L,1b1 / 3a<2a1 & NI/E/D & +{-}/+/{-} & F & W:2a,3a \\
C1488 & 1a<o0L 1b<o2 / 2a<1a1L,1b1 / 3a<2a1L & NI/E/D & +{-}/+/+ & F & W:2a,3a \\
C1489 & 1a<o0L 1b<o2 / 2a<1a1L,1b1 / 3a<2a2 & NI/E/D & +{-}/+/{-} & F & W:2a,3a \\
C1490 & 1a<o0L 1b<o2 / 2a<1a1L,1b1 / 3a<2a2L & NI/E/D & +{-}/+/+ & F & W:2a,3a \\
C1491 & 1a<o0L 1b<o2 / 2a<1a1L,1b2 / 3a<2a0L & NI/E/N & +{-}/+/+ & F & CW:2a \\
C1492 & 1a<o0L 1b<o2 / 2a<1a1L,1b2 / 3a<2a1 & NI/E/D & +{-}/+/{-} & F & CW:2a,3a \\
C1493 & 1a<o0L 1b<o2 / 2a<1a1L,1b2 / 3a<2a1L & NI/E/D & +{-}/+/+ & F & CW:2a,3a \\
C1494 & 1a<o0L 1b<o2 / 2a<1a1L,1b2 / 3a<2a2 & NI/E/D & +{-}/+/{-} & F & CW:2a,3a \\
C1495 & 1a<o0L 1b<o2 / 2a<1a1L,1b2 / 3a<2a2L & NI/E/D & +{-}/+/+ & F & CW:2a,3a \\
C1496 & 1a<o0L 1b<o2 / 2a<1a2,1b0L / 3a<2a0L & NI/N/N & +{-}/{-}/{-} & N &  \\
C1497 & 1a<o0L 1b<o2 / 2a<1a2,1b0L / 3a<2a1 & NI/N/N & +{-}/{-}/{-} & N &  \\
C1498 & 1a<o0L 1b<o2 / 2a<1a2,1b0L / 3a<2a1L & NI/N/N & +{-}/{-}/{-} & N &  \\
C1499 & 1a<o0L 1b<o2 / 2a<1a2,1b0L / 3a<2a2 & NI/N/N & +{-}/{-}/{-} & N &  \\
C1500 & 1a<o0L 1b<o2 / 2a<1a2,1b0L / 3a<2a2L & NI/N/N & +{-}/{-}/{-} & N &  \\
C1501 & 1a<o0L 1b<o2 / 2a<1a2,1b1 / 3a<2a0L & NI/E/N & +{-}/{-}/{-} & F & W:2a \\
C1502 & 1a<o0L 1b<o2 / 2a<1a2,1b1 / 3a<2a1 & NI/E/D & +{-}/{-}/{-} & F & W:2a,3a \\
C1503 & 1a<o0L 1b<o2 / 2a<1a2,1b1 / 3a<2a1L & NI/E/D & +{-}/{-}/{-} & F & W:2a,3a \\
C1504 & 1a<o0L 1b<o2 / 2a<1a2,1b1 / 3a<2a2 & NI/E/D & +{-}/{-}/{-} & F & W:2a,3a \\
C1505 & 1a<o0L 1b<o2 / 2a<1a2,1b1 / 3a<2a2L & NI/E/D & +{-}/{-}/{-} & F & W:2a,3a \\
C1506 & 1a<o0L 1b<o2 / 2a<1a2,1b1L / 3a<2a0L & NI/I/N & +{-}/{-}/{-} & F & CP:2a \\
C1507 & 1a<o0L 1b<o2 / 2a<1a2,1b1L / 3a<2a1 & NI/I/I & +{-}/{-}/{-} & F & CP:2a,3a \\
C1508 & 1a<o0L 1b<o2 / 2a<1a2,1b1L / 3a<2a1L & NI/I/I & +{-}/{-}/{-} & F & CP:2a,3a \\
C1509 & 1a<o0L 1b<o2 / 2a<1a2,1b1L / 3a<2a2 & NI/I/I & +{-}/{-}/{-} & F & CP:2a,3a \\
C1510 & 1a<o0L 1b<o2 / 2a<1a2,1b1L / 3a<2a2L & NI/I/I & +{-}/{-}/{-} & F & CP:2a,3a \\
C1511 & 1a<o0L 1b<o2 / 2a<1a2,1b2 / 3a<2a0L & NI/E/N & +{-}/{-}/{-} & F & W:2a \\
C1512 & 1a<o0L 1b<o2 / 2a<1a2,1b2 / 3a<2a1 & NI/E/D & +{-}/{-}/{-} & F & W:2a,3a \\
C1513 & 1a<o0L 1b<o2 / 2a<1a2,1b2 / 3a<2a1L & NI/E/D & +{-}/{-}/{-} & F & W:2a,3a \\
C1514 & 1a<o0L 1b<o2 / 2a<1a2,1b2 / 3a<2a2 & NI/E/D & +{-}/{-}/{-} & F & W:2a,3a \\
C1515 & 1a<o0L 1b<o2 / 2a<1a2,1b2 / 3a<2a2L & NI/E/D & +{-}/{-}/{-} & F & W:2a,3a \\
C1516 & 1a<o0L 1b<o2 / 2a<1a2,1b2L / 3a<2a0L & NI/I/N & +{-}/{-}/{-} & F & CP:2a \\
C1517 & 1a<o0L 1b<o2 / 2a<1a2,1b2L / 3a<2a1 & NI/I/I & +{-}/{-}/{-} & F & CP:2a,3a \\
C1518 & 1a<o0L 1b<o2 / 2a<1a2,1b2L / 3a<2a1L & NI/I/I & +{-}/{-}/{-} & F & CP:2a,3a \\
C1519 & 1a<o0L 1b<o2 / 2a<1a2,1b2L / 3a<2a2 & NI/I/I & +{-}/{-}/{-} & F & CP:2a,3a \\
C1520 & 1a<o0L 1b<o2 / 2a<1a2,1b2L / 3a<2a2L & NI/I/I & +{-}/{-}/{-} & F & CP:2a,3a \\
C1521 & 1a<o0L 1b<o2 / 2a<1a2L,1b1 / 3a<2a0L & NI/E/N & +{-}/+/+ & F & W:2a \\
C1522 & 1a<o0L 1b<o2 / 2a<1a2L,1b1 / 3a<2a1 & NI/E/D & +{-}/+/{-} & F & W:2a,3a \\
C1523 & 1a<o0L 1b<o2 / 2a<1a2L,1b1 / 3a<2a1L & NI/E/D & +{-}/+/+ & F & W:2a,3a \\
C1524 & 1a<o0L 1b<o2 / 2a<1a2L,1b1 / 3a<2a2 & NI/E/D & +{-}/+/{-} & F & W:2a,3a \\
C1525 & 1a<o0L 1b<o2 / 2a<1a2L,1b1 / 3a<2a2L & NI/E/D & +{-}/+/+ & F & W:2a,3a \\
C1526 & 1a<o0L 1b<o2 / 2a<1a2L,1b2 / 3a<2a0L & NI/E/N & +{-}/+/+ & F & W:2a \\
C1527 & 1a<o0L 1b<o2 / 2a<1a2L,1b2 / 3a<2a1 & NI/E/D & +{-}/+/{-} & F & W:2a,3a \\
C1528 & 1a<o0L 1b<o2 / 2a<1a2L,1b2 / 3a<2a1L & NI/E/D & +{-}/+/+ & F & W:2a,3a \\
C1529 & 1a<o0L 1b<o2 / 2a<1a2L,1b2 / 3a<2a2 & NI/E/D & +{-}/+/{-} & F & W:2a,3a \\
C1530 & 1a<o0L 1b<o2 / 2a<1a2L,1b2 / 3a<2a2L & NI/E/D & +{-}/+/+ & F & W:2a,3a \\
C1531 & 1a<o1 / 2a<1a0L 2b<1a1 / 3a<2a0L,2b1 & I/NI/I & {-}/{-}{-}/{-} & N &  \\
C1532 & 1a<o1 / 2a<1a0L 2b<1a1 / 3a<2a0L,2b2 & I/NI/I & {-}/{-}{-}/{-} & N &  \\
C1533 & 1a<o1 / 2a<1a0L 2b<1a1 / 3a<2a1,2b0L & I/NI/N & {-}/{-}{-}/{-} & N &  \\
C1534 & 1a<o1 / 2a<1a0L 2b<1a1 / 3a<2a1,2b1 & I/NI/E & {-}/{-}{-}/{-} & F & W:3a \\
C1535 & 1a<o1 / 2a<1a0L 2b<1a1 / 3a<2a1,2b1L & I/NI/I & {-}/{-}{-}/{-} & F & CP:3a \\
C1536 & 1a<o1 / 2a<1a0L 2b<1a1 / 3a<2a1,2b2 & I/NI/E & {-}/{-}{-}/{-} & F & CW:3a \\
C1537 & 1a<o1 / 2a<1a0L 2b<1a1 / 3a<2a1,2b2L & I/NI/I & {-}/{-}{-}/{-} & F & P:3a \\
C1538 & 1a<o1 / 2a<1a0L 2b<1a1 / 3a<2a1L,2b1 & I/NI/E & {-}/{-}{-}/{-} & F & W:3a \\
C1539 & 1a<o1 / 2a<1a0L 2b<1a1 / 3a<2a1L,2b2 & I/NI/E & {-}/{-}{-}/{-} & F & CW:3a \\
C1540 & 1a<o1 / 2a<1a0L 2b<1a1 / 3a<2a2,2b0L & I/NI/N & {-}/{-}{-}/{-} & N &  \\
C1541 & 1a<o1 / 2a<1a0L 2b<1a1 / 3a<2a2,2b1 & I/NI/E & {-}/{-}{-}/{-} & F & W:3a \\
C1542 & 1a<o1 / 2a<1a0L 2b<1a1 / 3a<2a2,2b1L & I/NI/I & {-}/{-}{-}/{-} & F & CP:3a \\
C1543 & 1a<o1 / 2a<1a0L 2b<1a1 / 3a<2a2,2b2 & I/NI/E & {-}/{-}{-}/{-} & F & W:3a \\
C1544 & 1a<o1 / 2a<1a0L 2b<1a1 / 3a<2a2,2b2L & I/NI/I & {-}/{-}{-}/{-} & F & CP:3a \\
C1545 & 1a<o1 / 2a<1a0L 2b<1a1 / 3a<2a2L,2b1 & I/NI/E & {-}/{-}{-}/{-} & F & W:3a \\
C1546 & 1a<o1 / 2a<1a0L 2b<1a1 / 3a<2a2L,2b2 & I/NI/E & {-}/{-}{-}/{-} & F & W:3a \\
C1547 & 1a<o1 / 2a<1a0L 2b<1a2 / 3a<2a0L,2b1 & I/NI/I & {-}/{-}{-}/{-} & N &  \\
C1548 & 1a<o1 / 2a<1a0L 2b<1a2 / 3a<2a0L,2b2 & I/NI/I & {-}/{-}{-}/{-} & N &  \\
C1549 & 1a<o1 / 2a<1a0L 2b<1a2 / 3a<2a1,2b0L & I/NI/N & {-}/{-}{-}/{-} & N &  \\
C1550 & 1a<o1 / 2a<1a0L 2b<1a2 / 3a<2a1,2b1 & I/NI/E & {-}/{-}{-}/{-} & F & W:3a \\
C1551 & 1a<o1 / 2a<1a0L 2b<1a2 / 3a<2a1,2b1L & I/NI/I & {-}/{-}{-}/{-} & F & CP:3a \\
C1552 & 1a<o1 / 2a<1a0L 2b<1a2 / 3a<2a1,2b2 & I/NI/E & {-}/{-}{-}/{-} & F & CW:3a \\
C1553 & 1a<o1 / 2a<1a0L 2b<1a2 / 3a<2a1,2b2L & I/NI/I & {-}/{-}{-}/{-} & F & P:3a \\
C1554 & 1a<o1 / 2a<1a0L 2b<1a2 / 3a<2a1L,2b1 & I/NI/E & {-}/{-}{-}/{-} & F & W:3a \\
C1555 & 1a<o1 / 2a<1a0L 2b<1a2 / 3a<2a1L,2b2 & I/NI/E & {-}/{-}{-}/{-} & F & CW:3a \\
C1556 & 1a<o1 / 2a<1a0L 2b<1a2 / 3a<2a2,2b0L & I/NI/N & {-}/{-}{-}/{-} & N &  \\
C1557 & 1a<o1 / 2a<1a0L 2b<1a2 / 3a<2a2,2b1 & I/NI/E & {-}/{-}{-}/{-} & F & W:3a \\
C1558 & 1a<o1 / 2a<1a0L 2b<1a2 / 3a<2a2,2b1L & I/NI/I & {-}/{-}{-}/{-} & F & CP:3a \\
C1559 & 1a<o1 / 2a<1a0L 2b<1a2 / 3a<2a2,2b2 & I/NI/E & {-}/{-}{-}/{-} & F & W:3a \\
C1560 & 1a<o1 / 2a<1a0L 2b<1a2 / 3a<2a2,2b2L & I/NI/I & {-}/{-}{-}/{-} & F & CP:3a \\
C1561 & 1a<o1 / 2a<1a0L 2b<1a2 / 3a<2a2L,2b1 & I/NI/E & {-}/{-}{-}/{-} & F & W:3a \\
C1562 & 1a<o1 / 2a<1a0L 2b<1a2 / 3a<2a2L,2b2 & I/NI/E & {-}/{-}{-}/{-} & F & W:3a \\
C1563 & 1a<o1 / 2a<1a1 2b<1a1 / 3a<2a0L,2b1 & I/EE/D & {-}/{-}{-}/{-} & E & W:1a \\
C1564 & 1a<o1 / 2a<1a1 2b<1a1 / 3a<2a0L,2b2 & I/EE/D & {-}/{-}{-}/{-} & E & W:1a \\
C1565 & 1a<o1 / 2a<1a1 2b<1a1 / 3a<2a1,2b1 & I/EE/D & {-}/{-}{-}/{-} & E & W:1a \\
C1566 & 1a<o1 / 2a<1a1 2b<1a1 / 3a<2a1,2b1L & I/EE/D & {-}/{-}{-}/{-} & E & W:1a \\
C1567 & 1a<o1 / 2a<1a1 2b<1a1 / 3a<2a1,2b2 & I/EE/D & {-}/{-}{-}/{-} & E & W:1a \\
C1568 & 1a<o1 / 2a<1a1 2b<1a1 / 3a<2a1,2b2L & I/EE/D & {-}/{-}{-}/{-} & E & W:1a \\
C1569 & 1a<o1 / 2a<1a1 2b<1a1 / 3a<2a1L,2b2 & I/EE/D & {-}/{-}{-}/{-} & E & W:1a \\
C1570 & 1a<o1 / 2a<1a1 2b<1a1 / 3a<2a2,2b2 & I/EE/D & {-}/{-}{-}/{-} & E & W:1a \\
C1571 & 1a<o1 / 2a<1a1 2b<1a1 / 3a<2a2,2b2L & I/EE/D & {-}/{-}{-}/{-} & E & W:1a \\
C1572 & 1a<o1 / 2a<1a1 2b<1a1L / 3a<2a0L,2b1 & I/LI/I & {-}/{-}{-}/{-} & B & CP:2b,3a W:1a,2b,3a \\
C1573 & 1a<o1 / 2a<1a1 2b<1a1L / 3a<2a0L,2b2 & I/LI/I & {-}/{-}{-}/{-} & B & CP:2b,3a W:1a,2b,3a \\
C1574 & 1a<o1 / 2a<1a1 2b<1a1L / 3a<2a1,2b0L & I/LI/D & {-}/{-}{-}/{-} & B & CP:2b W:1a,2b \\
C1575 & 1a<o1 / 2a<1a1 2b<1a1L / 3a<2a1,2b1 & I/LI/E & {-}/{-}{-}/{-} & BF & CP:2b W:1a,2b \\
C1576 & 1a<o1 / 2a<1a1 2b<1a1L / 3a<2a1,2b1L & I/LI/I & {-}/{-}{-}/{-} & BF & CP:2b,3a W:1a,2b,3a \\
C1577 & 1a<o1 / 2a<1a1 2b<1a1L / 3a<2a1,2b2 & I/LI/E & {-}/{-}{-}/{-} & BF & CP:2b W:1a,2b \\
C1578 & 1a<o1 / 2a<1a1 2b<1a1L / 3a<2a1,2b2L & I/LI/I & {-}/{-}{-}/{-} & BF & CP:2b,3a W:1a,2b,3a \\
C1579 & 1a<o1 / 2a<1a1 2b<1a1L / 3a<2a1L,2b1 & I/LI/E & {-}/{-}{-}/{-} & BF & CP:2b W:1a,2b \\
C1580 & 1a<o1 / 2a<1a1 2b<1a1L / 3a<2a1L,2b2 & I/LI/E & {-}/{-}{-}/{-} & BF & CP:2b W:1a,2b \\
C1581 & 1a<o1 / 2a<1a1 2b<1a1L / 3a<2a2,2b0L & I/LI/D & {-}/{-}{-}/{-} & B & CP:2b W:1a,2b \\
C1582 & 1a<o1 / 2a<1a1 2b<1a1L / 3a<2a2,2b1 & I/LI/E & {-}/{-}{-}/{-} & BF & CP:2b W:1a,2b \\
C1583 & 1a<o1 / 2a<1a1 2b<1a1L / 3a<2a2,2b1L & I/LI/I & {-}/{-}{-}/{-} & BF & CP:2b,3a W:1a,2b,3a \\
C1584 & 1a<o1 / 2a<1a1 2b<1a1L / 3a<2a2,2b2 & I/LI/E & {-}/{-}{-}/{-} & BF & CP:2b W:1a,2b \\
C1585 & 1a<o1 / 2a<1a1 2b<1a1L / 3a<2a2,2b2L & I/LI/I & {-}/{-}{-}/{-} & BF & CP:2b,3a W:1a,2b,3a \\
C1586 & 1a<o1 / 2a<1a1 2b<1a1L / 3a<2a2L,2b1 & I/LI/E & {-}/{-}{-}/{-} & BF & CP:2b W:1a,2b \\
C1587 & 1a<o1 / 2a<1a1 2b<1a1L / 3a<2a2L,2b2 & I/LI/E & {-}/{-}{-}/{-} & BF & CP:2b W:1a,2b \\
C1588 & 1a<o1 / 2a<1a1 2b<1a2 / 3a<2a0L,2b1 & I/EE/D & {-}/{-}{-}/{-} & E & C:2b,3a W:1a \\
C1589 & 1a<o1 / 2a<1a1 2b<1a2 / 3a<2a0L,2b2 & I/EE/D & {-}/{-}{-}/{-} & E & C:2b,3a W:1a \\
C1590 & 1a<o1 / 2a<1a1 2b<1a2 / 3a<2a1,2b0L & I/EE/D & {-}/{-}{-}/{-} & E & C:2b W:1a \\
C1591 & 1a<o1 / 2a<1a1 2b<1a2 / 3a<2a1,2b1 & I/EE/D & {-}/{-}{-}/{-} & E & C:2b W:1a \\
C1592 & 1a<o1 / 2a<1a1 2b<1a2 / 3a<2a1,2b1L & I/EE/D & {-}/{-}{-}/{-} & E & C:2b W:1a \\
C1593 & 1a<o1 / 2a<1a1 2b<1a2 / 3a<2a1,2b2 & I/EE/D & {-}/{-}{-}/{-} & E & C:2b,3a W:1a \\
C1594 & 1a<o1 / 2a<1a1 2b<1a2 / 3a<2a1,2b2L & I/EE/D & {-}/{-}{-}/{-} & E & C:2b,3a W:1a \\
C1595 & 1a<o1 / 2a<1a1 2b<1a2 / 3a<2a1L,2b1 & I/EE/D & {-}/{-}{-}/{-} & E & C:2b W:1a \\
C1596 & 1a<o1 / 2a<1a1 2b<1a2 / 3a<2a1L,2b2 & I/EE/D & {-}/{-}{-}/{-} & E & C:2b,3a W:1a \\
C1597 & 1a<o1 / 2a<1a1 2b<1a2 / 3a<2a2,2b0L & I/EE/D & {-}/{-}{-}/{-} & E & C:2b W:1a \\
C1598 & 1a<o1 / 2a<1a1 2b<1a2 / 3a<2a2,2b1 & I/EE/D & {-}/{-}{-}/{-} & E & C:2b W:1a \\
C1599 & 1a<o1 / 2a<1a1 2b<1a2 / 3a<2a2,2b1L & I/EE/D & {-}/{-}{-}/{-} & E & C:2b W:1a \\
C1600 & 1a<o1 / 2a<1a1 2b<1a2 / 3a<2a2,2b2 & I/EE/D & {-}/{-}{-}/{-} & E & C:2b W:1a \\
C1601 & 1a<o1 / 2a<1a1 2b<1a2 / 3a<2a2,2b2L & I/EE/D & {-}/{-}{-}/{-} & E & C:2b W:1a \\
C1602 & 1a<o1 / 2a<1a1 2b<1a2 / 3a<2a2L,2b1 & I/EE/D & {-}/{-}{-}/{-} & E & C:2b W:1a \\
C1603 & 1a<o1 / 2a<1a1 2b<1a2 / 3a<2a2L,2b2 & I/EE/D & {-}/{-}{-}/{-} & E & C:2b W:1a \\
C1604 & 1a<o1 / 2a<1a1 2b<1a2L / 3a<2a0L,2b1 & I/LI/I & {-}/{-}{-}/{-} & B & P:2b,3a W:1a,2b,3a \\
C1605 & 1a<o1 / 2a<1a1 2b<1a2L / 3a<2a0L,2b2 & I/LI/I & {-}/{-}{-}/{-} & B & P:2b,3a W:1a,2b,3a \\
C1606 & 1a<o1 / 2a<1a1 2b<1a2L / 3a<2a1,2b0L & I/LI/D & {-}/{-}{-}/{-} & B & P:2b W:1a,2b \\
C1607 & 1a<o1 / 2a<1a1 2b<1a2L / 3a<2a1,2b1 & I/LI/E & {-}/{-}{-}/{-} & BF & P:2b W:1a,2b \\
C1608 & 1a<o1 / 2a<1a1 2b<1a2L / 3a<2a1,2b1L & I/LI/I & {-}/{-}{-}/{-} & BF & C:3a P:2b,3a W:1a,2b,3a \\
C1609 & 1a<o1 / 2a<1a1 2b<1a2L / 3a<2a1,2b2 & I/LI/E & {-}/{-}{-}/{-} & BF & C:3a P:2b W:1a,2b \\
C1610 & 1a<o1 / 2a<1a1 2b<1a2L / 3a<2a1,2b2L & I/LI/I & {-}/{-}{-}/{-} & BF & P:2b,3a W:1a,2b,3a \\
C1611 & 1a<o1 / 2a<1a1 2b<1a2L / 3a<2a1L,2b1 & I/LI/E & {-}/{-}{-}/{-} & BF & P:2b W:1a,2b \\
C1612 & 1a<o1 / 2a<1a1 2b<1a2L / 3a<2a1L,2b2 & I/LI/E & {-}/{-}{-}/{-} & BF & C:3a P:2b W:1a,2b \\
C1613 & 1a<o1 / 2a<1a1 2b<1a2L / 3a<2a2,2b0L & I/LI/D & {-}/{-}{-}/{-} & B & P:2b W:1a,2b \\
C1614 & 1a<o1 / 2a<1a1 2b<1a2L / 3a<2a2,2b1 & I/LI/E & {-}/{-}{-}/{-} & BF & P:2b W:1a,2b \\
C1615 & 1a<o1 / 2a<1a1 2b<1a2L / 3a<2a2,2b1L & I/LI/I & {-}/{-}{-}/{-} & BF & C:3a P:2b,3a W:1a,2b,3a \\
C1616 & 1a<o1 / 2a<1a1 2b<1a2L / 3a<2a2,2b2 & I/LI/E & {-}/{-}{-}/{-} & BF & P:2b W:1a,2b \\
C1617 & 1a<o1 / 2a<1a1 2b<1a2L / 3a<2a2,2b2L & I/LI/I & {-}/{-}{-}/{-} & BF & C:3a P:2b,3a W:1a,2b,3a \\
C1618 & 1a<o1 / 2a<1a1 2b<1a2L / 3a<2a2L,2b1 & I/LI/E & {-}/{-}{-}/{-} & BF & P:2b W:1a,2b \\
C1619 & 1a<o1 / 2a<1a1 2b<1a2L / 3a<2a2L,2b2 & I/LI/E & {-}/{-}{-}/{-} & BF & P:2b W:1a,2b \\
C1620 & 1a<o1 / 2a<1a1L 2b<1a2 / 3a<2a0L,2b1 & I/IL/D & {-}/{-}{-}/{-} & B & C:2a,2b,3a P:2a W:1a,2a \\
C1621 & 1a<o1 / 2a<1a1L 2b<1a2 / 3a<2a0L,2b2 & I/IL/D & {-}/{-}{-}/{-} & B & C:2a,2b,3a P:2a W:1a,2a \\
C1622 & 1a<o1 / 2a<1a1L 2b<1a2 / 3a<2a1,2b0L & I/IL/I & {-}/{-}{-}/{-} & B & C:2a,2b,3a P:2a,3a W:1a,2a,3a \\
C1623 & 1a<o1 / 2a<1a1L 2b<1a2 / 3a<2a1,2b1 & I/IL/E & {-}/{-}{-}/{-} & BF & C:2a,2b P:2a W:1a,2a \\
C1624 & 1a<o1 / 2a<1a1L 2b<1a2 / 3a<2a1,2b1L & I/IL/E & {-}/{-}{-}/{-} & BF & C:2a,2b P:2a W:1a,2a \\
C1625 & 1a<o1 / 2a<1a1L 2b<1a2 / 3a<2a1,2b2 & I/IL/E & {-}/{-}{-}/{-} & BF & C:2a,2b,3a P:2a W:1a,2a \\
C1626 & 1a<o1 / 2a<1a1L 2b<1a2 / 3a<2a1,2b2L & I/IL/E & {-}/{-}{-}/{-} & BF & C:2a,2b,3a P:2a W:1a,2a \\
C1627 & 1a<o1 / 2a<1a1L 2b<1a2 / 3a<2a1L,2b1 & I/IL/I & {-}/{-}{-}/{-} & BF & C:2a,2b,3a P:2a,3a W:1a,2a,3a \\
C1628 & 1a<o1 / 2a<1a1L 2b<1a2 / 3a<2a1L,2b2 & I/IL/I & {-}/{-}{-}/{-} & BF & C:2a,2b P:2a,3a W:1a,2a,3a \\
C1629 & 1a<o1 / 2a<1a1L 2b<1a2 / 3a<2a2,2b0L & I/IL/I & {-}/{-}{-}/{-} & B & C:2a,2b,3a P:2a,3a W:1a,2a,3a \\
C1630 & 1a<o1 / 2a<1a1L 2b<1a2 / 3a<2a2,2b1 & I/IL/E & {-}/{-}{-}/{-} & BF & C:2a,2b P:2a W:1a,2a \\
C1631 & 1a<o1 / 2a<1a1L 2b<1a2 / 3a<2a2,2b1L & I/IL/E & {-}/{-}{-}/{-} & BF & C:2a,2b P:2a W:1a,2a \\
C1632 & 1a<o1 / 2a<1a1L 2b<1a2 / 3a<2a2,2b2 & I/IL/E & {-}/{-}{-}/{-} & BF & C:2a,2b P:2a W:1a,2a \\
C1633 & 1a<o1 / 2a<1a1L 2b<1a2 / 3a<2a2,2b2L & I/IL/E & {-}/{-}{-}/{-} & BF & C:2a,2b P:2a W:1a,2a \\
C1634 & 1a<o1 / 2a<1a1L 2b<1a2 / 3a<2a2L,2b1 & I/IL/I & {-}/{-}{-}/{-} & BF & C:2a,2b,3a P:2a,3a W:1a,2a,3a \\
C1635 & 1a<o1 / 2a<1a1L 2b<1a2 / 3a<2a2L,2b2 & I/IL/I & {-}/{-}{-}/{-} & BF & C:2a,2b,3a P:2a,3a W:1a,2a,3a \\
C1636 & 1a<o1 / 2a<1a2 2b<1a2 / 3a<2a0L,2b1 & I/EE/D & {-}/{-}{-}/{-} & E & W:1a \\
C1637 & 1a<o1 / 2a<1a2 2b<1a2 / 3a<2a0L,2b2 & I/EE/D & {-}/{-}{-}/{-} & E & W:1a \\
C1638 & 1a<o1 / 2a<1a2 2b<1a2 / 3a<2a1,2b1 & I/EE/D & {-}/{-}{-}/{-} & E & W:1a \\
C1639 & 1a<o1 / 2a<1a2 2b<1a2 / 3a<2a1,2b1L & I/EE/D & {-}/{-}{-}/{-} & E & W:1a \\
C1640 & 1a<o1 / 2a<1a2 2b<1a2 / 3a<2a1,2b2 & I/EE/D & {-}/{-}{-}/{-} & E & W:1a \\
C1641 & 1a<o1 / 2a<1a2 2b<1a2 / 3a<2a1,2b2L & I/EE/D & {-}/{-}{-}/{-} & E & W:1a \\
C1642 & 1a<o1 / 2a<1a2 2b<1a2 / 3a<2a1L,2b2 & I/EE/D & {-}/{-}{-}/{-} & E & W:1a \\
C1643 & 1a<o1 / 2a<1a2 2b<1a2 / 3a<2a2,2b2 & I/EE/D & {-}/{-}{-}/{-} & E & W:1a \\
C1644 & 1a<o1 / 2a<1a2 2b<1a2 / 3a<2a2,2b2L & I/EE/D & {-}/{-}{-}/{-} & E & W:1a \\
C1645 & 1a<o1 / 2a<1a2 2b<1a2L / 3a<2a0L,2b1 & I/LI/I & {-}/{-}{-}/{-} & B & CP:2b,3a W:1a,2b,3a \\
C1646 & 1a<o1 / 2a<1a2 2b<1a2L / 3a<2a0L,2b2 & I/LI/I & {-}/{-}{-}/{-} & B & CP:2b,3a W:1a,2b,3a \\
C1647 & 1a<o1 / 2a<1a2 2b<1a2L / 3a<2a1,2b0L & I/LI/D & {-}/{-}{-}/{-} & B & CP:2b W:1a,2b \\
C1648 & 1a<o1 / 2a<1a2 2b<1a2L / 3a<2a1,2b1 & I/LI/E & {-}/{-}{-}/{-} & BF & CP:2b W:1a,2b \\
C1649 & 1a<o1 / 2a<1a2 2b<1a2L / 3a<2a1,2b1L & I/LI/I & {-}/{-}{-}/{-} & BF & CP:2b,3a W:1a,2b,3a \\
C1650 & 1a<o1 / 2a<1a2 2b<1a2L / 3a<2a1,2b2 & I/LI/E & {-}/{-}{-}/{-} & BF & CP:2b W:1a,2b \\
C1651 & 1a<o1 / 2a<1a2 2b<1a2L / 3a<2a1,2b2L & I/LI/I & {-}/{-}{-}/{-} & BF & CP:2b,3a W:1a,2b,3a \\
C1652 & 1a<o1 / 2a<1a2 2b<1a2L / 3a<2a1L,2b1 & I/LI/E & {-}/{-}{-}/{-} & BF & CP:2b W:1a,2b \\
C1653 & 1a<o1 / 2a<1a2 2b<1a2L / 3a<2a1L,2b2 & I/LI/E & {-}/{-}{-}/{-} & BF & CP:2b W:1a,2b \\
C1654 & 1a<o1 / 2a<1a2 2b<1a2L / 3a<2a2,2b0L & I/LI/D & {-}/{-}{-}/{-} & B & CP:2b W:1a,2b \\
C1655 & 1a<o1 / 2a<1a2 2b<1a2L / 3a<2a2,2b1 & I/LI/E & {-}/{-}{-}/{-} & BF & CP:2b W:1a,2b \\
C1656 & 1a<o1 / 2a<1a2 2b<1a2L / 3a<2a2,2b1L & I/LI/I & {-}/{-}{-}/{-} & BF & CP:2b,3a W:1a,2b,3a \\
C1657 & 1a<o1 / 2a<1a2 2b<1a2L / 3a<2a2,2b2 & I/LI/E & {-}/{-}{-}/{-} & BF & CP:2b W:1a,2b \\
C1658 & 1a<o1 / 2a<1a2 2b<1a2L / 3a<2a2,2b2L & I/LI/I & {-}/{-}{-}/{-} & BF & CP:2b,3a W:1a,2b,3a \\
C1659 & 1a<o1 / 2a<1a2 2b<1a2L / 3a<2a2L,2b1 & I/LI/E & {-}/{-}{-}/{-} & BF & CP:2b W:1a,2b \\
C1660 & 1a<o1 / 2a<1a2 2b<1a2L / 3a<2a2L,2b2 & I/LI/E & {-}/{-}{-}/{-} & BF & CP:2b W:1a,2b \\
C1661 & 1a<o1 1b<o1 / 2a<1a0L,1b1 / 3a<2a0L & EE/D/N & {-}{-}/{-}/{-} & E &  \\
C1662 & 1a<o1 1b<o1 / 2a<1a0L,1b1 / 3a<2a1 & EE/D/D & {-}{-}/{-}/{-} & E &  \\
C1663 & 1a<o1 1b<o1 / 2a<1a0L,1b1 / 3a<2a1L & EE/D/D & {-}{-}/{-}/{-} & E &  \\
C1664 & 1a<o1 1b<o1 / 2a<1a0L,1b1 / 3a<2a2 & EE/D/D & {-}{-}/{-}/{-} & E &  \\
C1665 & 1a<o1 1b<o1 / 2a<1a0L,1b1 / 3a<2a2L & EE/D/D & {-}{-}/{-}/{-} & E &  \\
C1666 & 1a<o1 1b<o1 / 2a<1a0L,1b2 / 3a<2a0L & EE/D/N & {-}{-}/{-}/{-} & E &  \\
C1667 & 1a<o1 1b<o1 / 2a<1a0L,1b2 / 3a<2a1 & EE/D/D & {-}{-}/{-}/{-} & E &  \\
C1668 & 1a<o1 1b<o1 / 2a<1a0L,1b2 / 3a<2a1L & EE/D/D & {-}{-}/{-}/{-} & E &  \\
C1669 & 1a<o1 1b<o1 / 2a<1a0L,1b2 / 3a<2a2 & EE/D/D & {-}{-}/{-}/{-} & E &  \\
C1670 & 1a<o1 1b<o1 / 2a<1a0L,1b2 / 3a<2a2L & EE/D/D & {-}{-}/{-}/{-} & E &  \\
C1671 & 1a<o1 1b<o1 / 2a<1a1,1b1 / 3a<2a0L & EE/D/N & {-}{-}/{-}/{-} & E &  \\
C1672 & 1a<o1 1b<o1 / 2a<1a1,1b1 / 3a<2a1 & EE/D/D & {-}{-}/{-}/{-} & E &  \\
C1673 & 1a<o1 1b<o1 / 2a<1a1,1b1 / 3a<2a1L & EE/D/D & {-}{-}/{-}/{-} & E &  \\
C1674 & 1a<o1 1b<o1 / 2a<1a1,1b1 / 3a<2a2 & EE/D/D & {-}{-}/{-}/{-} & E &  \\
C1675 & 1a<o1 1b<o1 / 2a<1a1,1b1 / 3a<2a2L & EE/D/D & {-}{-}/{-}/{-} & E &  \\
C1676 & 1a<o1 1b<o1 / 2a<1a1,1b1L / 3a<2a0L & EE/D/N & {-}{-}/{-}/{-} & E &  \\
C1677 & 1a<o1 1b<o1 / 2a<1a1,1b1L / 3a<2a1 & EE/D/D & {-}{-}/{-}/{-} & E &  \\
C1678 & 1a<o1 1b<o1 / 2a<1a1,1b1L / 3a<2a1L & EE/D/D & {-}{-}/{-}/{-} & E &  \\
C1679 & 1a<o1 1b<o1 / 2a<1a1,1b1L / 3a<2a2 & EE/D/D & {-}{-}/{-}/{-} & E &  \\
C1680 & 1a<o1 1b<o1 / 2a<1a1,1b1L / 3a<2a2L & EE/D/D & {-}{-}/{-}/{-} & E &  \\
C1681 & 1a<o1 1b<o1 / 2a<1a1,1b2 / 3a<2a0L & EE/D/N & {-}{-}/{-}/{-} & E &  \\
C1682 & 1a<o1 1b<o1 / 2a<1a1,1b2 / 3a<2a1 & EE/D/D & {-}{-}/{-}/{-} & E &  \\
C1683 & 1a<o1 1b<o1 / 2a<1a1,1b2 / 3a<2a1L & EE/D/D & {-}{-}/{-}/{-} & E &  \\
C1684 & 1a<o1 1b<o1 / 2a<1a1,1b2 / 3a<2a2 & EE/D/D & {-}{-}/{-}/{-} & E &  \\
C1685 & 1a<o1 1b<o1 / 2a<1a1,1b2 / 3a<2a2L & EE/D/D & {-}{-}/{-}/{-} & E &  \\
C1686 & 1a<o1 1b<o1 / 2a<1a1,1b2L / 3a<2a0L & EE/D/N & {-}{-}/{-}/{-} & E &  \\
C1687 & 1a<o1 1b<o1 / 2a<1a1,1b2L / 3a<2a1 & EE/D/D & {-}{-}/{-}/{-} & E &  \\
C1688 & 1a<o1 1b<o1 / 2a<1a1,1b2L / 3a<2a1L & EE/D/D & {-}{-}/{-}/{-} & E &  \\
C1689 & 1a<o1 1b<o1 / 2a<1a1,1b2L / 3a<2a2 & EE/D/D & {-}{-}/{-}/{-} & E &  \\
C1690 & 1a<o1 1b<o1 / 2a<1a1,1b2L / 3a<2a2L & EE/D/D & {-}{-}/{-}/{-} & E &  \\
C1691 & 1a<o1 1b<o1 / 2a<1a1L,1b2 / 3a<2a0L & EE/D/N & {-}{-}/{-}/{-} & E &  \\
C1692 & 1a<o1 1b<o1 / 2a<1a1L,1b2 / 3a<2a1 & EE/D/D & {-}{-}/{-}/{-} & E &  \\
C1693 & 1a<o1 1b<o1 / 2a<1a1L,1b2 / 3a<2a1L & EE/D/D & {-}{-}/{-}/{-} & E &  \\
C1694 & 1a<o1 1b<o1 / 2a<1a1L,1b2 / 3a<2a2 & EE/D/D & {-}{-}/{-}/{-} & E &  \\
C1695 & 1a<o1 1b<o1 / 2a<1a1L,1b2 / 3a<2a2L & EE/D/D & {-}{-}/{-}/{-} & E &  \\
C1696 & 1a<o1 1b<o1 / 2a<1a2,1b2 / 3a<2a0L & EE/D/N & {-}{-}/{-}/{-} & E &  \\
C1697 & 1a<o1 1b<o1 / 2a<1a2,1b2 / 3a<2a1 & EE/D/D & {-}{-}/{-}/{-} & E &  \\
C1698 & 1a<o1 1b<o1 / 2a<1a2,1b2 / 3a<2a1L & EE/D/D & {-}{-}/{-}/{-} & E &  \\
C1699 & 1a<o1 1b<o1 / 2a<1a2,1b2 / 3a<2a2 & EE/D/D & {-}{-}/{-}/{-} & E &  \\
C1700 & 1a<o1 1b<o1 / 2a<1a2,1b2 / 3a<2a2L & EE/D/D & {-}{-}/{-}/{-} & E &  \\
C1701 & 1a<o1 1b<o1 / 2a<1a2,1b2L / 3a<2a0L & EE/D/N & {-}{-}/{-}/{-} & E &  \\
C1702 & 1a<o1 1b<o1 / 2a<1a2,1b2L / 3a<2a1 & EE/D/D & {-}{-}/{-}/{-} & E &  \\
C1703 & 1a<o1 1b<o1 / 2a<1a2,1b2L / 3a<2a1L & EE/D/D & {-}{-}/{-}/{-} & E &  \\
C1704 & 1a<o1 1b<o1 / 2a<1a2,1b2L / 3a<2a2 & EE/D/D & {-}{-}/{-}/{-} & E &  \\
C1705 & 1a<o1 1b<o1 / 2a<1a2,1b2L / 3a<2a2L & EE/D/D & {-}{-}/{-}/{-} & E &  \\
C1706 & 1a<o1 1b<o1L / 2a<1a0L,1b1 / 3a<2a0L & LI/I/N & {-}+/{-}/{-} & B & CPW:1b,2a \\
C1707 & 1a<o1 1b<o1L / 2a<1a0L,1b1 / 3a<2a1 & LI/I/I & {-}+/{-}/{-} & B & CPW:1b,2a,3a \\
C1708 & 1a<o1 1b<o1L / 2a<1a0L,1b1 / 3a<2a1L & LI/I/I & {-}+/{-}/{-} & B & CPW:1b,2a,3a \\
C1709 & 1a<o1 1b<o1L / 2a<1a0L,1b1 / 3a<2a2 & LI/I/I & {-}+/{-}/{-} & B & CPW:1b,2a,3a \\
C1710 & 1a<o1 1b<o1L / 2a<1a0L,1b1 / 3a<2a2L & LI/I/I & {-}+/{-}/{-} & B & CPW:1b,2a,3a \\
C1711 & 1a<o1 1b<o1L / 2a<1a0L,1b2 / 3a<2a0L & LI/I/N & {-}+/{-}/{-} & B & CPW:1b,2a \\
C1712 & 1a<o1 1b<o1L / 2a<1a0L,1b2 / 3a<2a1 & LI/I/I & {-}+/{-}/{-} & B & CPW:1b,2a,3a \\
C1713 & 1a<o1 1b<o1L / 2a<1a0L,1b2 / 3a<2a1L & LI/I/I & {-}+/{-}/{-} & B & CPW:1b,2a,3a \\
C1714 & 1a<o1 1b<o1L / 2a<1a0L,1b2 / 3a<2a2 & LI/I/I & {-}+/{-}/{-} & B & CPW:1b,2a,3a \\
C1715 & 1a<o1 1b<o1L / 2a<1a0L,1b2 / 3a<2a2L & LI/I/I & {-}+/{-}/{-} & B & CPW:1b,2a,3a \\
C1716 & 1a<o1 1b<o1L / 2a<1a1,1b0L / 3a<2a0L & LI/D/N & {-}+/+/+ & B & CPW:1b \\
C1717 & 1a<o1 1b<o1L / 2a<1a1,1b0L / 3a<2a1 & LI/D/D & {-}+/+/{-} & B & CPW:1b \\
C1718 & 1a<o1 1b<o1L / 2a<1a1,1b0L / 3a<2a1L & LI/D/D & {-}+/+/+ & B & CPW:1b \\
C1719 & 1a<o1 1b<o1L / 2a<1a1,1b0L / 3a<2a2 & LI/D/D & {-}+/+/{-} & B & CPW:1b \\
C1720 & 1a<o1 1b<o1L / 2a<1a1,1b0L / 3a<2a2L & LI/D/D & {-}+/+/+ & B & CPW:1b \\
C1721 & 1a<o1 1b<o1L / 2a<1a1,1b1 / 3a<2a0L & LI/E/N & {-}+/{-}/{-} & BF & CPW:1b \\
C1722 & 1a<o1 1b<o1L / 2a<1a1,1b1 / 3a<2a1 & LI/E/D & {-}+/{-}/{-} & BF & CPW:1b \\
C1723 & 1a<o1 1b<o1L / 2a<1a1,1b1 / 3a<2a1L & LI/E/D & {-}+/{-}/{-} & BF & CPW:1b \\
C1724 & 1a<o1 1b<o1L / 2a<1a1,1b1 / 3a<2a2 & LI/E/D & {-}+/{-}/{-} & BF & CPW:1b \\
C1725 & 1a<o1 1b<o1L / 2a<1a1,1b1 / 3a<2a2L & LI/E/D & {-}+/{-}/{-} & BF & CPW:1b \\
C1726 & 1a<o1 1b<o1L / 2a<1a1,1b1L / 3a<2a0L & LI/I/N & {-}+/+/+ & BF & CPW:1b,2a \\
C1727 & 1a<o1 1b<o1L / 2a<1a1,1b1L / 3a<2a1 & LI/I/I & {-}+/+/{-} & BF & CPW:1b,2a,3a \\
C1728 & 1a<o1 1b<o1L / 2a<1a1,1b1L / 3a<2a1L & LI/I/I & {-}+/+/+ & BF & CPW:1b,2a,3a \\
C1729 & 1a<o1 1b<o1L / 2a<1a1,1b1L / 3a<2a2 & LI/I/I & {-}+/+/{-} & BF & CPW:1b,2a,3a \\
C1730 & 1a<o1 1b<o1L / 2a<1a1,1b1L / 3a<2a2L & LI/I/I & {-}+/+/+ & BF & CPW:1b,2a,3a \\
C1731 & 1a<o1 1b<o1L / 2a<1a1,1b2 / 3a<2a0L & LI/E/N & {-}+/{-}/{-} & BF & CPW:1b \\
C1732 & 1a<o1 1b<o1L / 2a<1a1,1b2 / 3a<2a1 & LI/E/D & {-}+/{-}/{-} & BF & CPW:1b \\
C1733 & 1a<o1 1b<o1L / 2a<1a1,1b2 / 3a<2a1L & LI/E/D & {-}+/{-}/{-} & BF & CPW:1b \\
C1734 & 1a<o1 1b<o1L / 2a<1a1,1b2 / 3a<2a2 & LI/E/D & {-}+/{-}/{-} & BF & CPW:1b \\
C1735 & 1a<o1 1b<o1L / 2a<1a1,1b2 / 3a<2a2L & LI/E/D & {-}+/{-}/{-} & BF & CPW:1b \\
C1736 & 1a<o1 1b<o1L / 2a<1a1,1b2L / 3a<2a0L & LI/I/N & {-}+/+/+ & BF & CPW:1b,2a \\
C1737 & 1a<o1 1b<o1L / 2a<1a1,1b2L / 3a<2a1 & LI/I/I & {-}+/+/{-} & BF & CPW:1b,2a,3a \\
C1738 & 1a<o1 1b<o1L / 2a<1a1,1b2L / 3a<2a1L & LI/I/I & {-}+/+/+ & BF & CPW:1b,2a,3a \\
C1739 & 1a<o1 1b<o1L / 2a<1a1,1b2L / 3a<2a2 & LI/I/I & {-}+/+/{-} & BF & CPW:1b,2a,3a \\
C1740 & 1a<o1 1b<o1L / 2a<1a1,1b2L / 3a<2a2L & LI/I/I & {-}+/+/+ & BF & CPW:1b,2a,3a \\
C1741 & 1a<o1 1b<o1L / 2a<1a1L,1b1 / 3a<2a0L & LI/E/N & {-}+/{-}/{-} & BF & CPW:1b \\
C1742 & 1a<o1 1b<o1L / 2a<1a1L,1b1 / 3a<2a1 & LI/E/D & {-}+/{-}/{-} & BF & CPW:1b \\
C1743 & 1a<o1 1b<o1L / 2a<1a1L,1b1 / 3a<2a1L & LI/E/D & {-}+/{-}/{-} & BF & CPW:1b \\
C1744 & 1a<o1 1b<o1L / 2a<1a1L,1b1 / 3a<2a2 & LI/E/D & {-}+/{-}/{-} & BF & CPW:1b \\
C1745 & 1a<o1 1b<o1L / 2a<1a1L,1b1 / 3a<2a2L & LI/E/D & {-}+/{-}/{-} & BF & CPW:1b \\
C1746 & 1a<o1 1b<o1L / 2a<1a1L,1b2 / 3a<2a0L & LI/E/N & {-}+/{-}/{-} & BF & CPW:1b \\
C1747 & 1a<o1 1b<o1L / 2a<1a1L,1b2 / 3a<2a1 & LI/E/D & {-}+/{-}/{-} & BF & CPW:1b \\
C1748 & 1a<o1 1b<o1L / 2a<1a1L,1b2 / 3a<2a1L & LI/E/D & {-}+/{-}/{-} & BF & CPW:1b \\
C1749 & 1a<o1 1b<o1L / 2a<1a1L,1b2 / 3a<2a2 & LI/E/D & {-}+/{-}/{-} & BF & CPW:1b \\
C1750 & 1a<o1 1b<o1L / 2a<1a1L,1b2 / 3a<2a2L & LI/E/D & {-}+/{-}/{-} & BF & CPW:1b \\
C1751 & 1a<o1 1b<o1L / 2a<1a2,1b0L / 3a<2a0L & LI/D/N & {-}+/+/+ & B & CPW:1b \\
C1752 & 1a<o1 1b<o1L / 2a<1a2,1b0L / 3a<2a1 & LI/D/D & {-}+/+/{-} & B & CPW:1b \\
C1753 & 1a<o1 1b<o1L / 2a<1a2,1b0L / 3a<2a1L & LI/D/D & {-}+/+/+ & B & CPW:1b \\
C1754 & 1a<o1 1b<o1L / 2a<1a2,1b0L / 3a<2a2 & LI/D/D & {-}+/+/{-} & B & CPW:1b \\
C1755 & 1a<o1 1b<o1L / 2a<1a2,1b0L / 3a<2a2L & LI/D/D & {-}+/+/+ & B & CPW:1b \\
C1756 & 1a<o1 1b<o1L / 2a<1a2,1b1 / 3a<2a0L & LI/E/N & {-}+/{-}/{-} & BF & CPW:1b \\
C1757 & 1a<o1 1b<o1L / 2a<1a2,1b1 / 3a<2a1 & LI/E/D & {-}+/{-}/{-} & BF & CPW:1b \\
C1758 & 1a<o1 1b<o1L / 2a<1a2,1b1 / 3a<2a1L & LI/E/D & {-}+/{-}/{-} & BF & CPW:1b \\
C1759 & 1a<o1 1b<o1L / 2a<1a2,1b1 / 3a<2a2 & LI/E/D & {-}+/{-}/{-} & BF & CPW:1b \\
C1760 & 1a<o1 1b<o1L / 2a<1a2,1b1 / 3a<2a2L & LI/E/D & {-}+/{-}/{-} & BF & CPW:1b \\
C1761 & 1a<o1 1b<o1L / 2a<1a2,1b1L / 3a<2a0L & LI/I/N & {-}+/+/+ & BF & CPW:1b,2a \\
C1762 & 1a<o1 1b<o1L / 2a<1a2,1b1L / 3a<2a1 & LI/I/I & {-}+/+/{-} & BF & CPW:1b,2a,3a \\
C1763 & 1a<o1 1b<o1L / 2a<1a2,1b1L / 3a<2a1L & LI/I/I & {-}+/+/+ & BF & CPW:1b,2a,3a \\
C1764 & 1a<o1 1b<o1L / 2a<1a2,1b1L / 3a<2a2 & LI/I/I & {-}+/+/{-} & BF & CPW:1b,2a,3a \\
C1765 & 1a<o1 1b<o1L / 2a<1a2,1b1L / 3a<2a2L & LI/I/I & {-}+/+/+ & BF & CPW:1b,2a,3a \\
C1766 & 1a<o1 1b<o1L / 2a<1a2,1b2 / 3a<2a0L & LI/E/N & {-}+/{-}/{-} & BF & CPW:1b \\
C1767 & 1a<o1 1b<o1L / 2a<1a2,1b2 / 3a<2a1 & LI/E/D & {-}+/{-}/{-} & BF & CPW:1b \\
C1768 & 1a<o1 1b<o1L / 2a<1a2,1b2 / 3a<2a1L & LI/E/D & {-}+/{-}/{-} & BF & CPW:1b \\
C1769 & 1a<o1 1b<o1L / 2a<1a2,1b2 / 3a<2a2 & LI/E/D & {-}+/{-}/{-} & BF & CPW:1b \\
C1770 & 1a<o1 1b<o1L / 2a<1a2,1b2 / 3a<2a2L & LI/E/D & {-}+/{-}/{-} & BF & CPW:1b \\
C1771 & 1a<o1 1b<o1L / 2a<1a2,1b2L / 3a<2a0L & LI/I/N & {-}+/+/+ & BF & CPW:1b,2a \\
C1772 & 1a<o1 1b<o1L / 2a<1a2,1b2L / 3a<2a1 & LI/I/I & {-}+/+/{-} & BF & CPW:1b,2a,3a \\
C1773 & 1a<o1 1b<o1L / 2a<1a2,1b2L / 3a<2a1L & LI/I/I & {-}+/+/+ & BF & CPW:1b,2a,3a \\
C1774 & 1a<o1 1b<o1L / 2a<1a2,1b2L / 3a<2a2 & LI/I/I & {-}+/+/{-} & BF & CPW:1b,2a,3a \\
C1775 & 1a<o1 1b<o1L / 2a<1a2,1b2L / 3a<2a2L & LI/I/I & {-}+/+/+ & BF & CPW:1b,2a,3a \\
C1776 & 1a<o1 1b<o1L / 2a<1a2L,1b1 / 3a<2a0L & LI/E/N & {-}+/{-}/{-} & BF & CPW:1b \\
C1777 & 1a<o1 1b<o1L / 2a<1a2L,1b1 / 3a<2a1 & LI/E/D & {-}+/{-}/{-} & BF & CPW:1b \\
C1778 & 1a<o1 1b<o1L / 2a<1a2L,1b1 / 3a<2a1L & LI/E/D & {-}+/{-}/{-} & BF & CPW:1b \\
C1779 & 1a<o1 1b<o1L / 2a<1a2L,1b1 / 3a<2a2 & LI/E/D & {-}+/{-}/{-} & BF & CPW:1b \\
C1780 & 1a<o1 1b<o1L / 2a<1a2L,1b1 / 3a<2a2L & LI/E/D & {-}+/{-}/{-} & BF & CPW:1b \\
C1781 & 1a<o1 1b<o1L / 2a<1a2L,1b2 / 3a<2a0L & LI/E/N & {-}+/{-}/{-} & BF & CPW:1b \\
C1782 & 1a<o1 1b<o1L / 2a<1a2L,1b2 / 3a<2a1 & LI/E/D & {-}+/{-}/{-} & BF & CPW:1b \\
C1783 & 1a<o1 1b<o1L / 2a<1a2L,1b2 / 3a<2a1L & LI/E/D & {-}+/{-}/{-} & BF & CPW:1b \\
C1784 & 1a<o1 1b<o1L / 2a<1a2L,1b2 / 3a<2a2 & LI/E/D & {-}+/{-}/{-} & BF & CPW:1b \\
C1785 & 1a<o1 1b<o1L / 2a<1a2L,1b2 / 3a<2a2L & LI/E/D & {-}+/{-}/{-} & BF & CPW:1b \\
C1786 & 1a<o1 1b<o2 / 2a<1a0L,1b1 / 3a<2a0L & EE/D/N & {-}{-}/{-}/{-} & E & C:1b,2a \\
C1787 & 1a<o1 1b<o2 / 2a<1a0L,1b1 / 3a<2a1 & EE/D/D & {-}{-}/{-}/{-} & E & C:1b,2a,3a \\
C1788 & 1a<o1 1b<o2 / 2a<1a0L,1b1 / 3a<2a1L & EE/D/D & {-}{-}/{-}/{-} & E & C:1b,2a,3a \\
C1789 & 1a<o1 1b<o2 / 2a<1a0L,1b1 / 3a<2a2 & EE/D/D & {-}{-}/{-}/{-} & E & C:1b,2a,3a \\
C1790 & 1a<o1 1b<o2 / 2a<1a0L,1b1 / 3a<2a2L & EE/D/D & {-}{-}/{-}/{-} & E & C:1b,2a,3a \\
C1791 & 1a<o1 1b<o2 / 2a<1a0L,1b2 / 3a<2a0L & EE/D/N & {-}{-}/{-}/{-} & E & C:1b,2a \\
C1792 & 1a<o1 1b<o2 / 2a<1a0L,1b2 / 3a<2a1 & EE/D/D & {-}{-}/{-}/{-} & E & C:1b,2a,3a \\
C1793 & 1a<o1 1b<o2 / 2a<1a0L,1b2 / 3a<2a1L & EE/D/D & {-}{-}/{-}/{-} & E & C:1b,2a,3a \\
C1794 & 1a<o1 1b<o2 / 2a<1a0L,1b2 / 3a<2a2 & EE/D/D & {-}{-}/{-}/{-} & E & C:1b,2a,3a \\
C1795 & 1a<o1 1b<o2 / 2a<1a0L,1b2 / 3a<2a2L & EE/D/D & {-}{-}/{-}/{-} & E & C:1b,2a,3a \\
C1796 & 1a<o1 1b<o2 / 2a<1a1,1b0L / 3a<2a0L & EE/D/N & {-}{-}/{-}/{-} & E & C:1b \\
C1797 & 1a<o1 1b<o2 / 2a<1a1,1b0L / 3a<2a1 & EE/D/D & {-}{-}/{-}/{-} & E & C:1b \\
C1798 & 1a<o1 1b<o2 / 2a<1a1,1b0L / 3a<2a1L & EE/D/D & {-}{-}/{-}/{-} & E & C:1b \\
C1799 & 1a<o1 1b<o2 / 2a<1a1,1b0L / 3a<2a2 & EE/D/D & {-}{-}/{-}/{-} & E & C:1b \\
C1800 & 1a<o1 1b<o2 / 2a<1a1,1b0L / 3a<2a2L & EE/D/D & {-}{-}/{-}/{-} & E & C:1b \\
C1801 & 1a<o1 1b<o2 / 2a<1a1,1b1 / 3a<2a0L & EE/D/N & {-}{-}/{-}/{-} & E & C:1b \\
C1802 & 1a<o1 1b<o2 / 2a<1a1,1b1 / 3a<2a1 & EE/D/D & {-}{-}/{-}/{-} & E & C:1b \\
C1803 & 1a<o1 1b<o2 / 2a<1a1,1b1 / 3a<2a1L & EE/D/D & {-}{-}/{-}/{-} & E & C:1b \\
C1804 & 1a<o1 1b<o2 / 2a<1a1,1b1 / 3a<2a2 & EE/D/D & {-}{-}/{-}/{-} & E & C:1b \\
C1805 & 1a<o1 1b<o2 / 2a<1a1,1b1 / 3a<2a2L & EE/D/D & {-}{-}/{-}/{-} & E & C:1b \\
C1806 & 1a<o1 1b<o2 / 2a<1a1,1b1L / 3a<2a0L & EE/D/N & {-}{-}/{-}/{-} & E & C:1b \\
C1807 & 1a<o1 1b<o2 / 2a<1a1,1b1L / 3a<2a1 & EE/D/D & {-}{-}/{-}/{-} & E & C:1b \\
C1808 & 1a<o1 1b<o2 / 2a<1a1,1b1L / 3a<2a1L & EE/D/D & {-}{-}/{-}/{-} & E & C:1b \\
C1809 & 1a<o1 1b<o2 / 2a<1a1,1b1L / 3a<2a2 & EE/D/D & {-}{-}/{-}/{-} & E & C:1b \\
C1810 & 1a<o1 1b<o2 / 2a<1a1,1b1L / 3a<2a2L & EE/D/D & {-}{-}/{-}/{-} & E & C:1b \\
C1811 & 1a<o1 1b<o2 / 2a<1a1,1b2 / 3a<2a0L & EE/D/N & {-}{-}/{-}/{-} & E & C:1b,2a \\
C1812 & 1a<o1 1b<o2 / 2a<1a1,1b2 / 3a<2a1 & EE/D/D & {-}{-}/{-}/{-} & E & C:1b,2a,3a \\
C1813 & 1a<o1 1b<o2 / 2a<1a1,1b2 / 3a<2a1L & EE/D/D & {-}{-}/{-}/{-} & E & C:1b,2a,3a \\
C1814 & 1a<o1 1b<o2 / 2a<1a1,1b2 / 3a<2a2 & EE/D/D & {-}{-}/{-}/{-} & E & C:1b,2a,3a \\
C1815 & 1a<o1 1b<o2 / 2a<1a1,1b2 / 3a<2a2L & EE/D/D & {-}{-}/{-}/{-} & E & C:1b,2a,3a \\
C1816 & 1a<o1 1b<o2 / 2a<1a1,1b2L / 3a<2a0L & EE/D/N & {-}{-}/{-}/{-} & E & C:1b,2a \\
C1817 & 1a<o1 1b<o2 / 2a<1a1,1b2L / 3a<2a1 & EE/D/D & {-}{-}/{-}/{-} & E & C:1b,2a,3a \\
C1818 & 1a<o1 1b<o2 / 2a<1a1,1b2L / 3a<2a1L & EE/D/D & {-}{-}/{-}/{-} & E & C:1b,2a,3a \\
C1819 & 1a<o1 1b<o2 / 2a<1a1,1b2L / 3a<2a2 & EE/D/D & {-}{-}/{-}/{-} & E & C:1b,2a,3a \\
C1820 & 1a<o1 1b<o2 / 2a<1a1,1b2L / 3a<2a2L & EE/D/D & {-}{-}/{-}/{-} & E & C:1b,2a,3a \\
C1821 & 1a<o1 1b<o2 / 2a<1a1L,1b1 / 3a<2a0L & EE/D/N & {-}{-}/{-}/{-} & E & C:1b \\
C1822 & 1a<o1 1b<o2 / 2a<1a1L,1b1 / 3a<2a1 & EE/D/D & {-}{-}/{-}/{-} & E & C:1b \\
C1823 & 1a<o1 1b<o2 / 2a<1a1L,1b1 / 3a<2a1L & EE/D/D & {-}{-}/{-}/{-} & E & C:1b \\
C1824 & 1a<o1 1b<o2 / 2a<1a1L,1b1 / 3a<2a2 & EE/D/D & {-}{-}/{-}/{-} & E & C:1b \\
C1825 & 1a<o1 1b<o2 / 2a<1a1L,1b1 / 3a<2a2L & EE/D/D & {-}{-}/{-}/{-} & E & C:1b \\
C1826 & 1a<o1 1b<o2 / 2a<1a1L,1b2 / 3a<2a0L & EE/D/N & {-}{-}/{-}/{-} & E & C:1b,2a \\
C1827 & 1a<o1 1b<o2 / 2a<1a1L,1b2 / 3a<2a1 & EE/D/D & {-}{-}/{-}/{-} & E & C:1b,2a,3a \\
C1828 & 1a<o1 1b<o2 / 2a<1a1L,1b2 / 3a<2a1L & EE/D/D & {-}{-}/{-}/{-} & E & C:1b,2a,3a \\
C1829 & 1a<o1 1b<o2 / 2a<1a1L,1b2 / 3a<2a2 & EE/D/D & {-}{-}/{-}/{-} & E & C:1b,2a,3a \\
C1830 & 1a<o1 1b<o2 / 2a<1a1L,1b2 / 3a<2a2L & EE/D/D & {-}{-}/{-}/{-} & E & C:1b,2a,3a \\
C1831 & 1a<o1 1b<o2 / 2a<1a2,1b0L / 3a<2a0L & EE/D/N & {-}{-}/{-}/{-} & E & C:1b \\
C1832 & 1a<o1 1b<o2 / 2a<1a2,1b0L / 3a<2a1 & EE/D/D & {-}{-}/{-}/{-} & E & C:1b \\
C1833 & 1a<o1 1b<o2 / 2a<1a2,1b0L / 3a<2a1L & EE/D/D & {-}{-}/{-}/{-} & E & C:1b \\
C1834 & 1a<o1 1b<o2 / 2a<1a2,1b0L / 3a<2a2 & EE/D/D & {-}{-}/{-}/{-} & E & C:1b \\
C1835 & 1a<o1 1b<o2 / 2a<1a2,1b0L / 3a<2a2L & EE/D/D & {-}{-}/{-}/{-} & E & C:1b \\
C1836 & 1a<o1 1b<o2 / 2a<1a2,1b1 / 3a<2a0L & EE/D/N & {-}{-}/{-}/{-} & E & C:1b \\
C1837 & 1a<o1 1b<o2 / 2a<1a2,1b1 / 3a<2a1 & EE/D/D & {-}{-}/{-}/{-} & E & C:1b \\
C1838 & 1a<o1 1b<o2 / 2a<1a2,1b1 / 3a<2a1L & EE/D/D & {-}{-}/{-}/{-} & E & C:1b \\
C1839 & 1a<o1 1b<o2 / 2a<1a2,1b1 / 3a<2a2 & EE/D/D & {-}{-}/{-}/{-} & E & C:1b \\
C1840 & 1a<o1 1b<o2 / 2a<1a2,1b1 / 3a<2a2L & EE/D/D & {-}{-}/{-}/{-} & E & C:1b \\
C1841 & 1a<o1 1b<o2 / 2a<1a2,1b1L / 3a<2a0L & EE/D/N & {-}{-}/{-}/{-} & E & C:1b \\
C1842 & 1a<o1 1b<o2 / 2a<1a2,1b1L / 3a<2a1 & EE/D/D & {-}{-}/{-}/{-} & E & C:1b \\
C1843 & 1a<o1 1b<o2 / 2a<1a2,1b1L / 3a<2a1L & EE/D/D & {-}{-}/{-}/{-} & E & C:1b \\
C1844 & 1a<o1 1b<o2 / 2a<1a2,1b1L / 3a<2a2 & EE/D/D & {-}{-}/{-}/{-} & E & C:1b \\
C1845 & 1a<o1 1b<o2 / 2a<1a2,1b1L / 3a<2a2L & EE/D/D & {-}{-}/{-}/{-} & E & C:1b \\
C1846 & 1a<o1 1b<o2 / 2a<1a2,1b2 / 3a<2a0L & EE/D/N & {-}{-}/{-}/{-} & E & C:1b \\
C1847 & 1a<o1 1b<o2 / 2a<1a2,1b2 / 3a<2a1 & EE/D/D & {-}{-}/{-}/{-} & E & C:1b \\
C1848 & 1a<o1 1b<o2 / 2a<1a2,1b2 / 3a<2a1L & EE/D/D & {-}{-}/{-}/{-} & E & C:1b \\
C1849 & 1a<o1 1b<o2 / 2a<1a2,1b2 / 3a<2a2 & EE/D/D & {-}{-}/{-}/{-} & E & C:1b \\
C1850 & 1a<o1 1b<o2 / 2a<1a2,1b2 / 3a<2a2L & EE/D/D & {-}{-}/{-}/{-} & E & C:1b \\
C1851 & 1a<o1 1b<o2 / 2a<1a2,1b2L / 3a<2a0L & EE/D/N & {-}{-}/{-}/{-} & E & C:1b \\
C1852 & 1a<o1 1b<o2 / 2a<1a2,1b2L / 3a<2a1 & EE/D/D & {-}{-}/{-}/{-} & E & C:1b \\
C1853 & 1a<o1 1b<o2 / 2a<1a2,1b2L / 3a<2a1L & EE/D/D & {-}{-}/{-}/{-} & E & C:1b \\
C1854 & 1a<o1 1b<o2 / 2a<1a2,1b2L / 3a<2a2 & EE/D/D & {-}{-}/{-}/{-} & E & C:1b \\
C1855 & 1a<o1 1b<o2 / 2a<1a2,1b2L / 3a<2a2L & EE/D/D & {-}{-}/{-}/{-} & E & C:1b \\
C1856 & 1a<o1 1b<o2 / 2a<1a2L,1b1 / 3a<2a0L & EE/D/N & {-}{-}/{-}/{-} & E & C:1b \\
C1857 & 1a<o1 1b<o2 / 2a<1a2L,1b1 / 3a<2a1 & EE/D/D & {-}{-}/{-}/{-} & E & C:1b \\
C1858 & 1a<o1 1b<o2 / 2a<1a2L,1b1 / 3a<2a1L & EE/D/D & {-}{-}/{-}/{-} & E & C:1b \\
C1859 & 1a<o1 1b<o2 / 2a<1a2L,1b1 / 3a<2a2 & EE/D/D & {-}{-}/{-}/{-} & E & C:1b \\
C1860 & 1a<o1 1b<o2 / 2a<1a2L,1b1 / 3a<2a2L & EE/D/D & {-}{-}/{-}/{-} & E & C:1b \\
C1861 & 1a<o1 1b<o2 / 2a<1a2L,1b2 / 3a<2a0L & EE/D/N & {-}{-}/{-}/{-} & E & C:1b \\
C1862 & 1a<o1 1b<o2 / 2a<1a2L,1b2 / 3a<2a1 & EE/D/D & {-}{-}/{-}/{-} & E & C:1b \\
C1863 & 1a<o1 1b<o2 / 2a<1a2L,1b2 / 3a<2a1L & EE/D/D & {-}{-}/{-}/{-} & E & C:1b \\
C1864 & 1a<o1 1b<o2 / 2a<1a2L,1b2 / 3a<2a2 & EE/D/D & {-}{-}/{-}/{-} & E & C:1b \\
C1865 & 1a<o1 1b<o2 / 2a<1a2L,1b2 / 3a<2a2L & EE/D/D & {-}{-}/{-}/{-} & E & C:1b \\
C1866 & 1a<o1 1b<o2L / 2a<1a0L,1b1 / 3a<2a0L & LI/I/N & {-}+/{-}/{-} & B & PW:1b,2a \\
C1867 & 1a<o1 1b<o2L / 2a<1a0L,1b1 / 3a<2a1 & LI/I/I & {-}+/{-}/{-} & B & PW:1b,2a,3a \\
C1868 & 1a<o1 1b<o2L / 2a<1a0L,1b1 / 3a<2a1L & LI/I/I & {-}+/{-}/{-} & B & PW:1b,2a,3a \\
C1869 & 1a<o1 1b<o2L / 2a<1a0L,1b1 / 3a<2a2 & LI/I/I & {-}+/{-}/{-} & B & PW:1b,2a,3a \\
C1870 & 1a<o1 1b<o2L / 2a<1a0L,1b1 / 3a<2a2L & LI/I/I & {-}+/{-}/{-} & B & PW:1b,2a,3a \\
C1871 & 1a<o1 1b<o2L / 2a<1a0L,1b2 / 3a<2a0L & LI/I/N & {-}+/{-}/{-} & B & PW:1b,2a \\
C1872 & 1a<o1 1b<o2L / 2a<1a0L,1b2 / 3a<2a1 & LI/I/I & {-}+/{-}/{-} & B & PW:1b,2a,3a \\
C1873 & 1a<o1 1b<o2L / 2a<1a0L,1b2 / 3a<2a1L & LI/I/I & {-}+/{-}/{-} & B & PW:1b,2a,3a \\
C1874 & 1a<o1 1b<o2L / 2a<1a0L,1b2 / 3a<2a2 & LI/I/I & {-}+/{-}/{-} & B & PW:1b,2a,3a \\
C1875 & 1a<o1 1b<o2L / 2a<1a0L,1b2 / 3a<2a2L & LI/I/I & {-}+/{-}/{-} & B & PW:1b,2a,3a \\
C1876 & 1a<o1 1b<o2L / 2a<1a1,1b0L / 3a<2a0L & LI/D/N & {-}+/+/+ & B & PW:1b \\
C1877 & 1a<o1 1b<o2L / 2a<1a1,1b0L / 3a<2a1 & LI/D/D & {-}+/+/{-} & B & PW:1b \\
C1878 & 1a<o1 1b<o2L / 2a<1a1,1b0L / 3a<2a1L & LI/D/D & {-}+/+/+ & B & PW:1b \\
C1879 & 1a<o1 1b<o2L / 2a<1a1,1b0L / 3a<2a2 & LI/D/D & {-}+/+/{-} & B & PW:1b \\
C1880 & 1a<o1 1b<o2L / 2a<1a1,1b0L / 3a<2a2L & LI/D/D & {-}+/+/+ & B & PW:1b \\
C1881 & 1a<o1 1b<o2L / 2a<1a1,1b1 / 3a<2a0L & LI/E/N & {-}+/{-}/{-} & BF & PW:1b \\
C1882 & 1a<o1 1b<o2L / 2a<1a1,1b1 / 3a<2a1 & LI/E/D & {-}+/{-}/{-} & BF & PW:1b \\
C1883 & 1a<o1 1b<o2L / 2a<1a1,1b1 / 3a<2a1L & LI/E/D & {-}+/{-}/{-} & BF & PW:1b \\
C1884 & 1a<o1 1b<o2L / 2a<1a1,1b1 / 3a<2a2 & LI/E/D & {-}+/{-}/{-} & BF & PW:1b \\
C1885 & 1a<o1 1b<o2L / 2a<1a1,1b1 / 3a<2a2L & LI/E/D & {-}+/{-}/{-} & BF & PW:1b \\
C1886 & 1a<o1 1b<o2L / 2a<1a1,1b1L / 3a<2a0L & LI/I/N & {-}+/+/+ & BF & C:2a PW:1b,2a \\
C1887 & 1a<o1 1b<o2L / 2a<1a1,1b1L / 3a<2a1 & LI/I/I & {-}+/+/{-} & BF & C:2a,3a PW:1b,2a,3a \\
C1888 & 1a<o1 1b<o2L / 2a<1a1,1b1L / 3a<2a1L & LI/I/I & {-}+/+/+ & BF & C:2a,3a PW:1b,2a,3a \\
C1889 & 1a<o1 1b<o2L / 2a<1a1,1b1L / 3a<2a2 & LI/I/I & {-}+/+/{-} & BF & C:2a,3a PW:1b,2a,3a \\
C1890 & 1a<o1 1b<o2L / 2a<1a1,1b1L / 3a<2a2L & LI/I/I & {-}+/+/+ & BF & C:2a,3a PW:1b,2a,3a \\
C1891 & 1a<o1 1b<o2L / 2a<1a1,1b2 / 3a<2a0L & LI/E/N & {-}+/{-}/{-} & BF & C:2a PW:1b \\
C1892 & 1a<o1 1b<o2L / 2a<1a1,1b2 / 3a<2a1 & LI/E/D & {-}+/{-}/{-} & BF & C:2a,3a PW:1b \\
C1893 & 1a<o1 1b<o2L / 2a<1a1,1b2 / 3a<2a1L & LI/E/D & {-}+/{-}/{-} & BF & C:2a,3a PW:1b \\
C1894 & 1a<o1 1b<o2L / 2a<1a1,1b2 / 3a<2a2 & LI/E/D & {-}+/{-}/{-} & BF & C:2a,3a PW:1b \\
C1895 & 1a<o1 1b<o2L / 2a<1a1,1b2 / 3a<2a2L & LI/E/D & {-}+/{-}/{-} & BF & C:2a,3a PW:1b \\
C1896 & 1a<o1 1b<o2L / 2a<1a1,1b2L / 3a<2a0L & LI/I/N & {-}+/+/+ & BF & PW:1b,2a \\
C1897 & 1a<o1 1b<o2L / 2a<1a1,1b2L / 3a<2a1 & LI/I/I & {-}+/+/{-} & BF & PW:1b,2a,3a \\
C1898 & 1a<o1 1b<o2L / 2a<1a1,1b2L / 3a<2a1L & LI/I/I & {-}+/+/+ & BF & PW:1b,2a,3a \\
C1899 & 1a<o1 1b<o2L / 2a<1a1,1b2L / 3a<2a2 & LI/I/I & {-}+/+/{-} & BF & PW:1b,2a,3a \\
C1900 & 1a<o1 1b<o2L / 2a<1a1,1b2L / 3a<2a2L & LI/I/I & {-}+/+/+ & BF & PW:1b,2a,3a \\
C1901 & 1a<o1 1b<o2L / 2a<1a1L,1b1 / 3a<2a0L & LI/E/N & {-}+/{-}/{-} & BF & PW:1b \\
C1902 & 1a<o1 1b<o2L / 2a<1a1L,1b1 / 3a<2a1 & LI/E/D & {-}+/{-}/{-} & BF & PW:1b \\
C1903 & 1a<o1 1b<o2L / 2a<1a1L,1b1 / 3a<2a1L & LI/E/D & {-}+/{-}/{-} & BF & PW:1b \\
C1904 & 1a<o1 1b<o2L / 2a<1a1L,1b1 / 3a<2a2 & LI/E/D & {-}+/{-}/{-} & BF & PW:1b \\
C1905 & 1a<o1 1b<o2L / 2a<1a1L,1b1 / 3a<2a2L & LI/E/D & {-}+/{-}/{-} & BF & PW:1b \\
C1906 & 1a<o1 1b<o2L / 2a<1a1L,1b2 / 3a<2a0L & LI/E/N & {-}+/{-}/{-} & BF & C:2a PW:1b \\
C1907 & 1a<o1 1b<o2L / 2a<1a1L,1b2 / 3a<2a1 & LI/E/D & {-}+/{-}/{-} & BF & C:2a,3a PW:1b \\
C1908 & 1a<o1 1b<o2L / 2a<1a1L,1b2 / 3a<2a1L & LI/E/D & {-}+/{-}/{-} & BF & C:2a,3a PW:1b \\
C1909 & 1a<o1 1b<o2L / 2a<1a1L,1b2 / 3a<2a2 & LI/E/D & {-}+/{-}/{-} & BF & C:2a,3a PW:1b \\
C1910 & 1a<o1 1b<o2L / 2a<1a1L,1b2 / 3a<2a2L & LI/E/D & {-}+/{-}/{-} & BF & C:2a,3a PW:1b \\
C1911 & 1a<o1 1b<o2L / 2a<1a2,1b0L / 3a<2a0L & LI/D/N & {-}+/+/+ & B & PW:1b \\
C1912 & 1a<o1 1b<o2L / 2a<1a2,1b0L / 3a<2a1 & LI/D/D & {-}+/+/{-} & B & PW:1b \\
C1913 & 1a<o1 1b<o2L / 2a<1a2,1b0L / 3a<2a1L & LI/D/D & {-}+/+/+ & B & PW:1b \\
C1914 & 1a<o1 1b<o2L / 2a<1a2,1b0L / 3a<2a2 & LI/D/D & {-}+/+/{-} & B & PW:1b \\
C1915 & 1a<o1 1b<o2L / 2a<1a2,1b0L / 3a<2a2L & LI/D/D & {-}+/+/+ & B & PW:1b \\
C1916 & 1a<o1 1b<o2L / 2a<1a2,1b1 / 3a<2a0L & LI/E/N & {-}+/{-}/{-} & BF & PW:1b \\
C1917 & 1a<o1 1b<o2L / 2a<1a2,1b1 / 3a<2a1 & LI/E/D & {-}+/{-}/{-} & BF & PW:1b \\
C1918 & 1a<o1 1b<o2L / 2a<1a2,1b1 / 3a<2a1L & LI/E/D & {-}+/{-}/{-} & BF & PW:1b \\
C1919 & 1a<o1 1b<o2L / 2a<1a2,1b1 / 3a<2a2 & LI/E/D & {-}+/{-}/{-} & BF & PW:1b \\
C1920 & 1a<o1 1b<o2L / 2a<1a2,1b1 / 3a<2a2L & LI/E/D & {-}+/{-}/{-} & BF & PW:1b \\
C1921 & 1a<o1 1b<o2L / 2a<1a2,1b1L / 3a<2a0L & LI/I/N & {-}+/+/+ & BF & C:2a PW:1b,2a \\
C1922 & 1a<o1 1b<o2L / 2a<1a2,1b1L / 3a<2a1 & LI/I/I & {-}+/+/{-} & BF & C:2a,3a PW:1b,2a,3a \\
C1923 & 1a<o1 1b<o2L / 2a<1a2,1b1L / 3a<2a1L & LI/I/I & {-}+/+/+ & BF & C:2a,3a PW:1b,2a,3a \\
C1924 & 1a<o1 1b<o2L / 2a<1a2,1b1L / 3a<2a2 & LI/I/I & {-}+/+/{-} & BF & C:2a,3a PW:1b,2a,3a \\
C1925 & 1a<o1 1b<o2L / 2a<1a2,1b1L / 3a<2a2L & LI/I/I & {-}+/+/+ & BF & C:2a,3a PW:1b,2a,3a \\
C1926 & 1a<o1 1b<o2L / 2a<1a2,1b2 / 3a<2a0L & LI/E/N & {-}+/{-}/{-} & BF & PW:1b \\
C1927 & 1a<o1 1b<o2L / 2a<1a2,1b2 / 3a<2a1 & LI/E/D & {-}+/{-}/{-} & BF & PW:1b \\
C1928 & 1a<o1 1b<o2L / 2a<1a2,1b2 / 3a<2a1L & LI/E/D & {-}+/{-}/{-} & BF & PW:1b \\
C1929 & 1a<o1 1b<o2L / 2a<1a2,1b2 / 3a<2a2 & LI/E/D & {-}+/{-}/{-} & BF & PW:1b \\
C1930 & 1a<o1 1b<o2L / 2a<1a2,1b2 / 3a<2a2L & LI/E/D & {-}+/{-}/{-} & BF & PW:1b \\
C1931 & 1a<o1 1b<o2L / 2a<1a2,1b2L / 3a<2a0L & LI/I/N & {-}+/+/+ & BF & C:2a PW:1b,2a \\
C1932 & 1a<o1 1b<o2L / 2a<1a2,1b2L / 3a<2a1 & LI/I/I & {-}+/+/{-} & BF & C:2a,3a PW:1b,2a,3a \\
C1933 & 1a<o1 1b<o2L / 2a<1a2,1b2L / 3a<2a1L & LI/I/I & {-}+/+/+ & BF & C:2a,3a PW:1b,2a,3a \\
C1934 & 1a<o1 1b<o2L / 2a<1a2,1b2L / 3a<2a2 & LI/I/I & {-}+/+/{-} & BF & C:2a,3a PW:1b,2a,3a \\
C1935 & 1a<o1 1b<o2L / 2a<1a2,1b2L / 3a<2a2L & LI/I/I & {-}+/+/+ & BF & C:2a,3a PW:1b,2a,3a \\
C1936 & 1a<o1 1b<o2L / 2a<1a2L,1b1 / 3a<2a0L & LI/E/N & {-}+/{-}/{-} & BF & PW:1b \\
C1937 & 1a<o1 1b<o2L / 2a<1a2L,1b1 / 3a<2a1 & LI/E/D & {-}+/{-}/{-} & BF & PW:1b \\
C1938 & 1a<o1 1b<o2L / 2a<1a2L,1b1 / 3a<2a1L & LI/E/D & {-}+/{-}/{-} & BF & PW:1b \\
C1939 & 1a<o1 1b<o2L / 2a<1a2L,1b1 / 3a<2a2 & LI/E/D & {-}+/{-}/{-} & BF & PW:1b \\
C1940 & 1a<o1 1b<o2L / 2a<1a2L,1b1 / 3a<2a2L & LI/E/D & {-}+/{-}/{-} & BF & PW:1b \\
C1941 & 1a<o1 1b<o2L / 2a<1a2L,1b2 / 3a<2a0L & LI/E/N & {-}+/{-}/{-} & BF & PW:1b \\
C1942 & 1a<o1 1b<o2L / 2a<1a2L,1b2 / 3a<2a1 & LI/E/D & {-}+/{-}/{-} & BF & PW:1b \\
C1943 & 1a<o1 1b<o2L / 2a<1a2L,1b2 / 3a<2a1L & LI/E/D & {-}+/{-}/{-} & BF & PW:1b \\
C1944 & 1a<o1 1b<o2L / 2a<1a2L,1b2 / 3a<2a2 & LI/E/D & {-}+/{-}/{-} & BF & PW:1b \\
C1945 & 1a<o1 1b<o2L / 2a<1a2L,1b2 / 3a<2a2L & LI/E/D & {-}+/{-}/{-} & BF & PW:1b \\
C1946 & 1a<o1L / 2a<1a0L 2b<1a1 / 3a<2a0L,2b1 & I/NI/I & +/+{-}/+ & N &  \\
C1947 & 1a<o1L / 2a<1a0L 2b<1a1 / 3a<2a0L,2b2 & I/NI/I & +/+{-}/+ & N &  \\
C1948 & 1a<o1L / 2a<1a0L 2b<1a1 / 3a<2a1,2b0L & I/NI/N & +/+{-}/{-} & N &  \\
C1949 & 1a<o1L / 2a<1a0L 2b<1a1 / 3a<2a1,2b1 & I/NI/E & +/+{-}/{-} & F & W:3a \\
C1950 & 1a<o1L / 2a<1a0L 2b<1a1 / 3a<2a1,2b1L & I/NI/I & +/+{-}/{-} & F & CP:3a \\
C1951 & 1a<o1L / 2a<1a0L 2b<1a1 / 3a<2a1,2b2 & I/NI/E & +/+{-}/{-} & F & CW:3a \\
C1952 & 1a<o1L / 2a<1a0L 2b<1a1 / 3a<2a1,2b2L & I/NI/I & +/+{-}/{-} & F & P:3a \\
C1953 & 1a<o1L / 2a<1a0L 2b<1a1 / 3a<2a1L,2b1 & I/NI/E & +/+{-}/+ & F & W:3a \\
C1954 & 1a<o1L / 2a<1a0L 2b<1a1 / 3a<2a1L,2b2 & I/NI/E & +/+{-}/+ & F & CW:3a \\
C1955 & 1a<o1L / 2a<1a0L 2b<1a1 / 3a<2a2,2b0L & I/NI/N & +/+{-}/{-} & N &  \\
C1956 & 1a<o1L / 2a<1a0L 2b<1a1 / 3a<2a2,2b1 & I/NI/E & +/+{-}/{-} & F & W:3a \\
C1957 & 1a<o1L / 2a<1a0L 2b<1a1 / 3a<2a2,2b1L & I/NI/I & +/+{-}/{-} & F & CP:3a \\
C1958 & 1a<o1L / 2a<1a0L 2b<1a1 / 3a<2a2,2b2 & I/NI/E & +/+{-}/{-} & F & W:3a \\
C1959 & 1a<o1L / 2a<1a0L 2b<1a1 / 3a<2a2,2b2L & I/NI/I & +/+{-}/{-} & F & CP:3a \\
C1960 & 1a<o1L / 2a<1a0L 2b<1a1 / 3a<2a2L,2b1 & I/NI/E & +/+{-}/+ & F & W:3a \\
C1961 & 1a<o1L / 2a<1a0L 2b<1a1 / 3a<2a2L,2b2 & I/NI/E & +/+{-}/+ & F & W:3a \\
C1962 & 1a<o1L / 2a<1a0L 2b<1a2 / 3a<2a0L,2b1 & I/NI/I & +/+{-}/+ & N &  \\
C1963 & 1a<o1L / 2a<1a0L 2b<1a2 / 3a<2a0L,2b2 & I/NI/I & +/+{-}/+ & N &  \\
C1964 & 1a<o1L / 2a<1a0L 2b<1a2 / 3a<2a1,2b0L & I/NI/N & +/+{-}/{-} & N &  \\
C1965 & 1a<o1L / 2a<1a0L 2b<1a2 / 3a<2a1,2b1 & I/NI/E & +/+{-}/{-} & F & W:3a \\
C1966 & 1a<o1L / 2a<1a0L 2b<1a2 / 3a<2a1,2b1L & I/NI/I & +/+{-}/{-} & F & CP:3a \\
C1967 & 1a<o1L / 2a<1a0L 2b<1a2 / 3a<2a1,2b2 & I/NI/E & +/+{-}/{-} & F & CW:3a \\
C1968 & 1a<o1L / 2a<1a0L 2b<1a2 / 3a<2a1,2b2L & I/NI/I & +/+{-}/{-} & F & P:3a \\
C1969 & 1a<o1L / 2a<1a0L 2b<1a2 / 3a<2a1L,2b1 & I/NI/E & +/+{-}/+ & F & W:3a \\
C1970 & 1a<o1L / 2a<1a0L 2b<1a2 / 3a<2a1L,2b2 & I/NI/E & +/+{-}/+ & F & CW:3a \\
C1971 & 1a<o1L / 2a<1a0L 2b<1a2 / 3a<2a2,2b0L & I/NI/N & +/+{-}/{-} & N &  \\
C1972 & 1a<o1L / 2a<1a0L 2b<1a2 / 3a<2a2,2b1 & I/NI/E & +/+{-}/{-} & F & W:3a \\
C1973 & 1a<o1L / 2a<1a0L 2b<1a2 / 3a<2a2,2b1L & I/NI/I & +/+{-}/{-} & F & CP:3a \\
C1974 & 1a<o1L / 2a<1a0L 2b<1a2 / 3a<2a2,2b2 & I/NI/E & +/+{-}/{-} & F & W:3a \\
C1975 & 1a<o1L / 2a<1a0L 2b<1a2 / 3a<2a2,2b2L & I/NI/I & +/+{-}/{-} & F & CP:3a \\
C1976 & 1a<o1L / 2a<1a0L 2b<1a2 / 3a<2a2L,2b1 & I/NI/E & +/+{-}/+ & F & W:3a \\
C1977 & 1a<o1L / 2a<1a0L 2b<1a2 / 3a<2a2L,2b2 & I/NI/E & +/+{-}/+ & F & W:3a \\
C1978 & 1a<o1L / 2a<1a1 2b<1a1 / 3a<2a0L,2b1 & I/EE/D & +/{-}{-}/{-} & E & W:1a \\
C1979 & 1a<o1L / 2a<1a1 2b<1a1 / 3a<2a0L,2b2 & I/EE/D & +/{-}{-}/{-} & E & W:1a \\
C1980 & 1a<o1L / 2a<1a1 2b<1a1 / 3a<2a1,2b1 & I/EE/D & +/{-}{-}/{-} & E & W:1a \\
C1981 & 1a<o1L / 2a<1a1 2b<1a1 / 3a<2a1,2b1L & I/EE/D & +/{-}{-}/{-} & E & W:1a \\
C1982 & 1a<o1L / 2a<1a1 2b<1a1 / 3a<2a1,2b2 & I/EE/D & +/{-}{-}/{-} & E & W:1a \\
C1983 & 1a<o1L / 2a<1a1 2b<1a1 / 3a<2a1,2b2L & I/EE/D & +/{-}{-}/{-} & E & W:1a \\
C1984 & 1a<o1L / 2a<1a1 2b<1a1 / 3a<2a1L,2b2 & I/EE/D & +/{-}{-}/{-} & E & W:1a \\
C1985 & 1a<o1L / 2a<1a1 2b<1a1 / 3a<2a2,2b2 & I/EE/D & +/{-}{-}/{-} & E & W:1a \\
C1986 & 1a<o1L / 2a<1a1 2b<1a1 / 3a<2a2,2b2L & I/EE/D & +/{-}{-}/{-} & E & W:1a \\
C1987 & 1a<o1L / 2a<1a1 2b<1a1L / 3a<2a0L,2b1 & I/LI/I & +/{-}+/{-} & B & CP:2b,3a W:1a,2b,3a \\
C1988 & 1a<o1L / 2a<1a1 2b<1a1L / 3a<2a0L,2b2 & I/LI/I & +/{-}+/{-} & B & CP:2b,3a W:1a,2b,3a \\
C1989 & 1a<o1L / 2a<1a1 2b<1a1L / 3a<2a1,2b0L & I/LI/D & +/{-}+/+ & B & CP:2b W:1a,2b \\
C1990 & 1a<o1L / 2a<1a1 2b<1a1L / 3a<2a1,2b1 & I/LI/E & +/{-}+/{-} & BF & CP:2b W:1a,2b \\
C1991 & 1a<o1L / 2a<1a1 2b<1a1L / 3a<2a1,2b1L & I/LI/I & +/{-}+/+ & BF & CP:2b,3a W:1a,2b,3a \\
C1992 & 1a<o1L / 2a<1a1 2b<1a1L / 3a<2a1,2b2 & I/LI/E & +/{-}+/{-} & BF & CP:2b W:1a,2b \\
C1993 & 1a<o1L / 2a<1a1 2b<1a1L / 3a<2a1,2b2L & I/LI/I & +/{-}+/+ & BF & CP:2b,3a W:1a,2b,3a \\
C1994 & 1a<o1L / 2a<1a1 2b<1a1L / 3a<2a1L,2b1 & I/LI/E & +/{-}+/{-} & BF & CP:2b W:1a,2b \\
C1995 & 1a<o1L / 2a<1a1 2b<1a1L / 3a<2a1L,2b2 & I/LI/E & +/{-}+/{-} & BF & CP:2b W:1a,2b \\
C1996 & 1a<o1L / 2a<1a1 2b<1a1L / 3a<2a2,2b0L & I/LI/D & +/{-}+/+ & B & CP:2b W:1a,2b \\
C1997 & 1a<o1L / 2a<1a1 2b<1a1L / 3a<2a2,2b1 & I/LI/E & +/{-}+/{-} & BF & CP:2b W:1a,2b \\
C1998 & 1a<o1L / 2a<1a1 2b<1a1L / 3a<2a2,2b1L & I/LI/I & +/{-}+/+ & BF & CP:2b,3a W:1a,2b,3a \\
C1999 & 1a<o1L / 2a<1a1 2b<1a1L / 3a<2a2,2b2 & I/LI/E & +/{-}+/{-} & BF & CP:2b W:1a,2b \\
C2000 & 1a<o1L / 2a<1a1 2b<1a1L / 3a<2a2,2b2L & I/LI/I & +/{-}+/+ & BF & CP:2b,3a W:1a,2b,3a \\
C2001 & 1a<o1L / 2a<1a1 2b<1a1L / 3a<2a2L,2b1 & I/LI/E & +/{-}+/{-} & BF & CP:2b W:1a,2b \\
C2002 & 1a<o1L / 2a<1a1 2b<1a1L / 3a<2a2L,2b2 & I/LI/E & +/{-}+/{-} & BF & CP:2b W:1a,2b \\
C2003 & 1a<o1L / 2a<1a1 2b<1a2 / 3a<2a0L,2b1 & I/EE/D & +/{-}{-}/{-} & E & C:2b,3a W:1a \\
C2004 & 1a<o1L / 2a<1a1 2b<1a2 / 3a<2a0L,2b2 & I/EE/D & +/{-}{-}/{-} & E & C:2b,3a W:1a \\
C2005 & 1a<o1L / 2a<1a1 2b<1a2 / 3a<2a1,2b0L & I/EE/D & +/{-}{-}/{-} & E & C:2b W:1a \\
C2006 & 1a<o1L / 2a<1a1 2b<1a2 / 3a<2a1,2b1 & I/EE/D & +/{-}{-}/{-} & E & C:2b W:1a \\
C2007 & 1a<o1L / 2a<1a1 2b<1a2 / 3a<2a1,2b1L & I/EE/D & +/{-}{-}/{-} & E & C:2b W:1a \\
C2008 & 1a<o1L / 2a<1a1 2b<1a2 / 3a<2a1,2b2 & I/EE/D & +/{-}{-}/{-} & E & C:2b,3a W:1a \\
C2009 & 1a<o1L / 2a<1a1 2b<1a2 / 3a<2a1,2b2L & I/EE/D & +/{-}{-}/{-} & E & C:2b,3a W:1a \\
C2010 & 1a<o1L / 2a<1a1 2b<1a2 / 3a<2a1L,2b1 & I/EE/D & +/{-}{-}/{-} & E & C:2b W:1a \\
C2011 & 1a<o1L / 2a<1a1 2b<1a2 / 3a<2a1L,2b2 & I/EE/D & +/{-}{-}/{-} & E & C:2b,3a W:1a \\
C2012 & 1a<o1L / 2a<1a1 2b<1a2 / 3a<2a2,2b0L & I/EE/D & +/{-}{-}/{-} & E & C:2b W:1a \\
C2013 & 1a<o1L / 2a<1a1 2b<1a2 / 3a<2a2,2b1 & I/EE/D & +/{-}{-}/{-} & E & C:2b W:1a \\
C2014 & 1a<o1L / 2a<1a1 2b<1a2 / 3a<2a2,2b1L & I/EE/D & +/{-}{-}/{-} & E & C:2b W:1a \\
C2015 & 1a<o1L / 2a<1a1 2b<1a2 / 3a<2a2,2b2 & I/EE/D & +/{-}{-}/{-} & E & C:2b W:1a \\
C2016 & 1a<o1L / 2a<1a1 2b<1a2 / 3a<2a2,2b2L & I/EE/D & +/{-}{-}/{-} & E & C:2b W:1a \\
C2017 & 1a<o1L / 2a<1a1 2b<1a2 / 3a<2a2L,2b1 & I/EE/D & +/{-}{-}/{-} & E & C:2b W:1a \\
C2018 & 1a<o1L / 2a<1a1 2b<1a2 / 3a<2a2L,2b2 & I/EE/D & +/{-}{-}/{-} & E & C:2b W:1a \\
C2019 & 1a<o1L / 2a<1a1 2b<1a2L / 3a<2a0L,2b1 & I/LI/I & +/{-}+/{-} & B & P:2b,3a W:1a,2b,3a \\
C2020 & 1a<o1L / 2a<1a1 2b<1a2L / 3a<2a0L,2b2 & I/LI/I & +/{-}+/{-} & B & P:2b,3a W:1a,2b,3a \\
C2021 & 1a<o1L / 2a<1a1 2b<1a2L / 3a<2a1,2b0L & I/LI/D & +/{-}+/+ & B & P:2b W:1a,2b \\
C2022 & 1a<o1L / 2a<1a1 2b<1a2L / 3a<2a1,2b1 & I/LI/E & +/{-}+/{-} & BF & P:2b W:1a,2b \\
C2023 & 1a<o1L / 2a<1a1 2b<1a2L / 3a<2a1,2b1L & I/LI/I & +/{-}+/+ & BF & C:3a P:2b,3a W:1a,2b,3a \\
C2024 & 1a<o1L / 2a<1a1 2b<1a2L / 3a<2a1,2b2 & I/LI/E & +/{-}+/{-} & BF & C:3a P:2b W:1a,2b \\
C2025 & 1a<o1L / 2a<1a1 2b<1a2L / 3a<2a1,2b2L & I/LI/I & +/{-}+/+ & BF & P:2b,3a W:1a,2b,3a \\
C2026 & 1a<o1L / 2a<1a1 2b<1a2L / 3a<2a1L,2b1 & I/LI/E & +/{-}+/{-} & BF & P:2b W:1a,2b \\
C2027 & 1a<o1L / 2a<1a1 2b<1a2L / 3a<2a1L,2b2 & I/LI/E & +/{-}+/{-} & BF & C:3a P:2b W:1a,2b \\
C2028 & 1a<o1L / 2a<1a1 2b<1a2L / 3a<2a2,2b0L & I/LI/D & +/{-}+/+ & B & P:2b W:1a,2b \\
C2029 & 1a<o1L / 2a<1a1 2b<1a2L / 3a<2a2,2b1 & I/LI/E & +/{-}+/{-} & BF & P:2b W:1a,2b \\
C2030 & 1a<o1L / 2a<1a1 2b<1a2L / 3a<2a2,2b1L & I/LI/I & +/{-}+/+ & BF & C:3a P:2b,3a W:1a,2b,3a \\
C2031 & 1a<o1L / 2a<1a1 2b<1a2L / 3a<2a2,2b2 & I/LI/E & +/{-}+/{-} & BF & P:2b W:1a,2b \\
C2032 & 1a<o1L / 2a<1a1 2b<1a2L / 3a<2a2,2b2L & I/LI/I & +/{-}+/+ & BF & C:3a P:2b,3a W:1a,2b,3a \\
C2033 & 1a<o1L / 2a<1a1 2b<1a2L / 3a<2a2L,2b1 & I/LI/E & +/{-}+/{-} & BF & P:2b W:1a,2b \\
C2034 & 1a<o1L / 2a<1a1 2b<1a2L / 3a<2a2L,2b2 & I/LI/E & +/{-}+/{-} & BF & P:2b W:1a,2b \\
C2035 & 1a<o1L / 2a<1a1L 2b<1a2 / 3a<2a0L,2b1 & I/IL/D & +/+{-}/+ & B & C:2a,2b,3a P:2a W:1a,2a \\
C2036 & 1a<o1L / 2a<1a1L 2b<1a2 / 3a<2a0L,2b2 & I/IL/D & +/+{-}/+ & B & C:2a,2b,3a P:2a W:1a,2a \\
C2037 & 1a<o1L / 2a<1a1L 2b<1a2 / 3a<2a1,2b0L & I/IL/I & +/+{-}/{-} & B & C:2a,2b,3a P:2a,3a W:1a,2a,3a \\
C2038 & 1a<o1L / 2a<1a1L 2b<1a2 / 3a<2a1,2b1 & I/IL/E & +/+{-}/{-} & BF & C:2a,2b P:2a W:1a,2a \\
C2039 & 1a<o1L / 2a<1a1L 2b<1a2 / 3a<2a1,2b1L & I/IL/E & +/+{-}/{-} & BF & C:2a,2b P:2a W:1a,2a \\
C2040 & 1a<o1L / 2a<1a1L 2b<1a2 / 3a<2a1,2b2 & I/IL/E & +/+{-}/{-} & BF & C:2a,2b,3a P:2a W:1a,2a \\
C2041 & 1a<o1L / 2a<1a1L 2b<1a2 / 3a<2a1,2b2L & I/IL/E & +/+{-}/{-} & BF & C:2a,2b,3a P:2a W:1a,2a \\
C2042 & 1a<o1L / 2a<1a1L 2b<1a2 / 3a<2a1L,2b1 & I/IL/I & +/+{-}/+ & BF & C:2a,2b,3a P:2a,3a W:1a,2a,3a \\
C2043 & 1a<o1L / 2a<1a1L 2b<1a2 / 3a<2a1L,2b2 & I/IL/I & +/+{-}/+ & BF & C:2a,2b P:2a,3a W:1a,2a,3a \\
C2044 & 1a<o1L / 2a<1a1L 2b<1a2 / 3a<2a2,2b0L & I/IL/I & +/+{-}/{-} & B & C:2a,2b,3a P:2a,3a W:1a,2a,3a \\
C2045 & 1a<o1L / 2a<1a1L 2b<1a2 / 3a<2a2,2b1 & I/IL/E & +/+{-}/{-} & BF & C:2a,2b P:2a W:1a,2a \\
C2046 & 1a<o1L / 2a<1a1L 2b<1a2 / 3a<2a2,2b1L & I/IL/E & +/+{-}/{-} & BF & C:2a,2b P:2a W:1a,2a \\
C2047 & 1a<o1L / 2a<1a1L 2b<1a2 / 3a<2a2,2b2 & I/IL/E & +/+{-}/{-} & BF & C:2a,2b P:2a W:1a,2a \\
C2048 & 1a<o1L / 2a<1a1L 2b<1a2 / 3a<2a2,2b2L & I/IL/E & +/+{-}/{-} & BF & C:2a,2b P:2a W:1a,2a \\
C2049 & 1a<o1L / 2a<1a1L 2b<1a2 / 3a<2a2L,2b1 & I/IL/I & +/+{-}/+ & BF & C:2a,2b,3a P:2a,3a W:1a,2a,3a \\
C2050 & 1a<o1L / 2a<1a1L 2b<1a2 / 3a<2a2L,2b2 & I/IL/I & +/+{-}/+ & BF & C:2a,2b,3a P:2a,3a W:1a,2a,3a \\
C2051 & 1a<o1L / 2a<1a2 2b<1a2 / 3a<2a0L,2b1 & I/EE/D & +/{-}{-}/{-} & E & W:1a \\
C2052 & 1a<o1L / 2a<1a2 2b<1a2 / 3a<2a0L,2b2 & I/EE/D & +/{-}{-}/{-} & E & W:1a \\
C2053 & 1a<o1L / 2a<1a2 2b<1a2 / 3a<2a1,2b1 & I/EE/D & +/{-}{-}/{-} & E & W:1a \\
C2054 & 1a<o1L / 2a<1a2 2b<1a2 / 3a<2a1,2b1L & I/EE/D & +/{-}{-}/{-} & E & W:1a \\
C2055 & 1a<o1L / 2a<1a2 2b<1a2 / 3a<2a1,2b2 & I/EE/D & +/{-}{-}/{-} & E & W:1a \\
C2056 & 1a<o1L / 2a<1a2 2b<1a2 / 3a<2a1,2b2L & I/EE/D & +/{-}{-}/{-} & E & W:1a \\
C2057 & 1a<o1L / 2a<1a2 2b<1a2 / 3a<2a1L,2b2 & I/EE/D & +/{-}{-}/{-} & E & W:1a \\
C2058 & 1a<o1L / 2a<1a2 2b<1a2 / 3a<2a2,2b2 & I/EE/D & +/{-}{-}/{-} & E & W:1a \\
C2059 & 1a<o1L / 2a<1a2 2b<1a2 / 3a<2a2,2b2L & I/EE/D & +/{-}{-}/{-} & E & W:1a \\
C2060 & 1a<o1L / 2a<1a2 2b<1a2L / 3a<2a0L,2b1 & I/LI/I & +/{-}+/{-} & B & CP:2b,3a W:1a,2b,3a \\
C2061 & 1a<o1L / 2a<1a2 2b<1a2L / 3a<2a0L,2b2 & I/LI/I & +/{-}+/{-} & B & CP:2b,3a W:1a,2b,3a \\
C2062 & 1a<o1L / 2a<1a2 2b<1a2L / 3a<2a1,2b0L & I/LI/D & +/{-}+/+ & B & CP:2b W:1a,2b \\
C2063 & 1a<o1L / 2a<1a2 2b<1a2L / 3a<2a1,2b1 & I/LI/E & +/{-}+/{-} & BF & CP:2b W:1a,2b \\
C2064 & 1a<o1L / 2a<1a2 2b<1a2L / 3a<2a1,2b1L & I/LI/I & +/{-}+/+ & BF & CP:2b,3a W:1a,2b,3a \\
C2065 & 1a<o1L / 2a<1a2 2b<1a2L / 3a<2a1,2b2 & I/LI/E & +/{-}+/{-} & BF & CP:2b W:1a,2b \\
C2066 & 1a<o1L / 2a<1a2 2b<1a2L / 3a<2a1,2b2L & I/LI/I & +/{-}+/+ & BF & CP:2b,3a W:1a,2b,3a \\
C2067 & 1a<o1L / 2a<1a2 2b<1a2L / 3a<2a1L,2b1 & I/LI/E & +/{-}+/{-} & BF & CP:2b W:1a,2b \\
C2068 & 1a<o1L / 2a<1a2 2b<1a2L / 3a<2a1L,2b2 & I/LI/E & +/{-}+/{-} & BF & CP:2b W:1a,2b \\
C2069 & 1a<o1L / 2a<1a2 2b<1a2L / 3a<2a2,2b0L & I/LI/D & +/{-}+/+ & B & CP:2b W:1a,2b \\
C2070 & 1a<o1L / 2a<1a2 2b<1a2L / 3a<2a2,2b1 & I/LI/E & +/{-}+/{-} & BF & CP:2b W:1a,2b \\
C2071 & 1a<o1L / 2a<1a2 2b<1a2L / 3a<2a2,2b1L & I/LI/I & +/{-}+/+ & BF & CP:2b,3a W:1a,2b,3a \\
C2072 & 1a<o1L / 2a<1a2 2b<1a2L / 3a<2a2,2b2 & I/LI/E & +/{-}+/{-} & BF & CP:2b W:1a,2b \\
C2073 & 1a<o1L / 2a<1a2 2b<1a2L / 3a<2a2,2b2L & I/LI/I & +/{-}+/+ & BF & CP:2b,3a W:1a,2b,3a \\
C2074 & 1a<o1L / 2a<1a2 2b<1a2L / 3a<2a2L,2b1 & I/LI/E & +/{-}+/{-} & BF & CP:2b W:1a,2b \\
C2075 & 1a<o1L / 2a<1a2 2b<1a2L / 3a<2a2L,2b2 & I/LI/E & +/{-}+/{-} & BF & CP:2b W:1a,2b \\
C2076 & 1a<o1L 1b<o2 / 2a<1a0L,1b1 / 3a<2a0L & IL/D/N & +{-}/+/+ & B & C:1a,1b,2a PW:1a \\
C2077 & 1a<o1L 1b<o2 / 2a<1a0L,1b1 / 3a<2a1 & IL/D/D & +{-}/+/{-} & B & C:1a,1b,2a,3a PW:1a \\
C2078 & 1a<o1L 1b<o2 / 2a<1a0L,1b1 / 3a<2a1L & IL/D/D & +{-}/+/+ & B & C:1a,1b,2a,3a PW:1a \\
C2079 & 1a<o1L 1b<o2 / 2a<1a0L,1b1 / 3a<2a2 & IL/D/D & +{-}/+/{-} & B & C:1a,1b,2a,3a PW:1a \\
C2080 & 1a<o1L 1b<o2 / 2a<1a0L,1b1 / 3a<2a2L & IL/D/D & +{-}/+/+ & B & C:1a,1b,2a,3a PW:1a \\
C2081 & 1a<o1L 1b<o2 / 2a<1a0L,1b2 / 3a<2a0L & IL/D/N & +{-}/+/+ & B & C:1a,1b,2a PW:1a \\
C2082 & 1a<o1L 1b<o2 / 2a<1a0L,1b2 / 3a<2a1 & IL/D/D & +{-}/+/{-} & B & C:1a,1b,2a,3a PW:1a \\
C2083 & 1a<o1L 1b<o2 / 2a<1a0L,1b2 / 3a<2a1L & IL/D/D & +{-}/+/+ & B & C:1a,1b,2a,3a PW:1a \\
C2084 & 1a<o1L 1b<o2 / 2a<1a0L,1b2 / 3a<2a2 & IL/D/D & +{-}/+/{-} & B & C:1a,1b,2a,3a PW:1a \\
C2085 & 1a<o1L 1b<o2 / 2a<1a0L,1b2 / 3a<2a2L & IL/D/D & +{-}/+/+ & B & C:1a,1b,2a,3a PW:1a \\
C2086 & 1a<o1L 1b<o2 / 2a<1a1,1b0L / 3a<2a0L & IL/I/N & +{-}/{-}/{-} & B & C:1a,1b,2a PW:1a,2a \\
C2087 & 1a<o1L 1b<o2 / 2a<1a1,1b0L / 3a<2a1 & IL/I/I & +{-}/{-}/{-} & B & C:1a,1b,2a,3a PW:1a,2a,3a \\
C2088 & 1a<o1L 1b<o2 / 2a<1a1,1b0L / 3a<2a1L & IL/I/I & +{-}/{-}/{-} & B & C:1a,1b,2a,3a PW:1a,2a,3a \\
C2089 & 1a<o1L 1b<o2 / 2a<1a1,1b0L / 3a<2a2 & IL/I/I & +{-}/{-}/{-} & B & C:1a,1b,2a,3a PW:1a,2a,3a \\
C2090 & 1a<o1L 1b<o2 / 2a<1a1,1b0L / 3a<2a2L & IL/I/I & +{-}/{-}/{-} & B & C:1a,1b,2a,3a PW:1a,2a,3a \\
C2091 & 1a<o1L 1b<o2 / 2a<1a1,1b1 / 3a<2a0L & IL/E/N & +{-}/{-}/{-} & BF & C:1a,1b PW:1a \\
C2092 & 1a<o1L 1b<o2 / 2a<1a1,1b1 / 3a<2a1 & IL/E/D & +{-}/{-}/{-} & BF & C:1a,1b PW:1a \\
C2093 & 1a<o1L 1b<o2 / 2a<1a1,1b1 / 3a<2a1L & IL/E/D & +{-}/{-}/{-} & BF & C:1a,1b PW:1a \\
C2094 & 1a<o1L 1b<o2 / 2a<1a1,1b1 / 3a<2a2 & IL/E/D & +{-}/{-}/{-} & BF & C:1a,1b PW:1a \\
C2095 & 1a<o1L 1b<o2 / 2a<1a1,1b1 / 3a<2a2L & IL/E/D & +{-}/{-}/{-} & BF & C:1a,1b PW:1a \\
C2096 & 1a<o1L 1b<o2 / 2a<1a1,1b1L / 3a<2a0L & IL/E/N & +{-}/{-}/{-} & BF & C:1a,1b PW:1a \\
C2097 & 1a<o1L 1b<o2 / 2a<1a1,1b1L / 3a<2a1 & IL/E/D & +{-}/{-}/{-} & BF & C:1a,1b PW:1a \\
C2098 & 1a<o1L 1b<o2 / 2a<1a1,1b1L / 3a<2a1L & IL/E/D & +{-}/{-}/{-} & BF & C:1a,1b PW:1a \\
C2099 & 1a<o1L 1b<o2 / 2a<1a1,1b1L / 3a<2a2 & IL/E/D & +{-}/{-}/{-} & BF & C:1a,1b PW:1a \\
C2100 & 1a<o1L 1b<o2 / 2a<1a1,1b1L / 3a<2a2L & IL/E/D & +{-}/{-}/{-} & BF & C:1a,1b PW:1a \\
C2101 & 1a<o1L 1b<o2 / 2a<1a1,1b2 / 3a<2a0L & IL/E/N & +{-}/{-}/{-} & BF & C:1a,1b,2a PW:1a \\
C2102 & 1a<o1L 1b<o2 / 2a<1a1,1b2 / 3a<2a1 & IL/E/D & +{-}/{-}/{-} & BF & C:1a,1b,2a,3a PW:1a \\
C2103 & 1a<o1L 1b<o2 / 2a<1a1,1b2 / 3a<2a1L & IL/E/D & +{-}/{-}/{-} & BF & C:1a,1b,2a,3a PW:1a \\
C2104 & 1a<o1L 1b<o2 / 2a<1a1,1b2 / 3a<2a2 & IL/E/D & +{-}/{-}/{-} & BF & C:1a,1b,2a,3a PW:1a \\
C2105 & 1a<o1L 1b<o2 / 2a<1a1,1b2 / 3a<2a2L & IL/E/D & +{-}/{-}/{-} & BF & C:1a,1b,2a,3a PW:1a \\
C2106 & 1a<o1L 1b<o2 / 2a<1a1,1b2L / 3a<2a0L & IL/E/N & +{-}/{-}/{-} & BF & C:1a,1b,2a PW:1a \\
C2107 & 1a<o1L 1b<o2 / 2a<1a1,1b2L / 3a<2a1 & IL/E/D & +{-}/{-}/{-} & BF & C:1a,1b,2a,3a PW:1a \\
C2108 & 1a<o1L 1b<o2 / 2a<1a1,1b2L / 3a<2a1L & IL/E/D & +{-}/{-}/{-} & BF & C:1a,1b,2a,3a PW:1a \\
C2109 & 1a<o1L 1b<o2 / 2a<1a1,1b2L / 3a<2a2 & IL/E/D & +{-}/{-}/{-} & BF & C:1a,1b,2a,3a PW:1a \\
C2110 & 1a<o1L 1b<o2 / 2a<1a1,1b2L / 3a<2a2L & IL/E/D & +{-}/{-}/{-} & BF & C:1a,1b,2a,3a PW:1a \\
C2111 & 1a<o1L 1b<o2 / 2a<1a1L,1b1 / 3a<2a0L & IL/I/N & +{-}/+/+ & BF & C:1a,1b,2a PW:1a,2a \\
C2112 & 1a<o1L 1b<o2 / 2a<1a1L,1b1 / 3a<2a1 & IL/I/I & +{-}/+/{-} & BF & C:1a,1b,2a,3a PW:1a,2a,3a \\
C2113 & 1a<o1L 1b<o2 / 2a<1a1L,1b1 / 3a<2a1L & IL/I/I & +{-}/+/+ & BF & C:1a,1b,2a,3a PW:1a,2a,3a \\
C2114 & 1a<o1L 1b<o2 / 2a<1a1L,1b1 / 3a<2a2 & IL/I/I & +{-}/+/{-} & BF & C:1a,1b,2a,3a PW:1a,2a,3a \\
C2115 & 1a<o1L 1b<o2 / 2a<1a1L,1b1 / 3a<2a2L & IL/I/I & +{-}/+/+ & BF & C:1a,1b,2a,3a PW:1a,2a,3a \\
C2116 & 1a<o1L 1b<o2 / 2a<1a1L,1b2 / 3a<2a0L & IL/I/N & +{-}/+/+ & BF & C:1a,1b PW:1a,2a \\
C2117 & 1a<o1L 1b<o2 / 2a<1a1L,1b2 / 3a<2a1 & IL/I/I & +{-}/+/{-} & BF & C:1a,1b PW:1a,2a,3a \\
C2118 & 1a<o1L 1b<o2 / 2a<1a1L,1b2 / 3a<2a1L & IL/I/I & +{-}/+/+ & BF & C:1a,1b PW:1a,2a,3a \\
C2119 & 1a<o1L 1b<o2 / 2a<1a1L,1b2 / 3a<2a2 & IL/I/I & +{-}/+/{-} & BF & C:1a,1b PW:1a,2a,3a \\
C2120 & 1a<o1L 1b<o2 / 2a<1a1L,1b2 / 3a<2a2L & IL/I/I & +{-}/+/+ & BF & C:1a,1b PW:1a,2a,3a \\
C2121 & 1a<o1L 1b<o2 / 2a<1a2,1b0L / 3a<2a0L & IL/I/N & +{-}/{-}/{-} & B & C:1a,1b,2a PW:1a,2a \\
C2122 & 1a<o1L 1b<o2 / 2a<1a2,1b0L / 3a<2a1 & IL/I/I & +{-}/{-}/{-} & B & C:1a,1b,2a,3a PW:1a,2a,3a \\
C2123 & 1a<o1L 1b<o2 / 2a<1a2,1b0L / 3a<2a1L & IL/I/I & +{-}/{-}/{-} & B & C:1a,1b,2a,3a PW:1a,2a,3a \\
C2124 & 1a<o1L 1b<o2 / 2a<1a2,1b0L / 3a<2a2 & IL/I/I & +{-}/{-}/{-} & B & C:1a,1b,2a,3a PW:1a,2a,3a \\
C2125 & 1a<o1L 1b<o2 / 2a<1a2,1b0L / 3a<2a2L & IL/I/I & +{-}/{-}/{-} & B & C:1a,1b,2a,3a PW:1a,2a,3a \\
C2126 & 1a<o1L 1b<o2 / 2a<1a2,1b1 / 3a<2a0L & IL/E/N & +{-}/{-}/{-} & BF & C:1a,1b PW:1a \\
C2127 & 1a<o1L 1b<o2 / 2a<1a2,1b1 / 3a<2a1 & IL/E/D & +{-}/{-}/{-} & BF & C:1a,1b PW:1a \\
C2128 & 1a<o1L 1b<o2 / 2a<1a2,1b1 / 3a<2a1L & IL/E/D & +{-}/{-}/{-} & BF & C:1a,1b PW:1a \\
C2129 & 1a<o1L 1b<o2 / 2a<1a2,1b1 / 3a<2a2 & IL/E/D & +{-}/{-}/{-} & BF & C:1a,1b PW:1a \\
C2130 & 1a<o1L 1b<o2 / 2a<1a2,1b1 / 3a<2a2L & IL/E/D & +{-}/{-}/{-} & BF & C:1a,1b PW:1a \\
C2131 & 1a<o1L 1b<o2 / 2a<1a2,1b1L / 3a<2a0L & IL/E/N & +{-}/{-}/{-} & BF & C:1a,1b PW:1a \\
C2132 & 1a<o1L 1b<o2 / 2a<1a2,1b1L / 3a<2a1 & IL/E/D & +{-}/{-}/{-} & BF & C:1a,1b PW:1a \\
C2133 & 1a<o1L 1b<o2 / 2a<1a2,1b1L / 3a<2a1L & IL/E/D & +{-}/{-}/{-} & BF & C:1a,1b PW:1a \\
C2134 & 1a<o1L 1b<o2 / 2a<1a2,1b1L / 3a<2a2 & IL/E/D & +{-}/{-}/{-} & BF & C:1a,1b PW:1a \\
C2135 & 1a<o1L 1b<o2 / 2a<1a2,1b1L / 3a<2a2L & IL/E/D & +{-}/{-}/{-} & BF & C:1a,1b PW:1a \\
C2136 & 1a<o1L 1b<o2 / 2a<1a2,1b2 / 3a<2a0L & IL/E/N & +{-}/{-}/{-} & BF & C:1a,1b PW:1a \\
C2137 & 1a<o1L 1b<o2 / 2a<1a2,1b2 / 3a<2a1 & IL/E/D & +{-}/{-}/{-} & BF & C:1a,1b PW:1a \\
C2138 & 1a<o1L 1b<o2 / 2a<1a2,1b2 / 3a<2a1L & IL/E/D & +{-}/{-}/{-} & BF & C:1a,1b PW:1a \\
C2139 & 1a<o1L 1b<o2 / 2a<1a2,1b2 / 3a<2a2 & IL/E/D & +{-}/{-}/{-} & BF & C:1a,1b PW:1a \\
C2140 & 1a<o1L 1b<o2 / 2a<1a2,1b2 / 3a<2a2L & IL/E/D & +{-}/{-}/{-} & BF & C:1a,1b PW:1a \\
C2141 & 1a<o1L 1b<o2 / 2a<1a2,1b2L / 3a<2a0L & IL/E/N & +{-}/{-}/{-} & BF & C:1a,1b PW:1a \\
C2142 & 1a<o1L 1b<o2 / 2a<1a2,1b2L / 3a<2a1 & IL/E/D & +{-}/{-}/{-} & BF & C:1a,1b PW:1a \\
C2143 & 1a<o1L 1b<o2 / 2a<1a2,1b2L / 3a<2a1L & IL/E/D & +{-}/{-}/{-} & BF & C:1a,1b PW:1a \\
C2144 & 1a<o1L 1b<o2 / 2a<1a2,1b2L / 3a<2a2 & IL/E/D & +{-}/{-}/{-} & BF & C:1a,1b PW:1a \\
C2145 & 1a<o1L 1b<o2 / 2a<1a2,1b2L / 3a<2a2L & IL/E/D & +{-}/{-}/{-} & BF & C:1a,1b PW:1a \\
C2146 & 1a<o1L 1b<o2 / 2a<1a2L,1b1 / 3a<2a0L & IL/I/N & +{-}/+/+ & BF & C:1a,1b,2a PW:1a,2a \\
C2147 & 1a<o1L 1b<o2 / 2a<1a2L,1b1 / 3a<2a1 & IL/I/I & +{-}/+/{-} & BF & C:1a,1b,2a,3a PW:1a,2a,3a \\
C2148 & 1a<o1L 1b<o2 / 2a<1a2L,1b1 / 3a<2a1L & IL/I/I & +{-}/+/+ & BF & C:1a,1b,2a,3a PW:1a,2a,3a \\
C2149 & 1a<o1L 1b<o2 / 2a<1a2L,1b1 / 3a<2a2 & IL/I/I & +{-}/+/{-} & BF & C:1a,1b,2a,3a PW:1a,2a,3a \\
C2150 & 1a<o1L 1b<o2 / 2a<1a2L,1b1 / 3a<2a2L & IL/I/I & +{-}/+/+ & BF & C:1a,1b,2a,3a PW:1a,2a,3a \\
C2151 & 1a<o1L 1b<o2 / 2a<1a2L,1b2 / 3a<2a0L & IL/I/N & +{-}/+/+ & BF & C:1a,1b,2a PW:1a,2a \\
C2152 & 1a<o1L 1b<o2 / 2a<1a2L,1b2 / 3a<2a1 & IL/I/I & +{-}/+/{-} & BF & C:1a,1b,2a,3a PW:1a,2a,3a \\
C2153 & 1a<o1L 1b<o2 / 2a<1a2L,1b2 / 3a<2a1L & IL/I/I & +{-}/+/+ & BF & C:1a,1b,2a,3a PW:1a,2a,3a \\
C2154 & 1a<o1L 1b<o2 / 2a<1a2L,1b2 / 3a<2a2 & IL/I/I & +{-}/+/{-} & BF & C:1a,1b,2a,3a PW:1a,2a,3a \\
C2155 & 1a<o1L 1b<o2 / 2a<1a2L,1b2 / 3a<2a2L & IL/I/I & +{-}/+/+ & BF & C:1a,1b,2a,3a PW:1a,2a,3a \\
C2156 & 1a<o2 / 2a<1a0L 2b<1a1 / 3a<2a0L,2b1 & I/NI/I & {-}/{-}{-}/{-} & N &  \\
C2157 & 1a<o2 / 2a<1a0L 2b<1a1 / 3a<2a0L,2b2 & I/NI/I & {-}/{-}{-}/{-} & N &  \\
C2158 & 1a<o2 / 2a<1a0L 2b<1a1 / 3a<2a1,2b0L & I/NI/N & {-}/{-}{-}/{-} & N &  \\
C2159 & 1a<o2 / 2a<1a0L 2b<1a1 / 3a<2a1,2b1 & I/NI/E & {-}/{-}{-}/{-} & F & W:3a \\
C2160 & 1a<o2 / 2a<1a0L 2b<1a1 / 3a<2a1,2b1L & I/NI/I & {-}/{-}{-}/{-} & F & CP:3a \\
C2161 & 1a<o2 / 2a<1a0L 2b<1a1 / 3a<2a1,2b2 & I/NI/E & {-}/{-}{-}/{-} & F & CW:3a \\
C2162 & 1a<o2 / 2a<1a0L 2b<1a1 / 3a<2a1,2b2L & I/NI/I & {-}/{-}{-}/{-} & F & P:3a \\
C2163 & 1a<o2 / 2a<1a0L 2b<1a1 / 3a<2a1L,2b1 & I/NI/E & {-}/{-}{-}/{-} & F & W:3a \\
C2164 & 1a<o2 / 2a<1a0L 2b<1a1 / 3a<2a1L,2b2 & I/NI/E & {-}/{-}{-}/{-} & F & CW:3a \\
C2165 & 1a<o2 / 2a<1a0L 2b<1a1 / 3a<2a2,2b0L & I/NI/N & {-}/{-}{-}/{-} & N &  \\
C2166 & 1a<o2 / 2a<1a0L 2b<1a1 / 3a<2a2,2b1 & I/NI/E & {-}/{-}{-}/{-} & F & W:3a \\
C2167 & 1a<o2 / 2a<1a0L 2b<1a1 / 3a<2a2,2b1L & I/NI/I & {-}/{-}{-}/{-} & F & CP:3a \\
C2168 & 1a<o2 / 2a<1a0L 2b<1a1 / 3a<2a2,2b2 & I/NI/E & {-}/{-}{-}/{-} & F & W:3a \\
C2169 & 1a<o2 / 2a<1a0L 2b<1a1 / 3a<2a2,2b2L & I/NI/I & {-}/{-}{-}/{-} & F & CP:3a \\
C2170 & 1a<o2 / 2a<1a0L 2b<1a1 / 3a<2a2L,2b1 & I/NI/E & {-}/{-}{-}/{-} & F & W:3a \\
C2171 & 1a<o2 / 2a<1a0L 2b<1a1 / 3a<2a2L,2b2 & I/NI/E & {-}/{-}{-}/{-} & F & W:3a \\
C2172 & 1a<o2 / 2a<1a0L 2b<1a2 / 3a<2a0L,2b1 & I/NI/I & {-}/{-}{-}/{-} & N &  \\
C2173 & 1a<o2 / 2a<1a0L 2b<1a2 / 3a<2a0L,2b2 & I/NI/I & {-}/{-}{-}/{-} & N &  \\
C2174 & 1a<o2 / 2a<1a0L 2b<1a2 / 3a<2a1,2b0L & I/NI/N & {-}/{-}{-}/{-} & N &  \\
C2175 & 1a<o2 / 2a<1a0L 2b<1a2 / 3a<2a1,2b1 & I/NI/E & {-}/{-}{-}/{-} & F & W:3a \\
C2176 & 1a<o2 / 2a<1a0L 2b<1a2 / 3a<2a1,2b1L & I/NI/I & {-}/{-}{-}/{-} & F & CP:3a \\
C2177 & 1a<o2 / 2a<1a0L 2b<1a2 / 3a<2a1,2b2 & I/NI/E & {-}/{-}{-}/{-} & F & CW:3a \\
C2178 & 1a<o2 / 2a<1a0L 2b<1a2 / 3a<2a1,2b2L & I/NI/I & {-}/{-}{-}/{-} & F & P:3a \\
C2179 & 1a<o2 / 2a<1a0L 2b<1a2 / 3a<2a1L,2b1 & I/NI/E & {-}/{-}{-}/{-} & F & W:3a \\
C2180 & 1a<o2 / 2a<1a0L 2b<1a2 / 3a<2a1L,2b2 & I/NI/E & {-}/{-}{-}/{-} & F & CW:3a \\
C2181 & 1a<o2 / 2a<1a0L 2b<1a2 / 3a<2a2,2b0L & I/NI/N & {-}/{-}{-}/{-} & N &  \\
C2182 & 1a<o2 / 2a<1a0L 2b<1a2 / 3a<2a2,2b1 & I/NI/E & {-}/{-}{-}/{-} & F & W:3a \\
C2183 & 1a<o2 / 2a<1a0L 2b<1a2 / 3a<2a2,2b1L & I/NI/I & {-}/{-}{-}/{-} & F & CP:3a \\
C2184 & 1a<o2 / 2a<1a0L 2b<1a2 / 3a<2a2,2b2 & I/NI/E & {-}/{-}{-}/{-} & F & W:3a \\
C2185 & 1a<o2 / 2a<1a0L 2b<1a2 / 3a<2a2,2b2L & I/NI/I & {-}/{-}{-}/{-} & F & CP:3a \\
C2186 & 1a<o2 / 2a<1a0L 2b<1a2 / 3a<2a2L,2b1 & I/NI/E & {-}/{-}{-}/{-} & F & W:3a \\
C2187 & 1a<o2 / 2a<1a0L 2b<1a2 / 3a<2a2L,2b2 & I/NI/E & {-}/{-}{-}/{-} & F & W:3a \\
C2188 & 1a<o2 / 2a<1a1 2b<1a1 / 3a<2a0L,2b1 & I/EE/D & {-}/{-}{-}/{-} & E & W:1a \\
C2189 & 1a<o2 / 2a<1a1 2b<1a1 / 3a<2a0L,2b2 & I/EE/D & {-}/{-}{-}/{-} & E & W:1a \\
C2190 & 1a<o2 / 2a<1a1 2b<1a1 / 3a<2a1,2b1 & I/EE/D & {-}/{-}{-}/{-} & E & W:1a \\
C2191 & 1a<o2 / 2a<1a1 2b<1a1 / 3a<2a1,2b1L & I/EE/D & {-}/{-}{-}/{-} & E & W:1a \\
C2192 & 1a<o2 / 2a<1a1 2b<1a1 / 3a<2a1,2b2 & I/EE/D & {-}/{-}{-}/{-} & E & W:1a \\
C2193 & 1a<o2 / 2a<1a1 2b<1a1 / 3a<2a1,2b2L & I/EE/D & {-}/{-}{-}/{-} & E & W:1a \\
C2194 & 1a<o2 / 2a<1a1 2b<1a1 / 3a<2a1L,2b2 & I/EE/D & {-}/{-}{-}/{-} & E & W:1a \\
C2195 & 1a<o2 / 2a<1a1 2b<1a1 / 3a<2a2,2b2 & I/EE/D & {-}/{-}{-}/{-} & E & W:1a \\
C2196 & 1a<o2 / 2a<1a1 2b<1a1 / 3a<2a2,2b2L & I/EE/D & {-}/{-}{-}/{-} & E & W:1a \\
C2197 & 1a<o2 / 2a<1a1 2b<1a1L / 3a<2a0L,2b1 & I/LI/I & {-}/{-}{-}/{-} & B & CP:2b,3a W:1a,2b,3a \\
C2198 & 1a<o2 / 2a<1a1 2b<1a1L / 3a<2a0L,2b2 & I/LI/I & {-}/{-}{-}/{-} & B & CP:2b,3a W:1a,2b,3a \\
C2199 & 1a<o2 / 2a<1a1 2b<1a1L / 3a<2a1,2b0L & I/LI/D & {-}/{-}{-}/{-} & B & CP:2b W:1a,2b \\
C2200 & 1a<o2 / 2a<1a1 2b<1a1L / 3a<2a1,2b1 & I/LI/E & {-}/{-}{-}/{-} & BF & CP:2b W:1a,2b \\
C2201 & 1a<o2 / 2a<1a1 2b<1a1L / 3a<2a1,2b1L & I/LI/I & {-}/{-}{-}/{-} & BF & CP:2b,3a W:1a,2b,3a \\
C2202 & 1a<o2 / 2a<1a1 2b<1a1L / 3a<2a1,2b2 & I/LI/E & {-}/{-}{-}/{-} & BF & CP:2b W:1a,2b \\
C2203 & 1a<o2 / 2a<1a1 2b<1a1L / 3a<2a1,2b2L & I/LI/I & {-}/{-}{-}/{-} & BF & CP:2b,3a W:1a,2b,3a \\
C2204 & 1a<o2 / 2a<1a1 2b<1a1L / 3a<2a1L,2b1 & I/LI/E & {-}/{-}{-}/{-} & BF & CP:2b W:1a,2b \\
C2205 & 1a<o2 / 2a<1a1 2b<1a1L / 3a<2a1L,2b2 & I/LI/E & {-}/{-}{-}/{-} & BF & CP:2b W:1a,2b \\
C2206 & 1a<o2 / 2a<1a1 2b<1a1L / 3a<2a2,2b0L & I/LI/D & {-}/{-}{-}/{-} & B & CP:2b W:1a,2b \\
C2207 & 1a<o2 / 2a<1a1 2b<1a1L / 3a<2a2,2b1 & I/LI/E & {-}/{-}{-}/{-} & BF & CP:2b W:1a,2b \\
C2208 & 1a<o2 / 2a<1a1 2b<1a1L / 3a<2a2,2b1L & I/LI/I & {-}/{-}{-}/{-} & BF & CP:2b,3a W:1a,2b,3a \\
C2209 & 1a<o2 / 2a<1a1 2b<1a1L / 3a<2a2,2b2 & I/LI/E & {-}/{-}{-}/{-} & BF & CP:2b W:1a,2b \\
C2210 & 1a<o2 / 2a<1a1 2b<1a1L / 3a<2a2,2b2L & I/LI/I & {-}/{-}{-}/{-} & BF & CP:2b,3a W:1a,2b,3a \\
C2211 & 1a<o2 / 2a<1a1 2b<1a1L / 3a<2a2L,2b1 & I/LI/E & {-}/{-}{-}/{-} & BF & CP:2b W:1a,2b \\
C2212 & 1a<o2 / 2a<1a1 2b<1a1L / 3a<2a2L,2b2 & I/LI/E & {-}/{-}{-}/{-} & BF & CP:2b W:1a,2b \\
C2213 & 1a<o2 / 2a<1a1 2b<1a2 / 3a<2a0L,2b1 & I/EE/D & {-}/{-}{-}/{-} & E & C:2b,3a W:1a \\
C2214 & 1a<o2 / 2a<1a1 2b<1a2 / 3a<2a0L,2b2 & I/EE/D & {-}/{-}{-}/{-} & E & C:2b,3a W:1a \\
C2215 & 1a<o2 / 2a<1a1 2b<1a2 / 3a<2a1,2b0L & I/EE/D & {-}/{-}{-}/{-} & E & C:2b W:1a \\
C2216 & 1a<o2 / 2a<1a1 2b<1a2 / 3a<2a1,2b1 & I/EE/D & {-}/{-}{-}/{-} & E & C:2b W:1a \\
C2217 & 1a<o2 / 2a<1a1 2b<1a2 / 3a<2a1,2b1L & I/EE/D & {-}/{-}{-}/{-} & E & C:2b W:1a \\
C2218 & 1a<o2 / 2a<1a1 2b<1a2 / 3a<2a1,2b2 & I/EE/D & {-}/{-}{-}/{-} & E & C:2b,3a W:1a \\
C2219 & 1a<o2 / 2a<1a1 2b<1a2 / 3a<2a1,2b2L & I/EE/D & {-}/{-}{-}/{-} & E & C:2b,3a W:1a \\
C2220 & 1a<o2 / 2a<1a1 2b<1a2 / 3a<2a1L,2b1 & I/EE/D & {-}/{-}{-}/{-} & E & C:2b W:1a \\
C2221 & 1a<o2 / 2a<1a1 2b<1a2 / 3a<2a1L,2b2 & I/EE/D & {-}/{-}{-}/{-} & E & C:2b,3a W:1a \\
C2222 & 1a<o2 / 2a<1a1 2b<1a2 / 3a<2a2,2b0L & I/EE/D & {-}/{-}{-}/{-} & E & C:2b W:1a \\
C2223 & 1a<o2 / 2a<1a1 2b<1a2 / 3a<2a2,2b1 & I/EE/D & {-}/{-}{-}/{-} & E & C:2b W:1a \\
C2224 & 1a<o2 / 2a<1a1 2b<1a2 / 3a<2a2,2b1L & I/EE/D & {-}/{-}{-}/{-} & E & C:2b W:1a \\
C2225 & 1a<o2 / 2a<1a1 2b<1a2 / 3a<2a2,2b2 & I/EE/D & {-}/{-}{-}/{-} & E & C:2b W:1a \\
C2226 & 1a<o2 / 2a<1a1 2b<1a2 / 3a<2a2,2b2L & I/EE/D & {-}/{-}{-}/{-} & E & C:2b W:1a \\
C2227 & 1a<o2 / 2a<1a1 2b<1a2 / 3a<2a2L,2b1 & I/EE/D & {-}/{-}{-}/{-} & E & C:2b W:1a \\
C2228 & 1a<o2 / 2a<1a1 2b<1a2 / 3a<2a2L,2b2 & I/EE/D & {-}/{-}{-}/{-} & E & C:2b W:1a \\
C2229 & 1a<o2 / 2a<1a1 2b<1a2L / 3a<2a0L,2b1 & I/LI/I & {-}/{-}{-}/{-} & B & P:2b,3a W:1a,2b,3a \\
C2230 & 1a<o2 / 2a<1a1 2b<1a2L / 3a<2a0L,2b2 & I/LI/I & {-}/{-}{-}/{-} & B & P:2b,3a W:1a,2b,3a \\
C2231 & 1a<o2 / 2a<1a1 2b<1a2L / 3a<2a1,2b0L & I/LI/D & {-}/{-}{-}/{-} & B & P:2b W:1a,2b \\
C2232 & 1a<o2 / 2a<1a1 2b<1a2L / 3a<2a1,2b1 & I/LI/E & {-}/{-}{-}/{-} & BF & P:2b W:1a,2b \\
C2233 & 1a<o2 / 2a<1a1 2b<1a2L / 3a<2a1,2b1L & I/LI/I & {-}/{-}{-}/{-} & BF & C:3a P:2b,3a W:1a,2b,3a \\
C2234 & 1a<o2 / 2a<1a1 2b<1a2L / 3a<2a1,2b2 & I/LI/E & {-}/{-}{-}/{-} & BF & C:3a P:2b W:1a,2b \\
C2235 & 1a<o2 / 2a<1a1 2b<1a2L / 3a<2a1,2b2L & I/LI/I & {-}/{-}{-}/{-} & BF & P:2b,3a W:1a,2b,3a \\
C2236 & 1a<o2 / 2a<1a1 2b<1a2L / 3a<2a1L,2b1 & I/LI/E & {-}/{-}{-}/{-} & BF & P:2b W:1a,2b \\
C2237 & 1a<o2 / 2a<1a1 2b<1a2L / 3a<2a1L,2b2 & I/LI/E & {-}/{-}{-}/{-} & BF & C:3a P:2b W:1a,2b \\
C2238 & 1a<o2 / 2a<1a1 2b<1a2L / 3a<2a2,2b0L & I/LI/D & {-}/{-}{-}/{-} & B & P:2b W:1a,2b \\
C2239 & 1a<o2 / 2a<1a1 2b<1a2L / 3a<2a2,2b1 & I/LI/E & {-}/{-}{-}/{-} & BF & P:2b W:1a,2b \\
C2240 & 1a<o2 / 2a<1a1 2b<1a2L / 3a<2a2,2b1L & I/LI/I & {-}/{-}{-}/{-} & BF & C:3a P:2b,3a W:1a,2b,3a \\
C2241 & 1a<o2 / 2a<1a1 2b<1a2L / 3a<2a2,2b2 & I/LI/E & {-}/{-}{-}/{-} & BF & P:2b W:1a,2b \\
C2242 & 1a<o2 / 2a<1a1 2b<1a2L / 3a<2a2,2b2L & I/LI/I & {-}/{-}{-}/{-} & BF & C:3a P:2b,3a W:1a,2b,3a \\
C2243 & 1a<o2 / 2a<1a1 2b<1a2L / 3a<2a2L,2b1 & I/LI/E & {-}/{-}{-}/{-} & BF & P:2b W:1a,2b \\
C2244 & 1a<o2 / 2a<1a1 2b<1a2L / 3a<2a2L,2b2 & I/LI/E & {-}/{-}{-}/{-} & BF & P:2b W:1a,2b \\
C2245 & 1a<o2 / 2a<1a1L 2b<1a2 / 3a<2a0L,2b1 & I/IL/D & {-}/{-}{-}/{-} & B & C:2a,2b,3a P:2a W:1a,2a \\
C2246 & 1a<o2 / 2a<1a1L 2b<1a2 / 3a<2a0L,2b2 & I/IL/D & {-}/{-}{-}/{-} & B & C:2a,2b,3a P:2a W:1a,2a \\
C2247 & 1a<o2 / 2a<1a1L 2b<1a2 / 3a<2a1,2b0L & I/IL/I & {-}/{-}{-}/{-} & B & C:2a,2b,3a P:2a,3a W:1a,2a,3a \\
C2248 & 1a<o2 / 2a<1a1L 2b<1a2 / 3a<2a1,2b1 & I/IL/E & {-}/{-}{-}/{-} & BF & C:2a,2b P:2a W:1a,2a \\
C2249 & 1a<o2 / 2a<1a1L 2b<1a2 / 3a<2a1,2b1L & I/IL/E & {-}/{-}{-}/{-} & BF & C:2a,2b P:2a W:1a,2a \\
C2250 & 1a<o2 / 2a<1a1L 2b<1a2 / 3a<2a1,2b2 & I/IL/E & {-}/{-}{-}/{-} & BF & C:2a,2b,3a P:2a W:1a,2a \\
C2251 & 1a<o2 / 2a<1a1L 2b<1a2 / 3a<2a1,2b2L & I/IL/E & {-}/{-}{-}/{-} & BF & C:2a,2b,3a P:2a W:1a,2a \\
C2252 & 1a<o2 / 2a<1a1L 2b<1a2 / 3a<2a1L,2b1 & I/IL/I & {-}/{-}{-}/{-} & BF & C:2a,2b,3a P:2a,3a W:1a,2a,3a \\
C2253 & 1a<o2 / 2a<1a1L 2b<1a2 / 3a<2a1L,2b2 & I/IL/I & {-}/{-}{-}/{-} & BF & C:2a,2b P:2a,3a W:1a,2a,3a \\
C2254 & 1a<o2 / 2a<1a1L 2b<1a2 / 3a<2a2,2b0L & I/IL/I & {-}/{-}{-}/{-} & B & C:2a,2b,3a P:2a,3a W:1a,2a,3a \\
C2255 & 1a<o2 / 2a<1a1L 2b<1a2 / 3a<2a2,2b1 & I/IL/E & {-}/{-}{-}/{-} & BF & C:2a,2b P:2a W:1a,2a \\
C2256 & 1a<o2 / 2a<1a1L 2b<1a2 / 3a<2a2,2b1L & I/IL/E & {-}/{-}{-}/{-} & BF & C:2a,2b P:2a W:1a,2a \\
C2257 & 1a<o2 / 2a<1a1L 2b<1a2 / 3a<2a2,2b2 & I/IL/E & {-}/{-}{-}/{-} & BF & C:2a,2b P:2a W:1a,2a \\
C2258 & 1a<o2 / 2a<1a1L 2b<1a2 / 3a<2a2,2b2L & I/IL/E & {-}/{-}{-}/{-} & BF & C:2a,2b P:2a W:1a,2a \\
C2259 & 1a<o2 / 2a<1a1L 2b<1a2 / 3a<2a2L,2b1 & I/IL/I & {-}/{-}{-}/{-} & BF & C:2a,2b,3a P:2a,3a W:1a,2a,3a \\
C2260 & 1a<o2 / 2a<1a1L 2b<1a2 / 3a<2a2L,2b2 & I/IL/I & {-}/{-}{-}/{-} & BF & C:2a,2b,3a P:2a,3a W:1a,2a,3a \\
C2261 & 1a<o2 / 2a<1a2 2b<1a2 / 3a<2a0L,2b1 & I/EE/D & {-}/{-}{-}/{-} & E & W:1a \\
C2262 & 1a<o2 / 2a<1a2 2b<1a2 / 3a<2a0L,2b2 & I/EE/D & {-}/{-}{-}/{-} & E & W:1a \\
C2263 & 1a<o2 / 2a<1a2 2b<1a2 / 3a<2a1,2b1 & I/EE/D & {-}/{-}{-}/{-} & E & W:1a \\
C2264 & 1a<o2 / 2a<1a2 2b<1a2 / 3a<2a1,2b1L & I/EE/D & {-}/{-}{-}/{-} & E & W:1a \\
C2265 & 1a<o2 / 2a<1a2 2b<1a2 / 3a<2a1,2b2 & I/EE/D & {-}/{-}{-}/{-} & E & W:1a \\
C2266 & 1a<o2 / 2a<1a2 2b<1a2 / 3a<2a1,2b2L & I/EE/D & {-}/{-}{-}/{-} & E & W:1a \\
C2267 & 1a<o2 / 2a<1a2 2b<1a2 / 3a<2a1L,2b2 & I/EE/D & {-}/{-}{-}/{-} & E & W:1a \\
C2268 & 1a<o2 / 2a<1a2 2b<1a2 / 3a<2a2,2b2 & I/EE/D & {-}/{-}{-}/{-} & E & W:1a \\
C2269 & 1a<o2 / 2a<1a2 2b<1a2 / 3a<2a2,2b2L & I/EE/D & {-}/{-}{-}/{-} & E & W:1a \\
C2270 & 1a<o2 / 2a<1a2 2b<1a2L / 3a<2a0L,2b1 & I/LI/I & {-}/{-}{-}/{-} & B & CP:2b,3a W:1a,2b,3a \\
C2271 & 1a<o2 / 2a<1a2 2b<1a2L / 3a<2a0L,2b2 & I/LI/I & {-}/{-}{-}/{-} & B & CP:2b,3a W:1a,2b,3a \\
C2272 & 1a<o2 / 2a<1a2 2b<1a2L / 3a<2a1,2b0L & I/LI/D & {-}/{-}{-}/{-} & B & CP:2b W:1a,2b \\
C2273 & 1a<o2 / 2a<1a2 2b<1a2L / 3a<2a1,2b1 & I/LI/E & {-}/{-}{-}/{-} & BF & CP:2b W:1a,2b \\
C2274 & 1a<o2 / 2a<1a2 2b<1a2L / 3a<2a1,2b1L & I/LI/I & {-}/{-}{-}/{-} & BF & CP:2b,3a W:1a,2b,3a \\
C2275 & 1a<o2 / 2a<1a2 2b<1a2L / 3a<2a1,2b2 & I/LI/E & {-}/{-}{-}/{-} & BF & CP:2b W:1a,2b \\
C2276 & 1a<o2 / 2a<1a2 2b<1a2L / 3a<2a1,2b2L & I/LI/I & {-}/{-}{-}/{-} & BF & CP:2b,3a W:1a,2b,3a \\
C2277 & 1a<o2 / 2a<1a2 2b<1a2L / 3a<2a1L,2b1 & I/LI/E & {-}/{-}{-}/{-} & BF & CP:2b W:1a,2b \\
C2278 & 1a<o2 / 2a<1a2 2b<1a2L / 3a<2a1L,2b2 & I/LI/E & {-}/{-}{-}/{-} & BF & CP:2b W:1a,2b \\
C2279 & 1a<o2 / 2a<1a2 2b<1a2L / 3a<2a2,2b0L & I/LI/D & {-}/{-}{-}/{-} & B & CP:2b W:1a,2b \\
C2280 & 1a<o2 / 2a<1a2 2b<1a2L / 3a<2a2,2b1 & I/LI/E & {-}/{-}{-}/{-} & BF & CP:2b W:1a,2b \\
C2281 & 1a<o2 / 2a<1a2 2b<1a2L / 3a<2a2,2b1L & I/LI/I & {-}/{-}{-}/{-} & BF & CP:2b,3a W:1a,2b,3a \\
C2282 & 1a<o2 / 2a<1a2 2b<1a2L / 3a<2a2,2b2 & I/LI/E & {-}/{-}{-}/{-} & BF & CP:2b W:1a,2b \\
C2283 & 1a<o2 / 2a<1a2 2b<1a2L / 3a<2a2,2b2L & I/LI/I & {-}/{-}{-}/{-} & BF & CP:2b,3a W:1a,2b,3a \\
C2284 & 1a<o2 / 2a<1a2 2b<1a2L / 3a<2a2L,2b1 & I/LI/E & {-}/{-}{-}/{-} & BF & CP:2b W:1a,2b \\
C2285 & 1a<o2 / 2a<1a2 2b<1a2L / 3a<2a2L,2b2 & I/LI/E & {-}/{-}{-}/{-} & BF & CP:2b W:1a,2b \\
C2286 & 1a<o2 1b<o2 / 2a<1a0L,1b1 / 3a<2a0L & EE/D/N & {-}{-}/{-}/{-} & E &  \\
C2287 & 1a<o2 1b<o2 / 2a<1a0L,1b1 / 3a<2a1 & EE/D/D & {-}{-}/{-}/{-} & E &  \\
C2288 & 1a<o2 1b<o2 / 2a<1a0L,1b1 / 3a<2a1L & EE/D/D & {-}{-}/{-}/{-} & E &  \\
C2289 & 1a<o2 1b<o2 / 2a<1a0L,1b1 / 3a<2a2 & EE/D/D & {-}{-}/{-}/{-} & E &  \\
C2290 & 1a<o2 1b<o2 / 2a<1a0L,1b1 / 3a<2a2L & EE/D/D & {-}{-}/{-}/{-} & E &  \\
C2291 & 1a<o2 1b<o2 / 2a<1a0L,1b2 / 3a<2a0L & EE/D/N & {-}{-}/{-}/{-} & E &  \\
C2292 & 1a<o2 1b<o2 / 2a<1a0L,1b2 / 3a<2a1 & EE/D/D & {-}{-}/{-}/{-} & E &  \\
C2293 & 1a<o2 1b<o2 / 2a<1a0L,1b2 / 3a<2a1L & EE/D/D & {-}{-}/{-}/{-} & E &  \\
C2294 & 1a<o2 1b<o2 / 2a<1a0L,1b2 / 3a<2a2 & EE/D/D & {-}{-}/{-}/{-} & E &  \\
C2295 & 1a<o2 1b<o2 / 2a<1a0L,1b2 / 3a<2a2L & EE/D/D & {-}{-}/{-}/{-} & E &  \\
C2296 & 1a<o2 1b<o2 / 2a<1a1,1b1 / 3a<2a0L & EE/D/N & {-}{-}/{-}/{-} & E &  \\
C2297 & 1a<o2 1b<o2 / 2a<1a1,1b1 / 3a<2a1 & EE/D/D & {-}{-}/{-}/{-} & E &  \\
C2298 & 1a<o2 1b<o2 / 2a<1a1,1b1 / 3a<2a1L & EE/D/D & {-}{-}/{-}/{-} & E &  \\
C2299 & 1a<o2 1b<o2 / 2a<1a1,1b1 / 3a<2a2 & EE/D/D & {-}{-}/{-}/{-} & E &  \\
C2300 & 1a<o2 1b<o2 / 2a<1a1,1b1 / 3a<2a2L & EE/D/D & {-}{-}/{-}/{-} & E &  \\
C2301 & 1a<o2 1b<o2 / 2a<1a1,1b1L / 3a<2a0L & EE/D/N & {-}{-}/{-}/{-} & E &  \\
C2302 & 1a<o2 1b<o2 / 2a<1a1,1b1L / 3a<2a1 & EE/D/D & {-}{-}/{-}/{-} & E &  \\
C2303 & 1a<o2 1b<o2 / 2a<1a1,1b1L / 3a<2a1L & EE/D/D & {-}{-}/{-}/{-} & E &  \\
C2304 & 1a<o2 1b<o2 / 2a<1a1,1b1L / 3a<2a2 & EE/D/D & {-}{-}/{-}/{-} & E &  \\
C2305 & 1a<o2 1b<o2 / 2a<1a1,1b1L / 3a<2a2L & EE/D/D & {-}{-}/{-}/{-} & E &  \\
C2306 & 1a<o2 1b<o2 / 2a<1a1,1b2 / 3a<2a0L & EE/D/N & {-}{-}/{-}/{-} & E &  \\
C2307 & 1a<o2 1b<o2 / 2a<1a1,1b2 / 3a<2a1 & EE/D/D & {-}{-}/{-}/{-} & E &  \\
C2308 & 1a<o2 1b<o2 / 2a<1a1,1b2 / 3a<2a1L & EE/D/D & {-}{-}/{-}/{-} & E &  \\
C2309 & 1a<o2 1b<o2 / 2a<1a1,1b2 / 3a<2a2 & EE/D/D & {-}{-}/{-}/{-} & E &  \\
C2310 & 1a<o2 1b<o2 / 2a<1a1,1b2 / 3a<2a2L & EE/D/D & {-}{-}/{-}/{-} & E &  \\
C2311 & 1a<o2 1b<o2 / 2a<1a1,1b2L / 3a<2a0L & EE/D/N & {-}{-}/{-}/{-} & E &  \\
C2312 & 1a<o2 1b<o2 / 2a<1a1,1b2L / 3a<2a1 & EE/D/D & {-}{-}/{-}/{-} & E &  \\
C2313 & 1a<o2 1b<o2 / 2a<1a1,1b2L / 3a<2a1L & EE/D/D & {-}{-}/{-}/{-} & E &  \\
C2314 & 1a<o2 1b<o2 / 2a<1a1,1b2L / 3a<2a2 & EE/D/D & {-}{-}/{-}/{-} & E &  \\
C2315 & 1a<o2 1b<o2 / 2a<1a1,1b2L / 3a<2a2L & EE/D/D & {-}{-}/{-}/{-} & E &  \\
C2316 & 1a<o2 1b<o2 / 2a<1a1L,1b2 / 3a<2a0L & EE/D/N & {-}{-}/{-}/{-} & E &  \\
C2317 & 1a<o2 1b<o2 / 2a<1a1L,1b2 / 3a<2a1 & EE/D/D & {-}{-}/{-}/{-} & E &  \\
C2318 & 1a<o2 1b<o2 / 2a<1a1L,1b2 / 3a<2a1L & EE/D/D & {-}{-}/{-}/{-} & E &  \\
C2319 & 1a<o2 1b<o2 / 2a<1a1L,1b2 / 3a<2a2 & EE/D/D & {-}{-}/{-}/{-} & E &  \\
C2320 & 1a<o2 1b<o2 / 2a<1a1L,1b2 / 3a<2a2L & EE/D/D & {-}{-}/{-}/{-} & E &  \\
C2321 & 1a<o2 1b<o2 / 2a<1a2,1b2 / 3a<2a0L & EE/D/N & {-}{-}/{-}/{-} & E &  \\
C2322 & 1a<o2 1b<o2 / 2a<1a2,1b2 / 3a<2a1 & EE/D/D & {-}{-}/{-}/{-} & E &  \\
C2323 & 1a<o2 1b<o2 / 2a<1a2,1b2 / 3a<2a1L & EE/D/D & {-}{-}/{-}/{-} & E &  \\
C2324 & 1a<o2 1b<o2 / 2a<1a2,1b2 / 3a<2a2 & EE/D/D & {-}{-}/{-}/{-} & E &  \\
C2325 & 1a<o2 1b<o2 / 2a<1a2,1b2 / 3a<2a2L & EE/D/D & {-}{-}/{-}/{-} & E &  \\
C2326 & 1a<o2 1b<o2 / 2a<1a2,1b2L / 3a<2a0L & EE/D/N & {-}{-}/{-}/{-} & E &  \\
C2327 & 1a<o2 1b<o2 / 2a<1a2,1b2L / 3a<2a1 & EE/D/D & {-}{-}/{-}/{-} & E &  \\
C2328 & 1a<o2 1b<o2 / 2a<1a2,1b2L / 3a<2a1L & EE/D/D & {-}{-}/{-}/{-} & E &  \\
C2329 & 1a<o2 1b<o2 / 2a<1a2,1b2L / 3a<2a2 & EE/D/D & {-}{-}/{-}/{-} & E &  \\
C2330 & 1a<o2 1b<o2 / 2a<1a2,1b2L / 3a<2a2L & EE/D/D & {-}{-}/{-}/{-} & E &  \\
C2331 & 1a<o2 1b<o2L / 2a<1a0L,1b1 / 3a<2a0L & LI/I/N & {-}+/{-}/{-} & B & CPW:1b,2a \\
C2332 & 1a<o2 1b<o2L / 2a<1a0L,1b1 / 3a<2a1 & LI/I/I & {-}+/{-}/{-} & B & CPW:1b,2a,3a \\
C2333 & 1a<o2 1b<o2L / 2a<1a0L,1b1 / 3a<2a1L & LI/I/I & {-}+/{-}/{-} & B & CPW:1b,2a,3a \\
C2334 & 1a<o2 1b<o2L / 2a<1a0L,1b1 / 3a<2a2 & LI/I/I & {-}+/{-}/{-} & B & CPW:1b,2a,3a \\
C2335 & 1a<o2 1b<o2L / 2a<1a0L,1b1 / 3a<2a2L & LI/I/I & {-}+/{-}/{-} & B & CPW:1b,2a,3a \\
C2336 & 1a<o2 1b<o2L / 2a<1a0L,1b2 / 3a<2a0L & LI/I/N & {-}+/{-}/{-} & B & CPW:1b,2a \\
C2337 & 1a<o2 1b<o2L / 2a<1a0L,1b2 / 3a<2a1 & LI/I/I & {-}+/{-}/{-} & B & CPW:1b,2a,3a \\
C2338 & 1a<o2 1b<o2L / 2a<1a0L,1b2 / 3a<2a1L & LI/I/I & {-}+/{-}/{-} & B & CPW:1b,2a,3a \\
C2339 & 1a<o2 1b<o2L / 2a<1a0L,1b2 / 3a<2a2 & LI/I/I & {-}+/{-}/{-} & B & CPW:1b,2a,3a \\
C2340 & 1a<o2 1b<o2L / 2a<1a0L,1b2 / 3a<2a2L & LI/I/I & {-}+/{-}/{-} & B & CPW:1b,2a,3a \\
C2341 & 1a<o2 1b<o2L / 2a<1a1,1b0L / 3a<2a0L & LI/D/N & {-}+/+/+ & B & CPW:1b \\
C2342 & 1a<o2 1b<o2L / 2a<1a1,1b0L / 3a<2a1 & LI/D/D & {-}+/+/{-} & B & CPW:1b \\
C2343 & 1a<o2 1b<o2L / 2a<1a1,1b0L / 3a<2a1L & LI/D/D & {-}+/+/+ & B & CPW:1b \\
C2344 & 1a<o2 1b<o2L / 2a<1a1,1b0L / 3a<2a2 & LI/D/D & {-}+/+/{-} & B & CPW:1b \\
C2345 & 1a<o2 1b<o2L / 2a<1a1,1b0L / 3a<2a2L & LI/D/D & {-}+/+/+ & B & CPW:1b \\
C2346 & 1a<o2 1b<o2L / 2a<1a1,1b1 / 3a<2a0L & LI/E/N & {-}+/{-}/{-} & BF & CPW:1b \\
C2347 & 1a<o2 1b<o2L / 2a<1a1,1b1 / 3a<2a1 & LI/E/D & {-}+/{-}/{-} & BF & CPW:1b \\
C2348 & 1a<o2 1b<o2L / 2a<1a1,1b1 / 3a<2a1L & LI/E/D & {-}+/{-}/{-} & BF & CPW:1b \\
C2349 & 1a<o2 1b<o2L / 2a<1a1,1b1 / 3a<2a2 & LI/E/D & {-}+/{-}/{-} & BF & CPW:1b \\
C2350 & 1a<o2 1b<o2L / 2a<1a1,1b1 / 3a<2a2L & LI/E/D & {-}+/{-}/{-} & BF & CPW:1b \\
C2351 & 1a<o2 1b<o2L / 2a<1a1,1b1L / 3a<2a0L & LI/I/N & {-}+/+/+ & BF & CPW:1b,2a \\
C2352 & 1a<o2 1b<o2L / 2a<1a1,1b1L / 3a<2a1 & LI/I/I & {-}+/+/{-} & BF & CPW:1b,2a,3a \\
C2353 & 1a<o2 1b<o2L / 2a<1a1,1b1L / 3a<2a1L & LI/I/I & {-}+/+/+ & BF & CPW:1b,2a,3a \\
C2354 & 1a<o2 1b<o2L / 2a<1a1,1b1L / 3a<2a2 & LI/I/I & {-}+/+/{-} & BF & CPW:1b,2a,3a \\
C2355 & 1a<o2 1b<o2L / 2a<1a1,1b1L / 3a<2a2L & LI/I/I & {-}+/+/+ & BF & CPW:1b,2a,3a \\
C2356 & 1a<o2 1b<o2L / 2a<1a1,1b2 / 3a<2a0L & LI/E/N & {-}+/{-}/{-} & BF & CPW:1b \\
C2357 & 1a<o2 1b<o2L / 2a<1a1,1b2 / 3a<2a1 & LI/E/D & {-}+/{-}/{-} & BF & CPW:1b \\
C2358 & 1a<o2 1b<o2L / 2a<1a1,1b2 / 3a<2a1L & LI/E/D & {-}+/{-}/{-} & BF & CPW:1b \\
C2359 & 1a<o2 1b<o2L / 2a<1a1,1b2 / 3a<2a2 & LI/E/D & {-}+/{-}/{-} & BF & CPW:1b \\
C2360 & 1a<o2 1b<o2L / 2a<1a1,1b2 / 3a<2a2L & LI/E/D & {-}+/{-}/{-} & BF & CPW:1b \\
C2361 & 1a<o2 1b<o2L / 2a<1a1,1b2L / 3a<2a0L & LI/I/N & {-}+/+/+ & BF & CPW:1b,2a \\
C2362 & 1a<o2 1b<o2L / 2a<1a1,1b2L / 3a<2a1 & LI/I/I & {-}+/+/{-} & BF & CPW:1b,2a,3a \\
C2363 & 1a<o2 1b<o2L / 2a<1a1,1b2L / 3a<2a1L & LI/I/I & {-}+/+/+ & BF & CPW:1b,2a,3a \\
C2364 & 1a<o2 1b<o2L / 2a<1a1,1b2L / 3a<2a2 & LI/I/I & {-}+/+/{-} & BF & CPW:1b,2a,3a \\
C2365 & 1a<o2 1b<o2L / 2a<1a1,1b2L / 3a<2a2L & LI/I/I & {-}+/+/+ & BF & CPW:1b,2a,3a \\
C2366 & 1a<o2 1b<o2L / 2a<1a1L,1b1 / 3a<2a0L & LI/E/N & {-}+/{-}/{-} & BF & CPW:1b \\
C2367 & 1a<o2 1b<o2L / 2a<1a1L,1b1 / 3a<2a1 & LI/E/D & {-}+/{-}/{-} & BF & CPW:1b \\
C2368 & 1a<o2 1b<o2L / 2a<1a1L,1b1 / 3a<2a1L & LI/E/D & {-}+/{-}/{-} & BF & CPW:1b \\
C2369 & 1a<o2 1b<o2L / 2a<1a1L,1b1 / 3a<2a2 & LI/E/D & {-}+/{-}/{-} & BF & CPW:1b \\
C2370 & 1a<o2 1b<o2L / 2a<1a1L,1b1 / 3a<2a2L & LI/E/D & {-}+/{-}/{-} & BF & CPW:1b \\
C2371 & 1a<o2 1b<o2L / 2a<1a1L,1b2 / 3a<2a0L & LI/E/N & {-}+/{-}/{-} & BF & CPW:1b \\
C2372 & 1a<o2 1b<o2L / 2a<1a1L,1b2 / 3a<2a1 & LI/E/D & {-}+/{-}/{-} & BF & CPW:1b \\
C2373 & 1a<o2 1b<o2L / 2a<1a1L,1b2 / 3a<2a1L & LI/E/D & {-}+/{-}/{-} & BF & CPW:1b \\
C2374 & 1a<o2 1b<o2L / 2a<1a1L,1b2 / 3a<2a2 & LI/E/D & {-}+/{-}/{-} & BF & CPW:1b \\
C2375 & 1a<o2 1b<o2L / 2a<1a1L,1b2 / 3a<2a2L & LI/E/D & {-}+/{-}/{-} & BF & CPW:1b \\
C2376 & 1a<o2 1b<o2L / 2a<1a2,1b0L / 3a<2a0L & LI/D/N & {-}+/+/+ & B & CPW:1b \\
C2377 & 1a<o2 1b<o2L / 2a<1a2,1b0L / 3a<2a1 & LI/D/D & {-}+/+/{-} & B & CPW:1b \\
C2378 & 1a<o2 1b<o2L / 2a<1a2,1b0L / 3a<2a1L & LI/D/D & {-}+/+/+ & B & CPW:1b \\
C2379 & 1a<o2 1b<o2L / 2a<1a2,1b0L / 3a<2a2 & LI/D/D & {-}+/+/{-} & B & CPW:1b \\
C2380 & 1a<o2 1b<o2L / 2a<1a2,1b0L / 3a<2a2L & LI/D/D & {-}+/+/+ & B & CPW:1b \\
C2381 & 1a<o2 1b<o2L / 2a<1a2,1b1 / 3a<2a0L & LI/E/N & {-}+/{-}/{-} & BF & CPW:1b \\
C2382 & 1a<o2 1b<o2L / 2a<1a2,1b1 / 3a<2a1 & LI/E/D & {-}+/{-}/{-} & BF & CPW:1b \\
C2383 & 1a<o2 1b<o2L / 2a<1a2,1b1 / 3a<2a1L & LI/E/D & {-}+/{-}/{-} & BF & CPW:1b \\
C2384 & 1a<o2 1b<o2L / 2a<1a2,1b1 / 3a<2a2 & LI/E/D & {-}+/{-}/{-} & BF & CPW:1b \\
C2385 & 1a<o2 1b<o2L / 2a<1a2,1b1 / 3a<2a2L & LI/E/D & {-}+/{-}/{-} & BF & CPW:1b \\
C2386 & 1a<o2 1b<o2L / 2a<1a2,1b1L / 3a<2a0L & LI/I/N & {-}+/+/+ & BF & CPW:1b,2a \\
C2387 & 1a<o2 1b<o2L / 2a<1a2,1b1L / 3a<2a1 & LI/I/I & {-}+/+/{-} & BF & CPW:1b,2a,3a \\
C2388 & 1a<o2 1b<o2L / 2a<1a2,1b1L / 3a<2a1L & LI/I/I & {-}+/+/+ & BF & CPW:1b,2a,3a \\
C2389 & 1a<o2 1b<o2L / 2a<1a2,1b1L / 3a<2a2 & LI/I/I & {-}+/+/{-} & BF & CPW:1b,2a,3a \\
C2390 & 1a<o2 1b<o2L / 2a<1a2,1b1L / 3a<2a2L & LI/I/I & {-}+/+/+ & BF & CPW:1b,2a,3a \\
C2391 & 1a<o2 1b<o2L / 2a<1a2,1b2 / 3a<2a0L & LI/E/N & {-}+/{-}/{-} & BF & CPW:1b \\
C2392 & 1a<o2 1b<o2L / 2a<1a2,1b2 / 3a<2a1 & LI/E/D & {-}+/{-}/{-} & BF & CPW:1b \\
C2393 & 1a<o2 1b<o2L / 2a<1a2,1b2 / 3a<2a1L & LI/E/D & {-}+/{-}/{-} & BF & CPW:1b \\
C2394 & 1a<o2 1b<o2L / 2a<1a2,1b2 / 3a<2a2 & LI/E/D & {-}+/{-}/{-} & BF & CPW:1b \\
C2395 & 1a<o2 1b<o2L / 2a<1a2,1b2 / 3a<2a2L & LI/E/D & {-}+/{-}/{-} & BF & CPW:1b \\
C2396 & 1a<o2 1b<o2L / 2a<1a2,1b2L / 3a<2a0L & LI/I/N & {-}+/+/+ & BF & CPW:1b,2a \\
C2397 & 1a<o2 1b<o2L / 2a<1a2,1b2L / 3a<2a1 & LI/I/I & {-}+/+/{-} & BF & CPW:1b,2a,3a \\
C2398 & 1a<o2 1b<o2L / 2a<1a2,1b2L / 3a<2a1L & LI/I/I & {-}+/+/+ & BF & CPW:1b,2a,3a \\
C2399 & 1a<o2 1b<o2L / 2a<1a2,1b2L / 3a<2a2 & LI/I/I & {-}+/+/{-} & BF & CPW:1b,2a,3a \\
C2400 & 1a<o2 1b<o2L / 2a<1a2,1b2L / 3a<2a2L & LI/I/I & {-}+/+/+ & BF & CPW:1b,2a,3a \\
C2401 & 1a<o2 1b<o2L / 2a<1a2L,1b1 / 3a<2a0L & LI/E/N & {-}+/{-}/{-} & BF & CPW:1b \\
C2402 & 1a<o2 1b<o2L / 2a<1a2L,1b1 / 3a<2a1 & LI/E/D & {-}+/{-}/{-} & BF & CPW:1b \\
C2403 & 1a<o2 1b<o2L / 2a<1a2L,1b1 / 3a<2a1L & LI/E/D & {-}+/{-}/{-} & BF & CPW:1b \\
C2404 & 1a<o2 1b<o2L / 2a<1a2L,1b1 / 3a<2a2 & LI/E/D & {-}+/{-}/{-} & BF & CPW:1b \\
C2405 & 1a<o2 1b<o2L / 2a<1a2L,1b1 / 3a<2a2L & LI/E/D & {-}+/{-}/{-} & BF & CPW:1b \\
C2406 & 1a<o2 1b<o2L / 2a<1a2L,1b2 / 3a<2a0L & LI/E/N & {-}+/{-}/{-} & BF & CPW:1b \\
C2407 & 1a<o2 1b<o2L / 2a<1a2L,1b2 / 3a<2a1 & LI/E/D & {-}+/{-}/{-} & BF & CPW:1b \\
C2408 & 1a<o2 1b<o2L / 2a<1a2L,1b2 / 3a<2a1L & LI/E/D & {-}+/{-}/{-} & BF & CPW:1b \\
C2409 & 1a<o2 1b<o2L / 2a<1a2L,1b2 / 3a<2a2 & LI/E/D & {-}+/{-}/{-} & BF & CPW:1b \\
C2410 & 1a<o2 1b<o2L / 2a<1a2L,1b2 / 3a<2a2L & LI/E/D & {-}+/{-}/{-} & BF & CPW:1b \\
C2411 & 1a<o2L / 2a<1a0L 2b<1a1 / 3a<2a0L,2b1 & I/NI/I & +/+{-}/+ & N &  \\
C2412 & 1a<o2L / 2a<1a0L 2b<1a1 / 3a<2a0L,2b2 & I/NI/I & +/+{-}/+ & N &  \\
C2413 & 1a<o2L / 2a<1a0L 2b<1a1 / 3a<2a1,2b0L & I/NI/N & +/+{-}/{-} & N &  \\
C2414 & 1a<o2L / 2a<1a0L 2b<1a1 / 3a<2a1,2b1 & I/NI/E & +/+{-}/{-} & F & W:3a \\
C2415 & 1a<o2L / 2a<1a0L 2b<1a1 / 3a<2a1,2b1L & I/NI/I & +/+{-}/{-} & F & CP:3a \\
C2416 & 1a<o2L / 2a<1a0L 2b<1a1 / 3a<2a1,2b2 & I/NI/E & +/+{-}/{-} & F & CW:3a \\
C2417 & 1a<o2L / 2a<1a0L 2b<1a1 / 3a<2a1,2b2L & I/NI/I & +/+{-}/{-} & F & P:3a \\
C2418 & 1a<o2L / 2a<1a0L 2b<1a1 / 3a<2a1L,2b1 & I/NI/E & +/+{-}/+ & F & W:3a \\
C2419 & 1a<o2L / 2a<1a0L 2b<1a1 / 3a<2a1L,2b2 & I/NI/E & +/+{-}/+ & F & CW:3a \\
C2420 & 1a<o2L / 2a<1a0L 2b<1a1 / 3a<2a2,2b0L & I/NI/N & +/+{-}/{-} & N &  \\
C2421 & 1a<o2L / 2a<1a0L 2b<1a1 / 3a<2a2,2b1 & I/NI/E & +/+{-}/{-} & F & W:3a \\
C2422 & 1a<o2L / 2a<1a0L 2b<1a1 / 3a<2a2,2b1L & I/NI/I & +/+{-}/{-} & F & CP:3a \\
C2423 & 1a<o2L / 2a<1a0L 2b<1a1 / 3a<2a2,2b2 & I/NI/E & +/+{-}/{-} & F & W:3a \\
C2424 & 1a<o2L / 2a<1a0L 2b<1a1 / 3a<2a2,2b2L & I/NI/I & +/+{-}/{-} & F & CP:3a \\
C2425 & 1a<o2L / 2a<1a0L 2b<1a1 / 3a<2a2L,2b1 & I/NI/E & +/+{-}/+ & F & W:3a \\
C2426 & 1a<o2L / 2a<1a0L 2b<1a1 / 3a<2a2L,2b2 & I/NI/E & +/+{-}/+ & F & W:3a \\
C2427 & 1a<o2L / 2a<1a0L 2b<1a2 / 3a<2a0L,2b1 & I/NI/I & +/+{-}/+ & N &  \\
C2428 & 1a<o2L / 2a<1a0L 2b<1a2 / 3a<2a0L,2b2 & I/NI/I & +/+{-}/+ & N &  \\
C2429 & 1a<o2L / 2a<1a0L 2b<1a2 / 3a<2a1,2b0L & I/NI/N & +/+{-}/{-} & N &  \\
C2430 & 1a<o2L / 2a<1a0L 2b<1a2 / 3a<2a1,2b1 & I/NI/E & +/+{-}/{-} & F & W:3a \\
C2431 & 1a<o2L / 2a<1a0L 2b<1a2 / 3a<2a1,2b1L & I/NI/I & +/+{-}/{-} & F & CP:3a \\
C2432 & 1a<o2L / 2a<1a0L 2b<1a2 / 3a<2a1,2b2 & I/NI/E & +/+{-}/{-} & F & CW:3a \\
C2433 & 1a<o2L / 2a<1a0L 2b<1a2 / 3a<2a1,2b2L & I/NI/I & +/+{-}/{-} & F & P:3a \\
C2434 & 1a<o2L / 2a<1a0L 2b<1a2 / 3a<2a1L,2b1 & I/NI/E & +/+{-}/+ & F & W:3a \\
C2435 & 1a<o2L / 2a<1a0L 2b<1a2 / 3a<2a1L,2b2 & I/NI/E & +/+{-}/+ & F & CW:3a \\
C2436 & 1a<o2L / 2a<1a0L 2b<1a2 / 3a<2a2,2b0L & I/NI/N & +/+{-}/{-} & N &  \\
C2437 & 1a<o2L / 2a<1a0L 2b<1a2 / 3a<2a2,2b1 & I/NI/E & +/+{-}/{-} & F & W:3a \\
C2438 & 1a<o2L / 2a<1a0L 2b<1a2 / 3a<2a2,2b1L & I/NI/I & +/+{-}/{-} & F & CP:3a \\
C2439 & 1a<o2L / 2a<1a0L 2b<1a2 / 3a<2a2,2b2 & I/NI/E & +/+{-}/{-} & F & W:3a \\
C2440 & 1a<o2L / 2a<1a0L 2b<1a2 / 3a<2a2,2b2L & I/NI/I & +/+{-}/{-} & F & CP:3a \\
C2441 & 1a<o2L / 2a<1a0L 2b<1a2 / 3a<2a2L,2b1 & I/NI/E & +/+{-}/+ & F & W:3a \\
C2442 & 1a<o2L / 2a<1a0L 2b<1a2 / 3a<2a2L,2b2 & I/NI/E & +/+{-}/+ & F & W:3a \\
C2443 & 1a<o2L / 2a<1a1 2b<1a1 / 3a<2a0L,2b1 & I/EE/D & +/{-}{-}/{-} & E & W:1a \\
C2444 & 1a<o2L / 2a<1a1 2b<1a1 / 3a<2a0L,2b2 & I/EE/D & +/{-}{-}/{-} & E & W:1a \\
C2445 & 1a<o2L / 2a<1a1 2b<1a1 / 3a<2a1,2b1 & I/EE/D & +/{-}{-}/{-} & E & W:1a \\
C2446 & 1a<o2L / 2a<1a1 2b<1a1 / 3a<2a1,2b1L & I/EE/D & +/{-}{-}/{-} & E & W:1a \\
C2447 & 1a<o2L / 2a<1a1 2b<1a1 / 3a<2a1,2b2 & I/EE/D & +/{-}{-}/{-} & E & W:1a \\
C2448 & 1a<o2L / 2a<1a1 2b<1a1 / 3a<2a1,2b2L & I/EE/D & +/{-}{-}/{-} & E & W:1a \\
C2449 & 1a<o2L / 2a<1a1 2b<1a1 / 3a<2a1L,2b2 & I/EE/D & +/{-}{-}/{-} & E & W:1a \\
C2450 & 1a<o2L / 2a<1a1 2b<1a1 / 3a<2a2,2b2 & I/EE/D & +/{-}{-}/{-} & E & W:1a \\
C2451 & 1a<o2L / 2a<1a1 2b<1a1 / 3a<2a2,2b2L & I/EE/D & +/{-}{-}/{-} & E & W:1a \\
C2452 & 1a<o2L / 2a<1a1 2b<1a1L / 3a<2a0L,2b1 & I/LI/I & +/{-}+/{-} & B & CP:2b,3a W:1a,2b,3a \\
C2453 & 1a<o2L / 2a<1a1 2b<1a1L / 3a<2a0L,2b2 & I/LI/I & +/{-}+/{-} & B & CP:2b,3a W:1a,2b,3a \\
C2454 & 1a<o2L / 2a<1a1 2b<1a1L / 3a<2a1,2b0L & I/LI/D & +/{-}+/+ & B & CP:2b W:1a,2b \\
C2455 & 1a<o2L / 2a<1a1 2b<1a1L / 3a<2a1,2b1 & I/LI/E & +/{-}+/{-} & BF & CP:2b W:1a,2b \\
C2456 & 1a<o2L / 2a<1a1 2b<1a1L / 3a<2a1,2b1L & I/LI/I & +/{-}+/+ & BF & CP:2b,3a W:1a,2b,3a \\
C2457 & 1a<o2L / 2a<1a1 2b<1a1L / 3a<2a1,2b2 & I/LI/E & +/{-}+/{-} & BF & CP:2b W:1a,2b \\
C2458 & 1a<o2L / 2a<1a1 2b<1a1L / 3a<2a1,2b2L & I/LI/I & +/{-}+/+ & BF & CP:2b,3a W:1a,2b,3a \\
C2459 & 1a<o2L / 2a<1a1 2b<1a1L / 3a<2a1L,2b1 & I/LI/E & +/{-}+/{-} & BF & CP:2b W:1a,2b \\
C2460 & 1a<o2L / 2a<1a1 2b<1a1L / 3a<2a1L,2b2 & I/LI/E & +/{-}+/{-} & BF & CP:2b W:1a,2b \\
C2461 & 1a<o2L / 2a<1a1 2b<1a1L / 3a<2a2,2b0L & I/LI/D & +/{-}+/+ & B & CP:2b W:1a,2b \\
C2462 & 1a<o2L / 2a<1a1 2b<1a1L / 3a<2a2,2b1 & I/LI/E & +/{-}+/{-} & BF & CP:2b W:1a,2b \\
C2463 & 1a<o2L / 2a<1a1 2b<1a1L / 3a<2a2,2b1L & I/LI/I & +/{-}+/+ & BF & CP:2b,3a W:1a,2b,3a \\
C2464 & 1a<o2L / 2a<1a1 2b<1a1L / 3a<2a2,2b2 & I/LI/E & +/{-}+/{-} & BF & CP:2b W:1a,2b \\
C2465 & 1a<o2L / 2a<1a1 2b<1a1L / 3a<2a2,2b2L & I/LI/I & +/{-}+/+ & BF & CP:2b,3a W:1a,2b,3a \\
C2466 & 1a<o2L / 2a<1a1 2b<1a1L / 3a<2a2L,2b1 & I/LI/E & +/{-}+/{-} & BF & CP:2b W:1a,2b \\
C2467 & 1a<o2L / 2a<1a1 2b<1a1L / 3a<2a2L,2b2 & I/LI/E & +/{-}+/{-} & BF & CP:2b W:1a,2b \\
C2468 & 1a<o2L / 2a<1a1 2b<1a2 / 3a<2a0L,2b1 & I/EE/D & +/{-}{-}/{-} & E & C:2b,3a W:1a \\
C2469 & 1a<o2L / 2a<1a1 2b<1a2 / 3a<2a0L,2b2 & I/EE/D & +/{-}{-}/{-} & E & C:2b,3a W:1a \\
C2470 & 1a<o2L / 2a<1a1 2b<1a2 / 3a<2a1,2b0L & I/EE/D & +/{-}{-}/{-} & E & C:2b W:1a \\
C2471 & 1a<o2L / 2a<1a1 2b<1a2 / 3a<2a1,2b1 & I/EE/D & +/{-}{-}/{-} & E & C:2b W:1a \\
C2472 & 1a<o2L / 2a<1a1 2b<1a2 / 3a<2a1,2b1L & I/EE/D & +/{-}{-}/{-} & E & C:2b W:1a \\
C2473 & 1a<o2L / 2a<1a1 2b<1a2 / 3a<2a1,2b2 & I/EE/D & +/{-}{-}/{-} & E & C:2b,3a W:1a \\
C2474 & 1a<o2L / 2a<1a1 2b<1a2 / 3a<2a1,2b2L & I/EE/D & +/{-}{-}/{-} & E & C:2b,3a W:1a \\
C2475 & 1a<o2L / 2a<1a1 2b<1a2 / 3a<2a1L,2b1 & I/EE/D & +/{-}{-}/{-} & E & C:2b W:1a \\
C2476 & 1a<o2L / 2a<1a1 2b<1a2 / 3a<2a1L,2b2 & I/EE/D & +/{-}{-}/{-} & E & C:2b,3a W:1a \\
C2477 & 1a<o2L / 2a<1a1 2b<1a2 / 3a<2a2,2b0L & I/EE/D & +/{-}{-}/{-} & E & C:2b W:1a \\
C2478 & 1a<o2L / 2a<1a1 2b<1a2 / 3a<2a2,2b1 & I/EE/D & +/{-}{-}/{-} & E & C:2b W:1a \\
C2479 & 1a<o2L / 2a<1a1 2b<1a2 / 3a<2a2,2b1L & I/EE/D & +/{-}{-}/{-} & E & C:2b W:1a \\
C2480 & 1a<o2L / 2a<1a1 2b<1a2 / 3a<2a2,2b2 & I/EE/D & +/{-}{-}/{-} & E & C:2b W:1a \\
C2481 & 1a<o2L / 2a<1a1 2b<1a2 / 3a<2a2,2b2L & I/EE/D & +/{-}{-}/{-} & E & C:2b W:1a \\
C2482 & 1a<o2L / 2a<1a1 2b<1a2 / 3a<2a2L,2b1 & I/EE/D & +/{-}{-}/{-} & E & C:2b W:1a \\
C2483 & 1a<o2L / 2a<1a1 2b<1a2 / 3a<2a2L,2b2 & I/EE/D & +/{-}{-}/{-} & E & C:2b W:1a \\
C2484 & 1a<o2L / 2a<1a1 2b<1a2L / 3a<2a0L,2b1 & I/LI/I & +/{-}+/{-} & B & P:2b,3a W:1a,2b,3a \\
C2485 & 1a<o2L / 2a<1a1 2b<1a2L / 3a<2a0L,2b2 & I/LI/I & +/{-}+/{-} & B & P:2b,3a W:1a,2b,3a \\
C2486 & 1a<o2L / 2a<1a1 2b<1a2L / 3a<2a1,2b0L & I/LI/D & +/{-}+/+ & B & P:2b W:1a,2b \\
C2487 & 1a<o2L / 2a<1a1 2b<1a2L / 3a<2a1,2b1 & I/LI/E & +/{-}+/{-} & BF & P:2b W:1a,2b \\
C2488 & 1a<o2L / 2a<1a1 2b<1a2L / 3a<2a1,2b1L & I/LI/I & +/{-}+/+ & BF & C:3a P:2b,3a W:1a,2b,3a \\
C2489 & 1a<o2L / 2a<1a1 2b<1a2L / 3a<2a1,2b2 & I/LI/E & +/{-}+/{-} & BF & C:3a P:2b W:1a,2b \\
C2490 & 1a<o2L / 2a<1a1 2b<1a2L / 3a<2a1,2b2L & I/LI/I & +/{-}+/+ & BF & P:2b,3a W:1a,2b,3a \\
C2491 & 1a<o2L / 2a<1a1 2b<1a2L / 3a<2a1L,2b1 & I/LI/E & +/{-}+/{-} & BF & P:2b W:1a,2b \\
C2492 & 1a<o2L / 2a<1a1 2b<1a2L / 3a<2a1L,2b2 & I/LI/E & +/{-}+/{-} & BF & C:3a P:2b W:1a,2b \\
C2493 & 1a<o2L / 2a<1a1 2b<1a2L / 3a<2a2,2b0L & I/LI/D & +/{-}+/+ & B & P:2b W:1a,2b \\
C2494 & 1a<o2L / 2a<1a1 2b<1a2L / 3a<2a2,2b1 & I/LI/E & +/{-}+/{-} & BF & P:2b W:1a,2b \\
C2495 & 1a<o2L / 2a<1a1 2b<1a2L / 3a<2a2,2b1L & I/LI/I & +/{-}+/+ & BF & C:3a P:2b,3a W:1a,2b,3a \\
C2496 & 1a<o2L / 2a<1a1 2b<1a2L / 3a<2a2,2b2 & I/LI/E & +/{-}+/{-} & BF & P:2b W:1a,2b \\
C2497 & 1a<o2L / 2a<1a1 2b<1a2L / 3a<2a2,2b2L & I/LI/I & +/{-}+/+ & BF & C:3a P:2b,3a W:1a,2b,3a \\
C2498 & 1a<o2L / 2a<1a1 2b<1a2L / 3a<2a2L,2b1 & I/LI/E & +/{-}+/{-} & BF & P:2b W:1a,2b \\
C2499 & 1a<o2L / 2a<1a1 2b<1a2L / 3a<2a2L,2b2 & I/LI/E & +/{-}+/{-} & BF & P:2b W:1a,2b \\
C2500 & 1a<o2L / 2a<1a1L 2b<1a2 / 3a<2a0L,2b1 & I/IL/D & +/+{-}/+ & B & C:2a,2b,3a P:2a W:1a,2a \\
C2501 & 1a<o2L / 2a<1a1L 2b<1a2 / 3a<2a0L,2b2 & I/IL/D & +/+{-}/+ & B & C:2a,2b,3a P:2a W:1a,2a \\
C2502 & 1a<o2L / 2a<1a1L 2b<1a2 / 3a<2a1,2b0L & I/IL/I & +/+{-}/{-} & B & C:2a,2b,3a P:2a,3a W:1a,2a,3a \\
C2503 & 1a<o2L / 2a<1a1L 2b<1a2 / 3a<2a1,2b1 & I/IL/E & +/+{-}/{-} & BF & C:2a,2b P:2a W:1a,2a \\
C2504 & 1a<o2L / 2a<1a1L 2b<1a2 / 3a<2a1,2b1L & I/IL/E & +/+{-}/{-} & BF & C:2a,2b P:2a W:1a,2a \\
C2505 & 1a<o2L / 2a<1a1L 2b<1a2 / 3a<2a1,2b2 & I/IL/E & +/+{-}/{-} & BF & C:2a,2b,3a P:2a W:1a,2a \\
C2506 & 1a<o2L / 2a<1a1L 2b<1a2 / 3a<2a1,2b2L & I/IL/E & +/+{-}/{-} & BF & C:2a,2b,3a P:2a W:1a,2a \\
C2507 & 1a<o2L / 2a<1a1L 2b<1a2 / 3a<2a1L,2b1 & I/IL/I & +/+{-}/+ & BF & C:2a,2b,3a P:2a,3a W:1a,2a,3a \\
C2508 & 1a<o2L / 2a<1a1L 2b<1a2 / 3a<2a1L,2b2 & I/IL/I & +/+{-}/+ & BF & C:2a,2b P:2a,3a W:1a,2a,3a \\
C2509 & 1a<o2L / 2a<1a1L 2b<1a2 / 3a<2a2,2b0L & I/IL/I & +/+{-}/{-} & B & C:2a,2b,3a P:2a,3a W:1a,2a,3a \\
C2510 & 1a<o2L / 2a<1a1L 2b<1a2 / 3a<2a2,2b1 & I/IL/E & +/+{-}/{-} & BF & C:2a,2b P:2a W:1a,2a \\
C2511 & 1a<o2L / 2a<1a1L 2b<1a2 / 3a<2a2,2b1L & I/IL/E & +/+{-}/{-} & BF & C:2a,2b P:2a W:1a,2a \\
C2512 & 1a<o2L / 2a<1a1L 2b<1a2 / 3a<2a2,2b2 & I/IL/E & +/+{-}/{-} & BF & C:2a,2b P:2a W:1a,2a \\
C2513 & 1a<o2L / 2a<1a1L 2b<1a2 / 3a<2a2,2b2L & I/IL/E & +/+{-}/{-} & BF & C:2a,2b P:2a W:1a,2a \\
C2514 & 1a<o2L / 2a<1a1L 2b<1a2 / 3a<2a2L,2b1 & I/IL/I & +/+{-}/+ & BF & C:2a,2b,3a P:2a,3a W:1a,2a,3a \\
C2515 & 1a<o2L / 2a<1a1L 2b<1a2 / 3a<2a2L,2b2 & I/IL/I & +/+{-}/+ & BF & C:2a,2b,3a P:2a,3a W:1a,2a,3a \\
C2516 & 1a<o2L / 2a<1a2 2b<1a2 / 3a<2a0L,2b1 & I/EE/D & +/{-}{-}/{-} & E & W:1a \\
C2517 & 1a<o2L / 2a<1a2 2b<1a2 / 3a<2a0L,2b2 & I/EE/D & +/{-}{-}/{-} & E & W:1a \\
C2518 & 1a<o2L / 2a<1a2 2b<1a2 / 3a<2a1,2b1 & I/EE/D & +/{-}{-}/{-} & E & W:1a \\
C2519 & 1a<o2L / 2a<1a2 2b<1a2 / 3a<2a1,2b1L & I/EE/D & +/{-}{-}/{-} & E & W:1a \\
C2520 & 1a<o2L / 2a<1a2 2b<1a2 / 3a<2a1,2b2 & I/EE/D & +/{-}{-}/{-} & E & W:1a \\
C2521 & 1a<o2L / 2a<1a2 2b<1a2 / 3a<2a1,2b2L & I/EE/D & +/{-}{-}/{-} & E & W:1a \\
C2522 & 1a<o2L / 2a<1a2 2b<1a2 / 3a<2a1L,2b2 & I/EE/D & +/{-}{-}/{-} & E & W:1a \\
C2523 & 1a<o2L / 2a<1a2 2b<1a2 / 3a<2a2,2b2 & I/EE/D & +/{-}{-}/{-} & E & W:1a \\
C2524 & 1a<o2L / 2a<1a2 2b<1a2 / 3a<2a2,2b2L & I/EE/D & +/{-}{-}/{-} & E & W:1a \\
C2525 & 1a<o2L / 2a<1a2 2b<1a2L / 3a<2a0L,2b1 & I/LI/I & +/{-}+/{-} & B & CP:2b,3a W:1a,2b,3a \\
C2526 & 1a<o2L / 2a<1a2 2b<1a2L / 3a<2a0L,2b2 & I/LI/I & +/{-}+/{-} & B & CP:2b,3a W:1a,2b,3a \\
C2527 & 1a<o2L / 2a<1a2 2b<1a2L / 3a<2a1,2b0L & I/LI/D & +/{-}+/+ & B & CP:2b W:1a,2b \\
C2528 & 1a<o2L / 2a<1a2 2b<1a2L / 3a<2a1,2b1 & I/LI/E & +/{-}+/{-} & BF & CP:2b W:1a,2b \\
C2529 & 1a<o2L / 2a<1a2 2b<1a2L / 3a<2a1,2b1L & I/LI/I & +/{-}+/+ & BF & CP:2b,3a W:1a,2b,3a \\
C2530 & 1a<o2L / 2a<1a2 2b<1a2L / 3a<2a1,2b2 & I/LI/E & +/{-}+/{-} & BF & CP:2b W:1a,2b \\
C2531 & 1a<o2L / 2a<1a2 2b<1a2L / 3a<2a1,2b2L & I/LI/I & +/{-}+/+ & BF & CP:2b,3a W:1a,2b,3a \\
C2532 & 1a<o2L / 2a<1a2 2b<1a2L / 3a<2a1L,2b1 & I/LI/E & +/{-}+/{-} & BF & CP:2b W:1a,2b \\
C2533 & 1a<o2L / 2a<1a2 2b<1a2L / 3a<2a1L,2b2 & I/LI/E & +/{-}+/{-} & BF & CP:2b W:1a,2b \\
C2534 & 1a<o2L / 2a<1a2 2b<1a2L / 3a<2a2,2b0L & I/LI/D & +/{-}+/+ & B & CP:2b W:1a,2b \\
C2535 & 1a<o2L / 2a<1a2 2b<1a2L / 3a<2a2,2b1 & I/LI/E & +/{-}+/{-} & BF & CP:2b W:1a,2b \\
C2536 & 1a<o2L / 2a<1a2 2b<1a2L / 3a<2a2,2b1L & I/LI/I & +/{-}+/+ & BF & CP:2b,3a W:1a,2b,3a \\
C2537 & 1a<o2L / 2a<1a2 2b<1a2L / 3a<2a2,2b2 & I/LI/E & +/{-}+/{-} & BF & CP:2b W:1a,2b \\
C2538 & 1a<o2L / 2a<1a2 2b<1a2L / 3a<2a2,2b2L & I/LI/I & +/{-}+/+ & BF & CP:2b,3a W:1a,2b,3a \\
C2539 & 1a<o2L / 2a<1a2 2b<1a2L / 3a<2a2L,2b1 & I/LI/E & +/{-}+/{-} & BF & CP:2b W:1a,2b \\
C2540 & 1a<o2L / 2a<1a2 2b<1a2L / 3a<2a2L,2b2 & I/LI/E & +/{-}+/{-} & BF & CP:2b W:1a,2b \\
C2541 & 1a<o0L / 2a<1a0L / 3a<2a0L 3b*<2a1 & N/N/NN & +/+/+{-} & N &  \\
C2542 & 1a<o0L / 2a<1a0L / 3a<2a0L 3b*<2a2 & N/N/NN & +/+/+{-} & N &  \\
C2543 & 1a<o0L / 2a<1a0L / 3a<2a1 3b*<2a1 & N/N/NN & +/+/{-}{-} & N &  \\
C2544 & 1a<o0L / 2a<1a0L / 3a<2a1 3b*<2a2 & N/N/NN & +/+/{-}{-} & N &  \\
C2545 & 1a<o0L / 2a<1a0L / 3a<2a1L 3b*<2a1 & N/N/NN & +/+/+{-} & N &  \\
C2546 & 1a<o0L / 2a<1a0L / 3a<2a1L 3b*<2a2 & N/N/NN & +/+/+{-} & N &  \\
C2547 & 1a<o0L / 2a<1a0L / 3a<2a2 3b*<2a1 & N/N/NN & +/+/{-}{-} & N &  \\
C2548 & 1a<o0L / 2a<1a0L / 3a<2a2 3b*<2a2 & N/N/NN & +/+/{-}{-} & N &  \\
C2549 & 1a<o0L / 2a<1a0L / 3a<2a2L 3b*<2a1 & N/N/NN & +/+/+{-} & N &  \\
C2550 & 1a<o0L / 2a<1a0L / 3a<2a2L 3b*<2a2 & N/N/NN & +/+/+{-} & N &  \\
C2551 & 1a<o0L / 2a<1a0L 2b<1a1 / 3a*<2a1 & N/NN/N & +/+{-}/{-} & N &  \\
C2552 & 1a<o0L / 2a<1a0L 2b<1a1 / 3a*<2a2 & N/NN/N & +/+{-}/{-} & N &  \\
C2553 & 1a<o0L / 2a<1a0L 2b<1a1 / 3a*<2b1 & N/NN/N & +/+{-}/{-} & N &  \\
C2554 & 1a<o0L / 2a<1a0L 2b<1a1 / 3a*<2b2 & N/NN/N & +/+{-}/{-} & N &  \\
C2555 & 1a<o0L / 2a<1a0L 2b<1a2 / 3a*<2a1 & N/NN/N & +/+{-}/{-} & N &  \\
C2556 & 1a<o0L / 2a<1a0L 2b<1a2 / 3a*<2a2 & N/NN/N & +/+{-}/{-} & N &  \\
C2557 & 1a<o0L / 2a<1a0L 2b<1a2 / 3a*<2b1 & N/NN/N & +/+{-}/{-} & N &  \\
C2558 & 1a<o0L / 2a<1a0L 2b<1a2 / 3a*<2b2 & N/NN/N & +/+{-}/{-} & N &  \\
C2559 & 1a<o0L / 2a<1a0L 2b*<1a1 / 3a<2a0L & N/NN/N & +/+{-}/+ & N &  \\
C2560 & 1a<o0L / 2a<1a0L 2b*<1a1 / 3a<2a1 & N/NN/N & +/+{-}/{-} & N &  \\
C2561 & 1a<o0L / 2a<1a0L 2b*<1a1 / 3a<2a1L & N/NN/N & +/+{-}/+ & N &  \\
C2562 & 1a<o0L / 2a<1a0L 2b*<1a1 / 3a<2a2 & N/NN/N & +/+{-}/{-} & N &  \\
C2563 & 1a<o0L / 2a<1a0L 2b*<1a1 / 3a<2a2L & N/NN/N & +/+{-}/+ & N &  \\
C2564 & 1a<o0L / 2a<1a0L 2b*<1a1 / 3a<2b1 & N/NN/N & +/+{-}/{-} & N &  \\
C2565 & 1a<o0L / 2a<1a0L 2b*<1a1 / 3a<2b2 & N/NN/N & +/+{-}/{-} & N &  \\
C2566 & 1a<o0L / 2a<1a0L 2b*<1a2 / 3a<2a0L & N/NN/N & +/+{-}/+ & N &  \\
C2567 & 1a<o0L / 2a<1a0L 2b*<1a2 / 3a<2a1 & N/NN/N & +/+{-}/{-} & N &  \\
C2568 & 1a<o0L / 2a<1a0L 2b*<1a2 / 3a<2a1L & N/NN/N & +/+{-}/+ & N &  \\
C2569 & 1a<o0L / 2a<1a0L 2b*<1a2 / 3a<2a2 & N/NN/N & +/+{-}/{-} & N &  \\
C2570 & 1a<o0L / 2a<1a0L 2b*<1a2 / 3a<2a2L & N/NN/N & +/+{-}/+ & N &  \\
C2571 & 1a<o0L / 2a<1a0L 2b*<1a2 / 3a<2b1 & N/NN/N & +/+{-}/{-} & N &  \\
C2572 & 1a<o0L / 2a<1a0L 2b*<1a2 / 3a<2b2 & N/NN/N & +/+{-}/{-} & N &  \\
C2573 & 1a<o0L / 2a<1a1 / 3a<2a0L 3b*<2a1 & N/N/NN & +/{-}/{-}{-} & N &  \\
C2574 & 1a<o0L / 2a<1a1 / 3a<2a0L 3b*<2a2 & N/N/NN & +/{-}/{-}{-} & N &  \\
C2575 & 1a<o0L / 2a<1a1 / 3a<2a1 3b*<2a1 & N/N/NN & +/{-}/{-}{-} & N &  \\
C2576 & 1a<o0L / 2a<1a1 / 3a<2a1 3b*<2a2 & N/N/NN & +/{-}/{-}{-} & N &  \\
C2577 & 1a<o0L / 2a<1a1 / 3a<2a1L 3b*<2a1 & N/N/NN & +/{-}/{-}{-} & N &  \\
C2578 & 1a<o0L / 2a<1a1 / 3a<2a1L 3b*<2a2 & N/N/NN & +/{-}/{-}{-} & N &  \\
C2579 & 1a<o0L / 2a<1a1 / 3a<2a2 3b*<2a1 & N/N/NN & +/{-}/{-}{-} & N &  \\
C2580 & 1a<o0L / 2a<1a1 / 3a<2a2 3b*<2a2 & N/N/NN & +/{-}/{-}{-} & N &  \\
C2581 & 1a<o0L / 2a<1a1 / 3a<2a2L 3b*<2a1 & N/N/NN & +/{-}/{-}{-} & N &  \\
C2582 & 1a<o0L / 2a<1a1 / 3a<2a2L 3b*<2a2 & N/N/NN & +/{-}/{-}{-} & N &  \\
C2583 & 1a<o0L / 2a<1a1 2b<1a1 / 3a*<2a1 & N/NN/N & +/{-}{-}/{-} & N &  \\
C2584 & 1a<o0L / 2a<1a1 2b<1a1 / 3a*<2a2 & N/NN/N & +/{-}{-}/{-} & N &  \\
C2585 & 1a<o0L / 2a<1a1 2b<1a1L / 3a*<2a1 & N/NN/N & +/{-}+/{-} & N &  \\
C2586 & 1a<o0L / 2a<1a1 2b<1a1L / 3a*<2a2 & N/NN/N & +/{-}+/{-} & N &  \\
C2587 & 1a<o0L / 2a<1a1 2b<1a1L / 3a*<2b1 & N/NN/N & +/{-}+/{-} & N &  \\
C2588 & 1a<o0L / 2a<1a1 2b<1a1L / 3a*<2b2 & N/NN/N & +/{-}+/{-} & N &  \\
C2589 & 1a<o0L / 2a<1a1 2b<1a2 / 3a*<2a1 & N/NN/N & +/{-}{-}/{-} & N &  \\
C2590 & 1a<o0L / 2a<1a1 2b<1a2 / 3a*<2a2 & N/NN/N & +/{-}{-}/{-} & N &  \\
C2591 & 1a<o0L / 2a<1a1 2b<1a2 / 3a*<2b1 & N/NN/N & +/{-}{-}/{-} & N &  \\
C2592 & 1a<o0L / 2a<1a1 2b<1a2 / 3a*<2b2 & N/NN/N & +/{-}{-}/{-} & N &  \\
C2593 & 1a<o0L / 2a<1a1 2b<1a2L / 3a*<2a1 & N/NN/N & +/{-}+/{-} & N &  \\
C2594 & 1a<o0L / 2a<1a1 2b<1a2L / 3a*<2a2 & N/NN/N & +/{-}+/{-} & N &  \\
C2595 & 1a<o0L / 2a<1a1 2b<1a2L / 3a*<2b1 & N/NN/N & +/{-}+/{-} & N &  \\
C2596 & 1a<o0L / 2a<1a1 2b<1a2L / 3a*<2b2 & N/NN/N & +/{-}+/{-} & N &  \\
C2597 & 1a<o0L / 2a<1a1 2b*<1a1 / 3a<2a0L & N/NN/N & +/{-}{-}/{-} & N &  \\
C2598 & 1a<o0L / 2a<1a1 2b*<1a1 / 3a<2a1 & N/NN/N & +/{-}{-}/{-} & N &  \\
C2599 & 1a<o0L / 2a<1a1 2b*<1a1 / 3a<2a1L & N/NN/N & +/{-}{-}/{-} & N &  \\
C2600 & 1a<o0L / 2a<1a1 2b*<1a1 / 3a<2a2 & N/NN/N & +/{-}{-}/{-} & N &  \\
C2601 & 1a<o0L / 2a<1a1 2b*<1a1 / 3a<2a2L & N/NN/N & +/{-}{-}/{-} & N &  \\
C2602 & 1a<o0L / 2a<1a1 2b*<1a1 / 3a<2b1 & N/NN/N & +/{-}{-}/{-} & N &  \\
C2603 & 1a<o0L / 2a<1a1 2b*<1a1 / 3a<2b2 & N/NN/N & +/{-}{-}/{-} & N &  \\
C2604 & 1a<o0L / 2a<1a1 2b*<1a2 / 3a<2a0L & N/NN/N & +/{-}{-}/{-} & N &  \\
C2605 & 1a<o0L / 2a<1a1 2b*<1a2 / 3a<2a1 & N/NN/N & +/{-}{-}/{-} & N &  \\
C2606 & 1a<o0L / 2a<1a1 2b*<1a2 / 3a<2a1L & N/NN/N & +/{-}{-}/{-} & N &  \\
C2607 & 1a<o0L / 2a<1a1 2b*<1a2 / 3a<2a2 & N/NN/N & +/{-}{-}/{-} & N &  \\
C2608 & 1a<o0L / 2a<1a1 2b*<1a2 / 3a<2a2L & N/NN/N & +/{-}{-}/{-} & N &  \\
C2609 & 1a<o0L / 2a<1a1 2b*<1a2 / 3a<2b1 & N/NN/N & +/{-}{-}/{-} & N &  \\
C2610 & 1a<o0L / 2a<1a1 2b*<1a2 / 3a<2b2 & N/NN/N & +/{-}{-}/{-} & N &  \\
C2611 & 1a<o0L / 2a<1a1L / 3a<2a0L 3b*<2a1 & N/N/NN & +/+/+{-} & N &  \\
C2612 & 1a<o0L / 2a<1a1L / 3a<2a0L 3b*<2a2 & N/N/NN & +/+/+{-} & N &  \\
C2613 & 1a<o0L / 2a<1a1L / 3a<2a1 3b*<2a1 & N/N/NN & +/+/{-}{-} & N &  \\
C2614 & 1a<o0L / 2a<1a1L / 3a<2a1 3b*<2a2 & N/N/NN & +/+/{-}{-} & N &  \\
C2615 & 1a<o0L / 2a<1a1L / 3a<2a1L 3b*<2a1 & N/N/NN & +/+/+{-} & N &  \\
C2616 & 1a<o0L / 2a<1a1L / 3a<2a1L 3b*<2a2 & N/N/NN & +/+/+{-} & N &  \\
C2617 & 1a<o0L / 2a<1a1L / 3a<2a2 3b*<2a1 & N/N/NN & +/+/{-}{-} & N &  \\
C2618 & 1a<o0L / 2a<1a1L / 3a<2a2 3b*<2a2 & N/N/NN & +/+/{-}{-} & N &  \\
C2619 & 1a<o0L / 2a<1a1L / 3a<2a2L 3b*<2a1 & N/N/NN & +/+/+{-} & N &  \\
C2620 & 1a<o0L / 2a<1a1L / 3a<2a2L 3b*<2a2 & N/N/NN & +/+/+{-} & N &  \\
C2621 & 1a<o0L / 2a<1a1L 2b<1a2 / 3a*<2a1 & N/NN/N & +/+{-}/{-} & N &  \\
C2622 & 1a<o0L / 2a<1a1L 2b<1a2 / 3a*<2a2 & N/NN/N & +/+{-}/{-} & N &  \\
C2623 & 1a<o0L / 2a<1a1L 2b<1a2 / 3a*<2b1 & N/NN/N & +/+{-}/{-} & N &  \\
C2624 & 1a<o0L / 2a<1a1L 2b<1a2 / 3a*<2b2 & N/NN/N & +/+{-}/{-} & N &  \\
C2625 & 1a<o0L / 2a<1a1L 2b*<1a1 / 3a<2a0L & N/NN/N & +/+{-}/+ & N &  \\
C2626 & 1a<o0L / 2a<1a1L 2b*<1a1 / 3a<2a1 & N/NN/N & +/+{-}/{-} & N &  \\
C2627 & 1a<o0L / 2a<1a1L 2b*<1a1 / 3a<2a1L & N/NN/N & +/+{-}/+ & N &  \\
C2628 & 1a<o0L / 2a<1a1L 2b*<1a1 / 3a<2a2 & N/NN/N & +/+{-}/{-} & N &  \\
C2629 & 1a<o0L / 2a<1a1L 2b*<1a1 / 3a<2a2L & N/NN/N & +/+{-}/+ & N &  \\
C2630 & 1a<o0L / 2a<1a1L 2b*<1a1 / 3a<2b1 & N/NN/N & +/+{-}/{-} & N &  \\
C2631 & 1a<o0L / 2a<1a1L 2b*<1a1 / 3a<2b2 & N/NN/N & +/+{-}/{-} & N &  \\
C2632 & 1a<o0L / 2a<1a1L 2b*<1a2 / 3a<2a0L & N/NN/N & +/+{-}/+ & N &  \\
C2633 & 1a<o0L / 2a<1a1L 2b*<1a2 / 3a<2a1 & N/NN/N & +/+{-}/{-} & N &  \\
C2634 & 1a<o0L / 2a<1a1L 2b*<1a2 / 3a<2a1L & N/NN/N & +/+{-}/+ & N &  \\
C2635 & 1a<o0L / 2a<1a1L 2b*<1a2 / 3a<2a2 & N/NN/N & +/+{-}/{-} & N &  \\
C2636 & 1a<o0L / 2a<1a1L 2b*<1a2 / 3a<2a2L & N/NN/N & +/+{-}/+ & N &  \\
C2637 & 1a<o0L / 2a<1a1L 2b*<1a2 / 3a<2b1 & N/NN/N & +/+{-}/{-} & N &  \\
C2638 & 1a<o0L / 2a<1a1L 2b*<1a2 / 3a<2b2 & N/NN/N & +/+{-}/{-} & N &  \\
C2639 & 1a<o0L / 2a<1a2 / 3a<2a0L 3b*<2a1 & N/N/NN & +/{-}/{-}{-} & N &  \\
C2640 & 1a<o0L / 2a<1a2 / 3a<2a0L 3b*<2a2 & N/N/NN & +/{-}/{-}{-} & N &  \\
C2641 & 1a<o0L / 2a<1a2 / 3a<2a1 3b*<2a1 & N/N/NN & +/{-}/{-}{-} & N &  \\
C2642 & 1a<o0L / 2a<1a2 / 3a<2a1 3b*<2a2 & N/N/NN & +/{-}/{-}{-} & N &  \\
C2643 & 1a<o0L / 2a<1a2 / 3a<2a1L 3b*<2a1 & N/N/NN & +/{-}/{-}{-} & N &  \\
C2644 & 1a<o0L / 2a<1a2 / 3a<2a1L 3b*<2a2 & N/N/NN & +/{-}/{-}{-} & N &  \\
C2645 & 1a<o0L / 2a<1a2 / 3a<2a2 3b*<2a1 & N/N/NN & +/{-}/{-}{-} & N &  \\
C2646 & 1a<o0L / 2a<1a2 / 3a<2a2 3b*<2a2 & N/N/NN & +/{-}/{-}{-} & N &  \\
C2647 & 1a<o0L / 2a<1a2 / 3a<2a2L 3b*<2a1 & N/N/NN & +/{-}/{-}{-} & N &  \\
C2648 & 1a<o0L / 2a<1a2 / 3a<2a2L 3b*<2a2 & N/N/NN & +/{-}/{-}{-} & N &  \\
C2649 & 1a<o0L / 2a<1a2 2b<1a2 / 3a*<2a1 & N/NN/N & +/{-}{-}/{-} & N &  \\
C2650 & 1a<o0L / 2a<1a2 2b<1a2 / 3a*<2a2 & N/NN/N & +/{-}{-}/{-} & N &  \\
C2651 & 1a<o0L / 2a<1a2 2b<1a2L / 3a*<2a1 & N/NN/N & +/{-}+/{-} & N &  \\
C2652 & 1a<o0L / 2a<1a2 2b<1a2L / 3a*<2a2 & N/NN/N & +/{-}+/{-} & N &  \\
C2653 & 1a<o0L / 2a<1a2 2b<1a2L / 3a*<2b1 & N/NN/N & +/{-}+/{-} & N &  \\
C2654 & 1a<o0L / 2a<1a2 2b<1a2L / 3a*<2b2 & N/NN/N & +/{-}+/{-} & N &  \\
C2655 & 1a<o0L / 2a<1a2 2b*<1a1 / 3a<2a0L & N/NN/N & +/{-}{-}/{-} & N &  \\
C2656 & 1a<o0L / 2a<1a2 2b*<1a1 / 3a<2a1 & N/NN/N & +/{-}{-}/{-} & N &  \\
C2657 & 1a<o0L / 2a<1a2 2b*<1a1 / 3a<2a1L & N/NN/N & +/{-}{-}/{-} & N &  \\
C2658 & 1a<o0L / 2a<1a2 2b*<1a1 / 3a<2a2 & N/NN/N & +/{-}{-}/{-} & N &  \\
C2659 & 1a<o0L / 2a<1a2 2b*<1a1 / 3a<2a2L & N/NN/N & +/{-}{-}/{-} & N &  \\
C2660 & 1a<o0L / 2a<1a2 2b*<1a1 / 3a<2b1 & N/NN/N & +/{-}{-}/{-} & N &  \\
C2661 & 1a<o0L / 2a<1a2 2b*<1a1 / 3a<2b2 & N/NN/N & +/{-}{-}/{-} & N &  \\
C2662 & 1a<o0L / 2a<1a2 2b*<1a2 / 3a<2a0L & N/NN/N & +/{-}{-}/{-} & N &  \\
C2663 & 1a<o0L / 2a<1a2 2b*<1a2 / 3a<2a1 & N/NN/N & +/{-}{-}/{-} & N &  \\
C2664 & 1a<o0L / 2a<1a2 2b*<1a2 / 3a<2a1L & N/NN/N & +/{-}{-}/{-} & N &  \\
C2665 & 1a<o0L / 2a<1a2 2b*<1a2 / 3a<2a2 & N/NN/N & +/{-}{-}/{-} & N &  \\
C2666 & 1a<o0L / 2a<1a2 2b*<1a2 / 3a<2a2L & N/NN/N & +/{-}{-}/{-} & N &  \\
C2667 & 1a<o0L / 2a<1a2 2b*<1a2 / 3a<2b1 & N/NN/N & +/{-}{-}/{-} & N &  \\
C2668 & 1a<o0L / 2a<1a2 2b*<1a2 / 3a<2b2 & N/NN/N & +/{-}{-}/{-} & N &  \\
C2669 & 1a<o0L / 2a<1a2L / 3a<2a0L 3b*<2a1 & N/N/NN & +/+/+{-} & N &  \\
C2670 & 1a<o0L / 2a<1a2L / 3a<2a0L 3b*<2a2 & N/N/NN & +/+/+{-} & N &  \\
C2671 & 1a<o0L / 2a<1a2L / 3a<2a1 3b*<2a1 & N/N/NN & +/+/{-}{-} & N &  \\
C2672 & 1a<o0L / 2a<1a2L / 3a<2a1 3b*<2a2 & N/N/NN & +/+/{-}{-} & N &  \\
C2673 & 1a<o0L / 2a<1a2L / 3a<2a1L 3b*<2a1 & N/N/NN & +/+/+{-} & N &  \\
C2674 & 1a<o0L / 2a<1a2L / 3a<2a1L 3b*<2a2 & N/N/NN & +/+/+{-} & N &  \\
C2675 & 1a<o0L / 2a<1a2L / 3a<2a2 3b*<2a1 & N/N/NN & +/+/{-}{-} & N &  \\
C2676 & 1a<o0L / 2a<1a2L / 3a<2a2 3b*<2a2 & N/N/NN & +/+/{-}{-} & N &  \\
C2677 & 1a<o0L / 2a<1a2L / 3a<2a2L 3b*<2a1 & N/N/NN & +/+/+{-} & N &  \\
C2678 & 1a<o0L / 2a<1a2L / 3a<2a2L 3b*<2a2 & N/N/NN & +/+/+{-} & N &  \\
C2679 & 1a<o0L / 2a<1a2L 2b*<1a1 / 3a<2a0L & N/NN/N & +/+{-}/+ & N &  \\
C2680 & 1a<o0L / 2a<1a2L 2b*<1a1 / 3a<2a1 & N/NN/N & +/+{-}/{-} & N &  \\
C2681 & 1a<o0L / 2a<1a2L 2b*<1a1 / 3a<2a1L & N/NN/N & +/+{-}/+ & N &  \\
C2682 & 1a<o0L / 2a<1a2L 2b*<1a1 / 3a<2a2 & N/NN/N & +/+{-}/{-} & N &  \\
C2683 & 1a<o0L / 2a<1a2L 2b*<1a1 / 3a<2a2L & N/NN/N & +/+{-}/+ & N &  \\
C2684 & 1a<o0L / 2a<1a2L 2b*<1a1 / 3a<2b1 & N/NN/N & +/+{-}/{-} & N &  \\
C2685 & 1a<o0L / 2a<1a2L 2b*<1a1 / 3a<2b2 & N/NN/N & +/+{-}/{-} & N &  \\
C2686 & 1a<o0L / 2a<1a2L 2b*<1a2 / 3a<2a0L & N/NN/N & +/+{-}/+ & N &  \\
C2687 & 1a<o0L / 2a<1a2L 2b*<1a2 / 3a<2a1 & N/NN/N & +/+{-}/{-} & N &  \\
C2688 & 1a<o0L / 2a<1a2L 2b*<1a2 / 3a<2a1L & N/NN/N & +/+{-}/+ & N &  \\
C2689 & 1a<o0L / 2a<1a2L 2b*<1a2 / 3a<2a2 & N/NN/N & +/+{-}/{-} & N &  \\
C2690 & 1a<o0L / 2a<1a2L 2b*<1a2 / 3a<2a2L & N/NN/N & +/+{-}/+ & N &  \\
C2691 & 1a<o0L / 2a<1a2L 2b*<1a2 / 3a<2b1 & N/NN/N & +/+{-}/{-} & N &  \\
C2692 & 1a<o0L / 2a<1a2L 2b*<1a2 / 3a<2b2 & N/NN/N & +/+{-}/{-} & N &  \\
C2693 & 1a<o0L / 2a*<1a1 / 3a<2a1 3b<2a1 & N/N/NN & +/{-}/{-}{-} & N &  \\
C2694 & 1a<o0L / 2a*<1a1 / 3a<2a1 3b<2a2 & N/N/NN & +/{-}/{-}{-} & N &  \\
C2695 & 1a<o0L / 2a*<1a1 / 3a<2a2 3b<2a2 & N/N/NN & +/{-}/{-}{-} & N &  \\
C2696 & 1a<o0L / 2a*<1a2 / 3a<2a1 3b<2a1 & N/N/NN & +/{-}/{-}{-} & N &  \\
C2697 & 1a<o0L / 2a*<1a2 / 3a<2a1 3b<2a2 & N/N/NN & +/{-}/{-}{-} & N &  \\
C2698 & 1a<o0L / 2a*<1a2 / 3a<2a2 3b<2a2 & N/N/NN & +/{-}/{-}{-} & N &  \\
C2699 & 1a<o0L 1b<o1 / 2a<1a0L / 3a*<2a1 & NI/N/N & +{-}/+/{-} & N &  \\
C2700 & 1a<o0L 1b<o1 / 2a<1a0L / 3a*<2a2 & NI/N/N & +{-}/+/{-} & N &  \\
C2701 & 1a<o0L 1b<o1 / 2a<1a1 / 3a*<2a1 & NI/N/N & +{-}/{-}/{-} & N &  \\
C2702 & 1a<o0L 1b<o1 / 2a<1a1 / 3a*<2a2 & NI/N/N & +{-}/{-}/{-} & N &  \\
C2703 & 1a<o0L 1b<o1 / 2a<1a1L / 3a*<2a1 & NI/N/N & +{-}/+/{-} & N &  \\
C2704 & 1a<o0L 1b<o1 / 2a<1a1L / 3a*<2a2 & NI/N/N & +{-}/+/{-} & N &  \\
C2705 & 1a<o0L 1b<o1 / 2a<1a2 / 3a*<2a1 & NI/N/N & +{-}/{-}/{-} & N &  \\
C2706 & 1a<o0L 1b<o1 / 2a<1a2 / 3a*<2a2 & NI/N/N & +{-}/{-}/{-} & N &  \\
C2707 & 1a<o0L 1b<o1 / 2a<1a2L / 3a*<2a1 & NI/N/N & +{-}/+/{-} & N &  \\
C2708 & 1a<o0L 1b<o1 / 2a<1a2L / 3a*<2a2 & NI/N/N & +{-}/+/{-} & N &  \\
C2709 & 1a<o0L 1b<o1 / 2a<1b0L / 3a*<2a1 & NI/N/N & +{-}/{-}/{-} & N &  \\
C2710 & 1a<o0L 1b<o1 / 2a<1b0L / 3a*<2a2 & NI/N/N & +{-}/{-}/{-} & N &  \\
C2711 & 1a<o0L 1b<o1 / 2a<1b1 / 3a*<2a1 & NI/I/U & +{-}/{-}/{-} & N & W:3a \\
C2712 & 1a<o0L 1b<o1 / 2a<1b1 / 3a*<2a2 & NI/I/U & +{-}/{-}/{-} & N & W:3a \\
C2713 & 1a<o0L 1b<o1 / 2a<1b1L / 3a*<2a1 & NI/I/U & +{-}/{-}/{-} & N & W:3a \\
C2714 & 1a<o0L 1b<o1 / 2a<1b1L / 3a*<2a2 & NI/I/U & +{-}/{-}/{-} & N & W:3a \\
C2715 & 1a<o0L 1b<o1 / 2a<1b2 / 3a*<2a1 & NI/I/U & +{-}/{-}/{-} & N & W:3a \\
C2716 & 1a<o0L 1b<o1 / 2a<1b2 / 3a*<2a2 & NI/I/U & +{-}/{-}/{-} & N & W:3a \\
C2717 & 1a<o0L 1b<o1 / 2a<1b2L / 3a*<2a1 & NI/I/U & +{-}/{-}/{-} & N & W:3a \\
C2718 & 1a<o0L 1b<o1 / 2a<1b2L / 3a*<2a2 & NI/I/U & +{-}/{-}/{-} & N & W:3a \\
C2719 & 1a<o0L 1b<o1 / 2a*<1a1 / 3a<2a1 & NI/N/N & +{-}/{-}/{-} & N &  \\
C2720 & 1a<o0L 1b<o1 / 2a*<1a1 / 3a<2a2 & NI/N/N & +{-}/{-}/{-} & N &  \\
C2721 & 1a<o0L 1b<o1 / 2a*<1a2 / 3a<2a1 & NI/N/N & +{-}/{-}/{-} & N &  \\
C2722 & 1a<o0L 1b<o1 / 2a*<1a2 / 3a<2a2 & NI/N/N & +{-}/{-}/{-} & N &  \\
C2723 & 1a<o0L 1b<o1 / 2a*<1b1 / 3a<2a1 & NI/U/U & +{-}/{-}/{-} & N & W:2a,3a \\
C2724 & 1a<o0L 1b<o1 / 2a*<1b1 / 3a<2a2 & NI/U/U & +{-}/{-}/{-} & N & W:2a,3a \\
C2725 & 1a<o0L 1b<o1 / 2a*<1b2 / 3a<2a1 & NI/U/U & +{-}/{-}/{-} & N & W:2a,3a \\
C2726 & 1a<o0L 1b<o1 / 2a*<1b2 / 3a<2a2 & NI/U/U & +{-}/{-}/{-} & N & W:2a,3a \\
C2727 & 1a<o0L 1b<o2 / 2a<1a0L / 3a*<2a1 & NI/N/N & +{-}/+/{-} & N &  \\
C2728 & 1a<o0L 1b<o2 / 2a<1a0L / 3a*<2a2 & NI/N/N & +{-}/+/{-} & N &  \\
C2729 & 1a<o0L 1b<o2 / 2a<1a1 / 3a*<2a1 & NI/N/N & +{-}/{-}/{-} & N &  \\
C2730 & 1a<o0L 1b<o2 / 2a<1a1 / 3a*<2a2 & NI/N/N & +{-}/{-}/{-} & N &  \\
C2731 & 1a<o0L 1b<o2 / 2a<1a1L / 3a*<2a1 & NI/N/N & +{-}/+/{-} & N &  \\
C2732 & 1a<o0L 1b<o2 / 2a<1a1L / 3a*<2a2 & NI/N/N & +{-}/+/{-} & N &  \\
C2733 & 1a<o0L 1b<o2 / 2a<1a2 / 3a*<2a1 & NI/N/N & +{-}/{-}/{-} & N &  \\
C2734 & 1a<o0L 1b<o2 / 2a<1a2 / 3a*<2a2 & NI/N/N & +{-}/{-}/{-} & N &  \\
C2735 & 1a<o0L 1b<o2 / 2a<1a2L / 3a*<2a1 & NI/N/N & +{-}/+/{-} & N &  \\
C2736 & 1a<o0L 1b<o2 / 2a<1a2L / 3a*<2a2 & NI/N/N & +{-}/+/{-} & N &  \\
C2737 & 1a<o0L 1b<o2 / 2a<1b0L / 3a*<2a1 & NI/N/N & +{-}/{-}/{-} & N &  \\
C2738 & 1a<o0L 1b<o2 / 2a<1b0L / 3a*<2a2 & NI/N/N & +{-}/{-}/{-} & N &  \\
C2739 & 1a<o0L 1b<o2 / 2a<1b1 / 3a*<2a1 & NI/I/U & +{-}/{-}/{-} & N & W:3a \\
C2740 & 1a<o0L 1b<o2 / 2a<1b1 / 3a*<2a2 & NI/I/U & +{-}/{-}/{-} & N & W:3a \\
C2741 & 1a<o0L 1b<o2 / 2a<1b1L / 3a*<2a1 & NI/I/U & +{-}/{-}/{-} & N & W:3a \\
C2742 & 1a<o0L 1b<o2 / 2a<1b1L / 3a*<2a2 & NI/I/U & +{-}/{-}/{-} & N & W:3a \\
C2743 & 1a<o0L 1b<o2 / 2a<1b2 / 3a*<2a1 & NI/I/U & +{-}/{-}/{-} & N & W:3a \\
C2744 & 1a<o0L 1b<o2 / 2a<1b2 / 3a*<2a2 & NI/I/U & +{-}/{-}/{-} & N & W:3a \\
C2745 & 1a<o0L 1b<o2 / 2a<1b2L / 3a*<2a1 & NI/I/U & +{-}/{-}/{-} & N & W:3a \\
C2746 & 1a<o0L 1b<o2 / 2a<1b2L / 3a*<2a2 & NI/I/U & +{-}/{-}/{-} & N & W:3a \\
C2747 & 1a<o0L 1b<o2 / 2a*<1a1 / 3a<2a1 & NI/N/N & +{-}/{-}/{-} & N &  \\
C2748 & 1a<o0L 1b<o2 / 2a*<1a1 / 3a<2a2 & NI/N/N & +{-}/{-}/{-} & N &  \\
C2749 & 1a<o0L 1b<o2 / 2a*<1a2 / 3a<2a1 & NI/N/N & +{-}/{-}/{-} & N &  \\
C2750 & 1a<o0L 1b<o2 / 2a*<1a2 / 3a<2a2 & NI/N/N & +{-}/{-}/{-} & N &  \\
C2751 & 1a<o0L 1b<o2 / 2a*<1b1 / 3a<2a1 & NI/U/U & +{-}/{-}/{-} & N & W:2a,3a \\
C2752 & 1a<o0L 1b<o2 / 2a*<1b1 / 3a<2a2 & NI/U/U & +{-}/{-}/{-} & N & W:2a,3a \\
C2753 & 1a<o0L 1b<o2 / 2a*<1b2 / 3a<2a1 & NI/U/U & +{-}/{-}/{-} & N & W:2a,3a \\
C2754 & 1a<o0L 1b<o2 / 2a*<1b2 / 3a<2a2 & NI/U/U & +{-}/{-}/{-} & N & W:2a,3a \\
C2755 & 1a<o0L 1b*<o1 / 2a<1a0L / 3a<2a0L & NU/N/N & +{-}/+/+ & N & W:1b \\
C2756 & 1a<o0L 1b*<o1 / 2a<1a0L / 3a<2a1 & NU/N/N & +{-}/+/{-} & N & W:1b \\
C2757 & 1a<o0L 1b*<o1 / 2a<1a0L / 3a<2a1L & NU/N/N & +{-}/+/+ & N & W:1b \\
C2758 & 1a<o0L 1b*<o1 / 2a<1a0L / 3a<2a2 & NU/N/N & +{-}/+/{-} & N & W:1b \\
C2759 & 1a<o0L 1b*<o1 / 2a<1a0L / 3a<2a2L & NU/N/N & +{-}/+/+ & N & W:1b \\
C2760 & 1a<o0L 1b*<o1 / 2a<1a1 / 3a<2a0L & NU/N/N & +{-}/{-}/{-} & N & W:1b \\
C2761 & 1a<o0L 1b*<o1 / 2a<1a1 / 3a<2a1 & NU/N/N & +{-}/{-}/{-} & N & W:1b \\
C2762 & 1a<o0L 1b*<o1 / 2a<1a1 / 3a<2a1L & NU/N/N & +{-}/{-}/{-} & N & W:1b \\
C2763 & 1a<o0L 1b*<o1 / 2a<1a1 / 3a<2a2 & NU/N/N & +{-}/{-}/{-} & N & W:1b \\
C2764 & 1a<o0L 1b*<o1 / 2a<1a1 / 3a<2a2L & NU/N/N & +{-}/{-}/{-} & N & W:1b \\
C2765 & 1a<o0L 1b*<o1 / 2a<1a1L / 3a<2a0L & NU/N/N & +{-}/+/+ & N & W:1b \\
C2766 & 1a<o0L 1b*<o1 / 2a<1a1L / 3a<2a1 & NU/N/N & +{-}/+/{-} & N & W:1b \\
C2767 & 1a<o0L 1b*<o1 / 2a<1a1L / 3a<2a1L & NU/N/N & +{-}/+/+ & N & W:1b \\
C2768 & 1a<o0L 1b*<o1 / 2a<1a1L / 3a<2a2 & NU/N/N & +{-}/+/{-} & N & W:1b \\
C2769 & 1a<o0L 1b*<o1 / 2a<1a1L / 3a<2a2L & NU/N/N & +{-}/+/+ & N & W:1b \\
C2770 & 1a<o0L 1b*<o1 / 2a<1a2 / 3a<2a0L & NU/N/N & +{-}/{-}/{-} & N & W:1b \\
C2771 & 1a<o0L 1b*<o1 / 2a<1a2 / 3a<2a1 & NU/N/N & +{-}/{-}/{-} & N & W:1b \\
C2772 & 1a<o0L 1b*<o1 / 2a<1a2 / 3a<2a1L & NU/N/N & +{-}/{-}/{-} & N & W:1b \\
C2773 & 1a<o0L 1b*<o1 / 2a<1a2 / 3a<2a2 & NU/N/N & +{-}/{-}/{-} & N & W:1b \\
C2774 & 1a<o0L 1b*<o1 / 2a<1a2 / 3a<2a2L & NU/N/N & +{-}/{-}/{-} & N & W:1b \\
C2775 & 1a<o0L 1b*<o1 / 2a<1a2L / 3a<2a0L & NU/N/N & +{-}/+/+ & N & W:1b \\
C2776 & 1a<o0L 1b*<o1 / 2a<1a2L / 3a<2a1 & NU/N/N & +{-}/+/{-} & N & W:1b \\
C2777 & 1a<o0L 1b*<o1 / 2a<1a2L / 3a<2a1L & NU/N/N & +{-}/+/+ & N & W:1b \\
C2778 & 1a<o0L 1b*<o1 / 2a<1a2L / 3a<2a2 & NU/N/N & +{-}/+/{-} & N & W:1b \\
C2779 & 1a<o0L 1b*<o1 / 2a<1a2L / 3a<2a2L & NU/N/N & +{-}/+/+ & N & W:1b \\
C2780 & 1a<o0L 1b*<o1 / 2a<1b1 / 3a<2a0L & NU/U/N & +{-}/{-}/{-} & N & W:1b,2a \\
C2781 & 1a<o0L 1b*<o1 / 2a<1b1 / 3a<2a1 & NU/U/U & +{-}/{-}/{-} & N & W:1b,2a,3a \\
C2782 & 1a<o0L 1b*<o1 / 2a<1b1 / 3a<2a1L & NU/U/U & +{-}/{-}/{-} & N & W:1b,2a,3a \\
C2783 & 1a<o0L 1b*<o1 / 2a<1b1 / 3a<2a2 & NU/U/U & +{-}/{-}/{-} & N & W:1b,2a,3a \\
C2784 & 1a<o0L 1b*<o1 / 2a<1b1 / 3a<2a2L & NU/U/U & +{-}/{-}/{-} & N & W:1b,2a,3a \\
C2785 & 1a<o0L 1b*<o1 / 2a<1b2 / 3a<2a0L & NU/U/N & +{-}/{-}/{-} & N & W:1b,2a \\
C2786 & 1a<o0L 1b*<o1 / 2a<1b2 / 3a<2a1 & NU/U/U & +{-}/{-}/{-} & N & W:1b,2a,3a \\
C2787 & 1a<o0L 1b*<o1 / 2a<1b2 / 3a<2a1L & NU/U/U & +{-}/{-}/{-} & N & W:1b,2a,3a \\
C2788 & 1a<o0L 1b*<o1 / 2a<1b2 / 3a<2a2 & NU/U/U & +{-}/{-}/{-} & N & W:1b,2a,3a \\
C2789 & 1a<o0L 1b*<o1 / 2a<1b2 / 3a<2a2L & NU/U/U & +{-}/{-}/{-} & N & W:1b,2a,3a \\
C2790 & 1a<o0L 1b*<o2 / 2a<1a0L / 3a<2a0L & NU/N/N & +{-}/+/+ & N & W:1b \\
C2791 & 1a<o0L 1b*<o2 / 2a<1a0L / 3a<2a1 & NU/N/N & +{-}/+/{-} & N & W:1b \\
C2792 & 1a<o0L 1b*<o2 / 2a<1a0L / 3a<2a1L & NU/N/N & +{-}/+/+ & N & W:1b \\
C2793 & 1a<o0L 1b*<o2 / 2a<1a0L / 3a<2a2 & NU/N/N & +{-}/+/{-} & N & W:1b \\
C2794 & 1a<o0L 1b*<o2 / 2a<1a0L / 3a<2a2L & NU/N/N & +{-}/+/+ & N & W:1b \\
C2795 & 1a<o0L 1b*<o2 / 2a<1a1 / 3a<2a0L & NU/N/N & +{-}/{-}/{-} & N & W:1b \\
C2796 & 1a<o0L 1b*<o2 / 2a<1a1 / 3a<2a1 & NU/N/N & +{-}/{-}/{-} & N & W:1b \\
C2797 & 1a<o0L 1b*<o2 / 2a<1a1 / 3a<2a1L & NU/N/N & +{-}/{-}/{-} & N & W:1b \\
C2798 & 1a<o0L 1b*<o2 / 2a<1a1 / 3a<2a2 & NU/N/N & +{-}/{-}/{-} & N & W:1b \\
C2799 & 1a<o0L 1b*<o2 / 2a<1a1 / 3a<2a2L & NU/N/N & +{-}/{-}/{-} & N & W:1b \\
C2800 & 1a<o0L 1b*<o2 / 2a<1a1L / 3a<2a0L & NU/N/N & +{-}/+/+ & N & W:1b \\
C2801 & 1a<o0L 1b*<o2 / 2a<1a1L / 3a<2a1 & NU/N/N & +{-}/+/{-} & N & W:1b \\
C2802 & 1a<o0L 1b*<o2 / 2a<1a1L / 3a<2a1L & NU/N/N & +{-}/+/+ & N & W:1b \\
C2803 & 1a<o0L 1b*<o2 / 2a<1a1L / 3a<2a2 & NU/N/N & +{-}/+/{-} & N & W:1b \\
C2804 & 1a<o0L 1b*<o2 / 2a<1a1L / 3a<2a2L & NU/N/N & +{-}/+/+ & N & W:1b \\
C2805 & 1a<o0L 1b*<o2 / 2a<1a2 / 3a<2a0L & NU/N/N & +{-}/{-}/{-} & N & W:1b \\
C2806 & 1a<o0L 1b*<o2 / 2a<1a2 / 3a<2a1 & NU/N/N & +{-}/{-}/{-} & N & W:1b \\
C2807 & 1a<o0L 1b*<o2 / 2a<1a2 / 3a<2a1L & NU/N/N & +{-}/{-}/{-} & N & W:1b \\
C2808 & 1a<o0L 1b*<o2 / 2a<1a2 / 3a<2a2 & NU/N/N & +{-}/{-}/{-} & N & W:1b \\
C2809 & 1a<o0L 1b*<o2 / 2a<1a2 / 3a<2a2L & NU/N/N & +{-}/{-}/{-} & N & W:1b \\
C2810 & 1a<o0L 1b*<o2 / 2a<1a2L / 3a<2a0L & NU/N/N & +{-}/+/+ & N & W:1b \\
C2811 & 1a<o0L 1b*<o2 / 2a<1a2L / 3a<2a1 & NU/N/N & +{-}/+/{-} & N & W:1b \\
C2812 & 1a<o0L 1b*<o2 / 2a<1a2L / 3a<2a1L & NU/N/N & +{-}/+/+ & N & W:1b \\
C2813 & 1a<o0L 1b*<o2 / 2a<1a2L / 3a<2a2 & NU/N/N & +{-}/+/{-} & N & W:1b \\
C2814 & 1a<o0L 1b*<o2 / 2a<1a2L / 3a<2a2L & NU/N/N & +{-}/+/+ & N & W:1b \\
C2815 & 1a<o0L 1b*<o2 / 2a<1b1 / 3a<2a0L & NU/U/N & +{-}/{-}/{-} & N & W:1b,2a \\
C2816 & 1a<o0L 1b*<o2 / 2a<1b1 / 3a<2a1 & NU/U/U & +{-}/{-}/{-} & N & W:1b,2a,3a \\
C2817 & 1a<o0L 1b*<o2 / 2a<1b1 / 3a<2a1L & NU/U/U & +{-}/{-}/{-} & N & W:1b,2a,3a \\
C2818 & 1a<o0L 1b*<o2 / 2a<1b1 / 3a<2a2 & NU/U/U & +{-}/{-}/{-} & N & W:1b,2a,3a \\
C2819 & 1a<o0L 1b*<o2 / 2a<1b1 / 3a<2a2L & NU/U/U & +{-}/{-}/{-} & N & W:1b,2a,3a \\
C2820 & 1a<o0L 1b*<o2 / 2a<1b2 / 3a<2a0L & NU/U/N & +{-}/{-}/{-} & N & W:1b,2a \\
C2821 & 1a<o0L 1b*<o2 / 2a<1b2 / 3a<2a1 & NU/U/U & +{-}/{-}/{-} & N & W:1b,2a,3a \\
C2822 & 1a<o0L 1b*<o2 / 2a<1b2 / 3a<2a1L & NU/U/U & +{-}/{-}/{-} & N & W:1b,2a,3a \\
C2823 & 1a<o0L 1b*<o2 / 2a<1b2 / 3a<2a2 & NU/U/U & +{-}/{-}/{-} & N & W:1b,2a,3a \\
C2824 & 1a<o0L 1b*<o2 / 2a<1b2 / 3a<2a2L & NU/U/U & +{-}/{-}/{-} & N & W:1b,2a,3a \\
C2825 & 1a<o1 / 2a<1a0L / 3a<2a0L 3b*<2a1 & I/N/NN & {-}/{-}/{-}{-} & N &  \\
C2826 & 1a<o1 / 2a<1a0L / 3a<2a0L 3b*<2a2 & I/N/NN & {-}/{-}/{-}{-} & N &  \\
C2827 & 1a<o1 / 2a<1a0L / 3a<2a1 3b*<2a1 & I/N/NN & {-}/{-}/{-}{-} & N &  \\
C2828 & 1a<o1 / 2a<1a0L / 3a<2a1 3b*<2a2 & I/N/NN & {-}/{-}/{-}{-} & N &  \\
C2829 & 1a<o1 / 2a<1a0L / 3a<2a1L 3b*<2a1 & I/N/NN & {-}/{-}/{-}{-} & N &  \\
C2830 & 1a<o1 / 2a<1a0L / 3a<2a1L 3b*<2a2 & I/N/NN & {-}/{-}/{-}{-} & N &  \\
C2831 & 1a<o1 / 2a<1a0L / 3a<2a2 3b*<2a1 & I/N/NN & {-}/{-}/{-}{-} & N &  \\
C2832 & 1a<o1 / 2a<1a0L / 3a<2a2 3b*<2a2 & I/N/NN & {-}/{-}/{-}{-} & N &  \\
C2833 & 1a<o1 / 2a<1a0L / 3a<2a2L 3b*<2a1 & I/N/NN & {-}/{-}/{-}{-} & N &  \\
C2834 & 1a<o1 / 2a<1a0L / 3a<2a2L 3b*<2a2 & I/N/NN & {-}/{-}/{-}{-} & N &  \\
C2835 & 1a<o1 / 2a<1a0L 2b<1a1 / 3a*<2a1 & I/NI/N & {-}/{-}{-}/{-} & N &  \\
C2836 & 1a<o1 / 2a<1a0L 2b<1a1 / 3a*<2a2 & I/NI/N & {-}/{-}{-}/{-} & N &  \\
C2837 & 1a<o1 / 2a<1a0L 2b<1a1 / 3a*<2b1 & I/NI/U & {-}/{-}{-}/{-} & N & W:3a \\
C2838 & 1a<o1 / 2a<1a0L 2b<1a1 / 3a*<2b2 & I/NI/U & {-}/{-}{-}/{-} & N & W:3a \\
C2839 & 1a<o1 / 2a<1a0L 2b<1a2 / 3a*<2a1 & I/NI/N & {-}/{-}{-}/{-} & N &  \\
C2840 & 1a<o1 / 2a<1a0L 2b<1a2 / 3a*<2a2 & I/NI/N & {-}/{-}{-}/{-} & N &  \\
C2841 & 1a<o1 / 2a<1a0L 2b<1a2 / 3a*<2b1 & I/NI/U & {-}/{-}{-}/{-} & N & W:3a \\
C2842 & 1a<o1 / 2a<1a0L 2b<1a2 / 3a*<2b2 & I/NI/U & {-}/{-}{-}/{-} & N & W:3a \\
C2843 & 1a<o1 / 2a<1a0L 2b*<1a1 / 3a<2a0L & I/NU/N & {-}/{-}{-}/{-} & N & W:2b \\
C2844 & 1a<o1 / 2a<1a0L 2b*<1a1 / 3a<2a1 & I/NU/N & {-}/{-}{-}/{-} & N & W:2b \\
C2845 & 1a<o1 / 2a<1a0L 2b*<1a1 / 3a<2a1L & I/NU/N & {-}/{-}{-}/{-} & N & W:2b \\
C2846 & 1a<o1 / 2a<1a0L 2b*<1a1 / 3a<2a2 & I/NU/N & {-}/{-}{-}/{-} & N & W:2b \\
C2847 & 1a<o1 / 2a<1a0L 2b*<1a1 / 3a<2a2L & I/NU/N & {-}/{-}{-}/{-} & N & W:2b \\
C2848 & 1a<o1 / 2a<1a0L 2b*<1a1 / 3a<2b1 & I/NU/U & {-}/{-}{-}/{-} & N & W:2b,3a \\
C2849 & 1a<o1 / 2a<1a0L 2b*<1a1 / 3a<2b2 & I/NU/U & {-}/{-}{-}/{-} & N & W:2b,3a \\
C2850 & 1a<o1 / 2a<1a0L 2b*<1a2 / 3a<2a0L & I/NU/N & {-}/{-}{-}/{-} & N & W:2b \\
C2851 & 1a<o1 / 2a<1a0L 2b*<1a2 / 3a<2a1 & I/NU/N & {-}/{-}{-}/{-} & N & W:2b \\
C2852 & 1a<o1 / 2a<1a0L 2b*<1a2 / 3a<2a1L & I/NU/N & {-}/{-}{-}/{-} & N & W:2b \\
C2853 & 1a<o1 / 2a<1a0L 2b*<1a2 / 3a<2a2 & I/NU/N & {-}/{-}{-}/{-} & N & W:2b \\
C2854 & 1a<o1 / 2a<1a0L 2b*<1a2 / 3a<2a2L & I/NU/N & {-}/{-}{-}/{-} & N & W:2b \\
C2855 & 1a<o1 / 2a<1a0L 2b*<1a2 / 3a<2b1 & I/NU/U & {-}/{-}{-}/{-} & N & W:2b,3a \\
C2856 & 1a<o1 / 2a<1a0L 2b*<1a2 / 3a<2b2 & I/NU/U & {-}/{-}{-}/{-} & N & W:2b,3a \\
C2857 & 1a<o1 / 2a<1a1 / 3a<2a0L 3b*<2a1 & I/I/NU & {-}/{-}/{-}{-} & N & W:3b \\
C2858 & 1a<o1 / 2a<1a1 / 3a<2a0L 3b*<2a2 & I/I/NU & {-}/{-}/{-}{-} & N & W:3b \\
C2859 & 1a<o1 / 2a<1a1 / 3a<2a1 3b*<2a1 & I/I/IU & {-}/{-}/{-}{-} & N & W:1a,2a,3a \\
C2860 & 1a<o1 / 2a<1a1 / 3a<2a1 3b*<2a2 & I/I/IU & {-}/{-}/{-}{-} & N & W:1a,2a,3a \\
C2861 & 1a<o1 / 2a<1a1 / 3a<2a1L 3b*<2a1 & I/I/IU & {-}/{-}/{-}{-} & N & W:1a,2a,3a \\
C2862 & 1a<o1 / 2a<1a1 / 3a<2a1L 3b*<2a2 & I/I/IU & {-}/{-}/{-}{-} & N & W:1a,2a,3a \\
C2863 & 1a<o1 / 2a<1a1 / 3a<2a2 3b*<2a1 & I/I/IU & {-}/{-}/{-}{-} & N & W:1a,2a,3a \\
C2864 & 1a<o1 / 2a<1a1 / 3a<2a2 3b*<2a2 & I/I/IU & {-}/{-}/{-}{-} & N & W:1a,2a,3a \\
C2865 & 1a<o1 / 2a<1a1 / 3a<2a2L 3b*<2a1 & I/I/IU & {-}/{-}/{-}{-} & N & W:1a,2a,3a \\
C2866 & 1a<o1 / 2a<1a1 / 3a<2a2L 3b*<2a2 & I/I/IU & {-}/{-}/{-}{-} & N & W:1a,2a,3a \\
C2867 & 1a<o1 / 2a<1a1 2b<1a1 / 3a*<2a1 & I/EE/U & {-}/{-}{-}/{-} & E & W:1a \\
C2868 & 1a<o1 / 2a<1a1 2b<1a1 / 3a*<2a2 & I/EE/U & {-}/{-}{-}/{-} & E & W:1a \\
C2869 & 1a<o1 / 2a<1a1 2b<1a1L / 3a*<2a1 & I/LI/U & {-}/{-}{-}/{-} & B & CP:2b W:1a,2b \\
C2870 & 1a<o1 / 2a<1a1 2b<1a1L / 3a*<2a2 & I/LI/U & {-}/{-}{-}/{-} & B & CP:2b W:1a,2b \\
C2871 & 1a<o1 / 2a<1a1 2b<1a1L / 3a*<2b1 & I/LI/U & {-}/{-}{-}/{-} & B & CP:2b W:1a,2b \\
C2872 & 1a<o1 / 2a<1a1 2b<1a1L / 3a*<2b2 & I/LI/U & {-}/{-}{-}/{-} & B & CP:2b W:1a,2b \\
C2873 & 1a<o1 / 2a<1a1 2b<1a2 / 3a*<2a1 & I/EE/U & {-}/{-}{-}/{-} & E & C:2b W:1a \\
C2874 & 1a<o1 / 2a<1a1 2b<1a2 / 3a*<2a2 & I/EE/U & {-}/{-}{-}/{-} & E & C:2b W:1a \\
C2875 & 1a<o1 / 2a<1a1 2b<1a2 / 3a*<2b1 & I/EE/U & {-}/{-}{-}/{-} & E & C:2b W:1a \\
C2876 & 1a<o1 / 2a<1a1 2b<1a2 / 3a*<2b2 & I/EE/U & {-}/{-}{-}/{-} & E & C:2b W:1a \\
C2877 & 1a<o1 / 2a<1a1 2b<1a2L / 3a*<2a1 & I/LI/U & {-}/{-}{-}/{-} & B & P:2b W:1a,2b \\
C2878 & 1a<o1 / 2a<1a1 2b<1a2L / 3a*<2a2 & I/LI/U & {-}/{-}{-}/{-} & B & P:2b W:1a,2b \\
C2879 & 1a<o1 / 2a<1a1 2b<1a2L / 3a*<2b1 & I/LI/U & {-}/{-}{-}/{-} & B & P:2b W:1a,2b \\
C2880 & 1a<o1 / 2a<1a1 2b<1a2L / 3a*<2b2 & I/LI/U & {-}/{-}{-}/{-} & B & P:2b W:1a,2b \\
C2881 & 1a<o1 / 2a<1a1 2b*<1a1 / 3a<2a0L & I/IU/N & {-}/{-}{-}/{-} & N & W:1a,2a \\
C2882 & 1a<o1 / 2a<1a1 2b*<1a1 / 3a<2a1 & I/IU/I & {-}/{-}{-}/{-} & N & W:1a,2a,3a \\
C2883 & 1a<o1 / 2a<1a1 2b*<1a1 / 3a<2a1L & I/IU/I & {-}/{-}{-}/{-} & N & W:1a,2a,3a \\
C2884 & 1a<o1 / 2a<1a1 2b*<1a1 / 3a<2a2 & I/IU/I & {-}/{-}{-}/{-} & N & W:1a,2a,3a \\
C2885 & 1a<o1 / 2a<1a1 2b*<1a1 / 3a<2a2L & I/IU/I & {-}/{-}{-}/{-} & N & W:1a,2a,3a \\
C2886 & 1a<o1 / 2a<1a1 2b*<1a1 / 3a<2b1 & I/IU/U & {-}/{-}{-}/{-} & N & W:1a,2a \\
C2887 & 1a<o1 / 2a<1a1 2b*<1a1 / 3a<2b2 & I/IU/U & {-}/{-}{-}/{-} & N & W:1a,2a \\
C2888 & 1a<o1 / 2a<1a1 2b*<1a2 / 3a<2a0L & I/IU/N & {-}/{-}{-}/{-} & N & W:1a,2a \\
C2889 & 1a<o1 / 2a<1a1 2b*<1a2 / 3a<2a1 & I/IU/I & {-}/{-}{-}/{-} & N & W:1a,2a,3a \\
C2890 & 1a<o1 / 2a<1a1 2b*<1a2 / 3a<2a1L & I/IU/I & {-}/{-}{-}/{-} & N & W:1a,2a,3a \\
C2891 & 1a<o1 / 2a<1a1 2b*<1a2 / 3a<2a2 & I/IU/I & {-}/{-}{-}/{-} & N & W:1a,2a,3a \\
C2892 & 1a<o1 / 2a<1a1 2b*<1a2 / 3a<2a2L & I/IU/I & {-}/{-}{-}/{-} & N & W:1a,2a,3a \\
C2893 & 1a<o1 / 2a<1a1 2b*<1a2 / 3a<2b1 & I/IU/U & {-}/{-}{-}/{-} & N & W:1a,2a \\
C2894 & 1a<o1 / 2a<1a1 2b*<1a2 / 3a<2b2 & I/IU/U & {-}/{-}{-}/{-} & N & W:1a,2a \\
C2895 & 1a<o1 / 2a<1a1L / 3a<2a0L 3b*<2a1 & I/I/NU & {-}/{-}/{-}{-} & N & W:3b \\
C2896 & 1a<o1 / 2a<1a1L / 3a<2a0L 3b*<2a2 & I/I/NU & {-}/{-}/{-}{-} & N & W:3b \\
C2897 & 1a<o1 / 2a<1a1L / 3a<2a1 3b*<2a1 & I/I/IU & {-}/{-}/{-}{-} & N & W:1a,2a,3a \\
C2898 & 1a<o1 / 2a<1a1L / 3a<2a1 3b*<2a2 & I/I/IU & {-}/{-}/{-}{-} & N & W:1a,2a,3a \\
C2899 & 1a<o1 / 2a<1a1L / 3a<2a1L 3b*<2a1 & I/I/IU & {-}/{-}/{-}{-} & N & W:1a,2a,3a \\
C2900 & 1a<o1 / 2a<1a1L / 3a<2a1L 3b*<2a2 & I/I/IU & {-}/{-}/{-}{-} & N & W:1a,2a,3a \\
C2901 & 1a<o1 / 2a<1a1L / 3a<2a2 3b*<2a1 & I/I/IU & {-}/{-}/{-}{-} & N & W:1a,2a,3a \\
C2902 & 1a<o1 / 2a<1a1L / 3a<2a2 3b*<2a2 & I/I/IU & {-}/{-}/{-}{-} & N & W:1a,2a,3a \\
C2903 & 1a<o1 / 2a<1a1L / 3a<2a2L 3b*<2a1 & I/I/IU & {-}/{-}/{-}{-} & N & W:1a,2a,3a \\
C2904 & 1a<o1 / 2a<1a1L / 3a<2a2L 3b*<2a2 & I/I/IU & {-}/{-}/{-}{-} & N & W:1a,2a,3a \\
C2905 & 1a<o1 / 2a<1a1L 2b<1a2 / 3a*<2a1 & I/IL/U & {-}/{-}{-}/{-} & B & C:2a,2b P:2a W:1a,2a \\
C2906 & 1a<o1 / 2a<1a1L 2b<1a2 / 3a*<2a2 & I/IL/U & {-}/{-}{-}/{-} & B & C:2a,2b P:2a W:1a,2a \\
C2907 & 1a<o1 / 2a<1a1L 2b<1a2 / 3a*<2b1 & I/IL/U & {-}/{-}{-}/{-} & B & C:2a,2b P:2a W:1a,2a \\
C2908 & 1a<o1 / 2a<1a1L 2b<1a2 / 3a*<2b2 & I/IL/U & {-}/{-}{-}/{-} & B & C:2a,2b P:2a W:1a,2a \\
C2909 & 1a<o1 / 2a<1a1L 2b*<1a1 / 3a<2a0L & I/IU/N & {-}/{-}{-}/{-} & N & W:1a,2a \\
C2910 & 1a<o1 / 2a<1a1L 2b*<1a1 / 3a<2a1 & I/IU/I & {-}/{-}{-}/{-} & N & W:1a,2a,3a \\
C2911 & 1a<o1 / 2a<1a1L 2b*<1a1 / 3a<2a1L & I/IU/I & {-}/{-}{-}/{-} & N & W:1a,2a,3a \\
C2912 & 1a<o1 / 2a<1a1L 2b*<1a1 / 3a<2a2 & I/IU/I & {-}/{-}{-}/{-} & N & W:1a,2a,3a \\
C2913 & 1a<o1 / 2a<1a1L 2b*<1a1 / 3a<2a2L & I/IU/I & {-}/{-}{-}/{-} & N & W:1a,2a,3a \\
C2914 & 1a<o1 / 2a<1a1L 2b*<1a1 / 3a<2b1 & I/IU/U & {-}/{-}{-}/{-} & N & W:1a,2a \\
C2915 & 1a<o1 / 2a<1a1L 2b*<1a1 / 3a<2b2 & I/IU/U & {-}/{-}{-}/{-} & N & W:1a,2a \\
C2916 & 1a<o1 / 2a<1a1L 2b*<1a2 / 3a<2a0L & I/IU/N & {-}/{-}{-}/{-} & N & W:1a,2a \\
C2917 & 1a<o1 / 2a<1a1L 2b*<1a2 / 3a<2a1 & I/IU/I & {-}/{-}{-}/{-} & N & W:1a,2a,3a \\
C2918 & 1a<o1 / 2a<1a1L 2b*<1a2 / 3a<2a1L & I/IU/I & {-}/{-}{-}/{-} & N & W:1a,2a,3a \\
C2919 & 1a<o1 / 2a<1a1L 2b*<1a2 / 3a<2a2 & I/IU/I & {-}/{-}{-}/{-} & N & W:1a,2a,3a \\
C2920 & 1a<o1 / 2a<1a1L 2b*<1a2 / 3a<2a2L & I/IU/I & {-}/{-}{-}/{-} & N & W:1a,2a,3a \\
C2921 & 1a<o1 / 2a<1a1L 2b*<1a2 / 3a<2b1 & I/IU/U & {-}/{-}{-}/{-} & N & W:1a,2a \\
C2922 & 1a<o1 / 2a<1a1L 2b*<1a2 / 3a<2b2 & I/IU/U & {-}/{-}{-}/{-} & N & W:1a,2a \\
C2923 & 1a<o1 / 2a<1a2 / 3a<2a0L 3b*<2a1 & I/I/NU & {-}/{-}/{-}{-} & N & W:3b \\
C2924 & 1a<o1 / 2a<1a2 / 3a<2a0L 3b*<2a2 & I/I/NU & {-}/{-}/{-}{-} & N & W:3b \\
C2925 & 1a<o1 / 2a<1a2 / 3a<2a1 3b*<2a1 & I/I/IU & {-}/{-}/{-}{-} & N & W:1a,2a,3a \\
C2926 & 1a<o1 / 2a<1a2 / 3a<2a1 3b*<2a2 & I/I/IU & {-}/{-}/{-}{-} & N & W:1a,2a,3a \\
C2927 & 1a<o1 / 2a<1a2 / 3a<2a1L 3b*<2a1 & I/I/IU & {-}/{-}/{-}{-} & N & W:1a,2a,3a \\
C2928 & 1a<o1 / 2a<1a2 / 3a<2a1L 3b*<2a2 & I/I/IU & {-}/{-}/{-}{-} & N & W:1a,2a,3a \\
C2929 & 1a<o1 / 2a<1a2 / 3a<2a2 3b*<2a1 & I/I/IU & {-}/{-}/{-}{-} & N & W:1a,2a,3a \\
C2930 & 1a<o1 / 2a<1a2 / 3a<2a2 3b*<2a2 & I/I/IU & {-}/{-}/{-}{-} & N & W:1a,2a,3a \\
C2931 & 1a<o1 / 2a<1a2 / 3a<2a2L 3b*<2a1 & I/I/IU & {-}/{-}/{-}{-} & N & W:1a,2a,3a \\
C2932 & 1a<o1 / 2a<1a2 / 3a<2a2L 3b*<2a2 & I/I/IU & {-}/{-}/{-}{-} & N & W:1a,2a,3a \\
C2933 & 1a<o1 / 2a<1a2 2b<1a2 / 3a*<2a1 & I/EE/U & {-}/{-}{-}/{-} & E & W:1a \\
C2934 & 1a<o1 / 2a<1a2 2b<1a2 / 3a*<2a2 & I/EE/U & {-}/{-}{-}/{-} & E & W:1a \\
C2935 & 1a<o1 / 2a<1a2 2b<1a2L / 3a*<2a1 & I/LI/U & {-}/{-}{-}/{-} & B & CP:2b W:1a,2b \\
C2936 & 1a<o1 / 2a<1a2 2b<1a2L / 3a*<2a2 & I/LI/U & {-}/{-}{-}/{-} & B & CP:2b W:1a,2b \\
C2937 & 1a<o1 / 2a<1a2 2b<1a2L / 3a*<2b1 & I/LI/U & {-}/{-}{-}/{-} & B & CP:2b W:1a,2b \\
C2938 & 1a<o1 / 2a<1a2 2b<1a2L / 3a*<2b2 & I/LI/U & {-}/{-}{-}/{-} & B & CP:2b W:1a,2b \\
C2939 & 1a<o1 / 2a<1a2 2b*<1a1 / 3a<2a0L & I/IU/N & {-}/{-}{-}/{-} & N & W:1a,2a \\
C2940 & 1a<o1 / 2a<1a2 2b*<1a1 / 3a<2a1 & I/IU/I & {-}/{-}{-}/{-} & N & W:1a,2a,3a \\
C2941 & 1a<o1 / 2a<1a2 2b*<1a1 / 3a<2a1L & I/IU/I & {-}/{-}{-}/{-} & N & W:1a,2a,3a \\
C2942 & 1a<o1 / 2a<1a2 2b*<1a1 / 3a<2a2 & I/IU/I & {-}/{-}{-}/{-} & N & W:1a,2a,3a \\
C2943 & 1a<o1 / 2a<1a2 2b*<1a1 / 3a<2a2L & I/IU/I & {-}/{-}{-}/{-} & N & W:1a,2a,3a \\
C2944 & 1a<o1 / 2a<1a2 2b*<1a1 / 3a<2b1 & I/IU/U & {-}/{-}{-}/{-} & N & W:1a,2a \\
C2945 & 1a<o1 / 2a<1a2 2b*<1a1 / 3a<2b2 & I/IU/U & {-}/{-}{-}/{-} & N & W:1a,2a \\
C2946 & 1a<o1 / 2a<1a2 2b*<1a2 / 3a<2a0L & I/IU/N & {-}/{-}{-}/{-} & N & W:1a,2a \\
C2947 & 1a<o1 / 2a<1a2 2b*<1a2 / 3a<2a1 & I/IU/I & {-}/{-}{-}/{-} & N & W:1a,2a,3a \\
C2948 & 1a<o1 / 2a<1a2 2b*<1a2 / 3a<2a1L & I/IU/I & {-}/{-}{-}/{-} & N & W:1a,2a,3a \\
C2949 & 1a<o1 / 2a<1a2 2b*<1a2 / 3a<2a2 & I/IU/I & {-}/{-}{-}/{-} & N & W:1a,2a,3a \\
C2950 & 1a<o1 / 2a<1a2 2b*<1a2 / 3a<2a2L & I/IU/I & {-}/{-}{-}/{-} & N & W:1a,2a,3a \\
C2951 & 1a<o1 / 2a<1a2 2b*<1a2 / 3a<2b1 & I/IU/U & {-}/{-}{-}/{-} & N & W:1a,2a \\
C2952 & 1a<o1 / 2a<1a2 2b*<1a2 / 3a<2b2 & I/IU/U & {-}/{-}{-}/{-} & N & W:1a,2a \\
C2953 & 1a<o1 / 2a<1a2L / 3a<2a0L 3b*<2a1 & I/I/NU & {-}/{-}/{-}{-} & N & W:3b \\
C2954 & 1a<o1 / 2a<1a2L / 3a<2a0L 3b*<2a2 & I/I/NU & {-}/{-}/{-}{-} & N & W:3b \\
C2955 & 1a<o1 / 2a<1a2L / 3a<2a1 3b*<2a1 & I/I/IU & {-}/{-}/{-}{-} & N & W:1a,2a,3a \\
C2956 & 1a<o1 / 2a<1a2L / 3a<2a1 3b*<2a2 & I/I/IU & {-}/{-}/{-}{-} & N & W:1a,2a,3a \\
C2957 & 1a<o1 / 2a<1a2L / 3a<2a1L 3b*<2a1 & I/I/IU & {-}/{-}/{-}{-} & N & W:1a,2a,3a \\
C2958 & 1a<o1 / 2a<1a2L / 3a<2a1L 3b*<2a2 & I/I/IU & {-}/{-}/{-}{-} & N & W:1a,2a,3a \\
C2959 & 1a<o1 / 2a<1a2L / 3a<2a2 3b*<2a1 & I/I/IU & {-}/{-}/{-}{-} & N & W:1a,2a,3a \\
C2960 & 1a<o1 / 2a<1a2L / 3a<2a2 3b*<2a2 & I/I/IU & {-}/{-}/{-}{-} & N & W:1a,2a,3a \\
C2961 & 1a<o1 / 2a<1a2L / 3a<2a2L 3b*<2a1 & I/I/IU & {-}/{-}/{-}{-} & N & W:1a,2a,3a \\
C2962 & 1a<o1 / 2a<1a2L / 3a<2a2L 3b*<2a2 & I/I/IU & {-}/{-}/{-}{-} & N & W:1a,2a,3a \\
C2963 & 1a<o1 / 2a<1a2L 2b*<1a1 / 3a<2a0L & I/IU/N & {-}/{-}{-}/{-} & N & W:1a,2a \\
C2964 & 1a<o1 / 2a<1a2L 2b*<1a1 / 3a<2a1 & I/IU/I & {-}/{-}{-}/{-} & N & W:1a,2a,3a \\
C2965 & 1a<o1 / 2a<1a2L 2b*<1a1 / 3a<2a1L & I/IU/I & {-}/{-}{-}/{-} & N & W:1a,2a,3a \\
C2966 & 1a<o1 / 2a<1a2L 2b*<1a1 / 3a<2a2 & I/IU/I & {-}/{-}{-}/{-} & N & W:1a,2a,3a \\
C2967 & 1a<o1 / 2a<1a2L 2b*<1a1 / 3a<2a2L & I/IU/I & {-}/{-}{-}/{-} & N & W:1a,2a,3a \\
C2968 & 1a<o1 / 2a<1a2L 2b*<1a1 / 3a<2b1 & I/IU/U & {-}/{-}{-}/{-} & N & W:1a,2a \\
C2969 & 1a<o1 / 2a<1a2L 2b*<1a1 / 3a<2b2 & I/IU/U & {-}/{-}{-}/{-} & N & W:1a,2a \\
C2970 & 1a<o1 / 2a<1a2L 2b*<1a2 / 3a<2a0L & I/IU/N & {-}/{-}{-}/{-} & N & W:1a,2a \\
C2971 & 1a<o1 / 2a<1a2L 2b*<1a2 / 3a<2a1 & I/IU/I & {-}/{-}{-}/{-} & N & W:1a,2a,3a \\
C2972 & 1a<o1 / 2a<1a2L 2b*<1a2 / 3a<2a1L & I/IU/I & {-}/{-}{-}/{-} & N & W:1a,2a,3a \\
C2973 & 1a<o1 / 2a<1a2L 2b*<1a2 / 3a<2a2 & I/IU/I & {-}/{-}{-}/{-} & N & W:1a,2a,3a \\
C2974 & 1a<o1 / 2a<1a2L 2b*<1a2 / 3a<2a2L & I/IU/I & {-}/{-}{-}/{-} & N & W:1a,2a,3a \\
C2975 & 1a<o1 / 2a<1a2L 2b*<1a2 / 3a<2b1 & I/IU/U & {-}/{-}{-}/{-} & N & W:1a,2a \\
C2976 & 1a<o1 / 2a<1a2L 2b*<1a2 / 3a<2b2 & I/IU/U & {-}/{-}{-}/{-} & N & W:1a,2a \\
C2977 & 1a<o1 / 2a*<1a1 / 3a<2a1 3b<2a1 & I/U/UU & {-}/{-}/{-}{-} & N & W:1a \\
C2978 & 1a<o1 / 2a*<1a1 / 3a<2a1 3b<2a2 & I/U/UU & {-}/{-}/{-}{-} & N & W:1a \\
C2979 & 1a<o1 / 2a*<1a1 / 3a<2a2 3b<2a2 & I/U/UU & {-}/{-}/{-}{-} & N & W:1a \\
C2980 & 1a<o1 / 2a*<1a2 / 3a<2a1 3b<2a1 & I/U/UU & {-}/{-}/{-}{-} & N & W:1a \\
C2981 & 1a<o1 / 2a*<1a2 / 3a<2a1 3b<2a2 & I/U/UU & {-}/{-}/{-}{-} & N & W:1a \\
C2982 & 1a<o1 / 2a*<1a2 / 3a<2a2 3b<2a2 & I/U/UU & {-}/{-}/{-}{-} & N & W:1a \\
C2983 & 1a<o1 1b<o1 / 2a<1a0L / 3a*<2a1 & EE/N/N & {-}{-}/{-}/{-} & E &  \\
C2984 & 1a<o1 1b<o1 / 2a<1a0L / 3a*<2a2 & EE/N/N & {-}{-}/{-}/{-} & E &  \\
C2985 & 1a<o1 1b<o1 / 2a<1a1 / 3a*<2a1 & EE/D/U & {-}{-}/{-}/{-} & E &  \\
C2986 & 1a<o1 1b<o1 / 2a<1a1 / 3a*<2a2 & EE/D/U & {-}{-}/{-}/{-} & E &  \\
C2987 & 1a<o1 1b<o1 / 2a<1a1L / 3a*<2a1 & EE/D/U & {-}{-}/{-}/{-} & E &  \\
C2988 & 1a<o1 1b<o1 / 2a<1a1L / 3a*<2a2 & EE/D/U & {-}{-}/{-}/{-} & E &  \\
C2989 & 1a<o1 1b<o1 / 2a<1a2 / 3a*<2a1 & EE/D/U & {-}{-}/{-}/{-} & E &  \\
C2990 & 1a<o1 1b<o1 / 2a<1a2 / 3a*<2a2 & EE/D/U & {-}{-}/{-}/{-} & E &  \\
C2991 & 1a<o1 1b<o1 / 2a<1a2L / 3a*<2a1 & EE/D/U & {-}{-}/{-}/{-} & E &  \\
C2992 & 1a<o1 1b<o1 / 2a<1a2L / 3a*<2a2 & EE/D/U & {-}{-}/{-}/{-} & E &  \\
C2993 & 1a<o1 1b<o1 / 2a*<1a1 / 3a<2a1 & EE/U/U & {-}{-}/{-}/{-} & E &  \\
C2994 & 1a<o1 1b<o1 / 2a*<1a1 / 3a<2a2 & EE/U/U & {-}{-}/{-}/{-} & E &  \\
C2995 & 1a<o1 1b<o1 / 2a*<1a2 / 3a<2a1 & EE/U/U & {-}{-}/{-}/{-} & E &  \\
C2996 & 1a<o1 1b<o1 / 2a*<1a2 / 3a<2a2 & EE/U/U & {-}{-}/{-}/{-} & E &  \\
C2997 & 1a<o1 1b<o1L / 2a<1a0L / 3a*<2a1 & LI/N/N & {-}+/{-}/{-} & B & CPW:1b \\
C2998 & 1a<o1 1b<o1L / 2a<1a0L / 3a*<2a2 & LI/N/N & {-}+/{-}/{-} & B & CPW:1b \\
C2999 & 1a<o1 1b<o1L / 2a<1a1 / 3a*<2a1 & LI/D/U & {-}+/{-}/{-} & B & CPW:1b \\
C3000 & 1a<o1 1b<o1L / 2a<1a1 / 3a*<2a2 & LI/D/U & {-}+/{-}/{-} & B & CPW:1b \\
C3001 & 1a<o1 1b<o1L / 2a<1a1L / 3a*<2a1 & LI/D/U & {-}+/{-}/{-} & B & CPW:1b \\
C3002 & 1a<o1 1b<o1L / 2a<1a1L / 3a*<2a2 & LI/D/U & {-}+/{-}/{-} & B & CPW:1b \\
C3003 & 1a<o1 1b<o1L / 2a<1a2 / 3a*<2a1 & LI/D/U & {-}+/{-}/{-} & B & CPW:1b \\
C3004 & 1a<o1 1b<o1L / 2a<1a2 / 3a*<2a2 & LI/D/U & {-}+/{-}/{-} & B & CPW:1b \\
C3005 & 1a<o1 1b<o1L / 2a<1a2L / 3a*<2a1 & LI/D/U & {-}+/{-}/{-} & B & CPW:1b \\
C3006 & 1a<o1 1b<o1L / 2a<1a2L / 3a*<2a2 & LI/D/U & {-}+/{-}/{-} & B & CPW:1b \\
C3007 & 1a<o1 1b<o1L / 2a<1b0L / 3a*<2a1 & LI/N/N & {-}+/+/{-} & B & CPW:1b \\
C3008 & 1a<o1 1b<o1L / 2a<1b0L / 3a*<2a2 & LI/N/N & {-}+/+/{-} & B & CPW:1b \\
C3009 & 1a<o1 1b<o1L / 2a<1b1 / 3a*<2a1 & LI/I/U & {-}+/{-}/{-} & B & CPW:1b,2a \\
C3010 & 1a<o1 1b<o1L / 2a<1b1 / 3a*<2a2 & LI/I/U & {-}+/{-}/{-} & B & CPW:1b,2a \\
C3011 & 1a<o1 1b<o1L / 2a<1b1L / 3a*<2a1 & LI/I/U & {-}+/+/{-} & B & CPW:1b,2a \\
C3012 & 1a<o1 1b<o1L / 2a<1b1L / 3a*<2a2 & LI/I/U & {-}+/+/{-} & B & CPW:1b,2a \\
C3013 & 1a<o1 1b<o1L / 2a<1b2 / 3a*<2a1 & LI/I/U & {-}+/{-}/{-} & B & CPW:1b,2a \\
C3014 & 1a<o1 1b<o1L / 2a<1b2 / 3a*<2a2 & LI/I/U & {-}+/{-}/{-} & B & CPW:1b,2a \\
C3015 & 1a<o1 1b<o1L / 2a<1b2L / 3a*<2a1 & LI/I/U & {-}+/+/{-} & B & CPW:1b,2a \\
C3016 & 1a<o1 1b<o1L / 2a<1b2L / 3a*<2a2 & LI/I/U & {-}+/+/{-} & B & CPW:1b,2a \\
C3017 & 1a<o1 1b<o1L / 2a*<1a1 / 3a<2a1 & LI/U/U & {-}+/{-}/{-} & B & CPW:1b \\
C3018 & 1a<o1 1b<o1L / 2a*<1a1 / 3a<2a2 & LI/U/U & {-}+/{-}/{-} & B & CPW:1b \\
C3019 & 1a<o1 1b<o1L / 2a*<1a2 / 3a<2a1 & LI/U/U & {-}+/{-}/{-} & B & CPW:1b \\
C3020 & 1a<o1 1b<o1L / 2a*<1a2 / 3a<2a2 & LI/U/U & {-}+/{-}/{-} & B & CPW:1b \\
C3021 & 1a<o1 1b<o1L / 2a*<1b1 / 3a<2a1 & LI/U/U & {-}+/{-}/{-} & B & CPW:1b \\
C3022 & 1a<o1 1b<o1L / 2a*<1b1 / 3a<2a2 & LI/U/U & {-}+/{-}/{-} & B & CPW:1b \\
C3023 & 1a<o1 1b<o1L / 2a*<1b2 / 3a<2a1 & LI/U/U & {-}+/{-}/{-} & B & CPW:1b \\
C3024 & 1a<o1 1b<o1L / 2a*<1b2 / 3a<2a2 & LI/U/U & {-}+/{-}/{-} & B & CPW:1b \\
C3025 & 1a<o1 1b<o2 / 2a<1a0L / 3a*<2a1 & EE/N/N & {-}{-}/{-}/{-} & E & C:1b \\
C3026 & 1a<o1 1b<o2 / 2a<1a0L / 3a*<2a2 & EE/N/N & {-}{-}/{-}/{-} & E & C:1b \\
C3027 & 1a<o1 1b<o2 / 2a<1a1 / 3a*<2a1 & EE/D/U & {-}{-}/{-}/{-} & E & C:1b \\
C3028 & 1a<o1 1b<o2 / 2a<1a1 / 3a*<2a2 & EE/D/U & {-}{-}/{-}/{-} & E & C:1b \\
C3029 & 1a<o1 1b<o2 / 2a<1a1L / 3a*<2a1 & EE/D/U & {-}{-}/{-}/{-} & E & C:1b \\
C3030 & 1a<o1 1b<o2 / 2a<1a1L / 3a*<2a2 & EE/D/U & {-}{-}/{-}/{-} & E & C:1b \\
C3031 & 1a<o1 1b<o2 / 2a<1a2 / 3a*<2a1 & EE/D/U & {-}{-}/{-}/{-} & E & C:1b \\
C3032 & 1a<o1 1b<o2 / 2a<1a2 / 3a*<2a2 & EE/D/U & {-}{-}/{-}/{-} & E & C:1b \\
C3033 & 1a<o1 1b<o2 / 2a<1a2L / 3a*<2a1 & EE/D/U & {-}{-}/{-}/{-} & E & C:1b \\
C3034 & 1a<o1 1b<o2 / 2a<1a2L / 3a*<2a2 & EE/D/U & {-}{-}/{-}/{-} & E & C:1b \\
C3035 & 1a<o1 1b<o2 / 2a<1b0L / 3a*<2a1 & EE/N/N & {-}{-}/{-}/{-} & E & C:1b \\
C3036 & 1a<o1 1b<o2 / 2a<1b0L / 3a*<2a2 & EE/N/N & {-}{-}/{-}/{-} & E & C:1b \\
C3037 & 1a<o1 1b<o2 / 2a<1b1 / 3a*<2a1 & EE/D/U & {-}{-}/{-}/{-} & E & C:1b,2a \\
C3038 & 1a<o1 1b<o2 / 2a<1b1 / 3a*<2a2 & EE/D/U & {-}{-}/{-}/{-} & E & C:1b,2a \\
C3039 & 1a<o1 1b<o2 / 2a<1b1L / 3a*<2a1 & EE/D/U & {-}{-}/{-}/{-} & E & C:1b,2a \\
C3040 & 1a<o1 1b<o2 / 2a<1b1L / 3a*<2a2 & EE/D/U & {-}{-}/{-}/{-} & E & C:1b,2a \\
C3041 & 1a<o1 1b<o2 / 2a<1b2 / 3a*<2a1 & EE/D/U & {-}{-}/{-}/{-} & E & C:1b,2a \\
C3042 & 1a<o1 1b<o2 / 2a<1b2 / 3a*<2a2 & EE/D/U & {-}{-}/{-}/{-} & E & C:1b,2a \\
C3043 & 1a<o1 1b<o2 / 2a<1b2L / 3a*<2a1 & EE/D/U & {-}{-}/{-}/{-} & E & C:1b,2a \\
C3044 & 1a<o1 1b<o2 / 2a<1b2L / 3a*<2a2 & EE/D/U & {-}{-}/{-}/{-} & E & C:1b,2a \\
C3045 & 1a<o1 1b<o2 / 2a*<1a1 / 3a<2a1 & EE/U/U & {-}{-}/{-}/{-} & E & C:1b \\
C3046 & 1a<o1 1b<o2 / 2a*<1a1 / 3a<2a2 & EE/U/U & {-}{-}/{-}/{-} & E & C:1b \\
C3047 & 1a<o1 1b<o2 / 2a*<1a2 / 3a<2a1 & EE/U/U & {-}{-}/{-}/{-} & E & C:1b \\
C3048 & 1a<o1 1b<o2 / 2a*<1a2 / 3a<2a2 & EE/U/U & {-}{-}/{-}/{-} & E & C:1b \\
C3049 & 1a<o1 1b<o2 / 2a*<1b1 / 3a<2a1 & EE/U/U & {-}{-}/{-}/{-} & E & C:1b \\
C3050 & 1a<o1 1b<o2 / 2a*<1b1 / 3a<2a2 & EE/U/U & {-}{-}/{-}/{-} & E & C:1b \\
C3051 & 1a<o1 1b<o2 / 2a*<1b2 / 3a<2a1 & EE/U/U & {-}{-}/{-}/{-} & E & C:1b \\
C3052 & 1a<o1 1b<o2 / 2a*<1b2 / 3a<2a2 & EE/U/U & {-}{-}/{-}/{-} & E & C:1b \\
C3053 & 1a<o1 1b<o2L / 2a<1a0L / 3a*<2a1 & LI/N/N & {-}+/{-}/{-} & B & PW:1b \\
C3054 & 1a<o1 1b<o2L / 2a<1a0L / 3a*<2a2 & LI/N/N & {-}+/{-}/{-} & B & PW:1b \\
C3055 & 1a<o1 1b<o2L / 2a<1a1 / 3a*<2a1 & LI/D/U & {-}+/{-}/{-} & B & PW:1b \\
C3056 & 1a<o1 1b<o2L / 2a<1a1 / 3a*<2a2 & LI/D/U & {-}+/{-}/{-} & B & PW:1b \\
C3057 & 1a<o1 1b<o2L / 2a<1a1L / 3a*<2a1 & LI/D/U & {-}+/{-}/{-} & B & PW:1b \\
C3058 & 1a<o1 1b<o2L / 2a<1a1L / 3a*<2a2 & LI/D/U & {-}+/{-}/{-} & B & PW:1b \\
C3059 & 1a<o1 1b<o2L / 2a<1a2 / 3a*<2a1 & LI/D/U & {-}+/{-}/{-} & B & PW:1b \\
C3060 & 1a<o1 1b<o2L / 2a<1a2 / 3a*<2a2 & LI/D/U & {-}+/{-}/{-} & B & PW:1b \\
C3061 & 1a<o1 1b<o2L / 2a<1a2L / 3a*<2a1 & LI/D/U & {-}+/{-}/{-} & B & PW:1b \\
C3062 & 1a<o1 1b<o2L / 2a<1a2L / 3a*<2a2 & LI/D/U & {-}+/{-}/{-} & B & PW:1b \\
C3063 & 1a<o1 1b<o2L / 2a<1b0L / 3a*<2a1 & LI/N/N & {-}+/+/{-} & B & PW:1b \\
C3064 & 1a<o1 1b<o2L / 2a<1b0L / 3a*<2a2 & LI/N/N & {-}+/+/{-} & B & PW:1b \\
C3065 & 1a<o1 1b<o2L / 2a<1b1 / 3a*<2a1 & LI/I/U & {-}+/{-}/{-} & B & PW:1b,2a \\
C3066 & 1a<o1 1b<o2L / 2a<1b1 / 3a*<2a2 & LI/I/U & {-}+/{-}/{-} & B & PW:1b,2a \\
C3067 & 1a<o1 1b<o2L / 2a<1b1L / 3a*<2a1 & LI/I/U & {-}+/+/{-} & B & PW:1b,2a \\
C3068 & 1a<o1 1b<o2L / 2a<1b1L / 3a*<2a2 & LI/I/U & {-}+/+/{-} & B & PW:1b,2a \\
C3069 & 1a<o1 1b<o2L / 2a<1b2 / 3a*<2a1 & LI/I/U & {-}+/{-}/{-} & B & PW:1b,2a \\
C3070 & 1a<o1 1b<o2L / 2a<1b2 / 3a*<2a2 & LI/I/U & {-}+/{-}/{-} & B & PW:1b,2a \\
C3071 & 1a<o1 1b<o2L / 2a<1b2L / 3a*<2a1 & LI/I/U & {-}+/+/{-} & B & PW:1b,2a \\
C3072 & 1a<o1 1b<o2L / 2a<1b2L / 3a*<2a2 & LI/I/U & {-}+/+/{-} & B & PW:1b,2a \\
C3073 & 1a<o1 1b<o2L / 2a*<1a1 / 3a<2a1 & LI/U/U & {-}+/{-}/{-} & B & PW:1b \\
C3074 & 1a<o1 1b<o2L / 2a*<1a1 / 3a<2a2 & LI/U/U & {-}+/{-}/{-} & B & PW:1b \\
C3075 & 1a<o1 1b<o2L / 2a*<1a2 / 3a<2a1 & LI/U/U & {-}+/{-}/{-} & B & PW:1b \\
C3076 & 1a<o1 1b<o2L / 2a*<1a2 / 3a<2a2 & LI/U/U & {-}+/{-}/{-} & B & PW:1b \\
C3077 & 1a<o1 1b<o2L / 2a*<1b1 / 3a<2a1 & LI/U/U & {-}+/{-}/{-} & B & PW:1b \\
C3078 & 1a<o1 1b<o2L / 2a*<1b1 / 3a<2a2 & LI/U/U & {-}+/{-}/{-} & B & PW:1b \\
C3079 & 1a<o1 1b<o2L / 2a*<1b2 / 3a<2a1 & LI/U/U & {-}+/{-}/{-} & B & PW:1b \\
C3080 & 1a<o1 1b<o2L / 2a*<1b2 / 3a<2a2 & LI/U/U & {-}+/{-}/{-} & B & PW:1b \\
C3081 & 1a<o1 1b*<o1 / 2a<1a0L / 3a<2a0L & IU/N/N & {-}{-}/{-}/{-} & N & W:1a \\
C3082 & 1a<o1 1b*<o1 / 2a<1a0L / 3a<2a1 & IU/N/N & {-}{-}/{-}/{-} & N & W:1a \\
C3083 & 1a<o1 1b*<o1 / 2a<1a0L / 3a<2a1L & IU/N/N & {-}{-}/{-}/{-} & N & W:1a \\
C3084 & 1a<o1 1b*<o1 / 2a<1a0L / 3a<2a2 & IU/N/N & {-}{-}/{-}/{-} & N & W:1a \\
C3085 & 1a<o1 1b*<o1 / 2a<1a0L / 3a<2a2L & IU/N/N & {-}{-}/{-}/{-} & N & W:1a \\
C3086 & 1a<o1 1b*<o1 / 2a<1a1 / 3a<2a0L & IU/I/N & {-}{-}/{-}/{-} & N & W:1a,2a \\
C3087 & 1a<o1 1b*<o1 / 2a<1a1 / 3a<2a1 & IU/I/I & {-}{-}/{-}/{-} & N & W:1a,2a,3a \\
C3088 & 1a<o1 1b*<o1 / 2a<1a1 / 3a<2a1L & IU/I/I & {-}{-}/{-}/{-} & N & W:1a,2a,3a \\
C3089 & 1a<o1 1b*<o1 / 2a<1a1 / 3a<2a2 & IU/I/I & {-}{-}/{-}/{-} & N & W:1a,2a,3a \\
C3090 & 1a<o1 1b*<o1 / 2a<1a1 / 3a<2a2L & IU/I/I & {-}{-}/{-}/{-} & N & W:1a,2a,3a \\
C3091 & 1a<o1 1b*<o1 / 2a<1a1L / 3a<2a0L & IU/I/N & {-}{-}/{-}/{-} & N & W:1a,2a \\
C3092 & 1a<o1 1b*<o1 / 2a<1a1L / 3a<2a1 & IU/I/I & {-}{-}/{-}/{-} & N & W:1a,2a,3a \\
C3093 & 1a<o1 1b*<o1 / 2a<1a1L / 3a<2a1L & IU/I/I & {-}{-}/{-}/{-} & N & W:1a,2a,3a \\
C3094 & 1a<o1 1b*<o1 / 2a<1a1L / 3a<2a2 & IU/I/I & {-}{-}/{-}/{-} & N & W:1a,2a,3a \\
C3095 & 1a<o1 1b*<o1 / 2a<1a1L / 3a<2a2L & IU/I/I & {-}{-}/{-}/{-} & N & W:1a,2a,3a \\
C3096 & 1a<o1 1b*<o1 / 2a<1a2 / 3a<2a0L & IU/I/N & {-}{-}/{-}/{-} & N & W:1a,2a \\
C3097 & 1a<o1 1b*<o1 / 2a<1a2 / 3a<2a1 & IU/I/I & {-}{-}/{-}/{-} & N & W:1a,2a,3a \\
C3098 & 1a<o1 1b*<o1 / 2a<1a2 / 3a<2a1L & IU/I/I & {-}{-}/{-}/{-} & N & W:1a,2a,3a \\
C3099 & 1a<o1 1b*<o1 / 2a<1a2 / 3a<2a2 & IU/I/I & {-}{-}/{-}/{-} & N & W:1a,2a,3a \\
C3100 & 1a<o1 1b*<o1 / 2a<1a2 / 3a<2a2L & IU/I/I & {-}{-}/{-}/{-} & N & W:1a,2a,3a \\
C3101 & 1a<o1 1b*<o1 / 2a<1a2L / 3a<2a0L & IU/I/N & {-}{-}/{-}/{-} & N & W:1a,2a \\
C3102 & 1a<o1 1b*<o1 / 2a<1a2L / 3a<2a1 & IU/I/I & {-}{-}/{-}/{-} & N & W:1a,2a,3a \\
C3103 & 1a<o1 1b*<o1 / 2a<1a2L / 3a<2a1L & IU/I/I & {-}{-}/{-}/{-} & N & W:1a,2a,3a \\
C3104 & 1a<o1 1b*<o1 / 2a<1a2L / 3a<2a2 & IU/I/I & {-}{-}/{-}/{-} & N & W:1a,2a,3a \\
C3105 & 1a<o1 1b*<o1 / 2a<1a2L / 3a<2a2L & IU/I/I & {-}{-}/{-}/{-} & N & W:1a,2a,3a \\
C3106 & 1a<o1 1b*<o1 / 2a<1b1 / 3a<2a0L & IU/U/N & {-}{-}/{-}/{-} & N & W:1a \\
C3107 & 1a<o1 1b*<o1 / 2a<1b1 / 3a<2a1 & IU/U/U & {-}{-}/{-}/{-} & N & W:1a \\
C3108 & 1a<o1 1b*<o1 / 2a<1b1 / 3a<2a1L & IU/U/U & {-}{-}/{-}/{-} & N & W:1a \\
C3109 & 1a<o1 1b*<o1 / 2a<1b1 / 3a<2a2 & IU/U/U & {-}{-}/{-}/{-} & N & W:1a \\
C3110 & 1a<o1 1b*<o1 / 2a<1b1 / 3a<2a2L & IU/U/U & {-}{-}/{-}/{-} & N & W:1a \\
C3111 & 1a<o1 1b*<o1 / 2a<1b2 / 3a<2a0L & IU/U/N & {-}{-}/{-}/{-} & N & W:1a \\
C3112 & 1a<o1 1b*<o1 / 2a<1b2 / 3a<2a1 & IU/U/U & {-}{-}/{-}/{-} & N & W:1a \\
C3113 & 1a<o1 1b*<o1 / 2a<1b2 / 3a<2a1L & IU/U/U & {-}{-}/{-}/{-} & N & W:1a \\
C3114 & 1a<o1 1b*<o1 / 2a<1b2 / 3a<2a2 & IU/U/U & {-}{-}/{-}/{-} & N & W:1a \\
C3115 & 1a<o1 1b*<o1 / 2a<1b2 / 3a<2a2L & IU/U/U & {-}{-}/{-}/{-} & N & W:1a \\
C3116 & 1a<o1 1b*<o2 / 2a<1a0L / 3a<2a0L & IU/N/N & {-}{-}/{-}/{-} & N & W:1a \\
C3117 & 1a<o1 1b*<o2 / 2a<1a0L / 3a<2a1 & IU/N/N & {-}{-}/{-}/{-} & N & W:1a \\
C3118 & 1a<o1 1b*<o2 / 2a<1a0L / 3a<2a1L & IU/N/N & {-}{-}/{-}/{-} & N & W:1a \\
C3119 & 1a<o1 1b*<o2 / 2a<1a0L / 3a<2a2 & IU/N/N & {-}{-}/{-}/{-} & N & W:1a \\
C3120 & 1a<o1 1b*<o2 / 2a<1a0L / 3a<2a2L & IU/N/N & {-}{-}/{-}/{-} & N & W:1a \\
C3121 & 1a<o1 1b*<o2 / 2a<1a1 / 3a<2a0L & IU/I/N & {-}{-}/{-}/{-} & N & W:1a,2a \\
C3122 & 1a<o1 1b*<o2 / 2a<1a1 / 3a<2a1 & IU/I/I & {-}{-}/{-}/{-} & N & W:1a,2a,3a \\
C3123 & 1a<o1 1b*<o2 / 2a<1a1 / 3a<2a1L & IU/I/I & {-}{-}/{-}/{-} & N & W:1a,2a,3a \\
C3124 & 1a<o1 1b*<o2 / 2a<1a1 / 3a<2a2 & IU/I/I & {-}{-}/{-}/{-} & N & W:1a,2a,3a \\
C3125 & 1a<o1 1b*<o2 / 2a<1a1 / 3a<2a2L & IU/I/I & {-}{-}/{-}/{-} & N & W:1a,2a,3a \\
C3126 & 1a<o1 1b*<o2 / 2a<1a1L / 3a<2a0L & IU/I/N & {-}{-}/{-}/{-} & N & W:1a,2a \\
C3127 & 1a<o1 1b*<o2 / 2a<1a1L / 3a<2a1 & IU/I/I & {-}{-}/{-}/{-} & N & W:1a,2a,3a \\
C3128 & 1a<o1 1b*<o2 / 2a<1a1L / 3a<2a1L & IU/I/I & {-}{-}/{-}/{-} & N & W:1a,2a,3a \\
C3129 & 1a<o1 1b*<o2 / 2a<1a1L / 3a<2a2 & IU/I/I & {-}{-}/{-}/{-} & N & W:1a,2a,3a \\
C3130 & 1a<o1 1b*<o2 / 2a<1a1L / 3a<2a2L & IU/I/I & {-}{-}/{-}/{-} & N & W:1a,2a,3a \\
C3131 & 1a<o1 1b*<o2 / 2a<1a2 / 3a<2a0L & IU/I/N & {-}{-}/{-}/{-} & N & W:1a,2a \\
C3132 & 1a<o1 1b*<o2 / 2a<1a2 / 3a<2a1 & IU/I/I & {-}{-}/{-}/{-} & N & W:1a,2a,3a \\
C3133 & 1a<o1 1b*<o2 / 2a<1a2 / 3a<2a1L & IU/I/I & {-}{-}/{-}/{-} & N & W:1a,2a,3a \\
C3134 & 1a<o1 1b*<o2 / 2a<1a2 / 3a<2a2 & IU/I/I & {-}{-}/{-}/{-} & N & W:1a,2a,3a \\
C3135 & 1a<o1 1b*<o2 / 2a<1a2 / 3a<2a2L & IU/I/I & {-}{-}/{-}/{-} & N & W:1a,2a,3a \\
C3136 & 1a<o1 1b*<o2 / 2a<1a2L / 3a<2a0L & IU/I/N & {-}{-}/{-}/{-} & N & W:1a,2a \\
C3137 & 1a<o1 1b*<o2 / 2a<1a2L / 3a<2a1 & IU/I/I & {-}{-}/{-}/{-} & N & W:1a,2a,3a \\
C3138 & 1a<o1 1b*<o2 / 2a<1a2L / 3a<2a1L & IU/I/I & {-}{-}/{-}/{-} & N & W:1a,2a,3a \\
C3139 & 1a<o1 1b*<o2 / 2a<1a2L / 3a<2a2 & IU/I/I & {-}{-}/{-}/{-} & N & W:1a,2a,3a \\
C3140 & 1a<o1 1b*<o2 / 2a<1a2L / 3a<2a2L & IU/I/I & {-}{-}/{-}/{-} & N & W:1a,2a,3a \\
C3141 & 1a<o1 1b*<o2 / 2a<1b1 / 3a<2a0L & IU/U/N & {-}{-}/{-}/{-} & N & W:1a \\
C3142 & 1a<o1 1b*<o2 / 2a<1b1 / 3a<2a1 & IU/U/U & {-}{-}/{-}/{-} & N & W:1a \\
C3143 & 1a<o1 1b*<o2 / 2a<1b1 / 3a<2a1L & IU/U/U & {-}{-}/{-}/{-} & N & W:1a \\
C3144 & 1a<o1 1b*<o2 / 2a<1b1 / 3a<2a2 & IU/U/U & {-}{-}/{-}/{-} & N & W:1a \\
C3145 & 1a<o1 1b*<o2 / 2a<1b1 / 3a<2a2L & IU/U/U & {-}{-}/{-}/{-} & N & W:1a \\
C3146 & 1a<o1 1b*<o2 / 2a<1b2 / 3a<2a0L & IU/U/N & {-}{-}/{-}/{-} & N & W:1a \\
C3147 & 1a<o1 1b*<o2 / 2a<1b2 / 3a<2a1 & IU/U/U & {-}{-}/{-}/{-} & N & W:1a \\
C3148 & 1a<o1 1b*<o2 / 2a<1b2 / 3a<2a1L & IU/U/U & {-}{-}/{-}/{-} & N & W:1a \\
C3149 & 1a<o1 1b*<o2 / 2a<1b2 / 3a<2a2 & IU/U/U & {-}{-}/{-}/{-} & N & W:1a \\
C3150 & 1a<o1 1b*<o2 / 2a<1b2 / 3a<2a2L & IU/U/U & {-}{-}/{-}/{-} & N & W:1a \\
C3151 & 1a<o1L / 2a<1a0L / 3a<2a0L 3b*<2a1 & I/N/NN & +/+/+{-} & N &  \\
C3152 & 1a<o1L / 2a<1a0L / 3a<2a0L 3b*<2a2 & I/N/NN & +/+/+{-} & N &  \\
C3153 & 1a<o1L / 2a<1a0L / 3a<2a1 3b*<2a1 & I/N/NN & +/+/{-}{-} & N &  \\
C3154 & 1a<o1L / 2a<1a0L / 3a<2a1 3b*<2a2 & I/N/NN & +/+/{-}{-} & N &  \\
C3155 & 1a<o1L / 2a<1a0L / 3a<2a1L 3b*<2a1 & I/N/NN & +/+/+{-} & N &  \\
C3156 & 1a<o1L / 2a<1a0L / 3a<2a1L 3b*<2a2 & I/N/NN & +/+/+{-} & N &  \\
C3157 & 1a<o1L / 2a<1a0L / 3a<2a2 3b*<2a1 & I/N/NN & +/+/{-}{-} & N &  \\
C3158 & 1a<o1L / 2a<1a0L / 3a<2a2 3b*<2a2 & I/N/NN & +/+/{-}{-} & N &  \\
C3159 & 1a<o1L / 2a<1a0L / 3a<2a2L 3b*<2a1 & I/N/NN & +/+/+{-} & N &  \\
C3160 & 1a<o1L / 2a<1a0L / 3a<2a2L 3b*<2a2 & I/N/NN & +/+/+{-} & N &  \\
C3161 & 1a<o1L / 2a<1a0L 2b<1a1 / 3a*<2a1 & I/NI/N & +/+{-}/{-} & N &  \\
C3162 & 1a<o1L / 2a<1a0L 2b<1a1 / 3a*<2a2 & I/NI/N & +/+{-}/{-} & N &  \\
C3163 & 1a<o1L / 2a<1a0L 2b<1a1 / 3a*<2b1 & I/NI/U & +/+{-}/{-} & N & W:3a \\
C3164 & 1a<o1L / 2a<1a0L 2b<1a1 / 3a*<2b2 & I/NI/U & +/+{-}/{-} & N & W:3a \\
C3165 & 1a<o1L / 2a<1a0L 2b<1a2 / 3a*<2a1 & I/NI/N & +/+{-}/{-} & N &  \\
C3166 & 1a<o1L / 2a<1a0L 2b<1a2 / 3a*<2a2 & I/NI/N & +/+{-}/{-} & N &  \\
C3167 & 1a<o1L / 2a<1a0L 2b<1a2 / 3a*<2b1 & I/NI/U & +/+{-}/{-} & N & W:3a \\
C3168 & 1a<o1L / 2a<1a0L 2b<1a2 / 3a*<2b2 & I/NI/U & +/+{-}/{-} & N & W:3a \\
C3169 & 1a<o1L / 2a<1a0L 2b*<1a1 / 3a<2a0L & I/NU/N & +/+{-}/+ & N & W:2b \\
C3170 & 1a<o1L / 2a<1a0L 2b*<1a1 / 3a<2a1 & I/NU/N & +/+{-}/{-} & N & W:2b \\
C3171 & 1a<o1L / 2a<1a0L 2b*<1a1 / 3a<2a1L & I/NU/N & +/+{-}/+ & N & W:2b \\
C3172 & 1a<o1L / 2a<1a0L 2b*<1a1 / 3a<2a2 & I/NU/N & +/+{-}/{-} & N & W:2b \\
C3173 & 1a<o1L / 2a<1a0L 2b*<1a1 / 3a<2a2L & I/NU/N & +/+{-}/+ & N & W:2b \\
C3174 & 1a<o1L / 2a<1a0L 2b*<1a1 / 3a<2b1 & I/NU/U & +/+{-}/{-} & N & W:2b,3a \\
C3175 & 1a<o1L / 2a<1a0L 2b*<1a1 / 3a<2b2 & I/NU/U & +/+{-}/{-} & N & W:2b,3a \\
C3176 & 1a<o1L / 2a<1a0L 2b*<1a2 / 3a<2a0L & I/NU/N & +/+{-}/+ & N & W:2b \\
C3177 & 1a<o1L / 2a<1a0L 2b*<1a2 / 3a<2a1 & I/NU/N & +/+{-}/{-} & N & W:2b \\
C3178 & 1a<o1L / 2a<1a0L 2b*<1a2 / 3a<2a1L & I/NU/N & +/+{-}/+ & N & W:2b \\
C3179 & 1a<o1L / 2a<1a0L 2b*<1a2 / 3a<2a2 & I/NU/N & +/+{-}/{-} & N & W:2b \\
C3180 & 1a<o1L / 2a<1a0L 2b*<1a2 / 3a<2a2L & I/NU/N & +/+{-}/+ & N & W:2b \\
C3181 & 1a<o1L / 2a<1a0L 2b*<1a2 / 3a<2b1 & I/NU/U & +/+{-}/{-} & N & W:2b,3a \\
C3182 & 1a<o1L / 2a<1a0L 2b*<1a2 / 3a<2b2 & I/NU/U & +/+{-}/{-} & N & W:2b,3a \\
C3183 & 1a<o1L / 2a<1a1 / 3a<2a0L 3b*<2a1 & I/I/NU & +/{-}/{-}{-} & N & W:3b \\
C3184 & 1a<o1L / 2a<1a1 / 3a<2a0L 3b*<2a2 & I/I/NU & +/{-}/{-}{-} & N & W:3b \\
C3185 & 1a<o1L / 2a<1a1 / 3a<2a1 3b*<2a1 & I/I/IU & +/{-}/{-}{-} & N & W:1a,2a,3a \\
C3186 & 1a<o1L / 2a<1a1 / 3a<2a1 3b*<2a2 & I/I/IU & +/{-}/{-}{-} & N & W:1a,2a,3a \\
C3187 & 1a<o1L / 2a<1a1 / 3a<2a1L 3b*<2a1 & I/I/IU & +/{-}/{-}{-} & N & W:1a,2a,3a \\
C3188 & 1a<o1L / 2a<1a1 / 3a<2a1L 3b*<2a2 & I/I/IU & +/{-}/{-}{-} & N & W:1a,2a,3a \\
C3189 & 1a<o1L / 2a<1a1 / 3a<2a2 3b*<2a1 & I/I/IU & +/{-}/{-}{-} & N & W:1a,2a,3a \\
C3190 & 1a<o1L / 2a<1a1 / 3a<2a2 3b*<2a2 & I/I/IU & +/{-}/{-}{-} & N & W:1a,2a,3a \\
C3191 & 1a<o1L / 2a<1a1 / 3a<2a2L 3b*<2a1 & I/I/IU & +/{-}/{-}{-} & N & W:1a,2a,3a \\
C3192 & 1a<o1L / 2a<1a1 / 3a<2a2L 3b*<2a2 & I/I/IU & +/{-}/{-}{-} & N & W:1a,2a,3a \\
C3193 & 1a<o1L / 2a<1a1 2b<1a1 / 3a*<2a1 & I/EE/U & +/{-}{-}/{-} & E & W:1a \\
C3194 & 1a<o1L / 2a<1a1 2b<1a1 / 3a*<2a2 & I/EE/U & +/{-}{-}/{-} & E & W:1a \\
C3195 & 1a<o1L / 2a<1a1 2b<1a1L / 3a*<2a1 & I/LI/U & +/{-}+/{-} & B & CP:2b W:1a,2b \\
C3196 & 1a<o1L / 2a<1a1 2b<1a1L / 3a*<2a2 & I/LI/U & +/{-}+/{-} & B & CP:2b W:1a,2b \\
C3197 & 1a<o1L / 2a<1a1 2b<1a1L / 3a*<2b1 & I/LI/U & +/{-}+/{-} & B & CP:2b W:1a,2b \\
C3198 & 1a<o1L / 2a<1a1 2b<1a1L / 3a*<2b2 & I/LI/U & +/{-}+/{-} & B & CP:2b W:1a,2b \\
C3199 & 1a<o1L / 2a<1a1 2b<1a2 / 3a*<2a1 & I/EE/U & +/{-}{-}/{-} & E & C:2b W:1a \\
C3200 & 1a<o1L / 2a<1a1 2b<1a2 / 3a*<2a2 & I/EE/U & +/{-}{-}/{-} & E & C:2b W:1a \\
C3201 & 1a<o1L / 2a<1a1 2b<1a2 / 3a*<2b1 & I/EE/U & +/{-}{-}/{-} & E & C:2b W:1a \\
C3202 & 1a<o1L / 2a<1a1 2b<1a2 / 3a*<2b2 & I/EE/U & +/{-}{-}/{-} & E & C:2b W:1a \\
C3203 & 1a<o1L / 2a<1a1 2b<1a2L / 3a*<2a1 & I/LI/U & +/{-}+/{-} & B & P:2b W:1a,2b \\
C3204 & 1a<o1L / 2a<1a1 2b<1a2L / 3a*<2a2 & I/LI/U & +/{-}+/{-} & B & P:2b W:1a,2b \\
C3205 & 1a<o1L / 2a<1a1 2b<1a2L / 3a*<2b1 & I/LI/U & +/{-}+/{-} & B & P:2b W:1a,2b \\
C3206 & 1a<o1L / 2a<1a1 2b<1a2L / 3a*<2b2 & I/LI/U & +/{-}+/{-} & B & P:2b W:1a,2b \\
C3207 & 1a<o1L / 2a<1a1 2b*<1a1 / 3a<2a0L & I/IU/N & +/{-}{-}/{-} & N & W:1a,2a \\
C3208 & 1a<o1L / 2a<1a1 2b*<1a1 / 3a<2a1 & I/IU/I & +/{-}{-}/{-} & N & W:1a,2a,3a \\
C3209 & 1a<o1L / 2a<1a1 2b*<1a1 / 3a<2a1L & I/IU/I & +/{-}{-}/{-} & N & W:1a,2a,3a \\
C3210 & 1a<o1L / 2a<1a1 2b*<1a1 / 3a<2a2 & I/IU/I & +/{-}{-}/{-} & N & W:1a,2a,3a \\
C3211 & 1a<o1L / 2a<1a1 2b*<1a1 / 3a<2a2L & I/IU/I & +/{-}{-}/{-} & N & W:1a,2a,3a \\
C3212 & 1a<o1L / 2a<1a1 2b*<1a1 / 3a<2b1 & I/IU/U & +/{-}{-}/{-} & N & W:1a,2a \\
C3213 & 1a<o1L / 2a<1a1 2b*<1a1 / 3a<2b2 & I/IU/U & +/{-}{-}/{-} & N & W:1a,2a \\
C3214 & 1a<o1L / 2a<1a1 2b*<1a2 / 3a<2a0L & I/IU/N & +/{-}{-}/{-} & N & W:1a,2a \\
C3215 & 1a<o1L / 2a<1a1 2b*<1a2 / 3a<2a1 & I/IU/I & +/{-}{-}/{-} & N & W:1a,2a,3a \\
C3216 & 1a<o1L / 2a<1a1 2b*<1a2 / 3a<2a1L & I/IU/I & +/{-}{-}/{-} & N & W:1a,2a,3a \\
C3217 & 1a<o1L / 2a<1a1 2b*<1a2 / 3a<2a2 & I/IU/I & +/{-}{-}/{-} & N & W:1a,2a,3a \\
C3218 & 1a<o1L / 2a<1a1 2b*<1a2 / 3a<2a2L & I/IU/I & +/{-}{-}/{-} & N & W:1a,2a,3a \\
C3219 & 1a<o1L / 2a<1a1 2b*<1a2 / 3a<2b1 & I/IU/U & +/{-}{-}/{-} & N & W:1a,2a \\
C3220 & 1a<o1L / 2a<1a1 2b*<1a2 / 3a<2b2 & I/IU/U & +/{-}{-}/{-} & N & W:1a,2a \\
C3221 & 1a<o1L / 2a<1a1L / 3a<2a0L 3b*<2a1 & I/I/NU & +/+/+{-} & N & W:3b \\
C3222 & 1a<o1L / 2a<1a1L / 3a<2a0L 3b*<2a2 & I/I/NU & +/+/+{-} & N & W:3b \\
C3223 & 1a<o1L / 2a<1a1L / 3a<2a1 3b*<2a1 & I/I/IU & +/+/{-}{-} & N & W:1a,2a,3a \\
C3224 & 1a<o1L / 2a<1a1L / 3a<2a1 3b*<2a2 & I/I/IU & +/+/{-}{-} & N & W:1a,2a,3a \\
C3225 & 1a<o1L / 2a<1a1L / 3a<2a1L 3b*<2a1 & I/I/IU & +/+/+{-} & N & W:1a,2a,3a \\
C3226 & 1a<o1L / 2a<1a1L / 3a<2a1L 3b*<2a2 & I/I/IU & +/+/+{-} & N & W:1a,2a,3a \\
C3227 & 1a<o1L / 2a<1a1L / 3a<2a2 3b*<2a1 & I/I/IU & +/+/{-}{-} & N & W:1a,2a,3a \\
C3228 & 1a<o1L / 2a<1a1L / 3a<2a2 3b*<2a2 & I/I/IU & +/+/{-}{-} & N & W:1a,2a,3a \\
C3229 & 1a<o1L / 2a<1a1L / 3a<2a2L 3b*<2a1 & I/I/IU & +/+/+{-} & N & W:1a,2a,3a \\
C3230 & 1a<o1L / 2a<1a1L / 3a<2a2L 3b*<2a2 & I/I/IU & +/+/+{-} & N & W:1a,2a,3a \\
C3231 & 1a<o1L / 2a<1a1L 2b<1a2 / 3a*<2a1 & I/IL/U & +/+{-}/{-} & B & C:2a,2b P:2a W:1a,2a \\
C3232 & 1a<o1L / 2a<1a1L 2b<1a2 / 3a*<2a2 & I/IL/U & +/+{-}/{-} & B & C:2a,2b P:2a W:1a,2a \\
C3233 & 1a<o1L / 2a<1a1L 2b<1a2 / 3a*<2b1 & I/IL/U & +/+{-}/{-} & B & C:2a,2b P:2a W:1a,2a \\
C3234 & 1a<o1L / 2a<1a1L 2b<1a2 / 3a*<2b2 & I/IL/U & +/+{-}/{-} & B & C:2a,2b P:2a W:1a,2a \\
C3235 & 1a<o1L / 2a<1a1L 2b*<1a1 / 3a<2a0L & I/IU/N & +/+{-}/+ & N & W:1a,2a \\
C3236 & 1a<o1L / 2a<1a1L 2b*<1a1 / 3a<2a1 & I/IU/I & +/+{-}/{-} & N & W:1a,2a,3a \\
C3237 & 1a<o1L / 2a<1a1L 2b*<1a1 / 3a<2a1L & I/IU/I & +/+{-}/+ & N & W:1a,2a,3a \\
C3238 & 1a<o1L / 2a<1a1L 2b*<1a1 / 3a<2a2 & I/IU/I & +/+{-}/{-} & N & W:1a,2a,3a \\
C3239 & 1a<o1L / 2a<1a1L 2b*<1a1 / 3a<2a2L & I/IU/I & +/+{-}/+ & N & W:1a,2a,3a \\
C3240 & 1a<o1L / 2a<1a1L 2b*<1a1 / 3a<2b1 & I/IU/U & +/+{-}/{-} & N & W:1a,2a \\
C3241 & 1a<o1L / 2a<1a1L 2b*<1a1 / 3a<2b2 & I/IU/U & +/+{-}/{-} & N & W:1a,2a \\
C3242 & 1a<o1L / 2a<1a1L 2b*<1a2 / 3a<2a0L & I/IU/N & +/+{-}/+ & N & W:1a,2a \\
C3243 & 1a<o1L / 2a<1a1L 2b*<1a2 / 3a<2a1 & I/IU/I & +/+{-}/{-} & N & W:1a,2a,3a \\
C3244 & 1a<o1L / 2a<1a1L 2b*<1a2 / 3a<2a1L & I/IU/I & +/+{-}/+ & N & W:1a,2a,3a \\
C3245 & 1a<o1L / 2a<1a1L 2b*<1a2 / 3a<2a2 & I/IU/I & +/+{-}/{-} & N & W:1a,2a,3a \\
C3246 & 1a<o1L / 2a<1a1L 2b*<1a2 / 3a<2a2L & I/IU/I & +/+{-}/+ & N & W:1a,2a,3a \\
C3247 & 1a<o1L / 2a<1a1L 2b*<1a2 / 3a<2b1 & I/IU/U & +/+{-}/{-} & N & W:1a,2a \\
C3248 & 1a<o1L / 2a<1a1L 2b*<1a2 / 3a<2b2 & I/IU/U & +/+{-}/{-} & N & W:1a,2a \\
C3249 & 1a<o1L / 2a<1a2 / 3a<2a0L 3b*<2a1 & I/I/NU & +/{-}/{-}{-} & N & W:3b \\
C3250 & 1a<o1L / 2a<1a2 / 3a<2a0L 3b*<2a2 & I/I/NU & +/{-}/{-}{-} & N & W:3b \\
C3251 & 1a<o1L / 2a<1a2 / 3a<2a1 3b*<2a1 & I/I/IU & +/{-}/{-}{-} & N & W:1a,2a,3a \\
C3252 & 1a<o1L / 2a<1a2 / 3a<2a1 3b*<2a2 & I/I/IU & +/{-}/{-}{-} & N & W:1a,2a,3a \\
C3253 & 1a<o1L / 2a<1a2 / 3a<2a1L 3b*<2a1 & I/I/IU & +/{-}/{-}{-} & N & W:1a,2a,3a \\
C3254 & 1a<o1L / 2a<1a2 / 3a<2a1L 3b*<2a2 & I/I/IU & +/{-}/{-}{-} & N & W:1a,2a,3a \\
C3255 & 1a<o1L / 2a<1a2 / 3a<2a2 3b*<2a1 & I/I/IU & +/{-}/{-}{-} & N & W:1a,2a,3a \\
C3256 & 1a<o1L / 2a<1a2 / 3a<2a2 3b*<2a2 & I/I/IU & +/{-}/{-}{-} & N & W:1a,2a,3a \\
C3257 & 1a<o1L / 2a<1a2 / 3a<2a2L 3b*<2a1 & I/I/IU & +/{-}/{-}{-} & N & W:1a,2a,3a \\
C3258 & 1a<o1L / 2a<1a2 / 3a<2a2L 3b*<2a2 & I/I/IU & +/{-}/{-}{-} & N & W:1a,2a,3a \\
C3259 & 1a<o1L / 2a<1a2 2b<1a2 / 3a*<2a1 & I/EE/U & +/{-}{-}/{-} & E & W:1a \\
C3260 & 1a<o1L / 2a<1a2 2b<1a2 / 3a*<2a2 & I/EE/U & +/{-}{-}/{-} & E & W:1a \\
C3261 & 1a<o1L / 2a<1a2 2b<1a2L / 3a*<2a1 & I/LI/U & +/{-}+/{-} & B & CP:2b W:1a,2b \\
C3262 & 1a<o1L / 2a<1a2 2b<1a2L / 3a*<2a2 & I/LI/U & +/{-}+/{-} & B & CP:2b W:1a,2b \\
C3263 & 1a<o1L / 2a<1a2 2b<1a2L / 3a*<2b1 & I/LI/U & +/{-}+/{-} & B & CP:2b W:1a,2b \\
C3264 & 1a<o1L / 2a<1a2 2b<1a2L / 3a*<2b2 & I/LI/U & +/{-}+/{-} & B & CP:2b W:1a,2b \\
C3265 & 1a<o1L / 2a<1a2 2b*<1a1 / 3a<2a0L & I/IU/N & +/{-}{-}/{-} & N & W:1a,2a \\
C3266 & 1a<o1L / 2a<1a2 2b*<1a1 / 3a<2a1 & I/IU/I & +/{-}{-}/{-} & N & W:1a,2a,3a \\
C3267 & 1a<o1L / 2a<1a2 2b*<1a1 / 3a<2a1L & I/IU/I & +/{-}{-}/{-} & N & W:1a,2a,3a \\
C3268 & 1a<o1L / 2a<1a2 2b*<1a1 / 3a<2a2 & I/IU/I & +/{-}{-}/{-} & N & W:1a,2a,3a \\
C3269 & 1a<o1L / 2a<1a2 2b*<1a1 / 3a<2a2L & I/IU/I & +/{-}{-}/{-} & N & W:1a,2a,3a \\
C3270 & 1a<o1L / 2a<1a2 2b*<1a1 / 3a<2b1 & I/IU/U & +/{-}{-}/{-} & N & W:1a,2a \\
C3271 & 1a<o1L / 2a<1a2 2b*<1a1 / 3a<2b2 & I/IU/U & +/{-}{-}/{-} & N & W:1a,2a \\
C3272 & 1a<o1L / 2a<1a2 2b*<1a2 / 3a<2a0L & I/IU/N & +/{-}{-}/{-} & N & W:1a,2a \\
C3273 & 1a<o1L / 2a<1a2 2b*<1a2 / 3a<2a1 & I/IU/I & +/{-}{-}/{-} & N & W:1a,2a,3a \\
C3274 & 1a<o1L / 2a<1a2 2b*<1a2 / 3a<2a1L & I/IU/I & +/{-}{-}/{-} & N & W:1a,2a,3a \\
C3275 & 1a<o1L / 2a<1a2 2b*<1a2 / 3a<2a2 & I/IU/I & +/{-}{-}/{-} & N & W:1a,2a,3a \\
C3276 & 1a<o1L / 2a<1a2 2b*<1a2 / 3a<2a2L & I/IU/I & +/{-}{-}/{-} & N & W:1a,2a,3a \\
C3277 & 1a<o1L / 2a<1a2 2b*<1a2 / 3a<2b1 & I/IU/U & +/{-}{-}/{-} & N & W:1a,2a \\
C3278 & 1a<o1L / 2a<1a2 2b*<1a2 / 3a<2b2 & I/IU/U & +/{-}{-}/{-} & N & W:1a,2a \\
C3279 & 1a<o1L / 2a<1a2L / 3a<2a0L 3b*<2a1 & I/I/NU & +/+/+{-} & N & W:3b \\
C3280 & 1a<o1L / 2a<1a2L / 3a<2a0L 3b*<2a2 & I/I/NU & +/+/+{-} & N & W:3b \\
C3281 & 1a<o1L / 2a<1a2L / 3a<2a1 3b*<2a1 & I/I/IU & +/+/{-}{-} & N & W:1a,2a,3a \\
C3282 & 1a<o1L / 2a<1a2L / 3a<2a1 3b*<2a2 & I/I/IU & +/+/{-}{-} & N & W:1a,2a,3a \\
C3283 & 1a<o1L / 2a<1a2L / 3a<2a1L 3b*<2a1 & I/I/IU & +/+/+{-} & N & W:1a,2a,3a \\
C3284 & 1a<o1L / 2a<1a2L / 3a<2a1L 3b*<2a2 & I/I/IU & +/+/+{-} & N & W:1a,2a,3a \\
C3285 & 1a<o1L / 2a<1a2L / 3a<2a2 3b*<2a1 & I/I/IU & +/+/{-}{-} & N & W:1a,2a,3a \\
C3286 & 1a<o1L / 2a<1a2L / 3a<2a2 3b*<2a2 & I/I/IU & +/+/{-}{-} & N & W:1a,2a,3a \\
C3287 & 1a<o1L / 2a<1a2L / 3a<2a2L 3b*<2a1 & I/I/IU & +/+/+{-} & N & W:1a,2a,3a \\
C3288 & 1a<o1L / 2a<1a2L / 3a<2a2L 3b*<2a2 & I/I/IU & +/+/+{-} & N & W:1a,2a,3a \\
C3289 & 1a<o1L / 2a<1a2L 2b*<1a1 / 3a<2a0L & I/IU/N & +/+{-}/+ & N & W:1a,2a \\
C3290 & 1a<o1L / 2a<1a2L 2b*<1a1 / 3a<2a1 & I/IU/I & +/+{-}/{-} & N & W:1a,2a,3a \\
C3291 & 1a<o1L / 2a<1a2L 2b*<1a1 / 3a<2a1L & I/IU/I & +/+{-}/+ & N & W:1a,2a,3a \\
C3292 & 1a<o1L / 2a<1a2L 2b*<1a1 / 3a<2a2 & I/IU/I & +/+{-}/{-} & N & W:1a,2a,3a \\
C3293 & 1a<o1L / 2a<1a2L 2b*<1a1 / 3a<2a2L & I/IU/I & +/+{-}/+ & N & W:1a,2a,3a \\
C3294 & 1a<o1L / 2a<1a2L 2b*<1a1 / 3a<2b1 & I/IU/U & +/+{-}/{-} & N & W:1a,2a \\
C3295 & 1a<o1L / 2a<1a2L 2b*<1a1 / 3a<2b2 & I/IU/U & +/+{-}/{-} & N & W:1a,2a \\
C3296 & 1a<o1L / 2a<1a2L 2b*<1a2 / 3a<2a0L & I/IU/N & +/+{-}/+ & N & W:1a,2a \\
C3297 & 1a<o1L / 2a<1a2L 2b*<1a2 / 3a<2a1 & I/IU/I & +/+{-}/{-} & N & W:1a,2a,3a \\
C3298 & 1a<o1L / 2a<1a2L 2b*<1a2 / 3a<2a1L & I/IU/I & +/+{-}/+ & N & W:1a,2a,3a \\
C3299 & 1a<o1L / 2a<1a2L 2b*<1a2 / 3a<2a2 & I/IU/I & +/+{-}/{-} & N & W:1a,2a,3a \\
C3300 & 1a<o1L / 2a<1a2L 2b*<1a2 / 3a<2a2L & I/IU/I & +/+{-}/+ & N & W:1a,2a,3a \\
C3301 & 1a<o1L / 2a<1a2L 2b*<1a2 / 3a<2b1 & I/IU/U & +/+{-}/{-} & N & W:1a,2a \\
C3302 & 1a<o1L / 2a<1a2L 2b*<1a2 / 3a<2b2 & I/IU/U & +/+{-}/{-} & N & W:1a,2a \\
C3303 & 1a<o1L / 2a*<1a1 / 3a<2a1 3b<2a1 & I/U/UU & +/{-}/{-}{-} & N & W:1a \\
C3304 & 1a<o1L / 2a*<1a1 / 3a<2a1 3b<2a2 & I/U/UU & +/{-}/{-}{-} & N & W:1a \\
C3305 & 1a<o1L / 2a*<1a1 / 3a<2a2 3b<2a2 & I/U/UU & +/{-}/{-}{-} & N & W:1a \\
C3306 & 1a<o1L / 2a*<1a2 / 3a<2a1 3b<2a1 & I/U/UU & +/{-}/{-}{-} & N & W:1a \\
C3307 & 1a<o1L / 2a*<1a2 / 3a<2a1 3b<2a2 & I/U/UU & +/{-}/{-}{-} & N & W:1a \\
C3308 & 1a<o1L / 2a*<1a2 / 3a<2a2 3b<2a2 & I/U/UU & +/{-}/{-}{-} & N & W:1a \\
C3309 & 1a<o1L 1b<o2 / 2a<1a0L / 3a*<2a1 & IL/N/N & +{-}/+/{-} & B & C:1a,1b PW:1a \\
C3310 & 1a<o1L 1b<o2 / 2a<1a0L / 3a*<2a2 & IL/N/N & +{-}/+/{-} & B & C:1a,1b PW:1a \\
C3311 & 1a<o1L 1b<o2 / 2a<1a1 / 3a*<2a1 & IL/I/U & +{-}/{-}/{-} & B & C:1a,1b,2a PW:1a,2a \\
C3312 & 1a<o1L 1b<o2 / 2a<1a1 / 3a*<2a2 & IL/I/U & +{-}/{-}/{-} & B & C:1a,1b,2a PW:1a,2a \\
C3313 & 1a<o1L 1b<o2 / 2a<1a1L / 3a*<2a1 & IL/I/U & +{-}/+/{-} & B & C:1a,1b,2a PW:1a,2a \\
C3314 & 1a<o1L 1b<o2 / 2a<1a1L / 3a*<2a2 & IL/I/U & +{-}/+/{-} & B & C:1a,1b,2a PW:1a,2a \\
C3315 & 1a<o1L 1b<o2 / 2a<1a2 / 3a*<2a1 & IL/I/U & +{-}/{-}/{-} & B & C:1a,1b,2a PW:1a,2a \\
C3316 & 1a<o1L 1b<o2 / 2a<1a2 / 3a*<2a2 & IL/I/U & +{-}/{-}/{-} & B & C:1a,1b,2a PW:1a,2a \\
C3317 & 1a<o1L 1b<o2 / 2a<1a2L / 3a*<2a1 & IL/I/U & +{-}/+/{-} & B & C:1a,1b,2a PW:1a,2a \\
C3318 & 1a<o1L 1b<o2 / 2a<1a2L / 3a*<2a2 & IL/I/U & +{-}/+/{-} & B & C:1a,1b,2a PW:1a,2a \\
C3319 & 1a<o1L 1b<o2 / 2a<1b0L / 3a*<2a1 & IL/N/N & +{-}/{-}/{-} & B & C:1a,1b PW:1a \\
C3320 & 1a<o1L 1b<o2 / 2a<1b0L / 3a*<2a2 & IL/N/N & +{-}/{-}/{-} & B & C:1a,1b PW:1a \\
C3321 & 1a<o1L 1b<o2 / 2a<1b1 / 3a*<2a1 & IL/D/U & +{-}/{-}/{-} & B & C:1a,1b,2a PW:1a \\
C3322 & 1a<o1L 1b<o2 / 2a<1b1 / 3a*<2a2 & IL/D/U & +{-}/{-}/{-} & B & C:1a,1b,2a PW:1a \\
C3323 & 1a<o1L 1b<o2 / 2a<1b1L / 3a*<2a1 & IL/D/U & +{-}/{-}/{-} & B & C:1a,1b,2a PW:1a \\
C3324 & 1a<o1L 1b<o2 / 2a<1b1L / 3a*<2a2 & IL/D/U & +{-}/{-}/{-} & B & C:1a,1b,2a PW:1a \\
C3325 & 1a<o1L 1b<o2 / 2a<1b2 / 3a*<2a1 & IL/D/U & +{-}/{-}/{-} & B & C:1a,1b,2a PW:1a \\
C3326 & 1a<o1L 1b<o2 / 2a<1b2 / 3a*<2a2 & IL/D/U & +{-}/{-}/{-} & B & C:1a,1b,2a PW:1a \\
C3327 & 1a<o1L 1b<o2 / 2a<1b2L / 3a*<2a1 & IL/D/U & +{-}/{-}/{-} & B & C:1a,1b,2a PW:1a \\
C3328 & 1a<o1L 1b<o2 / 2a<1b2L / 3a*<2a2 & IL/D/U & +{-}/{-}/{-} & B & C:1a,1b,2a PW:1a \\
C3329 & 1a<o1L 1b<o2 / 2a*<1a1 / 3a<2a1 & IL/U/U & +{-}/{-}/{-} & B & C:1a,1b PW:1a \\
C3330 & 1a<o1L 1b<o2 / 2a*<1a1 / 3a<2a2 & IL/U/U & +{-}/{-}/{-} & B & C:1a,1b PW:1a \\
C3331 & 1a<o1L 1b<o2 / 2a*<1a2 / 3a<2a1 & IL/U/U & +{-}/{-}/{-} & B & C:1a,1b PW:1a \\
C3332 & 1a<o1L 1b<o2 / 2a*<1a2 / 3a<2a2 & IL/U/U & +{-}/{-}/{-} & B & C:1a,1b PW:1a \\
C3333 & 1a<o1L 1b<o2 / 2a*<1b1 / 3a<2a1 & IL/U/U & +{-}/{-}/{-} & B & C:1a,1b PW:1a \\
C3334 & 1a<o1L 1b<o2 / 2a*<1b1 / 3a<2a2 & IL/U/U & +{-}/{-}/{-} & B & C:1a,1b PW:1a \\
C3335 & 1a<o1L 1b<o2 / 2a*<1b2 / 3a<2a1 & IL/U/U & +{-}/{-}/{-} & B & C:1a,1b PW:1a \\
C3336 & 1a<o1L 1b<o2 / 2a*<1b2 / 3a<2a2 & IL/U/U & +{-}/{-}/{-} & B & C:1a,1b PW:1a \\
C3337 & 1a<o1L 1b*<o1 / 2a<1a0L / 3a<2a0L & IU/N/N & +{-}/+/+ & N & W:1a \\
C3338 & 1a<o1L 1b*<o1 / 2a<1a0L / 3a<2a1 & IU/N/N & +{-}/+/{-} & N & W:1a \\
C3339 & 1a<o1L 1b*<o1 / 2a<1a0L / 3a<2a1L & IU/N/N & +{-}/+/+ & N & W:1a \\
C3340 & 1a<o1L 1b*<o1 / 2a<1a0L / 3a<2a2 & IU/N/N & +{-}/+/{-} & N & W:1a \\
C3341 & 1a<o1L 1b*<o1 / 2a<1a0L / 3a<2a2L & IU/N/N & +{-}/+/+ & N & W:1a \\
C3342 & 1a<o1L 1b*<o1 / 2a<1a1 / 3a<2a0L & IU/I/N & +{-}/{-}/{-} & N & W:1a,2a \\
C3343 & 1a<o1L 1b*<o1 / 2a<1a1 / 3a<2a1 & IU/I/I & +{-}/{-}/{-} & N & W:1a,2a,3a \\
C3344 & 1a<o1L 1b*<o1 / 2a<1a1 / 3a<2a1L & IU/I/I & +{-}/{-}/{-} & N & W:1a,2a,3a \\
C3345 & 1a<o1L 1b*<o1 / 2a<1a1 / 3a<2a2 & IU/I/I & +{-}/{-}/{-} & N & W:1a,2a,3a \\
C3346 & 1a<o1L 1b*<o1 / 2a<1a1 / 3a<2a2L & IU/I/I & +{-}/{-}/{-} & N & W:1a,2a,3a \\
C3347 & 1a<o1L 1b*<o1 / 2a<1a1L / 3a<2a0L & IU/I/N & +{-}/+/+ & N & W:1a,2a \\
C3348 & 1a<o1L 1b*<o1 / 2a<1a1L / 3a<2a1 & IU/I/I & +{-}/+/{-} & N & W:1a,2a,3a \\
C3349 & 1a<o1L 1b*<o1 / 2a<1a1L / 3a<2a1L & IU/I/I & +{-}/+/+ & N & W:1a,2a,3a \\
C3350 & 1a<o1L 1b*<o1 / 2a<1a1L / 3a<2a2 & IU/I/I & +{-}/+/{-} & N & W:1a,2a,3a \\
C3351 & 1a<o1L 1b*<o1 / 2a<1a1L / 3a<2a2L & IU/I/I & +{-}/+/+ & N & W:1a,2a,3a \\
C3352 & 1a<o1L 1b*<o1 / 2a<1a2 / 3a<2a0L & IU/I/N & +{-}/{-}/{-} & N & W:1a,2a \\
C3353 & 1a<o1L 1b*<o1 / 2a<1a2 / 3a<2a1 & IU/I/I & +{-}/{-}/{-} & N & W:1a,2a,3a \\
C3354 & 1a<o1L 1b*<o1 / 2a<1a2 / 3a<2a1L & IU/I/I & +{-}/{-}/{-} & N & W:1a,2a,3a \\
C3355 & 1a<o1L 1b*<o1 / 2a<1a2 / 3a<2a2 & IU/I/I & +{-}/{-}/{-} & N & W:1a,2a,3a \\
C3356 & 1a<o1L 1b*<o1 / 2a<1a2 / 3a<2a2L & IU/I/I & +{-}/{-}/{-} & N & W:1a,2a,3a \\
C3357 & 1a<o1L 1b*<o1 / 2a<1a2L / 3a<2a0L & IU/I/N & +{-}/+/+ & N & W:1a,2a \\
C3358 & 1a<o1L 1b*<o1 / 2a<1a2L / 3a<2a1 & IU/I/I & +{-}/+/{-} & N & W:1a,2a,3a \\
C3359 & 1a<o1L 1b*<o1 / 2a<1a2L / 3a<2a1L & IU/I/I & +{-}/+/+ & N & W:1a,2a,3a \\
C3360 & 1a<o1L 1b*<o1 / 2a<1a2L / 3a<2a2 & IU/I/I & +{-}/+/{-} & N & W:1a,2a,3a \\
C3361 & 1a<o1L 1b*<o1 / 2a<1a2L / 3a<2a2L & IU/I/I & +{-}/+/+ & N & W:1a,2a,3a \\
C3362 & 1a<o1L 1b*<o1 / 2a<1b1 / 3a<2a0L & IU/U/N & +{-}/{-}/{-} & N & W:1a \\
C3363 & 1a<o1L 1b*<o1 / 2a<1b1 / 3a<2a1 & IU/U/U & +{-}/{-}/{-} & N & W:1a \\
C3364 & 1a<o1L 1b*<o1 / 2a<1b1 / 3a<2a1L & IU/U/U & +{-}/{-}/{-} & N & W:1a \\
C3365 & 1a<o1L 1b*<o1 / 2a<1b1 / 3a<2a2 & IU/U/U & +{-}/{-}/{-} & N & W:1a \\
C3366 & 1a<o1L 1b*<o1 / 2a<1b1 / 3a<2a2L & IU/U/U & +{-}/{-}/{-} & N & W:1a \\
C3367 & 1a<o1L 1b*<o1 / 2a<1b2 / 3a<2a0L & IU/U/N & +{-}/{-}/{-} & N & W:1a \\
C3368 & 1a<o1L 1b*<o1 / 2a<1b2 / 3a<2a1 & IU/U/U & +{-}/{-}/{-} & N & W:1a \\
C3369 & 1a<o1L 1b*<o1 / 2a<1b2 / 3a<2a1L & IU/U/U & +{-}/{-}/{-} & N & W:1a \\
C3370 & 1a<o1L 1b*<o1 / 2a<1b2 / 3a<2a2 & IU/U/U & +{-}/{-}/{-} & N & W:1a \\
C3371 & 1a<o1L 1b*<o1 / 2a<1b2 / 3a<2a2L & IU/U/U & +{-}/{-}/{-} & N & W:1a \\
C3372 & 1a<o1L 1b*<o2 / 2a<1a0L / 3a<2a0L & IU/N/N & +{-}/+/+ & N & W:1a \\
C3373 & 1a<o1L 1b*<o2 / 2a<1a0L / 3a<2a1 & IU/N/N & +{-}/+/{-} & N & W:1a \\
C3374 & 1a<o1L 1b*<o2 / 2a<1a0L / 3a<2a1L & IU/N/N & +{-}/+/+ & N & W:1a \\
C3375 & 1a<o1L 1b*<o2 / 2a<1a0L / 3a<2a2 & IU/N/N & +{-}/+/{-} & N & W:1a \\
C3376 & 1a<o1L 1b*<o2 / 2a<1a0L / 3a<2a2L & IU/N/N & +{-}/+/+ & N & W:1a \\
C3377 & 1a<o1L 1b*<o2 / 2a<1a1 / 3a<2a0L & IU/I/N & +{-}/{-}/{-} & N & W:1a,2a \\
C3378 & 1a<o1L 1b*<o2 / 2a<1a1 / 3a<2a1 & IU/I/I & +{-}/{-}/{-} & N & W:1a,2a,3a \\
C3379 & 1a<o1L 1b*<o2 / 2a<1a1 / 3a<2a1L & IU/I/I & +{-}/{-}/{-} & N & W:1a,2a,3a \\
C3380 & 1a<o1L 1b*<o2 / 2a<1a1 / 3a<2a2 & IU/I/I & +{-}/{-}/{-} & N & W:1a,2a,3a \\
C3381 & 1a<o1L 1b*<o2 / 2a<1a1 / 3a<2a2L & IU/I/I & +{-}/{-}/{-} & N & W:1a,2a,3a \\
C3382 & 1a<o1L 1b*<o2 / 2a<1a1L / 3a<2a0L & IU/I/N & +{-}/+/+ & N & W:1a,2a \\
C3383 & 1a<o1L 1b*<o2 / 2a<1a1L / 3a<2a1 & IU/I/I & +{-}/+/{-} & N & W:1a,2a,3a \\
C3384 & 1a<o1L 1b*<o2 / 2a<1a1L / 3a<2a1L & IU/I/I & +{-}/+/+ & N & W:1a,2a,3a \\
C3385 & 1a<o1L 1b*<o2 / 2a<1a1L / 3a<2a2 & IU/I/I & +{-}/+/{-} & N & W:1a,2a,3a \\
C3386 & 1a<o1L 1b*<o2 / 2a<1a1L / 3a<2a2L & IU/I/I & +{-}/+/+ & N & W:1a,2a,3a \\
C3387 & 1a<o1L 1b*<o2 / 2a<1a2 / 3a<2a0L & IU/I/N & +{-}/{-}/{-} & N & W:1a,2a \\
C3388 & 1a<o1L 1b*<o2 / 2a<1a2 / 3a<2a1 & IU/I/I & +{-}/{-}/{-} & N & W:1a,2a,3a \\
C3389 & 1a<o1L 1b*<o2 / 2a<1a2 / 3a<2a1L & IU/I/I & +{-}/{-}/{-} & N & W:1a,2a,3a \\
C3390 & 1a<o1L 1b*<o2 / 2a<1a2 / 3a<2a2 & IU/I/I & +{-}/{-}/{-} & N & W:1a,2a,3a \\
C3391 & 1a<o1L 1b*<o2 / 2a<1a2 / 3a<2a2L & IU/I/I & +{-}/{-}/{-} & N & W:1a,2a,3a \\
C3392 & 1a<o1L 1b*<o2 / 2a<1a2L / 3a<2a0L & IU/I/N & +{-}/+/+ & N & W:1a,2a \\
C3393 & 1a<o1L 1b*<o2 / 2a<1a2L / 3a<2a1 & IU/I/I & +{-}/+/{-} & N & W:1a,2a,3a \\
C3394 & 1a<o1L 1b*<o2 / 2a<1a2L / 3a<2a1L & IU/I/I & +{-}/+/+ & N & W:1a,2a,3a \\
C3395 & 1a<o1L 1b*<o2 / 2a<1a2L / 3a<2a2 & IU/I/I & +{-}/+/{-} & N & W:1a,2a,3a \\
C3396 & 1a<o1L 1b*<o2 / 2a<1a2L / 3a<2a2L & IU/I/I & +{-}/+/+ & N & W:1a,2a,3a \\
C3397 & 1a<o1L 1b*<o2 / 2a<1b1 / 3a<2a0L & IU/U/N & +{-}/{-}/{-} & N & W:1a \\
C3398 & 1a<o1L 1b*<o2 / 2a<1b1 / 3a<2a1 & IU/U/U & +{-}/{-}/{-} & N & W:1a \\
C3399 & 1a<o1L 1b*<o2 / 2a<1b1 / 3a<2a1L & IU/U/U & +{-}/{-}/{-} & N & W:1a \\
C3400 & 1a<o1L 1b*<o2 / 2a<1b1 / 3a<2a2 & IU/U/U & +{-}/{-}/{-} & N & W:1a \\
C3401 & 1a<o1L 1b*<o2 / 2a<1b1 / 3a<2a2L & IU/U/U & +{-}/{-}/{-} & N & W:1a \\
C3402 & 1a<o1L 1b*<o2 / 2a<1b2 / 3a<2a0L & IU/U/N & +{-}/{-}/{-} & N & W:1a \\
C3403 & 1a<o1L 1b*<o2 / 2a<1b2 / 3a<2a1 & IU/U/U & +{-}/{-}/{-} & N & W:1a \\
C3404 & 1a<o1L 1b*<o2 / 2a<1b2 / 3a<2a1L & IU/U/U & +{-}/{-}/{-} & N & W:1a \\
C3405 & 1a<o1L 1b*<o2 / 2a<1b2 / 3a<2a2 & IU/U/U & +{-}/{-}/{-} & N & W:1a \\
C3406 & 1a<o1L 1b*<o2 / 2a<1b2 / 3a<2a2L & IU/U/U & +{-}/{-}/{-} & N & W:1a \\
C3407 & 1a<o2 / 2a<1a0L / 3a<2a0L 3b*<2a1 & I/N/NN & {-}/{-}/{-}{-} & N &  \\
C3408 & 1a<o2 / 2a<1a0L / 3a<2a0L 3b*<2a2 & I/N/NN & {-}/{-}/{-}{-} & N &  \\
C3409 & 1a<o2 / 2a<1a0L / 3a<2a1 3b*<2a1 & I/N/NN & {-}/{-}/{-}{-} & N &  \\
C3410 & 1a<o2 / 2a<1a0L / 3a<2a1 3b*<2a2 & I/N/NN & {-}/{-}/{-}{-} & N &  \\
C3411 & 1a<o2 / 2a<1a0L / 3a<2a1L 3b*<2a1 & I/N/NN & {-}/{-}/{-}{-} & N &  \\
C3412 & 1a<o2 / 2a<1a0L / 3a<2a1L 3b*<2a2 & I/N/NN & {-}/{-}/{-}{-} & N &  \\
C3413 & 1a<o2 / 2a<1a0L / 3a<2a2 3b*<2a1 & I/N/NN & {-}/{-}/{-}{-} & N &  \\
C3414 & 1a<o2 / 2a<1a0L / 3a<2a2 3b*<2a2 & I/N/NN & {-}/{-}/{-}{-} & N &  \\
C3415 & 1a<o2 / 2a<1a0L / 3a<2a2L 3b*<2a1 & I/N/NN & {-}/{-}/{-}{-} & N &  \\
C3416 & 1a<o2 / 2a<1a0L / 3a<2a2L 3b*<2a2 & I/N/NN & {-}/{-}/{-}{-} & N &  \\
C3417 & 1a<o2 / 2a<1a0L 2b<1a1 / 3a*<2a1 & I/NI/N & {-}/{-}{-}/{-} & N &  \\
C3418 & 1a<o2 / 2a<1a0L 2b<1a1 / 3a*<2a2 & I/NI/N & {-}/{-}{-}/{-} & N &  \\
C3419 & 1a<o2 / 2a<1a0L 2b<1a1 / 3a*<2b1 & I/NI/U & {-}/{-}{-}/{-} & N & W:3a \\
C3420 & 1a<o2 / 2a<1a0L 2b<1a1 / 3a*<2b2 & I/NI/U & {-}/{-}{-}/{-} & N & W:3a \\
C3421 & 1a<o2 / 2a<1a0L 2b<1a2 / 3a*<2a1 & I/NI/N & {-}/{-}{-}/{-} & N &  \\
C3422 & 1a<o2 / 2a<1a0L 2b<1a2 / 3a*<2a2 & I/NI/N & {-}/{-}{-}/{-} & N &  \\
C3423 & 1a<o2 / 2a<1a0L 2b<1a2 / 3a*<2b1 & I/NI/U & {-}/{-}{-}/{-} & N & W:3a \\
C3424 & 1a<o2 / 2a<1a0L 2b<1a2 / 3a*<2b2 & I/NI/U & {-}/{-}{-}/{-} & N & W:3a \\
C3425 & 1a<o2 / 2a<1a0L 2b*<1a1 / 3a<2a0L & I/NU/N & {-}/{-}{-}/{-} & N & W:2b \\
C3426 & 1a<o2 / 2a<1a0L 2b*<1a1 / 3a<2a1 & I/NU/N & {-}/{-}{-}/{-} & N & W:2b \\
C3427 & 1a<o2 / 2a<1a0L 2b*<1a1 / 3a<2a1L & I/NU/N & {-}/{-}{-}/{-} & N & W:2b \\
C3428 & 1a<o2 / 2a<1a0L 2b*<1a1 / 3a<2a2 & I/NU/N & {-}/{-}{-}/{-} & N & W:2b \\
C3429 & 1a<o2 / 2a<1a0L 2b*<1a1 / 3a<2a2L & I/NU/N & {-}/{-}{-}/{-} & N & W:2b \\
C3430 & 1a<o2 / 2a<1a0L 2b*<1a1 / 3a<2b1 & I/NU/U & {-}/{-}{-}/{-} & N & W:2b,3a \\
C3431 & 1a<o2 / 2a<1a0L 2b*<1a1 / 3a<2b2 & I/NU/U & {-}/{-}{-}/{-} & N & W:2b,3a \\
C3432 & 1a<o2 / 2a<1a0L 2b*<1a2 / 3a<2a0L & I/NU/N & {-}/{-}{-}/{-} & N & W:2b \\
C3433 & 1a<o2 / 2a<1a0L 2b*<1a2 / 3a<2a1 & I/NU/N & {-}/{-}{-}/{-} & N & W:2b \\
C3434 & 1a<o2 / 2a<1a0L 2b*<1a2 / 3a<2a1L & I/NU/N & {-}/{-}{-}/{-} & N & W:2b \\
C3435 & 1a<o2 / 2a<1a0L 2b*<1a2 / 3a<2a2 & I/NU/N & {-}/{-}{-}/{-} & N & W:2b \\
C3436 & 1a<o2 / 2a<1a0L 2b*<1a2 / 3a<2a2L & I/NU/N & {-}/{-}{-}/{-} & N & W:2b \\
C3437 & 1a<o2 / 2a<1a0L 2b*<1a2 / 3a<2b1 & I/NU/U & {-}/{-}{-}/{-} & N & W:2b,3a \\
C3438 & 1a<o2 / 2a<1a0L 2b*<1a2 / 3a<2b2 & I/NU/U & {-}/{-}{-}/{-} & N & W:2b,3a \\
C3439 & 1a<o2 / 2a<1a1 / 3a<2a0L 3b*<2a1 & I/I/NU & {-}/{-}/{-}{-} & N & W:3b \\
C3440 & 1a<o2 / 2a<1a1 / 3a<2a0L 3b*<2a2 & I/I/NU & {-}/{-}/{-}{-} & N & W:3b \\
C3441 & 1a<o2 / 2a<1a1 / 3a<2a1 3b*<2a1 & I/I/IU & {-}/{-}/{-}{-} & N & W:1a,2a,3a \\
C3442 & 1a<o2 / 2a<1a1 / 3a<2a1 3b*<2a2 & I/I/IU & {-}/{-}/{-}{-} & N & W:1a,2a,3a \\
C3443 & 1a<o2 / 2a<1a1 / 3a<2a1L 3b*<2a1 & I/I/IU & {-}/{-}/{-}{-} & N & W:1a,2a,3a \\
C3444 & 1a<o2 / 2a<1a1 / 3a<2a1L 3b*<2a2 & I/I/IU & {-}/{-}/{-}{-} & N & W:1a,2a,3a \\
C3445 & 1a<o2 / 2a<1a1 / 3a<2a2 3b*<2a1 & I/I/IU & {-}/{-}/{-}{-} & N & W:1a,2a,3a \\
C3446 & 1a<o2 / 2a<1a1 / 3a<2a2 3b*<2a2 & I/I/IU & {-}/{-}/{-}{-} & N & W:1a,2a,3a \\
C3447 & 1a<o2 / 2a<1a1 / 3a<2a2L 3b*<2a1 & I/I/IU & {-}/{-}/{-}{-} & N & W:1a,2a,3a \\
C3448 & 1a<o2 / 2a<1a1 / 3a<2a2L 3b*<2a2 & I/I/IU & {-}/{-}/{-}{-} & N & W:1a,2a,3a \\
C3449 & 1a<o2 / 2a<1a1 2b<1a1 / 3a*<2a1 & I/EE/U & {-}/{-}{-}/{-} & E & W:1a \\
C3450 & 1a<o2 / 2a<1a1 2b<1a1 / 3a*<2a2 & I/EE/U & {-}/{-}{-}/{-} & E & W:1a \\
C3451 & 1a<o2 / 2a<1a1 2b<1a1L / 3a*<2a1 & I/LI/U & {-}/{-}{-}/{-} & B & CP:2b W:1a,2b \\
C3452 & 1a<o2 / 2a<1a1 2b<1a1L / 3a*<2a2 & I/LI/U & {-}/{-}{-}/{-} & B & CP:2b W:1a,2b \\
C3453 & 1a<o2 / 2a<1a1 2b<1a1L / 3a*<2b1 & I/LI/U & {-}/{-}{-}/{-} & B & CP:2b W:1a,2b \\
C3454 & 1a<o2 / 2a<1a1 2b<1a1L / 3a*<2b2 & I/LI/U & {-}/{-}{-}/{-} & B & CP:2b W:1a,2b \\
C3455 & 1a<o2 / 2a<1a1 2b<1a2 / 3a*<2a1 & I/EE/U & {-}/{-}{-}/{-} & E & C:2b W:1a \\
C3456 & 1a<o2 / 2a<1a1 2b<1a2 / 3a*<2a2 & I/EE/U & {-}/{-}{-}/{-} & E & C:2b W:1a \\
C3457 & 1a<o2 / 2a<1a1 2b<1a2 / 3a*<2b1 & I/EE/U & {-}/{-}{-}/{-} & E & C:2b W:1a \\
C3458 & 1a<o2 / 2a<1a1 2b<1a2 / 3a*<2b2 & I/EE/U & {-}/{-}{-}/{-} & E & C:2b W:1a \\
C3459 & 1a<o2 / 2a<1a1 2b<1a2L / 3a*<2a1 & I/LI/U & {-}/{-}{-}/{-} & B & P:2b W:1a,2b \\
C3460 & 1a<o2 / 2a<1a1 2b<1a2L / 3a*<2a2 & I/LI/U & {-}/{-}{-}/{-} & B & P:2b W:1a,2b \\
C3461 & 1a<o2 / 2a<1a1 2b<1a2L / 3a*<2b1 & I/LI/U & {-}/{-}{-}/{-} & B & P:2b W:1a,2b \\
C3462 & 1a<o2 / 2a<1a1 2b<1a2L / 3a*<2b2 & I/LI/U & {-}/{-}{-}/{-} & B & P:2b W:1a,2b \\
C3463 & 1a<o2 / 2a<1a1 2b*<1a1 / 3a<2a0L & I/IU/N & {-}/{-}{-}/{-} & N & W:1a,2a \\
C3464 & 1a<o2 / 2a<1a1 2b*<1a1 / 3a<2a1 & I/IU/I & {-}/{-}{-}/{-} & N & W:1a,2a,3a \\
C3465 & 1a<o2 / 2a<1a1 2b*<1a1 / 3a<2a1L & I/IU/I & {-}/{-}{-}/{-} & N & W:1a,2a,3a \\
C3466 & 1a<o2 / 2a<1a1 2b*<1a1 / 3a<2a2 & I/IU/I & {-}/{-}{-}/{-} & N & W:1a,2a,3a \\
C3467 & 1a<o2 / 2a<1a1 2b*<1a1 / 3a<2a2L & I/IU/I & {-}/{-}{-}/{-} & N & W:1a,2a,3a \\
C3468 & 1a<o2 / 2a<1a1 2b*<1a1 / 3a<2b1 & I/IU/U & {-}/{-}{-}/{-} & N & W:1a,2a \\
C3469 & 1a<o2 / 2a<1a1 2b*<1a1 / 3a<2b2 & I/IU/U & {-}/{-}{-}/{-} & N & W:1a,2a \\
C3470 & 1a<o2 / 2a<1a1 2b*<1a2 / 3a<2a0L & I/IU/N & {-}/{-}{-}/{-} & N & W:1a,2a \\
C3471 & 1a<o2 / 2a<1a1 2b*<1a2 / 3a<2a1 & I/IU/I & {-}/{-}{-}/{-} & N & W:1a,2a,3a \\
C3472 & 1a<o2 / 2a<1a1 2b*<1a2 / 3a<2a1L & I/IU/I & {-}/{-}{-}/{-} & N & W:1a,2a,3a \\
C3473 & 1a<o2 / 2a<1a1 2b*<1a2 / 3a<2a2 & I/IU/I & {-}/{-}{-}/{-} & N & W:1a,2a,3a \\
C3474 & 1a<o2 / 2a<1a1 2b*<1a2 / 3a<2a2L & I/IU/I & {-}/{-}{-}/{-} & N & W:1a,2a,3a \\
C3475 & 1a<o2 / 2a<1a1 2b*<1a2 / 3a<2b1 & I/IU/U & {-}/{-}{-}/{-} & N & W:1a,2a \\
C3476 & 1a<o2 / 2a<1a1 2b*<1a2 / 3a<2b2 & I/IU/U & {-}/{-}{-}/{-} & N & W:1a,2a \\
C3477 & 1a<o2 / 2a<1a1L / 3a<2a0L 3b*<2a1 & I/I/NU & {-}/{-}/{-}{-} & N & W:3b \\
C3478 & 1a<o2 / 2a<1a1L / 3a<2a0L 3b*<2a2 & I/I/NU & {-}/{-}/{-}{-} & N & W:3b \\
C3479 & 1a<o2 / 2a<1a1L / 3a<2a1 3b*<2a1 & I/I/IU & {-}/{-}/{-}{-} & N & W:1a,2a,3a \\
C3480 & 1a<o2 / 2a<1a1L / 3a<2a1 3b*<2a2 & I/I/IU & {-}/{-}/{-}{-} & N & W:1a,2a,3a \\
C3481 & 1a<o2 / 2a<1a1L / 3a<2a1L 3b*<2a1 & I/I/IU & {-}/{-}/{-}{-} & N & W:1a,2a,3a \\
C3482 & 1a<o2 / 2a<1a1L / 3a<2a1L 3b*<2a2 & I/I/IU & {-}/{-}/{-}{-} & N & W:1a,2a,3a \\
C3483 & 1a<o2 / 2a<1a1L / 3a<2a2 3b*<2a1 & I/I/IU & {-}/{-}/{-}{-} & N & W:1a,2a,3a \\
C3484 & 1a<o2 / 2a<1a1L / 3a<2a2 3b*<2a2 & I/I/IU & {-}/{-}/{-}{-} & N & W:1a,2a,3a \\
C3485 & 1a<o2 / 2a<1a1L / 3a<2a2L 3b*<2a1 & I/I/IU & {-}/{-}/{-}{-} & N & W:1a,2a,3a \\
C3486 & 1a<o2 / 2a<1a1L / 3a<2a2L 3b*<2a2 & I/I/IU & {-}/{-}/{-}{-} & N & W:1a,2a,3a \\
C3487 & 1a<o2 / 2a<1a1L 2b<1a2 / 3a*<2a1 & I/IL/U & {-}/{-}{-}/{-} & B & C:2a,2b P:2a W:1a,2a \\
C3488 & 1a<o2 / 2a<1a1L 2b<1a2 / 3a*<2a2 & I/IL/U & {-}/{-}{-}/{-} & B & C:2a,2b P:2a W:1a,2a \\
C3489 & 1a<o2 / 2a<1a1L 2b<1a2 / 3a*<2b1 & I/IL/U & {-}/{-}{-}/{-} & B & C:2a,2b P:2a W:1a,2a \\
C3490 & 1a<o2 / 2a<1a1L 2b<1a2 / 3a*<2b2 & I/IL/U & {-}/{-}{-}/{-} & B & C:2a,2b P:2a W:1a,2a \\
C3491 & 1a<o2 / 2a<1a1L 2b*<1a1 / 3a<2a0L & I/IU/N & {-}/{-}{-}/{-} & N & W:1a,2a \\
C3492 & 1a<o2 / 2a<1a1L 2b*<1a1 / 3a<2a1 & I/IU/I & {-}/{-}{-}/{-} & N & W:1a,2a,3a \\
C3493 & 1a<o2 / 2a<1a1L 2b*<1a1 / 3a<2a1L & I/IU/I & {-}/{-}{-}/{-} & N & W:1a,2a,3a \\
C3494 & 1a<o2 / 2a<1a1L 2b*<1a1 / 3a<2a2 & I/IU/I & {-}/{-}{-}/{-} & N & W:1a,2a,3a \\
C3495 & 1a<o2 / 2a<1a1L 2b*<1a1 / 3a<2a2L & I/IU/I & {-}/{-}{-}/{-} & N & W:1a,2a,3a \\
C3496 & 1a<o2 / 2a<1a1L 2b*<1a1 / 3a<2b1 & I/IU/U & {-}/{-}{-}/{-} & N & W:1a,2a \\
C3497 & 1a<o2 / 2a<1a1L 2b*<1a1 / 3a<2b2 & I/IU/U & {-}/{-}{-}/{-} & N & W:1a,2a \\
C3498 & 1a<o2 / 2a<1a1L 2b*<1a2 / 3a<2a0L & I/IU/N & {-}/{-}{-}/{-} & N & W:1a,2a \\
C3499 & 1a<o2 / 2a<1a1L 2b*<1a2 / 3a<2a1 & I/IU/I & {-}/{-}{-}/{-} & N & W:1a,2a,3a \\
C3500 & 1a<o2 / 2a<1a1L 2b*<1a2 / 3a<2a1L & I/IU/I & {-}/{-}{-}/{-} & N & W:1a,2a,3a \\
C3501 & 1a<o2 / 2a<1a1L 2b*<1a2 / 3a<2a2 & I/IU/I & {-}/{-}{-}/{-} & N & W:1a,2a,3a \\
C3502 & 1a<o2 / 2a<1a1L 2b*<1a2 / 3a<2a2L & I/IU/I & {-}/{-}{-}/{-} & N & W:1a,2a,3a \\
C3503 & 1a<o2 / 2a<1a1L 2b*<1a2 / 3a<2b1 & I/IU/U & {-}/{-}{-}/{-} & N & W:1a,2a \\
C3504 & 1a<o2 / 2a<1a1L 2b*<1a2 / 3a<2b2 & I/IU/U & {-}/{-}{-}/{-} & N & W:1a,2a \\
C3505 & 1a<o2 / 2a<1a2 / 3a<2a0L 3b*<2a1 & I/I/NU & {-}/{-}/{-}{-} & N & W:3b \\
C3506 & 1a<o2 / 2a<1a2 / 3a<2a0L 3b*<2a2 & I/I/NU & {-}/{-}/{-}{-} & N & W:3b \\
C3507 & 1a<o2 / 2a<1a2 / 3a<2a1 3b*<2a1 & I/I/IU & {-}/{-}/{-}{-} & N & W:1a,2a,3a \\
C3508 & 1a<o2 / 2a<1a2 / 3a<2a1 3b*<2a2 & I/I/IU & {-}/{-}/{-}{-} & N & W:1a,2a,3a \\
C3509 & 1a<o2 / 2a<1a2 / 3a<2a1L 3b*<2a1 & I/I/IU & {-}/{-}/{-}{-} & N & W:1a,2a,3a \\
C3510 & 1a<o2 / 2a<1a2 / 3a<2a1L 3b*<2a2 & I/I/IU & {-}/{-}/{-}{-} & N & W:1a,2a,3a \\
C3511 & 1a<o2 / 2a<1a2 / 3a<2a2 3b*<2a1 & I/I/IU & {-}/{-}/{-}{-} & N & W:1a,2a,3a \\
C3512 & 1a<o2 / 2a<1a2 / 3a<2a2 3b*<2a2 & I/I/IU & {-}/{-}/{-}{-} & N & W:1a,2a,3a \\
C3513 & 1a<o2 / 2a<1a2 / 3a<2a2L 3b*<2a1 & I/I/IU & {-}/{-}/{-}{-} & N & W:1a,2a,3a \\
C3514 & 1a<o2 / 2a<1a2 / 3a<2a2L 3b*<2a2 & I/I/IU & {-}/{-}/{-}{-} & N & W:1a,2a,3a \\
C3515 & 1a<o2 / 2a<1a2 2b<1a2 / 3a*<2a1 & I/EE/U & {-}/{-}{-}/{-} & E & W:1a \\
C3516 & 1a<o2 / 2a<1a2 2b<1a2 / 3a*<2a2 & I/EE/U & {-}/{-}{-}/{-} & E & W:1a \\
C3517 & 1a<o2 / 2a<1a2 2b<1a2L / 3a*<2a1 & I/LI/U & {-}/{-}{-}/{-} & B & CP:2b W:1a,2b \\
C3518 & 1a<o2 / 2a<1a2 2b<1a2L / 3a*<2a2 & I/LI/U & {-}/{-}{-}/{-} & B & CP:2b W:1a,2b \\
C3519 & 1a<o2 / 2a<1a2 2b<1a2L / 3a*<2b1 & I/LI/U & {-}/{-}{-}/{-} & B & CP:2b W:1a,2b \\
C3520 & 1a<o2 / 2a<1a2 2b<1a2L / 3a*<2b2 & I/LI/U & {-}/{-}{-}/{-} & B & CP:2b W:1a,2b \\
C3521 & 1a<o2 / 2a<1a2 2b*<1a1 / 3a<2a0L & I/IU/N & {-}/{-}{-}/{-} & N & W:1a,2a \\
C3522 & 1a<o2 / 2a<1a2 2b*<1a1 / 3a<2a1 & I/IU/I & {-}/{-}{-}/{-} & N & W:1a,2a,3a \\
C3523 & 1a<o2 / 2a<1a2 2b*<1a1 / 3a<2a1L & I/IU/I & {-}/{-}{-}/{-} & N & W:1a,2a,3a \\
C3524 & 1a<o2 / 2a<1a2 2b*<1a1 / 3a<2a2 & I/IU/I & {-}/{-}{-}/{-} & N & W:1a,2a,3a \\
C3525 & 1a<o2 / 2a<1a2 2b*<1a1 / 3a<2a2L & I/IU/I & {-}/{-}{-}/{-} & N & W:1a,2a,3a \\
C3526 & 1a<o2 / 2a<1a2 2b*<1a1 / 3a<2b1 & I/IU/U & {-}/{-}{-}/{-} & N & W:1a,2a \\
C3527 & 1a<o2 / 2a<1a2 2b*<1a1 / 3a<2b2 & I/IU/U & {-}/{-}{-}/{-} & N & W:1a,2a \\
C3528 & 1a<o2 / 2a<1a2 2b*<1a2 / 3a<2a0L & I/IU/N & {-}/{-}{-}/{-} & N & W:1a,2a \\
C3529 & 1a<o2 / 2a<1a2 2b*<1a2 / 3a<2a1 & I/IU/I & {-}/{-}{-}/{-} & N & W:1a,2a,3a \\
C3530 & 1a<o2 / 2a<1a2 2b*<1a2 / 3a<2a1L & I/IU/I & {-}/{-}{-}/{-} & N & W:1a,2a,3a \\
C3531 & 1a<o2 / 2a<1a2 2b*<1a2 / 3a<2a2 & I/IU/I & {-}/{-}{-}/{-} & N & W:1a,2a,3a \\
C3532 & 1a<o2 / 2a<1a2 2b*<1a2 / 3a<2a2L & I/IU/I & {-}/{-}{-}/{-} & N & W:1a,2a,3a \\
C3533 & 1a<o2 / 2a<1a2 2b*<1a2 / 3a<2b1 & I/IU/U & {-}/{-}{-}/{-} & N & W:1a,2a \\
C3534 & 1a<o2 / 2a<1a2 2b*<1a2 / 3a<2b2 & I/IU/U & {-}/{-}{-}/{-} & N & W:1a,2a \\
C3535 & 1a<o2 / 2a<1a2L / 3a<2a0L 3b*<2a1 & I/I/NU & {-}/{-}/{-}{-} & N & W:3b \\
C3536 & 1a<o2 / 2a<1a2L / 3a<2a0L 3b*<2a2 & I/I/NU & {-}/{-}/{-}{-} & N & W:3b \\
C3537 & 1a<o2 / 2a<1a2L / 3a<2a1 3b*<2a1 & I/I/IU & {-}/{-}/{-}{-} & N & W:1a,2a,3a \\
C3538 & 1a<o2 / 2a<1a2L / 3a<2a1 3b*<2a2 & I/I/IU & {-}/{-}/{-}{-} & N & W:1a,2a,3a \\
C3539 & 1a<o2 / 2a<1a2L / 3a<2a1L 3b*<2a1 & I/I/IU & {-}/{-}/{-}{-} & N & W:1a,2a,3a \\
C3540 & 1a<o2 / 2a<1a2L / 3a<2a1L 3b*<2a2 & I/I/IU & {-}/{-}/{-}{-} & N & W:1a,2a,3a \\
C3541 & 1a<o2 / 2a<1a2L / 3a<2a2 3b*<2a1 & I/I/IU & {-}/{-}/{-}{-} & N & W:1a,2a,3a \\
C3542 & 1a<o2 / 2a<1a2L / 3a<2a2 3b*<2a2 & I/I/IU & {-}/{-}/{-}{-} & N & W:1a,2a,3a \\
C3543 & 1a<o2 / 2a<1a2L / 3a<2a2L 3b*<2a1 & I/I/IU & {-}/{-}/{-}{-} & N & W:1a,2a,3a \\
C3544 & 1a<o2 / 2a<1a2L / 3a<2a2L 3b*<2a2 & I/I/IU & {-}/{-}/{-}{-} & N & W:1a,2a,3a \\
C3545 & 1a<o2 / 2a<1a2L 2b*<1a1 / 3a<2a0L & I/IU/N & {-}/{-}{-}/{-} & N & W:1a,2a \\
C3546 & 1a<o2 / 2a<1a2L 2b*<1a1 / 3a<2a1 & I/IU/I & {-}/{-}{-}/{-} & N & W:1a,2a,3a \\
C3547 & 1a<o2 / 2a<1a2L 2b*<1a1 / 3a<2a1L & I/IU/I & {-}/{-}{-}/{-} & N & W:1a,2a,3a \\
C3548 & 1a<o2 / 2a<1a2L 2b*<1a1 / 3a<2a2 & I/IU/I & {-}/{-}{-}/{-} & N & W:1a,2a,3a \\
C3549 & 1a<o2 / 2a<1a2L 2b*<1a1 / 3a<2a2L & I/IU/I & {-}/{-}{-}/{-} & N & W:1a,2a,3a \\
C3550 & 1a<o2 / 2a<1a2L 2b*<1a1 / 3a<2b1 & I/IU/U & {-}/{-}{-}/{-} & N & W:1a,2a \\
C3551 & 1a<o2 / 2a<1a2L 2b*<1a1 / 3a<2b2 & I/IU/U & {-}/{-}{-}/{-} & N & W:1a,2a \\
C3552 & 1a<o2 / 2a<1a2L 2b*<1a2 / 3a<2a0L & I/IU/N & {-}/{-}{-}/{-} & N & W:1a,2a \\
C3553 & 1a<o2 / 2a<1a2L 2b*<1a2 / 3a<2a1 & I/IU/I & {-}/{-}{-}/{-} & N & W:1a,2a,3a \\
C3554 & 1a<o2 / 2a<1a2L 2b*<1a2 / 3a<2a1L & I/IU/I & {-}/{-}{-}/{-} & N & W:1a,2a,3a \\
C3555 & 1a<o2 / 2a<1a2L 2b*<1a2 / 3a<2a2 & I/IU/I & {-}/{-}{-}/{-} & N & W:1a,2a,3a \\
C3556 & 1a<o2 / 2a<1a2L 2b*<1a2 / 3a<2a2L & I/IU/I & {-}/{-}{-}/{-} & N & W:1a,2a,3a \\
C3557 & 1a<o2 / 2a<1a2L 2b*<1a2 / 3a<2b1 & I/IU/U & {-}/{-}{-}/{-} & N & W:1a,2a \\
C3558 & 1a<o2 / 2a<1a2L 2b*<1a2 / 3a<2b2 & I/IU/U & {-}/{-}{-}/{-} & N & W:1a,2a \\
C3559 & 1a<o2 / 2a*<1a1 / 3a<2a1 3b<2a1 & I/U/UU & {-}/{-}/{-}{-} & N & W:1a \\
C3560 & 1a<o2 / 2a*<1a1 / 3a<2a1 3b<2a2 & I/U/UU & {-}/{-}/{-}{-} & N & W:1a \\
C3561 & 1a<o2 / 2a*<1a1 / 3a<2a2 3b<2a2 & I/U/UU & {-}/{-}/{-}{-} & N & W:1a \\
C3562 & 1a<o2 / 2a*<1a2 / 3a<2a1 3b<2a1 & I/U/UU & {-}/{-}/{-}{-} & N & W:1a \\
C3563 & 1a<o2 / 2a*<1a2 / 3a<2a1 3b<2a2 & I/U/UU & {-}/{-}/{-}{-} & N & W:1a \\
C3564 & 1a<o2 / 2a*<1a2 / 3a<2a2 3b<2a2 & I/U/UU & {-}/{-}/{-}{-} & N & W:1a \\
C3565 & 1a<o2 1b<o2 / 2a<1a0L / 3a*<2a1 & EE/N/N & {-}{-}/{-}/{-} & E &  \\
C3566 & 1a<o2 1b<o2 / 2a<1a0L / 3a*<2a2 & EE/N/N & {-}{-}/{-}/{-} & E &  \\
C3567 & 1a<o2 1b<o2 / 2a<1a1 / 3a*<2a1 & EE/D/U & {-}{-}/{-}/{-} & E &  \\
C3568 & 1a<o2 1b<o2 / 2a<1a1 / 3a*<2a2 & EE/D/U & {-}{-}/{-}/{-} & E &  \\
C3569 & 1a<o2 1b<o2 / 2a<1a1L / 3a*<2a1 & EE/D/U & {-}{-}/{-}/{-} & E &  \\
C3570 & 1a<o2 1b<o2 / 2a<1a1L / 3a*<2a2 & EE/D/U & {-}{-}/{-}/{-} & E &  \\
C3571 & 1a<o2 1b<o2 / 2a<1a2 / 3a*<2a1 & EE/D/U & {-}{-}/{-}/{-} & E &  \\
C3572 & 1a<o2 1b<o2 / 2a<1a2 / 3a*<2a2 & EE/D/U & {-}{-}/{-}/{-} & E &  \\
C3573 & 1a<o2 1b<o2 / 2a<1a2L / 3a*<2a1 & EE/D/U & {-}{-}/{-}/{-} & E &  \\
C3574 & 1a<o2 1b<o2 / 2a<1a2L / 3a*<2a2 & EE/D/U & {-}{-}/{-}/{-} & E &  \\
C3575 & 1a<o2 1b<o2 / 2a*<1a1 / 3a<2a1 & EE/U/U & {-}{-}/{-}/{-} & E &  \\
C3576 & 1a<o2 1b<o2 / 2a*<1a1 / 3a<2a2 & EE/U/U & {-}{-}/{-}/{-} & E &  \\
C3577 & 1a<o2 1b<o2 / 2a*<1a2 / 3a<2a1 & EE/U/U & {-}{-}/{-}/{-} & E &  \\
C3578 & 1a<o2 1b<o2 / 2a*<1a2 / 3a<2a2 & EE/U/U & {-}{-}/{-}/{-} & E &  \\
C3579 & 1a<o2 1b<o2L / 2a<1a0L / 3a*<2a1 & LI/N/N & {-}+/{-}/{-} & B & CPW:1b \\
C3580 & 1a<o2 1b<o2L / 2a<1a0L / 3a*<2a2 & LI/N/N & {-}+/{-}/{-} & B & CPW:1b \\
C3581 & 1a<o2 1b<o2L / 2a<1a1 / 3a*<2a1 & LI/D/U & {-}+/{-}/{-} & B & CPW:1b \\
C3582 & 1a<o2 1b<o2L / 2a<1a1 / 3a*<2a2 & LI/D/U & {-}+/{-}/{-} & B & CPW:1b \\
C3583 & 1a<o2 1b<o2L / 2a<1a1L / 3a*<2a1 & LI/D/U & {-}+/{-}/{-} & B & CPW:1b \\
C3584 & 1a<o2 1b<o2L / 2a<1a1L / 3a*<2a2 & LI/D/U & {-}+/{-}/{-} & B & CPW:1b \\
C3585 & 1a<o2 1b<o2L / 2a<1a2 / 3a*<2a1 & LI/D/U & {-}+/{-}/{-} & B & CPW:1b \\
C3586 & 1a<o2 1b<o2L / 2a<1a2 / 3a*<2a2 & LI/D/U & {-}+/{-}/{-} & B & CPW:1b \\
C3587 & 1a<o2 1b<o2L / 2a<1a2L / 3a*<2a1 & LI/D/U & {-}+/{-}/{-} & B & CPW:1b \\
C3588 & 1a<o2 1b<o2L / 2a<1a2L / 3a*<2a2 & LI/D/U & {-}+/{-}/{-} & B & CPW:1b \\
C3589 & 1a<o2 1b<o2L / 2a<1b0L / 3a*<2a1 & LI/N/N & {-}+/+/{-} & B & CPW:1b \\
C3590 & 1a<o2 1b<o2L / 2a<1b0L / 3a*<2a2 & LI/N/N & {-}+/+/{-} & B & CPW:1b \\
C3591 & 1a<o2 1b<o2L / 2a<1b1 / 3a*<2a1 & LI/I/U & {-}+/{-}/{-} & B & CPW:1b,2a \\
C3592 & 1a<o2 1b<o2L / 2a<1b1 / 3a*<2a2 & LI/I/U & {-}+/{-}/{-} & B & CPW:1b,2a \\
C3593 & 1a<o2 1b<o2L / 2a<1b1L / 3a*<2a1 & LI/I/U & {-}+/+/{-} & B & CPW:1b,2a \\
C3594 & 1a<o2 1b<o2L / 2a<1b1L / 3a*<2a2 & LI/I/U & {-}+/+/{-} & B & CPW:1b,2a \\
C3595 & 1a<o2 1b<o2L / 2a<1b2 / 3a*<2a1 & LI/I/U & {-}+/{-}/{-} & B & CPW:1b,2a \\
C3596 & 1a<o2 1b<o2L / 2a<1b2 / 3a*<2a2 & LI/I/U & {-}+/{-}/{-} & B & CPW:1b,2a \\
C3597 & 1a<o2 1b<o2L / 2a<1b2L / 3a*<2a1 & LI/I/U & {-}+/+/{-} & B & CPW:1b,2a \\
C3598 & 1a<o2 1b<o2L / 2a<1b2L / 3a*<2a2 & LI/I/U & {-}+/+/{-} & B & CPW:1b,2a \\
C3599 & 1a<o2 1b<o2L / 2a*<1a1 / 3a<2a1 & LI/U/U & {-}+/{-}/{-} & B & CPW:1b \\
C3600 & 1a<o2 1b<o2L / 2a*<1a1 / 3a<2a2 & LI/U/U & {-}+/{-}/{-} & B & CPW:1b \\
C3601 & 1a<o2 1b<o2L / 2a*<1a2 / 3a<2a1 & LI/U/U & {-}+/{-}/{-} & B & CPW:1b \\
C3602 & 1a<o2 1b<o2L / 2a*<1a2 / 3a<2a2 & LI/U/U & {-}+/{-}/{-} & B & CPW:1b \\
C3603 & 1a<o2 1b<o2L / 2a*<1b1 / 3a<2a1 & LI/U/U & {-}+/{-}/{-} & B & CPW:1b \\
C3604 & 1a<o2 1b<o2L / 2a*<1b1 / 3a<2a2 & LI/U/U & {-}+/{-}/{-} & B & CPW:1b \\
C3605 & 1a<o2 1b<o2L / 2a*<1b2 / 3a<2a1 & LI/U/U & {-}+/{-}/{-} & B & CPW:1b \\
C3606 & 1a<o2 1b<o2L / 2a*<1b2 / 3a<2a2 & LI/U/U & {-}+/{-}/{-} & B & CPW:1b \\
C3607 & 1a<o2 1b*<o1 / 2a<1a0L / 3a<2a0L & IU/N/N & {-}{-}/{-}/{-} & N & W:1a \\
C3608 & 1a<o2 1b*<o1 / 2a<1a0L / 3a<2a1 & IU/N/N & {-}{-}/{-}/{-} & N & W:1a \\
C3609 & 1a<o2 1b*<o1 / 2a<1a0L / 3a<2a1L & IU/N/N & {-}{-}/{-}/{-} & N & W:1a \\
C3610 & 1a<o2 1b*<o1 / 2a<1a0L / 3a<2a2 & IU/N/N & {-}{-}/{-}/{-} & N & W:1a \\
C3611 & 1a<o2 1b*<o1 / 2a<1a0L / 3a<2a2L & IU/N/N & {-}{-}/{-}/{-} & N & W:1a \\
C3612 & 1a<o2 1b*<o1 / 2a<1a1 / 3a<2a0L & IU/I/N & {-}{-}/{-}/{-} & N & W:1a,2a \\
C3613 & 1a<o2 1b*<o1 / 2a<1a1 / 3a<2a1 & IU/I/I & {-}{-}/{-}/{-} & N & W:1a,2a,3a \\
C3614 & 1a<o2 1b*<o1 / 2a<1a1 / 3a<2a1L & IU/I/I & {-}{-}/{-}/{-} & N & W:1a,2a,3a \\
C3615 & 1a<o2 1b*<o1 / 2a<1a1 / 3a<2a2 & IU/I/I & {-}{-}/{-}/{-} & N & W:1a,2a,3a \\
C3616 & 1a<o2 1b*<o1 / 2a<1a1 / 3a<2a2L & IU/I/I & {-}{-}/{-}/{-} & N & W:1a,2a,3a \\
C3617 & 1a<o2 1b*<o1 / 2a<1a1L / 3a<2a0L & IU/I/N & {-}{-}/{-}/{-} & N & W:1a,2a \\
C3618 & 1a<o2 1b*<o1 / 2a<1a1L / 3a<2a1 & IU/I/I & {-}{-}/{-}/{-} & N & W:1a,2a,3a \\
C3619 & 1a<o2 1b*<o1 / 2a<1a1L / 3a<2a1L & IU/I/I & {-}{-}/{-}/{-} & N & W:1a,2a,3a \\
C3620 & 1a<o2 1b*<o1 / 2a<1a1L / 3a<2a2 & IU/I/I & {-}{-}/{-}/{-} & N & W:1a,2a,3a \\
C3621 & 1a<o2 1b*<o1 / 2a<1a1L / 3a<2a2L & IU/I/I & {-}{-}/{-}/{-} & N & W:1a,2a,3a \\
C3622 & 1a<o2 1b*<o1 / 2a<1a2 / 3a<2a0L & IU/I/N & {-}{-}/{-}/{-} & N & W:1a,2a \\
C3623 & 1a<o2 1b*<o1 / 2a<1a2 / 3a<2a1 & IU/I/I & {-}{-}/{-}/{-} & N & W:1a,2a,3a \\
C3624 & 1a<o2 1b*<o1 / 2a<1a2 / 3a<2a1L & IU/I/I & {-}{-}/{-}/{-} & N & W:1a,2a,3a \\
C3625 & 1a<o2 1b*<o1 / 2a<1a2 / 3a<2a2 & IU/I/I & {-}{-}/{-}/{-} & N & W:1a,2a,3a \\
C3626 & 1a<o2 1b*<o1 / 2a<1a2 / 3a<2a2L & IU/I/I & {-}{-}/{-}/{-} & N & W:1a,2a,3a \\
C3627 & 1a<o2 1b*<o1 / 2a<1a2L / 3a<2a0L & IU/I/N & {-}{-}/{-}/{-} & N & W:1a,2a \\
C3628 & 1a<o2 1b*<o1 / 2a<1a2L / 3a<2a1 & IU/I/I & {-}{-}/{-}/{-} & N & W:1a,2a,3a \\
C3629 & 1a<o2 1b*<o1 / 2a<1a2L / 3a<2a1L & IU/I/I & {-}{-}/{-}/{-} & N & W:1a,2a,3a \\
C3630 & 1a<o2 1b*<o1 / 2a<1a2L / 3a<2a2 & IU/I/I & {-}{-}/{-}/{-} & N & W:1a,2a,3a \\
C3631 & 1a<o2 1b*<o1 / 2a<1a2L / 3a<2a2L & IU/I/I & {-}{-}/{-}/{-} & N & W:1a,2a,3a \\
C3632 & 1a<o2 1b*<o1 / 2a<1b1 / 3a<2a0L & IU/U/N & {-}{-}/{-}/{-} & N & W:1a \\
C3633 & 1a<o2 1b*<o1 / 2a<1b1 / 3a<2a1 & IU/U/U & {-}{-}/{-}/{-} & N & W:1a \\
C3634 & 1a<o2 1b*<o1 / 2a<1b1 / 3a<2a1L & IU/U/U & {-}{-}/{-}/{-} & N & W:1a \\
C3635 & 1a<o2 1b*<o1 / 2a<1b1 / 3a<2a2 & IU/U/U & {-}{-}/{-}/{-} & N & W:1a \\
C3636 & 1a<o2 1b*<o1 / 2a<1b1 / 3a<2a2L & IU/U/U & {-}{-}/{-}/{-} & N & W:1a \\
C3637 & 1a<o2 1b*<o1 / 2a<1b2 / 3a<2a0L & IU/U/N & {-}{-}/{-}/{-} & N & W:1a \\
C3638 & 1a<o2 1b*<o1 / 2a<1b2 / 3a<2a1 & IU/U/U & {-}{-}/{-}/{-} & N & W:1a \\
C3639 & 1a<o2 1b*<o1 / 2a<1b2 / 3a<2a1L & IU/U/U & {-}{-}/{-}/{-} & N & W:1a \\
C3640 & 1a<o2 1b*<o1 / 2a<1b2 / 3a<2a2 & IU/U/U & {-}{-}/{-}/{-} & N & W:1a \\
C3641 & 1a<o2 1b*<o1 / 2a<1b2 / 3a<2a2L & IU/U/U & {-}{-}/{-}/{-} & N & W:1a \\
C3642 & 1a<o2 1b*<o2 / 2a<1a0L / 3a<2a0L & IU/N/N & {-}{-}/{-}/{-} & N & W:1a \\
C3643 & 1a<o2 1b*<o2 / 2a<1a0L / 3a<2a1 & IU/N/N & {-}{-}/{-}/{-} & N & W:1a \\
C3644 & 1a<o2 1b*<o2 / 2a<1a0L / 3a<2a1L & IU/N/N & {-}{-}/{-}/{-} & N & W:1a \\
C3645 & 1a<o2 1b*<o2 / 2a<1a0L / 3a<2a2 & IU/N/N & {-}{-}/{-}/{-} & N & W:1a \\
C3646 & 1a<o2 1b*<o2 / 2a<1a0L / 3a<2a2L & IU/N/N & {-}{-}/{-}/{-} & N & W:1a \\
C3647 & 1a<o2 1b*<o2 / 2a<1a1 / 3a<2a0L & IU/I/N & {-}{-}/{-}/{-} & N & W:1a,2a \\
C3648 & 1a<o2 1b*<o2 / 2a<1a1 / 3a<2a1 & IU/I/I & {-}{-}/{-}/{-} & N & W:1a,2a,3a \\
C3649 & 1a<o2 1b*<o2 / 2a<1a1 / 3a<2a1L & IU/I/I & {-}{-}/{-}/{-} & N & W:1a,2a,3a \\
C3650 & 1a<o2 1b*<o2 / 2a<1a1 / 3a<2a2 & IU/I/I & {-}{-}/{-}/{-} & N & W:1a,2a,3a \\
C3651 & 1a<o2 1b*<o2 / 2a<1a1 / 3a<2a2L & IU/I/I & {-}{-}/{-}/{-} & N & W:1a,2a,3a \\
C3652 & 1a<o2 1b*<o2 / 2a<1a1L / 3a<2a0L & IU/I/N & {-}{-}/{-}/{-} & N & W:1a,2a \\
C3653 & 1a<o2 1b*<o2 / 2a<1a1L / 3a<2a1 & IU/I/I & {-}{-}/{-}/{-} & N & W:1a,2a,3a \\
C3654 & 1a<o2 1b*<o2 / 2a<1a1L / 3a<2a1L & IU/I/I & {-}{-}/{-}/{-} & N & W:1a,2a,3a \\
C3655 & 1a<o2 1b*<o2 / 2a<1a1L / 3a<2a2 & IU/I/I & {-}{-}/{-}/{-} & N & W:1a,2a,3a \\
C3656 & 1a<o2 1b*<o2 / 2a<1a1L / 3a<2a2L & IU/I/I & {-}{-}/{-}/{-} & N & W:1a,2a,3a \\
C3657 & 1a<o2 1b*<o2 / 2a<1a2 / 3a<2a0L & IU/I/N & {-}{-}/{-}/{-} & N & W:1a,2a \\
C3658 & 1a<o2 1b*<o2 / 2a<1a2 / 3a<2a1 & IU/I/I & {-}{-}/{-}/{-} & N & W:1a,2a,3a \\
C3659 & 1a<o2 1b*<o2 / 2a<1a2 / 3a<2a1L & IU/I/I & {-}{-}/{-}/{-} & N & W:1a,2a,3a \\
C3660 & 1a<o2 1b*<o2 / 2a<1a2 / 3a<2a2 & IU/I/I & {-}{-}/{-}/{-} & N & W:1a,2a,3a \\
C3661 & 1a<o2 1b*<o2 / 2a<1a2 / 3a<2a2L & IU/I/I & {-}{-}/{-}/{-} & N & W:1a,2a,3a \\
C3662 & 1a<o2 1b*<o2 / 2a<1a2L / 3a<2a0L & IU/I/N & {-}{-}/{-}/{-} & N & W:1a,2a \\
C3663 & 1a<o2 1b*<o2 / 2a<1a2L / 3a<2a1 & IU/I/I & {-}{-}/{-}/{-} & N & W:1a,2a,3a \\
C3664 & 1a<o2 1b*<o2 / 2a<1a2L / 3a<2a1L & IU/I/I & {-}{-}/{-}/{-} & N & W:1a,2a,3a \\
C3665 & 1a<o2 1b*<o2 / 2a<1a2L / 3a<2a2 & IU/I/I & {-}{-}/{-}/{-} & N & W:1a,2a,3a \\
C3666 & 1a<o2 1b*<o2 / 2a<1a2L / 3a<2a2L & IU/I/I & {-}{-}/{-}/{-} & N & W:1a,2a,3a \\
C3667 & 1a<o2 1b*<o2 / 2a<1b1 / 3a<2a0L & IU/U/N & {-}{-}/{-}/{-} & N & W:1a \\
C3668 & 1a<o2 1b*<o2 / 2a<1b1 / 3a<2a1 & IU/U/U & {-}{-}/{-}/{-} & N & W:1a \\
C3669 & 1a<o2 1b*<o2 / 2a<1b1 / 3a<2a1L & IU/U/U & {-}{-}/{-}/{-} & N & W:1a \\
C3670 & 1a<o2 1b*<o2 / 2a<1b1 / 3a<2a2 & IU/U/U & {-}{-}/{-}/{-} & N & W:1a \\
C3671 & 1a<o2 1b*<o2 / 2a<1b1 / 3a<2a2L & IU/U/U & {-}{-}/{-}/{-} & N & W:1a \\
C3672 & 1a<o2 1b*<o2 / 2a<1b2 / 3a<2a0L & IU/U/N & {-}{-}/{-}/{-} & N & W:1a \\
C3673 & 1a<o2 1b*<o2 / 2a<1b2 / 3a<2a1 & IU/U/U & {-}{-}/{-}/{-} & N & W:1a \\
C3674 & 1a<o2 1b*<o2 / 2a<1b2 / 3a<2a1L & IU/U/U & {-}{-}/{-}/{-} & N & W:1a \\
C3675 & 1a<o2 1b*<o2 / 2a<1b2 / 3a<2a2 & IU/U/U & {-}{-}/{-}/{-} & N & W:1a \\
C3676 & 1a<o2 1b*<o2 / 2a<1b2 / 3a<2a2L & IU/U/U & {-}{-}/{-}/{-} & N & W:1a \\
C3677 & 1a<o2L / 2a<1a0L / 3a<2a0L 3b*<2a1 & I/N/NN & +/+/+{-} & N &  \\
C3678 & 1a<o2L / 2a<1a0L / 3a<2a0L 3b*<2a2 & I/N/NN & +/+/+{-} & N &  \\
C3679 & 1a<o2L / 2a<1a0L / 3a<2a1 3b*<2a1 & I/N/NN & +/+/{-}{-} & N &  \\
C3680 & 1a<o2L / 2a<1a0L / 3a<2a1 3b*<2a2 & I/N/NN & +/+/{-}{-} & N &  \\
C3681 & 1a<o2L / 2a<1a0L / 3a<2a1L 3b*<2a1 & I/N/NN & +/+/+{-} & N &  \\
C3682 & 1a<o2L / 2a<1a0L / 3a<2a1L 3b*<2a2 & I/N/NN & +/+/+{-} & N &  \\
C3683 & 1a<o2L / 2a<1a0L / 3a<2a2 3b*<2a1 & I/N/NN & +/+/{-}{-} & N &  \\
C3684 & 1a<o2L / 2a<1a0L / 3a<2a2 3b*<2a2 & I/N/NN & +/+/{-}{-} & N &  \\
C3685 & 1a<o2L / 2a<1a0L / 3a<2a2L 3b*<2a1 & I/N/NN & +/+/+{-} & N &  \\
C3686 & 1a<o2L / 2a<1a0L / 3a<2a2L 3b*<2a2 & I/N/NN & +/+/+{-} & N &  \\
C3687 & 1a<o2L / 2a<1a0L 2b<1a1 / 3a*<2a1 & I/NI/N & +/+{-}/{-} & N &  \\
C3688 & 1a<o2L / 2a<1a0L 2b<1a1 / 3a*<2a2 & I/NI/N & +/+{-}/{-} & N &  \\
C3689 & 1a<o2L / 2a<1a0L 2b<1a1 / 3a*<2b1 & I/NI/U & +/+{-}/{-} & N & W:3a \\
C3690 & 1a<o2L / 2a<1a0L 2b<1a1 / 3a*<2b2 & I/NI/U & +/+{-}/{-} & N & W:3a \\
C3691 & 1a<o2L / 2a<1a0L 2b<1a2 / 3a*<2a1 & I/NI/N & +/+{-}/{-} & N &  \\
C3692 & 1a<o2L / 2a<1a0L 2b<1a2 / 3a*<2a2 & I/NI/N & +/+{-}/{-} & N &  \\
C3693 & 1a<o2L / 2a<1a0L 2b<1a2 / 3a*<2b1 & I/NI/U & +/+{-}/{-} & N & W:3a \\
C3694 & 1a<o2L / 2a<1a0L 2b<1a2 / 3a*<2b2 & I/NI/U & +/+{-}/{-} & N & W:3a \\
C3695 & 1a<o2L / 2a<1a0L 2b*<1a1 / 3a<2a0L & I/NU/N & +/+{-}/+ & N & W:2b \\
C3696 & 1a<o2L / 2a<1a0L 2b*<1a1 / 3a<2a1 & I/NU/N & +/+{-}/{-} & N & W:2b \\
C3697 & 1a<o2L / 2a<1a0L 2b*<1a1 / 3a<2a1L & I/NU/N & +/+{-}/+ & N & W:2b \\
C3698 & 1a<o2L / 2a<1a0L 2b*<1a1 / 3a<2a2 & I/NU/N & +/+{-}/{-} & N & W:2b \\
C3699 & 1a<o2L / 2a<1a0L 2b*<1a1 / 3a<2a2L & I/NU/N & +/+{-}/+ & N & W:2b \\
C3700 & 1a<o2L / 2a<1a0L 2b*<1a1 / 3a<2b1 & I/NU/U & +/+{-}/{-} & N & W:2b,3a \\
C3701 & 1a<o2L / 2a<1a0L 2b*<1a1 / 3a<2b2 & I/NU/U & +/+{-}/{-} & N & W:2b,3a \\
C3702 & 1a<o2L / 2a<1a0L 2b*<1a2 / 3a<2a0L & I/NU/N & +/+{-}/+ & N & W:2b \\
C3703 & 1a<o2L / 2a<1a0L 2b*<1a2 / 3a<2a1 & I/NU/N & +/+{-}/{-} & N & W:2b \\
C3704 & 1a<o2L / 2a<1a0L 2b*<1a2 / 3a<2a1L & I/NU/N & +/+{-}/+ & N & W:2b \\
C3705 & 1a<o2L / 2a<1a0L 2b*<1a2 / 3a<2a2 & I/NU/N & +/+{-}/{-} & N & W:2b \\
C3706 & 1a<o2L / 2a<1a0L 2b*<1a2 / 3a<2a2L & I/NU/N & +/+{-}/+ & N & W:2b \\
C3707 & 1a<o2L / 2a<1a0L 2b*<1a2 / 3a<2b1 & I/NU/U & +/+{-}/{-} & N & W:2b,3a \\
C3708 & 1a<o2L / 2a<1a0L 2b*<1a2 / 3a<2b2 & I/NU/U & +/+{-}/{-} & N & W:2b,3a \\
C3709 & 1a<o2L / 2a<1a1 / 3a<2a0L 3b*<2a1 & I/I/NU & +/{-}/{-}{-} & N & W:3b \\
C3710 & 1a<o2L / 2a<1a1 / 3a<2a0L 3b*<2a2 & I/I/NU & +/{-}/{-}{-} & N & W:3b \\
C3711 & 1a<o2L / 2a<1a1 / 3a<2a1 3b*<2a1 & I/I/IU & +/{-}/{-}{-} & N & W:1a,2a,3a \\
C3712 & 1a<o2L / 2a<1a1 / 3a<2a1 3b*<2a2 & I/I/IU & +/{-}/{-}{-} & N & W:1a,2a,3a \\
C3713 & 1a<o2L / 2a<1a1 / 3a<2a1L 3b*<2a1 & I/I/IU & +/{-}/{-}{-} & N & W:1a,2a,3a \\
C3714 & 1a<o2L / 2a<1a1 / 3a<2a1L 3b*<2a2 & I/I/IU & +/{-}/{-}{-} & N & W:1a,2a,3a \\
C3715 & 1a<o2L / 2a<1a1 / 3a<2a2 3b*<2a1 & I/I/IU & +/{-}/{-}{-} & N & W:1a,2a,3a \\
C3716 & 1a<o2L / 2a<1a1 / 3a<2a2 3b*<2a2 & I/I/IU & +/{-}/{-}{-} & N & W:1a,2a,3a \\
C3717 & 1a<o2L / 2a<1a1 / 3a<2a2L 3b*<2a1 & I/I/IU & +/{-}/{-}{-} & N & W:1a,2a,3a \\
C3718 & 1a<o2L / 2a<1a1 / 3a<2a2L 3b*<2a2 & I/I/IU & +/{-}/{-}{-} & N & W:1a,2a,3a \\
C3719 & 1a<o2L / 2a<1a1 2b<1a1 / 3a*<2a1 & I/EE/U & +/{-}{-}/{-} & E & W:1a \\
C3720 & 1a<o2L / 2a<1a1 2b<1a1 / 3a*<2a2 & I/EE/U & +/{-}{-}/{-} & E & W:1a \\
C3721 & 1a<o2L / 2a<1a1 2b<1a1L / 3a*<2a1 & I/LI/U & +/{-}+/{-} & B & CP:2b W:1a,2b \\
C3722 & 1a<o2L / 2a<1a1 2b<1a1L / 3a*<2a2 & I/LI/U & +/{-}+/{-} & B & CP:2b W:1a,2b \\
C3723 & 1a<o2L / 2a<1a1 2b<1a1L / 3a*<2b1 & I/LI/U & +/{-}+/{-} & B & CP:2b W:1a,2b \\
C3724 & 1a<o2L / 2a<1a1 2b<1a1L / 3a*<2b2 & I/LI/U & +/{-}+/{-} & B & CP:2b W:1a,2b \\
C3725 & 1a<o2L / 2a<1a1 2b<1a2 / 3a*<2a1 & I/EE/U & +/{-}{-}/{-} & E & C:2b W:1a \\
C3726 & 1a<o2L / 2a<1a1 2b<1a2 / 3a*<2a2 & I/EE/U & +/{-}{-}/{-} & E & C:2b W:1a \\
C3727 & 1a<o2L / 2a<1a1 2b<1a2 / 3a*<2b1 & I/EE/U & +/{-}{-}/{-} & E & C:2b W:1a \\
C3728 & 1a<o2L / 2a<1a1 2b<1a2 / 3a*<2b2 & I/EE/U & +/{-}{-}/{-} & E & C:2b W:1a \\
C3729 & 1a<o2L / 2a<1a1 2b<1a2L / 3a*<2a1 & I/LI/U & +/{-}+/{-} & B & P:2b W:1a,2b \\
C3730 & 1a<o2L / 2a<1a1 2b<1a2L / 3a*<2a2 & I/LI/U & +/{-}+/{-} & B & P:2b W:1a,2b \\
C3731 & 1a<o2L / 2a<1a1 2b<1a2L / 3a*<2b1 & I/LI/U & +/{-}+/{-} & B & P:2b W:1a,2b \\
C3732 & 1a<o2L / 2a<1a1 2b<1a2L / 3a*<2b2 & I/LI/U & +/{-}+/{-} & B & P:2b W:1a,2b \\
C3733 & 1a<o2L / 2a<1a1 2b*<1a1 / 3a<2a0L & I/IU/N & +/{-}{-}/{-} & N & W:1a,2a \\
C3734 & 1a<o2L / 2a<1a1 2b*<1a1 / 3a<2a1 & I/IU/I & +/{-}{-}/{-} & N & W:1a,2a,3a \\
C3735 & 1a<o2L / 2a<1a1 2b*<1a1 / 3a<2a1L & I/IU/I & +/{-}{-}/{-} & N & W:1a,2a,3a \\
C3736 & 1a<o2L / 2a<1a1 2b*<1a1 / 3a<2a2 & I/IU/I & +/{-}{-}/{-} & N & W:1a,2a,3a \\
C3737 & 1a<o2L / 2a<1a1 2b*<1a1 / 3a<2a2L & I/IU/I & +/{-}{-}/{-} & N & W:1a,2a,3a \\
C3738 & 1a<o2L / 2a<1a1 2b*<1a1 / 3a<2b1 & I/IU/U & +/{-}{-}/{-} & N & W:1a,2a \\
C3739 & 1a<o2L / 2a<1a1 2b*<1a1 / 3a<2b2 & I/IU/U & +/{-}{-}/{-} & N & W:1a,2a \\
C3740 & 1a<o2L / 2a<1a1 2b*<1a2 / 3a<2a0L & I/IU/N & +/{-}{-}/{-} & N & W:1a,2a \\
C3741 & 1a<o2L / 2a<1a1 2b*<1a2 / 3a<2a1 & I/IU/I & +/{-}{-}/{-} & N & W:1a,2a,3a \\
C3742 & 1a<o2L / 2a<1a1 2b*<1a2 / 3a<2a1L & I/IU/I & +/{-}{-}/{-} & N & W:1a,2a,3a \\
C3743 & 1a<o2L / 2a<1a1 2b*<1a2 / 3a<2a2 & I/IU/I & +/{-}{-}/{-} & N & W:1a,2a,3a \\
C3744 & 1a<o2L / 2a<1a1 2b*<1a2 / 3a<2a2L & I/IU/I & +/{-}{-}/{-} & N & W:1a,2a,3a \\
C3745 & 1a<o2L / 2a<1a1 2b*<1a2 / 3a<2b1 & I/IU/U & +/{-}{-}/{-} & N & W:1a,2a \\
C3746 & 1a<o2L / 2a<1a1 2b*<1a2 / 3a<2b2 & I/IU/U & +/{-}{-}/{-} & N & W:1a,2a \\
C3747 & 1a<o2L / 2a<1a1L / 3a<2a0L 3b*<2a1 & I/I/NU & +/+/+{-} & N & W:3b \\
C3748 & 1a<o2L / 2a<1a1L / 3a<2a0L 3b*<2a2 & I/I/NU & +/+/+{-} & N & W:3b \\
C3749 & 1a<o2L / 2a<1a1L / 3a<2a1 3b*<2a1 & I/I/IU & +/+/{-}{-} & N & W:1a,2a,3a \\
C3750 & 1a<o2L / 2a<1a1L / 3a<2a1 3b*<2a2 & I/I/IU & +/+/{-}{-} & N & W:1a,2a,3a \\
C3751 & 1a<o2L / 2a<1a1L / 3a<2a1L 3b*<2a1 & I/I/IU & +/+/+{-} & N & W:1a,2a,3a \\
C3752 & 1a<o2L / 2a<1a1L / 3a<2a1L 3b*<2a2 & I/I/IU & +/+/+{-} & N & W:1a,2a,3a \\
C3753 & 1a<o2L / 2a<1a1L / 3a<2a2 3b*<2a1 & I/I/IU & +/+/{-}{-} & N & W:1a,2a,3a \\
C3754 & 1a<o2L / 2a<1a1L / 3a<2a2 3b*<2a2 & I/I/IU & +/+/{-}{-} & N & W:1a,2a,3a \\
C3755 & 1a<o2L / 2a<1a1L / 3a<2a2L 3b*<2a1 & I/I/IU & +/+/+{-} & N & W:1a,2a,3a \\
C3756 & 1a<o2L / 2a<1a1L / 3a<2a2L 3b*<2a2 & I/I/IU & +/+/+{-} & N & W:1a,2a,3a \\
C3757 & 1a<o2L / 2a<1a1L 2b<1a2 / 3a*<2a1 & I/IL/U & +/+{-}/{-} & B & C:2a,2b P:2a W:1a,2a \\
C3758 & 1a<o2L / 2a<1a1L 2b<1a2 / 3a*<2a2 & I/IL/U & +/+{-}/{-} & B & C:2a,2b P:2a W:1a,2a \\
C3759 & 1a<o2L / 2a<1a1L 2b<1a2 / 3a*<2b1 & I/IL/U & +/+{-}/{-} & B & C:2a,2b P:2a W:1a,2a \\
C3760 & 1a<o2L / 2a<1a1L 2b<1a2 / 3a*<2b2 & I/IL/U & +/+{-}/{-} & B & C:2a,2b P:2a W:1a,2a \\
C3761 & 1a<o2L / 2a<1a1L 2b*<1a1 / 3a<2a0L & I/IU/N & +/+{-}/+ & N & W:1a,2a \\
C3762 & 1a<o2L / 2a<1a1L 2b*<1a1 / 3a<2a1 & I/IU/I & +/+{-}/{-} & N & W:1a,2a,3a \\
C3763 & 1a<o2L / 2a<1a1L 2b*<1a1 / 3a<2a1L & I/IU/I & +/+{-}/+ & N & W:1a,2a,3a \\
C3764 & 1a<o2L / 2a<1a1L 2b*<1a1 / 3a<2a2 & I/IU/I & +/+{-}/{-} & N & W:1a,2a,3a \\
C3765 & 1a<o2L / 2a<1a1L 2b*<1a1 / 3a<2a2L & I/IU/I & +/+{-}/+ & N & W:1a,2a,3a \\
C3766 & 1a<o2L / 2a<1a1L 2b*<1a1 / 3a<2b1 & I/IU/U & +/+{-}/{-} & N & W:1a,2a \\
C3767 & 1a<o2L / 2a<1a1L 2b*<1a1 / 3a<2b2 & I/IU/U & +/+{-}/{-} & N & W:1a,2a \\
C3768 & 1a<o2L / 2a<1a1L 2b*<1a2 / 3a<2a0L & I/IU/N & +/+{-}/+ & N & W:1a,2a \\
C3769 & 1a<o2L / 2a<1a1L 2b*<1a2 / 3a<2a1 & I/IU/I & +/+{-}/{-} & N & W:1a,2a,3a \\
C3770 & 1a<o2L / 2a<1a1L 2b*<1a2 / 3a<2a1L & I/IU/I & +/+{-}/+ & N & W:1a,2a,3a \\
C3771 & 1a<o2L / 2a<1a1L 2b*<1a2 / 3a<2a2 & I/IU/I & +/+{-}/{-} & N & W:1a,2a,3a \\
C3772 & 1a<o2L / 2a<1a1L 2b*<1a2 / 3a<2a2L & I/IU/I & +/+{-}/+ & N & W:1a,2a,3a \\
C3773 & 1a<o2L / 2a<1a1L 2b*<1a2 / 3a<2b1 & I/IU/U & +/+{-}/{-} & N & W:1a,2a \\
C3774 & 1a<o2L / 2a<1a1L 2b*<1a2 / 3a<2b2 & I/IU/U & +/+{-}/{-} & N & W:1a,2a \\
C3775 & 1a<o2L / 2a<1a2 / 3a<2a0L 3b*<2a1 & I/I/NU & +/{-}/{-}{-} & N & W:3b \\
C3776 & 1a<o2L / 2a<1a2 / 3a<2a0L 3b*<2a2 & I/I/NU & +/{-}/{-}{-} & N & W:3b \\
C3777 & 1a<o2L / 2a<1a2 / 3a<2a1 3b*<2a1 & I/I/IU & +/{-}/{-}{-} & N & W:1a,2a,3a \\
C3778 & 1a<o2L / 2a<1a2 / 3a<2a1 3b*<2a2 & I/I/IU & +/{-}/{-}{-} & N & W:1a,2a,3a \\
C3779 & 1a<o2L / 2a<1a2 / 3a<2a1L 3b*<2a1 & I/I/IU & +/{-}/{-}{-} & N & W:1a,2a,3a \\
C3780 & 1a<o2L / 2a<1a2 / 3a<2a1L 3b*<2a2 & I/I/IU & +/{-}/{-}{-} & N & W:1a,2a,3a \\
C3781 & 1a<o2L / 2a<1a2 / 3a<2a2 3b*<2a1 & I/I/IU & +/{-}/{-}{-} & N & W:1a,2a,3a \\
C3782 & 1a<o2L / 2a<1a2 / 3a<2a2 3b*<2a2 & I/I/IU & +/{-}/{-}{-} & N & W:1a,2a,3a \\
C3783 & 1a<o2L / 2a<1a2 / 3a<2a2L 3b*<2a1 & I/I/IU & +/{-}/{-}{-} & N & W:1a,2a,3a \\
C3784 & 1a<o2L / 2a<1a2 / 3a<2a2L 3b*<2a2 & I/I/IU & +/{-}/{-}{-} & N & W:1a,2a,3a \\
C3785 & 1a<o2L / 2a<1a2 2b<1a2 / 3a*<2a1 & I/EE/U & +/{-}{-}/{-} & E & W:1a \\
C3786 & 1a<o2L / 2a<1a2 2b<1a2 / 3a*<2a2 & I/EE/U & +/{-}{-}/{-} & E & W:1a \\
C3787 & 1a<o2L / 2a<1a2 2b<1a2L / 3a*<2a1 & I/LI/U & +/{-}+/{-} & B & CP:2b W:1a,2b \\
C3788 & 1a<o2L / 2a<1a2 2b<1a2L / 3a*<2a2 & I/LI/U & +/{-}+/{-} & B & CP:2b W:1a,2b \\
C3789 & 1a<o2L / 2a<1a2 2b<1a2L / 3a*<2b1 & I/LI/U & +/{-}+/{-} & B & CP:2b W:1a,2b \\
C3790 & 1a<o2L / 2a<1a2 2b<1a2L / 3a*<2b2 & I/LI/U & +/{-}+/{-} & B & CP:2b W:1a,2b \\
C3791 & 1a<o2L / 2a<1a2 2b*<1a1 / 3a<2a0L & I/IU/N & +/{-}{-}/{-} & N & W:1a,2a \\
C3792 & 1a<o2L / 2a<1a2 2b*<1a1 / 3a<2a1 & I/IU/I & +/{-}{-}/{-} & N & W:1a,2a,3a \\
C3793 & 1a<o2L / 2a<1a2 2b*<1a1 / 3a<2a1L & I/IU/I & +/{-}{-}/{-} & N & W:1a,2a,3a \\
C3794 & 1a<o2L / 2a<1a2 2b*<1a1 / 3a<2a2 & I/IU/I & +/{-}{-}/{-} & N & W:1a,2a,3a \\
C3795 & 1a<o2L / 2a<1a2 2b*<1a1 / 3a<2a2L & I/IU/I & +/{-}{-}/{-} & N & W:1a,2a,3a \\
C3796 & 1a<o2L / 2a<1a2 2b*<1a1 / 3a<2b1 & I/IU/U & +/{-}{-}/{-} & N & W:1a,2a \\
C3797 & 1a<o2L / 2a<1a2 2b*<1a1 / 3a<2b2 & I/IU/U & +/{-}{-}/{-} & N & W:1a,2a \\
C3798 & 1a<o2L / 2a<1a2 2b*<1a2 / 3a<2a0L & I/IU/N & +/{-}{-}/{-} & N & W:1a,2a \\
C3799 & 1a<o2L / 2a<1a2 2b*<1a2 / 3a<2a1 & I/IU/I & +/{-}{-}/{-} & N & W:1a,2a,3a \\
C3800 & 1a<o2L / 2a<1a2 2b*<1a2 / 3a<2a1L & I/IU/I & +/{-}{-}/{-} & N & W:1a,2a,3a \\
C3801 & 1a<o2L / 2a<1a2 2b*<1a2 / 3a<2a2 & I/IU/I & +/{-}{-}/{-} & N & W:1a,2a,3a \\
C3802 & 1a<o2L / 2a<1a2 2b*<1a2 / 3a<2a2L & I/IU/I & +/{-}{-}/{-} & N & W:1a,2a,3a \\
C3803 & 1a<o2L / 2a<1a2 2b*<1a2 / 3a<2b1 & I/IU/U & +/{-}{-}/{-} & N & W:1a,2a \\
C3804 & 1a<o2L / 2a<1a2 2b*<1a2 / 3a<2b2 & I/IU/U & +/{-}{-}/{-} & N & W:1a,2a \\
C3805 & 1a<o2L / 2a<1a2L / 3a<2a0L 3b*<2a1 & I/I/NU & +/+/+{-} & N & W:3b \\
C3806 & 1a<o2L / 2a<1a2L / 3a<2a0L 3b*<2a2 & I/I/NU & +/+/+{-} & N & W:3b \\
C3807 & 1a<o2L / 2a<1a2L / 3a<2a1 3b*<2a1 & I/I/IU & +/+/{-}{-} & N & W:1a,2a,3a \\
C3808 & 1a<o2L / 2a<1a2L / 3a<2a1 3b*<2a2 & I/I/IU & +/+/{-}{-} & N & W:1a,2a,3a \\
C3809 & 1a<o2L / 2a<1a2L / 3a<2a1L 3b*<2a1 & I/I/IU & +/+/+{-} & N & W:1a,2a,3a \\
C3810 & 1a<o2L / 2a<1a2L / 3a<2a1L 3b*<2a2 & I/I/IU & +/+/+{-} & N & W:1a,2a,3a \\
C3811 & 1a<o2L / 2a<1a2L / 3a<2a2 3b*<2a1 & I/I/IU & +/+/{-}{-} & N & W:1a,2a,3a \\
C3812 & 1a<o2L / 2a<1a2L / 3a<2a2 3b*<2a2 & I/I/IU & +/+/{-}{-} & N & W:1a,2a,3a \\
C3813 & 1a<o2L / 2a<1a2L / 3a<2a2L 3b*<2a1 & I/I/IU & +/+/+{-} & N & W:1a,2a,3a \\
C3814 & 1a<o2L / 2a<1a2L / 3a<2a2L 3b*<2a2 & I/I/IU & +/+/+{-} & N & W:1a,2a,3a \\
C3815 & 1a<o2L / 2a<1a2L 2b*<1a1 / 3a<2a0L & I/IU/N & +/+{-}/+ & N & W:1a,2a \\
C3816 & 1a<o2L / 2a<1a2L 2b*<1a1 / 3a<2a1 & I/IU/I & +/+{-}/{-} & N & W:1a,2a,3a \\
C3817 & 1a<o2L / 2a<1a2L 2b*<1a1 / 3a<2a1L & I/IU/I & +/+{-}/+ & N & W:1a,2a,3a \\
C3818 & 1a<o2L / 2a<1a2L 2b*<1a1 / 3a<2a2 & I/IU/I & +/+{-}/{-} & N & W:1a,2a,3a \\
C3819 & 1a<o2L / 2a<1a2L 2b*<1a1 / 3a<2a2L & I/IU/I & +/+{-}/+ & N & W:1a,2a,3a \\
C3820 & 1a<o2L / 2a<1a2L 2b*<1a1 / 3a<2b1 & I/IU/U & +/+{-}/{-} & N & W:1a,2a \\
C3821 & 1a<o2L / 2a<1a2L 2b*<1a1 / 3a<2b2 & I/IU/U & +/+{-}/{-} & N & W:1a,2a \\
C3822 & 1a<o2L / 2a<1a2L 2b*<1a2 / 3a<2a0L & I/IU/N & +/+{-}/+ & N & W:1a,2a \\
C3823 & 1a<o2L / 2a<1a2L 2b*<1a2 / 3a<2a1 & I/IU/I & +/+{-}/{-} & N & W:1a,2a,3a \\
C3824 & 1a<o2L / 2a<1a2L 2b*<1a2 / 3a<2a1L & I/IU/I & +/+{-}/+ & N & W:1a,2a,3a \\
C3825 & 1a<o2L / 2a<1a2L 2b*<1a2 / 3a<2a2 & I/IU/I & +/+{-}/{-} & N & W:1a,2a,3a \\
C3826 & 1a<o2L / 2a<1a2L 2b*<1a2 / 3a<2a2L & I/IU/I & +/+{-}/+ & N & W:1a,2a,3a \\
C3827 & 1a<o2L / 2a<1a2L 2b*<1a2 / 3a<2b1 & I/IU/U & +/+{-}/{-} & N & W:1a,2a \\
C3828 & 1a<o2L / 2a<1a2L 2b*<1a2 / 3a<2b2 & I/IU/U & +/+{-}/{-} & N & W:1a,2a \\
C3829 & 1a<o2L / 2a*<1a1 / 3a<2a1 3b<2a1 & I/U/UU & +/{-}/{-}{-} & N & W:1a \\
C3830 & 1a<o2L / 2a*<1a1 / 3a<2a1 3b<2a2 & I/U/UU & +/{-}/{-}{-} & N & W:1a \\
C3831 & 1a<o2L / 2a*<1a1 / 3a<2a2 3b<2a2 & I/U/UU & +/{-}/{-}{-} & N & W:1a \\
C3832 & 1a<o2L / 2a*<1a2 / 3a<2a1 3b<2a1 & I/U/UU & +/{-}/{-}{-} & N & W:1a \\
C3833 & 1a<o2L / 2a*<1a2 / 3a<2a1 3b<2a2 & I/U/UU & +/{-}/{-}{-} & N & W:1a \\
C3834 & 1a<o2L / 2a*<1a2 / 3a<2a2 3b<2a2 & I/U/UU & +/{-}/{-}{-} & N & W:1a \\
C3835 & 1a<o2L 1b*<o1 / 2a<1a0L / 3a<2a0L & IU/N/N & +{-}/+/+ & N & W:1a \\
C3836 & 1a<o2L 1b*<o1 / 2a<1a0L / 3a<2a1 & IU/N/N & +{-}/+/{-} & N & W:1a \\
C3837 & 1a<o2L 1b*<o1 / 2a<1a0L / 3a<2a1L & IU/N/N & +{-}/+/+ & N & W:1a \\
C3838 & 1a<o2L 1b*<o1 / 2a<1a0L / 3a<2a2 & IU/N/N & +{-}/+/{-} & N & W:1a \\
C3839 & 1a<o2L 1b*<o1 / 2a<1a0L / 3a<2a2L & IU/N/N & +{-}/+/+ & N & W:1a \\
C3840 & 1a<o2L 1b*<o1 / 2a<1a1 / 3a<2a0L & IU/I/N & +{-}/{-}/{-} & N & W:1a,2a \\
C3841 & 1a<o2L 1b*<o1 / 2a<1a1 / 3a<2a1 & IU/I/I & +{-}/{-}/{-} & N & W:1a,2a,3a \\
C3842 & 1a<o2L 1b*<o1 / 2a<1a1 / 3a<2a1L & IU/I/I & +{-}/{-}/{-} & N & W:1a,2a,3a \\
C3843 & 1a<o2L 1b*<o1 / 2a<1a1 / 3a<2a2 & IU/I/I & +{-}/{-}/{-} & N & W:1a,2a,3a \\
C3844 & 1a<o2L 1b*<o1 / 2a<1a1 / 3a<2a2L & IU/I/I & +{-}/{-}/{-} & N & W:1a,2a,3a \\
C3845 & 1a<o2L 1b*<o1 / 2a<1a1L / 3a<2a0L & IU/I/N & +{-}/+/+ & N & W:1a,2a \\
C3846 & 1a<o2L 1b*<o1 / 2a<1a1L / 3a<2a1 & IU/I/I & +{-}/+/{-} & N & W:1a,2a,3a \\
C3847 & 1a<o2L 1b*<o1 / 2a<1a1L / 3a<2a1L & IU/I/I & +{-}/+/+ & N & W:1a,2a,3a \\
C3848 & 1a<o2L 1b*<o1 / 2a<1a1L / 3a<2a2 & IU/I/I & +{-}/+/{-} & N & W:1a,2a,3a \\
C3849 & 1a<o2L 1b*<o1 / 2a<1a1L / 3a<2a2L & IU/I/I & +{-}/+/+ & N & W:1a,2a,3a \\
C3850 & 1a<o2L 1b*<o1 / 2a<1a2 / 3a<2a0L & IU/I/N & +{-}/{-}/{-} & N & W:1a,2a \\
C3851 & 1a<o2L 1b*<o1 / 2a<1a2 / 3a<2a1 & IU/I/I & +{-}/{-}/{-} & N & W:1a,2a,3a \\
C3852 & 1a<o2L 1b*<o1 / 2a<1a2 / 3a<2a1L & IU/I/I & +{-}/{-}/{-} & N & W:1a,2a,3a \\
C3853 & 1a<o2L 1b*<o1 / 2a<1a2 / 3a<2a2 & IU/I/I & +{-}/{-}/{-} & N & W:1a,2a,3a \\
C3854 & 1a<o2L 1b*<o1 / 2a<1a2 / 3a<2a2L & IU/I/I & +{-}/{-}/{-} & N & W:1a,2a,3a \\
C3855 & 1a<o2L 1b*<o1 / 2a<1a2L / 3a<2a0L & IU/I/N & +{-}/+/+ & N & W:1a,2a \\
C3856 & 1a<o2L 1b*<o1 / 2a<1a2L / 3a<2a1 & IU/I/I & +{-}/+/{-} & N & W:1a,2a,3a \\
C3857 & 1a<o2L 1b*<o1 / 2a<1a2L / 3a<2a1L & IU/I/I & +{-}/+/+ & N & W:1a,2a,3a \\
C3858 & 1a<o2L 1b*<o1 / 2a<1a2L / 3a<2a2 & IU/I/I & +{-}/+/{-} & N & W:1a,2a,3a \\
C3859 & 1a<o2L 1b*<o1 / 2a<1a2L / 3a<2a2L & IU/I/I & +{-}/+/+ & N & W:1a,2a,3a \\
C3860 & 1a<o2L 1b*<o1 / 2a<1b1 / 3a<2a0L & IU/U/N & +{-}/{-}/{-} & N & W:1a \\
C3861 & 1a<o2L 1b*<o1 / 2a<1b1 / 3a<2a1 & IU/U/U & +{-}/{-}/{-} & N & W:1a \\
C3862 & 1a<o2L 1b*<o1 / 2a<1b1 / 3a<2a1L & IU/U/U & +{-}/{-}/{-} & N & W:1a \\
C3863 & 1a<o2L 1b*<o1 / 2a<1b1 / 3a<2a2 & IU/U/U & +{-}/{-}/{-} & N & W:1a \\
C3864 & 1a<o2L 1b*<o1 / 2a<1b1 / 3a<2a2L & IU/U/U & +{-}/{-}/{-} & N & W:1a \\
C3865 & 1a<o2L 1b*<o1 / 2a<1b2 / 3a<2a0L & IU/U/N & +{-}/{-}/{-} & N & W:1a \\
C3866 & 1a<o2L 1b*<o1 / 2a<1b2 / 3a<2a1 & IU/U/U & +{-}/{-}/{-} & N & W:1a \\
C3867 & 1a<o2L 1b*<o1 / 2a<1b2 / 3a<2a1L & IU/U/U & +{-}/{-}/{-} & N & W:1a \\
C3868 & 1a<o2L 1b*<o1 / 2a<1b2 / 3a<2a2 & IU/U/U & +{-}/{-}/{-} & N & W:1a \\
C3869 & 1a<o2L 1b*<o1 / 2a<1b2 / 3a<2a2L & IU/U/U & +{-}/{-}/{-} & N & W:1a \\
C3870 & 1a<o2L 1b*<o2 / 2a<1a0L / 3a<2a0L & IU/N/N & +{-}/+/+ & N & W:1a \\
C3871 & 1a<o2L 1b*<o2 / 2a<1a0L / 3a<2a1 & IU/N/N & +{-}/+/{-} & N & W:1a \\
C3872 & 1a<o2L 1b*<o2 / 2a<1a0L / 3a<2a1L & IU/N/N & +{-}/+/+ & N & W:1a \\
C3873 & 1a<o2L 1b*<o2 / 2a<1a0L / 3a<2a2 & IU/N/N & +{-}/+/{-} & N & W:1a \\
C3874 & 1a<o2L 1b*<o2 / 2a<1a0L / 3a<2a2L & IU/N/N & +{-}/+/+ & N & W:1a \\
C3875 & 1a<o2L 1b*<o2 / 2a<1a1 / 3a<2a0L & IU/I/N & +{-}/{-}/{-} & N & W:1a,2a \\
C3876 & 1a<o2L 1b*<o2 / 2a<1a1 / 3a<2a1 & IU/I/I & +{-}/{-}/{-} & N & W:1a,2a,3a \\
C3877 & 1a<o2L 1b*<o2 / 2a<1a1 / 3a<2a1L & IU/I/I & +{-}/{-}/{-} & N & W:1a,2a,3a \\
C3878 & 1a<o2L 1b*<o2 / 2a<1a1 / 3a<2a2 & IU/I/I & +{-}/{-}/{-} & N & W:1a,2a,3a \\
C3879 & 1a<o2L 1b*<o2 / 2a<1a1 / 3a<2a2L & IU/I/I & +{-}/{-}/{-} & N & W:1a,2a,3a \\
C3880 & 1a<o2L 1b*<o2 / 2a<1a1L / 3a<2a0L & IU/I/N & +{-}/+/+ & N & W:1a,2a \\
C3881 & 1a<o2L 1b*<o2 / 2a<1a1L / 3a<2a1 & IU/I/I & +{-}/+/{-} & N & W:1a,2a,3a \\
C3882 & 1a<o2L 1b*<o2 / 2a<1a1L / 3a<2a1L & IU/I/I & +{-}/+/+ & N & W:1a,2a,3a \\
C3883 & 1a<o2L 1b*<o2 / 2a<1a1L / 3a<2a2 & IU/I/I & +{-}/+/{-} & N & W:1a,2a,3a \\
C3884 & 1a<o2L 1b*<o2 / 2a<1a1L / 3a<2a2L & IU/I/I & +{-}/+/+ & N & W:1a,2a,3a \\
C3885 & 1a<o2L 1b*<o2 / 2a<1a2 / 3a<2a0L & IU/I/N & +{-}/{-}/{-} & N & W:1a,2a \\
C3886 & 1a<o2L 1b*<o2 / 2a<1a2 / 3a<2a1 & IU/I/I & +{-}/{-}/{-} & N & W:1a,2a,3a \\
C3887 & 1a<o2L 1b*<o2 / 2a<1a2 / 3a<2a1L & IU/I/I & +{-}/{-}/{-} & N & W:1a,2a,3a \\
C3888 & 1a<o2L 1b*<o2 / 2a<1a2 / 3a<2a2 & IU/I/I & +{-}/{-}/{-} & N & W:1a,2a,3a \\
C3889 & 1a<o2L 1b*<o2 / 2a<1a2 / 3a<2a2L & IU/I/I & +{-}/{-}/{-} & N & W:1a,2a,3a \\
C3890 & 1a<o2L 1b*<o2 / 2a<1a2L / 3a<2a0L & IU/I/N & +{-}/+/+ & N & W:1a,2a \\
C3891 & 1a<o2L 1b*<o2 / 2a<1a2L / 3a<2a1 & IU/I/I & +{-}/+/{-} & N & W:1a,2a,3a \\
C3892 & 1a<o2L 1b*<o2 / 2a<1a2L / 3a<2a1L & IU/I/I & +{-}/+/+ & N & W:1a,2a,3a \\
C3893 & 1a<o2L 1b*<o2 / 2a<1a2L / 3a<2a2 & IU/I/I & +{-}/+/{-} & N & W:1a,2a,3a \\
C3894 & 1a<o2L 1b*<o2 / 2a<1a2L / 3a<2a2L & IU/I/I & +{-}/+/+ & N & W:1a,2a,3a \\
C3895 & 1a<o2L 1b*<o2 / 2a<1b1 / 3a<2a0L & IU/U/N & +{-}/{-}/{-} & N & W:1a \\
C3896 & 1a<o2L 1b*<o2 / 2a<1b1 / 3a<2a1 & IU/U/U & +{-}/{-}/{-} & N & W:1a \\
C3897 & 1a<o2L 1b*<o2 / 2a<1b1 / 3a<2a1L & IU/U/U & +{-}/{-}/{-} & N & W:1a \\
C3898 & 1a<o2L 1b*<o2 / 2a<1b1 / 3a<2a2 & IU/U/U & +{-}/{-}/{-} & N & W:1a \\
C3899 & 1a<o2L 1b*<o2 / 2a<1b1 / 3a<2a2L & IU/U/U & +{-}/{-}/{-} & N & W:1a \\
C3900 & 1a<o2L 1b*<o2 / 2a<1b2 / 3a<2a0L & IU/U/N & +{-}/{-}/{-} & N & W:1a \\
C3901 & 1a<o2L 1b*<o2 / 2a<1b2 / 3a<2a1 & IU/U/U & +{-}/{-}/{-} & N & W:1a \\
C3902 & 1a<o2L 1b*<o2 / 2a<1b2 / 3a<2a1L & IU/U/U & +{-}/{-}/{-} & N & W:1a \\
C3903 & 1a<o2L 1b*<o2 / 2a<1b2 / 3a<2a2 & IU/U/U & +{-}/{-}/{-} & N & W:1a \\
C3904 & 1a<o2L 1b*<o2 / 2a<1b2 / 3a<2a2L & IU/U/U & +{-}/{-}/{-} & N & W:1a \\
C3905 & 1a*<o1 / 2a<1a1 / 3a<2a0L 3b<2a1 & U/U/NU & {-}/{-}/{-}{-} & N & W:1a,2a,3b \\
C3906 & 1a*<o1 / 2a<1a1 / 3a<2a0L 3b<2a2 & U/U/NU & {-}/{-}/{-}{-} & N & W:1a,2a,3b \\
C3907 & 1a*<o1 / 2a<1a1 / 3a<2a1 3b<2a1 & U/U/UU & {-}/{-}/{-}{-} & N &  \\
C3908 & 1a*<o1 / 2a<1a1 / 3a<2a1 3b<2a1L & U/U/UU & {-}/{-}/{-}{-} & N &  \\
C3909 & 1a*<o1 / 2a<1a1 / 3a<2a1 3b<2a2 & U/U/UU & {-}/{-}/{-}{-} & N &  \\
C3910 & 1a*<o1 / 2a<1a1 / 3a<2a1 3b<2a2L & U/U/UU & {-}/{-}/{-}{-} & N &  \\
C3911 & 1a*<o1 / 2a<1a1 / 3a<2a1L 3b<2a2 & U/U/UU & {-}/{-}/{-}{-} & N &  \\
C3912 & 1a*<o1 / 2a<1a1 / 3a<2a2 3b<2a2 & U/U/UU & {-}/{-}/{-}{-} & N &  \\
C3913 & 1a*<o1 / 2a<1a1 / 3a<2a2 3b<2a2L & U/U/UU & {-}/{-}/{-}{-} & N &  \\
C3914 & 1a*<o1 / 2a<1a1 2b<1a1 / 3a<2a0L & U/UU/N & {-}/{-}{-}/{-} & N &  \\
C3915 & 1a*<o1 / 2a<1a1 2b<1a1 / 3a<2a1 & U/UU/U & {-}/{-}{-}/{-} & N &  \\
C3916 & 1a*<o1 / 2a<1a1 2b<1a1 / 3a<2a1L & U/UU/U & {-}/{-}{-}/{-} & N &  \\
C3917 & 1a*<o1 / 2a<1a1 2b<1a1 / 3a<2a2 & U/UU/U & {-}/{-}{-}/{-} & N &  \\
C3918 & 1a*<o1 / 2a<1a1 2b<1a1 / 3a<2a2L & U/UU/U & {-}/{-}{-}/{-} & N &  \\
C3919 & 1a*<o1 / 2a<1a1 2b<1a2 / 3a<2a0L & U/UU/N & {-}/{-}{-}/{-} & N &  \\
C3920 & 1a*<o1 / 2a<1a1 2b<1a2 / 3a<2a1 & U/UU/U & {-}/{-}{-}/{-} & N &  \\
C3921 & 1a*<o1 / 2a<1a1 2b<1a2 / 3a<2a1L & U/UU/U & {-}/{-}{-}/{-} & N &  \\
C3922 & 1a*<o1 / 2a<1a1 2b<1a2 / 3a<2a2 & U/UU/U & {-}/{-}{-}/{-} & N &  \\
C3923 & 1a*<o1 / 2a<1a1 2b<1a2 / 3a<2a2L & U/UU/U & {-}/{-}{-}/{-} & N &  \\
C3924 & 1a*<o1 / 2a<1a1 2b<1a2 / 3a<2b0L & U/UU/N & {-}/{-}{-}/{-} & N &  \\
C3925 & 1a*<o1 / 2a<1a1 2b<1a2 / 3a<2b1 & U/UU/U & {-}/{-}{-}/{-} & N &  \\
C3926 & 1a*<o1 / 2a<1a1 2b<1a2 / 3a<2b1L & U/UU/U & {-}/{-}{-}/{-} & N &  \\
C3927 & 1a*<o1 / 2a<1a1 2b<1a2 / 3a<2b2 & U/UU/U & {-}/{-}{-}/{-} & N &  \\
C3928 & 1a*<o1 / 2a<1a1 2b<1a2 / 3a<2b2L & U/UU/U & {-}/{-}{-}/{-} & N &  \\
C3929 & 1a*<o1 / 2a<1a2 / 3a<2a0L 3b<2a1 & U/U/NU & {-}/{-}/{-}{-} & N & W:1a,2a,3b \\
C3930 & 1a*<o1 / 2a<1a2 / 3a<2a0L 3b<2a2 & U/U/NU & {-}/{-}/{-}{-} & N & W:1a,2a,3b \\
C3931 & 1a*<o1 / 2a<1a2 / 3a<2a1 3b<2a1 & U/U/UU & {-}/{-}/{-}{-} & N &  \\
C3932 & 1a*<o1 / 2a<1a2 / 3a<2a1 3b<2a1L & U/U/UU & {-}/{-}/{-}{-} & N &  \\
C3933 & 1a*<o1 / 2a<1a2 / 3a<2a1 3b<2a2 & U/U/UU & {-}/{-}/{-}{-} & N &  \\
C3934 & 1a*<o1 / 2a<1a2 / 3a<2a1 3b<2a2L & U/U/UU & {-}/{-}/{-}{-} & N &  \\
C3935 & 1a*<o1 / 2a<1a2 / 3a<2a1L 3b<2a2 & U/U/UU & {-}/{-}/{-}{-} & N &  \\
C3936 & 1a*<o1 / 2a<1a2 / 3a<2a2 3b<2a2 & U/U/UU & {-}/{-}/{-}{-} & N &  \\
C3937 & 1a*<o1 / 2a<1a2 / 3a<2a2 3b<2a2L & U/U/UU & {-}/{-}/{-}{-} & N &  \\
C3938 & 1a*<o1 / 2a<1a2 2b<1a2 / 3a<2a0L & U/UU/N & {-}/{-}{-}/{-} & N &  \\
C3939 & 1a*<o1 / 2a<1a2 2b<1a2 / 3a<2a1 & U/UU/U & {-}/{-}{-}/{-} & N &  \\
C3940 & 1a*<o1 / 2a<1a2 2b<1a2 / 3a<2a1L & U/UU/U & {-}/{-}{-}/{-} & N &  \\
C3941 & 1a*<o1 / 2a<1a2 2b<1a2 / 3a<2a2 & U/UU/U & {-}/{-}{-}/{-} & N &  \\
C3942 & 1a*<o1 / 2a<1a2 2b<1a2 / 3a<2a2L & U/UU/U & {-}/{-}{-}/{-} & N &  \\
C3943 & 1a*<o2 / 2a<1a1 / 3a<2a0L 3b<2a1 & U/U/NU & {-}/{-}/{-}{-} & N & W:1a,2a,3b \\
C3944 & 1a*<o2 / 2a<1a1 / 3a<2a0L 3b<2a2 & U/U/NU & {-}/{-}/{-}{-} & N & W:1a,2a,3b \\
C3945 & 1a*<o2 / 2a<1a1 / 3a<2a1 3b<2a1 & U/U/UU & {-}/{-}/{-}{-} & N &  \\
C3946 & 1a*<o2 / 2a<1a1 / 3a<2a1 3b<2a1L & U/U/UU & {-}/{-}/{-}{-} & N &  \\
C3947 & 1a*<o2 / 2a<1a1 / 3a<2a1 3b<2a2 & U/U/UU & {-}/{-}/{-}{-} & N &  \\
C3948 & 1a*<o2 / 2a<1a1 / 3a<2a1 3b<2a2L & U/U/UU & {-}/{-}/{-}{-} & N &  \\
C3949 & 1a*<o2 / 2a<1a1 / 3a<2a1L 3b<2a2 & U/U/UU & {-}/{-}/{-}{-} & N &  \\
C3950 & 1a*<o2 / 2a<1a1 / 3a<2a2 3b<2a2 & U/U/UU & {-}/{-}/{-}{-} & N &  \\
C3951 & 1a*<o2 / 2a<1a1 / 3a<2a2 3b<2a2L & U/U/UU & {-}/{-}/{-}{-} & N &  \\
C3952 & 1a*<o2 / 2a<1a1 2b<1a1 / 3a<2a0L & U/UU/N & {-}/{-}{-}/{-} & N &  \\
C3953 & 1a*<o2 / 2a<1a1 2b<1a1 / 3a<2a1 & U/UU/U & {-}/{-}{-}/{-} & N &  \\
C3954 & 1a*<o2 / 2a<1a1 2b<1a1 / 3a<2a1L & U/UU/U & {-}/{-}{-}/{-} & N &  \\
C3955 & 1a*<o2 / 2a<1a1 2b<1a1 / 3a<2a2 & U/UU/U & {-}/{-}{-}/{-} & N &  \\
C3956 & 1a*<o2 / 2a<1a1 2b<1a1 / 3a<2a2L & U/UU/U & {-}/{-}{-}/{-} & N &  \\
C3957 & 1a*<o2 / 2a<1a1 2b<1a2 / 3a<2a0L & U/UU/N & {-}/{-}{-}/{-} & N &  \\
C3958 & 1a*<o2 / 2a<1a1 2b<1a2 / 3a<2a1 & U/UU/U & {-}/{-}{-}/{-} & N &  \\
C3959 & 1a*<o2 / 2a<1a1 2b<1a2 / 3a<2a1L & U/UU/U & {-}/{-}{-}/{-} & N &  \\
C3960 & 1a*<o2 / 2a<1a1 2b<1a2 / 3a<2a2 & U/UU/U & {-}/{-}{-}/{-} & N &  \\
C3961 & 1a*<o2 / 2a<1a1 2b<1a2 / 3a<2a2L & U/UU/U & {-}/{-}{-}/{-} & N &  \\
C3962 & 1a*<o2 / 2a<1a1 2b<1a2 / 3a<2b0L & U/UU/N & {-}/{-}{-}/{-} & N &  \\
C3963 & 1a*<o2 / 2a<1a1 2b<1a2 / 3a<2b1 & U/UU/U & {-}/{-}{-}/{-} & N &  \\
C3964 & 1a*<o2 / 2a<1a1 2b<1a2 / 3a<2b1L & U/UU/U & {-}/{-}{-}/{-} & N &  \\
C3965 & 1a*<o2 / 2a<1a1 2b<1a2 / 3a<2b2 & U/UU/U & {-}/{-}{-}/{-} & N &  \\
C3966 & 1a*<o2 / 2a<1a1 2b<1a2 / 3a<2b2L & U/UU/U & {-}/{-}{-}/{-} & N &  \\
C3967 & 1a*<o2 / 2a<1a2 / 3a<2a0L 3b<2a1 & U/U/NU & {-}/{-}/{-}{-} & N & W:1a,2a,3b \\
C3968 & 1a*<o2 / 2a<1a2 / 3a<2a0L 3b<2a2 & U/U/NU & {-}/{-}/{-}{-} & N & W:1a,2a,3b \\
C3969 & 1a*<o2 / 2a<1a2 / 3a<2a1 3b<2a1 & U/U/UU & {-}/{-}/{-}{-} & N &  \\
C3970 & 1a*<o2 / 2a<1a2 / 3a<2a1 3b<2a1L & U/U/UU & {-}/{-}/{-}{-} & N &  \\
C3971 & 1a*<o2 / 2a<1a2 / 3a<2a1 3b<2a2 & U/U/UU & {-}/{-}/{-}{-} & N &  \\
C3972 & 1a*<o2 / 2a<1a2 / 3a<2a1 3b<2a2L & U/U/UU & {-}/{-}/{-}{-} & N &  \\
C3973 & 1a*<o2 / 2a<1a2 / 3a<2a1L 3b<2a2 & U/U/UU & {-}/{-}/{-}{-} & N &  \\
C3974 & 1a*<o2 / 2a<1a2 / 3a<2a2 3b<2a2 & U/U/UU & {-}/{-}/{-}{-} & N &  \\
C3975 & 1a*<o2 / 2a<1a2 / 3a<2a2 3b<2a2L & U/U/UU & {-}/{-}/{-}{-} & N &  \\
C3976 & 1a*<o2 / 2a<1a2 2b<1a2 / 3a<2a0L & U/UU/N & {-}/{-}{-}/{-} & N &  \\
C3977 & 1a*<o2 / 2a<1a2 2b<1a2 / 3a<2a1 & U/UU/U & {-}/{-}{-}/{-} & N &  \\
C3978 & 1a*<o2 / 2a<1a2 2b<1a2 / 3a<2a1L & U/UU/U & {-}/{-}{-}/{-} & N &  \\
C3979 & 1a*<o2 / 2a<1a2 2b<1a2 / 3a<2a2 & U/UU/U & {-}/{-}{-}/{-} & N &  \\
C3980 & 1a*<o2 / 2a<1a2 2b<1a2 / 3a<2a2L & U/UU/U & {-}/{-}{-}/{-} & N &  \\
C3981 & 1a<o0L / 2a<1a0L 2b<1a1 / 3a*<2a1,2b1 & N/NN/N & +/+{-}/{-} & N &  \\
C3982 & 1a<o0L / 2a<1a0L 2b<1a1 / 3a*<2a1,2b2 & N/NN/N & +/+{-}/{-} & N &  \\
C3983 & 1a<o0L / 2a<1a0L 2b<1a1 / 3a*<2a2,2b1 & N/NN/N & +/+{-}/{-} & N &  \\
C3984 & 1a<o0L / 2a<1a0L 2b<1a1 / 3a*<2a2,2b2 & N/NN/N & +/+{-}/{-} & N &  \\
C3985 & 1a<o0L / 2a<1a0L 2b<1a2 / 3a*<2a1,2b1 & N/NN/N & +/+{-}/{-} & N &  \\
C3986 & 1a<o0L / 2a<1a0L 2b<1a2 / 3a*<2a1,2b2 & N/NN/N & +/+{-}/{-} & N &  \\
C3987 & 1a<o0L / 2a<1a0L 2b<1a2 / 3a*<2a2,2b1 & N/NN/N & +/+{-}/{-} & N &  \\
C3988 & 1a<o0L / 2a<1a0L 2b<1a2 / 3a*<2a2,2b2 & N/NN/N & +/+{-}/{-} & N &  \\
C3989 & 1a<o0L / 2a<1a0L 2b*<1a1 / 3a<2a0L,2b1 & N/NN/N & +/+{-}/+ & N &  \\
C3990 & 1a<o0L / 2a<1a0L 2b*<1a1 / 3a<2a0L,2b2 & N/NN/N & +/+{-}/+ & N &  \\
C3991 & 1a<o0L / 2a<1a0L 2b*<1a1 / 3a<2a1,2b1 & N/NN/N & +/+{-}/{-} & N &  \\
C3992 & 1a<o0L / 2a<1a0L 2b*<1a1 / 3a<2a1,2b2 & N/NN/N & +/+{-}/{-} & N &  \\
C3993 & 1a<o0L / 2a<1a0L 2b*<1a1 / 3a<2a1L,2b1 & N/NN/N & +/+{-}/+ & N &  \\
C3994 & 1a<o0L / 2a<1a0L 2b*<1a1 / 3a<2a1L,2b2 & N/NN/N & +/+{-}/+ & N &  \\
C3995 & 1a<o0L / 2a<1a0L 2b*<1a1 / 3a<2a2,2b1 & N/NN/N & +/+{-}/{-} & N &  \\
C3996 & 1a<o0L / 2a<1a0L 2b*<1a1 / 3a<2a2,2b2 & N/NN/N & +/+{-}/{-} & N &  \\
C3997 & 1a<o0L / 2a<1a0L 2b*<1a1 / 3a<2a2L,2b1 & N/NN/N & +/+{-}/+ & N &  \\
C3998 & 1a<o0L / 2a<1a0L 2b*<1a1 / 3a<2a2L,2b2 & N/NN/N & +/+{-}/+ & N &  \\
C3999 & 1a<o0L / 2a<1a0L 2b*<1a2 / 3a<2a0L,2b1 & N/NN/N & +/+{-}/+ & N &  \\
C4000 & 1a<o0L / 2a<1a0L 2b*<1a2 / 3a<2a0L,2b2 & N/NN/N & +/+{-}/+ & N &  \\
C4001 & 1a<o0L / 2a<1a0L 2b*<1a2 / 3a<2a1,2b1 & N/NN/N & +/+{-}/{-} & N &  \\
C4002 & 1a<o0L / 2a<1a0L 2b*<1a2 / 3a<2a1,2b2 & N/NN/N & +/+{-}/{-} & N &  \\
C4003 & 1a<o0L / 2a<1a0L 2b*<1a2 / 3a<2a1L,2b1 & N/NN/N & +/+{-}/+ & N &  \\
C4004 & 1a<o0L / 2a<1a0L 2b*<1a2 / 3a<2a1L,2b2 & N/NN/N & +/+{-}/+ & N &  \\
C4005 & 1a<o0L / 2a<1a0L 2b*<1a2 / 3a<2a2,2b1 & N/NN/N & +/+{-}/{-} & N &  \\
C4006 & 1a<o0L / 2a<1a0L 2b*<1a2 / 3a<2a2,2b2 & N/NN/N & +/+{-}/{-} & N &  \\
C4007 & 1a<o0L / 2a<1a0L 2b*<1a2 / 3a<2a2L,2b1 & N/NN/N & +/+{-}/+ & N &  \\
C4008 & 1a<o0L / 2a<1a0L 2b*<1a2 / 3a<2a2L,2b2 & N/NN/N & +/+{-}/+ & N &  \\
C4009 & 1a<o0L / 2a<1a1 2b<1a1 / 3a*<2a1,2b1 & N/NN/N & +/{-}{-}/{-} & N &  \\
C4010 & 1a<o0L / 2a<1a1 2b<1a1 / 3a*<2a1,2b2 & N/NN/N & +/{-}{-}/{-} & N &  \\
C4011 & 1a<o0L / 2a<1a1 2b<1a1 / 3a*<2a2,2b2 & N/NN/N & +/{-}{-}/{-} & N &  \\
C4012 & 1a<o0L / 2a<1a1 2b<1a1L / 3a*<2a1,2b1 & N/NN/N & +/{-}+/{-} & N &  \\
C4013 & 1a<o0L / 2a<1a1 2b<1a1L / 3a*<2a1,2b2 & N/NN/N & +/{-}+/{-} & N &  \\
C4014 & 1a<o0L / 2a<1a1 2b<1a1L / 3a*<2a2,2b1 & N/NN/N & +/{-}+/{-} & N &  \\
C4015 & 1a<o0L / 2a<1a1 2b<1a1L / 3a*<2a2,2b2 & N/NN/N & +/{-}+/{-} & N &  \\
C4016 & 1a<o0L / 2a<1a1 2b<1a2 / 3a*<2a1,2b1 & N/NN/N & +/{-}{-}/{-} & N &  \\
C4017 & 1a<o0L / 2a<1a1 2b<1a2 / 3a*<2a1,2b2 & N/NN/N & +/{-}{-}/{-} & N &  \\
C4018 & 1a<o0L / 2a<1a1 2b<1a2 / 3a*<2a2,2b1 & N/NN/N & +/{-}{-}/{-} & N &  \\
C4019 & 1a<o0L / 2a<1a1 2b<1a2 / 3a*<2a2,2b2 & N/NN/N & +/{-}{-}/{-} & N &  \\
C4020 & 1a<o0L / 2a<1a1 2b<1a2L / 3a*<2a1,2b1 & N/NN/N & +/{-}+/{-} & N &  \\
C4021 & 1a<o0L / 2a<1a1 2b<1a2L / 3a*<2a1,2b2 & N/NN/N & +/{-}+/{-} & N &  \\
C4022 & 1a<o0L / 2a<1a1 2b<1a2L / 3a*<2a2,2b1 & N/NN/N & +/{-}+/{-} & N &  \\
C4023 & 1a<o0L / 2a<1a1 2b<1a2L / 3a*<2a2,2b2 & N/NN/N & +/{-}+/{-} & N &  \\
C4024 & 1a<o0L / 2a<1a1 2b*<1a1 / 3a<2a0L,2b1 & N/NN/N & +/{-}{-}/{-} & N &  \\
C4025 & 1a<o0L / 2a<1a1 2b*<1a1 / 3a<2a0L,2b2 & N/NN/N & +/{-}{-}/{-} & N &  \\
C4026 & 1a<o0L / 2a<1a1 2b*<1a1 / 3a<2a1,2b1 & N/NN/N & +/{-}{-}/{-} & N &  \\
C4027 & 1a<o0L / 2a<1a1 2b*<1a1 / 3a<2a1,2b2 & N/NN/N & +/{-}{-}/{-} & N &  \\
C4028 & 1a<o0L / 2a<1a1 2b*<1a1 / 3a<2a1L,2b1 & N/NN/N & +/{-}{-}/{-} & N &  \\
C4029 & 1a<o0L / 2a<1a1 2b*<1a1 / 3a<2a1L,2b2 & N/NN/N & +/{-}{-}/{-} & N &  \\
C4030 & 1a<o0L / 2a<1a1 2b*<1a1 / 3a<2a2,2b1 & N/NN/N & +/{-}{-}/{-} & N &  \\
C4031 & 1a<o0L / 2a<1a1 2b*<1a1 / 3a<2a2,2b2 & N/NN/N & +/{-}{-}/{-} & N &  \\
C4032 & 1a<o0L / 2a<1a1 2b*<1a1 / 3a<2a2L,2b1 & N/NN/N & +/{-}{-}/{-} & N &  \\
C4033 & 1a<o0L / 2a<1a1 2b*<1a1 / 3a<2a2L,2b2 & N/NN/N & +/{-}{-}/{-} & N &  \\
C4034 & 1a<o0L / 2a<1a1 2b*<1a2 / 3a<2a0L,2b1 & N/NN/N & +/{-}{-}/{-} & N &  \\
C4035 & 1a<o0L / 2a<1a1 2b*<1a2 / 3a<2a0L,2b2 & N/NN/N & +/{-}{-}/{-} & N &  \\
C4036 & 1a<o0L / 2a<1a1 2b*<1a2 / 3a<2a1,2b1 & N/NN/N & +/{-}{-}/{-} & N &  \\
C4037 & 1a<o0L / 2a<1a1 2b*<1a2 / 3a<2a1,2b2 & N/NN/N & +/{-}{-}/{-} & N &  \\
C4038 & 1a<o0L / 2a<1a1 2b*<1a2 / 3a<2a1L,2b1 & N/NN/N & +/{-}{-}/{-} & N &  \\
C4039 & 1a<o0L / 2a<1a1 2b*<1a2 / 3a<2a1L,2b2 & N/NN/N & +/{-}{-}/{-} & N &  \\
C4040 & 1a<o0L / 2a<1a1 2b*<1a2 / 3a<2a2,2b1 & N/NN/N & +/{-}{-}/{-} & N &  \\
C4041 & 1a<o0L / 2a<1a1 2b*<1a2 / 3a<2a2,2b2 & N/NN/N & +/{-}{-}/{-} & N &  \\
C4042 & 1a<o0L / 2a<1a1 2b*<1a2 / 3a<2a2L,2b1 & N/NN/N & +/{-}{-}/{-} & N &  \\
C4043 & 1a<o0L / 2a<1a1 2b*<1a2 / 3a<2a2L,2b2 & N/NN/N & +/{-}{-}/{-} & N &  \\
C4044 & 1a<o0L / 2a<1a1L 2b<1a2 / 3a*<2a1,2b1 & N/NN/N & +/+{-}/{-} & N &  \\
C4045 & 1a<o0L / 2a<1a1L 2b<1a2 / 3a*<2a1,2b2 & N/NN/N & +/+{-}/{-} & N &  \\
C4046 & 1a<o0L / 2a<1a1L 2b<1a2 / 3a*<2a2,2b1 & N/NN/N & +/+{-}/{-} & N &  \\
C4047 & 1a<o0L / 2a<1a1L 2b<1a2 / 3a*<2a2,2b2 & N/NN/N & +/+{-}/{-} & N &  \\
C4048 & 1a<o0L / 2a<1a1L 2b*<1a1 / 3a<2a0L,2b1 & N/NN/N & +/+{-}/+ & N &  \\
C4049 & 1a<o0L / 2a<1a1L 2b*<1a1 / 3a<2a0L,2b2 & N/NN/N & +/+{-}/+ & N &  \\
C4050 & 1a<o0L / 2a<1a1L 2b*<1a1 / 3a<2a1,2b1 & N/NN/N & +/+{-}/{-} & N &  \\
C4051 & 1a<o0L / 2a<1a1L 2b*<1a1 / 3a<2a1,2b2 & N/NN/N & +/+{-}/{-} & N &  \\
C4052 & 1a<o0L / 2a<1a1L 2b*<1a1 / 3a<2a1L,2b1 & N/NN/N & +/+{-}/+ & N &  \\
C4053 & 1a<o0L / 2a<1a1L 2b*<1a1 / 3a<2a1L,2b2 & N/NN/N & +/+{-}/+ & N &  \\
C4054 & 1a<o0L / 2a<1a1L 2b*<1a1 / 3a<2a2,2b1 & N/NN/N & +/+{-}/{-} & N &  \\
C4055 & 1a<o0L / 2a<1a1L 2b*<1a1 / 3a<2a2,2b2 & N/NN/N & +/+{-}/{-} & N &  \\
C4056 & 1a<o0L / 2a<1a1L 2b*<1a1 / 3a<2a2L,2b1 & N/NN/N & +/+{-}/+ & N &  \\
C4057 & 1a<o0L / 2a<1a1L 2b*<1a1 / 3a<2a2L,2b2 & N/NN/N & +/+{-}/+ & N &  \\
C4058 & 1a<o0L / 2a<1a1L 2b*<1a2 / 3a<2a0L,2b1 & N/NN/N & +/+{-}/+ & N &  \\
C4059 & 1a<o0L / 2a<1a1L 2b*<1a2 / 3a<2a0L,2b2 & N/NN/N & +/+{-}/+ & N &  \\
C4060 & 1a<o0L / 2a<1a1L 2b*<1a2 / 3a<2a1,2b1 & N/NN/N & +/+{-}/{-} & N &  \\
C4061 & 1a<o0L / 2a<1a1L 2b*<1a2 / 3a<2a1,2b2 & N/NN/N & +/+{-}/{-} & N &  \\
C4062 & 1a<o0L / 2a<1a1L 2b*<1a2 / 3a<2a1L,2b1 & N/NN/N & +/+{-}/+ & N &  \\
C4063 & 1a<o0L / 2a<1a1L 2b*<1a2 / 3a<2a1L,2b2 & N/NN/N & +/+{-}/+ & N &  \\
C4064 & 1a<o0L / 2a<1a1L 2b*<1a2 / 3a<2a2,2b1 & N/NN/N & +/+{-}/{-} & N &  \\
C4065 & 1a<o0L / 2a<1a1L 2b*<1a2 / 3a<2a2,2b2 & N/NN/N & +/+{-}/{-} & N &  \\
C4066 & 1a<o0L / 2a<1a1L 2b*<1a2 / 3a<2a2L,2b1 & N/NN/N & +/+{-}/+ & N &  \\
C4067 & 1a<o0L / 2a<1a1L 2b*<1a2 / 3a<2a2L,2b2 & N/NN/N & +/+{-}/+ & N &  \\
C4068 & 1a<o0L / 2a<1a2 2b<1a2 / 3a*<2a1,2b1 & N/NN/N & +/{-}{-}/{-} & N &  \\
C4069 & 1a<o0L / 2a<1a2 2b<1a2 / 3a*<2a1,2b2 & N/NN/N & +/{-}{-}/{-} & N &  \\
C4070 & 1a<o0L / 2a<1a2 2b<1a2 / 3a*<2a2,2b2 & N/NN/N & +/{-}{-}/{-} & N &  \\
C4071 & 1a<o0L / 2a<1a2 2b<1a2L / 3a*<2a1,2b1 & N/NN/N & +/{-}+/{-} & N &  \\
C4072 & 1a<o0L / 2a<1a2 2b<1a2L / 3a*<2a1,2b2 & N/NN/N & +/{-}+/{-} & N &  \\
C4073 & 1a<o0L / 2a<1a2 2b<1a2L / 3a*<2a2,2b1 & N/NN/N & +/{-}+/{-} & N &  \\
C4074 & 1a<o0L / 2a<1a2 2b<1a2L / 3a*<2a2,2b2 & N/NN/N & +/{-}+/{-} & N &  \\
C4075 & 1a<o0L / 2a<1a2 2b*<1a1 / 3a<2a0L,2b1 & N/NN/N & +/{-}{-}/{-} & N &  \\
C4076 & 1a<o0L / 2a<1a2 2b*<1a1 / 3a<2a0L,2b2 & N/NN/N & +/{-}{-}/{-} & N &  \\
C4077 & 1a<o0L / 2a<1a2 2b*<1a1 / 3a<2a1,2b1 & N/NN/N & +/{-}{-}/{-} & N &  \\
C4078 & 1a<o0L / 2a<1a2 2b*<1a1 / 3a<2a1,2b2 & N/NN/N & +/{-}{-}/{-} & N &  \\
C4079 & 1a<o0L / 2a<1a2 2b*<1a1 / 3a<2a1L,2b1 & N/NN/N & +/{-}{-}/{-} & N &  \\
C4080 & 1a<o0L / 2a<1a2 2b*<1a1 / 3a<2a1L,2b2 & N/NN/N & +/{-}{-}/{-} & N &  \\
C4081 & 1a<o0L / 2a<1a2 2b*<1a1 / 3a<2a2,2b1 & N/NN/N & +/{-}{-}/{-} & N &  \\
C4082 & 1a<o0L / 2a<1a2 2b*<1a1 / 3a<2a2,2b2 & N/NN/N & +/{-}{-}/{-} & N &  \\
C4083 & 1a<o0L / 2a<1a2 2b*<1a1 / 3a<2a2L,2b1 & N/NN/N & +/{-}{-}/{-} & N &  \\
C4084 & 1a<o0L / 2a<1a2 2b*<1a1 / 3a<2a2L,2b2 & N/NN/N & +/{-}{-}/{-} & N &  \\
C4085 & 1a<o0L / 2a<1a2 2b*<1a2 / 3a<2a0L,2b1 & N/NN/N & +/{-}{-}/{-} & N &  \\
C4086 & 1a<o0L / 2a<1a2 2b*<1a2 / 3a<2a0L,2b2 & N/NN/N & +/{-}{-}/{-} & N &  \\
C4087 & 1a<o0L / 2a<1a2 2b*<1a2 / 3a<2a1,2b1 & N/NN/N & +/{-}{-}/{-} & N &  \\
C4088 & 1a<o0L / 2a<1a2 2b*<1a2 / 3a<2a1,2b2 & N/NN/N & +/{-}{-}/{-} & N &  \\
C4089 & 1a<o0L / 2a<1a2 2b*<1a2 / 3a<2a1L,2b1 & N/NN/N & +/{-}{-}/{-} & N &  \\
C4090 & 1a<o0L / 2a<1a2 2b*<1a2 / 3a<2a1L,2b2 & N/NN/N & +/{-}{-}/{-} & N &  \\
C4091 & 1a<o0L / 2a<1a2 2b*<1a2 / 3a<2a2,2b1 & N/NN/N & +/{-}{-}/{-} & N &  \\
C4092 & 1a<o0L / 2a<1a2 2b*<1a2 / 3a<2a2,2b2 & N/NN/N & +/{-}{-}/{-} & N &  \\
C4093 & 1a<o0L / 2a<1a2 2b*<1a2 / 3a<2a2L,2b1 & N/NN/N & +/{-}{-}/{-} & N &  \\
C4094 & 1a<o0L / 2a<1a2 2b*<1a2 / 3a<2a2L,2b2 & N/NN/N & +/{-}{-}/{-} & N &  \\
C4095 & 1a<o0L / 2a<1a2L 2b*<1a1 / 3a<2a0L,2b1 & N/NN/N & +/+{-}/+ & N &  \\
C4096 & 1a<o0L / 2a<1a2L 2b*<1a1 / 3a<2a0L,2b2 & N/NN/N & +/+{-}/+ & N &  \\
C4097 & 1a<o0L / 2a<1a2L 2b*<1a1 / 3a<2a1,2b1 & N/NN/N & +/+{-}/{-} & N &  \\
C4098 & 1a<o0L / 2a<1a2L 2b*<1a1 / 3a<2a1,2b2 & N/NN/N & +/+{-}/{-} & N &  \\
C4099 & 1a<o0L / 2a<1a2L 2b*<1a1 / 3a<2a1L,2b1 & N/NN/N & +/+{-}/+ & N &  \\
C4100 & 1a<o0L / 2a<1a2L 2b*<1a1 / 3a<2a1L,2b2 & N/NN/N & +/+{-}/+ & N &  \\
C4101 & 1a<o0L / 2a<1a2L 2b*<1a1 / 3a<2a2,2b1 & N/NN/N & +/+{-}/{-} & N &  \\
C4102 & 1a<o0L / 2a<1a2L 2b*<1a1 / 3a<2a2,2b2 & N/NN/N & +/+{-}/{-} & N &  \\
C4103 & 1a<o0L / 2a<1a2L 2b*<1a1 / 3a<2a2L,2b1 & N/NN/N & +/+{-}/+ & N &  \\
C4104 & 1a<o0L / 2a<1a2L 2b*<1a1 / 3a<2a2L,2b2 & N/NN/N & +/+{-}/+ & N &  \\
C4105 & 1a<o0L / 2a<1a2L 2b*<1a2 / 3a<2a0L,2b1 & N/NN/N & +/+{-}/+ & N &  \\
C4106 & 1a<o0L / 2a<1a2L 2b*<1a2 / 3a<2a0L,2b2 & N/NN/N & +/+{-}/+ & N &  \\
C4107 & 1a<o0L / 2a<1a2L 2b*<1a2 / 3a<2a1,2b1 & N/NN/N & +/+{-}/{-} & N &  \\
C4108 & 1a<o0L / 2a<1a2L 2b*<1a2 / 3a<2a1,2b2 & N/NN/N & +/+{-}/{-} & N &  \\
C4109 & 1a<o0L / 2a<1a2L 2b*<1a2 / 3a<2a1L,2b1 & N/NN/N & +/+{-}/+ & N &  \\
C4110 & 1a<o0L / 2a<1a2L 2b*<1a2 / 3a<2a1L,2b2 & N/NN/N & +/+{-}/+ & N &  \\
C4111 & 1a<o0L / 2a<1a2L 2b*<1a2 / 3a<2a2,2b1 & N/NN/N & +/+{-}/{-} & N &  \\
C4112 & 1a<o0L / 2a<1a2L 2b*<1a2 / 3a<2a2,2b2 & N/NN/N & +/+{-}/{-} & N &  \\
C4113 & 1a<o0L / 2a<1a2L 2b*<1a2 / 3a<2a2L,2b1 & N/NN/N & +/+{-}/+ & N &  \\
C4114 & 1a<o0L / 2a<1a2L 2b*<1a2 / 3a<2a2L,2b2 & N/NN/N & +/+{-}/+ & N &  \\
C4115 & 1a<o0L 1b<o1 / 2a<1a0L,1b1 / 3a*<2a1 & NI/I/U & +{-}/+/{-} & N & W:3a \\
C4116 & 1a<o0L 1b<o1 / 2a<1a0L,1b1 / 3a*<2a2 & NI/I/U & +{-}/+/{-} & N & W:3a \\
C4117 & 1a<o0L 1b<o1 / 2a<1a0L,1b2 / 3a*<2a1 & NI/I/U & +{-}/+/{-} & N & W:3a \\
C4118 & 1a<o0L 1b<o1 / 2a<1a0L,1b2 / 3a*<2a2 & NI/I/U & +{-}/+/{-} & N & W:3a \\
C4119 & 1a<o0L 1b<o1 / 2a<1a1,1b0L / 3a*<2a1 & NI/N/N & +{-}/{-}/{-} & N &  \\
C4120 & 1a<o0L 1b<o1 / 2a<1a1,1b0L / 3a*<2a2 & NI/N/N & +{-}/{-}/{-} & N &  \\
C4121 & 1a<o0L 1b<o1 / 2a<1a1,1b1 / 3a*<2a1 & NI/E/U & +{-}/{-}/{-} & F & W:2a,3a \\
C4122 & 1a<o0L 1b<o1 / 2a<1a1,1b1 / 3a*<2a2 & NI/E/U & +{-}/{-}/{-} & F & W:2a,3a \\
C4123 & 1a<o0L 1b<o1 / 2a<1a1,1b1L / 3a*<2a1 & NI/I/U & +{-}/{-}/{-} & F & CP:2a W:3a \\
C4124 & 1a<o0L 1b<o1 / 2a<1a1,1b1L / 3a*<2a2 & NI/I/U & +{-}/{-}/{-} & F & CP:2a W:3a \\
C4125 & 1a<o0L 1b<o1 / 2a<1a1,1b2 / 3a*<2a1 & NI/E/U & +{-}/{-}/{-} & F & C:2a W:2a,3a \\
C4126 & 1a<o0L 1b<o1 / 2a<1a1,1b2 / 3a*<2a2 & NI/E/U & +{-}/{-}/{-} & F & C:2a W:2a,3a \\
C4127 & 1a<o0L 1b<o1 / 2a<1a1,1b2L / 3a*<2a1 & NI/I/U & +{-}/{-}/{-} & F & P:2a W:3a \\
C4128 & 1a<o0L 1b<o1 / 2a<1a1,1b2L / 3a*<2a2 & NI/I/U & +{-}/{-}/{-} & F & P:2a W:3a \\
C4129 & 1a<o0L 1b<o1 / 2a<1a1L,1b1 / 3a*<2a1 & NI/E/U & +{-}/+/{-} & F & W:2a,3a \\
C4130 & 1a<o0L 1b<o1 / 2a<1a1L,1b1 / 3a*<2a2 & NI/E/U & +{-}/+/{-} & F & W:2a,3a \\
C4131 & 1a<o0L 1b<o1 / 2a<1a1L,1b2 / 3a*<2a1 & NI/E/U & +{-}/+/{-} & F & C:2a W:2a,3a \\
C4132 & 1a<o0L 1b<o1 / 2a<1a1L,1b2 / 3a*<2a2 & NI/E/U & +{-}/+/{-} & F & C:2a W:2a,3a \\
C4133 & 1a<o0L 1b<o1 / 2a<1a2,1b0L / 3a*<2a1 & NI/N/N & +{-}/{-}/{-} & N &  \\
C4134 & 1a<o0L 1b<o1 / 2a<1a2,1b0L / 3a*<2a2 & NI/N/N & +{-}/{-}/{-} & N &  \\
C4135 & 1a<o0L 1b<o1 / 2a<1a2,1b1 / 3a*<2a1 & NI/E/U & +{-}/{-}/{-} & F & W:2a,3a \\
C4136 & 1a<o0L 1b<o1 / 2a<1a2,1b1 / 3a*<2a2 & NI/E/U & +{-}/{-}/{-} & F & W:2a,3a \\
C4137 & 1a<o0L 1b<o1 / 2a<1a2,1b1L / 3a*<2a1 & NI/I/U & +{-}/{-}/{-} & F & CP:2a W:3a \\
C4138 & 1a<o0L 1b<o1 / 2a<1a2,1b1L / 3a*<2a2 & NI/I/U & +{-}/{-}/{-} & F & CP:2a W:3a \\
C4139 & 1a<o0L 1b<o1 / 2a<1a2,1b2 / 3a*<2a1 & NI/E/U & +{-}/{-}/{-} & F & W:2a,3a \\
C4140 & 1a<o0L 1b<o1 / 2a<1a2,1b2 / 3a*<2a2 & NI/E/U & +{-}/{-}/{-} & F & W:2a,3a \\
C4141 & 1a<o0L 1b<o1 / 2a<1a2,1b2L / 3a*<2a1 & NI/I/U & +{-}/{-}/{-} & F & CP:2a W:3a \\
C4142 & 1a<o0L 1b<o1 / 2a<1a2,1b2L / 3a*<2a2 & NI/I/U & +{-}/{-}/{-} & F & CP:2a W:3a \\
C4143 & 1a<o0L 1b<o1 / 2a<1a2L,1b1 / 3a*<2a1 & NI/E/U & +{-}/+/{-} & F & W:2a,3a \\
C4144 & 1a<o0L 1b<o1 / 2a<1a2L,1b1 / 3a*<2a2 & NI/E/U & +{-}/+/{-} & F & W:2a,3a \\
C4145 & 1a<o0L 1b<o1 / 2a<1a2L,1b2 / 3a*<2a1 & NI/E/U & +{-}/+/{-} & F & W:2a,3a \\
C4146 & 1a<o0L 1b<o1 / 2a<1a2L,1b2 / 3a*<2a2 & NI/E/U & +{-}/+/{-} & F & W:2a,3a \\
C4147 & 1a<o0L 1b<o1 / 2a*<1a1,1b1 / 3a<2a1 & NI/U/U & +{-}/{-}/{-} & N & W:2a,3a \\
C4148 & 1a<o0L 1b<o1 / 2a*<1a1,1b1 / 3a<2a2 & NI/U/U & +{-}/{-}/{-} & N & W:2a,3a \\
C4149 & 1a<o0L 1b<o1 / 2a*<1a1,1b2 / 3a<2a1 & NI/U/U & +{-}/{-}/{-} & N & W:2a,3a \\
C4150 & 1a<o0L 1b<o1 / 2a*<1a1,1b2 / 3a<2a2 & NI/U/U & +{-}/{-}/{-} & N & W:2a,3a \\
C4151 & 1a<o0L 1b<o1 / 2a*<1a2,1b1 / 3a<2a1 & NI/U/U & +{-}/{-}/{-} & N & W:2a,3a \\
C4152 & 1a<o0L 1b<o1 / 2a*<1a2,1b1 / 3a<2a2 & NI/U/U & +{-}/{-}/{-} & N & W:2a,3a \\
C4153 & 1a<o0L 1b<o1 / 2a*<1a2,1b2 / 3a<2a1 & NI/U/U & +{-}/{-}/{-} & N & W:2a,3a \\
C4154 & 1a<o0L 1b<o1 / 2a*<1a2,1b2 / 3a<2a2 & NI/U/U & +{-}/{-}/{-} & N & W:2a,3a \\
C4155 & 1a<o0L 1b<o2 / 2a<1a0L,1b1 / 3a*<2a1 & NI/I/U & +{-}/+/{-} & N & W:3a \\
C4156 & 1a<o0L 1b<o2 / 2a<1a0L,1b1 / 3a*<2a2 & NI/I/U & +{-}/+/{-} & N & W:3a \\
C4157 & 1a<o0L 1b<o2 / 2a<1a0L,1b2 / 3a*<2a1 & NI/I/U & +{-}/+/{-} & N & W:3a \\
C4158 & 1a<o0L 1b<o2 / 2a<1a0L,1b2 / 3a*<2a2 & NI/I/U & +{-}/+/{-} & N & W:3a \\
C4159 & 1a<o0L 1b<o2 / 2a<1a1,1b0L / 3a*<2a1 & NI/N/N & +{-}/{-}/{-} & N &  \\
C4160 & 1a<o0L 1b<o2 / 2a<1a1,1b0L / 3a*<2a2 & NI/N/N & +{-}/{-}/{-} & N &  \\
C4161 & 1a<o0L 1b<o2 / 2a<1a1,1b1 / 3a*<2a1 & NI/E/U & +{-}/{-}/{-} & F & W:2a,3a \\
C4162 & 1a<o0L 1b<o2 / 2a<1a1,1b1 / 3a*<2a2 & NI/E/U & +{-}/{-}/{-} & F & W:2a,3a \\
C4163 & 1a<o0L 1b<o2 / 2a<1a1,1b1L / 3a*<2a1 & NI/I/U & +{-}/{-}/{-} & F & CP:2a W:3a \\
C4164 & 1a<o0L 1b<o2 / 2a<1a1,1b1L / 3a*<2a2 & NI/I/U & +{-}/{-}/{-} & F & CP:2a W:3a \\
C4165 & 1a<o0L 1b<o2 / 2a<1a1,1b2 / 3a*<2a1 & NI/E/U & +{-}/{-}/{-} & F & C:2a W:2a,3a \\
C4166 & 1a<o0L 1b<o2 / 2a<1a1,1b2 / 3a*<2a2 & NI/E/U & +{-}/{-}/{-} & F & C:2a W:2a,3a \\
C4167 & 1a<o0L 1b<o2 / 2a<1a1,1b2L / 3a*<2a1 & NI/I/U & +{-}/{-}/{-} & F & P:2a W:3a \\
C4168 & 1a<o0L 1b<o2 / 2a<1a1,1b2L / 3a*<2a2 & NI/I/U & +{-}/{-}/{-} & F & P:2a W:3a \\
C4169 & 1a<o0L 1b<o2 / 2a<1a1L,1b1 / 3a*<2a1 & NI/E/U & +{-}/+/{-} & F & W:2a,3a \\
C4170 & 1a<o0L 1b<o2 / 2a<1a1L,1b1 / 3a*<2a2 & NI/E/U & +{-}/+/{-} & F & W:2a,3a \\
C4171 & 1a<o0L 1b<o2 / 2a<1a1L,1b2 / 3a*<2a1 & NI/E/U & +{-}/+/{-} & F & C:2a W:2a,3a \\
C4172 & 1a<o0L 1b<o2 / 2a<1a1L,1b2 / 3a*<2a2 & NI/E/U & +{-}/+/{-} & F & C:2a W:2a,3a \\
C4173 & 1a<o0L 1b<o2 / 2a<1a2,1b0L / 3a*<2a1 & NI/N/N & +{-}/{-}/{-} & N &  \\
C4174 & 1a<o0L 1b<o2 / 2a<1a2,1b0L / 3a*<2a2 & NI/N/N & +{-}/{-}/{-} & N &  \\
C4175 & 1a<o0L 1b<o2 / 2a<1a2,1b1 / 3a*<2a1 & NI/E/U & +{-}/{-}/{-} & F & W:2a,3a \\
C4176 & 1a<o0L 1b<o2 / 2a<1a2,1b1 / 3a*<2a2 & NI/E/U & +{-}/{-}/{-} & F & W:2a,3a \\
C4177 & 1a<o0L 1b<o2 / 2a<1a2,1b1L / 3a*<2a1 & NI/I/U & +{-}/{-}/{-} & F & CP:2a W:3a \\
C4178 & 1a<o0L 1b<o2 / 2a<1a2,1b1L / 3a*<2a2 & NI/I/U & +{-}/{-}/{-} & F & CP:2a W:3a \\
C4179 & 1a<o0L 1b<o2 / 2a<1a2,1b2 / 3a*<2a1 & NI/E/U & +{-}/{-}/{-} & F & W:2a,3a \\
C4180 & 1a<o0L 1b<o2 / 2a<1a2,1b2 / 3a*<2a2 & NI/E/U & +{-}/{-}/{-} & F & W:2a,3a \\
C4181 & 1a<o0L 1b<o2 / 2a<1a2,1b2L / 3a*<2a1 & NI/I/U & +{-}/{-}/{-} & F & CP:2a W:3a \\
C4182 & 1a<o0L 1b<o2 / 2a<1a2,1b2L / 3a*<2a2 & NI/I/U & +{-}/{-}/{-} & F & CP:2a W:3a \\
C4183 & 1a<o0L 1b<o2 / 2a<1a2L,1b1 / 3a*<2a1 & NI/E/U & +{-}/+/{-} & F & W:2a,3a \\
C4184 & 1a<o0L 1b<o2 / 2a<1a2L,1b1 / 3a*<2a2 & NI/E/U & +{-}/+/{-} & F & W:2a,3a \\
C4185 & 1a<o0L 1b<o2 / 2a<1a2L,1b2 / 3a*<2a1 & NI/E/U & +{-}/+/{-} & F & W:2a,3a \\
C4186 & 1a<o0L 1b<o2 / 2a<1a2L,1b2 / 3a*<2a2 & NI/E/U & +{-}/+/{-} & F & W:2a,3a \\
C4187 & 1a<o0L 1b<o2 / 2a*<1a1,1b1 / 3a<2a1 & NI/U/U & +{-}/{-}/{-} & N & W:2a,3a \\
C4188 & 1a<o0L 1b<o2 / 2a*<1a1,1b1 / 3a<2a2 & NI/U/U & +{-}/{-}/{-} & N & W:2a,3a \\
C4189 & 1a<o0L 1b<o2 / 2a*<1a1,1b2 / 3a<2a1 & NI/U/U & +{-}/{-}/{-} & N & W:2a,3a \\
C4190 & 1a<o0L 1b<o2 / 2a*<1a1,1b2 / 3a<2a2 & NI/U/U & +{-}/{-}/{-} & N & W:2a,3a \\
C4191 & 1a<o0L 1b<o2 / 2a*<1a2,1b1 / 3a<2a1 & NI/U/U & +{-}/{-}/{-} & N & W:2a,3a \\
C4192 & 1a<o0L 1b<o2 / 2a*<1a2,1b1 / 3a<2a2 & NI/U/U & +{-}/{-}/{-} & N & W:2a,3a \\
C4193 & 1a<o0L 1b<o2 / 2a*<1a2,1b2 / 3a<2a1 & NI/U/U & +{-}/{-}/{-} & N & W:2a,3a \\
C4194 & 1a<o0L 1b<o2 / 2a*<1a2,1b2 / 3a<2a2 & NI/U/U & +{-}/{-}/{-} & N & W:2a,3a \\
C4195 & 1a<o0L 1b*<o1 / 2a<1a0L,1b1 / 3a<2a0L & NU/U/N & +{-}/+/+ & N & W:1b,2a \\
C4196 & 1a<o0L 1b*<o1 / 2a<1a0L,1b1 / 3a<2a1 & NU/U/U & +{-}/+/{-} & N & W:1b,2a,3a \\
C4197 & 1a<o0L 1b*<o1 / 2a<1a0L,1b1 / 3a<2a1L & NU/U/U & +{-}/+/+ & N & W:1b,2a,3a \\
C4198 & 1a<o0L 1b*<o1 / 2a<1a0L,1b1 / 3a<2a2 & NU/U/U & +{-}/+/{-} & N & W:1b,2a,3a \\
C4199 & 1a<o0L 1b*<o1 / 2a<1a0L,1b1 / 3a<2a2L & NU/U/U & +{-}/+/+ & N & W:1b,2a,3a \\
C4200 & 1a<o0L 1b*<o1 / 2a<1a0L,1b2 / 3a<2a0L & NU/U/N & +{-}/+/+ & N & W:1b,2a \\
C4201 & 1a<o0L 1b*<o1 / 2a<1a0L,1b2 / 3a<2a1 & NU/U/U & +{-}/+/{-} & N & W:1b,2a,3a \\
C4202 & 1a<o0L 1b*<o1 / 2a<1a0L,1b2 / 3a<2a1L & NU/U/U & +{-}/+/+ & N & W:1b,2a,3a \\
C4203 & 1a<o0L 1b*<o1 / 2a<1a0L,1b2 / 3a<2a2 & NU/U/U & +{-}/+/{-} & N & W:1b,2a,3a \\
C4204 & 1a<o0L 1b*<o1 / 2a<1a0L,1b2 / 3a<2a2L & NU/U/U & +{-}/+/+ & N & W:1b,2a,3a \\
C4205 & 1a<o0L 1b*<o1 / 2a<1a1,1b1 / 3a<2a0L & NU/U/N & +{-}/{-}/{-} & N & W:1b,2a \\
C4206 & 1a<o0L 1b*<o1 / 2a<1a1,1b1 / 3a<2a1 & NU/U/U & +{-}/{-}/{-} & N & W:1b,2a,3a \\
C4207 & 1a<o0L 1b*<o1 / 2a<1a1,1b1 / 3a<2a1L & NU/U/U & +{-}/{-}/{-} & N & W:1b,2a,3a \\
C4208 & 1a<o0L 1b*<o1 / 2a<1a1,1b1 / 3a<2a2 & NU/U/U & +{-}/{-}/{-} & N & W:1b,2a,3a \\
C4209 & 1a<o0L 1b*<o1 / 2a<1a1,1b1 / 3a<2a2L & NU/U/U & +{-}/{-}/{-} & N & W:1b,2a,3a \\
C4210 & 1a<o0L 1b*<o1 / 2a<1a1,1b2 / 3a<2a0L & NU/U/N & +{-}/{-}/{-} & N & W:1b,2a \\
C4211 & 1a<o0L 1b*<o1 / 2a<1a1,1b2 / 3a<2a1 & NU/U/U & +{-}/{-}/{-} & N & W:1b,2a,3a \\
C4212 & 1a<o0L 1b*<o1 / 2a<1a1,1b2 / 3a<2a1L & NU/U/U & +{-}/{-}/{-} & N & W:1b,2a,3a \\
C4213 & 1a<o0L 1b*<o1 / 2a<1a1,1b2 / 3a<2a2 & NU/U/U & +{-}/{-}/{-} & N & W:1b,2a,3a \\
C4214 & 1a<o0L 1b*<o1 / 2a<1a1,1b2 / 3a<2a2L & NU/U/U & +{-}/{-}/{-} & N & W:1b,2a,3a \\
C4215 & 1a<o0L 1b*<o1 / 2a<1a1L,1b1 / 3a<2a0L & NU/U/N & +{-}/+/+ & N & W:1b,2a \\
C4216 & 1a<o0L 1b*<o1 / 2a<1a1L,1b1 / 3a<2a1 & NU/U/U & +{-}/+/{-} & N & W:1b,2a,3a \\
C4217 & 1a<o0L 1b*<o1 / 2a<1a1L,1b1 / 3a<2a1L & NU/U/U & +{-}/+/+ & N & W:1b,2a,3a \\
C4218 & 1a<o0L 1b*<o1 / 2a<1a1L,1b1 / 3a<2a2 & NU/U/U & +{-}/+/{-} & N & W:1b,2a,3a \\
C4219 & 1a<o0L 1b*<o1 / 2a<1a1L,1b1 / 3a<2a2L & NU/U/U & +{-}/+/+ & N & W:1b,2a,3a \\
C4220 & 1a<o0L 1b*<o1 / 2a<1a1L,1b2 / 3a<2a0L & NU/U/N & +{-}/+/+ & N & W:1b,2a \\
C4221 & 1a<o0L 1b*<o1 / 2a<1a1L,1b2 / 3a<2a1 & NU/U/U & +{-}/+/{-} & N & W:1b,2a,3a \\
C4222 & 1a<o0L 1b*<o1 / 2a<1a1L,1b2 / 3a<2a1L & NU/U/U & +{-}/+/+ & N & W:1b,2a,3a \\
C4223 & 1a<o0L 1b*<o1 / 2a<1a1L,1b2 / 3a<2a2 & NU/U/U & +{-}/+/{-} & N & W:1b,2a,3a \\
C4224 & 1a<o0L 1b*<o1 / 2a<1a1L,1b2 / 3a<2a2L & NU/U/U & +{-}/+/+ & N & W:1b,2a,3a \\
C4225 & 1a<o0L 1b*<o1 / 2a<1a2,1b1 / 3a<2a0L & NU/U/N & +{-}/{-}/{-} & N & W:1b,2a \\
C4226 & 1a<o0L 1b*<o1 / 2a<1a2,1b1 / 3a<2a1 & NU/U/U & +{-}/{-}/{-} & N & W:1b,2a,3a \\
C4227 & 1a<o0L 1b*<o1 / 2a<1a2,1b1 / 3a<2a1L & NU/U/U & +{-}/{-}/{-} & N & W:1b,2a,3a \\
C4228 & 1a<o0L 1b*<o1 / 2a<1a2,1b1 / 3a<2a2 & NU/U/U & +{-}/{-}/{-} & N & W:1b,2a,3a \\
C4229 & 1a<o0L 1b*<o1 / 2a<1a2,1b1 / 3a<2a2L & NU/U/U & +{-}/{-}/{-} & N & W:1b,2a,3a \\
C4230 & 1a<o0L 1b*<o1 / 2a<1a2,1b2 / 3a<2a0L & NU/U/N & +{-}/{-}/{-} & N & W:1b,2a \\
C4231 & 1a<o0L 1b*<o1 / 2a<1a2,1b2 / 3a<2a1 & NU/U/U & +{-}/{-}/{-} & N & W:1b,2a,3a \\
C4232 & 1a<o0L 1b*<o1 / 2a<1a2,1b2 / 3a<2a1L & NU/U/U & +{-}/{-}/{-} & N & W:1b,2a,3a \\
C4233 & 1a<o0L 1b*<o1 / 2a<1a2,1b2 / 3a<2a2 & NU/U/U & +{-}/{-}/{-} & N & W:1b,2a,3a \\
C4234 & 1a<o0L 1b*<o1 / 2a<1a2,1b2 / 3a<2a2L & NU/U/U & +{-}/{-}/{-} & N & W:1b,2a,3a \\
C4235 & 1a<o0L 1b*<o1 / 2a<1a2L,1b1 / 3a<2a0L & NU/U/N & +{-}/+/+ & N & W:1b,2a \\
C4236 & 1a<o0L 1b*<o1 / 2a<1a2L,1b1 / 3a<2a1 & NU/U/U & +{-}/+/{-} & N & W:1b,2a,3a \\
C4237 & 1a<o0L 1b*<o1 / 2a<1a2L,1b1 / 3a<2a1L & NU/U/U & +{-}/+/+ & N & W:1b,2a,3a \\
C4238 & 1a<o0L 1b*<o1 / 2a<1a2L,1b1 / 3a<2a2 & NU/U/U & +{-}/+/{-} & N & W:1b,2a,3a \\
C4239 & 1a<o0L 1b*<o1 / 2a<1a2L,1b1 / 3a<2a2L & NU/U/U & +{-}/+/+ & N & W:1b,2a,3a \\
C4240 & 1a<o0L 1b*<o1 / 2a<1a2L,1b2 / 3a<2a0L & NU/U/N & +{-}/+/+ & N & W:1b,2a \\
C4241 & 1a<o0L 1b*<o1 / 2a<1a2L,1b2 / 3a<2a1 & NU/U/U & +{-}/+/{-} & N & W:1b,2a,3a \\
C4242 & 1a<o0L 1b*<o1 / 2a<1a2L,1b2 / 3a<2a1L & NU/U/U & +{-}/+/+ & N & W:1b,2a,3a \\
C4243 & 1a<o0L 1b*<o1 / 2a<1a2L,1b2 / 3a<2a2 & NU/U/U & +{-}/+/{-} & N & W:1b,2a,3a \\
C4244 & 1a<o0L 1b*<o1 / 2a<1a2L,1b2 / 3a<2a2L & NU/U/U & +{-}/+/+ & N & W:1b,2a,3a \\
C4245 & 1a<o0L 1b*<o2 / 2a<1a0L,1b1 / 3a<2a0L & NU/U/N & +{-}/+/+ & N & W:1b,2a \\
C4246 & 1a<o0L 1b*<o2 / 2a<1a0L,1b1 / 3a<2a1 & NU/U/U & +{-}/+/{-} & N & W:1b,2a,3a \\
C4247 & 1a<o0L 1b*<o2 / 2a<1a0L,1b1 / 3a<2a1L & NU/U/U & +{-}/+/+ & N & W:1b,2a,3a \\
C4248 & 1a<o0L 1b*<o2 / 2a<1a0L,1b1 / 3a<2a2 & NU/U/U & +{-}/+/{-} & N & W:1b,2a,3a \\
C4249 & 1a<o0L 1b*<o2 / 2a<1a0L,1b1 / 3a<2a2L & NU/U/U & +{-}/+/+ & N & W:1b,2a,3a \\
C4250 & 1a<o0L 1b*<o2 / 2a<1a0L,1b2 / 3a<2a0L & NU/U/N & +{-}/+/+ & N & W:1b,2a \\
C4251 & 1a<o0L 1b*<o2 / 2a<1a0L,1b2 / 3a<2a1 & NU/U/U & +{-}/+/{-} & N & W:1b,2a,3a \\
C4252 & 1a<o0L 1b*<o2 / 2a<1a0L,1b2 / 3a<2a1L & NU/U/U & +{-}/+/+ & N & W:1b,2a,3a \\
C4253 & 1a<o0L 1b*<o2 / 2a<1a0L,1b2 / 3a<2a2 & NU/U/U & +{-}/+/{-} & N & W:1b,2a,3a \\
C4254 & 1a<o0L 1b*<o2 / 2a<1a0L,1b2 / 3a<2a2L & NU/U/U & +{-}/+/+ & N & W:1b,2a,3a \\
C4255 & 1a<o0L 1b*<o2 / 2a<1a1,1b1 / 3a<2a0L & NU/U/N & +{-}/{-}/{-} & N & W:1b,2a \\
C4256 & 1a<o0L 1b*<o2 / 2a<1a1,1b1 / 3a<2a1 & NU/U/U & +{-}/{-}/{-} & N & W:1b,2a,3a \\
C4257 & 1a<o0L 1b*<o2 / 2a<1a1,1b1 / 3a<2a1L & NU/U/U & +{-}/{-}/{-} & N & W:1b,2a,3a \\
C4258 & 1a<o0L 1b*<o2 / 2a<1a1,1b1 / 3a<2a2 & NU/U/U & +{-}/{-}/{-} & N & W:1b,2a,3a \\
C4259 & 1a<o0L 1b*<o2 / 2a<1a1,1b1 / 3a<2a2L & NU/U/U & +{-}/{-}/{-} & N & W:1b,2a,3a \\
C4260 & 1a<o0L 1b*<o2 / 2a<1a1,1b2 / 3a<2a0L & NU/U/N & +{-}/{-}/{-} & N & W:1b,2a \\
C4261 & 1a<o0L 1b*<o2 / 2a<1a1,1b2 / 3a<2a1 & NU/U/U & +{-}/{-}/{-} & N & W:1b,2a,3a \\
C4262 & 1a<o0L 1b*<o2 / 2a<1a1,1b2 / 3a<2a1L & NU/U/U & +{-}/{-}/{-} & N & W:1b,2a,3a \\
C4263 & 1a<o0L 1b*<o2 / 2a<1a1,1b2 / 3a<2a2 & NU/U/U & +{-}/{-}/{-} & N & W:1b,2a,3a \\
C4264 & 1a<o0L 1b*<o2 / 2a<1a1,1b2 / 3a<2a2L & NU/U/U & +{-}/{-}/{-} & N & W:1b,2a,3a \\
C4265 & 1a<o0L 1b*<o2 / 2a<1a1L,1b1 / 3a<2a0L & NU/U/N & +{-}/+/+ & N & W:1b,2a \\
C4266 & 1a<o0L 1b*<o2 / 2a<1a1L,1b1 / 3a<2a1 & NU/U/U & +{-}/+/{-} & N & W:1b,2a,3a \\
C4267 & 1a<o0L 1b*<o2 / 2a<1a1L,1b1 / 3a<2a1L & NU/U/U & +{-}/+/+ & N & W:1b,2a,3a \\
C4268 & 1a<o0L 1b*<o2 / 2a<1a1L,1b1 / 3a<2a2 & NU/U/U & +{-}/+/{-} & N & W:1b,2a,3a \\
C4269 & 1a<o0L 1b*<o2 / 2a<1a1L,1b1 / 3a<2a2L & NU/U/U & +{-}/+/+ & N & W:1b,2a,3a \\
C4270 & 1a<o0L 1b*<o2 / 2a<1a1L,1b2 / 3a<2a0L & NU/U/N & +{-}/+/+ & N & W:1b,2a \\
C4271 & 1a<o0L 1b*<o2 / 2a<1a1L,1b2 / 3a<2a1 & NU/U/U & +{-}/+/{-} & N & W:1b,2a,3a \\
C4272 & 1a<o0L 1b*<o2 / 2a<1a1L,1b2 / 3a<2a1L & NU/U/U & +{-}/+/+ & N & W:1b,2a,3a \\
C4273 & 1a<o0L 1b*<o2 / 2a<1a1L,1b2 / 3a<2a2 & NU/U/U & +{-}/+/{-} & N & W:1b,2a,3a \\
C4274 & 1a<o0L 1b*<o2 / 2a<1a1L,1b2 / 3a<2a2L & NU/U/U & +{-}/+/+ & N & W:1b,2a,3a \\
C4275 & 1a<o0L 1b*<o2 / 2a<1a2,1b1 / 3a<2a0L & NU/U/N & +{-}/{-}/{-} & N & W:1b,2a \\
C4276 & 1a<o0L 1b*<o2 / 2a<1a2,1b1 / 3a<2a1 & NU/U/U & +{-}/{-}/{-} & N & W:1b,2a,3a \\
C4277 & 1a<o0L 1b*<o2 / 2a<1a2,1b1 / 3a<2a1L & NU/U/U & +{-}/{-}/{-} & N & W:1b,2a,3a \\
C4278 & 1a<o0L 1b*<o2 / 2a<1a2,1b1 / 3a<2a2 & NU/U/U & +{-}/{-}/{-} & N & W:1b,2a,3a \\
C4279 & 1a<o0L 1b*<o2 / 2a<1a2,1b1 / 3a<2a2L & NU/U/U & +{-}/{-}/{-} & N & W:1b,2a,3a \\
C4280 & 1a<o0L 1b*<o2 / 2a<1a2,1b2 / 3a<2a0L & NU/U/N & +{-}/{-}/{-} & N & W:1b,2a \\
C4281 & 1a<o0L 1b*<o2 / 2a<1a2,1b2 / 3a<2a1 & NU/U/U & +{-}/{-}/{-} & N & W:1b,2a,3a \\
C4282 & 1a<o0L 1b*<o2 / 2a<1a2,1b2 / 3a<2a1L & NU/U/U & +{-}/{-}/{-} & N & W:1b,2a,3a \\
C4283 & 1a<o0L 1b*<o2 / 2a<1a2,1b2 / 3a<2a2 & NU/U/U & +{-}/{-}/{-} & N & W:1b,2a,3a \\
C4284 & 1a<o0L 1b*<o2 / 2a<1a2,1b2 / 3a<2a2L & NU/U/U & +{-}/{-}/{-} & N & W:1b,2a,3a \\
C4285 & 1a<o0L 1b*<o2 / 2a<1a2L,1b1 / 3a<2a0L & NU/U/N & +{-}/+/+ & N & W:1b,2a \\
C4286 & 1a<o0L 1b*<o2 / 2a<1a2L,1b1 / 3a<2a1 & NU/U/U & +{-}/+/{-} & N & W:1b,2a,3a \\
C4287 & 1a<o0L 1b*<o2 / 2a<1a2L,1b1 / 3a<2a1L & NU/U/U & +{-}/+/+ & N & W:1b,2a,3a \\
C4288 & 1a<o0L 1b*<o2 / 2a<1a2L,1b1 / 3a<2a2 & NU/U/U & +{-}/+/{-} & N & W:1b,2a,3a \\
C4289 & 1a<o0L 1b*<o2 / 2a<1a2L,1b1 / 3a<2a2L & NU/U/U & +{-}/+/+ & N & W:1b,2a,3a \\
C4290 & 1a<o0L 1b*<o2 / 2a<1a2L,1b2 / 3a<2a0L & NU/U/N & +{-}/+/+ & N & W:1b,2a \\
C4291 & 1a<o0L 1b*<o2 / 2a<1a2L,1b2 / 3a<2a1 & NU/U/U & +{-}/+/{-} & N & W:1b,2a,3a \\
C4292 & 1a<o0L 1b*<o2 / 2a<1a2L,1b2 / 3a<2a1L & NU/U/U & +{-}/+/+ & N & W:1b,2a,3a \\
C4293 & 1a<o0L 1b*<o2 / 2a<1a2L,1b2 / 3a<2a2 & NU/U/U & +{-}/+/{-} & N & W:1b,2a,3a \\
C4294 & 1a<o0L 1b*<o2 / 2a<1a2L,1b2 / 3a<2a2L & NU/U/U & +{-}/+/+ & N & W:1b,2a,3a \\
C4295 & 1a<o1 / 2a<1a0L 2b<1a1 / 3a*<2a1,2b1 & I/NI/U & {-}/{-}{-}/{-} & N & W:3a \\
C4296 & 1a<o1 / 2a<1a0L 2b<1a1 / 3a*<2a1,2b2 & I/NI/U & {-}/{-}{-}/{-} & N & W:3a \\
C4297 & 1a<o1 / 2a<1a0L 2b<1a1 / 3a*<2a2,2b1 & I/NI/U & {-}/{-}{-}/{-} & N & W:3a \\
C4298 & 1a<o1 / 2a<1a0L 2b<1a1 / 3a*<2a2,2b2 & I/NI/U & {-}/{-}{-}/{-} & N & W:3a \\
C4299 & 1a<o1 / 2a<1a0L 2b<1a2 / 3a*<2a1,2b1 & I/NI/U & {-}/{-}{-}/{-} & N & W:3a \\
C4300 & 1a<o1 / 2a<1a0L 2b<1a2 / 3a*<2a1,2b2 & I/NI/U & {-}/{-}{-}/{-} & N & W:3a \\
C4301 & 1a<o1 / 2a<1a0L 2b<1a2 / 3a*<2a2,2b1 & I/NI/U & {-}/{-}{-}/{-} & N & W:3a \\
C4302 & 1a<o1 / 2a<1a0L 2b<1a2 / 3a*<2a2,2b2 & I/NI/U & {-}/{-}{-}/{-} & N & W:3a \\
C4303 & 1a<o1 / 2a<1a0L 2b*<1a1 / 3a<2a0L,2b1 & I/NU/U & {-}/{-}{-}/{-} & N & W:2b,3a \\
C4304 & 1a<o1 / 2a<1a0L 2b*<1a1 / 3a<2a0L,2b2 & I/NU/U & {-}/{-}{-}/{-} & N & W:2b,3a \\
C4305 & 1a<o1 / 2a<1a0L 2b*<1a1 / 3a<2a1,2b1 & I/NU/U & {-}/{-}{-}/{-} & N & W:2b,3a \\
C4306 & 1a<o1 / 2a<1a0L 2b*<1a1 / 3a<2a1,2b2 & I/NU/U & {-}/{-}{-}/{-} & N & W:2b,3a \\
C4307 & 1a<o1 / 2a<1a0L 2b*<1a1 / 3a<2a1L,2b1 & I/NU/U & {-}/{-}{-}/{-} & N & W:2b,3a \\
C4308 & 1a<o1 / 2a<1a0L 2b*<1a1 / 3a<2a1L,2b2 & I/NU/U & {-}/{-}{-}/{-} & N & W:2b,3a \\
C4309 & 1a<o1 / 2a<1a0L 2b*<1a1 / 3a<2a2,2b1 & I/NU/U & {-}/{-}{-}/{-} & N & W:2b,3a \\
C4310 & 1a<o1 / 2a<1a0L 2b*<1a1 / 3a<2a2,2b2 & I/NU/U & {-}/{-}{-}/{-} & N & W:2b,3a \\
C4311 & 1a<o1 / 2a<1a0L 2b*<1a1 / 3a<2a2L,2b1 & I/NU/U & {-}/{-}{-}/{-} & N & W:2b,3a \\
C4312 & 1a<o1 / 2a<1a0L 2b*<1a1 / 3a<2a2L,2b2 & I/NU/U & {-}/{-}{-}/{-} & N & W:2b,3a \\
C4313 & 1a<o1 / 2a<1a0L 2b*<1a2 / 3a<2a0L,2b1 & I/NU/U & {-}/{-}{-}/{-} & N & W:2b,3a \\
C4314 & 1a<o1 / 2a<1a0L 2b*<1a2 / 3a<2a0L,2b2 & I/NU/U & {-}/{-}{-}/{-} & N & W:2b,3a \\
C4315 & 1a<o1 / 2a<1a0L 2b*<1a2 / 3a<2a1,2b1 & I/NU/U & {-}/{-}{-}/{-} & N & W:2b,3a \\
C4316 & 1a<o1 / 2a<1a0L 2b*<1a2 / 3a<2a1,2b2 & I/NU/U & {-}/{-}{-}/{-} & N & W:2b,3a \\
C4317 & 1a<o1 / 2a<1a0L 2b*<1a2 / 3a<2a1L,2b1 & I/NU/U & {-}/{-}{-}/{-} & N & W:2b,3a \\
C4318 & 1a<o1 / 2a<1a0L 2b*<1a2 / 3a<2a1L,2b2 & I/NU/U & {-}/{-}{-}/{-} & N & W:2b,3a \\
C4319 & 1a<o1 / 2a<1a0L 2b*<1a2 / 3a<2a2,2b1 & I/NU/U & {-}/{-}{-}/{-} & N & W:2b,3a \\
C4320 & 1a<o1 / 2a<1a0L 2b*<1a2 / 3a<2a2,2b2 & I/NU/U & {-}/{-}{-}/{-} & N & W:2b,3a \\
C4321 & 1a<o1 / 2a<1a0L 2b*<1a2 / 3a<2a2L,2b1 & I/NU/U & {-}/{-}{-}/{-} & N & W:2b,3a \\
C4322 & 1a<o1 / 2a<1a0L 2b*<1a2 / 3a<2a2L,2b2 & I/NU/U & {-}/{-}{-}/{-} & N & W:2b,3a \\
C4323 & 1a<o1 / 2a<1a1 2b<1a1 / 3a*<2a1,2b1 & I/EE/U & {-}/{-}{-}/{-} & E & W:1a \\
C4324 & 1a<o1 / 2a<1a1 2b<1a1 / 3a*<2a1,2b2 & I/EE/U & {-}/{-}{-}/{-} & E & W:1a \\
C4325 & 1a<o1 / 2a<1a1 2b<1a1 / 3a*<2a2,2b2 & I/EE/U & {-}/{-}{-}/{-} & E & W:1a \\
C4326 & 1a<o1 / 2a<1a1 2b<1a1L / 3a*<2a1,2b1 & I/LI/U & {-}/{-}{-}/{-} & B & CP:2b W:1a,2b \\
C4327 & 1a<o1 / 2a<1a1 2b<1a1L / 3a*<2a1,2b2 & I/LI/U & {-}/{-}{-}/{-} & B & CP:2b W:1a,2b \\
C4328 & 1a<o1 / 2a<1a1 2b<1a1L / 3a*<2a2,2b1 & I/LI/U & {-}/{-}{-}/{-} & B & CP:2b W:1a,2b \\
C4329 & 1a<o1 / 2a<1a1 2b<1a1L / 3a*<2a2,2b2 & I/LI/U & {-}/{-}{-}/{-} & B & CP:2b W:1a,2b \\
C4330 & 1a<o1 / 2a<1a1 2b<1a2 / 3a*<2a1,2b1 & I/EE/U & {-}/{-}{-}/{-} & E & C:2b W:1a \\
C4331 & 1a<o1 / 2a<1a1 2b<1a2 / 3a*<2a1,2b2 & I/EE/U & {-}/{-}{-}/{-} & E & C:2b W:1a \\
C4332 & 1a<o1 / 2a<1a1 2b<1a2 / 3a*<2a2,2b1 & I/EE/U & {-}/{-}{-}/{-} & E & C:2b W:1a \\
C4333 & 1a<o1 / 2a<1a1 2b<1a2 / 3a*<2a2,2b2 & I/EE/U & {-}/{-}{-}/{-} & E & C:2b W:1a \\
C4334 & 1a<o1 / 2a<1a1 2b<1a2L / 3a*<2a1,2b1 & I/LI/U & {-}/{-}{-}/{-} & B & P:2b W:1a,2b \\
C4335 & 1a<o1 / 2a<1a1 2b<1a2L / 3a*<2a1,2b2 & I/LI/U & {-}/{-}{-}/{-} & B & P:2b W:1a,2b \\
C4336 & 1a<o1 / 2a<1a1 2b<1a2L / 3a*<2a2,2b1 & I/LI/U & {-}/{-}{-}/{-} & B & P:2b W:1a,2b \\
C4337 & 1a<o1 / 2a<1a1 2b<1a2L / 3a*<2a2,2b2 & I/LI/U & {-}/{-}{-}/{-} & B & P:2b W:1a,2b \\
C4338 & 1a<o1 / 2a<1a1 2b*<1a1 / 3a<2a0L,2b1 & I/IU/U & {-}/{-}{-}/{-} & N & W:1a,2a \\
C4339 & 1a<o1 / 2a<1a1 2b*<1a1 / 3a<2a0L,2b2 & I/IU/U & {-}/{-}{-}/{-} & N & W:1a,2a \\
C4340 & 1a<o1 / 2a<1a1 2b*<1a1 / 3a<2a1,2b1 & I/IU/I & {-}/{-}{-}/{-} & N & W:1a,2a,3a \\
C4341 & 1a<o1 / 2a<1a1 2b*<1a1 / 3a<2a1,2b2 & I/IU/I & {-}/{-}{-}/{-} & N & W:1a,2a,3a \\
C4342 & 1a<o1 / 2a<1a1 2b*<1a1 / 3a<2a1L,2b1 & I/IU/I & {-}/{-}{-}/{-} & N & W:1a,2a,3a \\
C4343 & 1a<o1 / 2a<1a1 2b*<1a1 / 3a<2a1L,2b2 & I/IU/I & {-}/{-}{-}/{-} & N & W:1a,2a,3a \\
C4344 & 1a<o1 / 2a<1a1 2b*<1a1 / 3a<2a2,2b1 & I/IU/I & {-}/{-}{-}/{-} & N & W:1a,2a,3a \\
C4345 & 1a<o1 / 2a<1a1 2b*<1a1 / 3a<2a2,2b2 & I/IU/I & {-}/{-}{-}/{-} & N & W:1a,2a,3a \\
C4346 & 1a<o1 / 2a<1a1 2b*<1a1 / 3a<2a2L,2b1 & I/IU/I & {-}/{-}{-}/{-} & N & W:1a,2a,3a \\
C4347 & 1a<o1 / 2a<1a1 2b*<1a1 / 3a<2a2L,2b2 & I/IU/I & {-}/{-}{-}/{-} & N & W:1a,2a,3a \\
C4348 & 1a<o1 / 2a<1a1 2b*<1a2 / 3a<2a0L,2b1 & I/IU/U & {-}/{-}{-}/{-} & N & W:1a,2a \\
C4349 & 1a<o1 / 2a<1a1 2b*<1a2 / 3a<2a0L,2b2 & I/IU/U & {-}/{-}{-}/{-} & N & W:1a,2a \\
C4350 & 1a<o1 / 2a<1a1 2b*<1a2 / 3a<2a1,2b1 & I/IU/I & {-}/{-}{-}/{-} & N & W:1a,2a,3a \\
C4351 & 1a<o1 / 2a<1a1 2b*<1a2 / 3a<2a1,2b2 & I/IU/I & {-}/{-}{-}/{-} & N & W:1a,2a,3a \\
C4352 & 1a<o1 / 2a<1a1 2b*<1a2 / 3a<2a1L,2b1 & I/IU/I & {-}/{-}{-}/{-} & N & W:1a,2a,3a \\
C4353 & 1a<o1 / 2a<1a1 2b*<1a2 / 3a<2a1L,2b2 & I/IU/I & {-}/{-}{-}/{-} & N & W:1a,2a,3a \\
C4354 & 1a<o1 / 2a<1a1 2b*<1a2 / 3a<2a2,2b1 & I/IU/I & {-}/{-}{-}/{-} & N & W:1a,2a,3a \\
C4355 & 1a<o1 / 2a<1a1 2b*<1a2 / 3a<2a2,2b2 & I/IU/I & {-}/{-}{-}/{-} & N & W:1a,2a,3a \\
C4356 & 1a<o1 / 2a<1a1 2b*<1a2 / 3a<2a2L,2b1 & I/IU/I & {-}/{-}{-}/{-} & N & W:1a,2a,3a \\
C4357 & 1a<o1 / 2a<1a1 2b*<1a2 / 3a<2a2L,2b2 & I/IU/I & {-}/{-}{-}/{-} & N & W:1a,2a,3a \\
C4358 & 1a<o1 / 2a<1a1L 2b<1a2 / 3a*<2a1,2b1 & I/IL/U & {-}/{-}{-}/{-} & B & C:2a,2b P:2a W:1a,2a \\
C4359 & 1a<o1 / 2a<1a1L 2b<1a2 / 3a*<2a1,2b2 & I/IL/U & {-}/{-}{-}/{-} & B & C:2a,2b P:2a W:1a,2a \\
C4360 & 1a<o1 / 2a<1a1L 2b<1a2 / 3a*<2a2,2b1 & I/IL/U & {-}/{-}{-}/{-} & B & C:2a,2b P:2a W:1a,2a \\
C4361 & 1a<o1 / 2a<1a1L 2b<1a2 / 3a*<2a2,2b2 & I/IL/U & {-}/{-}{-}/{-} & B & C:2a,2b P:2a W:1a,2a \\
C4362 & 1a<o1 / 2a<1a1L 2b*<1a1 / 3a<2a0L,2b1 & I/IU/U & {-}/{-}{-}/{-} & N & W:1a,2a \\
C4363 & 1a<o1 / 2a<1a1L 2b*<1a1 / 3a<2a0L,2b2 & I/IU/U & {-}/{-}{-}/{-} & N & W:1a,2a \\
C4364 & 1a<o1 / 2a<1a1L 2b*<1a1 / 3a<2a1,2b1 & I/IU/I & {-}/{-}{-}/{-} & N & W:1a,2a,3a \\
C4365 & 1a<o1 / 2a<1a1L 2b*<1a1 / 3a<2a1,2b2 & I/IU/I & {-}/{-}{-}/{-} & N & W:1a,2a,3a \\
C4366 & 1a<o1 / 2a<1a1L 2b*<1a1 / 3a<2a1L,2b1 & I/IU/I & {-}/{-}{-}/{-} & N & W:1a,2a,3a \\
C4367 & 1a<o1 / 2a<1a1L 2b*<1a1 / 3a<2a1L,2b2 & I/IU/I & {-}/{-}{-}/{-} & N & W:1a,2a,3a \\
C4368 & 1a<o1 / 2a<1a1L 2b*<1a1 / 3a<2a2,2b1 & I/IU/I & {-}/{-}{-}/{-} & N & W:1a,2a,3a \\
C4369 & 1a<o1 / 2a<1a1L 2b*<1a1 / 3a<2a2,2b2 & I/IU/I & {-}/{-}{-}/{-} & N & W:1a,2a,3a \\
C4370 & 1a<o1 / 2a<1a1L 2b*<1a1 / 3a<2a2L,2b1 & I/IU/I & {-}/{-}{-}/{-} & N & W:1a,2a,3a \\
C4371 & 1a<o1 / 2a<1a1L 2b*<1a1 / 3a<2a2L,2b2 & I/IU/I & {-}/{-}{-}/{-} & N & W:1a,2a,3a \\
C4372 & 1a<o1 / 2a<1a1L 2b*<1a2 / 3a<2a0L,2b1 & I/IU/U & {-}/{-}{-}/{-} & N & W:1a,2a \\
C4373 & 1a<o1 / 2a<1a1L 2b*<1a2 / 3a<2a0L,2b2 & I/IU/U & {-}/{-}{-}/{-} & N & W:1a,2a \\
C4374 & 1a<o1 / 2a<1a1L 2b*<1a2 / 3a<2a1,2b1 & I/IU/I & {-}/{-}{-}/{-} & N & W:1a,2a,3a \\
C4375 & 1a<o1 / 2a<1a1L 2b*<1a2 / 3a<2a1,2b2 & I/IU/I & {-}/{-}{-}/{-} & N & W:1a,2a,3a \\
C4376 & 1a<o1 / 2a<1a1L 2b*<1a2 / 3a<2a1L,2b1 & I/IU/I & {-}/{-}{-}/{-} & N & W:1a,2a,3a \\
C4377 & 1a<o1 / 2a<1a1L 2b*<1a2 / 3a<2a1L,2b2 & I/IU/I & {-}/{-}{-}/{-} & N & W:1a,2a,3a \\
C4378 & 1a<o1 / 2a<1a1L 2b*<1a2 / 3a<2a2,2b1 & I/IU/I & {-}/{-}{-}/{-} & N & W:1a,2a,3a \\
C4379 & 1a<o1 / 2a<1a1L 2b*<1a2 / 3a<2a2,2b2 & I/IU/I & {-}/{-}{-}/{-} & N & W:1a,2a,3a \\
C4380 & 1a<o1 / 2a<1a1L 2b*<1a2 / 3a<2a2L,2b1 & I/IU/I & {-}/{-}{-}/{-} & N & W:1a,2a,3a \\
C4381 & 1a<o1 / 2a<1a1L 2b*<1a2 / 3a<2a2L,2b2 & I/IU/I & {-}/{-}{-}/{-} & N & W:1a,2a,3a \\
C4382 & 1a<o1 / 2a<1a2 2b<1a2 / 3a*<2a1,2b1 & I/EE/U & {-}/{-}{-}/{-} & E & W:1a \\
C4383 & 1a<o1 / 2a<1a2 2b<1a2 / 3a*<2a1,2b2 & I/EE/U & {-}/{-}{-}/{-} & E & W:1a \\
C4384 & 1a<o1 / 2a<1a2 2b<1a2 / 3a*<2a2,2b2 & I/EE/U & {-}/{-}{-}/{-} & E & W:1a \\
C4385 & 1a<o1 / 2a<1a2 2b<1a2L / 3a*<2a1,2b1 & I/LI/U & {-}/{-}{-}/{-} & B & CP:2b W:1a,2b \\
C4386 & 1a<o1 / 2a<1a2 2b<1a2L / 3a*<2a1,2b2 & I/LI/U & {-}/{-}{-}/{-} & B & CP:2b W:1a,2b \\
C4387 & 1a<o1 / 2a<1a2 2b<1a2L / 3a*<2a2,2b1 & I/LI/U & {-}/{-}{-}/{-} & B & CP:2b W:1a,2b \\
C4388 & 1a<o1 / 2a<1a2 2b<1a2L / 3a*<2a2,2b2 & I/LI/U & {-}/{-}{-}/{-} & B & CP:2b W:1a,2b \\
C4389 & 1a<o1 / 2a<1a2 2b*<1a1 / 3a<2a0L,2b1 & I/IU/U & {-}/{-}{-}/{-} & N & W:1a,2a \\
C4390 & 1a<o1 / 2a<1a2 2b*<1a1 / 3a<2a0L,2b2 & I/IU/U & {-}/{-}{-}/{-} & N & W:1a,2a \\
C4391 & 1a<o1 / 2a<1a2 2b*<1a1 / 3a<2a1,2b1 & I/IU/I & {-}/{-}{-}/{-} & N & W:1a,2a,3a \\
C4392 & 1a<o1 / 2a<1a2 2b*<1a1 / 3a<2a1,2b2 & I/IU/I & {-}/{-}{-}/{-} & N & W:1a,2a,3a \\
C4393 & 1a<o1 / 2a<1a2 2b*<1a1 / 3a<2a1L,2b1 & I/IU/I & {-}/{-}{-}/{-} & N & W:1a,2a,3a \\
C4394 & 1a<o1 / 2a<1a2 2b*<1a1 / 3a<2a1L,2b2 & I/IU/I & {-}/{-}{-}/{-} & N & W:1a,2a,3a \\
C4395 & 1a<o1 / 2a<1a2 2b*<1a1 / 3a<2a2,2b1 & I/IU/I & {-}/{-}{-}/{-} & N & W:1a,2a,3a \\
C4396 & 1a<o1 / 2a<1a2 2b*<1a1 / 3a<2a2,2b2 & I/IU/I & {-}/{-}{-}/{-} & N & W:1a,2a,3a \\
C4397 & 1a<o1 / 2a<1a2 2b*<1a1 / 3a<2a2L,2b1 & I/IU/I & {-}/{-}{-}/{-} & N & W:1a,2a,3a \\
C4398 & 1a<o1 / 2a<1a2 2b*<1a1 / 3a<2a2L,2b2 & I/IU/I & {-}/{-}{-}/{-} & N & W:1a,2a,3a \\
C4399 & 1a<o1 / 2a<1a2 2b*<1a2 / 3a<2a0L,2b1 & I/IU/U & {-}/{-}{-}/{-} & N & W:1a,2a \\
C4400 & 1a<o1 / 2a<1a2 2b*<1a2 / 3a<2a0L,2b2 & I/IU/U & {-}/{-}{-}/{-} & N & W:1a,2a \\
C4401 & 1a<o1 / 2a<1a2 2b*<1a2 / 3a<2a1,2b1 & I/IU/I & {-}/{-}{-}/{-} & N & W:1a,2a,3a \\
C4402 & 1a<o1 / 2a<1a2 2b*<1a2 / 3a<2a1,2b2 & I/IU/I & {-}/{-}{-}/{-} & N & W:1a,2a,3a \\
C4403 & 1a<o1 / 2a<1a2 2b*<1a2 / 3a<2a1L,2b1 & I/IU/I & {-}/{-}{-}/{-} & N & W:1a,2a,3a \\
C4404 & 1a<o1 / 2a<1a2 2b*<1a2 / 3a<2a1L,2b2 & I/IU/I & {-}/{-}{-}/{-} & N & W:1a,2a,3a \\
C4405 & 1a<o1 / 2a<1a2 2b*<1a2 / 3a<2a2,2b1 & I/IU/I & {-}/{-}{-}/{-} & N & W:1a,2a,3a \\
C4406 & 1a<o1 / 2a<1a2 2b*<1a2 / 3a<2a2,2b2 & I/IU/I & {-}/{-}{-}/{-} & N & W:1a,2a,3a \\
C4407 & 1a<o1 / 2a<1a2 2b*<1a2 / 3a<2a2L,2b1 & I/IU/I & {-}/{-}{-}/{-} & N & W:1a,2a,3a \\
C4408 & 1a<o1 / 2a<1a2 2b*<1a2 / 3a<2a2L,2b2 & I/IU/I & {-}/{-}{-}/{-} & N & W:1a,2a,3a \\
C4409 & 1a<o1 / 2a<1a2L 2b*<1a1 / 3a<2a0L,2b1 & I/IU/U & {-}/{-}{-}/{-} & N & W:1a,2a \\
C4410 & 1a<o1 / 2a<1a2L 2b*<1a1 / 3a<2a0L,2b2 & I/IU/U & {-}/{-}{-}/{-} & N & W:1a,2a \\
C4411 & 1a<o1 / 2a<1a2L 2b*<1a1 / 3a<2a1,2b1 & I/IU/I & {-}/{-}{-}/{-} & N & W:1a,2a,3a \\
C4412 & 1a<o1 / 2a<1a2L 2b*<1a1 / 3a<2a1,2b2 & I/IU/I & {-}/{-}{-}/{-} & N & W:1a,2a,3a \\
C4413 & 1a<o1 / 2a<1a2L 2b*<1a1 / 3a<2a1L,2b1 & I/IU/I & {-}/{-}{-}/{-} & N & W:1a,2a,3a \\
C4414 & 1a<o1 / 2a<1a2L 2b*<1a1 / 3a<2a1L,2b2 & I/IU/I & {-}/{-}{-}/{-} & N & W:1a,2a,3a \\
C4415 & 1a<o1 / 2a<1a2L 2b*<1a1 / 3a<2a2,2b1 & I/IU/I & {-}/{-}{-}/{-} & N & W:1a,2a,3a \\
C4416 & 1a<o1 / 2a<1a2L 2b*<1a1 / 3a<2a2,2b2 & I/IU/I & {-}/{-}{-}/{-} & N & W:1a,2a,3a \\
C4417 & 1a<o1 / 2a<1a2L 2b*<1a1 / 3a<2a2L,2b1 & I/IU/I & {-}/{-}{-}/{-} & N & W:1a,2a,3a \\
C4418 & 1a<o1 / 2a<1a2L 2b*<1a1 / 3a<2a2L,2b2 & I/IU/I & {-}/{-}{-}/{-} & N & W:1a,2a,3a \\
C4419 & 1a<o1 / 2a<1a2L 2b*<1a2 / 3a<2a0L,2b1 & I/IU/U & {-}/{-}{-}/{-} & N & W:1a,2a \\
C4420 & 1a<o1 / 2a<1a2L 2b*<1a2 / 3a<2a0L,2b2 & I/IU/U & {-}/{-}{-}/{-} & N & W:1a,2a \\
C4421 & 1a<o1 / 2a<1a2L 2b*<1a2 / 3a<2a1,2b1 & I/IU/I & {-}/{-}{-}/{-} & N & W:1a,2a,3a \\
C4422 & 1a<o1 / 2a<1a2L 2b*<1a2 / 3a<2a1,2b2 & I/IU/I & {-}/{-}{-}/{-} & N & W:1a,2a,3a \\
C4423 & 1a<o1 / 2a<1a2L 2b*<1a2 / 3a<2a1L,2b1 & I/IU/I & {-}/{-}{-}/{-} & N & W:1a,2a,3a \\
C4424 & 1a<o1 / 2a<1a2L 2b*<1a2 / 3a<2a1L,2b2 & I/IU/I & {-}/{-}{-}/{-} & N & W:1a,2a,3a \\
C4425 & 1a<o1 / 2a<1a2L 2b*<1a2 / 3a<2a2,2b1 & I/IU/I & {-}/{-}{-}/{-} & N & W:1a,2a,3a \\
C4426 & 1a<o1 / 2a<1a2L 2b*<1a2 / 3a<2a2,2b2 & I/IU/I & {-}/{-}{-}/{-} & N & W:1a,2a,3a \\
C4427 & 1a<o1 / 2a<1a2L 2b*<1a2 / 3a<2a2L,2b1 & I/IU/I & {-}/{-}{-}/{-} & N & W:1a,2a,3a \\
C4428 & 1a<o1 / 2a<1a2L 2b*<1a2 / 3a<2a2L,2b2 & I/IU/I & {-}/{-}{-}/{-} & N & W:1a,2a,3a \\
C4429 & 1a<o1 1b<o1 / 2a<1a0L,1b1 / 3a*<2a1 & EE/D/U & {-}{-}/{-}/{-} & E &  \\
C4430 & 1a<o1 1b<o1 / 2a<1a0L,1b1 / 3a*<2a2 & EE/D/U & {-}{-}/{-}/{-} & E &  \\
C4431 & 1a<o1 1b<o1 / 2a<1a0L,1b2 / 3a*<2a1 & EE/D/U & {-}{-}/{-}/{-} & E &  \\
C4432 & 1a<o1 1b<o1 / 2a<1a0L,1b2 / 3a*<2a2 & EE/D/U & {-}{-}/{-}/{-} & E &  \\
C4433 & 1a<o1 1b<o1 / 2a<1a1,1b1 / 3a*<2a1 & EE/D/U & {-}{-}/{-}/{-} & E &  \\
C4434 & 1a<o1 1b<o1 / 2a<1a1,1b1 / 3a*<2a2 & EE/D/U & {-}{-}/{-}/{-} & E &  \\
C4435 & 1a<o1 1b<o1 / 2a<1a1,1b1L / 3a*<2a1 & EE/D/U & {-}{-}/{-}/{-} & E &  \\
C4436 & 1a<o1 1b<o1 / 2a<1a1,1b1L / 3a*<2a2 & EE/D/U & {-}{-}/{-}/{-} & E &  \\
C4437 & 1a<o1 1b<o1 / 2a<1a1,1b2 / 3a*<2a1 & EE/D/U & {-}{-}/{-}/{-} & E &  \\
C4438 & 1a<o1 1b<o1 / 2a<1a1,1b2 / 3a*<2a2 & EE/D/U & {-}{-}/{-}/{-} & E &  \\
C4439 & 1a<o1 1b<o1 / 2a<1a1,1b2L / 3a*<2a1 & EE/D/U & {-}{-}/{-}/{-} & E &  \\
C4440 & 1a<o1 1b<o1 / 2a<1a1,1b2L / 3a*<2a2 & EE/D/U & {-}{-}/{-}/{-} & E &  \\
C4441 & 1a<o1 1b<o1 / 2a<1a1L,1b2 / 3a*<2a1 & EE/D/U & {-}{-}/{-}/{-} & E &  \\
C4442 & 1a<o1 1b<o1 / 2a<1a1L,1b2 / 3a*<2a2 & EE/D/U & {-}{-}/{-}/{-} & E &  \\
C4443 & 1a<o1 1b<o1 / 2a<1a2,1b2 / 3a*<2a1 & EE/D/U & {-}{-}/{-}/{-} & E &  \\
C4444 & 1a<o1 1b<o1 / 2a<1a2,1b2 / 3a*<2a2 & EE/D/U & {-}{-}/{-}/{-} & E &  \\
C4445 & 1a<o1 1b<o1 / 2a<1a2,1b2L / 3a*<2a1 & EE/D/U & {-}{-}/{-}/{-} & E &  \\
C4446 & 1a<o1 1b<o1 / 2a<1a2,1b2L / 3a*<2a2 & EE/D/U & {-}{-}/{-}/{-} & E &  \\
C4447 & 1a<o1 1b<o1 / 2a*<1a1,1b1 / 3a<2a1 & EE/U/U & {-}{-}/{-}/{-} & E &  \\
C4448 & 1a<o1 1b<o1 / 2a*<1a1,1b1 / 3a<2a2 & EE/U/U & {-}{-}/{-}/{-} & E &  \\
C4449 & 1a<o1 1b<o1 / 2a*<1a1,1b2 / 3a<2a1 & EE/U/U & {-}{-}/{-}/{-} & E &  \\
C4450 & 1a<o1 1b<o1 / 2a*<1a1,1b2 / 3a<2a2 & EE/U/U & {-}{-}/{-}/{-} & E &  \\
C4451 & 1a<o1 1b<o1 / 2a*<1a2,1b2 / 3a<2a1 & EE/U/U & {-}{-}/{-}/{-} & E &  \\
C4452 & 1a<o1 1b<o1 / 2a*<1a2,1b2 / 3a<2a2 & EE/U/U & {-}{-}/{-}/{-} & E &  \\
C4453 & 1a<o1 1b<o1L / 2a<1a0L,1b1 / 3a*<2a1 & LI/I/U & {-}+/{-}/{-} & B & CPW:1b,2a \\
C4454 & 1a<o1 1b<o1L / 2a<1a0L,1b1 / 3a*<2a2 & LI/I/U & {-}+/{-}/{-} & B & CPW:1b,2a \\
C4455 & 1a<o1 1b<o1L / 2a<1a0L,1b2 / 3a*<2a1 & LI/I/U & {-}+/{-}/{-} & B & CPW:1b,2a \\
C4456 & 1a<o1 1b<o1L / 2a<1a0L,1b2 / 3a*<2a2 & LI/I/U & {-}+/{-}/{-} & B & CPW:1b,2a \\
C4457 & 1a<o1 1b<o1L / 2a<1a1,1b0L / 3a*<2a1 & LI/D/U & {-}+/+/{-} & B & CPW:1b \\
C4458 & 1a<o1 1b<o1L / 2a<1a1,1b0L / 3a*<2a2 & LI/D/U & {-}+/+/{-} & B & CPW:1b \\
C4459 & 1a<o1 1b<o1L / 2a<1a1,1b1 / 3a*<2a1 & LI/E/U & {-}+/{-}/{-} & BF & CPW:1b \\
C4460 & 1a<o1 1b<o1L / 2a<1a1,1b1 / 3a*<2a2 & LI/E/U & {-}+/{-}/{-} & BF & CPW:1b \\
C4461 & 1a<o1 1b<o1L / 2a<1a1,1b1L / 3a*<2a1 & LI/I/U & {-}+/+/{-} & BF & CPW:1b,2a \\
C4462 & 1a<o1 1b<o1L / 2a<1a1,1b1L / 3a*<2a2 & LI/I/U & {-}+/+/{-} & BF & CPW:1b,2a \\
C4463 & 1a<o1 1b<o1L / 2a<1a1,1b2 / 3a*<2a1 & LI/E/U & {-}+/{-}/{-} & BF & CPW:1b \\
C4464 & 1a<o1 1b<o1L / 2a<1a1,1b2 / 3a*<2a2 & LI/E/U & {-}+/{-}/{-} & BF & CPW:1b \\
C4465 & 1a<o1 1b<o1L / 2a<1a1,1b2L / 3a*<2a1 & LI/I/U & {-}+/+/{-} & BF & CPW:1b,2a \\
C4466 & 1a<o1 1b<o1L / 2a<1a1,1b2L / 3a*<2a2 & LI/I/U & {-}+/+/{-} & BF & CPW:1b,2a \\
C4467 & 1a<o1 1b<o1L / 2a<1a1L,1b1 / 3a*<2a1 & LI/E/U & {-}+/{-}/{-} & BF & CPW:1b \\
C4468 & 1a<o1 1b<o1L / 2a<1a1L,1b1 / 3a*<2a2 & LI/E/U & {-}+/{-}/{-} & BF & CPW:1b \\
C4469 & 1a<o1 1b<o1L / 2a<1a1L,1b2 / 3a*<2a1 & LI/E/U & {-}+/{-}/{-} & BF & CPW:1b \\
C4470 & 1a<o1 1b<o1L / 2a<1a1L,1b2 / 3a*<2a2 & LI/E/U & {-}+/{-}/{-} & BF & CPW:1b \\
C4471 & 1a<o1 1b<o1L / 2a<1a2,1b0L / 3a*<2a1 & LI/D/U & {-}+/+/{-} & B & CPW:1b \\
C4472 & 1a<o1 1b<o1L / 2a<1a2,1b0L / 3a*<2a2 & LI/D/U & {-}+/+/{-} & B & CPW:1b \\
C4473 & 1a<o1 1b<o1L / 2a<1a2,1b1 / 3a*<2a1 & LI/E/U & {-}+/{-}/{-} & BF & CPW:1b \\
C4474 & 1a<o1 1b<o1L / 2a<1a2,1b1 / 3a*<2a2 & LI/E/U & {-}+/{-}/{-} & BF & CPW:1b \\
C4475 & 1a<o1 1b<o1L / 2a<1a2,1b1L / 3a*<2a1 & LI/I/U & {-}+/+/{-} & BF & CPW:1b,2a \\
C4476 & 1a<o1 1b<o1L / 2a<1a2,1b1L / 3a*<2a2 & LI/I/U & {-}+/+/{-} & BF & CPW:1b,2a \\
C4477 & 1a<o1 1b<o1L / 2a<1a2,1b2 / 3a*<2a1 & LI/E/U & {-}+/{-}/{-} & BF & CPW:1b \\
C4478 & 1a<o1 1b<o1L / 2a<1a2,1b2 / 3a*<2a2 & LI/E/U & {-}+/{-}/{-} & BF & CPW:1b \\
C4479 & 1a<o1 1b<o1L / 2a<1a2,1b2L / 3a*<2a1 & LI/I/U & {-}+/+/{-} & BF & CPW:1b,2a \\
C4480 & 1a<o1 1b<o1L / 2a<1a2,1b2L / 3a*<2a2 & LI/I/U & {-}+/+/{-} & BF & CPW:1b,2a \\
C4481 & 1a<o1 1b<o1L / 2a<1a2L,1b1 / 3a*<2a1 & LI/E/U & {-}+/{-}/{-} & BF & CPW:1b \\
C4482 & 1a<o1 1b<o1L / 2a<1a2L,1b1 / 3a*<2a2 & LI/E/U & {-}+/{-}/{-} & BF & CPW:1b \\
C4483 & 1a<o1 1b<o1L / 2a<1a2L,1b2 / 3a*<2a1 & LI/E/U & {-}+/{-}/{-} & BF & CPW:1b \\
C4484 & 1a<o1 1b<o1L / 2a<1a2L,1b2 / 3a*<2a2 & LI/E/U & {-}+/{-}/{-} & BF & CPW:1b \\
C4485 & 1a<o1 1b<o1L / 2a*<1a1,1b1 / 3a<2a1 & LI/U/U & {-}+/{-}/{-} & B & CPW:1b \\
C4486 & 1a<o1 1b<o1L / 2a*<1a1,1b1 / 3a<2a2 & LI/U/U & {-}+/{-}/{-} & B & CPW:1b \\
C4487 & 1a<o1 1b<o1L / 2a*<1a1,1b2 / 3a<2a1 & LI/U/U & {-}+/{-}/{-} & B & CPW:1b \\
C4488 & 1a<o1 1b<o1L / 2a*<1a1,1b2 / 3a<2a2 & LI/U/U & {-}+/{-}/{-} & B & CPW:1b \\
C4489 & 1a<o1 1b<o1L / 2a*<1a2,1b1 / 3a<2a1 & LI/U/U & {-}+/{-}/{-} & B & CPW:1b \\
C4490 & 1a<o1 1b<o1L / 2a*<1a2,1b1 / 3a<2a2 & LI/U/U & {-}+/{-}/{-} & B & CPW:1b \\
C4491 & 1a<o1 1b<o1L / 2a*<1a2,1b2 / 3a<2a1 & LI/U/U & {-}+/{-}/{-} & B & CPW:1b \\
C4492 & 1a<o1 1b<o1L / 2a*<1a2,1b2 / 3a<2a2 & LI/U/U & {-}+/{-}/{-} & B & CPW:1b \\
C4493 & 1a<o1 1b<o2 / 2a<1a0L,1b1 / 3a*<2a1 & EE/D/U & {-}{-}/{-}/{-} & E & C:1b,2a \\
C4494 & 1a<o1 1b<o2 / 2a<1a0L,1b1 / 3a*<2a2 & EE/D/U & {-}{-}/{-}/{-} & E & C:1b,2a \\
C4495 & 1a<o1 1b<o2 / 2a<1a0L,1b2 / 3a*<2a1 & EE/D/U & {-}{-}/{-}/{-} & E & C:1b,2a \\
C4496 & 1a<o1 1b<o2 / 2a<1a0L,1b2 / 3a*<2a2 & EE/D/U & {-}{-}/{-}/{-} & E & C:1b,2a \\
C4497 & 1a<o1 1b<o2 / 2a<1a1,1b0L / 3a*<2a1 & EE/D/U & {-}{-}/{-}/{-} & E & C:1b \\
C4498 & 1a<o1 1b<o2 / 2a<1a1,1b0L / 3a*<2a2 & EE/D/U & {-}{-}/{-}/{-} & E & C:1b \\
C4499 & 1a<o1 1b<o2 / 2a<1a1,1b1 / 3a*<2a1 & EE/D/U & {-}{-}/{-}/{-} & E & C:1b \\
C4500 & 1a<o1 1b<o2 / 2a<1a1,1b1 / 3a*<2a2 & EE/D/U & {-}{-}/{-}/{-} & E & C:1b \\
C4501 & 1a<o1 1b<o2 / 2a<1a1,1b1L / 3a*<2a1 & EE/D/U & {-}{-}/{-}/{-} & E & C:1b \\
C4502 & 1a<o1 1b<o2 / 2a<1a1,1b1L / 3a*<2a2 & EE/D/U & {-}{-}/{-}/{-} & E & C:1b \\
C4503 & 1a<o1 1b<o2 / 2a<1a1,1b2 / 3a*<2a1 & EE/D/U & {-}{-}/{-}/{-} & E & C:1b,2a \\
C4504 & 1a<o1 1b<o2 / 2a<1a1,1b2 / 3a*<2a2 & EE/D/U & {-}{-}/{-}/{-} & E & C:1b,2a \\
C4505 & 1a<o1 1b<o2 / 2a<1a1,1b2L / 3a*<2a1 & EE/D/U & {-}{-}/{-}/{-} & E & C:1b,2a \\
C4506 & 1a<o1 1b<o2 / 2a<1a1,1b2L / 3a*<2a2 & EE/D/U & {-}{-}/{-}/{-} & E & C:1b,2a \\
C4507 & 1a<o1 1b<o2 / 2a<1a1L,1b1 / 3a*<2a1 & EE/D/U & {-}{-}/{-}/{-} & E & C:1b \\
C4508 & 1a<o1 1b<o2 / 2a<1a1L,1b1 / 3a*<2a2 & EE/D/U & {-}{-}/{-}/{-} & E & C:1b \\
C4509 & 1a<o1 1b<o2 / 2a<1a1L,1b2 / 3a*<2a1 & EE/D/U & {-}{-}/{-}/{-} & E & C:1b,2a \\
C4510 & 1a<o1 1b<o2 / 2a<1a1L,1b2 / 3a*<2a2 & EE/D/U & {-}{-}/{-}/{-} & E & C:1b,2a \\
C4511 & 1a<o1 1b<o2 / 2a<1a2,1b0L / 3a*<2a1 & EE/D/U & {-}{-}/{-}/{-} & E & C:1b \\
C4512 & 1a<o1 1b<o2 / 2a<1a2,1b0L / 3a*<2a2 & EE/D/U & {-}{-}/{-}/{-} & E & C:1b \\
C4513 & 1a<o1 1b<o2 / 2a<1a2,1b1 / 3a*<2a1 & EE/D/U & {-}{-}/{-}/{-} & E & C:1b \\
C4514 & 1a<o1 1b<o2 / 2a<1a2,1b1 / 3a*<2a2 & EE/D/U & {-}{-}/{-}/{-} & E & C:1b \\
C4515 & 1a<o1 1b<o2 / 2a<1a2,1b1L / 3a*<2a1 & EE/D/U & {-}{-}/{-}/{-} & E & C:1b \\
C4516 & 1a<o1 1b<o2 / 2a<1a2,1b1L / 3a*<2a2 & EE/D/U & {-}{-}/{-}/{-} & E & C:1b \\
C4517 & 1a<o1 1b<o2 / 2a<1a2,1b2 / 3a*<2a1 & EE/D/U & {-}{-}/{-}/{-} & E & C:1b \\
C4518 & 1a<o1 1b<o2 / 2a<1a2,1b2 / 3a*<2a2 & EE/D/U & {-}{-}/{-}/{-} & E & C:1b \\
C4519 & 1a<o1 1b<o2 / 2a<1a2,1b2L / 3a*<2a1 & EE/D/U & {-}{-}/{-}/{-} & E & C:1b \\
C4520 & 1a<o1 1b<o2 / 2a<1a2,1b2L / 3a*<2a2 & EE/D/U & {-}{-}/{-}/{-} & E & C:1b \\
C4521 & 1a<o1 1b<o2 / 2a<1a2L,1b1 / 3a*<2a1 & EE/D/U & {-}{-}/{-}/{-} & E & C:1b \\
C4522 & 1a<o1 1b<o2 / 2a<1a2L,1b1 / 3a*<2a2 & EE/D/U & {-}{-}/{-}/{-} & E & C:1b \\
C4523 & 1a<o1 1b<o2 / 2a<1a2L,1b2 / 3a*<2a1 & EE/D/U & {-}{-}/{-}/{-} & E & C:1b \\
C4524 & 1a<o1 1b<o2 / 2a<1a2L,1b2 / 3a*<2a2 & EE/D/U & {-}{-}/{-}/{-} & E & C:1b \\
C4525 & 1a<o1 1b<o2 / 2a*<1a1,1b1 / 3a<2a1 & EE/U/U & {-}{-}/{-}/{-} & E & C:1b \\
C4526 & 1a<o1 1b<o2 / 2a*<1a1,1b1 / 3a<2a2 & EE/U/U & {-}{-}/{-}/{-} & E & C:1b \\
C4527 & 1a<o1 1b<o2 / 2a*<1a1,1b2 / 3a<2a1 & EE/U/U & {-}{-}/{-}/{-} & E & C:1b \\
C4528 & 1a<o1 1b<o2 / 2a*<1a1,1b2 / 3a<2a2 & EE/U/U & {-}{-}/{-}/{-} & E & C:1b \\
C4529 & 1a<o1 1b<o2 / 2a*<1a2,1b1 / 3a<2a1 & EE/U/U & {-}{-}/{-}/{-} & E & C:1b \\
C4530 & 1a<o1 1b<o2 / 2a*<1a2,1b1 / 3a<2a2 & EE/U/U & {-}{-}/{-}/{-} & E & C:1b \\
C4531 & 1a<o1 1b<o2 / 2a*<1a2,1b2 / 3a<2a1 & EE/U/U & {-}{-}/{-}/{-} & E & C:1b \\
C4532 & 1a<o1 1b<o2 / 2a*<1a2,1b2 / 3a<2a2 & EE/U/U & {-}{-}/{-}/{-} & E & C:1b \\
C4533 & 1a<o1 1b<o2L / 2a<1a0L,1b1 / 3a*<2a1 & LI/I/U & {-}+/{-}/{-} & B & PW:1b,2a \\
C4534 & 1a<o1 1b<o2L / 2a<1a0L,1b1 / 3a*<2a2 & LI/I/U & {-}+/{-}/{-} & B & PW:1b,2a \\
C4535 & 1a<o1 1b<o2L / 2a<1a0L,1b2 / 3a*<2a1 & LI/I/U & {-}+/{-}/{-} & B & PW:1b,2a \\
C4536 & 1a<o1 1b<o2L / 2a<1a0L,1b2 / 3a*<2a2 & LI/I/U & {-}+/{-}/{-} & B & PW:1b,2a \\
C4537 & 1a<o1 1b<o2L / 2a<1a1,1b0L / 3a*<2a1 & LI/D/U & {-}+/+/{-} & B & PW:1b \\
C4538 & 1a<o1 1b<o2L / 2a<1a1,1b0L / 3a*<2a2 & LI/D/U & {-}+/+/{-} & B & PW:1b \\
C4539 & 1a<o1 1b<o2L / 2a<1a1,1b1 / 3a*<2a1 & LI/E/U & {-}+/{-}/{-} & BF & PW:1b \\
C4540 & 1a<o1 1b<o2L / 2a<1a1,1b1 / 3a*<2a2 & LI/E/U & {-}+/{-}/{-} & BF & PW:1b \\
C4541 & 1a<o1 1b<o2L / 2a<1a1,1b1L / 3a*<2a1 & LI/I/U & {-}+/+/{-} & BF & C:2a PW:1b,2a \\
C4542 & 1a<o1 1b<o2L / 2a<1a1,1b1L / 3a*<2a2 & LI/I/U & {-}+/+/{-} & BF & C:2a PW:1b,2a \\
C4543 & 1a<o1 1b<o2L / 2a<1a1,1b2 / 3a*<2a1 & LI/E/U & {-}+/{-}/{-} & BF & C:2a PW:1b \\
C4544 & 1a<o1 1b<o2L / 2a<1a1,1b2 / 3a*<2a2 & LI/E/U & {-}+/{-}/{-} & BF & C:2a PW:1b \\
C4545 & 1a<o1 1b<o2L / 2a<1a1,1b2L / 3a*<2a1 & LI/I/U & {-}+/+/{-} & BF & PW:1b,2a \\
C4546 & 1a<o1 1b<o2L / 2a<1a1,1b2L / 3a*<2a2 & LI/I/U & {-}+/+/{-} & BF & PW:1b,2a \\
C4547 & 1a<o1 1b<o2L / 2a<1a1L,1b1 / 3a*<2a1 & LI/E/U & {-}+/{-}/{-} & BF & PW:1b \\
C4548 & 1a<o1 1b<o2L / 2a<1a1L,1b1 / 3a*<2a2 & LI/E/U & {-}+/{-}/{-} & BF & PW:1b \\
C4549 & 1a<o1 1b<o2L / 2a<1a1L,1b2 / 3a*<2a1 & LI/E/U & {-}+/{-}/{-} & BF & C:2a PW:1b \\
C4550 & 1a<o1 1b<o2L / 2a<1a1L,1b2 / 3a*<2a2 & LI/E/U & {-}+/{-}/{-} & BF & C:2a PW:1b \\
C4551 & 1a<o1 1b<o2L / 2a<1a2,1b0L / 3a*<2a1 & LI/D/U & {-}+/+/{-} & B & PW:1b \\
C4552 & 1a<o1 1b<o2L / 2a<1a2,1b0L / 3a*<2a2 & LI/D/U & {-}+/+/{-} & B & PW:1b \\
C4553 & 1a<o1 1b<o2L / 2a<1a2,1b1 / 3a*<2a1 & LI/E/U & {-}+/{-}/{-} & BF & PW:1b \\
C4554 & 1a<o1 1b<o2L / 2a<1a2,1b1 / 3a*<2a2 & LI/E/U & {-}+/{-}/{-} & BF & PW:1b \\
C4555 & 1a<o1 1b<o2L / 2a<1a2,1b1L / 3a*<2a1 & LI/I/U & {-}+/+/{-} & BF & C:2a PW:1b,2a \\
C4556 & 1a<o1 1b<o2L / 2a<1a2,1b1L / 3a*<2a2 & LI/I/U & {-}+/+/{-} & BF & C:2a PW:1b,2a \\
C4557 & 1a<o1 1b<o2L / 2a<1a2,1b2 / 3a*<2a1 & LI/E/U & {-}+/{-}/{-} & BF & PW:1b \\
C4558 & 1a<o1 1b<o2L / 2a<1a2,1b2 / 3a*<2a2 & LI/E/U & {-}+/{-}/{-} & BF & PW:1b \\
C4559 & 1a<o1 1b<o2L / 2a<1a2,1b2L / 3a*<2a1 & LI/I/U & {-}+/+/{-} & BF & C:2a PW:1b,2a \\
C4560 & 1a<o1 1b<o2L / 2a<1a2,1b2L / 3a*<2a2 & LI/I/U & {-}+/+/{-} & BF & C:2a PW:1b,2a \\
C4561 & 1a<o1 1b<o2L / 2a<1a2L,1b1 / 3a*<2a1 & LI/E/U & {-}+/{-}/{-} & BF & PW:1b \\
C4562 & 1a<o1 1b<o2L / 2a<1a2L,1b1 / 3a*<2a2 & LI/E/U & {-}+/{-}/{-} & BF & PW:1b \\
C4563 & 1a<o1 1b<o2L / 2a<1a2L,1b2 / 3a*<2a1 & LI/E/U & {-}+/{-}/{-} & BF & PW:1b \\
C4564 & 1a<o1 1b<o2L / 2a<1a2L,1b2 / 3a*<2a2 & LI/E/U & {-}+/{-}/{-} & BF & PW:1b \\
C4565 & 1a<o1 1b<o2L / 2a*<1a1,1b1 / 3a<2a1 & LI/U/U & {-}+/{-}/{-} & B & PW:1b \\
C4566 & 1a<o1 1b<o2L / 2a*<1a1,1b1 / 3a<2a2 & LI/U/U & {-}+/{-}/{-} & B & PW:1b \\
C4567 & 1a<o1 1b<o2L / 2a*<1a1,1b2 / 3a<2a1 & LI/U/U & {-}+/{-}/{-} & B & PW:1b \\
C4568 & 1a<o1 1b<o2L / 2a*<1a1,1b2 / 3a<2a2 & LI/U/U & {-}+/{-}/{-} & B & PW:1b \\
C4569 & 1a<o1 1b<o2L / 2a*<1a2,1b1 / 3a<2a1 & LI/U/U & {-}+/{-}/{-} & B & PW:1b \\
C4570 & 1a<o1 1b<o2L / 2a*<1a2,1b1 / 3a<2a2 & LI/U/U & {-}+/{-}/{-} & B & PW:1b \\
C4571 & 1a<o1 1b<o2L / 2a*<1a2,1b2 / 3a<2a1 & LI/U/U & {-}+/{-}/{-} & B & PW:1b \\
C4572 & 1a<o1 1b<o2L / 2a*<1a2,1b2 / 3a<2a2 & LI/U/U & {-}+/{-}/{-} & B & PW:1b \\
C4573 & 1a<o1 1b*<o1 / 2a<1a0L,1b1 / 3a<2a0L & IU/U/N & {-}{-}/{-}/{-} & N & W:1a \\
C4574 & 1a<o1 1b*<o1 / 2a<1a0L,1b1 / 3a<2a1 & IU/U/U & {-}{-}/{-}/{-} & N & W:1a \\
C4575 & 1a<o1 1b*<o1 / 2a<1a0L,1b1 / 3a<2a1L & IU/U/U & {-}{-}/{-}/{-} & N & W:1a \\
C4576 & 1a<o1 1b*<o1 / 2a<1a0L,1b1 / 3a<2a2 & IU/U/U & {-}{-}/{-}/{-} & N & W:1a \\
C4577 & 1a<o1 1b*<o1 / 2a<1a0L,1b1 / 3a<2a2L & IU/U/U & {-}{-}/{-}/{-} & N & W:1a \\
C4578 & 1a<o1 1b*<o1 / 2a<1a0L,1b2 / 3a<2a0L & IU/U/N & {-}{-}/{-}/{-} & N & W:1a \\
C4579 & 1a<o1 1b*<o1 / 2a<1a0L,1b2 / 3a<2a1 & IU/U/U & {-}{-}/{-}/{-} & N & W:1a \\
C4580 & 1a<o1 1b*<o1 / 2a<1a0L,1b2 / 3a<2a1L & IU/U/U & {-}{-}/{-}/{-} & N & W:1a \\
C4581 & 1a<o1 1b*<o1 / 2a<1a0L,1b2 / 3a<2a2 & IU/U/U & {-}{-}/{-}/{-} & N & W:1a \\
C4582 & 1a<o1 1b*<o1 / 2a<1a0L,1b2 / 3a<2a2L & IU/U/U & {-}{-}/{-}/{-} & N & W:1a \\
C4583 & 1a<o1 1b*<o1 / 2a<1a1,1b1 / 3a<2a0L & IU/I/N & {-}{-}/{-}/{-} & N & W:1a,2a \\
C4584 & 1a<o1 1b*<o1 / 2a<1a1,1b1 / 3a<2a1 & IU/I/I & {-}{-}/{-}/{-} & N & W:1a,2a,3a \\
C4585 & 1a<o1 1b*<o1 / 2a<1a1,1b1 / 3a<2a1L & IU/I/I & {-}{-}/{-}/{-} & N & W:1a,2a,3a \\
C4586 & 1a<o1 1b*<o1 / 2a<1a1,1b1 / 3a<2a2 & IU/I/I & {-}{-}/{-}/{-} & N & W:1a,2a,3a \\
C4587 & 1a<o1 1b*<o1 / 2a<1a1,1b1 / 3a<2a2L & IU/I/I & {-}{-}/{-}/{-} & N & W:1a,2a,3a \\
C4588 & 1a<o1 1b*<o1 / 2a<1a1,1b2 / 3a<2a0L & IU/I/N & {-}{-}/{-}/{-} & N & W:1a,2a \\
C4589 & 1a<o1 1b*<o1 / 2a<1a1,1b2 / 3a<2a1 & IU/I/I & {-}{-}/{-}/{-} & N & W:1a,2a,3a \\
C4590 & 1a<o1 1b*<o1 / 2a<1a1,1b2 / 3a<2a1L & IU/I/I & {-}{-}/{-}/{-} & N & W:1a,2a,3a \\
C4591 & 1a<o1 1b*<o1 / 2a<1a1,1b2 / 3a<2a2 & IU/I/I & {-}{-}/{-}/{-} & N & W:1a,2a,3a \\
C4592 & 1a<o1 1b*<o1 / 2a<1a1,1b2 / 3a<2a2L & IU/I/I & {-}{-}/{-}/{-} & N & W:1a,2a,3a \\
C4593 & 1a<o1 1b*<o1 / 2a<1a1L,1b1 / 3a<2a0L & IU/I/N & {-}{-}/{-}/{-} & N & W:1a,2a \\
C4594 & 1a<o1 1b*<o1 / 2a<1a1L,1b1 / 3a<2a1 & IU/I/I & {-}{-}/{-}/{-} & N & W:1a,2a,3a \\
C4595 & 1a<o1 1b*<o1 / 2a<1a1L,1b1 / 3a<2a1L & IU/I/I & {-}{-}/{-}/{-} & N & W:1a,2a,3a \\
C4596 & 1a<o1 1b*<o1 / 2a<1a1L,1b1 / 3a<2a2 & IU/I/I & {-}{-}/{-}/{-} & N & W:1a,2a,3a \\
C4597 & 1a<o1 1b*<o1 / 2a<1a1L,1b1 / 3a<2a2L & IU/I/I & {-}{-}/{-}/{-} & N & W:1a,2a,3a \\
C4598 & 1a<o1 1b*<o1 / 2a<1a1L,1b2 / 3a<2a0L & IU/I/N & {-}{-}/{-}/{-} & N & W:1a,2a \\
C4599 & 1a<o1 1b*<o1 / 2a<1a1L,1b2 / 3a<2a1 & IU/I/I & {-}{-}/{-}/{-} & N & W:1a,2a,3a \\
C4600 & 1a<o1 1b*<o1 / 2a<1a1L,1b2 / 3a<2a1L & IU/I/I & {-}{-}/{-}/{-} & N & W:1a,2a,3a \\
C4601 & 1a<o1 1b*<o1 / 2a<1a1L,1b2 / 3a<2a2 & IU/I/I & {-}{-}/{-}/{-} & N & W:1a,2a,3a \\
C4602 & 1a<o1 1b*<o1 / 2a<1a1L,1b2 / 3a<2a2L & IU/I/I & {-}{-}/{-}/{-} & N & W:1a,2a,3a \\
C4603 & 1a<o1 1b*<o1 / 2a<1a2,1b1 / 3a<2a0L & IU/I/N & {-}{-}/{-}/{-} & N & W:1a,2a \\
C4604 & 1a<o1 1b*<o1 / 2a<1a2,1b1 / 3a<2a1 & IU/I/I & {-}{-}/{-}/{-} & N & W:1a,2a,3a \\
C4605 & 1a<o1 1b*<o1 / 2a<1a2,1b1 / 3a<2a1L & IU/I/I & {-}{-}/{-}/{-} & N & W:1a,2a,3a \\
C4606 & 1a<o1 1b*<o1 / 2a<1a2,1b1 / 3a<2a2 & IU/I/I & {-}{-}/{-}/{-} & N & W:1a,2a,3a \\
C4607 & 1a<o1 1b*<o1 / 2a<1a2,1b1 / 3a<2a2L & IU/I/I & {-}{-}/{-}/{-} & N & W:1a,2a,3a \\
C4608 & 1a<o1 1b*<o1 / 2a<1a2,1b2 / 3a<2a0L & IU/I/N & {-}{-}/{-}/{-} & N & W:1a,2a \\
C4609 & 1a<o1 1b*<o1 / 2a<1a2,1b2 / 3a<2a1 & IU/I/I & {-}{-}/{-}/{-} & N & W:1a,2a,3a \\
C4610 & 1a<o1 1b*<o1 / 2a<1a2,1b2 / 3a<2a1L & IU/I/I & {-}{-}/{-}/{-} & N & W:1a,2a,3a \\
C4611 & 1a<o1 1b*<o1 / 2a<1a2,1b2 / 3a<2a2 & IU/I/I & {-}{-}/{-}/{-} & N & W:1a,2a,3a \\
C4612 & 1a<o1 1b*<o1 / 2a<1a2,1b2 / 3a<2a2L & IU/I/I & {-}{-}/{-}/{-} & N & W:1a,2a,3a \\
C4613 & 1a<o1 1b*<o1 / 2a<1a2L,1b1 / 3a<2a0L & IU/I/N & {-}{-}/{-}/{-} & N & W:1a,2a \\
C4614 & 1a<o1 1b*<o1 / 2a<1a2L,1b1 / 3a<2a1 & IU/I/I & {-}{-}/{-}/{-} & N & W:1a,2a,3a \\
C4615 & 1a<o1 1b*<o1 / 2a<1a2L,1b1 / 3a<2a1L & IU/I/I & {-}{-}/{-}/{-} & N & W:1a,2a,3a \\
C4616 & 1a<o1 1b*<o1 / 2a<1a2L,1b1 / 3a<2a2 & IU/I/I & {-}{-}/{-}/{-} & N & W:1a,2a,3a \\
C4617 & 1a<o1 1b*<o1 / 2a<1a2L,1b1 / 3a<2a2L & IU/I/I & {-}{-}/{-}/{-} & N & W:1a,2a,3a \\
C4618 & 1a<o1 1b*<o1 / 2a<1a2L,1b2 / 3a<2a0L & IU/I/N & {-}{-}/{-}/{-} & N & W:1a,2a \\
C4619 & 1a<o1 1b*<o1 / 2a<1a2L,1b2 / 3a<2a1 & IU/I/I & {-}{-}/{-}/{-} & N & W:1a,2a,3a \\
C4620 & 1a<o1 1b*<o1 / 2a<1a2L,1b2 / 3a<2a1L & IU/I/I & {-}{-}/{-}/{-} & N & W:1a,2a,3a \\
C4621 & 1a<o1 1b*<o1 / 2a<1a2L,1b2 / 3a<2a2 & IU/I/I & {-}{-}/{-}/{-} & N & W:1a,2a,3a \\
C4622 & 1a<o1 1b*<o1 / 2a<1a2L,1b2 / 3a<2a2L & IU/I/I & {-}{-}/{-}/{-} & N & W:1a,2a,3a \\
C4623 & 1a<o1 1b*<o2 / 2a<1a0L,1b1 / 3a<2a0L & IU/U/N & {-}{-}/{-}/{-} & N & W:1a \\
C4624 & 1a<o1 1b*<o2 / 2a<1a0L,1b1 / 3a<2a1 & IU/U/U & {-}{-}/{-}/{-} & N & W:1a \\
C4625 & 1a<o1 1b*<o2 / 2a<1a0L,1b1 / 3a<2a1L & IU/U/U & {-}{-}/{-}/{-} & N & W:1a \\
C4626 & 1a<o1 1b*<o2 / 2a<1a0L,1b1 / 3a<2a2 & IU/U/U & {-}{-}/{-}/{-} & N & W:1a \\
C4627 & 1a<o1 1b*<o2 / 2a<1a0L,1b1 / 3a<2a2L & IU/U/U & {-}{-}/{-}/{-} & N & W:1a \\
C4628 & 1a<o1 1b*<o2 / 2a<1a0L,1b2 / 3a<2a0L & IU/U/N & {-}{-}/{-}/{-} & N & W:1a \\
C4629 & 1a<o1 1b*<o2 / 2a<1a0L,1b2 / 3a<2a1 & IU/U/U & {-}{-}/{-}/{-} & N & W:1a \\
C4630 & 1a<o1 1b*<o2 / 2a<1a0L,1b2 / 3a<2a1L & IU/U/U & {-}{-}/{-}/{-} & N & W:1a \\
C4631 & 1a<o1 1b*<o2 / 2a<1a0L,1b2 / 3a<2a2 & IU/U/U & {-}{-}/{-}/{-} & N & W:1a \\
C4632 & 1a<o1 1b*<o2 / 2a<1a0L,1b2 / 3a<2a2L & IU/U/U & {-}{-}/{-}/{-} & N & W:1a \\
C4633 & 1a<o1 1b*<o2 / 2a<1a1,1b1 / 3a<2a0L & IU/I/N & {-}{-}/{-}/{-} & N & W:1a,2a \\
C4634 & 1a<o1 1b*<o2 / 2a<1a1,1b1 / 3a<2a1 & IU/I/I & {-}{-}/{-}/{-} & N & W:1a,2a,3a \\
C4635 & 1a<o1 1b*<o2 / 2a<1a1,1b1 / 3a<2a1L & IU/I/I & {-}{-}/{-}/{-} & N & W:1a,2a,3a \\
C4636 & 1a<o1 1b*<o2 / 2a<1a1,1b1 / 3a<2a2 & IU/I/I & {-}{-}/{-}/{-} & N & W:1a,2a,3a \\
C4637 & 1a<o1 1b*<o2 / 2a<1a1,1b1 / 3a<2a2L & IU/I/I & {-}{-}/{-}/{-} & N & W:1a,2a,3a \\
C4638 & 1a<o1 1b*<o2 / 2a<1a1,1b2 / 3a<2a0L & IU/I/N & {-}{-}/{-}/{-} & N & W:1a,2a \\
C4639 & 1a<o1 1b*<o2 / 2a<1a1,1b2 / 3a<2a1 & IU/I/I & {-}{-}/{-}/{-} & N & W:1a,2a,3a \\
C4640 & 1a<o1 1b*<o2 / 2a<1a1,1b2 / 3a<2a1L & IU/I/I & {-}{-}/{-}/{-} & N & W:1a,2a,3a \\
C4641 & 1a<o1 1b*<o2 / 2a<1a1,1b2 / 3a<2a2 & IU/I/I & {-}{-}/{-}/{-} & N & W:1a,2a,3a \\
C4642 & 1a<o1 1b*<o2 / 2a<1a1,1b2 / 3a<2a2L & IU/I/I & {-}{-}/{-}/{-} & N & W:1a,2a,3a \\
C4643 & 1a<o1 1b*<o2 / 2a<1a1L,1b1 / 3a<2a0L & IU/I/N & {-}{-}/{-}/{-} & N & W:1a,2a \\
C4644 & 1a<o1 1b*<o2 / 2a<1a1L,1b1 / 3a<2a1 & IU/I/I & {-}{-}/{-}/{-} & N & W:1a,2a,3a \\
C4645 & 1a<o1 1b*<o2 / 2a<1a1L,1b1 / 3a<2a1L & IU/I/I & {-}{-}/{-}/{-} & N & W:1a,2a,3a \\
C4646 & 1a<o1 1b*<o2 / 2a<1a1L,1b1 / 3a<2a2 & IU/I/I & {-}{-}/{-}/{-} & N & W:1a,2a,3a \\
C4647 & 1a<o1 1b*<o2 / 2a<1a1L,1b1 / 3a<2a2L & IU/I/I & {-}{-}/{-}/{-} & N & W:1a,2a,3a \\
C4648 & 1a<o1 1b*<o2 / 2a<1a1L,1b2 / 3a<2a0L & IU/I/N & {-}{-}/{-}/{-} & N & W:1a,2a \\
C4649 & 1a<o1 1b*<o2 / 2a<1a1L,1b2 / 3a<2a1 & IU/I/I & {-}{-}/{-}/{-} & N & W:1a,2a,3a \\
C4650 & 1a<o1 1b*<o2 / 2a<1a1L,1b2 / 3a<2a1L & IU/I/I & {-}{-}/{-}/{-} & N & W:1a,2a,3a \\
C4651 & 1a<o1 1b*<o2 / 2a<1a1L,1b2 / 3a<2a2 & IU/I/I & {-}{-}/{-}/{-} & N & W:1a,2a,3a \\
C4652 & 1a<o1 1b*<o2 / 2a<1a1L,1b2 / 3a<2a2L & IU/I/I & {-}{-}/{-}/{-} & N & W:1a,2a,3a \\
C4653 & 1a<o1 1b*<o2 / 2a<1a2,1b1 / 3a<2a0L & IU/I/N & {-}{-}/{-}/{-} & N & W:1a,2a \\
C4654 & 1a<o1 1b*<o2 / 2a<1a2,1b1 / 3a<2a1 & IU/I/I & {-}{-}/{-}/{-} & N & W:1a,2a,3a \\
C4655 & 1a<o1 1b*<o2 / 2a<1a2,1b1 / 3a<2a1L & IU/I/I & {-}{-}/{-}/{-} & N & W:1a,2a,3a \\
C4656 & 1a<o1 1b*<o2 / 2a<1a2,1b1 / 3a<2a2 & IU/I/I & {-}{-}/{-}/{-} & N & W:1a,2a,3a \\
C4657 & 1a<o1 1b*<o2 / 2a<1a2,1b1 / 3a<2a2L & IU/I/I & {-}{-}/{-}/{-} & N & W:1a,2a,3a \\
C4658 & 1a<o1 1b*<o2 / 2a<1a2,1b2 / 3a<2a0L & IU/I/N & {-}{-}/{-}/{-} & N & W:1a,2a \\
C4659 & 1a<o1 1b*<o2 / 2a<1a2,1b2 / 3a<2a1 & IU/I/I & {-}{-}/{-}/{-} & N & W:1a,2a,3a \\
C4660 & 1a<o1 1b*<o2 / 2a<1a2,1b2 / 3a<2a1L & IU/I/I & {-}{-}/{-}/{-} & N & W:1a,2a,3a \\
C4661 & 1a<o1 1b*<o2 / 2a<1a2,1b2 / 3a<2a2 & IU/I/I & {-}{-}/{-}/{-} & N & W:1a,2a,3a \\
C4662 & 1a<o1 1b*<o2 / 2a<1a2,1b2 / 3a<2a2L & IU/I/I & {-}{-}/{-}/{-} & N & W:1a,2a,3a \\
C4663 & 1a<o1 1b*<o2 / 2a<1a2L,1b1 / 3a<2a0L & IU/I/N & {-}{-}/{-}/{-} & N & W:1a,2a \\
C4664 & 1a<o1 1b*<o2 / 2a<1a2L,1b1 / 3a<2a1 & IU/I/I & {-}{-}/{-}/{-} & N & W:1a,2a,3a \\
C4665 & 1a<o1 1b*<o2 / 2a<1a2L,1b1 / 3a<2a1L & IU/I/I & {-}{-}/{-}/{-} & N & W:1a,2a,3a \\
C4666 & 1a<o1 1b*<o2 / 2a<1a2L,1b1 / 3a<2a2 & IU/I/I & {-}{-}/{-}/{-} & N & W:1a,2a,3a \\
C4667 & 1a<o1 1b*<o2 / 2a<1a2L,1b1 / 3a<2a2L & IU/I/I & {-}{-}/{-}/{-} & N & W:1a,2a,3a \\
C4668 & 1a<o1 1b*<o2 / 2a<1a2L,1b2 / 3a<2a0L & IU/I/N & {-}{-}/{-}/{-} & N & W:1a,2a \\
C4669 & 1a<o1 1b*<o2 / 2a<1a2L,1b2 / 3a<2a1 & IU/I/I & {-}{-}/{-}/{-} & N & W:1a,2a,3a \\
C4670 & 1a<o1 1b*<o2 / 2a<1a2L,1b2 / 3a<2a1L & IU/I/I & {-}{-}/{-}/{-} & N & W:1a,2a,3a \\
C4671 & 1a<o1 1b*<o2 / 2a<1a2L,1b2 / 3a<2a2 & IU/I/I & {-}{-}/{-}/{-} & N & W:1a,2a,3a \\
C4672 & 1a<o1 1b*<o2 / 2a<1a2L,1b2 / 3a<2a2L & IU/I/I & {-}{-}/{-}/{-} & N & W:1a,2a,3a \\
C4673 & 1a<o1L / 2a<1a0L 2b<1a1 / 3a*<2a1,2b1 & I/NI/U & +/+{-}/{-} & N & W:3a \\
C4674 & 1a<o1L / 2a<1a0L 2b<1a1 / 3a*<2a1,2b2 & I/NI/U & +/+{-}/{-} & N & W:3a \\
C4675 & 1a<o1L / 2a<1a0L 2b<1a1 / 3a*<2a2,2b1 & I/NI/U & +/+{-}/{-} & N & W:3a \\
C4676 & 1a<o1L / 2a<1a0L 2b<1a1 / 3a*<2a2,2b2 & I/NI/U & +/+{-}/{-} & N & W:3a \\
C4677 & 1a<o1L / 2a<1a0L 2b<1a2 / 3a*<2a1,2b1 & I/NI/U & +/+{-}/{-} & N & W:3a \\
C4678 & 1a<o1L / 2a<1a0L 2b<1a2 / 3a*<2a1,2b2 & I/NI/U & +/+{-}/{-} & N & W:3a \\
C4679 & 1a<o1L / 2a<1a0L 2b<1a2 / 3a*<2a2,2b1 & I/NI/U & +/+{-}/{-} & N & W:3a \\
C4680 & 1a<o1L / 2a<1a0L 2b<1a2 / 3a*<2a2,2b2 & I/NI/U & +/+{-}/{-} & N & W:3a \\
C4681 & 1a<o1L / 2a<1a0L 2b*<1a1 / 3a<2a0L,2b1 & I/NU/U & +/+{-}/+ & N & W:2b,3a \\
C4682 & 1a<o1L / 2a<1a0L 2b*<1a1 / 3a<2a0L,2b2 & I/NU/U & +/+{-}/+ & N & W:2b,3a \\
C4683 & 1a<o1L / 2a<1a0L 2b*<1a1 / 3a<2a1,2b1 & I/NU/U & +/+{-}/{-} & N & W:2b,3a \\
C4684 & 1a<o1L / 2a<1a0L 2b*<1a1 / 3a<2a1,2b2 & I/NU/U & +/+{-}/{-} & N & W:2b,3a \\
C4685 & 1a<o1L / 2a<1a0L 2b*<1a1 / 3a<2a1L,2b1 & I/NU/U & +/+{-}/+ & N & W:2b,3a \\
C4686 & 1a<o1L / 2a<1a0L 2b*<1a1 / 3a<2a1L,2b2 & I/NU/U & +/+{-}/+ & N & W:2b,3a \\
C4687 & 1a<o1L / 2a<1a0L 2b*<1a1 / 3a<2a2,2b1 & I/NU/U & +/+{-}/{-} & N & W:2b,3a \\
C4688 & 1a<o1L / 2a<1a0L 2b*<1a1 / 3a<2a2,2b2 & I/NU/U & +/+{-}/{-} & N & W:2b,3a \\
C4689 & 1a<o1L / 2a<1a0L 2b*<1a1 / 3a<2a2L,2b1 & I/NU/U & +/+{-}/+ & N & W:2b,3a \\
C4690 & 1a<o1L / 2a<1a0L 2b*<1a1 / 3a<2a2L,2b2 & I/NU/U & +/+{-}/+ & N & W:2b,3a \\
C4691 & 1a<o1L / 2a<1a0L 2b*<1a2 / 3a<2a0L,2b1 & I/NU/U & +/+{-}/+ & N & W:2b,3a \\
C4692 & 1a<o1L / 2a<1a0L 2b*<1a2 / 3a<2a0L,2b2 & I/NU/U & +/+{-}/+ & N & W:2b,3a \\
C4693 & 1a<o1L / 2a<1a0L 2b*<1a2 / 3a<2a1,2b1 & I/NU/U & +/+{-}/{-} & N & W:2b,3a \\
C4694 & 1a<o1L / 2a<1a0L 2b*<1a2 / 3a<2a1,2b2 & I/NU/U & +/+{-}/{-} & N & W:2b,3a \\
C4695 & 1a<o1L / 2a<1a0L 2b*<1a2 / 3a<2a1L,2b1 & I/NU/U & +/+{-}/+ & N & W:2b,3a \\
C4696 & 1a<o1L / 2a<1a0L 2b*<1a2 / 3a<2a1L,2b2 & I/NU/U & +/+{-}/+ & N & W:2b,3a \\
C4697 & 1a<o1L / 2a<1a0L 2b*<1a2 / 3a<2a2,2b1 & I/NU/U & +/+{-}/{-} & N & W:2b,3a \\
C4698 & 1a<o1L / 2a<1a0L 2b*<1a2 / 3a<2a2,2b2 & I/NU/U & +/+{-}/{-} & N & W:2b,3a \\
C4699 & 1a<o1L / 2a<1a0L 2b*<1a2 / 3a<2a2L,2b1 & I/NU/U & +/+{-}/+ & N & W:2b,3a \\
C4700 & 1a<o1L / 2a<1a0L 2b*<1a2 / 3a<2a2L,2b2 & I/NU/U & +/+{-}/+ & N & W:2b,3a \\
C4701 & 1a<o1L / 2a<1a1 2b<1a1 / 3a*<2a1,2b1 & I/EE/U & +/{-}{-}/{-} & E & W:1a \\
C4702 & 1a<o1L / 2a<1a1 2b<1a1 / 3a*<2a1,2b2 & I/EE/U & +/{-}{-}/{-} & E & W:1a \\
C4703 & 1a<o1L / 2a<1a1 2b<1a1 / 3a*<2a2,2b2 & I/EE/U & +/{-}{-}/{-} & E & W:1a \\
C4704 & 1a<o1L / 2a<1a1 2b<1a1L / 3a*<2a1,2b1 & I/LI/U & +/{-}+/{-} & B & CP:2b W:1a,2b \\
C4705 & 1a<o1L / 2a<1a1 2b<1a1L / 3a*<2a1,2b2 & I/LI/U & +/{-}+/{-} & B & CP:2b W:1a,2b \\
C4706 & 1a<o1L / 2a<1a1 2b<1a1L / 3a*<2a2,2b1 & I/LI/U & +/{-}+/{-} & B & CP:2b W:1a,2b \\
C4707 & 1a<o1L / 2a<1a1 2b<1a1L / 3a*<2a2,2b2 & I/LI/U & +/{-}+/{-} & B & CP:2b W:1a,2b \\
C4708 & 1a<o1L / 2a<1a1 2b<1a2 / 3a*<2a1,2b1 & I/EE/U & +/{-}{-}/{-} & E & C:2b W:1a \\
C4709 & 1a<o1L / 2a<1a1 2b<1a2 / 3a*<2a1,2b2 & I/EE/U & +/{-}{-}/{-} & E & C:2b W:1a \\
C4710 & 1a<o1L / 2a<1a1 2b<1a2 / 3a*<2a2,2b1 & I/EE/U & +/{-}{-}/{-} & E & C:2b W:1a \\
C4711 & 1a<o1L / 2a<1a1 2b<1a2 / 3a*<2a2,2b2 & I/EE/U & +/{-}{-}/{-} & E & C:2b W:1a \\
C4712 & 1a<o1L / 2a<1a1 2b<1a2L / 3a*<2a1,2b1 & I/LI/U & +/{-}+/{-} & B & P:2b W:1a,2b \\
C4713 & 1a<o1L / 2a<1a1 2b<1a2L / 3a*<2a1,2b2 & I/LI/U & +/{-}+/{-} & B & P:2b W:1a,2b \\
C4714 & 1a<o1L / 2a<1a1 2b<1a2L / 3a*<2a2,2b1 & I/LI/U & +/{-}+/{-} & B & P:2b W:1a,2b \\
C4715 & 1a<o1L / 2a<1a1 2b<1a2L / 3a*<2a2,2b2 & I/LI/U & +/{-}+/{-} & B & P:2b W:1a,2b \\
C4716 & 1a<o1L / 2a<1a1 2b*<1a1 / 3a<2a0L,2b1 & I/IU/U & +/{-}{-}/{-} & N & W:1a,2a \\
C4717 & 1a<o1L / 2a<1a1 2b*<1a1 / 3a<2a0L,2b2 & I/IU/U & +/{-}{-}/{-} & N & W:1a,2a \\
C4718 & 1a<o1L / 2a<1a1 2b*<1a1 / 3a<2a1,2b1 & I/IU/I & +/{-}{-}/{-} & N & W:1a,2a,3a \\
C4719 & 1a<o1L / 2a<1a1 2b*<1a1 / 3a<2a1,2b2 & I/IU/I & +/{-}{-}/{-} & N & W:1a,2a,3a \\
C4720 & 1a<o1L / 2a<1a1 2b*<1a1 / 3a<2a1L,2b1 & I/IU/I & +/{-}{-}/{-} & N & W:1a,2a,3a \\
C4721 & 1a<o1L / 2a<1a1 2b*<1a1 / 3a<2a1L,2b2 & I/IU/I & +/{-}{-}/{-} & N & W:1a,2a,3a \\
C4722 & 1a<o1L / 2a<1a1 2b*<1a1 / 3a<2a2,2b1 & I/IU/I & +/{-}{-}/{-} & N & W:1a,2a,3a \\
C4723 & 1a<o1L / 2a<1a1 2b*<1a1 / 3a<2a2,2b2 & I/IU/I & +/{-}{-}/{-} & N & W:1a,2a,3a \\
C4724 & 1a<o1L / 2a<1a1 2b*<1a1 / 3a<2a2L,2b1 & I/IU/I & +/{-}{-}/{-} & N & W:1a,2a,3a \\
C4725 & 1a<o1L / 2a<1a1 2b*<1a1 / 3a<2a2L,2b2 & I/IU/I & +/{-}{-}/{-} & N & W:1a,2a,3a \\
C4726 & 1a<o1L / 2a<1a1 2b*<1a2 / 3a<2a0L,2b1 & I/IU/U & +/{-}{-}/{-} & N & W:1a,2a \\
C4727 & 1a<o1L / 2a<1a1 2b*<1a2 / 3a<2a0L,2b2 & I/IU/U & +/{-}{-}/{-} & N & W:1a,2a \\
C4728 & 1a<o1L / 2a<1a1 2b*<1a2 / 3a<2a1,2b1 & I/IU/I & +/{-}{-}/{-} & N & W:1a,2a,3a \\
C4729 & 1a<o1L / 2a<1a1 2b*<1a2 / 3a<2a1,2b2 & I/IU/I & +/{-}{-}/{-} & N & W:1a,2a,3a \\
C4730 & 1a<o1L / 2a<1a1 2b*<1a2 / 3a<2a1L,2b1 & I/IU/I & +/{-}{-}/{-} & N & W:1a,2a,3a \\
C4731 & 1a<o1L / 2a<1a1 2b*<1a2 / 3a<2a1L,2b2 & I/IU/I & +/{-}{-}/{-} & N & W:1a,2a,3a \\
C4732 & 1a<o1L / 2a<1a1 2b*<1a2 / 3a<2a2,2b1 & I/IU/I & +/{-}{-}/{-} & N & W:1a,2a,3a \\
C4733 & 1a<o1L / 2a<1a1 2b*<1a2 / 3a<2a2,2b2 & I/IU/I & +/{-}{-}/{-} & N & W:1a,2a,3a \\
C4734 & 1a<o1L / 2a<1a1 2b*<1a2 / 3a<2a2L,2b1 & I/IU/I & +/{-}{-}/{-} & N & W:1a,2a,3a \\
C4735 & 1a<o1L / 2a<1a1 2b*<1a2 / 3a<2a2L,2b2 & I/IU/I & +/{-}{-}/{-} & N & W:1a,2a,3a \\
C4736 & 1a<o1L / 2a<1a1L 2b<1a2 / 3a*<2a1,2b1 & I/IL/U & +/+{-}/{-} & B & C:2a,2b P:2a W:1a,2a \\
C4737 & 1a<o1L / 2a<1a1L 2b<1a2 / 3a*<2a1,2b2 & I/IL/U & +/+{-}/{-} & B & C:2a,2b P:2a W:1a,2a \\
C4738 & 1a<o1L / 2a<1a1L 2b<1a2 / 3a*<2a2,2b1 & I/IL/U & +/+{-}/{-} & B & C:2a,2b P:2a W:1a,2a \\
C4739 & 1a<o1L / 2a<1a1L 2b<1a2 / 3a*<2a2,2b2 & I/IL/U & +/+{-}/{-} & B & C:2a,2b P:2a W:1a,2a \\
C4740 & 1a<o1L / 2a<1a1L 2b*<1a1 / 3a<2a0L,2b1 & I/IU/U & +/+{-}/+ & N & W:1a,2a \\
C4741 & 1a<o1L / 2a<1a1L 2b*<1a1 / 3a<2a0L,2b2 & I/IU/U & +/+{-}/+ & N & W:1a,2a \\
C4742 & 1a<o1L / 2a<1a1L 2b*<1a1 / 3a<2a1,2b1 & I/IU/I & +/+{-}/{-} & N & W:1a,2a,3a \\
C4743 & 1a<o1L / 2a<1a1L 2b*<1a1 / 3a<2a1,2b2 & I/IU/I & +/+{-}/{-} & N & W:1a,2a,3a \\
C4744 & 1a<o1L / 2a<1a1L 2b*<1a1 / 3a<2a1L,2b1 & I/IU/I & +/+{-}/+ & N & W:1a,2a,3a \\
C4745 & 1a<o1L / 2a<1a1L 2b*<1a1 / 3a<2a1L,2b2 & I/IU/I & +/+{-}/+ & N & W:1a,2a,3a \\
C4746 & 1a<o1L / 2a<1a1L 2b*<1a1 / 3a<2a2,2b1 & I/IU/I & +/+{-}/{-} & N & W:1a,2a,3a \\
C4747 & 1a<o1L / 2a<1a1L 2b*<1a1 / 3a<2a2,2b2 & I/IU/I & +/+{-}/{-} & N & W:1a,2a,3a \\
C4748 & 1a<o1L / 2a<1a1L 2b*<1a1 / 3a<2a2L,2b1 & I/IU/I & +/+{-}/+ & N & W:1a,2a,3a \\
C4749 & 1a<o1L / 2a<1a1L 2b*<1a1 / 3a<2a2L,2b2 & I/IU/I & +/+{-}/+ & N & W:1a,2a,3a \\
C4750 & 1a<o1L / 2a<1a1L 2b*<1a2 / 3a<2a0L,2b1 & I/IU/U & +/+{-}/+ & N & W:1a,2a \\
C4751 & 1a<o1L / 2a<1a1L 2b*<1a2 / 3a<2a0L,2b2 & I/IU/U & +/+{-}/+ & N & W:1a,2a \\
C4752 & 1a<o1L / 2a<1a1L 2b*<1a2 / 3a<2a1,2b1 & I/IU/I & +/+{-}/{-} & N & W:1a,2a,3a \\
C4753 & 1a<o1L / 2a<1a1L 2b*<1a2 / 3a<2a1,2b2 & I/IU/I & +/+{-}/{-} & N & W:1a,2a,3a \\
C4754 & 1a<o1L / 2a<1a1L 2b*<1a2 / 3a<2a1L,2b1 & I/IU/I & +/+{-}/+ & N & W:1a,2a,3a \\
C4755 & 1a<o1L / 2a<1a1L 2b*<1a2 / 3a<2a1L,2b2 & I/IU/I & +/+{-}/+ & N & W:1a,2a,3a \\
C4756 & 1a<o1L / 2a<1a1L 2b*<1a2 / 3a<2a2,2b1 & I/IU/I & +/+{-}/{-} & N & W:1a,2a,3a \\
C4757 & 1a<o1L / 2a<1a1L 2b*<1a2 / 3a<2a2,2b2 & I/IU/I & +/+{-}/{-} & N & W:1a,2a,3a \\
C4758 & 1a<o1L / 2a<1a1L 2b*<1a2 / 3a<2a2L,2b1 & I/IU/I & +/+{-}/+ & N & W:1a,2a,3a \\
C4759 & 1a<o1L / 2a<1a1L 2b*<1a2 / 3a<2a2L,2b2 & I/IU/I & +/+{-}/+ & N & W:1a,2a,3a \\
C4760 & 1a<o1L / 2a<1a2 2b<1a2 / 3a*<2a1,2b1 & I/EE/U & +/{-}{-}/{-} & E & W:1a \\
C4761 & 1a<o1L / 2a<1a2 2b<1a2 / 3a*<2a1,2b2 & I/EE/U & +/{-}{-}/{-} & E & W:1a \\
C4762 & 1a<o1L / 2a<1a2 2b<1a2 / 3a*<2a2,2b2 & I/EE/U & +/{-}{-}/{-} & E & W:1a \\
C4763 & 1a<o1L / 2a<1a2 2b<1a2L / 3a*<2a1,2b1 & I/LI/U & +/{-}+/{-} & B & CP:2b W:1a,2b \\
C4764 & 1a<o1L / 2a<1a2 2b<1a2L / 3a*<2a1,2b2 & I/LI/U & +/{-}+/{-} & B & CP:2b W:1a,2b \\
C4765 & 1a<o1L / 2a<1a2 2b<1a2L / 3a*<2a2,2b1 & I/LI/U & +/{-}+/{-} & B & CP:2b W:1a,2b \\
C4766 & 1a<o1L / 2a<1a2 2b<1a2L / 3a*<2a2,2b2 & I/LI/U & +/{-}+/{-} & B & CP:2b W:1a,2b \\
C4767 & 1a<o1L / 2a<1a2 2b*<1a1 / 3a<2a0L,2b1 & I/IU/U & +/{-}{-}/{-} & N & W:1a,2a \\
C4768 & 1a<o1L / 2a<1a2 2b*<1a1 / 3a<2a0L,2b2 & I/IU/U & +/{-}{-}/{-} & N & W:1a,2a \\
C4769 & 1a<o1L / 2a<1a2 2b*<1a1 / 3a<2a1,2b1 & I/IU/I & +/{-}{-}/{-} & N & W:1a,2a,3a \\
C4770 & 1a<o1L / 2a<1a2 2b*<1a1 / 3a<2a1,2b2 & I/IU/I & +/{-}{-}/{-} & N & W:1a,2a,3a \\
C4771 & 1a<o1L / 2a<1a2 2b*<1a1 / 3a<2a1L,2b1 & I/IU/I & +/{-}{-}/{-} & N & W:1a,2a,3a \\
C4772 & 1a<o1L / 2a<1a2 2b*<1a1 / 3a<2a1L,2b2 & I/IU/I & +/{-}{-}/{-} & N & W:1a,2a,3a \\
C4773 & 1a<o1L / 2a<1a2 2b*<1a1 / 3a<2a2,2b1 & I/IU/I & +/{-}{-}/{-} & N & W:1a,2a,3a \\
C4774 & 1a<o1L / 2a<1a2 2b*<1a1 / 3a<2a2,2b2 & I/IU/I & +/{-}{-}/{-} & N & W:1a,2a,3a \\
C4775 & 1a<o1L / 2a<1a2 2b*<1a1 / 3a<2a2L,2b1 & I/IU/I & +/{-}{-}/{-} & N & W:1a,2a,3a \\
C4776 & 1a<o1L / 2a<1a2 2b*<1a1 / 3a<2a2L,2b2 & I/IU/I & +/{-}{-}/{-} & N & W:1a,2a,3a \\
C4777 & 1a<o1L / 2a<1a2 2b*<1a2 / 3a<2a0L,2b1 & I/IU/U & +/{-}{-}/{-} & N & W:1a,2a \\
C4778 & 1a<o1L / 2a<1a2 2b*<1a2 / 3a<2a0L,2b2 & I/IU/U & +/{-}{-}/{-} & N & W:1a,2a \\
C4779 & 1a<o1L / 2a<1a2 2b*<1a2 / 3a<2a1,2b1 & I/IU/I & +/{-}{-}/{-} & N & W:1a,2a,3a \\
C4780 & 1a<o1L / 2a<1a2 2b*<1a2 / 3a<2a1,2b2 & I/IU/I & +/{-}{-}/{-} & N & W:1a,2a,3a \\
C4781 & 1a<o1L / 2a<1a2 2b*<1a2 / 3a<2a1L,2b1 & I/IU/I & +/{-}{-}/{-} & N & W:1a,2a,3a \\
C4782 & 1a<o1L / 2a<1a2 2b*<1a2 / 3a<2a1L,2b2 & I/IU/I & +/{-}{-}/{-} & N & W:1a,2a,3a \\
C4783 & 1a<o1L / 2a<1a2 2b*<1a2 / 3a<2a2,2b1 & I/IU/I & +/{-}{-}/{-} & N & W:1a,2a,3a \\
C4784 & 1a<o1L / 2a<1a2 2b*<1a2 / 3a<2a2,2b2 & I/IU/I & +/{-}{-}/{-} & N & W:1a,2a,3a \\
C4785 & 1a<o1L / 2a<1a2 2b*<1a2 / 3a<2a2L,2b1 & I/IU/I & +/{-}{-}/{-} & N & W:1a,2a,3a \\
C4786 & 1a<o1L / 2a<1a2 2b*<1a2 / 3a<2a2L,2b2 & I/IU/I & +/{-}{-}/{-} & N & W:1a,2a,3a \\
C4787 & 1a<o1L / 2a<1a2L 2b*<1a1 / 3a<2a0L,2b1 & I/IU/U & +/+{-}/+ & N & W:1a,2a \\
C4788 & 1a<o1L / 2a<1a2L 2b*<1a1 / 3a<2a0L,2b2 & I/IU/U & +/+{-}/+ & N & W:1a,2a \\
C4789 & 1a<o1L / 2a<1a2L 2b*<1a1 / 3a<2a1,2b1 & I/IU/I & +/+{-}/{-} & N & W:1a,2a,3a \\
C4790 & 1a<o1L / 2a<1a2L 2b*<1a1 / 3a<2a1,2b2 & I/IU/I & +/+{-}/{-} & N & W:1a,2a,3a \\
C4791 & 1a<o1L / 2a<1a2L 2b*<1a1 / 3a<2a1L,2b1 & I/IU/I & +/+{-}/+ & N & W:1a,2a,3a \\
C4792 & 1a<o1L / 2a<1a2L 2b*<1a1 / 3a<2a1L,2b2 & I/IU/I & +/+{-}/+ & N & W:1a,2a,3a \\
C4793 & 1a<o1L / 2a<1a2L 2b*<1a1 / 3a<2a2,2b1 & I/IU/I & +/+{-}/{-} & N & W:1a,2a,3a \\
C4794 & 1a<o1L / 2a<1a2L 2b*<1a1 / 3a<2a2,2b2 & I/IU/I & +/+{-}/{-} & N & W:1a,2a,3a \\
C4795 & 1a<o1L / 2a<1a2L 2b*<1a1 / 3a<2a2L,2b1 & I/IU/I & +/+{-}/+ & N & W:1a,2a,3a \\
C4796 & 1a<o1L / 2a<1a2L 2b*<1a1 / 3a<2a2L,2b2 & I/IU/I & +/+{-}/+ & N & W:1a,2a,3a \\
C4797 & 1a<o1L / 2a<1a2L 2b*<1a2 / 3a<2a0L,2b1 & I/IU/U & +/+{-}/+ & N & W:1a,2a \\
C4798 & 1a<o1L / 2a<1a2L 2b*<1a2 / 3a<2a0L,2b2 & I/IU/U & +/+{-}/+ & N & W:1a,2a \\
C4799 & 1a<o1L / 2a<1a2L 2b*<1a2 / 3a<2a1,2b1 & I/IU/I & +/+{-}/{-} & N & W:1a,2a,3a \\
C4800 & 1a<o1L / 2a<1a2L 2b*<1a2 / 3a<2a1,2b2 & I/IU/I & +/+{-}/{-} & N & W:1a,2a,3a \\
C4801 & 1a<o1L / 2a<1a2L 2b*<1a2 / 3a<2a1L,2b1 & I/IU/I & +/+{-}/+ & N & W:1a,2a,3a \\
C4802 & 1a<o1L / 2a<1a2L 2b*<1a2 / 3a<2a1L,2b2 & I/IU/I & +/+{-}/+ & N & W:1a,2a,3a \\
C4803 & 1a<o1L / 2a<1a2L 2b*<1a2 / 3a<2a2,2b1 & I/IU/I & +/+{-}/{-} & N & W:1a,2a,3a \\
C4804 & 1a<o1L / 2a<1a2L 2b*<1a2 / 3a<2a2,2b2 & I/IU/I & +/+{-}/{-} & N & W:1a,2a,3a \\
C4805 & 1a<o1L / 2a<1a2L 2b*<1a2 / 3a<2a2L,2b1 & I/IU/I & +/+{-}/+ & N & W:1a,2a,3a \\
C4806 & 1a<o1L / 2a<1a2L 2b*<1a2 / 3a<2a2L,2b2 & I/IU/I & +/+{-}/+ & N & W:1a,2a,3a \\
C4807 & 1a<o1L 1b<o2 / 2a<1a0L,1b1 / 3a*<2a1 & IL/D/U & +{-}/+/{-} & B & C:1a,1b,2a PW:1a \\
C4808 & 1a<o1L 1b<o2 / 2a<1a0L,1b1 / 3a*<2a2 & IL/D/U & +{-}/+/{-} & B & C:1a,1b,2a PW:1a \\
C4809 & 1a<o1L 1b<o2 / 2a<1a0L,1b2 / 3a*<2a1 & IL/D/U & +{-}/+/{-} & B & C:1a,1b,2a PW:1a \\
C4810 & 1a<o1L 1b<o2 / 2a<1a0L,1b2 / 3a*<2a2 & IL/D/U & +{-}/+/{-} & B & C:1a,1b,2a PW:1a \\
C4811 & 1a<o1L 1b<o2 / 2a<1a1,1b0L / 3a*<2a1 & IL/I/U & +{-}/{-}/{-} & B & C:1a,1b,2a PW:1a,2a \\
C4812 & 1a<o1L 1b<o2 / 2a<1a1,1b0L / 3a*<2a2 & IL/I/U & +{-}/{-}/{-} & B & C:1a,1b,2a PW:1a,2a \\
C4813 & 1a<o1L 1b<o2 / 2a<1a1,1b1 / 3a*<2a1 & IL/E/U & +{-}/{-}/{-} & BF & C:1a,1b PW:1a \\
C4814 & 1a<o1L 1b<o2 / 2a<1a1,1b1 / 3a*<2a2 & IL/E/U & +{-}/{-}/{-} & BF & C:1a,1b PW:1a \\
C4815 & 1a<o1L 1b<o2 / 2a<1a1,1b1L / 3a*<2a1 & IL/E/U & +{-}/{-}/{-} & BF & C:1a,1b PW:1a \\
C4816 & 1a<o1L 1b<o2 / 2a<1a1,1b1L / 3a*<2a2 & IL/E/U & +{-}/{-}/{-} & BF & C:1a,1b PW:1a \\
C4817 & 1a<o1L 1b<o2 / 2a<1a1,1b2 / 3a*<2a1 & IL/E/U & +{-}/{-}/{-} & BF & C:1a,1b,2a PW:1a \\
C4818 & 1a<o1L 1b<o2 / 2a<1a1,1b2 / 3a*<2a2 & IL/E/U & +{-}/{-}/{-} & BF & C:1a,1b,2a PW:1a \\
C4819 & 1a<o1L 1b<o2 / 2a<1a1,1b2L / 3a*<2a1 & IL/E/U & +{-}/{-}/{-} & BF & C:1a,1b,2a PW:1a \\
C4820 & 1a<o1L 1b<o2 / 2a<1a1,1b2L / 3a*<2a2 & IL/E/U & +{-}/{-}/{-} & BF & C:1a,1b,2a PW:1a \\
C4821 & 1a<o1L 1b<o2 / 2a<1a1L,1b1 / 3a*<2a1 & IL/I/U & +{-}/+/{-} & BF & C:1a,1b,2a PW:1a,2a \\
C4822 & 1a<o1L 1b<o2 / 2a<1a1L,1b1 / 3a*<2a2 & IL/I/U & +{-}/+/{-} & BF & C:1a,1b,2a PW:1a,2a \\
C4823 & 1a<o1L 1b<o2 / 2a<1a1L,1b2 / 3a*<2a1 & IL/I/U & +{-}/+/{-} & BF & C:1a,1b PW:1a,2a \\
C4824 & 1a<o1L 1b<o2 / 2a<1a1L,1b2 / 3a*<2a2 & IL/I/U & +{-}/+/{-} & BF & C:1a,1b PW:1a,2a \\
C4825 & 1a<o1L 1b<o2 / 2a<1a2,1b0L / 3a*<2a1 & IL/I/U & +{-}/{-}/{-} & B & C:1a,1b,2a PW:1a,2a \\
C4826 & 1a<o1L 1b<o2 / 2a<1a2,1b0L / 3a*<2a2 & IL/I/U & +{-}/{-}/{-} & B & C:1a,1b,2a PW:1a,2a \\
C4827 & 1a<o1L 1b<o2 / 2a<1a2,1b1 / 3a*<2a1 & IL/E/U & +{-}/{-}/{-} & BF & C:1a,1b PW:1a \\
C4828 & 1a<o1L 1b<o2 / 2a<1a2,1b1 / 3a*<2a2 & IL/E/U & +{-}/{-}/{-} & BF & C:1a,1b PW:1a \\
C4829 & 1a<o1L 1b<o2 / 2a<1a2,1b1L / 3a*<2a1 & IL/E/U & +{-}/{-}/{-} & BF & C:1a,1b PW:1a \\
C4830 & 1a<o1L 1b<o2 / 2a<1a2,1b1L / 3a*<2a2 & IL/E/U & +{-}/{-}/{-} & BF & C:1a,1b PW:1a \\
C4831 & 1a<o1L 1b<o2 / 2a<1a2,1b2 / 3a*<2a1 & IL/E/U & +{-}/{-}/{-} & BF & C:1a,1b PW:1a \\
C4832 & 1a<o1L 1b<o2 / 2a<1a2,1b2 / 3a*<2a2 & IL/E/U & +{-}/{-}/{-} & BF & C:1a,1b PW:1a \\
C4833 & 1a<o1L 1b<o2 / 2a<1a2,1b2L / 3a*<2a1 & IL/E/U & +{-}/{-}/{-} & BF & C:1a,1b PW:1a \\
C4834 & 1a<o1L 1b<o2 / 2a<1a2,1b2L / 3a*<2a2 & IL/E/U & +{-}/{-}/{-} & BF & C:1a,1b PW:1a \\
C4835 & 1a<o1L 1b<o2 / 2a<1a2L,1b1 / 3a*<2a1 & IL/I/U & +{-}/+/{-} & BF & C:1a,1b,2a PW:1a,2a \\
C4836 & 1a<o1L 1b<o2 / 2a<1a2L,1b1 / 3a*<2a2 & IL/I/U & +{-}/+/{-} & BF & C:1a,1b,2a PW:1a,2a \\
C4837 & 1a<o1L 1b<o2 / 2a<1a2L,1b2 / 3a*<2a1 & IL/I/U & +{-}/+/{-} & BF & C:1a,1b,2a PW:1a,2a \\
C4838 & 1a<o1L 1b<o2 / 2a<1a2L,1b2 / 3a*<2a2 & IL/I/U & +{-}/+/{-} & BF & C:1a,1b,2a PW:1a,2a \\
C4839 & 1a<o1L 1b<o2 / 2a*<1a1,1b1 / 3a<2a1 & IL/U/U & +{-}/{-}/{-} & B & C:1a,1b PW:1a \\
C4840 & 1a<o1L 1b<o2 / 2a*<1a1,1b1 / 3a<2a2 & IL/U/U & +{-}/{-}/{-} & B & C:1a,1b PW:1a \\
C4841 & 1a<o1L 1b<o2 / 2a*<1a1,1b2 / 3a<2a1 & IL/U/U & +{-}/{-}/{-} & B & C:1a,1b PW:1a \\
C4842 & 1a<o1L 1b<o2 / 2a*<1a1,1b2 / 3a<2a2 & IL/U/U & +{-}/{-}/{-} & B & C:1a,1b PW:1a \\
C4843 & 1a<o1L 1b<o2 / 2a*<1a2,1b1 / 3a<2a1 & IL/U/U & +{-}/{-}/{-} & B & C:1a,1b PW:1a \\
C4844 & 1a<o1L 1b<o2 / 2a*<1a2,1b1 / 3a<2a2 & IL/U/U & +{-}/{-}/{-} & B & C:1a,1b PW:1a \\
C4845 & 1a<o1L 1b<o2 / 2a*<1a2,1b2 / 3a<2a1 & IL/U/U & +{-}/{-}/{-} & B & C:1a,1b PW:1a \\
C4846 & 1a<o1L 1b<o2 / 2a*<1a2,1b2 / 3a<2a2 & IL/U/U & +{-}/{-}/{-} & B & C:1a,1b PW:1a \\
C4847 & 1a<o1L 1b*<o1 / 2a<1a0L,1b1 / 3a<2a0L & IU/U/N & +{-}/+/+ & N & W:1a \\
C4848 & 1a<o1L 1b*<o1 / 2a<1a0L,1b1 / 3a<2a1 & IU/U/U & +{-}/+/{-} & N & W:1a \\
C4849 & 1a<o1L 1b*<o1 / 2a<1a0L,1b1 / 3a<2a1L & IU/U/U & +{-}/+/+ & N & W:1a \\
C4850 & 1a<o1L 1b*<o1 / 2a<1a0L,1b1 / 3a<2a2 & IU/U/U & +{-}/+/{-} & N & W:1a \\
C4851 & 1a<o1L 1b*<o1 / 2a<1a0L,1b1 / 3a<2a2L & IU/U/U & +{-}/+/+ & N & W:1a \\
C4852 & 1a<o1L 1b*<o1 / 2a<1a0L,1b2 / 3a<2a0L & IU/U/N & +{-}/+/+ & N & W:1a \\
C4853 & 1a<o1L 1b*<o1 / 2a<1a0L,1b2 / 3a<2a1 & IU/U/U & +{-}/+/{-} & N & W:1a \\
C4854 & 1a<o1L 1b*<o1 / 2a<1a0L,1b2 / 3a<2a1L & IU/U/U & +{-}/+/+ & N & W:1a \\
C4855 & 1a<o1L 1b*<o1 / 2a<1a0L,1b2 / 3a<2a2 & IU/U/U & +{-}/+/{-} & N & W:1a \\
C4856 & 1a<o1L 1b*<o1 / 2a<1a0L,1b2 / 3a<2a2L & IU/U/U & +{-}/+/+ & N & W:1a \\
C4857 & 1a<o1L 1b*<o1 / 2a<1a1,1b1 / 3a<2a0L & IU/I/N & +{-}/{-}/{-} & N & W:1a,2a \\
C4858 & 1a<o1L 1b*<o1 / 2a<1a1,1b1 / 3a<2a1 & IU/I/I & +{-}/{-}/{-} & N & W:1a,2a,3a \\
C4859 & 1a<o1L 1b*<o1 / 2a<1a1,1b1 / 3a<2a1L & IU/I/I & +{-}/{-}/{-} & N & W:1a,2a,3a \\
C4860 & 1a<o1L 1b*<o1 / 2a<1a1,1b1 / 3a<2a2 & IU/I/I & +{-}/{-}/{-} & N & W:1a,2a,3a \\
C4861 & 1a<o1L 1b*<o1 / 2a<1a1,1b1 / 3a<2a2L & IU/I/I & +{-}/{-}/{-} & N & W:1a,2a,3a \\
C4862 & 1a<o1L 1b*<o1 / 2a<1a1,1b2 / 3a<2a0L & IU/I/N & +{-}/{-}/{-} & N & W:1a,2a \\
C4863 & 1a<o1L 1b*<o1 / 2a<1a1,1b2 / 3a<2a1 & IU/I/I & +{-}/{-}/{-} & N & W:1a,2a,3a \\
C4864 & 1a<o1L 1b*<o1 / 2a<1a1,1b2 / 3a<2a1L & IU/I/I & +{-}/{-}/{-} & N & W:1a,2a,3a \\
C4865 & 1a<o1L 1b*<o1 / 2a<1a1,1b2 / 3a<2a2 & IU/I/I & +{-}/{-}/{-} & N & W:1a,2a,3a \\
C4866 & 1a<o1L 1b*<o1 / 2a<1a1,1b2 / 3a<2a2L & IU/I/I & +{-}/{-}/{-} & N & W:1a,2a,3a \\
C4867 & 1a<o1L 1b*<o1 / 2a<1a1L,1b1 / 3a<2a0L & IU/I/N & +{-}/+/+ & N & W:1a,2a \\
C4868 & 1a<o1L 1b*<o1 / 2a<1a1L,1b1 / 3a<2a1 & IU/I/I & +{-}/+/{-} & N & W:1a,2a,3a \\
C4869 & 1a<o1L 1b*<o1 / 2a<1a1L,1b1 / 3a<2a1L & IU/I/I & +{-}/+/+ & N & W:1a,2a,3a \\
C4870 & 1a<o1L 1b*<o1 / 2a<1a1L,1b1 / 3a<2a2 & IU/I/I & +{-}/+/{-} & N & W:1a,2a,3a \\
C4871 & 1a<o1L 1b*<o1 / 2a<1a1L,1b1 / 3a<2a2L & IU/I/I & +{-}/+/+ & N & W:1a,2a,3a \\
C4872 & 1a<o1L 1b*<o1 / 2a<1a1L,1b2 / 3a<2a0L & IU/I/N & +{-}/+/+ & N & W:1a,2a \\
C4873 & 1a<o1L 1b*<o1 / 2a<1a1L,1b2 / 3a<2a1 & IU/I/I & +{-}/+/{-} & N & W:1a,2a,3a \\
C4874 & 1a<o1L 1b*<o1 / 2a<1a1L,1b2 / 3a<2a1L & IU/I/I & +{-}/+/+ & N & W:1a,2a,3a \\
C4875 & 1a<o1L 1b*<o1 / 2a<1a1L,1b2 / 3a<2a2 & IU/I/I & +{-}/+/{-} & N & W:1a,2a,3a \\
C4876 & 1a<o1L 1b*<o1 / 2a<1a1L,1b2 / 3a<2a2L & IU/I/I & +{-}/+/+ & N & W:1a,2a,3a \\
C4877 & 1a<o1L 1b*<o1 / 2a<1a2,1b1 / 3a<2a0L & IU/I/N & +{-}/{-}/{-} & N & W:1a,2a \\
C4878 & 1a<o1L 1b*<o1 / 2a<1a2,1b1 / 3a<2a1 & IU/I/I & +{-}/{-}/{-} & N & W:1a,2a,3a \\
C4879 & 1a<o1L 1b*<o1 / 2a<1a2,1b1 / 3a<2a1L & IU/I/I & +{-}/{-}/{-} & N & W:1a,2a,3a \\
C4880 & 1a<o1L 1b*<o1 / 2a<1a2,1b1 / 3a<2a2 & IU/I/I & +{-}/{-}/{-} & N & W:1a,2a,3a \\
C4881 & 1a<o1L 1b*<o1 / 2a<1a2,1b1 / 3a<2a2L & IU/I/I & +{-}/{-}/{-} & N & W:1a,2a,3a \\
C4882 & 1a<o1L 1b*<o1 / 2a<1a2,1b2 / 3a<2a0L & IU/I/N & +{-}/{-}/{-} & N & W:1a,2a \\
C4883 & 1a<o1L 1b*<o1 / 2a<1a2,1b2 / 3a<2a1 & IU/I/I & +{-}/{-}/{-} & N & W:1a,2a,3a \\
C4884 & 1a<o1L 1b*<o1 / 2a<1a2,1b2 / 3a<2a1L & IU/I/I & +{-}/{-}/{-} & N & W:1a,2a,3a \\
C4885 & 1a<o1L 1b*<o1 / 2a<1a2,1b2 / 3a<2a2 & IU/I/I & +{-}/{-}/{-} & N & W:1a,2a,3a \\
C4886 & 1a<o1L 1b*<o1 / 2a<1a2,1b2 / 3a<2a2L & IU/I/I & +{-}/{-}/{-} & N & W:1a,2a,3a \\
C4887 & 1a<o1L 1b*<o1 / 2a<1a2L,1b1 / 3a<2a0L & IU/I/N & +{-}/+/+ & N & W:1a,2a \\
C4888 & 1a<o1L 1b*<o1 / 2a<1a2L,1b1 / 3a<2a1 & IU/I/I & +{-}/+/{-} & N & W:1a,2a,3a \\
C4889 & 1a<o1L 1b*<o1 / 2a<1a2L,1b1 / 3a<2a1L & IU/I/I & +{-}/+/+ & N & W:1a,2a,3a \\
C4890 & 1a<o1L 1b*<o1 / 2a<1a2L,1b1 / 3a<2a2 & IU/I/I & +{-}/+/{-} & N & W:1a,2a,3a \\
C4891 & 1a<o1L 1b*<o1 / 2a<1a2L,1b1 / 3a<2a2L & IU/I/I & +{-}/+/+ & N & W:1a,2a,3a \\
C4892 & 1a<o1L 1b*<o1 / 2a<1a2L,1b2 / 3a<2a0L & IU/I/N & +{-}/+/+ & N & W:1a,2a \\
C4893 & 1a<o1L 1b*<o1 / 2a<1a2L,1b2 / 3a<2a1 & IU/I/I & +{-}/+/{-} & N & W:1a,2a,3a \\
C4894 & 1a<o1L 1b*<o1 / 2a<1a2L,1b2 / 3a<2a1L & IU/I/I & +{-}/+/+ & N & W:1a,2a,3a \\
C4895 & 1a<o1L 1b*<o1 / 2a<1a2L,1b2 / 3a<2a2 & IU/I/I & +{-}/+/{-} & N & W:1a,2a,3a \\
C4896 & 1a<o1L 1b*<o1 / 2a<1a2L,1b2 / 3a<2a2L & IU/I/I & +{-}/+/+ & N & W:1a,2a,3a \\
C4897 & 1a<o1L 1b*<o2 / 2a<1a0L,1b1 / 3a<2a0L & IU/U/N & +{-}/+/+ & N & W:1a \\
C4898 & 1a<o1L 1b*<o2 / 2a<1a0L,1b1 / 3a<2a1 & IU/U/U & +{-}/+/{-} & N & W:1a \\
C4899 & 1a<o1L 1b*<o2 / 2a<1a0L,1b1 / 3a<2a1L & IU/U/U & +{-}/+/+ & N & W:1a \\
C4900 & 1a<o1L 1b*<o2 / 2a<1a0L,1b1 / 3a<2a2 & IU/U/U & +{-}/+/{-} & N & W:1a \\
C4901 & 1a<o1L 1b*<o2 / 2a<1a0L,1b1 / 3a<2a2L & IU/U/U & +{-}/+/+ & N & W:1a \\
C4902 & 1a<o1L 1b*<o2 / 2a<1a0L,1b2 / 3a<2a0L & IU/U/N & +{-}/+/+ & N & W:1a \\
C4903 & 1a<o1L 1b*<o2 / 2a<1a0L,1b2 / 3a<2a1 & IU/U/U & +{-}/+/{-} & N & W:1a \\
C4904 & 1a<o1L 1b*<o2 / 2a<1a0L,1b2 / 3a<2a1L & IU/U/U & +{-}/+/+ & N & W:1a \\
C4905 & 1a<o1L 1b*<o2 / 2a<1a0L,1b2 / 3a<2a2 & IU/U/U & +{-}/+/{-} & N & W:1a \\
C4906 & 1a<o1L 1b*<o2 / 2a<1a0L,1b2 / 3a<2a2L & IU/U/U & +{-}/+/+ & N & W:1a \\
C4907 & 1a<o1L 1b*<o2 / 2a<1a1,1b1 / 3a<2a0L & IU/I/N & +{-}/{-}/{-} & N & W:1a,2a \\
C4908 & 1a<o1L 1b*<o2 / 2a<1a1,1b1 / 3a<2a1 & IU/I/I & +{-}/{-}/{-} & N & W:1a,2a,3a \\
C4909 & 1a<o1L 1b*<o2 / 2a<1a1,1b1 / 3a<2a1L & IU/I/I & +{-}/{-}/{-} & N & W:1a,2a,3a \\
C4910 & 1a<o1L 1b*<o2 / 2a<1a1,1b1 / 3a<2a2 & IU/I/I & +{-}/{-}/{-} & N & W:1a,2a,3a \\
C4911 & 1a<o1L 1b*<o2 / 2a<1a1,1b1 / 3a<2a2L & IU/I/I & +{-}/{-}/{-} & N & W:1a,2a,3a \\
C4912 & 1a<o1L 1b*<o2 / 2a<1a1,1b2 / 3a<2a0L & IU/I/N & +{-}/{-}/{-} & N & W:1a,2a \\
C4913 & 1a<o1L 1b*<o2 / 2a<1a1,1b2 / 3a<2a1 & IU/I/I & +{-}/{-}/{-} & N & W:1a,2a,3a \\
C4914 & 1a<o1L 1b*<o2 / 2a<1a1,1b2 / 3a<2a1L & IU/I/I & +{-}/{-}/{-} & N & W:1a,2a,3a \\
C4915 & 1a<o1L 1b*<o2 / 2a<1a1,1b2 / 3a<2a2 & IU/I/I & +{-}/{-}/{-} & N & W:1a,2a,3a \\
C4916 & 1a<o1L 1b*<o2 / 2a<1a1,1b2 / 3a<2a2L & IU/I/I & +{-}/{-}/{-} & N & W:1a,2a,3a \\
C4917 & 1a<o1L 1b*<o2 / 2a<1a1L,1b1 / 3a<2a0L & IU/I/N & +{-}/+/+ & N & W:1a,2a \\
C4918 & 1a<o1L 1b*<o2 / 2a<1a1L,1b1 / 3a<2a1 & IU/I/I & +{-}/+/{-} & N & W:1a,2a,3a \\
C4919 & 1a<o1L 1b*<o2 / 2a<1a1L,1b1 / 3a<2a1L & IU/I/I & +{-}/+/+ & N & W:1a,2a,3a \\
C4920 & 1a<o1L 1b*<o2 / 2a<1a1L,1b1 / 3a<2a2 & IU/I/I & +{-}/+/{-} & N & W:1a,2a,3a \\
C4921 & 1a<o1L 1b*<o2 / 2a<1a1L,1b1 / 3a<2a2L & IU/I/I & +{-}/+/+ & N & W:1a,2a,3a \\
C4922 & 1a<o1L 1b*<o2 / 2a<1a1L,1b2 / 3a<2a0L & IU/I/N & +{-}/+/+ & N & W:1a,2a \\
C4923 & 1a<o1L 1b*<o2 / 2a<1a1L,1b2 / 3a<2a1 & IU/I/I & +{-}/+/{-} & N & W:1a,2a,3a \\
C4924 & 1a<o1L 1b*<o2 / 2a<1a1L,1b2 / 3a<2a1L & IU/I/I & +{-}/+/+ & N & W:1a,2a,3a \\
C4925 & 1a<o1L 1b*<o2 / 2a<1a1L,1b2 / 3a<2a2 & IU/I/I & +{-}/+/{-} & N & W:1a,2a,3a \\
C4926 & 1a<o1L 1b*<o2 / 2a<1a1L,1b2 / 3a<2a2L & IU/I/I & +{-}/+/+ & N & W:1a,2a,3a \\
C4927 & 1a<o1L 1b*<o2 / 2a<1a2,1b1 / 3a<2a0L & IU/I/N & +{-}/{-}/{-} & N & W:1a,2a \\
C4928 & 1a<o1L 1b*<o2 / 2a<1a2,1b1 / 3a<2a1 & IU/I/I & +{-}/{-}/{-} & N & W:1a,2a,3a \\
C4929 & 1a<o1L 1b*<o2 / 2a<1a2,1b1 / 3a<2a1L & IU/I/I & +{-}/{-}/{-} & N & W:1a,2a,3a \\
C4930 & 1a<o1L 1b*<o2 / 2a<1a2,1b1 / 3a<2a2 & IU/I/I & +{-}/{-}/{-} & N & W:1a,2a,3a \\
C4931 & 1a<o1L 1b*<o2 / 2a<1a2,1b1 / 3a<2a2L & IU/I/I & +{-}/{-}/{-} & N & W:1a,2a,3a \\
C4932 & 1a<o1L 1b*<o2 / 2a<1a2,1b2 / 3a<2a0L & IU/I/N & +{-}/{-}/{-} & N & W:1a,2a \\
C4933 & 1a<o1L 1b*<o2 / 2a<1a2,1b2 / 3a<2a1 & IU/I/I & +{-}/{-}/{-} & N & W:1a,2a,3a \\
C4934 & 1a<o1L 1b*<o2 / 2a<1a2,1b2 / 3a<2a1L & IU/I/I & +{-}/{-}/{-} & N & W:1a,2a,3a \\
C4935 & 1a<o1L 1b*<o2 / 2a<1a2,1b2 / 3a<2a2 & IU/I/I & +{-}/{-}/{-} & N & W:1a,2a,3a \\
C4936 & 1a<o1L 1b*<o2 / 2a<1a2,1b2 / 3a<2a2L & IU/I/I & +{-}/{-}/{-} & N & W:1a,2a,3a \\
C4937 & 1a<o1L 1b*<o2 / 2a<1a2L,1b1 / 3a<2a0L & IU/I/N & +{-}/+/+ & N & W:1a,2a \\
C4938 & 1a<o1L 1b*<o2 / 2a<1a2L,1b1 / 3a<2a1 & IU/I/I & +{-}/+/{-} & N & W:1a,2a,3a \\
C4939 & 1a<o1L 1b*<o2 / 2a<1a2L,1b1 / 3a<2a1L & IU/I/I & +{-}/+/+ & N & W:1a,2a,3a \\
C4940 & 1a<o1L 1b*<o2 / 2a<1a2L,1b1 / 3a<2a2 & IU/I/I & +{-}/+/{-} & N & W:1a,2a,3a \\
C4941 & 1a<o1L 1b*<o2 / 2a<1a2L,1b1 / 3a<2a2L & IU/I/I & +{-}/+/+ & N & W:1a,2a,3a \\
C4942 & 1a<o1L 1b*<o2 / 2a<1a2L,1b2 / 3a<2a0L & IU/I/N & +{-}/+/+ & N & W:1a,2a \\
C4943 & 1a<o1L 1b*<o2 / 2a<1a2L,1b2 / 3a<2a1 & IU/I/I & +{-}/+/{-} & N & W:1a,2a,3a \\
C4944 & 1a<o1L 1b*<o2 / 2a<1a2L,1b2 / 3a<2a1L & IU/I/I & +{-}/+/+ & N & W:1a,2a,3a \\
C4945 & 1a<o1L 1b*<o2 / 2a<1a2L,1b2 / 3a<2a2 & IU/I/I & +{-}/+/{-} & N & W:1a,2a,3a \\
C4946 & 1a<o1L 1b*<o2 / 2a<1a2L,1b2 / 3a<2a2L & IU/I/I & +{-}/+/+ & N & W:1a,2a,3a \\
C4947 & 1a<o2 / 2a<1a0L 2b<1a1 / 3a*<2a1,2b1 & I/NI/U & {-}/{-}{-}/{-} & N & W:3a \\
C4948 & 1a<o2 / 2a<1a0L 2b<1a1 / 3a*<2a1,2b2 & I/NI/U & {-}/{-}{-}/{-} & N & W:3a \\
C4949 & 1a<o2 / 2a<1a0L 2b<1a1 / 3a*<2a2,2b1 & I/NI/U & {-}/{-}{-}/{-} & N & W:3a \\
C4950 & 1a<o2 / 2a<1a0L 2b<1a1 / 3a*<2a2,2b2 & I/NI/U & {-}/{-}{-}/{-} & N & W:3a \\
C4951 & 1a<o2 / 2a<1a0L 2b<1a2 / 3a*<2a1,2b1 & I/NI/U & {-}/{-}{-}/{-} & N & W:3a \\
C4952 & 1a<o2 / 2a<1a0L 2b<1a2 / 3a*<2a1,2b2 & I/NI/U & {-}/{-}{-}/{-} & N & W:3a \\
C4953 & 1a<o2 / 2a<1a0L 2b<1a2 / 3a*<2a2,2b1 & I/NI/U & {-}/{-}{-}/{-} & N & W:3a \\
C4954 & 1a<o2 / 2a<1a0L 2b<1a2 / 3a*<2a2,2b2 & I/NI/U & {-}/{-}{-}/{-} & N & W:3a \\
C4955 & 1a<o2 / 2a<1a0L 2b*<1a1 / 3a<2a0L,2b1 & I/NU/U & {-}/{-}{-}/{-} & N & W:2b,3a \\
C4956 & 1a<o2 / 2a<1a0L 2b*<1a1 / 3a<2a0L,2b2 & I/NU/U & {-}/{-}{-}/{-} & N & W:2b,3a \\
C4957 & 1a<o2 / 2a<1a0L 2b*<1a1 / 3a<2a1,2b1 & I/NU/U & {-}/{-}{-}/{-} & N & W:2b,3a \\
C4958 & 1a<o2 / 2a<1a0L 2b*<1a1 / 3a<2a1,2b2 & I/NU/U & {-}/{-}{-}/{-} & N & W:2b,3a \\
C4959 & 1a<o2 / 2a<1a0L 2b*<1a1 / 3a<2a1L,2b1 & I/NU/U & {-}/{-}{-}/{-} & N & W:2b,3a \\
C4960 & 1a<o2 / 2a<1a0L 2b*<1a1 / 3a<2a1L,2b2 & I/NU/U & {-}/{-}{-}/{-} & N & W:2b,3a \\
C4961 & 1a<o2 / 2a<1a0L 2b*<1a1 / 3a<2a2,2b1 & I/NU/U & {-}/{-}{-}/{-} & N & W:2b,3a \\
C4962 & 1a<o2 / 2a<1a0L 2b*<1a1 / 3a<2a2,2b2 & I/NU/U & {-}/{-}{-}/{-} & N & W:2b,3a \\
C4963 & 1a<o2 / 2a<1a0L 2b*<1a1 / 3a<2a2L,2b1 & I/NU/U & {-}/{-}{-}/{-} & N & W:2b,3a \\
C4964 & 1a<o2 / 2a<1a0L 2b*<1a1 / 3a<2a2L,2b2 & I/NU/U & {-}/{-}{-}/{-} & N & W:2b,3a \\
C4965 & 1a<o2 / 2a<1a0L 2b*<1a2 / 3a<2a0L,2b1 & I/NU/U & {-}/{-}{-}/{-} & N & W:2b,3a \\
C4966 & 1a<o2 / 2a<1a0L 2b*<1a2 / 3a<2a0L,2b2 & I/NU/U & {-}/{-}{-}/{-} & N & W:2b,3a \\
C4967 & 1a<o2 / 2a<1a0L 2b*<1a2 / 3a<2a1,2b1 & I/NU/U & {-}/{-}{-}/{-} & N & W:2b,3a \\
C4968 & 1a<o2 / 2a<1a0L 2b*<1a2 / 3a<2a1,2b2 & I/NU/U & {-}/{-}{-}/{-} & N & W:2b,3a \\
C4969 & 1a<o2 / 2a<1a0L 2b*<1a2 / 3a<2a1L,2b1 & I/NU/U & {-}/{-}{-}/{-} & N & W:2b,3a \\
C4970 & 1a<o2 / 2a<1a0L 2b*<1a2 / 3a<2a1L,2b2 & I/NU/U & {-}/{-}{-}/{-} & N & W:2b,3a \\
C4971 & 1a<o2 / 2a<1a0L 2b*<1a2 / 3a<2a2,2b1 & I/NU/U & {-}/{-}{-}/{-} & N & W:2b,3a \\
C4972 & 1a<o2 / 2a<1a0L 2b*<1a2 / 3a<2a2,2b2 & I/NU/U & {-}/{-}{-}/{-} & N & W:2b,3a \\
C4973 & 1a<o2 / 2a<1a0L 2b*<1a2 / 3a<2a2L,2b1 & I/NU/U & {-}/{-}{-}/{-} & N & W:2b,3a \\
C4974 & 1a<o2 / 2a<1a0L 2b*<1a2 / 3a<2a2L,2b2 & I/NU/U & {-}/{-}{-}/{-} & N & W:2b,3a \\
C4975 & 1a<o2 / 2a<1a1 2b<1a1 / 3a*<2a1,2b1 & I/EE/U & {-}/{-}{-}/{-} & E & W:1a \\
C4976 & 1a<o2 / 2a<1a1 2b<1a1 / 3a*<2a1,2b2 & I/EE/U & {-}/{-}{-}/{-} & E & W:1a \\
C4977 & 1a<o2 / 2a<1a1 2b<1a1 / 3a*<2a2,2b2 & I/EE/U & {-}/{-}{-}/{-} & E & W:1a \\
C4978 & 1a<o2 / 2a<1a1 2b<1a1L / 3a*<2a1,2b1 & I/LI/U & {-}/{-}{-}/{-} & B & CP:2b W:1a,2b \\
C4979 & 1a<o2 / 2a<1a1 2b<1a1L / 3a*<2a1,2b2 & I/LI/U & {-}/{-}{-}/{-} & B & CP:2b W:1a,2b \\
C4980 & 1a<o2 / 2a<1a1 2b<1a1L / 3a*<2a2,2b1 & I/LI/U & {-}/{-}{-}/{-} & B & CP:2b W:1a,2b \\
C4981 & 1a<o2 / 2a<1a1 2b<1a1L / 3a*<2a2,2b2 & I/LI/U & {-}/{-}{-}/{-} & B & CP:2b W:1a,2b \\
C4982 & 1a<o2 / 2a<1a1 2b<1a2 / 3a*<2a1,2b1 & I/EE/U & {-}/{-}{-}/{-} & E & C:2b W:1a \\
C4983 & 1a<o2 / 2a<1a1 2b<1a2 / 3a*<2a1,2b2 & I/EE/U & {-}/{-}{-}/{-} & E & C:2b W:1a \\
C4984 & 1a<o2 / 2a<1a1 2b<1a2 / 3a*<2a2,2b1 & I/EE/U & {-}/{-}{-}/{-} & E & C:2b W:1a \\
C4985 & 1a<o2 / 2a<1a1 2b<1a2 / 3a*<2a2,2b2 & I/EE/U & {-}/{-}{-}/{-} & E & C:2b W:1a \\
C4986 & 1a<o2 / 2a<1a1 2b<1a2L / 3a*<2a1,2b1 & I/LI/U & {-}/{-}{-}/{-} & B & P:2b W:1a,2b \\
C4987 & 1a<o2 / 2a<1a1 2b<1a2L / 3a*<2a1,2b2 & I/LI/U & {-}/{-}{-}/{-} & B & P:2b W:1a,2b \\
C4988 & 1a<o2 / 2a<1a1 2b<1a2L / 3a*<2a2,2b1 & I/LI/U & {-}/{-}{-}/{-} & B & P:2b W:1a,2b \\
C4989 & 1a<o2 / 2a<1a1 2b<1a2L / 3a*<2a2,2b2 & I/LI/U & {-}/{-}{-}/{-} & B & P:2b W:1a,2b \\
C4990 & 1a<o2 / 2a<1a1 2b*<1a1 / 3a<2a0L,2b1 & I/IU/U & {-}/{-}{-}/{-} & N & W:1a,2a \\
C4991 & 1a<o2 / 2a<1a1 2b*<1a1 / 3a<2a0L,2b2 & I/IU/U & {-}/{-}{-}/{-} & N & W:1a,2a \\
C4992 & 1a<o2 / 2a<1a1 2b*<1a1 / 3a<2a1,2b1 & I/IU/I & {-}/{-}{-}/{-} & N & W:1a,2a,3a \\
C4993 & 1a<o2 / 2a<1a1 2b*<1a1 / 3a<2a1,2b2 & I/IU/I & {-}/{-}{-}/{-} & N & W:1a,2a,3a \\
C4994 & 1a<o2 / 2a<1a1 2b*<1a1 / 3a<2a1L,2b1 & I/IU/I & {-}/{-}{-}/{-} & N & W:1a,2a,3a \\
C4995 & 1a<o2 / 2a<1a1 2b*<1a1 / 3a<2a1L,2b2 & I/IU/I & {-}/{-}{-}/{-} & N & W:1a,2a,3a \\
C4996 & 1a<o2 / 2a<1a1 2b*<1a1 / 3a<2a2,2b1 & I/IU/I & {-}/{-}{-}/{-} & N & W:1a,2a,3a \\
C4997 & 1a<o2 / 2a<1a1 2b*<1a1 / 3a<2a2,2b2 & I/IU/I & {-}/{-}{-}/{-} & N & W:1a,2a,3a \\
C4998 & 1a<o2 / 2a<1a1 2b*<1a1 / 3a<2a2L,2b1 & I/IU/I & {-}/{-}{-}/{-} & N & W:1a,2a,3a \\
C4999 & 1a<o2 / 2a<1a1 2b*<1a1 / 3a<2a2L,2b2 & I/IU/I & {-}/{-}{-}/{-} & N & W:1a,2a,3a \\
C5000 & 1a<o2 / 2a<1a1 2b*<1a2 / 3a<2a0L,2b1 & I/IU/U & {-}/{-}{-}/{-} & N & W:1a,2a \\
C5001 & 1a<o2 / 2a<1a1 2b*<1a2 / 3a<2a0L,2b2 & I/IU/U & {-}/{-}{-}/{-} & N & W:1a,2a \\
C5002 & 1a<o2 / 2a<1a1 2b*<1a2 / 3a<2a1,2b1 & I/IU/I & {-}/{-}{-}/{-} & N & W:1a,2a,3a \\
C5003 & 1a<o2 / 2a<1a1 2b*<1a2 / 3a<2a1,2b2 & I/IU/I & {-}/{-}{-}/{-} & N & W:1a,2a,3a \\
C5004 & 1a<o2 / 2a<1a1 2b*<1a2 / 3a<2a1L,2b1 & I/IU/I & {-}/{-}{-}/{-} & N & W:1a,2a,3a \\
C5005 & 1a<o2 / 2a<1a1 2b*<1a2 / 3a<2a1L,2b2 & I/IU/I & {-}/{-}{-}/{-} & N & W:1a,2a,3a \\
C5006 & 1a<o2 / 2a<1a1 2b*<1a2 / 3a<2a2,2b1 & I/IU/I & {-}/{-}{-}/{-} & N & W:1a,2a,3a \\
C5007 & 1a<o2 / 2a<1a1 2b*<1a2 / 3a<2a2,2b2 & I/IU/I & {-}/{-}{-}/{-} & N & W:1a,2a,3a \\
C5008 & 1a<o2 / 2a<1a1 2b*<1a2 / 3a<2a2L,2b1 & I/IU/I & {-}/{-}{-}/{-} & N & W:1a,2a,3a \\
C5009 & 1a<o2 / 2a<1a1 2b*<1a2 / 3a<2a2L,2b2 & I/IU/I & {-}/{-}{-}/{-} & N & W:1a,2a,3a \\
C5010 & 1a<o2 / 2a<1a1L 2b<1a2 / 3a*<2a1,2b1 & I/IL/U & {-}/{-}{-}/{-} & B & C:2a,2b P:2a W:1a,2a \\
C5011 & 1a<o2 / 2a<1a1L 2b<1a2 / 3a*<2a1,2b2 & I/IL/U & {-}/{-}{-}/{-} & B & C:2a,2b P:2a W:1a,2a \\
C5012 & 1a<o2 / 2a<1a1L 2b<1a2 / 3a*<2a2,2b1 & I/IL/U & {-}/{-}{-}/{-} & B & C:2a,2b P:2a W:1a,2a \\
C5013 & 1a<o2 / 2a<1a1L 2b<1a2 / 3a*<2a2,2b2 & I/IL/U & {-}/{-}{-}/{-} & B & C:2a,2b P:2a W:1a,2a \\
C5014 & 1a<o2 / 2a<1a1L 2b*<1a1 / 3a<2a0L,2b1 & I/IU/U & {-}/{-}{-}/{-} & N & W:1a,2a \\
C5015 & 1a<o2 / 2a<1a1L 2b*<1a1 / 3a<2a0L,2b2 & I/IU/U & {-}/{-}{-}/{-} & N & W:1a,2a \\
C5016 & 1a<o2 / 2a<1a1L 2b*<1a1 / 3a<2a1,2b1 & I/IU/I & {-}/{-}{-}/{-} & N & W:1a,2a,3a \\
C5017 & 1a<o2 / 2a<1a1L 2b*<1a1 / 3a<2a1,2b2 & I/IU/I & {-}/{-}{-}/{-} & N & W:1a,2a,3a \\
C5018 & 1a<o2 / 2a<1a1L 2b*<1a1 / 3a<2a1L,2b1 & I/IU/I & {-}/{-}{-}/{-} & N & W:1a,2a,3a \\
C5019 & 1a<o2 / 2a<1a1L 2b*<1a1 / 3a<2a1L,2b2 & I/IU/I & {-}/{-}{-}/{-} & N & W:1a,2a,3a \\
C5020 & 1a<o2 / 2a<1a1L 2b*<1a1 / 3a<2a2,2b1 & I/IU/I & {-}/{-}{-}/{-} & N & W:1a,2a,3a \\
C5021 & 1a<o2 / 2a<1a1L 2b*<1a1 / 3a<2a2,2b2 & I/IU/I & {-}/{-}{-}/{-} & N & W:1a,2a,3a \\
C5022 & 1a<o2 / 2a<1a1L 2b*<1a1 / 3a<2a2L,2b1 & I/IU/I & {-}/{-}{-}/{-} & N & W:1a,2a,3a \\
C5023 & 1a<o2 / 2a<1a1L 2b*<1a1 / 3a<2a2L,2b2 & I/IU/I & {-}/{-}{-}/{-} & N & W:1a,2a,3a \\
C5024 & 1a<o2 / 2a<1a1L 2b*<1a2 / 3a<2a0L,2b1 & I/IU/U & {-}/{-}{-}/{-} & N & W:1a,2a \\
C5025 & 1a<o2 / 2a<1a1L 2b*<1a2 / 3a<2a0L,2b2 & I/IU/U & {-}/{-}{-}/{-} & N & W:1a,2a \\
C5026 & 1a<o2 / 2a<1a1L 2b*<1a2 / 3a<2a1,2b1 & I/IU/I & {-}/{-}{-}/{-} & N & W:1a,2a,3a \\
C5027 & 1a<o2 / 2a<1a1L 2b*<1a2 / 3a<2a1,2b2 & I/IU/I & {-}/{-}{-}/{-} & N & W:1a,2a,3a \\
C5028 & 1a<o2 / 2a<1a1L 2b*<1a2 / 3a<2a1L,2b1 & I/IU/I & {-}/{-}{-}/{-} & N & W:1a,2a,3a \\
C5029 & 1a<o2 / 2a<1a1L 2b*<1a2 / 3a<2a1L,2b2 & I/IU/I & {-}/{-}{-}/{-} & N & W:1a,2a,3a \\
C5030 & 1a<o2 / 2a<1a1L 2b*<1a2 / 3a<2a2,2b1 & I/IU/I & {-}/{-}{-}/{-} & N & W:1a,2a,3a \\
C5031 & 1a<o2 / 2a<1a1L 2b*<1a2 / 3a<2a2,2b2 & I/IU/I & {-}/{-}{-}/{-} & N & W:1a,2a,3a \\
C5032 & 1a<o2 / 2a<1a1L 2b*<1a2 / 3a<2a2L,2b1 & I/IU/I & {-}/{-}{-}/{-} & N & W:1a,2a,3a \\
C5033 & 1a<o2 / 2a<1a1L 2b*<1a2 / 3a<2a2L,2b2 & I/IU/I & {-}/{-}{-}/{-} & N & W:1a,2a,3a \\
C5034 & 1a<o2 / 2a<1a2 2b<1a2 / 3a*<2a1,2b1 & I/EE/U & {-}/{-}{-}/{-} & E & W:1a \\
C5035 & 1a<o2 / 2a<1a2 2b<1a2 / 3a*<2a1,2b2 & I/EE/U & {-}/{-}{-}/{-} & E & W:1a \\
C5036 & 1a<o2 / 2a<1a2 2b<1a2 / 3a*<2a2,2b2 & I/EE/U & {-}/{-}{-}/{-} & E & W:1a \\
C5037 & 1a<o2 / 2a<1a2 2b<1a2L / 3a*<2a1,2b1 & I/LI/U & {-}/{-}{-}/{-} & B & CP:2b W:1a,2b \\
C5038 & 1a<o2 / 2a<1a2 2b<1a2L / 3a*<2a1,2b2 & I/LI/U & {-}/{-}{-}/{-} & B & CP:2b W:1a,2b \\
C5039 & 1a<o2 / 2a<1a2 2b<1a2L / 3a*<2a2,2b1 & I/LI/U & {-}/{-}{-}/{-} & B & CP:2b W:1a,2b \\
C5040 & 1a<o2 / 2a<1a2 2b<1a2L / 3a*<2a2,2b2 & I/LI/U & {-}/{-}{-}/{-} & B & CP:2b W:1a,2b \\
C5041 & 1a<o2 / 2a<1a2 2b*<1a1 / 3a<2a0L,2b1 & I/IU/U & {-}/{-}{-}/{-} & N & W:1a,2a \\
C5042 & 1a<o2 / 2a<1a2 2b*<1a1 / 3a<2a0L,2b2 & I/IU/U & {-}/{-}{-}/{-} & N & W:1a,2a \\
C5043 & 1a<o2 / 2a<1a2 2b*<1a1 / 3a<2a1,2b1 & I/IU/I & {-}/{-}{-}/{-} & N & W:1a,2a,3a \\
C5044 & 1a<o2 / 2a<1a2 2b*<1a1 / 3a<2a1,2b2 & I/IU/I & {-}/{-}{-}/{-} & N & W:1a,2a,3a \\
C5045 & 1a<o2 / 2a<1a2 2b*<1a1 / 3a<2a1L,2b1 & I/IU/I & {-}/{-}{-}/{-} & N & W:1a,2a,3a \\
C5046 & 1a<o2 / 2a<1a2 2b*<1a1 / 3a<2a1L,2b2 & I/IU/I & {-}/{-}{-}/{-} & N & W:1a,2a,3a \\
C5047 & 1a<o2 / 2a<1a2 2b*<1a1 / 3a<2a2,2b1 & I/IU/I & {-}/{-}{-}/{-} & N & W:1a,2a,3a \\
C5048 & 1a<o2 / 2a<1a2 2b*<1a1 / 3a<2a2,2b2 & I/IU/I & {-}/{-}{-}/{-} & N & W:1a,2a,3a \\
C5049 & 1a<o2 / 2a<1a2 2b*<1a1 / 3a<2a2L,2b1 & I/IU/I & {-}/{-}{-}/{-} & N & W:1a,2a,3a \\
C5050 & 1a<o2 / 2a<1a2 2b*<1a1 / 3a<2a2L,2b2 & I/IU/I & {-}/{-}{-}/{-} & N & W:1a,2a,3a \\
C5051 & 1a<o2 / 2a<1a2 2b*<1a2 / 3a<2a0L,2b1 & I/IU/U & {-}/{-}{-}/{-} & N & W:1a,2a \\
C5052 & 1a<o2 / 2a<1a2 2b*<1a2 / 3a<2a0L,2b2 & I/IU/U & {-}/{-}{-}/{-} & N & W:1a,2a \\
C5053 & 1a<o2 / 2a<1a2 2b*<1a2 / 3a<2a1,2b1 & I/IU/I & {-}/{-}{-}/{-} & N & W:1a,2a,3a \\
C5054 & 1a<o2 / 2a<1a2 2b*<1a2 / 3a<2a1,2b2 & I/IU/I & {-}/{-}{-}/{-} & N & W:1a,2a,3a \\
C5055 & 1a<o2 / 2a<1a2 2b*<1a2 / 3a<2a1L,2b1 & I/IU/I & {-}/{-}{-}/{-} & N & W:1a,2a,3a \\
C5056 & 1a<o2 / 2a<1a2 2b*<1a2 / 3a<2a1L,2b2 & I/IU/I & {-}/{-}{-}/{-} & N & W:1a,2a,3a \\
C5057 & 1a<o2 / 2a<1a2 2b*<1a2 / 3a<2a2,2b1 & I/IU/I & {-}/{-}{-}/{-} & N & W:1a,2a,3a \\
C5058 & 1a<o2 / 2a<1a2 2b*<1a2 / 3a<2a2,2b2 & I/IU/I & {-}/{-}{-}/{-} & N & W:1a,2a,3a \\
C5059 & 1a<o2 / 2a<1a2 2b*<1a2 / 3a<2a2L,2b1 & I/IU/I & {-}/{-}{-}/{-} & N & W:1a,2a,3a \\
C5060 & 1a<o2 / 2a<1a2 2b*<1a2 / 3a<2a2L,2b2 & I/IU/I & {-}/{-}{-}/{-} & N & W:1a,2a,3a \\
C5061 & 1a<o2 / 2a<1a2L 2b*<1a1 / 3a<2a0L,2b1 & I/IU/U & {-}/{-}{-}/{-} & N & W:1a,2a \\
C5062 & 1a<o2 / 2a<1a2L 2b*<1a1 / 3a<2a0L,2b2 & I/IU/U & {-}/{-}{-}/{-} & N & W:1a,2a \\
C5063 & 1a<o2 / 2a<1a2L 2b*<1a1 / 3a<2a1,2b1 & I/IU/I & {-}/{-}{-}/{-} & N & W:1a,2a,3a \\
C5064 & 1a<o2 / 2a<1a2L 2b*<1a1 / 3a<2a1,2b2 & I/IU/I & {-}/{-}{-}/{-} & N & W:1a,2a,3a \\
C5065 & 1a<o2 / 2a<1a2L 2b*<1a1 / 3a<2a1L,2b1 & I/IU/I & {-}/{-}{-}/{-} & N & W:1a,2a,3a \\
C5066 & 1a<o2 / 2a<1a2L 2b*<1a1 / 3a<2a1L,2b2 & I/IU/I & {-}/{-}{-}/{-} & N & W:1a,2a,3a \\
C5067 & 1a<o2 / 2a<1a2L 2b*<1a1 / 3a<2a2,2b1 & I/IU/I & {-}/{-}{-}/{-} & N & W:1a,2a,3a \\
C5068 & 1a<o2 / 2a<1a2L 2b*<1a1 / 3a<2a2,2b2 & I/IU/I & {-}/{-}{-}/{-} & N & W:1a,2a,3a \\
C5069 & 1a<o2 / 2a<1a2L 2b*<1a1 / 3a<2a2L,2b1 & I/IU/I & {-}/{-}{-}/{-} & N & W:1a,2a,3a \\
C5070 & 1a<o2 / 2a<1a2L 2b*<1a1 / 3a<2a2L,2b2 & I/IU/I & {-}/{-}{-}/{-} & N & W:1a,2a,3a \\
C5071 & 1a<o2 / 2a<1a2L 2b*<1a2 / 3a<2a0L,2b1 & I/IU/U & {-}/{-}{-}/{-} & N & W:1a,2a \\
C5072 & 1a<o2 / 2a<1a2L 2b*<1a2 / 3a<2a0L,2b2 & I/IU/U & {-}/{-}{-}/{-} & N & W:1a,2a \\
C5073 & 1a<o2 / 2a<1a2L 2b*<1a2 / 3a<2a1,2b1 & I/IU/I & {-}/{-}{-}/{-} & N & W:1a,2a,3a \\
C5074 & 1a<o2 / 2a<1a2L 2b*<1a2 / 3a<2a1,2b2 & I/IU/I & {-}/{-}{-}/{-} & N & W:1a,2a,3a \\
C5075 & 1a<o2 / 2a<1a2L 2b*<1a2 / 3a<2a1L,2b1 & I/IU/I & {-}/{-}{-}/{-} & N & W:1a,2a,3a \\
C5076 & 1a<o2 / 2a<1a2L 2b*<1a2 / 3a<2a1L,2b2 & I/IU/I & {-}/{-}{-}/{-} & N & W:1a,2a,3a \\
C5077 & 1a<o2 / 2a<1a2L 2b*<1a2 / 3a<2a2,2b1 & I/IU/I & {-}/{-}{-}/{-} & N & W:1a,2a,3a \\
C5078 & 1a<o2 / 2a<1a2L 2b*<1a2 / 3a<2a2,2b2 & I/IU/I & {-}/{-}{-}/{-} & N & W:1a,2a,3a \\
C5079 & 1a<o2 / 2a<1a2L 2b*<1a2 / 3a<2a2L,2b1 & I/IU/I & {-}/{-}{-}/{-} & N & W:1a,2a,3a \\
C5080 & 1a<o2 / 2a<1a2L 2b*<1a2 / 3a<2a2L,2b2 & I/IU/I & {-}/{-}{-}/{-} & N & W:1a,2a,3a \\
C5081 & 1a<o2 1b<o2 / 2a<1a0L,1b1 / 3a*<2a1 & EE/D/U & {-}{-}/{-}/{-} & E &  \\
C5082 & 1a<o2 1b<o2 / 2a<1a0L,1b1 / 3a*<2a2 & EE/D/U & {-}{-}/{-}/{-} & E &  \\
C5083 & 1a<o2 1b<o2 / 2a<1a0L,1b2 / 3a*<2a1 & EE/D/U & {-}{-}/{-}/{-} & E &  \\
C5084 & 1a<o2 1b<o2 / 2a<1a0L,1b2 / 3a*<2a2 & EE/D/U & {-}{-}/{-}/{-} & E &  \\
C5085 & 1a<o2 1b<o2 / 2a<1a1,1b1 / 3a*<2a1 & EE/D/U & {-}{-}/{-}/{-} & E &  \\
C5086 & 1a<o2 1b<o2 / 2a<1a1,1b1 / 3a*<2a2 & EE/D/U & {-}{-}/{-}/{-} & E &  \\
C5087 & 1a<o2 1b<o2 / 2a<1a1,1b1L / 3a*<2a1 & EE/D/U & {-}{-}/{-}/{-} & E &  \\
C5088 & 1a<o2 1b<o2 / 2a<1a1,1b1L / 3a*<2a2 & EE/D/U & {-}{-}/{-}/{-} & E &  \\
C5089 & 1a<o2 1b<o2 / 2a<1a1,1b2 / 3a*<2a1 & EE/D/U & {-}{-}/{-}/{-} & E &  \\
C5090 & 1a<o2 1b<o2 / 2a<1a1,1b2 / 3a*<2a2 & EE/D/U & {-}{-}/{-}/{-} & E &  \\
C5091 & 1a<o2 1b<o2 / 2a<1a1,1b2L / 3a*<2a1 & EE/D/U & {-}{-}/{-}/{-} & E &  \\
C5092 & 1a<o2 1b<o2 / 2a<1a1,1b2L / 3a*<2a2 & EE/D/U & {-}{-}/{-}/{-} & E &  \\
C5093 & 1a<o2 1b<o2 / 2a<1a1L,1b2 / 3a*<2a1 & EE/D/U & {-}{-}/{-}/{-} & E &  \\
C5094 & 1a<o2 1b<o2 / 2a<1a1L,1b2 / 3a*<2a2 & EE/D/U & {-}{-}/{-}/{-} & E &  \\
C5095 & 1a<o2 1b<o2 / 2a<1a2,1b2 / 3a*<2a1 & EE/D/U & {-}{-}/{-}/{-} & E &  \\
C5096 & 1a<o2 1b<o2 / 2a<1a2,1b2 / 3a*<2a2 & EE/D/U & {-}{-}/{-}/{-} & E &  \\
C5097 & 1a<o2 1b<o2 / 2a<1a2,1b2L / 3a*<2a1 & EE/D/U & {-}{-}/{-}/{-} & E &  \\
C5098 & 1a<o2 1b<o2 / 2a<1a2,1b2L / 3a*<2a2 & EE/D/U & {-}{-}/{-}/{-} & E &  \\
C5099 & 1a<o2 1b<o2 / 2a*<1a1,1b1 / 3a<2a1 & EE/U/U & {-}{-}/{-}/{-} & E &  \\
C5100 & 1a<o2 1b<o2 / 2a*<1a1,1b1 / 3a<2a2 & EE/U/U & {-}{-}/{-}/{-} & E &  \\
C5101 & 1a<o2 1b<o2 / 2a*<1a1,1b2 / 3a<2a1 & EE/U/U & {-}{-}/{-}/{-} & E &  \\
C5102 & 1a<o2 1b<o2 / 2a*<1a1,1b2 / 3a<2a2 & EE/U/U & {-}{-}/{-}/{-} & E &  \\
C5103 & 1a<o2 1b<o2 / 2a*<1a2,1b2 / 3a<2a1 & EE/U/U & {-}{-}/{-}/{-} & E &  \\
C5104 & 1a<o2 1b<o2 / 2a*<1a2,1b2 / 3a<2a2 & EE/U/U & {-}{-}/{-}/{-} & E &  \\
C5105 & 1a<o2 1b<o2L / 2a<1a0L,1b1 / 3a*<2a1 & LI/I/U & {-}+/{-}/{-} & B & CPW:1b,2a \\
C5106 & 1a<o2 1b<o2L / 2a<1a0L,1b1 / 3a*<2a2 & LI/I/U & {-}+/{-}/{-} & B & CPW:1b,2a \\
C5107 & 1a<o2 1b<o2L / 2a<1a0L,1b2 / 3a*<2a1 & LI/I/U & {-}+/{-}/{-} & B & CPW:1b,2a \\
C5108 & 1a<o2 1b<o2L / 2a<1a0L,1b2 / 3a*<2a2 & LI/I/U & {-}+/{-}/{-} & B & CPW:1b,2a \\
C5109 & 1a<o2 1b<o2L / 2a<1a1,1b0L / 3a*<2a1 & LI/D/U & {-}+/+/{-} & B & CPW:1b \\
C5110 & 1a<o2 1b<o2L / 2a<1a1,1b0L / 3a*<2a2 & LI/D/U & {-}+/+/{-} & B & CPW:1b \\
C5111 & 1a<o2 1b<o2L / 2a<1a1,1b1 / 3a*<2a1 & LI/E/U & {-}+/{-}/{-} & BF & CPW:1b \\
C5112 & 1a<o2 1b<o2L / 2a<1a1,1b1 / 3a*<2a2 & LI/E/U & {-}+/{-}/{-} & BF & CPW:1b \\
C5113 & 1a<o2 1b<o2L / 2a<1a1,1b1L / 3a*<2a1 & LI/I/U & {-}+/+/{-} & BF & CPW:1b,2a \\
C5114 & 1a<o2 1b<o2L / 2a<1a1,1b1L / 3a*<2a2 & LI/I/U & {-}+/+/{-} & BF & CPW:1b,2a \\
C5115 & 1a<o2 1b<o2L / 2a<1a1,1b2 / 3a*<2a1 & LI/E/U & {-}+/{-}/{-} & BF & CPW:1b \\
C5116 & 1a<o2 1b<o2L / 2a<1a1,1b2 / 3a*<2a2 & LI/E/U & {-}+/{-}/{-} & BF & CPW:1b \\
C5117 & 1a<o2 1b<o2L / 2a<1a1,1b2L / 3a*<2a1 & LI/I/U & {-}+/+/{-} & BF & CPW:1b,2a \\
C5118 & 1a<o2 1b<o2L / 2a<1a1,1b2L / 3a*<2a2 & LI/I/U & {-}+/+/{-} & BF & CPW:1b,2a \\
C5119 & 1a<o2 1b<o2L / 2a<1a1L,1b1 / 3a*<2a1 & LI/E/U & {-}+/{-}/{-} & BF & CPW:1b \\
C5120 & 1a<o2 1b<o2L / 2a<1a1L,1b1 / 3a*<2a2 & LI/E/U & {-}+/{-}/{-} & BF & CPW:1b \\
C5121 & 1a<o2 1b<o2L / 2a<1a1L,1b2 / 3a*<2a1 & LI/E/U & {-}+/{-}/{-} & BF & CPW:1b \\
C5122 & 1a<o2 1b<o2L / 2a<1a1L,1b2 / 3a*<2a2 & LI/E/U & {-}+/{-}/{-} & BF & CPW:1b \\
C5123 & 1a<o2 1b<o2L / 2a<1a2,1b0L / 3a*<2a1 & LI/D/U & {-}+/+/{-} & B & CPW:1b \\
C5124 & 1a<o2 1b<o2L / 2a<1a2,1b0L / 3a*<2a2 & LI/D/U & {-}+/+/{-} & B & CPW:1b \\
C5125 & 1a<o2 1b<o2L / 2a<1a2,1b1 / 3a*<2a1 & LI/E/U & {-}+/{-}/{-} & BF & CPW:1b \\
C5126 & 1a<o2 1b<o2L / 2a<1a2,1b1 / 3a*<2a2 & LI/E/U & {-}+/{-}/{-} & BF & CPW:1b \\
C5127 & 1a<o2 1b<o2L / 2a<1a2,1b1L / 3a*<2a1 & LI/I/U & {-}+/+/{-} & BF & CPW:1b,2a \\
C5128 & 1a<o2 1b<o2L / 2a<1a2,1b1L / 3a*<2a2 & LI/I/U & {-}+/+/{-} & BF & CPW:1b,2a \\
C5129 & 1a<o2 1b<o2L / 2a<1a2,1b2 / 3a*<2a1 & LI/E/U & {-}+/{-}/{-} & BF & CPW:1b \\
C5130 & 1a<o2 1b<o2L / 2a<1a2,1b2 / 3a*<2a2 & LI/E/U & {-}+/{-}/{-} & BF & CPW:1b \\
C5131 & 1a<o2 1b<o2L / 2a<1a2,1b2L / 3a*<2a1 & LI/I/U & {-}+/+/{-} & BF & CPW:1b,2a \\
C5132 & 1a<o2 1b<o2L / 2a<1a2,1b2L / 3a*<2a2 & LI/I/U & {-}+/+/{-} & BF & CPW:1b,2a \\
C5133 & 1a<o2 1b<o2L / 2a<1a2L,1b1 / 3a*<2a1 & LI/E/U & {-}+/{-}/{-} & BF & CPW:1b \\
C5134 & 1a<o2 1b<o2L / 2a<1a2L,1b1 / 3a*<2a2 & LI/E/U & {-}+/{-}/{-} & BF & CPW:1b \\
C5135 & 1a<o2 1b<o2L / 2a<1a2L,1b2 / 3a*<2a1 & LI/E/U & {-}+/{-}/{-} & BF & CPW:1b \\
C5136 & 1a<o2 1b<o2L / 2a<1a2L,1b2 / 3a*<2a2 & LI/E/U & {-}+/{-}/{-} & BF & CPW:1b \\
C5137 & 1a<o2 1b<o2L / 2a*<1a1,1b1 / 3a<2a1 & LI/U/U & {-}+/{-}/{-} & B & CPW:1b \\
C5138 & 1a<o2 1b<o2L / 2a*<1a1,1b1 / 3a<2a2 & LI/U/U & {-}+/{-}/{-} & B & CPW:1b \\
C5139 & 1a<o2 1b<o2L / 2a*<1a1,1b2 / 3a<2a1 & LI/U/U & {-}+/{-}/{-} & B & CPW:1b \\
C5140 & 1a<o2 1b<o2L / 2a*<1a1,1b2 / 3a<2a2 & LI/U/U & {-}+/{-}/{-} & B & CPW:1b \\
C5141 & 1a<o2 1b<o2L / 2a*<1a2,1b1 / 3a<2a1 & LI/U/U & {-}+/{-}/{-} & B & CPW:1b \\
C5142 & 1a<o2 1b<o2L / 2a*<1a2,1b1 / 3a<2a2 & LI/U/U & {-}+/{-}/{-} & B & CPW:1b \\
C5143 & 1a<o2 1b<o2L / 2a*<1a2,1b2 / 3a<2a1 & LI/U/U & {-}+/{-}/{-} & B & CPW:1b \\
C5144 & 1a<o2 1b<o2L / 2a*<1a2,1b2 / 3a<2a2 & LI/U/U & {-}+/{-}/{-} & B & CPW:1b \\
C5145 & 1a<o2 1b*<o1 / 2a<1a0L,1b1 / 3a<2a0L & IU/U/N & {-}{-}/{-}/{-} & N & W:1a \\
C5146 & 1a<o2 1b*<o1 / 2a<1a0L,1b1 / 3a<2a1 & IU/U/U & {-}{-}/{-}/{-} & N & W:1a \\
C5147 & 1a<o2 1b*<o1 / 2a<1a0L,1b1 / 3a<2a1L & IU/U/U & {-}{-}/{-}/{-} & N & W:1a \\
C5148 & 1a<o2 1b*<o1 / 2a<1a0L,1b1 / 3a<2a2 & IU/U/U & {-}{-}/{-}/{-} & N & W:1a \\
C5149 & 1a<o2 1b*<o1 / 2a<1a0L,1b1 / 3a<2a2L & IU/U/U & {-}{-}/{-}/{-} & N & W:1a \\
C5150 & 1a<o2 1b*<o1 / 2a<1a0L,1b2 / 3a<2a0L & IU/U/N & {-}{-}/{-}/{-} & N & W:1a \\
C5151 & 1a<o2 1b*<o1 / 2a<1a0L,1b2 / 3a<2a1 & IU/U/U & {-}{-}/{-}/{-} & N & W:1a \\
C5152 & 1a<o2 1b*<o1 / 2a<1a0L,1b2 / 3a<2a1L & IU/U/U & {-}{-}/{-}/{-} & N & W:1a \\
C5153 & 1a<o2 1b*<o1 / 2a<1a0L,1b2 / 3a<2a2 & IU/U/U & {-}{-}/{-}/{-} & N & W:1a \\
C5154 & 1a<o2 1b*<o1 / 2a<1a0L,1b2 / 3a<2a2L & IU/U/U & {-}{-}/{-}/{-} & N & W:1a \\
C5155 & 1a<o2 1b*<o1 / 2a<1a1,1b1 / 3a<2a0L & IU/I/N & {-}{-}/{-}/{-} & N & W:1a,2a \\
C5156 & 1a<o2 1b*<o1 / 2a<1a1,1b1 / 3a<2a1 & IU/I/I & {-}{-}/{-}/{-} & N & W:1a,2a,3a \\
C5157 & 1a<o2 1b*<o1 / 2a<1a1,1b1 / 3a<2a1L & IU/I/I & {-}{-}/{-}/{-} & N & W:1a,2a,3a \\
C5158 & 1a<o2 1b*<o1 / 2a<1a1,1b1 / 3a<2a2 & IU/I/I & {-}{-}/{-}/{-} & N & W:1a,2a,3a \\
C5159 & 1a<o2 1b*<o1 / 2a<1a1,1b1 / 3a<2a2L & IU/I/I & {-}{-}/{-}/{-} & N & W:1a,2a,3a \\
C5160 & 1a<o2 1b*<o1 / 2a<1a1,1b2 / 3a<2a0L & IU/I/N & {-}{-}/{-}/{-} & N & W:1a,2a \\
C5161 & 1a<o2 1b*<o1 / 2a<1a1,1b2 / 3a<2a1 & IU/I/I & {-}{-}/{-}/{-} & N & W:1a,2a,3a \\
C5162 & 1a<o2 1b*<o1 / 2a<1a1,1b2 / 3a<2a1L & IU/I/I & {-}{-}/{-}/{-} & N & W:1a,2a,3a \\
C5163 & 1a<o2 1b*<o1 / 2a<1a1,1b2 / 3a<2a2 & IU/I/I & {-}{-}/{-}/{-} & N & W:1a,2a,3a \\
C5164 & 1a<o2 1b*<o1 / 2a<1a1,1b2 / 3a<2a2L & IU/I/I & {-}{-}/{-}/{-} & N & W:1a,2a,3a \\
C5165 & 1a<o2 1b*<o1 / 2a<1a1L,1b1 / 3a<2a0L & IU/I/N & {-}{-}/{-}/{-} & N & W:1a,2a \\
C5166 & 1a<o2 1b*<o1 / 2a<1a1L,1b1 / 3a<2a1 & IU/I/I & {-}{-}/{-}/{-} & N & W:1a,2a,3a \\
C5167 & 1a<o2 1b*<o1 / 2a<1a1L,1b1 / 3a<2a1L & IU/I/I & {-}{-}/{-}/{-} & N & W:1a,2a,3a \\
C5168 & 1a<o2 1b*<o1 / 2a<1a1L,1b1 / 3a<2a2 & IU/I/I & {-}{-}/{-}/{-} & N & W:1a,2a,3a \\
C5169 & 1a<o2 1b*<o1 / 2a<1a1L,1b1 / 3a<2a2L & IU/I/I & {-}{-}/{-}/{-} & N & W:1a,2a,3a \\
C5170 & 1a<o2 1b*<o1 / 2a<1a1L,1b2 / 3a<2a0L & IU/I/N & {-}{-}/{-}/{-} & N & W:1a,2a \\
C5171 & 1a<o2 1b*<o1 / 2a<1a1L,1b2 / 3a<2a1 & IU/I/I & {-}{-}/{-}/{-} & N & W:1a,2a,3a \\
C5172 & 1a<o2 1b*<o1 / 2a<1a1L,1b2 / 3a<2a1L & IU/I/I & {-}{-}/{-}/{-} & N & W:1a,2a,3a \\
C5173 & 1a<o2 1b*<o1 / 2a<1a1L,1b2 / 3a<2a2 & IU/I/I & {-}{-}/{-}/{-} & N & W:1a,2a,3a \\
C5174 & 1a<o2 1b*<o1 / 2a<1a1L,1b2 / 3a<2a2L & IU/I/I & {-}{-}/{-}/{-} & N & W:1a,2a,3a \\
C5175 & 1a<o2 1b*<o1 / 2a<1a2,1b1 / 3a<2a0L & IU/I/N & {-}{-}/{-}/{-} & N & W:1a,2a \\
C5176 & 1a<o2 1b*<o1 / 2a<1a2,1b1 / 3a<2a1 & IU/I/I & {-}{-}/{-}/{-} & N & W:1a,2a,3a \\
C5177 & 1a<o2 1b*<o1 / 2a<1a2,1b1 / 3a<2a1L & IU/I/I & {-}{-}/{-}/{-} & N & W:1a,2a,3a \\
C5178 & 1a<o2 1b*<o1 / 2a<1a2,1b1 / 3a<2a2 & IU/I/I & {-}{-}/{-}/{-} & N & W:1a,2a,3a \\
C5179 & 1a<o2 1b*<o1 / 2a<1a2,1b1 / 3a<2a2L & IU/I/I & {-}{-}/{-}/{-} & N & W:1a,2a,3a \\
C5180 & 1a<o2 1b*<o1 / 2a<1a2,1b2 / 3a<2a0L & IU/I/N & {-}{-}/{-}/{-} & N & W:1a,2a \\
C5181 & 1a<o2 1b*<o1 / 2a<1a2,1b2 / 3a<2a1 & IU/I/I & {-}{-}/{-}/{-} & N & W:1a,2a,3a \\
C5182 & 1a<o2 1b*<o1 / 2a<1a2,1b2 / 3a<2a1L & IU/I/I & {-}{-}/{-}/{-} & N & W:1a,2a,3a \\
C5183 & 1a<o2 1b*<o1 / 2a<1a2,1b2 / 3a<2a2 & IU/I/I & {-}{-}/{-}/{-} & N & W:1a,2a,3a \\
C5184 & 1a<o2 1b*<o1 / 2a<1a2,1b2 / 3a<2a2L & IU/I/I & {-}{-}/{-}/{-} & N & W:1a,2a,3a \\
C5185 & 1a<o2 1b*<o1 / 2a<1a2L,1b1 / 3a<2a0L & IU/I/N & {-}{-}/{-}/{-} & N & W:1a,2a \\
C5186 & 1a<o2 1b*<o1 / 2a<1a2L,1b1 / 3a<2a1 & IU/I/I & {-}{-}/{-}/{-} & N & W:1a,2a,3a \\
C5187 & 1a<o2 1b*<o1 / 2a<1a2L,1b1 / 3a<2a1L & IU/I/I & {-}{-}/{-}/{-} & N & W:1a,2a,3a \\
C5188 & 1a<o2 1b*<o1 / 2a<1a2L,1b1 / 3a<2a2 & IU/I/I & {-}{-}/{-}/{-} & N & W:1a,2a,3a \\
C5189 & 1a<o2 1b*<o1 / 2a<1a2L,1b1 / 3a<2a2L & IU/I/I & {-}{-}/{-}/{-} & N & W:1a,2a,3a \\
C5190 & 1a<o2 1b*<o1 / 2a<1a2L,1b2 / 3a<2a0L & IU/I/N & {-}{-}/{-}/{-} & N & W:1a,2a \\
C5191 & 1a<o2 1b*<o1 / 2a<1a2L,1b2 / 3a<2a1 & IU/I/I & {-}{-}/{-}/{-} & N & W:1a,2a,3a \\
C5192 & 1a<o2 1b*<o1 / 2a<1a2L,1b2 / 3a<2a1L & IU/I/I & {-}{-}/{-}/{-} & N & W:1a,2a,3a \\
C5193 & 1a<o2 1b*<o1 / 2a<1a2L,1b2 / 3a<2a2 & IU/I/I & {-}{-}/{-}/{-} & N & W:1a,2a,3a \\
C5194 & 1a<o2 1b*<o1 / 2a<1a2L,1b2 / 3a<2a2L & IU/I/I & {-}{-}/{-}/{-} & N & W:1a,2a,3a \\
C5195 & 1a<o2 1b*<o2 / 2a<1a0L,1b1 / 3a<2a0L & IU/U/N & {-}{-}/{-}/{-} & N & W:1a \\
C5196 & 1a<o2 1b*<o2 / 2a<1a0L,1b1 / 3a<2a1 & IU/U/U & {-}{-}/{-}/{-} & N & W:1a \\
C5197 & 1a<o2 1b*<o2 / 2a<1a0L,1b1 / 3a<2a1L & IU/U/U & {-}{-}/{-}/{-} & N & W:1a \\
C5198 & 1a<o2 1b*<o2 / 2a<1a0L,1b1 / 3a<2a2 & IU/U/U & {-}{-}/{-}/{-} & N & W:1a \\
C5199 & 1a<o2 1b*<o2 / 2a<1a0L,1b1 / 3a<2a2L & IU/U/U & {-}{-}/{-}/{-} & N & W:1a \\
C5200 & 1a<o2 1b*<o2 / 2a<1a0L,1b2 / 3a<2a0L & IU/U/N & {-}{-}/{-}/{-} & N & W:1a \\
C5201 & 1a<o2 1b*<o2 / 2a<1a0L,1b2 / 3a<2a1 & IU/U/U & {-}{-}/{-}/{-} & N & W:1a \\
C5202 & 1a<o2 1b*<o2 / 2a<1a0L,1b2 / 3a<2a1L & IU/U/U & {-}{-}/{-}/{-} & N & W:1a \\
C5203 & 1a<o2 1b*<o2 / 2a<1a0L,1b2 / 3a<2a2 & IU/U/U & {-}{-}/{-}/{-} & N & W:1a \\
C5204 & 1a<o2 1b*<o2 / 2a<1a0L,1b2 / 3a<2a2L & IU/U/U & {-}{-}/{-}/{-} & N & W:1a \\
C5205 & 1a<o2 1b*<o2 / 2a<1a1,1b1 / 3a<2a0L & IU/I/N & {-}{-}/{-}/{-} & N & W:1a,2a \\
C5206 & 1a<o2 1b*<o2 / 2a<1a1,1b1 / 3a<2a1 & IU/I/I & {-}{-}/{-}/{-} & N & W:1a,2a,3a \\
C5207 & 1a<o2 1b*<o2 / 2a<1a1,1b1 / 3a<2a1L & IU/I/I & {-}{-}/{-}/{-} & N & W:1a,2a,3a \\
C5208 & 1a<o2 1b*<o2 / 2a<1a1,1b1 / 3a<2a2 & IU/I/I & {-}{-}/{-}/{-} & N & W:1a,2a,3a \\
C5209 & 1a<o2 1b*<o2 / 2a<1a1,1b1 / 3a<2a2L & IU/I/I & {-}{-}/{-}/{-} & N & W:1a,2a,3a \\
C5210 & 1a<o2 1b*<o2 / 2a<1a1,1b2 / 3a<2a0L & IU/I/N & {-}{-}/{-}/{-} & N & W:1a,2a \\
C5211 & 1a<o2 1b*<o2 / 2a<1a1,1b2 / 3a<2a1 & IU/I/I & {-}{-}/{-}/{-} & N & W:1a,2a,3a \\
C5212 & 1a<o2 1b*<o2 / 2a<1a1,1b2 / 3a<2a1L & IU/I/I & {-}{-}/{-}/{-} & N & W:1a,2a,3a \\
C5213 & 1a<o2 1b*<o2 / 2a<1a1,1b2 / 3a<2a2 & IU/I/I & {-}{-}/{-}/{-} & N & W:1a,2a,3a \\
C5214 & 1a<o2 1b*<o2 / 2a<1a1,1b2 / 3a<2a2L & IU/I/I & {-}{-}/{-}/{-} & N & W:1a,2a,3a \\
C5215 & 1a<o2 1b*<o2 / 2a<1a1L,1b1 / 3a<2a0L & IU/I/N & {-}{-}/{-}/{-} & N & W:1a,2a \\
C5216 & 1a<o2 1b*<o2 / 2a<1a1L,1b1 / 3a<2a1 & IU/I/I & {-}{-}/{-}/{-} & N & W:1a,2a,3a \\
C5217 & 1a<o2 1b*<o2 / 2a<1a1L,1b1 / 3a<2a1L & IU/I/I & {-}{-}/{-}/{-} & N & W:1a,2a,3a \\
C5218 & 1a<o2 1b*<o2 / 2a<1a1L,1b1 / 3a<2a2 & IU/I/I & {-}{-}/{-}/{-} & N & W:1a,2a,3a \\
C5219 & 1a<o2 1b*<o2 / 2a<1a1L,1b1 / 3a<2a2L & IU/I/I & {-}{-}/{-}/{-} & N & W:1a,2a,3a \\
C5220 & 1a<o2 1b*<o2 / 2a<1a1L,1b2 / 3a<2a0L & IU/I/N & {-}{-}/{-}/{-} & N & W:1a,2a \\
C5221 & 1a<o2 1b*<o2 / 2a<1a1L,1b2 / 3a<2a1 & IU/I/I & {-}{-}/{-}/{-} & N & W:1a,2a,3a \\
C5222 & 1a<o2 1b*<o2 / 2a<1a1L,1b2 / 3a<2a1L & IU/I/I & {-}{-}/{-}/{-} & N & W:1a,2a,3a \\
C5223 & 1a<o2 1b*<o2 / 2a<1a1L,1b2 / 3a<2a2 & IU/I/I & {-}{-}/{-}/{-} & N & W:1a,2a,3a \\
C5224 & 1a<o2 1b*<o2 / 2a<1a1L,1b2 / 3a<2a2L & IU/I/I & {-}{-}/{-}/{-} & N & W:1a,2a,3a \\
C5225 & 1a<o2 1b*<o2 / 2a<1a2,1b1 / 3a<2a0L & IU/I/N & {-}{-}/{-}/{-} & N & W:1a,2a \\
C5226 & 1a<o2 1b*<o2 / 2a<1a2,1b1 / 3a<2a1 & IU/I/I & {-}{-}/{-}/{-} & N & W:1a,2a,3a \\
C5227 & 1a<o2 1b*<o2 / 2a<1a2,1b1 / 3a<2a1L & IU/I/I & {-}{-}/{-}/{-} & N & W:1a,2a,3a \\
C5228 & 1a<o2 1b*<o2 / 2a<1a2,1b1 / 3a<2a2 & IU/I/I & {-}{-}/{-}/{-} & N & W:1a,2a,3a \\
C5229 & 1a<o2 1b*<o2 / 2a<1a2,1b1 / 3a<2a2L & IU/I/I & {-}{-}/{-}/{-} & N & W:1a,2a,3a \\
C5230 & 1a<o2 1b*<o2 / 2a<1a2,1b2 / 3a<2a0L & IU/I/N & {-}{-}/{-}/{-} & N & W:1a,2a \\
C5231 & 1a<o2 1b*<o2 / 2a<1a2,1b2 / 3a<2a1 & IU/I/I & {-}{-}/{-}/{-} & N & W:1a,2a,3a \\
C5232 & 1a<o2 1b*<o2 / 2a<1a2,1b2 / 3a<2a1L & IU/I/I & {-}{-}/{-}/{-} & N & W:1a,2a,3a \\
C5233 & 1a<o2 1b*<o2 / 2a<1a2,1b2 / 3a<2a2 & IU/I/I & {-}{-}/{-}/{-} & N & W:1a,2a,3a \\
C5234 & 1a<o2 1b*<o2 / 2a<1a2,1b2 / 3a<2a2L & IU/I/I & {-}{-}/{-}/{-} & N & W:1a,2a,3a \\
C5235 & 1a<o2 1b*<o2 / 2a<1a2L,1b1 / 3a<2a0L & IU/I/N & {-}{-}/{-}/{-} & N & W:1a,2a \\
C5236 & 1a<o2 1b*<o2 / 2a<1a2L,1b1 / 3a<2a1 & IU/I/I & {-}{-}/{-}/{-} & N & W:1a,2a,3a \\
C5237 & 1a<o2 1b*<o2 / 2a<1a2L,1b1 / 3a<2a1L & IU/I/I & {-}{-}/{-}/{-} & N & W:1a,2a,3a \\
C5238 & 1a<o2 1b*<o2 / 2a<1a2L,1b1 / 3a<2a2 & IU/I/I & {-}{-}/{-}/{-} & N & W:1a,2a,3a \\
C5239 & 1a<o2 1b*<o2 / 2a<1a2L,1b1 / 3a<2a2L & IU/I/I & {-}{-}/{-}/{-} & N & W:1a,2a,3a \\
C5240 & 1a<o2 1b*<o2 / 2a<1a2L,1b2 / 3a<2a0L & IU/I/N & {-}{-}/{-}/{-} & N & W:1a,2a \\
C5241 & 1a<o2 1b*<o2 / 2a<1a2L,1b2 / 3a<2a1 & IU/I/I & {-}{-}/{-}/{-} & N & W:1a,2a,3a \\
C5242 & 1a<o2 1b*<o2 / 2a<1a2L,1b2 / 3a<2a1L & IU/I/I & {-}{-}/{-}/{-} & N & W:1a,2a,3a \\
C5243 & 1a<o2 1b*<o2 / 2a<1a2L,1b2 / 3a<2a2 & IU/I/I & {-}{-}/{-}/{-} & N & W:1a,2a,3a \\
C5244 & 1a<o2 1b*<o2 / 2a<1a2L,1b2 / 3a<2a2L & IU/I/I & {-}{-}/{-}/{-} & N & W:1a,2a,3a \\
C5245 & 1a<o2L / 2a<1a0L 2b<1a1 / 3a*<2a1,2b1 & I/NI/U & +/+{-}/{-} & N & W:3a \\
C5246 & 1a<o2L / 2a<1a0L 2b<1a1 / 3a*<2a1,2b2 & I/NI/U & +/+{-}/{-} & N & W:3a \\
C5247 & 1a<o2L / 2a<1a0L 2b<1a1 / 3a*<2a2,2b1 & I/NI/U & +/+{-}/{-} & N & W:3a \\
C5248 & 1a<o2L / 2a<1a0L 2b<1a1 / 3a*<2a2,2b2 & I/NI/U & +/+{-}/{-} & N & W:3a \\
C5249 & 1a<o2L / 2a<1a0L 2b<1a2 / 3a*<2a1,2b1 & I/NI/U & +/+{-}/{-} & N & W:3a \\
C5250 & 1a<o2L / 2a<1a0L 2b<1a2 / 3a*<2a1,2b2 & I/NI/U & +/+{-}/{-} & N & W:3a \\
C5251 & 1a<o2L / 2a<1a0L 2b<1a2 / 3a*<2a2,2b1 & I/NI/U & +/+{-}/{-} & N & W:3a \\
C5252 & 1a<o2L / 2a<1a0L 2b<1a2 / 3a*<2a2,2b2 & I/NI/U & +/+{-}/{-} & N & W:3a \\
C5253 & 1a<o2L / 2a<1a0L 2b*<1a1 / 3a<2a0L,2b1 & I/NU/U & +/+{-}/+ & N & W:2b,3a \\
C5254 & 1a<o2L / 2a<1a0L 2b*<1a1 / 3a<2a0L,2b2 & I/NU/U & +/+{-}/+ & N & W:2b,3a \\
C5255 & 1a<o2L / 2a<1a0L 2b*<1a1 / 3a<2a1,2b1 & I/NU/U & +/+{-}/{-} & N & W:2b,3a \\
C5256 & 1a<o2L / 2a<1a0L 2b*<1a1 / 3a<2a1,2b2 & I/NU/U & +/+{-}/{-} & N & W:2b,3a \\
C5257 & 1a<o2L / 2a<1a0L 2b*<1a1 / 3a<2a1L,2b1 & I/NU/U & +/+{-}/+ & N & W:2b,3a \\
C5258 & 1a<o2L / 2a<1a0L 2b*<1a1 / 3a<2a1L,2b2 & I/NU/U & +/+{-}/+ & N & W:2b,3a \\
C5259 & 1a<o2L / 2a<1a0L 2b*<1a1 / 3a<2a2,2b1 & I/NU/U & +/+{-}/{-} & N & W:2b,3a \\
C5260 & 1a<o2L / 2a<1a0L 2b*<1a1 / 3a<2a2,2b2 & I/NU/U & +/+{-}/{-} & N & W:2b,3a \\
C5261 & 1a<o2L / 2a<1a0L 2b*<1a1 / 3a<2a2L,2b1 & I/NU/U & +/+{-}/+ & N & W:2b,3a \\
C5262 & 1a<o2L / 2a<1a0L 2b*<1a1 / 3a<2a2L,2b2 & I/NU/U & +/+{-}/+ & N & W:2b,3a \\
C5263 & 1a<o2L / 2a<1a0L 2b*<1a2 / 3a<2a0L,2b1 & I/NU/U & +/+{-}/+ & N & W:2b,3a \\
C5264 & 1a<o2L / 2a<1a0L 2b*<1a2 / 3a<2a0L,2b2 & I/NU/U & +/+{-}/+ & N & W:2b,3a \\
C5265 & 1a<o2L / 2a<1a0L 2b*<1a2 / 3a<2a1,2b1 & I/NU/U & +/+{-}/{-} & N & W:2b,3a \\
C5266 & 1a<o2L / 2a<1a0L 2b*<1a2 / 3a<2a1,2b2 & I/NU/U & +/+{-}/{-} & N & W:2b,3a \\
C5267 & 1a<o2L / 2a<1a0L 2b*<1a2 / 3a<2a1L,2b1 & I/NU/U & +/+{-}/+ & N & W:2b,3a \\
C5268 & 1a<o2L / 2a<1a0L 2b*<1a2 / 3a<2a1L,2b2 & I/NU/U & +/+{-}/+ & N & W:2b,3a \\
C5269 & 1a<o2L / 2a<1a0L 2b*<1a2 / 3a<2a2,2b1 & I/NU/U & +/+{-}/{-} & N & W:2b,3a \\
C5270 & 1a<o2L / 2a<1a0L 2b*<1a2 / 3a<2a2,2b2 & I/NU/U & +/+{-}/{-} & N & W:2b,3a \\
C5271 & 1a<o2L / 2a<1a0L 2b*<1a2 / 3a<2a2L,2b1 & I/NU/U & +/+{-}/+ & N & W:2b,3a \\
C5272 & 1a<o2L / 2a<1a0L 2b*<1a2 / 3a<2a2L,2b2 & I/NU/U & +/+{-}/+ & N & W:2b,3a \\
C5273 & 1a<o2L / 2a<1a1 2b<1a1 / 3a*<2a1,2b1 & I/EE/U & +/{-}{-}/{-} & E & W:1a \\
C5274 & 1a<o2L / 2a<1a1 2b<1a1 / 3a*<2a1,2b2 & I/EE/U & +/{-}{-}/{-} & E & W:1a \\
C5275 & 1a<o2L / 2a<1a1 2b<1a1 / 3a*<2a2,2b2 & I/EE/U & +/{-}{-}/{-} & E & W:1a \\
C5276 & 1a<o2L / 2a<1a1 2b<1a1L / 3a*<2a1,2b1 & I/LI/U & +/{-}+/{-} & B & CP:2b W:1a,2b \\
C5277 & 1a<o2L / 2a<1a1 2b<1a1L / 3a*<2a1,2b2 & I/LI/U & +/{-}+/{-} & B & CP:2b W:1a,2b \\
C5278 & 1a<o2L / 2a<1a1 2b<1a1L / 3a*<2a2,2b1 & I/LI/U & +/{-}+/{-} & B & CP:2b W:1a,2b \\
C5279 & 1a<o2L / 2a<1a1 2b<1a1L / 3a*<2a2,2b2 & I/LI/U & +/{-}+/{-} & B & CP:2b W:1a,2b \\
C5280 & 1a<o2L / 2a<1a1 2b<1a2 / 3a*<2a1,2b1 & I/EE/U & +/{-}{-}/{-} & E & C:2b W:1a \\
C5281 & 1a<o2L / 2a<1a1 2b<1a2 / 3a*<2a1,2b2 & I/EE/U & +/{-}{-}/{-} & E & C:2b W:1a \\
C5282 & 1a<o2L / 2a<1a1 2b<1a2 / 3a*<2a2,2b1 & I/EE/U & +/{-}{-}/{-} & E & C:2b W:1a \\
C5283 & 1a<o2L / 2a<1a1 2b<1a2 / 3a*<2a2,2b2 & I/EE/U & +/{-}{-}/{-} & E & C:2b W:1a \\
C5284 & 1a<o2L / 2a<1a1 2b<1a2L / 3a*<2a1,2b1 & I/LI/U & +/{-}+/{-} & B & P:2b W:1a,2b \\
C5285 & 1a<o2L / 2a<1a1 2b<1a2L / 3a*<2a1,2b2 & I/LI/U & +/{-}+/{-} & B & P:2b W:1a,2b \\
C5286 & 1a<o2L / 2a<1a1 2b<1a2L / 3a*<2a2,2b1 & I/LI/U & +/{-}+/{-} & B & P:2b W:1a,2b \\
C5287 & 1a<o2L / 2a<1a1 2b<1a2L / 3a*<2a2,2b2 & I/LI/U & +/{-}+/{-} & B & P:2b W:1a,2b \\
C5288 & 1a<o2L / 2a<1a1 2b*<1a1 / 3a<2a0L,2b1 & I/IU/U & +/{-}{-}/{-} & N & W:1a,2a \\
C5289 & 1a<o2L / 2a<1a1 2b*<1a1 / 3a<2a0L,2b2 & I/IU/U & +/{-}{-}/{-} & N & W:1a,2a \\
C5290 & 1a<o2L / 2a<1a1 2b*<1a1 / 3a<2a1,2b1 & I/IU/I & +/{-}{-}/{-} & N & W:1a,2a,3a \\
C5291 & 1a<o2L / 2a<1a1 2b*<1a1 / 3a<2a1,2b2 & I/IU/I & +/{-}{-}/{-} & N & W:1a,2a,3a \\
C5292 & 1a<o2L / 2a<1a1 2b*<1a1 / 3a<2a1L,2b1 & I/IU/I & +/{-}{-}/{-} & N & W:1a,2a,3a \\
C5293 & 1a<o2L / 2a<1a1 2b*<1a1 / 3a<2a1L,2b2 & I/IU/I & +/{-}{-}/{-} & N & W:1a,2a,3a \\
C5294 & 1a<o2L / 2a<1a1 2b*<1a1 / 3a<2a2,2b1 & I/IU/I & +/{-}{-}/{-} & N & W:1a,2a,3a \\
C5295 & 1a<o2L / 2a<1a1 2b*<1a1 / 3a<2a2,2b2 & I/IU/I & +/{-}{-}/{-} & N & W:1a,2a,3a \\
C5296 & 1a<o2L / 2a<1a1 2b*<1a1 / 3a<2a2L,2b1 & I/IU/I & +/{-}{-}/{-} & N & W:1a,2a,3a \\
C5297 & 1a<o2L / 2a<1a1 2b*<1a1 / 3a<2a2L,2b2 & I/IU/I & +/{-}{-}/{-} & N & W:1a,2a,3a \\
C5298 & 1a<o2L / 2a<1a1 2b*<1a2 / 3a<2a0L,2b1 & I/IU/U & +/{-}{-}/{-} & N & W:1a,2a \\
C5299 & 1a<o2L / 2a<1a1 2b*<1a2 / 3a<2a0L,2b2 & I/IU/U & +/{-}{-}/{-} & N & W:1a,2a \\
C5300 & 1a<o2L / 2a<1a1 2b*<1a2 / 3a<2a1,2b1 & I/IU/I & +/{-}{-}/{-} & N & W:1a,2a,3a \\
C5301 & 1a<o2L / 2a<1a1 2b*<1a2 / 3a<2a1,2b2 & I/IU/I & +/{-}{-}/{-} & N & W:1a,2a,3a \\
C5302 & 1a<o2L / 2a<1a1 2b*<1a2 / 3a<2a1L,2b1 & I/IU/I & +/{-}{-}/{-} & N & W:1a,2a,3a \\
C5303 & 1a<o2L / 2a<1a1 2b*<1a2 / 3a<2a1L,2b2 & I/IU/I & +/{-}{-}/{-} & N & W:1a,2a,3a \\
C5304 & 1a<o2L / 2a<1a1 2b*<1a2 / 3a<2a2,2b1 & I/IU/I & +/{-}{-}/{-} & N & W:1a,2a,3a \\
C5305 & 1a<o2L / 2a<1a1 2b*<1a2 / 3a<2a2,2b2 & I/IU/I & +/{-}{-}/{-} & N & W:1a,2a,3a \\
C5306 & 1a<o2L / 2a<1a1 2b*<1a2 / 3a<2a2L,2b1 & I/IU/I & +/{-}{-}/{-} & N & W:1a,2a,3a \\
C5307 & 1a<o2L / 2a<1a1 2b*<1a2 / 3a<2a2L,2b2 & I/IU/I & +/{-}{-}/{-} & N & W:1a,2a,3a \\
C5308 & 1a<o2L / 2a<1a1L 2b<1a2 / 3a*<2a1,2b1 & I/IL/U & +/+{-}/{-} & B & C:2a,2b P:2a W:1a,2a \\
C5309 & 1a<o2L / 2a<1a1L 2b<1a2 / 3a*<2a1,2b2 & I/IL/U & +/+{-}/{-} & B & C:2a,2b P:2a W:1a,2a \\
C5310 & 1a<o2L / 2a<1a1L 2b<1a2 / 3a*<2a2,2b1 & I/IL/U & +/+{-}/{-} & B & C:2a,2b P:2a W:1a,2a \\
C5311 & 1a<o2L / 2a<1a1L 2b<1a2 / 3a*<2a2,2b2 & I/IL/U & +/+{-}/{-} & B & C:2a,2b P:2a W:1a,2a \\
C5312 & 1a<o2L / 2a<1a1L 2b*<1a1 / 3a<2a0L,2b1 & I/IU/U & +/+{-}/+ & N & W:1a,2a \\
C5313 & 1a<o2L / 2a<1a1L 2b*<1a1 / 3a<2a0L,2b2 & I/IU/U & +/+{-}/+ & N & W:1a,2a \\
C5314 & 1a<o2L / 2a<1a1L 2b*<1a1 / 3a<2a1,2b1 & I/IU/I & +/+{-}/{-} & N & W:1a,2a,3a \\
C5315 & 1a<o2L / 2a<1a1L 2b*<1a1 / 3a<2a1,2b2 & I/IU/I & +/+{-}/{-} & N & W:1a,2a,3a \\
C5316 & 1a<o2L / 2a<1a1L 2b*<1a1 / 3a<2a1L,2b1 & I/IU/I & +/+{-}/+ & N & W:1a,2a,3a \\
C5317 & 1a<o2L / 2a<1a1L 2b*<1a1 / 3a<2a1L,2b2 & I/IU/I & +/+{-}/+ & N & W:1a,2a,3a \\
C5318 & 1a<o2L / 2a<1a1L 2b*<1a1 / 3a<2a2,2b1 & I/IU/I & +/+{-}/{-} & N & W:1a,2a,3a \\
C5319 & 1a<o2L / 2a<1a1L 2b*<1a1 / 3a<2a2,2b2 & I/IU/I & +/+{-}/{-} & N & W:1a,2a,3a \\
C5320 & 1a<o2L / 2a<1a1L 2b*<1a1 / 3a<2a2L,2b1 & I/IU/I & +/+{-}/+ & N & W:1a,2a,3a \\
C5321 & 1a<o2L / 2a<1a1L 2b*<1a1 / 3a<2a2L,2b2 & I/IU/I & +/+{-}/+ & N & W:1a,2a,3a \\
C5322 & 1a<o2L / 2a<1a1L 2b*<1a2 / 3a<2a0L,2b1 & I/IU/U & +/+{-}/+ & N & W:1a,2a \\
C5323 & 1a<o2L / 2a<1a1L 2b*<1a2 / 3a<2a0L,2b2 & I/IU/U & +/+{-}/+ & N & W:1a,2a \\
C5324 & 1a<o2L / 2a<1a1L 2b*<1a2 / 3a<2a1,2b1 & I/IU/I & +/+{-}/{-} & N & W:1a,2a,3a \\
C5325 & 1a<o2L / 2a<1a1L 2b*<1a2 / 3a<2a1,2b2 & I/IU/I & +/+{-}/{-} & N & W:1a,2a,3a \\
C5326 & 1a<o2L / 2a<1a1L 2b*<1a2 / 3a<2a1L,2b1 & I/IU/I & +/+{-}/+ & N & W:1a,2a,3a \\
C5327 & 1a<o2L / 2a<1a1L 2b*<1a2 / 3a<2a1L,2b2 & I/IU/I & +/+{-}/+ & N & W:1a,2a,3a \\
C5328 & 1a<o2L / 2a<1a1L 2b*<1a2 / 3a<2a2,2b1 & I/IU/I & +/+{-}/{-} & N & W:1a,2a,3a \\
C5329 & 1a<o2L / 2a<1a1L 2b*<1a2 / 3a<2a2,2b2 & I/IU/I & +/+{-}/{-} & N & W:1a,2a,3a \\
C5330 & 1a<o2L / 2a<1a1L 2b*<1a2 / 3a<2a2L,2b1 & I/IU/I & +/+{-}/+ & N & W:1a,2a,3a \\
C5331 & 1a<o2L / 2a<1a1L 2b*<1a2 / 3a<2a2L,2b2 & I/IU/I & +/+{-}/+ & N & W:1a,2a,3a \\
C5332 & 1a<o2L / 2a<1a2 2b<1a2 / 3a*<2a1,2b1 & I/EE/U & +/{-}{-}/{-} & E & W:1a \\
C5333 & 1a<o2L / 2a<1a2 2b<1a2 / 3a*<2a1,2b2 & I/EE/U & +/{-}{-}/{-} & E & W:1a \\
C5334 & 1a<o2L / 2a<1a2 2b<1a2 / 3a*<2a2,2b2 & I/EE/U & +/{-}{-}/{-} & E & W:1a \\
C5335 & 1a<o2L / 2a<1a2 2b<1a2L / 3a*<2a1,2b1 & I/LI/U & +/{-}+/{-} & B & CP:2b W:1a,2b \\
C5336 & 1a<o2L / 2a<1a2 2b<1a2L / 3a*<2a1,2b2 & I/LI/U & +/{-}+/{-} & B & CP:2b W:1a,2b \\
C5337 & 1a<o2L / 2a<1a2 2b<1a2L / 3a*<2a2,2b1 & I/LI/U & +/{-}+/{-} & B & CP:2b W:1a,2b \\
C5338 & 1a<o2L / 2a<1a2 2b<1a2L / 3a*<2a2,2b2 & I/LI/U & +/{-}+/{-} & B & CP:2b W:1a,2b \\
C5339 & 1a<o2L / 2a<1a2 2b*<1a1 / 3a<2a0L,2b1 & I/IU/U & +/{-}{-}/{-} & N & W:1a,2a \\
C5340 & 1a<o2L / 2a<1a2 2b*<1a1 / 3a<2a0L,2b2 & I/IU/U & +/{-}{-}/{-} & N & W:1a,2a \\
C5341 & 1a<o2L / 2a<1a2 2b*<1a1 / 3a<2a1,2b1 & I/IU/I & +/{-}{-}/{-} & N & W:1a,2a,3a \\
C5342 & 1a<o2L / 2a<1a2 2b*<1a1 / 3a<2a1,2b2 & I/IU/I & +/{-}{-}/{-} & N & W:1a,2a,3a \\
C5343 & 1a<o2L / 2a<1a2 2b*<1a1 / 3a<2a1L,2b1 & I/IU/I & +/{-}{-}/{-} & N & W:1a,2a,3a \\
C5344 & 1a<o2L / 2a<1a2 2b*<1a1 / 3a<2a1L,2b2 & I/IU/I & +/{-}{-}/{-} & N & W:1a,2a,3a \\
C5345 & 1a<o2L / 2a<1a2 2b*<1a1 / 3a<2a2,2b1 & I/IU/I & +/{-}{-}/{-} & N & W:1a,2a,3a \\
C5346 & 1a<o2L / 2a<1a2 2b*<1a1 / 3a<2a2,2b2 & I/IU/I & +/{-}{-}/{-} & N & W:1a,2a,3a \\
C5347 & 1a<o2L / 2a<1a2 2b*<1a1 / 3a<2a2L,2b1 & I/IU/I & +/{-}{-}/{-} & N & W:1a,2a,3a \\
C5348 & 1a<o2L / 2a<1a2 2b*<1a1 / 3a<2a2L,2b2 & I/IU/I & +/{-}{-}/{-} & N & W:1a,2a,3a \\
C5349 & 1a<o2L / 2a<1a2 2b*<1a2 / 3a<2a0L,2b1 & I/IU/U & +/{-}{-}/{-} & N & W:1a,2a \\
C5350 & 1a<o2L / 2a<1a2 2b*<1a2 / 3a<2a0L,2b2 & I/IU/U & +/{-}{-}/{-} & N & W:1a,2a \\
C5351 & 1a<o2L / 2a<1a2 2b*<1a2 / 3a<2a1,2b1 & I/IU/I & +/{-}{-}/{-} & N & W:1a,2a,3a \\
C5352 & 1a<o2L / 2a<1a2 2b*<1a2 / 3a<2a1,2b2 & I/IU/I & +/{-}{-}/{-} & N & W:1a,2a,3a \\
C5353 & 1a<o2L / 2a<1a2 2b*<1a2 / 3a<2a1L,2b1 & I/IU/I & +/{-}{-}/{-} & N & W:1a,2a,3a \\
C5354 & 1a<o2L / 2a<1a2 2b*<1a2 / 3a<2a1L,2b2 & I/IU/I & +/{-}{-}/{-} & N & W:1a,2a,3a \\
C5355 & 1a<o2L / 2a<1a2 2b*<1a2 / 3a<2a2,2b1 & I/IU/I & +/{-}{-}/{-} & N & W:1a,2a,3a \\
C5356 & 1a<o2L / 2a<1a2 2b*<1a2 / 3a<2a2,2b2 & I/IU/I & +/{-}{-}/{-} & N & W:1a,2a,3a \\
C5357 & 1a<o2L / 2a<1a2 2b*<1a2 / 3a<2a2L,2b1 & I/IU/I & +/{-}{-}/{-} & N & W:1a,2a,3a \\
C5358 & 1a<o2L / 2a<1a2 2b*<1a2 / 3a<2a2L,2b2 & I/IU/I & +/{-}{-}/{-} & N & W:1a,2a,3a \\
C5359 & 1a<o2L / 2a<1a2L 2b*<1a1 / 3a<2a0L,2b1 & I/IU/U & +/+{-}/+ & N & W:1a,2a \\
C5360 & 1a<o2L / 2a<1a2L 2b*<1a1 / 3a<2a0L,2b2 & I/IU/U & +/+{-}/+ & N & W:1a,2a \\
C5361 & 1a<o2L / 2a<1a2L 2b*<1a1 / 3a<2a1,2b1 & I/IU/I & +/+{-}/{-} & N & W:1a,2a,3a \\
C5362 & 1a<o2L / 2a<1a2L 2b*<1a1 / 3a<2a1,2b2 & I/IU/I & +/+{-}/{-} & N & W:1a,2a,3a \\
C5363 & 1a<o2L / 2a<1a2L 2b*<1a1 / 3a<2a1L,2b1 & I/IU/I & +/+{-}/+ & N & W:1a,2a,3a \\
C5364 & 1a<o2L / 2a<1a2L 2b*<1a1 / 3a<2a1L,2b2 & I/IU/I & +/+{-}/+ & N & W:1a,2a,3a \\
C5365 & 1a<o2L / 2a<1a2L 2b*<1a1 / 3a<2a2,2b1 & I/IU/I & +/+{-}/{-} & N & W:1a,2a,3a \\
C5366 & 1a<o2L / 2a<1a2L 2b*<1a1 / 3a<2a2,2b2 & I/IU/I & +/+{-}/{-} & N & W:1a,2a,3a \\
C5367 & 1a<o2L / 2a<1a2L 2b*<1a1 / 3a<2a2L,2b1 & I/IU/I & +/+{-}/+ & N & W:1a,2a,3a \\
C5368 & 1a<o2L / 2a<1a2L 2b*<1a1 / 3a<2a2L,2b2 & I/IU/I & +/+{-}/+ & N & W:1a,2a,3a \\
C5369 & 1a<o2L / 2a<1a2L 2b*<1a2 / 3a<2a0L,2b1 & I/IU/U & +/+{-}/+ & N & W:1a,2a \\
C5370 & 1a<o2L / 2a<1a2L 2b*<1a2 / 3a<2a0L,2b2 & I/IU/U & +/+{-}/+ & N & W:1a,2a \\
C5371 & 1a<o2L / 2a<1a2L 2b*<1a2 / 3a<2a1,2b1 & I/IU/I & +/+{-}/{-} & N & W:1a,2a,3a \\
C5372 & 1a<o2L / 2a<1a2L 2b*<1a2 / 3a<2a1,2b2 & I/IU/I & +/+{-}/{-} & N & W:1a,2a,3a \\
C5373 & 1a<o2L / 2a<1a2L 2b*<1a2 / 3a<2a1L,2b1 & I/IU/I & +/+{-}/+ & N & W:1a,2a,3a \\
C5374 & 1a<o2L / 2a<1a2L 2b*<1a2 / 3a<2a1L,2b2 & I/IU/I & +/+{-}/+ & N & W:1a,2a,3a \\
C5375 & 1a<o2L / 2a<1a2L 2b*<1a2 / 3a<2a2,2b1 & I/IU/I & +/+{-}/{-} & N & W:1a,2a,3a \\
C5376 & 1a<o2L / 2a<1a2L 2b*<1a2 / 3a<2a2,2b2 & I/IU/I & +/+{-}/{-} & N & W:1a,2a,3a \\
C5377 & 1a<o2L / 2a<1a2L 2b*<1a2 / 3a<2a2L,2b1 & I/IU/I & +/+{-}/+ & N & W:1a,2a,3a \\
C5378 & 1a<o2L / 2a<1a2L 2b*<1a2 / 3a<2a2L,2b2 & I/IU/I & +/+{-}/+ & N & W:1a,2a,3a \\
C5379 & 1a<o2L 1b*<o1 / 2a<1a0L,1b1 / 3a<2a0L & IU/U/N & +{-}/+/+ & N & W:1a \\
C5380 & 1a<o2L 1b*<o1 / 2a<1a0L,1b1 / 3a<2a1 & IU/U/U & +{-}/+/{-} & N & W:1a \\
C5381 & 1a<o2L 1b*<o1 / 2a<1a0L,1b1 / 3a<2a1L & IU/U/U & +{-}/+/+ & N & W:1a \\
C5382 & 1a<o2L 1b*<o1 / 2a<1a0L,1b1 / 3a<2a2 & IU/U/U & +{-}/+/{-} & N & W:1a \\
C5383 & 1a<o2L 1b*<o1 / 2a<1a0L,1b1 / 3a<2a2L & IU/U/U & +{-}/+/+ & N & W:1a \\
C5384 & 1a<o2L 1b*<o1 / 2a<1a0L,1b2 / 3a<2a0L & IU/U/N & +{-}/+/+ & N & W:1a \\
C5385 & 1a<o2L 1b*<o1 / 2a<1a0L,1b2 / 3a<2a1 & IU/U/U & +{-}/+/{-} & N & W:1a \\
C5386 & 1a<o2L 1b*<o1 / 2a<1a0L,1b2 / 3a<2a1L & IU/U/U & +{-}/+/+ & N & W:1a \\
C5387 & 1a<o2L 1b*<o1 / 2a<1a0L,1b2 / 3a<2a2 & IU/U/U & +{-}/+/{-} & N & W:1a \\
C5388 & 1a<o2L 1b*<o1 / 2a<1a0L,1b2 / 3a<2a2L & IU/U/U & +{-}/+/+ & N & W:1a \\
C5389 & 1a<o2L 1b*<o1 / 2a<1a1,1b1 / 3a<2a0L & IU/I/N & +{-}/{-}/{-} & N & W:1a,2a \\
C5390 & 1a<o2L 1b*<o1 / 2a<1a1,1b1 / 3a<2a1 & IU/I/I & +{-}/{-}/{-} & N & W:1a,2a,3a \\
C5391 & 1a<o2L 1b*<o1 / 2a<1a1,1b1 / 3a<2a1L & IU/I/I & +{-}/{-}/{-} & N & W:1a,2a,3a \\
C5392 & 1a<o2L 1b*<o1 / 2a<1a1,1b1 / 3a<2a2 & IU/I/I & +{-}/{-}/{-} & N & W:1a,2a,3a \\
C5393 & 1a<o2L 1b*<o1 / 2a<1a1,1b1 / 3a<2a2L & IU/I/I & +{-}/{-}/{-} & N & W:1a,2a,3a \\
C5394 & 1a<o2L 1b*<o1 / 2a<1a1,1b2 / 3a<2a0L & IU/I/N & +{-}/{-}/{-} & N & W:1a,2a \\
C5395 & 1a<o2L 1b*<o1 / 2a<1a1,1b2 / 3a<2a1 & IU/I/I & +{-}/{-}/{-} & N & W:1a,2a,3a \\
C5396 & 1a<o2L 1b*<o1 / 2a<1a1,1b2 / 3a<2a1L & IU/I/I & +{-}/{-}/{-} & N & W:1a,2a,3a \\
C5397 & 1a<o2L 1b*<o1 / 2a<1a1,1b2 / 3a<2a2 & IU/I/I & +{-}/{-}/{-} & N & W:1a,2a,3a \\
C5398 & 1a<o2L 1b*<o1 / 2a<1a1,1b2 / 3a<2a2L & IU/I/I & +{-}/{-}/{-} & N & W:1a,2a,3a \\
C5399 & 1a<o2L 1b*<o1 / 2a<1a1L,1b1 / 3a<2a0L & IU/I/N & +{-}/+/+ & N & W:1a,2a \\
C5400 & 1a<o2L 1b*<o1 / 2a<1a1L,1b1 / 3a<2a1 & IU/I/I & +{-}/+/{-} & N & W:1a,2a,3a \\
C5401 & 1a<o2L 1b*<o1 / 2a<1a1L,1b1 / 3a<2a1L & IU/I/I & +{-}/+/+ & N & W:1a,2a,3a \\
C5402 & 1a<o2L 1b*<o1 / 2a<1a1L,1b1 / 3a<2a2 & IU/I/I & +{-}/+/{-} & N & W:1a,2a,3a \\
C5403 & 1a<o2L 1b*<o1 / 2a<1a1L,1b1 / 3a<2a2L & IU/I/I & +{-}/+/+ & N & W:1a,2a,3a \\
C5404 & 1a<o2L 1b*<o1 / 2a<1a1L,1b2 / 3a<2a0L & IU/I/N & +{-}/+/+ & N & W:1a,2a \\
C5405 & 1a<o2L 1b*<o1 / 2a<1a1L,1b2 / 3a<2a1 & IU/I/I & +{-}/+/{-} & N & W:1a,2a,3a \\
C5406 & 1a<o2L 1b*<o1 / 2a<1a1L,1b2 / 3a<2a1L & IU/I/I & +{-}/+/+ & N & W:1a,2a,3a \\
C5407 & 1a<o2L 1b*<o1 / 2a<1a1L,1b2 / 3a<2a2 & IU/I/I & +{-}/+/{-} & N & W:1a,2a,3a \\
C5408 & 1a<o2L 1b*<o1 / 2a<1a1L,1b2 / 3a<2a2L & IU/I/I & +{-}/+/+ & N & W:1a,2a,3a \\
C5409 & 1a<o2L 1b*<o1 / 2a<1a2,1b1 / 3a<2a0L & IU/I/N & +{-}/{-}/{-} & N & W:1a,2a \\
C5410 & 1a<o2L 1b*<o1 / 2a<1a2,1b1 / 3a<2a1 & IU/I/I & +{-}/{-}/{-} & N & W:1a,2a,3a \\
C5411 & 1a<o2L 1b*<o1 / 2a<1a2,1b1 / 3a<2a1L & IU/I/I & +{-}/{-}/{-} & N & W:1a,2a,3a \\
C5412 & 1a<o2L 1b*<o1 / 2a<1a2,1b1 / 3a<2a2 & IU/I/I & +{-}/{-}/{-} & N & W:1a,2a,3a \\
C5413 & 1a<o2L 1b*<o1 / 2a<1a2,1b1 / 3a<2a2L & IU/I/I & +{-}/{-}/{-} & N & W:1a,2a,3a \\
C5414 & 1a<o2L 1b*<o1 / 2a<1a2,1b2 / 3a<2a0L & IU/I/N & +{-}/{-}/{-} & N & W:1a,2a \\
C5415 & 1a<o2L 1b*<o1 / 2a<1a2,1b2 / 3a<2a1 & IU/I/I & +{-}/{-}/{-} & N & W:1a,2a,3a \\
C5416 & 1a<o2L 1b*<o1 / 2a<1a2,1b2 / 3a<2a1L & IU/I/I & +{-}/{-}/{-} & N & W:1a,2a,3a \\
C5417 & 1a<o2L 1b*<o1 / 2a<1a2,1b2 / 3a<2a2 & IU/I/I & +{-}/{-}/{-} & N & W:1a,2a,3a \\
C5418 & 1a<o2L 1b*<o1 / 2a<1a2,1b2 / 3a<2a2L & IU/I/I & +{-}/{-}/{-} & N & W:1a,2a,3a \\
C5419 & 1a<o2L 1b*<o1 / 2a<1a2L,1b1 / 3a<2a0L & IU/I/N & +{-}/+/+ & N & W:1a,2a \\
C5420 & 1a<o2L 1b*<o1 / 2a<1a2L,1b1 / 3a<2a1 & IU/I/I & +{-}/+/{-} & N & W:1a,2a,3a \\
C5421 & 1a<o2L 1b*<o1 / 2a<1a2L,1b1 / 3a<2a1L & IU/I/I & +{-}/+/+ & N & W:1a,2a,3a \\
C5422 & 1a<o2L 1b*<o1 / 2a<1a2L,1b1 / 3a<2a2 & IU/I/I & +{-}/+/{-} & N & W:1a,2a,3a \\
C5423 & 1a<o2L 1b*<o1 / 2a<1a2L,1b1 / 3a<2a2L & IU/I/I & +{-}/+/+ & N & W:1a,2a,3a \\
C5424 & 1a<o2L 1b*<o1 / 2a<1a2L,1b2 / 3a<2a0L & IU/I/N & +{-}/+/+ & N & W:1a,2a \\
C5425 & 1a<o2L 1b*<o1 / 2a<1a2L,1b2 / 3a<2a1 & IU/I/I & +{-}/+/{-} & N & W:1a,2a,3a \\
C5426 & 1a<o2L 1b*<o1 / 2a<1a2L,1b2 / 3a<2a1L & IU/I/I & +{-}/+/+ & N & W:1a,2a,3a \\
C5427 & 1a<o2L 1b*<o1 / 2a<1a2L,1b2 / 3a<2a2 & IU/I/I & +{-}/+/{-} & N & W:1a,2a,3a \\
C5428 & 1a<o2L 1b*<o1 / 2a<1a2L,1b2 / 3a<2a2L & IU/I/I & +{-}/+/+ & N & W:1a,2a,3a \\
C5429 & 1a<o2L 1b*<o2 / 2a<1a0L,1b1 / 3a<2a0L & IU/U/N & +{-}/+/+ & N & W:1a \\
C5430 & 1a<o2L 1b*<o2 / 2a<1a0L,1b1 / 3a<2a1 & IU/U/U & +{-}/+/{-} & N & W:1a \\
C5431 & 1a<o2L 1b*<o2 / 2a<1a0L,1b1 / 3a<2a1L & IU/U/U & +{-}/+/+ & N & W:1a \\
C5432 & 1a<o2L 1b*<o2 / 2a<1a0L,1b1 / 3a<2a2 & IU/U/U & +{-}/+/{-} & N & W:1a \\
C5433 & 1a<o2L 1b*<o2 / 2a<1a0L,1b1 / 3a<2a2L & IU/U/U & +{-}/+/+ & N & W:1a \\
C5434 & 1a<o2L 1b*<o2 / 2a<1a0L,1b2 / 3a<2a0L & IU/U/N & +{-}/+/+ & N & W:1a \\
C5435 & 1a<o2L 1b*<o2 / 2a<1a0L,1b2 / 3a<2a1 & IU/U/U & +{-}/+/{-} & N & W:1a \\
C5436 & 1a<o2L 1b*<o2 / 2a<1a0L,1b2 / 3a<2a1L & IU/U/U & +{-}/+/+ & N & W:1a \\
C5437 & 1a<o2L 1b*<o2 / 2a<1a0L,1b2 / 3a<2a2 & IU/U/U & +{-}/+/{-} & N & W:1a \\
C5438 & 1a<o2L 1b*<o2 / 2a<1a0L,1b2 / 3a<2a2L & IU/U/U & +{-}/+/+ & N & W:1a \\
C5439 & 1a<o2L 1b*<o2 / 2a<1a1,1b1 / 3a<2a0L & IU/I/N & +{-}/{-}/{-} & N & W:1a,2a \\
C5440 & 1a<o2L 1b*<o2 / 2a<1a1,1b1 / 3a<2a1 & IU/I/I & +{-}/{-}/{-} & N & W:1a,2a,3a \\
C5441 & 1a<o2L 1b*<o2 / 2a<1a1,1b1 / 3a<2a1L & IU/I/I & +{-}/{-}/{-} & N & W:1a,2a,3a \\
C5442 & 1a<o2L 1b*<o2 / 2a<1a1,1b1 / 3a<2a2 & IU/I/I & +{-}/{-}/{-} & N & W:1a,2a,3a \\
C5443 & 1a<o2L 1b*<o2 / 2a<1a1,1b1 / 3a<2a2L & IU/I/I & +{-}/{-}/{-} & N & W:1a,2a,3a \\
C5444 & 1a<o2L 1b*<o2 / 2a<1a1,1b2 / 3a<2a0L & IU/I/N & +{-}/{-}/{-} & N & W:1a,2a \\
C5445 & 1a<o2L 1b*<o2 / 2a<1a1,1b2 / 3a<2a1 & IU/I/I & +{-}/{-}/{-} & N & W:1a,2a,3a \\
C5446 & 1a<o2L 1b*<o2 / 2a<1a1,1b2 / 3a<2a1L & IU/I/I & +{-}/{-}/{-} & N & W:1a,2a,3a \\
C5447 & 1a<o2L 1b*<o2 / 2a<1a1,1b2 / 3a<2a2 & IU/I/I & +{-}/{-}/{-} & N & W:1a,2a,3a \\
C5448 & 1a<o2L 1b*<o2 / 2a<1a1,1b2 / 3a<2a2L & IU/I/I & +{-}/{-}/{-} & N & W:1a,2a,3a \\
C5449 & 1a<o2L 1b*<o2 / 2a<1a1L,1b1 / 3a<2a0L & IU/I/N & +{-}/+/+ & N & W:1a,2a \\
C5450 & 1a<o2L 1b*<o2 / 2a<1a1L,1b1 / 3a<2a1 & IU/I/I & +{-}/+/{-} & N & W:1a,2a,3a \\
C5451 & 1a<o2L 1b*<o2 / 2a<1a1L,1b1 / 3a<2a1L & IU/I/I & +{-}/+/+ & N & W:1a,2a,3a \\
C5452 & 1a<o2L 1b*<o2 / 2a<1a1L,1b1 / 3a<2a2 & IU/I/I & +{-}/+/{-} & N & W:1a,2a,3a \\
C5453 & 1a<o2L 1b*<o2 / 2a<1a1L,1b1 / 3a<2a2L & IU/I/I & +{-}/+/+ & N & W:1a,2a,3a \\
C5454 & 1a<o2L 1b*<o2 / 2a<1a1L,1b2 / 3a<2a0L & IU/I/N & +{-}/+/+ & N & W:1a,2a \\
C5455 & 1a<o2L 1b*<o2 / 2a<1a1L,1b2 / 3a<2a1 & IU/I/I & +{-}/+/{-} & N & W:1a,2a,3a \\
C5456 & 1a<o2L 1b*<o2 / 2a<1a1L,1b2 / 3a<2a1L & IU/I/I & +{-}/+/+ & N & W:1a,2a,3a \\
C5457 & 1a<o2L 1b*<o2 / 2a<1a1L,1b2 / 3a<2a2 & IU/I/I & +{-}/+/{-} & N & W:1a,2a,3a \\
C5458 & 1a<o2L 1b*<o2 / 2a<1a1L,1b2 / 3a<2a2L & IU/I/I & +{-}/+/+ & N & W:1a,2a,3a \\
C5459 & 1a<o2L 1b*<o2 / 2a<1a2,1b1 / 3a<2a0L & IU/I/N & +{-}/{-}/{-} & N & W:1a,2a \\
C5460 & 1a<o2L 1b*<o2 / 2a<1a2,1b1 / 3a<2a1 & IU/I/I & +{-}/{-}/{-} & N & W:1a,2a,3a \\
C5461 & 1a<o2L 1b*<o2 / 2a<1a2,1b1 / 3a<2a1L & IU/I/I & +{-}/{-}/{-} & N & W:1a,2a,3a \\
C5462 & 1a<o2L 1b*<o2 / 2a<1a2,1b1 / 3a<2a2 & IU/I/I & +{-}/{-}/{-} & N & W:1a,2a,3a \\
C5463 & 1a<o2L 1b*<o2 / 2a<1a2,1b1 / 3a<2a2L & IU/I/I & +{-}/{-}/{-} & N & W:1a,2a,3a \\
C5464 & 1a<o2L 1b*<o2 / 2a<1a2,1b2 / 3a<2a0L & IU/I/N & +{-}/{-}/{-} & N & W:1a,2a \\
C5465 & 1a<o2L 1b*<o2 / 2a<1a2,1b2 / 3a<2a1 & IU/I/I & +{-}/{-}/{-} & N & W:1a,2a,3a \\
C5466 & 1a<o2L 1b*<o2 / 2a<1a2,1b2 / 3a<2a1L & IU/I/I & +{-}/{-}/{-} & N & W:1a,2a,3a \\
C5467 & 1a<o2L 1b*<o2 / 2a<1a2,1b2 / 3a<2a2 & IU/I/I & +{-}/{-}/{-} & N & W:1a,2a,3a \\
C5468 & 1a<o2L 1b*<o2 / 2a<1a2,1b2 / 3a<2a2L & IU/I/I & +{-}/{-}/{-} & N & W:1a,2a,3a \\
C5469 & 1a<o2L 1b*<o2 / 2a<1a2L,1b1 / 3a<2a0L & IU/I/N & +{-}/+/+ & N & W:1a,2a \\
C5470 & 1a<o2L 1b*<o2 / 2a<1a2L,1b1 / 3a<2a1 & IU/I/I & +{-}/+/{-} & N & W:1a,2a,3a \\
C5471 & 1a<o2L 1b*<o2 / 2a<1a2L,1b1 / 3a<2a1L & IU/I/I & +{-}/+/+ & N & W:1a,2a,3a \\
C5472 & 1a<o2L 1b*<o2 / 2a<1a2L,1b1 / 3a<2a2 & IU/I/I & +{-}/+/{-} & N & W:1a,2a,3a \\
C5473 & 1a<o2L 1b*<o2 / 2a<1a2L,1b1 / 3a<2a2L & IU/I/I & +{-}/+/+ & N & W:1a,2a,3a \\
C5474 & 1a<o2L 1b*<o2 / 2a<1a2L,1b2 / 3a<2a0L & IU/I/N & +{-}/+/+ & N & W:1a,2a \\
C5475 & 1a<o2L 1b*<o2 / 2a<1a2L,1b2 / 3a<2a1 & IU/I/I & +{-}/+/{-} & N & W:1a,2a,3a \\
C5476 & 1a<o2L 1b*<o2 / 2a<1a2L,1b2 / 3a<2a1L & IU/I/I & +{-}/+/+ & N & W:1a,2a,3a \\
C5477 & 1a<o2L 1b*<o2 / 2a<1a2L,1b2 / 3a<2a2 & IU/I/I & +{-}/+/{-} & N & W:1a,2a,3a \\
C5478 & 1a<o2L 1b*<o2 / 2a<1a2L,1b2 / 3a<2a2L & IU/I/I & +{-}/+/+ & N & W:1a,2a,3a \\
C5479 & 1a*<o1 / 2a<1a1 2b<1a1 / 3a<2a0L,2b1 & U/UU/U & {-}/{-}{-}/{-} & N &  \\
C5480 & 1a*<o1 / 2a<1a1 2b<1a1 / 3a<2a0L,2b2 & U/UU/U & {-}/{-}{-}/{-} & N &  \\
C5481 & 1a*<o1 / 2a<1a1 2b<1a1 / 3a<2a1,2b1 & U/UU/U & {-}/{-}{-}/{-} & N &  \\
C5482 & 1a*<o1 / 2a<1a1 2b<1a1 / 3a<2a1,2b1L & U/UU/U & {-}/{-}{-}/{-} & N &  \\
C5483 & 1a*<o1 / 2a<1a1 2b<1a1 / 3a<2a1,2b2 & U/UU/U & {-}/{-}{-}/{-} & N &  \\
C5484 & 1a*<o1 / 2a<1a1 2b<1a1 / 3a<2a1,2b2L & U/UU/U & {-}/{-}{-}/{-} & N &  \\
C5485 & 1a*<o1 / 2a<1a1 2b<1a1 / 3a<2a1L,2b2 & U/UU/U & {-}/{-}{-}/{-} & N &  \\
C5486 & 1a*<o1 / 2a<1a1 2b<1a1 / 3a<2a2,2b2 & U/UU/U & {-}/{-}{-}/{-} & N &  \\
C5487 & 1a*<o1 / 2a<1a1 2b<1a1 / 3a<2a2,2b2L & U/UU/U & {-}/{-}{-}/{-} & N &  \\
C5488 & 1a*<o1 / 2a<1a1 2b<1a2 / 3a<2a0L,2b1 & U/UU/U & {-}/{-}{-}/{-} & N &  \\
C5489 & 1a*<o1 / 2a<1a1 2b<1a2 / 3a<2a0L,2b2 & U/UU/U & {-}/{-}{-}/{-} & N &  \\
C5490 & 1a*<o1 / 2a<1a1 2b<1a2 / 3a<2a1,2b0L & U/UU/U & {-}/{-}{-}/{-} & N &  \\
C5491 & 1a*<o1 / 2a<1a1 2b<1a2 / 3a<2a1,2b1 & U/UU/U & {-}/{-}{-}/{-} & N &  \\
C5492 & 1a*<o1 / 2a<1a1 2b<1a2 / 3a<2a1,2b1L & U/UU/U & {-}/{-}{-}/{-} & N &  \\
C5493 & 1a*<o1 / 2a<1a1 2b<1a2 / 3a<2a1,2b2 & U/UU/U & {-}/{-}{-}/{-} & N &  \\
C5494 & 1a*<o1 / 2a<1a1 2b<1a2 / 3a<2a1,2b2L & U/UU/U & {-}/{-}{-}/{-} & N &  \\
C5495 & 1a*<o1 / 2a<1a1 2b<1a2 / 3a<2a1L,2b1 & U/UU/U & {-}/{-}{-}/{-} & N &  \\
C5496 & 1a*<o1 / 2a<1a1 2b<1a2 / 3a<2a1L,2b2 & U/UU/U & {-}/{-}{-}/{-} & N &  \\
C5497 & 1a*<o1 / 2a<1a1 2b<1a2 / 3a<2a2,2b0L & U/UU/U & {-}/{-}{-}/{-} & N &  \\
C5498 & 1a*<o1 / 2a<1a1 2b<1a2 / 3a<2a2,2b1 & U/UU/U & {-}/{-}{-}/{-} & N &  \\
C5499 & 1a*<o1 / 2a<1a1 2b<1a2 / 3a<2a2,2b1L & U/UU/U & {-}/{-}{-}/{-} & N &  \\
C5500 & 1a*<o1 / 2a<1a1 2b<1a2 / 3a<2a2,2b2 & U/UU/U & {-}/{-}{-}/{-} & N &  \\
C5501 & 1a*<o1 / 2a<1a1 2b<1a2 / 3a<2a2,2b2L & U/UU/U & {-}/{-}{-}/{-} & N &  \\
C5502 & 1a*<o1 / 2a<1a1 2b<1a2 / 3a<2a2L,2b1 & U/UU/U & {-}/{-}{-}/{-} & N &  \\
C5503 & 1a*<o1 / 2a<1a1 2b<1a2 / 3a<2a2L,2b2 & U/UU/U & {-}/{-}{-}/{-} & N &  \\
C5504 & 1a*<o1 / 2a<1a2 2b<1a2 / 3a<2a0L,2b1 & U/UU/U & {-}/{-}{-}/{-} & N &  \\
C5505 & 1a*<o1 / 2a<1a2 2b<1a2 / 3a<2a0L,2b2 & U/UU/U & {-}/{-}{-}/{-} & N &  \\
C5506 & 1a*<o1 / 2a<1a2 2b<1a2 / 3a<2a1,2b1 & U/UU/U & {-}/{-}{-}/{-} & N &  \\
C5507 & 1a*<o1 / 2a<1a2 2b<1a2 / 3a<2a1,2b1L & U/UU/U & {-}/{-}{-}/{-} & N &  \\
C5508 & 1a*<o1 / 2a<1a2 2b<1a2 / 3a<2a1,2b2 & U/UU/U & {-}/{-}{-}/{-} & N &  \\
C5509 & 1a*<o1 / 2a<1a2 2b<1a2 / 3a<2a1,2b2L & U/UU/U & {-}/{-}{-}/{-} & N &  \\
C5510 & 1a*<o1 / 2a<1a2 2b<1a2 / 3a<2a1L,2b2 & U/UU/U & {-}/{-}{-}/{-} & N &  \\
C5511 & 1a*<o1 / 2a<1a2 2b<1a2 / 3a<2a2,2b2 & U/UU/U & {-}/{-}{-}/{-} & N &  \\
C5512 & 1a*<o1 / 2a<1a2 2b<1a2 / 3a<2a2,2b2L & U/UU/U & {-}/{-}{-}/{-} & N &  \\
C5513 & 1a*<o2 / 2a<1a1 2b<1a1 / 3a<2a0L,2b1 & U/UU/U & {-}/{-}{-}/{-} & N &  \\
C5514 & 1a*<o2 / 2a<1a1 2b<1a1 / 3a<2a0L,2b2 & U/UU/U & {-}/{-}{-}/{-} & N &  \\
C5515 & 1a*<o2 / 2a<1a1 2b<1a1 / 3a<2a1,2b1 & U/UU/U & {-}/{-}{-}/{-} & N &  \\
C5516 & 1a*<o2 / 2a<1a1 2b<1a1 / 3a<2a1,2b1L & U/UU/U & {-}/{-}{-}/{-} & N &  \\
C5517 & 1a*<o2 / 2a<1a1 2b<1a1 / 3a<2a1,2b2 & U/UU/U & {-}/{-}{-}/{-} & N &  \\
C5518 & 1a*<o2 / 2a<1a1 2b<1a1 / 3a<2a1,2b2L & U/UU/U & {-}/{-}{-}/{-} & N &  \\
C5519 & 1a*<o2 / 2a<1a1 2b<1a1 / 3a<2a1L,2b2 & U/UU/U & {-}/{-}{-}/{-} & N &  \\
C5520 & 1a*<o2 / 2a<1a1 2b<1a1 / 3a<2a2,2b2 & U/UU/U & {-}/{-}{-}/{-} & N &  \\
C5521 & 1a*<o2 / 2a<1a1 2b<1a1 / 3a<2a2,2b2L & U/UU/U & {-}/{-}{-}/{-} & N &  \\
C5522 & 1a*<o2 / 2a<1a1 2b<1a2 / 3a<2a0L,2b1 & U/UU/U & {-}/{-}{-}/{-} & N &  \\
C5523 & 1a*<o2 / 2a<1a1 2b<1a2 / 3a<2a0L,2b2 & U/UU/U & {-}/{-}{-}/{-} & N &  \\
C5524 & 1a*<o2 / 2a<1a1 2b<1a2 / 3a<2a1,2b0L & U/UU/U & {-}/{-}{-}/{-} & N &  \\
C5525 & 1a*<o2 / 2a<1a1 2b<1a2 / 3a<2a1,2b1 & U/UU/U & {-}/{-}{-}/{-} & N &  \\
C5526 & 1a*<o2 / 2a<1a1 2b<1a2 / 3a<2a1,2b1L & U/UU/U & {-}/{-}{-}/{-} & N &  \\
C5527 & 1a*<o2 / 2a<1a1 2b<1a2 / 3a<2a1,2b2 & U/UU/U & {-}/{-}{-}/{-} & N &  \\
C5528 & 1a*<o2 / 2a<1a1 2b<1a2 / 3a<2a1,2b2L & U/UU/U & {-}/{-}{-}/{-} & N &  \\
C5529 & 1a*<o2 / 2a<1a1 2b<1a2 / 3a<2a1L,2b1 & U/UU/U & {-}/{-}{-}/{-} & N &  \\
C5530 & 1a*<o2 / 2a<1a1 2b<1a2 / 3a<2a1L,2b2 & U/UU/U & {-}/{-}{-}/{-} & N &  \\
C5531 & 1a*<o2 / 2a<1a1 2b<1a2 / 3a<2a2,2b0L & U/UU/U & {-}/{-}{-}/{-} & N &  \\
C5532 & 1a*<o2 / 2a<1a1 2b<1a2 / 3a<2a2,2b1 & U/UU/U & {-}/{-}{-}/{-} & N &  \\
C5533 & 1a*<o2 / 2a<1a1 2b<1a2 / 3a<2a2,2b1L & U/UU/U & {-}/{-}{-}/{-} & N &  \\
C5534 & 1a*<o2 / 2a<1a1 2b<1a2 / 3a<2a2,2b2 & U/UU/U & {-}/{-}{-}/{-} & N &  \\
C5535 & 1a*<o2 / 2a<1a1 2b<1a2 / 3a<2a2,2b2L & U/UU/U & {-}/{-}{-}/{-} & N &  \\
C5536 & 1a*<o2 / 2a<1a1 2b<1a2 / 3a<2a2L,2b1 & U/UU/U & {-}/{-}{-}/{-} & N &  \\
C5537 & 1a*<o2 / 2a<1a1 2b<1a2 / 3a<2a2L,2b2 & U/UU/U & {-}/{-}{-}/{-} & N &  \\
C5538 & 1a*<o2 / 2a<1a2 2b<1a2 / 3a<2a0L,2b1 & U/UU/U & {-}/{-}{-}/{-} & N &  \\
C5539 & 1a*<o2 / 2a<1a2 2b<1a2 / 3a<2a0L,2b2 & U/UU/U & {-}/{-}{-}/{-} & N &  \\
C5540 & 1a*<o2 / 2a<1a2 2b<1a2 / 3a<2a1,2b1 & U/UU/U & {-}/{-}{-}/{-} & N &  \\
C5541 & 1a*<o2 / 2a<1a2 2b<1a2 / 3a<2a1,2b1L & U/UU/U & {-}/{-}{-}/{-} & N &  \\
C5542 & 1a*<o2 / 2a<1a2 2b<1a2 / 3a<2a1,2b2 & U/UU/U & {-}/{-}{-}/{-} & N &  \\
C5543 & 1a*<o2 / 2a<1a2 2b<1a2 / 3a<2a1,2b2L & U/UU/U & {-}/{-}{-}/{-} & N &  \\
C5544 & 1a*<o2 / 2a<1a2 2b<1a2 / 3a<2a1L,2b2 & U/UU/U & {-}/{-}{-}/{-} & N &  \\
C5545 & 1a*<o2 / 2a<1a2 2b<1a2 / 3a<2a2,2b2 & U/UU/U & {-}/{-}{-}/{-} & N &  \\
C5546 & 1a*<o2 / 2a<1a2 2b<1a2 / 3a<2a2,2b2L & U/UU/U & {-}/{-}{-}/{-} & N &  \\
\end{casetable}


\end{document}
